\documentclass[11pt,a4paper]{article}
\usepackage[T1]{fontenc}
\usepackage{lmodern}
\usepackage{amsmath,amssymb,amsthm,mathtools,mathrsfs}
\usepackage[margin=25mm]{geometry}
\usepackage{microtype,booktabs,longtable,array,enumitem,hyperref,xurl,graphicx}
\hypersetup{colorlinks=true,linkcolor=blue,citecolor=blue,urlcolor=blue,
 pdftitle={LM4: Finite group actions, loop structures, and profinite extensions},
 pdfauthor={AI-produced by GPT-6 Astra (OpenAI, ChatGPT Pro mode)}}
\setlength{\emergencystretch}{1.5em}
\setlength{\parskip}{3pt}
\newcommand{\Z}{\mathbb Z}\newcommand{\Q}{\mathbb Q}
\newcommand{\R}{\mathbb R}\newcommand{\C}{\mathbb C}\newcommand{\F}{\mathbb F}
\DeclareMathOperator{\Aut}{Aut}\DeclareMathOperator{\Out}{Out}
\DeclareMathOperator{\Inn}{Inn}\DeclareMathOperator{\Hom}{Hom}
\DeclareMathOperator{\id}{id}\DeclareMathOperator{\im}{im}
\DeclareMathOperator{\coker}{coker}\DeclareMathOperator{\Spec}{Spec}
\DeclareMathOperator{\Gal}{Gal}\DeclareMathOperator{\Map}{Map}
\DeclareMathOperator{\Homeo}{Homeo}\DeclareMathOperator{\Diff}{Diff}
\DeclareMathOperator{\rank}{rank}\DeclareMathOperator{\Sec}{Sec}
\DeclareMathOperator{\GL}{GL}\DeclareMathOperator{\PGL}{PGL}
\newenvironment{claim}[3]{\par\medskip\phantomsection\label{cl:#1}%
 \noindent\textbf{#1. [#2] #3.}\ \ignorespaces}{\par\medskip}
\newcommand{\passage}[1]{\par\noindent{\footnotesize\emph{Manuscript:} #1.}\par}
\begin{document}
\begin{center}
{\Large\bfseries Finite group actions, loop structures,\par
and profinite extensions\par}\smallskip
{\large A mathematical development of \emph{La Longue Marche}: LM4\par}\medskip
\end{center}
\noindent AI-produced by GPT-6 Astra (OpenAI, ChatGPT Pro mode). No priority claimed.
\par\noindent 13 November 2026

\noindent\textbf{Working manuscript: INCOMPLETE.}
The numbered statements distinguish proved assertions, complete counterexamples,
source theorems with their reductions, and the exact obligations still under
investigation. This working record makes no assertion of completion of the task.

\smallskip
\noindent\textbf{Notation and its manuscript concordance.}
Local generators, base points, coefficients and comparison maps are defined at
their use. The table fixes the global symbols and the meaning of each variant.

{\small\renewcommand{\arraystretch}{1.13}\noindent
\begin{tabular}{@{}p{.16\textwidth}p{.49\textwidth}p{.29\textwidth}@{}}
\toprule
Our notation & Meaning and Grothendieck's notation & Introducing passage / scans\\
\midrule
$X,S,U=X-S$ & Compact surface, marked set, punctured surface; $X,S,U$ & \S12: 91--92; \S19: 204\\
$(g,\nu),I$ & Genus, number and index set of punctures; $(g,\nu),I$ & \S13 bis: 115--116\\
$\pi,\pi_0,P$ & Discrete surface group, chosen discretization, its completion; $\pi,\pi_0,\pi$ & \S12: 91; \S28: 352\\
$L_i,\kappa_i$ & Indexed peripheral subgroup and exterior source map; $L_i,\kappa_i$ & \S12: 91--99\\
$T=T_\pi$ & Covariant orientation line; dual of the earlier cohomological line & 91; 116; 130\\
$A,A^0$ & Surface homeomorphism group and identity component; $A_{g,\nu},A^0$ & \S19: 204--205\\
$M,\Gamma_{g,\nu}$ & Geometric mapping class group; $\mathcal T_{g,\nu},\Gamma_{g,\nu}$ & 193--203; 204--205\\
$!,+$ & Pointwise marked-point preservation; orientation preservation & 204--205\\
$G,\Gamma,E$ & Finite action, operator group, lift extension; his $\Gamma,G,E$ as locally specified & \S15: 137--139; 307--312\\
$\Phi,\mu$ & Maximal partial-section data and torsion multiplier; $\Phi,\mu$ & 150--160; 197--203\\
$SE,S\Gamma$ & Special actual and exterior automorphism groups; $SE,S\Gamma$ & \S25: 296--300\\
$B,\overline M$ & $\Out_{\rm lac}(P)$ and closure of a specified geometric group; $\widehat{\widehat{\mathcal T}},\widehat{\mathcal T}$ where identified & 318--320; 352--358\\
$\operatorname{Discr},\operatorname{Prediscr}$ & Actual, outer, or closed-stabilizer cosets as explicitly specified & \S28: 352--358\\
$GT_0,\lambda,f$ & The normalized centralizer of $S_3$, multiplier and commutator parameter & Two-group problem: 312, 316--317\\
\bottomrule
\end{tabular}}

\clearpage
\section*{Reading the proof record}
\addcontentsline{toc}{section}{Reading the proof record}
We work in the order A--F of the supplied task. A group with loops retains its
indexed exterior peripheral maps and its orientation line; this includes the
exceptional cases where the image subgroups alone lose that data. All
homomorphisms between profinite or pro-$p$ groups are continuous. A section
retains its prescribed quotient labels, and its equivalence is conjugacy by
the named kernel. Every passage citation gives the manuscript section, folio,
and scan; a folio followed by ``verso'' names its reverse side. A transcription
alternative is retained when mathematical evidence bears on it. Our corrections
are assertions about the displayed mathematics, not declarations that a
doubtful reading has been settled.

The results register contains the full numbered assertions and dependencies.
The readings register records the editorial consequences. The report gives
the obligation-by-obligation state, source checks and computation provenance.
An ``Open'' label records an unfinished obligation in this working manuscript;
it does not replace a requested proof. A ``Known'' label includes the
derivation from the precise source theorem. The code uses exact finite
arithmetic, and the proofs identify what the computations certify and what
requires a separate infinite-group argument.

\tableofcontents
\clearpage

\section{Groups with loops and their rigidity}
\label{sec:part-a}
We use $[x,y]=xyx^{-1}y^{-1}$. A peripheral subgroup is indexed even when two indices give the same subgroup. The notation $\Out_{\mathrm{lac}}$ means preservation of the indexed peripheral conjugacy classes, with their permitted permutations; in hyperbolic type these permutations are determined by the automorphism. We specify the additional data needed in the exceptional types below. All profinite homomorphisms and automorphisms in this section are continuous.

\begin{claim}{A.1}{Proved}{The presentation, ranks, and exceptional types}
Let $X$ be a connected closed oriented topological surface of genus $g$, let $S$ consist of $\nu$ distinct points, and put $U=X\setminus S$. A discrete group with loops of type $(g,\nu)$ consists of a group $\pi$, a finite index set $I$ of cardinality $\nu$, a free rank-one integral orientation module $T$, and exterior homomorphisms $\kappa_i:T\to\pi$. Here an exterior homomorphism is a conjugacy class, under conjugation in the target, of actual homomorphisms. After choosing a generator of $T$, representatives $c_i=\kappa_i(1)$, an ordering of $I$, and the handle generators, these data admit the presentation
\[
 \pi_{g,\nu}=\left\langle a_1,b_1,\ldots,a_g,b_g,c_1,\ldots,c_\nu\ \middle|\
 \prod_{j=1}^{g}[a_j,b_j]c_1\cdots c_\nu=1\right\rangle,
 \qquad L_i=\langle c_i\rangle.
\]
The profinite version uses a free rank-one $\widehat\Z$-module $T$, continuous exterior homomorphisms, and the corresponding presentation in profinite groups, namely the quotient of the free profinite group by the closed normal subgroup generated by this relator. When $I$ is empty one additionally identifies $T$ with the dual of $H^2(\pi;\Z)$, or its continuous profinite counterpart. This is the covariant orientation convention of scan 116. More precisely, an earlier cohomological marking $\kappa:H^2(\pi;\Z)\xrightarrow{\sim}T_{\mathrm{coh}}$ on scan 91 gives, by dualization, $\kappa^\vee:T_{\mathrm{coh}}^\vee\xrightarrow{\sim}H_2(\pi;\Z)$; the covariant module used here is $T=T_{\mathrm{coh}}^\vee$. The same typed dual map specifies the profinite convention. For an isomorphism of groups with loops the data are an isomorphism $f:\pi\to\pi_1$, a bijection $s:I\to I_1$, and an isomorphism $t:T\to T_1$ such that $f\kappa_i$ is conjugate to $\kappa_{s(i)}t$ for every $i$, with the corresponding dual-cohomology compatibility when $I$ is empty. An exterior group with loops regards two such $f$ as equivalent when they differ by target inner conjugation, with the same $s,t$. Thus its isomorphism sets are explicitly the quotient of the preceding isomorphism sets by $\Inn(\pi_1)$. In all nontrivial admitted types $L_i=\kappa_i(T)$ recovers its source by A.2; the one-puncture sphere requires the original source map.
For $\nu>0$ the group is free of rank $2g+\nu-1$, and its abelianization has that rank. For $\nu=0$ the abelianization has rank $2g$. The abelian types are exactly
\[
 (0,0),\quad(0,1),\quad(0,2),\quad(1,0),
\]
with groups $1,1,\Z,\Z^2$, respectively. The same list holds for their profinite completions. The peripheral maps from the source orientation group $\Z$ to $\pi$ are injective except at $(0,1)$, among the types admitted in scans 91--92. The closed sphere is explicitly excluded there.
\end{claim}
\passage{12; 44--48; 91--100; 13 bis; 56--57; 115--118}

\begin{proof}
Eliminate $c_\nu$ when $\nu>0$; the remaining generators have no relation. Abelianizing the displayed presentation leaves only $c_1+\cdots+c_\nu=0$, giving both ranks. A free group of rank at least two is nonabelian, as witnessed by its quotient $S_3$, generated by $(12)$ and $(123)$. A closed genus-$g$ group for $g\ge2$ maps onto the free group on $a_1,\ldots,a_g$ by killing all the $b_j$. The same finite quotient witnesses noncommutativity after completion. For $\nu\ge2$, each peripheral generator belongs to a free basis obtained by eliminating a different peripheral generator. For $\nu=1$, $g\ge1$, the nonempty cyclically reduced word $\prod[a_j,b_j]$ is nontrivial and has infinite order. Its image has infinite order in every pro-$p$ completion: map the first pair into the upper unitriangular $3\times3$ matrices over $\Z_p$, and all other pairs to $1$; the commutator is $1+E_{13}$. This also proves that its procyclic closure in the full completion is $\widehat\Z$, since every prime component injects. At $(0,1)$ the peripheral word is $1$, so the source map $\Z\to1$ is not injective. These calculations prove the correction $\nu\le2$ to the bracket $\nu\le3$ on scan 116. This is mathematical evidence of an error, not a changed reading of the visible numeral.
\end{proof}
\noindent\emph{Manuscript:} \S12, folios 44--48, scans 91--100; \S13 bis, folios 56--57, scans 115--118.

\begin{claim}{A.2}{Proved}{Canonical comparison of the orientation modules}
Assume $\nu>0$ and $\pi\ne1$. There is a canonical exact sequence
\[
0\longrightarrow T_\pi\longrightarrow\bigoplus_{i\in S}L_i
 \longrightarrow\pi^{\mathrm{ab}}\longrightarrow H_1(X;\Z)\longrightarrow0,
\qquad
T_\pi=\ker\left(\bigoplus_iL_i\longrightarrow\pi^{\mathrm{ab}}\right).
\]
Each projection $T_\pi\to L_i$ is an isomorphism of rank-one integral modules. If $M\subset U$ is obtained by removing small open discs at the punctures, the boundary morphism and evaluation give canonical isomorphisms
\[
 H_2(M,\partial M;\Z)\xrightarrow{\sim}T_\pi,
 \qquad
 T_\pi\xrightarrow{\sim}\Hom\bigl(H^2_c(U;\Z),\Z\bigr).
\]
For $\nu=0$, $g\ge1$, the corresponding formula is
$T_\pi=\Hom(H^2(\pi;\Z),\Z)=H_2(X;\Z)$.
Filling the discs identifies both orientation lines with $H_2(X;\Z)$, naturally for homeomorphisms of the marked surface. All these comparisons persist with $\widehat\Z$ and continuous relative group cohomology. At $(0,1)$ the kernel formed from the image subgroup $L_1=1$ is zero, although $H^2_c(\R^2;\Z)=\Z$; the source peripheral morphism is therefore essential. At $(0,0)$ group cohomology $H^2(1;\Z)=0$ cannot supply the surface orientation line.
\end{claim}
\passage{13 bis; 56--57; 116--118; 14; 64 verso; 130}

\begin{proof}
The relevant portion of the integral homology sequence of the compact pair is
\[
0\to H_2(M,\partial M)\xrightarrow{\partial}
H_1(\partial M)\to H_1(M)\to H_1(M,\partial M).
\]
The initial zero follows from $H_2(M)=0$. A relative fundamental class has boundary the sum of the consistently oriented boundary circles. Identifying each boundary circle with its peripheral subgroup, the boundary morphism identifies $H_2(M,\partial M)$ with the stated kernel and has degree $1$ on each factor. The cokernel is $H_1(X)$ by attaching the discs and using the homology sequence. Universal coefficients for the pair, whose homology is free, identify $H^2(M,\partial M)$ with the dual of $H_2(M,\partial M)$; excision identifies the former with $H^2_c(U)$. Let $D\subset X$ be the union of the filled discs. The canonical isomorphisms $H_2(X)\xrightarrow{\sim}H_2(X,D)\xrightarrow{\sim}H_2(M,\partial M)$ are, respectively, the pair map and excision. Their inverse is the filling comparison, sending a relative fundamental class to the closed fundamental class and proving compatibility without a chosen sign. Conjugating a peripheral subgroup changes its identification by an inner normalizer action, which is trivial by A.3; hence the sequence is independent of paths to the boundary. For the annulus the same conclusion follows directly because its group is abelian.

For clarity, the profinite compact-support notation is the relative continuous cochain complex
\[
 C^*_{\mathrm{rel}}(\widehat\pi;(\widehat L_i);\Z/n)
 =\operatorname{Cone}\left(C^*_{\mathrm{cont}}(\widehat\pi;\Z/n)
 \longrightarrow\bigoplus_iC^*_{\mathrm{cont}}(\widehat L_i;\Z/n)\right)[-1].
\]
For $\nu>0$, $\widehat\pi$ has cohomological dimension one and each $\widehat L_i\simeq\widehat\Z$; thus its degree-two group is the cokernel of restriction on $H^1$. In the displayed presentation the restriction matrix is the transpose of the integral peripheral inclusion matrix. Its cokernel is canonically $\Hom(T_\pi,\Z/n)$. Taking the inverse limit over $n$ gives the claimed duality. For a closed surface of genus at least two, the continuous duality pairing identifies $H^2(\widehat\pi;\widehat\Z)$ with the completed surface orientation cohomology; this is the surface Poincare duality statement in \cite[Proposition 1.3]{A_Mochizuki2007}. At genus one, for each prime $p$ the Koszul resolution of the trivial module over $\Z_p[[\Z_p^2]]$ has differentials given by the two generators minus $1$. Applying trivial coefficients $\Z/p^r$ makes both differentials zero, so the cup product identifies degree two with the exterior square of degree one, exactly as for the discrete torus. The other prime factors have zero positive-degree cohomology with these coefficients. The Chinese remainder theorem gives the comparison for every finite coefficient $\Z/n$. The transition maps on degree one are surjective, so their first derived inverse limit vanishes; consequently passing to inverse limits introduces no extra degree-two term. The canonical map from the discrete presentation to its completion induces the identity on the surface relation and on every peripheral source; consequently the finite-coefficient comparisons commute with filling. The two trivial-group exceptions follow by the explicit cohomology calculations in the statement.
\end{proof}
\noindent\emph{Manuscript:} \S13 bis, folios 56--57, scans 116--118; \S14, folio 64 verso, scan 130. We read $H^2_!$ on scan 130 as the relative complex just specified, rather than as ordinary cohomology of the underlying abstract profinite group.

\begin{claim}{A.3}{Proved}{Discrete peripheral rigidity in every hyperbolic type}
For $2g+\nu>2$, the group $\pi=\pi_{g,\nu}$ has trivial centre and $C_\pi(H)=1$ for every finite-index subgroup $H$. For every peripheral subgroup and every integer $n\ge1$,
\[
 N_\pi(L_i)=C_\pi(L_i)=C_\pi(L_i^n)=L_i.
\]
Distinct peripheral indices give distinct, noncommensurable conjugacy classes.
\end{claim}
\passage{12; 47 verso--48; 98--99; 13 bis; 57; 117--118}

\begin{proof}
For $\nu>0$, the free action of a free group on its Cayley tree shows that a nonidentity element has an axis and that its centralizer is the infinite cyclic group of translations generated by its primitive root. An element normalizing this cyclic group preserves the axis. A reversal of the axis fixes a vertex or an edge midpoint, which is incompatible with the free, inversion-free action; the normalizer is the same cyclic group. For $\nu\ge2$ each peripheral element is primitive in a free basis. For $\nu=1$, the cyclic word $\prod[a_j,b_j]$ contains the positive letter $a_1$ exactly once. A proper cyclic power has the number of occurrences of every signed letter divisible by its exponent, so this word is not a proper power. This proves the peripheral formulas including the marginal case of scan 117.

For the closed case, the hyperbolic polygon realization of \cite[Theorem 1.2]{A_FarbMargalit2012} embeds $\pi_g$ as a torsion-free cocompact group of orientation-preserving isometries of the hyperbolic plane. Every nonidentity element is hyperbolic. Its centralizer preserves its axis and acts discretely by translations, hence is cyclic. An axis reversal would be an order-two elliptic element and is absent. A finite-index subgroup is a surface group of genus $1+[\pi:H](g-1)\ge2$. In the punctured case the covering has free rank $1+[\pi:H](2g+\nu-2)\ge2$. Thus a finite-index subgroup cannot lie in a cyclic centralizer. This proves $C_\pi(H)=1$, including $H=\pi$.

For distinct punctures with $\nu\ge3$, their primitive abelianization vectors span distinct lines; for $\nu=2$, $g\ge1$, choose the free basis containing $c_1$. The cyclically reduced lengths of $c_1$ and $c_2=(\prod[a_j,b_j]c_1)^{-1}$ are $1$ and $4g+1$, and their abelianizations force any commensurability to match equal absolute exponents. They cannot be conjugate. This proves the final assertion; at $\nu=1$ it is vacuous.
\end{proof}
\noindent\emph{Manuscript:} \S12, folios 47 verso--48, scans 98--99; \S13 bis, folio 57, scans 117--118.

\begin{claim}{A.4}{Proved}{Exact exceptional-case table}
The discrete and profinite exceptional groups have the following properties; in the profinite column replace $\Z$ by $\widehat\Z$.
\[
\begin{array}{c|c|c|c|c}
(g,\nu)&\pi& Z(\pi)&C_\pi(H)\ (H\text{ finite index/open})&\text{peripherals}\\\hline
(0,0)&1&1&1&\text{none}\\
(0,1)&1&1&1&L_1=1\\
(0,2)&\Z&\Z&\Z&L_1=L_2=\pi\\
(1,0)&\Z^2&\Z^2&\Z^2&\text{none}
\end{array}
\]
Consequently the peripheral normalizer and power-centralizer equalities hold in the plane and annulus, while distinctness of the two peripheral classes fails on the annulus. The nonzero vector $1$ on the annulus and $(1,0)$ on the torus witness failure of both centre and finite-index-centralizer assertions. The plane's failure is the noninjective source map, not a nontrivial centre. For the standard base-element formula of scan 501, the annulus satisfies the equality since its basis element generates its whole group. On the torus $C(\langle(1,0)\rangle^n)=\Z^2$ for every $n\ge1$, and $(0,1)$ witnesses its failure to equal $\langle(1,0)\rangle$. The same vector works after completion. The sphere and plane have no nontrivial base elements to test.
\end{claim}
\passage{20; 111--115 verso; 222--230; 12; 47 verso; 98; 13 bis; 56 verso; 116}

\begin{proof}
A.1 identifies the groups and peripheral images. Every subgroup of an abelian group is centralized by the entire group. The assertions now follow by direct substitution. For the profinite version the listed elements remain nonzero, as their reductions modulo $2$ show.
\end{proof}
\noindent\emph{Manuscript:} \S20, folios 111--115 verso, scans 222--230; omissions on scans 98 and 116.

\begin{claim}{A.5}{Proved}{Peripheral decomposition groups and transporters}
Let $1\to\pi\to E\xrightarrow{p}\Gamma\to1$ arise from an action on a hyperbolic marked surface. Choose one peripheral subgroup $L_s$ and the corresponding peripheral lift group $E_s\subset E$; thus $E_s\cap\pi=L_s$ and $p(E_s)=\Gamma_s$, the stabilizer of that puncture. Define
\[
\operatorname{Trans}_E(L_s^d,L_s)
 =\{e\in E:eL_s^de^{-1}\subseteq L_s\}\qquad(d\ge1).
\]
Then $N_E(L_s)=\operatorname{Trans}_E(L_s^d,L_s)=E_s$. The same assertions hold for the corresponding profinite extension and closed peripheral groups. On the annulus these assertions fail for $z\mapsto1/z$: with $E=\Z\rtimes C_2$ and the involution acting by inversion, $E_s=\Z$ but $N_E(L_s)=E$ and every transporter is $E$. More generally for every annulus action both normalizer and transporter equal $E$, while $E_s$ is the inverse image of the end stabilizer. Thus their equality with $E_s$ holds exactly when the action fixes both ends. On the one-puncture sphere the peripheral mixed group maps onto $E=\Gamma$ with kernel $\Z$; using its image, both equalities hold, whereas identifying that mixed group itself with a subgroup is invalid. The closed sphere and torus have no peripheral statement.
\end{claim}
\passage{9; 33; 69; 12; 48; 99}

\begin{proof}
A normalizer element maps to $\Gamma_s$, since the peripheral conjugacy classes are distinct. Choose $e_s\in E_s$ with the same image. Then $e_s^{-1}e$ belongs to $N_\pi(L_s)=L_s$, proving the normalizer equality. For the transporter assertion, $eL_s^de^{-1}\subset L_s$ implies $L_s^d\subset e^{-1}L_se$. The latter is abelian, so it lies in $C_\pi(L_s^d)=L_s$, by A.3 or A.7. In the discrete case a peripheral generator cannot be a proper power, so the inclusion of peripheral groups $e^{-1}L_se\subset L_s$ is an equality. In the profinite case use compactness of $E$: a closed subgroup $L$ of a compact group cannot be properly contained in an inner conjugate. Indeed $aLa^{-1}\subset L$ gives $a^kLa^{-k}\subset L$ for all $k\ge1$, and a net of positive powers of $a$ converges to $a^{-1}$ inside its procyclic closure. Passing to the limit gives the reverse conjugation inclusion and hence equality. Thus the transporter normalizes $L_s$ in both cases. This argument only uses the exact power-centralizer statement and avoids replacing the one-sided edge statement in \cite[Proposition 1.2(i)]{A_Mochizuki2007} by a stronger unstated assertion about chosen subgroup representatives. The peripheral lift groups for the action on a punctured disc map onto the puncture stabilizer by construction, so the preceding choice of $e_s$ is legitimate. For the annulus the involution swaps the two punctures, whose inertia subgroups both equal $\pi$. Its lift is the explicit element in the normalizer outside $E_s$. Replacing $\Z$ with $\widehat\Z$ gives the completed counterexample.
\end{proof}
\noindent\emph{Manuscript:} \S9, folio 33, scan 69; \S12, folio 48, scan 99. In the plane, the map from the peripheral mixed group to $E$ has a kernel $\Z$; writing that mixed group as a subgroup would already discard essential data.

\begin{claim}{A.6}{Proved}{Centre-free extension reconstruction and its abelian correction}
For a centre-free discrete group $\pi$ and a homomorphism $\rho:\Gamma\to\Out(\pi)$, the extension inducing $\rho$ is, up to a unique isomorphism inducing the identity on its prescribed kernel and quotient,
\[
 E_\rho=\{(a,\gamma)\in\Aut(\pi)\times\Gamma:[a]=\rho(\gamma)\}.
\]
The same formula holds for a topologically finitely generated centre-free profinite $\pi$ and continuous $\rho$, with the profinite topology. For $A=\Z$ or $\Z^2$ and a prescribed action $\rho:\Gamma\to\Aut(A)$, extension classes with their kernel and quotient identified are instead classified by $H^2(\Gamma;A_\rho)$. Their automorphisms inducing the identity on both ends form $Z^1(\Gamma;A_\rho)$. For finite $\Gamma$ this statement holds with $A=\widehat\Z$ or $\widehat\Z^2$ and continuous cochains. Thus both the annulus and the torus require the abelian correction.
\end{claim}
\passage{12; 47 verso; 98; 17; 77--80; 155--160}

\begin{proof}
An extension $E$ maps to $E_\rho$ by $e\mapsto(\operatorname{conj}_e,p(e))$. Its kernel is $Z(\pi)=1$. Given $(a,\gamma)$, choose $e_0$ above $\gamma$. The automorphism $a\operatorname{conj}_{e_0}^{-1}$ is inner, implemented by a unique element of $\pi$, so the map is onto. An isomorphism inducing the identity on both ends fixes both components of this formula and is unique. In the profinite case this continuous bijection between compact Hausdorff groups is an isomorphism.

For abelian $A$, choose a set-theoretic section $t$ of the quotient with $t(1)=1$ and write $t(g)t(h)=c(g,h)t(gh)$. Associativity is exactly
$c(g,h)+c(gh,k)=\rho(g)c(h,k)+c(g,hk)$. Changing $t$ by a function $b:\Gamma\to A$ adds its coboundary. Conversely this identity makes $(a,g)(b,h)=(a+\rho(g)b+c(g,h),gh)$ associative. An automorphism fixing both ends has the form $(a,g)\mapsto(a+b(g),g)$, where the homomorphism equation says $b$ is a $1$-cocycle. For finite $\Gamma$ all choices of section are continuous. For example $\Gamma=C_2$, $A=\Z$ with trivial action has the two extension classes $\Z\times C_2$ and $\Z$, with kernel $\Z\xrightarrow{\times2}\Z$ in the latter. They induce the same outer action and are not isomorphic as extensions; reduction modulo $2$ detects the different value of $t(\sigma)^2$.
\end{proof}
\noindent\emph{Manuscript:} \S12, folio 47 verso, scan 98; \S17, folios 77--80, scans 155--160.

\begin{claim}{A.7}{Known}{The full profinite rigidity statements, with reductions}
For the full profinite completion $\Pi=\widehat\pi_{g,\nu}$ of a hyperbolic type, $\Pi$ is torsion-free, $C_\Pi(H)=1$ for every open subgroup $H$, and
\[
N_\Pi(\widehat L_i)=C_\Pi(\widehat L_i^n)=\widehat L_i
\quad(n\in\Z_{>0}).
\]
Distinct peripheral classes are noncommensurable. The same conclusions hold for the pro-$p$ completions, for every prime $p$, including $2$ and $3$. The four exceptional types have exactly the completed table of A.4, and A.5 supplies the extension and transporter assertions.
\end{claim}
\passage{13 bis; 57 verso; 118; 35; 253--254; 499--501}

\begin{proof}
Take the full formation of finite groups, respectively the formation of finite $p$-groups, in \cite[Remark 1.2.2 and Proposition 1.4]{A_MochizukiTamagawa2008}. A hyperbolic complex curve with the prescribed topological presentation identifies its geometric group with $\Pi$, respectively the pro-$p$ completion. Proposition 1.4 asserts precisely the vanishing of centralizers of open subgroups, with no properness assumption, so taking $H=\Pi$ gives centre-freeness. Remark 1.2.2 gives finite $\ell$-cohomological dimension for every relevant prime, hence excludes an element of order $\ell$: such an element would give a closed $C_\ell$ subgroup, whose $\ell$-cohomological dimension is infinite. Every nontrivial finite-order element has a prime-order power; this proves torsion-freeness.

For peripherals, \cite[Lemma 1.3.7]{A_Mochizuki2004} gives $\operatorname{Comm}_\Pi(\widehat L_i)=\widehat L_i$. A normalizer is contained in that commensurator. An element centralizing $\widehat L_i^n$ conjugates the finite-index subgroup $n\widehat L_i$ to itself, so it too belongs to the commensurator. Conversely the procyclic group $\widehat L_i$ centralizes every power. This proves both equalities. In the semi-graph of the smooth pointed curve, peripherals are the cusp edge groups. Proposition 1.2(i) of \cite{A_Mochizuki2007} says that an intersection open in one edge group identifies the underlying edge; hence distinct punctures cannot have commensurable inertia groups. This gives the needed distinction of peripheral indices; A.5 proves its transporter assertion separately using power-centralizers and compactness. No claim about centralizers of arbitrary non-peripheral elements is needed here.

For $\nu\ge2$ there is also a direct source for the power-centralizer assertion: each peripheral generator can be made a free basis element, and \cite[Lemma (2.1.2)]{A_Nakamura1994} gives the centralizer of every nonidentity element of its procyclic closure. For $\nu=1$ the commensurator argument covers precisely the case that this basis argument does not cover. For every $\nu\ge1$, each handle generator $a_j,b_j$ belongs to the free basis obtained by eliminating $c_\nu$. Applying that same lemma proves equation (4) of scan 501 for all these base elements, including the once-punctured case. The basis condition is material if the phrase ``systeme de generateurs on scan 500 is read literally without a cardinality restriction: in a free profinite group on $a,b$, the redundant generating list $a,b,a^2$ has $C(a^2)=\overline{\langle a\rangle}\supsetneq\overline{\langle a^2\rangle}$. The last index is $2$, detected in abelianization modulo $2$. For the pro-$p$ version use $a,b,a^p$ and index $p$. The stated source theorem applies to a member of a free basis; it does not supply the assertion for every redundant generating list.
\end{proof}
\noindent\emph{Manuscript:} \S13 bis, folio 57 verso, scan 118; \S35, folios 253--254, scans 499--501.

\begin{claim}{A.8}{Proved}{Lemme 5 from open-subgroup centralizers alone}
Let $P$ be a discrete or profinite group for which the centralizer of every finite-index, respectively open, subgroup is trivial. Let $H\subset P$ be such a subgroup and let $u\in\Aut(P)$ preserve $H$. An automorphism restricting to the identity on $H$ is the identity on $P$. An automorphism whose restriction equals $\operatorname{conj}_h$ for some $h\in H$ equals $\operatorname{conj}_h$ on $P$.
\end{claim}
\passage{35; 252 verso--254; 498--501}

\begin{proof}
Take the finite-index normal core $N=\bigcap_{x\in P}xHx^{-1}\subset H$. If $u|_H=\id$, then for $x\in P$ and $n\in N$,
\[
 u(x)nu(x)^{-1}=u(xnx^{-1})=xnx^{-1}.
\]
Thus $x^{-1}u(x)\in C_P(N)=1$ for every $x$. For the second assertion apply the first to $\operatorname{conj}_{h^{-1}}u$. In the profinite case the core is an intersection of finitely many open subgroups and is open. This proof uses neither peripheral generators nor their power-centralizers.
\end{proof}
\noindent\emph{Manuscript:} \S35, folio 252 verso, scan 498, Lemme 5; folios 253--254, scans 500--501, proposed proof.

\begin{claim}{A.9}{Proved}{Faithfulness ascends along a finite etale cover}
Let $k$ be a characteristic-zero field, $U/k$ a hyperbolic geometrically connected smooth curve, and $V\to U$ a finite etale geometrically connected covering defined over $k$. Put $\Gamma=\operatorname{Gal}(\overline k/k)$ and choose compatible geometric base points. Faithfulness of $\Gamma\to\Out(\pi_1(U_{\overline k}))$ implies faithfulness of $\Gamma\to\Out(\pi_1(V_{\overline k}))$.
\end{claim}
\passage{35; 253; 499--500}

\begin{proof}
The covering gives an open embedding $\Pi_V\hookrightarrow\Pi_U$ of arithmetic fundamental groups, compatible with their surjections onto the same $\Gamma$. Their geometric kernels are an open subgroup $H\subset P$. Let $\gamma\in\Gamma$ act trivially in $\Out(H)$ and choose a lift $e\in\Pi_V$. Conjugation by $e$ preserves $H$ and restricts there to conjugation by some $h\in H$. By A.7 and A.8 its action on $P$ is also conjugation by $h$. Thus $\gamma$ acts trivially in $\Out(P)$ and equals $1$. A change of compatible base points conjugates these homomorphisms and does not change their kernels. The proof does not use the structure of finite subgroups of $\Gamma$.
\end{proof}
\noindent\emph{Manuscript:} \S35, folio 253, scans 499--500, Corollaire 2.

\begin{claim}{A.10}{Known}{Closed-surface base elements: two separate source reductions}
Let $\Pi=\widehat\pi_g$ with $g\ge2$, let $x$ be any of the standard surface generators $a_i,b_i$, and let $\Lambda=\overline{\langle x\rangle}$. Then $\Lambda\simeq\widehat\Z$ and, for every positive integer $n$,
\[
 C_\Pi(\Lambda^n)=N_\Pi(\Lambda^n)=\Lambda.
\]
\end{claim}
\passage{35; 254; 501}

\begin{proof}
The standard generator maps to a primitive vector of $\Pi^{\mathrm{ab}}=\widehat\Z^{2g}$; hence its procyclic closure injects there and is $\widehat\Z$. For the first source reduction, $x$ is represented by a nonseparating simple closed curve, and is therefore algebraically simple in the terminology of \cite[Definition 1.3]{A_BoggiZalesskii2019}. Proposition 3.6 of that paper gives $N_\Pi(\Lambda^n)=\Lambda$. Centralizers lie in normalizers and contain $\Lambda$, proving the centralizer equality.

For the second reduction, pinch that curve on the genus-$g$ surface. The resulting stable one-node curve has a normalization of genus $g-1$ with two marked inverse branches; at $g=2$ this is a twice-marked genus-one component, still hyperbolic. Its associated semi-graph of anabelioids has one edge group whose image in the smoothing's profinite group is exactly $\Lambda$, up to the chosen base-path conjugation. Proposition 1.2(ii) of \cite{A_Mochizuki2007} makes this edge group commensurably terminal. An element normalizing $\Lambda^n$ commensurates $\Lambda$, proving the same normalizer equality. These are independent applications of the two stated source results, and neither replaces a closed surface group by a free group.
\end{proof}
\noindent\emph{Manuscript:} \S35, folio 254, scan 501, equation (4).

\begin{claim}{A.11}{Proved}{What the abstract profinite presentation specifies}
The full profinite presentation required by scan 318, equation (4), already determines a profinite group with loops isomorphic to $\widehat\pi_{g,\nu}$ with its standard loops. Hence A.7 and A.10 apply to every abstract object satisfying that presentation axiom. Merely requiring generators to satisfy the displayed equation would not suffice: $P=\widehat\Z^2$, $c_1=(1,0)$, $c_2=(0,1)$, $c_3=(-1,-1)$ satisfies the $(0,3)$ relation, has three injective procyclic peripheral subgroups and the rank-one orientation kernel, but has centre $P$.
\end{claim}
\passage{26; 161; 318--319}

\begin{proof}
The universal property gives a continuous homomorphism from the quotient of the free profinite group by the closed normal closure of the relator to $P$. The assertion that this is a defining relation makes this homomorphism an isomorphism; its inverse sends the specified base to the standard base. This gives exactly the comparison map in scan 319, equation (5). For the weaker condition the explicit abelian example has $c_1+c_2+c_3=0$ and generates $P$, but has the additional relation $[c_1,c_2]=1$. Its peripheral inclusion matrix has kernel generated by $(1,1,1)$, and its centre contains $(1,0)\ne0$. The finite reductions and the rank calculation are checked by the accompanying code.
\end{proof}
\noindent\emph{Manuscript:} \S26, folio 161, scans 318--319, equations (4)--(5). The words ``relation generatrice'' are indispensable; weakening them changes the objects by the explicit quotient above.

\begin{claim}{A.12}{Proved}{The discrete filling isomorphism and all compatibility maps}
Let $U=X\setminus S$ be hyperbolic, choose $b\in U$, and put $U'=U\setminus\{b\}$, $\pi'=\pi_1(U')$, and $\pi=\pi'/\langle\!\langle L_b\rangle\!\rangle$. Choose the peripheral representative $L_b\subset\pi'$. The normalization of representatives defines an isomorphism of discrete groups
\[
 d_b:\Out_{\mathrm{lac}}(\pi';b)\xrightarrow{\sim}\Aut_{\mathrm{lac}}(\pi).
\]
On both sides old punctures may be permuted and orientation may be reversed. It identifies the pure orientation-preserving subgroups, whose common index is $2\nu!$, where $\nu=|S|$. It commutes with forgetting the distinguished puncture, with the maps to $\Out_{\mathrm{lac}}(\pi)$, with orientation characters, with the permutation maps to $\mathfrak S_\nu$, and with point-pushing and inner conjugation on their kernels.
\end{claim}
\passage{18; 93--94; 186--189}

\begin{proof}
An outer class fixing $[L_b]$ has a representative $\alpha$ preserving $L_b$ itself: conjugate any representative by an element taking its peripheral image back to $L_b$. Two normalized representatives differ by conjugation by an element of $N_{\pi'}(L_b)=L_b$, by A.3. Such conjugation becomes the identity after killing $L_b$. Thus $\alpha$ induces a well-defined actual automorphism $d_b([\alpha])$ of the specified quotient $\pi$. Products of normalized representatives are normalized, so $d_b$ is a homomorphism. Its definition also proves compatibility with the permutation of the remaining loops and with the orientation module, using A.2's filling comparison.

We give the group-theoretic identification of this homomorphism with the topological one. The Dehn--Nielsen--Baer theorem identifies the extended mapping class group of a hyperbolic punctured surface with the outer automorphisms preserving its peripheral conjugacy classes, and in the closed case with $\Out(\pi)$; the precise versions are \cite[Theorems 8.1 and 8.8]{A_FarbMargalit2012}. The evaluation map of the homeomorphism group of $(X,S)$ at $b$ is a locally trivial fibration over $U$: over a sufficiently small disc, move $b$ to a nearby point by a homeomorphism supported in that disc, which constructs a local section. Its connecting map pushes $b$ round a loop. The identity component of this homeomorphism group is contractible in hyperbolic type, by the surface homeomorphism-group theorem in \cite[Lemma 2.1 and Theorem 1.1]{A_Yagasaki2000}; hence the homotopy sequence gives
\[
1\to\pi\xrightarrow{\mathrm{push}}\Out_{\mathrm{lac}}(\pi';b)
\xrightarrow{\mathrm{fill}}\Out_{\mathrm{lac}}(\pi)\to1.
\]
One may alternatively use the extended Birman exact sequence, \cite[Theorem 4.6 and the immediately following extended-group remark]{A_FarbMargalit2012}. Surjectivity follows directly by moving the image of $b$ back to $b$ along an arc avoiding $S$.

Choose the convention that the point-push associated to a based loop $\gamma$ moves the base-point track in the direction for which the induced based map is $\eta\mapsto\gamma\eta\gamma^{-1}$. This can also be checked on the boundary of the homotopy square traversed by $\eta$ under the isotopy. It proves $d_b\circ\mathrm{push}=\operatorname{conj}$. The filled mapping class gives the same outer automorphism as $d_b$ by construction. Thus the last sequence maps to
\[
1\to\pi\xrightarrow{\operatorname{conj}}\Aut_{\mathrm{lac}}(\pi)
\to\Out_{\mathrm{lac}}(\pi)\to1
\]
by the identity on the end groups. Injectivity and surjectivity of $d_b$ follow by the elementary exact-sequence argument: a kernel element belongs to point-pushing and maps trivially by the injective inner map, and any target element can first be matched in the outer quotient and then in the kernel.

Finally every permutation of old punctures is realized by half-twists supported in small discs along arcs, and an orientation-reversing homeomorphism can be chosen to fix the finite marked set pointwise. Hence the map to $\{\pm1\}\times\mathfrak S_\nu$ is onto; its kernel is precisely the pure orientation-preserving group. Its index is $2\nu!$. This also identifies the subgroups denoted $\mathcal T^{00}$ on scan 188 and proves the compatibility checks left on scan 189.
\end{proof}
\noindent\emph{Manuscript:} \S18, folios 93--94, scans 186--189. On scan 187 the inequality must be $2g+\nu'\ge4$ when $\nu'$ counts punctures before filling; the meaningful alternative $\le4$ contradicts both examples and the stated hyperbolicity requirement. The mathematical argument supports the first alternative without settling the handwriting.

\begin{claim}{A.13}{Proved}{Filling into the two abelian targets}
Exactly one of the drilled and filled groups is abelian in the two cases $(0,3)\to(0,2)$ and $(1,1)\to(1,0)$. In the first, the homomorphism on automorphisms of the underlying groups with peripheral subgroup classes is
\[
 C_2\times C_2=\Out_{\mathrm{lac}}(F_2;[\langle xy\rangle])
 \longrightarrow\Aut(\Z)=C_2,
\]
with kernel of order two. Its nontrivial element has representative $x\mapsto y^{-1}$, $y\mapsto x^{-1}$. Its restriction to the positive-orientation subgroup is an isomorphism. Retaining the two indexed peripheral source maps and their orientation line on the annulus enlarges the target automorphism group to $C_2\times C_2$ and restores the full isomorphism. In the second case the filling map is the isomorphism $\operatorname{GL}_2(\Z)\to\operatorname{GL}_2(\Z)$, with positive subgroup $\operatorname{SL}_2(\Z)$.
\end{claim}
\passage{18; 94--97; 189--194}

\begin{proof}
A.1 lists the two cases. For a three-punctured sphere the extended mapping class group is $\mathfrak S_3\times C_2$: a homeomorphism fixing the three punctures is isotopic, in each orientation class, to the identity or a fixed reflection, by cutting along three joining arcs and applying the disc isotopy. The stabilizer of the distinguished puncture is $C_2\times C_2$. On $F_2=\langle x,y\rangle$, the old-puncture transposition $x\leftrightarrow y$ and the reflection $x\mapsto x^{-1},y\mapsto y^{-1}$ both induce inversion after imposing $xy=1$. Their product is the displayed nontrivial kernel element. It takes the distinguished loop to its inverse and reverses orientation. Direct free-word and finite-group calculations appear in the attached program.

For an annulus with ordered end sources $\kappa_1(1)=t$ and $\kappa_2(1)=t^{-1}$, an automorphism includes a sign $\epsilon\in\{\pm1\}$ on $T$ and a permutation $s\in\mathfrak S_2$. Compatibility forces the underlying map $t\mapsto t^{\epsilon\operatorname{sgn}(s)}$. Thus there are four automorphisms of this full data, and the map forgetting $s,\epsilon$ is exactly the preceding quotient. This calculation specifies the information lost by retaining only subgroup images.

For the torus, every matrix in $\operatorname{GL}_2(\Z)$ is realized by a linear torus homeomorphism fixing the origin. A based torus homeomorphism inducing the identity matrix is isotopic relative to that point to the identity. Indeed the restriction of evaluation to the translation subgroup $\R^2/\Z^2$ is the identity of the torus, so evaluation is surjective on fundamental groups and the point-pushing kernel is zero. The unbased kernel is the identity component: isotope the images of two primitive curves intersecting once back to those curves, use the identity matrix to remove their remaining twists, and cut along them to a disc, where the Alexander isotopy finishes the argument. Removing the fixed point identifies this mapping class group with the peripheral-preserving outer group of the once-punctured torus by Dehn--Nielsen--Baer. Filling induces abelianization, hence the asserted identity map on matrices. Its orientation sign is the determinant.
\end{proof}
\noindent\emph{Manuscript:} \S18, folios 94--97, scans 189--194. The abelian cases in the margin of scan 194 refer to the filled types; the corresponding drilled groups are both $F_2$.

\begin{claim}{A.14}{Known}{Primitive loops in finite-step solvable quotients}
Let $\Sigma$ be a nonempty set of primes, $F$ a free pro-$\Sigma$ group of finite rank at least two, and $F^{[j]}$ its closed derived series. Put $F_m=F/F^{[m]}$, $m\ge2$. For any element $x$ of a free basis and any positive integer $n$,
\[
 C_{F_m}(x^n)=N_{F_m}(\overline{\langle x\rangle})
 =\overline{\langle x\rangle}\simeq\Z_\Sigma,
 \qquad \Z_\Sigma=\prod_{p\in\Sigma}\Z_p.
\]
Every open subgroup of $F_m$ has trivial centralizer. For $\nu\ge2$, all peripheral loops of the pro-$\Sigma$ group of type $(g,\nu)$ satisfy these formulas in every $F_m$. Their indexed conjugacy classes are pairwise distinct in all hyperbolic types.
\end{claim}
\passage{13 bis; 57--57 verso; 117--118; strengthened PartA(c)}

\begin{proof}
The power-centralizer formula is \cite[Theorem 1.5, finite-rank case]{A_Yamaguchi2026}; the source explicitly treats every nonzero integer power, including powers divisible by $2$ or $3$. We checked its finite-rank proof: Proposition 1.3 uses the finite semidirect product $(\Z/N)[G/\langle\bar x\rangle]\rtimes(G\times\langle\bar x\rangle)$ to force the centralizer into $\langle x\rangle F_m^{[m-1]}$; the coefficient at the distinguished coset forces the first $G$ coordinate into $\langle\bar x\rangle$. At $m=2$, the Blanchfield--Lyndon injection of $F'/F''$ into $\Z_\Sigma[[F^{\mathrm{ab}}]]^r$ and injectivity of multiplication by $\bar x^n-1$ remove the intersection with $F_m'$. That injectivity is proved by passing from $C_M$ to $C_{kM}$: invariant coefficients downstairs are divisible by every $\Sigma$-integer $k$, so are zero. The finite-rank induction passes to finite abelian covers where a positive power of $x$ is a free basis element. These are exactly the finite-rank hypotheses used here; no infinite-rank inverse-limit assertion is used.

If $u$ normalizes $\langle x\rangle$, write $uxu^{-1}=x^\lambda$, $\lambda\in\Z_\Sigma^\times$. The primitive abelianization coordinate gives $\lambda=1$, so the normalizer lies in the already computed centralizer. An open subgroup contains positive powers of two distinct free generators; their cyclic centralizers intersect trivially because their abelianization coordinates are independent. This proves the open-centralizer assertion directly.

Each peripheral generator for $\nu\ge2$ belongs to a free basis after eliminating another boundary generator. For $\nu\ge3$ distinct peripheral abelianization lines separate the conjugacy classes. When $\nu=2$ and $g\ge1$, kill $c_1$ and all but one handle pair, mapping onto a free metabelian group on $a,b$. The image of $c_1$ is $1$ and that of $c_2$ is $[a,b]^{-1}\ne1$; the latter has infinite order in the upper unitriangular $3\times3$ group over $\Z_\Sigma$. Hence they cannot be conjugate, or commensurable with finite indices. These maps factor through every $F_m$, $m\ge2$.
\end{proof}
\noindent\emph{Manuscript:} \S13 bis, scans 117--118, and the strengthened finite-step task in Part A(c).

\begin{claim}{A.15}{Proved}{The compact intersection calculation used in descent}
Let $(H_\alpha)$ be a decreasing directed family of closed subgroups of a profinite group $P$, and put $K=\bigcap_\alpha H_\alpha$. Then
\[
 K^{[j]}=\bigcap_\alpha H_\alpha^{[j]}\quad(j\ge0).
\]
For a closed subgroup $L\subset P$ and decreasing closed normal subgroups $N_\alpha$ of $P$, one also has $\bigcap_\alpha LN_\alpha=L(\bigcap_\alpha N_\alpha)$.
\end{claim}
\passage{Auxiliary compact-group lemma for finite-step PartA(c)}

\begin{proof}
For any open normal $V\subset P$, the images of $H_\alpha$ in the finite group $P/V$ eventually stabilize. Their eventual image equals that of $K$: each compatible coset fibre is a decreasing family of nonempty compact sets. Taking a fixed number of derived subgroups commutes with homomorphic images, and taking the inverse limit of all finite images recovers the closed derived subgroup. This proves the first equality. For the second equality, an element $z$ of every $LN_\alpha$ gives nonempty closed sets $L\cap zN_\alpha$; compactness makes their intersection nonempty. A member of that intersection gives the required single factorization of $z$.
\end{proof}

\begin{claim}{A.16}{Proved}{The one-puncture solvable threshold, with a finite certificate}
Let $F$ be the pro-$\Sigma$ completion of $\pi_{g,1}$, $g\ge1$, let $c=\prod_{j=1}^g[a_j,b_j]$ (the inverse peripheral generator), and write $c$ for its image in $F_m=F/F^{[m]}$. For every positive integer $n$,
\[
\begin{array}{ll}
m=2:& C_{F_2}(c^n)=N_{F_2}(\overline{\langle c\rangle})=F_2',\\[2pt]
m\ge3:&C_{F_m}(c^n)=N_{F_m}(\overline{\langle c\rangle})
 =\overline{\langle c\rangle}.
\end{array}
\]
In the first line the inclusion $\overline{\langle c\rangle}\subset F_2'$ is strict. Thus peripheral rigidity fails at the smallest solvable depth $m=2$ and is restored at every $m\ge3$, for the full profinite and every pro-$p$ version.
\end{claim}
\passage{13 bis; 57 verso; 118; strengthened PartA(c)}

\begin{proof}
Let $H\subset F$ run over open normal subgroups containing $F'$, with $|F/H|\ge3$. They form a family whose intersection is $F'$. The corresponding finite abelian covering of the once-punctured surface has exactly $d=|F/H|$ punctures, because $c\in F'$ has trivial image in the deck group. Choose the cusp above the original based peripheral loop. Its loop is $c$ itself; the other loops are its conjugates by deck representatives. The submodule of $H^{\mathrm{ab}}$ generated by these cusps is
\[
 \Z_\Sigma^{F/H}\big/\Z_\Sigma\!\left(\sum_{q\in F/H}e_q\right),
\]
with the deck group acting by left translation. This follows either from the relative surface homology sequence of A.2 or directly from the free-group peripheral presentation. For $d\ge3$ the lines generated by distinct $e_q$ in this module are distinct; moreover $ne_q=ae_1$, for $0\ne n\in\Z$ and $a\in\Z_\Sigma$, forces $q=1$ and $a=n$. To check this last assertion, lift the equality to the free module and compare a coordinate different from both $1$ and $q$; its coefficient forces the possible multiple of $\sum e_q$ to be zero, after which the coordinates at $q$ and $1$ give the result.

For $m\ge2$, $F^{[m]}\subset H'$, so the natural map $H\to H^{\mathrm{ab}}$ factors through the image of $H$ in $F_m$. Let $u\in F$ represent an element of $C_{F_m}(c^n)$. The relation $uc^nu^{-1}=c^n$ modulo $F^{[m]}$ implies $ne_{\bar u}=ne_1$ in the cusp submodule of $H^{\mathrm{ab}}$. Hence $u\in H$ for every such $H$, so its image belongs to $F_m'$. The same argument for a normalizer, using $uc u^{-1}=c^\lambda$ modulo $F^{[m]}$, $\lambda\in\Z_\Sigma^\times$, gives $u\in F'$ and $\lambda=1$. The subgroup $\overline{\langle c\rangle}$ is indeed $\Z_\Sigma$: the unitriangular $3\times3$ quotient used in A.1 proves injectivity of all its prime components. Thus its normalizer equals its centralizer.

When $m=2$, the derived subgroup is abelian, contains $c$, and centralizes it. This gives equality with $F_2'$. Now let $m\ge3$ and let $uF^{[m]}$ centralize $c^n$. We have proved $u\in F'\subset H$. Since $F'\subset H$,
$F^{[m]}=(F')^{[m-1]}\subset H^{[m-1]}$; therefore there is a well-defined continuous epimorphism
\[
 H/F^{[m]}\longrightarrow H/H^{[m-1]}.
\]
The group $H$ is free pro-$\Sigma$, and $c$ is a member of a free basis because this cover has at least two punctures. By A.14, applied at depth $m-1\ge2$, the image of $u$ in this quotient belongs to $\overline{\langle c\rangle}$. Consequently
\[
 u\in\bigcap_H\overline{\langle c\rangle}H^{[m-1]}
 =\overline{\langle c\rangle}\left(\bigcap_H H^{[m-1]}\right)
 =\overline{\langle c\rangle}(F')^{[m-1]}
 =\overline{\langle c\rangle}F^{[m]},
\]
where both interchanges are proved in A.15. This proves the second line.

For strictness in the first line, set $a=a_1$, $b=b_1$, and kill the remaining handle pairs. In $\operatorname{UT}_4(\F_p)$ take
\[
 x=1+E_{12}+E_{23},\qquad y=1+E_{34}.
\]
Their group is metabelian; direct multiplication gives $[x,y]=1+E_{24}$ and $x[x,y]x^{-1}=1+E_{24}+E_{14}$. The second matrix is not a power of the first. Thus $aca^{-1}\in F_2'$ centralizes $c$ but does not belong to its procyclic closure. For every $p\in\Sigma$ this quotient factors through $F_2$ and proves strictness; the attached program independently enumerates and verifies the cases $p=2,3,5$. In particular there is no excluded residue-characteristic caveat.
\end{proof}
\noindent\emph{Manuscript:} \S13 bis, scan 118, in the strengthened finite-step question of Part A(c). This assertion does not assert a new literature priority; the searched sources and the proof audit are listed in the report.

\begin{claim}{A.17}{Proved}{Centre and torsion in all finite-step surface quotients}
Let $P$ be the full profinite or pro-$\Sigma$ completion of a hyperbolic surface group, open or closed, and let $P_m=P/P^{[m]}$. For every $m\ge2$, $P_m$ is torsion-free and centre-free. This verifies the form of \cite[Theorem 2.9 and Corollary 2.10]{A_Yamaguchi2026} needed here by the following argument from finite covers and transfer.
\end{claim}
\passage{Strengthened PartA(c); Yamaguchi2026 Theorem2.9 and Corollary2.10}

\begin{proof}
Every open subgroup of $P$ is the pro-$\Sigma$ group of a finite covering surface, so its abelianization is a free $\Z_\Sigma$-module. A closed subgroup $K$ is the intersection of its open neighbourhood subgroups; by A.15 its abelianization is the inverse limit of their abelianizations. More explicitly, a compatible system in those abelianizations lifts to a compatible coset system in the compact group, and the kernel is the intersection of the closed derived subgroups. Thus $K^{\mathrm{ab}}$ is torsion-free. Applying this to each $P^{[j]}$ and inducting in
\[
1\to(P^{[m-1]})^{\mathrm{ab}}\to P_m\to P_{m-1}\to1
\]
proves torsion-freeness, starting with the torsion-free abelianization.

We also need the faithful action of a finite deck group $Q=H/N$ on $N^{\mathrm{ab}}$, for open $N\trianglelefteq H\subset P$. Choose $\ell\in\Sigma$. If a subgroup $Q_0$ acts trivially on $H^1$ of the covering surface with $\Q_\ell$ coefficients, transfer gives an isomorphism of that $H^1$ with the $H^1$ of the quotient by $Q_0$. The two surfaces are both proper or both nonproper, and their Euler characteristics are consequently equal. Multiplicativity of Euler characteristic gives $\chi(Y)=|Q_0|\chi(Y/Q_0)$. Both are negative, so $Q_0=1$. The comparison with $N^{\mathrm{ab}}\otimes_{\Z_\Sigma}\Q_\ell$ is the canonical duality between first homology and first cohomology. Thus the deck action on $N^{\mathrm{ab}}$ is faithful.

Put $K=P^{[m-1]}$, $A=K/K'$, and $Q=P/K'$. For an open normal subgroup $N\subset Q$ containing $A$, its preimage $\widetilde N\subset P$ contains $K$, hence $K'\subset\widetilde N'$. Therefore $N^{\mathrm{ab}}=\widetilde N^{\mathrm{ab}}$, and the preceding deck faithfulness implies that an element of $Z(Q)$ belongs to every such $N$. Their intersection is $A$.

We finish by proving $A^{P/K}=0$. For any open normal $V\subset P$ containing $K$, write $j:V^{\mathrm{ab}}\to P^{\mathrm{ab}}$ for inclusion and $t:P^{\mathrm{ab}}\to V^{\mathrm{ab}}$ for transfer. The composite $tj$ on the $P/V$-invariants is multiplication by $[P:V]$, which is injective because $V^{\mathrm{ab}}$ is torsion-free. Hence inclusion is injective on those invariants. Passing to the inverse limit over $V$ shows that $K^{\mathrm{ab},P/K}\to P^{\mathrm{ab}}$ is injective. But $K\subset P'$, so this map is zero. Therefore $A^{P/K}=0$. Since $Z(Q)\subset A$, the centre is zero, as claimed. This proof checks precisely the transfer and finite-cover steps used in the cited unrefereed result.
\end{proof}

\begin{claim}{A.18}{Proved}{The remaining affine solvable rigidity list and a closed bound}
For all hyperbolic types with $\nu>0$, and all $m\ge2$, the maximal $m$-step-solvable full profinite or pro-$\Sigma$ quotient has trivial centre and trivial centralizers of open subgroups, and is torsion-free. Peripheral conjugacy classes remain distinct. For $\nu\ge2$ all peripheral normalizer and integer-power-centralizer equalities hold at every $m\ge2$; for $\nu=1$ their exact values are those of A.16. In an extension obtained by quotienting a geometric action extension by the $m$th closed derived subgroup of its kernel, the equality $N_E(L_s)=E_s$ holds in the former cases and at $m\ge3$ in the latter. At $\nu=1,m=2$ its correct value is $N_E(L_s)=F_2'E_s$, with the intersection $F_2'\cap E_s=L_s$. In all these cases every transporter $\operatorname{Trans}_E(L_s^d,L_s)$, $d\ge1$, equals the same displayed normalizer. Since the kernels are centre-free by A.17, the extension reconstruction of A.6 also applies to all these solvable quotients. Standard free-basis generators satisfy the power-centralizer formula at every $m\ge2$ by A.14.
For a closed hyperbolic group $P$ and an open subgroup $H\subset P_m$, the proven universal bound is
\[
 C_{P_m}(H)\subseteq P_m^{[m-1]}.
\]
\end{claim}
\passage{13 bis; 57 verso; 118; 35; 252 verso--254; 498--501; strengthened PartA(c)}

\begin{proof}
The first assertions are A.14, A.16, and A.17. For the extension calculation, the quotient map is injective on the peripheral group, since the peripheral generator has infinite order in $P_2$ by A.1. The image of a geometric peripheral lift group still maps onto the puncture stabilizer. The kernel of the normalizer map to that stabilizer is the normalizer calculated in A.14 or A.16, so the same lifting argument as A.5 proves the equalities. In the metabelian one-puncture case there is just one puncture, so the stabilizer is the whole acting group and multiplication by $E_s$ gives the asserted $F_2'E_s$.

It remains to justify the one-sided transporter equality in these quotients. Suppose $eL_i^de^{-1}\subset L_i$. The action sends $L_i$ to a conjugate of some $L_j$ and its chosen generator to a conjugate of $c_j^\lambda$, with $\lambda\in\Z_\Sigma^\times$. For $\nu\ge3$, the primitive abelianization lines imply $j=i$. For $\nu=2$, $g\ge1$, the quotient killing $c_i$ and all but one handle pair sends $c_j$ to a nonzero commutator of infinite order in $\operatorname{UT}_3(\Z_\Sigma)$; thus $j\ne i$ would contradict the given containment. Choose an element of $E_i$ with the same image as $e$ in the acting group and divide it out. The resulting element $u$ lies in $P_m$ and satisfies $uc_i^du^{-1}=c_i^a$ for some $a\in\Z_\Sigma$. The primitive abelianization coordinate forces $a=d$, so A.14 puts $u$ in $L_i$. This proves the equality for $\nu\ge2$. For $\nu=1$, the acting group already fixes the only puncture. After the same division by $E_i$, apply the finite abelian cover calculation in A.16 to $uc^du^{-1}=c^a$. It forces $a=d$ and reduces the transporter to the integer-power centralizer calculated there. This gives $E_i$ at $m\ge3$ and $F_2'E_i$ at $m=2$.

The closed bound recorded here is strengthened to triviality in A.19 below. For this bound, put $A=P_m^{[m-1]}$ and choose any open normal $N\subset P_m$ containing $A$. The group $N^{\mathrm{ab}}$ is the torsion-free abelianization of its preimage in $P$, as in A.17. If $z$ centralizes $H$, its conjugation action fixes the image of $H\cap N$ in $N^{\mathrm{ab}}$. This image has finite index, so some nonzero integer multiple of every element is fixed. Torsion-freeness then makes the action trivial on all of $N^{\mathrm{ab}}$. Deck faithfulness from A.17 implies $z\in N$. Intersecting all these $N$ yields $z\in A$.
\end{proof}
\noindent\emph{Manuscript:} \S13 bis, scan 118; \S35, scans 498--501; additional finite-step task A(c).

\begin{claim}{A.19}{Proved}{Closed finite-step power centralizers and slimness}
Let $P$ be the pro-$\Sigma$ completion of a closed oriented surface group of genus $g\ge2$, let $m\ge2$, and put $P_m=P/P^{[m]}$. Let $x$ be represented by a nonseparating simple closed curve, in particular by any standard handle generator. Then for every positive integer $n$,
\[
 C_{P_m}(x^n)=N_{P_m}(\overline{\langle x^n\rangle})
 =\overline{\langle x\rangle}\simeq\Z_\Sigma.
\]
Consequently $C_{P_m}(H)=1$ for every open subgroup $H$. This includes the full profinite and every pro-$p$ quotient, without a restriction at $p=2$ or $3$.
\end{claim}
\passage{35; 254; 501; strengthened PartA(c)}

\begin{proof}
We first give the completed module construction required for the metabelian case. Let $F$ be a free pro-$\Sigma$ group on $z_1,\ldots,z_r$, let $A=F/F'$, and put $R=\Z_\Sigma[[A]]$. The continuous homomorphism
\[
 \mu:F\longrightarrow R^r\rtimes A,
 \qquad z_i\longmapsto(e_i,\bar z_i),
\]
exists by freeness; $A$ acts by multiplication on the coefficient module, so the product in the target is $(v,a)(w,b)=(v+aw,ab)$. Its kernel is exactly $F''$. The inclusion $F''\subset\ker\mu$ follows because the target is metabelian. To prove the reverse inclusion without assuming any exactness result about completed chain complexes, take an arbitrary finite metabelian quotient $G$ of $F$. Write $K=G'$, $Q=G/K$, and choose a section $t:Q\to G$ with $t(1)=1$. There is an injective homomorphism
\[
 G\longrightarrow K^Q\rtimes Q,
 \qquad g\longmapsto(f_g,\bar g),\qquad
 f_g(q)=t(q)g\,t(q\bar g)^{-1}.
\]
Indeed $f_{gh}(q)=f_g(q)f_h(q\bar g)$, and, when $\bar g=1$, evaluation at $q=1$ recovers $g$. Use additive notation for the finite abelian group $K$. The finite module $K^Q$ is an $R$-module through $A\to Q$ and right translation on $Q$. Sending $e_i$ to the function coordinate of the image of $z_i$ gives a continuous $R$-linear map $R^r\to K^Q$. The resulting homomorphism of semidirect products makes the homomorphism from $F$ factor through $\mu$, since it agrees on every free generator. Hence $\ker\mu$ maps trivially to every finite metabelian quotient of $F$. These quotients separate $F/F''$, proving the assertion. In particular the function coordinate gives an injective $R$-linear morphism
\[
 F'/F''\lhook\joinrel\longrightarrow R^r.
\]
Its image is a closed $R$-submodule: conjugation gives the $A$-action, integral powers give the scalar action, and continuity extends these to the completed group ring.

Apply this construction with $r=2g$ and
\[
 P=F/\langle\!\langle r_0\rangle\!\rangle,
 \qquad r_0=\prod_{i=1}^g[a_i,b_i],
 \qquad A=P^{\mathrm{ab}}=\Z_\Sigma^{2g}.
\]
The function coordinate of $r_0$ is
\[
 v=(1-\bar b_1,\bar a_1-1,\ldots,1-\bar b_g,\bar a_g-1)\in R^{2g}.
\]
This is obtained by multiplying the four factors of each commutator in $R^{2g}\rtimes A$. Since $r_0\in F'$, its closed normal closure maps in $F'/F''$ onto exactly $Rv$: its conjugates give the $A$-multiples of $v$, their products give their integral sums, and their closure is $Rv$. The last module is closed because it is the image of the compact ring $R$. Thus the natural quotient maps induce the injective comparison
\[
 B:=P'/P''\lhook\joinrel\longrightarrow R^{2g}/Rv.                 \tag{A.19.1}
\]
Indeed $P''$ is the image of $F''$ and the kernel in $F'/F''$ is precisely the just calculated module. This establishes the needed completed Fox-module comparison directly.

Choose a geometric basis with $x=a_1$, and set $h=\bar a_1^n-1$. We use two elementary facts about these completed group rings. First, multiplication by $h$ on $R$ is injective. Second, multiplication by $\bar a_2-1$ on $R/hR$ is injective. Here
\[
 R/hR\xrightarrow{\sim}
 \Z_\Sigma[[A/\overline{\langle\bar a_1^n\rangle}]]                \tag{A.19.2}
\]
is the map induced by the group quotient. Its kernel is the closed ideal generated by $h$; that ideal equals $hR$ since $hR$ is compact. To prove the injectivity assertions, one can use the following coefficient calculation. In a completed group ring of $\Z_\Sigma\times D$, consider an element invariant under the chosen infinite cyclic coordinate, or under its $n$th power. At finite coefficient level $\Z/k$ and finite coordinate level $C_M$, enlarge the coordinate level to $C_{kM}$, taking $M$ divisible by the $\Sigma$-part of $n$ in the second case. Invariance makes the coefficients constant on each orbit. The projection down to $C_M$ sums $k$ equal coefficients on each orbit fibre and therefore is zero modulo $k$. This works with every finite quotient of the remaining factor $D$. All finite images of the original element vanish, so the element is zero. Apply this argument first to the $a_1$ coordinate of $R$, and then to the independent $a_2$ coordinate in the group of (A.19.2). It applies to arbitrary $\Sigma$, because $k$ and the chosen coordinate orders have prime factors in $\Sigma$.

These ring facts show that multiplication by $h$ is injective on $R^{2g}/Rv$. If $hz=vt$, inspect the $b_2$ coordinate, which is $h z_{b_2}=(\bar a_2-1)t$. Reducing modulo $h$ and using the second injectivity gives $t=ht'$ for some $t'\in R$. The first injectivity then gives $z=vt'$ in every coordinate. Consequently (A.19.1) implies
\[
 B^{\langle x^n\rangle}=0.                                      \tag{A.19.3}
\]

For completeness, we next constrain the abelianization of an arbitrary centralizer element, at every depth $m\ge2$. Put $V=\Z_\Sigma^{2g}$ with its ordered handle basis and put $W=\bigwedge^2V/\Z_\Sigma\omega$, where $\omega=\sum_i a_i\wedge b_i$. For every $\Sigma$-integer $k$, form the finite class-two group with underlying set $(V/kV)\times(W/kW)$ and product
\[
 (q,w)(q',w')=(q+q',w+w'+\beta(q,q')),
 \qquad
 \beta(q,q')=\sum_{i<j}q_iq'_j\,z_i\wedge z_j.
\]
Its commutator is $(0,q\wedge q')$. The standard surface relation maps to $\omega=0$, so sending the handle generators to their coordinate vectors defines a quotient homomorphism from $P$. The finite group is metabelian and a $\Sigma$-group; the homomorphism therefore factors through every $P_m$, $m\ge2$. If $u\in P_m$ centralizes $a_1^n$ and has abelianization $q$, these homomorphisms give
$n(q\wedge a_1)=0$ in $W$. The module $W$ is free because $\omega$ is primitive, so multiplication by the nonzero integer $n$ is injective. Write $q\wedge a_1=t\omega$. The coefficient of $a_2\wedge b_2$ on the left is zero; on the right it is $t$. Hence $t=0$, and the remaining coordinates give $q=\lambda a_1$ for some $\lambda\in\Z_\Sigma$. The procyclic group generated by $a_1$ injects into $V$, so dividing $u$ by $a_1^\lambda$ puts it in $P_m'$ without changing its centralization property. At $m=2$, (A.19.3) makes this remaining element trivial, proving the centralizer formula.

We now induct on $m$, simultaneously for every genus at least two and every nonseparating simple closed curve. Suppose $m\ge3$, and divide a centralizer element as just explained so that it is represented by $u\in P'$. Let $H$ run over all open normal subgroups of $P$ containing $P'$, and let $d$ be the order of $x$ in $P/H$. The corresponding finite abelian covering of the closed surface has genus at least two. The based lift of the simple curve $x$ is represented by $x^d$, traversed once. This lift is a simple nonseparating curve: its integral homology class pushes forward to $d[x]\ne0$, whereas a separating curve has zero homology class. It is therefore a standard handle generator for some geometric basis of the covering surface group $H$.

Since $P'\subset H$, we have $P^{[m]}\subset H^{[m-1]}$. The quotient morphism
\[
 H/P^{[m]}\longrightarrow H/H^{[m-1]}
\]
is consequently defined. The image of $u$ centralizes $(x^d)^n$ in its target. By the induction hypothesis at depth $m-1$, that image belongs to $\overline{\langle x^d\rangle}$. Hence
\[
 u\in\bigcap_H\overline{\langle x\rangle}H^{[m-1]}
 =\overline{\langle x\rangle}\left(\bigcap_HH\right)^{[m-1]}
 =\overline{\langle x\rangle}P^{[m]},
\]
by A.15 and $\bigcap_HH=P'$. The intersection of $\overline{\langle x\rangle}$ with $P'$ is trivial, again by its primitive abelianization coordinate. Since $u\in P'$, its image in $P_m$ is trivial. This completes the induction.

If an element normalizes $\overline{\langle x^n\rangle}$, its conjugation sends $x^n$ to $x^{n\lambda}$ for a unit $\lambda\in\Z_\Sigma^\times$. Abelianization gives $n\lambda=n$, whence $\lambda=1$, so the normalizer equals the calculated centralizer. Finally an open subgroup contains positive integer powers of $a_1$ and $b_1$. Their calculated centralizers have trivial intersection because they inject into distinct independent abelianization coordinates. This proves slimness.
\end{proof}
\noindent\emph{Manuscript:} \S35, scan 501, equation (4), and the finite-step strengthening in Part A(c). The proof establishes the exact module comparison and all closed-group descent maps used; it does not rely on the unknown general slimness statement mentioned before Proposition 2.6 in \cite{A_Yamaguchi2026}, nor claim a literature priority.

\begin{claim}{A.20}{Proved}{The typed completion comparison for filling}
For the two hyperbolic groups in A.12, put $\Pi'=\widehat{\pi'}$ and $\Pi=\widehat\pi$. There is a commutative diagram of profinite groups
\[
\begin{array}{ccc}
 \widehat{\Out_{\mathrm{lac}}(\pi';b)}
 &\xrightarrow{\ \widehat{d_b}\ }&
 \widehat{\Aut_{\mathrm{lac}}(\pi)}\\
 \big\downarrow c_{\mathrm{out}}&&\big\downarrow c_{\mathrm{aut}}\\
 \Out_{\mathrm{lac}}(\Pi';b)
 &\xrightarrow{\ d_b^{\mathrm{cont}}\ }&
 \Aut_{\mathrm{lac}}(\Pi).
\end{array}
\]
The upper horizontal map is an isomorphism. The lower horizontal map is the normalized filling homomorphism on all continuous loop automorphisms. The vertical maps are induced by extending discrete automorphisms to group completions. The same construction with pro-$p$ completions of the surface groups still has the full profinite completions of the discrete automorphism groups in its upper row. These assertions specify the precise consequence of the discrete filling theorem; no isomorphism assertion for the lower horizontal or vertical arrows is used.
\end{claim}
\passage{18; 93--94; 186--189; comparison with profinite PartA(c)}

\begin{proof}
The presentation gives $\Pi'/\overline{\langle\!\langle L_b\rangle\!\rangle}=\Pi$. By A.7, $N_{\Pi'}(L_b)=L_b$. Normalize an outer representative to preserve $L_b$, then pass to this quotient. Two normalized representatives differ by conjugation by an element of $L_b$, which dies in the quotient. This proves that $d_b^{\mathrm{cont}}$ is a well-defined homomorphism. It is continuous: the closed subgroup of $\Aut(\Pi')$ preserving $L_b$ is compact, its quotient onto the indicated outer group has kernel the inner automorphisms by $L_b$, and the continuous quotient action factors through this quotient. The same proof works for the pro-$p$ groups using A.7.

Every discrete automorphism extends uniquely to an automorphism of its completion; extensions of inner automorphisms remain inner. Since the resulting target automorphism groups are profinite, the extension maps themselves extend to the full profinite completions of the discrete automorphism groups. Normalizing a discrete representative and then completing gives the same quotient automorphism as completing and then normalizing. Thus the square commutes on the dense discrete subgroup and therefore everywhere. A.12 implies that its upper arrow is an isomorphism. The quotient maps $\widehat\pi\twoheadrightarrow\pi^{(p)}$ and $\widehat{\pi'}\twoheadrightarrow(\pi')^{(p)}$ have characteristic kernels and commute with filling, giving the further comparison from the full square to its pro-$p$ version.

For odd $p$, replacing the upper groups by their pro-$p$ completions would in general fail to define these same maps: an orientation-reversing involution is killed by every homomorphism to a pro-$p$ group, whereas its action on the pro-$p$ orientation line is multiplication by $-1\ne1$. This is an explicit obstruction to that proposed factorization, rather than an identification of two differently typed completion operations.
\end{proof}

\section{Finite actions, sections, and rotation data}
\label{sec:partB}

In this section $U=X\setminus S$ is an orientable connected surface of finite type,
unless another hypothesis is explicitly stated. A universal covering
$p:\widetilde U\to U$ is fixed, and $\pi=\operatorname{Deck}(p)$ acts on the left.
For a faithful action of a finite group $G$ on $U$, we write $E$ for the group of
all lifts of its elements. Thus $1\to\pi\to E\to G\to1$ is an extension of
\emph{discrete} groups. A partial section means a finite subgroup of $E$ projecting
isomorphically onto a subgroup of $G$; its equivalence relation throughout this
section is conjugacy by $\pi$. We use $G^+$ for the orientation-preserving subgroup.
The claims below also specify the corrections required when $G^+=1$.

\subsection{The plane lemma and the section correspondence}

\begin{claim}{B.1}{Proved}{Finite plane actions and simultaneous local models}
Every faithful action of a finite group $H$ on $\R^2$ by homeomorphisms is
conjugate to an orthogonal action. When $H^+\ne1$, $H^+$ is cyclic, $H$ is cyclic
or dihedral, and $(\R^2)^H$ is a single point. A nontrivial orientation-reversing
involution has a fixed line. Every stabilizer in a finite action on a surface
with universal cover $\R^2$ has an invariant disc on which the whole stabilizer
acts by an orthogonal action. These conclusions supply the two local assertions
used in \S15, folios 70--74, scans 141--149, and \S33, scan 440.
\end{claim}
\passage{15; 70-74; 141-149; 33; 222; 440}

\begin{proof}
We use precisely Kolev's Theorem 1.2: a compact subgroup of
$\operatorname{Homeo}(S^2)$ is conjugate to a closed subgroup of $O(3)$
\cite{bKolev2006}. Each homeomorphism of $\R^2$ extends to its one-point
compactification, and the finite group fixes the added point $\infty$.
After the stated conjugacy the group fixes a point $q$ of the round sphere.
Its action on $q^\perp$ gives an injective homomorphism into $O(2)$.
Stereographic projection from $q$ is equivariant and identifies the action on
$S^2\setminus\{q\}$ with that linear action. Finite subgroups of $SO(2)$ are the
cyclic groups of roots of unity: an element of least positive angle generates
the group, by division with remainder of angles. An orientation-reversing
orthogonal matrix in dimension two is a reflection and conjugates every rotation
to its inverse. This proves the cyclic/dihedral assertion and all assertions about
fixed sets.

For the local statement, let $H$ stabilize $x\in U$ and choose
$\widetilde x\in p^{-1}(x)$. For each $h\in H$ there is a unique lift fixing
$\widetilde x$. Uniqueness makes the assignment a homomorphism, and it is
injective because its projection to $U$ is the original action. The preceding
plane argument linearizes these lifts simultaneously. Pull back a sufficiently
small round disc about the common fixed point; choose it inside an evenly
covered neighbourhood of $\widetilde x$. Its projection is the required
$H$-invariant disc. The same assertion for a sphere follows directly from
Kolev's theorem. In particular the fixed locus of an anti-involution is locally
a diameter, with the involution interchanging its two sides. No differentiability
of the original action has entered the argument. The single-element versions
are Constantin--Kolev, Lemma 2.5, Theorems 3.1 and 4.1, and Corollary 4.2
\cite{bConstantinKolev1994}.
\end{proof}

\begin{claim}{B.2}{Refuted}{The omitted orientation hypothesis in scans 146--147}
The displayed lemma on scan 147, taken without its preceding orientation
hypothesis, is false. Its correct form is: an orientation-preserving finite-order
homeomorphism $u$ of $\C$ commuting with $T(z)=z+1$ and satisfying $u(0)=0$ is
the identity. The quotient $\C/\Z$ provides a complete proof along the route
proposed in the manuscript.
\end{claim}
\passage{15; 72-73; 146-147}

\begin{proof}
Complex conjugation is a counterexample to the unqualified sentence: it has
order two, fixes zero, and commutes with $T$. For the corrected statement,
$u(k)=k$ for every integer $k$. B.1 gives at most one fixed point for a
nonidentity orientation-preserving periodic plane homeomorphism, so $u=1$.

Here is the quotient argument. Commutation with $T$ gives a finite-order
homeomorphism $\bar u$ of the annulus $A=\C/\Z$, inducing $+1$ on
$\pi_1(A)=\Z$. It preserves the orientation of $A$. The two end loops, oriented
as boundary loops of a compact subannulus, represent opposite homology classes.
Consequently an orientation-preserving map inducing $+1$ on $H_1(A;\Z)$ fixes
each end. Filling the ends extends $\bar u$ to an orientation-preserving periodic
homeomorphism of $S^2$ fixing both added points and $[0]$. By
Constantin--Kolev, Theorem 4.1, a nonidentity such map has exactly two fixed
points. Hence $\bar u=1$. Its lift $u$ is a deck transformation $T^k$; $u(0)=0$
forces $k=0$. On scan 146 the equation $G=G^\circ$ is exactly what makes the
lifted action orientation preserving. The mathematical repair therefore restores
an inherited hypothesis; it does not justify a spectral-sequence argument.
\end{proof}

\begin{claim}{B.3}{Proved}{Sections and isolated fixed points}
Let $\widetilde U\cong\R^2$. A faithful finite action on $U$ with $G^+\ne1$
has a bijection
\[
 U^G\longrightarrow \operatorname{Sec}(E\to G)/\pi.
\]
For every section subgroup $K$, $C_\pi(K)=1$. Existence of a section forces $G$
to be cyclic or dihedral. More generally, every nontrivial finite subgroup of
$E$ whose orientation-preserving part is nontrivial is contained in exactly one
maximal finite subgroup of $E$. For an orientation-preserving action the map
from points with nontrivial stabilizer to maximal partial-section classes is a
$G$-equivariant bijection. These are the theorem and corollary of \S15, folios
73--75, scans 148--151, with their stated orientation hypotheses.
\end{claim}
\passage{15; 73-75; 148-151}

\begin{proof}
For $x\in U^G$, the lifts of $G$ fixing a chosen point over $x$ form a section.
Changing that point conjugates the section by the corresponding element of
$\pi$. Conversely a section $K$ is a finite faithful group of plane
homeomorphisms. B.1 supplies a unique fixed point $\widetilde x$. Its image is
fixed by $G$, and the section fixing $\widetilde x$ is $K$. These constructions
are inverse: conjugate sections have deck-translated fixed points and hence
the same image in $U$. An element of $C_\pi(K)$ preserves the unique fixed point.
A deck transformation with a fixed point is the identity. This proves the
centralizer assertion. B.1 gives the possible finite groups.

Let $K\le E$ have nontrivial positive part and fix $\widetilde x$. The stabilizer
$E_{\widetilde x}$ is finite, since it intersects $\pi$ trivially and injects into
the finite group $G$. Any finite subgroup containing $K$ has a fixed point by
B.1; that point must be $\widetilde x$. It is therefore contained in
$E_{\widetilde x}$. This proves uniqueness as an actual subgroup, before any
conjugacy quotient. In the orientation-preserving case every nontrivial partial
section has this property. Applying the same construction to each stabilizer
$G_x$ gives the final bijection and its equivariance. The argument also applies
to a nonorientable $U$ whenever the finite lifted group has nontrivial positive
part; in particular the cardinality condition $|G|\ge3$ on scan 149 ensures it.
\end{proof}

\begin{claim}{B.4}{Proved}{Reflection sections, ovals, and punctured ovals}
Let $\sigma$ be an orientation-reversing involution of $U$ with
$\widetilde U\cong\R^2$, and let $E\to C_2$ be its lift extension. There is a
bijection
\[
 \pi_0(U^\sigma)\longrightarrow\operatorname{Sec}(E\to C_2)/\pi.
\]
For the section $s$ associated with a component $Z$, inclusion induces an
isomorphism $\pi_1(Z)\cong C_\pi(s)$. Thus the centralizer is $1$ for a line and
$\Z$ for an oval. This includes punctures fixed or permuted by $\sigma$ and proves
\S33, folios 222--225, scans 440--445, in the required noncompact cases.
\end{claim}
\passage{33; 222-225; 440-445}

\begin{proof}
Each order-two lift $s$ is a reflection of the plane by B.1. Its fixed set
$L$ is a line. The restriction $p:L\to U^\sigma$ is locally a homeomorphism
onto the real locus, by the local model in B.1. Its image is the entire component
$Z$ containing any of its points. To verify this without a properness assumption,
lift a path in $Z$ starting at $p(\widetilde x)$: applying $s$ gives another lift
with the same initial point, so uniqueness of path lifting fixes the lifted path.
The same argument for a homotopy shows that $L\to Z$ is its universal covering.
Every section has such a fixed line, so every section class arises. Two fixed
lines projecting to the same component contain lifts of the same point; a deck
transformation taking one lift to the other conjugates the associated involutions.
The two lifts fixing the resulting point agree, because lifts of a specified
base homeomorphism are unique once the value at one point is fixed.

A deck transformation preserving $L$ normalizes $\langle s\rangle$ and hence
centralizes $s$. Conversely a deck transformation centralizing $s$ preserves
$L$. Therefore the deck group of $L\to Z$ is exactly $C_\pi(s)$. Since $L\cong\R$,
this proves both the asserted isomorphism and injectivity of $\pi_1(Z)\to\pi$.
A finite-type real locus is a finite union of circles and open intervals after
its real punctures are removed. Each interval is homeomorphic to $\R$.
This argument is also the explicit reduction of Mochizuki's Theorem 3.6 and
Corollary 3.7 \cite{bMochizuki2003} to Grothendieck's formulation: his full
centralizer is $\pi_1(Z)\times C_2$, and intersecting it with $\pi$ gives the
centralizer used here.
\end{proof}

For $G=1$ the section set has one element and its centralizer is $\pi$; the
comparison is the bijection $\pi_0(U)\to\{*\}$. The map from individual points
is not injective. The hypothesis on the universal cover has content: on $S^2$
a nontrivial rotation has two fixed points while its lift extension, with kernel
$1$, has one section; the antipodal involution has no fixed point and still has
one section. These examples explain exactly why the same formula does not extend
to that excluded surface.

\begin{claim}{B.5}{Proved}{Recovering the incidence and cyclic order of real punctures}
For a hyperbolic $U=X\setminus S$ with anti-involution $\sigma$, its discrete
extension together with its indexed peripheral subgroups determines the following
finite graph. Its vertices are the section classes of B.4, its edges are the real
punctures, and each edge is attached to its two adjacent real half-branches.
Every punctured oval gives one cycle, allowing a loop or two parallel edges.
The graph recovers which punctures lie on each oval and their cyclic order up
to reversal. Unpunctured ovals are the vertices with infinite cyclic centralizer.
The claim uses the extension and peripheral incidence maps; a naked set of
section classes alone has no such incidence information.
\end{claim}
\passage{33; 221-225; 437,440-445}

\begin{proof}
A puncture is real exactly when its $C_2$-stabilizer is all of $C_2$. For such a
puncture its peripheral extension is
\[
 E_i=\langle t,s\mid s^2=1,;sts^{-1}=t^{-1}\rangle,
 \qquad L_i=\langle t\rangle.
\]
Indeed the local anti-involution of the filled disc is reflection. On the
logarithmic covering $z=\exp(2\pi i w)$ of the punctured disc its two types of
order-two lifts are $w\mapsto-\overline w$ and
$w\mapsto1-\overline w$. Algebraically the reflections $t^k s$ have two
$L_i$-conjugacy classes, distinguished by the parity of $k$, since conjugation
by $t^m$ adds $2m$. The two inclusions of these classes into the section set of
$E$ are precisely the two real half-branches. They can have the same image.
Attach the edge $i$ using these two maps. Walking through the intervals of a
punctured circle gives a cycle whose edge labels are its punctures. Conversely
these cycles determine the partition of the punctures and their circular order,
with the orientation of each circle unchosen. B.4 identifies every remaining
vertex with an unpunctured oval. A nonreal pair of punctures has trivial
stabilizer and supplies no edge. Finally any two configurations on circles with
the same circular order are related by a homeomorphism: map successive closed
arcs monotonically, and extend over collars. Thus this recovery specifies the
topological position of the punctures, including the single-puncture and
single-interval case, without specifying their analytic positions.
\end{proof}

\subsection{Finite automorphisms and the withdrawn assertions}

\begin{claim}{B.6}{Proved}{The fundamental loop lemma of \S18}
Let $\pi=\pi_1(U)$ be hyperbolic and discrete, with its indexed peripheral
structure and orientation module. A finite-order orientation-preserving
automorphism of $\pi$ fixing a peripheral index and normalizing its subgroup
$L_i$ is the identity. Consequently the peripheral extension attached to a
nontrivial finite-order positive outer automorphism fixing $i$ cannot split.
This proves the lemma and corollary of scans 193--194 without adding realizability
to the hypotheses.
\end{claim}
\passage{18; 96-98; 193-196}

\begin{proof}
The based Dehn--Nielsen--Baer correspondence identifies the given finite cyclic
group with a finite subgroup of the mapping class group with one additional
marked base point. Kerckhoff's realization theorem realizes this group by
homeomorphisms fixing that point and preserving the marked punctures. The
original paper explicitly treats cusps in \S IV, pp.~255--256, in addition to
the orientation-reversing convention on p.~235 \cite{C_Kerckhoff1983}; hence the
punctured case used here is included. The section determined by the fixed base
point identifies the induced genuine action on $\pi$ with the given one.

An orientation-preserving automorphism fixing $i$ acts as $+1$ on $L_i\cong\Z$,
so it fixes a generator $\ell$ of $L_i$. A nontrivial finite positive cyclic group
of plane lifts has one fixed point by B.1. Commutation with the nontrivial deck
transformation $\ell$ would give the second fixed point $\ell\widetilde x$,
a contradiction. Thus the automorphism is the identity. A splitting of the
peripheral extension gives a finite-order lift normalizing $L_i$; the result
just proved makes its outer class trivial.

The margin on scan 194 writes the \emph{filled} type as $(g,\nu)$ and the
acting group as type $(g,\nu+1)$. For $(g,\nu)=(0,2)$ or $(1,0)$ the acting
group is still a hyperbolic free group of rank two, so the same proof applies.
For a literal annulus, a positive automorphism fixing one of its two indexed
ends acts trivially on $\pi=\Z$. Forgetting the indices permits the positive
end interchange, which induces $-1$ on $\Z$ while normalizing its common
peripheral subgroup. For a literal unpunctured torus there is no peripheral
index, so the displayed loop lemma has no instance. These are the precise
comparisons required by the marginal remark, also reflected in A.13.
\end{proof}

\begin{claim}{B.7}{Refuted}{The unrestricted fixed-subgroup question, with its replacement}
The assertion $\pi^G=1$ for every finite $G\le\Aut_{\rm lac}(\pi)$ is false.
For a hyperbolic discrete $\pi$, a nontrivial finite genuine loop action has
fixed subgroup $1$ when its positive part is nontrivial. Its remaining possibility
is a reflection action of $C_2$; its fixed subgroup is $1$ or $\Z$, as in B.4.
For any finite-type faithful geometric action there are finitely many classes
of full and partial sections. This answers questions (c) and (e), scans 197--198.
\end{claim}
\passage{18; 98-99; 197-198}

\begin{proof}
A finite subgroup of $\Aut_{\rm lac}(\pi)$ intersects $\Inn(\pi)$ trivially:
$\pi$ is torsion-free and centre-free. Its outer image is therefore finite and
is realized by Kerckhoff. The canonical centre-free extension identifies the
given genuine action with a section of its geometric lift extension. B.3 and
B.4 now give the complete list of fixed subgroups. An explicit counterexample
is the double of an orientable one-holed torus, a genus-two surface, with the
reflection interchanging the two halves. Its fixed oval is essential and a
base point on it gives a genuine order-two automorphism with fixed subgroup
$\Z$.

For finiteness, compactify $U$ by its finitely many ends. Every nonidentity
positive periodic homeomorphism has isolated fixed points by B.1. The resulting
closed discrete subset of the compact surface is finite. A reflection has a
compact one-dimensional fixed submanifold, hence finitely many components;
removing finitely many punctures preserves finiteness. Apply B.3 or B.4 to each
of the finitely many subgroups of $G$. This proves the asserted finiteness.
For the abelian groups the genuine actions can be read directly: a nonidentity
positive finite-order element of $SL_2(\Z)$ fixes no nonzero lattice vector,
whereas a determinant $-1$ involution fixes a primitive rank-one sublattice.
For $\Z$, the nonidentity automorphism is $-1$ and its fixed subgroup is zero.
The trivial group of operators fixes all of $\pi$ in every type.
\end{proof}

\begin{claim}{B.8}{Refuted}{Unique maximal sections fail for reflections, already in type $(0,3)$}
The orientation-free extension of question (d), scan 198, and of the expectation
on scan 155 is false. Put
\[
 E=\langle s_0,s_1,s_2\mid s_j^2=1, (s_0s_1)^2=1, (s_1s_2)^3=1\rangle.
\]
There is a geometric extension $1\to\pi_{0,3}\to E\to S_3\times C_2\to1$.
The reflection subgroup $\langle s_1\rangle$ is contained in two distinct maximal
finite subgroups $\langle s_0,s_1\rangle\cong C_2^2$ and
$\langle s_1,s_2\rangle\cong S_3$, which are not $\pi_{0,3}$-conjugate.
The valid uniqueness statement is B.3, requiring a nontrivial positive part.
\end{claim}
\passage{15; 77; 155; 18; 99-101; 198,201-203}

\begin{proof}
Take a hyperbolic triangle with angles $\pi/2,\pi/3,0$ and reflect in its three
sides. Reflected triangles tile the plane; the local tilings at the two finite
vertices give exactly the displayed relations. Equivalently its group is the
amalgam $C_2^2*_{C_2}S_3$, with common subgroup generated by $s_1$.
Write $S_3=\langle r,a\mid r^3=a^2=1, ara=r^{-1}\rangle$, and let $\tau$ generate
the central $C_2$. Define the surjection
\[
 s_0\mapsto\tau,\qquad s_1\mapsto a\tau,\qquad s_2\mapsto ra\tau.
\]
It is injective on both vertex groups. A finite subgroup of this amalgam fixes
a vertex of its Bass--Serre tree: the finite orbit spans a finite tree, whose
centre is fixed (there are no edge inversions). Thus the kernel is torsion-free.
The character sending each $s_j$ to $-1$ factors through the exponent of $\tau$,
so the kernel preserves orientation. The group Euler characteristic is
$1/4+1/6-1/2=-1/12$, and the index is $12$; hence the kernel has Euler
characteristic $-1$. The stabilizer of the cusp is $\langle s_0,s_2\rangle$;
its image has order four. The kernel therefore has $12/4=3$ cusps, and its
quotient surface has genus zero. This constructs the asserted geometric extension.

Each vertex group is maximal finite, again by the tree action: a finite overgroup
would fix a vertex, and a vertex group of order greater than the edge group fixes
only its vertex. Their intersection is $\langle s_1\rangle$ by the normal form
for the amalgam. Their orders differ, and their generated group is all of the
infinite amalgam. This proves every claimed failure. The finite quotient, its
vertex and cusp images, and the Euler-characteristic calculation are checked in
\texttt{code/check\_part\_b.py}. Grothendieck's proof on scans 201--203 applies
when the finite overgroups are positive cyclic groups: they centralize their
intersection, and B.3 makes the kernel of that centralizer map trivial. A
reflection has a line of fixed points, so that argument supplies neither of
these requirements in the generalization just refuted.
\end{proof}

\begin{claim}{B.9}{Proved}{The valid normalizer assertion}
For a hyperbolic discrete group with loops, a nontrivial finite subgroup
$K\le\Aut^+_{\rm lac}(\pi)$ is cyclic and
\[
 N_{\Aut^+_{\rm lac}(\pi)}(K)=C_{\Aut^+_{\rm lac}(\pi)}(K).
\]
An orientation-reversing element of its full normalizer acts on $K$ by inversion.
Every finite positive overgroup of $K$ is cyclic. These are the valid assertions
within the passage labelled ``Ultra fort!'' on scans 198--199.
\end{claim}
\passage{18; 99-100; 198-200}

\begin{proof}
Realize the outer image of $K$ and use its given section. B.3 makes $K$ cyclic.
The boundary rotation character of B.12 is injective on $K$ and is natural for
loop automorphisms with their action on the orientation module. A positive
normalizer therefore preserves every value of that character, so fixes every
element of $K$. A negative normalizer changes its sign, so inverts every element.
The same realization argument proves cyclicity of each finite positive overgroup.
This proof compares a normalizing automorphism on the canonical oriented boundary;
it makes no assertion that an infinite normalizer acts as a finite isometry group.
\end{proof}

\begin{claim}{B.10}{Refuted}{The withdrawn uniqueness of topological roots}
The lemma marked ``faux!'' on scan 200 is false even for the punctured torus:
two different finite-order homeomorphisms fixing the base point can have the
same nontrivial square and induce the same automorphism of $\pi_1$.
\end{claim}
\passage{18; 100; 200}

\begin{proof}
On $U=\R^2/\Z^2\setminus\{0\}$ choose $p=(1/2,1/2)$ and let
$r(x,y)=(-y,x)$, $h=r^2$. Choose a small disc $D$ whose four $r$-translates are
disjoint and avoid both $0$ and $p$. Choose a nonidentity homeomorphism $f_0$
supported in $D$ and isotopic to the identity there relative to its boundary;
for example use a radial bump rotation in disc coordinates. Put
$f=f_0(hf_0h^{-1})$. Its two disjoint supports show that $fh=hf$; the same
supports show that $fr\ne rf$, since $r$ moves their union to the disjoint pair
$rD,r^3D$. Hence $v=frf^{-1}\ne r$, but $v^2=h$ and $v(p)=p$.
The isotopy of $f$ fixes $0$ and $p$, so $f_*=1$ on $\pi_1(U,p)$ and therefore
$v_*=r_*$. This verifies the exact data of the withdrawn statement. What survives
is uniqueness up to conjugacy by the identity component in the realization
comparison, not equality of the homeomorphisms.
\end{proof}

\begin{claim}{B.11}{Refuted}{Torsion in the centralizer need not be a subgroup}
The assertion following the normalizer assertion on scan 199, that the torsion
elements of the centralizer form one largest finite cyclic overgroup, is false
for the hyperbolic type $(1,1)$.
\end{claim}
\passage{18; 99-100; 198-200}

\begin{proof}
Keep $U,p,r,h$ as in B.10. The essential simple curve
$c=\{x+y=1/2\pmod1\}$ avoids $0,p$ and is invariant under $h$.
The based mapping class $T_c$ commutes with $h$, since conjugation of a positive
Dehn twist by an orientation-preserving map is the positive twist about the
image curve. Thus the two genuine automorphisms induced by $r$ and
$T_crT_c^{-1}$ have order four and common square $h_*$; both centralize $h_*$.
On $H_1(U;\Z)$, with a consistent choice of twist sign, their matrices are
\[
 S=\begin{pmatrix}0&-1\\1&0\end{pmatrix},\qquad
 TST^{-1}=\begin{pmatrix}-2&-5\\1&2\end{pmatrix},\qquad
 T=\begin{pmatrix}2&1\\-1&0\end{pmatrix}.
\]
Both square to $-I$, while their product has trace $-6$ and determinant one.
Its eigenvalues are $-3\pm\sqrt8$, so it has infinite order. The corresponding
product of automorphisms has infinite order as well. These integral equalities
are checked by the supplied program. Consequently the centralizer's torsion
set is not closed under multiplication, and there is no finite subgroup
containing all its finite overgroups of $\langle h_*\rangle$. This coexists with
the exact equality of normalizer and centralizer in B.9.
\end{proof}

\subsection{The multiplier and the abelian surfaces}

\begin{claim}{B.12}{Proved}{An intrinsic discrete formula for the multiplier}
For a geometric finite extension of a hyperbolic discrete surface group,
$E$ with its peripheral structure and orientation module $T$ determines
$\mu_E$ on all torsion elements. On $E^+$ it is the rotation number of its
canonical oriented boundary action. More explicitly, let
$\widetilde{\operatorname{Homeo}}^+(S^1)\to\operatorname{Homeo}^+(S^1)$ be the
central extension by the positively oriented deck translation $z$. For
$u\in E^+$ of order $n$ and any lift $\widetilde u$, write
$\widetilde u^{,n}=z^k$. Then
\[
 \mu_E(u)=\frac{k}{n}\pmod{\Z}\quad\hbox{in }T\otimes\Q/\Z.
\]
For a nonidentity orientation-reversing torsion element we put
$\mu_E(u)$ equal to the unique nonzero element of order two. With this
convention the formula is natural under isomorphisms preserving the loops and
transporting $T$, satisfies $\mu(u^m)=m\mu(u)$, and has the same order as $u$.
This answers the discrete intrinsicity question of \S15, folios 77--79,
scans 155--160.
\end{claim}
\passage{15; 77-79; 155-160}

\begin{proof}
Choose an oriented finite-area hyperbolic uniformization of $U$. Its boundary
at infinity is a circle. We explain the ambiguity in this choice. A marking
isomorphism of surface groups preserving their peripheral subgroups and the
orientation is induced by an orientation-preserving homeomorphism, by the
Dehn--Nielsen--Baer correspondence. We can choose that marking quasiconformal:
on a compact core smooth it, and on each cusp use standard cylindrical
coordinates and extend with bounded dilatation. Its lift extends to an oriented
boundary homeomorphism. This extension is uniquely determined by its
$\pi$-equivariance. Indeed it must carry the attracting and repelling endpoints
of every hyperbolic element to those of the corresponding element. These
endpoints are dense in the boundary circle of a finite-area surface group.
For completeness, density follows from its full limit set: a boundary arc
meets translates of a fundamental ideal polygon arbitrarily near infinity;
conjugating either of two hyperbolic elements with disjoint endpoints by the
corresponding translates gives endpoints in that arc. Thus two equivariant
boundary maps agree on a dense set and hence everywhere.

For each automorphism of $\pi$ preserving the peripheral structure,
Dehn--Nielsen--Baer provides the boundary map conjugating its deck action.
The preceding uniqueness makes composition exact and makes this circle,
together with its $E$-action, independent of the uniformization up to a unique
$\pi$-equivariant boundary isomorphism. The boundary orientation is recorded by
$T$, so changing orientation acts by $-1$. Realize the finite action and choose
an invariant conformal structure, using B.1's cyclic and dihedral local charts
on the quotient. Its hyperbolic metric makes a finite-order positive lift an
elliptic isometry. In the disc model centred at its fixed point it is the same
rotation on the tangent circle and on the circle at infinity. The displayed
formula therefore gives its local multiplier.

Replacing $\widetilde u$ by $z^j\widetilde u$ replaces $k$ by $k+nj$, so the
value is independent of the lift. An orientation-preserving finite-order circle
homeomorphism of order $n$ is conjugate to a rotation by $k/n$ with
$\gcd(k,n)=1$; this also follows by ordering the finite orbit on the circle
and extending a homeomorphism of one fundamental interval around that orbit.
This proves the order and power statements. A negative finite plane isometry
is a reflection, of order two, so its assigned value has exactly the required
power property and is unchanged by sign reversal.

The normalizer formula suggested in the margin of scan 159 is recovered as
follows. Drill the orbit of the fixed point of $u$. Above the new puncture the
positive peripheral normalizer is infinite cyclic, generated by a positively
oriented element $v$, with peripheral subgroup generated by $v^n$. A lift of
$u$ to this normalizer is $v^a$ modulo $\langle v^n\rangle$, and
$\mu(u)=a/n$. In a local coordinate this is the equality between the rotation
$z\mapsto e^{2\pi ia/n}z$ and monodromy in $z\mapsto z^n$. Harvey's Theorem 7
\cite{bHarvey1971} computes exactly this inverse-residue relation from the
elliptic generators of the orbifold epimorphism. His convention uses an inverse
deck transformation, accounting for the sign visible in his proof on
pp.~393--394. Our local-coordinate calculation also applies to cusps and to the
positive cyclic part of a dihedral stabilizer; B.1 supplies the reflection part.
\end{proof}

\begin{claim}{B.13}{Refuted}{Forgetting loops can change the multiplier}
An abstract extension, with its orientation line regarded as an additional
rank-one module, need not determine the multiplier after the peripheral
structure is forgotten. There are isomorphic extensions with kernel a free
group of rank two and quotient $C_3$ whose geometric multipliers disagree,
even after allowing either isomorphism between their orientation lines.
\end{claim}
\passage{15; 77-79; 155-160}

\begin{proof}
Let $E_0=\langle u,v\mid u^3=v^3=1\rangle=C_3*C_3$. Realize it as the
positive orbifold group of a sphere with one cusp and two cone points of order
three; choose $u,v$ as its positive local generators. The epimorphisms
\[
 q_+(u)=q_+(v)=1,\qquad q_-(u)=1,\quad q_-(v)=-1
 \quad (\text{values in }\Z/3)
\]
are injective on every finite subgroup, because finite subgroups of a free
product are conjugate into its factors (the tree argument in B.8 applies).
Their kernels are torsion-free of index three and Euler characteristic $-1$.
The cusp generator is $(uv)^{-1}$. Its image under $q_+$ is nonzero, so there
is one cusp and the kernel is a once-punctured torus group. Its image under
$q_-$ is zero, so there are three cusps and the kernel is a pair-of-pants group.
Both have free rank two.

The automorphism $\psi(u)=u,\ \psi(v)=v^{-1}$ satisfies $q_-\psi=q_+$ and
is an isomorphism of the two abstract extensions. The positive local generator
values in each realization are $\mu(u)=\mu(v)=1/3$. Thus $\psi$ preserves the
value at $u$ but changes the value at $v$ to $-1/3$. Reversing the isomorphism
of orientation lines would instead break the equality at $u$. No orientation-line
identification repairs both. The isomorphism sends the cusp word $uv$ to
$uv^{-1}$, displaying precisely which peripheral comparison has been discarded.
The finite monodromy calculation is included in the computational record.
\end{proof}

\begin{claim}{B.14}{Proved}{Finite annulus and torus actions, and their multipliers}
Every finite topological action on the torus is topologically conjugate to an
affine action on $V/\Lambda$, where $\Lambda\cong\Z^2$ and
$V=\Lambda\otimes\R$. For a fixed finite operator group $G$ such an action is
specified by a homomorphism $A:G\to GL(\Lambda)$ and a cocycle
$b\in Z^1(G,V/\Lambda)$, subject to faithfulness of
$g\mapsto(x\mapsto A_gx+b_g)$. Translation conjugacy corresponds exactly to a
coboundary, and
\[
 H^1(G,V/\Lambda)\xrightarrow{\sim}H^2(G,\Lambda)
\]
is the connecting homomorphism. Linear changes of the lattice give the further
conjugacies. Every finite annulus action has a simultaneous normal form
\[
 (x,t)\longmapsto(\epsilon_g x+a_g,\eta_g t),
 \qquad x\in\R/\Z,\quad t\in\R,\quad
 \epsilon_g,\eta_g\in\{1,-1\}.
\]
These descriptions determine the multiplier of every finite-order element in
the respective lift extensions.
\end{claim}
\passage{15; 78-79; 157-159; 17; 89-90; 179-181}

\begin{proof}
We first justify simultaneous conformalization. B.1 gives genuine cyclic or
dihedral charts for every finite stabilizer. The quotient is a compact
surface orbifold after the ends are added; its interior special points have
cyclic charts, and its mirror edges and corners have reflection and dihedral
charts. Choose a smooth metric on the quotient away from these charts, choose
round cone and mirror metrics in them, and patch in collars using an invariant
partition of unity on their finite local covers. The local charts make the
patched lifts smooth. Pullback is a smooth invariant conformal structure on
the original surface. This construction proves the required passage from the
original topological action; it does not average a merely topological metric
and call the result Riemannian.

Uniformization of the torus gives $\C/\Lambda$. Every holomorphic automorphism
lifts to an entire map $f$ with $f(z+\lambda)-f(z)\in\Lambda$. Its derivative
is periodic and hence bounded on $\C$, so it is constant; $f(z)=az+c$.
An antiholomorphic automorphism is obtained by composing such a map with
complex conjugation on its source. Hence every element is real affine on the
same lattice. The action law is precisely
$b_{gh}=b_g+A_gb_h$. Averaging cochains over the finite group gives
$H^q(G,V)=0$ for $q>0$. The exact sequence
$0\to\Lambda\to V\to V/\Lambda\to0$ therefore gives the asserted isomorphism.
In concrete terms, lifts $v_g\in V$ of $b_g$ give the extension cocycle
$v_g+A_gv_h-v_{gh}\in\Lambda$. Conjugating by translation adds a coboundary.
Moreover any topological conjugacy between two affine actions can be replaced
by an affine conjugacy. A lift of the conjugacy has form $Bx+f(x)$, with
$B\in GL_2(\Z)$ and $f$ periodic. Average its equivariance equation over the
torus with Haar measure. Affine lattice automorphisms preserve that measure,
so $Bx+\int f$ satisfies the same equation modulo $\Lambda$.

For the annulus, extend the action to its two-point sphere compactification
and apply Kolev. In the resulting conformal sphere choose a M\"obius coordinate
sending the two ends to zero and infinity. The maps are then of the forms
$z\mapsto az$, $a/z$, $a\overline z$, or $a/\overline z$. On the logarithmic
radial coordinate they act affinely with slopes $\eta=\pm1$. A finite group
of affine transformations of the real line fixes the average of any orbit;
translation to this common fixed point removes every radial constant. The
angular coordinate gives the claimed $\epsilon x+a$. This proves the displayed
simultaneous normal form. The action on $\pi_1$ is $\epsilon_g$ and its
orientation sign is $\epsilon_g\eta_g$.

For the promised computation on the torus, an element of the lift extension
is $x\mapsto Ax+v$. It has finite order exactly when $A$ has finite order,
say $n$, and $(1+A+\cdots+A^{n-1})v=0$. Averaging its orbit gives a fixed point,
and translation to that point conjugates it to $A$. A nonidentity positive
finite-order matrix has trace $-2,-1,0,1$, giving orders $2,3,4,6$, respectively.
For order two its multiplier is $1/2$. For the other orders it is $+1/n$ or
$-1/n$ according to the sign of $\det(w,Aw)$ for any nonzero $w$; this sign
is independent of $w$, as is seen after averaging an inner product and writing
$A$ as a rotation. A negative finite-order matrix is a reflection and receives
$1/2$. A lift with $A=I$ has finite order only when it is the identity.
For the annulus a lift is $(x,t)\mapsto(\epsilon x+a+k,\eta t)$. With
$\epsilon=-1$ its square is the identity; with $\epsilon=1$ finite order
requires $a+k=0$, leaving the identity or $t\mapsto-t$. Thus every nonidentity
torsion lift on an annulus has order two and multiplier $1/2$.
\end{proof}

\begin{claim}{B.15}{Refuted}{The ``only counterexample'' of scan 179}
The assertion that a faithful homotopically trivial finite annulus action
failing to be a translation action must be the single anti-involution
$z\mapsto1/\overline z$ is false. For the two non-simply-connected abelian
types, the exact replacement is: homotopically trivial torus actions are
translation actions; homotopically trivial annulus actions are the finite
subgroups of $SO(2)\times C_2$ in the normal form of B.14. Those contained in
$SO(2)$ are translations, while those with nontrivial second projection include
two different types of anti-involution and groups of arbitrarily large order.
\end{claim}
\passage{17; 89-90; 179-181}

\begin{proof}
On the torus, homotopic triviality is $A_g=I$ for every $g$, so B.14 proves the
translation statement and the implication doubted on scans 180--181.
On the annulus it is $\epsilon_g=1$, which gives the asserted subgroup.
For each $n\ge1$ the maps
\[
 r(z)=e^{2\pi i/n}z,\qquad s(z)=1/\overline z
\]
generate a faithful $C_n\times C_2$ action, and every element induces the
identity on $\pi_1(\C^*)$. The fixed circle of $s$ excludes a translation
realization: a nonidentity translation in any group is fixed-point free.
In addition $s_-(z)=-1/\overline z$ is an involution inducing $+1$ on $\pi_1$,
but it has no fixed point, since its fixed-point equation is $|z|^2=-1$.
It cannot be conjugate to $s$, and it cannot be a translation of a connected
topological group: every such translation is isotopic to the identity along
a path in that group and hence preserves orientation. These formulas give the required full
counterexamples, not merely a gap in the argument.

More explicitly let the rotation subgroup have order $n$ and generator
$(x,t)\mapsto(x+1/n,t)$. A second-coset generator can be chosen as
$(x,t)\mapsto(x+m/(2n),-t)$, with $m=0$ or $1$ after multiplying by a rotation.
Its square is the $m$th rotation. This exhausts the finite subgroups projecting
onto $C_2$, and includes both the reflection and glide cases. For the omitted
simply connected surfaces, finite plane actions are the orthogonal actions of
B.1 and are all homotopically trivial; positive sphere actions are likewise
homotopic to the identity. Nontrivial rotations fixing a point already refute
any extension of the translation formulation to those types.
\end{proof}

\subsection{Anti-involutions and the later tube construction}

\begin{claim}{B.16}{Proved}{Complete topological classification of anti-involutions}
Conjugacy classes of anti-involutions of a closed oriented genus-$g$ surface
are classified by the orientability of the quotient and its number $b$ of
boundary components. The possibilities are
\[
 \begin{array}{ll}
 \text{nonorientable quotient:}&0\le b\le g,\quad
                  h=g+1-b\ge1\text{ crosscaps};\\
 \text{orientable quotient:}&b\ge1,\quad g=2h+b-1,\quad h\ge0.
 \end{array}
\]
There are $(g+1)+\lfloor(g+2)/2\rfloor=\lfloor(3g+4)/2\rfloor$ classes.
For $\nu$ unlabelled invariant punctures the additional data are an integer
$r\ge0$ and an unordered list $(q_1,\ldots,q_b)$ of nonnegative integers with
$2r+\sum q_i=\nu$: $r$ is the number of interior punctures of the quotient,
and $q_i$ counts punctures on its $i$th boundary circle.
\end{claim}
\passage{33; 221-222; 437-440}

\begin{proof}
By B.1 the quotient $Y=X/\sigma$ is a compact surface with boundary
$X^\sigma$. The complement of the real locus covers its interior with degree
two; the Euler characteristic of the boundary is zero. Hence
$\chi(X)=2\chi(Y)$. If $Y$ is nonorientable with $h$ crosscaps,
$\chi(Y)=2-h-b$, yielding $g=h+b-1$; if orientable of genus $h$ it yields
$g=2h+b-1$. In the latter case $b$ must be positive, since the oriented double
of a closed orientable quotient would be disconnected. Conversely form the
orientation double over the interior of each allowed $Y$ and join the two
local sides at its boundary. The result is a connected closed orientable surface,
with the deck interchange extended to reflection along the boundary; its genus
is given by the displayed formula. The quotient recovers $Y$, so the classification
of compact surfaces proves both existence and uniqueness of these conjugacy classes.
This also proves the complete list underlying Gross--Harris, Proposition 3.1
\cite{bGrossHarris1981}, including the two corrections in the margins of scan 439.
Grothendieck's nonorientable parameter $\gamma$ is $h-1$; thus his formula
$g=\gamma+b$ agrees exactly with ours.

An anti-involution of a punctured finite-type surface extends to its end
compactification. Punctures in its quotient lie either in the interior,
with two-point inverse image, or on the boundary, with one-point inverse image.
Any two finite interior configurations of the same cardinality are related by
an isotopy supported away from the boundary. Boundary components may be
permuted, and any two unlabelled configurations of the same cardinality on one
circle are related by a circle homeomorphism, extended over a collar. These
facts prove completeness of the punctured invariants. With individually
labelled punctures one must instead retain their assignments and their circular
orders; B.5 provides the corresponding incidence reconstruction.
\end{proof}

\begin{claim}{B.17}{Proved}{Real moduli components, with the group and marking convention fixed}
Let $\mathcal M_{g,\nu}$ mean the real moduli stack of smooth proper
geometrically connected genus-$g$ curves with an \emph{unlabelled}
degree-$\nu$ reduced puncture divisor. We retain the stack and its
real forms also in the four unstable types. Its complex
mapping class group is $\mathcal T^+_{g,\nu}$, allowing permutations of the
punctures. Write $p_b(N)$ for the number of partitions of $N$ into at most
$b$ parts, with $p_0(0)=1$ and $p_0(N)=0$ for $N>0$. Then
\[
 \#\pi_0\bigl(\mathcal M_{g,\nu}(\R)\bigr)
 =\sum_{b=0}^{g}\ \sum_{r=0}^{\lfloor\nu/2\rfloor}p_b(\nu-2r)
  +\sum_{\substack{1\le b\le g+1\\ b\equiv g+1\ (2)}}
       \ \sum_{r=0}^{\lfloor\nu/2\rfloor}p_b(\nu-2r).
\]
For closed hyperbolic surfaces this is $\lfloor(3g+4)/2\rfloor$.
In the four unstable types $(0,0)$, $(0,1)$, $(0,2)$ and $(1,0)$ the counts are,
respectively, $2,1,3,3$. In the last case no origin is prescribed.
Distinct topological anti-involution classes remain distinct in the extended
mapping class group in these hyperbolic types. The last assertion fails for
$(0,0)$, $(0,2)$, and $(1,0)$.
\end{claim}
\passage{33; 222; 439-440}

\begin{proof}
C.5 and C.6, finite-action realization and uniqueness up to the identity component, identify
conjugacy classes of anti-involutions on the fixed punctured surface with
conjugacy classes of orientation-reversing involutions in
$\mathcal T_{g,\nu}$. One can use positive conjugacy here: an orientation-reversing
conjugator can be multiplied by the target involution to become positive.
Mochizuki, Corollary 3.8, identifies these conjugacy classes with the connected
components of the real moduli \emph{stack} \cite{bMochizuki2003}. His Example 3.3
uses the full modular group of a punctured surface, which is the unlabelled
convention just specified. B.16 supplies the list. For fixed $b,r$, unordered
nonnegative lists of length $b$ summing to $\nu-2r$ are exactly the partitions
counted by $p_b(\nu-2r)$. Summing over the two quotient families proves the
formula. The supplied code enumerates these numbers for $2g+\nu\le6$ and checks
the closed formula for $0\le g\le20$; the proof holds for all $g$.

We verify the unstable component counts without using mapping-class
injectivity. A real genus-zero curve is a split or nonsplit real conic;
analytically these are the sphere with $z\mapsto\overline z$ or with
$z\mapsto-1/\overline z$. To see these are the two possibilities,
conformalize its anti-involution as in B.14 and apply the two quotient
types in B.16; uniformization identifies the resulting conformal
double with the sphere. A real point exists exactly in the first type.
There is no complex moduli parameter in genus zero. With one puncture
the point is real, and $PGL_2(\R)$ acts transitively on the real
projective line, giving one component. With two unlabelled punctures
the possibilities are two real points on the split conic, one
conjugate pair on the split conic, or one conjugate pair on the
nonsplit conic. The first is transitive under $PGL_2(\R)$, the second
under its action on the upper half-plane, and the third under the
rotations commuting with the antipodal map. These give three
components. With no puncture there are the two conic components.

For an unpointed genus-one curve, B.14 and B.16 give the three
anti-involution types with zero, one, or two real ovals. Each type
has a connected parameter space of complex structures. Explicit
representatives are
\[
 \begin{array}{c|c|c}
 b&\Lambda_t&\sigma(z)\\\hline
 0&\Z+it\Z&\overline z+\tfrac12\\
 1&\Z+(\tfrac12+it)\Z&\overline z\\
 2&\Z+it\Z&\overline z
 \end{array}\qquad(t>0).
\]
To obtain this list from an invariant lattice, rotate a primitive
real period to $1$. Conjugation of a second period forces its real
part modulo $\Z$ to be $0$ or $1/2$. The translation cocycle of
B.14 is then removed by a change of origin, except for the
nonzero half-period class in the rectangular case. That class gives
the first row. Thus every real genus-one curve occurs in one of
these connected families. The number of real ovals is locally
constant in a smooth proper real family: local trivialization can
be chosen invariant under the involution by an invariant horizontal
metric and its path lifting. Distinct rows consequently belong to
different components. This proves the number three, with no origin
added to the moduli problem. Direct substitution of the four
unstable types into the displayed partition formula gives exactly
these counts.

The mapping-class exceptions can be checked directly. On $S^2$ reflection and the antipodal
map have different fixed sets, but both represent the orientation-reversing
mapping class. On $\R^2/\Z^2$ the maps $(x,y)\mapsto(x,-y)$ and
$(x,y)\mapsto(x+1/2,-y)$ have respectively two fixed circles and no fixed
points, but both induce $\operatorname{diag}(1,-1)$ and hence the same mapping
class. On the annulus, the reflection and glide anti-involutions of B.15 also have
the same mapping class: the path $(x,t)\mapsto(x+s/2,-t)$ joins them through
homeomorphisms preserving the two ends as a set. These examples establish the exact failure of the unrestricted
injectivity hoped for on scan 440. For labelled moduli spaces or a coarse
moduli space, the quotient and its real locus have to be specified separately;
the displayed formula is for the stack and puncture convention stated above.
\end{proof}

\begin{claim}{B.18}{Proved}{The two centralizer corollaries in both quotient families}
Let $g\ge2$, and let an anti-involution of $X_g$ have $b>0$ ovals, with outer
class $\sigma$. There are exactly $b$ classes of order-two genuine
loop automorphisms lifting $\sigma$, under conjugacy by $\pi_g$.
For each lift $u$, $\pi_g^u\cong\Z$. The positive centralizer of $\sigma$ acts
transitively on those $b$ classes. More precisely, writing
$A=\Aut_{\rm lac}(\pi_g)$ and $O=\Out_{\rm lac}(\pi_g)$, there is an exact sequence
\[
 1\longrightarrow\pi_g^u\longrightarrow C_{A^+}(u)
 \longrightarrow\operatorname{Stab}_{C_{O^+}(\sigma)}([u])\longrightarrow1,
\]
and the last group has index $b$ in $C_{O^+}(\sigma)$. The corresponding full
centralizer image also has index $b$. This proves Corollaries 1 and 2,
folios 225--226, scans 445--448, for both orientable and nonorientable quotients.
\end{claim}
\passage{33; 225-226; 445-448}

\begin{proof}
B.4 proves the count and the fixed subgroup. Any permutation of the boundary
components of a compact connected quotient surface is realized by a homeomorphism:
join two components by an embedded arc, interchange them in a neighbourhood of
that arc and the two collars, and use transpositions. Each quotient homeomorphism
lifts to the oriented double; among its two lifts choose the orientation-preserving
one. It commutes with the reflection. Consequently the positive centralizer is
transitive on the ovals, and hence on the section classes.

For exactness, the centre of $\pi_g$ is trivial, so $\Inn(\pi_g)$ is identified
with $\pi_g$. Its intersection with $C_A(u)$ is exactly $\pi_g^u$.
An outer class $\alpha$ centralizing $\sigma$ preserves $[u]$ precisely when
one, hence every, lift $\widetilde\alpha$ satisfies
$\widetilde\alpha u\widetilde\alpha^{-1}=vuv^{-1}$ for an inner automorphism
$v$. Then $v^{-1}\widetilde\alpha$ centralizes $u$, proving surjectivity onto
the displayed stabilizer. Orbit--stabilizer gives index $b$ in both versions.
For $b=0$ there are no order-two genuine lifts and these assertions about a
chosen lift have no instance. With punctures the same exact sequence holds,
but the orbit size records the boundary decorations from B.16; it need not be
the total number of ovals when different ovals carry different puncture data.
\end{proof}

\begin{claim}{B.19}{Proved}{The tube twists, their maximality property, and their normalizer}
For $g\ge2$, the elements $\rho_j$ of \S33, scans 453--470, are full Dehn twists
about a cut system $c_1,\ldots,c_g$: disjoint nonseparating curves with connected
complement. They generate $H\cong\Z^g$. This is not a maximal free abelian
subgroup of $\mathcal T_g$: it is contained in a multitwist subgroup of rank
$3g-3$. It has the precise maximality property of being a full multitwist group
of largest rank whose supporting multicurve has connected complement.
Furthermore
\[
 N_{\mathcal T_g}(H)=\operatorname{Stab}_{\mathcal T_g}(\{c_1,\ldots,c_g\}),
 \qquad N_{\mathcal T_g}(H)/C_{\mathcal T_g}(H)\cong S_g\times C_2.
\]
The action on $H$ permutes its basis and multiplies every basis vector by the
same orientation sign. Cutting and capping gives an exact sequence
\[
 1\longrightarrow H\longrightarrow N_{\mathcal T_g}(H)
 \longrightarrow\mathcal T^{\rm pairs}_{0,2g}\longrightarrow1,
\]
where the last extended sphere mapping class group preserves the partition of
its $2g$ punctures into the $g$ unordered gluing pairs.
\end{claim}
\passage{33; 229-237; 453-470}

\begin{proof}
In the coordinates of scan 461 the map is
$\rho_j(z,t)=(-z e^{\pi i t},t)$ on $S^1\times[-1,1]$ and is the identity outside
the tube. Its angular displacement increases from $0$ to $2\pi$, so its relative
boundary mapping class is one full Dehn twist, with sign fixed by that coordinate
orientation. The complement of the tubes in his construction is a sphere with
$2g$ boundary circles, establishing the cut-system assertion.

Twists in disjoint collars commute. They are independent: cut the surface along
the $c_j$, choose transverse curves dual to these cuts, and measure the integer
winding of the resulting arcs when the collars are glued back. A product
$\prod T_{c_j}^{n_j}$ changes the $j$th such winding by $n_j$; an isotopy cannot
change these relative-end winding differences. This proves $H\cong\Z^g$.
Complete the $2g$-holed sphere to a pants decomposition by adding $2g-3$
essential disjoint curves. The same winding argument gives a rank-$3g-3$
multitwist group containing $H$. In contrast, cutting $k$ disjoint curves while
keeping the complement connected reduces the genus to $g-k$, so $k\le g$.
This proves the stated maximality among connected-complement multitwist systems.

For the normalizer, conjugation satisfies
$fT_cf^{-1}=T_{f(c)}^{\epsilon(f)}$. A single twist can equal a product of twists
about the $c_j$ only by being one of these twists to the corresponding power.
Here is the needed support argument. Such a twist commutes with every $T_{c_j}$.
The elementary intersection formula
$i(T_d(c),c)=i(c,d)^2$ for a single twist forces $d$ to be disjoint from each
$c_j$; otherwise its conjugation changes $T_{c_j}$. Add $d$ to the disjoint
system if it is not already there and apply the preceding independent-winding
argument. The resulting nontrivial relation is impossible. Thus a normalizer
permutes precisely the $c_j$, proving the first equality. Conversely every such
permutation of curves normalizes $H$. Every permutation of the handles is
realized positively, and a reflection of the standard cut-system model reverses
the global orientation while preserving each curve. This gives the full image
$S_g\times C_2$. Its kernel is exactly the centralizer.

A representative preserving the multicurve can be isotoped to preserve its
collars. Its restriction to the cut surface, after capping each boundary to a
puncture, preserves the gluing pairs. Conversely a homeomorphism of the capped
sphere preserving those pairs can be made radial in small puncture discs and
then extended over the paired collars. Two such extensions differ by integer
collar twists. An isotopy on the capped sphere lifts after allowing exactly
these twists, by the isotopy extension theorem on collars. This proves both
surjectivity and the stated kernel, and types the comparison with the
complement suggested on scan 470.
\end{proof}

\begin{claim}{B.20}{Refuted}{Presentation (52) and its exact correction}
For $g\ge2$, presentation (52) on scan 469 does not present the stated
$\Z^g$-by-finite group. Put $r=\varepsilon_0\varepsilon_1$ and
$\rho_j=r^j\tau'\tau r^{-j}$, $0\le j<g$. The printed presentation gives
\[
 F_g\rtimes\bigl(D_{4g}\times\langle\tau\rangle\bigr),
\]
where $D_{4g}=\langle r,\varepsilon_1\mid r^{2g}=\varepsilon_1^2=1,
\varepsilon_1r\varepsilon_1=r^{-1}\rangle$ permutes the free generators through
its action on $\Z/g$, and $\tau$ inverts each free generator.
Adding the relations $[\rho_i,\rho_j]=1$ for $i<j$ gives exactly
$\Z^g\rtimes(D_{4g}\times C_2)$ with the manuscript's specified action.
\end{claim}
\passage{33; 237; 469}

\begin{proof}
The relations say that $\tau$ commutes with $D_{4g}$, that $\tau'$ commutes
with $\varepsilon_1$ and $r^g$, and that both $\tau,\tau'$ are involutions.
Set $\rho_0=\tau'\tau$. Then $\tau\rho_0\tau=\rho_0^{-1}$,
$\varepsilon_1\rho_0\varepsilon_1=\rho_0$, and
$r^g\rho_0r^{-g}=\rho_0$. Its conjugates are indexed by the $g$ cosets of
$\langle\varepsilon_1,r^g\rangle$ in $D_{4g}$, with no further relation among
them. To check the last assertion directly, form the displayed semidirect
product with a free basis $\rho_0,\ldots,\rho_{g-1}$ and define
$\tau'=\rho_0\tau$ in it. All printed relations hold. The reverse map sends
these free generators to $r^j\tau'\tau r^{-j}$; the two maps agree inversely
on all generators. This is a proof by explicit mutually inverse homomorphisms.
Abelianizing only the free factor adds exactly the stated commutators and gives
the corrected group.

There is also a finite certificate. For $g=2$ set
$\varepsilon_0=(13)$, $\varepsilon_1=\tau=1$, $\tau'=(12)$ in $S_3$.
All nine displayed relators evaluate to $1$, but
$\rho_0=(12)$ and $\rho_1=(23)$ do not commute. The program checks every one
of these evaluations. Relations (33), scan 454, already assert commutation
for the geometric twists; they are exactly the missing input to (52).
\end{proof}

\begin{claim}{B.21}{Proved}{The one-holed-torus subgroups and the proved part of the analogy}
For $g\ge2$, choose a one-holed torus $W_j$ containing $c_j$ and a transverse curve $d_j$
with $i(c_j,d_j)=1$. Its group fixing its boundary pointwise embeds by extension
by the identity into $\mathcal T_g^+$. It is
\[
 B_3=\langle A_j,B_j\mid A_jB_jA_j=B_jA_jB_j\rangle,
\]
and capping the boundary gives
\[
 1\longrightarrow\langle T_{\partial W_j}\rangle\cong\Z
 \longrightarrow B_3\longrightarrow SL_2(\Z)\longrightarrow1,
 \qquad (A_jB_j)^6=T_{\partial W_j}^{\pm1}.
\]
This extension has no section. For disjoint chosen $W_j$, their supported
subgroups commute and, together with $H$, preserve their boundary multicurve;
they do not generate $\mathcal T_g^+$. On explicit generators the tube subgroup
and these two twists act by
\[
 \begin{array}{lll}
 \rho_j:&a_j\mapsto a_j,&b_j\mapsto b_ja_j,\\
 B_j:&a_j\mapsto a_jb_j^{-1},&b_j\mapsto b_j,
 \end{array}
\]
fixing the other handle generators of
$\pi_g=\langle a_j,b_j\mid\prod_j[a_j,b_j]=1\rangle$.
\end{claim}
\passage{33; 236-237; 467-470}

\begin{proof}
Choose the based handle loops so that their commutator is the boundary loop.
The displayed substitutions preserve $[a_j,b_j]$ exactly, and are the effect
of crossing the appropriate twist collar. Twists in curves meeting once obey
the braid relation: in these generators substitution gives $ABA=BAB$.
The sixfold product $(AB)^6$ conjugates both generators by
$[a,b]^{-1}$, which is the boundary twist with this collar convention.
Both word equalities are checked by free reduction in
\texttt{code/free\_words.py}; the stated substitutions also give their direct
verification. Filling the boundary gives the standard two elementary matrices
generating $SL_2(\Z)$. Their action on primitive slopes, reduced by the Euclidean
algorithm, shows that the filled mapping class group is exactly $SL_2(\Z)$.
The kernel of filling is the boundary winding: cut along two based generating
arcs, isotope a kernel representative to fix them, and read its remaining
integer twist in the boundary collar. The kernel is therefore the infinite
cyclic group generated by the boundary twist. The braid group on three strands
has the same central extension: with $x=ABA,y=AB$ it is
$\langle x,y\mid x^2=y^3\rangle$, and its quotient by $y^6$ is $SL_2(\Z)$.
The comparison sends the kernel generator to the boundary twist, so it is an
isomorphism of these extensions.

Extension by the identity is injective because an isotopy to the identity on
the closed surface would fix the essential boundary and, after collar adjustment,
give an isotopy on the one-holed torus; a nonzero boundary twist remains
nontrivial by the winding argument in B.19. The braid group is torsion-free:
$\langle x,y\mid x^2=y^3\rangle$ is the amalgam of two infinite cyclic groups
along their subgroups of indices two and three; a finite subgroup fixes a vertex
in its tree and is conjugate into an infinite cyclic group. A section from
$SL_2(\Z)$ would inject its order-four element, contradicting this fact.

Disjoint supports commute. Each supported homeomorphism fixes the chosen
boundary curves, and so does every generator of $H$. An essential curve crossing
one of those boundaries has a twist which does not preserve it, by the
intersection formula in B.19. Hence the generated subgroup is proper.
The exact analogy proved here is the homomorphism to the usual rank-two
integral representation and its central boundary-winding extension, with braid
and disjoint-support relations. For overlapping choices of the strips one
must specify their intersections before computing the remaining relations;
no unspecified choice of strips supplies additional relations by itself.
\end{proof}

\subsection{Computations and the limits of the requested table}

\begin{claim}{B.22}{Refuted}{A finite subgroup alone does not index the requested abelian table}
For all types in Part B(f), a table indexed only by a finite subgroup of the
mapping class group cannot assign unique splitting numbers or $\Phi$.
The type $(1,0)$ and the single subgroup
$\langle\operatorname{diag}(1,-1)\rangle$ already have two different answers.
The additional extension datum is precisely the class in $H^2(G,\pi)$ in the
abelian cases of B.14.
\end{claim}
\passage{15; 78-79; 157-159; prompt B.f}

\begin{proof}
On the torus compare
$\sigma_0(x,y)=(x,-y)$ and
$\sigma_1(x,y)=(x+1/2,-y)$. Both are faithful involutions with the same
homomorphism to $GL_2(\Z)$, but the first fixes two circles and the second fixes
nothing. The corresponding extensions, with commuting lattice generators
$a,b$, have relations
\[
 \widetilde\sigma a\widetilde\sigma^{-1}=a,\qquad
 \widetilde\sigma b\widetilde\sigma^{-1}=b^{-1},\qquad
 \widetilde\sigma^2=a^\epsilon,\quad \epsilon=0\text{ or }1.
\]
The square of $a^mb^n\widetilde\sigma$ is $a^{2m+\epsilon}$. For $\epsilon=1$
there is no order-two lift. For $\epsilon=0$ such lifts require $m=0$,
and conjugation by $b$ changes $n$ by two. Thus there are exactly two section
classes in the first case and zero in the second. These are the two parity
classes of the extension cocycle. B.14 proves that retaining the extension
class gives the correct affine classification, with all translation conjugacies
accounted for. Thus the counterexample identifies the missing index and its
exact replacement.
\end{proof}

\begin{claim}{B.23}{Proved}{Complete finite-subgroup table in type $(0,3)$}
The ten conjugacy classes in the extended group
$\mathcal T_{0,3}\cong S_3\times C_2$, and their discrete section invariants,
are the following. In the table $F$ counts full-section classes, $P$ counts all
nontrivial partial-section classes with their subgroup images distinguished,
$M=|\Phi|$ counts maximal partial-section classes, and $O$ counts the group
orbits on $\Phi$.
\[
\begin{array}{lrrrrr}
 G&|G|&F&P&M&O\\\hline
 1&1&1&0&0&0\\
 C_2\text{ (positive)}&2&1&1&1&1\\
 C_2\text{ (horizontal mirror)}&2&3&3&3&3\\
 C_2\text{ (vertical mirror)}&2&1&1&1&1\\
 C_3&3&2&2&2&2\\
 C_2^2&4&1&6&3&2\\
 C_6&6&0&5&5&2\\
 S_3\text{ (positive)}&6&0&5&5&2\\
 S_3\text{ (mixed)}&6&2&7&2&2\\
 S_3\times C_2&12&0&16&5&2
\end{array}
\]
The attached CSV gives every maximal stabilizer; the JSON table gives every
partial section and all of its multiplier values. In particular it checks the
reflection failure of maximal-section uniqueness from B.8 inside this whole list.
\end{claim}
\passage{15; 75-77; 150-155; prompt B.f}

\begin{proof}
Place the three punctures at the vertices of an equatorial equilateral triangle.
The positive rotation group of that triangle is $S_3$: it is generated by the
$120$-degree vertical rotation $r$ and a half-turn $a$ around an axis through a
puncture. Reflection $\tau$ in the equator commutes with both. The mapping class
group map to permutations and orientation is an isomorphism. For its kernel,
fix three arcs forming a tree through the punctures by isotopy; its complement
is a disc, and Alexander's disc isotopy makes the remaining map isotopic to the
identity. This proves completeness of the ambient finite group and its stated
realization.

The program represents this group exactly by pairs $(p,e)\in S_3\times C_2$.
Starting with the trivial subgroup it repeatedly adjoins every group element
and takes closure, until no new subgroup occurs. Every subgroup is generated
by its elements, so this terminating procedure is exhaustive. It finds sixteen
subgroups and ten conjugacy classes, represented in the displayed order by
\[
 1,\ \langle a\rangle,\ \langle\tau\rangle,\
 \langle a\tau\rangle,\ \langle r\rangle,\ \langle a,\tau\rangle,\
 \langle r,\tau\rangle,\ \langle r,a\rangle,\
 \langle r,a\tau\rangle,\ \langle r,a,\tau\rangle.
\]
Here is the independent geometric input to its section calculation. Denote the
poles by $N,S$, and the three surviving opposite equatorial axis points by
$Q_i$; the other three axis points are the removed punctures. The fixed set of
$\langle r\rangle$ is $\{N,S\}$, and that of a positive half-turn is its single
surviving $Q_i$. The horizontal mirror has three open intervals, each containing
one $Q_i$. A vertical mirror has one real line after its single puncture is
removed, and this line contains $N,S,Q_i$. Intersect these fixed sets to find the
fixed sets of every subgroup. B.3 and B.4 turn their components into exactly
its section classes. A partial section on a mirror interval is contained in a
point stabilizer precisely when that point lies on the interval; this determines
all inclusions and hence all maximal classes. The code records and checks these
inclusions and the conjugation action of every element.

Every reflection and positive half-turn has multiplier $1/2$. For $r$, the
values at $N,S$ are respectively $1/3,-1/3$, and its powers give all remaining
values. This computes all entries in the attached JSON. The geometry also
verifies that all the centralizers occurring for nontrivial sections here are
trivial: every mirror component is a line, and every positive section fixes an
isolated point. The absence of a common finite overgroup for the two vertex
stabilizers along a vertical mirror is retained in the inclusion data.
\end{proof}

\begin{claim}{B.24}{Proved}{Completed finite-subgroup catalogue and exact indexing}
Every finite-subgroup conjugacy class in the requested range
$2g+\nu\le6$ has now been enumerated.  For the hyperbolic types,
the numbers of classes in the extended mapping class group are
\[
\begin{array}{c|r|l}
 (g,\nu)&\#\text{ classes}&\text{proofs and executed tables}\\\hline
 (0,3),(0,4),(0,5),(0,6)&10,30,20,57&B.23,B.25,B.26\\
 (1,1),(1,2),(1,3),(1,4)&13,40,38,115&B.25,B.27\\
 (2,0),(2,1),(2,2)&92,32,121&B.28,B.29\\
 (3,0)&292&B.33,B.37\text{--}B.41.
\end{array}                                             \tag{B.24.1}
\]
Each table includes the full and partial section classes, multipliers,
centralizers, extension incidences, maximal section classes $\Phi$,
and their group orbits and stabilizers.  In the four exceptional
abelian types, the finite mapping-class subgroup counts are
respectively $2,2,5,13$ for sphere, plane, annulus and torus.  Their
injectively realized finite-action counts are $3,2,7,17$; the exact
covering extension is retained as part of the section index in B.36,
as required by the explicit failures of extension-independent indexing
in B.22.  The families with a nontrivial kernel on passage to mapping
classes are described in B.14.
The exact tube words, stabilizer action, strip relations and arbitrary
transversal obstruction formerly required here are proved and executed
in B.30--B.32.
\end{claim}
\passage{\S{}18; folios 96--101; scans 193--203; \S{}33; folios 222--237; scans 440--470; Part B(e) and B(f) of the task}

\begin{proof}
The first three rows of (B.24.1) are precisely the complete catalogues
cited there.  We spell out the closed genus-three reconciliation to
record the coverage and the absence of double counting.
The positive portion has $55$ classes by B.33.  For a mixed finite
action let $H$ be its positive subgroup.  The positive quotient has
genus zero, one, two or three, by the positive classification.
The genus-zero case gives either a mirror disc ($145$ classes,
B.37) or a projective plane ($11$ classes, B.39).  The genus-one case
gives an annulus or M\"obius band ($57$ classes, B.40), or a Klein
bottle ($6$ classes, B.41).  The genus-two case is the free positive
$C_2$ action and gives $12$ mixed classes in B.38.  The genus-three
case has trivial positive part and gives its six negative involution
classes in B.38.  Thus the number of mixed classes is
\[
                   145+11+57+6+12+6=237,
                     \qquad55+237=292.                \tag{B.24.2}
\]
Every conjugacy descends to homeomorphisms of both the positive
quotient and the full quotient orbifold.  The positive quotient genus,
and the full quotient's boundary and orientability, separate the
listed portions.  Within each portion the cited proofs calculate
exactly the geometric kernel-normalizer orbits, so these counts are
irreducible counts of surface-action conjugacy classes.

B.34 supplies an independent hyperelliptic construction and is
counted through its explicit overlaps, not added a second time.
All $79$ of its rows are identified by monodromy isomorphisms:
$18$ with positive B.33 rows, $57$ with B.37 rows, and four with B.39
rows.  The executed program
\path{reconcile_genus3_catalogue.py} records these $79$ maps and
checks the corresponding section counts.  Its unified
\path{type_3_0_complete_catalogue.csv} has $292$ different
class identifiers, and \path{type_3_0_complete_sections.json}
contains the full finite group and section model for every row.
The exceptional replacement in B.36 is complete with its stated
realization index; it retains, in particular, the sphere antipodal
action's algebraic splitting without identifying it with a geometric
fixed point.  This proves the claimed coverage with exactly the
comparisons and indices used by the computations.
\end{proof}

\begin{claim}{B.25}{Proved}{Complete executed catalogues in types $(1,1)$ and $(0,4)$}
The extended mapping class groups of types $(1,1)$ and $(0,4)$ have,
respectively, $13$ and $30$ conjugacy classes of finite subgroups, including
the trivial subgroup.  The tables below give every class and the numbers of
full splittings, nontrivial partial splittings, maximal partial section
classes $\Phi$, and group orbits on $\Phi$.  The accompanying exact-arithmetic
files give every section, its image subgroup, its multiplier on every
element, its centralizer, all inclusion relations, the group stabilizers,
and the puncture stabilizers.  These are exhaustive catalogues for the two
specified surface types.
\end{claim}
\passage{\S{}15; folios 73--75; scans 148--151; \S{}15; folios 77--79; scans 155--160; \S{}18; folios 96--101; scans 193--203; Part B(f) of the task; no further manuscript assertion}

\begin{proof}
Put
\[
 S=\begin{pmatrix}0&-1\\1&0\end{pmatrix},\qquad
 R=\begin{pmatrix}0&-1\\1&1\end{pmatrix},\qquad
 F=\begin{pmatrix}0&1\\1&0\end{pmatrix}.
\]
For the punctured torus use their linear actions on
$\R^2/\Z^2\setminus\{0\}$.  The homology map identifies its extended
mapping class group with $\mathrm{GL}_2(\Z)$: surjectivity is supplied by
linear maps, and a map acting trivially on homology can be isotoped to fix
two transverse basis curves, after which the complementary punctured disc
is handled by the Alexander isotopy.  This is the punctured version of
\cite[Theorem 2.5 and pp.~54--55]{A_FarbMargalit2012}.

For the four-punctured sphere use the pillowcase
\[
 \bigl(\R^2/\Z^2\setminus (\tfrac12\Z^2/\Z^2)\bigr)/\{p\sim-p\}.
\]
Write $U(p)=p+(1/2,0)$ and $V(p)=p+(0,1/2)$.  The extended mapping
class group is
$V_4\rtimes\mathrm{PGL}_2(\Z)$ with these affine representatives.
The map to $\mathrm{PGL}_2(\Z)$ is obtained by lifting through the double
branched torus cover and acting on torus homology.  Its kernel consists of
the four half translations: a kernel element fixes all slope classes;
after a half translation it fixes the four punctures and two filling
curves, so the Alexander isotopy on their complementary discs kills it.
Linear maps split the homomorphism and act on half translations by
reduction modulo two.  This is the explicit morphism of
\cite[Proposition 2.7 and the extended statement on p.~57]{A_FarbMargalit2012}.

Here are finite amalgams that make enumeration complete.  The matrices
$S,R$ generate $\mathrm{SL}_2(\Z)$ by the Euclidean algorithm, since
$S^{-1}R=\left(\begin{smallmatrix}1&1\\0&1\end{smallmatrix}\right)$.
Modulo $\{\pm I\}$ their images have orders two and three, with no
further relation: on the real projective line, $\bar S$ sends the negative
half-line to the positive half-line, and both nonidentity powers of
$\bar R$ send the positive half-line into the negative half-line.  The
free-product ping-pong argument gives
$\mathrm{PSL}_2(\Z)=C_2*C_3$.
Conjugation by $F$ fixes $\bar S$ and inverts $\bar R$.  Consequently
\[
 \mathrm{PGL}_2(\Z)=V_4*_{C_2}S_3,
 \qquad
 \mathrm{GL}_2(\Z)=\langle S,F\rangle
       *_{\langle-I,F\rangle}\langle R,F\rangle
       =D_8*_{V_4}D_{12},
\]
where $D_k$ has order $k$.  The second identity follows also by lifting
the first through their common central subgroup $\langle-I\rangle$;
the normal form for the amalgam leaves no additional kernel.
Adding the normal group of half translations gives the pillowcase amalgam
\[
 (V_4\rtimes V_4)*_{V_4\rtimes C_2}(V_4\rtimes S_3),
\]
whose vertex groups have orders $16$ and $24$, and whose edge group has
order $8$.  The two vertex groups are $D_8\times C_2$ and $S_4$.

A finite subgroup of an amalgam fixes a vertex in its Bass--Serre tree.
Indeed the convex hull of a finite orbit is a finite tree; its centre is
fixed, and the action preserves the two vertex types, so an edge centre
also gives fixed endpoints.  Two vertex subgroups are conjugate in the
amalgam precisely through the equivalence generated by conjugacy in
each vertex group and identification of edge subgroups.  To prove the
nontrivial direction, a conjugacy makes the subgroup fix both endpoints
of the intervening path; it fixes every edge on that path, which
decomposes the conjugacy into those successive vertex and edge steps.
Thus finite enumeration in the two vertices, followed by these exact
fusion steps, is a complete conjugacy calculation in the infinite group.
The program finds $10$ and $16$ actual subgroups in the torus vertices,
and $35$ and $30$ in the pillowcase vertices.  Fusion leaves exactly
$13$ and $30$ classes.  All subgroup lists are saturated under adjoining
an arbitrary vertex-group element, so the enumeration has no search bound
that could omit a subgroup.

We next justify the computation of sections.  Let
$Q=(\tfrac1{12}\Z^2)/\Z^2$.  Every isolated fixed point needed for
either model lies in $Q$.  For the torus it solves $(A-I)p=0$; for the
pillowcase it solves
\[
 (A-\epsilon I)p=-t/2\pmod{\Z^2},\qquad
 \epsilon\in\{1,-1\},\quad t\in(\Z/2)^2.
\]
For every positive matrix in the finite vertices, the inverse matrices
and the particular solutions of the nonsingular equations have
denominators dividing $12$.  The program certifies this by exact
fractions for every matrix and translation, rather than inferring it
from a grid count.  The singular case for a nonidentity positive element
is a nontrivial translation and contributes no solutions to that sign.

For a reflection, the same equations are affine circle equations.  The
greatest common divisor of the nonzero matrix entries is one or two, so
every component meets $Q$.  On each such circle the primitive integral
kernel direction, divided by $12$, joins consecutive sampled points.
In the pillowcase the plus and minus circles meet only at two-torsion
points, all of which are included in $Q$.  Quotienting the sampled graph
by $p\sim-p$ therefore gives a subdivision of the true fixed circles.
Deleting the prescribed punctures gives exactly their connected
components.  A component adjacent to a deleted point is a line; the
others are ovals.  All rotations and all punctures are retained as
vertices, so no incidence needed for restriction of sections is lost.

B.3 and B.4 identify these fixed components with section classes.
For a partial section with a rotation use its fixed point; for a
reflection use its fixed line or oval.  Inclusion of image subgroups and
containment of the fixed point in the fixed component determine exactly
when one section extends another.  Applying each affine element gives
the exact action on this finite section poset.  This computes $\Phi$,
its orbits and its stabilizers without replacing pointwise section data
by conjugacy of image subgroups in the quotient.  Centralizers are $1$
for isolated points and lines and $\Z$ for ovals, by B.3--B.4.

At a torus fixed point the local positive derivative is $A$.  At a
nonbranch pillowcase fixed point it is $\epsilon A$, with the unique
sign specified by $Ap+t/2=\epsilon p$.  For these matrices the traces
$-2,-1,0,1$ give orders $2,3,4,6$; the half-turn has multiplier $1/2$,
and otherwise the sign of $\det(e_1,Ae_1)$ selects $+1/n$ or $-1/n$.
Reflections have multiplier $1/2$.  At a removed pillowcase branch
point the quotient coordinate squares a torus coordinate, so the
positive rotation angle doubles.  The program records these puncture
data separately.  It checks the order and power identities for every
element of every section and checks that every section containing a
rotation has exactly one maximal overgroup.  The reflection classes
with several maximal extensions are retained explicitly.

Finally, C.5--C.6 show that the chosen finite actions compute the action
attached to each mapping class subgroup up to conjugacy, so these
fixed-component computations depend on the listed subgroup class, not
on an unrecorded realization choice.
\end{proof}

In the tables $f$ is the number of full splitting classes, $p$ the number
of nontrivial partial splitting classes with image subgroup retained,
$\phi=|\Phi|$, and $o=|\Phi/\Gamma|$.  Juxtaposition in the generator
column means composition of the matrices or affine maps defined above.
For $(0,4)$ the letters $S,R,F$ mean their pillowcase classes.
The complete output files are
\texttt{type\_1\_1\_finite\_subgroups.csv},
\texttt{type\_0\_4\_finite\_subgroups.csv}, and the corresponding
\texttt{partial\_sections.json} files, produced by
\texttt{catalogue\_low\_complexity.py}.

\begin{center}Type $(1,1)$\end{center}
\begin{longtable}{lllrrrrr}
\toprule
Class & Group & Generators & $|\Gamma^+|$ & $f$ & $p$ & $\phi$ & $o$\\
\midrule\endhead
T01 & $1$ & $1$ & 1 & 1 & 0 & 0 & 0\\
T02 & $C_2$ & $SF$ & 1 & 2 & 2 & 2 & 2\\
T03 & $C_2$ & $SSF$ & 1 & 1 & 1 & 1 & 1\\
T04 & $C_2$ & $SS$ & 2 & 3 & 3 & 3 & 3\\
T05 & $C_3$ & $RR$ & 3 & 2 & 2 & 2 & 2\\
T06 & $C_2^2$ & $SS,SF$ & 2 & 3 & 10 & 3 & 3\\
T07 & $C_2^2$ & $SS,SSF$ & 2 & 1 & 6 & 3 & 2\\
T08 & $C_4$ & $S$ & 4 & 1 & 4 & 3 & 2\\
T09 & $S_3$ & $RR,RRF$ & 3 & 2 & 7 & 2 & 2\\
T10 & $S_3$ & $RR,RF$ & 3 & 0 & 5 & 5 & 2\\
T11 & $C_6$ & $R$ & 6 & 0 & 5 & 5 & 2\\
T12 & $D_8$ & $S,SF$ & 4 & 1 & 15 & 3 & 2\\
T13 & $D_{12}$ & $R,RRF$ & 6 & 0 & 16 & 5 & 2\\
\bottomrule
\end{longtable}

\begin{center}Type $(0,4)$\end{center}
\begin{longtable}{lllrrrrr}
\toprule
Class & Group & Generators & $|\Gamma^+|$ & $f$ & $p$ & $\phi$ & $o$\\
\midrule\endhead
S01 & $1$ & $1$ & 1 & 1 & 0 & 0 & 0\\
S02 & $C_2$ & $SF$ & 1 & 4 & 4 & 4 & 4\\
S03 & $C_2$ & $SFV$ & 1 & 1 & 1 & 1 & 1\\
S04 & $C_2$ & $SFUV$ & 1 & 0 & 0 & 0 & 0\\
S05 & $C_2$ & $F$ & 1 & 2 & 2 & 2 & 2\\
S06 & $C_2$ & $V$ & 2 & 2 & 2 & 2 & 2\\
S07 & $C_2$ & $S$ & 2 & 0 & 0 & 0 & 0\\
S08 & $C_3$ & $RR$ & 3 & 1 & 1 & 1 & 1\\
S09 & $C_2^2$ & $V,SF$ & 2 & 2 & 9 & 4 & 3\\
S10 & $C_2^2$ & $V,SFU$ & 2 & 0 & 3 & 3 & 2\\
S11 & $C_2^2$ & $UV,SF$ & 2 & 0 & 6 & 6 & 3\\
S12 & $C_2^2$ & $UV,SFV$ & 2 & 2 & 6 & 2 & 2\\
S13 & $C_2^2$ & $UV,F$ & 2 & 2 & 8 & 2 & 2\\
S14 & $C_2^2$ & $SF,F$ & 2 & 0 & 6 & 6 & 3\\
S15 & $C_2^2$ & $SFUV,F$ & 2 & 0 & 2 & 2 & 1\\
S16 & $C_4$ & $FU$ & 2 & 0 & 2 & 2 & 1\\
S17 & $C_2^2$ & $V,U$ & 4 & 0 & 6 & 6 & 3\\
S18 & $C_2^2$ & $UV,S$ & 4 & 0 & 2 & 2 & 1\\
S19 & $C_4$ & $SU$ & 4 & 2 & 4 & 2 & 2\\
S20 & $S_3$ & $RR,FR$ & 3 & 1 & 8 & 4 & 2\\
S21 & $C_2^3$ & $V,U,SF$ & 4 & 0 & 18 & 6 & 3\\
S22 & $C_2^3$ & $UV,SF,F$ & 4 & 0 & 12 & 6 & 2\\
S23 & $C_4\times C_2$ & $FU,SU$ & 4 & 0 & 8 & 6 & 2\\
S24 & $D_8$ & $FU,V$ & 4 & 0 & 12 & 6 & 2\\
S25 & $D_8$ & $SU,SFV$ & 4 & 2 & 16 & 2 & 2\\
S26 & $D_8$ & $FU,SFV$ & 4 & 0 & 6 & 2 & 1\\
S27 & $D_8$ & $SU,V$ & 8 & 0 & 8 & 6 & 2\\
S28 & $A_4$ & $RR,RUR$ & 12 & 0 & 10 & 10 & 2\\
S29 & $D_8\times C_2$ & $FU,SU,V$ & 8 & 0 & 28 & 6 & 2\\
S30 & $S_4$ & $FUR,URF$ & 12 & 0 & 32 & 10 & 2\\
\bottomrule
\end{longtable}

\begin{claim}{B.26}{Proved}{Complete executed spherical catalogues through six punctures}
For $\nu=3,4,5,6$, the extended mapping class group of the sphere with
$\nu$ unlabelled punctures has respectively
\[
                         10,\quad30,\quad20,\quad57
\]
conjugacy classes of finite subgroups, including the trivial subgroup.
The exact finite parameterization below gives every class once.  The
executed files \texttt{type\_0\_\nu\_orbifold\_catalogue.csv} and
\texttt{type\_0\_\nu\_orbifold\_sections.json}, for each of these four
values of $\nu$, give all partial-section classes, all their inclusion
relations, $\Phi$, its group action and stabilizers, the centralizers,
the multiplier on each section element, and the puncture stabilizers.
Thus these files settle the whole computation requested in B(f) for
the hyperbolic spherical types with $2g+\nu\le6$.
\end{claim}
\passage{\S{}15; folios 73--75; scans 148--151; \S{}15; folios 77--79; scans 155--160; \S{}18; folios 96--101; scans 193--203; \S{}33; folios 222--225; scans 440--445; Part B(f) of the task; no further manuscript assertion}

\begin{proof}
Extend an action over the punctures.  By B.1 and the sphere theorem of
\cite[Theorem 1.2]{bKolev2006}, a finite topological action on $S^2$ is
conjugate to an orthogonal action.  Write $Q=S^2/G$ with its local
orbifold groups.  A point has cyclic positive stabilizer, or a dihedral
stabilizer with mirrors; the mirror-only stabilizer has order two.
The covering formula is
\[
 \chi_{\rm orb}(Q)=\chi(|Q|)
  -\sum_{j=1}^r(1-1/m_j)
  -\frac12\sum_{j=1}^k(1-1/n_j)=\frac2N,
 \qquad N=|G|.                                      \tag{B.26.1}
\]
Here the $m_j$ are interior cone orders and the $n_j$ are corner orders.
One proves this formula directly by a triangulation whose vertices
include all cone and corner points: count each cell in the covering
with weight equal to the reciprocal of its stabilizer order.  Mirrors
are boundary strata; they introduce no additional term after the
boundary cells and their endpoints are counted together.

Since (B.26.1) is positive, the underlying compact surface is $S^2$,
$\mathbb{RP}^2$, or a disc.  The sphere has at most three cones.  A
single cone is impossible, since its local generator would be trivial
in the orbifold presentation.  Two cones have equal orders, since the
two local generators are inverses and both local cyclic groups inject.
For three cones the inequality $1/m_1+1/m_2+1/m_3>1$ gives precisely
$(2,2,n),(2,3,3),(2,3,4),(2,3,5)$.
The projective plane has at most one cone.  For a disc, positivity gives
$2r+k\le3$.  A single corner without a cone is impossible; two corners
without a cone have equal orders, by the same local-injectivity argument.
Three corners give the same four triples.  One cone and one corner
give $2/m+1/n>1$, hence $(m,n)=(2,n)$ or $(3,2)$.
Consequently the following table is exhaustive.  In its notation
$S(m_1,\ldots,m_r)$ is a sphere with cones, $P(m)$ is a projective
plane with one cone, $D(m_1,\ldots,m_r;n_1,\ldots,n_k)$ is a mirror
disc with the specified interior cones and cyclically ordered corners,
and an empty list means no such singularities.
\[
\begin{array}{c|c@{\qquad}c|c}
Q&N&Q&N\\\hline
S()&1&P()&2\\
S(n,n)&n&P(n)&2n\\
S(2,2,n)&2n&D(;)&2\\
S(2,3,3)&12&D(;n,n)&2n\\
S(2,3,4)&24&D(n;)&2n\\
S(2,3,5)&60&D(;2,2,n)&4n\\
&&D(2;n)&4n\\
&&D(;2,3,3)&24\\
&&D(;2,3,4)&48\\
&&D(;2,3,5)&120\\
&&D(3;2)&24
\end{array}                                                   \tag{B.26.2}
\]
Every integer $n$ in this table is at least two.  For $3\le\nu\le6$
only $2\le n\le6$ can occur.  Indeed, in every family depending on
$n$, the union of all point orbits of size strictly smaller than $n$
has at most two points; an invariant set of at least three punctures
therefore contains an orbit of size at least $n$.  This assertion also
follows immediately by dividing $N$ by each local group order in
(B.26.2).  In the two icosahedral rows every point orbit has at least
twelve points.  Those rows cannot occur in the present range.  These
are proved finite bounds, rather than bounds inferred from the output.

We now give the exact decoration rule used in the program.  A cone
orbit may be removed or retained; put $c_j\in\{0,1\}$ accordingly.
A corner orbit similarly has $d_j\in\{0,1\}$.  For each open mirror
edge let $e_i\ge0$ be the number of removed orbits on that edge; a
cornerless mirror circle is one edge for this purpose.  Finally let
$a\ge0$ be the number of regular interior orbits removed.  The single
equation is
\[
 \sum_jc_j\frac N{m_j}
 +\sum_jd_j\frac N{2n_j}
 +\sum_i e_i\frac N2+aN=\nu.                       \tag{B.26.3}
\]
On a sphere, permutations of cones of equal order identify decorations.
On a mirror disc, identify them under precisely the rotations and
reflections of the cyclic corner list that preserve its order labels.
The corresponding edge counts move with their edges.  For two corners
this rule includes the independent exchanges of the two corners and
the two edges.  For one or zero corners it introduces no further
identification of the single edge count.  There is at most one interior
cone in every remaining disc or projective-plane row.

For clarity, this rule also proves completeness of the conjugacy
classification.  A homeomorphism of an interval moves any two
unlabelled finite subsets of the same size into one another, and the
same assertion holds in the connected regular interior of each quotient
surface.  Such moves extend to the quotient surface by disjoint small
discs and strips.  The only remaining data are therefore the selected
singular strata and the cyclic mirror-edge counts in (B.26.3), modulo
the explicitly stated symmetries.  Conversely every quotient
homeomorphism preserves those strata, their orders, and their cyclic
incidence.  Since $S^2$ is the universal orbifold cover, a homeomorphism
between the decorated quotient orbifolds lifts to a conjugacy between
their covering actions; this follows by lifting paths and their local
cone or mirror charts.  A conjugacy of actions descends to such a
quotient homeomorphism.  Thus no extra equivalence of decorations is
possible.  Finally C.5 and C.6 identify these topological action classes
with conjugacy classes of finite subgroups in the extended mapping
class group of the punctured sphere.  The induced action there is
faithful: a nonidentity positive periodic map isotopic to the identity
would fix every puncture individually, whereas a nonidentity sphere
rotation has exactly two fixed points; a negative map cannot be
isotopic to the identity.  Here $\nu\ge3$ is essential.

Here are certificates for the finite group calculations.  The source
\texttt{catalogue\_spheres.py} constructs the cyclic and dihedral
families as regular permutation groups, the tetrahedral and octahedral
positive groups as $A_4,S_4$, their needed extensions by an explicitly
given central involution, and the full octahedral group as signed
permutations of three axes.  All generators, products, orders and signs
are recorded as finite permutations.  The program verifies each
orbifold relation, injection of every local cyclic or dihedral group,
and $N\chi_{\rm orb}(Q)=2$ using rational arithmetic.  These checks
prove that the permutation group is the entire orbifold group, as
follows.  The kernel of the displayed quotient of the orbifold
presentation produces a covering with no cone, corner or mirror
singularity because all local groups inject.  Its orientation
character is trivial, and its Euler characteristic is two.  It is
therefore a closed orientable sphere; its fundamental group is trivial.
This proves injectivity of the quotient presentation as well as its
already checked surjectivity.  In particular these certificates do
not merely exhibit a possibly proper finite quotient of an unexamined
presentation.

We describe how the section data are recovered without numerical
sampling.  Cone and corner lifts are cosets of their cyclic and
dihedral stabilizers.  In a disc presentation choose reflections
$r_0,\ldots,r_k$ with
\[
 (r_jr_{j+1})^{n_j}=1,\qquad
 r_k=e r_0e^{-1},\qquad x_1\cdots x_r e=1.            \tag{B.26.4}
\]
Indices of the corner orders are shifted here to start at zero.
The code stores $r_0,\ldots,r_{k-1}$ and derives the final reflection.
An edge lift is a coset of $\langle r_i\rangle$ and meets the adjacent
corner cosets; the closing incidence at edge zero is multiplied by
$e^{-1}$.  With no corners it instead joins the corresponding two
cosets of $\langle r_0\rangle$ by that same rule.  Insert the prescribed
removed vertices along each edge, retaining two auxiliary vertices in
each complementary segment.  These auxiliary vertices prevent a
one-edge loop from collapsing when represented as an ordinary finite
adjacency graph.  This is exactly the lifted mirror locus with a
finite subdivision.

For any reflection $r\in G$, keep the vertices whose stabilizers
contain $r$.  Before punctures are deleted the program verifies degree
two at every such vertex.  Its components are consequently the full
fixed circles of $r$.  Deleting the selected puncture vertices gives
all the fixed lines and ovals, with their exact incidence.  For a
subgroup whose positive part is nontrivial, the fixed points are
precisely the retained cone and corner lifts whose stabilizers contain
that subgroup.  A negative involution without a mirror has an empty
fixed set.  Thus B.3 and B.4 supply every section class, with centralizer
$1$ for a point or a line and $\Z$ for an oval.  A section over $H$
extends a section over $K<H$ exactly when its fixed component is
contained in the latter fixed component.  Hence the same finite
incidence calculation gives the entire partial-section poset and all
its maximal elements, including the reflection cases with several
maximal extensions.  Left multiplication on cosets gives the group
action and every stabilizer without a choice of orbit representatives.

The multiplier is exact as well.  Orient a local lift of each cone,
and at a corner use the positive generator $r_jr_{j+1}$ determined by
the oriented boundary order.  Write this generator as $u$, of order
$m$.  At the point represented by a coset $gD$, an element $h$ of its
positive stabilizer satisfies $g^{-1}hg=u^b$, and
\[
                   \mu(h)=\operatorname{sgn}(g)\frac bm
                                \pmod{\Z}.           \tag{B.26.5}
\]
Replacing $g$ by another coset representative leaves this value
unchanged: a positive local element commutes with $u$, while a
negative local element both reverses the sign and inverts $u$.
For a reflection the value is $1/2$, and for the identity it is zero.
Formula (B.26.5) is the local instance of B.12, and applies to all
sections in the files.  The code checks its order and every power
identity, as well as uniqueness of the maximal extension for every
section with nontrivial positive part.  All nontrivial subgroups of
each finite permutation group are enumerated by saturation under
adjoining each group element; this stops only when the subgroup set
is closed under every possible enlargement.

The following compact census is part of the executed output; each
entry is the number of conjugacy classes of that group order.  A dash
denotes zero.  Complete individual rows, rather than just this census,
are in the CSV files specified in the statement.
\[
\begin{array}{c|rrrrrrrrrrrrr|r}
\nu\backslash N&1&2&3&4&5&6&8&10&12&16&20&24&48&\text{total}\\\hline
3&1&3&1&1&-&3&-&-&1&-&-&-&-&10\\
4&1&6&1&11&-&1&7&-&1&1&-&1&-&30\\
5&1&4&1&4&1&3&1&3&1&-&1&-&-&20\\
6&1&7&1&14&1&9&8&1&9&1&-&4&1&57
\end{array}                                                    \tag{B.26.6}
\]
As an independent check, the program compares the multisets of all
six columns consisting of group order, positive order, full-section
count, nontrivial partial-section count, maximal-section count, and
group-orbit count, for $\nu=3,4$, against the previous permutation and
amalgam calculations in B.23 and B.25.  Both multisets agree, with
multiplicities.  The check and its inputs are attached; the completeness
proof above does not depend on that agreement.
\end{proof}

\begin{claim}{B.27}{Proved}{Complete executed punctured-torus catalogues through four punctures}
The extended mapping class groups of a torus with $\nu=1,2,3,4$
unlabelled punctures have, respectively,
\[
                         13,\quad40,\quad38,\quad115
\]
conjugacy classes of finite subgroups.  The files
\texttt{type\_1\_\nu\_affine\_catalogue.csv}, together with their
\texttt{affine\_sections.json} companions and
\texttt{torus\_affine\_normal\_forms.json}, list every class once.
They specify each group by its complete finite set of rational affine
transformations, its punctures, every full and nontrivial partial section
class, all extension incidences, $\Phi$ and its group orbits and stabilizers,
the centralizers, and $\mu$ on every section and puncture-stabilizer element.
Thus the torus portion of B(f) is complete for every hyperbolic type in the
requested range.
\end{claim}
\passage{\S{}15; folios 73--75; scans 148--151; \S{}15; folios 77--79; scans 155--160; \S{}18; folios 96--101; scans 193--203; Part B(f) of the task; no further manuscript assertion}

\begin{proof}
We give the finite classification underlying the program
\texttt{catalogue\_tori.py}; this also specifies the tables without relying
on abstract finite-group names.  By B.14 extend over the punctures and put
the action in affine form on $V/\Lambda=\R^2/\Z^2$.  Write
\[
 1\longrightarrow T\longrightarrow G\longrightarrow H\longrightarrow1,
 \qquad T\subset V/\Lambda,\quad H\subset\mathrm{GL}_2(\Z).
                                                        \tag{B.27.1}
\]
Here $T$ consists exactly of the translations in $G$.  It acts freely on
$V/\Lambda$, so each $T$-orbit on the puncture set has $|T|$ elements.
Consequently $|T|\mid\nu$ and in particular $|T|\le4$.  The complete list
of subgroups of $V/\Lambda$ of order at most four consists of the trivial
subgroup, the three cyclic subgroups of order two, the four cyclic
subgroups of order three, the six cyclic subgroups of order four, and
$(\frac12\Lambda)/\Lambda$.  Indeed a group of order at most four is
cyclic or $C_2^2$; a cyclic subgroup is generated by a point of that exact
order, and division by the number of generators gives precisely these
counts.  All fifteen groups lie in $(\frac1{12}\Lambda)/\Lambda$.
The thirteen possibilities for $H$, up to integral conjugacy, are exactly
those proved in B.25.  For each we retain only the $H$-invariant $T$.

We now enumerate the affine extensions, including the actions with no
common fixed point.  If $H$ contains a nonidentity positive element,
choose a generator $r$ of its cyclic positive subgroup.  Every affine
map with linear part $r$ has a fixed point, since $I-r$ is invertible
on $V$ and is therefore onto on $V/\Lambda$.  Translation of that point
to the origin makes this generator $(r,0)$.  A positive cyclic $H$
then gives exactly $G=\langle T,(r,0)\rangle$.  For a dihedral $H$
choose a negative generator $f$ and write its lift as $(f,b)$.  The two
relations in $H$ give exactly
\[
       (I-r)b\in T,\qquad (I+f)b\in T.                 \tag{B.27.2}
\]
Conversely these equations make the subgroup generated by
$T,(r,0),(f,b)$ a finite group with kernel $T$ and quotient $H$:
every word reduces to a translation followed by $r^j$ or $r^jf$, and
its linear part distinguishes these cosets.  Values of $b$ modulo $T$
are enough.  The finite set $(I-r)^{-1}T$ therefore supplies every
candidate, with no bound on words or finite quotients involved.
Translations of the origin preserving this normal form are exactly
$a$ with $(I-r)a\in T$; they change $b$ by $(I-f)a$ modulo $T$.

The remaining nontrivial linear groups have $H=\langle f\rangle$
of order two.  Put $W=(V/\Lambda)/T$.  Their extensions are represented
by $b\in W$ with $(I+f)b=0$, modulo $(I-f)W$.  Every such class has a
representative in $W[2]$.  To verify this assertion directly, choose a
lift $v\in V$ of $b$ for the enlarged lattice
$\Lambda_T=\{v\in V:v\bmod\Lambda\in T\}$.
Then $(I+f)v\in\Lambda_T$.  Write
$v=v_++v_-$ in the real $+1$ and $-1$ eigenspaces of $f$.
Subtracting $(I-f)(v_-/2)$ removes $v_-$, and the remaining $v_+$
satisfies $2v_+\in\Lambda_T$.  Thus the program uses precisely
\[
              2b\in T,\qquad (I+f)b\in T.             \tag{B.27.3}
\]
For two candidates $b,b'$ the translation-conjugacy test is also exact.
Let $w$ be a primitive integral row vector satisfying $wf=w$.
The image of $I-f$ on $V/\Lambda$ is the connected circle $\ker w$:
its real image is that eigenspace, and multiplication by any nonzero
integer is onto on a circle.  Hence the two candidates are equivalent
precisely when $w(b-b')\in w(T)$ in $\R/\Z$.
This includes both reflections and fixed-point-free glide reflections.

Integral changes of basis are handled without truncating an infinite
matrix group.  For $H=1$ or $\{\pm I\}$, use the generators $S,R,F$
of B.25 on the complete fifteen-element list of possible $T$'s.
Their orbit calculation is the full $\mathrm{GL}_2(\Z)$-conjugacy
calculation.  For a fixed $T$, a spanning tree of this finite orbit
and Schreier's formula give generators for its entire integral
stabilizer.  Explicitly, with a chosen representative $B_U$ carrying
$T$ to $U$, the generator associated to a matrix $A\in\{S,R,F,S^{-1},R^{-1}\}$
and the edge $U\mapsto AU$ is
\[
                       B_{AU}^{-1}AB_U.                \tag{B.27.4}
\]
For every other $H$ its integral normalizer is finite.  Here is an
exhaustiveness check for the small matrix enumeration in the code.
For a reflection conjugate to $\operatorname{diag}(-1,1)$, a commuting
integral unimodular matrix is diagonal with entries $\pm1$.
For a reflection conjugate to
$\left(\begin{smallmatrix}0&-1\\-1&0\end{smallmatrix}\right)$,
the equation is $A=\left(\begin{smallmatrix}a&b\\b&a\end{smallmatrix}\right)$,
with $(a-b)(a+b)=\pm1$, giving four solutions.  A normalizer of the
four-group generated by this reflection and $-I$ may also exchange
the two eigenspaces; the same calculation gives eight solutions.
A group with a rotation of order four has its normalizer in the eight
signed permutation matrices, since a commuting matrix has determinant
$a^2+b^2=1$.  A group with rotation of order three or six has its
normalizer in the twelve matrices for the hexagonal lattice: the
commuting equation gives the norm $a^2+ab+b^2=1$, and one reflection
supplies the other coset.  Testing these finite sets against $H$
gives the full normalizer.  Every resulting matrix has entries of
absolute value at most one, so the code's box of entries from $-2$
to $2$ contains every solution and introduces none outside the tested
normalizer equations.  Applying these matrices and the origin changes
above fuses all and only affine-conjugate candidates.
The result is $93$ affine action classes with $|T|\le4$ before choosing
punctures; some admit no invariant set of a given required size.

We explain why decorating these affine actions gives topological
conjugacy, including all permutations of the singular strata.
B.14 proves that a topological conjugacy between finite affine actions
can be replaced by an affine conjugacy.  The replacement also induces
the same permutation on the strata used here.  In detail write a lift
of the conjugacy as $h(x)=Ax+\delta(x)$ with $\delta$ periodic, and
put $\bar\delta=\int_{V/\Lambda}\delta$.  The equivariance equations
and invariance of Haar measure show that $x\mapsto Ax+\bar\delta$
has the same equivariance.  The straight interpolation of these two
lifts is an equivariant homotopy.  An isolated fixed point of a
nontrivial positive subgroup stays constant under this homotopy.
Each reflection circle stays in the same fixed-set component.
The degree on every such circle is fixed by the homology map $A$;
these circles are essential, since they lift to parallel lines.
Together with the images of the isolated vertices this determines
the permutation of the intervening open mirror intervals.  Thus the
affine normalizer gives exactly the possible stratum permutations.

To compute them, let $\mathcal V$ be the union of the fixed points of
all positive nontranslation elements.  It is a finite set.  The fixed
circles of the reflections, cut at $\mathcal V$, are open mirror
intervals or uncut mirror circles.  Take the $G$-orbits of the vertices
and of these mirror components; call their numbers $v$ and $e$.
There is one regular interior stratum in the quotient, since its
interior after removal of finitely many cone points is connected.
A vertex orbit may be removed once, a mirror-component orbit may
contain any nonnegative number of puncture orbits, and the regular
interior may contain any nonnegative number.  Thus the decoration is
\[
 (d_1,\ldots,d_v;c_1,\ldots,c_e;a),\qquad
 d_i\in\{0,1\},\ c_j,a\ge0,
\]
subject to the exact equation
\[
      \sum_i d_i|V_i|+\sum_j c_j\frac{|G|}{2}+a|G|=\nu.
                                                        \tag{B.27.5}
\]
Divide these finitely many decorations by the stratum permutations
of the affine normalizer computed above.  Counts are sufficient on
an open mirror interval or circle: an increasing interval map, or a
circle map of the prescribed degree, moves any two unlabelled finite
sets of the same size into each other.  The moves extend through
small strips in the quotient.  In the regular interior use disjoint
moving discs; its configuration space of finitely many unlabelled
points is connected.  Such quotient isotopies lift through the finite
orbifold covering.  They preserve all cone and corner charts and
therefore give equivariant homeomorphisms upstairs.  Conversely a
conjugacy preserves the displayed orbit sizes and the incidence of
these strata.  This proves completeness and injectivity of the
decorated classification.

All computations use exact integers on
$Q=(\frac1{144}\Z^2)/\Z^2$.  We record why this finite mesh loses no
component or conjugacy.  Coordinates of $T$ have denominator dividing
its order $k\le4$.  For the rotations of orders $2,3,4,6$, the inverses
of $I-r$ have denominators $2,3,2,1$, respectively.  Thus all solutions
needed for (B.27.2), all normalized origin changes, and all isolated
fixed points have denominators dividing $2k$ or $3k$.
The representatives (B.27.3) have denominator dividing $2k$.
For a reflection normalizer, the origin equation is a rank-one
integral equation with nonzero Smith invariant $1$ or $2$; every
connected solution component has a representative with denominator
dividing $4k$.  The integer $144$ is divisible by each of these
integers for $1\le k\le4$.  The code separately checks the inverse
matrix and particular-solution denominators for every positive affine
element.  On each reflection circle, moving by its primitive integral
direction divided by $144$ connects consecutive sampled points.
Every isolated intersection and every selected puncture is retained
as a vertex.  This graph is a subdivision of the actual fixed-circle
arrangement, so its components after removing punctures are exactly
the ovals and lines of the action.  The mesh has enough unused
representatives on each eligible stratum to place the at most four
punctures; the executed program explicitly verifies the prescribed
orbit counts and total cardinality for every decoration.

For each unremoved vertex $p$, enumerate all subgroups of its finite
cyclic or dihedral stabilizer that contain a positive nonidentity
element.  Each gives the section class with fixed point $p$.
For each reflection enumerate the components of its fixed circles
after deleting the punctures.  B.3--B.4 identify exactly these data
with all nontrivial partial-section classes, and identify their
centralizers as $1$ at points and lines and $\Z$ at ovals.
Containment of image subgroups and reverse containment of their fixed
components give all extension relations.  Applying every affine
matrix computes the group action on this poset and hence $\Phi$,
its orbit decomposition and every stabilizer.  At a positive fixed
point the local matrix is its actual linear part.  Its trace
$-2,-1,0,1$ gives order $2,3,4,6$, and the sign of
$\det(e_1,Ae_1)$ gives the sign of the angle except for a half-turn.
Reflections have multiplier $1/2$.  The program checks the order
and power identities on every element of every section.  It retains
rather than deletes the reflection section classes with several
maximal extensions.

Finally a nonidentity finite affine map acting trivially in the
punctured-torus mapping class group would have trivial homology
matrix and hence be a nonidentity translation.  It would also fix
each puncture, contradicting freeness of translations.  Thus each
action gives a faithful finite mapping class subgroup.  The realization
and uniqueness statements C.5--C.6 identify their conjugacy classes
with the topological action classes just enumerated.  The computation
returns the four counts in the statement.  For $\nu=1$, the multiset
of group orders, positive orders, full splitting counts and partial
splitting counts agrees, class for class, with the independent
amalgam calculation of B.25.  A first complete run with denominator
$72$ and the final run with denominator $144$ also give the same
four counts; the proof of completeness uses the latter mesh.
\end{proof}

\begin{claim}{B.28}{Proved}{Complete executed extended genus-two catalogue}
The extended mapping class group of the closed orientable genus-two
surface has exactly $92$ conjugacy classes of finite subgroups.  Of
these, $57$ contain the hyperelliptic involution and $35$ do not;
$21$ classes preserve orientation.  The files
\texttt{type\_2\_0\_extended\_catalogue.csv} and
\texttt{type\_2\_0\_extended\_sections.json} specify every group by
an exact finite permutation representation and give every full and
partial splitting, its multiplier and centralizer, all extension
incidences, $\Phi$, and its group orbits and stabilizers.
\end{claim}
\passage{\S{}15; folios 73--75; scans 148--151; \S{}15; folios 77--79; scans 155--160; \S{}18; folios 96--101; scans 193--203; \S{}33; folios 222--225; scans 440--445; Part B(f) of the task; no further manuscript assertion}

\begin{proof}
Let $\iota$ be the hyperelliptic mapping class.  The genus-two
Birman--Hilden homomorphism, with orientation reversal allowed, is
\[
 1\longrightarrow\langle\iota\rangle
 \longrightarrow\mathcal T_{2,0}
 \longrightarrow\mathcal T_{0,6}\longrightarrow1.       \tag{B.28.1}
\]
We recall its geometric content, to specify every comparison used here.
The involution is the deck map of the double cover of $S^2$ branched
at six points.  The covering is determined by the homomorphism from
the fundamental group of the six-times-punctured sphere to $C_2$
that sends each peripheral generator to the nonidentity element;
the product relation is respected because six is even.  Every
homeomorphism preserving the branch set preserves this homomorphism,
and hence has exactly two lifts, differing by $\iota$.
The genus-two hyperelliptic mapping class is central: the usual chain
of five nonseparating curves is invariant under the involution, their
Dehn twists generate the oriented genus-two mapping class group,
and a reflection of the standard branched-cover model commutes with
$\iota$ and supplies the negative coset.  The Birman--Hilden isotopy
property identifies the kernel in (B.28.1) with these two deck classes;
it also makes the lifting map onto.  This is the hyperelliptic
Birman--Hilden description used in B.25, in genus two instead of the
four-point pillowcase; see \cite[\S9.4.2, p.~257, with the proof in \S9.4.3]{A_FarbMargalit2012}.

For a finite $K\subset\mathcal T_{2,0}$ put
$E=K\langle\iota\rangle$ and let $H$ be its image in
$\mathcal T_{0,6}$.  The group $E$ is exactly the full inverse image
of $H$ in (B.28.1), and has order $2|H|$.  Thus B.26 supplies its
complete list of $57$ conjugacy classes.  A subgroup $K$ with this
fixed $E$ either equals $E$, or is a complement to
$\langle\iota\rangle$.  Every complement is normal in $E$ and is
uniquely the kernel of a homomorphism
\[
             \chi:E\longrightarrow\mathbf F_2,
             \qquad\chi(\iota)=1.                     \tag{B.28.2}
\]
Indeed $E=K\times\langle\iota\rangle$ for a complement, giving
this homomorphism; conversely the kernel of (B.28.2) has index two
and maps isomorphically onto $H$.  Because $\iota$ is central in the
whole mapping class group, conjugate $K$'s have conjugate associated
$E$'s.  For a fixed $E$, their further fusion is therefore exactly
the action of its normalizer on the characters (B.28.2).

Here are explicit presentations of all the finite groups $E$.
They are also used to calculate the sections, so the statement is
independent of the names assigned to abstract finite groups.  Let
$Q=S^2/H$ be one of the spherical orbifolds in B.26, with its
six-branch-point decoration.  Write $z=\iota$, $z^2=1$, and require
$z$ to commute with all generators.  An interior cone of order $m_j$
has a generator $x_j$ with
\[
                    x_j^{m_j}=z^{d_j},                 \tag{B.28.3}
\]
where $d_j=1$ exactly when its orbit is a branch orbit.  Let $a$ be
the number of regular interior branch orbits.  For a spherical
underlying quotient the remaining relation is
\[
                       \prod_jx_j=z^a.                 \tag{B.28.4}
\]
For a projective-plane quotient, with crosscap generator $u$, it is
$u^2\prod_jx_j=z^a$.

For a mirror disc, expand the cyclic corner list by inserting each
branch orbit on an open mirror edge as an extra entry $(m,d)=(1,1)$.
An original corner contributes $(m,d)$ with its original order $m$
and its branch flag $d$.  The distinction is retained as a color.
Let the resulting cyclic list be $(m_j,d_j)_{j=0}^{L-1}$.
There are reflections $R_0,\ldots,R_{L-1}$ and a boundary element
$e$, with
\[
 \begin{split}
 R_j^2&=1,\\
 \prod_i x_i\,e&=z^a,\\
 (R_jR_{j+1})^{m_j}&=z^{d_j},\\
 R_L&=eR_0e^{-1}.
 \end{split}
 \tag{B.28.5}
\]
When $L=0$ there is one reflection $R_0$ and the last relation is
$eR_0e^{-1}=R_0$.  A newly inserted mirror branch entry says exactly
$R_{j+1}=R_jz$.

These are presentations, rather than only relations satisfied by
some finite images.  To verify this, use the branched cover described
above.  Its full group of lifts of $H$ is a group of order $2|H|$.
A small positive circle about a branch point has deck monodromy $z$;
about an unbranched point it has trivial deck monodromy.  This gives
(B.28.3).  The product of the boundary loops gives (B.28.4), or its
crosscap or disc version.  At a branch point on a mirror, the local
map is $w\mapsto w^2$: the two reflections on the adjacent lifted
half-branches differ by $w\mapsto-w$.  This gives precisely the
inserted entry in (B.28.5).  At a corner the positive local cyclic
generator has its order doubled exactly when it is branched, giving
the other entries.  These local monodromies and $z$ generate the full
lift group.  Conversely quotienting the displayed presentation by
$z=1$ gives the original spherical orbifold presentation, which is
$H$ by the local-injectivity and Euler-characteristic certificate in
B.26.  Since $z$ is central of order at most two, the presented group
has at most $2|H|$ elements.  Its surjection onto the geometric lift
group proves equality and the presentation assertion.

We determine the entire normalizer action needed in (B.28.2), not
just conjugation inside the finite group.  A normalizing mapping
class descends through (B.28.1) and through the $H$-action to an
orbifold homeomorphism of $Q$ preserving its branch decoration.
Conversely any such orbifold homeomorphism lifts to $S^2$, because
$S^2$ is the universal orbifold cover of $Q$, and then to the double
branched cover, because the six peripheral mod-two monodromies are
preserved.  Thus its lifts give the whole normalizer action, with
inner automorphisms and the deck involution as the harmless lifting
ambiguities.  Both act trivially on characters to $\mathbf F_2$.

On a sphere, an orbifold homeomorphism sends each local cone
generator to a conjugate of a local generator with the same original
order and branch flag, or to its inverse when it reverses
orientation.  Characters ignore conjugation and inversion.  All
permutations among cones of the same pair $(m_j,d_j)$ are realized
by homeomorphisms of the sphere.  Point-pushing and pure braids
conjugate the local generators and therefore act trivially on these
characters.  The ordinary branch generators already equal the
central element $z$.  Consequently the exact action on (B.28.2)
is the permutation action on these colored cone coordinates.
The character equations are
\[
       m_j\chi(x_j)=d_j,\qquad
       \sum_j\chi(x_j)=a\pmod2.                       \tag{B.28.6}
\]
This includes, for example, the two elliptic complements in a
four-group lying over a sphere involution with no fixed branch
points: exchanging the two cone points exchanges those complements.
They must be fused, and the program does so.

On a mirror disc there is at most one interior cone.  The normalizer
preserves the cyclic order on its boundary, up to reversal.  It
therefore induces exactly the rotations and reversals of the
expanded colored corner list that preserve its colors.  Each such
symmetry is realized by a disc homeomorphism preserving the interior
decoration.  To see its action explicitly, extend the reflection
sequence by $R_{j+L}=eR_je^{-1}$.  A cyclic shift by $s$ sends
$R_j$ to $R_{j+s}$ and preserves $e$; a reversal sends $R_j$ to
$R_{s-j}$ and $e$ to $e^{-1}$.  The interior cone generator is sent
to a conjugate of itself or its inverse.  Modulo two these formulas
are literal permutations of the reflection coordinates and leave
the interior coordinate fixed.  They respect every relation in
(B.28.5).  Every remaining disc mapping class comes from moving the
interior marked points or from a boundary twist; its effect on local
generators is conjugation and hence is trivial on the characters.
The character equations are therefore
\[
 \begin{split}
 m_i\chi(x_i)&=d_i,\\
 \chi(e)&=a+\sum_i\chi(x_i),\\
 m_j\bigl(\chi(R_j)+\chi(R_{j+1})\bigr)&=d_j\pmod2.
 \end{split}
                                                        \tag{B.28.7}
\]
The program enumerates these bit vectors and divides by exactly the
colored dihedral permutations just proved.  For a projective-plane
quotient, the six-branch-point equation allows only the three cases
$P()$ with $a=3$, $P(2)$ with $(d,a)=(1,1)$, and $P(3)$ with
$(d,a)=(0,1)$.  Their full lifts are respectively $C_4,C_8,C_{12}$,
with $z$ the involution.  None has a character taking $z$ to one.
Thus no further complement fusion is needed in these cases.

Executing (B.28.6)--(B.28.7) gives $35$ character orbits, in addition
to the $57$ full groups.  The following order distribution is a
compact view of the complete $92$-row table:
\[
\begin{array}{c|rrrrrrrr}
 |K|&1&2&3&4&5&6&8&10\\\hline
 \#&1&7&1&13&1&11&14&3\\[2pt]
 |K|&12&16&20&24&32&48&96\\\hline
 \#&17&8&1&9&1&4&1
\end{array}                                             \tag{B.28.8}
\]
In particular its seven involution classes consist of the hyperelliptic
class, the elliptic positive class, and the five negative classes
predicted independently by B.16.

We finish by specifying the verification of all section data.
For each of the $57$ presentations, the program performs an exact
Todd--Coxeter enumeration and obtains the regular action on exactly
$2|H|$ cosets.  It evaluates every relator in these permutations and
checks that $z$ is nonidentity of order two.  In the quotient orbifold
for the genus-two action, an originally branched cone of order $m$
now has order $2m$, a regular branch orbit contributes a cone of
order two, and a branched corner has its order doubled.  Every
local cyclic or dihedral group injects into the computed finite
group.  The orbifold Euler characteristic multiplied by $2|H|$
is $-2$, and the covering subgroup is the kernel of a
monodromy homomorphism compatible with orientation.  Its torsion-free
covering surface is therefore orientable of genus two.
The program also checks that $z$ has exactly six section classes.

The cone and corner cosets, together with the lifted mirror coset
graph of B.26, give every fixed point and reflection oval.  Enumerating
subgroups of these local stabilizers and conjugating them by the
finite group is enough: every finite subgroup with a nonempty fixed
set is contained in such a local stabilizer.  The section multiplier
is read from the power of the corresponding positive local generator;
conjugating by a negative element reverses that angle.  For each
complement, retain exactly the section classes whose image subgroup
lies in its character kernel.  This restriction is exact because
the lift extension for the complement is the inverse image of that
subgroup in the full lift extension.  Recompute maximality in the
retained poset, intersect the section stabilizers with the complement,
and compute its actual orbits.  Thus no maximality or orbit count is
inferred merely by restricting the full-group summary.
All order and power identities for $\mu$ and all unique-maximal-extension
claims for sections containing a rotation are checked.  Finally the
$21$ positive rows agree with the independent primary-source
oriented genus-two catalogue used in B.33 on the complete multiset
of five invariants: group order, full sections, nontrivial partial
sections, $|\Phi|$, and $|\Phi/K|$.
\end{proof}

\begin{claim}{B.29}{Proved}{Complete executed punctured genus-two catalogues}
The extended mapping class groups of types $(2,1)$ and $(2,2)$ have
respectively $32$ and $121$ conjugacy classes of finite subgroups.
Their positive subgroups form $10$ and $23$ of these classes.
The files \texttt{type\_2\_\nu\_extended\_catalogue.csv} and
\texttt{type\_2\_\nu\_extended\_sections.json}, for $\nu=1,2$,
list every class with its exact group, puncture orbits and stabilizers,
all section classes and their incidences, $\Phi$, its group action
and stabilizers, centralizers, and all section and puncture multipliers.
\end{claim}
\passage{\S{}15; folios 73--75; scans 148--151; \S{}15; folios 77--79; scans 155--160; \S{}18; folios 96--101; scans 193--203; \S{}33; folios 222--225; scans 440--445; Part B(f) of the task; no further manuscript assertion}

\begin{proof}
Fill the punctures of a finite action on the genus-two surface.
The resulting finite subgroup in the closed mapping class group is
faithful: a finite homeomorphism isotopic to the identity on a closed
hyperbolic surface is the identity, by B.6.  Thus B.28 supplies the
complete list of closed actions that must be decorated.
For such an action $K$, keep the notation
$E=K\langle\iota\rangle$, $H=E/\langle\iota\rangle$, and
$Q=S^2/H$ from that proof.

The strata relevant to punctures are intrinsic to $K$.
Let $\mathcal V_K$ be the set of points whose stabilizer in $K$
contains a positive nonidentity element.  Take its $K$-orbits.
Cut the fixed circles of the reflections in $K$ at
$\mathcal V_K$ and take the $K$-orbits of the resulting open
intervals or uncut circles.  These are the mirror strata.
There is one connected regular interior in the quotient $S_2/K$.
A vertex orbit is selected zero or one time; a mirror stratum
receives any nonnegative number of puncture orbits; so does the
regular interior.  Their weights are, respectively,
\[
       |Kp|,\qquad |K|/2,\qquad |K|.                  \tag{B.29.1}
\]
The weighted sum must equal $\nu$.
Points fixed by $\iota$ but having no positive stabilizer in $K$
are ordinary points of these $K$-strata.  They introduce no extra
puncture label.  In particular punctures need not be preserved
by $\iota$ when $K$ is a complement in $E$.

We prove the normalizer permutations used to identify these
decorations.  Every mapping class normalizing $K$ normalizes $E$,
since $\iota$ is central in $\mathcal T_{2,0}$.  At the homeomorphism
level it may be represented by a normalizer of the chosen finite
action $E$ without changing its action on the intrinsic $K$-strata.
Indeed compare the two commuting involutions $\iota$ and
$h\iota h^{-1}$ after conjugating $K$ by $h$.  The finite actions of
the abstract extension $K\langle\iota\rangle$ have the same mapping
classes.  Uniqueness of realization in C.6 gives a conjugating
homeomorphism isotopic to the identity and commuting with $K$.
C.7 gives an isotopy through such commuting homeomorphisms.
It fixes the isolated $K$-fixed points individually and preserves
every reflection component.  Composing with it makes $h$ normalize
the chosen $E$ and leaves the required stratum permutation unchanged.
Thus it suffices, and is necessary, to compute the action of the
normalizer of $E$ that preserves $K$.

The full-group case $K=E$ is direct.  A cone in $Q$ has the color
consisting of its original order and its branch flag.
An ordinary branch orbit has its own color $(1,1)$.
The normalizer permits all permutations of identical interior
cone colors and exactly the dihedral symmetries of the expanded
colored boundary list of B.28.  Each is realized by a quotient
orbifold homeomorphism and hence lifts.  Pure interior braids
conjugate local generators, so they leave the corresponding
$E$-orbits unchanged.  These observations prove that this finite
permutation group is the whole normalizer action on the strata.

Suppose now $K=\ker\chi$ with $\chi(z)=1$.
For a local stabilizer $L\subset E$, its point orbit is $E/L$.
It is a single $K$-orbit when $\chi(L)=\mathbf F_2$.
When $\chi(L)=0$, it consists of two $K$-orbits, labelled
\[
                       b=\chi(g),\qquad gL\in E/L.    \tag{B.29.2}
\]
This is well defined precisely in the latter case.
It gives the exact orbit comparison, including the ambiguity in
choosing a lift of a quotient homeomorphism: composing with $z$
simultaneously changes every available sheet label $b$ to $b+1$.
The code includes this deck permutation.

An interior cone with generator $x$ contributes an intrinsic
$K$-vertex only when $K\cap\langle x\rangle$ is nontrivial.
For cone permutations, retain exactly the equal-color permutations
that also preserve the character values on their local generators.
On a sphere quotient, every split cone pair has an additional
independent interchange.  Here is the explicit lift proving that
assertion.  Choose a branch local generator $y$ with $\chi(y)=1$;
one exists because the branch decoration is nonempty.
An ordinary branch generator is $z$.  A branched old cone has
$y^m=z$ with $m$ odd whenever a character with $\chi(z)=1$ exists,
and therefore also has $\chi(y)=1$.
For a split cone generator $x$, we have $\chi(x)=0$.
A pure braid about these two quotient points acts by
\[
       x\longmapsto(xy)x(xy)^{-1},
       \qquad y\longmapsto xyx^{-1}.                 \tag{B.29.3}
\]
The other local generators are unchanged.
The character is preserved.  The conjugator for $x$ has character
one, so (B.29.2) interchanges its two $K$-orbits; the conjugator for
$y$ has character zero.  Thus it leaves the other split pairs
unchanged.  This proves the claimed independent interchanges.
Together with cone permutations they generate the whole stratum
action: pure braids generate the remaining mapping classes of a
punctured sphere, and their action on the two-element fibers can
only be these interchanges.

For a disc quotient there is at most one old interior cone.
An ordinary interior branch point permits exactly the same
interchange (B.29.3) of that split cone pair, supported away from
the boundary.  When there is no such interior branch point, the
disc has at most one interior marked cone; its pure braid group
introduces no additional permutation.  Boundary permutations are
computed as follows.  Let $L$ be the length of the expanded corner
list, write $R_{j+L}=eR_je^{-1}$, and put
$\epsilon=\chi(e)$.  A colored cyclic shift sends a corner index
$j$ to $s+j$ and an edge index $i$ to $s+i$.
A reversal sends them to $s-j-1$ and $s-i$, respectively.
For either image write the resulting integer as $qL+t$,
$0\le t<L$.  Its local generator is conjugated by $e^q$;
hence the sheet label changes by
\[
                          b\longmapsto b+q\epsilon.   \tag{B.29.4}
\]
Retain precisely the colored symmetries that preserve the reflection
character values, and include the global deck interchange.
These formulas give both vertex and mirror-component permutations.
For a cornerless boundary the same rule is its circular covering
with monodromy $\chi(e)$; the lifted mirror graph records whether
the two local sheet labels belong to one or two components.

There are no omitted boundary permutations.  An orbifold
homeomorphism of the disc preserves its cyclic boundary order
up to reversal.  After its colored dihedral permutation has been
removed, its remaining mapping class is an interior braid or a
boundary twist.  An interior braid has the local conjugators
already described and fixes the boundary.  A boundary twist has
the conjugator $e$, already represented by (B.29.4) and the deck
permutation.  Inner automorphisms coming from $K$ leave $K$-orbits
unchanged.  The projective-plane cases in B.28 have no complements;
their full-group decorations require only permutations of identical
interior cone colors.  Thus every case has an explicit complete
finite normalizer permutation group.

We spell out how this calculation handles smooth points of the
$K$-mirror locus that were corners for $E$.  The code starts with
the exact finite mirror coset graph of the full action $E$.
It finds the vertices with positive stabilizer in $K$, removes
only these vertices to form the open mirror strata, and computes
the fixed circles of each actual reflection in $K$.
Thus an $E$-corner that has only a reflection stabilizer in $K$
lies inside a single $K$-mirror stratum.  The sheet and boundary
formulas above are applied to an ordinary edge sample in each
such component.  They determine its whole component by
connectedness.  This is also a direct check that a puncture can
move through a noncritical Weierstrass point whenever the
$K$-stratum allows it.

Enumerate all vectors with the weights (B.29.1) summing to
$\nu=1$ or $2$, and divide by the proved finite permutation group.
These are precisely the topological action classes with punctures.
Indeed unlabelled finite subsets of an interval or circle with
the same cardinality are related by a homeomorphism of prescribed
degree, and the regular interior configuration space is connected.
Extend such moves by small discs and strips in the quotient and
lift the isotopies through the finite orbifold cover.  Conversely
an equivariant homeomorphism preserves the intrinsic strata,
their weights, and a stratum permutation in the normalizer just
computed.  C.5--C.6 then identify these action classes with the
finite-subgroup conjugacy classes in the punctured mapping class
group.

The executed program \texttt{catalogue\_genus2\_punctures.py}
returns $32$ and $121$ classes.  It places the punctures on the
exact coset graph, with one midpoint inserted in every edge before
any vertex is removed.  This subdivision guarantees that even
two punctures at formerly adjacent samples leave the intervening
open interval represented.  The program checks that an element
stabilizing an edge fixes both its endpoints, so the inserted
midpoint has exactly the edge stabilizer and creates no artificial
critical point.  It records regular puncture orbits by the free
$K$-set $K$, with trivial stabilizer.

All positive section classes from B.28 are retained exactly when
their fixed point has not been removed and their image subgroup
lies in $K$.  For every reflection, remove the punctures from its
actual fixed-circle graph and recompute the resulting ovals and
lines.  These operations give all partial-section classes by
B.3--B.4.  Their inclusion poset, maximal classes, $K$-action,
stabilizers and centralizers are then recomputed.
Local multipliers at punctures are retained from the positive
local generators or assigned $1/2$ to a reflection.
The positive portions contain $10$ and $23$ classes and agree
with the independent primary-source calculation in B.33 on the
full multiset of group order, full sections, partial sections,
$|\Phi|$, and $|\Phi/K|$.  The negative involutions provide
another independent check: there are $4$ classes with one
puncture and $11$ with two, exactly the decorations of the five
closed real types in B.16.
\end{proof}

\subsection{The full tube action and the strip subgroups}

\begin{claim}{B.30}{Proved}{Explicit outer words for the antipodally paired tube model}
For $g\ge2$, put $n=2g$ and label the equatorial holes by $j\in\Z/n$,
with $j$ glued by a tube to $j+g$. The following presentation and
substitutions compute the entire exterior action requested on scan 467:
\[
 \pi_g=\left\langle x_0,\ldots,x_{g-1},p_0,\ldots,p_{g-1}
 \ \middle|\ R=1\right\rangle,
 \qquad
 R=x_0\cdots x_{g-1}
   \prod_{j=0}^{g-1}(p_j^{-1}x_j^{-1}p_j).
\]
Extend the notation, with all subscripts read modulo $n$, by
\[
 \begin{gathered}
 x_{j+g}=p_j^{-1}x_j^{-1}p_j,\qquad p_{j+g}=p_j^{-1},\\
 u_j=(x_0\cdots x_j)^{-1},\quad u_{n-1}=1,\qquad
 q_j=u_j^{-1}p_j u_{j+g-1}.
 \end{gathered}
\]
Let $r=\varepsilon_0\varepsilon_1$ and $s=\varepsilon_1$ as in
B.20. Representatives of their outer automorphisms, and of $\tau$ and
$\rho_k$, are specified on the $2g$ displayed generators by
\[
\begin{array}{c|cc}
 &x_j&p_j\\ \hline
 r&x_{j+1}&p_{j+1}\\
 s&u_{-j}^{-1}x_{-j}u_{-j}&q_{-j}\\
 \tau&u_{j+g-1}^{-1}x_{j+g}^{-1}u_{j+g-1}&q_j^{-1}\\
 \rho_k&x_j&x_j^{\delta_{jk}}p_j.
\end{array}
\]
The map to $\operatorname{Out}(\pi_g)$ defined by these substitutions
is the geometric action of
$\Z^g\rtimes(D_{4g}\times C_2)$ of B.20. The finite factor acts by
\[
 r\rho_jr^{-1}=\rho_{j+1},\qquad
 s\rho_js^{-1}=\rho_{-j},\qquad
 \tau\rho_j\tau^{-1}=\rho_j^{-1},
\]
with indices in these three formulas taken modulo $g$.
\end{claim}
\passage{\S{}33; folios 234--236; scans 463--467}

\begin{proof}
We specify the base paths so that the formulas can be used as words,
rather than only as unbased pictures. Let $N,S$ be the north and south
poles in the sphere with the $n$ discs removed. Draw disjoint meridian
arcs $d_j:N\to S$ through the gap immediately following hole $j$.
The loop about that hole, based at $N$, is
$x_j=d_{j-1}d_j^{-1}$. Collapse the maximal-tree edge $d_{n-1}$ and
identify $N,S$ by that edge. The remaining arcs become
$u_j=(x_0\cdots x_j)^{-1}$, because the product telescopes to
$d_{n-1}d_j^{-1}$. In particular $\prod_{j=0}^{n-1}x_j=1$.

Let $p_j$ run from $N$ to the northern side of hole $j$, through the
tube to hole $j+g$, and back to $N$ on its northern side. The opposite
traversal is $p_{j+g}=p_j^{-1}$. Its boundary identification reads
$x_{j+g}=p_j^{-1}x_j^{-1}p_j$. Van Kampen therefore gives exactly the
presentation in the claim. The parallel path based at $S$, on the
southern side of the tube, is
$q_j=d_j^{-1}p_jd_{j+g-1}$ after inserting the northern paths, which
becomes the displayed formula on collapsing the tree edge.
These identities can also be taken as a two-vertex cellular description:
the sphere faces give the telescoping boundary relation and each tube
face gives the indicated boundary conjugation.

Rotation about the polar axis sends hole $j$ to $j+1$, preserves $N$,
and sends $p_j$ to $p_{j+1}$. This proves the $r$ row. The half-turn
$s$ about the equatorial axis through hole $0$ exchanges $N,S$ and
sends hole $j$ to $-j$. The southern boundary loop at hole $k$ is
$d_k^{-1}x_kd_k$, and the southern tube path is $q_k$, giving the
$s$ row. The antipodal map exchanges $N,S$, sends $j$ to $j+g$,
and reverses orientation. It therefore inverts the southern boundary
loop and reverses the corresponding tube path. The boundary-loop
formula written with $u_{j+g-1}$ equals that written with $u_{j+g}$,
since $u_{k-1}=x_ku_k$. This proves the $\tau$ row. Finally a positive
full twist of tube $k$ inserts $x_k$ into the crossing $p_k$ and
leaves all other chosen paths unchanged. This proves the $\rho_k$ row.
Each row comes from an actual homeomorphism, so it is an automorphism
of the stated group and has the prescribed orientation sign.

For $s$ and $\tau$ the chosen path from $N$ to the image of $N$ is
$d_{n-1}$. We consequently get actual automorphisms only after this
base-path choice. For example their squares are
\[
 s^2=\operatorname{Inn}(x_0),\qquad
 \tau^2=\operatorname{Inn}(x_0\cdots x_{g-1}),
\]
whereas $r^{2g}=1$ as an actual automorphism. These formulas follow
by concatenating the base path with its image: $s(d_{n-1})=d_0^{-1}$
and $\tau(d_{n-1})=d_{g-1}^{-1}$. Thus the displayed group relation
$s^2=\tau^2=1$ holds in the outer group with precisely the stated
inner ambiguity. The geometric symmetries satisfy $srs=r^{-1}$ and
commute with the antipodal involution on the surface. They preserve
and permute the tube collars; a positive symmetry preserves twist
sign and the antipodal map reverses it. This proves every relation
of the asserted semidirect product in the outer group. Changing the
maximal-tree edge conjugates each based comparison by its corresponding
base-path word and leaves this exterior action unchanged.

The accompanying programs check all relator images for $2\le g\le8$;
$r,s$ carry $R$ to a conjugate of $R$, and $\tau$ carries it to a
conjugate of $R^{-1}$. They additionally check all the finite-factor
and twist-conjugacy relations in the surface group, with the actual
inner conjugators above. For instance $srsr=1$ and
$\tau r\tau r^{-1}=\operatorname{Inn}(x_0\cdots x_g)$.
The word checker uses the symmetrized relator: its largest common
piece has length one, so the length-$4g$ presentation satisfies
$C'(1/6)$ for $g\ge2$, and the Dehn reduction is exact. These finite
certificates supplement the cellular proof for arbitrary $g$.
\end{proof}

\begin{claim}{B.31}{Proved}{The eight symmetries of one tube}
The stabilizer of tube $T_0$ in $D_{4g}\times\langle\tau\rangle$ is
$\langle r^g,s,\tau\rangle\cong C_2^3$. There are coordinates
$(\theta,t)\in(\R/2\pi\Z)\times[-1,1]$ on that tube in which its
three generators act as
\[
 r^g(\theta,t)=(-\theta,-t),\qquad
 s(\theta,t)=(\theta+\pi,t),\qquad
 \tau(\theta,t)=(\theta+\pi,-t).
\]
These formulas agree with the outer action in B.30 and with the
manuscript's free antipodal extension, after the explicit coordinate
change $\theta_{\rm new}=\theta_{\rm old}+\pi t/2$.
\end{claim}
\passage{\S{}33; folios 229--230; scans 453--454; \S{}33; folios 233--236; scans 461--467}

\begin{proof}
The stabilizer of the unordered pair $\{0,g\}$ under the polygon's
dihedral group is $\{1,r^g,s,r^gs\}$: a stabilizing rotation is either
$1$ or $r^g$, and every stabilizing reflection is one of these followed
by $s$. The independent central antipodal factor doubles it. In a
symmetric tube, $r^g$ interchanges its ends and reverses its angular
direction, while $s$ preserves both ends and rotates the cross section
by $\pi$. The displayed representatives commute, are involutions,
and have eight distinct combinations; the antipodal representative
is free. They extend the same boundary maps as the corresponding
sphere symmetries and hence give the desired tube extension.

The formula of scan 453, retained through scans 461--462, is
$\tau_{\rm old}(\theta,t)=(\theta+\pi+\pi t,-t)$.
Conjugating it by the stated change gives exactly
$(\theta+\pi,-t)$; this also transports both boundary identifications.
Thus it is an explicit isomorphism of the two actions. The full twist
is $(\theta,t)\mapsto(\theta+\pi+\pi t,t)$ in either coordinate.
Conjugation by the two positive generators preserves its class, while
conjugation by the negative generator inverts its class, agreeing with
B.30. All eight maps and their products are checked symbolically in
\texttt{code/verify\_tube\_local.py}.
\end{proof}

\begin{claim}{B.32}{Proved}{Relations for arbitrary strips and a uniform obstruction to generation}
Choose for each $c_j$ one simple transversal $d_j$, using a strip in
the complement of the other tubes, so that
$i(c_i,d_j)=\delta_{ij}$. Let $W_j$ be a regular neighbourhood of
$c_j\cup d_j$, a one-holed torus, and put
$A_j=T_{c_j}=\rho_j$, $B_j=T_{d_j}$. Its supported subgroup is the
central extension $\langle A_j,B_j\rangle\cong B_3$ of B.21, with
quotient $SL_2(\Z)$. The exact pairwise relations are
\[
 [A_i,A_j]=1,\quad [A_i,B_j]=1\ (i\ne j),\qquad
 A_jB_jA_j=B_jA_jB_j.
\]
For distinct $i,j$, the two supported subgroups commute precisely when
$i(d_i,d_j)=0$. For intersection one, $B_i,B_j$ obey the braid
relation and do not commute; for intersection at least two they freely
generate a free group of rank two. For every such choice of one strip
per tube and every $g\ge2$, these $g$ supported subgroups, together
with $H$, generate a proper subgroup of $\mathcal T_g^+$.
\end{claim}
\passage{\S{}33; folios 236--237; scans 467--470}

\begin{proof}
Each $W_j$ has genus one and one boundary component because the two
core curves meet transversely once. The extension-by-identity map and
its faithful boundary-twist kernel are the explicit maps of B.21.
The disjointness conditions give the first two commutation relations,
and the calculation on its two free generators in B.21 proves the
braid relation. For any two essential simple curves $a,b$, their
twists commute exactly when they are disjoint: conjugation sends
$T_b$ to $T_{T_a(b)}$, and
$i(T_a(b),b)=i(a,b)^2$. Hence intersecting curves cannot commute.
If $i(a,b)=1$, their regular neighbourhood is a one-holed torus, so
the same braid calculation applies.

For $k=i(a,b)\ge2$ the freedom assertion has an elementary ping-pong
proof. On isotopy classes of essential curves set
\[
 X_a=\{x:i(x,a)<i(x,b)\},\qquad
 X_b=\{x:i(x,b)<i(x,a)\}.
\]
Both are nonempty and disjoint. Counting crossings in a twist collar
and cancelling only crossings already present outside it gives
\[
 i(T_a^m(x),b)\ge |m|\,i(x,a)k-i(x,b).
\]
For $m\ne0$ and $x\in X_b$ this is strictly greater than
$i(x,a)=i(T_a^m(x),a)$, so $T_a^m(X_b)\subset X_a$.
The symmetric inclusion holds for powers of $T_b$. The ping-pong
argument on a reduced alternating word proves that it is nontrivial:
apply it to the appropriate one of $X_a,X_b$, inserting one further
nonzero power from the opposite factor when the first and last
factors agree. Thus $\langle T_a,T_b\rangle\cong\Z*\Z$.
These statements compute the required pairwise relations for arbitrary
intersections of the selected strips; no disjoint-strip convention is
being used to conclude the general generation assertion.

For that assertion work on $V=H_1(X_g;\F_2)$ with its nonsingular
intersection pairing. The $2g$ classes $[c_j],[d_j]$ form a basis:
pairing a linear relation first with the $[c_j]$ kills all coefficients
of the $[d_j]$, and pairing next with the $[d_j]$ kills the rest.
There is a unique quadratic refinement $q:V\to\F_2$ with
$q([c_j])=q([d_j])=1$. Explicitly, for this ordered basis $(v_i)$,
\[
 q\Bigl(\sum_i a_iv_i\Bigr)
 =\sum_i a_i+\sum_{i<j}a_ia_j\langle v_i,v_j\rangle.
\]
It obeys $q(v+w)=q(v)+q(w)+\langle v,w\rangle$ by expansion.
A twist about a curve of class $v$ acts by
$x\mapsto x+\langle x,v\rangle v$; substitution gives
\[
 q(T_vx)-q(x)=\langle x,v\rangle\bigl(q(v)+1\bigr).
\]
Every generator $A_j,B_j$ therefore preserves $q$, without any
restriction on the intersections among the $d_j$.

Join $c_1$ and $c_2$ by an embedded arc in their connected complement
and form a band sum $e$. This is a simple curve of nonzero class
$[e]=[c_1]+[c_2]$, so it is nonseparating. Since the two cores are
disjoint, $q([e])=0$. Taking $x=[d_1]$ gives
$\langle x,[e]\rangle=1$, so $T_e$ fails to preserve $q$.
Consequently $T_e$ belongs to $\mathcal T_g^+$ and lies outside the
group generated by the strip subgroups. This proves properness for
all choices under discussion. The quadratic calculation has been
exhaustively checked for every possible matrix of transversal
intersections modulo two in genera two, three and four; the general
identity is the displayed algebraic proof.

The rank-one comparison is thus the specified central quotient
$B_3\to SL_2(\Z)$ and its action on the homology of $W_j$, together
with the above braid and intersection relations. The subgroup they
generate fixes the explicit quadratic refinement. This supplies the
precise algebraic maps and the explicit obstruction governing the
proposed generation in scan 470.
\end{proof}

\begin{claim}{B.33}{Proved}{Executed positive catalogues in genera two and three}
For $(g,\nu)=(2,0),(2,1),(2,2),(3,0)$, the numbers of conjugacy
classes of finite orientation-preserving subgroups, allowing conjugacy
in the full extended mapping class group, are respectively
\[
                         21,\qquad10,\qquad23,\qquad55.
\]
The trivial subgroup is included. The files
\texttt{oriented\_higher\_genus\_catalogue.csv} and
\texttt{oriented\_higher\_genus\_sections.json} give each class,
its actual finite permutation group, every subgroup's number of
sections, all nontrivial partial-section classes and their inclusion
relations, the $G$-set $\Phi$, its stabilizers, and every value of
$\mu$ on every section. Punctures are unlabelled. These assertions
concern positive subgroups; they do not enumerate subgroups containing
orientation-reversing elements.
\end{claim}
\passage{\S{}15; folios73--75; scans148--151; \S{}17; folios77--79verso; scans155--160; \S{}18; folios96--101; scans193--203; Part B(f) of the task; positive higher-genus cases}

\begin{proof}
We specify the classification input and the computation using it.
Karab\'a\v{s}--Nedela--Skyvov\'a
\cite[Theorem 2.2, Corollary 2.3, Theorem 4.1 and \S8]{R-KNS}
identify the positive topological classes with the orbits of smooth
epimorphisms from the Fuchsian group
\[
 \Delta=\langle x_1,\ldots,x_r,a_1,b_1,\ldots,a_h,b_h\mid
 x_i^{m_i}=1,\ [a_1,b_1]\cdots[a_h,b_h]x_1\cdots x_r=1\rangle
\]
under target automorphisms and orientation-preserving orbifold
automorphisms. Here smooth means that every image of $x_i$ has
order exactly $m_i$. Their \S8 gives a complete genus-two list
of 21 classes and a genus-three list of 55 classes in 49
group/signature pairs, including the trivial actions. We opened
the accompanying data \cite{R-KNSData}, downloaded the two original
JSON files, and retained their attribution and SHA-256 digests.
The attached program reads those exact files. It verifies that
each generating vector generates the group of the stated order,
each elliptic generator has its specified order, every presentation
relation holds, and
\[
 |G|\left(2-2h-\sum_i(1-1/m_i)\right)=2-2g.
\]
The conversion from the source's right-acting permutations to the
left-acting permutations in our code is the isomorphism sending
each permutation to its inverse. For a handle word the two
commutator conventions are compared explicitly; replacing both
handle generators by their inverses converts $a^{-1}b^{-1}ab$
to the convention $aba^{-1}b^{-1}$. The files retain the actual
generating vector, so this convention does not affect a recorded
elliptic rotation.

The local-order checks imply that the kernel of each epimorphism
is torsion-free: every finite-order element of a Fuchsian group
is conjugate to a power of an elliptic generator. The covering
formula then identifies its quotient surface as genus $g$.
Conversely a positive finite action determines this epimorphism
by lifting paths in the quotient. The two constructions are
inverse modulo the stated source and target automorphisms.
Finally C.5--C.6 identify these action classes with conjugacy
classes of finite mapping-class subgroups. This supplies the
exact reduction from the classification input to the groups
asked for here.

We also recompute the equivalence relation whenever $h=0$,
rather than using only the published class count. For an ordered
generating vector $(u_1,\ldots,u_r)$ with product one, the braid
generator acts by
\[
 \sigma_i:(u_i,u_{i+1})\longmapsto
               (u_i u_{i+1}u_i^{-1},u_i).
\]
Use adjacent such generators when $m_i=m_{i+1}$ and all pure
braid generators, obtained by moving the $j$th entry next to the
$i$th, applying $\sigma_i^2$, and reversing the moves.
Theorem 4.1 of the cited paper proves that these transformations
induce the full orientation-preserving orbifold automorphism
action after inner automorphisms have been divided out. A
generating vector is compared modulo $\Aut(G)$ by its rooted,
edge-coloured Cayley graph: breadth-first enumeration with the
ordered generators gives a canonical transition table. Equality
of tables gives the explicit bijection of vertices preserving
every labelled multiplication edge; induction on words proves
that it is the unique target-group isomorphism carrying the
generating vectors to one another. Thus the canonicalization
does not approximate group automorphisms.

To allow conjugacy by negative homeomorphisms, we add the
reflection of the marked sphere. Its effect on the tuple is
\[
 (u_i)_{i=1}^r\longmapsto
 \left(p_{i-1}u_i^{-1}p_{i-1}^{-1}\right)_{i=1}^r,
 \qquad p_{i-1}=u_1\cdots u_{i-1}.
\]
The product of the new tuple is $(u_1\cdots u_r)^{-1}=1$,
as successive cancellation shows. The program computes the
orbits both before and after adding this transformation; it
recovers the published planar counts and finds no fusion.
All nonplanar group/signature pairs in these two genera have
only one positive class in the input classification. Reversing
orientation preserves the group and signature and hence fixes
that class. Consequently the positive counts remain 21 and 55
under conjugacy in the extended group.

For genus two we add all puncture configurations with
$\nu\leq2$. An invariant finite puncture set consists of selected
elliptic orbits and a number $a$ of ordinary orbits. Write
$c_i\in\{0,1\}$ for selection of the $i$th elliptic orbit.
The exact equation is
\[
                  \sum_i c_i\frac{|G|}{m_i}+a|G|=\nu.
                                                               \tag{B.33.1}
\]
The program enumerates every solution. It carries the bits
$c_i$ along under each braid move, exchanging the two bits
when the corresponding two elliptic entries are exchanged.
It then takes orbits of these decorated generating vectors
under the same coloured braids, target automorphisms, and
reflection. These are exactly the decorated topological
classes: an equivariant homeomorphism descends to the quotient
orbifold, and its action on the covering is the compatibility
of its induced automorphism with the epimorphism just described.
Conversely that compatibility lifts the quotient homeomorphism
to an equivariant homeomorphism. Ordinary unlabelled punctures
can be moved to any other such configuration through paths in
the regular part of the quotient; point-pushing the paths is
isotopic to the identity after these additional points are
forgotten and so preserves the original covering kernel.
Thus $a$ records all their remaining data.

There are just two nonplanar genus-two entries. For the trivial
group, there is one configuration for each $\nu$. The other
has group $C_2$ and signature $(1;2,2)$. Its two elliptic
images are the unique involution, and exchanging their small
discs preserves the monodromy after an adjustment of the two
handle paths. The four possible pairs of handle images in
$C_2^2$ are connected by point-pushing either branch point
around a handle: crossing the branch cut adds its nonzero
monodromy to the transverse handle image. These two pushes
change either handle coordinate independently. Hence the
exchange lifts and only the number of selected branch orbits
matters. These elementary moves account for all four vectors
in the supplied data and prove the nonplanar decoration step
used by the code. Equation (B.33.1) and the decorated orbits
give 10 and 23 classes at one and two punctures.

We now derive every section entry from the checked finite
groups. A retained elliptic orbit is the left $G$-set
$G/\langle u_i\rangle$. Its point at the coset
$v\langle u_i\rangle$ has stabilizer
$v\langle u_i\rangle v^{-1}$ and positive generator
$v u_i v^{-1}$. B.3 gives exactly one kernel-conjugacy class
of sections for every point fixed by a specified nontrivial
subgroup $K\leq G$. Such a subgroup is cyclic because it is
contained in the cyclic point stabilizer. Conversely every
such containment gives a section. Thus the exact count is
\[
 \#\Sec(E_K\to K)/\pi
  =\sum_{i:c_i=0}
       \#\{v\langle u_i\rangle:
                    K\leq v\langle u_i\rangle v^{-1}\}.
                                                               \tag{B.33.2}
\]
The formula is used for every subgroup, which the program
enumerates by saturation under adjoining each remaining group
element. It checks that no noncyclic subgroup has a section.
The trivial subgroup has one section, and it is not included
in the nontrivial partial-section totals.

The maximal partial sections are precisely these individual
elliptic points, with $G$ acting by left multiplication on the
cosets. A partial section lies below another exactly when
they have the same point and its subgroup is contained in the
other subgroup. This determines the whole inclusion relation
and proves uniqueness of a maximal extension in the positive
case. Each nontrivial section has centralizer one by B.3.
The total number of nontrivial partial-section classes is
independently checked against
\[
             \sum_{i:c_i=0}\frac{|G|}{m_i}\bigl(\tau(m_i)-1\bigr),
\]
where $\tau(m)$ is the number of positive divisors of $m$.
Finally, for $h=v u_i^b v^{-1}$ the value of $\mu$ at this
section is $b/m_i\pmod\Z$. This lists every power of every
local generator and is exactly the intrinsic multiplier
identified in B.12. Puncture stabilizers are the removed
cyclic groups or the trivial group for ordinary orbits.
All these data are computed with integer permutations and
rational fractions. The output has 109 individual action
rows and 59 independent planar braid-orbit checks.
\end{proof}

\begin{claim}{B.34}{Proved}{The complete genus-three subcatalogue containing a hyperelliptic involution}
There are exactly $79$ conjugacy classes of finite subgroups of the
extended mapping class group of a closed orientable genus-three surface
that contain a hyperelliptic involution.  Exactly $18$ preserve orientation
and $61$ contain orientation-reversing elements.  The latter $61$ classes
are listed below and are specified, together with every section and its
extension incidences, in
\texttt{type\_3\_0\_hyperelliptic\_mixed\_catalogue.csv} and
\texttt{type\_3\_0\_hyperelliptic\_sections.json}.
\end{claim}
\passage{\S{}18; folios 96--101; scans 193--203; \S{}33; folios 222--225; scans 440--445; Part B(f) of the task; the genus-three hyperelliptic subcatalogue is an added result, not a further manuscript assertion}

\begin{proof}
Fix a hyperelliptic involution $z$ and its branched quotient
$p:S_3\longrightarrow S^2$, with branch set $P$ of size eight.
On the unbranched complements this is the connected double cover
classified by the homomorphism
\[
 \pi_1(S^2\setminus P)\longrightarrow C_2
 \quad\hbox{sending each of the eight meridians to }1.
                                                        \tag{B.34.1}
\]
Their sum is zero modulo two.  Every homeomorphism of the sphere
preserving $P$ preserves (B.34.1), and therefore has exactly two lifts.
The Birman--Hilden isotopy property gives the exact sequence
\[
 1\longrightarrow\langle z\rangle\longrightarrow
 C_{\mathcal T_{3,0}}(z)\longrightarrow\mathcal T_{0,8}
 \longrightarrow1.                                    \tag{B.34.2}
\]
Here the left group is the centralizer of this specified involution;
the map is induced by the quotient $p$.  The lifting and isotopy argument
is precisely the fully ramified hyperelliptic case of
\cite[\S9.4]{A_FarbMargalit2012}.  Allowing orientation reversal adds the
negative lifts of the same sphere homeomorphisms, with the same kernel.

The involution $z$ is the unique hyperelliptic involution inside any
finite subgroup $K$ containing it.  To verify the assertion directly,
realize $K$ by isometries and use the resulting invariant complex
structure.  A hyperelliptic involution acts as $-1$ on $H^1(S_3;\mathbb R)$:
its invariant subspace is the pullback of $H^1(S^2;\mathbb R)=0$ under
its quotient, as follows by averaging cochains, or by the transfer
outside the branch points and filling their meridians.  If $z'$ were
another, the finite-order holomorphic map $zz'$ would act trivially on
$H^1$.  A nonidentity such map has isolated fixed points, each of
positive local Lefschetz index; a fixed-point-free map has Lefschetz
number zero.  But its Lefschetz number from homology is $2-6=-4$.
Thus $zz'=1$, proving uniqueness.
Consequently every element of $K$ centralizes $z$, and any conjugacy
between two such finite groups maps their hyperelliptic involutions
to one another.  Once $z$ is fixed, (B.34.2) therefore gives a bijection
\[
 \left\{\begin{array}{c}\text{finite subgroup conjugacy classes}\\
 \text{in }\mathcal T_{3,0}\text{ containing }z\end{array}\right\}
 \longleftrightarrow
 \left\{\begin{array}{c}\text{finite subgroup conjugacy classes}\\
 \text{in }\mathcal T_{0,8}\end{array}\right\}.
                                                        \tag{B.34.3}
\]
The inverse takes the full preimage, of order twice the sphere group.
The inclusion of every kernel element follows from the assumption
$z\in K$.  This also proves that no additional fusion occurs when the
chosen hyperelliptic involution is forgotten.

For completeness, B.26's decorated spherical-orbifold enumeration
applies verbatim to eight marked points: a selected cone orbit has
size $|H|/m$, a selected corner orbit has size $|H|/(2m)$, an orbit
on an open mirror edge has size $|H|/2$, and a regular orbit has size
$|H|$.  Their sum is eight.  The cyclic and dihedral rotation orders
are at most eight, since every non-polar point orbit has at least that
many points.  The tetrahedral and octahedral cases are included.
An icosahedral orbit has at least twelve points, so cannot occur.
The spherical classification and the homeomorphism classification
of these decorations in B.26 prove exhaustiveness and irredundancy.
Executing the integer weight equation gives $79$ decorated quotients,
$18$ spherical and $61$ with a negative element.

For each decorated quotient, the full lift is given by the exact
central extension presentation (B.28.3)--(B.28.5), now with eight
branch points.  Its presentation proof uses (B.34.1) in exactly the
same way: the displayed group surjects onto the full geometric lift
group of order $2|H|$, while its quotient by the central element $z$
is $H$, so it has at most $2|H|$ elements.  The two groups are
therefore isomorphic by that surjection.  For a projective-plane
quotient its generators include both $z$ and the crosscap element;
their relation is $u^2\prod x_i=z^a$.  In particular no cyclicity
assumption from the six-branch-point case is imposed.

The program \texttt{catalogue\_genus3\_hyperelliptic.py} was executed.
For every presentation it computes an exact regular permutation
representation of degree $2|H|$, verifies every defining relator,
checks that $z$ has order two, and checks injectivity on every local
cyclic or dihedral stabilizer.  The resulting orbifold Euler
characteristic is $-4/|K|$, so its oriented surface kernel has genus
three.  Exactly eight point sections have image $\langle z\rangle$,
providing an independent branch-number check.
The finite coset incidence graph of B.26 then computes every point
and oval fixed component and every partial section.  Each point
section has trivial centralizer in the surface group, each oval
section has infinite-cyclic centralizer, and the rotation multipliers,
extensions, maximal sections $\Phi$, group orbits and stabilizers
are recorded explicitly in the JSON file.  The calculation retains
all intersections of reflection ovals at dihedral stabilizers.
\end{proof}

In the following table, $b$ is both the G3H class suffix and the decorated sphere-eight row number
in \texttt{type\_3\_0\_hyperelliptic\_catalogue.csv}; $N$ is the group
order; $p$ counts all nontrivial partial splitting classes; and $m,o$
are respectively $|\Phi|$ and $|K\backslash\Phi|$.  The CSV gives the
cone, corner, mirror and ordinary branch decorations of every row,
and the JSON gives its finite group generators and all relators.
\begin{longtable}{r|r|r|r|r}
b&N&p&m&o\\\hline
\endhead
2&4&8&8&4\\
3&4&10&10&5\\
4&4&12&8&5\\
5&4&16&8&6\\
6&4&20&8&7\\
7&4&24&8&8\\
11&8&12&12&3\\
12&8&14&14&4\\
13&8&20&12&4\\
14&8&28&12&5\\
15&8&14&8&3\\
16&8&22&8&4\\
17&8&20&12&4\\
18&8&26&12&5\\
19&8&34&12&6\\
20&8&24&12&5\\
21&8&32&12&6\\
22&8&32&12&6\\
23&8&24&8&4\\
24&8&32&8&5\\
25&8&32&8&5\\
29&12&12&8&2\\
30&12&14&10&3\\
31&12&24&8&3\\
32&12&30&8&3\\
33&12&42&8&4\\
34&12&42&8&4\\
39&16&16&12&2\\
40&16&28&20&3\\
41&16&40&20&4\\
42&16&28&12&3\\
43&16&18&14&3\\
44&16&32&12&3\\
45&16&38&20&4\\
46&16&52&20&5\\
47&16&50&20&5\\
48&16&48&12&4\\
49&16&48&12&4\\
50&16&36&12&3\\
51&16&48&12&4\\
52&16&44&12&4\\
56&24&54&20&3\\
57&24&36&8&2\\
58&24&28&8&2\\
59&24&68&20&4\\
60&24&66&20&4\\
61&24&66&8&3\\
62&24&72&8&3\\
65&28&45&10&3\\
66&32&60&28&3\\
67&32&44&12&2\\
68&32&36&12&2\\
69&32&82&28&4\\
70&32&84&28&4\\
71&32&76&12&3\\
72&32&72&12&3\\
74&48&54&20&2\\
75&48&120&20&3\\
76&48&120&20&3\\
78&64&136&28&3\\
79&96&210&44&3\\
\end{longtable}
The mixed-class order distribution is
\[
\begin{array}{c|rrrrrrrrrr}
 |K|&4&8&12&16&24&28&32&48&64&96\\\hline
 \#&6&15&6&14&7&1&7&3&1&1.
\end{array}                                             \tag{B.34.4}
\]
The scope of (B.34.3) is the condition that the subgroup itself
contain a hyperelliptic involution.  A finite subgroup that only
centralizes an involution outside itself is represented by a
complement in a lift group; its conjugacy after forgetting the
external involution requires an additional calculation.  No such
complement is counted among these $79$ classes.

\begin{claim}{B.36}{Proved}{Complete exceptional catalogues and their realization fibers}
For the four exceptional types, the numbers of finite-subgroup
conjugacy classes in the extended mapping class group, and the
numbers of finite action classes inducing these subgroups
isomorphically, are as follows:
\[
\begin{array}{c|c|r|r}
(g,\nu)&\mathcal T_{g,\nu}&
 \text{finite subgroup classes}&\text{realized action classes}\\\hline
(0,0)&C_2&2&3\\
(0,1)&C_2&2&2\\
(0,2)&C_2^2&5&7\\
(1,0)&\mathrm{GL}_2(\Z)&13&17
\end{array}                                             \tag{B.36.1}
\]
The executed files
\path{exceptional_finite_subgroups_and_lifts.csv} and
\path{exceptional_sections.json} give all $29$ action rows over
these $22$ mapping class subgroup rows.  They retain the extension,
all full and partial section classes, $\Phi$ and its stabilizers,
centralizers, and every defined local multiplier.
The request in B(f) for one set of section data attached only to a
finite mapping class subgroup is false in these types: the displayed
realization fibers are required.  For the antipodal action on $S^2$,
the algebraic section exists and the fixed-point locus is empty;
there is no fixed-point rotation multiplier for that section.
\end{claim}
\passage{\S{}15; folios 68--75; scans 137--151; \S{}17; folios 88--90; scans 179--181; \S{}19; folios 111--115; scans 222--230; Part B(f) of the task; no further manuscript assertion}

\begin{proof}
We first identify the groups in (B.36.1) and the comparison being
classified.  Orientation is the component invariant for
$\operatorname{Homeo}(S^2)$ and for
$\operatorname{Homeo}(S^2,\{\infty\})$.
An orientation-preserving homeomorphism is isotopic to the identity:
fix a disc and its complementary disc by the Alexander isotopy,
with the marked point kept fixed in the second case.
For the annulus, the two component invariants are the sign
$\epsilon$ on its fundamental group and the sign $\eta$ of its
permutation of the two ends.  Their four representatives are
\[
       (x,t)\longmapsto(\epsilon x,\eta t),
       \qquad x\in\R/\Z,\quad t\in\R.                  \tag{B.36.2}
\]
The kernel of these invariants is connected: compactify the two
ends, isotope an arc joining them to a fixed arc, and apply the
Alexander isotopy on the complementary disc.  Winding of the arc
is absorbed by rotating a puncture neighborhood, since its end
is a puncture and has no prescribed boundary parametrization.
For the torus, homology gives
$\pi_0\operatorname{Homeo}(T^2)=\mathrm{GL}_2(\Z)$.
Surjectivity is given by linear maps.  A homologically trivial
homeomorphism is first isotoped to fix the origin, then the two
basis curves, and then their complementary disc; translations
lie in the identity component.  B.25 gives its thirteen finite
subgroup classes.

In this statement a realized action over a finite subgroup
$H\subset\pi_0\operatorname{Homeo}(U)$ means a finite subgroup
$G\subset\operatorname{Homeo}(U)$ for which the projection
$G\to\pi_0\operatorname{Homeo}(U)$ is an isomorphism onto a conjugate
of $H$.  Actions are identified by homeomorphism conjugacy.
This specifies the exact fiber of the comparison map.  The
additional finite actions having a nontrivial kernel in that
projection are the full affine families already classified in
B.14; their kernels are retained in that result and are not
silently identified with an injective realization.

On $S^2$, a negative involution is conjugate either to a reflection
or to the antipodal map.  Indeed the sphere theorem in B.1 conjugates
the finite action into $O(3)$.  An orthogonal involution has
eigenvalues $\pm1$; negative determinant leaves precisely the
patterns $(1,1,-1)$ and $(-1,-1,-1)$.  Their fixed sets are a
circle and the empty set.  Together with the trivial action this
gives three rows over the two mapping class subgroups.
On the plane, a negative involution is conjugate to a reflection
by B.1, giving two rows including the trivial action.
For both surfaces $\pi=1$, so the lift extension is simply
$E=G$.  For the negative $C_2$ there is exactly one full section,
exactly one nontrivial partial section, and one maximal class,
with centralizer $1$ in $\pi$ and stabilizer $C_2$.
For the trivial group there is one full section and no nontrivial
partial section.  These are algebraic assertions about the
extension and include the antipodal action.  In that action
$p=-p$ implies $p=0$, which has no solution on the unit sphere.
The program verifies $A^2=I$, $\det A=-1$, and $\det(I-A)=8$
for $A=-I_3$.  It therefore records zero fixed components next
to the one algebraic full section.

A reflection has the local multiplier $1/2$.  For the antipodal
section no fixed-point germ exists from which to read such a
multiplier; the table marks that entry as undefined.  It separately
records the formal value $1/2$ obtained by extending the
orientation-sign convention to every negative involution.
This convention has a stated source and is not substituted for
a nonexistent fixed-point comparison.  In particular the
circle fixed by a sphere reflection has trivial image in
$\pi_1(S^2)$, so the hyperbolic oval-centralizer formula $\Z$
does not transfer to this case.

For the annulus use the simultaneous normal form of B.14.
Let
\[
 a(x,t)=(-x,t),\qquad
 b(x,t)=(x,-t),\qquad
 b'(x,t)=(x+\tfrac12,-t),\qquad
 c(x,t)=(-x,-t).                                      \tag{B.36.3}
\]
The three nontrivial mapping class elements are represented by
$a,b,c$.  A lift of the class of $a$ or $c$ has an arbitrary
angular shift, removed by translation conjugacy.
A lift of the class of $b$ has angular shift $0$ or $1/2$,
because its square is the identity.  These give $b$ and $b'$;
their fixed sets distinguish them.
For the whole four-group, normalize its angular reflection to
$a$.  Commutation forces the radial generator's angular shift
to be $0$ or $1/2$, giving exactly
$\langle a,b\rangle$ and $\langle a,b'\rangle$.
Conjugation cannot exchange the three nontrivial mapping class
elements, since their component group is abelian.  The fixed
oval or its absence again distinguishes these two actions.
This proves the seven-row classification.

Here are all its section counts; $f,p,\phi,o$ mean full sections,
nontrivial partial sections, $|\Phi|$, and $|\Phi/G|$:
\[
\begin{array}{l|rrrr}
\text{action}&f&p&\phi&o\\\hline
1&1&0&0&0\\
\langle a\rangle&2&2&2&2\\
\langle b\rangle&1&1&1&1\\
\langle b'\rangle&0&0&0&0\\
\langle c\rangle&2&2&2&2\\
\langle a,b\rangle&2&7&2&2\\
\langle a,b'\rangle&0&4&4&2
\end{array}                                             \tag{B.36.4}
\]
To verify them, the fixed sets of $a,b,b',c$ are respectively
two lines, one oval, the empty set, and two points.
In the split four-group, each of its two global fixed points
extends the corresponding point section and the incident
reflection sections; the oval section has both maximal extensions.
In the glide four-group, the two angular-reflection lines and the
two fixed points of $ab'$ are disjoint.  They are all maximal,
and the group interchanges each pair.  Thus (B.36.4) also gives
the required partial-section distinction.
Every nonidentity torsion element in these sections has multiplier
$1/2$.  The centralizer in $\pi=\Z$ is $\Z$ for the radial oval
section and $1$ for every point or line section.

These calculations can also be read directly from the extension.
For an affine element write its chosen angular shift as
$a_g\in[0,1)$ and let $\epsilon_g$ be its action on $\pi=\Z$.
The program records the integer cocycle
\[
 c(g,h)=a_g+\epsilon_ga_h-a_{gh}
 \quad\hbox{and}\quad
 (n,g)(m,h)=(n+\epsilon_gm+c(g,h),gh).                 \tag{B.36.5}
\]
It checks the cocycle identity on every triple.
In the universal cover $\R^2$, a finite subgroup of the lift
extension fixes a point; its fixed set is affine and connected.
Passing to the deck quotient identifies its conjugacy class
with the corresponding fixed component just listed.
This proves the section calculation without using a
center-free hypothesis.

For the torus apply the exact affine classification in B.27 and
retain precisely the normal forms whose translation kernel has
order one.  The linear projection of such a finite group is
faithful, which is exactly the comparison stipulated above.
Conversely every injective realized mapping class subgroup has
this property and is included by B.14.
There are seventeen affine forms over the thirteen integral
linear classes.  In the order of the linear classes in B.27,
the fiber sizes are
\[
                   1,2,1,1,1,3,1,1,1,1,1,2,1.       \tag{B.36.6}
\]
The proof of B.27 enumerates all cocycles and all affine
conjugacies without a word bound.  Its averaging comparison in
B.14 proves that no further topological conjugacy is missing.
Run the same exact fixed-point and reflection-circle computation
with no punctures removed.  Every isolated point section has
centralizer $1$ in $\Z^2$; every reflection oval has centralizer
the primitive invariant copy of $\Z$.  The trivial section has
centralizer $\Z^2$.  Local rotation matrices give the multipliers,
and the complete section poset gives the full and partial counts,
$\Phi$, and every stabilizer.  The resulting seventeen rows and
their exact rational matrices are included in the JSON.

The executed program \texttt{catalogue\_exceptional.py} carries
out all these finite calculations, records the two annulus-end
stabilizers and the plane's puncture stabilizer, and verifies the
four mapping class counts and the four realization-fiber totals.
Consequently it supplies a complete corrected computational
answer for the exceptional types, including the precise fibers
where B.22 prevents attaching one extension or one section table
to a bare mapping class subgroup.
\end{proof}

\begin{claim}{B.37}{Proved}{Complete genus-three mirror-disc quotient catalogue}
There are exactly $145$ conjugacy classes of finite subgroups $K$ of
$\mathcal T_{3,0}$ for which $H=K\cap\mathcal T^+_{3,0}$ has spherical
quotient and the induced involution on $S_3/H$ has a fixed circle.
Exactly $57$ contain a hyperelliptic involution, and $88$ do not.
The files \path{type_3_0_reflection_quotient_catalogue.csv},
\path{type_3_0_reflection_quotient_certificates.json}, and
\path{type_3_0_reflection_quotient_sections.json} give every
class, its exact finite group, and all its section data.
\end{claim}
\passage{\S{}18; folios 96--101; scans 193--203; Part B(f) of the task; genus-three mirror-disc quotient subcatalogue}

\begin{proof}
The quotient by the induced involution is a mirror disc $Q$.  Let
$(m_1,\ldots,m_k)$ be its cyclic list of corner orders and
$(n_1,\ldots,n_l)$ its list of interior cone orders.  The orientation
double of $Q$ is the spherical orbifold with cone multiset
\[
               \{m_1,\ldots,m_k,n_1,n_1,\ldots,n_l,n_l\}.
                                                        \tag{B.37.1}
\]
The involution on a sphere with a fixed circle is conjugate to a
reflection: cut along that circle into two discs and extend a boundary
parametrization over either disc, then use the involution to define
the parametrization on the other disc.  Cone points on its circle
become the corners of $Q$, while its other cone points occur in
exchanged pairs.  Thus the corner and interior data in (B.37.1)
exhaust its quotient types.  The corner list is considered up to
cyclic permutation and reversal; the interior list is unordered.

Use the following presentation of $\Lambda=\pi_1^{\rm orb}(Q)$:
\[
 \left\langle x_1,\ldots,x_l,e,c_0,\ldots,c_k\ \left|
 \begin{array}{c}
 x_i^{n_i}=1,\quad x_1\cdots x_l e=1,\\
 c_j^2=1,\quad(c_{j-1}c_j)^{m_j}=1,\quad c_k=ec_0e^{-1}
 \end{array}\right.\right\rangle .                    \tag{B.37.2}
\]
For $k=0$ there is one reflection $c_0$, commuting with $e$.
The orientation kernel $\Delta=\Lambda^+$ is the orbifold group
of (B.37.1).  Put $t=c_0$, $Z_j=c_{j-1}c_j$,
$P_j=Z_1\cdots Z_j$, $U_i=x_{l+1-i}^{-1}$, and $V_i=tU_it$.
Schreier rewriting, or the doubled fundamental polygon, gives the
spherical generating vector
\[
 (Z_1,\ldots,Z_k,U_1,\ldots,U_l,V_l^{-1},\ldots,V_1^{-1}),
 \qquad P_k\,U_1\cdots U_l\,(V_1\cdots V_l)^{-1}=1.
                                                        \tag{B.37.3}
\]
For example $e=U_1\cdots U_l$ and
$P_k=t c_k=tet e^{-1}=(V_1\cdots V_l)e^{-1}$, which proves
the product relation with its stated order.  Conjugation by $t$
acts on these generators by
\[
 \psi(Z_j)=P_{j-1}Z_j^{-1}P_{j-1}^{-1},\qquad
 \psi(U_i)=V_i,\qquad\psi(V_i)=U_i.                    \tag{B.37.4}
\]
These formulas include every base-path conjugation.

The input \cite{R-KNSData}, whose completeness is established in
\cite[Theorem 2.2, Corollary 2.3, and the planar calculation]{R-KNS},
contains every $\operatorname{Aut}(H)$-class of smooth epimorphisms
$\theta:\Delta\to H$ in genus three.  Its $197$ vectors cover all
$49$ genus-three signature/group pairs, including those with positive
quotient genus.  Here we use precisely the spherical-quotient pairs.
The original input and its attribution are retained in
\texttt{research\_data/actions-genus3.json}; B.33 verifies its orders
and surface signatures.  A prescribed ordering (B.37.3) is obtained
from the published sorted ordering by explicit adjacent Hurwitz
moves.  This is a bijection of epimorphism sets, since each move has
its displayed inverse in the program.

For each vector the program tests whether its kernel is invariant
under (B.37.4).  More explicitly, it attempts to define
\[
              \alpha(\theta(a))=\theta(\psi(a))
                         \qquad(a\in\Delta)             \tag{B.37.5}
\]
by propagation on the finite Cayley graph.  Agreement on every edge
is exactly the condition that (B.37.5) be well defined.  Surjectivity
then makes it an automorphism and (B.37.4) gives $\alpha^2=1$.
The resulting full action group is explicitly
\[
           K=H\rtimes_\alpha\langle t\rangle,
                 \quad t^2=1.                         \tag{B.37.6}
\]
Its epimorphism from (B.37.2) sends $x_i$ to their recovered vector
entries and $c_j$ to $t\theta(P_j)$.  Its kernel is exactly
$\ker\theta\subset\Delta$.  In particular no choice of a further
extension class is omitted: the orbifold kernel determines
$\Lambda/\ker\theta$ and hence (B.37.6).  Conversely every action
in the claim supplies this invariant kernel and therefore occurs.

We now prove the precise fusion calculation.  Two epimorphisms from
the same orbifold group have the same kernel exactly when a unique
isomorphism between their finite images intertwines them: define
that isomorphism by $\theta(a)\mapsto\theta'(a)$, which is well defined
exactly when the kernels agree.  Thus the published
$\operatorname{Aut}(H)$-normalization records kernels without loss.
A conjugacy of the surface actions descends to a decorated-orbifold
homeomorphism of $Q$ carrying one kernel to the other.  Conversely
such a quotient homeomorphism lifts through the two regular
orbifold coverings because it carries their defining kernels to
one another; its lift is the required surface conjugacy.  This is
an exact equivalence, so it remains to calculate the full action of
these disc homeomorphisms on the finite kernel set.

A disc homeomorphism permutes its colored boundary corners cyclically,
or reverses their cyclic order.  After composing with that boundary
map it fixes the boundary pointwise up to isotopy.  The remaining
mapping class is a braid of the interior cone points.  To see this,
forget the interior points and isotope the boundary-fixed disc map
to the identity by the Alexander isotopy; the tracks of the points
are its braid, and an isotopy of those tracks gives precisely an
isotopy of the punctured disc map by isotopy extension.  The colored
braid subgroup is generated by pure braid generators and adjacent
exchanges among equal colors after the colors are sorted.  The
program uses the standard words
$\sigma_{j-1}\cdots\sigma_{i+1}\sigma_i^2
\sigma_{i+1}^{-1}\cdots\sigma_{j-1}^{-1}$ and those equal-color
exchanges.  On interior generators an adjacent positive move is
\[
                    (x_i,x_{i+1})\longmapsto
                    (x_ix_{i+1}x_i^{-1},x_i),           \tag{B.37.7}
\]
fixing $e$ and every boundary reflection.  The inverse is also used
when forming the pure braid words.

Here are the boundary moves, including their effect on the chosen
reflection.  Extend the sequence by $c_{j+k}=ec_je^{-1}$.  One cyclic
shift has
\[
 \begin{split}
 (Z_1,\ldots,Z_k)&\longmapsto
       (Z_2,\ldots,Z_k,eZ_1e^{-1}),\\
 (x_1,\ldots,x_l)&\longmapsto(x_1,\ldots,x_l),\\
 t&\longmapsto tZ_1,\qquad
 \alpha\longmapsto\operatorname{Inn}(Z_1^{-1})\alpha.
 \end{split}                                           \tag{B.37.8}
\]
A reversal has
\[
 \begin{split}
 (Z_1,\ldots,Z_k)&\longmapsto(Z_k^{-1},\ldots,Z_1^{-1}),\\
 (x_1,\ldots,x_l)&\longmapsto(x_l^{-1},\ldots,x_1^{-1}),\\
 t&\longmapsto tP_k,\qquad
 \alpha\longmapsto\operatorname{Inn}(P_k^{-1})\alpha.
 \end{split}                                           \tag{B.37.9}
\]
They follow by sending $c_j$ respectively to $c_{j+1}$ and
$c_{k-j}$ in (B.37.2); in the reversal $e$ goes to $e^{-1}$.
Only the shifts and shifted reversals preserving the corner-order
list are self-homeomorphisms of the chosen colored quotient.
A fixed braid restores the sorted interior order after a reversal.
Equations (B.37.7)--(B.37.9), together with the preceding braid
argument, prove both that every programmed fusion is realized and
that every actual conjugacy is included.

The program computes the finite orbits under these moves.  It
identifies two generating vectors under $\operatorname{Aut}(H)$ by
their rooted, edge-colored Cayley graphs: starting at the identity,
visit vertices in generator order and record each colored successor.
Equality defines a unique group isomorphism mapping each generator
to the corresponding generator, and such an isomorphism plainly
gives equality.  Thus this is an exact isomorphism test rather than
a coarse group invariant.  The computation gives $145$ orbits.

For each orbit the second program constructs the regular action of
(B.37.6) explicitly from
$(h,\epsilon)(h',\epsilon')=(h\alpha^\epsilon(h'),\epsilon+\epsilon')$.
It checks (B.37.2), every local stabilizer injection, and
$|K|\chi^{\rm orb}(Q)=-4$.  The resulting cover is therefore a closed
orientable genus-three surface.  The finite coset incidence graph
then gives every fixed point and oval, multiplier, section,
extension, maximal section, centralizer, orbit and stabilizer as in
B.26.  These checks passed for all $145$ classes.

Finally the overlap with B.34 is identified by actual monodromy
isomorphisms.  In each of B.34's $57$ mirror-disc groups the program
extracts the positive group and the generators $x_i$, $Z_j$ and
conjugation by $c_0$.  It applies (B.37.8)--(B.37.9) and the interior
braids to put the signature in the same ordering, and locates its
rooted colored Cayley graph in the computed kernel orbit.  The file
\texttt{genus3\_hyperelliptic\_reflection\_crosscheck.json} records
the resulting bijection with $57$ different rows and equality of
all eight reported section-count columns.  Thus these are exactly
the common classes, leaving $88$ additional ones.
\end{proof}

The $145$ classes have the following order distribution:
\[
\begin{array}{c|rrrrrrrr}
 |K|&4&6&8&12&14&16&18&24\\\hline
 \#&5&3&36&13&3&35&1&15\\[2pt]
 |K|&28&32&42&48&64&96&192&336\\\hline
 \#&1&13&1&11&2&4&1&1.
\end{array}                                             \tag{B.37.10}
\]
Together B.34 and B.37 therefore give $149$ distinct mixed classes.
The additional four are the projective-plane quotients in B.34.
The scope of B.37 requires both a spherical positive quotient and a
fixed-circle involution on it; the remaining quotient types are
listed in the accompanying computation report.

\begin{claim}{B.38}{Proved}{The free positive involution and the six negative involutions in genus three}
There are exactly twelve conjugacy classes of mixed finite subgroups
of $\mathcal T_{3,0}$ whose positive subgroup is a free involution.
Each is a four-group.  Their quotient signatures and numbers are
\[
\begin{array}{c|ccccc}
\text{underlying quotient}&S_{0,3}&S_{1,1}&N_{1,2}&N_{2,1}&N_{3,0}\\\hline
\text{number of classes}&4&1&4&2&1 ,
\end{array}                                             \tag{B.38.1}
\]
where $S_{h,b}$ is orientable of genus $h$ with $b$ mirror boundary
components, and $N_{q,b}$ is nonorientable with $q$ crosscaps and
$b$ mirror boundary components.  There are no cones or corners.
There are also exactly six mixed classes whose positive subgroup
is trivial: the negative involutions with quotient signatures
\[
 S_{0,4},\quad S_{1,2},\quad
 N_{1,3},\quad N_{2,2},\quad N_{3,1},\quad N_{4,0}.
                                                        \tag{B.38.2}
\]
All eighteen classes, their exact monodromies, complete section
posets, $\Phi$, stabilizers, centralizers and multipliers appear in
\path{type_3_0_free_positive_and_anti_catalogue.csv}
and its \texttt{sections.json} companion.
\end{claim}
\passage{\S{}15; folios 73--75; scans 148--151; \S{}15; folios 77--79; scans 155--160; \S{}18; folios 96--101; scans 193--203; \S{}33; folios 222--225; scans 440--445; Part B(f) of the task; no further manuscript assertion}

\begin{proof}
Let $X$ be the closed orientable genus-three surface.  First let
$G^+=\langle a\rangle$ have order two and act freely, and let
$G$ contain negative elements.  Then $|G|=4$.
A group of order four is cyclic or a four-group: an element of
order four gives the first case; otherwise $(xy)^2=1$ for all
$x,y$, which gives commutativity and the second case.
There is no nontrivial positive point stabilizer.
A negative point stabilizer is generated by a reflection, so
$Q=X/G$ has no cone or corner points.  Its orbifold Euler
characteristic is $-4/4=-1$.  The complete list of compact
underlying surfaces with this Euler characteristic is exactly
(B.38.1).

The cyclic possibility cannot occur.  A negative element of
order four has no fixed point, since a negative local finite
stabilizer element is a reflection of order two.  Its square
is free by assumption, so the whole action would be free.
The quotient would then be $N_{3,0}$.
Its fundamental group has the presentation
\[
 \langle u_1,u_2,u_3\mid u_1^2u_2^2u_3^2=1\rangle .
\]
An orientation-compatible monodromy to $C_4=\Z/4$ sends each
crosscap generator to an odd residue.  Each square has image
$2$, so the relator has image $6=2\pmod4$, a contradiction.
The program checks all eight assignments of odd residues.
Thus $G$ is the four-group
$V_4=\mathbf F_2^2$, with orientation given by its second
coordinate, and $a=(1,0)$.

We classify all four-group monodromies, including their full
topological equivalence relation.  The mirror orbifold group
of an orientable $S_{h,b}$ has generators
$a_i,b_i,e_j,c_j$ and relations
\[
 c_j^2=1,\qquad[e_j,c_j]=1,\qquad
 \prod_i[a_i,b_i]\prod_je_j=1.                         \tag{B.38.3}
\]
The negative $c_j$ is the reflection on the $j$th boundary,
and the positive $e_j$ traverses that boundary.
For $N_{q,b}$ replace the handles by crosscap generators
$u_i$, and replace the last relation by
\[
                         \prod_i u_i^2\prod_j e_j=1.  \tag{B.38.4}
\]
The orientation character $\varepsilon$ is zero on handles
and $e_j$, and one on crosscaps and $c_j$.
An orientation-compatible epimorphism to $V_4$ is exactly
\[
 q_\alpha:\Pi_Q\longrightarrow V_4,\qquad
 q_\alpha(x)=(\alpha(x),\varepsilon(x)),
 \quad
 \alpha\in H^1(\Pi_Q;\mathbf F_2)
       \setminus\{0,\varepsilon\}.                   \tag{B.38.5}
\]
In coordinates write the handle or crosscap values as $\eta_i$,
the boundary values as $\beta_j=\alpha(e_j)$, and the reflection
values as $\gamma_j=\alpha(c_j)$.
The only linear relation is
$\sum_j\beta_j=0$.
The two automorphisms of $V_4$ preserving orientation identify
$\alpha$ with $\alpha+\varepsilon$.

Every map (B.38.5) is injective on local reflection groups.
Its kernel is therefore torsion-free; one proves this either
from the universal hyperbolic orbifold cover or by cutting along
arcs to its reflection polygon, where every finite stabilizer
is a displayed local reflection group.
The kernel lies in the orientation kernel, and its Euler
characteristic is $4(-1)=-4$.  It is consequently the fundamental
group of a closed orientable genus-three surface.
Conversely every action under consideration yields exactly
this construction.  Two such actions are conjugate precisely
when an orbifold homeomorphism of $Q$ carries one monodromy
kernel to the other, or equivalently when it carries the
monodromy to the other after an automorphism of $V_4$.
This follows by lifting paths in the two covers, with the
local reflection charts giving the extension across mirrors.
Thus the following finite cohomology-orbit calculation is an
exact classification of actions.

We prove the entire mapping class action needed in this
calculation.  A boundary permutation permutes the pairs
$(\beta_j,\gamma_j)$.  The interior homology action preserves
the mod-two intersection pairing, its radical generated by
boundary classes, and the orientation character.
It introduces no reflection coordinate: an interior loop
is mapped to an interior loop, whereas each reflection generator
is mapped to a conjugate of the reflection at the corresponding
boundary.  Thus the $\gamma_j$ are affected only by boundary
permutations and by the global replacement
$\alpha\mapsto\alpha+\varepsilon$.

For $S_{0,3}$ the interior homology is generated by its boundary
classes, with their single sum relation, so boundary permutations
are the whole action.  When $\beta=(0,0,0)$, surjectivity leaves
one nonconstant reflection-color pattern up to permutation and
global color reversal.  When $\beta$ has two ones, distinguish
the boundary with value zero and use global color reversal to
make its reflection value zero.  The two remaining reflection
values have weights zero, one or two.  This gives $1+3=4$ classes.

For $S_{1,1}$ the boundary value is zero.
Global color reversal sets its reflection value to zero.
The remaining nonzero pair of handle values is a single orbit
under $\mathrm{Sp}_2(\mathbf F_2)$, generated by the two Dehn
twists in a pair of basis curves.  This gives one class.

For $N_{1,2}$ write $u$ for its one-sided core and $e$ for
either mod-two boundary class.  Its intersection form has
$u\cdot u=1$, with $e$ in the radical.
An automorphism preserving the boundary classes and the
orientation character sends $u$ either to $u$ or to $u+e$.
Both possibilities are realized.  Push one boundary disc once
around the one-sided core and return it as an unparametrized
disc; its boundary orientation may reverse.  A core arc crossing
the track once acquires one boundary circuit.  Thus, modulo two,
this homeomorphism sends $u$ to $u+e$ and leaves both reflection
coordinates unchanged.  The boundary exchange is realized away
from the core.  These maps generate every permitted action on
the displayed homology and boundary data.
After global color reversal sets $\eta=\alpha(u)$ to zero,
the case $\beta=(0,0)$ has two nonzero reflection patterns up
to exchange.  The case $\beta=(1,1)$ has two patterns: the
boundary push followed by global reversal interchanges
$\gamma=(0,0)$ and $(1,1)$, while $(0,1)$ and $(1,0)$ are
exchanged by the boundary permutation.  This gives four classes.
The boundary push is necessary for this count.

For $N_{2,1}$ the boundary value is zero.
Set its reflection value to zero by global reversal.
The intersection form on the two crosscap classes is the
identity matrix over $\mathbf F_2$; its orthogonal group has
exactly the identity and the exchange of those two classes.
They are realized by exchanging the two crosscaps.
The nonzero crosscap vectors therefore have two orbits,
represented by $(1,0)$ and $(1,1)$.

For $N_{3,0}$ the mod-two intersection form is the identity
matrix in three crosscap classes, and its characteristic
vector is $(1,1,1)$.  An orthonormal basis cannot contain that
characteristic vector, since it pairs to one with every vector
of odd self-intersection.  Its members must therefore be the
three standard basis vectors.  The full orthogonal group is
their permutation group, realized by crosscap permutations.
Global addition of $\varepsilon=(1,1,1)$ exchanges weights
one and two, leaving one orbit after $0,\varepsilon$ are
excluded.  This proves all five numbers in (B.38.1).

We now compute the sections explicitly.  A mirror boundary
with data $(\beta_j,\gamma_j)$ has local reflection
$c_j=(\gamma_j,1)$ and boundary monodromy
$e_j=(\beta_j,0)$.
Its preimage is described by the two cosets
$V_4/\langle c_j\rangle$; one trip around the boundary
permutes these cosets by $e_j$.
Consequently it gives two fixed ovals when $\beta_j=0$,
and one when $\beta_j=1$.
The two-oval case has an oval stabilizer of order two,
$\langle c_j\rangle$, and the one-oval case has stabilizer
$V_4=\langle c_j,e_j\rangle$.
Both are one $V_4$-orbit of section classes.
There are no point sections because the positive involution
is free, and there are no full sections because $V_4$ cannot
fix a point without giving a positive local stabilizer.
Every reflection section is maximal, has centralizer $\Z$
in the surface group, and has multiplier $1/2$ on its
nonidentity element.  Thus
\[
 |\Phi|=\#\{\text{nontrivial partial sections}\}
       =2b-\#\{j:\beta_j=1\},
 \qquad |\Phi/V_4|=b.                                 \tag{B.38.6}
\]
The program constructs the actual finite coset graph, with
three subdivisions per boundary traversal, computes its
components and the group action, and verifies every count
and stabilizer in this formula.

Finally a mixed group with trivial positive subgroup is $C_2$.
Its quotient has orbifold Euler characteristic $-2$ and no
cones or corners.  An orientable quotient must have mirror
boundary; otherwise its orientation character would be
trivial and could not give the required negative covering.
Solving $2-2h-b=-2$ with $b>0$, or $2-q-b=-2$ with $q>0$,
gives precisely (B.38.2).
For each quotient the monodromy is forced: it is the orientation
character, taking every crosscap and local reflection to the
nonidentity element and every handle and boundary traversal
to zero.  Its orientation double is connected, torsion-free,
and has Euler characteristic $-4$.  The classification of
compact quotient surfaces and path lifting prove that each
signature gives one conjugacy class and that distinct signatures
give distinct classes.
There are $b$ fixed ovals, hence $b$ full sections, $b$
nontrivial partial sections, and $b$ maximal classes, each
fixed by the whole $C_2$ and with surface-group centralizer
$\Z$.  The case $N_{4,0}$ has none.  These are the six
genus-three real types of B.16 in an executed form.

The code \texttt{catalogue\_genus3\_free\_positive.py} enumerates
every cohomology vector in (B.38.5), applies exactly the
proved generators of the equivalence relation, checks
surjectivity, all quotient relations, orientation compatibility,
and the covering Euler characteristic, and records every
orbit.  Its second calculation constructs all six orientation
doubles in (B.38.2).  This supplies the complete eighteen-row
catalogue and its section data.
\end{proof}

\begin{claim}{B.39}{Proved}{Complete genus-three projective-plane quotient catalogue}
There are exactly $11$ conjugacy classes of mixed finite genus-three
actions whose positive quotient is a sphere and whose induced
involution on that sphere is fixed-point-free.  Their full quotients
are projective planes with cone points.  Four classes are identified
with B.34 and seven are additional.  Consequently B.37 and the present
claim give exactly $156$ mixed classes with spherical positive quotient.
\end{claim}
\passage{\S{}18; folios 96--101; scans 193--203; Part B(f) of the task; genus-three projective-plane quotient subcatalogue}

\begin{proof}
A fixed-point-free sphere involution has projective-plane quotient:
its quotient is a closed surface of Euler characteristic one, and
its orientation cover is the given sphere.  Cone points are exchanged
in pairs.  Thus each multiplicity in the positive cone multiset must
be even.  Inspection of the complete input in B.33 leaves exactly
$\mathrm O3.4$, $\mathrm O3.9$, $\mathrm O3.11$, $\mathrm O3.15$,
$\mathrm O3.21$, $\mathrm O3.22$, and $\mathrm O3.25$.
Their paired cone orders are respectively
$(2,2,2,2)$, $(2,2,2)$, $(4,4)$, $(2,6)$, $(2,4)$, $(2,4)$,
and $(2,3)$.

For a projective plane with cone orders $n_1,\ldots,n_l$ use
\[
 \Lambda=\langle x_1,\ldots,x_l,d\mid
       x_i^{n_i}=1,\quad x_1\cdots x_l d^2=1\rangle .   \tag{B.39.1}
\]
Set $X_i=x_i$, $Y_i=d x_i d^{-1}$ and
$A=d^2=(X_1\cdots X_l)^{-1}$.  Its orientation double has the
spherical vector
\[
        v=(X_1,\ldots,X_l,Y_l^{-1},\ldots,Y_1^{-1}),
                \qquad\prod_iX_i=\prod_iY_i.            \tag{B.39.2}
\]
Conjugation by $d$ is
\[
       X_i\longmapsto Y_i,\qquad
       Y_i\longmapsto A X_iA^{-1}.                     \tag{B.39.3}
\]
For every published $\operatorname{Aut}(H)$-normalized vector in
these seven cases, the program tests that (B.39.3) descends to an
automorphism $\alpha$ of its finite positive image $H$.  The same
Cayley-graph propagation as in B.37 gives an exact kernel-invariance
test.  It also verifies
$\alpha^2=\operatorname{Inn}(A)$ and $\alpha(A)=A$.
The full finite image is then the set $H\times\{0,1\}$ with product
\[
 (h,\epsilon)(k,\delta)=
       \bigl(h\alpha^\epsilon(k) A^{\epsilon\delta},
                    \epsilon+\delta\bmod2\bigr).      \tag{B.39.4}
\]
The two identities just verified prove associativity; the element
$d=(1,1)$ has square $A$ and induces $\alpha$ on $H$.
Thus (B.39.4) is exactly $\Lambda/\ker\theta$, where
$\theta$ is the positive epimorphism (B.39.2).

For clarity, the full source fusion is as follows.  The colored
mapping class group of the punctured projective plane is generated
by the colored disc braids and by sliding each puncture once around
a one-sided loop.  One direct proof cuts the projective plane along
a projective line disjoint from the marked disc: a moving puncture
crossing the cut is separated off as a one-sided slide, after which
all the remaining tracks form a braid in the disc.  Equivalently,
forgetting punctures successively expresses their point-pushing
loops by peripheral loops and the one-sided generator.  The former
are the pure disc braids.  The unmarked projective-plane mapping
class group is trivial: isotope the image of a projective line to
that line, adjust its parametrization by a projective rotation, and
apply the disc Alexander isotopy after cutting.  This proves that
these point motions exhaust all mapping classes.  Color-preserving
permutations are supplied by the same equal-color braid generators
used in B.37.

The slide of the last puncture has the exact substitutions
\[
 x_i\longmapsto x_i\quad(i<l),\qquad
 x_l\longmapsto(x_l d)x_l^{-1}(x_l d)^{-1},\qquad
 d\longmapsto x_l d.                                   \tag{B.39.5}
\]
They are obtained by tracking the meridian through a neighborhood
of that puncture and the one-sided loop.  The two-cone case, including
its geometric puncture-slide interpretation, is verified independently
in \cite[Lemma 3.19]{C_GromadzkiSzepietowski2016}.  Other punctures
outside this neighborhood are fixed, so the same formula applies
to all $l$ used here.  It respects the product relation directly:
if $B=x_1\cdots x_{l-1}$ and $x=x_l$, then
\[
 B(x d x^{-1}d^{-1}x^{-1})(x d)^2=B x d^2=1.
\]
In the finite positive group its effect is
\[
 x_l\longmapsto x_l\alpha(x_l)^{-1}x_l^{-1},\qquad
 \alpha\longmapsto\operatorname{Inn}(x_l)\alpha,
 \qquad A\longmapsto x_l\alpha(x_l)A.                  \tag{B.39.6}
\]
To slide any other puncture, a specified sequence of adjacent braids
moves it to the last position, applies (B.39.6), and moves it back.
The full word preserves its color even when the intermediate braids
exchange different colors.  Equations (B.39.5)--(B.39.6) and the disc
braids therefore give exactly the source mapping-class action.
The kernel-to-cover conjugacy argument of B.37 applies without change.

Executing these finite orbits gives respectively $1,3,3,1,2,0,1$
classes in the seven signature/group pairs above.  For every class,
\texttt{catalogue\_genus3\_projective\_quotients.py} constructs the
regular representation of (B.39.4), checks every relation in
(B.39.1), injectivity of all local cyclic groups, and Euler
characteristic $-4$.  The finite point-coset calculation gives every
partial section and its multiplier, centralizer and extensions.
No negative element has a fixed point, since its image on the
positive quotient sphere is fixed-point-free.  Thus all nontrivial
sections in these rows have positive cyclic image and are point
sections; their surface-group centralizers are trivial.

Finally the four projective-plane groups in B.34 are compared using
their actual cone generators and conjugation by their crosscap
element.  After a fixed ordering braid, their doubled vectors are
identified in the finite kernel orbits by the rooted colored Cayley
graph isomorphism test.  The resulting exact identifications are
\[
 \mathrm{G3H}_{002}\leftrightarrow\mathrm{G3P}_{001},\quad
 \mathrm{G3H}_{011}\leftrightarrow\mathrm{G3P}_{002},\quad
 \mathrm{G3H}_{029}\leftrightarrow\mathrm{G3P}_{008},\quad
 \mathrm{G3H}_{039}\leftrightarrow\mathrm{G3P}_{010}.
                                                        \tag{B.39.7}
\]
All seven applicable section-count columns agree under this
bijection.  Hence precisely seven classes are additional to B.34.
The quotient involution on the positive sphere has a fixed circle
in B.37 and has empty fixed set here; a conjugacy induces a sphere
conjugacy carrying these fixed sets to each other.  Thus none of
these eleven rows can be conjugate to a B.37 row, proving the total
$145+11=156$.
\end{proof}

Here $p$ counts nontrivial partial section classes, and $m,o$ are
$|\Phi|$ and $|K\backslash\Phi|$.  The final column $i$ counts negative
involutions in the displayed finite group; all of them act freely.
\[
\begin{array}{c|r|l|r|r|r|r}
 \text{class}&|K|&\text{cone orders}&p&m&o&i\\\hline
 001&4&(2,2,2,2)&8&8&4&2\\
 002&8&(2,2,2)&12&12&3&0\\
 003&8&(2,2,2)&12&12&3&2\\
 004&8&(2,2,2)&12&12&3&4\\
 005&8&(4,4)&8&4&2&0\\
 006&8&(4,4)&8&4&2&2\\
 007&8&(4,4)&8&4&2&4\\
 008&12&(2,6)&12&8&2&2\\
 009&16&(2,4)&16&12&2&0\\
 010&16&(2,4)&16&12&2&0\\
 011&24&(2,3)&20&20&2&4
\end{array}                                             \tag{B.39.8}
\]
Every row has zero full splittings, since its full image contains a
negative element and no such element fixes a point.  The JSON
\texttt{type\_3\_0\_projective\_quotient\_sections.json} includes all
finite permutations, local generators, multipliers and section
incidences behind this table.

\begin{claim}{B.40}{Proved}{Complete annulus and M\"obius-band quotients in genus three}
There are exactly $57$ mixed finite genus-three action classes whose
full quotient is an annulus or a M\"obius band with its cone and mirror
corner structure.  There are $39$ annulus classes and $18$ M\"obius-band
classes.  Their positive quotient is a torus.  Every group, generating
vector, quotient signature and complete section poset is given in
\path{type_3_0_torus_mirrors_catalogue.csv},
\path{type_3_0_torus_mirrors_certificates.json}, and
\path{type_3_0_torus_mirrors_sections.json}.
\end{claim}
\passage{\S{}18; folios 96--101; scans 193--203; Part B(f) of the task; genus-three annulus and Mobius quotient cases}

\begin{proof}
An involution reversing the orientation of a torus has annulus,
M\"obius-band or Klein-bottle quotient.  Indeed its fixed set consists
of disjoint circles, and its quotient has Euler characteristic zero;
the orientation double is the torus.  The classification of compact
surfaces leaves the orientable annulus, the nonorientable M\"obius
band, and the closed nonorientable Klein bottle.  Their orientation
doubles give respectively two, one, and zero fixed circles.  Cone
points on a fixed circle become mirror corners, and other cone points
are paired to interior cones.  The Klein-bottle cases are B.41.

The complete positive input of B.33 leaves the following seven
signature/group pairs.  In each row $r$ is the number of positive
quotient cone points, all of order $m$; the last two columns are the
executed numbers of full annulus and M\"obius-band classes.
\[
\begin{array}{c|c|r|r|r|r}
 \text{positive row}&H&m&r&\text{annulus}&\text{M\"obius}\\\hline
 \mathrm O3.3&C_2&2&4&10&3\\
 \mathrm O3.5&C_3&3&2&5&4\\
 \mathrm O3.7&C_4&2&2&8&4\\
 \mathrm O3.8&C_2^2&2&2&9&4\\
 \mathrm O3.12&S_3&3&1&3&1\\
 \mathrm O3.17&D_8&2&1&3&1\\
 \mathrm O3.18&Q_8&2&1&1&1
\end{array}                                             \tag{B.40.1}
\]
Here $D_8$ has order eight.  For $l$ interior cones, the total number
of corners is $r-2l$.  On an annulus distribute them as
$k_1\le k_2$ with $k_1+k_2=r-2l$; on a M\"obius band place all
$k=r-2l$ on its single boundary.  All their orders are $m$, so these
integer choices already classify the colored quotient signatures.

We first give an exhaustive construction of the possible oriented
finite groups, rather than assuming a small-group name list.
Choose a negative coset representative $t$ in an index-two extension
of $H$.  Let $\alpha(h)=t h t^{-1}$ and $A=t^2$.  Its complete
extension data satisfy
\[
              \alpha^2=\operatorname{Inn}(A),\qquad
                          \alpha(A)=A.                 \tag{B.40.2}
\]
Conversely the product of (B.39.4), with these data, defines that
extension:
\[
 (h,\epsilon)(k,\delta)=
       (h\alpha^\epsilon(k)A^{\epsilon\delta},
                                  \epsilon+\delta\bmod2).
                                                        \tag{B.40.3}
\]
Every oriented group isomorphism restricts to some $\beta\in
\operatorname{Aut}(H)$ and sends $t$ to some $k t'$ in the negative
coset.  Therefore the entire isomorphism action on (B.40.2) is generated
by transport by $\beta$ and by
\[
       (\alpha,A)\longmapsto
       \bigl(\operatorname{Inn}(k)\alpha,\ k\alpha(k)A\bigr).
                                                        \tag{B.40.4}
\]
The converse follows by sending the displayed generators accordingly
in (B.40.3), so this is an exact fusion test.  The program enumerates
$\operatorname{Aut}(H)$ by assigning images to a generating set and
propagating on its complete Cayley graph, enumerates every pair
(B.40.2), and takes precisely the finite orbits (B.40.4).
For the seven $H$'s in (B.40.1), it obtains respectively
$2,2,4,3,1,4,4$ oriented extension classes.

Here are the full source presentations and enumeration conditions.
For an annulus, put $X=x_1\cdots x_l$ and use positive generators
$x_i,e_1,e_2$ with
\[
                x_i^m=1,\qquad X e_1e_2=1.             \tag{B.40.5}
\]
On its $j$th boundary there are reflections
$c_{j,0},\ldots,c_{j,k_j}$ satisfying
\[
 c_{j,s}^2=1,\quad(c_{j,s-1}c_{j,s})^m=1,\qquad
               c_{j,k_j}=e_j c_{j,0}e_j^{-1}.          \tag{B.40.6}
\]
For $k_j=0$, there is one reflection commuting with $e_j$.
For a M\"obius band there is one negative crosscap generator $d$,
one boundary with (B.40.6), and the relation
\[
                X e d^2=1,
                   \qquad e=X^{-1}d^{-2}.              \tag{B.40.7}
\]
In each finite extension the program assigns every possible positive
order-$m$ element to every $x_i$, every positive element to $e_1$
in the annulus, and every negative element to $d$ in the M\"obius
band.  It then assigns every sequence of negative involutions to
the mirrors, requiring exact order $m$ for each adjacent product
and the stated closing relation.  It retains precisely the vectors
generating the whole finite group.  Two involutions whose product
has exact order $m$ generate a dihedral group of order $2m$; thus
all local groups inject.  The orientation character takes the
specified signs on every generator, so the kernel is oriented and
contains no reflection.  The signature equation gives
$|K|\chi^{\rm orb}(Q)=-4$; its torsion-free surface kernel therefore
has genus three.  Conversely every action with this quotient yields
one of these enumerated epimorphisms.

Target automorphisms are removed by the exact rooted colored Cayley
graph test of B.37, now applied to the full NEC generating vector.
Its prescribed generator signs force any resulting isomorphism to
preserve the orientation character.  For source fusion, the following
are explicit complete generators.  The interior braids are
(B.37.7), with all other generators fixed.  A cyclic boundary move
sends $c_{j,s}$ to $c_{j,s+1}$, extending the reflection sequence
by (B.40.6), while fixing $x_i$, $e_j$, and the crosscap generator.
After $k_j$ moves it conjugates that boundary's reflection sequence
by $e_j$ and fixes the other generators, which supplies the boundary
Dehn twist where it is needed.

For the annulus, pushing the last interior point once around the
first boundary, with $a=x_l$ and $e=e_1$, has
\[
 \begin{aligned}
 a&\longmapsto (a e)a(a e)^{-1},\\
 e_1&\longmapsto a e_1a^{-1},\\
 c_{1,s}&\longmapsto a c_{1,s}a^{-1}.
 \end{aligned}                                           \tag{B.40.8}
\]
It fixes the other $x_i$, $e_2$ and the second reflection sequence.
This is the square of the Hurwitz exchange between the puncture
meridian and the first boundary loop; in particular it preserves
$a e_1$ exactly.  The interior braids conjugate this to pushes of
each other interior point.  When $k_1=k_2$, interchange of the two
boundary components is
\[
 \begin{aligned}
 e_1&\longmapsto e_1e_2e_1^{-1},&
 e_2&\longmapsto e_1,\\
 c_{1,s}&\longmapsto e_1c_{2,s}e_1^{-1},&
 c_{2,s}&\longmapsto c_{1,s},
 \end{aligned}                                           \tag{B.40.9}
\]
fixing the interior generators.  This is their orientation-preserving
boundary exchange, with the chosen connecting paths.

The remaining annulus generator reverses the orientation of its
fundamental region.  For the peripheral list
$(w_1,\ldots,w_{l+2})=(x_1,\ldots,x_l,e_1,e_2)$, write
$W_i=w_1\cdots w_i$.  It sends
\[
             w_i\longmapsto W_{i-1}w_i^{-1}W_{i-1}^{-1}.
                                                        \tag{B.40.10}
\]
For the first boundary its reflection sequence is reversed and
conjugated by $X$; for the second it is reversed and conjugated by
$X e_1$.  Reversing a cycle means
$(c_0,\ldots,c_{k-1})\mapsto(c_k,c_{k-1},\ldots,c_1)$,
with $c_k=e c_0e^{-1}$; an empty cycle retains its single reflection.
These substitutions verify (B.40.5)--(B.40.6) directly.

For the M\"obius band, a one-sided slide of the last interior point
is particularly explicit.  Set
\[
          a=x_l,\qquad a'=(a d)a^{-1}(a d)^{-1}.
\]
The substitution is
\[
 \begin{aligned}
  x_l&\longmapsto a',& d&\longmapsto a d,\\
  c_s&\longmapsto a'^{-1}a\,c_s\,a^{-1}a'.
 \end{aligned}                                           \tag{B.40.11}
\]
It fixes the other interior generators and determines $e$ by
(B.40.7).  To obtain it geometrically, regard the boundary temporarily
as the last peripheral loop in the projective-plane presentation:
exchange $(a,e)$ by its Hurwitz move, slide the now-last puncture
by (B.39.5), and undo that exchange.  This gives exactly
(B.40.11), including the displayed transport of every boundary
reflection.  Finally a reversal of the core has
\[
 \begin{aligned}
 (x_1,\ldots,x_l)&\longmapsto(x_l^{-1},\ldots,x_1^{-1}),\\
 d&\longmapsto d^{-1},\qquad
 c_s\longmapsto X c_{k-s}X^{-1}.
 \end{aligned}                                           \tag{B.40.12}
\]
Here $e$ goes to $X e^{-1}X^{-1}$, so the closing relation is
preserved.

We justify that this is the whole source action.  First prescribe
the cyclic permutations, and possible reversal or interchange, of
the marked boundary components.  The boundary moves above realize
exactly these permutations.  Upon forgetting the interior points,
a remaining annulus map is generated by its core twist, and a
remaining M\"obius-band map by reversal of its core.  Cutting along
a spanning arc gives this classification: the annulus permits the
integer winding of the arc, while cutting the M\"obius band reduces
a boundary-fixed map to a disc isotopy.  The boundary reversal is
(B.40.12).  The core twist is already among the full cyclic boundary
moves when both boundaries are marked; when a boundary is unmarked
its rotation isotopes it away in the unpunctured quotient.
The kernel of forgetting interior points consists of their braids.
An interior point in an annulus minus the other points has fundamental
group generated by the puncture meridians and the core loop; in a
M\"obius band the corresponding extra generator is one-sided.
Forgetting points one at a time, and using isotopy extension, therefore
shows that their braid groups are generated by the ordinary interior
braids and respectively (B.40.8) or (B.40.11).  This also includes
any twist that becomes trivial only after the interior points are
forgotten.  Thus every geometric source automorphism is included,
and each included automorphism is realized by the maps just described.
The kernel-to-cover lifting argument of B.37 proves that the computed
orbits are exactly the surface conjugacy classes.

Execution gives the $57$ classes in (B.40.1).  There is an independent
check on the source generators in every one-cone, no-corner case:
the complete primary automorphism lists in
\cite[Lemmas 3.7 and 3.8]{C_GromadzkiSzepietowski2016} give exactly the
same partitions of the entire target-normalized epimorphism sets.
The file \path{check_torus_mirror_generators_output.json}
records all six comparisons, comparing full orbit memberships and
not just their numbers.

The program then computes all section invariants.  Its mirror
incidence graph has a separate closed cycle for each quotient
boundary, with endpoint gluing by that boundary's actual element
$e_j$.  At every corner it uses the full local dihedral group, and
checks that the fixed subgraph of each reflection has degree two
at every vertex.  Consequently its connected components are exactly
the fixed ovals, including all their intersections with other
reflection fixed sets.  Cone and corner cosets give the isolated
fixed points.  The procedure of B.26 then supplies every section,
multiplier, centralizer, extension incidence, maximal section, group
orbit and stabilizer.  All $57$ presentations, local injections,
Euler characteristics and complete section calculations passed.
\end{proof}

The $57$ classes have orders $4,6,8,12,16$ with multiplicities
$13,9,25,4,6$, respectively.  The full quotient surfaces in this
claim have boundary, whereas those in B.41 are closed.  A conjugacy
of actions descends to a homeomorphism of their full quotient
orbifolds and therefore preserves boundary.  Their positive
quotients have genus one, whereas those in B.37 and B.39 have genus
zero and those in B.38 have genus two or three.  The induced
homeomorphism of positive quotients preserves genus.  These exact
quotient maps prove that the four portions contribute disjoint
conjugacy classes.

\begin{claim}{B.41}{Proved}{The six genus-three actions with Klein-bottle quotient}
There are exactly six mixed finite-subgroup classes in
$\mathcal T_{3,0}$ whose quotient orbifold has underlying
Klein bottle.  Their positive subgroups, full groups, cone orders,
and section counts are
\[
\begin{array}{c|c|c|rrrr}
G^+&G&\text{cone orders}&f&p&|\Phi|&|\Phi/G|\\\hline
C_2&C_4&2,2&0&4&4&2\\
C_2&C_2^2&2,2&0&4&4&2\\
C_3&C_6&3&0&2&2&1\\
C_4&C_8&2&0&4&4&1\\
C_4&C_4\times C_2&2&0&4&4&1\\
C_2^2&D_8&2&0&4&4&1
\end{array}                                             \tag{B.41.1}
\]
Here $f$ and $p$ denote full and nontrivial partial section classes,
and $D_8$ has order eight.  The files
\path{type_3_0_klein_quotient_catalogue.csv} and
\path{type_3_0_klein_quotient_sections.json} contain
every group operation, monodromy, local multiplier, section
stabilizer and centralizer.  All centralizers in the surface
group are trivial.  In the primary positive catalogue, these
rows lie over O3.3, O3.5, O3.7 and O3.8, respectively in
numbers $2,1,2,1$.
\end{claim}
\passage{\S{}15; folios 73--75; scans 148--151; \S{}15; folios 77--79; scans 155--160; \S{}18; folios 96--101; scans 193--203; Part B(f) of the task; no further manuscript assertion}

\begin{proof}
The orientation double of a Klein bottle is a torus.
An interior cone of order $m$ in the quotient has two lifts of
order $m$ in that double, and there are no mirrors or corners.
The complete positive catalogue in B.33 therefore leaves
precisely the positive actions
\[
 (C_2;2,2,2,2),\quad(C_3;3,3),\quad
 (C_4;2,2),\quad(C_2^2;2,2),                          \tag{B.41.2}
\]
with torus underlying quotient.
The other torus-positive rows have a single cone and cannot
arise by this double-cover comparison.
Thus the full groups have orders four, six or eight.

We enumerate all oriented group extensions, rather than
presupposing a list of abstract groups.
For a fixed positive group $H$ choose a negative representative
$t$, put $\alpha(h)=tht^{-1}$ and $A=t^2\in H$.
Associativity is equivalent to
\[
           \alpha^2=\operatorname{Inn}(A),
           \qquad\alpha(A)=A.                         \tag{B.41.3}
\]
Conversely the set $H\times\{0,1\}$ with normal form $ht^i$
and relations $th=\alpha(h)t$, $t^2=A$ is a group for every
pair in (B.41.3); checking the two possible negative factors
verifies associativity.  Changing $t$ to $kt$ changes the pair to
\[
 \bigl(\operatorname{Inn}(k)\alpha,\ k\alpha(k)A\bigr),
                                                        \tag{B.41.4}
\]
and changing the identification of $H$ by $\beta\in\operatorname{Aut}(H)$
changes it to
$(\beta\alpha\beta^{-1},\beta(A))$.
These are all orientation-preserving isomorphisms of the
extensions, because any such isomorphism restricts to an
automorphism of $H$ and sends a negative representative to
one negative coset representative.

For $H=C_2$ this gives $C_4$ and $C_2^2$.
For $H=C_3$, the two possible conjugation actions give $C_6$
and $S_3$.  Explicitly the identity action allows the square
of $t$ to be removed using (B.41.4), while inversion forces
$A=1$.
For $H=C_4=\langle a\rangle$, the identity action has two
square classes, $A=1$ or $a$, giving $C_4\times C_2$ and
$C_8$.  Inversion allows $A=1$ or $a^2$, giving $D_8$ and
$Q_8$.
For $H=C_2^2$, the identity action gives $C_2^3$ or
$C_4\times C_2$ according as $A=1$ or $A\ne1$.
Every nonidentity order-two automorphism of $H$ is conjugate
to a transvection.  Its fixed subgroup equals the image of
$1+\alpha$, so (B.41.4) removes either allowable square.
This gives one further oriented extension, $D_8$, with positive
subgroup a four-group.  These arguments prove all the extension
counts $2,2,4,3$ used in the computation.

First take a Klein bottle with one cone of order $m$.
Its orbifold group is
\[
 \Lambda_m=
 \langle d_1,d_2\mid(d_1^2d_2^2)^m=1\rangle.
                                                        \tag{B.41.5}
\]
Put $y=d_1$, $z=d_1d_2$.  Then $y$ is negative, $z$ is
positive, and a positive local cone generator is
\[
                     x=(yzy^{-1}z)^{-1}.              \tag{B.41.6}
\]
A smooth orientation-compatible epimorphism to an oriented
finite group $G$ is therefore exactly a pair $(Y,Z)$ such that
$Y$ is negative, $Z$ is positive, $\langle Y,Z\rangle=G$,
and the image of (B.41.6) has order $m$.
Gromadzki--Szepietowski \cite[Lemma 3.3]{C_GromadzkiSzepietowski2016}
prove that the full outer automorphism group of $\Lambda_m$
is generated by
\[
 (y,z)\longmapsto(yz,z),\qquad
 (y,z)\longmapsto(y^{-1},z),\qquad
 (y,z)\longmapsto(y,z^{-1}).                           \tag{B.41.7}
\]
Their proof identifies these maps with the Dehn twist and
puncture and crosscap slides of the one-pointed Klein bottle.
The automorphism assertion concerns $\Lambda_m$ itself and
does not impose their later nonorientable-kernel hypothesis.
We apply that assertion to the orientation-compatible
epimorphisms above; every generator in (B.41.7) preserves the
orientation character.  Quotienting the finite set of pairs
also by $\operatorname{Aut}(G)$ preserving orientation is
therefore exactly topological conjugacy of the surface actions.

For $H=C_3$ the cone order is three.
In the $S_3$ extension a negative element inverts every
positive element, so (B.41.6) is the identity.
In $C_6$ it equals $Z^{-2}$; $Z$ must generate the positive
$C_3$.  The first move in (B.41.7) runs through the entire
negative coset, and the last move exchanges the two choices
of $Z$.  Thus there is one class, with full group $C_6$.

For $H=C_4$ the cone order is two.
The negative elements of $D_8$ and $Q_8$ invert $H$, again
making (B.41.6) trivial.  In the two abelian extensions it is
$Z^{-2}$, so $Z$ must generate $C_4$.
For $C_4\times C_2$, multiplication of $Y$ by powers of $Z$
sets $Y$ to the order-two generator of the other factor.
For $C_8$, an automorphism sets the odd generator $Y$ to a
chosen generator, and inversion of $Z$ identifies its two
remaining possibilities.  Each gives exactly one class.

For $H=C_2^2$, the elementary abelian extension $C_2^3$
cannot be generated by two elements.
In the abelian extension $C_4\times C_2$, every $Z\in H$
has square one, so (B.41.6) is trivial.
Write the remaining extension as
\[
 D_8=\langle r,s\mid r^4=s^2=1,\ srs=r^{-1}\rangle,
 \qquad H=\langle r^2,s\rangle .
\]
The admissible $Z$ are $s,r^2s$, and every negative $Y$
gives the cone image $r^2$.
An orientation-preserving automorphism exchanges the two
choices of $Z$; the first move in (B.41.7) changes a negative
reflection to a negative rotation, and inversion then
identifies $r$ and $r^{-1}$.  All eight pairs are one orbit.
This gives the final row in (B.41.1).

It remains to handle the two order-two cones in the first case
of (B.41.2).  For either extension $G=C_4$ or $C_2^2$, let
$h$ be the nonidentity positive element.
Both cone generators must map to $h$.
The two crosscap generators may independently map to either
negative element.  Every one of these four assignments is
surjective and respects the relation
\[
                      d_1^2x_1d_2^2x_2=1.
                                                        \tag{B.41.8}
\]
They are all equivalent.  Here is a geometric source move with
its exact word comparison, so no unproved transitivity is needed.
In a one-crosscap neighborhood of a cone, slide that cone once
around the one-sided core.  On the free generators $(d,x)$,
with exterior boundary word $d^2x$, the move is
\[
        d\longmapsto dx,\qquad
        x\longmapsto (dx)^{-1}x^{-1}(dx).              \tag{B.41.9}
\]
It preserves $d^2x$ exactly and sends the cone generator to a
conjugate of its inverse.  It is an automorphism: first replace
the free basis $(d,x)$ by $(dx,x)$, then invert and conjugate
the second member by the first.  Cutting the crosscap
neighborhood along its core gives the stated cone slide;
the unchanged exterior boundary permits extension by the
identity to the other crosscap.
In either abelian target, where $x$ has order two, this
replaces the image of $d$ by its product with $h$ and leaves
the cone image unchanged.  Apply it independently at the two
cones to connect all four assignments.  This gives one class
for each of $C_4$ and $C_2^2$.  Their different exponents
distinguish the two full groups.

For each surviving monodromy the local cone groups inject.
The orbifold covering subgroup is therefore torsion-free,
and it lies in the orientation kernel by construction.
The Klein-bottle orbifold Euler characteristic is
$-\sum_j(1-1/m_j)$; multiplying it by $|G|$ gives $-4$ in
all six cases.  The covering surface is consequently closed,
orientable, and of genus three.  This proves existence and
faithfulness of every listed action, as well as completeness.

There are no mirrors, so no reflection sections.
For each cone generator $x_j$ of prime order $m_j=2$ or $3$,
the points of its critical orbit are exactly
$G/\langle x_j\rangle$.  At the coset represented by $g$,
the stabilizer is $g\langle x_j\rangle g^{-1}$.  Write
$\varepsilon:G\to\{\pm1\}$ for orientation.  The multiplier
map $\mu_{j,g}$ on the section represented by that point is
\[
      \mu_{j,g}(gx_j^kg^{-1})=
                  \varepsilon(g)\frac{k}{m_j}\pmod{\Z}.
                                                        \tag{B.41.10}
\]
Every such point supplies one nontrivial partial section,
which is maximal; its centralizer in the surface group is
trivial by B.3.  No whole $G$ is a local stabilizer, so there
are no full sections.  The resulting coset counts are
exactly (B.41.1).  The program records the permutations of
these cosets, their stabilizers and the multiplier of every
local element.  In the $C_6$ row the two point generators
have multipliers $1/3$ and $-1/3$ under the fixed positive
generator; the negative coset exchanges the two.

The executed standard-library program
\path{catalogue_genus3_klein.py} independently enumerates
all pairs (B.41.3), checks associativity of their finite
multiplication tables and the orientation homomorphism,
enumerates every smooth generating monodromy, and computes
its orbit under the complete source and target automorphism
actions.  It retains the unsuccessful oriented extensions
with their empty monodromy lists, so the exclusions above
are part of the reproducible computation.
\end{proof}

\section{C. Equivariant comparison, drilling, and the topology of special groups}
\label{sec:part-c}

Throughout this part, $U=X\setminus S$ is connected, $X$ is a closed
orientable surface, and $S$ is a finite set of cardinality $\nu$. We put
$A=\operatorname{Homeo}(X,S)$ with its compact-open topology and
$K=\pi_0(A)$. Thus $A$ permits both permutations of $S$ and orientation
reversal; $A^0$ fixes $S$ pointwise and preserves orientation. We call the
type hyperbolic when $2g-2+\nu>0$. A superscript $u$ means the centralizer
of the image of the specified homomorphism $u:\Gamma\to A$. In particular,
$(A^0)^u=A^0\cap A^u$; it is not, by definition, the identity component
of $A^u$. We distinguish these two groups until connectedness is proved.

\subsection{The categories and their exact fidelity conditions}

\begin{claim}{C.1}{Refuted}{The unrestricted full target has too many arrows}
The literal comparison from surfaces and homeomorphisms to the full
$2$-category of peripheral groupoid diagrams and arbitrary functors is
not locally essentially surjective, already for $\Gamma=1$ and type
$(0,3)$. The correct target for a realization biequivalence is the
$2$-groupoid of such diagrams with equivalences as $1$-arrows.
\end{claim}
\passage{\S{}23; folios 134--136; scans 267--270}

\begin{proof}
For the three-punctured sphere put $\pi=\langle x,y\rangle$ and take the
peripheral diagram
\[
 \coprod_{i=1}^3 B\Z\longrightarrow B\pi,
 \qquad 1\longmapsto x,\ y,\ (xy)^{-1}.
\]
The identity on its three-object index set, the trivial homomorphism on
each copy of $\Z$, and the trivial homomorphism $\pi\to\pi$ define a
strictly commuting endomorphism of this diagram. Its map on
$\pi_{\rm ab}=\Z^2$ is zero. A homeomorphism induces an automorphism on
this abelianization. Natural isomorphisms between functors induce inner
conjugations, which act trivially on abelianization. Consequently this
endomorphism cannot become isomorphic to the image of a homeomorphism.
This is a counterexample to the literal fullness assertion, not to its
core version.

Here is the precise core version we use. An admissible surface is a
finite disjoint union of finite-type orientable surfaces, each genus-zero
component having at least one puncture, as on scan 264, folio 133. The
source $2$-groupoid has strict $\Gamma$-actions as objects, equivariant
homeomorphisms as $1$-arrows, and endpoint-preserving homotopy classes of
equivariant isotopies as $2$-arrows. An object of the target is a diagram
of groupoids over $B\Gamma$
\[
 D_\Gamma\longrightarrow U_\Gamma\longrightarrow B\Gamma.
\]
Its homotopy fibre over the base object is a peripheral diagram
$\coprod_{i\in I}B\Z\to B\pi$ with the surface presentation
$\prod_{j=1}^{g}[a_j,b_j]\prod_{i=1}^{\nu}c_i=1$; the $i$th map sends
$1$ to $c_i$. The fibre condition includes the common orientation module
and permits the noninjective map in type $(0,1)$ and the repeated image
in type $(0,2)$. This is the condition $M_{\rm lac}$. For $I=\varnothing$
we use the closed surface group, with $g\geq1$.

A target $1$-arrow is a pair of groupoid equivalences over $B\Gamma$
together with the invertible natural transformation expressing
commutativity of the peripheral square. A $2$-arrow is a pair of natural
isomorphisms, with identity projection to $B\Gamma$, satisfying that
square's compatibility. Composition is composition of functors with
pasting of these specified natural transformations. Equivalently these
are equivalences of the corresponding opfibrations over
$[1]\times B\Gamma$ that preserve the displayed fibres. The topological
comparison takes the fundamental groupoids of the homotopy quotients of
a $\Gamma$-invariant punctured collar system and of $U$. Two collar
systems have a common smaller system; the induced equivalences and
natural transformations specify the independence of this choice.
No arbitrary functor is silently replaced by an equivalence.
\end{proof}

\begin{claim}{C.2}{Proved}{The coset category and the quotient correction}
For a hyperbolic peripheral pair $(\pi;(L_i)_{i\in I})$, the category
$\mathscr C$ with objects $u,p_i$, automorphism groups
$\Aut(u)=\pi$, $\Aut(p_i)=1$, and
$\Hom(p_i,u)=\pi/L_i$ determines the pair up to an isomorphism and
simultaneous inner conjugation. Its equivalence groupoid has components
$\Out_{\rm lac}(\pi)$ and trivial automorphism groups. For an action with
extension $1\to\pi\to E\to\Gamma\to1$, the corresponding category over
$B\Gamma$ uses the bimodule $E/L_i$, with right action of
$E_i/L_i=\Gamma_i$, not $E/E_i$.
\end{claim}
\passage{\S{}23; folios 133--136; scans 264--270}

\begin{proof}
The stabilizer in $\Aut(u)$ of a selected arrow $p_i\to u$ is $L_i$.
Changing that arrow to $aL_i$ conjugates its stabilizer by $a$. An
equivalence fixes the unique object with nontrivial noncyclic
endomorphism group and permutes the other objects. Its restriction at
$u$ is an automorphism $\alpha$ of $\pi$. Equivariance of its maps on
Hom-sets says precisely that $\alpha(L_i)$ is conjugate to
$L_{\sigma(i)}$. Such a map of transitive sets is determined by a coset
whose stabilizer has this property. The equality
$N_\pi(L_i)=L_i$ makes that coset unique. A natural transformation between
two such equivalences has its component $a\in\pi$ at $u$, its identity
components at the $p_i$, and imposes $\alpha'=\operatorname{Int}(a)\alpha$
and the corresponding multiplication of all cosets by $a$. An
automorphism of an equivalence therefore lies in $Z(\pi)=1$.
This proves the assertion about the equivalence groupoid and gives both
comparison constructions explicitly.

For the geometric exit interpretation, the space of exits from $p_i$ to
$u$ is the homotopy fibre of the link map $BL_i\to B\pi$. Its components
are $\pi/L_i$ and its fundamental groups are the kernels of
$L_i\to\pi$; its higher homotopy groups vanish. The peripheral injection
makes those components contractible. Thus the ordinary category above
has exactly these exit mapping spaces. This also identifies the exact
failure of this ordinary model for the plane: the fibre of
$B\Z\to *$ is $B\Z$, not a discrete space. More explicitly the
Grothendieck category of $BL_i\to B\pi$ has automorphism groups
$L_i,\pi$ and forward Hom-set $\pi$. Its comparison with the exit
category is identity at the generic group, trivial at the special
endomorphism group, and $\pi\to\pi/L_i$ on forward arrows. Conversely
the translation groupoid of the generic action on the exit Hom-set
is $(\pi/L_i)//\pi\simeq BL_i$, with its map to $B\pi$. Thus the
peripheral diagram is recovered by this explicit link construction;
these two displayed categories are not silently identified. On
classifying spaces the forward comparison collapses each link to its
special point, realizing the filling map
$B\pi\to B\pi\cup_{\coprod BL_i}\coprod *$.

In the action category take one special object for each $\Gamma$-orbit
of punctures, choosing $i$ in that orbit. The generic object has
automorphism group $E$, the special object has automorphism group
$\Gamma_i$, and composition
on $E/L_i$ is left multiplication by $E$ and right multiplication by
$E_i/L_i$. These actions commute. The homotopy fibre over the base object of
$B\Gamma$ has special objects indexed by $\Gamma/\Gamma_i$ for each
chosen representative, and hence recovers every upstairs puncture. An
equivalent unskeletonized model uses all punctures and the nonempty
$\Gamma$-transporter sets between those in the same orbit. The
projection of a morphism $eL_i$
to $B\Gamma$ is the image $\bar e\in\Gamma$; it is well-defined because
$L_i\subset\pi$. It is not well-defined on $E/E_i$ when
$\Gamma_i\ne1$: representatives $e$ and $ea$, $a\in E_i$, can have
different images in $\Gamma$.

For a geometric example take $U=\mathbf P^1\setminus\{0,\infty,1,-1\}$,
$\Gamma=C_2$ acting by $z\mapsto-z$, and quotient coordinate $w=z^2$.
Then $E=\pi_1(\mathbf P^1\setminus\{0,1,\infty\})=F_2$. With $a$ a
loop about $w=0$, the upstairs inertia is $L_0=\langle a^2\rangle$
and $E_0=\langle a\rangle$. The right $C_2$ action on
$E/\langle a^2\rangle$ is nontrivial, whereas replacing it by
$E/\langle a\rangle$ erases that action. This is the exact morphism
obstruction to the quotient-category inference in the prompt.
\end{proof}

\begin{claim}{C.3}{Proved}{The four levels and the coefficient ladder}
Let $A$ be a locally path-connected topological group whose components
are $1$-types. For the comparison from strict actions to weak actions in
the categorical group $\Pi_1A$, the first three fidelity conditions are,
on each comparison $M\to N$ of transport spaces,
\[
\begin{array}{c|l}
0&\pi_1M\to\pi_1N\text{ injective at every point},\\
1&\pi_0M\to\pi_0N\text{ injective and }\pi_1M\to\pi_1N\text{ isomorphic},\\
2&\pi_0M\to\pi_0N\text{ bijective and }\pi_1M\to\pi_1N\text{ isomorphic}.
\end{array}
\]
Level $3$ adds essential surjectivity on action objects. When $A^0$ is
simply connected, these give exactly scans 279--280: vanishing of
$\pi_1((A^0)^u)$; connectedness of $(A^0)^u$; conjugacy of two lifts of the
same $\Gamma\to K$ by an element of $A^0$; and existence of a lift.
\end{claim}
\passage{\S{}23; folios 136--137; scans 271--273; \S{}23; folios 140--141; scans 279--281}

\begin{proof}
The source Hom-groupoid is $\Pi_1M$. A functor of groupoids is faithful
exactly when it is injective on the automorphism group at each object.
It is fully faithful exactly when that condition is an isomorphism and
two source objects whose images are connected are already connected.
It is an equivalence exactly when it is fully faithful and every target
component is hit. Applying these three elementary tests proves the
table, including the cumulative meaning of each level.

When $\pi_1(A^0)=0$, the categorical group $\Pi_1A$ is equivalent to the
discrete group $K$. For strict homomorphisms $u,v:\Gamma\to A$ the map
of transport spaces is
\[
 \{a\in A:a u(\gamma)a^{-1}=v(\gamma)\ (\gamma\in\Gamma)\}
 \longrightarrow
 \{k\in K:k\bar u(\gamma)k^{-1}=\bar v(\gamma)\}.
\]
A nonempty source is a torsor under $A^u$. Its fibres over points of the
target are torsors under $(A^0)^u$. The first two lines of the table are
therefore the first two conditions stated. For the third, given a target
transporter $k$, choose $a\in A$ representing $k$. The two actions
$a u a^{-1}$ and $v$ have the same projection to $K$; conjugating them
by an element of $A^0$ produces the prescribed transporter. Conversely,
apply transporter surjectivity to the identity element of $K$. This
also verifies the stronger marginal correction on scan 280.

For comparison with the GALOIS ladder, first take the $1$-truncations
of the transport spaces; fixed spaces are not assumed to inherit the
ambient group's higher truncation. On a map of $1$-types $M\to N$, set-valued constant coefficients test bijectivity on $\pi_0$;
full faithfulness of restriction tests that bijectivity together with
surjectivity on $\pi_1$; equivalence tests bijectivity on $\pi_0$ and an
isomorphism on $\pi_1$. For groupoid-valued coefficients all three tests
coincide with a weak equivalence, because $\pi_iM=\pi_iN=0$ for $i\geq2$;
the same is true for $2$-type coefficients. These assertions follow by
applying the long exact sequence to the homotopy fibres: an $r$-connected
fibre gives isomorphisms through degree $r$ and surjectivity in degree
$r+1$. Thus level $2$ is the $r=0$ equivalence test on the transport
spaces. Levels $0$ and $1$ do not equal a rung in that ladder: their
$\pi_1$-injectivity and $\pi_0$-injectivity tests have been written out
above, whereas each rung of the coefficient ladder demands a component
bijection. Level $3$ concerns objects rather than this map of Hom-spaces.
This is a comparison of the specified maps, not an identification based
only on the words ``faithful'' and ``equivalence''.
\end{proof}

\begin{claim}{C.4}{Proved}{Coherent actions and rectification on spaces}
For a discrete group $P$, the space of self-homotopy equivalences of
$BP$ has components indexed by $\Out(P)$, each with homotopy type
$BZ(P)$. In particular, for the centre-free surface groups the map from
outer actions to coherent actions is an equivalence of action groupoids.
Every coherent action of a discrete group on a space has a strict
realization on a space of the same homotopy type. For $P=\Z$ or $\Z^2$,
the obstruction class in degree $3$ is zero and the coherent actions
with a specified linear action are classified by $H^2(\Gamma,P)$.
\end{claim}
\passage{\S{}23; folios 131--132; scans 261--263; \S{}23; folios 135--141; scans 268--281}

\begin{proof}
The groupoid of functors $BP\to BP$ has homomorphisms $P\to P$ as
objects and conjugating elements of $P$ as arrows. Its equivalence
subgroupoid has automorphisms as objects. The stabilizer of an
automorphism is $Z(P)$, and the orbits are $\Out(P)$. The nerve computes
the mapping space because nerves of groupoids represent $1$-types.
Delooping the resulting categorical group has $\pi_1=\Out(P)$,
$\pi_2=Z(P)$ and no higher homotopy groups. Its Postnikov class is the
usual extension obstruction in $H^3(\Out(P);Z(P))$. For $P$ abelian the
map $\Aut(P)\to\Out(P)$ is the identity and gives a strict section;
the categorical group is $BP\rtimes\Aut(P)$, so this Postnikov class
vanishes. A choice of coherence is then a $2$-cocycle with values in the
specified $\Gamma$-module $P$; changing it by a coherent equivalence
changes the cocycle by a coboundary. This proves the stated
classification without asserting a nonzero obstruction for either
abelian surface group.

For rectification, represent a coherent action by a fibration
$p:Y\to B\Gamma$ with the prescribed homotopy fibre. Pull it back along
$E\Gamma\to B\Gamma$. The pullback $Y'=E\Gamma\times_{B\Gamma}Y$ has a
strict deck action of $\Gamma$. Its projection to the contractible
space $E\Gamma$ has fibre the original homotopy type, so the inclusion
of any chosen fibre is a homotopy equivalence. Its homotopy quotient
recovers $Y\to B\Gamma$. This construction applies simultaneously to
peripheral diagrams by the diagrammatic pullback. It produces a space
model; the additional demand that the model be the original surface is
exactly the realization problem addressed below. For continuous
profinite actions the objects and the homotopy quotient must instead be
specified in a continuous/profinite homotopy theory; the preceding
construction makes no untyped transfer to that setting.
\end{proof}

\subsection{Finite actions on hyperbolic surfaces}

\begin{claim}{C.5}{Proved}{Existence, with punctures and orientation reversal}
For a hyperbolic type and every finite group $\Gamma$, each
homomorphism $\rho:\Gamma\to K$ lifts to a homomorphism
$u:\Gamma\to A$. Its image can be chosen to consist of isometries of a
complete finite-area hyperbolic metric on $U$.
\end{claim}
\passage{\S{}23; folios 136, 140--141; scans 270--271, 279--280; \S{}24; folio 147; scan 293}

\begin{proof}
Let $H=\rho(\Gamma)$, a finite subgroup of the extended mapping class
group. Kerckhoff's Theorem 5, together with the explicit extension in
\S IV, pp.~255--256, of the original paper \cite{C_Kerckhoff1983}, gives
a fixed hyperbolic structure for $H$. That extension explicitly permits
cusps; the paper's mapping-class convention explicitly includes
orientation-reversing maps. We use the original paper, rather than
inferring the punctured case from the announcement.

Choose the isometry in each class of $H$. It is unique. Indeed an
isometry isotopic to the identity has a lift to $\mathbb H^2$ commuting
with the surface deck group. That deck group contains two hyperbolic
elements with different endpoint pairs, so an isometry commuting with
both fixes at least three points of the ideal boundary and is the
identity. The same argument covers a finite-area punctured surface;
its non-elementary deck group again contains those elements. Products
of the chosen isometries represent products in $H$, and uniqueness
therefore makes the choices a homomorphism $H\to A$. Composing with
$\rho$ proves the claim. The isometries extend to the end compactification
and permute $S$, so the target really is the specified group $A$.
\end{proof}

\begin{claim}{C.6}{Proved}{Uniqueness up to the identity component}
Let $u,v:\Gamma\to A$ be homomorphisms with $\Gamma$ finite, $U$
hyperbolic, and $\bar u=\bar v:\Gamma\to K$. There is an
$h\in A^0$ such that $h u(\gamma)h^{-1}=v(\gamma)$ for every $\gamma$.
This holds for topological actions without a smoothness assumption.
\end{claim}
\passage{\S{}23; folios 136, 140--141; scans 270--271, 279--280}

\begin{proof}
We first justify the analytic model for either action. For a point
$x\in U$ and its finite stabilizer $H$, lift each element of $H$ to the
universal cover with the chosen lift of $x$ fixed. Uniqueness of a lift
with a specified value makes these lifts a finite group. Compactifying
$\widetilde U\simeq\R^2$ by one point gives a finite action on $S^2$
with two common fixed points. Kolev's compact-sphere theorem
\cite[Theorem 1.2]{bKolev2006} conjugates that action into $O(3)$.
A subgroup of $O(3)$ fixing two distinct points is contained in the
orthogonal group of the perpendicular plane; choosing a sufficiently
small round disc at the first point and pulling it back yields a common
invariant disc with an $O(2)$ action. Choosing it small enough makes its
translates by distinct cosets of $H$ disjoint. At a puncture we apply the
same finite-stabilizer disc argument on the compactification, as in the
finite-disc theorem proved from the same source.

The quotient consequently has ordinary charts, cone charts, mirror
charts, and corner charts with their finite cyclic or dihedral local
groups. A finite triangulation respecting these strata supplies a
smooth orbifold structure: smooth each triangle, use collar charts
along its edges, and use the displayed linear quotient charts at its
vertices. Pulling the charts to $X$ makes the original finite action
smooth for a compatible smooth structure. Average a metric and use its
conformal class. Uniformization of the punctured conformal surface gives
a finite-type structure for which each operator is conformal or
anticonformal. Thus we have conformal structures $c_u,c_v$ on the same
marked topological surface.

The extremal quasiconformal map $h:(U,c_u)\to(U,c_v)$ in the identity
isotopy class is unique, by the finite-type Teichm\"uller theorem
\cite[\S3(A)--(D)]{C_Bers1981}. For each $\gamma$, the map
$v(\gamma)^{-1}h u(\gamma)$ is orientation preserving: $u(\gamma)$ and
$v(\gamma)$ have the same orientation sign. Pre- and postcomposition by
these conformal or anticonformal maps leaves the maximal dilatation
unchanged. The equality of the two mapping-class homomorphisms says
that this map is isotopic to the identity, with the punctures retained.
Uniqueness therefore gives $v(\gamma)^{-1}h u(\gamma)=h$. The
quasiconformal map extends across the finite punctures; hence
$h\in A^0$ and has the asserted conjugacy property.
\end{proof}

\begin{claim}{C.7}{Proved}{The finite centralizer component, including mirrors}
For a finite action $u$ on a hyperbolic $U$, the group $(A^0)^u$ is
contractible. In particular it is the identity component of $A^u$.
\end{claim}
\passage{\S{}\S{}23--24; folios 140--144; scans 279--286}

\begin{proof}
Replace $\Gamma$ by its image $H$. The preceding local construction
makes $H$ a group of hyperbolic isometries. Let $\pi$ be the surface deck
group and let $E$ be the full group of lifts of $H$ to $\mathbb H^2$.
It is a discrete group of isometries, possibly with reflections; its
quotient $Q$ is a finite-type hyperbolic orbifold. Write $D$ for descent
of a centralizing homeomorphism to the labelled underlying surface of
$Q$, with cone points, corner points, mirror boundary, and punctures
retained.

We give the injectivity argument at the place where an
orientation-preserving-only Birman--Hilden theorem would not suffice.
For $f\in(A^0)^u$, the identity isotopy supplies a lift $\widetilde f$
that commutes with $\pi$. Since $f$ commutes with $H$, conjugation by
$\widetilde f$ induces an automorphism $\alpha$ of $E$ that is the
identity on $\pi$. For $e\in E$ and $a\in\pi$,
\[
 \alpha(e)a\alpha(e)^{-1}=\alpha(eae^{-1})=eae^{-1}.
\]
Thus $e^{-1}\alpha(e)$ centralizes $\pi$, and hence is the identity by
the distinct-endpoint argument in C.5. Consequently $\alpha=\id_E$.
The orbifold Dehn--Nielsen--Baer injectivity theorem, with its mirror
and corner structure, makes $D(f)$ isotopic to the identity through
homeomorphisms preserving those structures; the relevant precise form
is \cite[Theorem 2 and Proposition 3]{C_KolbeEvans2021}. At the algebraic
step in their Lemma 3 we have used the calculation just written:
centralizers of individual reflections need not be cyclic, but the
centralizer of the non-elementary surface subgroup is trivial. Thus
that inapplicable sentence of their proof is not being used. There is an orientation hypothesis in that source when the labelled
underlying quotient is orientable, and we verify it here. The canonical
$\pi$-equivariant lift $\widetilde f$ extends by the identity on
$\partial_\infty\mathbb H^2$. For a compact surface this follows from
bounded displacement on a compact fundamental domain. For a punctured
surface use a compact truncation and disjoint invariant horodiscs:
displacement is bounded on the truncated part, and preservation of each
end sends sufficiently deep cusp horodiscs into the corresponding cusp
horodiscs. A sequence escaping through distinct horodiscs has the same
ideal limit before and after the map, since their visual diameters tend
to zero; a sequence in one horodisc has its cusp endpoint as its only
ideal limit. These observations prove the claimed continuous extension.
Because $\alpha=\id$, $\widetilde f$ preserves each reflection axis.
It fixes that axis's two ideal endpoints individually, so preserves its
linear orientation, and, being orientation preserving on $\mathbb H^2$,
preserves its two sides. It therefore preserves every mirror chamber.
On such a chamber the quotient orientation is the one induced from
$\mathbb H^2$; hence $D(f)$ preserves it. Without mirrors an orientable
quotient is covered orientably and the same assertion is immediate.
Thus the theorem is applied inside its stated homeomorphism group.

Let $\mathcal H(Q)^0$ denote the identity component of that labelled
homeomorphism group. Every path starting at its identity lifts uniquely,
starting at the identity upstairs, to a path centralizing $H$. Lift
first on the regular covering locus. The local cyclic and dihedral
models extend the lift over cone points, corners and mirrors; continuity
of the extended lift follows in the charts from the finite quotient
map. The same construction for a parameter family proves the path and
homotopy lifting properties. The only ambiguity for an endpoint is a
central deck transformation. Such a deck transformation cannot belong
to $A^0$ unless it is the identity, by C.5's isometry argument. It follows
that descent identifies $(A^0)^u$ with $\mathcal H(Q)^0$, as soon as the
latter is simply connected; the construction is continuous in local
isotopy charts, so this is an isomorphism of topological groups.

We check that this group is contractible, including the small underlying
quotients. End-compactify the quotient and retain its finitely many
interior and boundary marked points. Yagasaki's compact-pair Lemma 2.1
\cite{A_Yagasaki2000} gives contractibility outside its explicit
exceptional pairs. A marked point on a boundary component is covered
by Lemma 2.1(ii). With no such point, the exceptional disk has at most
one interior marked point, the annulus and M\"obius band have none, the
sphere has at most two, the projective plane at most one, and the torus
and Klein bottle have none. Each of these has nonnegative orbifold
Euler characteristic: an interior cone subtracts $1-1/m$, an interior
puncture subtracts $1$, and a corner or boundary puncture would already
be a boundary marked point. But
$\chi_{\rm orb}(Q)=\chi(U)/|H|<0$. None occurs. The homeomorphism group
with labels has the same identity component as the group fixing all
these finite marks pointwise. Yagasaki's theorem therefore applies.
This proves contractibility, and also supplies the simple connectedness
used in the lifting construction.
\end{proof}

\begin{claim}{C.8}{Proved}{The finite core biequivalence and homotopy fixed points}
On hyperbolic surfaces, for every finite $\Gamma$ the corrected
comparison of C.1 has all four fidelity levels and is a biequivalence.
For a specified action $u$, the comparison
\[
 A^u\longrightarrow A^{h\Gamma}
      =\operatorname{Map}_\Gamma(E\Gamma,A),
 \qquad a\longmapsto(\text{constant map with value }a),
\]
where $\Gamma$ acts on $A$ by conjugation through $u$, is a weak homotopy
equivalence. In standard CW models it is a homotopy equivalence.
\end{claim}
\passage{\S{}23; folios 136--141; scans 270--281}

\begin{proof}
The nonequivariant peripheral comparison identifies $K$ with
$\Out_{\rm lac}(\pi)$ and has trivial $2$-automorphism groups, by the
peripheral Dehn--Nielsen--Baer theorem
\cite[Theorems 8.1 and 8.8]{A_FarbMargalit2012} and $Z(\pi)=1$.
C.7 gives levels $0$ and $1$, C.6 gives level $2$ by the transporter
argument of C.3, and C.5 gives level $3$. These arguments also treat a
nonfaithful homomorphism: its finite image is the action group, and a
finite-order homeomorphism in $A^0$ is the identity. Disconnected
admissible surfaces reduce to representatives of the finitely many
orbits of connected components and their stabilizers, with the stated
transporters providing the arrows between choices. This is the desired
biequivalence on its hyperbolic part.

The map $A\to K$ is $\Gamma$-equivariant. Its fibres are contractible by
the same homeomorphism-group theorem with $H=1$. Taking derived
homotopy fixed points gives a weak equivalence
$A^{h\Gamma}\to K^\Gamma=C_K(\bar u(\Gamma))$, the latter a discrete
space. The constant-map comparison followed by this map is exactly
$\pi_0(A^u)\to C_K(\bar u(\Gamma))$. Its surjectivity follows by applying
C.6 to $a u a^{-1}$ and $u$ for any representative $a$ of a centralizing
mapping class; its injectivity and its contractible fibres are C.7.
Thus the displayed map is a weak equivalence. The source components
have CW homotopy type by the ANR statement for quotient homeomorphism
groups in \cite{A_Yagasaki2000}; taking the usual CW model of the
derived target turns it into a homotopy equivalence. We state the model
convention because scan 281 writes ``homotopisme'' without specifying
an infinite-dimensional mapping-space model.
\end{proof}

\subsection{The exceptional types and infinite operator groups}

\begin{claim}{C.9}{Proved}{Corrected abelian tests and a torus calculation}
For $\pi_1(A^0)\ne0$, replacement of the four tests by connectedness and
uniqueness of lifts of $\Gamma\to K$ loses the coherent action data.
The valid tests are C.3 with weak transport spaces. At a specified
strict action the target automorphism space is the $1$-type
$A^{h\Gamma}$ and
\[
 \pi_1(A^{h\Gamma},1)=\pi_1(A,1)^\Gamma.
\]
For $A=\operatorname{Homeo}(T^2)$ and $u(C_2)=\{1,-1\}$, the group
$(A^0)^u$ has four components, all contractible, whereas the right-hand
invariant fundamental group is zero. Thus the naive connectedness test
fails; the corrected comparison does not fail in this example.
\end{claim}
\passage{\S{}23; folio 141; scans 280--281}

\begin{proof}
The homotopy-fixed-point groupoid of a $1$-type has loops equal to the
fixed automorphisms of its base object, proving the displayed equality.
Its components include the $H^1(\Gamma;\pi_1A)$ coherence data, with
possible obstruction for an individual component of $A$. Hence it is
not legitimate to test only $\pi_0A$ when $\pi_1A\ne0$.

For the torus, evaluation at $0$ gives
$A^0\cong T^2\times A^0_{\{0\}}$ by
$f\mapsto(f(0),t_{-f(0)}f)$. The second factor is contractible by the
homeomorphism-group theorem for a torus with a marked point. In the
centralizer of $-1$, $f(0)$ is one of the four points of $T^2[2]$.
Each is realized by translation. In the fibre $f(0)=0$, regarding the
torus as once punctured makes it hyperbolic; the fibre is contractible
by C.7. This proves the four-component assertion. Conjugation by $-1$
acts as $-I$ on $\pi_1(A^0)=\Z^2$, so its invariants vanish.
The inclusion of translations $T^2\hookrightarrow A^0$ is equivariant
and a weak equivalence; its homotopy fixed points have components
$H^1(C_2;\Z^2_{-})=(\Z/2)^2$ and zero higher homotopy groups. The four
translations just identified realize exactly these four components.

For the sphere, the peripheral groupoid diagram is empty over $*$ and
cannot encode its homotopy type; $\pi_2(S^2)=\Z$ already records the
failure of asphericity. The plane, annulus and torus have
$A^0\simeq S^1,S^1,T^2$, respectively; in the plane the peripheral
map $B\Z\to *$ retains the circle of automorphism $2$-arrows, and in the
annulus the two ends remain indexed even though their subgroup images
coincide. These exact maps explain why collapsing to a bare outer group
would erase data. The complete finite-action comparison for these models, including
its parametric centralizer proof, is C.10.
\end{proof}

\begin{claim}{C.10}{Proved}{Every finite action in the exceptional aspherical types}
For $U=\R^2$, $S^1\times\R$, or $T^2$, and every finite strict action,
the fixed-to-homotopy-fixed comparison in C.8 is a weak equivalence
and is a homotopy equivalence on CW models. The corrected peripheral
$2$-groupoid comparison has all four fidelity levels for these types.
The tests are those of C.3; ordinary outer actions alone do not retain
the required coherence classes.
\end{claim}
\passage{\S{}23; folios 136--141; scans 270--281}

\begin{proof}
We prove the parametric assertion, not merely classification of individual
actions. Write respectively
\[
 G=O(2),\qquad
 G=(\R/\Z)\rtimes(C_2\times C_2),\qquad
 G=(\R^2/\Z^2)\rtimes\operatorname{GL}_2(\Z).
\]
The annulus element $(a,\epsilon,\eta)$ sends
$(\theta,t)$ to $(\epsilon\theta+a,\eta t)$; the first sign acts on the
circle factor and the second on the two ends. In every case $G\to A$
is a weak equivalence. Here is a direct check of the identity components.
On the torus, evaluation at $0$ splits $A^0$ as $T^2$ times its
point-fixing subgroup. An identity isotopy can be made point fixing by
postcomposing at time $s$ with translation by the negative of the trace
of $0$. Thus that subgroup is its own identity component, and is
contractible by the punctured-torus case of the homeomorphism-group
theorem \cite[Theorem 1.1 and Lemma 2.1]{A_Yagasaki2000}. The same adjustment on the annulus uses angular rotation and
translation of its real coordinate. Its point-fixing subgroup is the
identity group of a hyperbolic three-ended sphere and is contractible,
so $A^0\simeq S^1$. On the plane first fix $0$ by translations and
then fix $1$ using multiplication by a nonzero complex number. The
remaining identity group fixes $0,1,\infty$ and is contractible; hence
$A^0\simeq S^1$ with generator the rotation circle. The component
groups are respectively $C_2$, $C_2\times C_2$, and
$\operatorname{GL}_2(\Z)$: orientation and end signs give the first
two, and the torus action on its lattice gives the third. Cutting the
plane or annulus along a ray or spanning arc verifies the stated
kernels; the torus case is the genus-one Dehn--Nielsen--Baer theorem.

By the simultaneous finite normal forms proved in B.1 and B.14,
conjugate the given finite image $H$ into $G$. We will prove
\[
                    C_G(H)\longrightarrow C_A(H)                 \tag{C.10.1}
\]
is a weak equivalence by examining its fibres. Conjugation by the chosen
homeomorphism transports both the fixed and derived-fixed comparisons,
so this reduction is an isomorphism of the actual diagrams.

\emph{Torus displacement.}
For $f\in C_A(H)$ let $M\in\operatorname{GL}_2(\Z)$ be its lattice
action and lift it as $F(x)=Mx+d(x)$, with $d$ periodic. Put
$b=\int_{T^2}d(x)\,dx$ and $g(x)=Mx+b$ modulo $\Z^2$.
Changing the lift changes $b$ by an integral vector and leaves $g$
unchanged. For $\gamma(x)=A_\gamma x+t_\gamma$ the commutation equation
is
\[
 d(A_\gamma x+t_\gamma)-A_\gamma d(x)
       =(I-M)t_\gamma+k_\gamma,\qquad k_\gamma\in\Z^2.
\]
Integration preserves its left-hand change of variables, because the
affine operator preserves Haar measure. It gives the same equation
with $d$ replaced by $b$. Thus $g$ commutes with $H$ and
$g^{-1}f\in A^0$ has a unique lift of zero mean displacement.
Consequently multiplication is a homeomorphism of spaces
\[
 C_G(H)\times\mathcal K\ \xrightarrow{\sim}\ C_A(H),
\]
where $\mathcal K$ consists of these zero-mean maps. The averaging
map need not be a group homomorphism; only this explicit space
factorization is used. A zero-mean lift commutes exactly with every
lift of $H$: the possible integral commutator above has zero average
and is therefore zero.

First divide by the finite translation subgroup $N=H\cap T^2$.
If $a$ represents a member of $N$ of order $m$, commutation implies
$F(x+a)=F(x)+a+k$. Iterating $m$ times and using $ma\in\Z^2$
forces $k=0$ for a zero-mean identity lift. Hence its displacement is
periodic for the enlarged covering lattice. Descent and lift identify
$\mathcal K$ with the same zero-mean space on $T'=T^2/N$ for $H/N$.
This quotient group has faithful linear part, a finite subgroup $P$
of $\operatorname{GL}_2(\Z)$.

When $P$ contains a nonidentity positive operator, its positive subgroup
$K$ is cyclic, generated by an affine rotation with $I-A$ invertible.
It has a nonempty finite fixed set $R$ on $T'$, stable under $H/N$.
A zero-mean lift commutes with each lift of that rotation and fixes its
unique Euclidean fixed point; thus it fixes every point of $R$.
Choose $p\in R$. The point-fixing argument above places our map in
$\operatorname{Homeo}(T',p)^0$. C.7, applied to the cyclic action on
$T'\setminus\{p\}$, supplies an equivariant identity isotopy. This
isotopy fixes all of $R$, since $R$ is finite. Our map therefore belongs
to $\operatorname{Homeo}(T',R)^0$. Apply C.7 again, now to the full
$H/N$ action on the hyperbolic $T'\setminus R$. Its identity centralizer
is exactly $\mathcal K$: conversely an equivariant identity isotopy
relative to $R$ has commuting identity lifts, whose mean is fixed by
$A$ and hence is zero. This proves contractibility of $\mathcal K$.

The remaining linear groups are $P=1$ and a group of order two generated
by a reflection. For $P=1$, the space factorization on the identity
component is $\operatorname{Homeo}(T')^0\cong T'\times\mathcal K$.
The first factor is a section of the mean map. Since the left side has
the homotopy type $T'$, the mean map is an isomorphism on $\pi_1$ and
$\mathcal K$ is weakly contractible. It is a retract of the ANR on the
left, so it has CW homotopy type and is contractible.

For a reflected or glide-reflected torus the underlying quotient is an
annulus, a M\"obius band, or a Klein bottle. This exhausts the cases:
the action reverses orientation, has no isolated rotation point, and
the quotient has Euler characteristic zero; its boundary consists of
mirrors. A zero-mean lift $F$ commutes exactly with the Euclidean orbifold
group $E$. Each lifted reflection line is preserved, and bounded
displacement prevents reversal of its infinite directions. In
particular the compact annulus quotient's two boundary circles cannot
be interchanged. The explicit straight homotopy
\[
                  F_s(x)=(1-s)F(x)+sx
\]
is $E$-equivariant, because every element of $E$ is affine. On each
reflection line it stays on that same line. It therefore descends to
a homotopy of the quotient as a surface with boundary, preserving its
boundary throughout. Its quotient domain is compact, so it is proper.
Epstein's Theorem 6.3, with the topological Appendix
\cite{C_Epstein1966}, applies with exactly this proper,
boundary-preserving homotopy and gives an identity isotopy. This also
handles the boundary-free Klein bottle; Theorem 5.8 of that source
separately verifies the M\"obius-band case. Thus no unverified inference
from an abstract outer automorphism to an exceptional-surface isotopy
is needed here.

Descent consequently puts $\mathcal K$ in the identity group of the
labelled quotient. Every identity isotopy there lifts; among the two
lifts exactly one preserves the orientation upstairs. This identifies
the identity component of $C_{A^0}(H/N)$ with the quotient identity
homeomorphism group. By Yagasaki's actual exception list its homotopy
type is $S^1$ in all three cases. The mean factorization on that component
is $T'^P_0\times\mathcal K$, with $T'^P_0\cong S^1$. Here
$P$ acts linearly on the translation torus; this fixed translation
circle is present even for an affine glide with no fixed point.
The section by translations makes the mean map surjective on $\pi_1$;
a surjection $\Z\to\Z$ is an isomorphism. Hence $\mathcal K$ is
weakly contractible, and is contractible by the same ANR-retract argument.
This completes (C.10.1) for the torus, including free glide reflections.

\emph{Annulus displacement.}
Lift $f\in C_A(H)$ to
$F(\theta,t)=(\epsilon_f\theta+d(\theta,t),F_2(\theta,t))$ and put
$b=\int_0^1d(\theta,0)\,d\theta$. Define
$g(\theta,t)=(\epsilon_f\theta+b,\eta_f t)$, where $\eta_f$ is its
end sign. The commutation equation, integrated on the invariant core
$t=0$, proves that $g$ commutes with $H$. It gives the same continuous
factorization $C_A(H)\cong C_G(H)\times\mathcal K$, with zero-mean
identity lifts in $\mathcal K$. Divide first by the finite rotation
subgroup. The resulting faithful sign group is a subgroup of
$C_2\times C_2$.

When it contains $(-1,-1)$, that positive involution has two fixed
points. Its subgroup is normal. A zero-mean lift fixes each of the
corresponding fixed points in $\R\times\R$. The same two applications
of C.7 as in the torus argument, first marking one fixed point and then
the whole invariant fixed set, identify $\mathcal K$ with a contractible
centralizer on a hyperbolic punctured annulus. An equivariant identity
lift in that punctured model has mean zero because its angular linear
sign is $-1$, proving equality in both directions.

A subgroup without $(-1,-1)$ is trivial, generated by $(-1,1)$, or
generated by $(1,-1)$. In the first case the factorization and
$A^0\simeq S^1$ give contractibility as above. An angular reflection
has a strip as quotient, with its two mirror boundary lines and two
ends. Exact commutation of a zero-mean lift preserves both lines and
both ends. The same straight homotopy in $\R\times\R$ is equivariant
for all the affine covering operators and preserves the boundary
lines. It is proper on the strip quotient: for $t\to+\infty$ both
$t$ and $F_2(\theta,t)$ tend uniformly to $+\infty$ on a compact
angular fundamental interval, and for $t\to-\infty$ both tend to
$-\infty$. Every convex combination has the same uniform property.
Epstein's proper boundary-homotopy theorem therefore places the
descended map in the identity group. That group is contractible by
Yagasaki, so $\mathcal K$ is contractible.
Finally a radial reflection or radial glide reflection has quotient a
half-open annulus or a punctured projective plane, respectively. The
straight homotopy is again equivariant; it preserves the fixed mirror
line in the reflection case and has the just-verified properness at
the quotient's end. Epstein's Theorem 6.3, or Theorem 5.8 for the open
M\"obius band, gives the required identity isotopy. Descent and isotopy
lifting identify the upstairs centralizer identity component with that
group. Yagasaki gives $S^1$ in both cases. The mean factorization has a
section by the fixed angular circle, hence again forces
$\mathcal K$ to be contractible. These cases include every affine
finite action and retain the end permutation throughout.

\emph{Plane centralizers.}
For a cyclic rotation group of order $n>1$, a commuting map fixes $0$.
Descent through $z\mapsto z^n$ is a covering of the orientation-preserving
centralizer onto $\operatorname{Homeo}^+(\R^2,0)$, with deck group
$C_n$. The base is connected with homotopy type $S^1$. A loop of one
rotation downstairs lifts from the identity to the deck rotation by
$2\pi/n$. Its monodromy is therefore onto $C_n$, so the covering
centralizer is connected and the upstairs rotation circle induces the
isomorphism on $\pi_1$; all higher homotopy groups vanish. A negative
map centralizes the rotation only for $n=2$, as its action on the
punctured-plane fundamental group inverts the deck generator. In that
case conjugation supplies the second component. These are precisely
the components and circle types of $C_{O(2)}(C_n)$.

For a reflection, descent has half-plane quotient. A commuting map has
two independent signs: its orientation on the fixed line and its
orientation on the plane. The four diagonal orthogonal maps realize
them. Each fibre is the identity homeomorphism group of the half-plane,
which is contractible by Yagasaki. For a dihedral group with rotation
order $n>2$, a centralizing map preserves plane orientation and each
reflection line. Preservation of cyclic order of their rays permits
only the identity permutation or simultaneous reversal of every ray;
these are realized by $I$ and $-I$. For $n=2$ the two perpendicular
reflection lines permit four independent ray-sign choices, realized by
the four diagonal orthogonal maps. After dividing by the matching map,
descent is to a sector homeomorphism fixing its vertex and both boundary
rays. This group is contractible: the sector is a half-plane with a
marked boundary point, and its identity group is covered by Yagasaki's
boundary-point case. The boundary rays also make the lift unique.
The trivial group case was verified at the start. This proves (C.10.1)
for every finite subgroup of $O(2)$.

It remains to identify the derived comparison, including all components.
The following elementary torus calculation does so for each of our
three $G$. Let $T=V/\Lambda$, with $V$ real and $\Lambda$ a lattice,
and let a finite group act on a $T$-torsor by affine maps. Its translation
part is a cocycle $b\in Z^1(\Gamma,T)$, where $T$ and $\Lambda$
have the specified linear $\Gamma$-module structures. The homotopy-fixed space is
the section space of its Borel fibration. Its obstruction is the class
$\delta[b]\in H^2(\Gamma,\Lambda)$ of the extension obtained by lifting
the affine maps to $V$. Averaging cochains gives
$H^j(\Gamma,V)=0$ for $j>0$, so
$\delta:H^1(\Gamma,T)\to H^2(\Gamma,\Lambda)$ is an isomorphism.
Thus a homotopy fixed point exists exactly when $b$ is a coboundary,
which is exactly existence of an actual fixed point. After choosing one,
\[
 \pi_0(T^{h\Gamma})=H^1(\Gamma,\Lambda),\qquad
 \pi_1(T^{h\Gamma})=\Lambda^\Gamma,
\]
and there are no higher groups. The exact sequence
\[
 0\to\Lambda^\Gamma\to V^\Gamma\to T^\Gamma
       \to H^1(\Gamma,\Lambda)\to0
\]
identifies the actual fixed components and their fundamental groups
with exactly these groups. The constant-map comparison induces these
identifications. Apply this to every invariant component of $G$, each
of which is just such an affine torus torsor. Noninvariant components
contain neither fixed nor homotopy-fixed points. It follows that
$G^\Gamma\to G^{h\Gamma}$ is a weak equivalence. The equivariant weak
equivalence $G\to A$, together with (C.10.1), now proves the claimed
comparison for $A$ by the commutative square of fixed and derived-fixed
spaces.

Finally $BG$ is the split $2$-type with fundamental group
$\pi_0G$, second group $\Lambda=\pi_1G$, and zero Postnikov invariant.
Coherent finite actions with a specified homomorphism to $\pi_0G$ are
therefore classified by $H^2(\Gamma,\Lambda)$. The same connecting
isomorphism realizes each by a $T$-valued cocycle, hence by a strict
homomorphism $\Gamma\to G$. It also identifies coherent equivalences
in the identity component with coboundaries, namely conjugacies by
translations. The fixed-space calculation supplies their full
transport-groupoid comparison; a prescribed nonidentity transporter
component is treated by composing with its representative in $G$.
Thus the corrected four fidelity levels all hold. We finish by checking
that this is the actual peripheral comparison, not just an abstract
replacement of $A$. For the plane the diagram is $B\Z\to *$.
Its equivalence $2$-group has component group $\{\pm1\}$,
$2$-automorphism group $\Z$, and sign action. For the annulus use
$B\Z\amalg B\Z\to B\Z$ with the two maps $+1,-1$.
An equivalence is specified on components by a generic sign $\epsilon$
and an end permutation of sign $\eta$. Its peripheral sign is forced
to be $\epsilon\eta$. The two square conjugators can be removed by
peripheral natural transformations; the surviving automorphism of
such a diagram equivalence is one diagonal $\Z$, on which the generic
sign $\epsilon$ acts. Thus its $2$-group has exactly the two-sign
component group and lattice action of the annulus $G$ displayed above.
For the torus, the empty peripheral diagram over $B\Z^2$ has
component group $\operatorname{GL}_2(\Z)$ and $2$-automorphism group
$\Z^2$ with its standard action. In all three cases strict diagram
functors with zero square-conjugators give a section, so the
Postnikov class is zero. The actual comparison maps a full plane or
annulus rotation to the generator traced by a peripheral or generic
basepoint, and torus translation loops to the two lattice generators.
It is therefore an isomorphism on these specified $\pi_1$ groups,
as well as on components. This proves that the models and comparisons
used throughout the proof are the ones induced by C.1's peripheral
functor. In particular the annulus end permutation has never been
forgotten.
\end{proof}

\begin{claim}{C.11}{Proved}{Infinite cyclic failures at levels one and two}
For every hyperbolic surface containing a disc, there is a strict
$\Z$-action for which $(A^0)^u$ is disconnected. There also are two
strict $\Z$-actions with the same mapping-class homomorphism that are
not conjugate. Thus finite-group hypotheses cannot be removed from
levels $1$ and $2$.
\end{claim}
\passage{\S{}23; folios 140--141; scans 279--281}

\begin{proof}
In a coordinate disc choose two disjoint small round discs centered at
$q_\pm=\pm1$. Define $f$ to rotate each small disc by angle
$\sqrt2(1-r/r_0)$ at radius $r<r_0$, and to be the identity elsewhere.
This is a homeomorphism isotopic to the identity through the same
formula with the angle multiplied by a parameter. The two centres are
isolated fixed points; all other fixed points lie outside the interiors
of these discs. At radius $r_0/2$ the rotation angle is
$\sqrt2/2$, whose ratio to $2\pi$ is irrational, so $f$ has infinite
order. Let $h$ rotate the larger disc by angle $\pi$ on a concentric disc
containing the two small discs, tapering continuously to angle zero
before the outer boundary. It is also isotopic to the identity.
It exchanges the two small discs by rigid rotation and therefore
commutes with $f$.

Set $u(1)=f$. Every homeomorphism commuting with $f$ preserves its fixed
set and sends isolated fixed points to isolated fixed points. Along a
path in the centralizer the induced permutation of $\{q_-,q_+\}$ is
constant. The identity fixes both and $h$ exchanges them. They therefore
belong to different components of $(A^0)^u$. This proves the level-$1$
failure by the exact test C.3. The action $u$ and the trivial $\Z$-action
both project to the trivial homomorphism into $K$, but a conjugate of
$f\ne\id$ cannot be the identity. This proves the level-$2$ failure.
The formulas are checked on an explicit grid in the supplied program;
the isolation and irrationality arguments prove the assertions beyond
that grid.
\end{proof}

\begin{claim}{C.12}{Open}{The requested infinite-group zero-level example}
For the admissible surfaces and the exact comparison of C.1--C.3,
we have not produced an infinite operator group whose level-$0$
comparison has a nontrivial kernel, nor proved universal injectivity.
For a hyperbolic surface the required kernel is
$\pi_1(C_{A^0}(u(\Gamma)),1)$, since the target is zero.
The disconnected centralizer in C.11 proves level-$1$ failure.
The enlarged sphere comparison in C.25 has its stated nontrivial
kernel, but the unpunctured sphere is excluded from C.1's object class.
The requested admissible infinite-group zero-level decision remains
unsettled here.
C.73 proves the actual prime-end comparison for arbitrary simply connected frontiers, including its dense image and the explicit proper-image and inverse-topology failures for a non-locally-connected frontier. C.74 gives a transitive actual disk loop with full proper local centralizer $\mathbb R$ and proves that the specified change of principal section makes its coordinate interpolation fail to extend at every original boundary point; the original loop contracts. C.75 classifies every actual orbit strip for a one-dimensional proper operator action and computes its finite-endpoint translation characters, with the source stabilizer and all component operators retained. C.76--C.79 construct the controlled lift across all interior frontier lines and the original circle of the actual closed disk. In particular every based loop in the group of homeomorphisms of that disk fixing its boundary pointwise and commuting with the given closed locally compact free proper one-dimensional interior operator group has a based filling in that same relative centralizer. The proof retains a shared increasing base map, both independent endpoint characters, all component operators and the actual inverse topology. C.80 proves the same exact relative filling for every based sphere family and hence the vanishing of all positive based homotopy groups of this specified disk centralizer. The general comparison (C.12.1), arbitrary-frontier descent and nonfixed-frontier pasting remain unresolved.
C.81 supplies the exact invariant metric and metric covering for a discrete proper operator group with an actual transitive boundary-fixed commuting subgroup. Compact families with a fixed compact quotient support extend jointly to the original circle, every based sphere family supported over a compact quotient subsurface fills, and fixing any finite configuration of operator orbits preserves every positive based homotopy group. Its explicit escaping-disk sequence proves discontinuity of the actual descent map into the Whitney topology. C.82 constructs a quotient triangulation whose lifted closed stars form a null family in the original disk. Restriction and equivariant pasting identify the pointwise full-skeleton stabilizer with a product of boundary-relative triangle groups, and the coordinatewise Alexander deformation contracts it in the original topology. A controlled deformation of an arbitrary original family into that stabilizer, and the general frontier and descent comparisons, remain to be constructed.
C.83 gives an explicit null-homotopic loop in the actual infinite evaluation image, with all operator labels and a null lifted triangulation retained, which has no continuous lift to the original closed-disk group. This records the obstruction to that specified infinite-evaluation route; the general C.12 comparison remains unresolved.
C.84 identifies the exact radial discrete-translation centralizer with its weighted sublinear-displacement topology and contracts the entire specified vertically monotone subspace. C.85 constructs a continuous cutoff of arbitrary actual families to common compact quotient support and proves vanishing of all positive homotopy groups for the radial model. C.86 gives inverse continuous group maps to the entire logarithmic-spiral relative centralizer of C.72 and transfers those sphere fillings, with its original boundary topology retained. C.87 uses the exact zero-twist annulus lift, a continuous support height and explicit original-disk normalizers to contract each entire radial and spiral relative centralizer while fixing the identity. C.88 identifies the original general discrete-operator centralizer exactly with interior equivariant maps whose actual invariant-metric displacement vanishes at infinity on the quotient. Restriction and identity extension are inverse isometric group maps; the compact-parameter criterion retains joint interior continuity and a uniform quotient tail. The controlled deformation of arbitrary original families to common compact quotient support remains unresolved.
C.89 characterizes the exact restriction image on the compact lifted skeleton and original circle. Three-point-normalized Jordan-domain maps give a global continuous extension section; its displayed product with the C.82 stabilizer gives a strong deformation retraction onto that section image. C.90 proves the vanishing of every positive based homotopy group for every compact quotient in the actual C.81 setting, including nonorientable quotients, by lifting the pointed compact-surface contraction and applying the original-boundary criterion. C.91 strengthens this to a continuous identity-fixed contraction of the entire original centralizer for each such compact quotient. Its literal descent and unique point-fixed lift identify the whole point stabilizer with the pointed identity component downstairs; an explicit global evaluation section supplies the product and its inverse. The noncompact quotient tail, the deformation of the full skeleton embedding image for general quotients, and the general frontier comparisons remain unresolved.

C.92 proves that every element of the entire original point stabilizer descends into the pointed identity path component for an arbitrary actual quotient. A locally finite based homotopy and the precise orientation comparison meet Epstein's based isotopy theorem. The unique point-fixed lift gives the exact subgroup whose actual lift displacement belongs to $C_0(N)$, with an isometric comparison to the original topology and the exact compact-family uniform-tail criterion. An explicit radial-disk sequence converges downstairs in the compact-open topology while its unique pointed lifts stay more than $2/5$ from the identity. A noncompact deformation with that uniform tail and the general frontier comparisons remain unresolved.
\end{claim}
\passage{\S{}23; folios 140--141; scans 279--281}

The attempted construction using a supported radial rotation produces
isolated fixed points and a component obstruction, not a nonzero loop.
Adding finitely many marked fixed points still gives aspherical
configuration spaces and cannot by that argument produce the required
$\pi_2$ of an evaluation base. No implication from the examples for
levels $1$ and $2$ to level $0$ is asserted. The independent infinite-group
non-realization theorem C.23 gives level-$3$ failure and does not close
this missing zero-level test. In particular, the possibility that level
$0$ holds for all operator groups has not been ruled out by this return.

\paragraph{Scope of the remaining infinite-group obligation.}
The source fixed in C.1 is the admissible-surface category of scan~264,
folio~133: every genus-zero compact component meets the finite puncture
set. The unpunctured sphere is therefore absent. The last sentence of
Part~C(b) asks for an infinite-group example at each fidelity level
within this specified comparison. An example on one of its admitted
abelian types would meet that existential request; a hyperbolic example
would settle the more specific question in C.12. The examples C.11 and
C.23 supply the independent obstructions at levels~$1$, $2$, and~$3$.
The level-$0$ example inside the admitted source remains unprovided.

Here is the exact kernel that is required. For an arbitrary discrete
operator group and a strict action $u:\Gamma\to A$, put
$H=u(\Gamma)$ and let $N_u$ be its weak automorphism groupoid in
$\Pi_1A$. At its identity object, a $2$-automorphism is precisely a
loop class in $A$ invariant under conjugation by every $u(\gamma)$.
The comparison is thus the inclusion-induced map
\[
 \pi_1(C_A(H),1)\longrightarrow
 \pi_1(N_u,1)=\pi_1(A,1)^\Gamma.                 \tag{C.12.1}
\]
Pointwise commutation proves that the image lies in these invariants.
Every based loop and every based homotopy in the source lies in $A^0$,
so its source can equally be written $\pi_1(C_{A^0}(H),1)$.
For hyperbolic types the target vanishes, giving exactly C.12's test.
For the plane, annulus, and torus the target is the invariant subgroup
of $\Z$, $\Z$, or $\Z^2$, respectively; the required obstruction is a
nonzero kernel of (C.12.1).

C.39 proves that every class in this kernel represented by a loop
whose values commute pairwise is zero. Its contraction is global on
the plane, including arbitrary changes and accumulation of orbit types,
and its explicitly lifted filling disk descends to every connected
admissible surface. Thus a remaining example must involve a class
without a commuting-valued representative. C.40 constructs such a representative for every loop whose generated
closure, in the actual compact-open topology, is locally compact.
No representative reduction for arbitrary loops is proved here. C.41 proves the exact
inclusion injective for every operator group with a dense orbit,
without a local compactness assumption on either group. C.36--C.37 separately
provide the stated contraction on the closed union of dimension-two
proper-image open orbits for general loops. C.42 identifies the actual
boundary quotient as an open map and proves that every component of
every real lifted fibre is a properly embedded line; it also supplies
an exact example with disconnected whole circle fibres. C.43 proves
that the real lifted fibre is connected and supplies actual product
coordinates. C.44 gives the five downstairs covering forms, with
source a plane or annulus. C.45 gives the exact relation between two
quotients and its three planar coordinate domains. C.46 constructs a
simultaneous intrinsic contraction when the full set of lifted
boundary fibre partitions has exactly two members, retaining every
operator and quotient relation. C.47 proves that three independent
actual quotients force the full compatible group on the source plane
to be locally compact and Lie, with its actual topology. C.48 gives
the simultaneous intrinsic contraction for the full one-foliation
family, including every component of the proper operator group and
the dense lattice case. C.49 identifies the actual closed lifted
generated group in the three-quotient case and contracts each whole
orbit closure, with exact overlap agreement and finite pasting.
C.50 supplies an exact uniform word bound and extends these
contractions through arbitrary accumulation, giving a joint
contraction on $X_*=\overline{U_2\cup U_3}$. C.51 extends
the intrinsic contractions through the entire frontier when
$O/Q_O$ is compact and the full boundary-quotient family has
at most two members, with the exact median boundary value
and finite pasting. For an empty family every frontier point
is then fixed. C.52 supplies exact agreement with the
prescribed exponential values on $X_*$ by a natural hybrid
target rule and diagonal frontier continuity. C.53 enlarges
that region to $X^\dagger$ using all global nonopen quotient
maps, and controls every accumulation of all $H$-translates
of finitely many remaining cocompact representatives.
C.54 proves actual exponential closure contractions whenever
there are two metabelian global quotient images, including
all diagonal affine weights and the product affine group.
C.55 replaces the proper operator image by the full local
centralizer, with direct equivariant-map limits and the
larger prescribed region $X^\ddagger$. C.56 gives a uniform
word bound for all diagonal one-parameter affine groups and
extends their fractions through arbitrary accumulation on
$X^\sharp$, retaining the full-local relative extension.
C.57 determines the exact full local centralizers and
quotient spaces for every full two-partition family.
C.58 proves a uniform word bound for the full product affine group by an exact antisymmetric commutator correction, and extends every C.54 source through arbitrary accumulation on the larger closed $X^\flat$, with the same relative finite-type extension. C.59 extends localized quotient-preserving maps over all operator translates, with a continuous restriction inverse, and proves that a based two-partition quotient-kernel path is trivial. C.60 contracts every based sphere family in the relative one-partition quotient kernel through arbitrary accumulation, using exact forward and inverse prefix-track bounds; it requires no compact quotient or separate ambient-null assumption. C.61 excludes open-orbit accumulation for equal-level real-fiber intervals. C.62 places every such interval in an exact compact three-word family, excludes ambient escape, and confines all limiting points to the actual orbit of the source limit; at a fixed limit point the intervals shrink uniformly. C.63 constructs a relative continuous lift of the moving principal base by a compact family of small original-path increments and interlaced convex selection. It contracts the localized loop on every invariant union of open orbits with one-dimensional full local centralizer and G-fixed frontier, including arbitrary accumulation and the empty global nonopen-target family. C.64 constructs one compact labelled family controlling every selected open orbit and embeds the full product of locally finite disjoint operator disk groups across G-fixed frontiers, with an explicit Alexander contraction. C.65 proves that finite-word paths fixing the entire principal base act trivially on every nonmetabelian nonopen frontier orbit. The full principal kernel has a central discrete cyclic image there; its actual graph-image quotient is a covering, followed by a continuous bijection onto the principal-base image whose inverse continuity is unproved. C.65 also proves that individually shrinking null evaluation loops have fillings inside the original orbit and small ambient disk. C.66 proves the global filling for every original loop whose nonopen plane locus is zero-dimensional, with all operator dimensions included. Its exact end-extension isomorphism retains the original frontier, infinity, and every operator and deck label. C.67 identifies the full group fixing the complement of a bounded simply connected locally connected-frontier domain with the boundary-fixed disk group through the actual prime-end quotient. It proves contractibility of that full relative group and exactly identifies all original operator centralizers. C.68 extends the exact comparison to spherical and unbounded plane domains and proves relative centralizer fillings for the specified closed-disk first-kind Fuchsian case. C.69 proves the same relative filling for every specified closed-disk conjugate of a torsion-free Fuchsian action of either kind, without finite-generation assumptions; its partial compactification, boundary correction and compact-parameter extension are explicit. C.70 gives exactly operator-equivariant individual evaluation fillings in compact-family carriers, uniformly shrinking at every fixed frontier point. C.71 proves that every based sphere family in the actual boundary-fixed centralizer of a two-dimensional proper operator group contracts, using a single compact three-word set to retain the original boundary topology. It also computes a cyclic rectangular compactification whose relative centralizer contracts despite failure of the proposed partial-boundary restriction. C.72 gives a spiral compactification with an actual transitive loop-generated group and its full cyclic local centralizer: an explicit interior centralizer map fails to extend at every boundary point, while the same map extends after the displayed interior Fuchsian conjugacy. The original loop in that example has its displayed contraction. Joint continuity in the evaluated point, simultaneous configuration disjointness and arbitrary infinite-domain pasting remain unproved. The remaining requirements include selecting prescribed values on nonfixed limit orbits and lifting nontrivial boundary deformations, untreated noncocompact quotients,
arbitrarily many $H$-inequivalent remaining components, and
extension through nonopen regions outside these closed sets.

C.25 proves such a kernel for the separately enlarged
$\Pi_1\operatorname{Homeo}(S^2)$ comparison. Its relationship with the
admitted punctured comparisons can be computed explicitly. Write
$p_\pm$ for the poles of its irrational rotation. In the identity
component its centralizer $C$ fixes both poles, and it is the same
subgroup in the unmarked, once-marked, and twice-marked sphere
homeomorphism groups. Forgetting marks induces
\[
\begin{gathered}
 \pi_1C=\Z
 \xrightarrow{\;1\;}\pi_1\operatorname{Homeo}(S^2,\{p_+,p_-\})=\Z
 \xrightarrow{\;1\;}\pi_1\operatorname{Homeo}(S^2,\{p_+\})=\Z,\\
 \pi_1\operatorname{Homeo}(S^2,\{p_+\})=\Z
 \xrightarrow{\bmod 2}\pi_1\operatorname{Homeo}(S^2)=\Z/2.
\end{gathered}                                           \tag{C.12.2}
\]
The same full axis-rotation loop generates the first three groups by
the plane and annulus comparisons of C.10; its inclusion into $SO(3)$
gives the last map. Conjugation by the irrational rotation is based
homotopic to the identity in each ambient group, so these are also the
weak-transporter target groups in (C.12.1). Thus puncturing one or both
poles removes C.25's kernel. Every nonpolar point has an infinite orbit,
so this action admits no other finite invariant puncture set.

Keeping the peripheral-diagram recipe while adding the sphere gives
yet another explicitly related target: its diagram is
$\varnothing\to *$, whose equivalence $2$-group is terminal. The
comparison factors through the functor
$\Pi_1\operatorname{Homeo}(S^2)\to\operatorname{Eq}(\varnothing\to*)$,
which sends the identity $2$-automorphism group $\Z/2$ to zero. For the
irrational action its fundamental-group map is the composite
$\Z\to\Z/2\to0$; C.25 computes the first arrow. Neither enlargement
changes the admitted object class fixed in C.1. Consequently C.25 is
retained with its exact stated comparison, while the admissible
infinite-group zero-level obligation remains recorded here.

\paragraph{Exact lift comparison for the remaining zero-level question.}
The compact-open topology agrees with the marked compactification
topology used in C.1. More precisely, restriction
$\operatorname{Homeo}(X,S)\to\operatorname{Homeo}(X\setminus S)$
is a topological-group isomorphism. A homeomorphism extends uniquely
to the finite end compactification, giving the algebraic inverse.
For its continuity at $f_0$, choose disjoint connected target end
neighborhoods $W_j$ of small diameter and complementary compact
core $C$. Choose a source core $K$ with
$f_0^{-1}(C)\subset\operatorname{int}K$, and connected source tails
$V_i$ whose $f_0$-images lie in the prescribed $W_j$.
The conditions $f^{-1}(C)\subset\operatorname{int}K$ and
$f(z_i)\in W_j$ for one point $z_i\in V_i$ are open:
inversion is continuous in the compact-open homeomorphism group
of a locally compact, locally connected surface.
Connectedness forces $f(V_i)\subset W_j$, since this image avoids
$C$ and meets $W_j$. This controls the extended maps uniformly on
all tails. Compact-open control on the remaining compact core
controls them on the rest of $X$, proving continuity of the inverse.

Let $U$ be the hyperbolic surface of C.12, let $A^0=\operatorname{Homeo}(U)^0$, let $H$ be the image of the specified operator group, and let $p:\widetilde U\to U$ be the universal covering with deck group $\pi$. Write $\widetilde H$ for the group of all lifts of elements of $H$, regarded as actual homeomorphisms of $\widetilde U$. There is a continuous injective homomorphism
\[
 j:A^0\longrightarrow\operatorname{Homeo}(\widetilde U),
 \qquad p\,j(f)=f\,p,
\]
characterized by $j(f)\delta=\delta j(f)$ for every $\delta\in\pi$. It identifies the precise centralizers:
\[
 j\bigl(C_{A^0}(H)\bigr)
   =j(A^0)\cap C_{\operatorname{Homeo}(\widetilde U)}(\widetilde H).
\]
In particular every based loop in the left-hand centralizer lifts to a closed loop of homeomorphisms commuting with the entire lift group.
\begin{proof}
Lift an isotopy from $\id$ to $f$ starting at $\id_{\widetilde U}$. Conjugation on the discrete deck group is constant during the lift, so its endpoint commutes with $\pi$. Two such endpoints differ by a deck transformation centralizing $\pi$; they are equal because $Z(\pi)=1$ by Part A. Uniqueness also makes $j$ a homomorphism and an injection. For continuity, fix $\widetilde u\in\widetilde U$. The evaluation map $A^0\to U$ lifts continuously to $\widetilde U$, since $A^0$ is contractible and locally path connected by the surface homeomorphism-group theorem. This lift gives $j(f)(\widetilde u)$, and the local lifting property of a covering gives continuous dependence of the lifted homeomorphism in the compact-open topology.

Suppose $f$ centralizes $H$ and $\widetilde h$ is any lift of $h\in H$. Then $[j(f),\widetilde h]$ is a deck transformation. Its conjugation on $\pi$ is the commutator of the identity automorphism with the automorphism induced by $\widetilde h$, hence is trivial. Centre-freeness makes this deck transformation the identity. This proves one inclusion in the displayed equality; descent proves the other. Finally $j(\id)=\id$, so a based loop has the same endpoint after lifting.
\end{proof}

This comparison does not supply a contraction inside the right-hand centralizer. Choosing a hyperbolic metric and interpolating from $j(f)(x)$ to $x$ along geodesics has two unproved properties: the arbitrary lift group $\widetilde H$ need not preserve these geodesics, and an interpolated map need not be a homeomorphism. The first prevents asserting equivariance; the second prevents asserting that the interpolation lies in $j(A^0)$. Neither property follows from the displayed centralizer equality. The universal vanishing of $\pi_1(C_{A^0}(H))$, and the requested counterexample to it, therefore remain unresolved by this argument.

\paragraph{Current orbitwise reduction and the remaining compatibility problem.}
C.26 settles the globally transitive centralizer case on a closed
hyperbolic surface and also gives an explicit closed subgroup with an
ambient-null but intrinsically nontrivial loop; its latter example is
not asserted to be a centralizer. For the canonical plane lift above,
C.28 identifies the orbits of the group generated by the loop:
points, intrinsic lines, circles, and open surface orbits.
C.27 gives an order-natural contraction on each intrinsic line.
C.29 proves that every circle restriction has winding zero by the
actual plane-loop evaluation map, and contracts that restriction
in its circle centralizer.
C.31 supplies inverse and track control, and C.34 proves joint
continuity of the contractions and their inverses on the entire
closed union $Y$ of the nonopen orbits. This is a based contraction
through homeomorphisms of $Y$, commuting with every homeomorphism
that commutes with the original loop. In particular C.34 settles
the hyperbolic centralizer loop whenever its lifted generated
group has no open orbit.

The open orbits have also been treated intrinsically. C.30 gives
the exact restricted operator group, including its restriction
kernel, and proves that its closed image acts freely and properly.
C.32 handles a discrete image, and C.33 treats every locally
compact closed image supplied by C.30. In each case the restricted
loop has a based null-homotopy inside the centralizer of the full
restricted operator group. These results give contractions on
all individual open orbits and on the whole nonopen locus.
The remaining step is to choose these contractions compatibly
across the frontiers of the open components, with joint
continuity of the resulting homeomorphisms and their inverses.

C.35 proves that this step cannot require extending the existing
median prescription unchanged. Its explicit plane action has
the boundary restriction map
\[
 \R^2\longrightarrow\R^3,\qquad
                  (A,B)\longmapsto(A,B,A-B).
\]
At the displayed parameter the prescribed three edge medians
are $(1,0,0)$, which lies outside its image because
$r_1-r_2-r_3=1$. Thus that prescribed homeomorphism of $Y$
has no equivariant extension over an adjacent open orbit.
The original loop in the same example has the explicit global
equivariant contraction given in (C.35.5).
C.36 constructs compatible choices on closures of open orbits
whose closed proper operator image has dimension two, using the
actual open quotient onto their simply transitive Lie
centralizers. Scaling the logarithm of each based universal-cover
lift respects the exact common boundary action. C.37 supplies a
uniform three-word estimate and proves joint continuity, including
inverses, on
\[
 X_2=\overline{\bigcup_{\dim Q_O=2}O},
\]
including every infinite accumulation of those open orbits.
Its bounded residual equations determine the unique logarithmic
limit on line orbits. The actual null residual word isotopies
and the two-point configuration map determine the real lifted
limit on circle orbits, including its winding.

C.49--C.50 subsequently enlarge this closed set to $X_*$,
including all open orbits with at least three actual boundary
quotient foliations. C.51 proves whole-closure extension for
the remaining open orbits with compact operator quotient,
using the median frontier prescription and finite pasting.
C.52 makes these extensions exactly relative to $X_*$.
C.53 enlarges the exponential region to $X^\dagger$ and
proves continuity across all accumulations of the operator
translates of finitely many remaining cocompact representatives.
C.54--C.56 enlarge the treated region through all
three-dimensional diagonal affine source groups, with
uniform accumulation control, and replace the compact
operator quotient condition by compactness of the full
local centralizer quotient. C.57 gives the exact plane,
open-cylinder, line and open-M\"obius quotient alternatives
in the remaining noncocompact two-family cases.
The outstanding problem includes untreated noncocompact local
quotients, infinitely many $H$-inequivalent remaining
representatives, and extension through exterior regions
consisting of nonopen orbits. The hybrid target prescription
has been proved continuous on the treated closed sets;
its transverse continuity on all of the nonopen locus has
not been proved. C.35 requires allowing changes to the
original median prescription. These exact remaining
requirements, and the general kernel in (C.12.1), remain
unresolved here.


\subsection{Configuration invariants and drilling}

\begin{claim}{C.13}{Proved}{The invariant braid group and nonabelian cohomology}
Let a finite group act faithfully and orientation preservingly on a
hyperbolic $U$, and let $R=U^!$ be its finite set of points with nontrivial
stabilizer. Put $V=U\setminus R$,
$B=\{f\in A^0:f|_R=\id\}$, and $B^0=\operatorname{Homeo}(V)^0$.
With the action by conjugation, we have
\[
 (B/B^0)^\Gamma=1,\qquad H^1(\Gamma,B^0)=1,
 \qquad (A^0)^\Gamma=B^\Gamma\cong
 \operatorname{Homeo}(V/\Gamma)^0.
\]
Here the quotient group identifies canonically with
$\pi_1(\operatorname{Conf}_R(U),R)$ by evaluation, and its invariant
subgroup is the $H^0$ asked for on scan 285.
\end{claim}
\passage{\S{}24; folios 142--144; scans 282--286}

\begin{proof}
Evaluation $A^0\to\operatorname{Conf}_R(U)$ is a locally trivial
principal $B$-bundle, by the local extension theorem for embeddings of
a finite set. Its total space is contractible. Its homotopy exact
sequence gives $\pi_0B\cong\pi_1\operatorname{Conf}_R(U)$, with the
specified base configuration, and its identity component is the
identity component after retaining the new punctures. The construction
is equivariant for $\Gamma$, so this is the stated action on the braid
group.

An invariant component $b\in B/B^0$ is a mapping class of $V$ in the
kernel of filling $R$ that centralizes the action of $\Gamma$ on $V$.
By C.6 and its transporter version, choose a representative $f$ on $V$
commuting with $\Gamma$. It extends across the punctures and induces the
identity mapping class on $U$. Thus $f\in(A^0)^\Gamma$, and C.7 joins
it to the identity by an equivariant isotopy on $U$. Such an isotopy fixes
each point of $R$: it preserves each finite critical-point set and
starts with the identity permutation. Restriction to $V$ proves $b=1$.

For a nonabelian $1$-cocycle $a_\gamma\in B^0$, the maps
$v(\gamma)=a_\gamma u(\gamma)$ form another finite action on $V$.
The cocycle equation is exactly the homomorphism equation for $v$.
Both actions induce the same homomorphism into the mapping class group
of $V$. By C.6 they are conjugate by an element $b\in B^0$; equivalently
$a_\gamma=b u(\gamma)b^{-1}u(\gamma)^{-1}$. Hence the cocycle is a
coboundary and $H^1(\Gamma,B^0)=1$.

Finally C.7 gives connectedness. Descent to the quotient with its branch
points retained is precisely descent to $V/\Gamma$ with its punctures
retained; the identity lift construction in C.7 is the displayed
isomorphism. In particular the unprimed $V$ in the transcription of the
identity-component line on scan 286 must be compared with the quotient
$V'$ of scans 282--285. The two homeomorphism groups need not be
isomorphic; the descent map identifies the quotient one.
\end{proof}

\begin{claim}{C.14}{Refuted}{Finite critical configurations do not include mirrors}
The finite-set configuration formulation of C.13 does not extend by
replacing an orientation-preserving action by an arbitrary action and
keeping $R=U^!$ equal to all points with nontrivial stabilizer. A reflection
has a one-dimensional fixed locus, and deleting finitely many points
cannot make that action free.
\end{claim}
\passage{\S{}24; folios 142--148; scans 282--294}

\begin{proof}
Take complex conjugation on
$U=\mathbf P^1(\C)\setminus\{0,1,\infty\}$. It acts faithfully on this
hyperbolic surface. Its fixed set is
$\mathbf P^1(\R)\setminus\{0,1,\infty\}$, a union of three open arcs.
Every point of these arcs has nontrivial stabilizer. A finite deletion
leaves points in each arc. At any remaining fixed point, the lift of
conjugation fixing a selected universal-cover point is an element of
order two in the lift extension. Therefore the drilled extension still
has torsion. This explicit counterexample refutes both the finiteness
premise and a torsion-free finite-drilling assertion for such an action.
The corrected geometry retains mirror boundary in the quotient or cuts
along the mirror locus; it is not the operation of adding finitely many
puncture loop classes. The equations C.13 were stated on scans 282--286
after the explicit restriction $\Gamma=\Gamma^+$.
\end{proof}

\begin{claim}{C.15}{Proved}{The oriented drilling equivalence and the bijection}
Fix a finite group $\Gamma$ and faithful orientation-preserving outer
actions on discrete hyperbolic groups with loops. Let $\Phi(\pi,\Gamma)$
be the set of $\pi$-conjugacy classes of nontrivial maximal finite
subgroups of the canonical extension. Filling a specified
$\Gamma$-stable set $J$ of punctures with nontrivial stabilizers induces
an equivalence between the category of such drilled actions and the
category of filled actions with a specified subset
$J\subset\Phi(\pi,\Gamma)$. In particular
\[
 \Phi(\pi',\Gamma)\amalg J\xrightarrow{\ \sim\ }\Phi(\pi,\Gamma)
\]
is a canonical $\Gamma$-equivariant bijection. Drilling all of $\Phi$
produces a torsion-free extension, uniquely up to the unique exterior
isomorphism over the specified filled object.
\end{claim}
\passage{\S{}24; folios 144--147; scans 286--293}

\begin{proof}
Realize the action by C.5. A finite nontrivial subgroup of the lift group
acts on $\mathbb H^2$ by orientation-preserving isometries and has a
unique fixed point. Its stabilizer at that point is the unique maximal
finite subgroup containing it. Conversely every point of $U$ with
nontrivial stabilizer gives that subgroup after choosing a lift.
Changing the lift conjugates by $\pi$. This identifies $\Phi$ with the
finite critical-point set of the realized action, including its
$\Gamma$ action. It also proves the assertion about maximality, rather
than assuming it for arbitrary finite subgroups containing reflections.

Deleting the points of $J$ gives the drilled surface. Remaining critical
points and deleted points partition the original critical-point set,
which proves the displayed bijection. The group map is the epimorphism
$\pi'\to\pi$ obtained by killing the normal closure of exactly the new
peripheral subgroups. Local punctured discs show that its induced map
on lift groups sends a cusp stabilizer to the corresponding finite
point stabilizer. Thus the geometric bijection is the algebraic map
of scan 290, with the same conjugacy ambiguities.

To prove fullness, realize an equivariant outer isomorphism of the
filled peripheral groups by an equivariant homeomorphism, using C.6's
transporter conclusion. It preserves the selected subset of $\Phi$ by
the preceding construction. Its restriction to the complements gives
the desired drilled outer isomorphism. To prove faithfulness, represent
a drilled equivariant outer automorphism in the kernel by an equivariant
homeomorphism of the drilled surface and extend it across the new
punctures. It is isotopic to the identity on the filled surface. By C.7
there is an equivariant identity isotopy there. Every point of $J$ stays
fixed throughout, since it is a member of the finite critical set, so
the isotopy restricts to the drilled surface. The kernel is therefore
trivial. These two arguments also prove uniqueness over a fixed filled
isomorphism, not merely uniqueness of isomorphism classes.

After all critical points are removed, the finite action is free. Its
lift extension acts freely on $\mathbb H^2$ and is torsion-free, since
a nontrivial finite-order isometry fixes a point. The operations of
deleting and filling are inverse by the just-proved full faithfulness.
The only use of geometric realization in this proof is C.5 together
with the equivariant homeomorphism and isotopy results C.6--C.7; the
comparison map on groups and the kernel argument have been specified.
\end{proof}

\begin{claim}{C.16}{Refuted}{Nonfaithful outer actions obstruct torsion-free drilling}
The unrestricted existence claim on scans 286--288 fails for the
trivial outer action of $C_2$ on the group of type $(0,3)$. The faithful,
orientation-preserving version is C.15.
\end{claim}
\passage{\S{}24; folios 144--145; scans 286--288}

\begin{proof}
Suppose a finite-puncture drilling of that action had torsion-free
extension. Realize its finite outer action on the drilled hyperbolic
surface by C.5. Filling the specified punctures extends these
homeomorphisms to a finite action on the filled three-punctured sphere.
The induced mapping class of every element is trivial. The finite-order
rigidity argument of C.5 makes every filled homeomorphism the identity.
Its restriction to the dense drilled surface is therefore the identity.
The upstairs outer action is trivial as well. For a centre-free kernel
its canonical extension is $\pi'\times C_2$, which contains
$\{1\}\times C_2$. This contradicts torsion-freeness. The argument proves
the same obstruction for every nonfaithful outer action after restricting
to a nontrivial element of its finite kernel.
\end{proof}

\begin{claim}{C.17}{Proved}{Torsion-free finite extensions with peripheral data}
Let $\pi$ be a discrete hyperbolic surface group, possibly punctured,
and let $1\to\pi\to E\to\Gamma\to1$ be a torsion-free extension by a
finite group whose outer action preserves the loop structure. Then $E$
is the fundamental group of a finite-type surface, possibly
nonorientable. Its peripheral subgroups are the centralizers in $E$ of
the peripheral subgroups of $\pi$. With their conjugacy classes it is a
surface group with loops in this extended sense. The associated compact
surface and boundary groups form a Poincar\'e duality pair of dimension
two with its orientation local system.
\end{claim}
\passage{\S{}24; folios 147--148; scans 292--294; \S{}28; folio 178; scan 353}

\begin{proof}
The centre-free extension is the pullback of
$\Aut(\pi)\to\Out(\pi)$. The centralizer of $\pi$ in $E$ maps
isomorphically to the kernel of the outer action: a representative
inducing an inner automorphism can be multiplied by the inverse inner
element to centralize $\pi$, and this representative is unique because
$Z(\pi)=1$. This centralizer is finite, so torsion-freeness forces it
to be trivial. Hence the outer action is faithful.

Apply C.5 and identify $E$ with the group of lifts of the realized
action, by uniqueness of the centre-free extension. A nontrivial
stabilizer of a point gives a finite subgroup of $E$, so the action on
$U$ is free. Its quotient $W=U/\Gamma$ is an ordinary surface and
$\mathbb H^2\to W$ is its universal covering. This gives the required
isomorphism $E\cong\pi_1(W)$, respecting the specified subgroup
$\pi$ and quotient $\Gamma$.

A stabilizer of an end of $U$ cannot contain an orientation-reversing
operator. The local finite-group model at its compactifying point would
then contain a reflection and fixed points arbitrarily close to that
point, contrary to freeness on $U$. Its local action is therefore cyclic
rotation. Consequently an upstairs peripheral subgroup is a finite-index
subgroup of the downstairs peripheral infinite cyclic group. In the
hyperbolic universal cover its centralizer is exactly that full cusp
group. These are all the cusps, since every end downstairs has a
preimage upstairs. Thus the stated peripheral classes, including their
full subgroups rather than merely inclusions, are obtained.

Choose disjoint cusp neighbourhoods in $W$ and truncate them to obtain
a compact surface $M$ with boundary. The inclusions of boundary circles
induce precisely these cyclic groups. Lift a triangulation to the
contractible universal cover. The cellular cap-product map with the
relative orientation class gives the chain equivalence between
cochains and complementary relative chains; the local sign is the
orientation character of $W$. This is the Poincar\'e--Lefschetz duality
pair, and its group-pair chain model is the relative cellular complex
because both the universal cover and the lifted boundary components
are contractible. This proves the duality-pair assertion rather than
inferring it from virtual freeness alone. The exact geometric input is
the finite realization theorem; a separate proof from abstract duality
axioms is not substituted for this argument.
\end{proof}

\begin{claim}{C.18}{Proved}{The discrete oriented realization criterion}
For a faithful orientation-preserving finite outer action on a discrete
hyperbolic group with loops, the criterion of scan 294 is valid with
the word admissible meaning the full-faithfulness and uniqueness
conditions of scan 293. The drilling in question is at the canonical
set $\Phi(\pi,\Gamma)$, and its extension carries the quotient loop
structure described in C.17.
\end{claim}
\passage{\S{}24; folios 147--148; scans 293--294}

\begin{proof}
A realized action can be drilled at its critical points. C.15 proves
that the resulting action is admissible and its extension torsion-free;
C.17 identifies that extension with a group with loops. This verifies
every requirement of the stated criterion.

Conversely let $\pi'\to\pi$ be the stipulated filling map and let
$E'$ be its extension, with the group-with-loops structure required in
the criterion. Realize $E'$ by a compact surface with its boundary
circles, then take its covering corresponding to $\pi'$. The comparison
of peripheral subgroups in the definition ensures that this cover has
exactly the prescribed upstairs loop classes. Its deck group is
$\Gamma$, acting freely before filling. The cyclic end stabilizers
extend across each new puncture by the local map $z\mapsto z^m$.
Filling the selected ends thus gives a $\Gamma$-action on the surface
whose peripheral group is $\pi$. The canonical map of lift extensions
identifies its outer action with the original one. Admissibility makes
this construction independent, up to the uniquely specified exterior
isomorphism, of the chosen drilled object. This proves the converse
by the typed covering and filling maps. The mirror counterexample
C.14 explains the necessary restriction of this formulation.
\end{proof}

\subsection{Special homeomorphisms: the topology and the extension class}

\begin{claim}{C.19}{Refuted}{Escaping support invalidates the compact-open version}
Let $SA\subset A$ consist of homeomorphisms equal to the identity on
some neighbourhood of every puncture, and give it the subspace
compact-open topology. The canonical special-group comparison of
scan 304 does not identify $\pi_0(SA)$ with the boundary-framed mapping
class group. A peripheral boundary twist is connected to the identity
inside this $SA$. In particular, the annulus conclusion
$\pi_0(SA)=\Z$ on scan 304 is false in this topology.
\end{claim}
\passage{\S{}25; folios 151--153; scans 301--304}

\begin{proof}
Use a coordinate $z=re^{i\theta}$ on a disc about a puncture at $0$.
Choose a continuous function $\eta:[0,\infty)\to[0,1]$ equal to $0$ on
$[0,1/3]$ and to $1$ on $[2/3,\infty)$. For $0<s\leq1$ define
\[
 h_s(re^{i\theta})=r\exp\bigl(i\theta+2\pi i\eta(r/s)\bigr),
 \qquad h_0=\id,
\]
and extend by the identity outside the coordinate disc. Each $h_s$ is
identity near the puncture and outside the annulus
$s/3<r<2s/3$. It is a boundary twist when viewed on a fixed compact
surface obtained by truncating the puncture. Nevertheless
$\sup_z|h_s(z)-z|\leq4s/3$, and the same bound holds for its inverse.
Thus $s\mapsto h_s$ is a continuous path in the stated subspace topology
from the identity to that twist.

For the annulus every compactly supported orientation-preserving
homeomorphism fixes both ends and is represented on a compact annulus
by an integer twist; cutting along a spanning arc and applying the
Alexander isotopy of a disc proves this integer classification. The
isotopy to the corresponding integer twist can be chosen in that fixed
compact annulus. The path above then kills that twist in $SA$. Hence
$SA$ is path connected in the compact-open topology, proving the
explicit contradiction. The programme checks the formula and its
uniform bound; the bound itself proves convergence for the full space,
not just on a finite sample.

The fixed-boundary formulation below retains a single compact surface
throughout the isotopy. That is the required correction. No change of
topology is silently made in this refutation.
\end{proof}

\begin{claim}{C.20}{Proved}{The fixed-boundary repair and its class}
Let $M$ be an oriented compact surface obtained by replacing the
$\nu>0$ punctures of a hyperbolic type by boundary circles. Put
$A_\partial=\operatorname{Homeo}^+(M\ {\rm rel}\ \partial M)$.
Then $A_\partial^0$ is contractible,
$\pi_0(A_\partial)\cong S\Gamma^{!+}_{g,\nu}$ for the framed group
of E.1--E.2, and there is a central extension
\[
 1\longrightarrow\Z^\nu\longrightarrow\pi_0(A_\partial)
 \longrightarrow\Gamma^{!+}_{g,\nu}\longrightarrow1.
\]
The $i$th kernel generator is the positive boundary twist. Its
extension class is the vector of Euler classes of the oriented tangent
circle lines at the marked points, with the sign fixed by that generator.
For the annulus the kernel is the quotient of $\Z^2$ by the diagonal
rotation relation; for the once-punctured sphere it is zero.
\end{claim}
\passage{\S{}25; folios 152--154; scans 302--306}

\begin{proof}
Allow boundary rotations first, retaining each boundary component
setwise. Restriction to the boundary is a locally trivial fibration
onto $\prod_i\operatorname{Homeo}^+(S^1)$: extend a small boundary
isotopy across collars. Its fibre is $A_\partial$.
Each boundary homeomorphism group deformation retracts onto rotations
by lifting to increasing maps of $\R$ and linearly interpolating the
periodic displacement. Thus the base has homotopy type $(S^1)^\nu$.
Capping each boundary by a disc with its centre marked identifies the
free-boundary mapping classes with the pure oriented punctured mapping
classes. The same collar-extension construction gives the homotopy
comparison on identity components.

The identity component of the free-boundary group is contractible in
a hyperbolic type by Yagasaki's compact-pair theorem. The homotopy
sequence therefore injects $\pi_1((S^1)^\nu)=\Z^\nu$ into the
components of the fibre and identifies the quotient as stated. Lifting
one rotation in the $i$th collar gives exactly the boundary twist.
Its conjugation action is trivial because the boundary circles are
retained individually and their orientations are fixed. Choose a point
on each boundary circle. A boundary-fixing homeomorphism acts on the
framed peripheral diagram of E.2, and an isotopy relative to the
boundary gives its prescribed natural isomorphism. This defines
$\pi_0(A_\partial)\to S\Gamma^{!+}_{g,\nu}$. On the quotient it is
the peripheral Dehn--Nielsen--Baer isomorphism. A positive boundary
twist changes the path from the generic basepoint to that boundary
point by one peripheral turn; hence the induced map on kernels sends
the $i$th boundary generator to the $i$th generator of $T^I$.
It is an isomorphism on both kernel and quotient, so the two short
exact sequences give the asserted isomorphism of extensions.
Contractibility of $A_\partial^0$ follows directly from Lemma 2.1(ii),
with the nonempty fixed boundary. The connecting homomorphism is consequently the
fundamental-group extension of the bundle of tangent-circle framings
at the marked points. Obstruction to a section of that circle bundle
is its Euler class, proving the class identification with its sign
convention. In complex notation this is the tangent-line first Chern
class; the usual cotangent $\psi_i$ has its opposite sign.

For the annulus the identity component of the free-boundary group is
$S^1$. Its fundamental-group map to the two boundary-rotation groups
is the diagonal when both circles use the same angular coordinate,
or the antidiagonal when both use boundary orientations. The cokernel
is $\Z$. For the disc this map is an isomorphism $\Z\to\Z$, so the
cokernel is zero. These are the exact boundary models underlying the
two abelian computations on scans 303--304.
\end{proof}

\begin{claim}{C.21}{Proved}{The four small extension classes}
With the positive boundary-twist convention of C.20, the special
fixed-boundary extension classes for $(0,3)$ and $(0,4)$ are zero. For
$(1,1)$ the class is a generator $e$ of
$H^2(\operatorname{SL}_2(\Z);\Z)=\Z/12$. For $(1,2)$ the vector class
is $(q^*e,q^*e)$, where
$q:\Gamma^{!+}_{1,2}\to\operatorname{SL}_2(\Z)$ forgets the second
marked point; its order is $12$.
\end{claim}
\passage{\S{}25; folios 153--154; scans 304--306}

\begin{proof}
The pure oriented mapping class group of a three-punctured sphere is
trivial. Filling one puncture of a four-punctured sphere gives the
point-pushing group $F_2$, since the target pure group is trivial and
the three-punctured surface is hyperbolic. A graph is a classifying space
for $F_2$, so $H^2(F_2;\Z^4)=0$. This proves the first two class
computations, including the noncanonical splitting in type $(0,4)$.

The integral matrices
\[
 S=\begin{pmatrix}0&-1\\1&0\end{pmatrix},\qquad
 R=\begin{pmatrix}0&-1\\1&1\end{pmatrix}
\]
have orders $4$ and $6$ and satisfy $S^2=R^3=-I$.
Their action on the Farey tessellation gives
$\operatorname{PSL}_2(\Z)=C_2*C_3$: the dual bipartite tree has one
edge as fundamental domain, vertex stabilizers $C_2,C_3$, and trivial
edge stabilizer. Connectivity and the absence of circuits follow from
the Euclidean algorithm for adjacent reduced fractions. Lifting the
central element $-I$ gives
$\operatorname{SL}_2(\Z)=C_4*_{C_2}C_6$. The cohomology sequence of
this tree of groups identifies $H^2$ with
\[
 \{(r,s)\in\Z/4\oplus\Z/6:r\equiv s\pmod2\}:
\]
the preceding $H^1(C_2;\Z)$ and $H^1(C_4;\Z)\oplus
H^1(C_6;\Z)$ vanish, while the restrictions in degree two reduce
modulo $2$. This group has twelve elements and is cyclic.
The tangent-circle bundle restricts on the square torus fixed by $S$
to the rotation representation by $\pi/2$, and on the hexagonal torus
fixed by $R$ to the representation by $\pi/3$. Lifting these circle
rotations to $\R$ gives fourth and sixth powers equal to one positive
full turn. By the convention of C.20 the Euler class therefore
restricts to $(1,1)$, an element of order $12$. This proves its generator
assertion without an unproved presentation for the boundary mapping
class group. The supplied code verifies the matrix orders and the
finite abelian-group calculation.

In a complex torus translation canonically identifies the tangent
spaces at any two points. These identifications vary with the marked
complex torus and commute with all isomorphisms. Hence the two tangent
circle bundles on the two-pointed moduli orbifold are isomorphic, and
both are pulled back from the one-pointed tangent bundle. This proves
the vector formula. It has order dividing $12$. A square torus with
marked points $0$ and $(1+i)/2$ admits the order-four rotation preserving
both; a hexagonal torus admits an order-three rotation with at least
two marked fixed points. The restrictions of the class to these cyclic
subgroups are generators, as can be seen from their nontrivial local
rotations, or from C.22 below. Its order is divisible by $4$ and $3$,
and thus is exactly $12$. The existence of the displayed fixed points
is also verified in the integral torus matrices supplied with the code.
\end{proof}

\begin{claim}{C.22}{Proved}{Restriction over a finite puncture stabilizer}
For a hyperbolic type, a nontrivial finite subgroup
$G\subset\Gamma_i^+$ is cyclic. The extension by the $i$th copy of
$T=\Z$ obtained from the special group is nonsplit over $G$; its
preimage is infinite cyclic and the kernel is the subgroup of index
$|G|$.
\end{claim}
\passage{\S{}25; folios 153--154; scans 305--306}

\begin{proof}
Realize $G$ by C.5. It fixes the marked point $i$ and preserves
orientation. The finite local action is a faithful rotation action:
an element trivial on a small disc is the identity as a periodic
surface homeomorphism, by the local classification and connectedness.
Thus $G$ embeds in the circle rotations and is cyclic. Delete that
point and lift a small punctured quotient disc to its universal cover.
Its group is $\Z$, while the upstairs punctured disc has subgroup
$|G|\Z$. The exact normalizer interpretation proved in E.3 identifies
this with the extension pushed out by the $i$th projection. An infinite
cyclic group contains no nontrivial finite subgroup, so the extension
cannot split. This proves the remark on scan 306 with its stated
orientation and puncture-stabilizer hypotheses.
\end{proof}


\begin{claim}{C.23}{Proved}{An infinite-group failure of existence}
For every closed orientable surface $X_g$ of genus $g\geq2$, the
identity homomorphism from the discrete group
$\Gamma=\operatorname{Mod}^+(X_g)$ to its subgroup in $\pi_0A_g$
has no homomorphic lift to $A_g$. Thus level $3$ fails for an infinite
operator group already on each such closed hyperbolic surface.
\end{claim}
\passage{\S{}23; folios 140--141; scans 279--280}

\begin{proof}
Markovi\'c--\v Sari\'c's Theorem 1.1
\cite{C_MarkovicSaric2008} states that the projection
$\operatorname{Homeo}^+(X_g)\to\operatorname{Mod}^+(X_g)$ has no group
section for $g\geq2$. The theorem is about homeomorphisms, not merely
diffeomorphisms; its first-page statement and definitions have been
checked in the original paper. A lift of the displayed homomorphism
would take every element to an orientation-preserving homeomorphism,
since orientation is invariant under isotopy. Its composite with that
projection would be the identity, giving the forbidden section. This
is the exact reduction to the source theorem. The group is infinite:
a twist about a nonseparating simple curve acts on $H_1(X_g;\Z)$ by
an integral transvection, whose distinct powers are distinct.
\end{proof}

\subsection{The integral sphere comparison}

\begin{claim}{C.24}{Refuted}{The sphere comparison already fails for an order-three rotation}
Let $A=\operatorname{Homeo}(S^2)$ with its compact-open topology, and let
$\Gamma=C_3$ act on $S^2$ by rotation through $2\pi/3$ about a fixed axis.
Give $A$ the resulting conjugation action and base both fixed-point spaces
at the identity. The actual comparison has
\[
 \pi_1(A^\Gamma,1)=\Z,\qquad
 \pi_1(A^{h\Gamma},1)=\Z_3\times\Z/2,
 \qquad n\longmapsto(n,n\bmod2).
\]
Consequently the requested integral homotopy equivalence fails in type
$(0,0)$. This failure concerns the actual homeomorphism mapping space;
it is independent of the earlier loss of the sphere's $\pi_2$ in its
fundamental groupoid. For every odd prime $p$, the same construction with
$C_p$ gives $\Z_p\times\Z/2$ and the same comparison formula with $p$
in place of $3$.
\end{claim}
\passage{\S{}23; folio 141; scans 280--281}

\begin{proof}
We give the comparison maps and compute their effect, including the integral
rather than only completed homotopy group. Write $G=SO(3)$ and
$\rho:C_p\hookrightarrow G$ for the rotation subgroup, with $p$ odd.
The inclusion $G\hookrightarrow A^+=\operatorname{Homeo}^+(S^2)$ is
$C_p$-equivariant and a weak equivalence. Indeed evaluation on three
ordered distinct points is a locally trivial fibration onto
$\operatorname{Conf}_3(S^2)$, by isotopy extension. Its fibre is the
pure homeomorphism group of the three-punctured sphere. Its component
group is trivial, and its identity component is contractible by the
surface homeomorphism-group theorem used in C.7. The unique M\"obius
transformation carrying one ordered triple to another identifies this
configuration space with $PSL_2(\C)$; polar decomposition retracts this
group onto $SO(3)$. The resulting evaluation diagram proves the stated
weak equivalence for the actual inclusion. Homotopy fixed points preserve
this objectwise weak equivalence, after applying singular complexes. A
path through the identity in $A^{hC_p}$ stays in $A^+$, because $EC_p$
is connected and orientation takes values in the discrete group $C_2$.
It therefore suffices to calculate $(G^{hC_p},1)$.

The identity component of the actual centralizer fixes the two poles
individually. Quotient by $z\mapsto z^p$ identifies its descent map to
$\operatorname{Homeo}^+(S^2;0,\infty)^0$ as a covering of degree $p$.
Every map of the quotient fixing both poles lifts: on the complementary
annulus its action on the meridian is $+1$, so it preserves the subgroup
$p\Z$ of the covering group. Each lift commutes with the deck group,
since the induced automorphism of $\Z/p$ is $+1$; it extends uniquely
over the two poles. The usual local lifting neighbourhoods establish
that descent is a covering in the compact-open topology. Its total space
is connected: the quotient group is connected, and all deck lifts lie
on the connected circle of axis rotations. The annulus calculation of
C.10 identifies the base with a space of homotopy type $S^1$. The
rotation circle maps to its quotient rotation circle by $z\mapsto z^p$;
its inclusion in the total centralizer therefore induces an isomorphism
on every homotopy group. In particular the actual fundamental group is
$\Z$, with generator one complete axis rotation.

Set $Y=BG=BSO(3)$. The adjoint bundle
$EC_p\times_{C_p}G\to BC_p$, with action through $\rho$, has section
space $G^{hC_p}$. It is also the bundle of gauge automorphisms of the
principal $G$-bundle classified by $B\rho$. The universal bundle gives
the natural weak equivalence
\[
 G^{hC_p}\simeq
 \Omega_{B\rho}\operatorname{Map}(BC_p,Y).
\]
One way to verify the map is to send a gauge automorphism to the loop
of its classifying map obtained by gluing the bundle over the two ends
of $BC_p\times[0,1]$. Conversely a loop of classifying maps gives a
bundle over this cylinder with its specified endpoint identifications;
trivialization in the interval direction gives the gauge automorphism.
The spaces of such trivializations are contractible. This construction
is natural for the inclusion $C_G(\rho(C_p))\hookrightarrow G^{hC_p}$.

We now compute $\pi_2$ of this mapping space using the arithmetic square
of the simply connected finite-type space $Y$. Let
\[
 \widehat Y=\prod_qY_q^{\wedge},\qquad
 V=\Bigl(\prod_q\Z_q\Bigr)\otimes\Q.
\]
Here $q$ runs over primes and superscript $\wedge$ means homotopical
$q$-completion, throughout this proof. We have
$Y_{\Q}\simeq K(\Q,4)$ and
$(\widehat Y)_{\Q}\simeq K(V,4)$. This follows from
$SO(3)\simeq\mathbb {RP}^3$: its rational homotopy is concentrated in
degree three, and each of its higher homotopy groups is finite. The
square
\[
 \begin{matrix}
 Y&\longrightarrow&\widehat Y\\
 \downarrow&&\downarrow\\
 K(\Q,4)&\longrightarrow&K(V,4)
 \end{matrix}
\]
is a homotopy pullback. For completeness its homotopy-group verification
reduces on each free summand to
\[
 0\longrightarrow\Z\longrightarrow
 \Q\oplus\prod_q\Z_q\longrightarrow V\longrightarrow0,
 \quad (a,(b_q))\longmapsto(b_q-a)_q.
\]
Surjectivity follows by clearing a common denominator and choosing its
finitely many required residues by the Chinese remainder theorem; the
kernel consists of rational numbers integral at every prime and is
$\Z$. On finite summands completion gives the product of their primary
parts and rationalization is zero. The long exact sequence of the
homotopy pullback and Whitehead's theorem prove the square, also giving
the usual arithmetic-square construction in the exact form needed here.

Apply $\operatorname{Map}(BC_p,-)$ to this homotopy pullback and take the
component containing $B\rho$. For every rational vector space $W$,
positive-degree cohomology $H^k(C_p;W)$ vanishes by averaging. The two
rational mapping spaces consequently have vanishing $\pi_2$ and
$\pi_3$. The homotopy-pullback exact sequence therefore gives
\[
 \pi_2\operatorname{Map}(BC_p,Y)_{B\rho}
 \cong\prod_q
 \pi_2\operatorname{Map}(BC_p,Y_q^{\wedge})_{B\rho}.
\]
For $q\ne p$, evaluation is a weak equivalence from this mapping-space
component to $Y_q^{\wedge}$. To check it, use the Postnikov tower:
positive-degree $C_p$-cohomology of every $q$-complete homotopy group
vanishes since multiplication by $p$ is invertible; at each tower stage
only constant maps contribute, and homotopy inverse limits preserve
the comparison. Thus only $q=2$ contributes to $\pi_2$ among these
primes, giving $\pi_2(Y_2^{\wedge})=\Z/2$.

At $q=p$ we use the exact finite-$p$-group mapping-space theorem:
for a finite $p$-group $P$, a compact Lie group $K$, and
$\alpha:P\to K$, the map induced by the centralizer identifies
\[
 \pi_n\operatorname{Map}(BP,(BK)_p^{\wedge})_{B\alpha}
  =\bigl(\pi_n B C_K(\alpha(P))\bigr)_p^{\wedge}
  \quad(n\ge2).
\]
This is the Dwyer--Zabrodsky--Notbohm formula, stated and derived in
\cite[Theorem 1.9 and its proof, using Theorem 3.2(iii) of their earlier
paper]{C_JMO1995}. Its hypotheses apply to the displayed finite cyclic
group and $SO(3)$. Since a nontrivial odd-order rotation has centralizer
exactly the axis circle $SO(2)$, the contribution is $\Z_p$.
This proves $\pi_1(G^{hC_p},1)=\Z_p\times\Z/2$.

Naturality of the centralizer map makes the $p$-component of the map
from the rotation circle the completion $\Z\to\Z_p$. At the prime
$2$ it is the map $\pi_1SO(2)\to\pi_1SO(3)$, namely reduction modulo
$2$: lifting a full rotation to unit quaternions gives the endpoint
$-1$. This proves the stated formula for the exact comparison.
For an explicit missing element at $p=3$, put
$\beta=\sum_{j\ge1}3^{j!}\in\Z_3$. Its base-three digits have
infinitely many zeros and infinitely many ones, whereas an ordinary
integer has eventually only zeros or eventually only twos. Hence
$(\beta,0)$ is outside the image. This also proves nonsurjectivity
without a cardinality argument; for general odd $p$ the same factorial
series has digits zero and one, rather than eventually zero or $p-1$. The explicit action, its finite order, the axis
centralizer and the finite reductions of this comparison are checked in
\texttt{code/verify\_sphere\_comparison.py}. The inverse-limit argument
above, rather than a finite computation, proves the nonsurjectivity.
\end{proof}

\begin{claim}{C.25}{Refuted}{Zero-level fidelity for a faithful infinite cyclic sphere action}
Let $\Gamma=\Z$ act faithfully on $S^2$ by an irrational axis rotation,
$u(1)(\theta,t)=(\theta+2\pi\alpha,t)$ with
$\alpha\in\R\setminus\Q$ and the usual collapse of the latitude
circles at $t=\pm1$. Let $M=A^\Gamma$ be its actual automorphism
space, based at the identity, and let $N$ be its weak automorphism
space in the categorical group $\Pi_1A$, where
$A=\operatorname{Homeo}(S^2)$. The induced transporter map has
\[
 \pi_1(M,1)=\Z\longrightarrow\pi_1(N,1)=\Z/2,
 \qquad n\longmapsto n\bmod2.
\]
Its kernel is $2\Z$, so level $0$ fails for this faithful action of an
infinite group. This is the exact fundamental-groupoid comparison;
C.24 separately retains the higher homotopy of $A$ when computing the
untruncated homotopy-fixed space. The hyperbolic zero-level question
recorded in C.12 has its stated, stronger surface hypothesis.
\end{claim}
\passage{\S{}23; folios 140--141; scans 279--281}

\begin{proof}
Commutation with an irrational rotation implies commutation with every
axis rotation: its powers are dense in $S^1$, and composition and
inversion in the compact-open homeomorphism group are continuous.
An equivariant homeomorphism preserves the two poles as a set and
sends each other rotation orbit to another orbit. In the orientation-
preserving identity component it has exactly the form
\[
 f(\theta,t)=(\theta+a(t),b(t)),\quad
 a\in C((-1,1),\R/2\pi\Z),\quad
 b\in\operatorname{Homeo}^+([-1,1]).
\]
To prove that this parametrization is onto, equivariance forces the
angular function on each orbit to be translation by a unique $a(t)$;
the orbit quotient supplies $b$. Orientation forces $b$ to be increasing.
Conversely every such pair extends continuously at both poles, because
those circle orbits collapse there, and the same argument applies to
its inverse. Thus the parametrization is a homeomorphism of spaces
with the compact-open topologies. Continuity at a pole imposes no
limit on $a(t)$; on any compact subinterval of $(-1,1)$ it is exactly
ordinary compact-open convergence.

The homotopy $a_s(t)=a((1-s)t)$ contracts the angular-function factor
to its evaluation circle at $0$. The convex interpolation
$b_s(t)=(1-s)b(t)+st$ contracts the increasing-homeomorphism factor
to the identity: it remains strictly increasing and fixes both endpoints.
These homotopies are continuous in the stated topologies and extend
continuously to the sphere. The actual identity centralizer therefore
has homotopy type $S^1$, with its axis-rotation circle giving the
specified generator of $\pi_1M$.

In the weak transporter groupoid, an automorphism of its identity
object is a loop in $A$ fixed by the conjugation action of $\Gamma$.
This is exactly the groupoid calculation of C.9, now applied to
$\Pi_1A$ rather than asserting any higher truncation of $A$ itself.
By the evaluation and polar-decomposition comparison in C.24,
$\pi_1(A,1)=\pi_1SO(3)=\Z/2$. Conjugation by an axis rotation acts
trivially on this group: joining that rotation to the identity in
$SO(3)$ gives a based homotopy of its conjugation map to the identity.
Hence $\pi_1N=\Z/2$.

The transporter comparison sends a loop of equivariant homeomorphisms
to that same loop in $A$. Its effect is consequently the usual
$\pi_1SO(2)\to\pi_1SO(3)$, reduction modulo two. More explicitly,
a loop of two complete rotations has degree two under the evaluation
$a\mapsto a(0)$ in the actual centralizer, so it is nonzero there.
Its lift to the quaternion double cover $S^3\to SO(3)$ is the closed
loop $t\mapsto\cos(2\pi t)+\mathbf i\sin(2\pi t)$, which contracts
in the simply connected $S^3$. It is therefore killed by the weak
transporter map. This proves the exact nontrivial kernel and the
zero-level failure, with no inference from disconnectedness.
\end{proof}

\begin{claim}{C.26}{Proved}{Transitive centralizers on a closed hyperbolic surface}
Let $X$ be a closed connected orientable surface of genus $g\geq2$,
and give $A=\operatorname{Homeo}(X)$ the compact-open topology.
For every subgroup $L\leq A$ acting transitively on $X$, the centralizer
$C_A(L)$ is finite. Consequently, for an arbitrary strict action
$u:\Gamma\to A$ and a based loop $f_t$ in $C_{A^0}(u(\Gamma))$,
the loop is equivariantly null-homotopic whenever the subgroup generated
by its values acts transitively on $X$.
There also exists a closed torsion-free Polish subgroup
$G\leq\operatorname{Homeo}(D^2,\partial D^2)$ with $\pi_1G=\Z$;
extension by the identity embeds it closedly in every hyperbolic
$A^0$, where its displayed generator becomes null-homotopic. This
last subgroup is not asserted to be a centralizer.
\end{claim}
\passage{23; folios 140--141; scans 279--281; Part C(b) infinite-operator request; exact partial reduction for C.12}

\begin{proof}
Put $K=\overline L$ and $N=C_A(K)$. Continuity of composition gives
$C_A(L)=C_A(K)$, and $K$ is a closed Polish subgroup of $A$ that
still acts transitively. Fix $x\in X$ and write $K_x$ for its stabilizer.
The open mapping principle for Polish actions gives a homeomorphism
\[
 K/K_x\longrightarrow X,\qquad kK_x\longmapsto kx.
                                                        \tag{C.26.1}
\]
Precisely, we use the ordinary-action case of
\cite[Theorem~3.5]{C_GomezPinedoUzcategui2017}: the acting group is
Polish, $X$ is nonmeagre Hausdorff, the action is continuous and
transitive, and the domain $K^x$ of the partial-action notation is all
of $K$, hence open. The source proves this open mapping statement;
no local section of (C.26.1) is being assumed.

Every $n\in N$ is determined by $n(x)$, since
$n(kx)=k n(x)$. Its possible values form exactly the closed subset
\[
 N_K(K_x)/K_x\ \subseteq\ K/K_x=X.                      \tag{C.26.2}
\]
Indeed, $n(x)=kx$ forces $K_x=K_{kx}=kK_xk^{-1}$.
Conversely a normalizing $k$ defines the homeomorphism
$gK_x\mapsto gkK_x$, commuting with $K$. The normalizer of the closed
subgroup $K_x$ is closed: both inclusions in
$kK_xk^{-1}=K_x$ are intersections of closed conditions. Its image
under (C.26.1) is closed because its inverse image under the quotient
map is that same normalizer. Thus (C.26.2) is compact.

Evaluation is a homeomorphism of $N$ onto (C.26.2). Only inverse
continuity needs proof. Fix $\varepsilon>0$. For every $y=g_yx$ choose
an identity neighborhood $W_y$ in $K$ such that each
$g_ywg_y^{-1}$, $w\in W_y$, moves every point of $X$ by less than
$\varepsilon/3$. The open sets $g_yW_yx$ cover $X$; take a finite
subcover with elements $g_1,\ldots,g_m$. When $n(x)$ is sufficiently
close to $n_0(x)$, each $g_i n(x)$ is within $\varepsilon/3$ of
$g_i n_0(x)$. At $z=g_iwx$ in the corresponding set, commutation gives
\[
 n(z)=g_iw n(x),\qquad n_0(z)=g_iw n_0(x).
\]
The two uniformly small movements by $g_iwg_i^{-1}$ and the middle
$\varepsilon/3$ bound imply $d(n(z),n_0(z))<\varepsilon$. This is
uniform in $z$ and proves inverse continuity in the compact-open
topology. In particular $N$ is compact.

An element of $N$ fixing one point fixes every point, by transitivity
and commutation. Every element of $A^0$ has Lefschetz number
$\chi(X)=2-2g\ne0$, hence has a fixed point. Therefore
\[
                       N\cap A^0=1.
\]
Projection to the discrete group $\pi_0A$ is thus a continuous
injection on $N$. Its image is compact and discrete, hence finite,
as asserted.

For the stated loop, put $L=\langle f_t:0\leq t\leq1\rangle$.
Every $u(\gamma)$ commutes with $L$, so the transitivity hypothesis
and the first assertion place $u(\Gamma)$ inside a finite group.
Write $H=u(\Gamma)$. The equality
$C_{A^0}(u(\Gamma))=C_{A^0}(H)$ is equality of subspaces, irrespective
of the cardinality of $\Gamma$. C.7 applied to the finite group $H$
contracts the loop in this exact centralizer. The assertion therefore
does not replace the operator group by a different action without a
comparison map.
\end{proof}

\paragraph{A proved obstruction to a closed-subgroup-only argument.}
There is a closed torsion-free Polish subgroup of
$\operatorname{Homeo}(D^2,\partial D^2)$ with fundamental group
$\Z$. We give the embeddings, because this example explains exactly
why closedness and absence of compact circles do not by themselves
settle C.12. It is not asserted to be a centralizer.

Let
\[
 E=\{f\in L^1([0,1],\R):f\text{ is integer-valued almost everywhere}\}.
\]
An $L^1$-convergent sequence has an almost-everywhere convergent
subsequence; hence $E$ is closed and Polish. For $0<s\leq1$ define
$D_sf(t)=f(t/s)$ on $[0,s]$ and zero on $(s,1]$, and put $D_0f=0$.
These maps remain in $E$ and satisfy $\|D_sf\|_1=s\|f\|_1$.
Strong continuity of dilation follows first for step functions and
then by $L^1$ approximation. Thus $(s,f)\mapsto D_sf$ contracts $E$
and each of its norm balls at zero. In particular $E$ is contractible
and locally contractible.

The constant integer functions form a closed discrete subgroup $\Z$.
The quotient $G=E/\Z$ is Polish, and $E\to G$ is a covering: translates
of a norm ball of radius less than $1/3$ by distinct integers are
disjoint. It is the universal cover, and
\[
 \pi_1(G)=\Z,\qquad
 \gamma(s)=[\mathbf1_{[0,s]}]\quad(0\leq s\leq1)          \tag{C.26.3}
\]
is a generator, because its lift joins zero to the constant function
$1$. Nevertheless $G$ is torsion-free: $n[f]=0$ says that $nf$ is
a constant integer, so the integer-valued $f$ is constant and $[f]=0$.

There is a closed isometric group embedding
\[
 G\longrightarrow V:=L^1([0,1],\R)/\R,
       \qquad[f]_{\Z}\longmapsto[f]_{\R},                 \tag{C.26.4}
\]
where the second denominator consists of all constant real functions.
For integer-valued $h$, the convex function $c\mapsto\int|h-c|$ is
linear between consecutive integers, tends to infinity at both ends,
and attains a minimum at an integer. Hence
\[
 \inf_{c\in\R}\|h-c\|_1=\inf_{n\in\Z}\|h-n\|_1,
\]
proving the isometry. The quotient metric on $G$ is complete: choose
a Cauchy subsequence with successive quotient distances below $2^{-j}$
and adjust each representative by an integer so that the successive
$L^1$ distances have that bound. The representatives converge in $E$,
and the entire original Cauchy sequence then converges in $G$.
This proves closedness in the Banach space $V$.

We spell out a closed linear isometric embedding
$T:V\hookrightarrow C_0([0,1],\R)$, where the subscript here means
vanishing at both endpoints. The weak-star unit ball $B$ of $V^*$ is
compact metrizable, and $v\mapsto(\ell\mapsto\ell(v))$ embeds $V$
isometrically into $C(B)$. A continuous surjection from the Cantor
set to $B$ exists by nested finite covers of decreasing diameters.
Pullback followed by linear interpolation across each complementary
interval of the Cantor set embeds $C(B)$ linearly and isometrically
in $C([0,1])$. Place this interval in the middle third and interpolate
to zero at the two outer endpoints. The resulting map is the required
$T$; its image is closed because $V$ is complete.

For $F\in C_0([0,1])$, the formula
\[
 \Phi_F(x,y)=
   \bigl(\tanh(\operatorname{artanh}(x)+F(y)),y\bigr)
          \qquad(|x|<1)
                                                        \tag{C.26.5}
\]
on $[-1,1]\times[0,1]$, extended by the identity on the vertical
sides, fixes the entire boundary and satisfies
$\Phi_F\Phi_Q=\Phi_{F+Q}$. It is a topological-group embedding:
$\tanh$ is Lipschitz, giving continuity, and evaluation at $(0,y)$
recovers $\tanh(F(y))$, from which $F$ is recovered uniformly near
any given profile. Its image is closed. A homeomorphism limit fixes
the boundary and preserves the horizontal leaves; the image of the
compact segment $\{0\}\times[0,1]$ stays a positive distance from the
vertical sides. Applying $\operatorname{artanh}$ therefore recovers
uniformly bounded, uniformly convergent profiles.

Composing (C.26.4), $T$, and (C.26.5) embeds $G$ as a closed subgroup
of $\operatorname{Homeo}(D^2,\partial D^2)$. Extension by the identity
outside a coordinate disc is a closed topological-group embedding
into $A^0$ for any hyperbolic surface: restriction recovers the maps,
and being the identity off the disc is a closed condition. Its loop
(C.26.3) is ambient-null but nontrivial in $G$.
The precise enlargement killing it is already visible: if
$F_s=T([\mathbf1_{[0,s]}]_{\R})$, then $F_0=F_1=0$ and
$(r,s)\mapsto\Phi_{rF_s}$ is a based contraction in the larger real
shear group. No operator group cutting out a centralizer that excludes
this contraction has been constructed here. Thus this example refutes
the auxiliary assertion of fundamental-group injectivity for all
closed subgroups, while C.12 retains its stated centralizer question.

\begin{claim}{C.27}{Proved}{An order-natural contraction of line centralizer loops}
Give $\operatorname{Homeo}(\R)$ its compact-open topology, let
$H\leq\operatorname{Homeo}(\R)$ be an arbitrary subgroup, and put
\[
              C=C_{\operatorname{Homeo}^+(\R)}(H).
\]
Every based loop in $C$ has an explicit based null-homotopy in $C$.
Orientation-reversing elements of $H$ are allowed. The construction
commutes with conjugation by every homeomorphism of the line and
requires no finiteness, differentiability, or local compactness of $H$.
\end{claim}
\passage{23; folios 140--141; scans 279--281; Part C(b) infinite-operator request; exact partial reduction for C.12}

\begin{proof}
Let $t\mapsto f_t$ be a compact-open continuous loop in $C$, with
$f_0=f_1=\id$. For $0\leq r\leq1$ define
\[
 m_r(x)=\min_{0\leq u\leq r}f_u(x),\qquad
 M_r(x)=\max_{0\leq u\leq r}f_u(x).                    \tag{C.27.1}
\]
Both are increasing homeomorphisms. Indeed, for $x<y$, compactness
gives
\[
 \min_{0\leq u\leq1}\bigl(f_u(y)-f_u(x)\bigr)>0,
\]
which proves strict increase for each minimum and maximum. Their
properness is uniform in $r$: compactness of the inverse family
$\{f_u^{-1}:0\leq u\leq1\}$ gives, for every bounded target interval,
a common bound on all its inverse images. Continuity of (C.27.1),
jointly in $r$ and $x$, follows from uniform continuity on compact
parameter rectangles. Proper continuous strictly increasing maps of
$\R$ are homeomorphisms.

Set
\[
 F_{s,t}(x)=\operatorname{med}
      \bigl(m_{1-s}(x),f_t(x),M_{1-s}(x)\bigr)
   =\max\{m_{1-s}(x),\min\{f_t(x),M_{1-s}(x)\}\}.
                                                        \tag{C.27.2}
\]
Finite minima and maxima of increasing homeomorphisms are again
increasing homeomorphisms, by the same strict-increase and properness
arguments. Formula (C.27.2) is continuous in $(s,t,x)$, uniformly on
compact $x$-sets; it is therefore a continuous family in the
compact-open homeomorphism group. At $s=0$, the prefix interval is
the entire loop, so $m_1\leq f_t\leq M_1$ and $F_{0,t}=f_t$.
At $s=1$, $m_0=M_0=\id$, giving $F_{1,t}=\id$. Since $u=0$ occurs
in every prefix interval, $m_r\leq\id\leq M_r$, and hence
$F_{s,0}=F_{s,1}=\id$ for every $s$.

It remains to check the exact equivariance. An increasing
$h\in H$ fixes both envelopes under conjugation, because conjugation
commutes with their pointwise minimum and maximum and fixes every
$f_u$. A decreasing $h\in H$ interchanges $m_r$ and $M_r$ instead.
Every order-preserving or order-reversing bijection preserves the
median of three points. Thus both cases give
\[
                   h F_{s,t}h^{-1}=F_{s,t}
       \qquad(h\in H, 0\leq s,t\leq1).
\]
This proves that the entire based contraction lies in $C$. Replacing
the loop by its conjugate by an arbitrary line homeomorphism gives
the conjugate of the same formula, proving the asserted naturality.
\end{proof}

The surface problem in C.12 still requires a continuous construction
across different orbit strata. Applying (C.27.2) separately after
choosing coordinates on individual line orbits does not establish
continuity on the surface. No such gluing statement follows from
C.27, and none is used here.

\begin{claim}{C.28}{Proved}{Orbit types of a loop-generated plane group}
Let $t\mapsto f_t$ be a compact-open continuous based loop in
$\operatorname{Homeo}^+(\R^2)$, with $f_0=f_1=\id$, and put
\[
 L=\{f_t:0\leq t\leq1\},\qquad P=L\cup L^{-1},
 \qquad G=\langle P\rangle.
\]
Every $G$-orbit is one of the following: an open subset of the plane;
a point; a Jordan curve; or the image of a continuous injective map
$\eta:\R\to\R^2$ that restricts to a homeomorphism on every compact
interval. In the last case its intrinsic arc order identifies it with
the line, but no global homeomorphism from $\R$ to the orbit with its
subspace topology is asserted. The loop restricts continuously in the
compact-open topology to the intrinsically ordered line. Every
homeomorphism commuting with $G$ carries these orbit types to the same
types and induces a homeomorphism between the corresponding intrinsic
lines, allowing reversal of their orders.
\end{claim}
\passage{23; folios 140--141; scans 279--281; Part C(b) infinite-operator request; exact partial reduction for C.12}

\begin{proof}
Fix $x\in\R^2$ and put
\[
                  K_n=P^n x,\qquad O=Gx=\bigcup_{n\geq1}K_n.
                                                        \tag{C.28.1}
\]
The sets $K_n$ are nested compact Peano continua. Indeed, $P$ is a
continuous image of the wedge of two circles, using the given loop
and its inverse; it is compact, connected and locally connected.
A finite product of such spaces has the same properties, and the
word-evaluation map has image $K_n$. To recall why local connectedness
passes through this compact surjection, for an open neighborhood of a
point in the image cover its compact inverse fibre by finitely many
connected open subsets of its inverse image. Their images all meet
that point and lie in its component. The complement of the image of
the remaining compact set is a smaller open neighborhood in that
component. Thus components of open sets in the image are open.
In particular $K_n$ is arcwise connected. The identity lies in $P$,
which proves the nesting in (C.28.1). Every element of $G$ can be
joined to the identity by a path in $G$, by writing it as a finite word
and using the specified paths to its factors.

Suppose first that $O$ contains a Jordan curve $J$. For any $g\in G$
choose a path $g_s$ from $\id$ to $g$, so that $g_sJ\subseteq O$.
Let $D_s$ be the bounded Jordan domain of $g_sJ$. If $gJ\ne J$, then
$D_0\mathbin\triangle D_1$ contains a nonempty open set away from the
two boundary curves. Every point $q$ in this open set lies on some
$g_sJ$: otherwise the winding number of the continuously varying curve
$g_sJ$ around $q$ would be constant, whereas its endpoint values are
$0$ and $1$. Hence $O$ contains a nonempty open set. Transitivity of
$G$ on $O$ carries this open set to a neighborhood of each point of
$O$, so $O$ is open. If every $g$ preserves $J$, then for a point
$y\in J$ its orbit is contained in $J$; since $Gy=O$ and $J\subseteq O$,
we obtain $O=J$. Thus a nonopen orbit containing a Jordan curve is
itself that Jordan curve.

We may now consider a nonopen orbit containing no Jordan curve.
Each $K_n$ is a dendrite: a compact, connected, locally connected
metric space containing no Jordan curve. Arcs between two points are
unique, since two distinct arcs would contain a Jordan curve. This
also proves uniqueness of arcs in all of $O$: two proposed arcs in
$O$ would give the same forbidden curve. Existence follows by putting
their endpoints into one $K_n$.

Each $K_n$ has countably many branch points. Choose a countable dense
subset $D_n$. Three distinct components of the complement of a branch
point are open there, so choose a point of $D_n$ in each. The branch
point is the unique common point of the three pairwise connecting
arcs for this triple: two branches meet only at their point of
separation. Thus each branch point is determined by one of the
countably many triples from $D_n$.

A triod point of $O$ is a branch point of some $K_n$. Put the three
outer endpoints of its triod into one $K_n$; the unique connecting
arcs in $K_n$ must be the arcs of that triod. The triod points of $O$
are consequently countable. Their set is invariant under every element
of $G$. If it were nonempty, transitivity would make every point of
$O$ a triod point. A nontrivial arcwise connected orbit contains an
arc and is uncountable, giving a contradiction. Thus a nontrivial
$O$ has no triod points, and each nondegenerate $K_n$ is an arc.

The nested arcs give $O$ a consistent linear arc order, unique up to
reversal. To see the absence of additional branches directly, fix two
distinct points $a,b$. Every other point is on $[a,b]$, beyond $a$,
or beyond $b$, because any fourth alternative would give a triod.
There are at most two endpoints of this ordered union. The endpoint
set is invariant under $G$, so transitivity and uncountability exclude
endpoints. We can therefore pass to a cofinal sequence of arcs $A_j$
such that $A_j$ lies in the relative interior of $A_{j+1}$: extend
both endpoints in $O$ and then put the two extensions and the next
$K_n$ into a later $K_m$.

Choose compatible homeomorphisms $[-j,j]\to A_j$, extending at each
step along the two added arc intervals. Their union is a continuous
injective map
\[
                 \eta:\R\longrightarrow\R^2,
                        \qquad\eta(\R)=O,              \tag{C.28.2}
\]
whose restriction to every compact interval is a homeomorphism onto
its image. These assertions follow on each finite interval from one
of the compact arc parametrizations. This supplies the stated
intrinsic line structure. It does not prove that the inverse of
(C.28.2) is continuous for the global subspace topology; accumulations
from far along an orbit have not been excluded.

The restricted loop is continuous for that intrinsic line topology.
A compact arc in $O$ is contained in some $K_m$, and
$f_t(K_m)\subseteq K_{m+1}$ for all $t$. The latter is a compact arc,
on which the intrinsic and subspace topologies coincide. Joint
continuity of $(t,z)\mapsto f_t(z)$ on each such compact arc therefore
gives compact-open continuity of
$\eta^{-1}f_t\eta\in\operatorname{Homeo}(\R)$.
Its order sign is positive throughout the loop: images of two distinct
points vary continuously in the line and cannot cross. Finally a
commuting homeomorphism $h$ satisfies the exact equality
\[
                  h(K_n(x))=K_n(hx).
\]
It therefore maps compact arcs homeomorphically to compact arcs and
induces a homeomorphism of the intrinsic ordered lines; its inverse
has the same property. It also preserves the point, Jordan-curve and
open alternatives, completing the proof.
\end{proof}

\paragraph{Exact application and remaining step.}
The canonical lift in the paragraph after C.12 sends a hyperbolic
centralizer loop to a closed based loop in
$\operatorname{Homeo}^+(\widetilde U)$, where
$\widetilde U\cong\R^2$. The entire group $\widetilde H$ of lifts of
the operator image commutes with it and hence with the $G$ of C.28.
Thus C.28 applies to this actual lift group, and C.27 provides an
order-natural contraction on each intrinsic line orbit after
restriction to its stabilizer in $\widetilde H$.
These are typed restriction and lift maps, but they do not yet give a
homotopy through homeomorphisms of the whole plane. Continuity across
distinct orbit strata, compatibility with their frontiers, and a
relative contraction on proper open orbits remain unproved.
C.26 treats global transitivity on a closed hyperbolic surface; a
proper open orbit in C.28 does not satisfy that hypothesis. The
admissible infinite-group level-$0$ question therefore remains C.12.

\begin{claim}{C.29}{Proved}{Circle centralizers and the winding of a plane loop}
For every subgroup $H\leq\operatorname{Homeo}(S^1)$, put
$C=C_{\operatorname{Homeo}^+(S^1)}(H)$, with the compact-open topology.
The inclusion induces an injection
\[
 \pi_1(C,\id)\longrightarrow
 \pi_1(\operatorname{Homeo}^+(S^1),\id)=\Z,              \tag{C.29.1}
\]
where the integer is the winding of the image of one point. Its image
lies in $\Z^H$, with $H$ acting on $\Z$ by its orientation character.
A based loop of winding zero has an explicit contraction in $C$ that
is natural under all homeomorphisms between circles.
In particular, a based null-homotopic loop in
$\operatorname{Homeo}^+(\R^2)$ preserving a Jordan curve restricts to
a winding-zero loop on that curve. Its restriction therefore has this
natural contraction against every commuting circle homeomorphism.
\end{claim}
\passage{23; folios 140--141; scans 279--281; Part C(b) infinite-operator request; exact partial reduction for C.12}

\begin{proof}
Write the circle as $\R/\Z$, and set
\[
 \widetilde A=\{F\in\operatorname{Homeo}^+(\R):
                        F(x+1)=F(x)+1\}.
\]
The quotient map $\widetilde A\to\operatorname{Homeo}^+(S^1)$ is a
covering with deck transformations $F\mapsto F+n$, $n\in\Z$.
Indeed, requiring $F(0)$ to lie in a fixed interval of length less
than one gives continuous unique lifts over the corresponding
evaluation neighborhood. The homotopy
$(r,F)\mapsto(1-r)F+r\id$ contracts $\widetilde A$: strict
increase, the translation equation, and properness hold at every
parameter. Thus the target fundamental group in (C.29.1) is precisely
$\Z$. For a based loop $f_t$, its unique lift $F_t$ starting at
$\id$ ends at $x\mapsto x+n$, where $n$ is its point winding.

Suppose that $f_t\in C$ and $n=0$. Every lift $\widetilde h$ of an
element $h\in H$, including an orientation-reversing one, commutes
with every $F_t$. The commutator projects to the identity on the
circle, so it is an integer translation depending continuously on
$t$; it is the identity at $t=0$, and hence at all $t$.
Consequently $F_t$ is a based loop in the line centralizer of the
entire lift group of $H$. This group also contains the integer
translations. Apply C.27, explicitly
\[
 \begin{split}
 m_r(x)&=\min_{0\leq u\leq r}F_u(x),\qquad
 M_r(x)=\max_{0\leq u\leq r}F_u(x),\\
 Q_{s,t}(x)&=\operatorname{med}
       \bigl(m_{1-s}(x),F_t(x),M_{1-s}(x)\bigr).
 \end{split}                                                   \tag{C.29.2}
\]
Every $Q_{s,t}$ commutes with all lifts and with $x\mapsto x+1$,
so (C.29.2) descends to a contraction through elements of $C$.
This proves injectivity in (C.29.1).

The winding of a conjugated loop $h f_t h^{-1}$ is the winding of
$f_t$ multiplied by the orientation sign of $h$. Since these loops
agree for $h\in H$, the image is contained in $\Z^H$. For naturality,
let $q$ be a homeomorphism between two circles and choose any lift
$\widetilde q$ to their universal lines. The based lift of the
conjugated loop is $\widetilde q F_t\widetilde q^{-1}$.
An increasing $\widetilde q$ carries the two extrema in (C.29.2)
to the corresponding extrema; a decreasing one interchanges them.
Every line homeomorphism preserves the median of three points.
Thus (C.29.2) is carried to exactly the same formula for the
conjugated loop. Changing the lift of $q$ by an integer translation
does not change the descended contraction.

It remains to prove the stated plane winding comparison. Conjugate
by a Schoenflies homeomorphism so that the preserved Jordan curve is
the unit circle. Choose two distinct points $a,b$ on it. Evaluation
and the difference direction give the continuous maps
\[
 \operatorname{Homeo}^+(\R^2)
 \longrightarrow\operatorname{Conf}_2(\R^2)
 \longrightarrow S^1,
 \qquad
 f\longmapsto(f(a),f(b))
 \longmapsto\frac{f(b)-f(a)}{|f(b)-f(a)|}.                 \tag{C.29.3}
\]
For the given circle-preserving loop choose continuous angular lifts
$\alpha(t),\beta(t)$ with
$\alpha(t)<\beta(t)<\alpha(t)+2\pi$.
They both acquire the increment $2\pi n$: their difference remains
in $(0,2\pi)$ and returns to its initial value modulo $2\pi$.
The identity
\[
 e^{i\beta(t)}-e^{i\alpha(t)}
 =2i\,e^{i(\alpha(t)+\beta(t))/2}
            \sin\frac{\beta(t)-\alpha(t)}2
\]
shows that the winding of (C.29.3) is exactly $n$.
A null-homotopy of the ambient plane loop maps under (C.29.3) to a
null-homotopy in $S^1$, forcing $n=0$. This proves the exact
comparison and the last assertion.
\end{proof}

\paragraph{Application to C.12.}
The canonical lift of a hyperbolic centralizer loop is null-homotopic
as a loop of plane homeomorphisms: it is the image under the continuous
canonical lift homomorphism of a loop in the contractible $A^0$.
Every Jordan-curve orbit in C.28 is preserved throughout this loop.
Its restriction therefore has winding zero by (C.29.3), and (C.29.2)
contracts it naturally against its stabilizer in the full operator
lift group. Commuting operators carrying one circle orbit to another
carry these contractions to each other. This gives a circle-orbit
contraction with the exact required equivariance. C.34 subsequently
proves joint continuity, including inverses, on the full nonopen
locus. C.35 shows that this particular prescription need not admit
an equivariant extension over an adjacent open orbit, and C.36
constructs compatible alternatives on finite unions of closures
with two-dimensional proper operator image. The remaining global
compatibility requirements are stated in C.12, which remains open.

\begin{claim}{C.30}{Proved}{Properness of a transitive-action centralizer}
Let $V$ be a connected second-countable surface without boundary,
and let $K$ be a closed subgroup of $\operatorname{Homeo}(V)$ acting
transitively on $V$, with the compact-open topology. Put
$N=C_{\operatorname{Homeo}(V)}(K)$. Then $N$ is locally compact,
its action on $V$ is free and proper, and for every $x\in V$
evaluation is a homeomorphism onto the closed subset
\[
 N\xrightarrow{\ \mathrm{ev}_x\ }
 N_K(K_x)/K_x\ \subseteq\ K/K_x\cong V.                  \tag{C.30.1}
\]
Under the normalizer-quotient group structure, this is an
anti-isomorphism of topological groups. For an open orbit $O$ of the
loop-generated group in C.28, the image in $\operatorname{Homeo}(O)$
of its operator stabilizer lies in such an $N$; the closure of that
image acts freely and properly on $O$.
\end{claim}
\passage{23; folios 140--141; scans 279--281; Part C(b) infinite-operator request; exact partial reduction for C.12}

\begin{proof}
The group $K$ is Polish. The open mapping theorem
\cite[Theorem~3.5]{C_GomezPinedoUzcategui2017}, with the ordinary
continuous transitive action on the nonmeagre Hausdorff surface $V$,
identifies $K/K_x$ homeomorphically with $V$. In particular the
orbit map $k\mapsto kx$ is open. We will use the following exact
consequence, requiring no local continuous section: for every
sequence $x_j\to x$, there are $k_j\to1$ in $K$ with $k_jx=x_j$.
To prove it, choose a decreasing identity-neighborhood basis $W_m$.
The open sets $W_mx$ contain $x$, so successive tails of the sequence
lie in them. Choose $k_j$ from the corresponding $W_{m(j)}$, with
$m(j)\to\infty$.

An element of $N$ is determined by its value at $x$, since
$n(kx)=kn(x)$. If $n(x)=ax$, equality of stabilizers gives
$aK_xa^{-1}=K_x$. Conversely, for $a$ in this normalizer the map
$kK_x\mapsto kaK_x$ is a homeomorphism commuting with $K$.
The normalizer is closed: both inclusions defining
$aK_xa^{-1}=K_x$ are intersections of closed conditions.
It is saturated for the quotient by $K_x$, so its image is a closed
subset of $K/K_x=V$. This proves the claimed set-theoretic image
and closedness in (C.30.1). If $n(x)=ax$ and $m(x)=bx$, then
$(nm)(x)=bax$, proving the stated anti-isomorphism.

We prove inverse continuity of evaluation, including the issue of
noncompact $V$. Suppose that $n_j(x)\to n(x)$, with $n_j,n\in N$.
For any compact set $B\subset V$, a failure of uniform convergence
on $B$ would provide a subsequence and $y_j\in B$ with
$y_j\to y$ witnessing that failure. Choose $a\in K$ with $ax=y$,
and then $b_j\to1$ with $b_jy=y_j$, using open evaluation at $y$.
Commutation gives
\[
 n_j(y_j)=n_j(b_jax)=b_ja\,n_j(x)\longrightarrow a n(x)=n(y).
                                                               \tag{C.30.2}
\]
Also $n(y_j)\to n(y)$, a contradiction. Hence $n_j\to n$ uniformly
on every compact set, which is compact-open convergence.
Inversion can be checked directly as well. Put $c_j=n^{-1}n_j$,
so $c_jx\to x$, and choose $a_j\to1$ with $a_jx=c_jx$.
Because $c_j$ commutes with $a_j$,
$c_j^{-1}x=a_j^{-1}x\to x$; applying (C.30.2) to $c_j^{-1}$
gives convergence of the inverses. All the spaces are metrizable,
so this sequential argument proves inverse continuity.
Thus (C.30.1) is a homeomorphism, and $N$ is locally compact as a
closed subspace of the locally compact surface.

If $n\in N$ fixes one point, transitivity and commutation make it
fix every point, proving freeness. To prove properness, consider
sequences with
\[
                  x_j\longrightarrow x,\qquad
                  n_jx_j\longrightarrow y.
\]
Choose $k_j\to1$ with $k_jx=x_j$. Then
\[
                n_jx=k_j^{-1}n_jx_j\longrightarrow y.    \tag{C.30.3}
\]
The image of evaluation at $x$ is closed, so $y=nx$ for a unique
$n\in N$. Its inverse continuity gives $n_j\to n$.
For any compact subset of $V\times V$, every sequence in its
inverse image under $(n,z)\mapsto(z,nz)$ therefore has a convergent
subsequence, with limit in the same inverse image. Metrizability
turns this sequential compactness into compactness, proving that
the action map is proper. Every closed subgroup of $N$ inherits
the free proper action.

Finally let $G$ be the group in C.28 and let $O$ be one of its open
orbits. It is connected because every element of $G$ is joined to
the identity by a path in $G$. The closure
\[
        K=\overline{\{g|_O:g\in G\}}
                       ^{\operatorname{Homeo}(O)}
\]
is a closed Polish group still transitive on $O$. Let $\widetilde H$
be any group commuting with $G$, and let $\widetilde H_O$ be its
setwise stabilizer of $O$. The precise restriction homomorphism is
\[
 \widetilde H_O\longrightarrow\operatorname{Homeo}(O),\qquad
 h\longmapsto h|_O,
 \quad
 \ker=\{h\in\widetilde H_O:h|_O=\id\}.                 \tag{C.30.4}
\]
Its image commutes with the restriction of $G$ and hence with $K$.
The closure of this image is a closed subgroup of $N$, and therefore
acts freely and properly on $O$, as claimed.
\end{proof}

\paragraph{Precise remaining use.}
The closure in the last assertion has its subspace topology from
$\operatorname{Homeo}(O)$. Formula (C.30.4) does not assert that the
original discrete operator group acts freely, since its kernel may
be nontrivial; nor does it equip a possibly dense abstract image with
the discrete topology and claim properness for that topology.
The result supplies the closed locally compact transformation group
actually present on an open orbit. No Lie-group classification or
quotient-bundle contraction is invoked here. A contraction relative
to orbit frontiers, and its continuous assembly with C.27 and C.29
on the one-dimensional strata, remain unproved. Thus C.30 does not
change the unresolved status of C.12.

\begin{claim}{C.31}{Proved}{Track control and continuity at the fixed stratum}
Let $f_t$ be a based null-homotopic loop in
$\operatorname{Homeo}^+(\R^2)$, let $G$ be its generated group as in
C.28, and let $Y$ be the union of the nonopen $G$-orbits. This is a
closed subset of the plane. Define $Q_{s,t}$ on $Y$ by C.27 on each
intrinsic line orbit, by the descended formula of C.29 on each circle
orbit, and by the identity on point orbits. Then every $Q_{s,t}$ is
a bijection of $Y$ preserving each orbit, and
\[
 Q_{s,t}^{-1}=Q_{s,t}[f^{-1}],\qquad
 Q_{s,t}(x)\in\{f_u(x):0\leq u\leq1\},\qquad
 Q_{s,t}^{-1}(x)\in\{f_u^{-1}(x):0\leq u\leq1\}.
                                                        \tag{C.31.1}
\]
Here $Q[f^{-1}]$ means the same construction applied to the based
inverse loop $u\mapsto f_u^{-1}$, which generates the same group.
Both evaluation maps $(s,t,x)\mapsto Q_{s,t}^{\pm1}(x)$ are continuous
at every point of $[0,1]^2\times Z$, where
$Z=\{x:f_u(x)=x\text{ for all }u\}$ is the point-orbit stratum.
The construction commutes with every homeomorphism commuting with
all $f_u$. No continuity at a nonfixed orbit is asserted here.
\end{claim}
\passage{\S{}23; folios140--141; scans279--281; Part C(b) infinite-operator request; fixed-stratum gluing reduction for C.12}

\begin{proof}
The union of open orbits is open, proving that $Y$ is closed.
C.28 supplies exactly the remaining point, circle and intrinsic-line
types. The circle restrictions have zero winding by C.29, since
the ambient loop is null-homotopic. Thus the formulas defining $Q$
are available on every orbit in $Y$, and are independent of the
chosen line or circle coordinates by their proved naturality.
Each formula gives a homeomorphism of its single intrinsic line
or circle, so together they give a bijection of $Y$.

We prove the inverse identity before making any continuity claim.
For increasing homeomorphisms of a line, pointwise inversion
reverses order and exchanges finite minimum and maximum:
\[
 (a\wedge b)^{-1}=a^{-1}\vee b^{-1},\qquad
 (a\vee b)^{-1}=a^{-1}\wedge b^{-1}.
                                                        \tag{C.31.2}
\]
For example, at a point $x$ the first inverse is the larger of
$a^{-1}(x)$ and $b^{-1}(x)$: at this larger point one of $a,b$
equals $x$, and the other is at least $x$.
The same argument works for compact families, whose extrema
and inverse extrema are attained. Therefore, with the notation
of C.27,
\[
 m_r^{-1}(x)=\max_{0\leq u\leq r} f_u^{-1}(x),\qquad
 M_r^{-1}(x)=\min_{0\leq u\leq r} f_u^{-1}(x).
                                                        \tag{C.31.3}
\]
The median of three line homeomorphisms is the lattice expression
$(a\wedge b)\vee(a\wedge c)\vee(b\wedge c)$; its dual expression
$(a\vee b)\wedge(a\vee c)\wedge(b\vee c)$ has the same value.
Applying (C.31.2) thus proves
\[
 \operatorname{med}(a,b,c)^{-1}
       =\operatorname{med}(a^{-1},b^{-1},c^{-1}).
                                                        \tag{C.31.4}
\]
Combining (C.31.3)--(C.31.4), and using symmetry of the median
in its two envelope entries, gives the first identity in (C.31.1)
on every line orbit. On a circle, the inverse of the unique based
lift $F_u$ is exactly the unique based lift of $f_u^{-1}$.
The same calculation on the universal line descends to the circle.
It also holds on point orbits.

At any point the median chooses one of its three entries. Each
envelope entry is $f_u(x)$ for a parameter attaining its extremum,
and the middle entry is $f_t(x)$ itself. This proves the second
assertion of (C.31.1) on lines; projection of the same statement
for the based lifts proves it on circles. The inverse identity
then proves the third assertion.

Fix $p\in Z$ and an open neighborhood $V$ of $p$ in the plane.
Joint continuity of the action, compactness of $[0,1]$, and
$f_u(p)=p$ give a neighborhood $W$ of $p$ such that
$f_u(W)\subseteq V$ for every $u$. The compact inverse family
gives the same assertion for all $f_u^{-1}$ after shrinking $W$.
By (C.31.1),
\[
 Q_{s,t}^{\pm1}(W\cap Y)\subseteq V\cap Y
                   \quad\text{for all }s,t\in[0,1].       \tag{C.31.5}
\]
Since $Q_{s,t}^{\pm1}(p)=p$, this proves joint continuity at every
$(s,t,p)$, uniformly in the two homotopy parameters.

A homeomorphism commuting with all $f_u$ permutes the $G$-orbits,
preserves their four types, and maps intrinsic lines by line
homeomorphisms by C.28. Naturality in C.27 and C.29 gives
$hQ_{s,t}h^{-1}=Q_{s,t}$ on all of $Y$. This proves the asserted
equivariance without any continuity assertion at the remaining
points.
\end{proof}

\paragraph{Scope of the gluing step.}
Formula (C.31.5) resolves continuity at the fixed stratum for the
orbitwise construction on $Y$. It does not yet make the bijections
$Q_{s,t}:Y\to Y$ homeomorphisms, because continuity at the nonfixed
line and circle strata remains unproved. Nor does it control a
contraction on the omitted open orbits. These are the remaining
global assembly problems for this route to C.12.

\begin{claim}{C.32}{Proved}{Discrete proper operator groups on open orbits}
Let $O$ be a connected noncompact second-countable surface without
boundary, and let a discrete group $Q\leq\operatorname{Homeo}(O)$ act
freely and properly. With compact-open topology on the homeomorphism
groups, inclusion induces an injection
\[
 \pi_1\!\left(C_{\operatorname{Homeo}(O)}(Q),\id\right)
       \longrightarrow \pi_1(\operatorname{Homeo}(O),\id).
                                                        \tag{C.32.1}
\]
Moreover, a based null-homotopic loop in
$\operatorname{Homeo}^+(\R^2)$ preserving a connected open subset
$O\subseteq\R^2$ restricts to a null-homotopic loop in
$\operatorname{Homeo}(O)$.
Consequently, for a hyperbolic centralizer loop and an open orbit in
C.28 whose closed restricted operator image in C.30 is discrete, its
restriction has a based equivariant null-homotopy on that orbit.
The original restriction kernel is retained; no extension of this
null-homotopy to the orbit frontier is asserted.
\end{claim}
\passage{Section 23, folios 140-141, scans 279-281; Part C(b) infinite-operator request; discrete-closure open-orbit reduction for C.12}

\begin{proof}
First consider the free proper discrete action on $O$. It gives an
actual covering
\[
                       p:O\longrightarrow V=O/Q.          \tag{C.32.2}
\]
Here discreteness refers to the topology in the proper-action
hypothesis. For completeness, a compact neighborhood of a point
meets only finitely many of its translates, by properness. Shrink
an open neighborhood inside it to avoid the nonidentity translates
among this finite set, using freeness. Distinct translates of the
result are disjoint, which gives the evenly covered neighborhoods
for (C.32.2). The quotient $V$ is Hausdorff by properness, is a
surface without boundary by these local charts, and is second
countable because images of a countable open basis of $O$ form an
open basis. The connected covering is regular, with deck group
exactly $Q$.

We record the homeomorphism-group input precisely. Every
second-countable surface admits a PL triangulation. Applied with
the compact subpolyhedron consisting of one point $v$, the compact
case of \cite[Lemma~2.1]{A_Yagasaki2000} and the noncompact case of
\cite[Theorem~1.1]{A_Yagasaki2000} give
\[
 \operatorname{Homeo}(V,v)^0\simeq *
 \quad\text{for }V\not\cong S^2,\ \R P^2,\ \R^2.          \tag{C.32.3}
\]
These are all exceptions relevant to a surface without boundary
with that one-point subpolyhedron. Evaluation
$\operatorname{Homeo}(V)^0\to V$ is a locally trivial bundle:
in a coordinate disc a continuously chosen homeomorphism supported
in that disc moves its center to a nearby point, giving a local
section. Its identity fibre component is exactly the group in
(C.32.3). The homotopy exact sequence therefore gives an injection
\[
 \pi_1(\operatorname{Homeo}(V),\id)
          \longrightarrow\pi_1(V,v)
       \quad(V\not\cong S^2,\R P^2,\R^2).                \tag{C.32.4}
\]
The base $V$ of (C.32.2) cannot be $S^2$ or $\R P^2$:
their fundamental groups are finite, so every connected covering
has finitely many sheets and compact total space, contrary to
noncompactness of $O$. If $V=\R^2$, its connected covering is
one-sheeted, so $Q=1$ and (C.32.1) is the identity homomorphism.

Let $f_t$ be a based loop in the centralizer that is null-homotopic
in $\operatorname{Homeo}(O)$. It descends to a based loop $b_t$
on $V$, defined by the exact equation
\[
                         b_t p=p f_t.                    \tag{C.32.5}
\]
The family is continuous: joint continuity descends through the
open quotient map $[0,1]\times O\to[0,1]\times V$, and the
compact-open exponential law applies since $V$ is locally compact.
For $x\in O$, the evaluation loop $f_t(x)$ is null-homotopic in
$O$. Thus $b_t(p(x))=p(f_t(x))$ is null-homotopic in $V$.
By (C.32.4), $b_t$ is a null-homotopic based loop of homeomorphisms
of $V$. Choose a continuous filling
$B:D^2\to\operatorname{Homeo}(V)^0$ with the prescribed boundary
$b_t$, and a point $z_*\in\partial D^2$ corresponding to the based
endpoint, so that $B_{z_*}=\id$.

The lift is a filling through homeomorphisms and through the exact
centralizer; we verify both assertions. Form
\[
 F:O\times D^2\longrightarrow V,\qquad
                       F(x,z)=B_z(p(x)).                 \tag{C.32.6}
\]
At $(x_*,z_*)$, the induced fundamental-group image is
$p_*\pi_1(O,x_*)$, since $D^2$ is simply connected and
$F(x,z_*)=p(x)$. The covering lifting criterion gives a unique
continuous lift $\widetilde B:O\times D^2\to O$ with
$\widetilde B(x,z_*)=x$.
Apply the identical construction to
$F^-(x,z)=B_z^{-1}(p(x))$, giving a continuous lift
$\widetilde C$ that is the identity at $z_*$.
The maps
$(x,z)\mapsto\widetilde B_z\widetilde C_z(x)$ and
$(x,z)\mapsto x$ lift the same map $(x,z)\mapsto p(x)$
and agree at the basepoint. Uniqueness of lifts makes them equal;
the reverse composition is equal for the same reason.
Hence every $\widetilde B_z$ is a homeomorphism, with inverse
$\widetilde C_z$, and the family is compact-open continuous.

For $q\in Q$, the two maps
$\widetilde B_z q$ and $q\widetilde B_z$ lift the same map
$B_zp$ and agree at $z_*$, so they are equal for every $z$.
On $O\times\partial D^2$, both $\widetilde B_z$ and the specified
$f_t$ lift (C.32.5) and agree at the based endpoint. They are
therefore equal on that entire connected boundary product.
Thus $\widetilde B$ fills the original loop in its centralizer.
This proves (C.32.1).

Now let $O$ be a connected open subset of the plane preserved by
a based ambient-null loop. Its restriction is a continuous loop
in $\operatorname{Homeo}(O)^0$. By
\cite[Theorem~1.1]{A_Yagasaki2000}, with the empty subpolyhedron,
this identity group is homotopically trivial unless $O$ is the
plane or the open annulus; the nonorientable and boundary
exceptions in the source cannot occur for $O\subseteq\R^2$.
In either remaining case its fundamental group is $\Z$.
Choose model coordinates on $O$, a coordinate Jordan circle
$J\subset O$, and two distinct points $a,b$ of $J$. A Schoenflies
homeomorphism $h$ of the ambient plane sends $J$ to the unit circle.
The map
\[
 \Theta:\operatorname{Homeo}(O)\longrightarrow S^1,\qquad
 k\longmapsto
       \frac{h(k(b))-h(k(a))}
            {|h(k(b))-h(k(a))|}
                                                        \tag{C.32.7}
\]
has induced homomorphism $\pi_1\operatorname{Homeo}(O)\to\Z$
that is an isomorphism. Indeed, the full rotation loop in the
chosen plane or annulus model preserves $J$ and has restriction
winding one. The explicit computation in C.29 shows that its
image under (C.32.7) also has winding one. A homomorphism
$\Z\to\Z$ whose image contains $1$ is an isomorphism.
Formula (C.32.7) extends as a continuous map on all ambient
plane homeomorphisms by the same two evaluations. The
ambient-null loop therefore has zero image under it, so its
restricted fundamental-group class is zero. This proves the
second assertion in all cases.

Finally take an open orbit from C.28 and let $Q$ be the closed
restricted operator image supplied by C.30. When this group is
discrete, C.30 gives a free proper action with precisely that
discrete topology, so (C.32.2) applies. The restricted canonical
plane loop commutes with $Q$ and is null-homotopic in
$\operatorname{Homeo}(O)$ by the second assertion. Its
null-homotopy supplied by (C.32.1) commutes with $Q$, hence with
the image of the operator stabilizer. The kernel in (C.30.4)
acts identically on $O$ throughout. This proves the stated
equivariant conclusion with the actual restriction homomorphism.
\end{proof}

\paragraph{Transport and frontier control.}
One may choose the obtained contraction on one representative of
an operator orbit of open $G$-orbits and transport it to the others
by conjugation. Equivariance under the setwise stabilizer makes
this independent of the transporting operator. This establishes
well-defined orbitwise families. The covering lift above supplies
no uniform estimate at orbit frontiers. C.33 subsequently treats
all positive-dimensional proper closed operator images
intrinsically, and C.34 proves joint continuity on the full
nonopen locus. C.35 records the obstruction to extending its
particular median prescription, while C.36 constructs compatible
alternatives on finite unions of closures with two-dimensional
proper operator image. The remaining frontier and accumulation
requirements are stated in C.12, which remains open.

\begin{claim}{C.33}{Proved}{All proper operator groups on an open orbit}
Let $O$ be a connected noncompact second-countable surface without
boundary. Let $Q$ be a closed locally compact subgroup of
$\operatorname{Homeo}(O)$, with its compact-open topology, acting
freely and properly. Then inclusion induces an injection
\[
 \pi_1\!\left(C_{\operatorname{Homeo}(O)}(Q),\id\right)
       \longrightarrow \pi_1(\operatorname{Homeo}(O),\id).
                                                        \tag{C.33.1}
\]
For every open orbit of the canonically lifted hyperbolic
centralizer loop in C.28, its restriction therefore has a based
null-homotopy commuting with the entire restricted operator
stabilizer. This conclusion retains the restriction kernel in
C.30 and imposes no finiteness or discreteness hypothesis on the
original operator group.
It supplies an intrinsic null-homotopy on each open orbit;
continuity across orbit frontiers remains unproved.
\end{claim}
\passage{Section 23, folios 140-141, scans 279-281; Part C(b) infinite-operator request; arbitrary proper-closure open-orbit reduction for C.12}

\begin{proof}
Write $N=C_{\operatorname{Homeo}(O)}(Q)$.
The topological group $Q$ acts faithfully on the connected
three-manifold $O\times\R$ by
\[
                         q(x,s)=(qx,s).
                                                        \tag{C.33.2}
\]
This action is continuous: on any compact subset of $O\times\R$
the first coordinate is controlled on its compact projection to
$O$. Pardon proves that a locally compact group acting faithfully
on a connected three-manifold is a Lie group
\cite[Theorem~1.5 and the reduction preceding it]{C_Pardon2013}.
Thus $Q$ is a Lie group. Since it is a subspace of the
second-countable homeomorphism group, it is second countable.
Set $A=Q^0$ and $d=\dim A$.
Properness and freeness make $q\mapsto qx$ a proper embedding of
$A$ in $O$ for each $x$. Invariance of dimension gives $d\leq2$.
The case $d=0$ is precisely C.32.

If $d=2$, invariance of domain makes every $A$-orbit open.
Connectedness gives a single orbit, and evaluation $A\to O$
is a homeomorphism. For $q\in Q$, choose $a\in A$ with $qx=ax$.
Then $a^{-1}q$ fixes $x$, so freeness gives $q=a$ and $Q=A$.
Under $O\cong A$ with its left translation action, the elements of
$N$ are exactly the right translations. The correspondence
is an anti-isomorphism of topological groups, and evaluation
$N\to O$ is a homeomorphism. A loop in $N$ null-homotopic
in $\operatorname{Homeo}(O)$ has a null evaluation loop in $O$,
and hence is null-homotopic in $N$. This proves (C.33.1) in
dimension two.

It remains to treat $d=1$. The connected one-dimensional Lie
group $A$ is $\R$ or $\R/\Z$, written additively.
Its restricted action is still free and proper.
Palais' slice theorem gives a locally trivial principal bundle
\[
                 p_0:O\longrightarrow B=O/A.
                                                        \tag{C.33.3}
\]
Here all hypotheses of the theorem hold: the surface is completely
regular, $A$ is Lie, and a relatively compact neighborhood has
relatively compact transporter by properness, which is the
Cartan condition. A slice at a point has trivial stabilizer.
The resulting local trivialization is the map
$(a,s)\mapsto as$ from $A\times S$ onto the open saturated
neighborhood of that slice. This is the free-action specialization
of \cite[Theorem~2.3.3]{C_Palais1961}; its maps and inverse are
also given in \cite[Theorem~3.7]{C_KramerMuratGarcia2026}.

We verify the dimension and local topology of the base rather
than assuming differentiability of this action. The elementary
fact needed is that a metrizable space $S$ for which
$S\times(-1,1)$ is a surface is a one-manifold without boundary.
Indeed, projecting path-connected open neighborhoods in the
product shows that $S$ is locally path connected. It has no
isolated point, since an isolated point would give an open
interval in a surface. Every embedded open arc in $S$ is open:
its product with $(-1,1)$ maps injectively into a surface, and
invariance of domain applies.

Fix $p\in S$. Local path connectedness and the absence of isolated
points give an embedded arc $\alpha:[0,1]\to S$ with
$\alpha(0)=p$. Choose a path-connected open neighborhood $W$
of $p$ avoiding $\alpha(1)$. If $W$ contains a point outside
$\alpha([0,1])$, join $p$ to that point by an embedded arc in
$W$. This second arc cannot enter and subsequently leave
$\alpha((0,1))$: that open arc has boundary contained in
$\{p,\alpha(1)\}$, and an embedded arc in $W$ cannot revisit
$p$ or pass through $\alpha(1)$. Thus the two arcs meet only
at $p$, and together contain an arc having $p$ as an interior
point. Such an open arc is an open neighborhood of $p$ by the
preceding paragraph.
In the other case $W\subseteq\alpha([0,1))$, and a smaller
neighborhood of $p$ is a half interval. Its product with
$(-1,1)$ has zero local second homology at $(p,0)$,
whereas a surface without boundary has local second homology
$\Z$. This is impossible. The elementary fact follows.
Applied to the slices in (C.33.3), it proves that $B$ is a
one-manifold without boundary. Properness gives a Hausdorff
quotient, the open quotient map gives second countability, and
connectedness is inherited from $O$. Consequently $B$ is
$\R$ or $\R/\Z$.

Every principal $A$-bundle over either of these one-manifolds
is trivial. One can see the actual section as follows.
Over an interval, local trivializations are joined successively
across overlapping subintervals, adjusting one section by its
continuous $A$-valued transition function. A locally finite
interval exhaustion gives a section on $\R$.
Over the circle, first cut at one point and choose a section
on the resulting closed interval. Its endpoint discrepancy is
an element of the connected group $A$; multiplying the section
by a path from $0$ to the negative of that element makes the
endpoints agree, so that it descends to the circle.
The map $(a,b)\mapsto a\,s(b)$, with inverse given in each
bundle chart, therefore identifies
\[
                      O\cong A\times B,\qquad
             c(a,b)=(c+a,b)\quad(c\in A).                \tag{C.33.4}
\]
The pair $A=B=\R/\Z$ is excluded because $O$ is noncompact.

Let $\Delta=Q/A$. This discrete group acts freely and properly
on $B$. Freeness follows directly: an element of $Q$ fixing a
point of $B$ sends any point of its fibre to an $A$-translate;
after that translate is removed, it fixes a point of $O$ and
is the identity.
For properness, suppose $b_j\to b$ and
$\delta_j b_j\to b'$ in $B$. Local sections give $x_j\to x$
above $b_j$ and $y_j\to y$ above $\delta_j b_j$.
Choose a representative $q_j$ of $\delta_j$ and $a_j\in A$
such that $a_jq_jx_j=y_j$. Properness of the $Q$-action gives
a convergent subsequence of $a_jq_j$, so their components
$\delta_j$ are eventually constant on that subsequence.
This proves properness of the discrete action on the locally
compact metrizable space $B$.

We spell out this discrete action. On $\R$, an orientation
reverser has a fixed point and is excluded by freeness.
If $\Delta$ is nontrivial, its orbit of a point is closed and
discrete, with a nearest point to the right. The element
taking the chosen point to that next point generates
$\Delta$: its iterates are unbounded in both directions, and
any intervening orbit point would contradict the choice of
the nearest one. Choosing coordinates on one fundamental
interval and extending by its iterates identifies the action
with
\[
                   \Delta=\Z,\qquad \beta(b)=b+1.
                                                        \tag{C.33.5}
\]
On $\R/\Z$, properness makes $\Delta$ finite and freeness
again excludes orientation reversers. The action on one
finite orbit preserves its cyclic order, so it identifies
$\Delta$ with a cyclic group $C_n$. Coordinates on the
successive complementary intervals identify a generator
with
\[
                   \beta(b)=b+1/n\quad\text{in }\R/\Z.
                                                        \tag{C.33.6}
\]
The trivial group is also allowed on either base.

Suppose $\Delta$ is nontrivial and choose a representative
$q\in Q$ of the indicated generator. Conjugation by $q$ on
$A$ is multiplication by $\lambda$, where
$\lambda\in\R^\times$ for $A=\R$, and
$\lambda\in\{1,-1\}$ for $A=\R/\Z$.
Equation (C.33.4) therefore expresses its action exactly as
\[
                 q(a,b)=(\lambda a+\tau(b),\beta(b))
                                                        \tag{C.33.7}
\]
for a continuous function $\tau:B\to A$.
We next give a coordinate change removing this translation
term, with all component relations retained.

For $B=\R$ and $\Delta=\Z$, choose a continuous function
$c:[0,1]\to A$ whose endpoints satisfy
$c(1)=\lambda c(0)+\tau(0)$, and extend it in both directions
by
\[
                      c(b+1)=\lambda c(b)+\tau(b).
                                                        \tag{C.33.8}
\]
The endpoint equation makes the extension continuous at every
integer; invertibility of multiplication by $\lambda$ gives
the extension towards negative integers. In the new fibre
coordinate $a'=a-c(b)$, (C.33.7) becomes
$q(a',b)=(\lambda a',b+1)$.

For $B=\R/\Z$ and $\Delta=C_n$, necessarily $A=\R$.
Write $q^n=T_{a_0}$, where $T_{a_0}$ is translation by $a_0$
on the fibre. Conjugation gives $\lambda^n=1$ and
$\lambda a_0=a_0$.
If $\lambda=1$, replacing $q$ by $T_{-a_0/n}q$ makes
$q^n=1$ and does not change the group generated with $A$.
If $\lambda=-1$, then $n$ is even and $a_0=0$, so $q^n=1$
already. Thus the negative value of $\lambda$ is retained.
For this finite-order representative define the real-valued
function
\[
 c(b)=\frac1n\sum_{k=0}^{n-1}
      \operatorname{pr}_A\!\left(q^k(0,\beta^{-k}b)\right).
                                                        \tag{C.33.9}
\]
The affine nature of (C.33.7), followed by cyclic
reindexing using $q^n=1$, gives
$c(\beta b)=\lambda c(b)+\tau(b)$.
The same coordinate change $a'=a-c(b)$ now gives
$q(a',b)=(\lambda a',\beta b)$.
We have proved, by explicit changes of coordinates, that
the nonidentity components can be represented as
\[
                         q(a,b)=(\lambda a,\beta b).
                                                       \tag{C.33.10}
\]
When $\Delta$ is trivial no such additional equation is needed.

Commuting with all translations in (C.33.4) forces a
homeomorphism of $A\times B$ to have the form
\[
       h(a,b)=(a+\phi(b),\psi(b)),\qquad
       \phi\in C(B,A),\quad\psi\in\operatorname{Homeo}(B).
                                                       \tag{C.33.11}
\]
The inverse is
$(a,b)\mapsto(a-\phi(\psi^{-1}b),\psi^{-1}b)$.
These formulas and evaluation at the zero section give
mutually inverse continuous maps for the compact-open
topologies. In the nontrivial-component case, commuting
with (C.33.10) is exactly the pair of equations
\[
             \psi\beta=\beta\psi,\qquad
             \phi(\beta b)=\lambda\phi(b).
                                                       \tag{C.33.12}
\]
Thus the map (C.33.11) identifies $N$, as a space, with
the product of the function space defined by the second
equation and the homeomorphism centralizer defined by the
first. In particular the two constraints can be contracted
separately without changing the group $Q$ being centralized.

Let $h_t$ be a based loop in $N$ null-homotopic in
$\operatorname{Homeo}(O)$, and let $(\phi_t,\psi_t)$ be its
coordinates. The path $\psi_t$ preserves the orientation of
the one-manifold $B$.
When $A=\R$, multiplication of $\phi_t$ by $1-s$ contracts
the function-coordinate loop to zero and preserves its
equation in (C.33.12), including $\lambda=-1$.
For $B=\R$, C.27 contracts $\psi_t$ in the required
centralizer, with the trivial operator group when
$\Delta=1$. For $B=\R/\Z$, the evaluation loop
$h_t(0,b)=(\phi_t(b),\psi_t(b))$ is null-homotopic in
$\R\times(\R/\Z)$; hence $\psi_t(b)$ has winding zero.
C.29 contracts $\psi_t$ in the centralizer of the finite
rotation group in (C.33.6).

When $A=\R/\Z$, necessarily $B=\R$.
The based family $\phi_t(b)$ has a unique continuous lift
$\varphi:[0,1]\times\R\to\R$ with $\varphi_0=0$.
Evaluation of $h_t$ in the annulus $A\times\R$ shows that
$\varphi_1(b)=0$: this endpoint is the winding number of
the first coordinate of the null evaluation loop.
If $\Delta=\Z$, equation (C.33.12) says that
\[
           \varphi_t(b+1)-\lambda\varphi_t(b)\in\Z.
                                                       \tag{C.33.13}
\]
Its left side is a continuous function of $(t,b)$, equal
to zero at $t=0$, and is therefore identically zero.
Multiplying $\varphi$ by $1-s$ and descending to $\R/\Z$
contracts $\phi_t$ while preserving exact equivariance.
The same scaling applies when $\Delta=1$, without (C.33.13).
C.27 contracts the remaining $\psi_t$ on the line.
All these contractions keep $t=0,1$ fixed. Equations
(C.33.11)--(C.33.12), including the inverse formula, show
that their product is a continuous based contraction in
$N$. This completes the proof of (C.33.1).

Finally, for an open loop-group orbit from C.28, C.30
supplies the actual closed locally compact free proper
operator image $Q$. The restricted canonical plane loop
commutes with $Q$. It is null-homotopic in
$\operatorname{Homeo}(O)$ by the second assertion of
C.32, so (C.33.1) contracts it in the exact centralizer.
Every original operator in the setwise stabilizer acts
through its restriction to $Q$, and the restriction
kernel acts identically on $O$. The stated conclusion
follows with that kernel retained.
\end{proof}

\paragraph{What the open-orbit result provides.}
Equations (C.33.11)--(C.33.13) treat every component of the
proper operator group, including a circle fibre with a
nontrivial $\Z$-action on its line base. There is no
replacement of the base-translation centralizer by the
homeomorphism group of its quotient circle.
As in C.32, the obtained contraction can be transported to
other open orbits in the same operator orbit; exact
stabilizer equivariance makes the transport independent
of choices. This proof supplies no uniform control near
the boundary of the open orbit. C.34 subsequently proves
joint continuity, including inverses, on the entire closed
union of nonopen orbits. C.35 exhibits an exact obstruction
to extending that particular median prescription over an
adjacent open orbit. C.36 constructs compatible contractions
on closures with two-dimensional proper operator image and
proves continuity on every finite union of those closures.
C.37 supplies the uniform three-word estimate and gives a
continuous contraction on the closed union $X_2$ of all those
open orbits, including infinite accumulation. Extension across
$X_2$ to its exterior and frontier-compatible choices for the
dimension-zero and dimension-one operator images remain
additional requirements. C.12 retains its unresolved status.

\begin{claim}{C.34}{Proved}{The median contraction on all nonopen orbits}
Let $f_t$ be a based null-homotopic loop in
$\operatorname{Homeo}^+(\R^2)$, and use $P,G,Y$ from
C.28 and C.31. The orbitwise bijections $Q_{s,t}:Y\to Y$
defined in C.31 are homeomorphisms. Both maps
\[
 [0,1]^2\times Y\longrightarrow Y,\qquad
                  (s,t,x)\longmapsto Q_{s,t}^{\pm1}(x)
                                                        \tag{C.34.1}
\]
are continuous. Thus these formulas give a based
compact-open null-homotopy on the entire closed nonopen
locus, commuting with every homeomorphism commuting
with all $f_t$.
In particular, a hyperbolic centralizer loop whose
canonically lifted generated group has no open orbit
is null-homotopic in the original surface centralizer.
No extension over the omitted open orbits is asserted.
\end{claim}
\passage{Section 23, folios 140-141, scans 279-281; Part C(b) infinite-operator request; nonopen-stratum assembly in C.12}

\begin{proof}
We prove continuity at a nonfixed orbit. Continuity at
the common fixed set, bijectivity on $Y$, equivariance,
and the exact inverse identity have already been proved
in C.31.

\emph{Short paths moving in either direction.}
Let $q$ be a point of a nonopen nontrivial orbit. It lies
on an intrinsic line or a circle by C.28. There are
paths $g_s$ from the identity, formed from two factors
in $P$, which move $q$ a nonzero amount in either
chosen intrinsic direction while keeping its whole
track inside any prescribed sufficiently small
intrinsic arc about $q$.
To see this, $f_t(q)$ is not identically $q$; otherwise
every generator, and hence every element of $G$, would
fix $q$. Choose a component $(a,b)$ of
$\{t:f_t(q)\ne q\}$ and a sufficiently short initial
subinterval $[a,a+\eta]$. The path
\[
          h_s=f_{a+s\eta}f_a^{-1},\qquad 0\leq s\leq1,
                                                        \tag{C.34.2}
\]
starts at the identity, and $h_s(q)=f_{a+s\eta}(q)$.
C.28 gives compact-open continuity on an intrinsic
line; the circle restriction is continuous as well.
Consequently $h_s$ and $h_s^{-1}$ can be made uniformly
close to the identity on an intrinsic neighborhood of
$q$. Their endpoint displacements at $q$ have opposite
intrinsic signs. One of these two paths has the chosen
sign, and its track stays in the prescribed arc.
The paths also become uniformly small on each fixed
ambient compact set as $\eta\to0$.
Every such $g_s$ is a product of two elements of $P$.
For any $z\in K_n(x)=P^nx$,
\[
                        g_s(z)\in K_{n+2}(x).
                                                        \tag{C.34.3}
\]
In particular, the same compact-word argument as in
C.28 proves intrinsic continuity of its track on
every nonopen orbit, including when $z$ varies along
a compact intrinsic subarc.

\emph{A small cap cannot persist at a moving limit point.}
We give the local argument used below. Straighten a
compact arc of a limit orbit, extending it a little
at both ends, so that it is an interval of the
horizontal axis in a plane chart. Suppose compact
intrinsic arcs $A_j$ in nearby nonopen orbits lie in
strips of height tending to zero around this interval.
Let $q_j\in A_j$ attain the minimum $m_j$ of the
horizontal coordinate on $A_j$, suppose
$q_j\to q=(m,0)$, and suppose both endpoints of
$A_j$ have horizontal coordinates at least $m+\delta$
for a fixed $\delta>0$.
Take a rectangle
\[
          U=(m-\rho,m+\rho)\times(-\varepsilon,\varepsilon),
          \qquad 0<\rho<\delta/3,
                                                        \tag{C.34.4}
\]
whose closure stays in the chart. For large $j$,
$q_j$ lies in $U$, all of $A_j$ has vertical coordinate
of absolute value less than $\varepsilon/2$, and its
horizontal minimum is greater than $m-\rho/2$.
Starting at $q_j$ in each intrinsic direction towards
the two endpoints of $A_j$, take the first point
outside $U$. Both first-exit points lie on its right
edge. They are distinct points of the embedded arc.
Let $J_j$ be the closed subarc between these points.
Its interior lies in $U$, and every point of $J_j$
has horizontal coordinate at least $m_j$.

Use (C.34.2), taking the inverse when needed, to choose
a path $g_s$ whose $q$-track stays in $U$ and whose
endpoint is strictly to the left of $q$. This choice
is possible because $q$ is nonfixed. The track is a
compact subset of $U$, so joint continuity gives
$g_s(q_j)\in U$ for every $s$ and every sufficiently
large $j$.
Its intrinsic continuity from (C.34.3) now matters:
a continuous path in the intrinsic line or circle
starting at an interior point of $J_j$ cannot leave
$J_j$ without meeting one of its two endpoints.
Both endpoints lie outside the open rectangle $U$.
Thus the entire $g_s(q_j)$-track stays in $J_j$.
At $s=1$ its horizontal coordinate is at least $m_j$,
whereas
\[
        g_1(q_j)\longrightarrow g_1(q),\qquad
        \operatorname{pr}_1(g_1(q))<m=\lim_jm_j.
                                                        \tag{C.34.5}
\]
This contradiction excludes the cap.
The argument with horizontal maxima and a path moving
to the right is identical. The two first-exit points
are the explicit obstruction to escaping the cap
along another part of the same orbit.

The same argument excludes a whole Jordan-curve orbit
contained in such a shrinking strip and having a
nondegenerate interval as its limit: use a global
minimum of the horizontal coordinate on the circle,
and in its two directions take first exits towards
a point whose coordinate is bounded above that
minimum. The resulting subarc $J_j$ again has
two distinct endpoints outside $U$.

\emph{No excess in subarcs near an intrinsic line.}
Fix a point $p$ on an intrinsic line orbit and
take $x_j\to p$ in $Y$. For each fixed $n$, compactness
of $P^n$ and continuity of word evaluation give
\[
                         K_n(x_j)\longrightarrow K_n(p)
              \quad\hbox{in Hausdorff distance}.         \tag{C.34.6}
\]
The limit is a nondegenerate compact arc. Straighten
a slightly larger compact arc of the limit orbit.
Then (C.34.6) places the nearby $K_n(x_j)$ in shrinking
horizontal strips in that chart. Each $K_n(x_j)$ is
an arc, a point, or a whole circle, by C.28 and
connectedness. The whole-circle possibility is
excluded for all sufficiently large $j$ by the
preceding cap argument.

Suppose $u_j,v_j\in K_n(x_j)$ tend to $u,v\in K_n(p)$.
Write $[u_j,v_j]_j$ for the subarc contained in
$K_n(x_j)$. Then
\[
 \begin{gathered}
 z_j\in[u_j,v_j]_j,\quad z_j\longrightarrow z
             \quad\Longrightarrow\quad z\in[u,v].
                                                        \tag{C.34.7}
 \end{gathered}
\]
Indeed, an alleged limit point strictly to the left
of both $u,v$ produces a horizontal minimum of
$[u_j,v_j]_j$ whose two endpoints stay a fixed
distance to its right. A subsequence of these
minimum points converges, by (C.34.6), to a point
of the nonfixed limit orbit. This is the forbidden
cap. A point to the right is excluded by the
maximum version. The assertion includes $u=v$:
then all of the intervening subarc tends to $u$.
This proves the needed continuity of subarcs using
the path action; Hausdorff convergence alone has
not been used to infer it.

It follows that medians of three points in $K_n(x_j)$
converge to the median of their limits. Each median
lies in all three pairwise connecting subarcs.
Every subsequential limit therefore lies, by
(C.34.7), in all three corresponding limit intervals.
Their intersection is the single limit median,
including when two or three of the limiting points
coincide.

\emph{A circle limit: the precise local covering.}
Let the orbit of $p$ be a Jordan curve $C$. Choose a
collar annulus $W$ of $C$, with coordinates
$(\R/\Z)\times(-1,1)$ sending $C$ to the middle circle,
and write
\[
 \pi_W:\widetilde W=\R\times(-1,1)\longrightarrow W
                                                        \tag{C.34.8}
\]
for its universal covering. The compact family
$f_t$ preserves $C$. Therefore a sufficiently small
open disc $V$ about $p$ satisfies
$f_t(V)\subset W$ for all $t$.
The actual map
\[
 [0,1]\times V\longrightarrow W,\qquad
                         (t,x)\longmapsto f_t(x)
                                                        \tag{C.34.9}
\]
has a unique continuous lift $\widehat f_t(x)$ after
choosing a lift of the initial inclusion $V\to W$.
Its endpoint differs from its starting point by
an integer deck translation depending continuously
on $x\in V$. It is zero at $p$, by C.29's exact
plane-to-circle winding comparison, so it is zero
on all of the connected disc $V$.
Thus these are closed lifted point tracks.
For $x_j\to p$ in $V\cap Y$ put
\[
 T_j(r)=\{\widehat f_u(x_j):0\leq u\leq r\},\qquad
 T(r)=\{\widehat f_u(p):0\leq u\leq r\}.
                                                        \tag{C.34.10}
\]
Joint continuity on the compact parameter interval
shows $T_j(r_j)\to T(r)$ in Hausdorff distance
whenever $r_j\to r$. The set $T(1)$ is a
nondegenerate compact interval on the line
$\R\times\{0\}$.

We verify the intrinsic meaning of the lifted
nearby tracks. An intrinsic line orbit, or a
proper component of a circle orbit inside $W$,
has an intrinsic interval as its component through
$x_j$. Its continuous point track stays in that
component and lifts to an injectively parametrized
interval in $\widetilde W$, with embedded compact
subintervals. If the entire circle orbit lies in
$W$, it is a Jordan curve. Its winding in $W$
is $0$ or $\pm1$: the planar Jordan winding
around the complementary inner disc has exactly
these values. Its lifted components are therefore
Jordan curves in the first case and intrinsic
lines in the second.
Consequently $T_j(1)$ is a compact arc, a point,
or a whole lifted Jordan curve, with the
corresponding intrinsic topology.

The cap argument applies in this covering chart
with the original group paths. Here is the exact
local map needed to justify that use. At any
limiting extremum $q\in T(1)$ choose the small
rectangle $U$ of (C.34.4) with horizontal width
less than one, so that $\pi_W|_U$ is injective.
Choose a path (C.34.2) at $\pi_W(q)$ with the
required intrinsic sign and sufficiently small
that its images of a compact neighborhood of
$\pi_W(q)$ stay in $W$.
The map
\[
 (s,z)\longmapsto g_s(\pi_W(z))
                                                        \tag{C.34.11}
\]
on that small lifted neighborhood has a unique
continuous lift starting with $z\mapsto z$.
Its $q$-track lies in $U$ and has the required
horizontal displacement. For $q_j\to q$ its
whole $q_j$-track also stays in $U$, by compact
parameter continuity. Downstairs it is an
intrinsically continuous path in the original
orbit, by (C.34.3). Its lift cannot leave the
cap subarc without meeting one of its first-exit
endpoints on $\partial U$. Thus (C.34.5) applies
unchanged. No action of $G$ on the entire
annulus or its covering has been assumed.

This excludes a whole lifted Jordan curve for
large $j$ and proves (C.34.7), and hence median
continuity, for compact subarcs of $T_j(1)$.
These intrinsic interval operations are exactly
the ones used by C.27 and C.29. For a line or
a proper circle component the lift of its
interval is order preserving or reversing.
For an essential circle the annular covering
is its universal line covering. For a
contractible circle, the now proper arc
$T_j(1)$ lies in a single interval of its
universal line covering; a path leaving that
interval would have to cover the whole circle.
In all cases the based circle lift or intrinsic
line formula, followed by projection through
(C.34.8), gives the original $Q_{s,t}(x_j)$.
Median invariance under reversal removes any
choice of the interval orientation.

\emph{Continuity of the prefix-clamping formula.}
We can now handle a sequence
$(s_j,t_j,x_j)\to(s,t,p)$ with $p$ nonfixed.
For a line limit use the actual point tracks
$T_j(r)=\{f_u(x_j):0\leq u\leq r\}$ inside
$K_1(x_j)$. For a circle limit use (C.34.10).
In either case $T_j(1)$ is eventually an
intrinsic compact arc, its subarcs satisfy
(C.34.7), and its prefix tracks have their
stated Hausdorff limits.
Set $r_j=1-s_j$ and $r=1-s$.
Let $a_j,b_j$ be the two unordered endpoints
of the prefix arc $T_j(r_j)$, allowing equality.
Every subsequence has a further subsequence
on which both endpoints converge, say to
$a,b\in T(r)$. By (C.34.7), the Hausdorff
limit of their intervening arc is contained
in $[a,b]$. It is also exactly $T(r)$ by
the prefix-track continuity. Since $a,b$
belong to $T(r)$, their interval is contained
in $T(r)$. Therefore
\[
                         T(r)=[a,b].                    \tag{C.34.12}
\]
In particular the limiting unordered pair is
precisely the endpoint pair of the prefix
interval, including a degenerate prefix.

The formula of C.27 or its C.29 lift is
the median of $a_j,b_j$ and the current
track point at $t_j$. The median continuity
just proved and (C.34.12) give its limit as
the median of the limiting prefix endpoints
and the current point at $t$. This is
$Q_{s,t}(p)$, with projection through
(C.34.8) in the circle case.
Every subsequence has this same limit.
The domains are metrizable, so this proves
joint continuity on $[0,1]^2$ times the
nonfixed part of $Y$. C.31 supplies the
remaining fixed points.

Apply the same argument to the based inverse
loop $u\mapsto f_u^{-1}$. It is ambient-null,
generates the same group $G$, and has the
same closed locus $Y$. The exact identity
$Q[f]^{-1}=Q[f^{-1}]$ in C.31 proves joint
continuity of the inverse family in
(C.34.1). Thus all the orbitwise bijections
are homeomorphisms of $Y$. The compact-open
continuity follows from the evaluation
exponential law, since $Y$ is a closed,
locally compact subspace of the plane.
The based endpoints and equivariance are
those already proved orbitwise in C.31.

Finally, when the canonical lifted group of
a hyperbolic surface loop has no open orbit,
$Y=\widetilde U\cong\R^2$.
The obtained plane homeomorphisms commute
with the full operator lift group, including
all deck transformations. They descend by
the exact equations
\[
                  \overline Q_{s,t}\,p_U=p_UQ_{s,t}
                 \quad(p_U:\widetilde U\to U)
                                                        \tag{C.34.13}
\]
to homeomorphisms of $U$ commuting with every
original operator. Their joint continuity
and that of their inverses descend through
the open covering quotient, giving
compact-open continuity downstairs.
The endpoints in (C.34.13) are the original
loop and the identity, with both based
edges fixed. Hence this is its null-homotopy
inside the original centralizer.
At every parameter it lies in the identity
component by these same paths.
Equivalently, its lift is the canonical
one: two deck-commuting lifts of such a
map differ by a central deck transformation,
and the hyperbolic surface group has trivial
center. This proves the stated corollary
with the actual lift and descent maps.
\end{proof}

\paragraph{Remaining global extension.}
C.34 establishes the full nonopen-locus
homeomorphism family and its inverse.
C.33 provides intrinsic equivariant
contractions on each open orbit.
C.35 proves that this particular median
family need not extend equivariantly over
an adjacent open orbit. C.36 and C.37
construct compatible alternatives on the
closed union $X_2$ of open orbits with
two-dimensional proper operator image,
including every infinite accumulation.
Matching that boundary family to the
exterior of $X_2$, and treating the other
proper operator images at their frontiers,
remain the assembly requirements in
C.12, whose general status remains open.

\begin{claim}{C.35}{Proved}{An obstruction to extending the fixed median prescription}
There is a continuous $\R^2$-action on the plane and a based
ambient-null loop in its centralizer such that the C.34 median
homeomorphisms of the closed nonopen locus do not extend
equivariantly over an adjacent open orbit at one parameter.
The closed restricted operator image on that open orbit is
the free proper translation group $\R^2$.
The original loop nevertheless has an explicit equivariant
null-homotopy on the entire plane.
Thus this example refutes extension of the fixed C.34 median
prescription; it is not a failure of the level-$0$ comparison.
\end{claim}
\passage{Section 23, folios 140-141, scans 279-281; Part C(b) infinite-operator request; obstruction to extending a fixed contraction prescription in the C.12 proof route}

\begin{proof}
Use the closed simplex
\[
 \Delta=\{(x,y,z)\in\R_{\geq0}^3:x+y+z=1\}.
\]
For $(a,b)\in\R^2$ define
\[
 T_{a,b}(x,y,z)
   =\frac{(e^a x,e^b y,z)}{e^a x+e^b y+z}.
                                                        \tag{C.35.1}
\]
The denominator is positive on $\Delta$.
Direct substitution gives
$T_{a,b}T_{c,d}=T_{a+c,b+d}$ and
$T_{a,b}^{-1}=T_{-a,-b}$, so (C.35.1) is a continuous
action by homeomorphisms, fixing the three vertices and
preserving each open edge.
On the interior the coordinate homeomorphism
\[
 (x,y,z)\longmapsto(u,v)=\bigl(\log(x/z),\log(y/z)\bigr),
 \qquad
 (u,v)\longmapsto
       \frac{(e^u,e^v,1)}{1+e^u+e^v}
                                                        \tag{C.35.2}
\]
turns this action into $(u,v)\mapsto(u+a,v+b)$.
Its restricted image is therefore the closed free proper
translation group on $\R^2$.

Double $\Delta$ along its entire boundary. Since $\Delta$
is a closed disc, this double is a sphere. Apply (C.35.1)
on both copies; the formulas agree on their identified
boundary and hence give a continuous sphere action.
Remove the fixed vertex $(0,0,1)$. The remaining surface
$M$ is a plane, and the action restricts to it.
It is orientation preserving: each map preserves the
orientation on either triangle interior, and preserves
the two sheets of the oriented double.

Write $Q=\{T_{a,b}:(a,b)\in\R^2\}$ for this operator
image. Its orbits in $M$ are the two triangle interiors,
the three open edges, and the remaining two vertices.
Use the following coordinates on the edges:
\[
\begin{array}{c|c|c}
 \text{edge}&\text{coordinate in }\R&
                      \text{translation by }T_{a,b}\\ \hline
 y=0&\log(x/z)&a\\
 x=0&\log(y/z)&b\\
 z=0&\log(x/y)&a-b .
\end{array}                                               \tag{C.35.3}
\]
Thus the union $Y$ of the three edges and the two
vertices is exactly the closed nonopen orbit locus.

Let $w(t)=(a(t),b(t))$ be the piecewise linear based
loop with successive values
\[
       (0,0),\quad(1,1),\quad(1,0),\quad(0,0)
           \quad\text{at }t=0,\tfrac13,\tfrac23,1,
                                                        \tag{C.35.4}
\]
and put $f_t=T_{a(t),b(t)}$.
All its values commute with $Q$, and
\[
                  F_{\sigma,t}
                  =T_{(1-\sigma)a(t),(1-\sigma)b(t)}
                                                        \tag{C.35.5}
\]
is a based equivariant null-homotopy on the entire plane.
Continuity, including that of the inverse family, follows
from (C.35.1), which is already continuous on the compact
double and preserves the removed point.
The subgroup generated by the loop is exactly $Q$:
its first segment contains all diagonal translations
$(r,r)$ for $0\leq r\leq1$, and the differences of
points of the second segment contain all $(0,r)$
for $-1\leq r\leq1$. Integer sums and inverses give
both full coordinate axes and hence $\R^2$.
Consequently C.34 uses precisely the orbit locus $Y$
just described.

At $s=2/3$ the prefix in the C.27 formula ends at
$r=1-s=1/3$. By (C.35.3)--(C.35.4), the three
prefix translation intervals are, in their displayed
order,
\[
                         [0,1],\quad[0,1],\quad\{0\}.
                                                        \tag{C.35.6}
\]
At $t=2/3$, the current edge translations are
$(1,0,1)$. Taking the prescribed medians with
the prefix endpoints gives
\[
       \left.Q^{\mathrm{med}}_{2/3,\,2/3}\right|_Y:
           \quad\text{edge translations }(1,0,0).
                                                        \tag{C.35.7}
\]
The vertices stay fixed. These edge maps fit
continuously at the vertices, as also guaranteed
by C.34.

Every $Q$-equivariant homeomorphism from one triangle
interior to itself is a translation in (C.35.2).
Indeed, commutation with all translations gives
$h(u,v)=(u,v)+h(0,0)$. If it maps one interior to
the other, the same equation holds in the identical
coordinates on the target sheet. The sheet exchange
fixes the common boundary, so in either case any
continuous boundary extension has edge translation
parameters
\[
  R:\R^2\longrightarrow\R^3,\qquad
                    R(A,B)=(A,B,A-B).                  \tag{C.35.8}
\]
This is the exact restriction morphism.
Its image is the kernel of
\[
       \ell:\R^3\longrightarrow\R,\qquad
                         \ell(r_1,r_2,r_3)=r_1-r_2-r_3.
                                                        \tag{C.35.9}
\]
But $\ell(1,0,0)=1$, so (C.35.7) does not lie in
that image. It cannot extend equivariantly even
over a single adjacent interior. This proves
all assertions, while (C.35.5) gives the explicit
global contraction of the original loop.
\end{proof}

\paragraph{Executed endpoint certificate and scope.}
The standard-library rational calculation
\texttt{code/part\_b/c35\_median\_extension\_certificate.py}
computes the prefix extrema, the three median
values, the restriction matrix in (C.35.8),
and the nonzero obstruction (C.35.9).
Its JSON output records $(1,0,0)$ and residual $1$.
This certificate verifies the finite endpoint
calculation; the preceding simplex, doubling,
and restriction arguments prove the topological
statement.

C.34 is a homeomorphism construction on the whole
nonopen locus, and C.33 contracts the open-orbit
restrictions. C.35 proves that a global argument
must permit changing the nonopen-locus contraction
to satisfy the actual boundary restriction
morphisms of the open components. The remaining
problem is therefore the existence of compatible
equivariant contractions, not extension of the
particular median family without modification.
C.36 and C.37 meet the actual boundary equations on
the closed union $X_2$ of open orbits with
two-dimensional proper operator image, including
infinite accumulation. Extending this family
compatibly to the exterior of $X_2$ and the other
proper operator cases remains required.
The general C.12 comparison remains unresolved.

\begin{claim}{C.36}{Proved}{Compatible contractions on two-dimensional proper closures}
Let $f_t$ be a based ambient-null plane loop, let $G$ be
its generated group, and let $H$ be any group of plane
homeomorphisms commuting with every $f_t$.
For an open $G$-orbit $O$, let $Q_O$ be the closure of
the restricted setwise $H$-stabilizer, as in C.30.
For every such $O$ with $\dim Q_O=2$, there is a based
null-homotopy on the entire closure $\overline O$
through homeomorphisms of $\overline O$ that commute
with its operator stabilizer.
These closure contractions agree on every intersection
$\overline O_1\cap\overline O_2$ and are carried to one
another by $H$. They therefore give a jointly continuous
family, with jointly continuous inverse, on every
finite union of these closures.
No continuity across an infinite accumulating family
of closures is asserted.
\end{claim}
\passage{Section 23, folios 140-141, scans 279-281; Part C(b) infinite-operator request; frontier-compatible dimension-two operator closures in the C.12 proof route}

\begin{proof}
\emph{The actual quotient action on the closure.}
C.30 makes $Q_O$ a closed locally compact free proper
group of homeomorphisms of $O$, and C.33 makes it Lie.
Its dimension-two argument shows that $Q_O$ is
connected and simply transitive on $O$, and that
\[
 N_O=C_{\operatorname{Homeo}(O)}(Q_O)
                    \cong Q_O^{\mathrm{op}}             \tag{C.36.1}
\]
is also a connected, simply transitive two-dimensional
Lie group. Evaluation $N_O\to O$ is a homeomorphism.
The image of $G$ in $N_O$ is all of $N_O$: for any
$n\in N_O$, transitivity of $G$ supplies $g\in G$
with $gx=nx$ at one point $x\in O$, and the
injectivity of evaluation on $N_O$ then gives
$g|_O=n$.

Put $L=\overline G$ in $\operatorname{Homeo}^+(\R^2)$.
This is a Polish group. It preserves every open
$G$-orbit $O$. Indeed, if $g_j\to l$ with $g_j\in G$,
the identities $g_j(\partial O)=\partial O$ and
the convergent inverse family give
$l(\partial O)=\partial O$. For $x\in O$,
$lx\in\overline O$ by convergence, and
$lx\notin\partial O$ because $l$ preserves that
boundary and is injective. Thus $lx\in O$;
applying the same argument to $l^{-1}$ proves
$l(O)=O$.
The elements of $L$ still commute with $H$.
Restriction consequently gives a continuous
surjective homomorphism
\[
                   r_O:L\longrightarrow N_O.            \tag{C.36.2}
\]
It is open by the open mapping theorem for Polish
groups. Equivalently, apply the open orbit-map
principle from
\cite[Theorem~3.5]{C_GomezPinedoUzcategui2017}
to the transitive action of $L$ on the Baire
space $N_O$ through (C.36.2).

The kernel of (C.36.2) fixes $O$ pointwise, and
hence fixes $\overline O$ pointwise. Therefore
there is a well-defined action
\[
 N_O\times\overline O\longrightarrow\overline O,
       \qquad(n,z)\longmapsto lz
                 \quad\text{for any }l\text{ with }r_O(l)=n.
                                                        \tag{C.36.3}
\]
This action is jointly continuous: $r_O\times\id$
is an open quotient map, and the continuous
$L$-action factors through it. No continuous
section of (C.36.2) is required.
In fact each $n$ has a preimage in $G$, as proved
above. Hence (C.36.3) preserves every original
$G$-orbit contained in $\overline O$, and its
image on such an orbit is exactly the image of
$G$. This conclusion does not require a claim
that $L$ preserves every other $G$-orbit in the
plane.

\emph{The two-dimensional Lie contraction.}
We record the Lie groups and the maps being used.
A two-dimensional real Lie algebra is either
abelian or has a basis $X,Y$ with $[X,Y]=Y$:
in the nonabelian case its one-dimensional
derived algebra is an ideal, and scaling a
complementary basis vector gives that bracket.
The simply connected groups are respectively
$\R^2$ and the positive affine group of the
line. In coordinates $(a,b)$, the latter has
product
\[
           (a,b)(a',b')=(a+a',\,b+e^a b').
                                                        \tag{C.36.4}
\]
Its exponential map is
\[
 \exp(u,v)=\left(u,v\,\frac{e^u-1}{u}\right),
 \quad\text{with }\frac{e^u-1}{u}=1\text{ at }u=0.
                                                        \tag{C.36.5}
\]
The scalar factor is positive for every real
$u$, so this is a global diffeomorphism.
For $\R^2$ the exponential is the identity.
The affine group has trivial center, so its
connected groups with this Lie algebra have
no further covering quotients. In the abelian
case the connected group is a quotient of
$\R^2$ by a discrete central subgroup.
The compact torus is excluded here because
$N_O\cong O$ is noncompact.
Thus in every case the universal covering
group $\widetilde N_O$ has a globally defined
inverse $\log$ to its exponential.

Let $n_O(t)=f_t|_O\in N_O$.
By the restriction assertion of C.32 this loop
is null-homotopic in $\operatorname{Homeo}(O)$.
Evaluation on $O$ and the homeomorphism
$N_O\to O$ show that it is null-homotopic
in $N_O$. Its unique based lift
$\widetilde n_O(t)$ to $\widetilde N_O$
is therefore a closed loop.
With $p_O:\widetilde N_O\to N_O$ the covering
homomorphism, define
\[
 R^O_{s,t}
   =p_O\!\left(\exp\bigl((1-s)\log\widetilde n_O(t)\bigr)\right)
       \quad\text{acting on }\overline O
       \text{ by (C.36.3)}.                             \tag{C.36.6}
\]
This is a continuous family of homeomorphisms
with continuous inverse, starts with $f_t$ on
$\overline O$, ends with the identity, and
keeps both based edges fixed.
On $O$ it lies in the exact centralizer of
$Q_O$. Its commutation with the original
setwise operator stabilizer extends to
$\overline O$ by continuity and density.

\emph{The common boundary image and its topology.}
Take two of these closures and a $G$-orbit
$B$ in their intersection. It is a point,
an intrinsic line, or a circle: an open orbit
cannot lie in the boundary of another open
orbit. The point case is fixed throughout.
For a nontrivial $B$, each $N_i=N_{O_i}$
acts transitively on $B$ by (C.36.3), and
both action images are exactly $G|_B$.
We verify that this common image has a single
Lie topology, rather than identify abstract
quotients without their topology.

For $p\in B$, its stabilizer in $N_i$ is
closed. The homogeneous space
$N_i/(N_i)_p$ maps continuously and
bijectively to $B$. Its dimension is one:
dimension zero would make the connected
orbit a point, and dimension two would make
$B$ open in the plane by invariance of
domain. Its topology agrees with the
intrinsic topology of $B$.
Indeed, a compact interval in this
homogeneous one-manifold maps to an
embedded compact arc. In an intrinsic
line orbit that arc is the unique
connecting arc supplied by C.28; on a
circle the same conclusion holds locally.
The interiors of these arcs give
intrinsic interval neighborhoods, so the
bijective orbit map is a local, and
therefore global, intrinsic homeomorphism.
Thus the action is continuous also in
the intrinsic topology.

Divide $N_i$ by its closed kernel on $B$.
The resulting connected Lie group $M_i$
acts faithfully and transitively on that
one-manifold, and has dimension at most
two. In dimension one it is abelian, so
a point stabilizer is normal and is
therefore the kernel; faithfulness makes
the stabilizer trivial. Evaluation then
identifies its action with the translation
group of $\R$ or the rotation group of
$\R/\Z$.
An abelian group of dimension two would
give the same freeness conclusion and
an orbit of dimension two, so cannot
occur faithfully here.
In the remaining case $M_i$ is the
positive affine group (C.36.4).
The identity component of a point
stabilizer is a one-dimensional subgroup.
It cannot be the translation subgroup,
which is normal and would lie in the
action kernel. Every other such subgroup
is conjugate to the dilation subgroup.
Its normalizer is that dilation subgroup
itself, as is seen by conjugating its
unique fixed point on the line.
The full point stabilizer is consequently
the dilation subgroup, and the homogeneous
action is the usual positive affine
action on $\R$.

These descriptions prove topological
embedding of each actual image in
$\operatorname{Homeo}(B)$ with its
intrinsic compact-open topology.
For translations or rotations one
evaluation recovers the parameter.
For the affine group, the images of
two distinct ordered points recover
the positive slope and the intercept
continuously. They also prove that
the image groups are closed in that
homeomorphism group.
Since both images are the same set
$G|_B$, we have obtained a common
topological Lie group $M_B$, and
continuous quotient homomorphisms
\[
                       \rho_i:N_i\longrightarrow M_B.
                                                        \tag{C.36.7}
\]
The universal cover of $M_B$ is
$\R$ or the positive affine group,
and its exponential is again a
global diffeomorphism.

\emph{Equality on overlaps.}
The homomorphisms (C.36.7) lift
uniquely to based homomorphisms
\[
               \widetilde\rho_i:
                    \widetilde N_i\longrightarrow
                    \widetilde M_B.
                                                        \tag{C.36.8}
\]
The two paths
$\widetilde\rho_i(\widetilde n_i(t))$
are based lifts of the identical
loop $f_t|_B$, so covering uniqueness
makes them equal to one path
$\widetilde b(t)$.
Naturality of the Lie exponential
and its global injectivity give
\[
 \begin{split}
 d\widetilde\rho_i(\log\widetilde n_i(t))
      &=\log\widetilde b(t),\\
 \widetilde\rho_i
   \!\left(\exp\bigl((1-s)\log\widetilde n_i(t)\bigr)\right)
      &=\exp\bigl((1-s)\log\widetilde b(t)\bigr).
 \end{split}                                             \tag{C.36.9}
\]
After projection to $M_B$, this says
precisely that $R^{O_1}_{s,t}$ and
$R^{O_2}_{s,t}$ agree on $B$.
The intersection of the two closures
is a union of such $G$-orbits, so
they agree everywhere on the
intersection. This includes circle
orbits with their actual based
universal-cover lifts; no choice
of an unrecorded integer rotation
is made.

For $h\in H$, conjugation carries
$Q_O$ onto $Q_{hO}$ and $N_O$
onto $N_{hO}$, as follows from
the exact restriction of the
setwise stabilizers. This is a
topological Lie-group isomorphism.
It carries $n_O(t)$ to
$n_{hO}(t)$, since $h$ commutes
with $f_t$. Its based lift and
exponential naturality show
\[
                    R^{hO}_{s,t}
                         =hR^O_{s,t}h^{-1}.
                                                        \tag{C.36.10}
\]
Thus the compatible families have
the asserted operator equivariance.

Finally, on a finite union of
these closures, their continuous
maps agree on every overlap.
The finite closed-set pasting
lemma gives joint continuity.
Each intersection is preserved,
because its $G$-orbits are
preserved; hence the inverse
maps also agree there.
Applying the same pasting
argument to the inverse families
gives the stated family of
homeomorphisms on the finite
union. This proves the claim.
\end{proof}

\paragraph{Exact scope of the closure construction.}
The quotients (C.36.2)--(C.36.3)
provide continuity on each
individual closure, and
(C.36.7)--(C.36.9) enforce the
actual common boundary restriction
relations. In particular they
respect the three dependent edge
characters in C.35.
The present proof uses finite
closed-set pasting. C.37
subsequently proves a uniform
three-word estimate and extends
these contractions continuously,
with their inverses, to the
closed union $X_2$ of all such
open orbits, including infinite
accumulation. Extension across
$X_2$ to its exterior, including
frontier contractions for
$\dim Q_O=0,1$ and compatible
choices elsewhere on $Y$,
remains required for a global
conclusion in C.12.

\begin{claim}{C.37}{Proved}{Uniform control on all two-dimensional operator closures}
Let $f_t$, $0\leq t\leq1$, be an ambient-null based loop in
$\operatorname{Homeo}^+(\R^2)$, let $G=\langle f_t:0\leq t\leq1\rangle$,
and let $H$ be any subgroup commuting with every $f_t$.
For an open $G$-orbit $O$, let $Q_O$ be the actual closed restricted
operator image of C.30. Put
\[
 U_2=\bigcup_{\dim Q_O=2}O,\qquad X_2=\overline{U_2}.
                                                        \tag{C.37.1}
\]
The contractions of C.36 extend to a jointly continuous based
contraction through homeomorphisms of $X_2$, with jointly continuous
inverses, commuting with $H$. This includes arbitrary accumulation
of the open orbits in (C.37.1).

The uniform estimate used in the proof is the following exact
word statement. In either simply connected two-dimensional Lie
group $\R^2$ or $\operatorname{Aff}^+(\R)$, every continuous path
$k:[0,t]\to N$ from $1$ to $g$ satisfies
\[
 \exp(r\log g)
   =k(b)k(a)^{-1}
 \quad\hbox{or}\quad
 \exp(r\log g)=g\,k(b)k(a)^{-1}
 \quad(0\leq r\leq1)                                   \tag{C.37.2}
\]
for suitable $a,b\in[0,t]$. The same assertion is used in the
universal covering group when the actual group is $\R\times S^1$.
Consequently every contraction value on an open orbit in $U_2$
is represented by a word of length at most three in the original
loop and its inverse, with equality of the specified based
universal-cover lifts.
\end{claim}
\passage{Section 23, folios 140-141, scans 279-281; Part C(b) infinite-operator request; uniform dimension-two open-orbit closure step for C.12}


\begin{proof}
\emph{The three-word estimate.}
For a continuous path $\gamma:[0,1]\to\R^2$ with
$\gamma(0)=0$, $\gamma(1)=v$, one has
\[
 rv=\gamma(b)-\gamma(a)
 \quad\hbox{or}\quad
 rv=v+\gamma(b)-\gamma(a) \qquad(0\leq r\leq1).           \tag{C.37.3}
\]
Here is an exact measure-theoretic proof that applies to arbitrary
paths. For a polygonal path apply the signed-measure splitting
theorem of Grandoni--Ravi--Singh--Zenklusen
\cite[Theorem 4.7 and \S4.3.2]{C_GrandoniRaviSinghZenklusen2014}
to the two signed measures with densities $\gamma_1'$ and
$\gamma_2'$. A single circular interval has fraction $r$ of both
total measures. Its integral is either a path difference or the
total vector plus a path difference, giving (C.37.3).
Approximate an arbitrary continuous path uniformly by polygonal
paths with the same endpoints. Compactness of the two endpoints
and of the two alternatives passes the equation to the limit.
The planar case of the cited splitting proof uses at most two
cuts of an interval; the two equal-mass halves together have at
most three pieces, so one is a single interval. This is the
one-circular-interval count required here.

For $N=\R^2$, (C.37.2) is (C.37.3).
Write the affine group in the exact coordinates of C.36:
\[
 (a,b)(a',b')=(a+a',\,b+e^a b').
                                                        \tag{C.37.4}
\]
Every nonidentity $g=(a,b)$ lies on its unique exponential
one-parameter subgroup. If $a\ne0$, conjugation by
$(0,b/(e^a-1))$ sends $g$ to $(a,0)$.
The global coordinate map
\[
 (a',b')\longmapsto (a'/a,\,e^{-a'}b')
                                                        \tag{C.37.5}
\]
then sends left multiplication by $g^r$ to translation by
$(r,0)$, and sends $1,g$ to $(0,0),(1,0)$.
If $g=(0,b)$ with $b\ne0$, use $(a',b')\mapsto(b'/b,a')$
instead. Apply (C.37.3) in these coordinates.
A difference $(r,0)$ says $k(b)=g^r k(a)$; the other
alternative says $k(a)=g^{1-r}k(b)$. These are exactly
(C.37.2). The identity case uses the identity word.
This proves the estimate also for the simply connected
covering group of every group in C.36.

\emph{The compact family of actual word isotopies.}
Use $r=1-s$ for the contraction parameter, and set
\[
 R^O_{r,t}=p_O\exp\bigl(r\log\widetilde n_O(t)\bigr),
 \qquad n_O(t)=f_t|_O,                                  \tag{C.37.6}
\]
as in C.36. Consider the following compact family of actual
plane isotopies, with the compact-open topology on maps and
inverses:
\[
 W_{\epsilon,a,b,t}(u)
       = f_{ut}^{\,\epsilon} f_{ub}f_{ua}^{-1},
 \quad 0\leq a,b\leq t\leq1,\quad
       \epsilon\in\{0,1\},\quad 0\leq u\leq1 .
                                                        \tag{C.37.7}
\]
Compactness follows from continuity of the original loop and
of multiplication and inversion. Formula (C.37.2), applied to
the based lift of its prefix $[0,t]$, represents every
$R^O_{r,t}$ by the endpoint of one of (C.37.7).
More precisely, the based lift of that word isotopy in $N_O$
has endpoint
$\exp(r\log\widetilde n_O(t))$. Thus the representation retains
the lift, not only the projected homeomorphism.

\emph{The exact groups on nonopen limit orbits.}
Each $N_O$ in (C.37.6) is metabelian. Therefore the identity
\[
          [[g_1,g_2],[g_3,g_4]]=1
                                                        \tag{C.37.8}
\]
holds on $U_2$ for all $g_i\in G$, and hence on $X_2$ by
continuity. The set $X_2$ is $G$-invariant, so it is a union
of whole $G$-orbits. Its open $G$-orbits are precisely those
already in $U_2$: any other open orbit has an open
neighborhood disjoint from $U_2$.

Let $B\subset X_2$ be a nontrivial nonopen orbit, endowed
with the intrinsic line or circle topology of C.28.
The image $G_B$ in $\operatorname{Homeo}^+(B)$ is transitive,
metabelian, and path connected. For the last assertion,
every finite word has a path from the identity obtained from
the original paths and their inverses; C.28 gives intrinsic
compact-open continuity of all these finite-word paths.
We now prove the exact classification needed at $B$:
\[
 G_B\cong
 \begin{cases}
   (\R,+)\ \hbox{or }\operatorname{Aff}^+(\R),& B\text{ a line},\\
   S^1,& B\text{ a circle},
 \end{cases}                                            \tag{C.37.9}
\]
with the usual transitive actions after conjugacy.

First, a transitive abelian subgroup $A$ of
$\operatorname{Homeo}^+(B)$ has abelian closure $\overline A$,
and this closure acts freely: an element fixing one point
commutes with $A$, so fixes its whole transitive orbit.
The closure is Polish and transitive. The open orbit-map
theorem used in C.26 identifies it homeomorphically with $B$
by evaluation, since its point stabilizer is trivial.
In particular it is locally compact. Its faithful product
action on the three-manifold $B\times\R^2$ makes it a Lie
group by the same Hilbert--Smith input as C.33.
It is the one-dimensional connected group $\R$ or $S^1$,
and evaluation conjugates its action to translation.
Since $A$ was already transitive, injectivity of evaluation
also proves $A=\overline A$.

If $G_B$ is nonabelian, its derived subgroup $D$ is
nontrivial, abelian, normal, and path connected.
Path connectedness follows by taking commutator paths and
finite products of them. The common fixed set of $D$ is
$G_B$-invariant. Transitivity and nontriviality make it empty.
Every $D$-orbit is therefore a connected nontrivial interval
or circle. Such an orbit has no included endpoint, since
that endpoint would be preserved by its transitive
orientation-preserving action. All these orbits are open in
$B$; connectedness of $B$ forces $D$ to be transitive.
The preceding paragraph identifies $D$ with the full
translation group.

The orientation-preserving normalizer of translations on
the line consists exactly of $x\mapsto cx+d$, $c>0$.
Indeed normalization says $h(x+a)=h(x)+\chi(a)$;
continuity and additivity give $\chi(a)=ca$.
The multiplier image of the path-connected group $G_B$
is a connected subgroup of $\R_{>0}$. It is nontrivial in
the present case, so it is all $\R_{>0}$; because $G_B$
contains every translation, it is the full affine group.
On the circle, the orientation-preserving normalizer of
the full rotation group is just that rotation group:
the induced continuous automorphism of $S^1$ has degree
$+1$. This would contradict nonabelianness. These
arguments prove (C.37.9), including the topology of its
image in the intrinsic homeomorphism group.

For a line $B$, define $R^B_{r,t}$ by the logarithmic formula
in this simply connected group. For a circle, use the
unique based real lift of $f_t|_B$ to the cover of its
rotation group and scale that lift by $r$.
Its endpoint is zero by the exact winding comparison C.29.
On a point orbit define $R^B_{r,t}=1$.
These definitions are independent of the conjugacies in
(C.37.9), since Lie homomorphisms commute with exponential
maps and with based universal-cover lifts.
They agree with C.36 on every boundary orbit to which
that construction already applies.

\emph{Limits from the dense open union.}
Take
\[
 x_i\in O_i\subset U_2,\quad x_i\longrightarrow x\in X_2,
 \qquad (r_i,t_i)\longrightarrow(r,t).
                                                        \tag{C.37.10}
\]
Choose word isotopies $W_i$ from (C.37.7) representing
$R^{O_i}_{r_i,t_i}$ with their specified based lifts.
After a subsequence all word parameters converge.
Then $W_i$ and their inverses converge uniformly on compact
sets, uniformly in $u$, to a word isotopy $W$ from the
same family. Write $w=W(1)$ and $g=f_t$.
If $x$ is in an open orbit of $X_2$, eventually $x_i$ is
in that same open orbit, where (C.37.6) is continuous.
If $x$ is a point orbit, $w(x)=x$ for every possible
limiting word, which gives the required limit immediately.
It remains to identify $w$ on the intrinsic line or circle
$B=Gx$.

We will repeatedly use the following exact passage of
identities. For every fixed $h\in G$,
$hx_i\in O_i$ and $hx_i\to hx$. An identity of word
homeomorphisms on $O_i$, whose word isotopies converge,
therefore passes to every point $hx$ of $B$.
In particular
\[
                        [w|_B,g|_B]=1,                 \tag{C.37.11}
\]
because a fractional power commutes with its endpoint
on $O_i$.

Fix a positive integer $n$ and set $m_i=\lfloor nr_i\rfloor$.
In the universal covering group of $N_{O_i}$ the exact
residual is
\[
 \widetilde w_i^{\,n}\,
       \widetilde n_{O_i}(t_i)^{-m_i}
    =\exp\bigl((nr_i-m_i)
                  \log\widetilde n_{O_i}(t_i)\bigr).
                                                        \tag{C.37.12}
\]
Its exponent fraction lies in $[0,1)$.
By (C.37.2), the right side is again represented by one of
the compact family (C.37.7), with equality of based lifts.
Passing to a further subsequence for this fixed $n$,
the integer $m_i$ is constant, say $m_n$, and the residual
word isotopies converge to $Z_n$. Hence
\[
 (w|_B)^n(g|_B)^{-m_n}=z_n,\qquad
 |m_n-nr|\leq1,\qquad z_n=Z_n(1)|_B.                    \tag{C.37.13}
\]
All $z_n$ belong to the same compact set of restricted
three-word endpoints, independently of $n$.
No simultaneous subsequence for all $n$ is required.

For a line, (C.37.9) and (C.37.11) allow logarithms in
$\R$ or $\operatorname{Aff}^+(\R)$.
In the latter group the centralizer of every nonidentity
element is its exponential one-parameter subgroup, as is
seen directly from (C.37.4). Thus logarithms add for
commuting elements in both groups. The compact set in
(C.37.13) has bounded logarithms, giving
\[
 n\log(w|_B)-m_n\log(g|_B)=\log z_n,\qquad
 \sup_n\|\log z_n\|<\infty.
                                                        \tag{C.37.14}
\]
Divide by $n$ and let $n$ tend to infinity.
It follows that $\log(w|_B)=r\log(g|_B)$.
This includes $g|_B=1$, in which case boundedness forces
$w|_B=1$. The limiting word therefore has exactly the
required fractional action on the line.

\emph{The lifted residual bound on a circle.}
Identify the action in (C.37.9) on a circle with rotations,
and let $\theta(u)$ be the continuous real angle of
$f_u|_B$, normalized by $\theta(0)=0$.
The normalization uses rotations $\R/\Z$.
C.29 gives $\theta(1)=0$.
Every isotopy in (C.37.7), restricted to $B$, has a
canonical real endpoint angle
\[
 \epsilon\theta(t)+\theta(b)-\theta(a).
                                                        \tag{C.37.15}
\]
Let $A$ be this angle for $W$, let $T=\theta(t)$,
and let $C_n$ be the angle for $Z_n$.
Writing $M=\max_{0\leq u\leq1}|\theta(u)|$ gives
$|C_n|\leq3M$.
We must prove the real, rather than merely circle-valued,
identity
\[
                     nA-m_nT=C_n.                      \tag{C.37.16}
\]

The word isotopy
\[
 \Lambda_i(u)=W_i(u)^n f_{u t_i}^{-m_i}Z_i(u)^{-1}
                                                        \tag{C.37.17}
\]
restricts to a based loop in $N_{O_i}$.
Equation (C.37.12), with the specified lifts of $W_i$ and
$Z_i$, says that its universal-cover lift ends at the
identity. Hence that loop is null-homotopic in $N_{O_i}$.
Choose $h\in G$ with $hx\ne x$, and evaluate (C.37.17)
at the two points $x_i,hx_i$ of $O_i$.
For large $i$ they are distinct. The resulting loop in
$\operatorname{Conf}_2(\R^2)$ has difference-direction
winding zero. The word isotopies converge uniformly on
compact sets to
$\Lambda(u)=W(u)^n f_{ut}^{-m_n}Z_n(u)^{-1}$.
The two-point loops therefore converge uniformly to its
evaluation at $x,hx$. This limiting pair stays distinct,
so winding is preserved.
On the Jordan circle $B$ its winding, by C.29, is exactly
$nA-m_nT-C_n$ (up to a common orientation sign).
It is zero, proving (C.37.16).

Using (C.37.13) and $|C_n|\leq3M$ now gives
\[
        |A-rT|\leq\frac{3M+|T|}{n}\quad(n\geq1).
                                                        \tag{C.37.18}
\]
Consequently $A=rT$. This determines the based lifted
fractional action even when $f_t|_B=1$ but $\theta(t)$
is a nonzero integer. Thus no hidden winding survives
the limit from the open orbits.

\emph{Joint continuity, inverses, and equivariance.}
We have shown that every convergent subsequence of word
parameters in (C.37.10) gives the same endpoint
$R^B_{r,t}(x)$. Compactness of (C.37.7) proves
\[
          R^{O_i}_{r_i,t_i}(x_i)
                  \longrightarrow R^B_{r,t}(x).
                                                        \tag{C.37.19}
\]
This holds for every sequence from the dense set $U_2$
approaching every point of $X_2$, with parameters also
varying. It proves joint continuity on all of $X_2$.
Explicitly, for a sequence of points of $X_2$, approximate
each by a point of $U_2$ so closely that both its distance
and its image distance for that fixed parameter are
less than $1/i$. Such approximations exist by
(C.37.19) with parameters fixed. Applying (C.37.19) to
the chosen approximating sequence gives the full
continuity assertion.

The inverse loop $f_t^{-1}$ has the same group $G$,
the same sets $U_2,X_2$, and the same actual operator
images. Its construction gives precisely
$(R^B_{r,t})^{-1}$, by inversion of the based lifts.
Applying the preceding argument to it proves joint
continuity of the inverse family as well. Each member
is an orbitwise bijection, hence a homeomorphism of
$X_2$. The identities $R_{0,t}=1$, $R_{1,t}=f_t|_{X_2}$,
and $R_{r,0}=R_{r,1}=1$ give the based contraction.

Finally every $h\in H$ permutes the open $G$-orbits and
conjugates their actual closed operator images.
It therefore preserves $U_2$ and $X_2$.
C.36 gives exact transported equivariance on $U_2$.
The established joint continuity extends that equality
to its closure $X_2$. This proves all assertions.
\end{proof}

\paragraph{The remaining global boundary requirement.}
C.37 treats the whole closed set $X_2$, including infinite
accumulation of its open orbits. Its boundary contraction
is the based Lie-group fractional action forced by
(C.37.14) and (C.37.16). Extending this family to the
complement of $X_2$, compatibly with the dimension-zero
and dimension-one proper operator cases, remains a
separate requirement. Even where that complement
contains only nonopen orbits, its median contraction
from C.34 has not been identified with this boundary
family. C.35 proves that imposing the fixed median
prescription can violate the actual restriction map.
Thus the uniform estimate proved here closes the
dimension-two accumulation step, while the global
centralizer conclusion in C.12 remains open.

\begin{claim}{C.38}{Proved}{Discrete cyclic operators on nonabelian line and circle orbits}
Let $B$ be a line or circle and let
$T\leq\operatorname{Homeo}^+(B)$ be transitive and
nonabelian. Its full compact-open centralizer
\[
                  N_B=C_{\operatorname{Homeo}(B)}(T)
                                                        \tag{C.38.1}
\]
is discrete and acts freely and properly. On a line it is
trivial or infinite cyclic; on a circle it is finite cyclic.

In the plane-loop setting of C.28, let $B$ be a nonopen
orbit whose intrinsic image $G|_B$ is nonabelian, and let
$H_B=\{h\in H:hB=B\}$ be the operator stabilizer.
Its actual restriction has the exact sequence
\[
 1\longrightarrow K_B\longrightarrow H_B
       \xrightarrow{\ \rho_B\ } P_B\longrightarrow1,
 \quad
 K_B=\{h\in H_B:h|_B=1\},\quad P_B=\rho_B(H_B)\leq N_B.
                                                        \tag{C.38.2}
\]
In the intrinsic compact-open topology, $P_B$ is already
closed and discrete, acts freely and properly, and
\[
                    p_B:B\longrightarrow B/P_B
                                                        \tag{C.38.3}
\]
is an actual regular covering. Its deck transformations
are exactly $P_B$. The quotient is a line for trivial
$P_B$ on a line, and otherwise a circle. On a circle
$P_B$ is finite and the covering degree is $|P_B|$.
The kernel $K_B$ acts identically on $B$ throughout this
comparison.
\end{claim}
\passage{Section 23, folios 140-141, scans 279-281; Part C(b) infinite-operator request; exact one-dimensional operator-image reduction at the boundary remaining in C.12}


\begin{proof}
Let $L$ be the closure of $T$ in the compact-open
homeomorphism group of $B$, choose $b\in B$, and put
$J=L_b$. The group $L$ is Polish and transitive, and
$C(L)=C(T)$ by continuity of composition.
The open orbit-map theorem used in C.26 gives
$L/J\cong B$. The normalizer $N_L(J)$ is closed, and
the same exact evaluation construction as in C.30 gives
\[
 \operatorname{ev}_b:N_B
       \xrightarrow{\ \cong\ }N_L(J)/J\ \subset B.
                                                        \tag{C.38.4}
\]
The right side is closed. The inverse sends the coset
represented by $k$ to the map $lb\mapsto lkb$.
The normalizer equation makes this well defined.
Its compact-open continuity follows by the same
microtransitivity argument as C.30: convergent points
of $B$ lift locally through the open orbit map to
elements of $L$ tending to the identity.
No surface-dimensional step occurs in that argument.
The construction also proves freeness, local compactness,
and properness. Explicitly, if
$n_i b_i=c_i$, $b_i\to b$, and $c_i\to c$, choose
$l_i\to1$ with $l_i b=b_i$. Then
$n_i b=l_i^{-1}c_i\to c$, and the closed image and
homeomorphism (C.38.4) give a convergent subsequence
in $N_B$. This is the required properness test.

Every orientation-reversing homeomorphism of a line
or circle has a fixed point. Freeness therefore places
$N_B$ in $\operatorname{Homeo}^+(B)$.
Its faithful product action on the three-manifold
$B\times\R^2$ makes it a Lie group by the verified
Hilbert--Smith theorem used in C.33.
Equation (C.38.4) bounds its dimension by one.
If that dimension were one, every orbit of $N_B^0$
would be open in $B$: its one-dimensional orbit map
is injective and hence open by invariance of domain.
Connectedness of $B$ would make $N_B^0$ transitive.
Freeness would then give $N_B=N_B^0$, since an element
of $N_B^0$ with the same evaluation as an element of
$N_B$ equals it. This is the translation group $\R$
or the rotation group $S^1$ in evaluation coordinates.
Its centralizer is that same abelian group of right
translations. But $T\leq C(N_B)$, contradicting the
assumed nonabelianness of $T$.
Consequently $N_B$ is a zero-dimensional Lie group,
so is discrete.

A free proper discrete action supplies evenly covered
neighborhoods, and its orbit space is a connected
one-manifold without boundary. For a line this
quotient is a line or a circle. Since the line is
simply connected, the first case has trivial deck
group and the second has deck group $\Z$.
For a circle the quotient is a circle and the
covering has finitely many sheets; its deck group
is finite cyclic. This proves the first assertion.

In the stated plane-loop application, C.28 proves
that each operator stabilizing $B$ induces an actual
homeomorphism for its intrinsic line or circle
topology. Commutation with $G$ puts its image in
(C.38.1). Every subgroup of this discrete group is
closed; restriction of a proper free action to a
closed subgroup is proper and free. These facts
give (C.38.3). Its fibers are precisely the
$P_B$-orbits. A deck transformation is determined
by one point in the connected covering, and each
point of a fiber is obtained by exactly one
element of $P_B$, so the whole deck group is
exactly $P_B$. The same one-manifold covering
classification gives the stated quotient and
degree. Equation (C.38.2) retains the original
stabilizer and its restriction kernel.
\end{proof}

\paragraph{Topology and the remaining relative extension.}
The topology asserted for $P_B$ is its intrinsic
compact-open image topology. The operator group may
be regarded as discrete as in the original problem;
no assertion about continuity of the restriction
map from an inherited plane topology to a recurrent
intrinsic line topology is needed.
C.38 gives exact covering quotients on every
nonabelian one-dimensional orbit. It supplies no
transverse control between different orbits and no
relative extension of the C.37 family from $X_2$
to its exterior. Those requirements remain in C.12.

\begin{claim}{C.39}{Proved}{A global contraction for loops with commuting values}
Let $f_t$, $0\leq t\leq1$, be a based loop in
$\operatorname{Homeo}^+(\R^2)$ with
$f_s f_t=f_t f_s$ for every $s,t$.
Suppose the loop is null-homotopic in that ambient group.
For every subgroup $H$ commuting with all $f_t$, the loop
has a based null-homotopy in
$C_{\operatorname{Homeo}^+(\R^2)}(H)$.
The homotopy and its inverse family are jointly continuous
on the whole plane, including every change and accumulation
of orbit types.

Consequently, on every connected admissible surface $U$
of C.1, an element of the kernel of the exact inclusion map
\[
 \pi_1 C_{\operatorname{Homeo}(U)}(H)
      \longrightarrow\pi_1\operatorname{Homeo}(U)
                                                        \tag{C.39.1}
\]
that is represented by a loop with pairwise commuting values
is zero. On a hyperbolic surface every such based loop in
$C_{\operatorname{Homeo}(U)^0}(H)$ is therefore null-homotopic
in that centralizer.
\end{claim}
\passage{Section 23, folios 140-141, scans 279-281; Part C(b) infinite-operator request; complete global commuting-value subcase of the exact C.12 level-zero comparison}


\begin{proof}
\emph{The actual image on every orbit.}
Set $G=\langle f_t\rangle$. This is an abelian group.
The orbit classification C.28 applies.
On an open orbit $O$, the closure $K_O$ of $G|_O$ in
$\operatorname{Homeo}(O)$ is Polish, abelian, connected,
and transitive. It acts freely: an element fixing a point
commutes with the transitive subgroup $G|_O$ and hence fixes
every point. The Polish open orbit-map theorem identifies
evaluation $K_O\to O$ homeomorphically, with trivial
stabilizer. Thus $K_O$ is locally compact. Its faithful
action on $O\times\R$ and the Hilbert--Smith input of C.33
make it a two-dimensional abelian Lie group.
Since $O$ is a noncompact open subset of the plane, the
possibilities are $\R^2$ and $\R\times S^1$.
The original image is already the whole group:
for each $k\in K_O$, transitivity of $G|_O$ gives an
element with the same evaluation, and freeness makes
the two elements equal.

On an intrinsic line or circle orbit $B$, C.28 gives
compact-open continuity of the restricted original paths.
The same abelian closure and free evaluation argument,
now in the intrinsic topology, identifies $G|_B$ with
the full translation group $\R$ or rotation group $S^1$.
This is also the abelian case of (C.37.9).
On a point orbit the image is trivial.
Denote the resulting actual Lie image on any orbit $B$
by $K_B$. Its simply connected covering group is a vector
group of dimension at most two.

The original restricted loop has a closed based lift
$\widetilde n_B(t)$ to this vector group.
For an open orbit, C.32 gives its null-homotopy in
$\operatorname{Homeo}(O)$; evaluation and the
homeomorphism $K_O\cong O$ then make its class in
$\pi_1K_O$ zero. For a line there is no covering kernel.
For a circle the based lift closes by the exact winding
comparison C.29. Define on each orbit
\[
              R^B_{r,t}
                =p_B\bigl(r\,\widetilde n_B(t)\bigr),
                    \qquad 0\leq r,t\leq1.
                                                        \tag{C.39.2}
\]
Here $p_B$ is the Lie-group covering homomorphism,
with the trivial interpretation for a point orbit.
The action of the right side on $B$ is its actual
homeomorphism action. Its endpoints are the identity
at $r=0$ and the original loop at $r=1$; both based
edges $t=0,1$ are fixed.

\emph{One compact word family for all orbit types.}
In dimension two, C.37 represents every value of
(C.39.2) by a word isotopy
\[
 W_{\epsilon,a,b,t}(u)
        =f_{ut}^{\,\epsilon}f_{ub}f_{ua}^{-1},
 \quad \epsilon\in\{0,1\},\quad 0\leq a,b\leq t\leq1,
                                                        \tag{C.39.3}
\]
with equality of its specified based universal-cover
endpoint. In dimension one the intermediate value
theorem already finds $b\leq t$ with
$\widetilde n_B(b)=r\widetilde n_B(t)$.
This is a member of (C.39.3) with $\epsilon=0,a=0$.
The identity word handles point orbits.
Thus the same compact family of original plane isotopies
works for every orbit, with the based lifts retained.

Take arbitrary points $x_i\to x$ in the plane and
parameters $(r_i,t_i)\to(r,t)$, and write $B_i=Gx_i$,
$B=Gx$. Choose representing word isotopies $W_i$ for
$R^{B_i}_{r_i,t_i}$. Compactness of the parameter set
in (C.39.3) gives a subsequence converging, together
with inverses and uniformly in $u$, to a word isotopy
$W$. Put $w=W(1)$.
If $B$ is open, all sufficiently large $x_i$ are in
that same orbit and continuity follows from its
Lie action. If $B$ is a point, every possible word
limit fixes $x$, giving the required convergence.
It remains to identify the limiting word on a line
or circle. This argument explicitly allows the $B_i$
to have different types.

For each fixed positive integer $n$, set
$m_i=\lfloor nr_i\rfloor$. In the vector covering
group of the actual image $K_{B_i}$, one has
\[
 n\,\widetilde w_i
       -m_i\,\widetilde n_{B_i}(t_i)
       =(nr_i-m_i)\widetilde n_{B_i}(t_i).
                                                        \tag{C.39.4}
\]
The right side is represented by another word isotopy
$Z_i$ from (C.39.3), with equality of based lifted
endpoints. For this fixed $n$, pass to a further
subsequence so that $m_i=m_n$ is constant and
$Z_i\to Z_n$. Then $|m_n-nr|\leq1$.
Identities on $B_i$ pass to the whole orbit $B$ by
evaluation at $hx_i\to hx$, for any fixed $h\in G$.
In particular, with $g=f_t$ and $z_n=Z_n(1)$,
\[
             (w|_B)^n(g|_B)^{-m_n}=z_n|_B.
                                                        \tag{C.39.5}
\]

For a line, let $\theta(u)$ be the real translation
parameter of $f_u|_B$, and put
$M=\max_u|\theta(u)|$.
The real parameters $A$ of $w|_B$ and $C_n$ of
$z_n|_B$ come from their word paths:
each has the form
$\epsilon\theta(t)+\theta(b)-\theta(a)$.
Consequently $|C_n|\leq3M$, and (C.39.5) is the
real identity $nA-m_n\theta(t)=C_n$.
Division by $n$ gives $A=r\theta(t)$, as required.

For a circle, use the based real rotation lift
$\theta$ of the original loop, so $\theta(0)=\theta(1)=0$.
The word paths again specify real endpoint angles
$A,C_n$ with $|C_n|\leq3M$, but (C.39.5) initially
gives their relation only modulo integers.
The exact lift in (C.39.4) removes this ambiguity.
Indeed,
\[
       \Lambda_i(u)
         =W_i(u)^n f_{ut_i}^{-m_i}Z_i(u)^{-1}
                                                        \tag{C.39.6}
\]
is a based loop in the actual group $K_{B_i}$ and
has a closed based universal-cover lift.
It is therefore null-homotopic in that group.
Choose $h\in G$ with $hx\ne x$.
For large $i$ the two points $x_i,hx_i$ of $B_i$
are distinct. Evaluation of (C.39.6) at them has
difference-direction winding zero in
$\operatorname{Conf}_2(\R^2)$, regardless of whether
$B_i$ is an open, line, or circle orbit.
This follows from the actual null-homotopy in
$K_{B_i}$ and its continuous evaluation into the
plane; C.28 supplies that evaluation for intrinsic
line orbits.

The word isotopies and the two evaluation points
converge. Their limiting pair stays distinct.
The winding therefore remains zero. On the limiting
Jordan circle its value, by C.29, is
$nA-m_n\theta(t)-C_n$, up to a common orientation
sign. Thus the real identity
\[
 nA-m_n\theta(t)=C_n,\qquad
 |A-r\theta(t)|
       \leq\frac{3M+|\theta(t)|}{n}
                                                        \tag{C.39.7}
\]
holds for every $n$. It forces $A=r\theta(t)$.
This also treats an identity endpoint rotation
whose based angle $\theta(t)$ is a nonzero integer.
All possible subsequential word limits therefore
give the prescribed value (C.39.2).

Compactness of (C.39.3) now proves joint continuity
of (C.39.2) on the entire plane.
The inverse loop has the same orbits and image
groups, and its construction gives precisely the
inverse of (C.39.2); the same proof gives joint
inverse continuity. Each map is an orbitwise
bijection, so these are plane homeomorphisms.
Every commuting operator carries the actual
abelian Lie image on one orbit isomorphically to
that on the transported orbit. These are continuous
Lie-group isomorphisms in the specified intrinsic
or open-orbit topology. Their based covering lifts
are linear, and hence commute with multiplication
by $r$. Formula (C.39.2) is therefore exactly
$H$-equivariant. Taking $r=1-s$ proves the plane
assertion.

\emph{The precise surface lift and descent.}
Every connected admissible surface in C.1 has
universal covering $p:\widetilde U\to U$ with
$\widetilde U\cong\R^2$. Let $f_t$ represent an
element of the kernel in (C.39.1), with its
commuting-value property. Choose an ambient
filling $A:D^2\to\operatorname{Homeo}(U)$ whose
boundary is this loop and whose specified based
boundary point maps to the identity.
The actual map
\[
 D^2\times\widetilde U\longrightarrow U,\qquad
                  (z,x)\longmapsto A(z)(p(x))
                                                        \tag{C.39.8}
\]
has a unique lift equal to the identity on the
specified based fiber. Its domain is simply
connected. Lifting the inverse family as well,
and using uniqueness, proves that its slices are
homeomorphisms with jointly continuous inverses.
It is an ambient filling of a closed based
plane loop $\widetilde f_t$.

For every lift $\widetilde h$ of an operator,
$[\widetilde f_t,\widetilde h]$ is a deck
transformation. It depends continuously on $t$
and starts at the identity, so it is identically
the identity. Similarly
$[\widetilde f_s,\widetilde f_t]$ is a deck
transformation depending continuously on $(s,t)$;
it is the identity on the axes, and hence
everywhere. Thus the lifted loop has commuting
values and centralizes the full group of
operator lifts, including the deck group.
Apply the plane assertion. Its equivariant
contraction descends through $p$, with both
joint continuities, to the original centralizer
on $U$. Its slices lie in the identity component
by the displayed paths. This proves (C.39.1).
For a hyperbolic $U$ the ambient identity
component is contractible, so the ambient
kernel condition is automatic.
\end{proof}

\paragraph{Scope of the global subcase.}
This theorem treats loops with commuting values on
every orbit at once, including open orbits whose
operator image has dimension zero or one.
It does not identify every fundamental-group class
of an arbitrary centralizer with such a loop.
No deformation of a general centralizer loop to
this subcase has been constructed here.
C.12 remains open for general loops.

\paragraph{The earlier abelian closed subgroup cannot be a centralizer.}
Fix the specific embedding $j:B\hookrightarrow A^0$ of C.26, where
$B=L^1([0,1],\Z)/\Z$, $A=\operatorname{Homeo}(U)$ and $U$ is
hyperbolic. For every $H\leq A$ commuting with $j(B)$, the induced map
$\pi_1(B)\to\pi_1(C_A(H),1)$ is zero. Indeed the displayed generator
of $\pi_1(B)=\Z$ has pairwise commuting values and is ambient-null,
so C.39 contracts its image; the induced homomorphism therefore
vanishes on the whole group. The based contraction lies in $A^0$,
so the same conclusion holds with target $\pi_1(C_{A^0}(H),1)$.
Consequently $j(B)$ cannot equal either of these centralizers,
since its topological-group identification with $B$ induces an
isomorphism on the nonzero fundamental group. This rules out that
exact realization as a zero-level counterexample.

\begin{claim}{C.40}{Proved}{The locally compact loop-closure case}
Let $U$ be any connected admissible surface of C.1, let
$A=\operatorname{Homeo}(U)$ with its compact-open group
topology, and let $H\leq A$ be arbitrary. Put $C=C_A(H)$.
For a based loop $f:[0,1]\to C$, form the actual closed
subgroup
\[
 L_f=\overline{\langle f_t:0\leq t\leq1\rangle}^{\,A}
\]
with the topology induced from $A$. Every class in the
kernel of
\[
 i_*:\pi_1(C,1)\longrightarrow\pi_1(A,1)^H
                                                        \tag{C.40.1}
\]
that has a representative with locally compact $L_f$
is zero. Here $i:C\hookrightarrow A$ is the inclusion;
its induced map lands in the displayed fixed subgroup
because each $h\in H$ fixes the entire representative
loop under conjugation.

In particular, the same conclusion holds for loops
contained in any closed locally compact subgroup of
$C$. If $C$ itself is locally compact, (C.40.1) is
injective. On a hyperbolic surface every based loop
in $C_{\operatorname{Homeo}(U)^0}(H)$ with locally
compact $L_f$ is null-homotopic in that centralizer.
\end{claim}
\passage{Section 23, folios 140-141, scans 279-281; Part C(b) infinite-operator request; exact inclusion map from the centralizer fundamental group to the H-invariant ambient fundamental group}


\begin{proof}
\emph{The exact subgroup and its Lie topology.}
Centralizers are closed in the Hausdorff topological
group $A$. Consequently $L=L_f$ is a closed subgroup
of $C$. The subgroup generated by the loop is
path connected: each finite product of loop values
and their inverses can be joined to the identity by
the corresponding product of paths. Its closure $L$
is connected. Thus no identity-component replacement
or change of topology is being made.
The topology on $L$ is the subspace topology from
$A$, and its action on $U$ is continuous and faithful.

Suppose $L$ is locally compact. The product action
\[
 L\times(U\times\R)\longrightarrow U\times\R,
       \qquad (l,(x,a))\longmapsto(lx,a)
                                                        \tag{C.40.2}
\]
is a faithful continuous action on a connected
three-manifold. The verified Hilbert--Smith theorem
used in C.33, \cite[Theorem 1.5 and the introductory
reduction]{C_Pardon2013}, therefore makes $L$ a Lie
group with precisely this topology. It is a connected
Lie group. All homotopies constructed inside $L$ below
are consequently homotopies through actual elements
of $C$, continuous in its given topology.

\emph{A based deformation into a compact subgroup.}
Borel's principal theorem gives a connected maximal
compact subgroup $K$ and closed one-parameter
subgroups $H_1,\ldots,H_d$, each isomorphic to $\R$,
such that multiplication is a homeomorphism
\[
 K\times\R^d\longrightarrow L,\qquad
 (k,a_1,\ldots,a_d)\longmapsto
       k\,h_1(a_1)\cdots h_d(a_d).
                                                        \tag{C.40.3}
\]
This is the actual product assertion in
\cite[\S1, p.~271; \S\S3--4, pp.~277--278]{C_Borel1952}:
the product notation is explicitly defined there
by uniqueness and continuous dependence of all
factors. The maps $h_j:\R\to H_j$ are topological
group isomorphisms with $h_j(0)=1$.
The identity has the unique coordinates $(1,0,\ldots,0)$.
Thus
\[
 D_s\bigl(kh_1(a_1)\cdots h_d(a_d)\bigr)
       =k h_1((1-s)a_1)\cdots h_d((1-s)a_d)
                                                        \tag{C.40.4}
\]
is a continuous strong deformation retraction
$D:L\times[0,1]\to L$ onto $K$, fixing the identity.
In particular $D_s(f_t)$ is a based homotopy in
$C$ from $f$ to a based loop $\beta$ in $K$.

\emph{A commuting representative of the same based class.}
Average an inner product on the Lie algebra of $K$
over its compact adjoint action. The resulting
positive definite inner product is adjoint
invariant, so left translation gives a bi-invariant
Riemannian metric on $K$. Its pullback to the
universal covering group
$p:\widetilde K\to K$ is also bi-invariant.
For left-invariant vector fields its Levi-Civita
connection satisfies
\[
               \nabla_XY=\tfrac12[X,Y].
                                                        \tag{C.40.5}
\]
Indeed the Koszul formula and adjoint invariance
give this identity directly. Hence all geodesics
through the identity are the one-parameter
subgroups, and their left translates are defined
for all real time. This proves geodesic completeness
of the lifted metric. The Hopf--Rinow theorem
therefore joins the identity to every point of
$\widetilde K$ by a geodesic.

Lift $\beta$ from the identity to a path
$\widetilde\beta$ with endpoint $z\in\ker p$.
Choose $X$ in the Lie algebra of $\widetilde K$
with $\exp_{\widetilde K}(X)=z$, using the geodesic
just obtained. The paths $\widetilde\beta(t)$
and $\exp_{\widetilde K}(tX)$ have the same
endpoints in the simply connected space
$\widetilde K$, so they are homotopic relative
to both endpoints. Projection through $p$ is a
based homotopy in $K$ from $\beta$ to
\[
             \gamma(t)=p\bigl(\exp_{\widetilde K}(tX)\bigr).
                                                        \tag{C.40.6}
\]
This is a based loop, since $p(z)=1$.
It has pairwise commuting values because it is
the restriction of a one-parameter subgroup.
Every step from $f$ to $\gamma$ stays in $C$.

Now let $[f]$ be in the kernel of (C.40.1).
The inclusion of the fixed subgroup
$\pi_1(A,1)^H$ into $\pi_1(A,1)$ is injective,
so $f$ is ambient-null. The based centralizer
homotopy just constructed makes $\gamma$
ambient-null as well. C.39 applies on this
same admissible surface to this actual commuting
loop and the same operator subgroup $H$.
It supplies a based contraction of $\gamma$
inside $C$. Concatenating the two based
homotopies contracts $f$ inside $C$.

A loop contained in a closed locally compact
subgroup $M\leq C$ has $L_f\subseteq M$ closed;
it is therefore locally compact. Taking $M=C$
proves the assertion about a locally compact
centralizer. In the hyperbolic case
$\operatorname{Homeo}(U)^0$ is contractible by
the earlier surface theorem, so the ambient
kernel condition holds for every based loop.
The constructed contraction lies in the identity
path component because it is a based homotopy;
it consequently lies in the displayed
centralizer inside $\operatorname{Homeo}(U)^0$.
\end{proof}

\paragraph{Exact remaining scope.}
C.40 proves a reduction to commuting values for
the specified locally compact closed subgroup.
No such based homotopy has been constructed for
a general loop whose generated closure is not
locally compact. The unrestricted comparison
C.12 remains open. In particular, the proof
does not replace the induced topology on that
closure by a different locally compact topology.

\begin{claim}{C.41}{Proved}{A dense operator orbit settles the zero-level kernel}
Let $U$ be a connected admissible surface of C.1, let
$A=\operatorname{Homeo}(U)$ with its compact-open group
topology, and let $H\leq A$ have a dense orbit in $U$.
Then the exact inclusion-induced comparison
\[
 \pi_1(C_A(H),1)\longrightarrow\pi_1(A,1)^H
                                                        \tag{C.41.1}
\]
is injective. In particular, on a hyperbolic surface
$\pi_1(C_{A^0}(H),1)=0$. No local compactness assumption
is imposed on $H$ or on the subgroup generated by
the values of the loop.
\end{claim}
\passage{Section 23, folios 140-141, scans 279-281; Part C(b) infinite-operator request; exact centralizer inclusion on fundamental groups for a dense operator orbit}


\begin{proof}
\emph{The plane argument.}
Let $f_t$ be a based ambient-null loop of plane
homeomorphisms centralizing a subgroup $H$ with
a dense orbit $D$. Put $G=\langle f_t\rangle$.
Since the loop is based, its maps preserve
orientation, and the orbit results C.28--C.37
apply. Every $h\in H$ transports $G$-orbits
according to the exact identity
\[
                       h(Gx)=G(hx).
                                                        \tag{C.41.2}
\]
If there is no open $G$-orbit, C.34 already
gives the required global equivariant contraction.

Otherwise let $V$ be the nonempty union of the
open $G$-orbits. It is open and $H$-invariant.
Density makes $D$ meet $V$, and invariance then
gives $D\subseteq V$. Consequently $V$ is dense
in the plane.

Fix an open $G$-orbit $O$ and choose $y\in D\cap O$.
Let $H_O=\{h\in H:hO=O\}$ be its operator
stabilizer. The actual orbit equality is
\[
                     H y\cap O=H_O y.
                                                        \tag{C.41.3}
\]
Indeed, $hy\in O$ implies that the two $G$-orbits
$hO=G(hy)$ and $O$ meet, hence are equal. The
reverse inclusion is immediate. Since $Hy=D$
is dense, (C.41.3) is dense in $O$.

Retain the exact restriction homomorphism
$H_O\to\operatorname{Homeo}(O)$, with its
pointwise-fixing kernel. Let $Q_O$ be the
closure of its image in that homeomorphism
group. C.30 proves that $Q_O$ is locally
compact and acts freely and properly on $O$;
C.33 makes it a Lie group of dimension at
most two. Properness makes every $Q_O$-orbit
closed. In particular $Q_Oy$ is a closed
subset of $O$ containing the dense set
$H_Oy$, so
\[
                         Q_Oy=O.
                                                        \tag{C.41.4}
\]
Thus this free proper action is transitive.
Its orbit map $Q_O\to O$ is a homeomorphism,
and $\dim Q_O=2$.

This holds for every open $G$-orbit. In the
notation of C.37, therefore,
\[
 U_2=V,\qquad
 X_2=\overline{U_2}=\R^2.
                                                        \tag{C.41.5}
\]
The contraction proved there is defined on
the whole plane, with jointly continuous
inverses, and centralizes every element of
the original $H$. This settles both alternatives
of the plane argument.

\emph{The actual surface comparison.}
Let a based loop $f_t$ represent a kernel class
in (C.41.1). The inclusion of the fixed
subgroup in $\pi_1(A,1)$ is injective, so
choose an ambient filling
$F:D^2\to A$ of that loop.
For the universal covering
$p:\widetilde U\to U$, every admissible
$\widetilde U$ is a plane. Lift the actual map
\[
 D^2\times\widetilde U\longrightarrow U,\qquad
           (z,x)\longmapsto F(z)(p(x))
                                                        \tag{C.41.6}
\]
with the identity prescribed on the based
boundary fiber. Its domain is simply connected.
As in C.39, lifting the inverse family and
using uniqueness gives an ambient filling by
plane homeomorphisms with jointly continuous
inverses. Its boundary is a closed based
plane loop $\widetilde f_t$.

Let $\widetilde H$ be the group of all lifts
of elements of $H$, including the deck group.
For every $\widetilde h\in\widetilde H$,
the commutator
$[\widetilde f_t,\widetilde h]$ is a deck
transformation depending continuously on $t$
and equal to the identity at $t=0$.
It is therefore always the identity.
The lifted loop centralizes precisely the
same operator lifts required for descent.

Choose $x\in U$ with dense $Hx$ and
$\widetilde x\in p^{-1}(x)$. One has the
exact equality
\[
       \widetilde H\,\widetilde x=p^{-1}(Hx).
                                                        \tag{C.41.7}
\]
One inclusion follows by projection. For the
other, choose a lift of the specified
$h\in H$ and then a deck transformation
taking its value at $\widetilde x$ to the
specified point over $hx$.
The right side is dense because $p$ is open:
every nonempty open set upstairs projects
to a nonempty open set meeting $Hx$.

Apply the plane result to $\widetilde H$.
Its contraction commutes with all deck
transformations and all operator lifts.
The equation
$\overline R_{s,t}p=pR_{s,t}$ therefore
defines a based contraction downstairs
inside $C_A(H)$. The open covering quotient
and its local sections give joint continuity
of these maps and their inverses.
It contracts the original class in the
source of (C.41.1), proving injectivity.
For a hyperbolic $U$, the ambient identity
component is contractible, so its based
centralizer loops all satisfy the kernel
condition. The based contraction stays in
$A^0$.
\end{proof}

\paragraph{Remaining scope.}
C.41 settles every operator action with a
dense orbit. It does not assert that a
general operator group has such an orbit
or supply an equivariant decomposition into
subsets on which this argument applies.
The general operator problem without a dense orbit
in C.12 remains open.

\begin{claim}{C.42}{Proved}{The actual boundary quotient has proper line components upstairs}
Let $G$ be generated by a based plane loop as in C.28,
let $O$ be an open $G$-orbit, and put
$L=\overline G\leq\operatorname{Homeo}(\R^2)$.
Fix $x\in O$ and $J=L_x$. For
$z\in\operatorname{Fix}(J)$ whose $G$-orbit $B=Gz$
is a nontrivial intrinsic line or a circle, the formula
\[
 \Phi_z(lx)=lz
\]
defines the actual continuous open quotient
$O=L/J\to B=L/L_z$, with the intrinsic topology on $B$.
For the universal coverings $p:\widetilde O\to O$
and $q:\R\to B$, every lift
$F:\widetilde O\to\R$ of $\Phi_zp$ is open and onto,
and every connected component of every fiber of $F$
is a properly embedded line in the covering plane
$\widetilde O$.
Every nondegenerate compact connected component of
a fiber of $\Phi_z$ is a circle.
Whole fibers of $\Phi_z$ can be disconnected, including
for quotients obtained as limits of the original
commuting operators.
\end{claim}
\passage{Section 23, folios 140-141, scans 279-281; Part C(b) infinite-operator request; exact boundary quotients in the remaining centralizer extension problem}


\begin{proof}
\emph{The quotient and its topology.}
Every element of $L$ preserves $O$: forward and inverse
limits preserve its boundary, while a forward limit
sends an interior point into $\overline O$ and cannot
send it to that boundary. The Polish orbit-map theorem
used in C.26 and C.30 therefore makes $l\mapsto lx$
an open quotient $L\to O$.
If $l_1x=l_2x$, then $l_2^{-1}l_1\in J$ fixes $z$,
so $\Phi_z$ is well defined and continuous.
More precisely, the continuous action map
\[
 \operatorname{Fix}(J)\times L\longrightarrow\R^2,
 \qquad (z,l)\longmapsto lz
\]
descends through the open quotient
$\operatorname{id}\times(L\to O)$.
Consequently
\[
 \operatorname{Map}_L(O,\R^2)
       \xrightarrow[\cong]{\,\operatorname{ev}_x\,}
       \operatorname{Fix}(J)
                                                        \tag{C.42.1}
\]
is a homeomorphism in the compact-open topology:
equivariance proves bijectivity, and the displayed
joint action proves continuity of the inverse.
For $l\in L$, choose $g\in G$ with $gx=lx$.
Then $g^{-1}l\in J$, hence $lz=gz$.
Thus the range is exactly $Lz=Gz=B$.

The map into an intrinsic line is continuous in its
intrinsic topology. Indeed, the image of a compact
coordinate disk in $O$ is a compact Peano continuum
in $B$. By C.28 it contains neither a Jordan curve
nor a triod, and hence is a point or a compact arc;
on such an arc the intrinsic and ambient topologies
agree. These disks give neighborhoods at every point
of $O$. For a circle the two topologies already agree.
Each $l\in L$ restricts to an intrinsic homeomorphism
of $B$, because compact arcs are carried to compact
arcs. For $b=gz$ one has
\[
                  lb=\Phi_z(lgx).
                                                        \tag{C.42.2}
\]
This makes every point evaluation of the restriction
$L\to\operatorname{Homeo}(B)$ continuous.
Finite ordered partitions of compact intervals show
that pointwise convergence of monotone homeomorphisms
to a homeomorphism is uniform on compact intervals;
apply the same argument to inverses. Finite cyclic
partitions give the circle assertion. Thus restriction
is continuous in the compact-open group topology.
All these restrictions preserve orientation, since
the original loop does and the restriction map is
continuous.

Apply the Polish orbit-map theorem to this transitive
action on the Polish one-manifold $B$. It identifies
$B$ with $L/L_z$ in its stated topology. Since
$J\subseteq L_z$, the map $\Phi_z$ is the open quotient
$lJ\mapsto lL_z$.
In particular, if operators $h_n$ commute with $G$
and $h_nx\to z$, then they commute with $L$ and
(C.42.1) gives $h_n|_O\to\Phi_z$ uniformly on compact
subsets of $O$.

\emph{Both covering maps and the compatible group.}
The universal cover of the connected open planar
surface $O$ is a plane. Fix a lift
\[
 F:\widetilde O\longrightarrow\R,\qquad
                      qF=\Phi_zp.
                                                        \tag{C.42.3}
\]
For a line $B$, take $q$ to be an identification.
Local covering charts and openness of $\Phi_zp$
show that $F$ is open.
Consider the following subgroup, with the topology
induced from the indicated product:
\[
\begin{split}
\mathcal L=\{(l,a,b)\in
 L\times\operatorname{Homeo}(\widetilde O)
       \times\operatorname{Homeo}^+(\R):\quad
 pa&=(l|_O)p,\\
 qb&=(l|_B)q,\qquad Fa=bF\}.
\end{split}
                                                        \tag{C.42.4}
\]
The three equations are closed conditions, as may
be checked by evaluation on countable dense sets.
Hence $\mathcal L$ is a Polish group, and its
actions by $a$ and by $b$ are continuous.
Every lift $a$ of $l|_O$ has a unique compatible
lift $b$ of $l|_B$: the maps $Fa$ and $bF$ lift
the same map to $B$, and the target deck translation
is chosen to make them agree at one point.
Uniqueness of lifts then gives equality everywhere.
For a line target there is no target deck ambiguity.

The action of $\mathcal L$ on $\widetilde O$ is
transitive. Choose $l$ carrying the projections of
two specified points to one another and adjust a
source lift by a deck transformation to carry the
specified points themselves.
The image of $F$ is a nonempty open interval $I$.
It is invariant under every compatible $b$.
If $B$ is a circle and $I$ has a finite endpoint,
every increasing $b$ fixes that endpoint, so its
image under $q$ would be fixed by every element
of the transitive $L$-action on $B$. This is
impossible. Therefore $I=\R$; for a line target
surjectivity is already immediate from (C.42.3).

\emph{A component that separates the plane.}
Fix $t\in\R$ and write $E=F^{-1}(t)$.
Choose $u,v$ with $F(u)<t<F(v)$. The closed set
$E$ separates $u$ from $v$. Frankel's planar
separation lemma \cite[Lemma 2.7, p.~1258]{C_Frankel2015}
gives a connected component $C$ of $E$ separating
them. The lemma applies to unbounded closed sets.
For completeness, take the complementary domain
$U$ of $E$ containing $u$, and then the component
$V$ of $\widetilde O\setminus\overline U$ containing
$v$. The planar Brouwer property makes
$\partial V$ connected. It lies in $E$ and
separates $u,v$, so the component of $E$ containing
$\partial V$ has the required property.

Let $\mathcal K$ be the closed subgroup of
$\mathcal L$ whose $b$-coordinate fixes $t$.
It acts transitively on $E$: transitivity on the
covering plane finds an element carrying any
two points of $E$ to one another, and
$Fa=bF$ forces $b(t)=t$.
Thus all components of $E$ are ambiently
homeomorphic to $C$, and all separate the plane.
The stabilizer $\mathcal K_C$ is a closed Polish
subgroup and acts transitively on $C$.
Closedness follows by applying forward and inverse
limits to the closed set $C$.

\emph{Accessibility from each fixed complementary domain.}
Let $D$ be any component of
$\widetilde O\setminus C$.
Its stabilizer in $\mathcal K_C$ is open:
for a fixed $d\in D$, every sufficiently small
group element sends $d$ into $D$ and therefore
sends the whole complementary domain onto $D$.
The open orbit-map theorem for
$\mathcal K_C\curvearrowright C$ makes every
orbit of this open stabilizer open in $C$.
Since $C$ is connected, the stabilizer acts
transitively on $C$.

There is an accessible point of $\partial D$:
an interior point and a closest boundary point
give a straight end-cut in a coordinate copy of
the plane. Its boundary endpoint belongs to $C$.
Transporting this end-cut by the domain
stabilizer proves that every point of $C$
is accessible from this same $D$.
This holds for every complementary domain.

There are exactly two such domains.
There are at least two because $C$ separates.
For three domains, choose three distinct points
$c_1,c_2,c_3\in C$. Inside each domain, join
end-cuts to these points to an embedded tree
with these three boundary leaves and one
interior branch point. Such a tree is obtained
by joining two end-cuts to an arc and stopping
a third connecting arc at its first intersection.
The three trees have disjoint interiors and
share just their specified boundary leaves.
Their union embeds $K_{3,3}$ in the plane,
contradicting the bipartite planar Euler
inequality $9\leq2\cdot6-4$.

\emph{The proper line conclusion.}
If $C$ is unbounded, put
$C^*=C\cup\{\infty\}$ in the sphere; if $C$
is bounded, put $C^*=C$.
This is a compact connected proper subset
with exactly two complementary domains.
Every point except possibly the added
point $\infty$ is accessible from both.
Choose a stereographic pole in one complementary
domain. Deleting that pole leaves that domain
connected, and accessibility at its boundary
is unaffected, since end-cuts can be chosen
to avoid the pole. In the resulting plane
both inaccessible subsets of $C^*$ have
at most one point.
Polulyakh's converse Jordan theorem
\cite[Theorem 1.1 and Corollary 1.1]{C_Polulyakh2000}
therefore makes $C^*$ a Jordan curve.
The bounded alternative is impossible:
the real-valued open map $F$ is constant
on that Jordan curve, and its restriction
to the bounded disk either has an interior
extremum or is constant on a nonempty
open set. Both contradict openness.
Thus $C$ is a Jordan curve with the
compactification point removed, a properly
embedded line. Transitivity of $\mathcal K$
transfers this conclusion to every component
of every fiber.

Downstairs, a fiber of $\Phi_z$ is homogeneous
under $L_z$, and each of its component
stabilizers is transitive by the same
orbit argument. A nondegenerate compact
component is therefore homogeneous by
ambient neighborhood homeomorphisms.
Brechner--Erkama
\cite[Theorem 2]{C_BrechnerErkama1979}
makes it a Jordan curve.

\emph{An exact disconnected-fiber example.}
On $r>1$ in the plane put $s=\log(r-1)$.
Let $\beta_b$ be the time-$b$ flow of
$\sin\delta\,\partial_\delta$ on $\R$.
It fixes $\pi\mathbb Z$, is $2\pi$-periodic
as a lifted circle map, and
$|\beta_b(\delta)-\delta|\leq\pi$.
Define
\[
\begin{aligned}
 A_a(r,\theta)&=(r,\theta+a) &&(r\leq1),\\
 A_a(r,\theta)&=(1+e^a(r-1),\theta+a)&&(r>1),\\
 B_b|_{\{r\leq1\}}&=\operatorname{id},\\
 B_b(s,\theta)&=(\theta+\beta_b(s-\theta),\theta)&&(r>1).
\end{aligned}
                                                        \tag{C.42.5}
\]
The last first coordinate is $s$.
Its expression is independent of the chosen
representative of $\theta$.
The bound on $\beta_b-\operatorname{id}$
gives continuity at $r=1$ and properness at
infinity for $B_b$ and $B_{-b}$.
The $A$-formulas have these properties directly.
They are continuous orientation-preserving
one-parameter actions and commute, since
$A_a$ preserves $s-\theta$.

The open spiral strip
\[
 (t,\delta)\longmapsto
  (1+e^{t+\delta},\,t\bmod2\pi),
 \qquad t\in\R,\quad0<\delta<\pi,
                                                        \tag{C.42.6}
\]
is an open orbit $O$. This parametrization is
injective because an angular ambiguity is
$2\pi k$, while two $\delta$-values differ
by less than $\pi$.
In these coordinates $A_a$ translates $t$
and $B_b$ sends $\delta$ to $\beta_b(\delta)$.
The resulting group is $\R^2$, simply
transitive on $O$ and closed in the plane
homeomorphism group. Indeed, the open-orbit
preservation argument at the start of the proof
shows that any homeomorphic limit preserves
$O$, using both the forward and inverse limits.
Its evaluation at $(0,\pi/2)$ therefore remains
inside $O$ and recovers finite parameters
$a=t$ and $b=\log\tan(\delta/2)$ continuously.
The parameter limits recover the limiting
plane homeomorphism, proving closedness.
A triangular parameter loop
$0\to(1,0)\to(0,1)\to0$ generates this
group and contracts by scaling both parameters.

For the actual commuting operator subgroup
$H=\langle A_{-2\pi}\rangle$, its powers
$h_n=A_{-2\pi n}$ converge on compact
subsets of $O$ to
\[
 \Phi(t,\delta)=(1,t\bmod2\pi).
                                                        \tag{C.42.7}
\]
This is exactly the quotient associated to
the limit of $h_n(0,\pi/2)$.
Its fiber over $\theta_0$ is the disjoint
union of the lines $t=\theta_0+2\pi k$,
$k\in\mathbb Z$. Its real lift is $t$.
All maps and the ambient contraction have
been specified, so this example concerns
disconnected whole fibers rather than a
nonzero centralizer fundamental-group class.
The executed symbolic certificate
\path{code/part_b/c42_spiral_fibre_certificate.py}
checks the flow law and differential equation
in the exact half-angle coordinate,
the $A$--$B$ commutation identities, the
integer-labelled whole fibers and their real
lift, and the based triangular parameter loop.
It uses no numerical approximations; the
topological conclusions are proved above.
\end{proof}

\paragraph{Scope of the reduction.}
C.42 proves the component statement for the
specified real-valued lift. It does not
identify a lifted component with a whole
downstairs connected component without
checking the actual covering restriction.
It does not yet give local product charts
for the component decomposition, simultaneous
charts for several boundary quotients, or
the frontier-compatible contraction needed
in C.12. Those global extension questions
remain open.

\begin{claim}{C.43}{Proved}{The lifted boundary quotient is a product projection}
For the actual quotient and covering maps in C.42, every
fiber of the real lift
$F:\widetilde O\to\R$ is connected.
There is a homeomorphism
\[
 T:\R\times\R\longrightarrow\widetilde O,
 \qquad F(T(u,t))=t.                                    \tag{C.43.1}
\]
In particular the components in C.42 form a foliation
with genuine local product charts. The original map
$\Phi_z:O\to B$ also has local product charts with its
specified target coordinate. The construction treats
one quotient map at a time; its auxiliary flow and
the homeomorphism $T$ carry no asserted equivariance
under the original operator family.
\end{claim}
\passage{Section 23, folios 140-141, scans 279-281; Part C(b) infinite-operator request; actual boundary-quotient maps in the remaining relative centralizer extension}


\begin{proof}
Write $X=\widetilde O$ and let $\mathcal L$ be the
closed Polish group (C.42.4), acting transitively on
the covering plane $X$. For $v\in X$, let $C_v$ be
the connected component of $F^{-1}(F(v))$ containing
$v$. By C.42 these sets partition $X$ into properly
embedded lines. The equation $Fa=bF$ for an element
$(l,a,b)\in\mathcal L$ implies
\[
                         a(C_v)=C_{a(v)}.                \tag{C.43.2}
\]
Thus this is the actual component partition preserved
by the given compatible action.

\emph{Regularity and local products.}
Let $I\subset C_v$ be a compact arc containing $v$,
and let $\epsilon>0$. In a fixed metric on the
covering plane, the condition
\[
       \sup_{w\in I}d(a(w),w)<\epsilon
\]
defines an identity neighborhood in $\mathcal L$.
The Polish open orbit-map theorem makes its image
at $v$ a neighborhood of $v$. Consequently, for
every sufficiently near $v'$, some such $a$ sends
$v$ to $v'$. Restricting $a$ to $I$ gives an
$\epsilon$-close embedding of $I$ into $C_{v'}$
carrying $v$ to $v'$. This verifies precisely the
regularity condition of Aarts--Oversteegen
\cite[Definition 5.3]{C_AartsOversteegen1991}.
Their local-section and flowbox theorems
\cite[Theorems 5.6--5.7]{C_AartsOversteegen1991}
therefore provide open product neighborhoods
$S\times(-1,1)\to X$, with each vertical plaque
contained in one $C_v$.

The transverse factor is an interval locally.
Indeed it is metrizable, and its product with
an interval is a surface without boundary.
The elementary product-cancellation argument
proved in C.33 applies: local path connectedness
descends to $S$, embedded open arcs in $S$ are
open by invariance of domain, and an endpoint
would have zero local second homology in the
product. Thus $S$ is locally a one-manifold
without boundary. Shrink to an interval and
write one of these charts as
\[
 \chi:J\times(-1,1)\longrightarrow V\subset X.
\]
Because every plaque lies in a fiber,
$F\chi(s,u)=f(s)$ for a continuous function $f$.
For each open $W\subset J$, the set
$\chi(W\times(-1,1))$ is open in $X$, so the
openness of $F$ proves that $f(W)$ is open.
A continuous open map between intervals is
strictly monotone: failure of strict monotonicity
produces a local maximum, local minimum, or a
constant interval, each contradicting openness.
It follows that $f$ is a homeomorphism onto
an open interval. Replacing the transverse
coordinate by $f(s)$ gives actual charts
in which $F$ is the projection.

\emph{The component quotient and its topology.}
Let $Y$ be the set of components $C_v$, equipped
with the quotient topology, and let
$\pi:X\to Y$ be the quotient map.
This map is open. To verify this directly,
let $U\subset X$ be open and let $v$ belong
to the saturation of $U$. Choose
$w\in C_v\cap U$ and the compact arc in $C_v$
joining $v$ to $w$. Apply the regularity
argument above with an error small enough
that the image of $w$ remains in $U$.
For every $v'$ in a neighborhood of $v$
the resulting image arc joins $v'$ to a
point of $U$ in $C_{v'}$. The saturation
of $U$ is therefore open, which proves
the assertion about $\pi$.

In a chart where $F$ is projection, distinct
plaques have distinct $F$-values. They therefore
belong to distinct global components.
The transversal maps homeomorphically onto
$\pi(V)$: its inverse is the restriction
of the continuous factor map $\overline F$
in
\[
 X\xrightarrow{\ \pi\ }Y
   \xrightarrow{\ \overline F\ }\R,
 \qquad F=\overline F\pi.                               \tag{C.43.3}
\]
Here $\pi(V)$ is open by the preceding
paragraph, and the forward inverse in
the chosen target coordinate is the
continuous transversal section.
Thus $Y$ is a one-manifold without boundary
in the sense that it is locally an interval;
no Hausdorff property has been assumed.
The same charts show that $\overline F$
is a local homeomorphism.

The quotient $Y$ is connected and second
countable: an open quotient carries a
countable base to a countable base.
Every element of $\mathcal L$ induces a
homeomorphism of $Y$ by (C.43.2), and
transitivity on $X$ makes $Y$ homogeneous.
Baillif--Gabard's theorem
\cite[Theorem 1.1, proved in Section 4]
{C_BaillifGabard2008} now proves that
$Y$ is Hausdorff.
The topological hypotheses here are exact:
second countability gives Lindel\"ofness,
and the actual quotient charts give the
local Euclidean structure required in
that theorem.

One can also see the separation mechanism
directly from their proof. In a locally
Hausdorff space, Zorn's lemma supplies
a maximal Hausdorff open set through
each point; maximality makes it dense.
A countable subcover exists by
Lindel\"ofness. A manifold is locally
Baire and hence Baire, so those dense
open sets have a common point. That
point is separated from every other
point inside one member of the subcover.
Homogeneity transports this property
to every point.

A connected second-countable Hausdorff
one-manifold without boundary is a line
or a circle. The circle alternative is
impossible because $\overline F$ is open
and has image $\R$ by C.42. Thus $Y$
is a line. Its local homeomorphism
$\overline F:Y\to\R$ is globally
strictly monotone: the local direction
cannot change on a connected line.
It is onto, hence a homeomorphism.
Each fiber of $F$ is consequently one
of the proper lines $C_v$.

\emph{An exact global product for this map.}
Orient each leaf using the orientation
of the covering plane together with
the increasing transverse coordinate
$F$. On an overlap of two
orientation-preserving charts $(u,F)$,
the leaf-coordinate transition is
increasing. This supplies the coherent
orientation required for the
Aarts--Oversteegen global-flow theorem
\cite[Main Theorem 0.5, its proof in
Section 4, and Remark 4.4]
{C_AartsOversteegen1991}.
It gives a continuous flow
\[
 R:\R\times X\longrightarrow X
\]
whose orbits are exactly the specified
leaves. The flow is free, since those
leaves are lines rather than circles,
and a nontrivial period would make
the corresponding orbit a circle.

Take any small transversal section
$\sigma:I\to X$ with $F\sigma(t)=t$.
The map
\[
 \Theta_\sigma:\R\times I\longrightarrow F^{-1}(I),
 \qquad (u,t)\longmapsto R(u,\sigma(t))                 \tag{C.43.4}
\]
is continuous and bijective: every
fiber is exactly one flow orbit, and
freeness gives uniqueness of $u$.
Both spaces are surfaces, so invariance
of domain makes (C.43.4) a homeomorphism.
These are actual local trivializations
of a principal $\R$-bundle over $\R$.

For completeness, its global section
can be constructed with the group
coordinate retained. Choose a locally
finite interval cover $(I_i)$ with
sections $\sigma_i$ and a subordinate
partition of unity $(\rho_i)$ whose
supports lie in $I_i$. On overlaps,
let the continuous real transition
functions be defined by
\[
             \sigma_j(t)=R(a_{ij}(t),\sigma_i(t)).
\]
They obey $a_{ij}=a_{ik}+a_{kj}$.
Set $b_i(t)=\sum_j\rho_j(t)a_{ij}(t)$
on $I_i$, extending the individual
terms by zero outside their supports.
The sum is locally finite and
$b_i=a_{ik}+b_k$ on overlaps.
Consequently
$\sigma(t)=R(b_i(t),\sigma_i(t))$
is a well-defined continuous global
section with $F\sigma(t)=t$.
The same map as (C.43.4), now over
all $\R$, gives (C.43.1).

Finally, the identities
$qF=\Phi_zp$ from C.42, together with
local charts for the two covering
maps $p$ and $q$, carry the projection
charts to the original map $\Phi_z$.
This proves its stated local product
property with the actual target
coordinate.
\end{proof}

\paragraph{Scope of the reduction.}
The auxiliary flow in the proof is used
to trivialize one specified quotient.
It has not been shown to commute with
the original operators, preserve another
boundary quotient, or extend across the
frontier of $O$. C.43 therefore strengthens
the structure of the necessary boundary
maps without completing the simultaneous
relative extension in C.12.

\begin{claim}{C.44}{Proved}{Exact downstairs forms of a nontrivial boundary quotient}
For the actual quotient $\Phi_z:O\to B$ of C.42,
the open planar orbit $O$ is a plane or an annulus.
Choose a coordinate $B=\R$ when $B$ is a line,
or a coordinate $B=\R/\Z$ when it is a circle,
and use the corresponding target covering in
the definition of $F$. There are source
homeomorphisms giving one of the following
normal forms for the actual map:
\[
\begin{array}{c|c|c|c}
 O & B & \Phi_z & \text{whole fiber}\\ \hline
 \R\times\R & \R & (u,t)\mapsto t & \R\\
 \R\times\R & \R/\Z & (u,t)\mapsto t+\Z
       & \coprod_{\Z}\R\\
 (\R/\Z)\times\R & \R & ([u],t)\mapsto t
       & \R/\Z\\
 (\R/\Z)\times\R & \R/\Z & ([u],t)\mapsto t+\Z
       & \coprod_{\Z}(\R/\Z)\\
 \R\times(\R/m\Z) & \R/\Z & (u,[t]_m)\mapsto[t]_1
       & \coprod_{j=0}^{m-1}\R,\quad m\geq1 .
\end{array}                                             \tag{C.44.1}
\]
In the last row the actual induced image
$\Phi_{z*}(\pi_1O)$ is $m\Z\subset\pi_1B=\Z$;
in all other rows that image is zero.
Every downstairs fiber component is a
properly embedded line or a circle, and
the source covering restricts onto each
such component as the line covering
specified by these formulas.
\end{claim}
\passage{Section 23, folios 140-141, scans 279-281; Part C(b) infinite-operator request; exact quotient and covering constraints in the remaining frontier-extension problem}


\begin{proof}
Let $\Gamma$ be the deck group of
$p:\widetilde O\to O$.
Use C.43 to identify the covering plane
with $\R_u\times\R_t$ so that $F(u,t)=t$.
The equation $qF=\Phi_zp$ implies, for
each $\gamma\in\Gamma$,
\[
 F\gamma=F+\eta(\gamma),                                \tag{C.44.2}
\]
where $\eta(\gamma)=0$ for a line target
and $\eta(\gamma)\in\Z$ for a circle
target. For the circle, the difference
is a continuous integer-valued function
on the connected covering plane and
is therefore constant. Composition
shows that $\eta:\Gamma\to\Z$ is a
homomorphism. Its value is exactly the
winding number of the image under
$\Phi_z$ of the source loop whose lift
ends at $\gamma$ times its initial
point. Thus its image is the actual
induced fundamental-group image.

Write $K=\ker\eta$. Every deck map
preserves the orientation pulled back
from the planar surface $O$. It has
the form
\[
 \gamma(u,t)=
       (h_{\gamma,t}(u),\,t+\eta(\gamma)),                \tag{C.44.3}
\]
and each $h_{\gamma,t}$ is increasing:
the base coordinate is increasing,
and reversing the fiber direction
would reverse the surface orientation.
The functions in (C.44.3) and their
fiberwise inverses are jointly
continuous, by the deck map and its
inverse.

\emph{The exact kernel on a real fiber.}
The action of $K$ on any fiber
$\R\times\{t\}$ is free and proper.
Freeness follows from deck freeness.
For properness, compact subsets of
this fiber are compact in the
covering plane, and the deck action
has finite transporters of compact
pairs. The restriction is faithful,
since a deck transformation fixing
one point is the identity.
An increasing free proper group on
a line is trivial or infinite cyclic.
Here is the direct argument.
Its orbit of a point is closed and
discrete. In the nontrivial case
choose the element taking that point
to its nearest orbit point on the
right. Its iterates tend to the
two ends, because a finite limiting
point would be fixed. Every other
orbit point is one of these iterates:
an intervening point could be
translated back between the original
point and its nearest right neighbor.
Freeness then identifies every
element with the corresponding power.
Consequently
\[
                          K=1\quad\hbox{or}\quad K\cong\Z.
                                                               \tag{C.44.4}
\]

Suppose $K=\langle k\rangle\cong\Z$,
choosing the generator to move one
fiber point to the right.
Since the displacement
$h_{k,t}(u)-u$ is continuous and
nowhere zero on the connected plane,
it is positive everywhere.
Put $k_t=h_{k,t}$.
For $v=n+r$, where $n\in\Z$ and
$0\leq r\leq1$, define
\[
 H(v,t)=\bigl(k_t^n(r\,k_t(0)),\,t\bigr).
                                                               \tag{C.44.5}
\]
At integral $v$ the two formulas
agree. The intervals with endpoints
$k_t^n(0)$ exhaust the line in both
directions, so $H$ is a continuous
bijection of the plane. Invariance
of domain makes it a homeomorphism.
It preserves $F$ and satisfies
\[
                        H(v+1,t)=kH(v,t).
                                                               \tag{C.44.6}
\]
Thus the actual intermediate quotient
is
\[
 \widetilde O/K\cong(\R/\Z)\times\R,
 \qquad F_K([v],t)=t.                                  \tag{C.44.7}
\]

\emph{The compactness obstruction retaining the quotient map.}
The case $K\cong\Z$ cannot coexist
with a nonzero image of $\eta$.
Indeed that image would be $m\Z$
for some $m\geq1$. The quotient
$\Gamma/K$ is then infinite cyclic,
and its chosen generator acts on
(C.44.7) by a circle homeomorphism
in the first coordinate and by
$t\mapsto t+m$ in the second.
The compact strip
$(\R/\Z)\times[0,m]$ maps onto the
full quotient by $\Gamma/K$:
each point can be translated into
that strip using its second coordinate.
It follows that
\[
     O=\widetilde O/\Gamma
       =(\widetilde O/K)/(\Gamma/K)
\]
is compact, contrary to its being
a nonempty open subset of the plane.
Therefore $K\cong\Z$ forces
$\eta=0$ and $\Gamma=K$.
Equations (C.44.5)--(C.44.7) give
the third and fourth rows of
(C.44.1), according to the target.

If $K=1$ and $\eta=0$, then
$\Gamma=1$, so C.43 gives the
first two rows.

\emph{The remaining nonzero covering degree.}
It remains that $K=1$ and
$\eta(\Gamma)=m\Z$, with $m\geq1$.
Then $\Gamma$ itself is infinite
cyclic. Choose its generator
$\gamma$ with $\eta(\gamma)=m$,
and write
\[
                   \gamma(u,t)=(h_t(u),t+m).
                                                               \tag{C.44.8}
\]
We give a coordinate change
preserving the real target
coordinate and removing $h_t$.
For $0\leq t\leq m$ set
\[
 \rho_t(u)=
          \left(1-\frac{t}{m}\right)u
                   +\frac{t}{m}h_0(u).
                                                               \tag{C.44.9}
\]
Each $\rho_t$ is an increasing
homeomorphism of the line, with
$\rho_0=\id$ and $\rho_m=h_0$.
Extend to all real $t$ in both
directions by the exact recursion
\[
                     \rho_{t+m}=h_t\circ\rho_t.
                                                               \tag{C.44.10}
\]
The endpoint identity makes the
extension jointly continuous at
every multiple of $m$.
The map
$H_1(u,t)=(\rho_t(u),t)$ is a
continuous bijection of the plane
and hence a homeomorphism.
Equations (C.44.8)--(C.44.10) give
\[
              H_1^{-1}\gamma H_1(u,t)=(u,t+m).
                                                               \tag{C.44.11}
\]
The quotient is therefore
$\R\times(\R/m\Z)$, and the
original map, still satisfying
$qF=\Phi_zp$, is precisely
$(u,[t]_m)\mapsto[t]_1$.
This is the final row of (C.44.1).

All changes of source coordinates
preserved $F$; the target was
changed only by its initially
specified coordinate.
The whole-fiber descriptions now
follow directly. For a circle
target, levels $F=t_0+j$ are
indexed by $j\in\Z$. The deck
group identifies their labels
by $\eta(\Gamma)$ and acts
within each level by $K$.
Thus the component labels are
$\Z/\eta(\Gamma)$ and each
component is $\R/K$ in the
explicit coordinates above.
For a line target there is
only one real level and the
component is again $\R/K$.
These formulas also prove
surjectivity of the restricted
source covering onto each
downstairs component and give
its deck group. In particular,
they establish the component
assertion without assuming
component-surjectivity for an
arbitrary restricted covering.
\end{proof}

\paragraph{Scope of the reduction.}
C.44 classifies each of the actual
quotient maps separately, including
its whole fibers and covering
labels. Several boundary points
can impose several such maps on
the same open orbit. No common
product coordinates or homotopy
preserving all those maps have
been constructed here. That
simultaneous compatibility and
frontier extension remain part
of C.12.

\begin{claim}{C.45}{Proved}{Two actual boundary quotients give explicit planar coordinates}
Let $\Phi_1:O\to B_1$ and $\Phi_2:O\to B_2$
be two actual nontrivial boundary quotients
of the same open orbit as in C.42.
Use the same source universal covering
$p:X\to O$ and the real target coverings
$q_i:\R\to B_i$, and choose lifts
$q_iF_i=\Phi_i p$.
Exactly one of the following alternatives holds.
Either there is a homeomorphism
$f:\R\to\R$ with
\[
                         F_2=fF_1,                     \tag{C.45.1}
\]
or the pair $(F_1,F_2)$ is a homeomorphism
onto an open domain $\Omega\subset\R^2$
whose horizontal and vertical sections
are intervals and whose two projections
are onto $\R$.
In the second alternative, explicit
independent changes of the two real
coordinates, and a possible exchange
of them, carry $\Omega$ onto one of
\[
 \R^2,\qquad
 \{(x,y):x<y\},\qquad
 \{(x,y):x<y<x+1\}.                                    \tag{C.45.2}
\]
The full increasing coordinate-pair
groups preserving these three domains
are, respectively,
\[
 \operatorname{Homeo}^+(\R)^2,\qquad
 \{(g,g):g\in\operatorname{Homeo}^+(\R)\},\qquad
 \{(g,g):g\in\operatorname{Homeo}^+(\R),\
                         g(x+1)=g(x)+1\}.               \tag{C.45.3}
\]
The actual compatible lifted action
maps into the corresponding group.
These statements retain both original
covering equations and make no
identification of their downstairs
component labels.
\end{claim}
\passage{Section 23, folios 140-141, scans 279-281; Part C(b) infinite-operator request; simultaneous actual boundary-quotient constraints in the remaining extension problem}


\begin{proof}
\emph{A common compatible group.}
Use the closed subgroup
\[
\begin{split}
 \mathcal P=\{(l,a,b_1,b_2)\in
 L\times\operatorname{Homeo}(X)
       \times\operatorname{Homeo}^+(\R)^2:\qquad&\\
 pa&=(l|_O)p,\\
 q_i b_i&=(l|_{B_i})q_i,\qquad
 F_i a=b_iF_i\quad(i=1,2)\}.
\end{split}                                             \tag{C.45.4}
\]
The equations are closed in the
product of the indicated Polish
groups. For every lift $a$ of
$l|_O$, each target lift $b_i$
exists uniquely after its deck
translation is chosen at one point,
exactly as in C.42. Consequently
$\mathcal P$ acts transitively on
the source covering plane.
Both $F_i$ are open onto $\R$
with connected proper-line fibers,
by C.43.

Fix $t\in\R$ and set
$C=F_1^{-1}(t)$.
The closed subgroup of $\mathcal P$
with $b_1(t)=t$ acts transitively
on this whole line. Its action
on the interval $I=F_2(C)$ is
also transitive, by the equality
$F_2a=b_2F_2$. All these target
maps are increasing.
A nondegenerate interval admitting
such a transitive action has no
included endpoint, so it is an
open interval.
The Polish open orbit-map theorem,
applied to the actions on $C$
and $I$, makes
$F_2|_C:C\to I$ the open map
between their homogeneous-space
quotients: the point stabilizer
in $C$ is contained in the
stabilizer of its image in $I$.
It is therefore a strictly
monotone homeomorphism.

If instead $I$ is a point,
transitivity on the source
plane transports this constancy
to every $F_1$-fiber. There is
then a function $f$ satisfying
(C.45.1). Since $F_1$ is an
open quotient, $f$ is continuous.
For every open $W\subset\R$,
\[
              f(W)=F_2(F_1^{-1}(W))
\]
is open. Thus $f$ is open and
onto, hence a homeomorphism
of the line. This is the first
alternative.

In the nonconstant case,
$(F_1,F_2)$ is injective,
since $F_2$ is injective on
each $F_1$-fiber. Invariance
of domain makes it a
homeomorphism onto an open
domain $\Omega$.
Each vertical section is the
open interval $I$ just proved.
Interchanging the two indices
gives the horizontal statement;
the constant alternative in
that direction would contradict
the already proved two-dimensional
open image.
Surjectivity of the two $F_i$
gives both full coordinate
projections.

\emph{The endpoint maps are actual homogeneous quotients.}
The coordinate action
$(b_1,b_2)$ of $\mathcal P$
preserves $\Omega$ and is
transitive on it.
Write its vertical intervals as
\[
 \Omega_x=(\alpha(x),\beta(x)),
 \quad -\infty\leq\alpha(x)<\beta(x)\leq+\infty.
                                                               \tag{C.45.5}
\]
The two endpoints satisfy the
exact equations
\[
 \alpha(b_1x)=b_2(\alpha(x)),\qquad
 \beta(b_1x)=b_2(\beta(x)),                               \tag{C.45.6}
\]
where increasing homeomorphisms
fix the appropriate infinite
endpoints.
Transitivity of the first
coordinate action implies that
$\alpha$ is either everywhere
finite or identically $-\infty$,
and that $\beta$ is either
everywhere finite or identically
$+\infty$.

Each finite endpoint map is a
homeomorphism of the real line.
For example, choose $x_0$ with
$\alpha(x_0)\in\R$. The formula
\[
 \alpha(b_1x_0)=b_2(\alpha(x_0))
\]
descends the continuous second
evaluation through the open
first orbit map
$\mathcal P\to\R$.
It proves continuity of
$\alpha$ in its actual
topology. Its image is
nonempty and invariant under
the transitive second coordinate
action, hence is all $\R$.
The stabilizer of $x_0$ in
the first action fixes
$\alpha(x_0)$ in the second
by (C.45.6). The two open
orbit maps identify $\alpha$
with the associated open
coset-space map.
It follows that $\alpha$ is
continuous, open, and onto,
and therefore strictly monotone
and a homeomorphism.
The same proof applies to
finite $\beta$.
This verifies the endpoint
topology without assuming
regularity of a boundary graph.

\emph{The three domains and the coordinate changes.}
If both endpoints are infinite,
$\Omega=\R^2$.
If precisely one endpoint is
finite, use its inverse as
the second coordinate.
The defining inequality becomes
$x<y$ or $y<x$, according to
which endpoint and monotonicity
occur. In the latter case
exchange the two coordinates.

Suppose both endpoints are
finite. They have the same
monotonicity: an increasing
homeomorphism and a decreasing
homeomorphism cannot satisfy
$\alpha(x)<\beta(x)$ for all
real $x$, by their limits at
the two ends.
When both increase, replace
$y$ by $\alpha^{-1}(y)$ and put
$\tau=\alpha^{-1}\beta$.
When both decrease, replace
$y$ by $\beta^{-1}(y)$ and put
$\tau=\beta^{-1}\alpha$.
In either case the domain is
now exactly
\[
                        x<y<\tau(x),
 \qquad \tau\in\operatorname{Homeo}^+(\R),\
                          \tau(x)>x.                    \tag{C.45.7}
\]
Choose $x_0$ and any increasing
homeomorphism
$h:[x_0,\tau(x_0)]\to[0,1]$.
Extend it by
$h(\tau^n x)=h(x)+n$.
The iterates of $\tau$ exhaust
the two ends, since $\tau$
has no fixed point, and the
endpoint formulas agree.
Thus $h$ is an increasing
homeomorphism of $\R$ with
$h\tau h^{-1}(x)=x+1$.
Applying $h$ to both
coordinates in (C.45.7)
gives the final domain in
(C.45.2).
These are the asserted
explicit source coordinate
maps; a reversal in one
coordinate still conjugates
increasing coordinate actions
to increasing actions.

Finally, an increasing
coordinate pair preserving
$x<y$ must preserve its
boundary graph $y=x$,
which says $g=h$.
For $x<y<x+1$, its lower
and upper endpoint graphs
are separately preserved,
which says
$h(x)=g(x)$ and
$h(x+1)=g(x)+1$.
Conversely these identities
preserve the indicated
domains exactly.
The unrestricted plane case
has no additional equation.
This proves (C.45.3).

All induced actions used
above arise from (C.45.4).
An element acting trivially
on the covering plane may
still act elsewhere under
its $l$-coordinate; no
faithfulness of that global
coordinate has been imposed.
The original equations
$q_iF_i=\Phi_i p$ remain
in force, so the separate
deck shifts and component
labels of C.44 are retained.
\end{proof}

\paragraph{Scope of the reduction.}
C.45 gives the relation between
two specified actual boundary
maps, including all coordinate
changes. A third quotient can
impose an additional equation
on the group in (C.45.3).
Preserving all such equations
and obtaining a global
frontier-compatible contraction
are not consequences of this
pairwise classification.

\begin{claim}{C.46}{Proved}{Two boundary foliations admit a compatible intrinsic contraction}
Use the actual based null-homotopic plane loop $f_t$, its generated group $G$, the
open orbit $O$, and $L,x,J$ of C.42. For every
$z\in\operatorname{Fix}(J)\cap\partial O$ with nontrivial
intrinsic line or circle orbit, lift its quotient $\Phi_z$
to the common source universal covering $p:X\to O$ as in
C.42. Let $\mathscr F$ be the set of fiber partitions of
these real lifts; changing their target coordinates or
their lift constants does not change this set.
For $|\mathscr F|=2$, the loop $f_t|_O$ has a based
compact-open null-homotopy
\[
 K_{s,t}\in\operatorname{Homeo}(O),\qquad
 K_{0,t}=f_t|_O,\quad K_{1,t}=K_{s,0}=K_{s,1}=\id.
                                                        \tag{C.46.1}
\]
It commutes with every plane homeomorphism commuting
with all $f_t$ and preserving $O$. For every indicated
point $z$, it also satisfies the exact quotient relation
\[
                 \Phi_z K_{s,t}=Q^{B_z}_{s,t}\Phi_z,
                  \qquad B_z=Gz,                       \tag{C.46.2}
\]
where $Q^{B_z}$ is the natural median contraction of the
original restricted loop on that intrinsic line or
circle, as in C.27--C.29. Thus all the specified
boundary quotient relations hold simultaneously.
The asserted homotopy is intrinsic to $O$; continuity
at its ambient frontier is a further assertion.
\end{claim}
\passage{Section 23, folios 140-141, scans 279-281; Part C(b), infinite-operator zero-level comparison; simultaneous boundary relations in C12}


\begin{proof}
Choose representatives $F_1,F_2:X\to\R$ of the two
partitions. C.45 supplies a homeomorphism from $X$
onto one of its three domains $D$, with these two
coordinates changed only by independent line
homeomorphisms and a possible exchange. Lift the
isotopy $f_t|_O$ from the identity to $a_t$ on $X$.
For each original boundary target, lift $f_t|_{B_i}$
to its based real path $b_{i,t}$. Covering uniqueness
gives
\[
                       F_i a_t=b_{i,t}F_i.              \tag{C.46.3}
\]
Each $b_{i,t}$ is increasing. For a line target its
endpoint is the identity. For a circle target C.29
identifies its endpoint translation with the winding
of the original ambient-null plane loop and proves
that translation zero. Hence both $b_{i,1}$ are the
identity. Since the pair $(F_1,F_2)$ is injective,
$a_1=\id$. We have obtained an actual based loop in
the corresponding coordinate-pair group (C.45.3),
without discarding a source or target deck shift.

In these coordinates write $a_t=(g_{1,t},g_{2,t})$.
Apply C.27 to each coordinate loop and set
\[
\begin{split}
 m_{i,r}(u)&=\min_{0\leq v\leq r}g_{i,v}(u),\qquad
 M_{i,r}(u)=\max_{0\leq v\leq r}g_{i,v}(u),\\
 R_{i,s,t}(u)&=\operatorname{med}
       \bigl(m_{i,1-s}(u),g_{i,t}(u),M_{i,1-s}(u)\bigr),\\
 \widetilde K_{s,t}(u,v)&=(R_{1,s,t}(u),R_{2,s,t}(v)).
\end{split}                                             \tag{C.46.4}
\]
The coordinate maps and their inverses are jointly
continuous by C.27 and C.31. They preserve the actual
domain $D$. For $D=\R^2$ this follows directly from
their being line homeomorphisms. For $D=\{u<v\}$,
C.45 gives $g_{1,t}=g_{2,t}$, so the two new coordinate
maps are identical and preserve the strict inequality.
For $D=\{u<v<u+1\}$ the same equality holds and each
$g_{i,t}$ commutes with translation by one. Minima,
maxima and the median preserve that equation, so the
new equal coordinate maps preserve both inequalities.
The inverse family has the same property. Formula
(C.46.4) is therefore a based null-homotopy in
$\operatorname{Homeo}(X)$ with the asserted endpoint
identities.

We establish the operator comparison before descending.
Let $h$ commute with all $f_t$ and preserve $O$.
It commutes with $L$ by continuity of composition and
inversion. The stabilizers of $x$ and $hx$ in $L$
are equal: an element fixes one precisely when it
fixes the other, by commutation with $h$ and $h^{-1}$.
Choose $a_h\in L$ with $a_hx=hx$. The exact stabilizer
identity gives $a_hJa_h^{-1}=J$. For every specified
$z$ and every $l\in L$, evaluation gives
\[
 \Phi_z h(lx)=\Phi_z(l a_hx)=l a_hz
                         =\Phi_{a_hz}(lx).             \tag{C.46.5}
\]
The point $a_hz$ lies in $\operatorname{Fix}(J)\cap
\partial O$: normalizing $J$ proves the first property,
and $a_h$ preserves $O$ and its boundary. Its orbit
is the same $Lz=Gz$ as that of $z$, so it has the
same intrinsic type. Thus pullback by $h$ permutes
the full set $\mathscr F$. The inverse $h^{-1}$
gives the inverse permutation. After lifting to $X$,
the two sides of (C.46.5) differ at most by the
specified target deck translation, which leaves
their fiber partitions identical.

Every lift $\widetilde h$ of $h|_O$ consequently
permutes the two coordinate foliations of $D$.
It is a coordinate pair of line homeomorphisms,
possibly followed by exchange of the coordinates.
For clarity, preservation of a fiber partition
defines its map on the corresponding real quotient;
the open quotient projection proves continuity of
that map and of its inverse. Surjectivity of both
coordinate projections then makes it a line
homeomorphism. It may be increasing or decreasing.
These quotient maps give exactly the stated
coordinate form of $\widetilde h$.

Moreover $\widetilde h$ commutes with $a_t$.
Their commutator projects to the identity on $O$
and is a continuous deck-valued function of $t$,
equal to the identity at $t=0$. The discrete deck
group makes it the identity for every $t$.
C.27's naturality under every line homeomorphism,
together with the identical construction in both
coordinates, now gives
\[
      \widetilde h\widetilde K_{s,t}\widetilde h^{-1}
                           =\widetilde K_{s,t}.         \tag{C.46.6}
\]
In particular this holds for every source deck
transformation, which is a lift of $h=\id$.
Consequently the homotopy descends to $K_{s,t}$
on $O$. The evaluation maps and their inverse
maps descend continuously through the open covering
quotient. This proves compact-open continuity,
the endpoint identities, and the asserted
commutation with the original operators.

Finally, take any other specified $z$. Its lifted
fiber partition equals one of the two chosen ones.
The first alternative of C.45 identifies its real
map as a line homeomorphism composed with that
coordinate. Naturality of the median formula under
this exact reparametrization proves
\[
 F_z\widetilde K_{s,t}=\widetilde Q^{B_z}_{s,t}F_z.
\]
For a circle target the formula commutes with its
integer deck translations, so it descends; for a
line target it already is the stated target map.
Using $q_zF_z=\Phi_zp$ proves (C.46.2).
No ambient limit of points approaching $\partial O$
has been inferred from these quotient equations.
\end{proof}

\begin{claim}{C.47}{Proved}{Three independent boundary quotients force a Lie action}
Let $\Phi_i:O\to B_i$, $1\leq i\leq3$, be three actual boundary
quotients of the same open orbit as in C.42. Use their common source
universal covering and real target coverings, with
\[
 p:X\longrightarrow O,\qquad
 q_iF_i=\Phi_i p,\qquad F_i:X\longrightarrow\R .
                                                        \tag{C.47.1}
\]
Suppose the three lifted fiber partitions are pairwise distinct;
equivalently, none of the pairs satisfies the first alternative
of C.45. The full group
\[
 \mathcal W=
 \{a\in\operatorname{Homeo}(X):
       F_i a=b_iF_i\text{ for some }
       b_i\in\operatorname{Homeo}^+(\R),\ 1\leq i\leq3\}
                                                        \tag{C.47.2}
\]
is a closed, locally compact subgroup in the actual compact-open
topology. It is a Lie group of dimension two or three and acts
transitively on $X$. If any pair $(F_i,F_j)$ has image different
from $\R^2$, this action is simply transitive and has dimension two.
In particular, the closure in $\operatorname{Homeo}(X)$ of the
compatible lifted loop group is locally compact.
All three covering equations in (C.47.1) and the kernel of the
global-to-restricted action are retained. The conclusion concerns
the action on the source covering plane; it does not assert an
extension of an intrinsic homotopy across the frontier of $O$.
\end{claim}
\passage{Section 23, folios 140-141, scans 279-281; Part C(b) infinite-operator request; simultaneous actual boundary-quotient constraints in the remaining frontier extension problem}


\begin{proof}
\emph{The actual group topology.}
Include all three target lifts in the single closed Polish group
of tuples $(l,a,b_1,b_2,b_3)$ satisfying
\[
 pa=(l|_O)p,\qquad q_i b_i=(l|_{B_i})q_i,\qquad
 F_i a=b_iF_i \quad(1\leq i\leq3).
                                                        \tag{C.47.3}
\]
For each source lift, the three compatible target lifts are
obtained exactly as in C.42--C.45. This group acts transitively
on $X$, so its image in $\mathcal W$ does also.
Each $b_i$ in (C.47.2) is unique because $F_i$ is onto.
C.43 supplies a continuous section
$\sigma_i:\R\to X$ of each $F_i$. Consequently
\[
              b_i=F_i a\sigma_i,\qquad
              b_i^{-1}=F_i a^{-1}\sigma_i.               \tag{C.47.4}
\]
These formulas are continuous for uniform convergence on compact
sets, for maps and inverses.

They also prove closedness. If $a_n\in\mathcal W$ converges to a
homeomorphism $a$ together with its inverse, (C.47.4) gives locally
uniform limits $b_i$ and $d_i$ of the target maps and their
inverses. The identities $b_{i,n}d_{i,n}=d_{i,n}b_{i,n}=\id$
pass to the limit, since the images of each compact set remain
in a common compact set. Thus $d_i=b_i^{-1}$; each $b_i$ is an
increasing homeomorphism, and $F_i a=b_iF_i$. Hence $\mathcal W$
is a closed Polish subgroup of $\operatorname{Homeo}(X)$.

C.45 identifies $X$ by any two of the $F_i$ with an open domain
having interval coordinate sections. In coordinates $(x,y)=(F_1,F_2)$,
write $F_3=w(x,y)$. Every restriction of $w$ to a horizontal or
vertical section is strictly monotone. The signs are constant
over the connected domain: for two fixed nearby values in one
coordinate, the continuous difference of their $w$-values cannot
change sign. The coordinate representation of $\mathcal W$ is
therefore exactly the triples $(g_1,g_2,k)$ satisfying
\[
 (g_1,g_2)\Omega=\Omega,\qquad
 w(g_1x,g_2y)=k(w(x,y)).                                 \tag{C.47.5}
\]
The map from $\mathcal W$ to these triples and its inverse are
continuous, by (C.47.4) and the pair-coordinate homeomorphism.
All topological group assertions below use this topology.

\emph{A strictly interior mean makes an ordered-pair stabilizer trivial.}
We first prove an elementary fact used in the non-plane cases.
Let
\[
 D=\{(a,b):a<b\}
 \quad\hbox{or}\quad
 D=\{(a,b):a<b<a+1\}.
                                                        \tag{C.47.6}
\]
Suppose a group of increasing line homeomorphisms preserves $D$
diagonally and preserves a continuous function $M:D\to\R$ with
\[
 a<M(a,b)<b,
 \qquad M\text{ injective in each argument separately}.   \tag{C.47.7}
\]
A group element fixing one pair $(a,b)\in D$ is the identity.
Indeed, it fixes all the points obtained by repeatedly inserting
$M$ between consecutive fixed endpoints. These subdivisions
have mesh tending to zero. To verify this, fix $\varepsilon>0$.
On the compact set of pairs
\[
 a\leq u<v\leq b,\qquad v-u\geq\varepsilon,
 \]
both $M(u,v)-u$ and $v-M(u,v)$ have a positive minimum.
Every such pair is still in $D$, including in the band case
because $v\leq b<a+1\leq u+1$.
Every subdivision of an interval of length at least
$\varepsilon$ therefore shortens each child by a fixed positive
amount. No sufficiently deep subdivision can have that length.
Continuity now makes the homeomorphism fix $[a,b]$ pointwise.

Let $I$ be the maximal interval of fixed points containing
$[a,b]$. A finite right endpoint $c$ belongs to the fixed set.
Choose $x<c$ in the interior of $I$ with $(x,c)\in D$.
For $y>c$ sufficiently close to $c$, the pair $(x,y)$ remains in
$D$ and $M(x,y)$ remains in the interior of $I$. Equivariance
gives
\[
 M(x,gy)=M(gx,gy)=gM(x,y)=M(x,y).
 \]
Injectivity in the second argument forces $gy=y$, extending
$I$ past $c$. The analogous argument, with a fixed $y>c$ and
$x<c$ close to a finite left endpoint, uses injectivity in the
first argument. Both finite endpoints are impossible.
Thus $I=\R$ and $g=\id$.

\emph{A non-plane pair and two full remaining pairs.}
Suppose some pair image is not $\R^2$. Relabel and apply the
explicit coordinate changes of C.45 so that
$\Omega=D$ is one of (C.47.6).
The first two target actions then coincide:
$g_1=g_2=g$; in the band case $g(x+1)=g(x)+1$.
First suppose the images of $(F_1,F_3)$ and $(F_2,F_3)$ are both
the full plane. Each $w$-level is then the graph
\[
                   y=r_z(x),\qquad z\in\R,               \tag{C.47.8}
\]
where $(x,z)\mapsto r_z(x)$ is continuous and every $r_z$ is
a homeomorphism of $\R$ onto $\R$.
It is increasing, since $r_z(x)>x$ for all $x$, which a
decreasing real-line homeomorphism cannot satisfy.
Reverse the $z$-coordinate when necessary. Then
$z\mapsto r_z(x)$ is strictly increasing, and its range is
$(x,+\infty)$ in the half-plane case and $(x,x+1)$ in the
band case. The sign is the same for every $x$ by connectedness.
The action equation is exactly
\[
                      g r_z g^{-1}=r_{kz}.               \tag{C.47.9}
\]

For $(a,b)\in D$, there is a unique $s=s(a,b)$ with
\[
                   r_s^2(a)=b.
 \quad\hbox{Define}\quad M(a,b)=r_s(a).                  \tag{C.47.10}
\]
Here are the endpoint and continuity checks. The left side is
continuous and strictly increasing in $s$.
As $s\to-\infty$, $r_s(x)\to x$, and hence $r_s^2(a)\to a$.
This last limit follows by bounding $r_s(a)$ between $a$ and
$a+\delta$ and then using monotonicity of $r_s$ at these two
fixed points. As $s\to+\infty$ the square tends to $+\infty$
in the half-plane case. In the band case $r_s(x)\to x+1$,
and the same fixed-point bracketing gives
$r_s^2(a)\to a+2$. The required value $b$ lies strictly between
the endpoint limits. Strict monotonicity gives uniqueness,
and two nearby $s$-values bracketing a solution continue to
bracket it when $(a,b)$ varies slightly; this proves continuity.

The value in (C.47.10) lies strictly between $a$ and $b$.
It is strictly increasing in $b$ for fixed $a$.
For fixed $b$ and $a_1<a_2$, the corresponding parameters satisfy
$s_1>s_2$, because $r_s^2$ increases in its argument.
If $r_{s_1}(a_1)\geq r_{s_2}(a_2)$, applying these increasing
maps once more and using $s_1>s_2$ gives
\[
 b=r_{s_1}(r_{s_1}(a_1))
   >r_{s_2}(r_{s_2}(a_2))=b,
 \]
a contradiction. Thus $M$ is strictly increasing in its first
argument as well.
Equation (C.47.9) and uniqueness in (C.47.10) prove
$M(ga,gb)=gM(a,b)$.
The preceding mean argument therefore makes every source
point stabilizer trivial: once $g=\id$, surjectivity of $w$
in (C.47.5) also gives $k=\id$.

\emph{A further non-plane pair identifies all three line actions.}
It remains, with $\Omega=D$, to consider the case where at least
one of the other pair images is not the full plane. One of its
finite endpoint homeomorphisms from C.45 gives an exact equation
\[
                         f g=k f.                       \tag{C.47.11}
\]
Replace the third coordinate by $f^{-1}F_3$.
Now all three target actions are the same $g$; this remains true
when $f$ is decreasing. Denote the new third coordinate again
by $w(x,y)$. An equality $w=x$ or $w=y$ at one point would
hold everywhere by transitivity and contradict independence.
The relative order of the three values is therefore constant.

In the half-plane, the projected group of $g$'s acts transitively
on all ordered pairs. Its centralizer in
$\operatorname{Homeo}(\R)$ is trivial: the stabilizer of a point
$x$ acts transitively on each of $(-\infty,x)$ and $(x,+\infty)$,
so its only common fixed point is $x$.
A commuting map must send $x$ to that same fixed point.
It follows that every finite endpoint of either remaining
pair domain, now an equivariant self-homeomorphism of the
same line, is the identity.

When $x<w(x,y)<y$, take $M=w$.
When $w(x,y)>y$, the $(x,w)$ image is exactly
$\{(x,z):x<z\}$: its finite lower endpoint must be the identity,
and a finite upper endpoint would be the same map and leave
an empty interval. For each $a<b$, define $M(a,b)=c$ by
\[
                         w(a,c)=b.                      \tag{C.47.12}
\]
The restriction in $c$ is increasing and maps $(a,+\infty)$
onto $(a,+\infty)$; its increase also follows from
$w(a,c)>c$ and the limit as $c\to+\infty$.
The unique solution satisfies $a<c<b$.
When $w(x,y)<x$, the $(y,w)$ image is exactly
$\{(y,z):z<y\}$; define instead
\[
                         w(c,b)=a.                      \tag{C.47.13}
\]
The restriction in $c$ is increasing and maps
$(-\infty,b)$ onto $(-\infty,b)$, so again $a<c<b$.
Both inverse constructions are continuous by the pair-coordinate
homeomorphisms. They are injective in each argument, by the
strict injectivity of the corresponding restrictions of $w$.
Equation (C.47.5), with all three actions equal, proves their
equivariance.

In the band, $g$ commutes with unit translation.
The equalities $w=x+n$ and $w=y+n$, for any integer $n$, are
excluded by the same transitivity and independence argument.
There is consequently one integer $n$, constant on the connected
domain, such that
\[
                         x<w(x,y)-n<x+1.
 \]
Either $w(x,y)-n<y$ everywhere, or the reverse inequality holds
everywhere. In the first case take $M(x,y)=w(x,y)-n$.
In the second case use the exact domain homeomorphism
\[
          (u,v)=(y,x+1),\qquad
          M(u,v)=w(v-1,u)-n.                            \tag{C.47.14}
\]
It sends the band onto itself, and $u<M(u,v)<v$.
Both constructions are continuous and injective in each argument.
They are equivariant because every $g$ commutes with the
integer shifts. We have obtained a mean satisfying
(C.47.7) in all cases.

The mean argument proves simple transitivity of $\mathcal W$
whenever a pair image is not the full plane.
The Polish open orbit-map theorem makes evaluation
$\mathcal W\to X$ a homeomorphism in this case.
Thus $\mathcal W$ is locally compact and is a topological
two-manifold. Its faithful product action on $X\times\R$
and the Hilbert--Smith theorem in dimension three
\cite[Theorem~1.5 and the introductory reduction]{C_Pardon2013}
make it a two-dimensional Lie group.

\emph{Three full pair images give a global ordered loop.}
Suppose now that all three pair images are $\R^2$.
The function $w:\R^2\to\R$ is then continuous, and each partial
map is a homeomorphism onto $\R$. Reverse input coordinates
as necessary to make it increasing in both arguments.
At a chosen source point $(a,b)$ put $c=w(a,b)$ and use the
explicit coordinates
\[
 u=w(x,b)-c,\qquad v=w(a,y)-c,\qquad z=w(x,y)-c.
                                                        \tag{C.47.15}
\]
They give a continuous multiplication $m:\R^2\to\R$ with
\[
 m(u,0)=u,\qquad m(0,v)=v,                              \tag{C.47.16}
\]
strictly increasing in both variables, each partial map onto
$\R$. Its two division maps are continuous: solving a strictly
increasing continuous equation is continuous by fixed
bracketing on either side of the solution.

Let $S$ be the stabilizer of $(0,0)$ in $\mathcal W$.
Setting either argument equal to zero in its action equation
shows that the three target maps coincide. Thus, in its actual
subspace topology, $S$ is exactly
\[
 \operatorname{Aut}^+(m)=
 \{f\in\operatorname{Homeo}^+(\R):
       f(0)=0,\quad f(m(x,y))=m(fx,fy)\}.                 \tag{C.47.17}
\]
This is a closed subgroup. We next prove its local compactness
by uniform estimates for both maps and inverses.

Fix $a>0$ and $0<A<B$. Set
\[
 \mathcal F_{A,B}=\{f\in\operatorname{Aut}^+(m):
                            A\leq f(a)\leq B\}.
                                                        \tag{C.47.18}
\]
Define $D(x)=m(x,x)$ and $T_c(x)=m(c,x)$.
The map $D$ is an increasing homeomorphism,
$D(x)>x$ for $x>0$, and $D(x)<x$ for $x<0$.
The numbers
\[
 v_n=D^{-n}(a)>0,\qquad
 u_n\text{ defined by }m(u_n,v_n)=0                      \tag{C.47.19}
\]
therefore converge to zero, with $u_n<0$.
Every automorphism respects these formulas, and hence
\[
 D^{-n}(A)\leq f(v_n)\leq D^{-n}(B).
                                                        \tag{C.47.20}
\]
The two bounds tend to zero. If $\ell(t)$ is the unique
solution of $m(\ell(t),t)=0$, then
$f(u_n)=\ell(f(v_n))$; continuity of $\ell$ makes this tend
to zero uniformly over $\mathcal F_{A,B}$ as well.

For $c>0$, $T_c$ is increasing, onto, and satisfies $T_c(x)>x$.
Its integer iterates of zero tend to both ends of $\R$.
For every integer $j$,
\[
                       f(T_a^j(0))=T_{f(a)}^j(0).       \tag{C.47.21}
\]
To bound $f$ on a given compact interval, enclose it between
two such iterates with $j=-N,N$.
The right sides of (C.47.21) range over a compact set as
$f(a)$ ranges over $[A,B]$; the negative iterates depend
continuously on $c$ by the division-map continuity.
Monotonicity then bounds all $f$ uniformly on that interval.

For equicontinuity, let $I$ be a compact interval.
For fixed $n$, both gaps in
\[
                 m(u_n,x)<x<m(v_n,x)\qquad(x\in I)
 \]
have positive minima on $I$. Thus a single $\delta_n>0$
ensures
\[
 |y-x|<\delta_n,\ x\in I
 \quad\Longrightarrow\quad
                   m(u_n,x)<y<m(v_n,x).
 \]
Applying $f$ bounds $f(y)$ between
$m(f(u_n),f(x))$ and $m(f(v_n),f(x))$.
The already proved compact bound on $f(I)$, the uniform limits
$f(u_n),f(v_n)\to0$, and continuity of $m$ show that both
bounds approach $f(x)$ uniformly in $x$ and $f$.
This is equicontinuity on every compact interval.

The inverse family has the same estimates with new fixed
positive bounds. Choose $n$ so large that $D^{-n}(B)<a$.
Then $f(v_n)<a$ and $f^{-1}(a)>v_n>0$.
Choose $N$ with $T_A^N(0)>a$. Since $T_c^N(0)$ increases
with $c>0$, (C.47.21) gives
$f(T_a^N(0))\geq T_A^N(0)>a$, and hence
\[
                     v_n<f^{-1}(a)<T_a^N(0).            \tag{C.47.22}
\]
Applying the same compact bounds and equicontinuity argument
to these positive bounds controls all inverses uniformly.

Arzel\`a--Ascoli on a countable compact exhaustion now gives,
from any sequence in $\mathcal F_{A,B}$, locally uniform
limits $f$ and $h$ of a common subsequence of maps and inverses.
Their compact image bounds permit passage to the composition
identities, so $fh=hf=\id$.
The multiplication equation and the bounds in (C.47.18)
pass to the limit. Thus $\mathcal F_{A,B}$ is compact in the
homeomorphism-group topology.
Taking $A<a<B$ gives a compact neighborhood of the identity
in $S$, and proves its local compactness.

Evaluation $f\mapsto f(a)$ is also injective.
An automorphism fixing $a$ fixes every $v_n$ and therefore every
$T_{v_n}^j(0)$, for $j\in\Z$.
For fixed $x$, choose the first such iterate at or to the right
of $x$. It lies in $[x,T_{v_n}(x)]$, whose length tends to
zero since $v_n\to0$. These fixed points tend to $x$, so
the automorphism fixes every $x$.
The compactness of (C.47.18) proves that this continuous
evaluation is a proper embedding
\[
                              S\lhook\joinrel\longrightarrow(0,\infty).
                                                        \tag{C.47.23}
\]

\emph{From the stabilizer to the full group, without changing topology.}
The locally compact group $S$ acts faithfully and continuously
on $\R^3$ by its action on the first coordinate.
Pardon's theorem makes $S$ a Lie group.
Its embedding (C.47.23) gives $\dim S\leq1$.
The full group $\mathcal W$ is still the closed Polish group
already defined; it has not yet been assumed locally compact.

Consider the free right action of $S$ on $\mathcal W$.
It is a Cartan action in its actual topology.
Indeed, choose an identity neighborhood $U$ in $\mathcal W$
such that $U\cap S$ has compact closure in $S$, using the
proved local compactness and the subspace topology of $S$.
Choose an identity neighborhood $V$ with $V^{-1}V\subset U$.
Then
\[
       \{s\in S:Vs\cap V\ne\varnothing\}
                   \subseteq V^{-1}V\cap S
 \]
has compact closure in $S$, and translates provide the
required neighborhoods at every point.
The space $\mathcal W$ is metrizable and completely regular.
The free-action form of Palais' slice theorem therefore gives
local sections and local trivializations for
\[
                  \mathcal W\longrightarrow\mathcal W/S
 \quad\cite[Theorem~2.3.3]{C_Palais1961}.
 \]
For clarity, the slice map has the form
$\sigma(u)s\mapsto(u,s)$; its inverse is multiplication.
If $\theta$ is the equivariant slice function to $S$, the
inverse coordinates of $w$ are
$(w\theta(w)^{-1},\theta(w))$, as in the explicit free-action
statement \cite[Theorem~3.7]{C_KramerMuratGarcia2026}.
The open orbit-map theorem identifies
$\mathcal W/S$ homeomorphically with $X=\R^2$.
These actual local products of the locally compact spaces
$S$ and an open subset of $X$ prove that $\mathcal W$ is
locally compact. Its faithful action on $X\times\R$ then
makes it a Lie group by Pardon again.
The same local products give
$\dim\mathcal W=2+\dim S$, which is two or three.

Finally, every compatible lifted loop value belongs to
$\mathcal W$, and its closure in $\operatorname{Homeo}(X)$
is a closed subgroup of this locally compact group.
It is therefore locally compact with precisely the induced
homeomorphism-group topology.
None of the arguments replaces the $l$-coordinate of
(C.47.3) by its action on $X$: the kernel of that restriction
can still act elsewhere.
The coordinate changes act on the specified real target lifts
and leave the original equations (C.47.1) available throughout.
This proves the asserted intrinsic reduction and leaves the
frontier extension step explicit.
\end{proof}

\begin{claim}{C.48}{Proved}{One boundary foliation admits a compatible intrinsic contraction}
Use the based ambient-null plane loop and the full family
$\mathscr F$ of lifted boundary fiber partitions in C.46.
For an open orbit $O$ with $|\mathscr F|=1$, there is a based
compact-open null-homotopy $K_{s,t}$ of $f_t|_O$ which
commutes with its entire original operator stabilizer and
satisfies, simultaneously for every specified boundary point,
\[
                 \Phi_zK_{s,t}=Q^{B_z}_{s,t}\Phi_z.
                                                        \tag{C.48.1}
\]
Here $Q^{B_z}$ is precisely the natural median homotopy of
C.27--C.29. The original source and target covering maps are
retained. This is an intrinsic assertion on $O$; ambient
continuity at its frontier is an additional requirement.
\end{claim}
\passage{Section 23, folios 140-141, scans 279-281; Part C(b), infinite-operator zero-level comparison and actual boundary quotient compatibility}


\begin{proof}
\emph{The actual lifted operator group.}
Let $p:X\to O$ be the universal covering and choose one
real quotient $F:X\to\R$ representing the unique partition.
By C.43, $X$ is a plane, $F$ has connected proper-line fibers,
and a product chart identifies it with the second projection.
Let $Q$ be the closure, in $\operatorname{Homeo}(O)$, of the
restrictions of the original $O$-stabilizing operators.
C.30 proves that this exact closed group acts freely and
properly. Define its group of all lifts by
\[
 \widetilde Q=
 \{a\in\operatorname{Homeo}(X):pa=qp\text{ for some }q\in Q\}.
                                                        \tag{C.48.2}
\]
It includes the source deck group. Its topology is the
subspace homeomorphism-group topology, with the maps and
their inverses controlled on compact subsets.

We give the needed covering comparisons explicitly.
A lift is determined by its value at one point once its
downstairs map is specified. For a convergent sequence of
downstairs maps and a convergent compatible choice of those
values, the lifts converge on compact subsets: cover a
compact path from that point by finitely many covering
charts and continue the chosen sheet successively; the
same argument on finitely many such neighborhoods controls
a compact set. Apply it to the inverse maps as well.
Conversely a convergent sequence of lifts gives convergence
downstairs on each compact set by finitely many local
sections of $p$. These observations prove that (C.48.2) is
closed and that its projection to $Q$ has the stated
topology of lifts.

The action of $\widetilde Q$ on $X$ is free. A lift fixing
a point has downstairs map fixing a point, hence equal to
the identity; the resulting deck map fixing a point is the
identity too. It is proper in this topology. Indeed, from
$x_n\to x$ and $a_nx_n\to y$, properness downstairs gives
a subsequence $q_n\to q$. There is a unique lift $a$ of $q$
with $ax=y$, and the covering-chart argument gives
$a_n\to a$. The same argument makes evaluation at a point
a homeomorphism onto a closed subset of $X$: convergence
of its values supplies convergence of the corresponding
downstairs maps and then of the selected lifts.
Thus $\widetilde Q$ is locally compact. Its faithful product
action on $X\times\R$ makes it a Lie group by the already
checked Hilbert--Smith theorem of C.33. Its dimension is
at most two, since each orbit is properly embedded in $X$.

The original operators preserve the unique partition by
(C.46.5). This property passes to their restrictions'
closure and to all lifts. More explicitly, choose lifts
of an approximating sequence whose values at one point
converge to the chosen lift. Equality of $F$-values on a
fiber passes to the limit, and the inverse maps give
equality of the resulting fiber partitions. Since $F$ is
an open quotient, every $a\in\widetilde Q$ therefore has a
unique induced line homeomorphism $\bar a$ with
\[
                         Fa=\bar aF.                   \tag{C.48.3}
\]
A continuous section of $F$ gives $\bar a=Fa\sigma$, and
the inverse formula uses $a^{-1}$; hence this induced
homomorphism is continuous in the actual topologies.

Lift the original restricted isotopy from the identity
to $a_t$ on $X$. C.32 proves its null evaluation loop in
$O$, so $a_1=\id$. Every element of (C.48.2) commutes with
$a_t$: its commutator is a continuous deck-valued function
of $t$, initially the identity. Write
\[
                         Fa_t=b_tF.                    \tag{C.48.4}
\]
These are the based target lifts, and $b_0=b_1=\id$ by
C.29 for a circle target and directly for a line target.
The group generated by the paths $b_t$ acts transitively
on $\R$. Its orbits are intervals, because every word is
joined to the identity by a path of words. A finite
endpoint of a proper orbit interval would be fixed by
every word. Such a common fixed point would project to a
fixed point of the nontrivial transitive boundary orbit,
which is impossible. Thus the orbit interval is all of
$\R$.

\emph{The base operator image and its kernel.}
Take the closed transitive group generated on this real
line. The evaluation proof of C.30 applies to the line
without change and shows that its full centralizer is
locally compact, free and proper. Its product action on
$\R^3$ makes it Lie. No decreasing element is possible:
it would have a fixed point, contrary to freeness.
In dimension one its identity component acts transitively
on the line, and freeness excludes additional components;
it is the real translation group. In dimension zero the
nearest-orbit-point argument of C.44 gives either the
identity group or a cyclic translation group. Choose the
resulting line coordinate. We have thereby proved, with
the comparison (C.48.3), that every base operator is a
translation. Write
\[
 Fa=F+\alpha(a),\qquad
 \alpha:\widetilde Q\longrightarrow(\R,+),\qquad
 K=\ker\alpha,
                                                        \tag{C.48.5}
\]
after this specified change of target coordinate.

Restriction of $K$ to any fiber is faithful, free and
proper. Evaluation gives it the same topology as its
closed orbit in that line. The preceding one-dimensional
classification therefore gives exactly
\[
                         K=1,\quad K\cong\Z,
                         \quad\text{or}\quad K\cong\R.
                                                        \tag{C.48.6}
\]
All coordinate changes below are homeomorphisms of the
source and reparametrizations of this same target.
We construct them and keep the action of every component.

\emph{The real kernel.}
For $K\cong\R$, its orbits are the entire fibers. A global
section of its principal bundle over the line, constructed
as in C.33, gives coordinates $(u,r)$ with $F(u,r)=r$ and
\[
                         k_c(u,r)=(u+c,r).              \tag{C.48.7}
\]
The quotient $\widetilde Q/K$ acts freely and properly on
$X/K=\R$, by the exact local-section argument in C.33.
Its image under $\alpha$ is consequently closed, and is
$0$, $d\Z$ with $d>0$, or $\R$.
For image zero (C.48.7) is the complete operator action.
For image $d\Z$, choose a label-$d$ element. Its action is
\[
                 q(u,r)=(\lambda u+\tau(r),r+d),
                         \lambda\in\R^\times.          \tag{C.48.8}
\]
The scalar is constant because conjugation on the actual
group $K$ is a single continuous group automorphism.
Choose $c$ on $[0,d]$ with
$c(d)=\lambda c(0)+\tau(0)$ and extend it by
$c(r+d)=\lambda c(r)+\tau(r)$. In the coordinates
$u'=u-c(r)$ the complete additional operator is
\[
                         q(u',r)=(\lambda u',r+d).
                                                        \tag{C.48.9}
\]
This retains negative $\lambda$ as well.

For image $\R$, the group $\widetilde Q$ acts transitively
and freely on $X$. Both its kernel and quotient are
connected lines, so it is connected. A Lie-algebra element
mapping to $1$ integrates to a section $s_r$ with
$\alpha(s_r)=r$. Conjugation on $K$ has the form
$s_rk_us_r^{-1}=k_{e^{cr}u}$ for a real constant $c$.
Evaluation gives the exact group coordinates and product
\[
 X\cong\widetilde Q\cong\R^2,\qquad
 (u,r)(v,z)=(u+e^{cr}v,r+z),\qquad F(u,r)=r.
                                                        \tag{C.48.10}
\]
The formula proves directly that these coordinates are
global; $k_us_r$ is the unique element with the specified
base and fiber values. The operator action is left
multiplication, and its full centralizer consists of
right multiplications.

\emph{The cyclic kernel.}
For $K\cong\Z$, the construction (C.44.5) supplies
coordinates with $F(u,r)=r$ and generator
$k(u,r)=(u+1,r)$. The quotient is the actual cylinder
$X/K=(\R/\Z)\times\R$. The action of
$\widetilde Q/K$ on it is free and proper: lift convergent
point sequences by local sections and correct by elements
of $K$, then apply properness upstairs, exactly as in
C.33. The compact circle fibers show that its base image
is closed. Indeed, a sequence of base translations in a
compact interval sends one fixed point into a compact
cylinder strip; properness gives a convergent subsequence
of the quotient elements. Thus again the image is
$0$, $d\Z$, or $\R$.

For image $d\Z$, the additional generator has the form
\[
 q(u,r)=(g_r(u),r+d),\qquad
 g_r(u+1)=g_r(u)+\epsilon,\quad\epsilon\in\{1,-1\}.
                                                        \tag{C.48.11}
\]
Here $g_r$ is increasing for $\epsilon=1$ and decreasing
for $\epsilon=-1$. Choose an increasing unit-translation
equivariant path $\rho_r$ on $0\leq r\leq d$ from
$\rho_0=\id$ to $\rho_d=g_0\circ(\epsilon\id)$.
The convex interpolation of these two increasing maps
is such a path. Extend it in both directions by
\[
                 \rho_{r+d}=g_r\rho_r(\epsilon\id).
                                                        \tag{C.48.12}
\]
The endpoint equation makes this extension continuous.
Every $\rho_r$ remains increasing and commutes with unit
translation. The homeomorphism $(u,r)\mapsto(\rho_r(u),r)$
therefore changes the two generators into exactly
\[
                 k(u,r)=(u+1,r),\qquad
                 q(u,r)=(\epsilon u,r+d).              \tag{C.48.13}
\]

For base image $\R$, the projection with discrete kernel
is a covering homomorphism of Lie groups. Its identity
component maps isomorphically onto the simply connected
line. It centralizes $K$, since its conjugation map to
$\operatorname{Aut}(\Z)=\{1,-1\}$ is continuous and is
initially the identity. Thus
$\widetilde Q=\Z\times\R$. Use this real subgroup to
transport the fiber $F^{-1}(0)$ to every other fiber,
and use the nearest-orbit-point coordinate for $K$ on
that first fiber. These are global coordinates in which
\[
           k(u,r)=(u+1,r),\qquad s_z(u,r)=(u,r+z).
                                                        \tag{C.48.14}
\]
For base image zero only $k$ remains.

\emph{Trivial kernel and the possible dense base image.}
For $K=1$, the group $\widetilde Q$ embeds algebraically
in the translation group and is abelian. A positive
dimensional identity component maps isomorphically onto
$\R$; injectivity excludes other components. The same
transport from $F^{-1}(0)$ gives just $s_z$ in (C.48.14).
In dimension zero, $X\to X/\widetilde Q$ is a regular
surface covering with simply connected source. Its deck
group is the actual fundamental group of the quotient.
A noncompact surface has free fundamental group, and
an abelian free group is trivial or infinite cyclic.
For a compact quotient the surface classification gives
only the torus with abelian torsion-free fundamental
group: the sphere and projective plane do not have plane
universal cover, the Klein-bottle group is nonabelian,
and every remaining closed surface group is nonabelian.
Consequently the discrete possibilities are exactly
$1$, $\Z$, and $\Z^2$.

The cyclic case has a nonzero base translation; choose
its positive generator. The construction
(C.48.11)--(C.48.13), now without the unit-translation
condition, changes it to $(u,r)\mapsto(\epsilon u,r+d)$.
The trivial case has no equations.
For $\widetilde Q\cong\Z^2$ the quotient is a torus,
and its base image is dense: an injective homomorphism
from $\Z^2$ into $\R$ has two generators with irrational
ratio. The foliation descends to an oriented minimal
topological foliation of that torus. To verify minimality,
the saturation of any lifted fiber under the deck group
is $F^{-1}(r+\alpha(\widetilde Q))$, which is dense because
$F$ is open and the displayed set of levels is dense.
The quotient torus is orientable, and its deck maps
preserve both the surface and the transverse orientation,
so the leaf orientation descends as asserted.

We give the elementary foliation conjugacy needed here,
including its maps. It is the minimal-torus conclusion
recorded in \cite[Proposition~2.5]{C_MionMouton2025}; no
bi-foliation classification or asymptotic-cycle statement
is needed. The oriented local flowboxes give a nonsingular
continuous flow with these leaves by the construction
used in C.43. Choose a small closed transverse interval,
with endpoints off the orbit of an interior point, and a
first-return arc from that point to a different point of
the interval. A flowbox excludes sufficiently short
positive returns; minimality supplies a later return,
and compactness gives a first one. The endpoints chosen
off the orbit put that return in the interval's interior.
The arc's interior avoids the interval, and its endpoint
differs from its initial point because there is no
periodic orbit. Its union with the intervening
transverse subinterval is a simple closed curve.
In a thin flow strip along the arc, tilt the arc strictly
in the transverse direction pointing from the return
point back to the initial point. The tilt may be chosen
arbitrarily small, because the two endpoint transverse
coordinates are distinct; choose it so that the new
terminal point remains strictly between the old
endpoints in the return interval. Join the new endpoints
inside the transverse interval. Both pieces now have
the same strict transverse direction; straightening
their two corners in flowboxes gives an embedded closed
transversal $\gamma$.

Every leaf meets $\gamma$, since it is dense and a
flowbox around $\gamma$ is open. Forward and backward
returns exist: on a compact minimal flow each forward
and backward limit set is a nonempty closed invariant
set, hence the entire torus. Flowboxes and avoidance
of $\gamma$ on a compact interior return arc show that
the first return time is continuous and positive.
The return map $h:\gamma\to\gamma$ is a minimal
orientation-preserving circle homeomorphism. The map
which follows each point for a fraction of its return
time gives the leaf-preserving homeomorphism from its
mapping torus
\[
 [0,1]\times\gamma/((1,z)\sim(0,hz))\longrightarrow X/\widetilde Q.
                                                        \tag{C.48.15}
\]
Surjectivity uses the preceding returns; uniqueness of
the last intersection gives injectivity, and the compact
source and Hausdorff target give the inverse continuity.

For completeness, choose a probability measure invariant
under $h$ by taking a weak limit of finite orbit averages.
Minimality gives full support, and the absence of periodic
points gives no atoms. Its cumulative distribution is
therefore a circle homeomorphism. It conjugates $h$ to a
Lebesgue-measure-preserving increasing homeomorphism,
hence to a rotation: preservation of the measure of
every interval makes the displacement constant.
The angle $\theta$ is irrational by minimality.
After this conjugacy, the explicit map
\[
              [s,z]\longmapsto[(s,z+s\theta)]
                         \quad\text{in }\R^2/\Z^2
                                                        \tag{C.48.16}
\]
is well defined at the identified endpoints and sends
the suspended leaves to the lines of slope $\theta$.
This proves the required foliation conjugacy.

Lift (C.48.15)--(C.48.16) to the universal covers. It
intertwines the original deck action with translations
by $\Z^2$, through the induced isomorphism of their
fundamental groups. Each original $F$-fiber becomes a
whole affine line of slope $\theta$. A linear transverse
coordinate $r$ and a complementary linear coordinate
$u$ thus give
\[
       q(u,r)=(u+\beta(q),r+\alpha'(q))\quad(q\in\widetilde Q),
                                                        \tag{C.48.17}
\]
where these translation vectors form a lattice, and
$\alpha'(\widetilde Q)$ is dense. The original function
$F$ is a line homeomorphism composed with $r$: both
are open quotients with exactly the same connected
fibers. This gives the exact target reparametrization,
rather than identifying the two quotient functions
without a comparison map.

\emph{Contraction in the established normal forms.}
Except for the separately displayed transitive case
(C.48.10), write the actual lifted loop as
\[
                 a_t(u,r)=(h_t(u,r),b_t(r)).             \tag{C.48.18}
\]
Each $h_t(\,\cdot\,,r)$ is an increasing line
homeomorphism. Its endpoints are the identity family.
Let $B_{s,t}$ be the natural median contraction of
$b_t$ from C.27, and set
\[
 H_{s,t}(u,r)=
       \bigl((1-s)h_t(u,r)+su,\ B_{s,t}(r)\bigr).
                                                        \tag{C.48.19}
\]
We check every operator equation rather than changing
the operator group. For $k(u,r)=(u+1,r)$ the equation
is $h_t(u+1,r)=h_t(u,r)+1$, which (C.48.19) preserves.
For $q(u,r)=(\epsilon u,r+d)$ it is
\[
 h_t(\epsilon u,r+d)=\epsilon h_t(u,r),\qquad
 b_t(r+d)=b_t(r)+d.
                                                        \tag{C.48.20}
\]
Convex interpolation preserves the first identity and
median naturality preserves the second.
For the real kernel (C.48.7), commutation says
$h_t(u,r)=u+\phi_t(r)$; with (C.48.9) the further
equation is $\phi_t(r+d)=\lambda\phi_t(r)$.
Scaling $\phi_t$ by $1-s$ preserves it, including
negative $\lambda$.
For the real base translations in (C.48.14), the
equations are that $h_t$ is independent of $r$ and
$b_t(r)=r+c_t$. Both are preserved by (C.48.19).

For the lattice (C.48.17), the dense set of base
translations and continuity force $b_t(r)=r+c_t$.
The fiber equation is exactly
\[
 h_t(u+\beta(q),r+\alpha'(q))=h_t(u,r)+\beta(q).
                                                        \tag{C.48.21}
\]
It is preserved by convex interpolation with $u$.
The median of translations is again a translation,
so the base equation is preserved as well.
The trivial operator case has no extra equation.
These cases exhaust (C.48.6).

In (C.48.10), write the loop as right multiplication
by $(v_t,c_t)$. Its action is
$(u,r)\mapsto(u+e^{cr}v_t,r+c_t)$.
The median base contraction is translation by some
$C_{s,t}$. Right multiplication by
\[
                         ((1-s)v_t,C_{s,t})             \tag{C.48.22}
\]
is the required contraction and still commutes with
every left multiplication.

Every map in (C.48.19) is a homeomorphism: its base
map is a line homeomorphism, and its fiber map is
an increasing convex combination of two increasing
onto maps. The maps are jointly continuous. Their
inverses are jointly continuous as well: on compact
parameter sets, monotonicity and two fixed bracketing
fiber values bound the inverse, and continuity and
uniqueness identify every limit. This is the inverse
continuity argument of C.31 applied to the displayed
fiber family and the continuous inverse base family.
Formula (C.48.22) has the same properties directly
from the group law. All endpoint identities are
preserved. We have therefore obtained a based
null-homotopy commuting with all of $\widetilde Q$.

In particular it commutes with the original source
deck group, so it descends to a compact-open
null-homotopy $K_{s,t}$ on $O$, with continuous
inverse family. Its commutation with every original
operator follows from (C.48.2), retaining the
restriction kernel, which acts identically on $O$.
Finally each other boundary real quotient is an
exact reparametrization of $F$ by the first case
of C.45. Median naturality under that line
homeomorphism gives its required quotient identity.
For a circle target it commutes with the target
integer translations and descends. The equations
$q_zF_z=\Phi_zp$ now give (C.48.1) for every $z$.
No limit across the ambient frontier has been
deduced from this intrinsic construction.
\end{proof}

\begin{claim}{C.49}{Proved}{Lie contractions on closures with three boundary foliations}
Let $f_t$ be an ambient-null based loop in
$\operatorname{Homeo}^+(\R^2)$, let $G$ be its generated group,
and let $H$ commute with every $f_t$.
For every open $G$-orbit $O$ having at least three distinct actual
lifted boundary-quotient fiber partitions of C.42, the group
generated by the based lifts of $f_t|_O$ to the universal covering
$X\to O$ is closed in $\operatorname{Homeo}(X)$. With this actual
topology it is one of
\[
 \R^2,\qquad \operatorname{Aff}^+(\R),\qquad
 \mathcal A_2=\R^2\rtimes\R ,
                                                        \tag{C.49.1}
\]
where the last factor acts on $\R^2$ by positive scalar
dilations. Each group is simply connected and has a globally
invertible exponential map.

There is a jointly continuous based contraction on the whole
$\overline O$, with jointly continuous inverse, obtained from
these actual lifted groups by fractional powers.
These contractions preserve every actual boundary-quotient
equation, agree on intersections of such orbit closures and
with the contraction of C.37 on $X_2$, and are transported to
one another by $H$.
They consequently give a continuous homeomorphism family on
every finite union of these closures and $X_2$.
The conclusion does not assert continuity across an arbitrary
infinite accumulating family of the new three-dimensional
cases.
\end{claim}
\passage{Section 23, folios 140-141, scans 279-281; Part C(b) infinite-operator request; actual simultaneous boundary-quotient and frontier-extension obligations}


\begin{proof}
\emph{The actual based source lifts are closed loops and act transitively.}
Choose three partitions and lifts $F_i:X\to\R$ as in C.47.
Lift the isotopy $f_t|_O$ from the identity to
$\widetilde f_t\in\operatorname{Homeo}(X)$.
The covering homotopy property gives continuity for maps and
inverses in their compact-open topology.
Every $\widetilde f_t$ commutes with every deck transformation:
the commutator is a continuous deck-valued function of $t$
and is the identity at $t=0$.
Its target lifts satisfy
\[
                    F_i\widetilde f_t=b_{i,t}F_i .
                                                        \tag{C.49.2}
\]
For a real target, $b_{i,1}=\id$ directly. For a circle target,
the exact two-point winding comparison of C.29 makes the
based real lift return to the identity, since the original
plane loop is ambient-null. Thus $b_{i,1}=\id$ for all $i$.
The injectivity of $(F_1,F_2)$ from C.45 now gives
$\widetilde f_1=\id$.

Let $N_{\mathrm{alg}}=\langle\widetilde f_t:0\leq t\leq1\rangle$.
This actual subgroup is path connected: each finite word is
joined to the identity by simultaneously shortening its
factor paths. Put
\[
 P=\{\widetilde f_t:0\leq t\leq1\}
       \cup\{\widetilde f_t^{-1}:0\leq t\leq1\}.
 \]
It is compact and contains the identity.
The source deck group $\Delta$ is countable, because a fiber
of a covering of a second-countable surface is a discrete
countable subset of its second-countable covering plane.
For $x\in X$, transitivity of $G$ on $O$ gives
\[
                       X=\bigcup_{\delta\in\Delta}
                            \ \bigcup_{n\geq1}\delta P^n x .
                                                        \tag{C.49.3}
\]
All sets in this countable union are compact.
Baire's theorem gives interior to one of them, hence to
$N_{\mathrm{alg}}x$. Transitivity on that orbit then makes
the entire orbit open.
Every other orbit is a deck translate of it, since the
generated lifted group commutes with $\Delta$.
These disjoint open orbits partition the connected plane;
there is only one. Thus $N_{\mathrm{alg}}$ is transitive
on $X$, not just on its projection $O$.

\emph{The connected three-web groups and their exact actions.}
Let $\mathcal W$ be the full group in C.47 and write
$W=\mathcal W^0$. The group $W$ is open in $\mathcal W$.
The open orbit map makes every $W$-orbit open in $X$;
normality of $W$ and connectedness of $X$ imply transitivity.
For a point $x$, let $J=W_x$.
The Lie principal bundle $W\to W/J=X=\R^2$ shows that
$J$ is connected: the exact sequence on components has
$\pi_1(X)=0$ and $\pi_0(W)=0$.

If $\dim W=2$, then $J$ is discrete and connected, hence
trivial. Evaluation identifies $W$ homeomorphically with
the simply connected plane. Its Lie algebra and group are
the two explicit possibilities $\R^2$ and
$\operatorname{Aff}^+(\R)$ in C.36.
Since $N_{\mathrm{alg}}\subset W$ is already transitive
and this action of $W$ is simply transitive, one has
$N_{\mathrm{alg}}=W$.

Suppose $\dim W=3$. Here C.47 puts us in the case where all
three pair images are full planes.
The connected one-dimensional isotropy $J$ is a subgroup
of the ordered-loop automorphism group in that proof.
It cannot be a circle. Indeed, every compact group of
increasing line homeomorphisms is trivial: the compact orbit
of any point has a largest element, which all group elements
must fix, so transitivity on that orbit makes it a singleton.
Thus $J\cong\R$. The principal bundle over the contractible
plane now gives $\pi_1(W)=0$.

We identify its Lie algebra and homogeneous action explicitly.
Let $\mathfrak w=\operatorname{Lie}(W)$ and
$\mathfrak j=\operatorname{Lie}(J)$.
For each $i$, the stabilizer $K_i$ of $F_i(x)$ in the
induced real target action is a closed subgroup containing
$J$. The open orbit map identifies $W/K_i$ with the actual
real target. The same principal-bundle component argument
makes $K_i$ connected, and its Lie algebra $\mathfrak k_i$
has dimension two.
The three $\mathfrak k_i$ are distinct: otherwise the
corresponding connected subgroups, and hence the quotient
fibers through $x$, would coincide, contradicting C.45.

For $0\ne Z\in\mathfrak j$, the induced linear map
$\operatorname{ad}Z$ on $\mathfrak w/\mathfrak j$ preserves
the three distinct lines $\mathfrak k_i/\mathfrak j$.
A real linear map on a two-dimensional space preserving
three distinct lines is scalar; write that scalar $\lambda$.
If $\lambda=0$, then
$[\mathfrak j,\mathfrak w]\subset\mathfrak j$, so
$\mathfrak j$ is an ideal.
Its connected subgroup $J$ is then normal in $W$.
As a normal point stabilizer in a transitive action it fixes
all of $X$, contrary to the faithful inclusion
$W\subset\operatorname{Homeo}(X)$.
Therefore $\lambda\ne0$.

The adjoint map on $\mathfrak w$ has the eigenspace
decomposition
\[
 \mathfrak w=\mathfrak j\oplus V,\qquad
 \dim V=2,\qquad [Z,v]=\lambda v\quad(v\in V).
 \]
For $v,w\in V$, the derivation identity gives
$[Z,[v,w]]=2\lambda[v,w]$.
The value $2\lambda$ is absent from the spectrum
$\{0,\lambda\}$, so $[V,V]=0$.
Replacing $Z$ by $Z/\lambda$ identifies the Lie algebra
with scalar homotheties:
\[
              [Z,v]=v,\qquad [V,V]=0.                  \tag{C.49.4}
\]
Since $W$ is simply connected, the corresponding group
is exactly $\mathcal A_2$, with multiplication
\[
        (a,v)(b,w)=(a+b,\ v+e^a w),
        \qquad a,b\in\R,\quad v,w\in\R^2.               \tag{C.49.5}
\]
Under this isomorphism $J$ is the dilation subgroup
$\{(a,0)\}$. The orbit identification
$W/J\to X$ therefore gives an actual equivariant
homeomorphism of the source plane with the standard action
\[
                         z\longmapsto v+e^a z .
                                                        \tag{C.49.6}
\]

\emph{The generated subgroup is closed, rather than merely analytic.}
Yamabe's theorem makes the path-connected subgroup
$N_{\mathrm{alg}}\subset W$ an analytic Lie subgroup
\cite[pp.~13--14]{C_Yamabe1950}.
This statement initially gives its analytic subgroup
structure; it does not identify that structure with the
induced topology. We now establish the latter point by
classifying the transitive subgroups of (C.49.6).

A transitive analytic subgroup has dimension at least two.
In dimension two its Lie algebra is either the translation
plane $V$, or has a nonzero projection to the dilation
coordinate. In the latter case it is
\[
                    \R(Z+v_0)\oplus L,                 \tag{C.49.7}
\]
where $L\subset V$ is a one-dimensional translation subspace.
On the quotient line $V/L$, the generator $Z+v_0$ has flow
\[
                 \bar z\longmapsto
                       e^t\bar z+(e^t-1)\bar v_0 .
 \]
It fixes $-\bar v_0$, while $L$ acts trivially there.
Thus the analytic subgroup in (C.49.7) preserves the proper
affine line $-v_0+L$ in the source plane, contradicting
transitivity. The remaining two-dimensional subgroup is
exactly all translations $V$.
A three-dimensional analytic subgroup is open and equals
the connected group $\mathcal A_2$.
Hence, in this case,
$N_{\mathrm{alg}}=V$ or $N_{\mathrm{alg}}=\mathcal A_2$.

Both are closed with their actual homeomorphism-group
topology. One evaluation recovers a translation vector.
For a homothety, evaluation at two fixed distinct points
recovers its positive scale and translation continuously;
a limit that is still a homeomorphism has positive scale.
Together with the dimension-two case already treated,
this proves that $N=N_{\mathrm{alg}}$ is closed and is
one of (C.49.1), with its induced topology.
It also proves that choosing a different triple of
partitions does not change $N$: it is the same subgroup
generated by the same based lifted loop.

For the semidirect products in (C.49.1), with vector
dimension one or two, the exponential and its inverse
are given by
\[
 \exp(u,v)=
   \left(u,\frac{e^u-1}{u}\,v\right),
 \qquad \frac{e^u-1}{u}=1\quad(u=0).                    \tag{C.49.8}
\]
The displayed scalar factor is positive everywhere.
For $\R^2$ the exponential is the identity.
This proves the asserted global exponential coordinates
without an unspecified choice of logarithm.

\emph{The exact Polish quotient extending the action to the closure.}
Put $L=\overline G$ in the plane homeomorphism group.
It preserves $O$ by the forward-and-inverse boundary argument
of C.36. Define the closed Polish subgroup
\[
 E_O=\{(l,n)\in L\times N:
                  p n=(l|_O)p\}.                       \tag{C.49.9}
\]
Projection $E_O\to N$ is onto: each $n\in N$ is an actual
finite word in the based lifted loop, so the identical
word in the original plane loop supplies its $l\in G$.
This continuous surjective homomorphism of Polish groups
is open.
Its kernel consists of pairs $(l,1)$ with $l$ fixing $O$
pointwise, and therefore fixing $\overline O$ pointwise.
Consequently
\[
 N\times\overline O\longrightarrow\overline O,\qquad
     (n,z)\longmapsto lz\quad\bigl((l,n)\in E_O\bigr)     \tag{C.49.10}
\]
is well defined and jointly continuous.
Indeed, the product of the open projection with
$\id_{\overline O}$ is an open quotient, and the continuous
$E_O$-action factors through it.
This argument keeps the exact source covering and the
possibly nontrivial global restriction kernel.

Every $n$ has its chosen preimage in $G$, so (C.49.10)
preserves every original $G$-orbit in $\overline O$ and
has on that orbit exactly the image of $G$.
It commutes with the original setwise $H$-stabilizer of
$O$, since those preimages commute with that stabilizer.
All source deck transformations commute with $N$ because
they commute with its generating based lifts, so the
restriction of (C.49.10) to $O$ is also the exact descended
source action.

The required closure homotopy is now
\[
 R^O_{s,t}
   =\exp\bigl((1-s)\log\widetilde f_t\bigr)
           \quad\hbox{acting by (C.49.10)}.              \tag{C.49.11}
\]
The source loop is closed and based, as proved above.
Thus (C.49.11) starts with $f_t|_{\overline O}$, ends
with the identity, and fixes both based edges.
The action and group inversion give joint continuity
of the family and its inverse.

\emph{All actual boundary quotients and common boundary images.}
Consider any nontrivial nonopen $G$-orbit
$B\subset\overline O$, with its intrinsic line or circle
topology from C.28. The action (C.49.10) is transitive
on its underlying set. Its point stabilizer $N_z$ is closed.
The continuous bijective orbit map $N/N_z\to B$ has
one-dimensional domain.
Dimension zero is excluded by connectedness and nontriviality.
If its dimension were at least two, restricting to a small
two-dimensional coordinate disk and applying invariance
of domain in the ambient plane would give interior to
$B$, making this whole transitive orbit open.

The bijection agrees with the intrinsic topology.
A compact interval in the homogeneous one-manifold maps
to an embedded arc. On a non-circular orbit that arc is
the unique connecting arc of C.28; on a circle the same
statement holds in a sufficiently small interval.
Its interior is therefore an intrinsic interval
neighborhood. This proves that the orbit identification
is a homeomorphism, and that the $N$-action is continuous
in the intrinsic topology as well.

Every group in (C.49.1) is metabelian.
Its image on $B$ is path connected and equals the actual
$G|_B$. The exact metabelian image argument of C.37
therefore identifies this common image, with its induced
intrinsic compact-open topology, as
\[
 M_B=
 \begin{cases}
   \R\text{ acting by translations, or }
        \operatorname{Aff}^+(\R),&B\text{ a line},\\
   \R/\Z\text{ acting by rotations},&B\text{ a circle}.
 \end{cases}                                           \tag{C.49.12}
\]
These groups are closed in the intrinsic homeomorphism
group: one evaluation recovers the translation or
rotation, and two recover the affine parameters.
The action supplies a continuous surjective Lie
homomorphism
\[
                            \rho_B:N\to M_B.            \tag{C.49.13}
\]
In particular every actual boundary quotient retains its
original equivariance equation under the closure action.
To verify this directly for
$\Phi_z(lx)=lz$, take the preimage $g\in G$ of $n\in N$.
The equality $\Phi_z(gy)=g\Phi_z(y)$ follows from the
defining quotient formula; replacing $g$ by another
preimage changes neither side because its restriction
kernel fixes $\overline O$. Thus
\[
            \Phi_z(R^O_{s,t}y)
                   = R^B_{s,t}(\Phi_z(y))               \tag{C.49.14}
\]
with the same original line or circle target and its
covering equation.

Since $N$ is simply connected, (C.49.13) has a unique
based lift $\widetilde\rho_B:N\to\widetilde M_B$.
The latter group is $\R$ or $\operatorname{Aff}^+(\R)$,
whose exponential is globally invertible.
The path $\widetilde\rho_B(\widetilde f_t)$ is precisely
the based lift $\widetilde b_t$ of the original
$f_t|_B$. Naturality and injectivity of the exponential
give the exact identity
\[
 \widetilde\rho_B
       \left(\exp\bigl((1-s)\log\widetilde f_t\bigr)\right)
       =\exp\bigl((1-s)\log\widetilde b_t\bigr).          \tag{C.49.15}
\]
This defines $R^B_{s,t}$ in (C.49.14) after projection.
On a point orbit the action is the identity.

\emph{Intersections, operators, and the finite-union conclusion.}
For two orbit closures, their intersection consists of
nonopen $G$-orbits. On each nontrivial one, both images
(C.49.12) are the identical group $G|_B$ with the identical
intrinsic topology. The two based target-cover paths
are the same lift of $f_t|_B$. Formula (C.49.15)
therefore makes the closure contractions equal there.
It is also exactly the boundary formula in C.36--C.37,
so it agrees on every nonopen orbit in the intersection
with $X_2$.
If $O$ itself belongs to $U_2$, evaluation identifies the
simply transitive group of C.36 with $O$ and its universal
cover with $X$. The lifted translations are the unique
based source lifts used here; their transitive generated
group is exactly $N$. The two fractional-power formulas
are consequently identical on $O$ and its closure.
An open orbit not in $U_2$ is disjoint from $X_2$, because
it is an open neighborhood disjoint from $U_2$.
These observations prove agreement on the entire
intersection with $X_2$.

For $h\in H$ with $hO=O'$, lift $h|_O$ to a source-cover
homeomorphism $\widetilde h:X\to X'$.
Covering uniqueness for the two based isotopies gives
\[
              \widetilde h\,\widetilde f_t\,
                 \widetilde h^{-1}=\widetilde f'_t .
                                                        \tag{C.49.16}
\]
Conjugation is consequently a continuous group isomorphism
$N\to N'$. It intertwines their exponential maps and
the unique logarithms in (C.49.8).
It therefore carries (C.49.11) to the corresponding
homotopy on $O'$. Equality extends to the closures.
Changing the lift of $h$ by a deck map does not change
this identity, since deck maps commute with the generated
lifted groups.
The actual boundary-quotient families are transported
at the same time: with basepoint $hx$, their exact
formula is
$\Phi^{O'}_{hz}(lhx)=h\Phi^O_z(lx)$.
Thus the class of open orbits under consideration is
$H$-invariant; no arbitrary selected triple is assumed
to be preserved by the operators.

Every closure and $X_2$ is closed in the plane.
The matching families, and separately their matching
inverse families, paste continuously over a finite union.
Each preserves its member closures and is inverted by
the pasted inverse, so it is a homeomorphism of that union.
These assertions include all original covering and
operator labels. Uniform control for infinitely many
new homothety orbits, and extension to the remaining
one- and two-foliation regions, are not supplied by
this finite-pasting argument.
\end{proof}

\begin{claim}{C.50}{Proved}{Uniform contraction through all the new Lie orbit closures}
In the setting of C.49, let $U_3$ be the union of the open
$G$-orbits having at least three distinct actual lifted
boundary-quotient fiber partitions, and retain $U_2$ from C.37.
Set
\[
                  U_*=U_2\cup U_3,\qquad X_*=\overline{U_*}.
                                                        \tag{C.50.1}
\]
The contractions of C.36 and C.49 extend to a jointly
continuous based contraction through homeomorphisms of $X_*$,
with jointly continuous inverse, commuting with $H$.
This includes arbitrary accumulation of the three-dimensional
homothety orbits.

The uniform algebraic bound is exact. For every continuous
path $k:[0,t]\to\mathcal A_2$ with $k(0)=1$, put
$P=k([0,t])\cup k([0,t])^{-1}$ and $g=k(t)$.
Then
\[
        \{\exp(r\log g):0\leq r\leq1\}
                                      \subseteq P^{1913}.
                                                        \tag{C.50.2}
\]
The identity belongs to $P$, so shorter words can be padded.
For $\R^2$ and $\operatorname{Aff}^+(\R)$ the three-letter
bound of C.37 applies. All word equalities used on an orbit
hold in its specified simply connected source group.
The bound does not identify an additive vector sum with a
noncommutative product: the latter's complete error is
corrected explicitly in the proof.
\end{claim}
\passage{Section 23, folios 140-141, scans 279-281; Part C(b) infinite-operator request; uniform frontier control for the actual three-quotient Lie reductions}


\begin{proof}
\emph{Two additive path identities.}
The planar identity of C.37 represents any fraction of the
endpoint of a based path in $\R^2$ by at most three signed
path values. We also require its three-dimensional version:
for a continuous $\gamma:[0,t]\to\R^3$ from $0$ to $v$,
\[
              rv=\sum_{j=1}^{m}\varepsilon_j\gamma(s_j),
 \qquad m\leq5,\quad
       \varepsilon_j\in\{1,-1\},\quad 0\leq s_j\leq t .
                                                        \tag{C.50.3}
\]
For a polygonal path, apply the signed-measure theorem
\cite[Theorem~4.7 and \S4.3.2]{C_GrandoniRaviSinghZenklusen2014}
to its three coordinate derivatives.
The chosen set is a union of at most two circular intervals.
Its integral is a sum of two path differences, with at most
one additional total vector from an interval crossing the
identified endpoint. This gives five signed path values.
Zero terms $\gamma(0)$ permit padding to a nonempty word
when necessary.
The interval count in the halving proof is the same corrected
count used in C.37: two intervals cut at at most three points
give at most five pieces, so one half uses at most two.
For the whole interval, three cuts give at most four pieces,
so one of the two halves uses at most two intervals; these
are also circular intervals.
Complementation, halving, and compactness give every
fraction. Uniform polygonal approximation with the same
endpoints, followed by compactness of the five parameters
and the finite sign set, proves (C.50.3) for arbitrary
continuous paths.

\emph{A maximizing anchor and its exact translation commutators.}
Use the coordinates
\[
        (a,v)(b,w)=(a+b,v+e^a w),\qquad
              v,w\in\R^2,
 \]
and write $k(s)=(a_s,v_s)$.
Let $A=\max_{0\leq s\leq t}|a_s|$.
When $A=0$, the path lies in the translation plane and
the three-letter bound already proves (C.50.2).
For $A>0$, choose $u$ attaining the maximum and let
$h=k(u)$ or $k(u)^{-1}$, choosing its sign so that its
logarithmic scale is $A$.
Thus $h$ is one letter from $P$, not a power of unbounded
length. Conjugate all coordinates by the translation
which sends its unique fixed point to zero.
This is one group automorphism applied to both sides
of every forthcoming identity; it adds no letters to
any word. In the new coordinates $h=(A,0)$ and still
$|a_s|\leq A$. Put
\[
 D=e^A-1,\qquad
 C_s=[h,k(s)]
       =h k(s)h^{-1}k(s)^{-1}
       =(0,Dv_s).
                                                        \tag{C.50.4}
\]
The path $C_s$ is continuous in the translation plane
and begins at the identity.
Every one of its values and inverses uses at most four
letters of $P$.

For any fixed $s$ and $|\beta|\leq B$, where $B$ is a
positive integer, the translation $C_s^\beta$ uses at
most $12B$ letters of $P$.
Indeed, write it as $B$ copies of
$C_s^{\beta/B}$.
For nonnegative $\beta/B\leq1$, the planar identity
applied to the prefix $C|_{[0,s]}$ gives at most three
signed commutator values. For a negative fraction invert
the resulting word. Substituting (C.50.4) costs at most
four letters per commutator.
No continuous choice of the parameters in these words
is required.

\emph{The large-scale case.}
Suppose $A\geq\log2$, so $D\geq1$.
For $0\leq q\leq1$, the translation $T_{qv_s}$ equals
$C_s^{q/D}$ and therefore uses at most twelve letters.
Writing $D_a=(a,0)$ gives
\[
                    D_{a_s}=T_{v_s}^{-1}k(s),
                                                        \tag{C.50.5}
\]
which uses at most thirteen.
Let $a=a_t$ and $v=v_t$. Define
\[
 \theta(a,r)=
 \begin{cases}
  (e^{ra}-1)/(e^a-1),&a\ne0,\\
  r,&a=0.
 \end{cases}
 \qquad
       g^r=(ra,\theta(a,r)v).                           \tag{C.50.6}
\]
The function $\theta$ is continuous and lies in $[0,1]$.
The intermediate value theorem supplies $s$ with
$a_s=ra$, since the path starts at logarithmic scale zero
and ends at $a$.
Then
$g^r=T_{\theta(a,r)v}D_{a_s}$ uses at most
$12+13=25$ letters.

\emph{The small-scale case and its full noncommutative correction.}
Now suppose $0<A\leq\log2$.
Apply (C.50.3) to
$\gamma(s)=(a_s,v_s)$, obtaining
\[
       ra=\sum_{j=1}^{m}\varepsilon_j a_{s_j},\qquad
       rv=\sum_{j=1}^{m}\varepsilon_j v_{s_j},
                         \qquad m\leq5.                 \tag{C.50.7}
\]
Form the actual product
$w=k(s_1)^{\varepsilon_1}\cdots k(s_m)^{\varepsilon_m}$.
For an inverse letter, its translation is
$-e^{-a_{s_j}}v_{s_j}$.
Consequently, with
\[
 p_j=
 \begin{cases}
  \sum_{\ell<j}\varepsilon_\ell a_{s_\ell},
                                       &\varepsilon_j=1,\\
  \sum_{\ell\leq j}\varepsilon_\ell a_{s_\ell},
                                       &\varepsilon_j=-1,
 \end{cases}
                                                        \tag{C.50.8}
\]
the product is exactly
\[
 \begin{split}
 w&=\left(ra,\sum_{j=1}^{m}
                    \varepsilon_j e^{p_j}v_{s_j}\right),\\
 \operatorname{pr}_v(w)
   &=rv+\sum_{j=1}^{m}
             \varepsilon_j(e^{p_j}-1)v_{s_j}.
 \end{split}                                            \tag{C.50.9}
\]
Every $|p_j|$ is at most $mA$.
The pure translation which changes (C.50.9) to (C.50.6)
by left multiplication is
\[
 c=(\theta(a,r)-r)v
       -\sum_{j=1}^{m}
                  \varepsilon_j(e^{p_j}-1)v_{s_j}.
                                                        \tag{C.50.10}
\]
Using the translation commutators (C.50.4), this is the
exact group equation
\[
 g^r=C_t^{\beta_0}
          \left(\prod_{j=1}^{m}C_{s_j}^{\beta_j}\right)w,
 \quad
 \beta_0=\frac{\theta(a,r)-r}{D},\qquad
 \beta_j=-\varepsilon_j\frac{e^{p_j}-1}{D}.
                                                        \tag{C.50.11}
\]
The powers before $w$ commute because they are translations.

Here are uniform bounds, including $a=0$.
The integral formula
\[
        \theta(a,r)=
          \frac{\int_0^r e^{a q}\,dq}
               {\int_0^1 e^{a q}\,dq}                  \tag{C.50.12}
\]
holds at zero as well as away from it.
For $|a|\leq A$, the inequality
$|e^{aq}-1|\leq D$ holds on $0\leq q\leq1$, and the
denominator is at least $e^{-A}\geq1/2$.
Subtracting $r$ in (C.50.12), the numerator is the
difference of an integral over $[0,r]$ and $r$ times
an integral over $[0,1]$, each with integrand
$e^{aq}-1$. Its absolute value is at most $2D$.
Therefore $|\beta_0|\leq4$.
For every other coefficient,
\[
 \frac{|e^{p_j}-1|}{D}
 \leq\frac{e^{mA}-1}{e^A-1}
   =\sum_{\ell=0}^{m-1}e^{\ell A}
   \leq 2^m-1\leq31.                                  \tag{C.50.13}
\]
The same bound covers negative $p_j$ because
$1-e^{-mA}\leq e^{mA}-1$.
The word construction following (C.50.4) now applies with
$B=4$ and $B=31$. Equation (C.50.11) uses at most
\[
                      5+12(4+5\cdot31)=1913             \tag{C.50.14}
\]
original letters. This proves (C.50.2) in all regimes.
The executable certificate accompanying this claim checks
the inverse-letter exponents in (C.50.8), the complete
product correction, the commutator and conjugacy formulas,
and the integer count by exact symbolic arithmetic.
The measure selection and the analytic coefficient bounds
are proved above, independently of that certificate.

\emph{One compact family of actual word isotopies.}
Set $L_0=1913$.
For the original plane loop, all words just proved belong
to the compact parameter family
\[
 W_{\boldsymbol\varepsilon,\boldsymbol s}(u)
       =\prod_{j=1}^{L_0}f_{u s_j}^{\varepsilon_j},
 \qquad
 \varepsilon_j\in\{1,-1\},\quad
 0\leq s_j\leq t\leq1,\quad0\leq u\leq1.
                                                        \tag{C.50.15}
\]
Identity letters are obtained by $s_j=0$.
Continuity of finite multiplication and inversion in the
plane homeomorphism group gives compactness for maps and
inverses, uniformly in $u$.
On a C.49 orbit, the word equality is in its simply
connected actual group $N$.
On a C.36 orbit, use the based path in its specified
universal covering group and the three-letter C.37 bound,
then pad it. Thus the endpoint of the based lift of
(C.50.15) is exactly the required fractional power,
not just a homeomorphism with an unspecified deck shift.

Use $r=1-s$ for the contraction parameter and write
$R^O_{r,t}$ for these fractional actions on an open
orbit in $U_*$. C.36 and C.49 already prove their
individual closure continuity and finite compatibility.
Their groups are metabelian, so
$[[g_1,g_2],[g_3,g_4]]=1$ on $U_*$ for every
$g_i\in G$, and hence on $X_*$.
Every nontrivial nonopen orbit $B\subset X_*$ therefore
has the exact intrinsic image
$M_B=\R$, $\operatorname{Aff}^+(\R)$, or the circle
rotation group, by the path-connected metabelian
argument of C.37.
Define $R^B_{r,t}$ by its unique line-group logarithm
or its based real rotation lift, and use the identity
on point orbits. The based circle lift ends at zero
by C.29. These are the same boundary prescriptions
as in C.36 and C.49.

\emph{Limits, with the bounded residual in the specified cover.}
Take
\[
 x_i\in O_i\subset U_*,\quad x_i\to x\in X_*,
       \qquad (r_i,t_i)\to(r,t).
                                                        \tag{C.50.16}
\]
Choose representations (C.50.15) for
$R^{O_i}_{r_i,t_i}$ and pass to convergent word parameters.
Let $W$ be the limiting plane isotopy, $w=W(1)$,
and $g=f_t$.
If $x$ belongs to an open orbit of $X_*$, then that
orbit already lies in $U_*$ and eventually all $x_i$
lie in it. Its intrinsic continuity applies.
If $x$ is fixed by $G$, every limiting word fixes it.
For a remaining intrinsic line or circle orbit $B=Gx$,
identities on $O_i$ pass to every point $hx\in B$
by evaluating them at $hx_i$ for fixed $h\in G$.
In particular $w|_B$ commutes with $g|_B$.

For a fixed positive integer $n$, set
$m_i=\lfloor nr_i\rfloor$. In each specified simply
connected group, the residual is exactly
\[
       \widetilde w_i^{\,n}\widetilde g_i^{-m_i}
           =\exp\bigl((nr_i-m_i)\log\widetilde g_i\bigr).
                                                        \tag{C.50.17}
\]
Its exponent fraction is in $[0,1)$, so it has another
representation from the same compact word family.
For this fixed $n$, pass to a subsequence on which
$m_i=m_n$ and the residual word isotopies converge
to $Z_n$. Then
\[
       (w|_B)^n(g|_B)^{-m_n}=z_n,\qquad
       |m_n-nr|\leq1,\qquad z_n=Z_n(1)|_B,              \tag{C.50.18}
\]
and all $z_n$ belong to the same compact family of
intrinsic restricted word endpoints.
For a line, logarithms of this compact family are bounded.
In the affine group commuting nonidentity elements belong
to the same exponential one-parameter subgroup; the
identity case gives the same additive logarithm equation.
Thus
\[
          n\log(w|_B)-m_n\log(g|_B)=\log z_n .
 \]
Division by $n$ proves
$\log(w|_B)=r\log(g|_B)$.

For a circle, let $\theta(u)$ be the based real angle
of $f_u|_B$, and put $M=\max_u|\theta(u)|$.
Let $A_W$ be the real endpoint angle of $W|_B$,
$T=\theta(t)$, and $C_n$ that of $Z_n|_B$.
The fixed word length gives $|C_n|\leq L_0M$.
The real, rather than merely circle-valued, equality is
\[
                         nA_W-m_nT=C_n.                \tag{C.50.19}
\]
To prove it, use the plane word isotopies
$W_i(u)^n f_{u t_i}^{-m_i}Z_i(u)^{-1}$.
Their restrictions lift to loops ending at the identity
by (C.50.17). The specified source groups are simply
connected, so these loops are null-homotopic there.
Evaluate at the two points $x_i,hx_i$, where $hx\ne x$.
Their difference-direction winding is zero.
The resulting loops converge uniformly to the analogous
two-point loop at $x,hx$, which stays off the diagonal;
its winding is also zero.
The exact circle comparison of C.29 identifies that
winding, up to the common orientation sign, with
$nA_W-m_nT-C_n$. This proves (C.50.19).
Consequently
\[
                  |A_W-rT|
                       \leq\frac{L_0M+|T|}{n}
                       \quad(n\geq1),
 \]
so $A_W=rT$.
This also covers an identity endpoint on the circle
whose specified real angle $T$ is a nonzero integer.

\emph{Joint continuity and the remaining scope.}
Every convergent subsequence of the compact word
parameters in (C.50.16) has therefore the prescribed
endpoint action at $x$. This proves convergence to
$R^B_{r,t}(x)$ for the entire sequence.
It gives joint continuity from the dense set $U_*$.
For a general sequence in $X_*$, approximate its
$i$th point by a point of $U_*$ with both point distance
and image distance at that fixed parameter less than
$1/i$. The just-proved dense-set limit supplies those
approximations and proves full joint continuity.

Apply the same proof to $f_t^{-1}$.
It has the same generated group, the same open-orbit
families, and the same boundary partitions.
Its fractional actions are the orbitwise inverses,
so the inverse family is jointly continuous as well.
Every member is an orbitwise bijection inverted by
this family and is therefore a homeomorphism of $X_*$.
The equations
$R_{0,t}=1$, $R_{1,t}=f_t|_{X_*}$, and
$R_{r,0}=R_{r,1}=1$ give the based contraction.
The operator transport statements of C.36 and C.49
make $U_*$ and $X_*$ invariant under $H$ and give
commutation on $U_*$. Continuity extends that equality
to $X_*$.

This settles infinite accumulation for the orbit closures
treated by C.49. The remaining global step is a compatible
extension from $X_*$ through the regions with fewer
actual boundary-quotient foliations and their frontiers.
The particular median choice of C.34 cannot be imposed
there without proving compatibility, as C.35 demonstrates.
The general conclusion in C.12 consequently remains open.
\end{proof}

\begin{claim}{C.51}{Proved}{Cocompact proper operators control the entire frontier}
Let $f_t$ be the based ambient-null plane loop of C.49,
let $G=\langle f_t\rangle$, and let $H$ commute with all
its values. For an open $G$-orbit $O$, put
\[
 H_O=\{h\in H:hO=O\},\qquad
 Q_O=\overline{\{h|_O:h\in H_O\}}^{\operatorname{Homeo}(O)}.
                                                        \tag{C.51.1}
\]
Consider all such $O$ for which $O/Q_O$ is compact and
the full family $\mathscr F_O$ of actual lifted boundary
quotient fiber partitions in C.46 has at most two members.
For every one of these orbits the restricted loop has
a based contraction on its entire closure,
\[
 \overline K^O_{s,t}\in\operatorname{Homeo}(\overline O),
 \qquad
 \overline K^O_{0,t}=f_t|_{\overline O},\qquad
 \overline K^O_{1,t}=\overline K^O_{s,0}
                         =\overline K^O_{s,1}=\id,
                                                        \tag{C.51.2}
\]
jointly continuous with its inverse. It commutes with
$H_O$ and equals exactly the C.34 median contraction
on every nonopen orbit in $\partial O$.
When $\mathscr F_O$ is empty, every point of
$\partial O$ is fixed by $G$.

The contractions can be chosen with exact transport
under $H$. They agree on all common frontiers and with
C.34, and hence paste over the closed nonopen locus
together with any finite collection of these closures.
This assertion does not prescribe their values on the
closed set $X_*$ of C.50: the comparison there with
its exponential prescription remains to be proved.
\end{claim}
\passage{Section 23, folios 140-141, scans 279-281; Part C(b), infinite-operator zero-level comparison and ambient frontier extension}


\begin{proof}
Write $Q=Q_O$ and $L=\overline G$ in the plane
homeomorphism group. C.30 proves that the actual group
$Q$ in (C.51.1) is closed, locally compact, free and
proper on $O$, with the specified compact-open topology.
The restriction kernel of $H_O\to Q$ acts trivially
on $O$ and therefore on its closure.
We retain that map; we do not replace an element of
$Q$ by an unspecified global extension.

\emph{A compact set meeting every operator orbit.}
The quotient map $O\to O/Q$ is open. For each point
of $O$, choose a relatively compact open neighborhood
inside $O$. The images of these neighborhoods cover
the compact quotient. Choose finitely many whose
images still cover it, and let $C\subset O$ be the
union of their closures. Then $C$ is compact and
\[
                              Q C=O.                   \tag{C.51.3}
\]
This argument supplies the exact compact fundamental
set required below; it uses no fundamental-domain
boundary or continuous choice of orbit representatives.

\emph{Every frontier sequence produces an actual quotient.}
Take a sequence $x_n\in O$ converging in the plane to
$z\in\partial O$. By (C.51.3), write
$x_n=q_ny_n$ with $q_n\in Q$ and $y_n\in C$.
Pass to a subsequence with $y_n\to y\in C$.
For the moment choose $h_n\in H_O$ with
\[
                    |h_ny_n-q_ny_n|<1/n.               \tag{C.51.4}
\]
Such a choice follows from the definition of $Q$ as
the closure of the actual restrictions in (C.51.1).
Here and below $|\cdot|$ is the Euclidean distance
in the original ambient plane.

Every $h_n$ commutes with $L$: it commutes with all
generators of $G$, and commutation passes to its
closure in the plane homeomorphism group.
The orbit map
\[
                      L\longrightarrow O,\quad l\mapsto ly
                                                        \tag{C.51.5}
\]
is open by the same Polish orbit-map theorem used in
C.30 and C.42. Thus there are $l_n\to1$ in $L$ with
$l_ny=y_n$.
For clarity, take a decreasing countable identity
neighborhood basis $V_k$ in $L$. Each $V_k y$ is a
neighborhood of $y$, so $y_n\in V_k y$ eventually.
Choose an increasing sequence of such eventual
indices and then $l_n\in V_{k(n)}$ with
$k(n)\to\infty$. This proves the stated selection.
Commutation now gives the exact equality
\[
                       h_ny=l_n^{-1}h_ny_n\longrightarrow z.
                                                        \tag{C.51.6}
\]
Convergence follows from (C.51.4), $x_n\to z$, and
compact-open convergence of $l_n^{-1}$ to the identity
on a compact neighborhood of $z$.

Let $J_y=L_y$. For $j\in J_y$,
$j h_ny=h_nj y=h_ny$, so (C.51.6) gives $jz=z$.
Consequently $z\in\operatorname{Fix}(J_y)$, and
C.42's actual evaluation comparison gives
\[
 h_n|_O\longrightarrow\Phi_z^y
       \quad\hbox{uniformly on compact subsets of }O,
 \qquad
 \Phi_z^y:O\longrightarrow\R^2,\quad \Phi_z^y(ly)=lz.
                                                        \tag{C.51.7}
\]
This uses the homeomorphism
$\operatorname{Map}_L(O,\R^2)\to\operatorname{Fix}(J_y)$
of (C.42.1), including its compact-open topology.
In particular its assertion applies when $Gz$ is
a single point, in which case $\Phi_z^y$ is constant.
There is no convergence claim for arbitrary proper
operator maps outside their actual global approximants.

We compare (C.51.7) with the quotient family defined
at the fixed basepoint $x$ in C.46. Choose $a\in L$
with $ax=y$, and put $z_0=a^{-1}z$.
Then $z_0\in\operatorname{Fix}(L_x)\cap\partial O$.
For every $l\in L$ the equality
\[
             \Phi_{z_0}^x(ly)
                =\Phi_{z_0}^x(lax)
                =la z_0=lz=\Phi_z^y(ly)                \tag{C.51.8}
\]
identifies the actual maps, with the same target orbit.
The orbit is contained in $\partial O$, since $L$
preserves $O$ and its boundary. It is nonopen; an
open orbit meeting $\partial O$ would be an open
neighborhood disjoint from $O$, contradicting the
definition of its closure. Thus C.28 makes a
nontrivial such orbit an intrinsic line or circle.
For that case (C.51.8) is exactly one of the original
boundary quotients whose lifts define $\mathscr F_O$.
It follows in particular that an empty $\mathscr F_O$
forces every frontier orbit to be a point: apply
the construction to a sequence approaching any
prescribed frontier point.

\emph{The intrinsic family and every boundary equation.}
For $|\mathscr F_O|=1$ use the contraction $K^O_{s,t}$
of C.48, and for $|\mathscr F_O|=2$ use C.46.
Both commute with $H_O$, hence with its restricted
closure $Q$, and both have jointly continuous inverses.
They retain every quotient equation
\[
                    \Phi K^O_{s,t}=Q^B_{s,t}\Phi,
                                                        \tag{C.51.9}
\]
where $Q^B$ is the C.34 median family on the actual
boundary orbit $B$, with its based real covering
lift when that orbit is a circle.
For an empty family, C.32 makes $f_t|_O$ null in
$\operatorname{Homeo}(O)$ and C.33 supplies a based
null-homotopy $K^O$ in its exact $Q$-centralizer.
All frontier orbits are then points, as just proved,
so (C.51.9) for each constant $\Phi$ is automatic.

Define $\overline K^O$ to be $K^O$ on $O$ and
the C.34 family on $\partial O$. We prove its full
joint continuity, rather than continuity only for
fixed parameters. Let
\[
       x_n\in O,\quad x_n\to z\in\partial O,
       \qquad(s_n,t_n)\to(s,t).
                                                        \tag{C.51.10}
\]
Use $x_n=q_ny_n$ and a subsequence $y_n\to y$ as
above. The set
\[
 S=C\ \cup\!\bigcup_{(u,v)\in[0,1]^2}
                    \bigl(K^O_{u,v}(C)\cup(K^O_{u,v})^{-1}(C)\bigr)
                                                        \tag{C.51.11}
\]
is compact in $O$, by joint continuity on a compact
parameter space. Refine the choice of actual operators
in (C.51.4) to satisfy
\[
                       \sup_{w\in S}|h_nw-q_nw|<1/n.
                                                        \tag{C.51.12}
\]
For each $n$ this is a compact-open neighborhood
condition on $q_n\in Q$, so density of the specified
restrictions supplies $h_n\in H_O$.
The previous argument still proves (C.51.6)--(C.51.8).

Commutation with $Q$ and (C.51.12) give
\[
 \begin{split}
 K^O_{s_n,t_n}(x_n)&=q_nK^O_{s_n,t_n}(y_n),\\
 \bigl|K^O_{s_n,t_n}(x_n)
            -h_nK^O_{s_n,t_n}(y_n)\bigr|&<1/n.
 \end{split}                                            \tag{C.51.13}
\]
The points $K^O_{s_n,t_n}(y_n)$ converge in $O$ to
$K^O_{s,t}(y)$ and lie in the compact set $S$.
Uniform convergence (C.51.7), followed by the exact
quotient identity (C.51.9), therefore yields
\[
 \begin{split}
 h_nK^O_{s_n,t_n}(y_n)
      &\longrightarrow\Phi_z^y K^O_{s,t}(y)\\
      &=Q^{Gz}_{s,t}(\Phi_z^y(y))
        =Q^{Gz}_{s,t}(z).
 \end{split}                                            \tag{C.51.14}
\]
For a point orbit the same computation uses the
constant map and gives $z$.
Every subsequence in (C.51.10) has a further subsequence
of this form because $C$ is compact. The value on the
right is independent of all choices. Thus (C.51.14)
holds for the entire sequence and proves continuity
from $O$ to the prescribed frontier value.
Continuity for sequences contained in $\partial O$
is precisely C.34. Splitting a sequence according
to its two pieces proves joint continuity everywhere
on $[0,1]^2\times\overline O$.

The inverse maps have the same argument, using
the inverse images already included in $S$.
Equation (C.51.9) implies
$\Phi(K^O_{s,t})^{-1}=(Q^B_{s,t})^{-1}\Phi$
by composing with these actual inverse maps.
The inverse boundary family is jointly continuous
by C.34. This proves the inverse version of
(C.51.14), and hence joint continuity of the
piecewise inverse on the entire closure.
The maps preserve $O$ and each frontier orbit and
are bijections there; the piecewise inverses are
their actual inverses. They are homeomorphisms.
The endpoint and based-edge identities and
commutation with $H_O$ hold on both pieces.

\emph{Transport and finite pasting.}
The class of orbits in the claim is $H$-invariant:
conjugation by $h\in H$ identifies the two exact
groups $Q_O,Q_{hO}$, their quotient spaces, and
their full boundary-quotient families by C.46.
Choose one orbit from each $H$-orbit of these open
components, make the construction there, and set
\[
             \overline K^{hO}_{s,t}
                   =h\overline K^O_{s,t}h^{-1}.
                                                        \tag{C.51.15}
\]
The equality with every stabilizing operator makes
this independent of the choice of $h$ carrying the
selected representative to a given component.
Median naturality in C.27--C.34 makes its frontier
restriction the identical prescribed C.34 family.

Distinct open-orbit closures meet only in the
closed nonopen locus $Y$. Every one of the new
families and its inverse equals the corresponding
C.34 map there. The finite closed-set pasting
lemma thus gives the asserted maps and inverse
maps on $Y$ together with finitely many closures.
For every finite collection $\mathcal I$, the
homeomorphism $h$ carries this union to the union
indexed by $h\mathcal I$ and intertwines their
homotopies by (C.51.15). This is the exact
operator-transport comparison on those finite unions.

The proof uses one fixed compact fundamental set
inside each treated orbit. It supplies no uniform
control over infinitely many different such sets,
and it does not identify the median boundary value
with the exponential value forced on $X_*$ by C.50.
These are explicit remaining comparisons in C.12.
\end{proof}

\begin{claim}{C.52}{Proved}{Cocompact frontier extensions relative to the exponential region}
Use the actual based ambient-null plane loop $f_t$, the groups
$G=\langle f_t\rangle$ and $H$, and the closed set $X_*$ with its
contraction $E_{s,t}$ from C.50.
For every nonopen $G$-orbit $B$, define a based family
$R^B_{s,t}$ as follows:
\[
 R^B_{s,t}=
 \begin{cases}
 \id,&B\text{ a point},\\
 \text{the exponential fraction with parameter }1-s,
       &G|_B\text{ metabelian},\\
 \text{the C.27--C.29 median family},
       &G|_B\text{ nonmetabelian}.
 \end{cases}                                             \tag{C.52.1}
\]
The exponential fraction uses the actual intrinsic Lie image
classified in C.37. On a circle it uses the entire based lift
of the original loop to the covering group, including its
real angle at each intermediate time.

For every open orbit $O$ disjoint from $X_*$ whose full family
$\mathscr F_O$ of actual lifted boundary partitions has
cardinality at most two, there is an intrinsic based
null-homotopy $K^O_{s,t}$ commuting with its original operator
stabilizer and satisfying every actual boundary quotient
equation
\[
                 \Phi_zK^O_{s,t}=R^{Gz}_{s,t}\Phi_z.      \tag{C.52.2}
\]
For each such $O$ with compact proper-operator quotient
$O/Q_O$, this family extends to a based contraction through
homeomorphisms of $\overline O$, jointly continuous with
its inverse, and equal to (C.52.1) on its entire frontier.
It agrees with $E_{s,t}$ on $\overline O\cap X_*$.

These extensions can be chosen with exact transport under
$H$. They agree on common frontiers. Consequently they
paste with the prescribed C.50 contraction on $X_*$ over
every finite union
\[
                  X_*\ \cup\!
                  \bigcup_{O\in\mathcal I}\overline O,    \tag{C.52.3}
\]
where $\mathcal I$ is a finite collection of the
cocompact orbits just specified.
No continuity of (C.52.1) on the entire nonopen locus,
and no extension over an infinite collection or over the
other noncocompact open orbits, is asserted.
\end{claim}
\passage{Section 23, folios 140-141, scans 279-281; Part C(b) infinite-operator request; relative frontier extension preserving the already prescribed exponential region}


\begin{proof}
\emph{The precise intrinsic target family.}
For a nontrivial intrinsic line or circle orbit, $G|_B$ is
path connected and transitive. The proof of (C.37.9)
classifies every such metabelian image, with its actual
intrinsic homeomorphism-group topology: on a line it is
conjugate to all translations or to the full positive
affine group, and on a circle to all rotations.
That proof uses only these three stated properties of
the image, so it applies to every metabelian orbit in
(C.52.1), without requiring that orbit to lie in $X_*$.

For a real lift $F:X\to\R$ of an actual boundary quotient,
let $b_t$ be its based target path, characterized by
$Fa_t=b_tF$ for the based source lift $a_t$.
The exact source and target coverings are those of
C.42--C.44. The path $b_t$ is a based loop: for a circle
target this is the zero-winding conclusion of C.29,
and for a line target it follows from the identity
endpoint of $f_t$.
Its generated group
\[
                         T=\langle b_t:0\leq t\leq1\rangle
                                                               \tag{C.52.4}
\]
is transitive on $\R$. Indeed its orbits are intervals,
as each word is joined to the identity by a path of
words. A finite endpoint of such an orbit interval is
a common fixed point, which would project to a common
fixed point of the nontrivial transitive target action.
Thus no proper orbit interval is possible.

For a line target, (C.52.4) is the conjugate of the
actual image $G|_B$. For a circle, projection gives a
surjection $T\to G|_B$. When the latter image is
metabelian, its rotation conjugacy lifts to a real
coordinate in which every $b_t$ is a translation.
Transitivity then makes $T$ the full translation group.
Conversely a metabelian $T$ has metabelian quotient
$G|_B$. This proves the exact equivalence used in the
case selection below.

Define, in the specified real lift,
\[
 B_{s,t}=
 \begin{cases}
 \exp_T\bigl((1-s)\log_T b_t\bigr),&T\text{ metabelian},\\
 \text{the C.27 median of the path }b_t,&T\text{ nonmetabelian}.
 \end{cases}                                             \tag{C.52.5}
\]
In the first case $T$ is a simply connected translation
or positive affine Lie group with a global exponential
and logarithm. Formula (C.52.5) and its inverse are
jointly continuous, with all four required based-edge
identities. They commute with every target homeomorphism
commuting with the original path: in the exponential
case each new value lies in $T$, and in the median case
this is C.27.

This rule is natural under any line reparametrization,
including a decreasing one. Such conjugation identifies
the actual groups (C.52.4) topologically, preserves
metabelianness, and, in the Lie cases, intertwines
exponential maps and their unique logarithms. The
median case has the same naturality by C.27.
The source of a real lift of a circle may also be
changed by a covering conjugacy; the same argument
retains the specified based target path.
In particular the resulting map on $B$ is exactly
(C.52.1), independent of all these coordinates.

For clarity, at a time when $f_t|_B=\id$ on a circle,
its based real lift can still be a nonzero integer
translation. Formula (C.52.5) scales that real
translation. It does not replace it by zero.
The endpoint $t=1$ is zero by C.29. This is the same
covering convention as C.50.

\emph{Replacement in the two-foliation construction.}
For $|\mathscr F_O|=2$, retain exactly the source cover,
the two real quotient coordinates, and the domain $D$
of C.45--C.46. Replace each of the two coordinate
median families in (C.46.4) by (C.52.5).
On $D=\R^2$ this is a product of increasing line
homeomorphisms, with a jointly continuous inverse.
On $D=\{u<v\}$ the original two coordinate loops are
identical. Their generated groups and their selected
families are therefore identical, so the strict
inequality is preserved. On $D=\{u<v<u+1\}$ the same
argument applies, and the selected maps commute with
unit translation. In the exponential case this holds
because that translation centralizes $T$; in the
median case it is the original median naturality.
Thus both inequalities are preserved, as are they
by the inverse family.

Every original stabilizing operator, and every one
of its source lifts, permutes the full two-element
family of partitions by the exact comparison
(C.46.5). In these coordinates it is a pair of line
homeomorphisms, possibly followed by exchange.
Commutation with the based source loop identifies
the corresponding target paths by these same
conjugacies. Rule (C.52.5) is natural under them
and makes the same choice on conjugate paths.
Consequently the modified source family commutes
with every such lift, including all source deck maps.
It descends to $O$ with continuous inverse.

Every other actual real quotient is a line
homeomorphism composed with one of the two selected
coordinates, by C.45. The preceding naturality
therefore proves its exact lifted equation too.
Descending through $q_zF_z=\Phi_zp$ proves (C.52.2),
retaining the original source and target covers.

\emph{Replacement in the one-foliation construction.}
For $|\mathscr F_O|=1$, keep the actual closed proper
operator group, its group of all source lifts, and
every coordinate change of C.48.
In (C.48.19) replace only the median base map by
(C.52.5), leaving the fiber interpolation as displayed:
\[
 (u,r)\longmapsto
       \bigl((1-s)h_t(u,r)+su,\ B_{s,t}(r)\bigr).        \tag{C.52.6}
\]
Here the choice of target coordinate is the exact
reparametrization already established in C.48.
We check why all of its operator equations survive
this replacement.

For the unit fiber translation, the fiber equation
$h_t(u+1,r)=h_t(u,r)+1$ is unchanged by the base
replacement. For an operator $(\epsilon u,r+d)$,
the equations are
$h_t(\epsilon u,r+d)=\epsilon h_t(u,r)$ and
$b_t(r+d)=b_t(r)+d$. The first is preserved by
the same convex interpolation; the second says
that translation by $d$ centralizes $T$, so it
holds for (C.52.5).
For the real fiber kernel and a further
$(\lambda u,r+d)$, the fiber function is
$h_t(u,r)=u+\phi_t(r)$ with
$\phi_t(r+d)=\lambda\phi_t(r)$. Scaling $\phi_t$
preserves this equality for either sign of $\lambda$,
and the base comparison is again the same one.

In the normal forms with all real base translations
or a dense lattice base image, the original base
loop is $b_t(r)=r+c_t$. Its selected family is
$B_{s,t}(r)=r+(1-s)c_t$ and commutes with every base
operator. The fiber equations, respectively
independence of $r$ or the exact affine lattice
relation (C.48.21), are preserved by (C.52.6).
The trivial operator case imposes no extra equation.
These are the C.48 normal forms outside its
transitive two-dimensional case. That latter case
already lies in $U_2\subset X_*$ and is absent from
the open orbits of the present claim.

Thus the modified family commutes with the exact
group of all operator lifts, with its original
restriction kernel retained downstairs.
The fiber and base inverse-continuity argument
following (C.48.22) applies unchanged to (C.52.6):
the base family has continuous inverses by
(C.52.5), and fixed bracketing fiber values bound
the increasing convex-interpolation inverses.
The homotopy descends to $O$, with the asserted
endpoint identities. All other actual quotients
are reparametrizations of the one selected
partition; naturality proves (C.52.2) for them.

For an empty family, C.32--C.33 supply a based
null-homotopy in the exact proper-operator
centralizer. There is then no nontrivial actual
quotient equation to check. This finishes the
intrinsic assertion for all three cardinalities.

\emph{Cocompactness and the entire frontier.}
Now assume $O/Q_O$ compact and put $Q=Q_O$.
The first part of C.51 supplies a compact
$C\subset O$ with $QC=O$. It also proves that
every nontrivial frontier orbit gives an actual
quotient in the specified family, after the
explicit basepoint comparison (C.51.8).
In particular an empty family makes all frontier
points fixed.

Define $\overline K^O$ by the just-constructed
$K^O$ on $O$ and by (C.52.1) on $\partial O$.
We first prove its limit along points of $O$;
we do not assume transverse continuity of
(C.52.1). Let
\[
 x_n\in O,\quad x_n\to z\in\partial O,\qquad
                       (s_n,t_n)\to(s,t).              \tag{C.52.7}
\]
Write $x_n=q_ny_n$ with $y_n\in C$ and pass to
a subsequence $y_n\to y\in C$.
The compact set consisting of $C$ and all its
images and inverse images under $K^O$ is exactly
the type of set used in (C.51.11).
Approximate $q_n$ uniformly on this set to
accuracy $1/n$ by restrictions of actual
$h_n\in H_O$, as in (C.51.12).
The open evaluation map of $L=\overline G$ gives
$l_n\to1$ with $l_ny=y_n$. Exact commutation
gives $h_ny=l_n^{-1}h_ny_n\to z$.

The C.42 evaluation homeomorphism, with the
basepoint comparison (C.51.8), now gives
$h_n|_O\to\Phi_z^y$ uniformly on compact subsets.
For a nontrivial target this is an actual
quotient covered by (C.52.2); for a fixed target
it is constant. Commutation with $Q$ therefore
gives precisely
\[
 \begin{split}
 K^O_{s_n,t_n}(x_n)
   &=q_nK^O_{s_n,t_n}(y_n)\\
   &\longrightarrow
       \Phi_z^y K^O_{s,t}(y)
        =R^{Gz}_{s,t}(z). 
 \end{split}                                             \tag{C.52.8}
\]
The approximation error in the second line is
less than $1/n$ before passage to the limit.
Compactness of $C$ and uniqueness of the value
give (C.52.8) for the entire sequence.
This repeats C.51's exact operator comparison
with the new quotient equation; no element of
the restricted closure is assigned an
unproved global extension.

\emph{The diagonal argument at frontier points.}
Take now $z_n\in\partial O$ with $z_n\to z$
and $(s_n,t_n)\to(s,t)$.
For each fixed triple $(s_n,t_n,z_n)$,
(C.52.8) lets us choose $x_n\in O$ such that
\[
 |x_n-z_n|<1/n,\qquad
 \bigl|K^O_{s_n,t_n}(x_n)
            -R^{Gz_n}_{s_n,t_n}(z_n)\bigr|<1/n.         \tag{C.52.9}
\]
Then $x_n\to z$, so applying (C.52.8) to this
new sequence proves
$R^{Gz_n}_{s_n,t_n}(z_n)\to R^{Gz}_{s,t}(z)$.
This establishes the needed joint continuity
on this entire frontier, even across different
target types in (C.52.1).
Splitting a sequence into its interior and
frontier terms proves joint continuity on
$[0,1]^2\times\overline O$.

The same proof applies to the inverse family.
The compact approximation set already contains
all inverse images, and (C.52.2) implies the
exact inverse equation
$\Phi_z(K^O_{s,t})^{-1}=(R^{Gz}_{s,t})^{-1}\Phi_z$.
Use this equation in (C.52.8), then repeat the
diagonal choice (C.52.9) for inverses.
Thus the piecewise inverse is jointly continuous.
All maps are bijections on $O$ and on each
frontier orbit, with these actual inverses;
they are homeomorphisms of the closure.

\emph{Agreement with the fixed exponential region.}
Every group image on an open orbit of $U_2\cup U_3$
is metabelian, by C.36 and C.49. Hence, for all
$g_1,g_2,g_3,g_4\in G$, the plane homeomorphism
$[[g_1,g_2],[g_3,g_4]]$ fixes $U_2\cup U_3$
pointwise. Continuity makes it fix $X_*$ pointwise.
The set $X_*$ is $G$-invariant. Therefore every
nonopen orbit $B\subset X_*$ has metabelian
actual image $G|_B$.

An open orbit outside $U_2\cup U_3$ is disjoint
from $X_*$, because it is an open neighborhood
disjoint from that union. Thus the intersections
$\overline O\cap X_*$ occurring here contain only
nonopen orbits. On a point orbit both
contractions are the identity. On every other
such orbit, (C.52.1) and C.50 use the same
actual intrinsic Lie image and the same
based target-cover path. Exponential naturality
therefore gives equality, including the
intermediate real circle angles described above.

Distinct open-orbit closures meet only in the
nonopen locus. The present construction assigns
there the same rule (C.52.1), so their contractions
and inverses agree. Conjugation by $h\in H$
identifies the actual target images and their
based paths, preserves metabelianness, and
intertwines both branches of (C.52.1).
Choose the intrinsic construction on one
representative of each $H$-orbit of components
and transport it by $h$. Commutation with the
full stabilizer makes this choice independent
of the transporting element, exactly as in
(C.51.15).

Finite closed-set pasting now gives the
homeomorphism families and inverse families
on (C.52.3). A general $h$ intertwines this
union with the union indexed by $h\mathcal I$;
when it preserves the union, it commutes with
the displayed contraction. The family on
$X_*$ is exactly the prescribed one throughout.
The diagonal proof concerns a single fixed
closure. It supplies no uniform control over
infinitely many different compact fundamental
sets, and no compact fundamental set for the
remaining noncocompact quotients. These are
still retained parts of C.12.
\end{proof}

\begin{claim}{C.53}{Proved}{Global quotients and finitely many operator representatives}
Use the ambient-null based plane loop $f_t$, its generated group
$G$, the commuting operator group $H$, and $L=\overline G$ of
C.42. For every open $G$-orbit $O$, choose $x\in O$ and
put $J=L_x$. Define $\mathscr F_O^{\rm all}$ to be the family
of fiber partitions on the common universal covering of $O$
of all real lifts of
\[
 \Phi_z^x:O\longrightarrow Gz,\qquad \Phi_z^x(lx)=lz,
 \quad z\in\operatorname{Fix}(J),\quad
                  Gz\text{ nonopen and nontrivial}.
                                                        \tag{C.53.1}
\]
The target points need not belong to $\partial O$.
Let $U^\dagger$ be the union of all open orbits with either
$\dim Q_O=2$ or $|\mathscr F_O^{\rm all}|\geq3$, and set
$X^\dagger=\overline{U^\dagger}$.
There is a jointly continuous based contraction on
$X^\dagger$, with jointly continuous inverse and exact
$H$-commutation, extending the prescribed contraction on
$X_*$ in C.50. On every nonopen orbit it has the
corresponding exponential value of C.52.1.

Now choose finitely many open orbits
$O_1,\ldots,O_m$ disjoint from $X^\dagger$, each with compact
$O_j/Q_{O_j}$, and put
\[
 U=\bigcup_{j=1}^m\ \bigcup_{h\in H}hO_j,
 \qquad X=\overline U.
                                                        \tag{C.53.2}
\]
The original loop has a based contraction through
homeomorphisms of $X^\dagger\cup X$, jointly continuous
with its inverse, commuting with $H$ and equal to the
specified exponential contraction on $X^\dagger$.
On every nonopen orbit in $X$ its value is exactly the
hybrid prescription $R^B_{s,t}$ of C.52.1.
This assertion allows infinitely many different open
components in (C.53.2), with arbitrary accumulation.
It retains the stated finite number of $H$-representatives
and compactness of their actual operator quotients.
\end{claim}
\passage{Section23, folios140-141, scans279-281; Part C(b), infinite-operator zero-level comparison, global quotient maps and arbitrary accumulation under finitely many operator representatives}


\begin{proof}
\emph{The global quotient family and its exact comparisons.}
C.42 is already stated for every $z\in\operatorname{Fix}(J)$
with nontrivial line or circle orbit; it imposes no
frontier condition on $z$. It gives the actual continuous
open map (C.53.1), its intrinsic target topology, the
common covering construction, and its proper fiber lines.
C.43--C.45 use these maps and their transitive compatible
action to give the connected real fibers, product
coordinates and the exact three pairwise coordinate
domains. Thus each of those conclusions applies to
(C.53.1) with the same maps.

Changing the source basepoint also preserves the complete
family. If $y=ax$ for $a\in L$, then
$L_y=aJa^{-1}$ and, for $z\in\operatorname{Fix}(L_y)$,
\[
                \Phi_z^y(ly)=lz
                     =\Phi_{a^{-1}z}^x(lax).
                                                        \tag{C.53.3}
\]
The targets are the same actual $G$-orbit, because
$Lz=Gz$ for each such quotient, as proved in C.42.
This proves the asserted identification of partitions,
including the original target coverings.

For $h\in H$ stabilizing $O$, choose $a_h\in L$ with
$a_hx=hx$. Exact commutation gives
$L_{hx}=L_x$, hence $a_hJa_h^{-1}=J$, and
\[
                   \Phi_z^x h=\Phi_{a_hz}^x.             \tag{C.53.4}
\]
The point $a_hz$ is fixed by $J$ and lies on the same
nonopen orbit as $z$. Therefore (C.53.4) permutes the
entire family $\mathscr F_O^{\rm all}$. For $hO=O'$, use
basepoint $hx$ in $O'$; then
\[
               \Phi_{hz}^{hx}h=h\Phi_z^x.              \tag{C.53.5}
\]
These equalities identify the families under operator
transport and prove that their cardinalities are
well defined and $H$-invariant.

\emph{Three global quotients give the same actual Lie group.}
Choose three distinct partitions in
$\mathscr F_O^{\rm all}$. The full three-foliation group
of C.47 is constructed exclusively from the real maps
in C.43--C.45 and their common transitive action.
Its proper-pair or ordered-loop argument therefore
applies to these exact maps and gives the same locally
compact Lie group of dimension two or three.
The based target lifts in C.49.2 are loops: for circle
targets C.29 uses the ambient-null plane loop and two
points of that circle, independently of the position
of $O$; real target endpoints are the identity.
The Baire transitivity proof C.49.3, its isotropy
calculation C.49.4--C.49.7 and its closedness argument
now apply to the same based source lifts. Their actual
generated group $N$ is consequently one of
$\mathbb R^2$, $\operatorname{Aff}^+(\mathbb R)$, or
$\mathcal A_2$, with exactly the topology, global
exponential and unique logarithm proved in C.49.

We verify the action comparison for every global
quotient, rather than only for frontier targets.
Use the closed Polish group
\[
 E_O=\{(l,n)\in L\times N:pn=(l|_O)p\}.
                                                        \tag{C.53.6}
\]
Its projection to $N$ is onto, since every $n$ is an
actual word in the lifted loop; the identical global
word supplies $l\in G$. It is an open homomorphism
of Polish groups. If $(l,1)$ is in its kernel, $l$
fixes $O$ pointwise. Equivariance of (C.53.1) then
makes $l$ fix its entire image $B=Gz$ pointwise:
\[
                         l\Phi_z^x(u)=\Phi_z^x(lu)
                                      =\Phi_z^x(u).
                                                        \tag{C.53.7}
\]
C.42 proves continuity of $L\to\operatorname{Homeo}(B)$
in the intrinsic compact-open topology. Its restriction
to $E_O$ therefore factors continuously through the
open quotient to an actual action of $N$ on $B$.
That image is exactly $G|_B$: actual loop words prove
both inclusions. It is metabelian and hence is the
translation, affine, or rotation Lie image of C.37.
The resulting continuous homomorphism and its unique
based covering lift give exactly
\[
 \Phi_z^x K^O_{s,t}=R^B_{s,t}\Phi_z^x,
 \qquad K^O_{s,t}=
        \exp_N\bigl((1-s)\log_N\widetilde f_t\bigr)
                       \text{ descended to }O.         \tag{C.53.8}
\]
Here $R^B$ is the exponential branch of C.52.1,
with the full based real angle for a circle.
The original source covering, every original quotient,
and every restriction kernel are retained.

\emph{The enlarged closed exponential region.}
For a dimension-two proper-operator orbit use the
based lifted contraction of C.36--C.37. For an orbit
with three global partitions use (C.53.8).
Where both constructions apply, their source groups
are the same actual group generated by the same
based lifted loop, as in C.49, so the maps agree.
Operator conjugation transports these groups and
intertwines their exponential maps, by (C.53.5) and
the based covering identity C.49.16.

Every source group just specified has the uniform
fractional-word bound $1913$: C.37 gives the stronger
bound three for the first two groups, and C.50.1--C.50.14
give the stated bound for $\mathcal A_2$.
Each word is a word in the original based path and
its inverse. In particular the full argument
C.50.15--C.50.19 applies to this enlarged union with
precisely the same compact set of actual global word
isotopies. We spell out the transfer of its remaining
inputs. All the source images are metabelian, so
each global commutator-of-commutators fixes
$U^\dagger$, hence its closure. At a nonopen limit
orbit the actual transitive image is therefore the
same translation, affine or rotation group classified
in C.37. For every integer $n$, the residual after
removing $\lfloor n(1-s)\rfloor$ endpoint powers is a
fractional power in the specified simply connected
source group, and thus has the same uniform word
bound. On a line the bounded residual logarithms
force the exponential limit. On a circle the source
residual word loop is null in that same simply
connected group; its exact two-point winding passes
to the limiting orbit, giving the real equality
C.50.19, including its integer part. These are all
inputs of C.50's limit identification and joint
continuity proof. Applying it also to the inverse
loop gives a jointly continuous inverse on
$X^\dagger$.

The original $U_2\cup U_3$ lies in $U^\dagger$.
On that original union the maps are the same
source fractional powers. Continuity extends this
equality to $X_*$. The operator identities also
extend from the dense union. This proves the first
assertion and fixes its prescription for the rest
of the proof.

\emph{Compatible intrinsic maps on the remaining representatives.}
Each $O_j$ in (C.53.2) has
$|\mathscr F_{O_j}^{\rm all}|\leq2$ and
$\dim Q_{O_j}\leq1$, by its disjointness from
$X^\dagger$. For a family of size one or two,
apply the construction of C.52 to this full global
family. Here is the exact transfer. Its source
coordinates and operator normal forms use
C.42--C.45, already verified above for (C.53.1).
Its full-partition invariance is (C.53.4).
The base-family replacement C.52.5 is natural under
every target reparametrization and exchange, and
retains every affine operator equation C.52.6.
No step of those intrinsic constructions uses the
position of the target point on a frontier.
They therefore give a based intrinsic contraction
$K^j_{s,t}$, jointly continuous with its inverse,
commuting with $H_{O_j}$ and with its actual
restricted closure $Q_{O_j}$, and satisfying
\[
                    \Phi_z K^j_{s,t}
                             =R^{Gz}_{s,t}\Phi_z
                 \quad\text{for every map in (C.53.1)}.
                                                        \tag{C.53.9}
\]
For an empty family, C.32--C.33 give the based
intrinsic contraction in that exact proper-operator
centralizer. Every constant equivariant quotient
also satisfies (C.53.9), since its target is fixed.

Discard repetitions among the finitely many
$H$-orbits represented by $O_j$, and transport
$K^j$ to each $hO_j$ by $hK^jh^{-1}$.
Commutation with the full stabilizer makes this
well defined. Its union gives a jointly continuous
homeomorphism family $K$ on the open set $U$;
components are open, and each carries its intrinsic
family and inverse. It commutes with $H$ exactly.

\emph{One compact comparison set for every accumulating sequence.}
For each representative, compactness of
$O_j/Q_{O_j}$ supplies a compact $C_j\subset O_j$
with $Q_{O_j}C_j=O_j$, by C.51's finite-cover
argument. Suppose
\[
 x_n\in U,\quad x_n\longrightarrow z\in X\setminus U,
                  \qquad(s_n,t_n)\longrightarrow(s,t).
                                                        \tag{C.53.10}
\]
The orbit $Gz$ is nonopen. Indeed an open orbit
containing $z$ would contain every sufficiently
late $x_n$, and so would itself be one of the
components in $U$, contradicting $z\notin U$.
There are only finitely many representatives.
On a subsequence, write
\[
                  x_n=h_nq_ny_n,
 \quad h_n\in H,\quad q_n\in Q_{O_j},\quad
                  y_n\in C_j,\quad y_n\longrightarrow y\in C_j
                                                        \tag{C.53.11}
\]
for one fixed $j$. Fix the compact subset of $O_j$
\[
 S_j=C_j\ \cup\!
       \bigcup_{(u,v)\in[0,1]^2}
       \bigl(K^j_{u,v}(C_j)\cup(K^j_{u,v})^{-1}(C_j)\bigr).
                                                        \tag{C.53.12}
\]
Both joint continuity assertions prove compactness.

Choose actual $a_n\in H_{O_j}$ so that
\[
 \sup_{w\in S_j}
          |h_na_nw-h_nq_nw|<1/n,
 \qquad g_n=h_na_n\in H.
                                                        \tag{C.53.13}
\]
For each fixed $n$, the first condition is an
open neighborhood of $q_n$ in
$\operatorname{Homeo}(O_j)$: composition with the
fixed ambient homeomorphism $h_n$ and restriction
to the compact set $S_j$ are continuous.
The defining density of the actual restrictions
$H_{O_j}\to Q_{O_j}$ supplies $a_n$.
Thus the estimate is made \emph{after} the possibly
distorting operator $h_n$; no uniform modulus of
continuity for the sequence $h_n$ is assumed.

The open $L$-orbit map supplies $l_n\to1$ with
$l_ny=y_n$, exactly as in C.51. Since $g_n$
commutes with $L$, (C.53.10)--(C.53.13) give
\[
                       g_ny=l_n^{-1}g_ny_n\longrightarrow z.
                                                        \tag{C.53.14}
\]
Every element of $L_y$ fixes this limit.
The evaluation homeomorphism C.42.1 therefore gives
\[
                g_n|_{O_j}\longrightarrow\Phi_z^y
                     \quad\text{uniformly on compact sets}.
                                                        \tag{C.53.15}
\]
The exact basepoint change (C.53.3) identifies this
map with one of the global quotients in (C.53.1)
when its target is nontrivial; otherwise it is
constant. This is the reason the enlarged family
is required in this argument.

Transport and commutation with $Q_{O_j}$ now yield
\[
 \begin{split}
 K_{s_n,t_n}(x_n)&=h_nq_nK^j_{s_n,t_n}(y_n),\\
 \bigl|K_{s_n,t_n}(x_n)
              -g_nK^j_{s_n,t_n}(y_n)\bigr|&<1/n.
 \end{split}                                            \tag{C.53.16}
\]
The inner points converge in $S_j$ to $K^j_{s,t}(y)$.
Equations (C.53.15) and (C.53.9) consequently give
\[
                   K_{s_n,t_n}(x_n)\longrightarrow
                         \Phi_z^yK^j_{s,t}(y)
                           =R^{Gz}_{s,t}(z).            \tag{C.53.17}
\]
Every subsequence admits a further subsequence of
this form, by the finite number of $j$ and the
compactness of their $C_j$. The value in (C.53.17)
is independent of every such choice. Hence the
entire sequence has that limit. If a representative
has an empty global family, (C.53.15) also proves
that every limit reached through its translates
is a fixed point.

Define the extension on $X\setminus U$ to be
$R^{Gz}_{s,t}(z)$. The diagonal approximation
C.52.9 now applies with the dense set $U$ in $X$:
for each varying boundary point and fixed parameter
pair choose a point of $U$ close to it and with
close image, using (C.53.17). Applying (C.53.17)
to the resulting sequence proves joint continuity
on all of $X$. This argument does not require
prior transverse continuity of the hybrid rule.

Use the inverse images already present in $S_j$
and the exact inverse version of (C.53.9) to
repeat (C.53.13)--(C.53.17) for inverse maps.
The same diagonal approximation then proves
joint inverse continuity on $X$.
The maps preserve $U$ and each nonopen limit
orbit, with the displayed maps as mutual
inverses. Thus they are homeomorphisms of $X$.
Their based-edge identities and $H$-commutation
extend from $U$ by continuity; the restriction
to every nonopen orbit is the prescribed $R^B$.

Finally, $U$ is disjoint from $X^\dagger$.
Every point of $X\cap X^\dagger$ lies on a
nonopen orbit; otherwise its open orbit would
meet $U$ and be a component of $U$.
The metabelian identity on $X^\dagger$ and the
first part of the proof show that both
contractions there use precisely the same
exponential branch of C.52.1, with the same
based covering path. They and their inverses
agree on the full intersection. Pasting the two
closed sets $X$ and $X^\dagger$ gives the
asserted contraction on their union, preserving
the prescribed family on $X^\dagger$.

The finite-representative argument does not
supply compact comparison sets for arbitrary
noncocompact quotients or infinitely many
$H$-inequivalent representatives. Nor does it
extend the family through an exterior region
of nonopen orbits. These remain exact parts
of the general C.12 obligation.
\end{proof}

\begin{claim}{C.54}{Proved}{Two metabelian quotient images give exponential closure contractions}
Use the actual global quotient family of C.53 for the
ambient-null based plane loop $f_t$.
Let $O$ be an open $G$-orbit admitting two distinct
partitions in that family whose actual intrinsic target
images are metabelian. On the universal covering
$p:X\to O$, the group $N$ generated by the based source
lifts is closed in $\operatorname{Homeo}(X)$.
With that actual topology, it is one of the following
simply connected groups:
\[
\begin{gathered}
 \R^2,\qquad \operatorname{Aff}^+(\R),\qquad
 N_{a,b}=\R^2\rtimes\R\quad((a,b)\ne(0,0)),\\
 \operatorname{Aff}^+(\R)\times\operatorname{Aff}^+(\R).
\end{gathered}                                           \tag{C.54.1}
\]
In $N_{a,b}$ the real parameter $t$ acts on the two
translation coordinates by
$\operatorname{diag}(e^{at},e^{bt})$.
Thus zero, equal, unequal, and opposite nonzero
weights are retained.
Each group has the explicit globally invertible
exponential map given below.

The based source exponential fractions descend and
extend to a jointly continuous based contraction
through homeomorphisms of the entire $\overline O$,
with a jointly continuous inverse.
They satisfy every actual global quotient equation,
are independent of the selected pair of metabelian
quotients, and are transported exactly by $H$.
They agree on common orbit frontiers and with the
prescribed contraction of C.53 on $X^\dagger$.
Consequently they paste over $X^\dagger$ together
with any finite collection of these closures.
No compactness of $O/Q_O$ is required.
Uniform continuity over an arbitrary infinite
collection of the new closures is not asserted here.
\end{claim}
\passage{Section 23, folios 140-141, scans 279-281; Part C(b), infinite-operator zero-level comparison; actual noncocompact two-quotient Lie reductions}


\begin{proof}
\emph{The actual target groups and the two coordinate domains.}
Choose two distinct real lifted quotient maps
$F_1,F_2:X\to\R$, with their original target covers
retained. Their based target paths $b_{i,t}$ are
loops, by C.29 for circle targets and directly
for line targets. As in C.52.4, the generated
groups $T_i=\langle b_{i,t}\rangle$ are transitive
on the real lines. The exact real-target
comparison in C.52 makes each $T_i$ a full
translation group or a full positive affine
group after a line conjugacy. These are the
actual groups with their intrinsic
homeomorphism-group topologies.

C.45 identifies the pair of quotient coordinates
with one of the three domains
$\R^2$, $\{x<y\}$, or $\{x<y<x+1\}$, by the
specified independent line reparametrizations.
The based source path has the coordinate form
$(b_{1,t},b_{2,t})$ and is a closed loop, since
the target endpoints are the identity and the
pair map is injective.

In the half-plane domain the two coordinate
paths are identical, and their generated group
is the same $T$. A transitive abelian line
group is conjugate to translations. Its
diagonal action cannot be transitive on an
open two-dimensional domain: in translation
coordinates the difference of the two
coordinates is constant along each orbit.
Thus $T$ is the full affine group.
Conjugate both coordinates by the same line
homeomorphism. The source action becomes
\[
          (x,y)\longmapsto(e^t x+u,e^t y+u),
                         \qquad x<y.                  \tag{C.54.2}
\]
The exact source change
$(x,y)\mapsto(x,\log(y-x))$ identifies this with
the simply transitive plane action
\[
                    (x,r)\longmapsto(e^t x+u,r+t).
                                                               \tag{C.54.3}
\]
Its actual generated group is the full
$\operatorname{Aff}^+(\R)$.

The band domain is impossible under the
present metabelian-target hypothesis.
Here again the two paths and their group $T$
coincide, and $T$ commutes with unit translation.
The full affine line group has trivial
centralizer: a map commuting with every
translation is a translation, and commuting
with a nontrivial dilation makes that
translation the identity.
Thus $T$ cannot be affine.
If $T$ is conjugate to translations, the
commuting unit translation is conjugate to
translation by a fixed nonzero constant.
The band therefore becomes a constant-width
band in those same coordinates, and the
diagonal translation action again preserves
the coordinate difference. It is not
transitive on the band, a contradiction.

It remains to classify the full-plane case.
Conjugate the two target groups independently
to their standard translation or affine
actions. The actual source group is then a
path-connected subgroup of
\[
 A=\operatorname{Aff}^+(\R)^2
   =\{(s_1,s_2,u_1,u_2):
        (x_1,x_2)\mapsto
             (e^{s_1}x_1+u_1,e^{s_2}x_2+u_2)\}.
                                                               \tag{C.54.4}
\]
We do not assume that this subgroup is closed.

The based source group is transitive on $X$.
Indeed its compact path-and-inverse set $P$
and the countable deck group $\Delta$ satisfy
$X=\bigcup_{\delta\in\Delta,n\geq1}\delta P^n x$,
as in C.49.3. Baire gives an open orbit.
Every other orbit is a deck translate,
because the based source path commutes with
every source deck map. Hence all orbits are
open, and connectedness of the covering plane
gives transitivity. Yamabe's theorem now
identifies this path-connected subgroup with
the analytic subgroup of its Lie algebra
\cite[pp.~13--14]{C_Yamabe1950}. We establish
closedness and the induced topology separately.

\emph{All transitive analytic subgroups of the product affine action.}
Write the Lie algebra of (C.54.4) as
$\mathfrak a=D\ltimes V$, where $D=\R^2$ is
the pair of log-scale directions, $V=\R^2$
is the translation plane, and
\[
 [(d,v),(e,w)]
       =(0,(d_1w_1-e_1v_1,\ d_2w_2-e_2v_2)).
                                                               \tag{C.54.5}
\]
Let $\mathfrak n$ be the generated group's
Lie algebra and put
\[
 S=\operatorname{pr}_D(\mathfrak n),\qquad
 W=\mathfrak n\cap V.
                                                               \tag{C.54.6}
\]
The subspace $W$ is invariant under the
diagonal action of $S$.
Transitivity makes the tangent evaluation
of $\mathfrak n$ onto the source plane
surjective at every point.

For $\dim S=0$, this forces $W=V$ and gives
all translations.
For $\dim S=1$, choose a lift
$Z=((a,b),v_0)$ of a nonzero generator of $S$.
The case $W=0$ has dimension one and cannot
be transitive. If $W=V$, subtracting $v_0$
from $Z$ gives exactly the Lie algebra and
connected group $N_{a,b}$ in (C.54.1).

Suppose instead $\dim W=1$. If $a=b\ne0$,
the induced affine vector field on $V/W$
has nonzero scalar linear part and hence
a fixed point. The corresponding affine
line in $V$ is invariant under the entire
connected subgroup, contradicting
transitivity.
If $a\ne b$, diagonal invariance makes $W$
one of the two coordinate axes.
For $W=\R(1,0)$, the second-coordinate
vector field of $Z$ is $by+(v_0)_2$.
If $b\ne0$, its fixed level is an invariant
affine line, again impossible.
If $b=0$ and $(v_0)_2=0$, the second
coordinate is constant on every orbit.
The remaining possibility is
\[
                  b=0,\qquad a\ne0,\qquad (v_0)_2\ne0.
                                                               \tag{C.54.7}
\]
A translation of the first source coordinate
removes $(v_0)_1$, and replacing the second
coordinate by $ay/(v_0)_2$ gives exactly
(C.54.3). The case $W=\R(0,1)$ is its
coordinate exchange. These are all
one-dimensional $S$ cases.

For $\dim S=2$, invariance under all diagonal
matrices makes $W$ either zero, one coordinate
axis, or $V$. If $W\ne V$, choose a coordinate
$i$ on which every vector of $W$ is zero.
For a linear section of
$\mathfrak n\to S=D$, let $c_i(d)$ be its
translation coordinate there. This is
independent of the chosen section modulo $W$.
Equation (C.54.5), whose bracket must belong
to $W$, gives
\[
                  d_i c_i(e)-e_i c_i(d)=0
                                  \quad(d,e\in D).
                                                               \tag{C.54.8}
\]
Choosing $e_i=1$ shows
$c_i(d)=c\,d_i$ for one constant $c$.
Every Lie-algebra vector field then preserves
the level $x_i=-c$. Its flows, and hence the
connected analytic subgroup that they
generate, preserve that affine line.
This contradicts transitivity.
Thus $W=V$, and $\mathfrak n=\mathfrak a$,
giving the full product affine group.

\emph{Closed groups, actual topology, and logarithms.}
The preceding list consists of explicit closed
subgroups of $A$. For $N_{a,b}$, the scale
vector lies on the closed line
$\R(a,b)$ and both translations are arbitrary.
The simply transitive affine case (C.54.3)
has, in its displayed source coordinates,
the closed equations $s_2=0$ and $u_2=s_1$.
The translation plane and $A$ itself are
closed as well.
Each listed group is connected, with exactly
the Lie algebra just obtained. Uniqueness
of the connected analytic subgroup with a
specified Lie algebra therefore identifies
the original generated subgroup with that
explicit closed group.

The topology of $A$ in $\operatorname{Homeo}(\R^2)$
is its ordinary parameter topology.
Evaluation at $(0,0),(1,0),(0,1)$ recovers
the translations and the two positive scale
factors continuously. A limit which is
still a homeomorphism cannot have either
scale zero. Conversely convergent parameters
give compact-open convergence of the maps
and their inverses by the affine formulas.
Thus $A$ is closed in this homeomorphism
group and has the asserted actual topology.
The source conjugacies used above preserve
compact-open group topology. This proves
the same assertions for the actual $N$ on
$X$, rather than only for its analytic
Lie-group structure.

Here are the promised exponential maps.
Set
\[
 \sigma(z)=\int_0^1 e^{rz}\,dr
   =
   \begin{cases}(e^z-1)/z,&z\ne0,\\1,&z=0.\end{cases}
                                                               \tag{C.54.9}
\]
This function is positive for every real
$z$. In $N_{a,b}$ the exact group law is
\[
 (t,u,v)(t',u',v')
     =(t+t',u+e^{at}u',v+e^{bt}v'),
                                                               \tag{C.54.10}
\]
and
\[
 \begin{split}
 \exp(t,u,v)&=(t,\sigma(at)u,\sigma(bt)v),\\
 \log(t,u,v)&=(t,u/\sigma(at),v/\sigma(bt)).
 \end{split}                                             \tag{C.54.11}
\]
The positive affine group has the same
formula with one translation coordinate
and weight one. For the full product $A$,
\[
 \begin{split}
 \exp(s_1,s_2,u_1,u_2)
      &=(s_1,s_2,\sigma(s_1)u_1,\sigma(s_2)u_2),\\
 \log(s_1,s_2,u_1,u_2)
      &=(s_1,s_2,u_1/\sigma(s_1),u_2/\sigma(s_2)).
 \end{split}                                             \tag{C.54.12}
\]
The translation exponential is the identity.
These formulas follow by integrating the
displayed affine vector fields, or directly
by differentiating their one-parameter
subgroups. Positivity of $\sigma$ proves
that each is a global continuous inverse
pair, including every zero weight.
The displayed group parameter spaces are
Euclidean spaces, so every group in
(C.54.1) is simply connected.

\emph{The exact action on the closure and on all global quotients.}
Use the original global closure $L$ and
the specified source covering to form
\[
 E_O=\{(l,n)\in L\times N:pn=(l|_O)p\}.
                                                               \tag{C.54.13}
\]
This is a closed Polish group.
Every $n\in N$ is an actual finite word
in the based source lifts; the identical
word in the original plane loop gives a
preimage in $E_O$. Projection to $N$ is
therefore onto and open.
Its kernel fixes $O$ pointwise and, by
continuity, fixes $\overline O$ pointwise.
Consequently the continuous action
$(l,n,z)\mapsto lz$ descends through the
open product quotient
$E_O\times\overline O\to N\times\overline O$.
This is the actual jointly continuous
$N$-action on the closure. Every value
has a global word preimage, so its orbits
are precisely the original $G$-orbits
there, with the original labels retained.

Define the source family by
\[
          n_{s,t}=\exp_N\bigl((1-s)\log_N\widetilde f_t\bigr).
                                                               \tag{C.54.14}
\]
The source deck group commutes with every
generator, hence with $N$ itself.
Formula (C.54.14) descends through $p$
and acts on $\overline O$ through the
just-proved quotient action. It and its
inverse are jointly continuous and satisfy
all based-edge identities.

For any actual global quotient
$\Phi_z:O\to B$, equivariance makes the
kernel of (C.54.13) fix its entire image:
$l\Phi_z(w)=\Phi_z(lw)=\Phi_z(w)$.
The intrinsic restriction continuity
proved in C.42 therefore descends through
the same open quotient to a continuous
$N$-action on $B$ with image exactly
$G|_B$, as in C.53.7.
All groups in (C.54.1) are metabelian.
Thus this target image is one of the
actual Lie groups in C.37.9.
The continuous homomorphism, and its
unique based covering lift for a circle,
intertwine (C.54.14) with the exponential
branch of C.52.1. Hence every actual
global quotient equation holds, not only
the two equations selected at the start.

\emph{Boundary comparison, choices, and finite pasting.}
For a nontrivial nonopen $G$-orbit
$B\subset\overline O$, the continuous
$N$-action just constructed has exactly
that orbit. For $z\in B$, its orbit map
gives a continuous bijection
$N/N_z\to B$ with the ambient topology
on the image. Its manifold dimension is
one: dimension zero would give a point
because $N$ is connected, while dimension
at least two would give an injective
map of a small two-dimensional disk into
the plane with interior image, by
invariance of domain. This would make
$B$ contain an ambient open subset,
contrary to its being nonopen.

On each compact interval of this
one-manifold the orbit map is an arc.
C.28 identifies its intrinsic interval
with that same compact arc. Consequently
the orbit bijection is a homeomorphism
onto the intrinsic line or circle $B$.
The induced action is continuous in
that intrinsic topology, and its
compact-open group topology, by the
continuous-action exponential law.
Its image is exactly $G|_B$ and is
metabelian. C.37.9 and the same based
covering comparison therefore identify
the boundary value of (C.54.14) with
the precise exponential branch of
C.52.1. Fixed boundary points remain
fixed.

The group $N$ was defined by the actual
based source lifts, independently of
the chosen quotient pair. The proof
has identified its actual topology and
its unique global exponential and
logarithm. A different quotient pair
therefore gives the identical family.
Changing the source covering conjugates
the based paths, with their specified
identity initial lift, and hence
conjugates the identical formula.
An operator $h\in H$ gives this same
comparison between the source covers
of $O$ and $hO$. Its lifted conjugacy
identifies the actual generated groups
and their paths. Lie homomorphism
naturality and uniqueness of logarithms
give exact transport of (C.54.14).

Distinct open-orbit closures meet only
in nonopen orbits. Their intrinsic
images and original based target
paths are the same $G|_B$, so the
exponential values just proved agree
there, including the full real
circle lifts.
On an open orbit already in
$U^\dagger$, the earlier construction
uses this same actual group generated
by the based source lifts.
Thus the two formulas agree on that
open orbit and its closure. For a
new orbit disjoint from $X^\dagger$,
its intersection with $X^\dagger$
consists only of nonopen orbits,
where both values are again the
same exponential prescription.
Finite closed-set pasting applies
to the maps and their joint inverses.
This proves the stated finite relative
contraction without a cocompactness
assumption. No uniform word bound for
the new three- or four-dimensional
groups has been used or asserted.
\end{proof}

\begin{claim}{C.55}{Proved}{The full local centralizer and its exact global limits}
Retain the ambient-null based plane loop $f_t$, its generated
group $G$, its commuting operator group $H$, and
$L=\overline G$ from C.42. On every open $G$-orbit define
the closed topological group
\[
 \mathcal C_O=C_{\operatorname{Homeo}(O)}(G|_O),
 \qquad H_O\longrightarrow\mathcal C_O,\quad h\longmapsto h|_O.
                                                        \tag{C.55.1}
\]
The topology is compact-open, including inverse convergence.
The kernel of the displayed homomorphism consists exactly
of operators fixing $O$ pointwise. Its image closure is
the previously used $Q_O$, a closed subgroup of
$\mathcal C_O$.

Let $U^\ddagger$ be the union of open orbits with
$\dim\mathcal C_O=2$ or
$|\mathscr F_O^{\rm all}|\geq3$, using C.53's full global
quotient family, and put $X^\ddagger=\overline{U^\ddagger}$.
There is a jointly continuous based contraction on
$X^\ddagger$, with jointly continuous inverse and
commutation with $H$, extending the specified contraction
on $X^\dagger$ in C.53. Its nonopen-orbit values are the
exponential branch of C.52.1.

Choose finitely many open orbits $O_1,\ldots,O_m$
disjoint from $X^\ddagger$ for which
$O_j/\mathcal C_{O_j}$ is compact. Set
\[
 U=\bigcup_{j=1}^{m}\ \bigcup_{h\in H}hO_j,
 \qquad X=\overline U.                                  \tag{C.55.2}
\]
There is a based contraction on $X^\ddagger\cup X$,
jointly continuous with its inverse, commuting with $H$,
and equal to the prescribed contraction on $X^\ddagger$.
On every nonopen orbit in $X$ it has the exact hybrid
value of C.52.1. The intrinsic contraction on each
representative commutes with all of $\mathcal C_{O_j}$.
Thus these conclusions apply with the full local
centralizer quotient, without requiring that its
elements extend to the plane.
\end{claim}
\passage{Section23; folios140-141; scans279-281; Part C(b), infinite-operator zero-level comparison}


\begin{proof}
\emph{The local group, its topology, and its global comparisons.}
The centralizer of $G|_O$ is also the centralizer of its
closure in $\operatorname{Homeo}(O)$. The latter is a
closed Polish group transitive on $O$. C.30 therefore
proves that $\mathcal C_O$ is locally compact, with a
free proper action. C.33's surface Lie argument applies
to this very action and makes $\mathcal C_O$ a Lie group
of dimension at most two. These conclusions use the
restriction image and its topology, not a discrete
topology on the original operators. Formula (C.55.1)
and its kernel follow directly from commutation on the
plane. Closure in $\operatorname{Homeo}(O)$ proves the
assertion about $Q_O$.

The global group $L$ preserves $O$, by the boundary
and inverse-limit argument of C.36. Its restriction is
continuous into $\operatorname{Homeo}(O)$: compact subsets
of $O$ and their inverse images remain compact subsets
of the plane. Hence every $q\in\mathcal C_O$ satisfies
\[
                      q(lu)=lq(u)\quad(l\in L,\ u\in O).
                                                        \tag{C.55.3}
\]
For example, choose $g_n\in G$ converging to $l$ globally,
restrict on compact subsets of $O$, and pass
$qg_nu=g_nqu$ to the limit. This proves (C.55.3)
without a plane extension of $q$.

For any $h\in H$, conjugation defines the topological
group isomorphism
\[
 \mathcal C_O\longrightarrow\mathcal C_{hO},\qquad
                         q\longmapsto hqh^{-1}.
                                                        \tag{C.55.4}
\]
The domain and codomain here are local homeomorphism
groups; the displayed conjugation uses only the actual
global operator $h$. It preserves dimension and the
compactness of the corresponding orbit quotient.

Fix $x\in O$, put $J=L_x$, and choose $a_q\in L$ with
$a_qx=qx$, using transitivity. Equation (C.55.3) and its
inverse give $L_{qx}=J$, and therefore
$a_qJa_q^{-1}=J$. For every global quotient in C.53.1,
\[
 \begin{split}
 \Phi_z^x q(lx)&=\Phi_z^x(la_qx)=la_qz,\\
                  \Phi_z^x q&=\Phi_{a_qz}^x.
 \end{split}                                            \tag{C.55.5}
\]
The target $a_qz$ is fixed by $J$ and lies in the same
actual nonopen orbit, since C.42 gives $Lz=Gz$.
Changing $a_q$ by its stabilizer ambiguity changes no
map in (C.55.5). Applying the formula also to $q^{-1}$
proves that the full global family is permuted by
$\mathcal C_O$. The exact transport between different
orbits is still C.53.5. All source and target coverings
in that family are therefore the ones already used
in C.43--C.48.

\emph{The enlarged exponential region.}
For $\dim\mathcal C_O=2$, the proof of C.36 applies with
$\mathcal C_O$ as its proper group. More explicitly, its
identity-component orbits are open by the free
dimension-two action and closed by properness; since
$O$ is connected, that component is transitive. Freeness
then gives $\mathcal C_O=\mathcal C_O^0$. The evaluation
map identifies it with $O$. Its full centralizer
\[
 N_O=C_{\operatorname{Homeo}(O)}(\mathcal C_O)
                            \cong\mathcal C_O^{\rm op}
                                                        \tag{C.55.6}
\]
is a simply transitive two-dimensional Lie group.
The actual restriction image $G|_O$ equals $N_O$:
transitivity supplies a word with any prescribed value
at one point, and evaluation is injective on $N_O$.

The resulting surjection $L\to N_O$ is continuous and
open. Its kernel fixes $\overline O$ pointwise, and it
fixes the image of every global quotient $\Phi_z$
pointwise by $l\Phi_z(u)=\Phi_z(lu)$. The two actions
therefore factor through the open quotient, with their
actual topology, exactly as in C.36.3 and C.53.6--C.53.7.
The universal covering group is $\mathbb R^2$ or
$\operatorname{Aff}^+(\mathbb R)$, with global exponential
and logarithm. The original based path lifts to a
closed path: C.32 gives its null-homotopy in
$\operatorname{Homeo}(O)$, and the evaluation
homeomorphism gives the identical statement in $N_O$.
Fractional powers in that specified cover define the
same closure contraction as C.36.6. On each global
nonopen quotient they give the exponential branch of
C.52.1, including the based real circle angle.

For $|\mathscr F_O^{\rm all}|\geq3$, retain exactly the
actual group and fractional powers of C.53.6--C.53.8.
They also commute with every $q\in\mathcal C_O$.
Indeed choose any lift $\widetilde q$ to the universal
cover of $O$. The path
$\widetilde q\widetilde f_t\widetilde q^{-1}$ is the
based lift of the same path $f_t|_O$, so uniqueness of
path lifting identifies it with $\widetilde f_t$.
It follows that $\widetilde q$ commutes with its entire
actual generated source group and with the displayed
fractional powers. Descent proves the assertion.
Where the two constructions apply together, they
use the same based path in the same actual simply
connected Lie group and agree.

All these source groups have the uniform fractional
word bound $1913$ of C.50; the two-dimensional groups
have the stronger bound three. Their images are
metabelian. The compact global word-isotopy family
C.50.15 consequently applies to their entire union.
For each positive integer $n$, its fractional residual
is in the same specified simply connected group and
has the same bound. On a limiting line orbit, the
bounded logarithm argument C.50.17--C.50.18 determines
the fractional limit. On a limiting circle orbit,
the null residual word loop has zero two-point winding,
so C.50.19 retains its real angle, including the
integer part. On fixed points every such word is the
identity. These are precisely the inputs of C.50's
dense-set limit and diagonal arguments. They prove
joint continuity on $X^\ddagger$ and, on applying them
to the inverse loop, joint inverse continuity.
Equation (C.55.4) and uniqueness of the based lift
give the required operator transport and commutation.

Since $Q_O$ is a closed subgroup of $\mathcal C_O$,
$\dim Q_O=2$ implies $\dim\mathcal C_O=2$. Thus
$U^\dagger\subset U^\ddagger$. The two contractions
agree on $U^\dagger$ by their actual source formulas,
and hence on $X^\dagger$ by continuity. This proves
the first extension assertion with its exact
prescription.

\emph{Intrinsic maps commuting with the full local group.}
Each remaining representative in (C.55.2) has
$\dim\mathcal C_{O_j}\leq1$ and at most two global
partitions. Apply the intrinsic construction of C.52
to the proper Lie group $\mathcal C_{O_j}$ and to this
full family. Every required operator property is now
proved: freeness, properness and the topology follow
from C.30; the Lie property follows from C.33; and
permutation of all quotient partitions is (C.55.5).
In the two-partition case, C.46 and C.52 use the exact
coordinate changes and target conjugacies. In the
one-partition case, C.48's proper-group normal forms
and C.52's fiber and base formulas retain every
relation of this very local group, including deck
relations. With an empty family, C.32--C.33 give the
intrinsic centralizer contraction using the same
ambient-null restriction. These intrinsic arguments
do not use a global extension of an operator.

They give a based family $K^j_{s,t}$ on $O_j$,
jointly continuous with its inverse, commuting with
$\mathcal C_{O_j}$, and satisfying
\[
              \Phi_z K^j_{s,t}=R^{Gz}_{s,t}\Phi_z
                  \quad\text{for every global quotient.}
                                                        \tag{C.55.7}
\]
Constant equivariant maps satisfy the identical
equation, with the fixed-point target prescription.
Remove duplicate $H$-representatives and set
$K=hK^jh^{-1}$ on $hO_j$. Its independence from the
chosen $h$ follows from $H_{O_j}|_{O_j}\subset
\mathcal C_{O_j}$. The union is a jointly continuous
homeomorphism family on the open set $U$, with
jointly continuous inverse and exact $H$-commutation.

\emph{Limits of the actual local equivariant maps.}
Compactness of $O_j/\mathcal C_{O_j}$ supplies a compact
$C_j\subset O_j$ with $\mathcal C_{O_j}C_j=O_j$.
For completeness, the quotient projection is open.
Relatively compact open neighborhoods in $O_j$ have
open quotient images covering the compact quotient;
a finite subcover and the union of their compact
closures give $C_j$.

Take $x_n\in U$ converging to $z\in X\setminus U$,
and $(s_n,t_n)\to(s,t)$. The orbit $Gz$ is nonopen,
because an open orbit containing $z$ would contain
all late $x_n$ and would be a component of $U$.
After a subsequence, the finite representative set
and compactness give one $j$ and
\[
 x_n=h_nq_ny_n,
 \quad h_n\in H,\ q_n\in\mathcal C_{O_j},\ y_n\in C_j,
 \qquad y_n\longrightarrow y\in C_j.
                                                        \tag{C.55.8}
\]
The maps
\[
 m_n:O_j\longrightarrow\mathbb R^2,
                  \qquad m_n(u)=h_nq_nu
                                                        \tag{C.55.9}
\]
are actual continuous $L$-equivariant maps by
(C.55.3) and global commutation of $h_n$ with $L$.
The open $L$-orbit map supplies $l_n\to1$ in $L$
with $l_ny=y_n$. Consequently
\[
                    m_n(y)=l_n^{-1}x_n\longrightarrow z.
                                                        \tag{C.55.10}
\]
Every element of $L_y$ fixes $z$. C.42's evaluation
homeomorphism, applied to the space of all continuous
$L$-equivariant maps $O_j\to\mathbb R^2$, gives
\[
                  m_n\longrightarrow\Phi_z^y
                         \quad\text{uniformly on compact sets}.
                                                        \tag{C.55.11}
\]
This is a direct use of the actual maps (C.55.9).
No element $q_n$ has been extended to the plane or
approximated by restrictions of global operators.

The points $K^j_{s_n,t_n}(y_n)$ converge to
$K^j_{s,t}(y)$ in a compact subset of $O_j$, by joint
continuity and compactness of the parameter square
times $C_j$. Transport, commutation, (C.55.11), the
basepoint comparison C.53.3, and (C.55.7) give exactly
\[
 \begin{split}
 K_{s_n,t_n}(x_n)
   &=h_nq_nK^j_{s_n,t_n}(y_n)\\
   &\longrightarrow\Phi_z^y K^j_{s,t}(y)
                      =R^{Gz}_{s,t}(z).
 \end{split}                                            \tag{C.55.12}
\]
Every subsequence has a further subsequence of this
form, and its limit is the same prescribed value.
Thus (C.55.12) holds for the whole sequence.
An empty global family makes each such limit orbit
a fixed point, since otherwise (C.55.11) itself
would supply a member of that family.

Define $K$ on $X\setminus U$ by this prescribed
hybrid value. At a varying boundary point and fixed
parameter, (C.55.12) supplies a nearby point of $U$
with nearby image. Choosing both distances less
than $1/n$ and applying (C.55.12) once more proves
joint continuity for boundary-only sequences.
Combining interior and boundary subsequences proves
joint continuity on the whole square times $X$.
The inverse form of (C.55.7) and the compact set of
all $(K^j_{s,t})^{-1}(C_j)$ give the identical argument
for inverses. Every map preserves the treated orbits
and has this displayed inverse; it is therefore a
homeomorphism of $X$. The based-edge identities and
$H$-commutation extend from the dense open set $U$.

Finally $X\cap X^\ddagger$ contains only nonopen
orbits. The metabelian identity holds on
$X^\ddagger$, so the hybrid value there is exactly
its prescribed exponential value, with the same
based covering lift. The two families and their
inverses agree on the intersection. Finite closed-set
pasting gives the contraction on $X^\ddagger\cup X$.

This argument still requires finitely many
$H$-representatives and compact quotients by their
full local centralizers. It supplies no extension
through arbitrary noncompact local quotients,
infinitely many inequivalent representatives, or
the remaining exterior nonopen region. Those parts
of C.12 remain under investigation.
\end{proof}

\begin{claim}{C.56}{Proved}{A uniform word bound for diagonal one-parameter affine groups}
For real weights $\alpha,\beta$, use the group
\[
 N_{\alpha,\beta}=\mathbb R\ltimes\mathbb R^2,\qquad
 (\tau,v)(\tau',w)
 =(\tau+\tau',v+\operatorname{diag}(e^{\alpha\tau},
                                     e^{\beta\tau})w).
                                                        \tag{C.56.1}
\]
For every continuous path $k:[0,t]\to N_{\alpha,\beta}$
with $k(0)=1$, put $P=k([0,t])\cup k([0,t])^{-1}$ and
$g=k(t)$. With the explicit uniform integer
\[
                         M=13\,363\,205,                \tag{C.56.2}
\]
every exponential fraction $\exp(r\log g)$, $0\leq r\leq1$,
belongs to $P^M$. The bound is independent of the
weights and the path. Zero, equal and opposite
weights are included.

In the plane-loop setting of C.54--C.55, let $U^\sharp$
be the union of $U^\ddagger$ and every C.54 orbit whose
actual generated source group has dimension at most
three, and put $X^\sharp=\overline{U^\sharp}$.
The specified exponential contractions extend to a
jointly continuous based contraction on $X^\sharp$,
with jointly continuous inverse and exact
$H$-commutation. It extends the prescribed family on
$X^\ddagger$ and retains the exponential value on
every nonopen orbit.
The finite-representative extension of C.55 also
holds relative to $X^\sharp$, for representatives
outside this larger set with compact quotients
by their full local centralizers.
\end{claim}
\passage{Section23; folios140-141; scans279-281; Part C(b), infinite-operator zero-level comparison}


\begin{proof}
\emph{The anchor and the exact product error.}
Write $k(s)=(\tau_s,v_s)$ and
$A=\max_{0\leq s\leq t}|\tau_s|$.
For $A=0$, the path is in the translation plane,
and C.37's three-letter additive path identity applies.
For $A>0$, choose a path value or its inverse
$h\in P$ whose first coordinate is $A$.
Conjugate all coordinates by a single translation
so that every nonzero-weight coordinate of $h$
has translation part zero. Its scale in that
coordinate differs from one, so this is possible.
Leave any zero-weight translation coordinate
unchanged. This common conjugacy adds no letters
to a word and preserves all group identities.

Put $\alpha_1=\alpha$, $\alpha_2=\beta$,
$x_i=\alpha_i A$, and $D_i=e^{x_i}-1$.
Direct multiplication gives the translation path
\[
 C_s=[h,k(s)]=(0,(D_1(v_s)_1,D_2(v_s)_2)).               \tag{C.56.3}
\]
For a zero weight both contributions to its
commutator coordinate vanish, including the
unchanged translation part of $h$.
Every $C_s$ or inverse costs at most four letters
of $P$. Conjugation by $h^j$ multiplies the
$i$th translation coordinate by $e^{jx_i}$
for every integer $j$.

Apply C.50.3's three-dimensional additive path
identity to $s\mapsto(\tau_s,v_s)$. It supplies
$m\leq5$, signs $\varepsilon_j\in\{1,-1\}$ and
times $s_j\in[0,t]$ with
\[
 r\tau_t=\sum_{j=1}^{m}\varepsilon_j\tau_{s_j},
 \qquad rv_t=\sum_{j=1}^{m}\varepsilon_jv_{s_j}.
                                                        \tag{C.56.4}
\]
Form the actual product $w=\prod_{j=1}^{m}k(s_j)^{\varepsilon_j}$.
For $\varepsilon_j=1$ put
$p_j=\sum_{\ell<j}\varepsilon_\ell\tau_{s_\ell}$;
for $\varepsilon_j=-1$ put
$p_j=\sum_{\ell\leq j}\varepsilon_\ell\tau_{s_\ell}$.
The full inverse formula in (C.56.1) proves
\[
 w=\left(r\tau_t,
       \left(\sum_{j=1}^{m}\varepsilon_j
           e^{\alpha_i p_j}(v_{s_j})_i\right)_{i=1,2}\right),
              \qquad |p_j|\leq5A.                     \tag{C.56.5}
\]
In particular, the additive identity has not
been substituted for a noncommutative product.

Define
\[
 \theta(z,r)=
 \begin{cases}(e^{rz}-1)/(e^z-1),&z\ne0,\\r,&z=0.\end{cases}
                                                        \tag{C.56.6}
\]
C.54.11 gives
$g^r=(r\tau_t,(\theta(\alpha_i\tau_t,r)(v_t)_i)_i)$.
Put $q=\tau_t/A$, $q_j=p_j/A$, so $|q|\leq1$ and
$|q_j|\leq5$. For a real variable $x$, define the
smooth functions
\[
 \begin{split}
 F_0(x)&=\frac{\theta(qx,r)-r}{e^x-1},\\
 F_j(x)&=-\varepsilon_j\frac{e^{q_jx}-1}{e^x-1}
                                      \quad(1\leq j\leq m),
 \end{split}                                            \tag{C.56.7}
\]
with their continuous values at zero.
By (C.56.4)--(C.56.5), the translation taking
$w$ to $g^r$ by left multiplication has
$i$th coordinate exactly
\[
 F_0(x_i)D_i(v_t)_i+
          \sum_{j=1}^{m}F_j(x_i)D_i(v_{s_j})_i.          \tag{C.56.8}
\]
Terms with $D_i=0$ vanish, just as the original
corrections in that coordinate vanish.

\emph{Uniform analytic bounds at every weight.}
Use $\sigma(x)=\int_0^1e^{ux}\,du$ from C.54.9.
For every real $x$,
\[
 \sigma(x)\geq e^{-|x|},\qquad
                       0\leq\sigma'(x)/\sigma(x)\leq1. \tag{C.56.9}
\]
For $|b|\leq5$,
$(e^{bx}-1)/(e^x-1)=b\sigma(bx)/\sigma(x)$
is smooth at zero, has absolute value at most
$5e^{6|x|}$, and has derivative bounded by
$30e^{6|x|}$.
Indeed its logarithmic derivative when nonzero
is the difference of a number of absolute
value at most $|b|$ and one in $[0,1]$.
The zero function for $b=0$ satisfies the
same bounds directly.

For $F_0$, regard $\theta(z,r)$ as the
probability of $[0,r]$ for the density
$e^{zu}/\sigma(z)$ on $[0,1]$.
Differentiating the two integrals gives
\[
 |\partial_z\theta(z,r)|\leq1,\qquad
 |\partial_z^2\theta(z,r)|\leq6.                         \tag{C.56.10}
\]
To check the second bound explicitly, its
expression in expectations for this density is
\[
 \begin{split}
 &\mathbb E(\mathbf1_{[0,r]}u^2)
 -2\mathbb E(\mathbf1_{[0,r]}u)\mathbb E(u)\\
 &\qquad
 -\mathbb E(\mathbf1_{[0,r]})\mathbb E(u^2)
 +2\mathbb E(\mathbf1_{[0,r]})\mathbb E(u)^2 .
 \end{split}
\]
All expectations lie in $[0,1]$, so the
absolute sum is at most six.
The first derivative is a difference of
two numbers in $[0,1]$.
Consequently
\[
 F_0(x)=
 \frac{q\int_0^1\partial_z\theta(uqx,r)\,du}{\sigma(x)}
                                                        \tag{C.56.11}
\]
has absolute value at most $e^{|x|}$.
The derivative of its numerator is at most
$6q^2\int_0^1u\,du\leq3$ in absolute value.
Equation (C.56.9) therefore bounds
$|F_0'(x)|$ by $4e^{|x|}$.
Every function in (C.56.7) thus satisfies
\[
                 |F(x)|\leq5e^{6|x|},\qquad
                 |F'(x)|\leq30e^{6|x|}.                \tag{C.56.12}
\]

\emph{Bounded interpolation at the two diagonal nodes.}
Set
\[
 \begin{split}
 A_8(x)&=(e^x+e^{-x})^8,\\
 B_8(x)&=(e^x-e^{-x})(e^x+e^{-x})^7,\qquad
 u(x)=\tanh x=B_8(x)/A_8(x).
 \end{split}                                            \tag{C.56.13}
\]
For $f(x)=F(x)/A_8(x)$, the estimates just proved give
\[
 |f(x)|\leq5,\quad |f'(x)|\leq70e^{-2|x|},\quad
        u'(x)=\operatorname{sech}^2x\geq e^{-2|x|}.
                                                        \tag{C.56.14}
\]
We used $A_8(x)\geq e^{8|x|}$ and
$A_8'(x)/A_8(x)=8\tanh x$.
Thus $f$ as a function of $u\in(-1,1)$ is
Lipschitz with constant seventy.

For distinct $x_1,x_2$ put
\[
 b=\frac{f(x_1)-f(x_2)}{u(x_1)-u(x_2)},\qquad
                         a=f(x_1)-bu(x_1).
\]
For equal nodes put $b=0$, $a=f(x_1)$.
In either case,
\[
 |a|\leq75,\quad |b|\leq70,\qquad
 F(x_i)=aA_8(x_i)+bB_8(x_i)\quad(i=1,2).                \tag{C.56.15}
\]
This includes a zero node and coincident weights.
The two node values come from the same function
$F$ in (C.56.7), as follows from the single
logarithmic scale parameter $\tau$.
No uncontrolled divided difference is used.

\emph{Realization by two explicit translation paths.}
Expand the Laurent polynomials
\[
 (z+z^{-1})^8=\sum_{j=-8}^{8}a_jz^j,\qquad
 (z-z^{-1})(z+z^{-1})^7=\sum_{j=-8}^{8}b_jz^j.           \tag{C.56.16}
\]
Their coefficients are integers with
$\sum|a_j|=256$ and $\sum|b_j|\leq256$.
Define continuous translation paths
\[
 \mathcal A_s=\prod_{j=-8}^{8}h^jC_s^{a_j}h^{-j},
 \qquad
 \mathcal B_s=\prod_{j=-8}^{8}h^jC_s^{b_j}h^{-j}.         \tag{C.56.17}
\]
They start at the identity.
Their $i$th translation coordinates are
$A_8(x_i)D_i(v_s)_i$ and $B_8(x_i)D_i(v_s)_i$,
respectively. All factors are translations, so
their order of multiplication is immaterial.

A conjugated copy of $C_s^{\pm1}$ costs at most
$2|j|+4$ letters of $P$. Repetition realizes
each integer coefficient. Hence each value or
inverse of either path costs at most
\[
                       (2\cdot8+4)\,256=5120.          \tag{C.56.18}
\]
The planar additive path identity expresses a
fraction in $[0,1]$ of $\mathcal A_s$ by three
signed values of $\mathcal A|_{[0,s]}$, and
likewise for $\mathcal B$.
For $|a|\leq75$ use seventy-five copies of the
fraction $|a|/75$, with inverse for a negative
sign. For $|b|\leq70$ use seventy copies.
Thus the translation
$\mathcal A_s^a\mathcal B_s^b$ realizes (C.56.15)
times the vector $(D_i(v_s)_i)_i$ and costs at most
\[
                       3(75+70)\,5120=2\,227\,200.     \tag{C.56.19}
\]
The word parameters need not be chosen
continuously: the translation paths to which
the additive identity applies are the actual
continuous paths (C.56.17).

Apply this to $F_0$ at $s=t$ and to every
$F_j$ at $s=s_j$. At most six such translations
give exactly (C.56.8). Their product followed
by $w$ is $g^r$, with length at most
\[
                       5+6\cdot2\,227\,200
                                         =13\,363\,205. \tag{C.56.20}
\]
Identity path values pad shorter words.
This proves (C.56.2), including all zero weights.

\emph{Exact algebraic verification.}
The executed program
\path{code/root/research_data/c56_diagonal_word_certificate.py}
expands the Laurent polynomials, verifies the
interpolation identities and coefficient bounds,
and checks the complete product correction with
rational arithmetic for every signed pattern
of length one through five in six weight regimes.
It also verifies the commutator and all Laurent
conjugation sums in those cases.
The regimes include equal, unequal, opposite and
zero weights. The analytic estimates and signed
path-selection result have their proofs above
and in C.50.3; these finite checks are not used
as their proof.

\emph{Accumulation and the relative extension.}
On $U^\ddagger$, C.55 gives the prescribed source
fractions with bound $1913$. For a new C.54 orbit
of dimension at most three, its actual source is
$\mathbb R^2$, $\operatorname{Aff}^+(\mathbb R)$,
or (C.56.1). The bounds three and (C.56.2)
therefore apply in its specified simply connected
source group. C.54 already proves its individual
closure action and every global quotient equation.
Operator conjugation transports these groups and
their based paths. Consequently the fractions
agree wherever definitions overlap.

Use $M$ in the compact global word-isotopy family
C.50.15. It is a finite union of continuous
images of a compact finite cube, for maps and
inverses. The metabelian identity holds on
$U^\sharp$ and hence on $X^\sharp$.
Every nonopen transitive image there is therefore
one of the line translation, line affine, or
circle rotation groups classified in C.37.9.

For every positive integer $n$, a sequence of
source fractions has the exact residual
\[
 \exp\bigl((nr_i-\lfloor nr_i\rfloor)
                         \log\widetilde f_{t_i}\bigr)
\]
in the same specified simply connected source.
Its exponent lies in $[0,1)$, so the residual
has the same uniform bound $M$.
Compact word parameters give bounded target
residual logarithms on a line.
On a circle, the source residual loop is null
in that same source group; its two-point winding
passes to the limit and is zero.
The real equation is exactly C.50.19, with $M$
in its bound. Division by $n$ identifies the
limiting fraction, including the integer part
of the based circle angle. Every word fixes
a fixed target point.
These are all inputs of C.50's dense-set limit
and diagonal arguments. They prove joint
continuity on $X^\sharp$.

Apply the same proof to the inverse loop.
The resulting family is the orbitwise inverse,
so it is jointly continuous and each map is
a homeomorphism of $X^\sharp$.
The based-edge and operator identities extend
from the dense $U^\sharp$.
The fractions agree with the old family on
$U^\ddagger$; continuity proves agreement
on $X^\ddagger$.

Finally, representatives outside $X^\sharp$
are outside $X^\ddagger$, so the intrinsic
construction and direct local-equivariant-map
limit C.55.7--C.55.12 apply with the same
compact fundamental sets and full global
quotient equations.
Their closure meets $X^\sharp$ only in
nonopen orbits. The actual image on
$X^\sharp$ is metabelian, so the hybrid
value on the intersection is exactly its
prescribed based exponential fraction.
Finite closed-set pasting of the maps and
inverses proves the last assertion.

The proof above uses one logarithmic scale parameter. C.58 proves the full product affine bound by a separate exact commutator correction in logarithmic coordinates. Arbitrary remaining
local quotients, infinitely many inequivalent
representatives and the exterior nonopen-region
extension remain parts of C.12.
\end{proof}

\begin{claim}{C.57}{Proved}{The full local centralizer with exactly two global quotients}
Use C.53's full global family and suppose
$|\mathscr F_O^{\rm all}|=2$ for an open $G$-orbit $O$.
Let $p:X\to O$ be its original universal covering,
$\Delta$ its actual deck group, and
$\mathcal C_O$ the full local centralizer in C.55.
Its group of all source lifts is
\[
 \widetilde{\mathcal C}_O
   =\{q\in\operatorname{Homeo}(X):
                   pq=\bar q p\text{ for some }\bar q\in\mathcal C_O\}.
                                                               \tag{C.57.1}
\]
In the exact pair-coordinate domains of C.45 the
following list is complete.

For the full-plane domain, put
$T_i=\langle b_{i,t}\rangle$ and
$D_i=C_{\operatorname{Homeo}(\R)}(T_i)$.
Each $D_i$ is, after a specified line conjugacy,
the identity group, the integer translation group,
or the real translation group. In those coordinates,
\[
 \widetilde{\mathcal C}_O=D_1\times D_2,\qquad
 \mathcal C_O\cong(D_1\times D_2)/\Delta,\qquad
 O/\mathcal C_O\cong\R^2/(D_1\times D_2).               \tag{C.57.2}
\]
All these are topological-group or quotient-space
identifications through the original covering.
Up to exchange of coordinates, the final quotient is
\[
\begin{array}{c|c|c}
 D_1\times D_2& O/\mathcal C_O &\text{compact}\\ \hline
 1&\R^2&\text{no}\\
 \Z\times1&S^1\times\R&\text{no}\\
 \R\times1&\R&\text{no}\\
 \Z^2&S^1\times S^1&\text{yes}\\
 \R\times\Z&S^1&\text{yes}\\
 \R^2&\{\ast\}&\text{yes}.
\end{array}                                              \tag{C.57.3}
\]

For the half-plane domain, $\Delta=1$ and
$\mathcal C_O=1$, so its quotient is a plane.
For the band domain $D=\{x<y<x+1\}$, put
$e(x,y)=(y,x+1)$. Then
\[
 \widetilde{\mathcal C}_O=\langle e\rangle\cong\Z,\qquad
 \Delta=1\ \text{or}\ \langle e^{2m}\rangle\ (m\geq1),
 \qquad
 \mathcal C_O\cong\langle e\rangle/\Delta.               \tag{C.57.4}
\]
Its quotient $O/\mathcal C_O$ is always the open
M\"obius band. Thus every noncocompact case in this
two-family list has full lifted local centralizer
$1$, $\Z$, or $\R$, with quotient a plane, open
cylinder, line, or open M\"obius band.
This identifies the groups and quotient maps;
it asserts no further contraction at their frontiers.
\end{claim}
\passage{Section 23, folios 140-141, scans 279-281; Part C(b), infinite-operator zero-level comparison; exact local groups and remaining noncocompact two-quotient domains}


\begin{proof}
\emph{The lift group and its one-dimensional factors.}
Every element of (C.57.1) commutes with the based
source path. Indeed its conjugation of that path
is another lift of the same downstairs path,
with the same identity initial value.
The path and its inverse have compact image.
Their generated source group is transitive by
the compact-word, countable-deck Baire argument
of C.49.3; that argument does not require the
generated group to be Lie.
The two based target paths are loops by C.29,
and their groups $T_i$ are transitive on their
real lines, as in C.52.4.

Equation (C.55.5) proves that every local
centralizer element permutes the full two-element
family of global quotient partitions.
Its source lifts consequently have coordinate-pair
or coordinate-exchange form, as proved in C.46.
The actual source deck maps preserve each selected
partition separately: the equation
$q_iF_i=\Phi_i p$ says that $F_i\delta$ and $F_i$
differ by the specified target deck map.
This retains the original deck group inside
(C.57.1).

For any transitive $T\leq\operatorname{Homeo}^+(\R)$,
its full centralizer $D$ acts freely. An element
fixing one point fixes its entire transitive
$T$-orbit. A decreasing real-line homeomorphism
has a fixed point, so no such element belongs
to $D$. The line version of C.30, used with
the closed transitive group $\overline T$,
makes $D$ locally compact, free and proper
in its actual topology.
The C.48 argument makes it a Lie group
of dimension at most one.

In dimension one its identity component
has open orbits, hence is transitive on
the connected line. Freeness excludes
additional components. Evaluation gives
the standard real translation action.
In dimension zero the orbit of a point
is a closed discrete ordered set.
A nonidentity free increasing action
has points of that orbit on both sides.
The element carrying the point to its
nearest successor generates the group:
its positive iterates cannot have a
finite accumulation point, by properness,
and there is no orbit point between
consecutive iterates. Freeness identifies
every group element with the unique
iterate having its value at that point.
An interpolation on one fundamental
interval and its translates conjugates
the action to integer translations.
The remaining case is the identity group.
This proves the asserted exact alternatives
for every $D_i$, including their topology.

\emph{The plane: the graph obstruction and every descending pair.}
Use the C.45 coordinate map $X\to\R^2$.
A coordinate exchange has the form
$q(x,y)=(a(y),b(x))$. Commutation with every
source generator would imply
\[
                   a b_{2,t}=b_{1,t}a
                                      \quad(0\leq t\leq1).
                                                               \tag{C.57.5}
\]
The proper nonempty graph
$\{(a(y),y):y\in\R\}$ is then preserved
by every generator, hence by the
transitive generated source group.
This is impossible. The argument allows
either orientation of $a$; no exchange
is omitted.

Every lift in (C.57.1) is therefore a
coordinate pair $(a(x),b(y))$.
Commutation with the path is exactly
$a\in D_1$, $b\in D_2$.
Thus the lift group is contained in
$P=D_1\times D_2$. In particular
$\Delta\subset P$.
The preceding one-dimensional
classification makes $P$ abelian.
Conversely every pair in $P$ commutes
with the original lifted path and
with $\Delta$. It therefore descends
through the original $p$ to a
homeomorphism of $O$ commuting with
every $f_t|_O$. Hence it belongs to
(C.57.1). This proves equality with $P$.

The source covering gives
$O\cong\R^2/\Delta$ and the exact
surjection $P\to\mathcal C_O$ has kernel
$\Delta$. Its quotient topology is
the actual compact-open topology:
the covering-chart comparison for
maps and inverses in C.48 proves
this for the group of all lifts,
and the group here is exactly that
group by (C.57.1).
The two maps from $X$ to the orbit
quotient have exactly the same fibers:
\[
      X\longrightarrow O\longrightarrow O/\mathcal C_O,
                   \qquad X\longrightarrow X/P.        \tag{C.57.6}
\]
Both are open quotient maps. Their
induced bijection is consequently a
homeomorphism, proving (C.57.2).
For completeness, the deck group
in these formulas is $1$ or $\Z$:
C.44 already proves that the source
of an actual nontrivial global quotient
is a plane or an open annulus.
Its particular embedding in $P$
remains the original deck embedding,
and is not discarded in the formula
for $\mathcal C_O$.

Choose the line conjugacies just
constructed for $D_1,D_2$ independently.
They give the six rows of (C.57.3).
The quotient maps, in the order
displayed there, are
$(x,y)$, $([x],y)$, $y$,
$([x],[y])$, $[y]$, and the constant
map. Here brackets mean the quotient
by the specified integer translations.
These formulas explicitly identify
the quotient spaces and their
compactness, independently of which
allowed source deck subgroup occurs.

\emph{The half-plane: no pair and no exchange survives.}
The normalized source path is
$(g_t(x),g_t(y))$ on $x<y$.
Write $T=\langle g_t\rangle$.
It acts transitively on ordered pairs
in this half-plane.
For any fixed $x$, its stabilizer
$T_x$ is transitive on each of
the two open rays about $x$:
use ordered pairs $(x,y)$ on the
right and $(y,x)$ on the left.
Consequently
$\operatorname{Fix}(T_x)=\{x\}$.

The coordinate functions of any
pair or exchange commuting with
this diagonal path centralize $T$.
They must be increasing by the
free-centralizer argument above.
In particular a decreasing exchanged
pair, which could otherwise preserve
the ordered half-plane, is excluded
by commutation and transitivity.
An increasing exchanged pair
$(a(y),b(x))$ cannot preserve the
half-plane: for fixed $x$, taking
$y\to+\infty$ makes $a(y)>b(x)$.

For a nonexchanged increasing pair
$(a(x),b(y))$, bijection of the
vertical ray over $x$ onto the
vertical ray over $a(x)$ gives
$b(x)=a(x)$.
Thus it is $(c,c)$ with $c$ centralizing
$T$. The point $c(x)$ is fixed by
$T_x$, so $c(x)=x$ for every $x$.
The whole lift group is therefore
trivial. It contains $\Delta$,
which is also trivial.
The source is the half-plane itself,
homeomorphic to a plane by
$(x,y)\mapsto(x,\log(y-x))$, and
its full local centralizer is trivial.

\emph{The band: the primitive period and the exact quotient.}
Now the source path is diagonal
on $x<y<x+1$, and $T$ commutes with
the unit translation $\tau(x)=x+1$.
Transitivity on the band makes
$T_x$ transitive on $(x,x+1)$.
Commutation with $\tau$ gives
transitivity on each
$(x+n,x+n+1)$, while every $x+n$
is fixed. Hence
\[
                         \operatorname{Fix}(T_x)=x+\Z.
                                                               \tag{C.57.7}
\]
For $c$ in the full line centralizer
of $T$, the point $c(x)$ belongs
to (C.57.7), so $c(x)=x+n$.
Both $c$ and $\tau^n$ commute with
the transitive group $T$.
Agreement at one point makes them
equal everywhere. Thus the full
line centralizer is exactly
$\langle\tau\rangle$.
In particular the given unit period
is primitive; a smaller positive
translation would contradict the
transitivity of $T_x$ on $(x,x+1)$.

The two coordinate functions of a
lifted local centralizer element
are therefore integer translations.
A nonexchanged pair
$(x+m,y+n)$ preserves the band
exactly when $m=n$, giving $e^{2m}$.
An exchanged pair
$(y+m,x+n)$ preserves it exactly
when $n=m+1$, giving $e^{2m+1}$.
Decreasing coordinate functions
are excluded by the same
free-centralizer argument.
This proves
$\widetilde{\mathcal C}_O\subset\langle e\rangle$.

Every source deck map preserves
the plane orientation induced
from the orientable open surface
$O$. The map $e$ reverses orientation,
whereas $e^2(x,y)=(x+1,y+1)$.
Thus $\Delta\subset\langle e^2\rangle$,
and it has one of the forms stated
in (C.57.4).
Every power of $e$ commutes with
the diagonal path and with this
deck group, so it descends.
The reverse inclusion in (C.57.4)
and the actual topological quotient
follow exactly as in the plane case.

Finally the explicit coordinates
\[
          (u,v)=(x+y,\ y-x-\tfrac12)
        \in\R\times(-\tfrac12,\tfrac12)                 \tag{C.57.8}
\]
send $e$ to $(u,v)\mapsto(u+1,-v)$.
Their inverse is
$x=(u-v-\tfrac12)/2$,
$y=(u+v+\tfrac12)/2$.
The quotient by $\langle e\rangle$
is therefore the open M\"obius
band, through this exact
glide-reflection quotient map.
For $\Delta=\langle e^{2m}\rangle$,
the source quotient identifies
$u$ modulo $2m$ and is the open
annulus; its local centralizer
is the cyclic group of order
$2m$. With $\Delta=1$ the source
is a plane and that centralizer
is infinite cyclic.
In both cases the final operator
quotient is the same open
M\"obius band.

The table and the two other
domains now give the last
classification in the claim.
All these coordinate comparisons
identify existing groups, covers,
and quotients. They do not
identify an arbitrary intrinsic
isotopy with an ambient extension
at their noncompact ends.
\end{proof}

\begin{claim}{C.58}{Proved}{A uniform word bound for the full product affine group}
Let
\[
 A=\operatorname{Aff}^+(\R)\times\operatorname{Aff}^+(\R).
\]
For a continuous path $k:[0,t]\to A$ with $k(0)=1$, put
$P=k([0,t])\cup k([0,t])^{-1}$ and $g=k(t)$.
Every exponential fraction satisfies the exact inclusion
\[
 \{\exp(r\log g):0\leq r\leq1\}\subseteq P^M,
 \qquad M=172\,496\,624\,025\,607.                       \tag{C.58.1}
\]
The integer is independent of the path and all scale coordinates,
including zero coordinates.

In the plane-loop setting of C.54--C.56, let $U^\flat$
be the union of $U^\sharp$ and all the remaining C.54
open orbits, whose actual generated source is $A$, and set
$X^\flat=\overline{U^\flat}$.
The prescribed source exponential fractions extend to a
jointly continuous based contraction through homeomorphisms
of $X^\flat$, with jointly continuous inverse and exact
$H$-commutation. They agree with the contraction on
$X^\sharp$ and retain the based exponential prescription
on every nonopen orbit in $X^\flat$.
The finite-representative extension of C.55 also holds
relative to $X^\flat$, for representatives outside that
set having compact quotients by their full local centralizers.
\end{claim}
\passage{Section 23, folios 140-141, scans 279-281; Part C(b), infinite-operator zero-level comparison and uniform frontier control}


\begin{proof}
\emph{Seven signed values in logarithmic coordinates.}
For every continuous based path $\gamma:[0,t]\to\R^4$
and $0\leq r\leq1$, the signed-measure result
\cite[Theorem~4.7 and \S4.3.2]{C_GrandoniRaviSinghZenklusen2014}
gives
\[
 r\gamma(t)=\sum_{j=1}^{m}\varepsilon_j\gamma(s_j),
 \quad m\leq7,\quad \varepsilon_j\in\{1,-1\},\quad
                         0\leq s_j\leq t.              \tag{C.58.2}
\]
For a polygonal path, apply it to the four coordinate
derivatives. Three circular intervals contribute six
signed endpoint values, and at most one wrap contributes
the total endpoint. The required halving count is explicit:
three intervals cut at four points give at most seven
pieces, so one half uses at most three intervals.
For the whole interval four cuts give at most five pieces,
so one half uses at most two. Complementation, halving
and compactness supply all fractions.
Uniform polygonal approximation and compactness of seven
times and their signs give (C.58.2) for continuous paths.

Apply (C.58.2) to the actual path $\log k(s)$, using the
global logarithm of C.54. Pad by $\log k(0)=0$ so that
$m=7$, and put
\[
 X_j=\varepsilon_j\log k(s_j),\qquad
 \ell_j=\exp X_j=k(s_j)^{\varepsilon_j},\qquad
 w=\ell_1\cdots\ell_7.
                                                               \tag{C.58.3}
\]
Thus $\sum_jX_j=r\log g$. All inverse letters are included
in the signed Lie coordinates $X_j$; the equality
$\exp(-X)=(\exp X)^{-1}$ is the one-parameter group identity.

\emph{The full product error and its antisymmetric identity.}
In one affine factor use logarithmic scale and group
translation coordinates
\[
 (a,v)(b,z)=(a+b,v+e^az),\qquad
 \sigma(a)=\int_0^1e^{ua}\,du,\qquad
                  \exp(a,u)=(a,\sigma(a)u).             \tag{C.58.4}
\]
The function $\sigma$ is positive and smooth, with
$\sigma(0)=1$. Write $X_j=(a_j,u_j)$ in this factor,
and set $T=\sum_j a_j$ and $P_j=\sum_{\ell<j}a_\ell$.
The pure translation taking $w$ to $\exp(\sum_jX_j)$
by \emph{left} multiplication is
\[
 c=\sum_{j=1}^7 F_j(a)u_j,\qquad
 F_j(a)=\sigma(T)-e^{P_j}\sigma(a_j),\qquad
                         a=(a_1,\ldots,a_7).           \tag{C.58.5}
\]
Indeed the product translation is
$\sum_je^{P_j}\sigma(a_j)u_j$, whereas the target
translation is $\sigma(T)\sum_ju_j$.
The exponential telescoping identity proves
\[
 \sum_j a_jF_j(a)
 =T\sigma(T)-\sum_je^{P_j}(e^{a_j}-1)=0.               \tag{C.58.6}
\]

For all $i,j$ define the smooth antisymmetric functions
\[
 f_{ij}(a)=\int_0^1 s\,
       (\partial_iF_j-\partial_jF_i)(sa)\,ds.
                                                               \tag{C.58.7}
\]
Differentiating (C.58.6) at $sa$ gives
$\sum_i a_i\partial_jF_i(sa)=-F_j(sa)/s$ for $s>0$.
Consequently
\[
 s\sum_i a_i(\partial_iF_j-\partial_jF_i)(sa)
                   =\frac{d}{ds}\{sF_j(sa)\}.
 \]
The endpoint at zero vanishes by continuity.
Integration therefore gives $F_j=\sum_i a_if_{ij}$ and
\[
 c=\sum_{i<j}f_{ij}(a)(a_i u_j-a_j u_i).                \tag{C.58.8}
\]
An explicit form, including its continuous interpretation
at $a_j=0$, is
\[
 \begin{split}
 f_{ij}(a)
   &=-\int_0^1s e^{sP_j}\sigma(sa_j)\,ds\\
   &=-\frac{\sigma(P_j+a_j)-\sigma(P_j)}{a_j}
          \quad(i<j),\qquad
 f_{ij}\big|_{a_j=0}=-\sigma'(P_j).
 \end{split}                                            \tag{C.58.9}
\]
Here the two derivatives of the common term $\sigma(T)$
cancel, and only $\partial_i(e^{P_j})$ remains.
Formula (C.58.9) also verifies (C.58.8) by telescoping.

With $[x,y]=xyx^{-1}y^{-1}$, direct multiplication gives
\[
 [\ell_i,\ell_j]
  =\bigl(0,\sigma(a_i)\sigma(a_j)
                         (a_i u_j-a_j u_i)\bigr).
                                                               \tag{C.58.10}
\]
Thus, putting
$G_{ij}(a)=f_{ij}(a)/(\sigma(a_i)\sigma(a_j))$,
the exact error in (C.58.5) is a sum of actual
commutator translations multiplied by $G_{ij}$.
In the product group the two factors evaluate
\emph{this same function} at their two scale vectors
$a^{(1)},a^{(2)}\in\R^7$.

\emph{Bounds smooth across every zero scale.}
Let $R=\sum_\ell|a_\ell|$.
The integral expression for $\sigma$ bounds it and its
first two derivatives at $x$ by $e^{|x|}$.
Each $F_j$ and each of its first and second partial
derivatives therefore have absolute value at most $2e^R$.
Equation (C.58.7) yields
\[
 |f_{ij}(a)|\leq2e^R,\qquad
 |\partial_\ell f_{ij}(a)|\leq\tfrac43e^R.
 \]
Also
$\sigma(x)\geq e^{-|x|}$ and
$0\leq\sigma'(x)/\sigma(x)\leq1$.
Since $i\ne j$, differentiation in one coordinate
affects at most one of the two denominator factors.
It follows that
\[
 |G_{ij}(a)|\leq2e^{2R},\qquad
                 |\partial_\ell G_{ij}(a)|<4e^{2R}.
                                                               \tag{C.58.11}
\]
No nonzero scale assumption occurs in these estimates.

\emph{Bounded interpolation at the two seven-dimensional nodes.}
For any one of the functions $G=G_{ij}$ put
\[
 \begin{split}
 \mathcal A(a)&=\prod_{\ell=1}^7
                         (e^{a_\ell}+e^{-a_\ell})^4,\\
 \mathcal B_q(a)&=\mathcal A(a)\tanh a_q
                                  \quad(1\leq q\leq7),\\
 z_\ell&=\tanh a_\ell,\qquad
                         \phi(z)=G(a)/\mathcal A(a).
 \end{split}                                            \tag{C.58.12}
\]
The map $a\mapsto z$ is a diffeomorphism onto $(-1,1)^7$.
We have $\mathcal A\geq e^{4R}$ and
$\partial_\ell\log\mathcal A=4\tanh a_\ell$.
Equation (C.58.11) gives
\[
 |\phi|\leq2,\quad
 |\partial_{a_\ell}(G/\mathcal A)|\leq12e^{-2R},
 \quad |\partial_{z_\ell}\phi|\leq12.                  \tag{C.58.13}
\]
For the last inequality use
$\operatorname{sech}^2a_\ell\geq e^{-2|a_\ell|}$.

For unequal nodes $z^{(1)},z^{(2)}$, choose $q$ maximizing
$|z_q^{(1)}-z_q^{(2)}|$, and define
\[
 d=\frac{\phi(z^{(1)})-\phi(z^{(2)})}
           {z_q^{(1)}-z_q^{(2)}},\qquad
                         b=\phi(z^{(1)})-d z_q^{(1)}.
 \]
Integration along the segment between the nodes in the
convex cube $(-1,1)^7$ gives $|d|\leq12\cdot7=84$
and $|b|\leq86$.
For equal nodes take $d=0$ and $b=\phi(z^{(1)})$,
with any $q$. In both cases the exact two equations are
\[
 G(a^{(\nu)})=
     b\mathcal A(a^{(\nu)})+d\mathcal B_q(a^{(\nu)})
 \quad(\nu=1,2),\qquad |b|\leq86,\quad |d|\leq84.
                                                               \tag{C.58.14}
\]
In particular, independent affine scales cause no
unbounded divided difference: the complete scale vectors
are the nodes of the same smooth function.

\emph{Realization by actual translation paths.}
For a fixed pair $i<j$, the path
\[
 C_{ij}(v)=
 [k(vs_i)^{\varepsilon_i},k(vs_j)^{\varepsilon_j}],
                    \qquad 0\leq v\leq1,               \tag{C.58.15}
\]
starts at the identity and lies in the translation
plane of $A$. Its values and inverses use at most
four letters of $P$.
Put $Z_\ell=e^{a_\ell}$ in (C.58.12).
Both $\mathcal A$ and $\mathcal B_q$ are integer Laurent
polynomials in seven variables.
Their exponents have absolute value at most four in
each variable, and both sums of absolute coefficients
are at most $2^{28}$.
Indeed each factor of $\mathcal A$ has coefficient sum
$2^4$, while the distinguished factor of $\mathcal B_q$
is $(Z_q-Z_q^{-1})(Z_q+Z_q^{-1})^3$ with coefficient
norm at most $2^4$.

For a Laurent monomial $Z_1^{n_1}\cdots Z_7^{n_7}$ use
the conjugator $h_n=\ell_1^{n_1}\cdots\ell_7^{n_7}$.
It uses at most $28$ letters.
Conjugation by $h_n$ multiplies a translation's
$\nu$th coordinate by
$\exp(\sum_\ell n_\ell a_\ell^{(\nu)})$, irrespective
of the conjugator's own translation.
Expand either polynomial and multiply the corresponding
integer numbers of conjugated copies of $C_{ij}(v)$
or its inverse.
The resulting continuous based translation paths
$\mathcal A_{ij}(v)$ and $\mathcal B_{ij,q}(v)$
multiply the $\nu$th coordinate of $C_{ij}(v)$ by
$\mathcal A(a^{(\nu)})$ and
$\mathcal B_q(a^{(\nu)})$, respectively.
All their factors are translations and commute.
Every value or inverse has the uniform cost
\[
                     L=(2\cdot28+4)2^{28}=60\cdot2^{28}.
                                                               \tag{C.58.16}
\]

The planar signed-path identity of C.37 represents a
fraction of either translation endpoint by three signed
values of its actual path.
Use $86$ copies of the fraction $|b|/86$, with inverse
for negative $b$, and $84$ copies for $|d|/84$.
The resulting translation
$\mathcal A_{ij}(1)^b\mathcal B_{ij,q}(1)^d$ has cost
at most $3(86+84)L$ and, by (C.58.14), has exactly the
two coordinate corrections
$G_{ij}(a^{(\nu)})\,\operatorname{pr}_{v_\nu}C_{ij}(1)$.
There are at most $\binom72=21$ pairs.
Their translations commute, and their product followed
by $w$ is exactly $\exp(\sum_jX_j)=g^r$.
Thus the total cost is
\[
 7+21\cdot3(86+84)\cdot60\cdot2^{28}
                         =172\,496\,624\,025\,607.      \tag{C.58.17}
\]
Zero path values pad shorter words.
The cuts and interpolation choices need not vary
continuously; all paths used to obtain their words are
the actual continuous paths just specified.
The equality is in the original group $A$, with the
full product error corrected on the left.

\emph{Exact executable checks.}
The accompanying program
\path{code/part_b/c58_affine_product_word_certificate.py}
verifies the product, inverse and commutator formulas,
the explicit radial coefficient and its telescoping
identity, the zero-scale values, the two-node equations,
the Laurent coefficient norms and the integer bound.
It also checks the entire corrected product for every
signed pattern of lengths one through seven in several
two-factor scale regimes, retaining zero and canceling
prefix scales by exact formal Laurent arithmetic.
The signed-measure selection, analytic estimates and
topological conclusions have the proofs in this text;
the finite checks verify their computable algebra.

\emph{Arbitrary accumulation and the remaining relative statement.}
C.54 supplies the actual closed simply connected source
group on each new orbit, its entire closure action, its
operator transport and every global quotient equation.
C.56 supplies the old contraction on $X^\sharp$.
The bound in (C.58.17) exceeds all the earlier bounds.
Therefore all source fractions on $U^\flat$ belong to
one compact family of actual global word isotopies,
obtained by replacing the length in C.50.15 by $M$.
This is a finite union of continuous images of the
compact cube of word parameters; maps and inverses
are both included. Each endpoint equality holds in
the specified simply connected source, including the
based lifts in the older two-dimensional cases.

All these sources are metabelian. Their group identity
passes from $U^\flat$ to $X^\flat$, so the nontrivial
nonopen intrinsic images there are precisely the line
translation, line affine and circle rotation images
of C.37. On them use the unique line logarithm or the
specified based real circle angle, and fix point orbits.

For a converging sequence of source fractions
$w_i=\exp(r_i\log g_i)$, every positive integer $n$
gives the exact source residual
\[
 w_i^n g_i^{-\lfloor nr_i\rfloor}
       =\exp\bigl((nr_i-\lfloor nr_i\rfloor)\log g_i\bigr).
                                                               \tag{C.58.18}
\]
The residual fraction lies in $[0,1)$ and has the
same fixed bound $M$.
The compact word-parameter argument of C.50 now applies
without a change of topology or a new boundary prescription.
On a line the limit of the residual logarithms is bounded,
and the limiting $w$ commutes with $g$, by passing the
source identity to every point of the target orbit.
Dividing the logarithmic residual equation by $n$
forces $\log w=r\log g$.

For completeness, the circle comparison retains the
real lift. If $A_W,T,C_n$ are the real endpoint angles
of the limiting word, original path and residual word,
respectively, the equality is
$nA_W-m_nT=C_n$, with $|m_n-nr|\leq1$ and
$|C_n|\leq M\max_u|\theta(u)|$.
The source loop
$W_i(u)^n f_{ut_i}^{-m_i}Z_i(u)^{-1}$
is null in its specified simply connected source,
by (C.58.18). Evaluate it at two points converging to
distinct points of the target circle.
Its difference-direction winding is zero and stays
zero in the limit. C.29 identifies that winding
with the integer ambiguity in the displayed real
angle equation. Thus $A_W=rT$, including an identity
circle endpoint with nonzero specified real angle.
At a fixed point every limiting word is the identity.

These are exactly the dense-set limits required by
C.50's diagonal approximation on the closure.
They prove joint continuity on $X^\flat$.
The inverse loop has the same sources and quotient
families, and its fractions are the orbitwise inverses,
so it supplies a jointly continuous inverse family.
The based-edge equations and operator commutation pass
from the dense union to its closure.
Agreement on $U^\sharp$ gives agreement on $X^\sharp$.

Finally, an outside representative in the last
assertion is also outside $X^\ddagger$.
C.55's direct equivariant-map evaluation and its
compact fundamental set therefore give the same
whole-closure limit for its intrinsic contraction.
Its intersection with $X^\flat$ consists of nonopen
orbits. The actual group image on that intersection
is metabelian, so its hybrid target value is exactly
the exponential prescription just proved.
Finite closed-set pasting of both families and their
inverses gives the asserted relative extension.

This includes all C.54 source groups through arbitrary
accumulation. It does not establish the remaining
relative extension through the noncompact local
quotients outside this region, through infinitely
many unrelated residual representatives, or through
the exterior nonopen locus. C.12 remains open.
\end{proof}

\begin{claim}{C.59}{Proved}{Localized quotient-kernel maps extend over every operator translate}
Retain the plane-loop setting, $G$, $L=\overline G$,
the original commuting group $H$, an open $G$-orbit $O$,
and the full local centralizer $\mathcal C_O$ of C.55.
Fix $x\in O$, $J=L_x$, and a compact set $D\subset O$.
Let $\mathcal V_O[D]$ be the subgroup of
$\operatorname{Homeo}(O)$ consisting of maps $a$ satisfying
\[
 \begin{gathered}
 aq=qa\quad(q\in\mathcal C_O),\qquad
 \Phi_z^x a=\Phi_z^x
 \quad\bigl(z\in\operatorname{Fix}(J),\
                   Gz\text{ nonopen and nontrivial}\bigr),\\
 \operatorname{Mov}(a):=\{u:a(u)\ne u\}
                                  \subseteq\mathcal C_O D .
 \end{gathered}                                           \tag{C.59.1}
\]
Use its actual compact-open group topology, including
inverse convergence. There is a continuous injective
group homomorphism
\[
 \begin{split}
 E_D:\mathcal V_O[D]&\longrightarrow
                C_{\operatorname{Homeo}(\R^2)}(H),\\
 E_D(a)|_{hO}&=h a h^{-1}\quad(h\in H),\\
 E_D(a)|_{\R^2\setminus\bigcup_{h\in H}hO}&=1 .
 \end{split}                                             \tag{C.59.2}
\]
Restriction to $O$ is its continuous inverse onto its
image. Thus every continuous compact parameter family
in $\mathcal V_O[D]$ extends jointly with its inverse,
even when $O/\mathcal C_O$ is noncompact and infinitely
many operator translates accumulate.

The identity path component of the kernel of all the
quotient equations in (C.59.1) is trivial when the
global family has two distinct fiber partitions.
The localized perturbations in that component can
therefore contribute only in the zero-partition and
one-partition cases. The assertion does not identify
the entire quotient kernel with its identity path component.
\end{claim}
\passage{Section 23, folios 140-141, scans 279-281; Part C(b), infinite-operator zero-level comparison and relative frontier extension}


\begin{proof}
\emph{The actual group and the transported maps.}
The set $\mathcal C_OD$ is closed: the proper action
in C.55 has a Hausdorff quotient, the image of $D$
in it is compact and closed, and its full inverse
image is $\mathcal C_OD$.
The support condition is exactly pointwise fixing of
the complement. It is closed in the local homeomorphism
group, as are the centralizer and quotient conditions.
They define a closed subgroup. In particular,
$\operatorname{Mov}(a^{-1})=\operatorname{Mov}(a)$.
The inverse satisfies every quotient equation by
composing $\Phi_z^x a=\Phi_z^x$ with $a^{-1}$.

Distinct $H$-translates of $O$ are disjoint open
$G$-orbits, because each operator commutes with $G$.
If $h_1O=h_2O$, the restriction of $h_2^{-1}h_1$
lies in the actual image
$H_O\to\mathcal C_O$ in C.55.1. Equation (C.59.1)
makes its conjugation of $a$ trivial.
Thus the first rule in (C.59.2) is independent of
the chosen operator. The kernel of $H_O\to\mathcal C_O$,
which fixes $O$ pointwise, causes no additional
identification. Algebraically these rules give a
bijection of the plane, inverted by the same rules
for $a^{-1}$. They respect products and commute with
every original operator, with its original label.

\emph{The compact-set limit at every frontier point.}
Let $K$ be a compact metric parameter space and
$\lambda\mapsto a_\lambda$ a continuous map
$K\to\mathcal V_O[D]$. Write $U=\bigcup_{h\in H}hO$.
Joint continuity inside $U$ follows on each of its
open orbits. Consider
\[
 \lambda_n\to\lambda,\qquad
       u_n\in U,\qquad u_n\to z\notin U.                \tag{C.59.3}
\]
Terms fixed by $E_D(a_{\lambda_n})$ already have the
required limit. For the remaining terms choose
\[
 u_n=h_nq_n y_n,\qquad h_n\in H,\quad
             q_n\in\mathcal C_O,\quad y_n\in D .
                                                               \tag{C.59.4}
\]
After a subsequence, $y_n\to y\in D$.
The open orbit evaluation $L\to O$ supplies
$l_n\to1$ with $l_ny=y_n$.
By C.55.3 each local $q_n$ commutes with $L|_O$,
so
\[
 m_n=h_nq_n:O\longrightarrow\R^2
 \quad\hbox{is $L$-equivariant},\qquad
                     m_n(y)=l_n^{-1}u_n\to z.          \tag{C.59.5}
\]
The evaluation homeomorphism C.42.1, now based at $y$,
therefore gives
$m_n\to\Phi_z^y$ uniformly on compact subsets of $O$.
Its range is exactly $Gz$.
That orbit is nonopen: an open orbit containing $z$
would be a neighborhood disjoint from every distinct
open orbit in (C.59.3).

For a point orbit the limit map is constant.
Otherwise it is one of the actual global quotient maps,
after the exact change of source basepoint.
Namely choose $l\in L$ with $lx=y$; then
\[
                 \Phi_z^y=\Phi_{l^{-1}z}^x.             \tag{C.59.6}
\]
The point $l^{-1}z$ is fixed by $J$ and has the same
nonopen orbit type, so (C.59.1) applies.
The points $a_{\lambda_n}(y_n)$ converge to $a_\lambda(y)$
and lie in the compact image of $K\times D$.
Consequently (C.59.4)--(C.59.6) imply
\[
 \begin{split}
 E_D(a_{\lambda_n})(u_n)
     &=h_nq_n a_{\lambda_n}(y_n)\\
     &\longrightarrow\Phi_z^y(a_\lambda(y))
       =\Phi_z^y(y)=z .
 \end{split}                                             \tag{C.59.7}
\]
Every subsequence admits the same limit argument,
so this proves convergence for the entire sequence.
Sequences outside $U$ are fixed.
This proves joint continuity on the whole plane.

The inverse family is continuous locally and has
the identical moving-point saturation by (C.59.1).
The same proof gives joint continuity of its transported
family. Thus the algebraic inverse above is a continuous
inverse, and every transported map is a homeomorphism.
Compact parameter families give continuous maps into
the ambient compact-open group.
For continuity of $E_D$ itself, apply this conclusion
to a convergent sequence and its limit in
$\mathcal V_O[D]$, which is a compact metrizable
parameter set. Both homeomorphism groups are metrizable,
so sequential continuity proves continuity.
Restriction to $O$ is continuous and recovers $a$,
proving the claimed topological embedding.
No compactness of the entire quotient was used;
only the one compact set $D$ in (C.59.4) was used.

\emph{The exact two-partition restriction.}
Let $a_s$, $0\leq s\leq1$, be a path from the identity
fixing all global quotient maps, and choose two maps
with distinct partitions.
Retain their original source and target covers
$p:\widetilde O\to O$ and $q_i:\R\to B_i$.
The corresponding real lifts $F_i$ satisfy
$q_iF_i=\Phi_i p$.
C.45 proves that
$(F_1,F_2)$ is injective on $\widetilde O$.
Lift the isotopy $a_s p$ to the based source isotopy
$\widetilde a_s$, with $\widetilde a_0=1$.
For each $i$ the quotient equation gives
\[
                q_iF_i\widetilde a_s=q_iF_i .          \tag{C.59.8}
\]
For a line target this is equality of the two real
maps. For a circle their difference is an integer
deck translation, continuous in $(s,u)$ and zero
at $s=0$. It is therefore zero throughout.
Thus both real coordinates of $\widetilde a_s(u)$
equal those of $u$. Pair injectivity makes
$\widetilde a_s=1$, and descent gives $a_s=1$.
This retains the discrete deck ambiguity for maps
outside the identity path component.

For representatives outside the invariant closed set
$X^\flat$, (C.59.2) fixes $X^\flat$ pointwise and gives
relative perturbations there. It does not fragment
the original loop: on a transitive open orbit the
union of that loop's moving sets is the whole orbit,
so imposing one compact quotient support on it would
force a compact quotient. The new assertion supplies
the exact extension for localized factors; existence
and convergence of a factorization sufficient for
C.12 are still unproved.
\end{proof}

\begin{claim}{C.60}{Proved}{All positive homotopy groups of the relative one-partition kernel vanish}
Use the original plane loop, its generated group $G$,
the commuting operator group $H$, and the full global
quotient families of C.53.
Let $W$ be any $H$-invariant union of open $G$-orbits,
each having exactly one partition in that family.
Define $\mathcal K_W$ to be the subgroup of
$C_{\operatorname{Homeo}(\R^2)}(H)$ consisting of maps
which fix $\R^2\setminus W$ pointwise, preserve each
of those open orbits, and satisfy
\[
                 \Phi_z a|_O=\Phi_z
 \quad\hbox{for every actual global quotient of }O\subset W.
                                                               \tag{C.60.1}
\]
Give it the induced compact-open group topology.
Then
\[
                 \pi_n(\mathcal K_W,1)=0
                         \qquad(n\geq1).               \tag{C.60.2}
\]
More precisely, every based loop $a_t$ in this group
has a jointly continuous based contraction $R_{r,t}$
in the same group, with jointly continuous inverse,
and with the pointwise control
\[
 \begin{split}
 R_{r,t}(x)&\in\{a_s(x):0\leq s\leq t\},\\
 R_{r,t}^{-1}(x)&\in\{a_s^{-1}(x):0\leq s\leq t\}
 \qquad(0\leq r,t\leq1).
 \end{split}                                             \tag{C.60.3}
\]
Here $R_{0,t}=1$, $R_{1,t}=a_t$, and
$R_{r,0}=R_{r,1}=1$.
Every based sphere family has a based contraction
with the analogous control by its entire compact
family of point images and inverse images.
No compactness of any local quotient, finiteness of
operator-orbit representatives, or separate ambient-null
hypothesis is required.
\end{claim}
\passage{Section 23, folios 140-141, scans 279-281; Part C(b), infinite-operator zero-level comparison and relative one-partition frontier extension}


\begin{proof}
\emph{The exact winding input and the source covers.}
For $W=\varnothing$ the group is trivial.
Otherwise choose a global quotient of one of its
orbits. Its target is a nontrivial nonopen $G$-orbit,
so it lies outside $W$ and contains two distinct
points. Every $a_t$ fixes these two points.
The difference-direction map on
$\operatorname{Homeo}^+(\R^2)$ induces an isomorphism
on fundamental groups, as proved in C.32.7.
Its value on this loop is constant. Thus $a_t$ is
automatically null-homotopic in the ambient plane
homeomorphism group.
C.32 consequently makes its restriction to every
preserved open orbit null-homotopic in
$\operatorname{Homeo}(O)$.

Retain the original source universal covering
$p:X\to O$ and one original target covering $q$,
with $qF=\Phi p$ as in C.43--C.44.
The based lift $\widetilde a_t$ of $a_t|_O$ is a
closed loop of homeomorphisms of $X$:
the restricted null-homotopy just proved lifts,
so its endpoint is the identity deck map.
Equation (C.60.1) gives
$qF\widetilde a_t=qF$.
For a circle target the difference between the real
lifts is a continuous integer deck translation,
initially zero; for a line it is zero directly.
In both cases
\[
                         F\widetilde a_t=F.             \tag{C.60.4}
\]
This retains both original covers and their labels.
The construction below uses only this closed based
source lift and (C.60.4) as its winding input.

\emph{The actual affine fiber coordinates and operator equations.}
Put $Q_O=\overline{\operatorname{im}
(H_O\to\operatorname{Homeo}(O))}$, with its actual
topology and the restriction kernel of C.30.
The loop $a_t|_O$ commutes with $Q_O$.
Its based lift commutes with every element of the
all-lifts group $\widetilde Q_O$ of C.48.2:
the commutator is deck-valued, continuous in $t$,
and initially the identity.

Apply the established C.48 normal forms to this
proper group and to the unique actual partition.
The transitive base group used there is still the
one supplied by the original $G$-action; it has
not been replaced by the kernel loop $a_t$.
The resulting exact product chart $X\cong\R^2$
has $F(u,v)=v$, and every lifted operator is affine
in its fiber coordinate:
\[
 \widetilde q(u,v)=
   \bigl(\alpha_{\widetilde q}(v)u+
         \beta_{\widetilde q}(v),\,\overline q(v)\bigr),
                  \qquad \alpha_{\widetilde q}(v)\ne0. \tag{C.60.5}
\]
This displayed form includes all the cases
C.48.7--C.48.17: real and cyclic fiber kernels,
discrete and real base images, the affine
two-dimensional group, and the normalized lattice.
In particular it includes every source deck map.
No affine assertion is made for the original
ambient coordinates.

By (C.60.4) write
$\widetilde a_t(u,v)=(h_t(u,v),v)$.
For each $t,v$ the function $h_t(\,\cdot\,,v)$ is
an increasing onto line homeomorphism, and
$h_0(u,v)=h_1(u,v)=u$.
Define
\[
 \widetilde R_{r,t}(u,v)
                   =\bigl((1-r)u+r h_t(u,v),\,v\bigr). \tag{C.60.6}
\]
These are homeomorphisms with jointly continuous
inverses by strict monotonicity, properness, and
the two-point inverse bracketing argument in C.48.
Their operator equations are exact. Namely
commutation of $\widetilde a_t$ with (C.60.5) reads
\[
 h_t(\alpha(v)u+\beta(v),\overline q(v))
                         =\alpha(v)h_t(u,v)+\beta(v).
 \]
Taking the displayed convex combination with the
identity proves the same equation for
$\widetilde R_{r,t}$, including negative $\alpha$.
It commutes with the original deck group and
descends to $R^O_{r,t}$ on $O$.
Every other global real quotient is an exact
line reparametrization of $F$, by the dependent
case of C.45. It too is fixed, and descent through
its specified target cover proves (C.60.1).

Choose the chart on one representative of each
$H$-orbit of open orbits and put
$R^{hO}_{r,t}=hR^O_{r,t}h^{-1}$.
The representative construction commutes with
$Q_O$, hence with the actual image of $H_O$.
This makes transport independent of the operator
chosen, including its restriction kernel.
Set $R_{r,t}=1$ outside $W$.

\emph{Actual prefix tracks control all frontier limits.}
Fix $(u,v)$ and $t$. The real path
$s\mapsto h_s(u,v)$, $0\leq s\leq t$, begins at
$u$ and ends at $h_t(u,v)$.
Its image contains the whole closed interval
between these endpoints.
The first coordinate in (C.60.6) lies in that interval.
After projection this proves the first inclusion
in (C.60.3) on $O$.

For the inverse, write
$g(u)=(1-r)u+r h_t(u,v)$ with $v$ fixed.
Its strict monotonicity shows that $g^{-1}(u)$
lies between $u$ and $h_t^{-1}(u,v)$:
the values of $g$ at these two endpoints bracket $u$.
The continuous inverse prefix
$s\mapsto h_s^{-1}(u,v)$ contains that interval.
Projection proves the second inclusion in (C.60.3).
Operator transport preserves both statements
because the original loop commutes with $H$.
No continuous selection of the prefix parameter
$s$ is needed.

For any sequence $x_n\to x\notin W$, joint
continuity of the original loop and compactness
of its parameter interval give
\[
 \sup_{0\leq s\leq1}|a_s(x_n)-x|\longrightarrow0,\qquad
 \sup_{0\leq s\leq1}|a_s^{-1}(x_n)-x|\longrightarrow0.
                                                               \tag{C.60.7}
\]
Here $a_s(x)=a_s^{-1}(x)=x$ throughout.
The two track inclusions therefore prove joint
continuity of the constructed maps and inverses
at $x$, uniformly in $r,t$, regardless of which
open orbits contain $x_n$.
Inside $W$ a point has its open-orbit neighborhood,
where the intrinsic construction is already
jointly continuous.
Thus both global families are jointly continuous.
They are inverse orbitwise bijections and so
are homeomorphisms of the plane.
They commute with every original operator,
fix the exterior, preserve each orbit and every
global quotient, and satisfy all based-edge
equalities. This proves (C.60.2) for $n=1$.

\emph{Based sphere families in every higher degree.}
Let $n\geq2$ and let
$a:S^n\to\mathcal K_W$ be continuous, with
$a_{\lambda_*}=1$.
For each original source covering $p:X\to O$ the map
\[
 X\times S^n\longrightarrow O,\qquad
                  (u,\lambda)\longmapsto a_\lambda(pu)
                                                               \tag{C.60.8}
\]
has simply connected, locally path-connected domain.
The covering lifting criterion gives its unique
continuous lift with value $u_*$ at $(u_*,\lambda_*)$.
On $X\times\{\lambda_*\}$ that lift is the identity.
Lift the inverse family from the same basepoint.
Composing the two lifts gives a lift of $p$ on the
connected product and agrees with the identity at
the basepoint; uniqueness proves both inverse
identities. Each lifted map is therefore a
homeomorphism and both families are jointly continuous.

Commutation with every lifted operator and equality
of the real quotient coordinate follow as before:
the corresponding deck elements are continuous in
the connected parameter space and trivial at
$\lambda_*$. Apply (C.60.6) with $h_\lambda$ in place
of $h_t$. It is a based sphere contraction, preserves
all the exact operator and quotient equations, and
transports over $H$ just as above.

For any $\lambda$, choose a path in $S^n$ from
$\lambda_*$ to $\lambda$. Its lifted fiber evaluation
contains the interval between $u$ and
$h_\lambda(u,v)$. The forward interpolation and
inverse bracketing consequently give
\[
 \begin{split}
 R_{r,\lambda}(x)&\in
                       \{a_\mu(x):\mu\in S^n\},\\
 R_{r,\lambda}^{-1}(x)&\in
                       \{a_\mu^{-1}(x):\mu\in S^n\}.
 \end{split}                                             \tag{C.60.9}
\]
No continuous choice of those parameter paths is
used. Compactness of the whole sphere family
replaces the time interval in (C.60.7), proving
joint frontier continuity of maps and inverses
uniformly in $r,\lambda$.
The resulting homotopy is based at $\lambda_*$
and constant at $r=0$, so it descends to a
continuous filling on the cone $CS^n\cong D^{n+1}$.
This proves (C.60.2) in all positive degrees.

The conclusion supplies based contractions of every
sphere family in the relative one-partition quotient kernel,
including arbitrary accumulation and noncompact
local quotients. It does not supply a lift of a
nontrivial deformation of the quotient actions
or a decomposition of the original loop into
such a lift and this kernel. No contraction of the
entire topological group is asserted. Those existence
requirements remain part of C.12.
\end{proof}

\begin{claim}{C.61}{Proved}{Compact prefix families prevent fiber intervals from returning to an open orbit}
Retain the original based plane loop $f_t$, its
generated group $G$, and the closed set
\[
 Y=\R^2\setminus\bigcup\{O:O\text{ is an open }G\text{-orbit}\}.
\]
For each open orbit under consideration, choose one
actual global quotient, its original universal
source cover $p_O:X_O\to O$, and one real target
lift $F_O:X_O\to\R$ as in C.43--C.44.
Fix these choices on each orbit.
Write $\widetilde f_t^O$ for the based source lift
of $f_t|_O$.

Suppose $x_n\in O_n$ converges to $x\in Y$, and
$s_n,t_n\in[0,1]$. Choose any source lift
$\widetilde x_n$ of each $x_n$ and require the
exact real equality
\[
 F_{O_n}(\widetilde f_{s_n}^{O_n}\widetilde x_n)
   =F_{O_n}(\widetilde f_{t_n}^{O_n}\widetilde x_n).
                                                               \tag{C.61.1}
\]
Let $I_n$ be the unique compact interval in that
proper real fiber joining those two lifted
endpoints, and put $J_n=p_{O_n}(I_n)$.
Then every compact set $K\subset\R^2\setminus Y$
is eventually disjoint from $J_n$:
\[
                  J_n\cap K=\varnothing
                      \quad\hbox{for all sufficiently large }n.
                                                               \tag{C.61.2}
\]
In particular every finite ambient accumulation
point of these interval images lies in $Y$.
The statement allows arbitrary variation of the
open orbit, either source covering type, and the
original target deck labels. It requires equality
in (C.61.1), rather than equality modulo a circle
target's deck group.
\end{claim}
\passage{Section 23, folios 140-141, scans 279-281; Part C(b), infinite-operator zero-level comparison and compact-strip frontier control}


\begin{proof}
C.43 gives connected proper real-line fibers and
an actual product chart
$X_O\cong\R^2$ in which $F_O(u,v)=v$.
Every based source lift preserves this partition:
its real quotient equation has the form
$F_O\widetilde f_t^O=b_tF_O$ with $b_t$ increasing.
The lift commutes with each original source deck
map $\delta$. Indeed,
$\widetilde f_t^O\delta(\widetilde f_t^O)^{-1}\delta^{-1}$
is a continuous deck-valued path starting at the
identity, and hence is the identity.
Also C.44 retains the exact equation
\[
                 F_O\delta=F_O+\eta_O(\delta),           \tag{C.61.3}
\]
where the label is zero for a line target and is
an integer for the specified circle cover.
Thus a deck map carries the interval in (C.61.1)
to the corresponding interval in its translated
real fiber, with no change of endpoint ambiguity.

For any one fixed $O$, the actual family
\[
 \mathcal P_O=
 \{\widetilde f_t^O(\widetilde f_s^O)^{-1}:
                                      s,t\in[0,1]\}
                                                               \tag{C.61.4}
\]
is compact in $\operatorname{Homeo}(X_O)$, with
inverse convergence. This follows from continuity
of the original lifted path and multiplication,
on the compact parameter square.
In the product chart its elements have the form
\[
            g(u,v)=(h_g(u,v),b_g(v)),                    \tag{C.61.5}
\]
where $h_g(\,\cdot\,,v)$ is increasing and onto.
For a convergent sequence $g_n\to g$ in
$\mathcal P_O$ and $v_n\to v$, these fiber maps
are uniformly proper along $v_n$. Explicitly,
given $M$, choose $A$ with $h_g(A,v)>M+1$.
Then $h_{g_n}(A,v_n)>M$ for all sufficiently
large $n$, and monotonicity gives
$h_{g_n}(u_n,v_n)>M$ whenever $u_n>A$.
The lower bound follows by choosing $A$ with
$h_g(A,v)<-M-1$.
Consequently
\[
 \begin{split}
 u_n\longrightarrow+\infty
   &\ \Longrightarrow\ h_{g_n}(u_n,v_n)\longrightarrow+\infty,\\
 u_n\longrightarrow-\infty
   &\ \Longrightarrow\ h_{g_n}(u_n,v_n)\longrightarrow-\infty.
 \end{split}                                             \tag{C.61.6}
\]
These estimates use the whole compact family
of lifted homeomorphisms and their actual
increasing fiber maps.

Suppose (C.61.2) fails. After a subsequence there
are $y_n\in J_n\cap K$ with $y_n\to y\in K$.
The point $y$ lies in some open orbit $O$.
Eventually $y_n\in O$, while $y_n\in O_n$.
Distinct $G$-orbits are disjoint, so $O_n=O$
on this tail. All source covers and the chosen
real quotient are now the same fixed maps.
Pass to $s_n\to s$ and $t_n\to t$.

Choose $\widetilde y_n\in I_n$ over $y_n$.
A covering neighborhood at $y$ supplies lifts
$\widehat y_n\to\widehat y$ in one compact source
chart. Choose the actual deck map $\delta_n$
with $\delta_n\widetilde y_n=\widehat y_n$.
Put
\[
 \begin{split}
 a_n&=\delta_n\widetilde f_{s_n}^O\widetilde x_n,\\
 c_n&=\delta_n\widetilde f_{t_n}^O\widetilde x_n,\qquad
 g_n=\widetilde f_{t_n}^O(\widetilde f_{s_n}^O)^{-1}.
 \end{split}                                             \tag{C.61.7}
\]
Commutation with the very same deck map
$\delta_n$ gives the exact equality $c_n=g_na_n$.
By (C.61.3), the interval $\delta_n I_n$ joins
$a_n$ to $c_n$ in one real fiber and contains
$\widehat y_n$.
Write
$a_n=(u_n,v_n)$ and
$\widehat y_n=(w_n,v_n)$.
Then $v_n\to v$ and $w_n\to w$ are finite.
Equation (C.61.1) also gives $b_{g_n}(v_n)=v_n$.

The first endpoints escape every compact subset
of $X_O$. Indeed,
\[
                 p_O(a_n)=f_{s_n}(x_n)\longrightarrow f_s(x).
 \]
The set $Y$ is $G$-invariant, so $f_s(x)\in Y$.
A convergent source subsequence would instead
project to a point of $O$.
Since $v_n$ is bounded, this proves
$|u_n|\to\infty$. Select either the positive
or the negative infinite subsequence.
Compactness of (C.61.4) gives $g_n\to g$,
and (C.61.6) forces the first coordinate of
$c_n=g_na_n$ to escape to the same end.
The real interval between these two coordinates
cannot contain the bounded $w_n$.
This contradicts
$\widehat y_n\in\delta_n I_n$ and proves
(C.61.2).

No uniform choice of charts on different
orbits was made: an alleged open accumulation
point forces a single actual orbit before
any chart is used.
The result does not bound accumulation along
$Y$, nor exclude escape to ambient infinity.
In particular it does not yet give the required
control near one specified nonopen orbit or
a lift of the nontrivial quotient contraction
left open after C.60.
\end{proof}

\begin{claim}{C.62}{Proved}{Every equal-level prefix interval is contained in a compact three-word image}
Retain the original based plane loop $f_t$ and its generated
group $G$. On every open orbit under consideration choose one
actual global quotient, the original universal source cover
$p_O:X_O\to O$, and the specified real target lift
$F_O:X_O\to\R$ of C.43--C.44. Write
$\widetilde f_t^O$ for the based source lift.
For $x\in O$, choose $\widetilde x\in p_O^{-1}(x)$ and
$s,t\in[0,1]$ satisfying the exact real equality
\[
 F_O(\widetilde f_s^O\widetilde x)
   =F_O(\widetilde f_t^O\widetilde x).
                                                               \tag{C.62.1}
\]
Let $I$ be the compact interval in that proper real fiber
joining the two displayed points. Then
\[
 \begin{split}
 I&\subset
 \{\widetilde f_u^O(\widetilde f_v^O)^{-1}
                   \widetilde f_a^O\widetilde x:
                          0\leq a\leq v\leq u\leq1\},\\
 p_O(I)&\subset\mathcal W_3(x),\qquad
 \mathcal W_3(x)=
 \{f_u f_v^{-1}f_a(x):0\leq a\leq v\leq u\leq1\}.
 \end{split}                                                   \tag{C.62.2}
\]
The same assertion holds for the original inverse path,
with projected word $f_u^{-1}f_vf_a^{-1}$.

Consequently, let $x_n\in O_n$ tend to any $x\in\R^2$,
allow the open orbit and the chosen actual quotient to vary,
and let $J_n=p_{O_n}(I_n)$ be any intervals obtained from
(C.62.1). These intervals remain in one compact subset of
the plane and satisfy
\[
 \sup_{y\in J_n}\operatorname{dist}
                (y,\mathcal W_3(x))\longrightarrow0.
                                                               \tag{C.62.3}
\]
Every accumulation point therefore lies in the same
$G$-orbit as $x$. In particular, at a point fixed by $G$,
\[
                    \sup_{y\in J_n}|y-x|\longrightarrow0.
                                                               \tag{C.62.4}
\]
All source and target covering labels in these statements
are the original ones. In particular, (C.62.1) cannot be
replaced by equality modulo the deck group of a circle
target.
\end{claim}
\passage{Section 23, folios 140-141, scans 279-281; Part C(b), infinite-operator zero-level comparison and full prefix-fiber interval control}


\begin{proof}
\emph{The actual source chart and its two sides.}
C.43 provides a homeomorphism
\[
 \chi:X_O\longrightarrow\R^2,
       \qquad F_O=\operatorname{pr}_2\chi.
                                                               \tag{C.62.5}
\]
We use this exact chart for the proof and retain $p_O$.
Thus a real fiber is a complete horizontal line, and its
two complementary components are precisely its strict
sublevel and superlevel sets. Every original source lift
preserves the fiber partition. Its induced real target
map is increasing, because it is joined to the identity
by the based lifted path. The following argument uses
these actual homeomorphisms and their actual inverse
maps; it does not assume that the auxiliary product flow
from C.43 commutes with them.

Swap $s,t$ if necessary so that $s\leq t$, and put
\[
 \gamma(r)=\widetilde f_r^O\widetilde x,
       \qquad c=F_O\gamma(s)=F_O\gamma(t),
       \qquad
 Z=\{r\in[s,t]:F_O\gamma(r)=c\}.
                                                               \tag{C.62.6}
\]
The set $Z$ is compact and contains both endpoints.
Every complementary interval $(a,b)$ has endpoints
in $Z$, and the continuous function
$F_O\gamma(r)-c$ has one strict sign throughout it.

\emph{A single excursion is filled by an actual word triangle.}
Fix such an interval $(a,b)$. Suppress the superscript
$O$ and define
\[
 A=\widetilde f_a\widetilde x,\qquad
 h_r=\widetilde f_r\widetilde f_a^{-1}
       \quad(a\leq r\leq b),\qquad
 g=h_b,\qquad B=gA.
                                                               \tag{C.62.7}
\]
In particular $h_a=\id$, $A,B\in F_O^{-1}(c)$, and
$h_rA=\gamma(r)$.
Write $\beta_r$ for the increasing target map induced
by $h_r$. Suppose first that the excursion is positive.
Then
\[
 \beta_r(c)>c\quad(a<r<b),\qquad
 \beta_r^{-1}(c)<c,\qquad \beta_b(c)=c.
                                                               \tag{C.62.8}
\]
Since $\beta_b$ is increasing and fixes $c$, the two
paths
\[
       r\longmapsto h_rA,\qquad
       r\longmapsto g h_r^{-1}A
                                                               \tag{C.62.9}
\]
lie strictly on opposite sides of the fiber for
$a<r<b$. The first goes from $A$ to $B$ and the second
from $B$ to $A$. A negative excursion reverses their
two sides and changes nothing else in the argument.

For completeness, these side conditions give a precise
winding calculation even when either path has
self-intersections. Let $q$ be any point in the open
fiber interval between $A$ and $B$, with $A\ne B$.
In the chart (C.62.5), each path can be homotoped,
relative to its endpoints, to a two-segment path in
its own strict half-plane: use straight interpolation
with the polygonal path whose middle vertex is
directly above or below the midpoint of $A,B$,
at unit distance from the fiber on the specified
side. Throughout the interpolation
all interior parameter values stay strictly on that
side and hence avoid $q$; the endpoints also avoid
$q$. The resulting two polygonal paths form the
boundary of a diamond containing $q$. Their
concatenation therefore has winding $1$ or $-1$
about $q$. This proves the same nonzero winding for
(C.62.9), without imposing simplicity on the original
paths.

Now take the closed triangle
\[
 T_{a,b}=\{(u,v):a\leq v\leq u\leq b\}
\]
and the continuous map
\[
 \begin{split}
 \Psi(u,v)&=h_u h_v^{-1}A\\
          &=\widetilde f_u\widetilde f_v^{-1}
                                      \widetilde f_a\widetilde x.
 \end{split}                                                   \tag{C.62.10}
\]
On its edge $v=a$ this is the first path in (C.62.9);
on its edge $u=b$ it is the second path, in the order
specified there; and on the diagonal $u=v$ it is
the constant $A$. Thus its boundary has the nonzero
winding just computed about every such $q$.
If $q$ were missing from $\Psi(T_{a,b})$, this very
triangle would null-homotope that boundary in
$X_O\setminus\{q\}$, contradicting the winding.
The endpoints are already images of vertices.
When $A=B$ the interval is a singleton and the
conclusion is immediate. We have proved in all
cases that
\[
 [\gamma(a),\gamma(b)]_{F_O=c}
       \subset \Psi(T_{a,b}).
                                                               \tag{C.62.11}
\]
The equality in (C.62.10) follows by cancellation
of the displayed original lifts, so its three
letters have exactly the labels claimed in
(C.62.2).

\emph{The full zero set introduces no missing fiber intervals.}
Write $\chi\gamma(r)=(\xi(r),\zeta(r))$.
On $Z$ retain the real function $\xi|_Z$.
On every gap $(a,b)$ of $Z$, extend it by the
linear function joining $\xi(a)$ to $\xi(b)$;
denote the resulting function on $[s,t]$ by $e$.
This extension is continuous, including at
accumulation points of gaps. Indeed, on a gap of
length less than $\delta$, uniform continuity of
$\xi$ bounds the difference between $e(r)$ and
$\xi(r)$ by its modulus of continuity at
$\delta$. There are only finitely many gaps of
length at least $\delta$, and the linear extension
is continuous at both endpoints of each of
those gaps. Letting $\delta$ decrease proves the
assertion at every point of $Z$.

The continuous image $e([s,t])$ contains the
entire real interval between $\xi(s)$ and
$\xi(t)$. At a parameter in $Z$ its value gives
the original point $\gamma(r)$ on the fiber.
On a gap its value lies between the two endpoint
fiber coordinates, and so lies in the interval
already covered by (C.62.11). Consequently the
full interval $I$ is covered by the zero-level
track and these excursion intervals.
A zero-level track point is itself a three-letter
value, using $a=v=u=r$. All excursion triangles
use $0\leq a\leq v\leq u\leq1$.
This proves the first inclusion in (C.62.2).

\emph{Projection and the original deck equations.}
For each source lift we have
$p_O\widetilde f_r^O=f_rp_O$, and the same
equation holds for its inverse. Projecting
(C.62.10) therefore gives exactly
$f_u f_v^{-1}f_a(x)$, proving the second
inclusion in (C.62.2).
No sheet selection enters this comparison.
More explicitly, each original deck map $\delta$
commutes with every $\widetilde f_r^O$, because
the deck-valued commutator path starts at the
identity. C.44 supplies the retained target
equation
\[
                   F_O\delta=F_O+\eta_O(\delta).
                                                               \tag{C.62.12}
\]
Replacing $\widetilde x$ by $\delta\widetilde x$
therefore carries the entire interval and every
word value to their $\delta$-images, with the
same translated real level. Projection removes
that same deck map. Equality only modulo a
circle period would not give the single
horizontal fiber or the two sides used in
(C.62.8), which explains the exact hypothesis
(C.62.1).
Applying the already proved argument to the
actual continuous path $(\widetilde f_r^O)^{-1}$
gives the stated inverse-path inclusion and its
projected word $f_u^{-1}f_vf_a^{-1}$.

\emph{Uniform ambient control, across arbitrary orbits.}
Let
\[
 \Delta_3=\{(a,v,u)\in[0,1]^3:a\leq v\leq u\},\qquad
 W((a,v,u),y)=f_u f_v^{-1}f_a(y).
                                                               \tag{C.62.13}
\]
This is one jointly continuous global map,
independent of all orbit charts and source covers.
The set $E=\{x\}\cup\{x_n:n\geq1\}$ is compact.
Hence $W(\Delta_3\times E)$ is compact and
contains every $J_n$ by (C.62.2).
Continuity on this compact product also gives
\[
 \sup_{\lambda\in\Delta_3}
      |W(\lambda,x_n)-W(\lambda,x)|\longrightarrow0.
                                                               \tag{C.62.14}
\]
The distance in (C.62.3) is bounded by this
quantity. Every value $W(\lambda,x)$ belongs
to $Gx$, so (C.62.3) confines all accumulation
to that same orbit and excludes ambient escape.
At a $G$-fixed point, every such value is $x$;
this proves (C.62.4). The identical compactness
argument applies to the inverse-path words.

This strengthens C.61 to compact containment
near the specified nonopen orbit and proves
the full interval estimate at fixed frontier
points. At a nonfixed orbit, the compact set
$\mathcal W_3(x)$ may contain more than one
point. The statement does not yet select the
specified median or exponential boundary
value, nor construct a jointly continuous
lift of the moving quotient deformation.
Those two remaining requirements are retained
in C.12.
\end{proof}

\begin{claim}{C.63}{Proved}{A controlled moving-base lift on principal orbits with fixed frontier}
Let $f_t$ be the original based ambient-null plane loop in
the centralizer of an arbitrary operator group $H$, and let
$G$ be the group generated by its values. Let $W$ be an
$H$-invariant union of open $G$-orbits. For every $O\subset W$
suppose that the full local centralizer
\[
 C_O=C_{\operatorname{Homeo}(O)}(G|_O)
\]
has dimension one, with its actual compact-open topology.
Suppose also that every point of $\partial W$ is fixed by $G$.
Then the actual global loop
\[
 \widehat f_t(x)=
 \begin{cases}f_t(x),&x\in W,\\x,&x\notin W\end{cases}
                                                               \tag{C.63.1}
\]
has a based null-homotopy in the $H$-centralizer, fixing the
complement of $W$ throughout. Both this homotopy and its
inverse are jointly continuous across every accumulation
of open orbits.

On each $O$, the induced contraction of the actual principal
quotient $O/C_O^0$ is the median contraction of its original
based real target lift. The original source cover, every
deck label, and the exact $H_O$ restriction kernel are
retained. No nontrivial global quotient to a nonopen orbit
is required, so the assertion includes the empty global
target family. In particular it contracts the original
loop when $\R^2\setminus W$ is fixed by $G$.
\end{claim}
\passage{Section 23, folios 140-141, scans 279-281; Part C(b), infinite-operator zero-level comparison and controlled moving-base extension at fixed frontiers}


\begin{proof}
\emph{The actual principal coordinates and source labels.}
C.30 and C.33 make $C_O$ a free proper Lie transformation
group. Put $A=C_O^0$. Their actual principal bundle and
section give
\[
 \Theta:A\times B\xrightarrow{\ \cong\ }O,
 \qquad\Theta(a,b)=a\,s(b),\qquad B=O/A,
                                                               \tag{C.63.2}
\]
where each of $A,B$ is a line or a circle. The pair of
circles is excluded because $O$ is noncompact.
Let $\eta_A:\R\to A$ and $\eta_B:\R\to B$ be the identity
on a line factor and the specified quotient by $\Z$ on
a circle factor. The map
\[
               p_*(u,v)=\Theta(\eta_A(u),\eta_B(v))
                                                               \tag{C.63.3}
\]
is a universal covering. With basepoints fixed, its unique
covering isomorphism with the original $p_O:X_O\to O$ is
a homeomorphism $\chi_O:X_O\to\R^2$ satisfying
$p_O=p_*\chi_O$. The inverse is the based lift in the
opposite direction; uniqueness of based lifts proves
both inverse identities. Hence
\[
 F_O=\operatorname{pr}_2\chi_O,
 \qquad
 \chi_O\delta\chi_O^{-1}(u,v)
       =(u+m_A(\delta),v+m_B(\delta))
                                                               \tag{C.63.4}
\]
for each original deck map $\delta$. A line factor has
zero label, and a circle factor has the original integer
label under this based covering comparison. At most one
factor is a circle.

Every element of $C_O$ normalizes its identity component
$A$ and preserves its orbit partition. C.33 proves that
$C_O/A$ acts freely and properly on $B$, so no induced
base map is decreasing: a decreasing line or circle
homeomorphism has a fixed point. Retain the complete
affine fiber changes constructed in C.33. On the original
universal cover, the full group of lifts of $C_O$ is
generated by all fiber translations and, when its real
base image is nontrivial, one map of the form
\[
 T_c(u,v)=(u+c,v),\qquad
 q(u,v)=(\lambda u+\kappa,v+d),
 \quad d>0,\quad\lambda\ne0.                            \tag{C.63.5}
\]
This also includes the source deck group. For a circle
base its lift labels supply the base period; for a
circle fiber its integer deck maps are among the
$T_c$. The constant $\kappa$ retains the choice of
component representative and lift. C.33's fiber gauge
removes the nonconstant translation term; a different
original lift only changes the constant term. The
negative values of $\lambda$ remain allowed. If the
base image is trivial, only $T_c$ occurs.

The based source lift of the original loop commutes
with every original lift of an operator in $C_O$.
Indeed its commutator with such a lift is deck-valued
and is initially the identity. Commutation with all
$T_c$ gives the exact expression
\[
 \widetilde f_t^O(u,v)=(u+\phi_t^O(v),\beta_t^O(v)),
                                                               \tag{C.63.6}
\]
where $\beta_t^O$ is increasing. Here and below this
formula is written through $\chi_O$. By C.32 the
ambient-null loop restricts to an ambient-null loop
on $O$, so its based source lift closes. Thus
\[
 \phi_0^O=\phi_1^O=0,
       \qquad\beta_0^O=\beta_1^O=\id.                 \tag{C.63.7}
\]
When (C.63.5) has the additional generator, its exact
commutation equations are
\[
 \phi_t^O(v+d)=\lambda\phi_t^O(v),\qquad
 \beta_t^O(v+d)=\beta_t^O(v)+d.                         \tag{C.63.8}
\]
The constant $\kappa$ cancels from this comparison,
rather than being discarded from the operator group.

The base maps $\beta_t^O$ have no common fixed point.
A common fixed point would project to a point of $B$
fixed by every original generator, contradicting the
transitivity of $G$ on $O$ and hence on $B$.
For any $a\in\R$ and any $\delta>0$, the compact
increment family
\[
 \{\beta_s^O(\beta_t^O)^{-1}:|s-t|\leq\delta\}
                                                               \tag{C.63.9}
\]
sends $a$ to points on both sides of $a$. To check
this for arbitrarily short increments, the path
$t\mapsto(\beta_t^O)^{-1}(a)$ is nonconstant. It
therefore has unequal values at arbitrarily close
times $s,t$; otherwise it is locally constant and
then constant on the connected parameter interval.
The corresponding increment moves $a$, and its
inverse moves $a$ to the opposite side. The
parameter set $|s-t|\leq\delta$ is convex, contains
the diagonal, and gives a path-connected increment
family. Its evaluation image consequently contains
a neighborhood of $a$. This uses inverse evaluation
at the specified point, not a displacement at an
unrelated moving point.

\emph{One compact family of global maps and original lifts.}
There are at most countably many open orbits, since
each contains a member of a fixed countable basis
disjoint from all the others. Consider the product
of the global plane homeomorphism group and the
homeomorphism groups of all the selected original
source covers. In this product take the closed
subgroup $\mathcal E$ of compatible tuples
\[
             (k,(\widetilde k^O)_{O\subset W}),
       \qquad p_O\widetilde k^O=k p_O,                 \tag{C.63.10}
\]
which commute globally with $H$ and in each cover
with every retained lift of $C_O$.
All these conditions are closed pointwise
equalities, including the inverse equalities.
Compatibility implies $k(O)=O$, using the
surjectivity of $p_O$ and the inverse tuple.
Alternatively each original word preserves the
closed complement of every $O$; limits of both
maps and inverses preserve that complement too.
Thus no limiting homeomorphism sends an open
orbit into its frontier.

Each factor is Polish in its compact-open
homeomorphism topology, with inverses controlled
on compact sets. The countable product and its
closed subgroup $\mathcal E$ are Polish. Choose
a complete compatible metric on $\mathcal E$
controlling these maps and inverses. For example,
after choosing a complete compatible metric,
add its value on the inverse pair; the resulting
metric is still complete and compatible.
The original path
\[
             L_t=(f_t,(\widetilde f_t^O)_O)
                                                               \tag{C.63.11}
\]
is a continuous based loop in this exact group.

We construct positive numbers $\delta_n$ so that
all the products
\[
 L_{s_n}L_{t_n}^{-1}\cdots
 L_{s_1}L_{t_1}^{-1}L_r,
       \qquad |s_j-t_j|\leq\delta_j,\quad r\in[0,1],
                                                               \tag{C.63.12}
\]
converge uniformly over their entire parameter
product as $n\to\infty$. Suppose the first $n-1$
choices have been made; their product image is
compact. Continuity of left multiplication,
uniform on this compact image, gives an identity
neighborhood whose multiplication changes every
such product by less than $2^{-n}$ in the chosen
metric. Uniform continuity of
$(s,t)\mapsto L_sL_t^{-1}$ near its compact
diagonal gives $\delta_n$ placing all the next
increments in that neighborhood. The successive
product maps are therefore uniformly Cauchy.
Completeness gives their uniform limit in
$\mathcal E$, including every global inverse
and every original lifted inverse.

Put
\[
 T=[0,1]\times\prod_{n\geq1}
       \{(s_n,t_n)\in[0,1]^2:|s_n-t_n|\leq\delta_n\}.
                                                               \tag{C.63.13}
\]
This is compact and path connected; coordinatewise
straight segments give paths with their exact
labels. The uniform limit just constructed is a
continuous map $T\to\mathcal E$. Denote its
labelled image by $\mathcal K$. Let $\mathcal K_n$
be the family of products through the first $n$
increments, with $\mathcal K_0=\{L_r\}$.
Setting all later increments to the identity
shows $\mathcal K_n\subset\mathcal K$ and
$\mathcal K_n\subset\mathcal K_{n+1}$.
Every global value of $\mathcal K$ commutes with
$H$ and fixes every point fixed by $G$.
The source labels in $\mathcal K$ are part of
the construction even when distinct parameters
have the same global homeomorphism.

The C.62 triangle works for any path in such a
labelled family. Its proof uses the relative
maps $h_r=k_rk_a^{-1}$, so it does not require
the selected path itself to start at the identity.
Connect any two chosen labels by a path in $T$.
Their equal-real-level fiber interval is contained
in the corresponding image of
\[
             \mathcal K\mathcal K^{-1}\mathcal K.
                                                               \tag{C.63.14}
\]
The projected family in (C.63.14) is compact in
the global plane homeomorphism group. The same
argument for the inverse labelled family uses
$\mathcal K^{-1}\mathcal K\mathcal K^{-1}$.
No continuous choice of connecting paths or
endpoint labels is needed for these uniform
containments.

\emph{Enlarging the allowed fiber intervals.}
Fix one source orbit and suppress $O$ from the
notation. Let $B_{\rho,t}$ be C.27's median of
the real path $\beta_t$, using prefix length
$\rho$. Thus $B_{0,t}=\id$, $B_{1,t}=\beta_t$,
and $B_{\rho,0}=B_{\rho,1}=\id$.
For every $v$, its value lies between $v$ and
$\beta_t(v)$, since every running minimum and
maximum contains the initial value $v$ between
them. Hence some original prefix time between
$0$ and $t$ attains $B_{\rho,t}(v)$.

Each labelled source map $k\in\mathcal K$ has
the form
\[
                    k(u,v)=(u+c_k(v),b_k(v)),
                                                               \tag{C.63.15}
\]
with its jointly continuous real functions in
the label and $v$. For $p=(\rho,t,v)$ define
the nonempty compact interval
\[
 A_n(p)=\operatorname{conv}
       \{c_k(v):k\in\mathcal K_n,
                          \ b_k(v)=B_{\rho,t}(v)\}.
                                                               \tag{C.63.16}
\]
Nonemptiness follows from the preceding prefix
observation, and compactness follows from the
closed equality in the compact labelled
parameter set. The identical definition with
$\mathcal K$ gives a compact interval
$A_\infty(p)$ containing every $A_n(p)$.

Here is the precise lower-limit property used
for selection:
\[
             A_n(p)\subseteq
                    \liminf_{p'\to p}A_{n+1}(p').       \tag{C.63.17}
\]
First take an endpoint value $c_k(v)$ before
convexification, with its fixed label
$k\in\mathcal K_n$ and $b_k(v)=a=B_{\rho,t}(v)$.
At nearby parameters $p'$, both $b_k(v')$ and
$B_{\rho',t'}(v')$ tend to $a$. By (C.63.9),
increments with arbitrarily small time
differences move $a$ to both sides of itself.
Choose two such increments within the
$(n+1)$st allowed family, as close to the
identity as desired. Their strict inequalities
continue to bracket the desired target at
all sufficiently nearby $p'$. A path in that
small increment family and the intermediate
value theorem therefore give an increment
$e_{p'}$ with
\[
          b_{e_{p'}k}(v')=B_{\rho',t'}(v').
                                                               \tag{C.63.18}
\]
As the chosen increment neighborhood decreases
to the identity, its effect on the fiber
coordinate also tends uniformly to zero near
the specified source point. Thus the value
$c_k(v)$ is approximated by allowed values
of $\mathcal K_{n+1}$ at every nearby parameter.
For a point between two allowed endpoint
values, perform this construction for both
endpoints and take the same convex combination.
This proves (C.63.17). It does not assert
lower semicontinuity of $A_n$ itself.

The operator equations are retained by doing
the selection in their actual line bundle.
If (C.63.5) has a base generator, all members
of $\mathcal K$ satisfy
\[
 c_k(v+d)=\lambda c_k(v),\qquad
 b_k(v+d)=b_k(v)+d.                                    \tag{C.63.19}
\]
The median has the same base equation.
Thus the intervals (C.63.16) define convex
fiber sets in the real line bundle
\[
 (v,z)\sim(v+d,\lambda z)
       \quad\hbox{over }\R/d\Z,                       \tag{C.63.20}
\]
pulled back over the parameter square.
Its continuous fiber norm is
$\|[v,z]\|=|\lambda|^{-v/d}|z|$.
For trivial base image use the ordinary real
line bundle over $[0,1]^2\times\R$.
Both bases are metrizable and paracompact.
Negative $\lambda$ gives the corresponding
nonorientable line bundle and preserves its
convex fiber sets just as in (C.63.19).

\emph{The relative continuous-selection construction.}
On the closed union $D$ of the four boundary
faces of the parameter square prescribe
\[
 c(0,t,v)=0,\quad c(1,t,v)=\phi_t(v),\quad
 c(\rho,0,v)=c(\rho,1,v)=0.                            \tag{C.63.21}
\]
This is a continuous section on $D$, including
the corners by (C.63.7), and belongs to every
allowed interval there. Replace $A_n$ by this
prescribed singleton on $D$ and leave it
unchanged elsewhere; denote the new fiber
sets by $F_n$. Property (C.63.17) survives.
At a point of $D$ the required inclusion for
approaching points off $D$ is the already
proved inclusion for its prescribed value,
and for approaching points in $D$ it is
continuity of (C.63.21). Away from $D$ a
neighborhood avoids $D$.

We give the elementary approximation proof
needed for these interlaced intervals. It
also applies in the line bundle (C.63.20).
Choose positive $\epsilon_n=2^{-n}$.
At each parameter point choose one value
of $F_0$ and extend it as a constant local
section in a bundle chart. By the lower-limit
inclusion, a sufficiently small neighborhood
makes this local section have distance less
than $\epsilon_1$ from $F_1$.
A locally finite partition of unity subordinate
to these neighborhoods combines the sections
to a continuous section $a_1$ with that same
distance bound. Indeed, at a given point the
corresponding nearby elements of $F_1$ have
their convex combination in $F_1$, so the
norm of the error is bounded by the same
weighted error bound.

Suppose $a_n$ has distance less than
$\epsilon_n$ from $F_n$. At each parameter
point choose a value in $F_n$ that close to
$a_n$, and extend it in a local bundle chart.
Shrink its neighborhood so that the extended
section remains within $2\epsilon_n$ of
$a_n$ and within $\epsilon_{n+1}$ of
$F_{n+1}$; the latter requirement is exactly
the lower-limit inclusion. Another locally
finite partition of unity gives $a_{n+1}$
with
\[
 \|a_{n+1}-a_n\|<2\epsilon_n,
       \qquad \operatorname{dist}(a_{n+1},F_{n+1})
                                      <\epsilon_{n+1}.
                                                               \tag{C.63.22}
\]
The summable first bounds give a limit section
$c$, uniformly in the fiber norm and hence
locally uniformly in every bundle chart.
It is continuous. At each fixed parameter,
every $F_n$ lies in the closed compact interval
$A_\infty$, and the second bound tends to zero;
therefore $c\in A_\infty$ at that parameter.
On $D$ the sets $F_n$ are the same prescribed
singleton, so $c$ equals (C.63.21) exactly.
This proves the required relative selection
without assuming that the original prefix
hulls themselves admit a continuous selection.

\emph{The intrinsic homeomorphisms and all operator equations.}
In the original real source chart set
\[
 K_{\rho,t}^O(u,v)
       =(u+c(\rho,t,v),B_{\rho,t}(v)).                 \tag{C.63.23}
\]
This is a jointly continuous family of
homeomorphisms, with inverse
\[
 (u,w)\longmapsto
 \bigl(u-c(\rho,t,B_{\rho,t}^{-1}(w)),
                         B_{\rho,t}^{-1}(w)\bigr).
                                                               \tag{C.63.24}
\]
The inverse family is jointly continuous
because the base median has a jointly
continuous inverse. The selected section
of (C.63.20) satisfies the first equation
of (C.63.19), and the base median satisfies
the second. Equations (C.63.5) therefore
commute exactly with (C.63.23); all fiber
translations commute as well. Consequently
every original source deck map commutes,
and (C.63.23) descends through $p_O$.
It commutes with the entire $C_O$ downstairs.
Equations (C.63.21) give the based
null-homotopy endpoint conditions.

Choose this construction on one representative
of each $H$-orbit of open components in $W$
and transport it by conjugation. The exact
restriction $H_O\to C_O$ has the kernel
specified in C.30.4. Commutation with its
image makes transport independent of the
chosen representative operator; its kernel
acts identically on $O$. Conjugation carries
$C_O^0$ to $C_{hO}^0$ and hence transports
the actual principal fibers. Original based
source lifts intertwine under a lift of $h$
by uniqueness at time zero. Their finite
products and the labelled limits above
intertwine too. Thus the original covering
labels are retained under this transport.

\emph{The compact-family bounds and fixed-frontier pasting.}
At every source point, membership
$c(\rho,t,v)\in A_\infty(\rho,t,v)$ places
the value (C.63.23) in the real fiber interval
between two labelled $\mathcal K$-values at
that same target level. The endpoints exist
because the defining level set in the compact
label space is compact. By (C.63.14), after
projection the value belongs to
\[
                   \mathcal K\mathcal K^{-1}\mathcal K(x).
                                                               \tag{C.63.25}
\]
For the inverse, put $y=(K_{\rho,t}^O)^{-1}x$.
The point $x$ is between two labelled values
$k_1y,k_2y$ in one real fiber. Since both
fiber maps are increasing, $y$ lies between
$k_1^{-1}x$ and $k_2^{-1}x$ in their common
real input fiber. The inverse version of
(C.63.14) therefore gives
\[
        (K_{\rho,t}^O)^{-1}(x)
             \in\mathcal K^{-1}\mathcal K\mathcal K^{-1}(x).
                                                               \tag{C.63.26}
\]
These are inclusions in one compact family
of actual global maps, independent of the
source orbit, the selected section, and all
homotopy parameters. They survive $H$
transport because every global member of
$\mathcal K$ commutes with $H$.

Define the global $K_{\rho,t}$ by these
transported maps on $W$ and the identity
outside $W$. It is a bijection: each open
orbit is preserved and the complement is
fixed. Let $x_n\to z\in\partial W$ and
allow the source orbit and both parameters
to vary arbitrarily. Every global map in
either compact family of (C.63.25)--(C.63.26)
fixes $z$, because $z$ is $G$-fixed.
Joint continuity of the compact global
evaluation family makes all its values at
$x_n$ tend uniformly to $z$.
The two inclusions thus prove continuity
of both the pasted map and its inverse at
every such limit. Interior continuity was
already proved, and points of the exterior
are fixed. Hence the global family consists
of homeomorphisms and is jointly continuous
with its inverse. It commutes with $H$ by
the transported construction and the
$H$-invariance of $W$. Its four endpoint
faces are precisely those required in
(C.63.1).

This establishes an actual moving-base lift
and global contraction at fixed frontiers,
including arbitrarily many open components
and their accumulation. At a nonfixed
frontier orbit, a compact global family need
not have one evaluation value; the argument
does not select the prescribed contraction
there. Open components with zero-dimensional
full local centralizer are also outside the
stated principal-fiber assertion. These
remaining cases of C.12 are retained.
\end{proof}

\begin{claim}{C.64}{Proved}{Compact local control and equivariant disk localization at a fixed frontier}
Let $f_t$ and $G$ be the original based ambient-null plane loop and
its generated group, and let $H$ commute with every $f_t$. Let $W$
be an $H$-invariant union of open $G$-orbits. There is a compact
path-connected labelled family $\mathcal K$ of actual global
$H$-centralizing homeomorphisms, containing the identity, such that
every member preserves each original open orbit and fixes the
$G$-fixed set, and
\[
 x\in\operatorname{int}_{O}
       \bigl(\mathcal K\mathcal K^{-1}(x)\bigr)
       \qquad(x\in O\subset W).                         \tag{C.64.1}
\]
All original source-cover labels are retained in this family.

Now suppose that $\partial W$ is $G$-fixed and that the full local
centralizer on every $O\subset W$ has dimension zero. One can
choose arbitrarily small closed embedded disks $D_i\subset O_i$,
with points $x_i\in D_i$, for which
\[
 D_i\subset\mathcal K\mathcal K^{-1}(x_i),               \tag{C.64.2}
\]
and every effective stabilizer translate of $D_i$ is disjoint
from it. Take any countable such family, with the $D_i$
representing pairwise distinct $H$-orbits of disks, whose
distinct $H$-translates are pairwise disjoint and locally
finite in $W$.
There is an actual continuous group embedding, with continuous
inverse on its image,
\[
 \mathfrak L:\prod_i\operatorname{Homeo}(D_i\mathbin{\mathrm{rel}}\partial D_i)
 \longrightarrow
 C_{\operatorname{Homeo}(\R^2)}(H),                       \tag{C.64.3}
\]
obtained by conjugating each disk map over its $H$-translates
and using the identity elsewhere. Its image is a contractible
subgroup. Both the maps and inverses are jointly continuous,
including at every accumulation on $\partial W$, uniformly
over this whole product of disk maps. In particular every
based loop in this subgroup has an explicitly prescribed
based contraction in the same global centralizer.
\end{claim}
\passage{Section 23, folios 140-141, scans 279-281; Part C(b), infinite-operator zero-level comparison and controlled localization at fixed frontiers}


\begin{proof}
\emph{The exact path group.}
Use every original universal cover $p_O:X_O\to O$ of the selected
open orbits, with the original based lifted loop. The open orbits
are countable, by a countable basis of the plane. The closed
compatible tuple construction (C.63.10) therefore gives a Polish
group of global homeomorphisms and their labelled original lifts.
It imposes $p_O\widetilde k^O=kp_O$, the inverse equations,
global commutation with $H$, and commutation with the retained
lifts of the full local centralizers. This construction uses
no assumption on their dimension. The original loop belongs
to this group: its commutators with fixed operator lifts are
deck-valued and initially trivial, and its based lifts close
by C.32. Let $L$ be the closure of the subgroup generated by
that labelled loop. Its global values preserve every selected
open orbit and fix every $G$-fixed point, since these are closed
conditions for maps and inverses.

Write $\mathscr P(L)$ for the group of continuous paths
$\sigma:[0,1]\to L$ with $\sigma(0)=1$, with pointwise
multiplication and the compact-open path topology. A complete
compatible metric on $L$ gives a complete uniform path metric.
Uniform limits are continuous and based. The space is separable:
uniform-continuity scales and finitely many values in a countable
dense set give a countable neighborhood coding; choose one actual
path from each nonempty code. Multiplication and inversion are
continuous uniformly on compact parameter sets. Thus
$\mathscr P(L)$ is a Polish group.

For $x\in O$, the continuous action by endpoint evaluation is
transitive. A finite original word carrying $x$ to a given point
of $O$ is the endpoint of a based word path: replace every
original time $t_j$ by $s t_j$, preserving the order and inverse
sign of every letter. These are paths in $L$, with their original
source lifts. The precise orbit-open-map theorem used in C.30,
\cite[Theorem~3.5]{C_GomezPinedoUzcategui2017}, consequently makes
\[
 e_x:\mathscr P(L)\longrightarrow O,\qquad
                 e_x(\sigma)=\sigma(1)x                 \tag{C.64.4}
\]
an open surjection. Its source is Polish and its target is the
nonmeagre Hausdorff surface $O$, as required there.

\emph{Compact lifting of a compact target.}
We give the needed elementary consequence of this open map.
Let $p:X\to Y$ be a continuous open surjection from a completely
metrizable space to a metric space. For an open $V\subset X$
and a compact $D\subset p(V)$, there is a compact
$S\subset V$ with $p(S)=D$.

Here is the full construction, including the equality of images.
For each $y\in D$ choose a lift in $V$, and an open metric ball
about it whose closure lies in $V$. Images of these balls cover
$D$. A finite subcover and a closed shrinking give compact
subsets $D_i$ covering $D$, each contained in the image of its
chosen ball $U_i$. A closed shrinking is obtained directly by
taking sufficiently small metric neighborhoods in the finite
cover of the compact space and then their closures in $D$.
Make each ball have diameter at most $1/2$, and each $D_i$
have diameter at most $1/2$, refining the finite cover when
needed.

Repeat inside every such pair $(U_i,D_i)$. Lift every point
of $D_i$ inside $U_i$; choose child balls with closures in
$U_i$ and diameter at most $2^{-2}$; use their open images
and a finite closed shrinking to cover $D_i$ by compact child
sets contained in those images and of diameter at most
$2^{-2}$. Continue with bound $2^{-n}$ at level $n$.
Discard empty child sets. This is a finitely branching tree.
Along each infinite branch the closed source balls are nested
and have diameters tending to zero. Completeness gives one
source limit, lying in the initial $V$. The nested nonempty
compact target sets similarly give one target limit in $D$.
At level $n$ that target limit has a lift in the corresponding
source ball, because the whole target set lies in its image.
These lifts converge to the source limit. Continuity of $p$
identifies its image with the target limit.

The infinite-branch space is compact, and the source-limit
map from it is continuous by the source diameter bounds.
Its image $S$ is therefore compact and maps into $D$.
Every $y\in D$ belongs to a nested sequence of child sets,
so some branch limit maps to $y$. This proves $p(S)=D$,
with no continuous endpoint section assumed.

\emph{One compact family, over all selected orbits.}
Fix a complete compatible metric on $L$, and use its uniform
path metric on $\mathscr P(L)$. Choose a compact exhaustion
$E_n$ of the open set $W$, whose interiors cover $W$.
Let $V_n$ be the open neighborhood of the constant identity
path consisting of paths uniformly within $2^{-n}$ of it.
For each $x\in E_n$, openness of (C.64.4) supplies a small
closed coordinate disk $D_{n,x}$ in $e_x(V_n)$ with $x$ in
its interior. These interiors cover $E_n$; retain finitely
many, with centers denoted $x_{n,j}$. The preceding compact
lifting construction gives compact sets
$S_{n,j}\subset V_n$ with
\[
 \{\sigma(1)x_{n,j}:\sigma\in S_{n,j}\}=D_{n,j}.          \tag{C.64.5}
\]
Take all labelled values $\sigma(s)$, for $0\leq s\leq1$
and $\sigma$ in these compact sets, and add the identity.
Their union $\mathcal K$ is compact in $L$: for each $n$
there are finitely many compact value families, and all
families with $n\to\infty$ converge uniformly to the identity.
The same observation proves sequential compactness directly
in the chosen metric. It is path connected, since every
one of its values is connected to the identity by the
corresponding initial segment of its path, entirely in
$\mathcal K$. Compactness and these paths hold for the
actual labelled tuples, including all source lifts and
inverse maps.

For $x$ in the interior of $D_{n,j}$ choose an endpoint
$k_x\in\mathcal K$ with $x=k_xx_{n,j}$. Every other point
of that disk is $k_yx_{n,j}$ for some $k_y\in\mathcal K$.
Consequently
\[
       D_{n,j}\subset\mathcal K\mathcal K^{-1}(x).
                                                               \tag{C.64.6}
\]
Their interiors cover $W$, proving (C.64.1). The projection
of $\mathcal K$ to global plane homeomorphisms is compact
and retains global commutation with $H$ and the fixed set.
We use the same notation for that projected family in the
pointwise estimates below.

\emph{Disks and the exact operator ambiguity.}
For an open orbit with zero-dimensional full local centralizer,
C.30 and C.33 make that centralizer a discrete Lie group
acting freely and properly. The effective image
\[
 \overline H_O=\{h|_O:h\in H_O\}
\]
is a subgroup of this discrete group, and hence is closed
and acts freely and properly. Its actual kernel is exactly
$\{h\in H_O:h|_O=\id\}$, as in (C.30.4).
For any $x\in O$, properness and freeness give a neighborhood
whose closure is disjoint from every nonidentity effective
translate. To see the shrinking explicitly, start with a
compact neighborhood. Properness permits only finitely
many translates meeting it. None of the nonidentity elements
in this finite set fixes $x$; disjoint neighborhoods of their
distinct images allow a further shrinking with disjoint
closures. Intersect also with the interior neighborhood in
(C.64.1). A sufficiently small closed coordinate disk $D$
about $x$ then satisfies (C.64.2) and the disjointness assertion.

Use any family specified in the claim. Distinct representatives
are taken modulo their whole $H$-orbits. An operator carrying
one representative disk to itself stabilizes its source orbit;
the preceding disjointness forces its effective restriction
to be the identity. Thus it fixes that disk pointwise.
For $\psi=(\psi_i)_i$ define
\[
 \mathfrak L(\psi)|_{hD_i}=h\psi_i h^{-1},\qquad
 \mathfrak L(\psi)(x)=x
 \quad\text{off the union of the disk interiors}.        \tag{C.64.7}
\]
Changing $h$ without changing its disk changes it by an
operator acting identically on the representative orbit.
Formula (C.64.7) is therefore unambiguous. All disk boundaries
are fixed; different disk interiors are disjoint. The formula
is a bijection with the identical formula for $\psi_i^{-1}$
as inverse, and it commutes with every element of $H$.

\emph{Uniform control at every fixed frontier.}
Put
\[
          Q=\mathcal K\mathcal K^{-1},\qquad
          P=Q^2=\mathcal K\mathcal K^{-1}
                         \mathcal K\mathcal K^{-1}.      \tag{C.64.8}
\]
These are compact global families, $Q=Q^{-1}$, and every
member fixes the $G$-fixed set and commutes with $H$.
For $x,y\in D_i$ choose $q_x,q_y\in Q$ with
$x=q_xx_i$ and $y=q_yx_i$. Then
$y=q_yq_x^{-1}x\in P(x)$. Commutation with $H$ gives
the same containment for every translated disk. Hence,
including the fixed exterior points by $1\in P$,
\[
 \mathfrak L(\psi)(x)\in P(x),\qquad
 \mathfrak L(\psi)^{-1}(x)\in P(x)                        \tag{C.64.9}
\]
for all disk maps $\psi$ and all $x\in\R^2$.

Let $x_m\to z\in\partial W$, allowing the disk index,
operator translate and every disk map to vary arbitrarily.
Compactness of $P$ and joint continuity of evaluation give
\[
             \sup_{p\in P}|p(x_m)-z|\longrightarrow0,    \tag{C.64.10}
\]
because $p(z)=z$ for every $p\in P$. Equations
(C.64.9)--(C.64.10) prove joint continuity of the pasted
maps and inverses at such points, uniformly over the
entire product. At an interior point of $W$, the stipulated
local finiteness leaves finitely many disk maps in a
neighborhood. Their ordinary disk evaluation and inverse
evaluation are jointly continuous, and they agree with
the identity on their boundaries. The same assertions hold
at exterior interior points, where all maps are the identity.
Thus (C.64.7) gives actual global homeomorphisms and joint
continuity of both families.

Since the plane is locally compact, joint continuity on a
compact set, followed by a finite product-neighborhood
subcover, gives continuity into its compact-open map space.
Apply this also to inverses. This proves continuity of
(C.64.3) for the product topology in its source and the
actual homeomorphism topology in its target. Composition
is componentwise on every disk, so it is a homomorphism.
Restriction to the representative disks is its continuous
inverse on the image, and in particular it is injective.

\emph{The actual contraction.}
Choose once an actual disk chart
$\theta_i:\overline{\mathbb D}\to D_i$. On the closed unit
disk, for $a\in\operatorname{Homeo}(\overline{\mathbb D}
\mathbin{\mathrm{rel}}S^1)$ set
\[
 \mathcal A_s(a)(z)=
 \begin{cases}
 s\,a(z/s),& |z|\leq s,\quad s>0,\\
 z,& |z|\geq s,
 \end{cases}
 \qquad \mathcal A_0(a)=\id.                            \tag{C.64.11}
\]
The two formulas agree on $|z|=s$ because $a$ fixes $S^1$.
Its inverse is $\mathcal A_s(a^{-1})$. At $s=0$ the
displacement of either map is at most $2s$, uniformly in
$a$; at positive $s$ ordinary joint evaluation proves
continuity, also across the joining circle. Therefore
(C.64.11) is a continuous contraction of the disk
homeomorphism group, with $\mathcal A_1(a)=a$ and
$\mathcal A_s(\id)=\id$.
Conjugate this formula by every $\theta_i$, apply it in
each factor with the same $s$, and then apply $\mathfrak L$.
Continuity in the product topology and the proved continuity
of $\mathfrak L$ give the stated contraction of its image.
For a based loop the endpoint identities are preserved by
$\mathcal A_s(\id)=\id$. Formula (C.64.9) controls the
whole contraction and its inverse at every frontier.

This proves actual equivariant disk localization and a
contractible global subgroup, with the original stabilizer
kernel retained. No factorization of an arbitrary original
loop into this subgroup is proved here. Such a controlled
factorization or a simultaneous graph deformation, and
extension at nonfixed frontier points, remain in C.12.
\end{proof}

\begin{claim}{C.65}{Proved}{Principal-kernel boundary comparison and controlled planar fillings}
Let $f_t,G$ be the original plane loop and its generated
group, and let $O$ be an open $G$-orbit with one-dimensional
full local centralizer. Retain the actual principal quotient
of C.33,
\[
 \pi_0:O\longrightarrow B_0=O/C_O^0,
 \qquad
 \beta:G\longrightarrow\operatorname{Homeo}(B_0).
                                                               \tag{C.65.1}
\]
Let $N$ consist of endpoints $h_1$ of paths
\[
\begin{gathered}
 h_s=\prod_{j=1}^{m}f_{a_j(s)}^{\epsilon_j},
 \qquad a_j:[0,1]\to[0,1]\ \text{continuous},\\
 \epsilon_j\in\{1,-1\},
 \qquad h_0=\id,\qquad\beta(h_s)=\id\ \text{for every }s.
\end{gathered}
                                                               \tag{C.65.2}
\]
The finite length and signs may depend on the path.
For every nontrivial nonopen $G$-orbit $B\subset\partial O$
whose actual intrinsic restriction image $\rho_B(G)$
is nonmetabelian,
\[
              \rho_B(N)=1,\qquad
 \rho_B:G\longrightarrow\operatorname{Homeo}(B).
                                                               \tag{C.65.3}
\]
In particular two actual finite-word paths $p_s,q_s$
starting at the identity and satisfying
$\beta(p_s)=\beta(q_s)$ for every $s$ have exactly the
same restriction to each such $B$ at every time.

For the full principal kernel $K_0=\ker\beta$, the
image $D_B=\rho_B(K_0)$ is central in $\rho_B(G)$
and is a closed discrete cyclic subgroup of the
intrinsic homeomorphism group of $B$; it is finite
when $B$ is a circle. Give
\[
 E_B=\{(\beta(g),\rho_B(g)):g\in G\}
       \subset\operatorname{Homeo}(B_0)
                    \times\operatorname{Homeo}(B)
\]
the actual induced product topology, and put
$\Delta_B=\{1\}\times D_B$. There are exact maps
\[
 E_B\longrightarrow E_B/\Delta_B
             \longrightarrow\beta(G).
\]
The first map is a covering with deck group
$\Delta_B$; the second is a continuous bijective
homomorphism onto the actual induced image
$\beta(G)$. No continuity of its inverse is asserted.

There is also the following controlled filling fact for
the same planar setting. For any open planar set $V$,
any open Euclidean disk $D$, and any connected component
$U$ of $V\cap D$, the actual inclusion induces an
injection
\[
                        \pi_1(U,x)\longrightarrow\pi_1(V,x)
                    \qquad(x\in U).                    \tag{C.65.4}
\]
Thus a based loop in $U$ null-homotopic in $V$ has a
based disk filling entirely in $U$. Applied to shrinking
null evaluation loops near a fixed frontier point,
this gives fillings inside the original orbit with
vanishing ambient diameter.
\end{claim}
\passage{Section 23, folios 140-141, scans 279-281; Part C(b), infinite-operator zero-level comparison, nonfixed-frontier path comparison and controlled planar fillings}


\begin{proof}
\emph{The actual kernel paths and their intrinsic topology.}
Every $g\in G$ commutes with $C_O^0$, by the definition
of the full local centralizer. It therefore induces
the actual base homeomorphism in (C.65.1). This
defines a homomorphism, with no change of quotient.
The set $N$ in (C.65.2) is a subgroup: products
and inverses of two displayed paths are again
finite-word paths in the exact kernel, and
identity letters can be inserted if needed.
Conjugation by a fixed original word preserves
that kernel and adds only finitely many constant
letters. Hence $N$ is normal in $G$.
Every intermediate $h_s$ also belongs to $N$,
by reparametrizing the path on the interval from
$0$ to $s$.

We verify the topology used on $B$. On a circular
orbit the intrinsic and embedded circle topologies
agree. On an intrinsic line orbit, C.28 puts every
fixed finite-word parameter track in a compact
intrinsic interval, embedded as a compact arc in
the plane. More generally, a compact intrinsic
interval $J\subset B$ lies in one of C.28's
finite-word compact sets: choose its two endpoints
there, and use the unique connecting arc in that
set. Applying a word of fixed length to every
point of $J$ remains in one larger finite-word
compact set, hence in one compact embedded
intrinsic interval. The ambient joint evaluation
map therefore gives joint intrinsic evaluation
on the path-parameter interval times $J$.
The same argument applies to the inverse word.
Consequently each path in (C.65.2) restricts to
a continuous path in the actual intrinsic
homeomorphism group of $B$, with inverse
convergence on compact intervals. All original
generators preserve the intrinsic orientation,
since their paths start at the identity; the
same is true of every word.

It follows in particular that $\rho_B(N)$ is
path connected in that intrinsic group: the
whole restricting path stays in $\rho_B(N)$
by the preceding reparametrization. This argument
uses the exact finite-word paths. It makes no
assertion about intrinsic convergence for an
arbitrary globally convergent sequence of
unbounded words.

\emph{The actual boundary image of the kernel is abelian.}
C.33 supplies a principal product chart
$O\cong A\times B_0$, where $A=C_O^0$ is the
additive line or circle. An element of
$\ker\beta$ commutes with all translations
by $A$, so its restriction has exactly the form
\[
                 (a,b)\longmapsto(a+\varphi(b),b).
                                                               \tag{C.65.5}
\]
Two such restrictions commute, including when
$A$ is a circle. Their actual global commutator
is therefore identity on $O$. Continuity makes
it identity on $\overline O$, and hence on $B$.
Thus $\rho_B(N)$ is abelian. It is normal in
$\rho_B(G)$ by the already proved normality of
$N$. No abstract source kernel has been
identified with a boundary kernel; the
commutator equality on $O$ and its extension
to $\overline O$ are the exact comparison.

\emph{A nontrivial path-connected abelian normal image forces
a metabelian action.}
Put $T=\rho_B(N)$ and $K=\rho_B(G)$, and
suppose $T\ne1$. The common fixed set of
$T$ is $K$-invariant because $T$ is normal.
Transitivity of $K$ on $B$ makes this fixed
set either all of $B$ or empty. It cannot
be all of $B$ when $T\ne1$, so it is empty.

Each $T$-orbit is path connected and nontrivial.
A nonconstant path in a line or circle contains
a nonempty open arc. Thus each orbit contains
such an arc. If an interior point of that
arc is the translate of the starting point
by an element of $T$, its inverse carries
the arc to a neighborhood of the starting
point within the same orbit. Hence every
$T$-orbit is open. The disjoint open orbits
partition the connected one-manifold $B$,
so $T$ is transitive.

An abelian transitive group acts freely:
an element fixing one point fixes every
translate of that point. Let $\overline T$
be the closure in the actual intrinsic
homeomorphism group. It is still abelian,
connected and transitive, and the same
argument makes it free. Its full centralizer
is exactly $\overline T$: a commuting map
agrees at one point with some member of the
transitive group, and that agreement
propagates to every point.

The evaluation argument of C.30, applied
in dimension one exactly as in C.48,
therefore identifies $\overline T$
homeomorphically with $B$ and makes it
locally compact. Its faithful product
action on $B\times\R^2$ is a continuous
action on a connected three-manifold.
The Hilbert--Smith input already verified
in C.33 makes it a Lie group.
It is connected of dimension one, hence
is the additive line or circle. Evaluation
is the actual conjugacy of its action
with all translations or all rotations,
respectively.

Every element of $K$ normalizes this
closure: it normalizes $T$, and conjugation
is continuous in the intrinsic
homeomorphism-group topology. On a line,
write the retained translations as $T_a$.
For a normalizing homeomorphism $g$,
\[
                 gT_ag^{-1}=T_{\theta_g(a)}
                                                               \tag{C.65.6}
\]
defines a continuous additive automorphism
of $\R$, so $\theta_g(a)=\lambda a$.
Its orientation gives $\lambda>0$.
Evaluation at zero yields
$g(a)=\lambda a+g(0)$.
Thus $K$ is a subgroup of the positive
affine group. On a circle the same
evaluation identity identifies the
induced continuous circle-group
automorphism with the degree of $g$.
That degree is $+1$, so the automorphism
is the identity and $g$ is a rotation.
In either case $K$ is metabelian.
This contradicts the hypothesis in
(C.65.3). We conclude that $T=1$,
which proves that equation.

For the claimed path comparison take
$h_s=p_sq_s^{-1}$. It is a finite-word
path starting at identity and satisfies
$\beta(h_s)=\id$ as an equality of whole
base homeomorphisms. Each $h_s$ belongs
to $N$; applying (C.65.3) gives
$\rho_B(p_s)=\rho_B(q_s)$ for every $s$.
Equality only at one varying base point
would not place this path in $\ker\beta$
and is not used here.

\emph{The full kernel and its exact discrete ambiguity.}
Let $k\in K_0$ and $g\in G$. Express $g$ by
one finite original word and replace each of its
time parameters $a_j$ by $s a_j$. This gives
an actual finite-word path $g_s$ from identity
to $g$. The path
\[
               [g_s,k]=g_s k g_s^{-1}k^{-1}
                                                               \tag{C.65.10}
\]
starts at identity, is again a single finite-word
path, and has principal-base image identity
at every time because $\beta(k)=1$.
Thus every value of (C.65.10) lies in $N$.
Equation (C.65.3) gives
$[\rho_B(g),\rho_B(k)]=1$.
Consequently $D_B$ is central in $\rho_B(G)$
and lies in its full intrinsic centralizer.

The action $\rho_B(G)$ is transitive and
nonmetabelian, hence nonabelian. C.38 identifies
that full centralizer, in its actual intrinsic
compact-open topology, as a closed discrete
group which is trivial or infinite cyclic
on a line, and finite cyclic on a circle.
Every subgroup of this closed discrete group
is closed in the intrinsic homeomorphism
group. This proves the asserted closedness,
discreteness, and cyclic description of $D_B$.
It does not require a topology on the
abstract kernel $K_0$.

Projection onto the first coordinate of
$E_B$ is a continuous surjective
homomorphism
\[
       p_B:E_B\longrightarrow\beta(G),
       \qquad \ker p_B=\Delta_B.                       \tag{C.65.11}
\]
Here both image topologies are the induced
topologies specified in the statement.
Since $D_B$ is closed and discrete in the
second factor, $\Delta_B$ is closed and
discrete in $E_B$. The preceding
centralization makes it central in $E_B$.

Write $q_B:E_B\to E_B/\Delta_B$ for the
topological group quotient. This map is
open: the inverse image of the image of
an open set $V$ is the open set
$V\Delta_B$. Choose an identity
neighborhood $U$ in $E_B$ with
$U^{-1}U\cap\Delta_B=\{1\}$, using
discreteness and continuity of multiplication
and inversion. The sets $U\delta$,
$\delta\in\Delta_B$, are pairwise disjoint,
and
\[
 q_B^{-1}(q_B(U))
       =\coprod_{\delta\in\Delta_B}U\delta.             \tag{C.65.12}
\]
The restriction of $q_B$ to each displayed
open set is a homeomorphism onto $q_B(U)$:
it is continuous and bijective, and its
openness follows from the just proved
openness of the quotient map. Translates
of these neighborhoods prove the covering
assertion.

The space $E_B$ is path connected. Indeed
the finite-word path $g_s$ used above is
continuous in both displayed image groups.
For the frontier factor this was proved
with the compact intrinsic tracks. For
the principal factor use the actual local
sections of $\pi_0$ in the product chart:
evaluation on such a section, followed
by $\pi_0$, is jointly continuous on
compact base sets, and the inverse word
gives the same fact for inverses.
This proves the required compact-open
path continuity in $\operatorname{Homeo}(B_0)$.
A deck transformation of a connected
covering is determined by its value at
one point: the agreement set of two
deck transformations is open and closed,
by evenly covered neighborhoods. Right
translation by $\Delta_B$ is transitive
on each fiber. It follows that the
entire deck group of $q_B$ is exactly
$\Delta_B$.

Finally (C.65.11) induces the continuous
bijective homomorphism
$\overline p_B:E_B/\Delta_B\to\beta(G)$.
Continuity is exactly the universal
property of the quotient topology.
The topological comparison stops at
this continuous bijection: no open-map
or inverse-continuity assertion for
$\overline p_B$ has been used or proved.

\emph{A filling inside the original domain and a prescribed
ambient disk.}
Let $\gamma$ be a based loop in $U$
which is null-homotopic in $V$.
Its compact image has an open
neighborhood in $U$. A sufficiently
fine polygonal approximation $\gamma'$
and the straight homotopies on the
small parameter intervals stay there,
fixing the basepoint. Choose finitely
many vertices generically if necessary,
so that its image is a finite planar
graph.

Take a small compact connected planar
regular neighborhood $M$ of this graph
inside $U$. For example, use small
vertex disks and thin edge rectangles
and then round their corners. Its
boundary has one outer Jordan curve
and finitely many inner Jordan curves.
Write $D_i$ for its bounded
complementary Jordan disks.
The whole closed outer Jordan disk
lies in $D$. Indeed its bounded
interior lies in the convex hull of
its boundary: a point outside that
convex hull is separated by a line
and can be joined to infinity without
meeting the boundary. The convex
hull of a compact subset of the
open Euclidean disk $D$ remains
inside $D$. In particular every
$D_i$ lies in $D$.

Fill into $M$ precisely the inner
disks $D_i$ which are entirely
contained in $V$, and call the
resulting compact planar surface $M'$.
Their boundaries already lie in $M$,
so the entire closed disk is added.
We have the actual inclusions
\[
              \gamma'\subset M\subset M'\subset V\cap D.
                                                               \tag{C.65.7}
\]
The set $M'$ is connected and meets
$U$, hence lies in that same component.
For each remaining inner disk choose
one point $z_i\in D_i\setminus V$.

The actual inclusion
\[
                  M'\hookrightarrow\R^2\setminus\{z_i\}
                                                               \tag{C.65.8}
\]
is a homotopy equivalence with an
explicit strong deformation retraction.
For each remaining closed inner Jordan
disk choose a Schoenflies disk coordinate
taking $z_i$ to zero and its boundary
to the unit circle. On the punctured
disk use
\[
        (r,\theta)\longmapsto((1-\tau)r+\tau,\theta),
                   \qquad 0<r\leq1.                    \tag{C.65.9}
\]
On the exterior of the outer Jordan
disk use its exterior Schoenflies
coordinate and the same formula for
$r\geq1$. All these homotopies fix
their boundary circles pointwise.
With the identity on $M'$, they
paste on the finite closed cover
of the punctured plane to a strong
deformation retraction onto $M'$.
Every basepoint in $M'$ stays fixed.

Since each $z_i$ lies outside $V$,
the given null-homotopy of $\gamma'$
in $V$ is a null-homotopy in the
punctured plane of (C.65.8).
Applying the displayed retraction
gives a based filling in $M'$.
Prepending the based polygonal
approximation homotopy gives a
filling of $\gamma$ in $U$.
This proves (C.65.4).

For original ambient-null evaluation loops
converging to a $G$-fixed frontier
point, compactness of the original
time interval makes their entire
tracks shrink to that point.
Choose open Euclidean disks
containing those tracks with
radii tending to zero. Each track
lies in one connected component
of its original orbit intersected
with that disk, and C.32 supplies
its null-homotopy in the original
orbit. The just proved injection
gives a filling in that same
component, with the stated
vanishing ambient diameter.

The first part compares specified
finite-word kernel paths and records
the central discrete ambiguity of the
full kernel through an actual covering.
It does not construct a new base lift
or identify the final continuous
bijection with a homeomorphism.
It also does not factor arbitrary
source germs through a frontier action.
The second part chooses each
filling separately; it does not
give a continuous choice of
fillings, simultaneous configuration
or edge isotopy extension, or
fragmentation of an arbitrary
original loop. These remaining
requirements of C.12 are retained.
\end{proof}

\begin{claim}{C.66}{Proved}{Zero-dimensional frontiers give a global centralizer filling}
Let $f_t$ be a based ambient-null loop in
$\operatorname{Homeo}^+(\R^2)$, let $G$ be the group
generated by its values, and let $H$ be any group
of plane homeomorphisms commuting with every $f_t$.
Let $Y$ be the closed union of the nonopen $G$-orbits.
When $Y$ is zero-dimensional, the original loop
has a based disk filling in
$C_{\operatorname{Homeo}(\R^2)}(H)$.
Every map of the filling fixes $Y$ pointwise.
This includes the empty frontier and imposes no
discreteness or local compactness assumption on $H$.

The topological comparison used here is the following
actual isomorphism. For every nonempty compact
zero-dimensional set $Z\subset S^2$, put $M=S^2\setminus Z$.
With compact-open group topologies, including inverse
convergence, restriction is an isomorphism of topological
groups
\[
 R_Z:\operatorname{Homeo}(S^2;Z)
             \xrightarrow{\ \cong\ }\operatorname{Homeo}(M),
 \qquad \operatorname{Homeo}(S^2;Z)
       =\{a:a(Z)=Z\}.                                  \tag{C.66.1}
\]
Its inverse is the unique end extension.
A connected parameter family containing the identity
extends to a family fixing every point of $Z$.

Consequently, for a connected admissible surface of C.1,
an ambient-null loop in the centralizer of an arbitrary
specified operator action is null-homotopic in that
same centralizer whenever the nonopen locus of its
based plane lift is zero-dimensional. All original
operator and deck labels are retained.
\end{claim}
\passage{Section 23, folios 140-141, scans 279-281; Part C(b), infinite-operator zero-level comparison for the zero-dimensional nonopen locus}


\begin{proof}
\emph{Finite disjoint Jordan covers of the deleted set.}
We first prove the planar compactification statement
with all the topology needed for parameter families.
End extension for zero-dimensional deleted sets is
also used in \cite[Proposition~1.5 and Appendix~A.1]
{C_BeguinCrovisierLeRoux2020}; the following construction
supplies the actual group topology comparison.

Let $Z\subset S^2$ be compact and zero-dimensional.
It is a proper subset. Choose a stereographic pole
outside $Z$, so $Z$ becomes a compact subset of the
coordinate plane. Given an open neighborhood $V$ of
$Z$ and $\varepsilon>0$, there is a finite family
of closed Jordan disks $D_1,\ldots,D_r$ such that
\[
 \begin{gathered}
 D_i\cap D_j=\varnothing\quad(i\ne j),\qquad
 Z\subset\bigcup_j\operatorname{int}D_j,\qquad
 D_j\subset V,\qquad D_j\cap Z\ne\varnothing,\\
              \operatorname{diam}_{S^2}D_j<\varepsilon .
 \end{gathered}                                         \tag{C.66.2}
\]
Here and below a diameter refers to one fixed compatible
metric on the compact sphere.

To construct the disks, partition $Z$ into finitely
many disjoint compact relatively clopen sets of
arbitrarily small coordinate diameter. Such a partition
exists by the clopen basis and compactness: take a finite
clopen cover of the required mesh and disjointize it by
successive differences. Distinct nonempty pieces have
positive distance. Cover them by a finite sufficiently
fine collection of coordinate grid cells and take a
small polygonal regular neighborhood of that union.
One obtains the neighborhood by separating corner
contacts in tiny vertex neighborhoods and rounding
the remaining boundary corners. Choose the cell mesh
and rounding size below the separation of the clopen
pieces. Each connected component of this regular
neighborhood then lies in a small neighborhood of one
partition piece. Delete any component not meeting $Z$.
Its other components are compact connected planar
surfaces, with finitely many disjoint Jordan boundary
curves, and they still contain $Z$ in their interiors.

Fill each component up to its outer Jordan curve.
The resulting disk lies in the convex hull of that
component: a point outside the convex hull is separated
by a line and belongs to the unbounded component of
the complement of its outer curve. Its coordinate
diameter is therefore bounded by that of the original
component. The outer curves of different components
are disjoint. Their closed Jordan disks are either
disjoint or nested; keep only the maximal disks.
Their closures are pairwise disjoint, and they cover
all the original components and hence $Z$.
The partition mesh and neighborhood sizes were arbitrary.
Choose them small enough that each retained disk,
which contains a point of $Z$, lies in $V$ and has
spherical diameter below $\varepsilon$. This proves
(C.66.2), without a local-connectivity assumption on $Z$.

The complement of $Z$ is path connected. For two
points outside $Z$, apply (C.66.2) with a neighborhood
avoiding both points. The complement of the interiors
of the finitely many disjoint disks is a connected
sphere with finitely many disks removed, and it
contains both points. A path in this complement
avoids $Z$. The same argument inside a Jordan disk
$D$ with $\partial D\cap Z=\varnothing$ shows that
$\operatorname{int}D\setminus Z$ is path connected:
cover the compact set $Z\cap D$ by finitely many
smaller disks in its interior, avoiding the two
chosen points, and join them in the remaining
finitely connected disk.

\emph{The exact end compactification.}
Apply (C.66.2) successively to obtain disk families
$D_{n,j}$ of mesh less than $2^{-n}$, with every
disk at stage $n+1$ contained in the interior of
one disk at stage $n$. Every disk meets $Z$.
Set
\[
 K_n=S^2\setminus\bigcup_j\operatorname{int}D_{n,j},
 \qquad
 U_{n,j}=\operatorname{int}D_{n,j}\setminus Z .
                                                               \tag{C.66.3}
\]
The sets $K_n$ are compact connected planar surfaces
contained in $M$, satisfy $K_n\subset\operatorname{int}K_{n+1}$,
and exhaust $M$. Indeed each disk meets $Z$ and its
diameter tends to zero, so every point outside $Z$
eventually lies outside all the disks. The preceding
connectivity argument shows that the $U_{n,j}$ are
exactly the components of $M\setminus K_n$.
Each is non-relatively-compact in $M$, since its
disk meets $Z$.

An end is a nested compatible choice of these
components. The corresponding nested closed disks
have diameters tending to zero. Their intersection
is one point of $Z$: each contains a point of $Z$,
and compactness and nesting give an intersection
with $Z$. Conversely every point of $Z$ determines
one such choice. Thus the map
\[
 \Xi:\widehat M\longrightarrow S^2,\qquad
 \Xi|_M=\id,\qquad
 \Xi((U_{n,j(n)})_n)=\bigcap_nD_{n,j(n)}
                                                               \tag{C.66.4}
\]
is a bijection from the end compactification.
A basic end neighborhood consisting of $U_{n,j}$
and the ends it contains maps to
$\operatorname{int}D_{n,j}$. These disks form a
neighborhood basis at their limiting point, by
their vanishing diameters. Equation (C.66.4)
and its inverse are therefore continuous,
and the topology on $M$ is unchanged.

Every $a\in\operatorname{Homeo}(M)$ is proper.
Given a target $K_m$, the compact set $a^{-1}(K_m)$
lies inside some $K_n$. The connected set
$a(U_{n,j})$ then lies in a unique component
of $M\setminus K_m$. Following a nested sequence
of source components defines its image end.
The same construction for $a^{-1}$ gives the
inverse map, and the basic neighborhoods just
described prove continuity. This is the unique
extension $\widehat a$ to $S^2$: $M$ is dense
because a zero-dimensional subset has empty
interior. The component construction, or
uniqueness on this dense set, gives
$\widehat{ab}=\widehat a\,\widehat b$ and
$\widehat{a^{-1}}=(\widehat a)^{-1}$.
The extension preserves $Z$ setwise.
Conversely a sphere homeomorphism preserving
$Z$ restricts to $M$ and is its unique extension.
This proves the algebraic bijection in (C.66.1).

\emph{The compact-open and inverse topology.}
Restriction in (C.66.1) is continuous on compact
subsets of $M$, for maps and inverses. We prove
continuity of extension at the identity; the
group homomorphism then gives continuity everywhere.
Both groups are metrizable, so it suffices to take
a sequence $a_i\to\id$ in the actual homeomorphism
group of $M$, including $a_i^{-1}\to\id$ on compact sets.

Fix $\varepsilon>0$ and choose $m$ so all disks
$D_{m,j}$ have diameter below $\varepsilon$.
The set
\[
             K_m\ \cup\ \bigcup_i a_i^{-1}(K_m)
                                                               \tag{C.66.5}
\]
is compact in $M$. To see this exactly, the
convergent sequence $\{\id,a_i^{-1}\}$ is a compact
parameter space, and its evaluation on the compact
set $K_m$ is a compact image. Choose $N\ge m$
with (C.66.5) contained in $K_N$.
Every $a_i(U_{N,j})$ consequently avoids $K_m$.
It is connected and hence lies in one component
of $M\setminus K_m$.

Choose a point $p_j$ in each of the finitely many
$U_{N,j}$. This set lies inside one parent
$U_{m,\ell(j)}$. Since $a_i(p_j)\to p_j$,
for all sufficiently large $i$ the component
containing $a_i(U_{N,j})$ is exactly that parent.
The end-extension construction then also sends
every end in $D_{N,j}$ into $D_{m,\ell(j)}$.
Thus each point of $\operatorname{int}D_{N,j}$,
including its deleted points, and its image
under $\widehat a_i$ lie in the same disk
$D_{m,\ell(j)}$ of diameter below $\varepsilon$.
On the remaining compact set $K_N$, the
original maps $a_i$ converge uniformly to identity.
It follows that $\widehat a_i\to\id$ uniformly
on the whole sphere. Apply the identical argument
to the sequence $a_i^{-1}$ to obtain uniform
inverse convergence. This proves (C.66.1) as
an isomorphism of the stated topological groups.

In particular every continuous parameter family
extends jointly with its inverse. When its parameter
space is connected and contains a parameter at which
the map is identity, evaluation at any $z\in Z$
is a continuous map into the zero-dimensional
space $Z$ with connected image. That image is
the singleton $\{z\}$. The entire family fixes
every deleted point individually. This is the
pointwise conclusion used below; the full group
in (C.66.1) preserves $Z$ only setwise.

\emph{Application to the original plane loop.}
The group $G$ is path connected: any finite word
in loop values is joined to identity by scaling
each of its time parameters to zero. Its orbits
are therefore path connected. A nonopen orbit
contained in the zero-dimensional set $Y$
must be a point, so $G$ fixes $Y$ pointwise.
Let
\[
                  Z=Y\cup\{\infty\}\subset S^2 .
                                                               \tag{C.66.6}
\]
This is compact and zero-dimensional. For
completeness, a closed zero-dimensional subset
of the plane is locally compact and has a basis
of compact relatively open-and-closed sets:
inside a compact neighborhood choose a clopen
basis element with closure still in that
neighborhood. Finite unions of such sets give
an increasing compact clopen exhaustion of $Y$.
Their complements in $Z$ form a clopen basis
at the point at infinity. The old clopen basis
at each finite point is unchanged. The empty
$Y$ gives the singleton $Z=\{\infty\}$.

The complement $O=\R^2\setminus Y=S^2\setminus Z$
is connected by the disk-cover argument.
It is a union of disjoint open $G$-orbits.
Connectedness makes it one such orbit.
Every $h\in H$ transports $G$-orbits to $G$-orbits
with their actual labels, since it commutes
with $G$. It preserves $Y$ and the unique open
orbit $O$.

Let $P$ be the actual restriction image of
$H\to\operatorname{Homeo}(O)$, retaining its
kernel $\{h:h|_O=1\}$, and let $Q$ be the
closure of $P$ in that homeomorphism group.
C.30 makes the full local centralizer of
$G|_O$ a closed locally compact group acting
freely and properly on $O$. The closed subgroup
$Q$ of this full centralizer has these same
properties. The restricted original loop
commutes with $Q$.

C.32 proves that this restriction is null-homotopic
in $\operatorname{Homeo}(O)$, using the specified
ambient-null input. C.33 therefore gives an actual
based disk filling
\[
 A:D^2\longrightarrow
          C_{\operatorname{Homeo}(O)}(Q),
       \qquad A|_{\partial D^2}=f_t|_O .
                                                               \tag{C.66.7}
\]
Apply the inverse of (C.66.1) to this whole family.
The disk is connected and contains the based
identity parameter, so every extended map fixes
all of $Z$, including the specified point
$\infty$. Restricting these sphere maps to the
original plane gives a jointly continuous family
of plane homeomorphisms and inverses, fixing $Y$.
On $O$ it commutes with each restriction of
each original $h\in H$. The two compositions
are continuous on the plane and agree on the
dense set $O$, so they agree everywhere.
The boundary family is the original $f_t$ by
the same density argument. This proves the
claimed filling with every operator label retained.

\emph{The actual lift and descent for an admissible surface.}
Let $p:\R^2\to U$ be the universal covering of a
connected admissible surface, and start with a
based centralizer loop which is null-homotopic
in its ambient homeomorphism group.
Lift that ambient disk family after composition
with $p$, choosing the identity at its based
parameter. The covering lifting criterion and
the corresponding inverse family give actual
plane homeomorphisms. Its boundary is a closed
ambient-null based plane loop $\widetilde f_t$.

Take $H$ to consist of every original lift of
every specified surface operator, including
all deck transformations. For each such lift
$h$, the maps $\widetilde f_t h$ and
$h\widetilde f_t$ lift the same surface family
and agree at the based parameter; uniqueness
of the covering lift makes them equal for
every $t$. Under the zero-dimensional locus
assumption the plane result supplies a disk
family $K_w$ commuting with every one of
these original lifts. Commutation with the
deck transformations gives the unique descent
\[
                    \overline K_w\,p=pK_w .
                                                               \tag{C.66.8}
\]
The open covering quotient makes the descended
map and inverse families jointly continuous.
Commutation with each lifted operator descends
to commutation with its specified surface
operator. The family has the original boundary.

For $U=X\setminus S$ with nonempty finite $S$,
the same end-extension topology argument, now
using disjoint coordinate disks about the
finitely many punctures in $X$, extends this
connected based family to $X$ fixing $S$
pointwise. In the compact case there is no
extension step. Thus (C.66.8) is a disk filling
in the original group $A^0\cap A^u$, with
its original operator action. No replacement
of that group by an abstract quotient is made.

The proof handles zero-dimensional $Y$ through
its exact end compactification. It does not
extend an arbitrary intrinsic filling across
a nontrivial line or circle frontier, or across
an arbitrary positive-dimensional fixed set.
Those remaining extensions in C.12 are unchanged.
\end{proof}

\begin{claim}{C.67}{Proved}{The exact relative disk group for a locally connected frontier}
Let $O\subset\R^2$ be a nonempty bounded simply
connected open set with locally connected frontier.
Choose a Riemann map $\phi:\mathbb D\to O$ and
retain its continuous Carath\'eodory extension
\[
                  q:\overline{\mathbb D}\longrightarrow\overline O.
                                                               \tag{C.67.1}
\]
Put
\[
 \begin{split}
 \mathcal D_\partial
     &=\{a\in\operatorname{Homeo}(\overline{\mathbb D}):
                                      a|_{S^1}=\id\},\\
 \mathcal K_O
     &=\{b\in\operatorname{Homeo}(\R^2):
                                b|_{\R^2\setminus O}=\id\}.
 \end{split}                                             \tag{C.67.2}
\]
The actual map
\[
 T_q:\mathcal D_\partial\longrightarrow\mathcal K_O,
 \qquad
 T_q(a)|_O=\phi a\phi^{-1},\qquad
 T_q(a)|_{\R^2\setminus O}=\id
                                                               \tag{C.67.3}
\]
is an isomorphism of topological groups with their
compact-open and inverse topologies.
In particular $\mathcal K_O$ is contractible.
The frontier need not be a Jordan curve.

For the actual stabilizer $H_O$ of $O$ in any
specified group of ambient operators, let
\[
 \eta_q:H_O\longrightarrow
                   \operatorname{Homeo}(\overline{\mathbb D})
                                                               \tag{C.67.4}
\]
be its original prime-end action in the chosen
Riemann coordinate. Its kernel is exactly
$\{h\in H_O:h|_O=\id\}$. The isomorphism (C.67.3)
restricts to the exact comparison
\[
 C_{\mathcal D_\partial}(\eta_q(H_O))
   \xrightarrow{\ \cong\ }
       \mathcal K_O\cap
               C_{\operatorname{Homeo}(\R^2)}(H_O).
                                                               \tag{C.67.5}
\]
Both directions have the actual induced group
topologies. The contraction asserted above is
for the full relative group $\mathcal K_O$;
no equivariance of the Alexander contraction
under a nontrivial operator is asserted.
\end{claim}
\passage{Section 23, folios 140-141, scans 279-281; Part C(b), infinite-operator zero-level comparison and relative extension across a locally connected fixed frontier}


\begin{proof}
\emph{The retained quotient and prime-end maps.}
The Carath\'eodory theorem for locally connected
boundaries gives (C.67.1); it is a continuous
surjection, equals $\phi$ on the open disk,
and maps the boundary circle onto $\partial O$.
The same conformal coordinate identifies the
prime-end compactification of $O$ with the
closed disk. We use precisely the crosscut-chain
compactification and the extension theorem
stated in \cite[Sections~1 and~3.1]
{C_BinderRojasYampolsky2014}.
The map $q$ may have several preimages for a
frontier point. No inverse of $q$ on the
frontier is used.

An ambient homeomorphism $h$ preserving $O$
induces a homeomorphism of this prime-end
compactification. Here is the relevant
functorial construction. It sends a crosscut,
with its two distinct frontier endpoints,
to the crosscut with their actual images.
It sends the corresponding complementary
crosscut regions to their actual images.
On the compact set $\overline O$, $h$ is
uniformly continuous. Thus disjoint crosscut
closures remain disjoint and diameters
tending to zero still tend to zero.
Nested fundamental chains are sent to nested
fundamental chains. The inclusion relations
defining equivalence of chains are preserved,
as are those relations under $h^{-1}$.
The image-chain rule is therefore a bijection
on prime ends. Basic crosscut-region
neighborhoods and the inverse rule show that
its extension is a homeomorphism.
This construction is independent of the
auxiliary interior basepoint used to select
crosscut sides: a chain tends to the frontier
and eventually excludes either chosen
interior point.

Conjugating this extension by the fixed
prime-end Riemann coordinate gives
$\widehat h=\eta_q(h)$, with the exact equations
\[
 \widehat h|_{\mathbb D}=\phi^{-1}h\phi,\qquad
                    q\widehat h=hq .
                                                               \tag{C.67.6}
\]
The second equation holds first on the
dense open disk and then everywhere by
continuity. Composition of image chains,
or uniqueness on the dense disk, proves
that $\eta_q$ is a homomorphism.
It is identity exactly when $h|_O$ is
identity. These are the asserted original
restriction labels and kernel.

\emph{A frontier-fixing map fixes the entire prime-end circle.}
Let $b\in\mathcal K_O$ and let $\widehat b$
be its just constructed prime-end extension.
It preserves orientation: $b$ is identity
on the nonempty open exterior of the
bounded set $\overline O$, so its global
orientation is positive.
Write $r=\widehat b|_{S^1}$ and
$q_\partial=q|_{S^1}$. Since $b$ fixes
$\partial O$ pointwise, (C.67.6) gives
\[
                         q_\partial r=q_\partial.
                                                               \tag{C.67.7}
\]
We prove that this exact equation forces
$r=\id$.

First, $q_\partial$ is constant on no
nonempty open arc. Otherwise, for a point
$p$ equal to that constant value, the
holomorphic map $\phi-p$ is continuous
up to an open circle arc and is zero
there. In a local conformal coordinate
straightening that arc to a real segment,
Schwarz reflection extends it holomorphically
across the segment. The identity theorem
then makes $\phi-p$ identically zero,
contrary to injectivity of $\phi$.

Fix $x_*\in O$. The circle-valued map
\[
 v(\zeta)=\frac{q_\partial(\zeta)-x_*}
                    {|q_\partial(\zeta)-x_*|}
                  \quad(\zeta\in S^1)                 \tag{C.67.8}
\]
has degree one. Indeed $\phi(sS^1)$,
for $s<1$ sufficiently close to one,
is a positively oriented Jordan curve
enclosing $x_*$. Uniform convergence of
$\phi(s\zeta)$ to $q_\partial(\zeta)$
and the positive distance from $x_*$
to $\partial O$ preserve that winding
number in the limit.

Use $\zeta=e^{2\pi i t}$ and choose
a real lift $R$ of $v$, so
$R(t+1)=R(t)+1$. For an increasing
lift $P$ of $r$, equation (C.67.7)
implies that $R(P(t))-R(t)$ is one
constant integer. Replacing $P$ by
the unique integer translate which
subtracts that integer gives
\[
                    R(P(t))=R(t)\quad(t\in\R).
                                                               \tag{C.67.9}
\]
The function $R(t)-t$ is periodic and
bounded. Thus the full orbit
$\{P^n(t):n\in\Z\}$ is bounded for
each $t$, since its $R$-value is
constant. An increasing real
homeomorphism with such an orbit
has a fixed point: a nonconstant
forward monotone orbit has a finite
limit, and continuity fixes that
limit. Hence $P$, and therefore
$r$, has a fixed point.
Moreover, if $r(e^{2\pi i t})=e^{2\pi i t}$,
then $P(t)=t+k$ for an integer $k$.
Periodicity would give $P^n(t)=t+nk$;
the established boundedness forces $k=0$.
Thus the fixed set of $r$ is exactly
the projection of $\operatorname{Fix}(P)$.

If $r$ were nonidentity, choose
a component of the complement
of its fixed set. Using the
periodic fixed set of $P$, lift
that component to an interval
between consecutive boundary
points of $\operatorname{Fix}(P)$.
The interval is bounded, its
endpoints are fixed, and $P$
has no fixed point inside.
Iterations of every interior
point therefore tend to one
endpoint. Equation (C.67.7)
and continuity make $q_\partial$
constant on that open arc.
This contradicts the no-constant-arc
property. We conclude that
$r=\id$ as claimed.

\emph{The algebraic comparison and its topology.}
For $a\in\mathcal D_\partial$, the
continuous map $qa$ on the closed
disk is constant on every fiber
of $q$. Interior fibers are single
points; on a boundary fiber $a$
fixes every point. The compact
quotient $q$ therefore gives a
continuous map on $\overline O$,
equal to the identity on its
frontier. The identical construction
with $a^{-1}$ gives its inverse.
Pasting this map with identity on
the closed complement of $O$
defines the actual homeomorphism
(C.67.3). The formulas on $O$
give injectivity and preservation
of products.

Conversely, for $b\in\mathcal K_O$,
the prime-end extension $\widehat b$
fixes $S^1$ by the preceding argument.
Thus $\widehat b\in\mathcal D_\partial$,
and its interior formula gives
$T_q(\widehat b)=b$ everywhere.
This proves bijectivity.

For continuity let $\omega_q$ be
the modulus of uniform continuity
of $q$ on the compact closed disk.
The actual forward maps satisfy
\[
 \sup_{x\in\R^2}|T_q(a)(x)-T_q(a')(x)|
       \leq
       \omega_q\!\left(
        \sup_{\zeta\in\overline{\mathbb D}}
                                      |a(\zeta)-a'(\zeta)|
                    \right).                           \tag{C.67.10}
\]
Outside $O$ both sides of the
comparison of maps are identity.
The same estimate applies to
their inverses. Hence $T_q$
is continuous in the actual
group topologies.

The source $\mathcal D_\partial$
is a closed subgroup of the
Polish homeomorphism group of
the compact disk. The target
$\mathcal K_O$ is a closed
subgroup of the Polish plane
homeomorphism group: it is
defined by the pointwise
equations on $\R^2\setminus O$.
The source acts continuously
and transitively on this Polish
target by $(a,b)\mapsto T_q(a)b$.
The Polish open orbit-map
theorem verified in C.26 and
used in C.30 makes its orbit
map at identity open.
That orbit map is exactly
$T_q$. The continuous bijection
is therefore open and has a
continuous inverse. This proves
the full topological group
isomorphism, without assuming
injectivity of the boundary
quotient $q$.

\emph{Contraction of the full relative group.}
The Alexander family on
$\mathcal D_\partial$ is
\[
 A_s(a)(\zeta)=
 \begin{cases}
       s\,a(\zeta/s),&|\zeta|\leq s,\quad s>0,\\
       \zeta,&|\zeta|\geq s ,
 \end{cases}
 \qquad A_0(a)=\id.                                    \tag{C.67.11}
\]
The boundary-fixing condition
gives agreement on $|\zeta|=s$.
Every map has inverse
$A_s(a^{-1})$, and fixes $S^1$.
For $s>0$ joint uniform
continuity follows from uniform
continuity on the compact disk;
at zero both the map and its
inverse have displacement at
most $2s$. Also
$\|A_s(a)-A_s(a')\|_\infty
\leq s\|a-a'\|_\infty$.
These observations give joint
continuity in $s$ and $a$,
including inverses.
We have $A_1(a)=a$,
$A_0(a)=\id$, and
$A_s(\id)=\id$.
Transporting (C.67.11) by the
proved isomorphism contracts
the full group $\mathcal K_O$.

\emph{Every original operator relation.}
For $h\in H_O$, (C.67.6)
implies on the open disk,
and therefore by continuity
on the relative plane maps,
\[
 T_q(\widehat h\,a\,\widehat h^{-1})
       =h\,T_q(a)\,h^{-1}.
                                                               \tag{C.67.12}
\]
Outside $O$ both sides are
identity, since $h$ preserves
$O$ and its complement.
Conjugation by $\widehat h$
preserves $\mathcal D_\partial$
even when $h$ reverses orientation.
Equation (C.67.12) proves both
directions of the exact
centralizer comparison (C.67.5).
Restriction of a homeomorphism
of topological groups to these
subgroups retains their
induced topologies.
An original operator in the
kernel of $H_O\to\operatorname{Homeo}(O)$
has $\widehat h=\id$ and
conjugates every map of
$\mathcal K_O$ trivially;
that actual kernel has been
retained.

The same identity also keeps
the labels of operators which
transport $O$ to another such
domain $O'$. With a chosen
Riemann map $\phi'$ and extension
$q'$ there, the prime-end
transport is the unique disk
homeomorphism extending
$\phi'^{-1}h\phi$, and satisfies
$q'\widehat h=hq$.
Thus
\[
 T_{q'}(\widehat h\,a\,\widehat h^{-1})
        =h\,T_q(a)\,h^{-1}.                             \tag{C.67.13}
\]
Composing two original operators
composes these exact transports,
by their uniqueness on the
open disk.

Equations (C.67.5), (C.67.12),
and (C.67.13) provide the
operator-preserving disk models.
The Alexander family was proved
to contract the full relative
group. Its preservation of the
operator centralizers in
(C.67.5) has not been established.
Nor does this one-domain
topology comparison supply
uniform control for contractions
over infinitely many transported
domains. These are retained
extension requirements in C.12.
\end{proof}

\begin{claim}{C.68}{Proved}{Relative first-kind disk centralizers and spherical frontier quotients}
Let $O\subset S^2$ be a simply connected domain whose
frontier is a nondegenerate locally connected continuum.
Retain a Riemann coordinate and its actual continuous
Carath\'eodory quotient
\[
 \phi:\mathbb D\longrightarrow O,\qquad
 q:\overline{\mathbb D}\longrightarrow\overline O .
                                                               \tag{C.68.1}
\]
Writing $\mathcal D_\partial$ as in C.67, put
\[
 \mathcal K_O^+
  =\{b\in\operatorname{Homeo}^+(S^2):
                         b|_{S^2\setminus O}=\id\}.
                                                               \tag{C.68.2}
\]
The map $T_q$, equal to $\phi a\phi^{-1}$ on $O$
and to the identity off $O$, is an isomorphism of
topological groups
\[
                  T_q:\mathcal D_\partial
                              \xrightarrow{\ \cong\ }\mathcal K_O^+.
                                                               \tag{C.68.3}
\]
The full group $\mathcal K_O^+$ is contractible.
No injectivity of $q|_{S^1}$ is assumed.

There is also the following operator comparison.
Let $A\leq\operatorname{Homeo}(\overline{\mathbb D})$,
and retain its actual restriction closure
\[
 Q=\overline{\{a|_{\mathbb D}:a\in A\}}
             \ \leq\operatorname{Homeo}(\mathbb D).
                                                               \tag{C.68.4}
\]
Suppose that $Q$ is locally compact and acts freely
and properly on $\mathbb D$. Retain a specified
$k\in\operatorname{Homeo}(\overline{\mathbb D})$ and
a torsion-free Fuchsian group $\Lambda$ of the first
kind, with the actual subgroup inclusion
\[
       (k|_{\mathbb D})^{-1}\Lambda(k|_{\mathbb D})\leq Q .
                                                               \tag{C.68.5}
\]
Here first kind means that the limit set of the
specified disk action is all of $S^1$.
Restriction is an isomorphism of topological groups
\[
 \operatorname{res}:
       C_{\mathcal D_\partial}(A)
       \xrightarrow{\ \cong\ }
       C_{\operatorname{Homeo}(\mathbb D)}(Q).
                                                               \tag{C.68.6}
\]
Every based loop in the left-hand group has a based
filling in that same group.
Neither finite generation nor cocompactness of
$\Lambda$ is required.

For original ambient operators stabilizing $O$,
take $A$ to be their actual prime-end image under
(C.68.1). Equations (C.68.3) and (C.68.6) then
give a based filling of every loop in the corresponding
relative ambient centralizer whenever the displayed
data (C.68.4)--(C.68.5) hold. This retains each original
operator and its restriction kernel.
For the plane application, retain the specified
point $\infty\notin O$; the entire filling fixes
that point and restricts to a filling through
plane homeomorphisms.
\end{claim}
\passage{Section 23, folios 140-141, scans 279-281; Part C(b), infinite-operator zero-level comparison and relative frontier extension}


\begin{proof}
\emph{The actual spherical boundary quotient.}
The locally connected Carath\'eodory theorem and
the prime-end compactification in
\cite[Sections~1 and~3.1]{C_BinderRojasYampolsky2014}
give (C.68.1). They apply in a spherical coordinate:
the conformal map may be meromorphic in a chosen
plane chart, but is a homeomorphism from the open
disk onto $O$. The boundary map
$q_\partial=q|_{S^1}$ is onto $\partial O$.

This boundary map is constant on no nonempty open
arc. Indeed, a constant value $p$ on such an arc
admits a target complex coordinate finite near $p$.
Continuity of $q$ gives a neighborhood of a smaller
closed subarc on which that coordinate of $\phi$
is holomorphic and continuous up to the arc.
Subtract its constant boundary value and straighten
the source circle locally. Schwarz reflection across
the resulting real segment gives a holomorphic
extension with a segment of zeros. The identity
theorem makes $\phi$ constant on an open subset,
and then on the connected disk, a contradiction.

We now prove the needed precise property of its
possibly nontrivial fibers. There cannot be four
cyclically ordered points $a,b,c,d$ in $S^1$ with
\[
            q(a)=q(c)=p,\qquad q(b)=q(d)=r,\qquad p\ne r .
                                                               \tag{C.68.7}
\]
Choose an embedded closed crosscut $\alpha$ joining
$a$ to $c$ and otherwise contained in $\mathbb D$.
The map $q$ embeds the interior of this crosscut,
and its interior image lies in $O$. Its only
coincident endpoint images are the point $p$.
Consequently
\[
           J=q(\alpha)\text{ is a Jordan curve in }S^2,
           \qquad J\cap\partial O=\{p\}.
                                                               \tag{C.68.8}
\]
The two components of $\mathbb D\setminus\alpha$
have connected images disjoint from $J$.
Near an interior point of $\alpha$, the map $\phi$
is a local homeomorphism and sends its two local
sides to the two local sides of $J$.
Thus the two entire component images lie in the
different complementary Jordan domains.
The cyclic order in (C.68.7) places $b$ and $d$
on the respective boundary arcs of the two source
components. Approaching them from those components
shows that $r$ lies in the closures of both
Jordan domains, and hence in $J$.
Since $r\in\partial O$, (C.68.8) gives $r=p$,
contrary to (C.68.7). This proves that distinct
fibers of $q_\partial$ do not have alternating
pairs of points.

Let $u\in\operatorname{Homeo}^+(S^1)$ satisfy
\[
                            q_\partial u=q_\partial .
                                                               \tag{C.68.9}
\]
It follows that $u=\id$. To see this, choose a
moved point $\zeta$ and an increasing real lift
in which $\zeta<u\zeta<\zeta+1$.
For every $\eta>\zeta$ sufficiently close to
$\zeta$, continuity and strict increase give
\[
                       \zeta<\eta<u\zeta<u\eta<\zeta+1 .
                                                               \tag{C.68.10}
\]
The absence of a constant boundary arc permits
such an $\eta$ with $q_\partial\eta\ne q_\partial\zeta$.
Equation (C.68.9) and (C.68.10) then give the
forbidden alternating pairs in two distinct
fibers. This proves the assertion.

\emph{The relative spherical group, its topology, and its operators.}
Every original sphere homeomorphism $h$ preserving
$O$ acts on its prime ends. This is precisely the
null-crosscut-chain functor proved in C.67:
uniform continuity on the compact sphere preserves
null chains, their inclusion relations, and their
equivalence, and the inverse gives the inverse
prime-end map. In the retained coordinate write
this map as $\widehat h$. It satisfies
\[
 \widehat h|_{\mathbb D}=\phi^{-1}h\phi,\qquad
                             q\widehat h=hq .
                                                               \tag{C.68.11}
\]
For $h\in\mathcal K_O^+$, the boundary restriction
of $\widehat h$ preserves orientation and satisfies
(C.68.9); it is therefore the identity.

For $a\in\mathcal D_\partial$, the continuous
map $qa$ is constant on every fiber of $q$:
interior fibers are singletons, and all points
in each boundary fiber are fixed by $a$.
The compact quotient $q$ therefore descends
$qa$ to a continuous self-map of $\overline O$.
The same construction for $a^{-1}$ is its inverse.
It is the identity on $\partial O$, so pasting
with identity on $S^2\setminus O$ gives $T_q(a)$.
It preserves orientation because its restriction
to $O$ does. Equations on $O$ prove injectivity
and the homomorphism identity. Equation (C.68.11)
and boundary rigidity give its inverse, proving
the algebraic bijection (C.68.3).

Uniform continuity of $q$ gives forward continuity
of $T_q$: uniform changes of $a$ and $a^{-1}$
on the closed disk give uniform changes of
$T_q(a)$ and its inverse on the sphere.
Both groups in (C.68.3) are closed Polish subgroups
of their homeomorphism groups. The left action
of $\mathcal D_\partial$ on $\mathcal K_O^+$
through $T_q$ is continuous and transitive.
The verified Polish open orbit-map theorem used
in C.26 and C.30 makes this bijective homomorphism
open. Its inverse is continuous, as claimed.
Conjugating the explicit Alexander contraction
of $\mathcal D_\partial$ in C.67 by this
isomorphism proves contractibility of the full
relative group.

For any original operator $h$ preserving $O$,
the exact interior conjugation and the exterior
identity give
\[
          T_q(\widehat h\,a\,\widehat h^{-1})
                        =h\,T_q(a)\,h^{-1}.
                                                               \tag{C.68.12}
\]
Orientation reversal by $h$ is allowed.
The kernel of $h\mapsto\widehat h$ is exactly
the kernel of restriction to $O$.
Thus (C.68.3) identifies the actual centralizers,
with the actual induced topologies and operator
labels. No equivariance of the Alexander formula
under these operators is used.

\emph{First-kind extension with the specified disk conjugacy.}
We use the original first-kind extension theorem
in \cite[Section~3.B, Proposition~3.5 and
Corollary~3.5.1, pp.~31--32]{C_Tukia1972}.
It states that a homeomorphism of the open disk
inducing a geometric isomorphism between
first-kind groups extends homeomorphically to
the closed disk. Its boundary map is determined
by the images of the attracting fixed points
of infinite-order elements. The proof constructs
the boundary map from nested intervals with
hyperbolic-axis endpoints and proves continuity
of the actual interior homeomorphism at their
limiting boundary points. The hypotheses in
that section impose no finite generation or
cocompactness requirement.

In particular, an interior homeomorphism $g$
commuting with $\Lambda$ induces its identity
group isomorphism. The boundary map supplied
by that theorem fixes every one of the stated
attracting points; they are dense for a
first-kind group. Its boundary extension is
therefore the identity:
\[
      g\in C_{\operatorname{Homeo}(\mathbb D)}(\Lambda)
        \quad\Longrightarrow\quad
      \overline g\in\mathcal D_\partial .
                                                               \tag{C.68.13}
\]
This is an extension of the given homeomorphism,
not a replacement by another representative.
Uniqueness follows from density of the open disk.

Now let $g\in C_{\operatorname{Homeo}(\mathbb D)}(Q)$.
By the actual inclusion (C.68.5), the map
$k g k^{-1}$ on the open disk commutes with
$\Lambda$. Equation (C.68.13) extends it to
the closed disk, fixing the boundary.
Conjugating this extension by the specified
closed-disk homeomorphism $k^{-1}$ gives the
unique extension $\overline g$ of $g$, still
fixing the boundary. For each $a\in A$,
the two continuous maps $\overline g\,a$
and $a\,\overline g$ agree on the dense open
disk, and hence everywhere.
Thus $\overline g\in C_{\mathcal D_\partial}(A)$.

Conversely, a member of
$C_{\mathcal D_\partial}(A)$ restricts to a
homeomorphism commuting with all restrictions
of $A$, and then with their actual closure
$Q$, by continuity of group operations.
This proves surjectivity and injectivity
of (C.68.6). Restriction is continuous.
The source and target are closed Polish
subgroups: commuting with each specified
operator is a closed condition, and fixing
the boundary is closed.
Applying the same open orbit-map theorem
to the transitive action through restriction
proves continuity of its inverse.
Consequently every compact parameter family,
including its inverse family, extends
continuously through this exact inverse map.
No separate assumption about a boundary
restriction fibration is needed for this
particular isomorphism.

\emph{Based fillings and the original frontier.}
Let $t\mapsto f_t$ be a based loop in
$C_{\mathcal D_\partial}(A)$.
As a loop in the full group
$\mathcal D_\partial$, it is null-homotopic
by the Alexander contraction in C.67.
Restriction makes it an ambient-null loop
in $\operatorname{Homeo}(\mathbb D)$.
The proper-action injection C.33 applies
to the exact group $Q$ in (C.68.4).
It supplies a based filling
\[
       B:D^2\longrightarrow
                    C_{\operatorname{Homeo}(\mathbb D)}(Q)
\]
of that restricted loop. Composing with
the continuous inverse in (C.68.6) gives
a based filling of the original $f_t$
through boundary-fixed closed-disk
homeomorphisms commuting with every
original $a\in A$.
This proves the asserted loop statement.

For the spherical domain, use
$A=\{\widehat h:h\in H_O\}$ from (C.68.11).
Apply (C.68.6) to the loop obtained through
$T_q^{-1}$, and return the resulting filling
through $T_q$. Equation (C.68.12) proves
commutation with every original operator
in $H_O$ throughout the filling; the
restriction kernel acts identically on
$O$ throughout. The complement of $O$
is fixed pointwise.
This passage uses the quotient $q$ and
never uses an inverse boundary value of
that quotient. It therefore includes slit
frontiers and other locally connected
frontiers with multiple prime-end accesses.

In the plane application $\infty\notin O$,
so $\infty$ belongs to the fixed complement.
Restriction from the sphere group fixing
$\infty$ gives the original compact-open
and inverse topology on plane homeomorphisms;
the converse continuity is the one-end
compact-core comparison proved in C.66.
Hence this is a filling in the required
plane group with its specified infinity.
\end{proof}

\emph{Scope of the extension.}
The closed-disk homeomorphism $k$ and the
first-kind subgroup inclusion (C.68.5) are
part of the proved comparison.
An arbitrary interior uniformization of
$O/Q$ does not supply this $k$ on the
original prime-end circle.
The argument makes no such substitution.
It also supplies no uniform estimate for
pasting contractions over infinitely many
different transported domains.
The remaining general frontier and
moving-quotient extension in C.12 stays open.

\begin{claim}{C.69}{Proved}{Relative disk centralizers for Fuchsian actions of either kind}
Let $\Gamma$ be a torsion-free discrete group of
orientation-preserving M\"obius transformations of
$\mathbb D$, with its actual action on
$\overline{\mathbb D}$. No finite-generation
assumption is imposed. Every based loop in
\[
                 C_{\mathcal D_\partial}(\Gamma)
                                                               \tag{C.69.1}
\]
has a based filling in that same group.
The same assertion holds for the actual action
$k^{-1}\Gamma k$, with any specified
$k\in\operatorname{Homeo}(\overline{\mathbb D})$.

More precisely, let $\mathcal L_\Gamma$ be the
limit set, and retain the actual spaces and
restriction map
\[
 \Omega_\Gamma=S^1\setminus\mathcal L_\Gamma,\qquad
 D_\Gamma=\mathbb D\cup\Omega_\Gamma,\qquad
 \operatorname{res}_\Gamma:
 C_{\mathcal D_\partial}(\Gamma)\longrightarrow
 C_{\operatorname{Homeo}(D_\Gamma,\Omega_\Gamma)}(\Gamma).
                                                               \tag{C.69.2}
\]
The notation on the right means pointwise fixing
$\Omega_\Gamma$. The map
$\operatorname{res}_\Gamma$ is an isomorphism
of topological groups with their compact-open
and inverse topologies. This isomorphism
extends the full filling, including its inverse
family, across the original limit set.

The proof also gives the following exact relative
surface comparison. For every second-countable
surface $N$ with boundary, inclusion induces
an injection
\[
 \pi_1\bigl(\operatorname{Homeo}(N,\partial N),\id\bigr)
       \longrightarrow
 \pi_1\bigl(\operatorname{Homeo}(N),\id\bigr).
                                                               \tag{C.69.3}
\]
The proof supplies an explicit correction of an
ambient parameter-disk filling on a proper collar.
It permits infinitely many boundary components
and preserves the original parameter-boundary
loop exactly.

Through the actual spherical quotient and
operator maps of C.68, (C.69.1) gives a relative
equivariant filling on a locally connected
simply connected domain whenever the original
prime-end action is the specified closed-disk
conjugate $k^{-1}\Gamma k$. Its original
restriction kernel is retained. In the plane
application the specified infinity lies outside
the domain and is fixed throughout.
\end{claim}
\passage{Section 23, folios 140-141, scans 279-281; Part C(b), infinite-operator zero-level comparison and exact relative frontier extension}


\begin{proof}
\emph{The original partial compactification and its proper action.}
The group $\Gamma=1$ has
$D_\Gamma=\overline{\mathbb D}$, and (C.69.1)
is the Alexander assertion in C.67.
We first record facts that also cover every
nontrivial $\Gamma$.

For any sequence of distinct elements
$\gamma_n\in\Gamma$, a subsequence has points
$a,b\in\mathcal L_\Gamma$ with
\[
 \gamma_n(z)\longrightarrow a
 \quad\hbox{uniformly on compact subsets of }
                        \overline{\mathbb D}\setminus\{b\}.
                                                               \tag{C.69.4}
\]
Here is the explicit disk calculation.
Write
\[
 \gamma_n(z)=e^{i\theta_n}
                  \frac{z-b_n}{1-\overline b_n z},
                    \qquad b_n=\gamma_n^{-1}(0).
                                                               \tag{C.69.5}
\]
Discreteness prevents a sequence of distinct
$\gamma_n$ from remaining in a compact subset
of the M\"obius group. Taking a subsequence
gives $b_n\to b\in S^1$ and
$e^{i\theta_n}\to e^{i\theta}$.
For $z\ne b$, the numerator at the limit
equals $-b(1-\overline b z)$.
Thus (C.69.4) follows with
$a=-e^{i\theta}b$; the denominator is bounded
away from zero on each stated compact set.
Both $a$ and $b$ are in the limit set,
by applying the definition to
$\gamma_n(0)$ and $\gamma_n^{-1}(0)$.

The limit set is $\Gamma$-invariant: applying
an element to a convergent orbit sequence
gives another orbit sequence, and its
inverse gives equality. Thus the displayed
action preserves $D_\Gamma$.
The space $D_\Gamma$ is a second-countable
surface with boundary $\Omega_\Gamma$.
At a point of $\Omega_\Gamma$ use a small
closed-disk boundary chart disjoint from
the closed set $\mathcal L_\Gamma$.
It is contractible: $(z,t)\mapsto(1-t)z$
is an actual contraction in $D_\Gamma$.

The $\Gamma$-action on this space is free
and proper. A nonidentity element is
hyperbolic or parabolic, because a discrete
elliptic element has finite order and
$\Gamma$ is torsion free. All its fixed
points lie in $\mathcal L_\Gamma$, so
it fixes no point of $D_\Gamma$.
For compact $K\subset D_\Gamma$, a sequence
of distinct translates meeting $K$
would contradict (C.69.4): the compact
$K$ avoids the repelling point $b$,
its images converge uniformly to
$a\in\mathcal L_\Gamma$, and $K$ has
positive distance from that closed
limit set. Consequently the exact quotient
\[
                p:D_\Gamma\longrightarrow
                           N_\Gamma=D_\Gamma/\Gamma
                                                               \tag{C.69.6}
\]
is a regular universal covering of
surfaces with boundary, with precisely
the deck group $\Gamma$.
Evenly covered neighborhoods, Hausdorffness,
and second countability follow exactly
as in C.32, using boundary half disks
when required. Its boundary is
$\Omega_\Gamma/\Gamma$, and its interior
is $\mathbb D/\Gamma$.

\emph{A dense set of controlled boundary accesses.}
Let $E_\Gamma$ be the union of
$\Omega_\Gamma$ and the attracting fixed
points of the hyperbolic elements of
$\Gamma$. Then
\[
                          \overline{E_\Gamma}=S^1.
                                                               \tag{C.69.7}
\]
For completeness this also follows from
(C.69.4). If $\Gamma$ contains a
hyperbolic element $\gamma$, use both
of its endpoints, the attracting
points of $\gamma$ and $\gamma^{-1}$.
For $\zeta\in\mathcal L_\Gamma$ choose
distinct $\gamma_n$ with
$\gamma_n(0)\to\zeta$, and pass to
(C.69.4). At least one of the two
chosen endpoints differs from $b$.
Its images under $\gamma_n$ tend to
$\zeta$ and are the attracting points
of the actual conjugates of $\gamma$
or $\gamma^{-1}$.
Thus the hyperbolic fixed points are
dense in the limit set.

If a nontrivial $\Gamma$ has no
hyperbolic element, all its elements
are parabolic. Two with distinct fixed
points would generate a hyperbolic
element: conjugate one to $z\mapsto z+1$
in the upper half plane, write the
other as a parabolic matrix with
nonzero lower-left entry $c$, and
observe that the trace of its product
with the $n$th power of the first
is its original trace plus $nc$.
Its absolute value exceeds two for
some integer $n$.
Thus all fixed points coincide.
After conjugating that point to
infinity, every element is a real
translation. Discreteness makes
this group cyclic, and its limit
set is a single point.
Its $\Omega_\Gamma$ is dense in
$S^1$. The trivial group has
$\Omega_\Gamma=S^1$.
These cases prove (C.69.7).

\emph{The precise cap-side extension.}
Let
\[
             h\in
 C_{\operatorname{Homeo}(D_\Gamma,\Omega_\Gamma)}(\Gamma).
                                                               \tag{C.69.8}
\]
At a point $p\in\Omega_\Gamma$, it and
its inverse already extend continuously
with value $p$.
At an attracting fixed point $p$ of a
hyperbolic $\gamma$, choose a terminal
ray of the hyperbolic axis ending at $p$.
Partition this ray into the actual
translates $\gamma^n A$ of one compact
fundamental segment $A$.
Its image under $h$ is partitioned
into $\gamma^n h(A)$, by commutation.
The compact set $h(A)$ has iterates
converging uniformly to $p$ under
the hyperbolic M\"obius map.
Thus this image ray ends at the
same $p$. The same argument applies
to $h^{-1}$. These are accesses
with the exact original axis labels.

Fix $\zeta\in S^1$ and $\epsilon>0$.
Choose two points $p,q\in E_\Gamma$
on opposite sides of $\zeta$ in
a sufficiently short circular arc
$I$ containing $\zeta$. Choose an
open crosscut $\alpha\subset\mathbb D$
whose closed arc has endpoints
$p$ and $q$, as follows. Near an
endpoint in $\Omega_\Gamma$, use
a short radial tail. Near a
hyperbolic fixed endpoint use
a terminal axis tail as above.
Truncate those two disjoint tails
and join their interior endpoints
by an embedded arc in a small
convex disk cap around $\zeta$.
The tails may be chosen in
disjoint small endpoint neighborhoods;
the joining arc avoids them except
at its two endpoints.
This gives a crosscut whose closure
and the closed boundary arc
$\overline I$ lie in a convex
neighborhood of diameter less than
$\epsilon$.

Let $V$ be the component of
$\mathbb D\setminus\alpha$ adjacent
to $I$. The Jordan curve
$\overline\alpha\cup\overline I$
bounds this cap. Its filled region
lies in the convex hull of that
curve, so
\[
                           \operatorname{diam} V<\epsilon .
                                                               \tag{C.69.9}
\]
The image $\beta=h^{-1}(\alpha)$
has a continuous embedded closure
with the same two endpoints $p,q$,
by the access calculation.
It is therefore another actual
crosscut. Let $U$ be its complementary
component adjacent to the same
boundary arc $I$.

We verify the side, rather than
inferring it from the endpoint set.
Choose $r\in E_\Gamma\cap I$.
There are points $z_n\in\mathbb D$
with both $z_n\to r$ and
$h^{-1}(z_n)\to r$.
For $r\in\Omega_\Gamma$ this follows
from continuity at $r$; for a
hyperbolic fixed point take
$z_n=\gamma^n z_0$ for its attracting
map and use
$h^{-1}(z_n)=\gamma^n h^{-1}(z_0)$.
Eventually $z_n\in V$ and
$h^{-1}(z_n)\in U$.
Since $h$ maps complementary
components of $\beta$ bijectively
to those of $\alpha$, it follows
that
\[
                               h(U)=V .
                                                               \tag{C.69.10}
\]
The set $U\cup I$ is a neighborhood
of $\zeta$ in the closed disk.
Equations (C.69.9)--(C.69.10)
prove that $h(z)\to\zeta$ as
$z\to\zeta$ from the open disk.
The same proof applies to the
inverse. Thus $h$ extends uniquely
to a boundary-fixed homeomorphism
of $\overline{\mathbb D}$.
Commutation with every original
$\gamma$ survives by density.

This proves algebraic surjectivity
of (C.69.2); injectivity follows
from density of the open disk.
Restriction is continuous. Both
groups are closed Polish groups
in their indicated compact-open
homeomorphism groups, since the
commutator equations and pointwise
boundary constraints are closed.
The Polish open orbit-map argument
of C.26/C.30 applied to this
continuous bijective homomorphism
proves continuity of its inverse.
In particular this gives joint
extension for every compact
parameter family and its inverse
family, uniformly on the closed
disk. It is the full topology
comparison, not merely pointwise
extension of the individual maps.

\emph{The relative parameter-disk correction on an arbitrary surface.}
We prove (C.69.3) independently.
Let $f_t$ be a based loop fixing
$\partial N$ pointwise, and let
$B_z$, $z\in D^2$, be a based
ambient filling with that boundary
loop.
Fix a proper closed collar
\[
                     c:\partial N\times[0,1]\hookrightarrow N .
                                                               \tag{C.69.11}
\]
We recall explicitly that such
a collar can be made proper.
A locally finite surface
triangulation with boundary
subcomplex gives a collar by
the adjacent boundary-edge
rectangles and the finite
interval fans at boundary vertices.
Their normal parameters agree
on the successive common edges
and give a product collar.
Start with any such collar
$c_0$ and choose a compatible
proper metric on $N$.
A positive continuous function
$\delta(b)<1/2$ can be chosen
so that
$d(c_0(b,s),b)<1$ for
$0\leq s\leq\delta(b)$:
local continuity gives the
admissible neighborhoods and
a locally finite partition
of unity gives a positive
continuous minorant.
The collar
$c(b,t)=c_0(b,\delta(b)t)$
is proper. Indeed the $b$
coordinates over a compact
target set lie in its closed
unit neighborhood, which is
compact, and the inverse
image is closed in that
compact set times $[0,1]$.

Each boundary component is
a line or a circle. The
connected based family $B_z$
preserves each component and
its orientation. Put
$\psi_z=B_z|_{\partial N}$.
On a line component choose
a real coordinate and denote
its actual increasing map
by $\Psi_z$.
On a circle component, lift
the joint map
$D^2\times\mathbb R\to S^1$
to a real map $\Psi_z$,
normalized by
$\Psi_{z_*}(u)=u$ at the
specified based parameter
$z_*\in\partial D^2$.
The parameter disk is simply
connected, so this is an
actual continuous lift.
It satisfies
\[
 \Psi_z(u+1)=\Psi_z(u)+1,\qquad
             \Psi_z(u)=u\quad(z\in\partial D^2).
                                                               \tag{C.69.12}
\]
For the second equality,
the original boundary loop
fixes the boundary circle;
the integer discrepancy
is continuous on the connected
parameter boundary and is
zero at $z_*$.
Every $\Psi_z$ is increasing.
The same normalized lifting
of $\psi_z^{-1}$ gives its
inverse, by uniqueness.

Choose a continuous function
$\chi:[0,1]\to[0,1]$ with
$\chi(0)=1$ and $\chi(t)=0$
for $t\geq2/3$.
On each real boundary
coordinate put
\[
             E_{z,t}(u)
                  =(1-\chi(t))u+\chi(t)\Psi_z(u).
                                                               \tag{C.69.13}
\]
This is an increasing proper
homeomorphism. On circle
components it commutes with
integer translation and
therefore descends.
The families and their
inverses are continuous:
strict increase is preserved
by the convex combination,
and compact parameter sets
give common properness bounds
from the compact inverse family
of $\Psi_z$.
Extend to $N$ by
\[
 E_z(c(b,t))=c(E_{z,t}(b),t),\qquad
                  E_z=\id\text{ outside the collar}.
                                                               \tag{C.69.14}
\]
The notation in (C.69.14)
uses the just specified
line or circle coordinates.
The formula is identity
near the outer collar edge,
so it pastes to a homeomorphism.

This is also a continuous
family for infinitely many
boundary components.
Properness of the collar
reduces each compact target
set to a compact subset
of $\partial N\times[0,1]$.
Only finitely many components
meet that subset, since
boundary components are
open in the boundary
one-manifold. All formulas
on those components have
the proved compact-parameter
and inverse continuity.
Thus (C.69.14) is continuous
in the original compact-open
and inverse topology.

We have
$E_z|_{\partial N}=\psi_z$
and $E_z=\id$ on the entire
parameter boundary. Consequently
\[
                            B'_z=E_z^{-1}B_z
                                                               \tag{C.69.15}
\]
is a based filling fixing
$\partial N$ pointwise,
with exactly the prescribed
loop $f_t$ on $\partial D^2$.
This proves (C.69.3) without
assuming a boundary-restriction
fibration or a global choice
of circle lifts over its
whole image.

\emph{Filling on the exact quotient and lifting every label.}
Return to a nontrivial
$\Gamma$ and a based loop
$f_t$ in (C.69.1).
Restriction through (C.69.2)
and descent through (C.69.6)
give a based loop $b_t$
of homeomorphisms of
$N_\Gamma$ fixing its
boundary pointwise, with
the exact equation
\[
                            b_t p=p f_t .
                                                               \tag{C.69.16}
\]
Joint continuity and inverse
continuity descend through
the open covering quotient,
as in C.32.

Choose $x\in\mathbb D$ and
$v=p(x)$ in the interior
of $N_\Gamma$.
By \cite[Theorem~1.1 and
Lemma~2.1]{A_Yagasaki2000},
the identity component of
$\operatorname{Homeo}(N_\Gamma,v)$
is contractible.
Here all exceptions are
excluded explicitly:
the only pointed exceptions
for this assertion are the
plane, the compact disk,
the sphere, and the projective
plane. The actual universal
cover (C.69.6) has nontrivial
infinite fundamental group
$\pi_1(N_\Gamma)=\Gamma$
downstairs, so none occurs.
An interior coordinate disk
gives a continuous local
section of evaluation by
homeomorphisms supported in
that disk. Hence evaluation
is a locally trivial bundle
onto the interior, and its
homotopy exact sequence gives
\[
 \pi_1(\operatorname{Homeo}(N_\Gamma),\id)
             \hookrightarrow\pi_1(\operatorname{Int}N_\Gamma,v).
                                                               \tag{C.69.17}
\]
The evaluation loop $b_t(v)$
is null-homotopic: it is
the projection of the closed
loop $f_t(x)$ in the
contractible open disk.
Thus (C.69.17) supplies an
ambient based filling of
$b_t$, and (C.69.15)
corrects it to a filling
$B'_z$ fixing the whole
boundary of $N_\Gamma$.

Lift the actual map
$(x,z)\mapsto B'_z(p(x))$
on $D_\Gamma\times D^2$
through (C.69.6), with the
identity lift at $z_*$.
The product is simply
connected, so the lift
exists and is unique.
Lift the inverse family
with the same normalization.
Their compositions lift
the projection $p$ and
equal identity at $z_*$,
so uniqueness makes them
mutual inverses everywhere.
For every $\gamma\in\Gamma$,
the two compositions of
the lifted family with
$\gamma$ lift the same
map and agree at the
based parameter. They
therefore agree for all
parameters: every deck
label is retained.

For $x\in\Omega_\Gamma$,
the lifted value remains
in the discrete fiber over
$p(x)$, because $B'_z$
fixes that boundary point.
Connectedness of $D^2$
and the identity value
at $z_*$ force it to
remain exactly $x$.
On the parameter boundary,
uniqueness of lifts over
the connected product
$D_\Gamma\times\partial D^2$
recovers the original
$f_t$ exactly.
This gives the desired
continuous filling in
the right side of (C.69.2).
Its continuous inverse
restriction extends the
entire filling to (C.69.1).

The elementary nontrivial
groups require no omitted
case. For a parabolic
generator, the upper-half-plane
coordinate $z\mapsto z+1$
identifies $D_\Gamma$ with
the closed upper half plane
and $N_\Gamma$ with
$S^1\times[0,\infty)$.
For a hyperbolic generator
$z\mapsto\lambda z$,
$\lambda>1$, logarithmic
polar coordinates
$(s,\theta)\in\mathbb R\times[0,\pi]$
give the deck translation
$s\mapsto s+\log\lambda$
and the compact annular
quotient. In both cases
the pointed group used in
(C.69.17) is contractible
by precisely the quoted
exception lists.
The trivial case was
handled at the start.

\emph{The original conjugacy and the frontier transport.}
Conjugation by the specified
closed-disk homeomorphism
$k$ is a topological group
isomorphism carrying the
centralizer of $k^{-1}\Gamma k$
to (C.69.1).
Transport the proved filling
by its inverse. Every
conjugate operator retains
its original label, and
the boundary is fixed
throughout.

For a spherical domain,
compose with the actual
$T_q$ and prime-end maps
of C.68. The exact
conjugation equation
(C.68.12) proves the
claimed ambient commutation,
and the original restriction
kernel acts identically
on the domain throughout.
The complement is fixed.
When the specified infinity
lies in that complement,
the entire filling descends
to the original plane
homeomorphism group with
the topology comparison
already proved in C.66.
\end{proof}

\emph{Remaining scope.}
The original closed-disk
conjugacy is retained in
this theorem. The construction
does not assert that a
uniformizing conjugacy of
an arbitrary interior proper
action extends to its
original prime-end circle.
Nor does it assert uniform
pasting over infinitely many
different domain conjugacies.
Those parts of the general
C.12 extension problem
remain unresolved.

\begin{claim}{C.70}{Proved}{Operator-compatible carriers for the original evaluation fillings}
Let $f_t$, $G$, $H$ and the $H$-invariant union $W$ of
original open $G$-orbits be as in C.64. There is one compact
path-connected family $P$ of actual global $H$-centralizing
homeomorphisms, containing the identity and every $f_t$,
preserving each original open orbit and fixing every
$G$-fixed point, with the following properties.

For a nonempty compact set $K\subset\R^2$, define its
filled hull by
\[
 \operatorname{Fill}(K)
   =\R^2\setminus U_\infty(K),
 \qquad U_\infty(K)\text{ the unbounded component of }\R^2\setminus K.
                                                               \tag{C.70.1}
\]
For $x\in W$, put $K_x=P^2(x)$. There is a unique
component $V_x$ of $\operatorname{int}\operatorname{Fill}(K_x)$
containing the entire set $P(x)$. It is simply connected.
Write $O_x$ for the original orbit containing $x$, and let
$U_x$ be the component of $V_x\cap O_x$ containing $P(x)$.
These actual domains satisfy
\[
 P(x)\subset U_x\subset O_x\cap\operatorname{Fill}(P^2(x)),
 \qquad hU_x=U_{hx}\quad(h\in H).
                                                               \tag{C.70.2}
\]
Inclusion induces an injection
\[
                  \pi_1(U_x,x)\longrightarrow\pi_1(O_x,x).
                                                               \tag{C.70.3}
\]
The original evaluation loop $t\mapsto f_t(x)$ therefore
has a based disk filling $b_x:D^2\to U_x$.
The fillings can be chosen simultaneously as a set-indexed
family satisfying the exact operator equations
\[
               b_{hx}(z)=h(b_x(z))
                    \qquad(x\in W,\ h\in H,\ z\in D^2).
                                                               \tag{C.70.4}
\]
Every choice of such fillings obeys the fixed-frontier control
\[
 \sup_{z\in D^2}|b_{x_n}(z)-y|\longrightarrow0
       \quad\text{for }x_n\to y\text{ with }y\in\operatorname{Fix}(G).
                                                               \tag{C.70.5}
\]
This is a statement about the specified individual evaluation
fillings and their simultaneous operator labels. Joint
continuity in $x$, avoidance of collisions between distinct
evaluation points, and extension to a family of homeomorphisms
are not established here.
\end{claim}
\passage{Section 23, folios 140-141, scans 279-281; Part C(b), infinite-operator zero-level comparison and controlled fixed-frontier evaluation fillings}

\begin{proof}
\emph{The compact family and its actual orbit neighborhoods.}
Let $\mathcal K$ be C.64's compact path-connected labelled
family, and retain its projection to actual global maps.
Put
\[
 P=\mathcal K\mathcal K^{-1}
        \ \cup\ \{f_t:0\leq t\leq1\}
        \ \cup\ \{f_t^{-1}:0\leq t\leq1\}.
                                                               \tag{C.70.6}
\]
Each of these three compact sets is path connected and
contains the identity; their union has the same properties.
Composition and inversion take place in the original global
homeomorphism group. Every displayed map commutes with each
original $h\in H$, preserves $O_x$, and fixes $\operatorname{Fix}(G)$.
The source-cover labels, when required, are the corresponding
products and inverses of C.64's retained labels and of the
original lifted loop.

By C.64.1, $x$ is in the plane interior of $P(x)$:
the orbit is an open subset of the plane. For each $p\in P$,
the image under the original homeomorphism $p$ of that
neighborhood lies in $pP(x)\subset P^2(x)$.
Consequently
\[
                   P(x)\subset\operatorname{int}P^2(x).
                                                               \tag{C.70.7}
\]
The set $P(x)$ is path connected and contains $x$.
The set $K_x=P^2(x)$ is also compact and path connected,
since it is the continuous evaluation image of $P\times P$.

\emph{Filled hulls and their simply connected interior components.}
For any compact $K$, the complement has exactly one
unbounded component: outside a sufficiently large closed
disk it is connected, and every unbounded component meets
that exterior. The filled hull is closed and bounded, hence
compact. It is contained in the Euclidean convex hull of $K$.
Indeed, a point outside that convex hull is separated from
it by a line and can be joined to infinity within the
corresponding open half-plane, without meeting $K$.

Every component $V$ of $\operatorname{int}\operatorname{Fill}(K)$
is simply connected. To prove this directly, take a Jordan
curve $J\subset V$ and its bounded Jordan disk $D_J$.
The connected unbounded set $U_\infty(K)$ is disjoint from
$J$, so it lies entirely on its unbounded side. Hence
$D_J\subset\operatorname{Fill}(K)$. The interior of
$D_J$ is open and is consequently in
$\operatorname{int}\operatorname{Fill}(K)$; it belongs
to the same component as $J$. Thus the whole disk lies in $V$.
For an arbitrary loop in $V$, polygonal approximation
inside $V$ and a finite planar regular neighborhood of
the resulting graph express it as a product of based
conjugates of its inner boundary Jordan curves. Each
bounds a disk in $V$ by what we just proved. This gives
a based null-homotopy of the original loop in $V$.

Apply this with $K=K_x$. Equation (C.70.7) places the
connected set $P(x)$ in a single component $V_x$ of
the interior of the filled hull, proving existence and
uniqueness of that component. It is simply connected
by the preceding argument. Since $P(x)\subset O_x$,
it is contained in one component $U_x$ of $V_x\cap O_x$.
This proves the inclusions in (C.70.2).

\emph{The actual inclusion on fundamental groups.}
The proof of C.65.4 works with a simply connected open
plane domain in place of its open Euclidean disk. We give
the comparison in this form. Let $O,V$ be open plane
domains with $V$ simply connected, and let $U$ be a
component of $O\cap V$. Take a based loop $\gamma$ in
$U$ which is null-homotopic in $O$.
Approximate it by a based polygonal loop within $U$,
and take a compact connected planar regular neighborhood
$M\subset U$ of its finite graph. The bounded disk
inside its outer boundary lies in $V$, by simple
connectedness and Jordan separation. In particular all
its inner complementary Jordan disks lie in $V$.

Fill precisely those inner disks wholly contained in $O$,
obtaining a compact connected planar surface $M'\subset U$.
From each remaining inner disk choose a point $z_i$ outside
$O$. The inclusion
\[
                   M'\hookrightarrow\R^2\setminus\{z_i\}
                                                               \tag{C.70.8}
\]
has the explicit strong deformation retraction of C.65.9:
choose a Schoenflies coordinate on each inner disk taking
its chosen $z_i$ to zero, apply the radial formula
$(r,\theta)\mapsto((1-\tau)r+\tau,\theta)$ there, and
apply that formula in the exterior Schoenflies coordinate
outside the outer boundary. The identity on $M'$ completes
the finite closed pasting. Every boundary and the chosen
basepoint stay fixed.
The original null-homotopy in $O$ avoids all $z_i$;
compose it with this retraction and append the polygonal
approximation homotopy. This gives a based filling in $U$.
The actual inclusion $U\hookrightarrow O$ is therefore
injective on fundamental groups. Taking $O=O_x$ and
$V=V_x$ proves (C.70.3).

By C.32, the original homeomorphism loop restricted
to $O_x$ is null-homotopic in its homeomorphism group.
Evaluation at the original point $x$ is continuous and
gives a based null-homotopy of $t\mapsto f_t(x)$ in
$O_x$. That entire loop lies in $P(x)\subset U_x$.
Equation (C.70.3) therefore gives the asserted filling
$b_x$ with image in this actual $U_x$.

\emph{All original operators and uniform control at fixed points.}
A plane homeomorphism sends the unbounded component
of the complement of a compact set to the unbounded
component of the image complement. To see unboundedness,
choose a sequence in that component escaping every
compact set. Properness of the inverse prevents its
image from lying in a compact set. It follows that
\[
          h\operatorname{Fill}(K)=\operatorname{Fill}(hK).
                                                               \tag{C.70.9}
\]
Every member of $P$ commutes with $h\in H$, so
$hP^2(x)=P^2(hx)$ and $hP(x)=P(hx)$.
Equation (C.70.9) sends the specified component $V_x$
to $V_{hx}$. Commutation with the original generated
group also gives $hO_x=O_{hx}$. Taking the indicated
components now proves $hU_x=U_{hx}$ exactly.

Choose one representative $x$ of each $H$-orbit in
$W$ and one of its based fillings in $U_x$. Define
the filling at $hx$ by $h\circ b_x$. This is independent
of the choice of $h$: an operator fixing $x$ fixes
every point of $O_x$, since each such point is $gx$
for some $g\in G$ and $h(gx)=g(hx)=gx$.
Its image therefore fixes the entire chosen filling.
The boundary identity follows from $hf_t(x)=f_t(hx)$.
This proves (C.70.4) for every original operator,
without any choice of an effective quotient representative.

Finally let $x_n\to y$ with $y\in\operatorname{Fix}(G)$.
The compact family $P^2$ fixes $y$ pointwise. Continuity
of evaluation on
$P^2\times(\{y\}\cup\{x_n:n\geq1\})$ gives
\[
             \sup_{p\in P^2}|p(x_n)-y|\longrightarrow0.
                                                               \tag{C.70.10}
\]
The Euclidean ball centered at $y$ with this radius
contains the convex hull of $P^2(x_n)$, and therefore
contains its filled hull and $U_{x_n}$. Every point of
every chosen filling lies in that same ball. This
proves (C.70.5), uniformly over the disk parameter
and independently of every filling choice.

The family just selected is equivariant as a family
indexed by points. Its proof imposes no continuity
between fillings at distinct source points. It also
imposes no configuration-space disjointness. Those
requirements, needed for an ambient homeomorphism
filling, remain part of C.12.
\end{proof}

\begin{claim}{C.71}{Proved}{Actual disk comparisons for two-dimensional operators and a cyclic compactification}
Let $Q$ be a closed locally compact subgroup of
$\operatorname{Homeo}(\mathbb D)$, with its actual compact-open
and inverse topology, acting freely and properly on $\mathbb D$.
When $\dim Q=2$, put
\[
 E_Q=\{h\in\mathcal D_\partial:
                    h|_{\mathbb D}\in
                    C_{\operatorname{Homeo}(\mathbb D)}(Q)\}.
                                                               \tag{C.71.1}
\]
Then every based sphere family $S^n\to E_Q$, $n\geq1$, has a
based contraction in $E_Q$. In particular every based loop
fills there. No extension of the $Q$-action to the boundary,
and no closed-disk conjugacy for that action, is assumed.
If $A\leq\operatorname{Homeo}(\overline{\mathbb D})$ has
actual restricted closure $Q$, then
$E_Q=C_{\mathcal D_\partial}(A)$ by continuity on the
interior and density in the closed disk.

There is also an explicit cyclic action on the actual closed
disk whose interior image is closed, discrete, free and proper,
and which admits no closed-disk
conjugacy to a torsion-free cyclic Fuchsian action. For its
specified invariant partial compactification, the analogue of
the restriction isomorphism (C.69.2) is not surjective. The
proof gives the actual maps witnessing both assertions.
Nevertheless its boundary-fixed centralizer is contractible:
it is identified topologically with an explicitly described
group of continuous families of increasing real lifts, and
linear interpolation supplies the contraction. Its interior
orbits are precisely the displayed vertical lines, so this
example does not supply an open transitive orbit for that
boundary-fixed centralizer.
\end{claim}
\passage{Section 23, folios 140-141, scans 279-281; Part C(b), admissible infinite-operator level-zero comparison and relative disk centralizers}


\begin{proof}
\emph{The actual two-dimensional transformation group.}
The Lie conclusion and principal orbit comparison of C.33
apply to $Q$. The identity component has open orbits because
the free orbit dimension is two. All these orbits are open;
connectedness of $\mathbb D$ leaves just one. Given $q\in Q$,
transitivity supplies $q_0\in Q^0$ with $q_0x=qx$ at one
specified $x$. Freeness gives $q=q_0$. Thus $Q=Q^0$, and
evaluation
\[
                 \epsilon_x:Q\longrightarrow\mathbb D,
                       \qquad q\longmapsto qx
                                                               \tag{C.71.2}
\]
is a homeomorphism. The inverse is the principal orbit
coordinate from C.33; alternatively properness makes the
continuous orbit bijection a homeomorphism. In particular
$Q$ is simply connected.

Its two-dimensional Lie algebra is either abelian or has
a basis $X,Y$ with $[X,Y]=Y$: in the nonabelian case the
derived algebra is a line and rescaling a complement gives
that bracket. The simply connected groups are respectively
$\mathbb R^2$ and
\[
 (a,b)(a',b')=(a+a',b+e^a b'),
                   \qquad(a,b)\in\mathbb R^2 .
                                                               \tag{C.71.3}
\]
This uses the actual Lie-group isomorphisms integrating those
bases. Their inverses are the integrations in the opposite
direction; the uniqueness for simply connected groups gives
both inverse identities. In each case exponential is a
global diffeomorphism. In (C.71.3) its explicit formula is
$\exp(a,v)=(a,v(e^a-1)/a)$, with the quotient at $a=0$
defined to be $1$.

Through (C.71.2), a map commuting with all left translations
is precisely a right translation. Indeed its value $c$ at
the identity forces its value at every $q$ to be $qc$.
Right translations and their inverses depend continuously
on $c$ in the compact-open topology, and evaluation gives
the inverse comparison. Thus the actual interior centralizer
is topologically the opposite group
\[
 N=C_{\operatorname{Homeo}(\mathbb D)}(Q)\cong Q^{\rm op}.
                                                               \tag{C.71.4}
\]
Inversion is an isomorphism $Q^{\rm op}\to Q$, so $N$ has
the same two specified simply connected exponential models.

\emph{Compact words retain the given disk boundary.}
Let $a:S^n\to E_Q$ be based, with $n\geq1$, and let
$g_z=a(z)|_{\mathbb D}\in N$. Define the interior maps
\[
               R_{r,z}=\exp_N(r\log_N g_z),
                     \qquad 0\leq r\leq1 .
                                                               \tag{C.71.5}
\]
They are continuous in the actual interior homeomorphism
topology. They equal the identity for $r=0$ and at the
sphere basepoint, and equal $g_z$ for $r=1$.

Put
\[
 P=a(S^n)\cup a(S^n)^{-1}\subset\mathcal D_\partial .
                                                               \tag{C.71.6}
\]
This is compact and contains the identity. For each $z$,
choose a path in $S^n$ from its basepoint to $z$.
Applying the exact two- or three-letter formula (C.37.2)
to its restricted path in the simply connected group $N$
places $R_{r,z}$ in the restriction of $P^3$.
No continuous choice of these sphere paths is used.
The resulting restriction map
\[
                  P^3\longrightarrow
                         \operatorname{Homeo}(\mathbb D)
                                                               \tag{C.71.7}
\]
is continuous and injective: two closed-disk maps agreeing
on the dense interior agree everywhere. Its source is
compact and its target is Hausdorff, so it is a
homeomorphism onto its image. Consequently (C.71.5)
has a unique continuous lift
$\overline R_{r,z}\in P^3$ through (C.71.7).
Continuity here includes the inverse maps, since the
source and target carry their homeomorphism-group
topologies. All these lifted maps fix the original
circle pointwise and centralize $Q$ on the interior.
They give the claimed based contraction. This argument
uses compact word control to prove boundary continuity;
it does not identify the full relative group topology
with the topology of an arbitrary subgroup of $N$.

\emph{The cyclic action and its exact disk coordinate.}
Write
\[
 \kappa(x)=\frac{x}{1+|x|},\qquad
 \kappa^{-1}(s)=\frac{s}{1-|s|}\quad(-1<s<1),
 \qquad S=[-1,1]^2.
                                                               \tag{C.71.8}
\]
The extension of $\kappa$ takes $-\infty,+\infty$ to
$-1,+1$. An explicit homeomorphism $J:S\to
\overline{\mathbb D}$ and its inverse are
\[
 \begin{aligned}
 J(v)&=\frac{\|v\|_\infty}{\|v\|_2}v\quad(v\ne0),
                          &J(0)&=0,\\
 J^{-1}(w)&=\frac{\|w\|_2}{\|w\|_\infty}w\quad(w\ne0),
                          &J^{-1}(0)&=0 .
 \end{aligned}                                                   \tag{C.71.9}
\]
The formulas are inverse on each ray and continuous at
zero. They take the square boundary onto the unit circle.
Retain the actual interior coordinate
$p=J\circ(\kappa\times\kappa):\mathbb R^2\to\mathbb D$.

On the real plane take $T(x,y)=(x,y+1)$.
Its extension on $S$ is
\[
 \overline T(s,t)=(s,\tau(t)),\qquad
 \tau(t)=\kappa(\kappa^{-1}(t)+1)\quad(|t|<1),
       \qquad\tau(-1)=-1,\quad\tau(1)=1.
                                                               \tag{C.71.10}
\]
The specified closed-disk generator is
$J\overline T J^{-1}$. The group
$Q=p\langle T\rangle p^{-1}$ is free and proper on
the interior: a compact real-plane set has bounded
$y$-coordinate and meets only finitely many of its
integer translates. It is closed and discrete in the
actual interior homeomorphism group, since evaluation
at $(0,0)$ sends $T^m$ to $(0,m)$; a convergent sequence
of such maps has eventually constant integer parameter.
Conjugation by $p$ transports these assertions, with
its inverse giving the opposite comparison.

The fixed set of $\overline T$ on $\partial S$ is
exactly the union of the two closed horizontal edges.
Indeed $\tau(t)>t$ for $-1<t<1$. Under a purported
closed-disk conjugacy $k$ to a generator $\gamma$ of
a torsion-free cyclic Fuchsian action, the equation
$k(J\overline T J^{-1})=\gamma k$ would imply
\[
       kJ\bigl(\operatorname{Fix}(\overline T)
                         \cap\partial S\bigr)
                  =\operatorname{Fix}(\gamma)\cap S^1.
                                                               \tag{C.71.11}
\]
The left side consists of two nondegenerate arcs,
whereas the right side has one point for a parabolic
element or two points for a hyperbolic element.
Since $kJ$ is injective, (C.71.11) is impossible.

There is a direct failure of the analogous partial
restriction comparison as well. Set
\[
 S_*=[-1,1]\times(-1,1),\qquad
 \Omega_*=\{-1,1\}\times(-1,1).
                                                               \tag{C.71.12}
\]
The $\langle\overline T\rangle$ action on $S_*$ is
free and proper. Through $(s,y)\mapsto(s,\kappa(y))$,
its actual quotient is the closed annulus
$[-1,1]\times(\mathbb R/\mathbb Z)$, with deck map
$y\mapsto y+1$ and boundary $\Omega_*/\langle\overline
T\rangle$. The horizontal translation $L(x,y)=(x+1,y)$
extends on $S$ as
$\overline L(s,t)=(\tau(s),t)$. It commutes with
$\overline T$, and its restriction to $S_*$ fixes
$\Omega_*$ pointwise. It does not fix the horizontal
edges pointwise. Thus
\[
 C_{\operatorname{Homeo}(S,\partial S)}(\overline T)
       \longrightarrow
 C_{\operatorname{Homeo}(S_*,\Omega_*)}(\overline T)
                                                               \tag{C.71.13}
\]
is not surjective: any preimage of $\overline L|_{S_*}$
would equal $\overline L$ on the dense set $S_*$ and
hence on $S$. The maps $J$ and $J^{-1}$ transport this
same counterexample to the specified disk.

\emph{The actual boundary-fixed centralizer of this action.}
Let $\mathcal V$ consist of continuous functions
$F:\mathbb R^2\to\mathbb R$, written $F_x(y)$, such that
each $F_x$ is an increasing homeomorphism of $\mathbb R$,
and
\[
 F_x(y+1)=F_x(y)+1,\qquad
 \lim_{|x|\to\infty}\sup_{0\leq y\leq1}
                         |F_x(y)-y|=0 .
                                                               \tag{C.71.14}
\]
Its product and inverse are the pointwise fiber
composition and inverse. Give it the metric topology
\[
 d_{\mathcal V}(F,G)=
  \sup_{x,y}|F_x(y)-G_x(y)|
   +\sup_{x,y}|F_x^{-1}(y)-G_x^{-1}(y)|.
                                                               \tag{C.71.15}
\]
Both suprema are finite. The displacements are
continuous and periodic in $y$, and vanish as
$|x|\to\infty$, so they are bounded globally.
The inverse displacements have the same suprema
at each $x$ and the same vanishing property.
Joint continuity of $(x,y)\mapsto F_x^{-1}(y)$
follows directly: on a compact $x$-interval the
periodic displacements are uniformly bounded,
which confines all inverse arguments to a compact
interval; strict monotonicity identifies every
subsequential limit.

There is an exact group isomorphism
\[
 \begin{split}
 \mathcal V&\longrightarrow
       C_{\operatorname{Homeo}(S,\partial S)}(\overline T),\\
 F&\longmapsto H_F,\qquad
 H_F(\kappa(x),\kappa(y))
          =(\kappa(x),\kappa(F_x(y))),
 \end{split}                                                   \tag{C.71.16}
\]
with the identity on $\partial S$. To prove surjectivity,
write the real-plane restriction of a map on the right
as $(X(x,y),Y(x,y))$. Commutation gives
$X(x,y+m)=X(x,y)$ and
$Y(x,y+m)=Y(x,y)+m$ for each integer $m$.
Letting $m\to+\infty$ and using the actual fixed top
edge forces $X(x,y)=x$. Bijectivity now gives a
homeomorphism $F_x(y)=Y(x,y)$ of each fiber, and the
integer equation forces it to be increasing.
Continuity at the two fixed vertical edges,
uniformly for $y\in[0,1]$, is exactly the vanishing
condition in (C.71.14).

Conversely (C.71.14) defines a homeomorphism of
the real plane with the stated inverse. For
bounded $x$ its displacement is uniformly bounded,
so it extends as the identity at the horizontal
edges. For $|x|\to\infty$ its displacement tends
uniformly to zero, giving the identity at the
vertical edges and at the corners. These assertions
also apply to the inverse. They prove the claimed
extension and bijection.

The isomorphism (C.71.16) is topological. In the
forward direction $\kappa$ is globally Lipschitz,
so each supremum in (C.71.15) controls the uniform
square displacement of the corresponding maps or
inverses. Conversely suppose $H_{F_j}$ and their
inverses converge uniformly to $H_F$ and its inverse.
For real $y\in[0,1]$, the values $F_x(y)$ lie in
one fixed compact real interval, uniformly in $x$.
Uniform convergence of their $\kappa$-coordinates
keeps all $F_{j,x}(y)$ in a slightly larger compact
real interval. Uniform continuity of $\kappa^{-1}$
there gives uniform convergence of these actual
values. Periodicity extends it to every real $y$.
Applying the same argument to the inverse maps
gives exactly (C.71.15).

Finally define
\[
       \mathcal H_r(F)_x(y)
                      =(1-r)F_x(y)+ry,\qquad 0\leq r\leq1 .
                                                               \tag{C.71.17}
\]
It is strictly increasing, proper, and has the
same integer equivariance and vanishing condition.
Forward continuity in (C.71.15) is immediate.
For inverse continuity, each fixed element of
$\mathcal V$ and its inverse is uniformly continuous
in the fiber variable, uniformly in $x$: its
displacement is continuous on
$\mathbb R\times(\mathbb R/\mathbb Z)$ and vanishes
at the two $x$-ends. The elementary increasing-map
inverse estimate therefore turns uniform forward
convergence to a member of $\mathcal V$ into uniform
inverse convergence. This applies to (C.71.17)
when $(r,F)$ varies. Hence it is a continuous
contraction of the full topological group to the
identity, fixing the identity throughout.
Conjugation by $J$ gives the stated disk contraction.

Every such map preserves the real $x$-coordinate.
For each specified $x$ and $c\in\mathbb R$, a
continuous compactly supported function $b$ with
$b(x)=c$ gives $F_u(y)=y+b(u)\in\mathcal V$.
Thus the exact group orbits in the interior are
the vertical lines. This records the transitivity
scope of the example explicitly.

The exact certificate
\path{code/part_b/c71_disk_compactification_certificate.py}
checks the rational compactification and inverse,
the translation conjugacies and commutation,
the displayed boundary and quotient maps, and
nonlinear periodic fiber models together with
their interpolation and inverse formulas.
The full topological comparisons and contractions
are proved above. This claim leaves the general
zero- and one-dimensional proper-operator disk
comparison unresolved.
\end{proof}

\begin{claim}{C.72}{Proved}{A spiral boundary obstruction with the full discrete local centralizer}
There is an explicit based loop of closed-disk homeomorphisms
fixing the boundary pointwise whose generated group $G$ is
transitive on the interior and has full interior centralizer
equal to a specified closed discrete free proper cyclic group
$Q=\langle T_*\rangle$. Every forward $Q$-orbit has the whole
boundary circle as its limit set. Nevertheless the actual
restriction map
\[
 C_{\mathcal D_\partial}(Q)\longrightarrow
              C_{\operatorname{Homeo}(\mathbb D)}(Q)
                                                               \tag{C.72.1}
\]
is not surjective: the proof displays an interior centralizer
map with no continuous extension at any boundary point.
Here $Q$ on the left means its specified closed-disk action.

The proof also supplies an explicit interior conjugacy
$j$ from $Q$ to a parabolic Fuchsian group. Under that
conjugacy the same obstructed map does extend as a
boundary-fixed homeomorphism. Thus the map
$h\mapsto jhj^{-1}$ between the two interior centralizers
does not identify their actual boundary-relative subgroups.
The exact failed inclusion is witnessed below; no
comparison of unnamed compactifications is used.

The original loop in this example has an explicit
boundary-fixed contraction. The result therefore
refutes the displayed extension comparison, including
the versions supplemented by transitivity and the
full-circle orbit limit set; it supplies no failure of
the level-zero loop comparison in C.12.
\end{claim}
\passage{Section 23, folios 140-141, scans 279-281; Part C(b), admissible infinite-operator level-zero comparison and the actual relative disk extension problem}


\begin{proof}
\emph{The actual spiral coordinate and its inverse.}
Regard the real plane as $\mathbb C$, with
$z=x+iy$, and define
\[
 \begin{split}
 p(z)&=\frac{z}{1+|z|}
                \exp\bigl(i\log(1+|z|)\bigr),\\
 p^{-1}(w)&=\frac{w}{1-|w|}
                \exp\bigl(i\log(1-|w|)\bigr)
                            \qquad(w\in\mathbb D).
 \end{split}                                                   \tag{C.72.2}
\]
Both maps take zero to zero. The radial coordinates
are $r\mapsto r/(1+r)$ and
$\rho\mapsto\rho/(1-\rho)$, which are inverse.
The angular changes in the two displayed formulas
are opposite because $1+r=(1-\rho)^{-1}$.
They prove both inverse identities and continuity,
including at zero. Thus (C.72.2) is an actual
homeomorphism of the plane with the disk interior.
It is used to specify the action on the given
closed disk; no boundary extension of $p$ is asserted.

The following elementary estimate retains its boundary
topology. Suppose $|a|\leq M$ and $r=|z|>2M+1$.
Write $r'=|z+a|$. The reverse triangle inequality
gives $|r'-r|\leq M$, and
\[
 \begin{split}
 \left|\frac{z+a}{r'}-\frac z r\right|
                     &\leq\frac{2M}{r-M},\\
 |\log(1+r')-\log(1+r)|
                     &\leq\frac{M}{1+r-M},\\
 \left|\frac{r'}{1+r'}-\frac r{1+r}\right|
                     &\leq\frac{M}{(1+r-M)(1+r)} .
 \end{split}                                                   \tag{C.72.3}
\]
The first inequality follows by adding and subtracting
$z/r'$, and the second by the mean value theorem.
Combining them with $|e^{iu}-e^{iv}|\leq|u-v|$ gives
\[
                       |p(z+a)-p(z)|
                            \leq\frac{4M}{r-M}.
                                                               \tag{C.72.4}
\]
In particular, any plane homeomorphism $h$ for
which both $h(z)-z$ and $h^{-1}(z)-z$ are bounded
gives a closed-disk homeomorphism
\[
 \overline h(w)=
 \begin{cases}
       ph p^{-1}(w),&w\in\mathbb D,\\
       w,&w\in S^1 .
 \end{cases}                                                   \tag{C.72.5}
\]
For a compact continuous parameter family with a
common bound on both displacements, (C.72.4) gives
joint boundary continuity of the maps and inverses.
Interior continuity is just the actual conjugation
in (C.72.2). These observations prove continuity
in the full closed-disk homeomorphism topology.

\emph{The actual cyclic and translation actions.}
For $a\in\mathbb C$ put $L_a(z)=z+a$ and denote
its extension (C.72.5) by $\overline L_a$.
The equations
\[
 \overline L_a\overline L_b=\overline L_{a+b},
       \qquad\overline L_a^{-1}=\overline L_{-a}
                                                               \tag{C.72.6}
\]
hold first on the interior and then on the boundary.
The parameter map $a\mapsto\overline L_a$ is continuous
by (C.72.4), and its inverse on its image is
$h\mapsto p^{-1}(h(0))$. Its image is closed:
if $\overline L_{a_n}$ converges to a closed-disk
homeomorphism $h$, then $h(0)$ is interior and
$a_n\to p^{-1}(h(0))$. Thus this is a specified
closed translation subgroup, topologically $\mathbb R^2$,
acting transitively on $\mathbb D$ and fixing $S^1$
pointwise.

Set
\[
             T(x,y)=(x,y+1),\qquad
             T_*=\overline L_i,\qquad Q=\langle T_*\rangle.
                                                               \tag{C.72.7}
\]
Its interior image is closed and discrete: evaluating
the conjugate $T^m$ at the real origin gives $(0,m)$.
It is free and proper there, since compact real-plane
sets have bounded vertical coordinates and meet only
finitely many integer translates. These assertions
are transported by the actual $p,p^{-1}$.
Every element of $Q$ fixes the entire boundary.

Every forward orbit of $T_*$ has every point of
$S^1$ as a limit. Fix $z=x+iy$ and
$\zeta=e^{i\alpha}$. For large positive integers $k$
choose an integer $n_k$ nearest to
$\exp(2\pi k+\alpha-\pi/2)$. Then
\[
 \begin{split}
 \arg(x+i(y+n_k))&\longrightarrow\pi/2,\\
 \log(1+|x+i(y+n_k)|)
                -(2\pi k+\alpha-\pi/2)&\longrightarrow0,\\
 |p(x+i(y+n_k))|&\longrightarrow1.
 \end{split}                                                   \tag{C.72.8}
\]
The second limit follows since the integer rounding
error is at most $1/2$ and
$|x+i(y+n_k)|/n_k\to1$. Substitution in (C.72.2)
gives $p(T^{n_k}z)\to\zeta$.
Conversely $|T^n z|\to\infty$, so every accumulation
point is on the circle. This proves the stated
full-circle forward limit set for each original point.

\emph{A based loop with precisely this full local centralizer.}
Use the continuous functions
\[
 \varphi(y)=\operatorname{dist}(y,\mathbb Z),
                 \qquad\psi(x)=\max(0,1-|x|).
                                                               \tag{C.72.9}
\]
The first has zero set exactly $\mathbb Z$ and
period one; the second is nonzero and compactly
supported. For $0\leq u\leq1$ set
\[
 \begin{array}{ll}
 P_1(u)(x,y)=(x+u,y),&
 P_2(u)(x,y)=(x,y+u),\\
 P_3(u)(x,y)=(x+u\varphi(y),y),&
 P_4(u)(x,y)=(x,y+u\psi(x)).
 \end{array}                                                   \tag{C.72.10}
\]
Each is a homeomorphism, with inverse obtained by
replacing $u$ by $-u$ in the same formula. They
all commute with $T$. Their displacements and those
of their inverses have a common bound of one, so
their extensions $\overline P_j(u)$ in (C.72.5)
are continuous paths in
$C_{\mathcal D_\partial}(Q)$.

For $(j-1)/4\leq t\leq j/4$ define the triangular
amplitude
\[
 a_j(t)=
   \min\{\,8(t-(j-1)/4),\,2-8(t-(j-1)/4)\,\},
 \qquad f_t=\overline P_j(a_j(t)).
                                                               \tag{C.72.11}
\]
At every joining parameter the amplitude is zero,
so these are an actual continuous based loop.
Replacing $a_j(t)$ by $(1-r)a_j(t)$ gives its
explicit based contraction in the same relative
centralizer. In particular its ambient class is
zero, with no appeal to a different group.

Let $G=\langle f_t:0\leq t\leq1\rangle$.
It contains all real horizontal and vertical
translations in the $p$-coordinate: the intervals
$[0,1]$ in the first two paths generate the two
additive real groups, using products and inverses.
Thus $G$ is transitive on $\mathbb D$.
An interior homeomorphism centralizing $G$ must
commute with all these translations. Conjugation
by $p^{-1}$ and evaluation at zero then force it
to be a translation $L_{(a,b)}$.
Its commutation with $P_3(1)$ says
$\varphi(y+b)=\varphi(y)$ for every $y$.
At $y=0$ this gives $b\in\mathbb Z$.
Commutation with $P_4(1)$ says
$\psi(x+a)=\psi(x)$ for all $x$.
If $a\ne0$, iteration would give
$\psi(na)=\psi(0)=1$ for arbitrarily large $|na|$,
contradicting its compact support.
Hence $a=0$. Conversely these integer vertical
translations commute with every path (C.72.10).
We have proved the exact equality
\[
             C_{\operatorname{Homeo}(\mathbb D)}(G|_{\mathbb D})
                       =Q|_{\mathbb D}.
                                                               \tag{C.72.12}
\]
In particular the effective full local operator
group has dimension zero and is the named closed
discrete proper group, in precisely the topology
of the original comparison. The larger
$C_{\mathcal D_\partial}(Q)$ is also transitive,
since it contains the same translation subgroup.

\emph{The interior centralizer map that cannot extend.}
The real-plane map
\[
                        A(x,y)=(2x,y)
                                                               \tag{C.72.13}
\]
is a homeomorphism and commutes with $T$. Thus
$A_*=pAp^{-1}$ belongs to the target of (C.72.1).
Fix again $\zeta=e^{i\alpha}$ and, for sufficiently
large $k$, put
\[
 \begin{split}
 z_k&=(e^{2\pi k+\alpha}-1,0),\\
 z'_k&=(0,e^{2\pi k+\alpha-\pi/2}-1).
 \end{split}                                                   \tag{C.72.14}
\]
Both $p(z_k)$ and $p(z'_k)$ approach $\zeta$,
with the angular value exactly $\alpha$ modulo
$2\pi$. On the vertical sequence $A$ is the
identity, so $p(Az'_k)\to\zeta$. On the horizontal
sequence its image angle is
\[
 2\pi k+\alpha+
            \log(2-e^{-(2\pi k+\alpha)}).
                                                               \tag{C.72.15}
\]
Its radius tends to one, so
$p(Az_k)\to e^{i\log2}\zeta$.
These limits differ, since $0<\log2<1<2\pi$.
Thus $A_*$ has no continuous extension at $\zeta$.
As $\zeta$ was arbitrary, it has no continuous
extension at any boundary point. This proves the
failure of surjectivity in (C.72.1).

\emph{The specified interior Fuchsian comparison.}
Let $\mathbb H=\{u+iv:v>0\}$ and retain the actual
homeomorphisms
\[
 \begin{aligned}
 B(x,y)&=y+i e^x,& B^{-1}(u+iv)&=(\log v,u),\\
 C(\zeta)&=\frac{\zeta-i}{\zeta+i},&
 C^{-1}(w)&=i\frac{1+w}{1-w},\\
 j&=CBp^{-1}:\mathbb D\longrightarrow\mathbb D .
 \end{aligned}                                                   \tag{C.72.16}
\]
The displayed inverse formulas give
$j^{-1}=pB^{-1}C^{-1}$.
Direct substitution yields
$BTB^{-1}(\zeta)=\zeta+1$.
Hence $j$ conjugates $Q|_{\mathbb D}$ to the
parabolic Fuchsian group generated by
$C(\zeta\mapsto\zeta+1)C^{-1}$.
This interior conjugacy cannot extend to a
closed-disk homeomorphism: such an extension
would take the whole fixed boundary circle of
$T_*$ bijectively onto the single boundary
fixed point of that parabolic generator.
This is the same direct fixed-set equation as
(C.71.11), with both maps specified here.

For the obstructed map itself, the same
comparison gives
\[
                     BAB^{-1}(u+iv)=u+i v^2 .
                                                               \tag{C.72.17}
\]
This homeomorphism of $\mathbb H$ extends as
the identity on $\mathbb R\cup\{\infty\}$.
At finite real points this follows from
$v^2\to0$, and at infinity from
$|u+i v^2|\to\infty$ whenever
$|u+iv|\to\infty$. The inverse
$u+iv\mapsto u+i\sqrt v$ has the same property.
Consequently $jA_*j^{-1}$ is in the actual
boundary-fixed centralizer of the parabolic
closed-disk action, although $A_*$ is absent
from the original boundary-relative group.
Conjugation by $j$ is an isomorphism of the
interior centralizers, with inverse conjugation
by $j^{-1}$, but it does not carry the original
relative subgroup onto the parabolic one.
This proves the precise comparison failure
claimed in the statement.

All paths and their explicit contraction fix
the circle, so a specified closed-disk embedding
in a surface transports them and extends them
by the identity outside that disk. The inverse
families have the same extension. This places
the example in the original local-orbit setting
without changing its operator or loop labels.
The contraction in (C.72.11) continues to be
the actual contraction there.

The exact certificate
\path{code/part_b/c72_spiral_comparison_certificate.py}
checks the radial inverse and phase formulas,
the two displayed asymptotic phases symbolically,
the specified interior conjugacy and centralizer
map, and the periodic shear, inverse, and loop
identities with rational arithmetic.
The uniform boundary estimates, complete orbit
limit sets, full local-centralizer equality and
topological comparison failures are proved above.
\end{proof}

\begin{claim}{C.73}{Proved}{The exact prime-end comparison for arbitrary simply connected frontiers}
Let $O\subset S^2$ be a simply connected domain with
nondegenerate frontier, and fix a conformal map
$\phi:\mathbb D\to O$. Retain the actual groups
\[
 K_O=\{h\in\operatorname{Homeo}^+(S^2):
                    h|_{S^2\setminus O}=\id\},
             \qquad \mathcal D_\partial
                 =\operatorname{Homeo}(\overline{\mathbb D},S^1).
                                                               \tag{C.73.1}
\]
Every $h\in K_O$ induces the identity on the prime-end
circle of $O$. Its actual prime-end extension defines a
continuous injective homomorphism
\[
 \begin{split}
 \iota_O:K_O&\longrightarrow\mathcal D_\partial,\\
 \iota_O(h)|_{\mathbb D}&=\phi^{-1}h\phi .
 \end{split}                                                   \tag{C.73.2}
\]
All groups have the compact-open and inverse topologies.
The image of (C.73.2) is dense. For locally connected
frontier it is the topological isomorphism of C.68.
For every non-locally-connected frontier its image is
proper, and its inverse on that image is not continuous.
The proof supplies the actual sequences and a disk map
witnessing these two failures.

Every actual ambient operator preserving $O$ induces
its own prime-end map, and (C.73.2) intertwines
conjugation by the original operator with conjugation
by that prime-end map. In particular it identifies
an ambient relative centralizer with the intersection
of its image and the corresponding disk centralizer,
with the original operator labels and restriction
kernel retained. This is a continuous injective
comparison; no injection on fundamental groups or
descent of a disk null-homotopy is asserted.
\end{claim}
\passage{Section 23, folios 140-141, scans 279-281; Part C(b), actual prime-end and ambient relative centralizer comparisons}


\begin{proof}
\emph{The prime-end space and the precise access data.}
Let $O^{\rm P}$ denote the prime-end compactification.
The exact extension of $\phi$ is a homeomorphism
$\phi^{\rm P}:\overline{\mathbb D}\to O^{\rm P}$.
The crosscut-chain definition and the neighborhood
basis are those in
\cite[\S3.1]{C_BinderRojasYampolsky2014}.
We use the spherical version, with diameters taken
in the compact sphere. The explicit spherical
definition and functorial action are also recorded
in \cite[\S2, pp.~1571--1572]{C_OrtegaRuizDelPortal2011}.
No continuous landing map on the whole prime-end
circle is assumed.

Here are the exact crosscut facts needed below.
A genuine crosscut is an arc whose interior lies
in $O$ and whose two distinct endpoints lie on
$\partial O$. In the prime-end disk its closure
is an arc with two distinct prime-end endpoints;
their principal sets are the respective singleton
endpoint sets. Each side of the crosscut determines
the corresponding open circle interval.
An access arc landing at an actual frontier point
has a unique endpoint in the prime-end disk, and
that endpoint has the singleton principal set
consisting of the landing point. Equivalently the
radial image under $\phi$ lands at that point.
These precise statements appear in
\cite[p.~1575 and p.~1581, before the proof proper
of Lemma~7]{C_OrtegaRuizDelPortal2011}; they apply to
spherical frontier points, including infinity.

Let $A\subset S^1$ be the set of accessible prime
ends in the specified coordinate. For $a\in A$,
write $\ell(a)$ for its unique principal point,
so
\[
                  \lim_{r\uparrow1}\phi(ra)=\ell(a).
                                                               \tag{C.73.3}
\]
Every nonempty open circle interval $I$ contains
two members of $A$ with distinct values of $\ell$.
To see this without postulating continuity of
$\ell$ on $I$, first choose a prime-end open
neighborhood whose closure meets the circle
only inside $I$. The crosscut-neighborhood basis
supplies a crosscut side contained in this
neighborhood. The two prime-end endpoints of its
closure lie in $I$, and the preceding genuine
crosscut statement gives them distinct actual
landing points. In particular $A$ is dense,
and $\ell|_{A\cap I}$ is never constant.

An ambient homeomorphism $h$ preserving $O$
takes null chains to null chains: it preserves
crosscuts, their sides and cofinal inclusions,
and uniform continuity on $S^2$ preserves
vanishing diameters. Its inverse gives the
inverse action. This is the actual induced
homeomorphism $h^{\rm P}$ of $O^{\rm P}$.
For $h\in K_O$, let $u$ be its action on the
circle through $\phi^{\rm P}$. It preserves
orientation. Indeed its interior action preserves
the oriented relative class of the closed disk;
naturality of the boundary isomorphism to the
oriented circle class gives degree $+1$ on the
boundary. The access statements give
\[
                    u(A)=A,\qquad \ell(ua)=\ell(a).
                                                               \tag{C.73.4}
\]
The second equality uses the actual fact that
$h$ fixes every point of $\partial O$.

\emph{The unlinked access fibres and circle rigidity.}
There cannot be four distinct accessible prime
ends $a,b,c,d$ in cyclic order such that
\[
                   \ell(a)=\ell(c)=p,\qquad
                   \ell(b)=\ell(d)=q,\qquad p\ne q.
                                                               \tag{C.73.5}
\]
For a complete separation argument, take radial
tails at $a$ and $c$ in the closed unit disk
and join their interior endpoints by a simple
arc inside the disk, disjoint from those tails
except at the two joining endpoints. For example,
an arc on a smaller concentric circle gives
such a joining arc. Their union is a genuine
disk crosscut $\alpha$ from $a$ to $c$.
By (C.73.3), $\phi(\alpha\cap\mathbb D)$
closes at both ends to the same point $p$.
It is injective on its interior and does not
contain $p$ there. The resulting map from the
circle obtained by identifying the two endpoints
is a continuous bijection onto its image.
Compactness makes this image a Jordan curve
$J\subset S^2$, and
\[
                         J\cap\partial O=\{p\}.
                                                               \tag{C.73.6}
\]
The two sides of $\mathbb D\setminus\alpha$
map to opposite components of $S^2\setminus J$.
This can be checked at any interior point of
$\alpha$, where the open-set homeomorphism
$\phi$ takes its two local sides to the two
local Jordan sides; connectedness then gives
the same assignment for each entire component.
The radial tails at $b$ and $d$ lie on opposite
disk sides. Their landing point $q$ therefore
lies in the closures of both Jordan domains,
whose intersection is $J$. Equation (C.73.6)
would force $q=p$. This proves the impossibility
of (C.73.5), without any assertion about other
impressions.

Suppose $u$ is nonidentity. Its moved set is
open and meets $A$, so choose a moved $a\in A$.
Choose real representatives so that
$a<u(a)<a+1$. There is a small open interval
$I$ just after $a$ such that every $b\in I$
satisfies
$a<b<u(a)<u(b)<a+1$.
The nonconstancy proved above supplies
$b\in A\cap I$ with $\ell(b)\ne\ell(a)$.
Equation (C.73.4) now gives exactly the forbidden
four-tuple $a,b,u(a),u(b)$.
Consequently $u=\id$.
This proves the prime-end rigidity assertion
and defines (C.73.2) on every element of $K_O$.
The group equations follow from the interior
conjugacy and density of the interior. Injectivity
also follows from that conjugacy and the specified
identity outside $O$.

\emph{Continuity in the actual ambient topology.}
The group $K_O$ is a closed subgroup of
$\operatorname{Homeo}^+(S^2)$, hence Polish;
$\mathcal D_\partial$ is Polish as well.
The graph
\[
       \mathcal G=
       \{(h,g)\in K_O\times\mathcal D_\partial:
                       g|_{\mathbb D}=\phi^{-1}h\phi\}
                                                               \tag{C.73.7}
\]
is closed. Namely, convergence of both coordinates
preserves the displayed equality at every interior
point; density and the fixed circle then determine
$g$ uniquely. Thus $\mathcal G$ is a closed
Polish subgroup. Its projection to $K_O$ is a
continuous bijective homomorphism, by the existence
and uniqueness already proved. The Polish-group
open-mapping consequence of the orbit theorem
used in C.30 makes that projection open: apply
the theorem to the transitive left action of
$\mathcal G$ on the Polish group $K_O$ through
the projection. Its inverse is therefore
continuous. Composing that inverse with the
second projection proves continuity of $\iota_O$,
including for compact families of maps and
inverses. This proves no continuity of its
inverse on its image.

Every compactly supported disk homeomorphism
$g$ belongs to the image. Its interior conjugate
$\phi g\phi^{-1}$ has compact support in $O$,
so extension by the identity gives an element
of $K_O$. Such disk maps are dense in
$\mathcal D_\partial$: for $0<r<1$, define
\[
 g^{[r]}(z)=
 \begin{cases}
      r\,g(z/r),& |z|\leq r,\\
      z,& |z|\geq r .
 \end{cases}                                                   \tag{C.73.8}
\]
The boundary-fixed equation makes these formulas
agree. Their inverse has the same formula with
$g^{-1}$. Uniform continuity of $g,g^{-1}$
on the closed disk gives
$g^{[r]}\to g$ with inverses as $r\uparrow1$.
Thus the image in (C.73.2) is dense.

\emph{A nontrivial impression gives both failures of comparison.}
In the fixed coordinate the impression of
$\zeta\in S^1$ can be written
\[
 I(\zeta)=\bigcap_{\epsilon>0}
       \overline{\phi\bigl(\mathbb D\cap
                                  B(\zeta,\epsilon)\bigr)}
                       \quad\text{in }S^2 .
                                                               \tag{C.73.9}
\]
This is the usual impression: the disk
neighborhoods and the crosscut neighborhoods
are cofinal bases in $O^{\rm P}$.
The sets in (C.73.9) are nested compact
connected sets, and their intersection is
nonempty and contained in $\partial O$.

A non-locally-connected frontier has a
nontrivial impression. Here is also a direct
verification of this consequence. If every
$I(\zeta)$ were a singleton, assign that value
to $\phi(\zeta)$. Any interior sequence tending
to $\zeta$ would have all its spherical cluster
points in that singleton, so would converge
to it. For a sequence of boundary points
$\zeta_n\to\zeta$, approximate each assigned
value by an interior value at a point within
$1/n$ of $\zeta_n$. The same interior-sequence
argument then proves continuity on the circle
and on the whole closed disk. Every frontier
point is a value on the circle, by compactness
applied to an interior approximating sequence.
The frontier would therefore be a continuous
image of the compact locally connected circle,
hence locally connected, by
\cite[Lemma~3.2]{C_BinderRojasYampolsky2014}.
This contradiction proves the assertion, and
also keeps the spherical topology explicit.

Fix distinct $a,b\in I(\zeta)$.
Choose successively positive numbers
$\epsilon_n\to0$ and points
$u_n,v_n\in\mathbb D\cap B(\zeta,\epsilon_n)$
such that
\[
                    \phi(u_n)\longrightarrow a,\qquad
                    \phi(v_n)\longrightarrow b .
                                                               \tag{C.73.10}
\]
Their existence follows from (C.73.9).
The lens $\mathbb D\cap B(\zeta,\epsilon_n)$
is convex and open, so a small closed
neighborhood of the segment $[u_n,v_n]$
is a convex topological disk $D_n$ contained
in that lens, with both points in its interior.
After choosing $D_n$, choose the next
$\epsilon_{n+1}$ smaller than its distance
from $\zeta$ and the distances of all earlier
disks from $\zeta$. This retains (C.73.10)
and makes the $D_n$ pairwise disjoint.
They form a locally finite family in
$\mathbb D$, with sole possible accumulation
point $\zeta$ in the closed disk, and
$\operatorname{diam}D_n\leq2\epsilon_n$.

Choose a homeomorphism $g_n$ supported in
$D_n$ taking $u_n$ to $v_n$. There is a
concrete such map: on each segment from
$u_n$ to a boundary point $z\in\partial D_n$,
send $u_n+t(z-u_n)$ to
$v_n+t(z-v_n)$, $0\leq t\leq1$.
Convexity and interiority give the unique
ray coordinates and the inverse formula
with $u_n,v_n$ interchanged. It fixes
$\partial D_n$ and extends by the identity
outside $D_n$. Therefore
\[
 \sup|g_n-\id|\leq2\epsilon_n,\qquad
 \sup|g_n^{-1}-\id|\leq2\epsilon_n .
                                                               \tag{C.73.11}
\]
Let $h_n\in K_O$ be its compactly supported
ambient conjugate. Then $\iota_O(h_n)=g_n$
tends to the identity in $\mathcal D_\partial$,
but $h_n$ does not tend to the identity:
it moves $\phi(u_n)$ to $\phi(v_n)$, whose
limiting points in (C.73.10) are distinct.
This proves the claimed failure of inverse
continuity on the actual image.

There is also one map witnessing nonsurjectivity.
Set $g=g_n$ on each $D_n$ and set $g=\id$
elsewhere on $\overline{\mathbb D}$.
Local finiteness gives a homeomorphism in the
interior, and the shrinking supports give
continuity at $\zeta$ and every other boundary
point. The same reasoning applies to the
piecewise inverses, so $g\in\mathcal D_\partial$.
If $g=\iota_O(h)$ for an $h\in K_O$, the
interior conjugacy would give
$h(\phi(u_n))=\phi(v_n)$.
Continuity at the fixed frontier point $a$
would then imply $a=b$, a contradiction.
Thus this specified $g$ is outside the image.
Together with density this proves the stated
proper-image assertion. For locally connected
frontier, C.68 gives the full opposite
comparison and its topology; no such
comparison is inferred here at a nontrivial
impression.

\emph{Original operators and the remaining filling problem.}
Let $H_O$ be any actual ambient operator group
preserving $O$. Denote its same-label
prime-end action in the chosen coordinate by
\[
       \widehat h=(\phi^{\rm P})^{-1}
                           h^{\rm P}\phi^{\rm P}
               \quad(h\in H_O).
                                                               \tag{C.73.12}
\]
An operator may reverse orientation. Its
conjugation still preserves $K_O$, and the
interior equations and density give exactly
\[
       \iota_O(hkh^{-1})
           =\widehat h\,\iota_O(k)\,\widehat h^{-1}
                       \quad(k\in K_O,\ h\in H_O).
                                                               \tag{C.73.13}
\]
Consequently
\[
 \iota_O\bigl(C_{K_O}(H_O)\bigr)
       =\iota_O(K_O)\cap
                 C_{\mathcal D_\partial}(\widehat H_O).
                                                               \tag{C.73.14}
\]
For the reverse implication, the disk equation
first gives the ambient commutation on $O$;
on its complement both relative maps are
the identity. This proves the original
same-label equation everywhere. An operator
acting trivially on $O$ induces the identity
prime-end map and already commutes with every
element of $K_O$, so the restriction kernel
is retained by (C.73.13).

The density proved above concerns
$\iota_O(K_O)$, not an unproved density
statement for the centralizer intersection.
Likewise a disk contraction need not stay
in this image, and (C.73.11) shows why
continuity of a putative inverse transfer
requires its own argument. These exact
comparisons leave the general equivariant
frontier-filling obligation of C.12 open.
\end{proof}

\begin{claim}{C.74}{Proved}{An exact obstruction to principal-coordinate interpolation}
There is an explicit based loop
$f_\tau\in\mathcal D_\partial$ whose generated group
$G$ is transitive on $\mathbb D$ and whose full
interior centralizer is a closed free proper group
$Q\cong\mathbb R$. The loop has an explicit based
contraction in $\mathcal D_\partial$ commuting with
this actual $Q$.

There are two specified principal coordinates for
the same $Q$-action, with their exact transition
map and inverse. In the second coordinate, write
the loop as
\[
        \widetilde f_\tau(s,t)
                  =(s+u_\tau(t),v_\tau(t)).
\]
The coordinate centralizer homotopy
\[
 (s,t)\longmapsto
       \bigl(s+r u_\tau(t),(1-r)t+r v_\tau(t)\bigr)
             \qquad(0\leq r\leq1)                       \tag{C.74.1}
\]
is a continuous family of interior homeomorphisms,
but its map at $r=1/2,\tau=1/2$ has no continuous
extension at any point of the original boundary
circle. Thus this literal coordinate interpolation
does not give the required relative contraction.

The same map admits no uniform bound by filled
orbits of any compact family of actual
boundary-fixed maps: for every compact
$P\subset\mathcal D_\partial$ and every positive
integer $N$, it fails the inclusion
\[
       A_*(x)\in\operatorname{Fill}\bigl(P^N(x)\bigr)
                         \quad\hbox{for all }x\in\mathbb D,
                                                               \tag{C.74.2}
\]
where $\operatorname{Fill}$ is the actual planar
filled hull of C.70 after the standard disk is
embedded in the plane. The comparison concerns
the stated interpolation and the original maps;
the displayed loop itself remains contractible.
\end{claim}
\passage{Section 23, folios 140-141, scans 279-281; Part C(b), exact actual-boundary comparison for one-dimensional proper operator groups}


\begin{proof}
\emph{The actual disk coordinate and three loop paths.}
Use the exact spiral coordinate of C.72, now
writing its real source variables as $(w,t)$:
\[
 \begin{split}
 p(z)&=\frac{z}{1+|z|}
                    \exp\bigl(i\log(1+|z|)\bigr),\\
 p^{-1}(\xi)&=\frac{\xi}{1-|\xi|}
                    \exp\bigl(i\log(1-|\xi|)\bigr),
                    \qquad z=w+it .
 \end{split}                                                   \tag{C.74.3}
\]
These are inverse homeomorphisms between the
real plane and $\mathbb D$ by their polar
formulas, including at zero.

Recall the precise compact-family boundary
control proved in C.72. If $|z'-z|\leq M$ and
$|z|>2M+1$, its direction, radial and logarithm
estimates give
\[
                  |p(z')-p(z)|\leq
                              \frac{4M}{|z|-M}.         \tag{C.74.4}
\]
Consequently every plane homeomorphism with
bounded displacement, and with bounded inverse
displacement, has a unique conjugate extension
fixing $S^1$ pointwise. A compact parameter
family with a common bound and jointly
continuous interior maps and inverses has a
jointly continuous extension with inverses.
The proof is the uniform estimate (C.74.4)
near the circle and the actual interior
conjugacy elsewhere.

Put $\psi(t)=\max\{0,1-|t|\}$ and define
\[
 \begin{split}
 P_{1,a}(w,t)&=(w+a,t),\\
 P_{2,a}(w,t)&=(w,t+a),\\
 P_{3,a}(w,t)&=(w+a\psi(t),t)
                         \qquad(0\leq a\leq1).
 \end{split}                                                   \tag{C.74.5}
\]
Their inverses have the same respective
formulas with $-a$. Their forward and inverse
displacements are at most one. Denote their
actual extensions through (C.74.3) by
$\overline P_{j,a}\in\mathcal D_\partial$.
For $\tau\in[(j-1)/3,j/3]$, set
\[
 \alpha_j(\tau)=
   \min\{6(\tau-(j-1)/3),\,2-6(\tau-(j-1)/3)\},
 \qquad
 f_\tau=\overline P_{j,\alpha_j(\tau)}.                \tag{C.74.6}
\]
Each amplitude is zero at the two endpoints
and one at the midpoint. The three pieces
therefore form one continuous based loop,
with its inverse continuous as well.

There is already an exact relative contraction:
replace each amplitude in (C.74.6) by
$r\alpha_j(\tau)$. The common displacement
bound one still applies. At $r=0$ this is
the identity loop and at $r=1$ it is the
original loop; all joins and the loop
basepoint stay the identity. All three
families commute with horizontal
translations, so this contraction retains
those actual operator equations.

\emph{Transitivity and the full actual interior centralizer.}
The values of the first path generate every
horizontal translation, using products and
inverses of parameters in $[0,1]$.
The values of the second generate every
vertical translation in the same way.
Thus $G$ contains all translations of the
$(w,t)$-plane and is transitive.

Any homeomorphism commuting with all plane
translations is itself a translation:
if its value at zero is $(a,b)$, commutation
forces its value at $(w,t)$ to be
$(w+a,t+b)$. Commutation with $P_{3,1}$
then says
\[
                         \psi(t+b)=\psi(t)
                                      \quad(t\in\mathbb R).
                                                               \tag{C.74.7}
\]
At $t=0$, this gives $\psi(b)=1$, whose unique
solution is $b=0$. Conversely every horizontal
translation commutes with all three families.
Consequently the full interior centralizer,
including possible orientation-reversing
candidates, is exactly
\[
 C_{\operatorname{Homeo}(\mathbb D)}(G|_{\mathbb D})
       =Q=\{p\,q_a\,p^{-1}:a\in\mathbb R\},
       \qquad q_a(w,t)=(w+a,t).                       \tag{C.74.8}
\]
There are no orientation-reversing elements
in this explicitly computed group.

The map $a\mapsto q_a$ is a topological
isomorphism onto its actual interior image:
evaluation at zero gives the inverse, and
the formulas give continuity of the maps
and inverses on every compact set.
This image is closed, since a homeomorphic
limit evaluated at zero has a finite real
parameter and hence is the corresponding
translation. The action is free. It is
proper because two compact source sets
have bounded first coordinates, which
bound every translation parameter taking
one set into the other.
Conjugation by $p$ and its displayed inverse
preserve these assertions. Thus $Q$ has
dimension exactly one and the topology
required in C.33. Its actual extensions to
the closed disk are themselves
boundary-fixed, by (C.74.4).

\emph{Two actual principal sections and their transition.}
The quotient by $Q$ is the real coordinate
$t$. A first section is $\sigma_0(t)=p(0,t)$,
with the principal coordinate
\[
               \Theta_0(w,t)=p(w,t).
\]
A second section of the same quotient is
$\sigma_1(t)=p(-t^3,t)$, and its principal
coordinate is
\[
 \begin{split}
 \Theta_1(s,t)&=p(s-t^3,t),\\
 H=\Theta_0^{-1}\Theta_1 &: (s,t)\longmapsto(s-t^3,t),\\
 H^{-1}&:(w,t)\longmapsto(w+t^3,t).
 \end{split}                                                   \tag{C.74.9}
\]
These are the actual coordinate maps and
their inverse identities. In both,
the same labelled element $p q_a p^{-1}$
acts by addition of $a$ to the first
coordinate. In particular the second
coordinate satisfies every C.33 principal
bundle requirement; no operator has
been changed by the section choice.

At the midpoint $\tau=1/2$ of the second
loop path, the original map is
$pP_{2,1}p^{-1}$. More generally its
vertical-translation path in the second
coordinate is
\[
 F_a(s,t)=
       \bigl(s+(t+a)^3-t^3,t+a\bigr)
             =H^{-1}P_{2,a}H(s,t).                   \tag{C.74.10}
\]
Its literal interpolation with the identity
in that coordinate is
\[
 M_{r,a}(s,t)=
       \bigl(s+r((t+a)^3-t^3),t+ra\bigr).
                                                               \tag{C.74.11}
\]
It is an interior homeomorphism commuting
with every horizontal translation. The
corresponding calculation in the first
coordinate is exact:
\[
 H M_{r,a}H^{-1}(w,t)
  =\bigl(w+3r(1-r)a^2t+r(1-r^2)a^3,\ t+ra\bigr).
                                                               \tag{C.74.12}
\]
In particular $r=1/2,a=1$ gives
\[
 A(w,t)=\left(w+\frac34t+\frac38,t+\frac12\right),
 \qquad
 A^{-1}(w,t)=\left(w-\frac34t,t-\frac12\right).
                                                               \tag{C.74.13}
\]
The actual candidate on the disk is
$A_*=pAp^{-1}$. Formula (C.74.13) is also
an explicit inverse calculation; it
establishes the interior assertion before
any boundary comparison.

\emph{Two sequences at each original boundary point.}
Fix $\zeta=e^{i\theta}\in S^1$ and take
sufficiently large positive integers $n$.
Define
\[
 \begin{split}
 R_n&=e^{2\pi n+\theta}-1,\qquad X_n=(R_n,0),\\
 S_n&=e^{2\pi n+\theta-\pi/2}-1,\qquad Y_n=(0,S_n).
 \end{split}                                                   \tag{C.74.14}
\]
The polar formula (C.74.3) gives
$p(X_n)\to\zeta$ and $p(Y_n)\to\zeta$.
For the horizontal sequence,
$A(X_n)-X_n=(3/8,1/2)$ is bounded, so
(C.74.4) gives $p(A(X_n))\to\zeta$.

For the vertical sequence there is the
exact factorization
\[
        A(Y_n)=\left(S_n+\frac12\right)
                                   \left(\frac34,1\right).
                                                               \tag{C.74.15}
\]
Its argument is $\pi/2-\arctan(3/4)$
and its norm is $\frac54(S_n+\frac12)$.
Thus the logarithmic radial change tends
to $\log(5/4)$, and
\[
       p(A(Y_n))\longrightarrow e^{i\delta}\zeta,
          \qquad
       \delta=\log(5/4)-\arctan(3/4).                  \tag{C.74.16}
\]
The two limits are distinct without any
numerical approximation. The integral
formulas give
$0<\log(5/4)<1/4$ and
$\arctan(3/4)> (3/4)/(1+(3/4)^2)=12/25$.
Consequently
\[
                   -\frac{\pi}{2}<\delta<
                         -\frac{23}{100}<0.           \tag{C.74.17}
\]
In particular $e^{i\delta}\ne1$.
The two sequences (C.74.14) prove that
$A_*$ has no continuous extension at
this specified $\zeta$. Since $\zeta$
was arbitrary, this holds at every
point of the original circle.

\emph{The exact obstruction to the proposed filled-word bound.}
Let $P\subset\mathcal D_\partial$ be any
compact family, and fix $N\geq1$.
The product family $P^N$ is compact.
Joint continuity of evaluation on this
compact family, and its pointwise-fixed
boundary, imply
\[
 \sup_{g\in P^N}|g(p(Y_n))-\zeta|\longrightarrow0.
                                                               \tag{C.74.18}
\]
Every planar filled hull is contained
in the Euclidean convex hull of its
compact set, as proved in C.70.
Therefore every point of
$\operatorname{Fill}(P^N(p(Y_n)))$
also tends uniformly to $\zeta$.
Its proposed member $A_*(p(Y_n))$
instead tends to $e^{i\delta}\zeta$.
This disproves (C.74.2), for these
actual maps and for every such
compact family and finite word length.
Extending each original boundary-fixed
map by the identity outside the closed
disk makes the same assertion in the
original plane homeomorphism group.

The original contraction (C.74.6)
with scaled amplitudes remains valid
and keeps the same operators.
The failed construction is exactly
the interpolation (C.74.1) in the
specified changed principal section.
A section-adapted filling or a
controlled selection still requires
its own boundary argument; this
example does not furnish a nontrivial
centralizer fundamental group.

The executable
\path{code/part_b/c74_principal_interpolation_certificate.py}
checks the coordinate and inverse
identities, the group equations,
the exact interpolation formula,
the two ray calculations and the
rational constants in (C.74.17).
The limit and compact-family
arguments above are topological
proofs, with no floating-point
verification used.
\end{proof}

\begin{claim}{C.75}{Proved}{The actual orbit strips for a one-dimensional proper operator action}
Let $O$ be a plane, including an actual open disk, and
let $Q\leq\operatorname{Homeo}(O)$ be closed, locally
compact, free and proper, with $\dim Q=1$.
Let $f_\tau$ be a continuous based loop in its
centralizer, and put $G=\langle f_\tau\rangle$.
Retain the actual principal coordinate of C.33,
with its inverse, in which
\[
 \begin{split}
 O&\cong\mathbb R_s\times\mathbb R_t,\\
 f_\tau(s,t)&=(s+u_\tau(t),\beta_\tau(t)),\\
 A&=\{t:\beta_\tau(t)=t\text{ for every }\tau\}.
 \end{split}                                                   \tag{C.75.1}
\]
The set $A$ is closed. On every component $I$ of
$\mathbb R\setminus A$, precisely one of the
following two descriptions holds.
Either $\mathbb R\times I$ is one open $G$-orbit,
or there is a continuous function $b_I:I\to\mathbb R$
such that all the orbits in that strip are exactly
\[
                   \{(c+b_I(t),t):t\in I\},
                                  \qquad c\in\mathbb R.        \tag{C.75.2}
\]
In the latter case the actual coordinate
$(s,t)\mapsto(s-b_I(t),t)$, with inverse
$(w,t)\mapsto(w+b_I(t),t)$, changes every restricted
generator to $(w,t)\mapsto(w,\beta_\tau(t))$.
The effective restricted action and its base image
are topologically isomorphic in their actual
compact-open and inverse topologies.

At each $a\in A$ the orbits on $\mathbb R\times\{a\}$
are individual points when every $u_\tau(a)$ vanishes,
and otherwise that entire horizontal line is one
orbit. This lists all orbits. In particular every
finite interior frontier line of an open orbit is
the actual proper line at an endpoint of its
base interval. There are no circular orbits.
The graph in (C.75.2) is proper in its indicated
strip; extension of $b_I$ to an endpoint is
not asserted.

There is a further exact comparison for a finite
endpoint $a\in A$ of a base interval $I$.
Write $g(s,t)=(s+u_g(t),\beta_g(t))$ for the
original element $g\in G$, and put
\[
 c_a(g)=u_g(a),\qquad
 G_t=\{g\in G:\beta_g(t)=t\}\quad(t\in I).
                                                               \tag{C.75.3}
\]
The character $c_a$ is continuous. There are
exactly three alternatives: $c_a=0$; or
$c_a\ne0$ and $c_a(G_t)=0$; or
$c_a(G_t)=\mathbb R$. In the second alternative
there is an actual homeomorphism
$b:I\to\mathbb R$ satisfying
\[
                 b(\beta_g(t))=b(t)+c_a(g)
                                  \qquad(g\in G,\ t\in I).
                                                               \tag{C.75.4}
\]
In the third alternative every real endpoint
displacement is realized by an actual source
element fixing the chosen base point $t$.
All component-operator equations and the kernel
of the comparison action used below are retained.
No continuous choice of these correcting source
elements, or frontier filling, is inferred.
\end{claim}
\passage{Section 23, folios 140-141, scans 279-281; Part C(b), exact one-dimensional proper-operator orbit and frontier reduction}


\begin{proof}
\emph{The principal coordinate and the original group words.}
C.33 gives a Lie group with identity component
$Q^0$, and its actual product principal bundle.
Because $O$ is a plane, neither factor in that
product can be a circle: a circle factor gives
a nontrivial fundamental-group factor, whereas
the plane has trivial fundamental group.
Thus $Q^0=\mathbb R$, its base is $\mathbb R$,
and an actual section gives the coordinate
and inverse in (C.75.1).
Every labelled element of $Q^0$ acts by
$(s,t)\mapsto(s+c,t)$.

Commutation with these translations forces
every $g\in G$ to have the stated form.
The base homeomorphism is increasing because
each generator is joined to the identity.
The inverse map is exactly
\[
 (s,t)\longmapsto
       \bigl(s-u_g(\beta_g^{-1}(t)),\beta_g^{-1}(t)\bigr).
                                                               \tag{C.75.5}
\]
For the composition convention $gh=g\circ h$,
\[
 \beta_{gh}=\beta_g\beta_h,\qquad
 u_{gh}(t)=u_h(t)+u_g(\beta_h(t)).
                                                               \tag{C.75.6}
\]
The functions and their inverse maps are jointly
continuous on every compact parameter family.

Put $P=\{f_\tau\}\cup\{f_\tau^{-1}\}$, retaining
the actual maps. It is compact, path connected,
and contains the identity. Every $P^n$ is compact
and path connected, and $\bigcup_nP^n=G$.
More explicitly an element
$g=f_{\tau_1}^{\epsilon_1}\cdots
f_{\tau_m}^{\epsilon_m}$ has the actual path
\[
 r\longmapsto
 f_{r\tau_1}^{\epsilon_1}\cdots
 f_{r\tau_m}^{\epsilon_m},\qquad 0\leq r\leq1,
                                                               \tag{C.75.7}
\]
from the identity to $g$, entirely in $P^m$.
No path in a larger closure is substituted.

The fixed set of all base generators is
exactly $A$, and it is also the fixed set
of their generated base group.
The base orbit of any $t\notin A$ is an
open interval. Indeed choose a based path
of base maps moving that point. That
path and its inverse give points on both
sides of it; their evaluation intervals
give a neighborhood in the orbit.
The same argument at each point of the
orbit makes the orbit open. It is an
interval because it is the increasing
union of the connected compact evaluation
images of $P^n$.
A finite endpoint of this orbit is fixed
by every increasing base map preserving
the orbit. Hence the orbit is exactly
the component $I$ of $\mathbb R\setminus A$
containing $t$.

\emph{The precise interval consequence of the word triangle.}
For a continuous path $h_r$ of actual
maps of the form in (C.75.1), starting at
the identity, suppose that $h_1(s,t)$
has second coordinate $t$.
The entire horizontal interval joining
$(s,t)$ to $h_1(s,t)$ is contained in
the orbit of $(s,t)$ under the group
generated by that path.
This is the exact C.62 word-triangle
argument with second-coordinate quotient.
Its hypotheses here can be checked directly:
the source is the actual product plane,
the real fibers are proper horizontal
lines, the induced base maps are
increasing, and the endpoint equality
is real equality.
No original-open-orbit hypothesis
enters that triangle proof.

More explicitly, subdivide the path by
the closed set of parameters at which
its second-coordinate evaluation equals
$t$. On an excursion $(a,b)$, put
$A_0=h_a(s,t)$ and $k_r=h_rh_a^{-1}$.
The paths $k_rA_0$ and
$k_bk_r^{-1}A_0$ have opposite strict
second-coordinate signs relative to
the original fiber. Their concatenation
has winding $1$ or $-1$ about every
interior point of the interval between
their two endpoints, as in C.62.9.
The triangle $(v,u)\mapsto k_uk_v^{-1}A_0$
therefore contains that interval.
Interpolating on the closed zero set
and these excursion gaps, with uniform
continuity on their compact parameter
set, gives the complete endpoint
interval, exactly as in C.62.
Every displayed value is an actual
three-word value from the path.
For (C.75.7) it consequently belongs
to $P^{3m}(s,t)$.

\emph{The two possible strip actions.}
For each $t\in I$, define the actual
additive subgroup
\[
       D_t=\{u_g(t):g\in G_t\}\leq\mathbb R.
                                                               \tag{C.75.8}
\]
Additivity and inverses follow from
(C.75.5--6). If $h$ takes the base point
$t$ to $t'$, conjugation by that actual
map takes a shift $c\in D_t$ to the
same shift $c$ in $D_{t'}$:
its first-coordinate translations
cancel. Thus $D_t=D_{t'}$ throughout
this base orbit.

If $D_t$ contains a nonzero $c$, take
its actual source element $g$ and the
path (C.75.7). The preceding triangle
puts every point of the interval from
$(0,t)$ to $(c,t)$ in the original orbit.
Each such orbit element fixes the
base point, so $D_t$ contains the
full interval from zero to $c$.
An additive subgroup containing a
nondegenerate interval is all of
$\mathbb R$. Transitivity on the
base interval, followed by these
arbitrary fiber translations, proves
that $\mathbb R\times I$ is one orbit.

Suppose instead that $D_t=0$.
Choose $t_0\in I$ and $x_0=(0,t_0)$.
Every fiber over $I$ meets $Gx_0$
exactly once: existence follows
from base transitivity, and two
such points would give a nonzero
element of $D_t$ after composing
their source words.
Write that point as $(b_I(t),t)$.

This function is continuous with
its actual, not a postulated,
graph topology. The compact set
$P^n(x_0)$ projects continuously
and bijectively to a compact
interval $I_n$. A compact-to-Hausdorff
bijection is a homeomorphism.
The intervals $I_n$ increase and
exhaust $I$. For each $t\in I$,
choose $t_-<t<t_+$ in $I$.
Both belong to some $I_n$, so
$[t_-,t_+]\subset I_n$.
Thus the inverse projection,
and hence $b_I$, is continuous
on a neighborhood of every $t$.

Every horizontal translation
commutes with $G$. Its translates
of $Gx_0$ are all the other
orbits, proving (C.75.2).
Invariance of this actual graph
under $g$ gives
\[
           u_g(t)=b_I(\beta_g(t))-b_I(t).
                                                               \tag{C.75.9}
\]
This verifies both conjugacy
directions of the asserted
coordinate change.
The graph is closed in
$\mathbb R\times I$, and its
projection is a homeomorphism
onto $I$, so it is proper
in that strip.

Under this coordinate, the
effective restricted group
is exactly $\{\,\id_{\mathbb R}
\times\beta:\beta\in\beta(G)|_I\,\}$.
Projection to the base has
the explicit continuous inverse
$\beta\mapsto\id\times\beta$,
also on inverse maps, in the
compact-open topologies.
Evaluation on the section
$\{0\}\times I$ proves
continuity in the other direction.
This proves the stated topology
of the effective images. It
makes no quotient-topology
claim for restriction from
the original global group.

For $a\in A$, the first-coordinate
shifts form the image of the
path-connected group $G$ under
the continuous evaluation
$g\mapsto u_g(a)$.
This is a connected additive
subgroup of $\mathbb R$, hence
zero or $\mathbb R$.
The zero alternative is exactly
the vanishing of every generator
$u_\tau(a)$. This proves the
remaining orbit description.
The frontier inside the product
plane of an open strip is
exactly the horizontal lines
at its finite endpoints.

\emph{The endpoint character with all source labels retained.}
Fix the indicated finite
endpoint $a$ and interval $I$.
By (C.75.6), $c_a$ is an
additive character, continuous
in the original group topology.
Form the labelled comparison
action on a second product:
\[
 \begin{split}
 T_a:G&\longrightarrow
                 \operatorname{Homeo}(\mathbb R_z\times I),\\
 T_a(g)(z,t)&=(z+c_a(g),\beta_g(t)),\\
 \ker T_a&=\{g:c_a(g)=0,\ \beta_g|_I=\id\}.
 \end{split}                                                   \tag{C.75.10}
\]
The displayed kernel is kept;
no faithfulness or closedness
of this image is assumed.
Every finite word and every
path (C.75.7) has precisely
its displayed image under $T_a$.

The strip proof just given
therefore applies to this
actual image action. Its
fiber stabilizer at $t$ is
exactly $c_a(G_t)$, so it is
zero or all of $\mathbb R$.
In the latter case the same
triangle supplies actual
source words with every
intermediate real displacement,
and their finite products
give every real displacement.
They fix this specified base
point; a continuous choice
of these words has not been
constructed.

In the zero case the actual
graph function from that
argument satisfies
\[
           b(\beta_g(t))=b(t)+c_a(g).
                                                               \tag{C.75.11}
\]
Suppose that $c_a$ is nonzero.
Its path-connected additive
image is then $\mathbb R$,
so $b(I)=\mathbb R$ by base
transitivity and (C.75.11).
We check that $b$ is injective.
If two distinct arguments
had the same value, either
$b$ would be constant on
the interval between them,
or it would have an interior
local maximum or minimum
there.
An open interval of constancy
is transported by the
transitive base action to
neighborhoods of every point,
which would make $b$ constant.
A local maximum is similarly
transported to every point.
A continuous real function
having a local maximum at
every point has countable
image: each of its values
is the maximum on one
rational closed interval
inside a local maximum
neighborhood. Its connected
image must therefore be
a singleton. The identical
argument with minima handles
the other case.
All alternatives contradict
$b(I)=\mathbb R$.
Thus $b$ is a continuous
injective surjection from
an interval to $\mathbb R$,
hence a strictly monotone
homeomorphism. This proves
(C.75.4) with its exact
inverse coordinate.

\emph{Component operators and the remaining boundary step.}
C.33 retains either no
nontrivial component operator
or one labelled generator
\[
       q(s,t)=(\lambda s+\kappa,t+d),
                    \qquad d>0,\quad\lambda\ne0 .
                                                               \tag{C.75.12}
\]
The original commutation
equations, with the constant
$\kappa$ retained, are
\[
 \beta_g(t+d)=\beta_g(t)+d,\qquad
 u_g(t+d)=\lambda u_g(t).
                                                               \tag{C.75.13}
\]
They give $A+d=A$, carry
each strip and its orbit
description to the translated
strip, and send its endpoint
character to $\lambda c_a$.
The actual maps $q$ and
$q^{-1}$ prove both transport
directions. Negative $\lambda$
is included.

All finite interior frontier
strata in this reduction
are consequently the named
horizontal lines, with the
original translation
characters just computed.
An actual disk coordinate
can still have complicated
limits at its original
outer circle, as C.74
demonstrates. Neither the
graph coordinate nor the
endpoint correction
surjectivity supplies a
jointly continuous filling
at that circle or across
the nonfixed interior
frontier lines. Those
extension steps remain
part of C.12.
\end{proof}

\begin{claim}{C.76}{Proved}{Simultaneous endpoint prescriptions with compact original-boundary control}
Let $Q$ be the actual closed, locally compact, free proper
one-dimensional group on $O=\operatorname{int}\overline D$
of C.75. Let $f_\tau$ be a based loop of homeomorphisms of
the actual closed disk, fixing its boundary pointwise and
commuting with $Q$ on $O$. Use the actual principal
coordinate and its inverse from C.75, with coordinates
$(s,t)$, common base fixed set $A$, and generated group $G$.

Retain the following components $I$ of $\mathbb R\setminus A$:
those for which every finite-endpoint character $c_a$ is zero,
and those having a finite endpoint $a$ with
\[
       c_a\ne0,\qquad c_a(G_t)=0\quad(t\in I).
                                                               \tag{C.76.1}
\]
The first prescription includes an interval without finite
endpoints. Let $Z\subset\overline D$ consist of the original
boundary circle, the actual lines over $A$, and all retained
strips. This is a closed subset, preserved by the loop and
by every actual interior operator wherever that operator
is defined. The restricted loop has a based contraction
through homeomorphisms of $Z$, with jointly continuous
inverses, fixing the original circle and commuting with
every $q\in Q$ on $Z\cap O$.

On every line over $a\in A$, the contraction is exactly
\[
                  (s,a)\longmapsto(s+r u_\tau(a),a),
                           \qquad 0\leq r\leq1.                \tag{C.76.2}
\]
On a strip retained by (C.76.1), choose C.75's actual
homeomorphism $b:I\to\mathbb R$ for one such endpoint.
The base contraction is
\[
 B_{r,\tau}(t)=b^{-1}\bigl(b(t)+r c_a(f_\tau)\bigr).
                                                               \tag{C.76.3}
\]
Both finite endpoint values, when there are two, are
prescribed simultaneously by (C.76.2). All original
component operators, including negative fiber multipliers,
are retained. One compact family of original closed-disk
homeomorphisms controls the forward and inverse evaluations
over all retained strips and all parameters. Hence the
construction includes infinitely many strips accumulating
at an interior line or at the original circle.

More precisely, let $\mathscr R_Z$ be the subgroup of the
original relative $Q$-centralizer preserving $Z$, and let
\[
 \begin{split}
 \mathscr A_Z={}&\{k\in\operatorname{Homeo}(Z):
          k|_{\partial\overline D}=\id,
          \ kq=qk\text{ on }Z\cap O\},\\
 \rho_Z:\mathscr R_Z&\longrightarrow\mathscr A_Z,
                         \qquad \rho_Z(h)=h|_Z.
 \end{split}                                                   \tag{C.76.4}
\]
These are the actual compact-open and inverse topologies,
and $\rho_Z$ is a continuous homomorphism. The contraction
just asserted fills the actual loop $\rho_Z(f_\tau)$ in
$\mathscr A_Z$. When every complementary interval is
retained, $Z=\overline D$ and this is a contraction in the
original relative centralizer itself. No lifting of the
constructed filling through $\rho_Z$ is asserted for a
proper $Z$.
\end{claim}
\passage{Section 23, folios 140-141, scans 279-281; Part C(b), actual endpoint-prescribed and original-boundary-controlled continuation of the dimension-one proper-operator reduction}


\begin{proof}
\emph{The endpoint image and its actual short increments.}
For any one strip with $k\in\{1,2\}$ finite endpoints,
write $c=(c_{a_1},\ldots,c_{a_k})$ and $V=c(G)\subset\mathbb R^k$.
This is the image of the actual group $G$, before taking
any closure. It is a real vector subspace. Indeed, every
$g\in G$ has the finite-word path (C.75.7) from the identity.
Its character image is a path from $0$ to $c(g)$ in $V$.
The planar three-word statement (C.37.3), using a zero
second coordinate for $k=1$, places every real fraction
$r c(g)$, $0\leq r\leq1$, in $V$. Addition and negation
then give all real multiples. This proves the assertion
and, in particular, that $V$ is closed in $\mathbb R^k$.
It is the real span of $\{c(f_\tau):0\leq\tau\leq1\}$:
one inclusion follows from generation, and the other from
the vector-space property just proved.

For $\delta>0$, retain the actual compact family
\[
 \begin{split}
 J_\delta&=\{f_s f_t^{-1}:s,t\in[0,1],\ |s-t|\leq\delta\},\\
 E_\delta&=J_\delta^6,
 \qquad D_\delta=c(J_\delta).
 \end{split}                                                   \tag{C.76.5}
\]
Both parameter families are compact and path connected,
contain the identity, and include all smaller choices of
$\delta$. Their global maps and inverses tend uniformly
to the identity as $\delta\downarrow0$, in the actual
closed-disk homeomorphism topology. The set $D_\delta$
is compact, symmetric and path connected. It spans $V$:
divide $[0,\tau]$ into intervals of length at most $\delta$
and telescope $c(f_\tau)$ into their character increments.

Choose a basis $v_1,\ldots,v_d$ of $V$ in $D_\delta$,
where $d\leq2$. For each $i$ a path in $D_\delta$ joins
$0$ to $v_i$. Again (C.37.3) gives
$[-v_i,v_i]\subset D_\delta+D_\delta+D_\delta$;
symmetry supplies the negative fractions. Consequently
\[
 \Bigl\{\sum_{i=1}^d a_i v_i:|a_i|\leq1\Bigr\}
       \subset c(E_\delta).
                                                               \tag{C.76.6}
\]
For $d<2$ use identity factors to reach six increments.
For $d=0$ the assertion is the identity image.
Thus $c(E_\delta)$ contains a neighborhood of $0$ in the
actual vector image $V$, for every $\delta>0$. This is a
statement about specified compact sets of actual source
words. It does not choose a continuous inverse to $c$.

Put $P=\{f_\tau\}\cup\{f_\tau^{-1}\}$. The same three-word
statement, now on the original prefix character path,
gives the uniform inclusion
\[
               r c(f_\tau)\in c(P^3)
                           \quad(0\leq r,\tau\leq1).
                                                               \tag{C.76.7}
\]
The single family $P^3$ works for every strip and both
endpoint characters, including varying ranks of $V$.

\emph{One compact family in the original relative group.}
Take the closure $L$ of $G$ in the group of actual
closed-disk homeomorphisms fixing the circle and commuting
with $Q$ on the interior. This is a Polish group: the
ambient compact-open homeomorphism topology controls maps
and inverses, and the stated pointwise equations are closed.
Every member of $L$ still has the principal form
\[
                  k(s,t)=(s+u_k(t),\beta_k(t)),
                                                               \tag{C.76.8}
\]
with $\beta_k$ increasing, fixing $A$ pointwise and
preserving every complementary interval. These are closed
conditions for maps and inverses. Each finite-endpoint
evaluation $c_a(k)=u_k(a)$ extends continuously to $L$.
For any particular endpoint vector its image on $L$ is
still $V$: density gives inclusion in the closed space
$V$, and the original image already equals $V$.
The translation equation (C.75.4) also passes to $L$ on
each retained strip by pointwise continuity. No assertion
about a quotient topology on $c(G)$ is used.

Choose a complete compatible metric on $L$. Exactly as
in the compact-product argument of C.63, choose
$\delta_n>0$ so that left multiplication by any member of
$E_{\delta_n}$ changes every previously formed product
by less than $2^{-n}$. This is possible by compactness
of the previous product family and the uniform identity
limit in (C.76.5). The families
\[
       K_n=E_{\delta_n}\cdots E_{\delta_1}P^3,
                       \qquad K_0=P^3,
                                                               \tag{C.76.9}
\]
are nested, and their infinite products converge uniformly
over the entire product of their compact parameter spaces.
Completeness gives a compact limit family $K\subset L$
containing every $K_n$. It is path connected: parameterize
$P$ by concatenating the original loop and its inverse,
and use the convex time bands in (C.76.5). The product
of these parameter spaces is path connected, and its
uniform product map into $L$ is continuous.
This includes convergence of the actual inverse maps.

Every member of $K$ fixes the original circle and commutes
with every labelled $q\in Q$. It preserves all principal
strips, whether retained or not. On a given strip the
functions $u_k(t),\beta_k(t)$ in (C.76.8) are jointly
continuous in its compact label and its real coordinate.

\emph{Base prescriptions and nonempty fiber intervals.}
On a strip with all finite-endpoint characters zero, use
C.27's based median contraction $B_{r,\tau}$ of the
original base loop. Its value is between $t$ and
$\beta_\tau(t)$, and hence is attained by an original
prefix value. Define
\[
 A_n(r,\tau,t)=\operatorname{conv}
       \{u_k(t):k\in K_n,\ \beta_k(t)=B_{r,\tau}(t)\}
                                                               \tag{C.76.10}
\]
on this strip. These are nonempty compact intervals.

On a strip satisfying (C.76.1), choose one of its eligible
endpoints $a$ and its actual coordinate $b$. Use (C.76.3),
and include \emph{all} finite endpoints in the vector $c$.
Define instead
\[
 A_n(r,\tau,t)=\operatorname{conv}
       \{u_k(t):k\in K_n,\ c(k)=r c(f_\tau)\}.
                                                               \tag{C.76.11}
\]
Equation (C.76.7) proves nonemptiness. For every label
occurring in (C.76.11), the $a$ coordinate and the exact
translation relation give
\[
                       \beta_k(t)=B_{r,\tau}(t).
                                                               \tag{C.76.12}
\]
Thus (C.76.11) really is an interval of displacements
in one specified principal fiber. Every one of its
labels has the prescribed displacement at \emph{both}
finite endpoints. The base map (C.76.3) lies between
$t$ and $\beta_\tau(t)$ even when $b$ is decreasing.
It is an increasing homeomorphism of $I$ with jointly
continuous inverse and all four based parameter faces.

In both definitions use $K$ in place of $K_n$ to obtain
a nonempty compact interval $A_\infty$. The four desired
fiber prescriptions are
\[
 \begin{split}
 U(0,\tau,t)&=0,\qquad U(1,\tau,t)=u_\tau(t),\\
 U(r,0,t)&=U(r,1,t)=0.
 \end{split}                                                   \tag{C.76.13}
\]
They belong to every $A_n$ on the respective faces.

\emph{The exact lower-limit comparison and relative selection.}
For each of (C.76.10) and (C.76.11) one has
\[
 A_n(p)\subset\liminf_{p'\to p}A_{n+1}(p'),
                     \qquad p=(r,\tau,t).
                                                               \tag{C.76.14}
\]
For (C.76.10), the short-increment argument of C.63 applies
directly. The base group is transitive on $I$ by C.75.
At every base point, arbitrarily short original increments
and their inverses move that point to both sides. The
inverse-evaluation argument (C.63.9) proves this assertion
for the actual specified point. These increments occur
in $E_{\delta_{n+1}}$, with five identity factors added.
The intermediate value theorem gives a correction of the
nearby target level, as close to the identity as desired.

For (C.76.11), take a label $k\in K_n$ with
$c(k)=r c(f_\tau)$. At a nearby parameter the required
character correction is the actual vector
\[
        v(p')=r'c(f_{\tau'})-c(k)\in V,
                              \qquad v(p')\longrightarrow0.
                                                               \tag{C.76.15}
\]
Choose any arbitrarily small $\delta'\leq\delta_{n+1}$.
By (C.76.6), every sufficiently nearby $p'$ has an actual
$e_{p'}\in E_{\delta'}$ with $c(e_{p'})=v(p')$.
Then $e_{p'}k\in K_{n+1}$ has precisely the new required
character vector. Equation (C.76.12) supplies its base
level at $t'$ automatically. As $\delta'\downarrow0$,
its fiber displacement at $t'$ tends uniformly to that
of $k$ near the chosen source point. This proves the
lower-limit inclusion for every unconvexified value.
Applying it to two endpoints and retaining their convex
coefficient proves (C.76.14) for the whole interval.
No continuous choice of $e_{p'}$ was used.

Replace the allowed sets by the prescribed singletons
(C.76.13) on the closed parameter boundary. The
lower-limit property survives by the same two comparisons
as in C.63: off that boundary it is (C.76.14), and along
the boundary it is continuity of the prescribed values.
Here are the precise selection ingredients. Choose at
each parameter a value of the $n$th interval, extend it
locally as a constant section, and use (C.76.14) to make
it arbitrarily close to the next interval. A locally
finite partition of unity combines these approximations,
since every fiber set is convex. Starting with error
$2^{-1}$, inductively choose continuous approximations
$U_n$ with distance less than $2^{-n}$ from the $n$th
modified interval and with
$|U_{n+1}-U_n|<2^{1-n}$. Their locally uniform limit is
continuous, belongs to the closed compact interval
$A_\infty$ at every parameter, and equals each prescribed
singleton exactly. This is C.63's relative selection
construction with the lower-limit hypothesis now proved
for the simultaneous character constraints.

\emph{Every original component operator.}
If $Q$ is connected, all its operators are the original
fiber translations and the selected maps commute with
them. Otherwise retain C.33's actual component generator
\[
                 q(s,t)=(\lambda s+\kappa,t+d),
                         \quad d>0,\quad\lambda\ne0.
                                                               \tag{C.76.16}
\]
All members of $L$, and hence all labels in $K$, obey
\[
       u_k(t+d)=\lambda u_k(t),\qquad
       \beta_k(t+d)=\beta_k(t)+d.
                                                               \tag{C.76.17}
\]
Consequently $A+d=A$, and the retained intervals are
permuted with all their exact endpoint criteria: the
character at $a+d$ is $\lambda c_a$.
For a proper interval $I\ne\mathbb R$, no nonzero power
of the base translation preserves $I$. Choose the
construction on one representative of each interval
orbit and transport it by the actual powers of $q$.
For the endpoint coordinate the transported choice is
$b'(t+d)=\lambda b(t)$, with any additive constant
permitted; (C.76.3) then transports exactly. The fiber
displacement transports by $U'(r,\tau,t+d)=\lambda U(r,\tau,t)$.
An operator stabilizing that proper interval belongs
to $Q^0$ and is an actual fiber translation, so this
transport is independent of its representative.

The only interval which can stabilize under a nonzero
component power is $I=\mathbb R$, when $A$ is empty.
It is retained by the zero-endpoint rule. Perform its
relative selection in the line bundle
$(t,z)\sim(t+d,\lambda z)$ exactly as in (C.63.20).
Its continuous fiber norm is $|\lambda|^{-t/d}|z|$;
the preceding approximation bounds are taken in this
norm. The convex fiber sets and their lower-limit
property are well-defined there by (C.76.17), and the
median commutes with the actual base translation.
This covers the negative values of $\lambda$ as well.
The constant $\kappa$ cancels in the commutation equation
and has not been removed from the operator label.

\emph{The homeomorphisms and their compact evaluation bounds.}
On each retained strip put
\[
 \begin{split}
 K_{r,\tau}(s,t)&=(s+U(r,\tau,t),B_{r,\tau}(t)),\\
 K_{r,\tau}^{-1}(s,v)&=
 \bigl(s-U(r,\tau,B_{r,\tau}^{-1}(v)),B_{r,\tau}^{-1}(v)\bigr).
 \end{split}
                                                               \tag{C.76.18}
\]
These are exact inverse homeomorphisms, jointly continuous
on each strip. Use (C.76.2) on lines over $A$ and the
identity on the original circle. The formulas preserve
each piece of $Z$ and have the four required based faces.
Their operator commutation follows from (C.76.16)--(C.76.17)
and the transported construction.

At every source point of a retained strip the selected
value lies between two $K$-values in one real fiber,
by (C.76.10)--(C.76.12). Connect their labels by a path
in the compact parameter space for $K$. The C.62 word
triangle, in precisely the form (C.63.14), gives
\[
       K_{r,\tau}(x)\in K K^{-1}K(x),\qquad
       K_{r,\tau}^{-1}(x)\in K^{-1}K K^{-1}(x).
                                                               \tag{C.76.19}
\]
For the inverse inclusion, the two increasing fiber
maps bracket the inverse point by their own inverses
in their common input fiber, as in (C.63.26).
On a line over $a\in A$, the real path
$\sigma\mapsto u_\sigma(a)$ joins $0$ to $u_\tau(a)$
on the prefix $[0,\tau]$. It attains $r u_\tau(a)$.
Thus the forward value in (C.76.2) is an original
prefix value, and its inverse is an original inverse
prefix value. Since $P\subset K$, the two bounds
(C.76.19) hold there as well. The bounds persist under
all component transports because each actual member
of $K$ commutes with $q$ on $O$.

\emph{All finite endpoint limits, including varying strips.}
Suppose $t_n\to a\in A$, with $t_n$ in retained strips
$I_n$, and allow both homotopy parameters to vary.
Choose a finite endpoint $a_n$ of $I_n$ lying between
$t_n$ and $a$. Then $a_n\to a$. For a zero-character
strip every $k\in K$ satisfies $u_k(a_n)=0$.
For a character-prescribed strip every label used
in its defining interval satisfies
\[
                     u_k(a_n)=r_n u_{\tau_n}(a_n),
                                                               \tag{C.76.20}
\]
because the vector constraint includes both finite
endpoints. This same equality holds for the zero-character
strip, whose right-hand side is zero.
Compactness of $K$ in the actual group makes
$u_k(t_n)-u_k(a_n)\to0$ uniformly in $k$: the principal
coordinate and its inverse are continuous on a compact
neighborhood of the finite line point under consideration.
Convexification preserves this estimate. It follows that
\[
         U(r_n,\tau_n,t_n)-r_n u_{\tau_n}(a_n)
                              \longrightarrow0.
                                                               \tag{C.76.21}
\]
The original compact loop makes its latter endpoint
values jointly continuous. Both base prescriptions
are between $t$ and $\beta_\tau(t)$, so their base
values tend to $a$ uniformly in the parameters.
Equations (C.76.20)--(C.76.21) give exactly (C.76.2)
at the limiting line, also when $a$ is not itself an
endpoint of one fixed interval.

For inverses, a base map lying pointwise between the
identity and $\beta_\tau$ has inverse lying pointwise
between the identity and $\beta_\tau^{-1}$. This follows
by applying the two increasing maps to the inequalities
at its inverse argument. Hence an inverse base argument
approaching $a$ also tends to $a$ uniformly. Applying
(C.76.21) at that argument proves the inverse limit
with displacement $-r u_\tau(a)$. Along $A$ itself,
both limits follow directly from (C.76.2).

\emph{The original circle and the exact restriction map.}
The two word families in (C.76.19) are compact families
of the original closed-disk maps, and every member
fixes the original circle pointwise. For $x_n\to\zeta$
on that circle, their evaluation values at $x_n$ tend
uniformly to $\zeta$. Thus (C.76.19) proves joint
continuity of both the constructed maps and their
inverses at every original boundary point, with no
restriction on the escaping principal coordinates
or on the varying strip index. The circle is fixed
throughout. The set $Z$ is closed because its complement
in the disk is exactly the union of the excluded open
strips. The preceding bijections and their continuous
inverses are therefore actual homeomorphisms of this
specified compact space, and the compact-parameter
evaluation criterion gives continuity into its actual
homeomorphism group.

Restriction to a closed invariant subset is continuous
for maps and inverses, so (C.76.4) is the asserted
continuous group homomorphism. Its source contains the
original loop. The formulas construct the claimed
filling of that loop's restriction; no assertion that
every intermediate map extends through an excluded
strip is made. The excluded intervals have a nonzero
finite-endpoint character but none satisfying (C.76.1).
C.75 identifies their nonzero endpoint-stabilizer images
as the full real line. Choosing a continuous,
boundary-controlled lift there remains a separate step
in C.12. In particular no source-word surjectivity alone
has been substituted for that lift.
\end{proof}

\begin{claim}{C.77}{Proved}{All abelian base strips and the exact remaining character action}
Keep the actual closed disk, proper one-dimensional
operator group $Q$, loop, principal coordinate, and
common base fixed set $A$ of C.76. For a component
$I$ of $\mathbb R\setminus A$, let
\[
        B_I=\{\beta_g|_I:g\in G\}
                    \leq\operatorname{Homeo}^+(I)
                                                               \tag{C.77.1}
\]
be its actual effective base image, with its induced
compact-open and inverse topology. Retain every strip
whose finite-endpoint characters are all zero, and
every strip for which $B_I$ is abelian. Let $Z_{\mathrm{ab}}$
be their union with all lines over $A$ and the original
circle. The restriction of the original loop to this
closed set has a based contraction through its actual
homeomorphisms, fixing the original circle and commuting
with $Q$ on its interior part. The inverse family is
jointly continuous. One compact family of original
closed-disk maps controls all evaluations, including
arbitrary accumulation of strips.

For an abelian base image there are an actual increasing
homeomorphism $b_I:I\to\mathbb R$ and a continuous
source character $d_I:G\to\mathbb R$ such that
\[
 b_I(\beta_g(t))=b_I(t)+d_I(g),\qquad
                         d_I(G)=\mathbb R.
                                                               \tag{C.77.2}
\]
The prescribed base contraction is
$b_I^{-1}(b_I(t)+r d_I(f_\tau))$, and each finite endpoint
has the simultaneous displacement $r u_\tau(a)$.
This includes abelian base strips with full-real
endpoint-stabilizer character images. The exact
restriction homomorphism is the map (C.76.4) with
$Z_{\mathrm{ab}}$ in place of $Z$; no lifting through
excluded strips is asserted. The set $Z$ of C.76 is
contained in $Z_{\mathrm{ab}}$ by the actual translation
comparison (C.75.4).

There is also an exact description of the remaining
character comparison. On a strip with nonabelian $B_I$,
write $c:G\to V\subset\mathbb R^k$ for the vector of its
finite-endpoint characters, $0\leq k\leq2$. Then
\[
 c(G_t)=V\quad(t\in I),
 \qquad
 T(g)(z,t)=(z+c(g),\beta_g(t))
                                                               \tag{C.77.3}
\]
defines a transitive action on the actual space $V\times I$.
Let $L$ be the actual closed relative group used in C.76,
and let $\mathscr P(L)$ be its Polish based path group.
For each $t\in I$, endpoint evaluation gives an open
surjection and the indicated homeomorphism
\[
 \begin{aligned}
 e_t:\mathscr P(L)&\longrightarrow V\times I,\\
 e_t(\sigma)&=(c(\sigma(1)),\beta_{\sigma(1)}(t)),\\
 \mathscr P(L)/J_t&\xrightarrow{\ \cong\ }V\times I,\\
 \sigma J_t&\longmapsto e_t(\sigma).
 \end{aligned}                                                 \tag{C.77.4}
\]
where $J_t=e_t^{-1}(0,t)$ is the actual stabilizer.
Every sufficiently small compact target set has a compact
preimage in any specified open path neighborhood whose
image contains it. No continuous section, nor uniform
compact preimage for targets escaping towards a strip
endpoint, is inferred.
\end{claim}
\passage{Section 23, folios 140-141, scans 279-281; Part C(b), enlargement of the actual one-dimensional proper-operator filling and exact remaining character/base quotient}


\begin{proof}
\emph{The actual topology of an abelian base image.}
C.75 proves that $B_I$ is transitive on $I$. Take its
closure $\overline B_I$ in the actual homeomorphism group
of $I$. It is abelian and still transitive. Its action
is free: an element fixing one point commutes with the
transitive original image and consequently fixes every
point. The verified Polish orbit-open-map theorem used
in C.30 identifies evaluation
$\overline B_I\to I$ homeomorphically, since its
stabilizer is trivial. The closure is locally compact.
The one-dimensional group comparison in (C.37.9) now
identifies its actual action, through this evaluation
homeomorphism, with the translation action of $\mathbb R$.
This is exactly C.37's transitive abelian line argument,
including its faithful action and topology hypotheses.
Because $B_I$ was already transitive, freeness of the
closure shows $B_I=\overline B_I$: the original element
with a specified evaluation value is the only element
of the closure with that value.

Choose the orientation of its real translation coordinate
so that $b_I$ is increasing. Then
\[
              d_I(g)=b_I(\beta_g(t_0))-b_I(t_0)
                                                               \tag{C.77.5}
\]
for any fixed $t_0\in I$ gives (C.77.2).
It is a character, is independent of the chosen point
$t_0$, and is continuous on the actual source group
by the displayed evaluation formula. Its kernel is
exactly $\{g:\beta_g|_I=\id\}$. Its image is all of
$\mathbb R$, by transitivity. Thus the continuous source
map and the topological isomorphism of its effective
base image have both been specified; a quotient-topology
claim for the original source restriction is unnecessary.

\emph{Three actual characters and their compact representatives.}
For an abelian base strip put
\[
       a_I(g)=(d_I(g),c_{a_1}(g),\ldots,c_{a_k}(g))
                    \in\mathbb R^{k+1},\qquad k\leq2,
                                                               \tag{C.77.6}
\]
including every finite endpoint, and write $V_I=a_I(G)$.
The three-dimensional additive path identity (C.50.3)
applies after adding zero coordinates when needed.
For every actual finite-word path from $1$ to $g$, it
places each real fraction of $a_I(g)$ in $V_I$ using
at most five signed path values. Hence $V_I$ is a real
vector subspace of dimension at most three, before
passing to any closure. It is the real span of the
original character path and is closed in this Euclidean
space. With $P=\{f_\tau\}\cup\{f_\tau^{-1}\}$ one has
\[
             r a_I(f_\tau)\in a_I(P^5)
                               \quad(0\leq r,\tau\leq1).
                                                               \tag{C.77.7}
\]

Retain the actual short-increment family $J_\delta$
of (C.76.5), and now put $E_\delta=J_\delta^{15}$.
Telescoping the original character path proves that
$a_I(J_\delta)$ spans $V_I$ for every $\delta>0$.
It is compact, symmetric and path connected.
Choose a basis of at most three vectors in this set.
The identity (C.50.3) puts the full segment from minus
to plus each basis vector in five signed values from
that set; symmetry absorbs the signs. Adding the at
most three segments therefore proves
\[
      a_I(E_\delta)\text{ contains a neighborhood of }0
                         \text{ in the actual }V_I.
                                                               \tag{C.77.8}
\]
Identity factors pad shorter products. The numbers
five and fifteen do not depend on the strip, the
dimension of its image, or the parameter values.
All original maps and inverses in $E_\delta$ tend
uniformly to the identity as $\delta\downarrow0$.

The character $d_I$ extends continuously to $L$ by
(C.77.5). The full translation equation is a closed
pointwise relation and still holds there. Each endpoint
character extends by its actual line evaluation.
Density and closedness of $V_I$ show $a_I(L)=V_I$,
with no replacement of the original image by a larger
vector space. Construct the one compact path-connected
global family $K$ from
\[
           K_0=P^5,\qquad
           K_n=E_{\delta_n}\cdots E_{\delta_1}P^5
                                                               \tag{C.77.9}
\]
as in (C.76.9). Complete compatible metrics, summable
left-product errors, compact parameter products, and
inverse control apply verbatim to these actual finite
families. Every element of $K$ is a homeomorphism of
the original closed disk fixing its circle and commuting
with the original operators.

\emph{The selection, every endpoint, and the original circle.}
On an abelian base strip use the base prescription
of the claim. Define the allowed fiber intervals by
\[
 A_n(r,\tau,t)=\operatorname{conv}
 \{u_k(t): k\in K_n,\ a_I(k)=r a_I(f_\tau)\}.
                                                               \tag{C.77.10}
\]
Equation (C.77.7) proves nonemptiness. Every label in
this set has exactly the prescribed base map by its
$d_I$ coordinate, and exactly the prescribed displacement
at every finite endpoint by the other coordinates.
For nearby parameters, the needed correction is
$r'a_I(f_{\tau'})-a_I(k)\in V_I$.
Equation (C.77.8), with an arbitrarily smaller
$\delta'\leq\delta_{n+1}$, realizes that correction
by an actual arbitrarily globally small increment.
Thus the complete lower-limit proof (C.76.14)--(C.76.15)
holds with $a_I$ in place of $c$. It uses no selected
continuous inverse of $a_I$.

On a retained nonabelian strip, all finite-endpoint
characters are zero by its definition. Use the median
base prescription and (C.76.10). The original single
short increments still lie in the new fifteenfold
family, so its lower-limit proof is unchanged.
The four parameter-face values (C.76.13) belong to
every defining interval. The relative partition-of-unity
approximation and summable limiting selection of C.76
therefore give a continuous fiber displacement $U$.
It lies in the compact limiting interval and has those
four faces exactly. Use the same explicit homeomorphism
and inverse formulas (C.76.18).

Every actual component operator still has the form
$q(s,t)=(\lambda s+\kappa,t+d)$ of C.75.
Transport the chosen abelian base coordinate by
$b_{I+d}(t+d)=b_I(t)$. This makes the transported
base character equal to $d_I$, while each transported
endpoint character is $\lambda c_a$.
Equations (C.77.2) and (C.77.10) therefore commute
with the actual $q$, including negative $\lambda$.
For proper intervals choose orbit representatives and
transport the fiber selections by multiplication by
$\lambda$, as in C.76. No nonzero component power
stabilizes a proper interval. If $I=\mathbb R$ and
the base image is abelian, the actual base translation
$t\mapsto t+d$ commutes with every $\beta_g$.
In the $b_I$ coordinate it therefore commutes with
all real translations and is itself a translation:
$b_I(t+d)=b_I(t)+e$ for one $e>0$.
The fractional base prescription consequently satisfies
$B_{r,\tau}(t+d)=B_{r,\tau}(t)+d$.
The vector $a_I$ then has only its $d_I$ coordinate,
independent of the evaluation point. Its intervals
(C.77.10) obey the exact multiplier equation
$A_n(r,\tau,t+d)=\lambda A_n(r,\tau,t)$.
Perform the selection in the actual line bundle
(C.63.20). This retains the claimed fractional base
prescription and every component equation, including
negative $\lambda$. Thus every interval and every
original operator label are covered.

At a finite endpoint $a$, every label defining an
abelian-strip fiber interval has
$u_k(a)=r u_\tau(a)$; on a zero-character strip the
same equality has both sides zero. These are exactly
the hypotheses used in (C.76.20)--(C.76.21).
Compactness of the single global $K$ gives the same
uniform endpoint estimate over varying intervals.
Every base prescription lies between the identity
and the original base map; their inverses have the
corresponding inverse bounds. Consequently the
homeomorphisms and inverses glue to (C.76.2) on all
lines over $A$, including accumulation points which
are not endpoints of any one fixed interval.

Each selected fiber value is between two $K$-values
at the same real base level. The labelled triangle
comparison therefore gives the same bounds
\[
         K_{r,\tau}(x)\in K K^{-1}K(x),\qquad
         K_{r,\tau}^{-1}(x)\in K^{-1}K K^{-1}(x).
                                                               \tag{C.77.11}
\]
On the lines over $A$, these follow from the original
prefix path as in C.76. Compactness of these actual
original families, all fixing the original circle,
proves joint boundary continuity and inverse continuity
uniformly over all strips. This verifies the claimed
contraction and the continuous restriction map in the
actual topology, by the final argument of C.76.

Every interval retained by (C.76.1) has an actual base
image conjugate to all real translations by (C.75.4)
and the nonzero character's full real image. Its base
image is therefore abelian, proving the specified
inclusion $Z\subset Z_{\mathrm{ab}}$. On the subgroup
of $\operatorname{Homeo}(Z_{\mathrm{ab}})$ preserving
$Z$, restriction to $Z$ is a continuous homomorphism
for maps and inverses. Every constructed map belongs
to that subgroup, since it preserves each strip and
each line over $A$. This actual restriction sends
the new filling to a filling of the same original
restricted loop. No identity between the two
independently selected interior homotopies is
required or asserted.

\emph{The nonabelian base and the actual transitive kernel.}
Now let $B_I$ be nonabelian, and put $D=[G,G]$.
This is a path-connected subgroup of the actual
source group: every $g,h\in G$ has the finite-word
paths (C.75.7), so their commutator has a commutator
path from the identity; finite products give paths
to every element of $D$. Its base image is nontrivial,
normal in $B_I$, and path connected. Its common
fixed set is $B_I$-invariant, and hence is either
empty or all of $I$, since $B_I$ is transitive.
Nontriviality excludes the latter.

For every point of $I$, an actual commutator-product
path moves it to a different point; its inverse
path supplies movement to the other side. The
evaluation interval therefore contains a neighborhood
of that point. Each $\beta(D)$-orbit is path connected
and open. These disjoint orbits partition the
connected interval $I$, so there is exactly one.
Thus $\beta(D)$ is transitive. All endpoint characters
vanish on $D$, because they are additive characters.

For $v\in V=c(G)$ choose an actual $g\in G$ with
$c(g)=v$. For any desired $t'\in I$, choose $h\in D$
with $\beta_h(\beta_g(t))=t'$. Then $hg$ has character
$v$ and the desired base value. Taking $t'=t$ proves
$c(G_t)=V$, and arbitrary $t'$ proves (C.77.3).
The inverse of its displayed map is exactly
$(z,t)\mapsto(z-c(g),\beta_g^{-1}(t))$.
Both map and inverse depend continuously on the
actual source element on every compact target set,
by continuity of the character and of base restriction.
Its action kernel is exactly
\[
             \{g\in G:c(g)=0, \beta_g|_I=\id\}.
                                                               \tag{C.77.12}
\]
This proves the statement for the actual source words,
not only for their closure.

The vector space $V$ and the endpoint characters
extend to $L$ as in C.76. Endpoint action gives a
continuous action of the Polish based path group
$\mathscr P(L)$ on $V\times I$. Every original word
$hg$ just used is reached by its finite-word path,
so this path-group action is transitive. The target
is the nonmeagre Hausdorff manifold $V\times I$.
The exact orbit-open-map theorem already verified
in C.30 makes $e_t$ in (C.77.4) open. Its fibers
are the actual cosets of its closed stabilizer
$J_t$. The displayed quotient map is consequently
a continuous open bijection and hence a homeomorphism.
The kernel of the whole path-group action consists
of the paths whose endpoints satisfy (C.77.12),
now read in $L$; it has not been suppressed.

Apply C.64's proved compact-covering lemma to this
open map, using the complete compatible path metric.
For every compact target set lying in the image of
an assigned open path neighborhood, it provides
a compact preimage entirely inside that neighborhood.
The union of the values of these compact paths is
also compact in the original relative group, including
the inverse maps. This proves the final compactness
assertion. The argument supplies no single compact
family for a collection of target points escaping
to a finite endpoint or to infinity in the base
coordinate. The original-boundary-controlled selection
on these remaining nonabelian strips is therefore
still a precise unresolved step in C.12.
\end{proof}

\begin{claim}{C.78}{Proved}{Rank-one endpoint images with a controlled joint character and base lift}
Keep the actual disk, proper one-dimensional interior
operator group $Q$, original relative loop, principal
coordinate, and set $A$ of C.76--C.77. For each
component $I$ of $\mathbb R\setminus A$, let $B_I$
be its actual base image, and let $V_I\subset\mathbb R^k$,
$0\leq k\leq2$, be the actual vector image of all
finite-endpoint characters. Retain the strips satisfying
\[
                  B_I\text{ is abelian}
                       \quad\hbox{or}\quad\dim V_I\leq1.
                                                               \tag{C.78.1}
\]
Let $Z_1$ consist of these strips, all lines over $A$,
and the original boundary circle. The original loop
restricted to this closed set has a based contraction
through its actual homeomorphisms, fixing the circle
and commuting with each original interior operator.
Both the contraction and its inverse are jointly
continuous. On every line over $a\in A$ its value
is exactly $(s,a)\mapsto(s+r u_\tau(a),a)$.

The construction includes arbitrary accumulation of
strips at interior lines and at the original circle.
It gives an explicit continuous restriction homomorphism
as in (C.76.4), with $Z_1$ as target subset. When
every strip is retained it contracts the original
loop in its original relative centralizer. For a
proper $Z_1$ no lift through the excluded strips is
asserted. Every excluded strip is a bounded interval
$I=(a,b)$ with nonabelian base image and
\[
                    (c_a,c_b)(G)=\mathbb R^2.
                                                               \tag{C.78.2}
\]

Two exact constructions supply the new assertion.
Compact convex averages of actual source maps give
a continuous increasing base deformation. The C.62
triangle in the scalar-character/base comparison plane
realizes every such value by one compact three-word
family of actual original maps. A second selection
uses compact local pair transports for the exact
Polish path-group actions. No continuous section of
an orbit or character map is assumed.
\end{claim}
\passage{Section 23, folios 140-141, scans 279-281; Part C(b), exact rank-one endpoint-image filling and remaining two-character bounded-strip obstruction}


\begin{proof}
\emph{The first compact source family.}
Use the original closed relative group $L$ of C.76,
which is the closure of the group generated by the
original loop in its actual closed-disk topology.
Every member has the principal form
\[
                    k(s,t)=(s+u_k(t),\beta_k(t)),
                                                               \tag{C.78.3}
\]
preserves each interval and each line over $A$,
fixes the original circle, and commutes with $Q$.
Its maps and inverse maps are controlled by a
complete compatible metric. Put
$P=\{f_\tau\}\cup\{f_\tau^{-1}\}$.
Construct the compact path-connected family $K$
and its nested compact families $K_n$ exactly as
in (C.77.9), starting with $K_0=P^5$ and using
fifteen short increments at each stage. Their
compact parameter products and uniform limits
take place in the actual group $L$.

The newly retained strips have nonabelian $B_I$
and $\dim V_I=1$. Choose one nonzero endpoint
character $c=c_a$ on such a strip. Every other
finite-endpoint character is a real multiple of
$c$: this is exactly the assertion that the
actual vector image $V_I$ is a line. The relation
holds on $L$ by continuity. The scalar image is
$c(G)=c(L)=\mathbb R$.
For $p=(r,\tau)\in[0,1]^2$, put
\[
       F_n(p)=\{k\in K_n:c(k)=r c(f_\tau)\},\qquad
       F_\infty(p)=\{k\in K:c(k)=r c(f_\tau)\}.
                                                               \tag{C.78.4}
\]
These are nonempty compact sets. Nonemptiness
already follows from the original prefix, since
the real path $c(f_\sigma)$ joins $0$ to $c(f_\tau)$.
All labels in a given set have the prescribed
fraction of \emph{every} finite endpoint character.

\emph{A compact space of positive averages, with its exact topology.}
Choose a countable dense sequence $\varphi_j$ in
the unit ball of $C(K,\mathbb R)$ and define
\[
       \chi(k)=(2^{-j}\varphi_j(k))_{j\geq1}\in\ell^2,
       \qquad
       \mathcal P(F)=\overline{\operatorname{conv}}\chi(F)
                              \quad(F\subset K\text{ compact}).
                                                               \tag{C.78.5}
\]
The map $\chi$ is continuous and injective: the
dense continuous functions separate distinct points.
Its inverse on its image is continuous by compactness.
The closed convex hull $\mathcal P(K)$ is compact.
Indeed, a finite $\epsilon$-net in $\chi(K)$ places
its convex hull within $\epsilon$ of the convex
hull of a finite set. That latter hull is a compact
finite-dimensional polytope. Thus the full hull
is totally bounded; its closure in the complete
Hilbert space is compact. The same proof applies
to every $\mathcal P(F)$.

Every $q\in\mathcal P(K)$ defines a positive
average $\Lambda_q(h)$ of a continuous function
$h\in C(K)$. To give the exact definition, approximate
$q$ by finite convex combinations of the vectors
$\chi(k)$ and take the corresponding averages of
$h(k)$. This has a unique limit: approximate
$h/\max(1,\|h\|_\infty)$ uniformly by one of the
$\varphi_j$, whose average is the specified
$j$th coordinate multiplied by $2^j$.
The uniform approximation bounds the discrepancy
for any two approximating convex combinations.
It also proves independence of the approximating
sequence. The resulting functional is linear,
positive, sends the constant $1$ to $1$, and has
$|\Lambda_q(h)|\leq\|h\|_\infty$.
The same approximation proves joint continuity
in $q$ and in $h$ with its uniform norm.
For $q\in\mathcal P(F)$ one moreover has
\[
              \min_F h\leq\Lambda_q(h)\leq\max_F h.
                                                               \tag{C.78.6}
\]
Thus no representation theorem for measures or
additional selection theorem is needed.

There is a continuous choice
\[
 q_p\in\mathcal P(F_\infty(p)),\qquad
 q_{0,\tau}=\chi(1),\quad q_{1,\tau}=\chi(f_\tau),\quad
 q_{r,0}=q_{r,1}=\chi(1).
                                                               \tag{C.78.7}
\]
Here is its full lower-limit input. If $k\in F_n(p)$,
the correction $r'c(f_{\tau'})-c(k)$ tends to zero.
C.77's short-increment neighborhood statement realizes
it by an actual arbitrarily small increment $e$ in
the next allowed family. The same increment corrects
every label of $F_n(p)$, since those labels have
the same character. For a finite convex combination
of their $\chi$-images, replace every $k$ by $ek$.
Uniform continuity of $\chi$ on the compact source
family makes the new combination arbitrarily close.
Approximation by finite convex combinations therefore
proves
\[
       \mathcal P(F_n(p))\subset
          \liminf_{p'\to p}\mathcal P(F_{n+1}(p')).
                                                               \tag{C.78.8}
\]
Replace these convex sets on the four parameter
faces by the prescribed singletons in (C.78.7).
The lower-limit inclusion survives by continuity
of those prescribed values. Apply the elementary
relative approximation argument of (C.63.22), now
in the Hilbert space: extend one selected vector
locally as a constant, use a partition of unity,
and make successive errors less than $2^{-n}$
with successive differences less than $2^{1-n}$.
Every approximation stays in the convex compact
set $\mathcal P(K)$. Completeness gives a continuous
limit. At each parameter all approximating allowed
sets lie in the compact closed
$\mathcal P(F_\infty(p))$, so the limit lies there;
on the parameter faces the error to the prescribed
singleton tends to zero. This proves (C.78.7).

\emph{The base average and its actual three-word realization.}
On the new strip define
\[
                B_p(t)=\Lambda_{q_p}\bigl(k\mapsto\beta_k(t)\bigr).
                                                               \tag{C.78.9}
\]
It is jointly continuous: on each compact real
coordinate interval, $t\mapsto(k\mapsto\beta_k(t))$
is continuous into $C(K)$ in its uniform norm.
For $t_1<t_2$, compactness gives
\[
       \min_{k\in K}\bigl(\beta_k(t_2)-\beta_k(t_1)\bigr)>0.
                                                               \tag{C.78.10}
\]
Positivity of the average makes $B_p$ strictly
increasing. It is proper onto $I$: for each
$v\in I$, the compact set
$\{\beta_k^{-1}(v):k\in K\}$ is bounded away
from both endpoints of $I$. Beyond its upper
bound every $\beta_k(t)$ exceeds $v$, and below
its lower bound every $\beta_k(t)$ is less than
$v$. The same statements hold for their average.
Thus $B_p$ is an increasing homeomorphism of $I$.
Strict increase and joint continuity give a
jointly continuous inverse, by bracketing its
value between two fixed nearby input points.
The four faces in (C.78.7) give the required
identity and original base maps exactly.

Use the actual comparison-plane action
\[
       \begin{split}
       T_I(k)(z,t)&=(z+c(k),\beta_k(t)),\\
       T_I(k)^{-1}(z,t)&=(z-c(k),\beta_k^{-1}(t)).
       \end{split}
                                                               \tag{C.78.11}
\]
The source $\mathbb R\times I$ is a plane; its
proper real fibers for the target $z$ are the
vertical copies of $I$. Every displayed map
preserves that partition and induces an increasing
translation on its real target. By (C.78.6), the
value $B_p(t)$ lies between the smallest and
largest values $\beta_k(t)$ for $k\in F_\infty(p)$.
Both extrema have actual compact labels in $K$,
and those endpoints have the same $z$ coordinate
$r c(f_\tau)$ in (C.78.11).

Connect their labels by a path in $K$ and apply
the C.62 triangle and zero-set argument in this
comparison plane, using $z$ as the exact real
target. Its proof only requires proper real
fibers and increasing target maps; both have
just been checked. It does not require affine
maps along the fibers. The resulting three-letter
word retains its actual labels in $L$ and realizes
the desired comparison point. Thus, with
\[
                     M_0=K K^{-1}K,
                                                               \tag{C.78.12}
\]
for every $p,t$ there is an actual $w\in M_0$
such that $c(w)=r c(f_\tau)$ and
$\beta_w(t)=B_p(t)$. This is a pointwise source
realization in one compact family. No continuous
choice of $w$ has been asserted.

\emph{All comparison spaces and compact local pair transports.}
For the strips retained already in C.77 use its
base prescription. An abelian base strip uses
the joint character $a_I$ of (C.77.6) and target
space $a_I(G)$ with its translation action.
Its source point is $0$ and desired target is
$r a_I(f_\tau)$; $P^5\subset M_0$ realizes it.
A nonabelian strip with zero endpoint image uses
target space $I$, its actual base action, source
point $t$, and C.27's median target $B_p(t)$;
an original prefix in $P$ realizes that target.
The new rank-one strips use (C.78.11), source
$(0,t)$, and target $(r c(f_\tau),B_p(t))$.
Their compact initial realization is (C.78.12).

There are countably many such spaces, one for
each retained interval. Denote them by $X_i$.
Each is a Euclidean vector space or a product
of a line with an interval, or an interval.
The actual based path group $\mathscr P(L)$
acts transitively on each: the vector action
uses the original character image, the interval
action uses C.75, and the comparison-plane
action uses (C.77.3)--(C.77.4).
All their orbit maps are open by the same
verified Polish orbit theorem. Every action
uses the original endpoint character or base
map; any action kernel remains its literal
kernel in $\mathscr P(L)$.

We require the following exact strengthening
of C.64's compact lifting construction. Given
an identity neighborhood $U$ in $L$, there is
a compact path-connected family $R\subset U$
containing $1$ with this property: for every
$y\in X_i$ and every identity neighborhood $V$
in $L$, there is a neighborhood $N$ of $y$
such that
\[
 \text{for }y_1,y_2\in N\text{ there is }e\in R\cap V
                       \text{ with }e\cdot y_1=y_2.
                                                               \tag{C.78.13}
\]
No continuity of the choice $e$ is part of the
assertion.

To prove it, choose an identity neighborhood
$U_0$ with $U_0U_0^{-1}\subset U$.
Exhaust the disjoint countable union of the
$X_i$ by nested compact sets $E_m$ whose
interiors cover it. Use based path neighborhoods
whose values lie in $U_0$ and within $\epsilon_m$
of $1$, where $\epsilon_m\downarrow0$ in a
compatible metric on $L$. At each point of
$E_m$, openness of its actual orbit map supplies
a small closed coordinate ball $D_{m,j}$ within
the image of that path neighborhood. A finite
collection of their interiors covers $E_m$.
C.64's proved compact-covering lemma gives a
compact family of these small paths whose
endpoints map the chosen center exactly onto
$D_{m,j}$. Let $S$ be the union of all their
path values, together with $1$.

This $S$ is compact and path connected: at each
level there are finitely many compact path-value
families, and the families with $m\to\infty$
converge uniformly to $1$. It lies in $U_0$.
Set $R=SS^{-1}$, so $R$ is compact, path connected,
and contained in $U$. Given $y,V$, take $m$
sufficiently large that $y\in E_m$ and that
products of two level-$m$ values, one inverted,
lie in $V$. Choose a ball $D_{m,j}$ whose
interior contains $y$. For two points in that
interior choose endpoint labels $s_1,s_2$ from
the same compact path family mapping its one
center to those two points. Then
$e=s_2s_1^{-1}\in R\cap V$ maps the first to
the second. This proves (C.78.13) simultaneously
for every comparison space. Both maps and
inverses are controlled in the original group.

\emph{The second compact product and the actual fiber lift.}
Starting with $M_0$ in (C.78.12), choose an
identity neighborhood $U_n$ so small that its
left multiplication changes every previously
formed product by less than $2^{-n}$ in a
complete compatible metric on $L$.
Choose $R_n\subset U_n$ by (C.78.13), and put
\[
                   M_n=R_n\cdots R_1M_0.
                                                               \tag{C.78.14}
\]
Uniform Cauchy convergence on the product of
the compact parameter families gives a compact
path-connected limit family $M\subset L$
containing every $M_n$. Inverse convergence is
included by the actual group topology.

On every retained strip define the fiber interval
using precisely the appropriate target equality:
\[
 A_n(p,t)=\operatorname{conv}
 \{u_k(t):k\in M_n,
          \ k\cdot x_i(t)=y_i(p,t)\}.
                                                               \tag{C.78.15}
\]
Here $(x_i(t),y_i(p,t))$ are the actual source
and target points specified above: respectively
the character-vector pair, the base pair, or
the scalar-character/base pair.
These sets are nonempty and compact by the
initial realizations in $M_0$. Their defining
equality ensures the same actual target base
level $B_p(t)$ in each case.

For a fixed allowed $k\in M_n$ at $(p,t)$,
the two comparison points
$k\cdot x_i(t')$ and $y_i(p',t')$ tend to the
same point as $(p',t')\to(p,t)$.
Property (C.78.13), applied to $R_{n+1}$ and
an arbitrarily small $V$, supplies an actual
$e\in R_{n+1}\cap V$ carrying the first point
to the second. Then $ek\in M_{n+1}$ has exactly
the new target equality. Its fiber displacement
approaches that of $k$ as $V$ decreases to the
identity, uniformly near the specified source
point. Using the same convex coefficient for
two corrected endpoint values proves
\[
                   A_n(p,t)\subset
              \liminf_{(p',t')\to(p,t)}A_{n+1}(p',t').
                                                               \tag{C.78.16}
\]
The four prescribed fiber values are again
(C.76.13), represented by $1$ or $f_\tau$ in
$M_0$. The full relative interval selection
argument of C.76 therefore supplies a continuous
$U(p,t)$ in the limiting compact fiber interval,
with those faces exactly. Use (C.76.18) for the
actual homeomorphism and its displayed inverse.

\emph{Operator transport and every finite endpoint.}
Retain the original component map
$q(s,t)=(\lambda s+\kappa,t+d)$ of C.75,
when it exists. Choose all constructions on
representatives of its orbits of proper intervals.
For a new rank-one strip the chosen endpoint
character transports to $\lambda c$. Its
level sets (C.78.4) for the transported strip
are therefore identical subsets of the same
global $K$. Transport the selected positive
average by using the same vector $q_p$ in
$\mathcal P(K)$. Equation (C.76.17), linearity
and preservation of the constant function by
$\Lambda_{q_p}$ give
$B'_p(t+d)=B_p(t)+d$ exactly.
The allowed actual fiber intervals transform
by multiplication by $\lambda$.
Transport their selection by that multiplier;
the additive $\kappa$ cancels in the actual
commutation equation. For the older retained
strips use precisely the character or median
transport verified in C.77. A nonzero component
power never stabilizes a proper interval.
The interval $I=\mathbb R$ has no endpoint
characters and is one of those older cases;
its possible stabilizer is handled by the
actual line-bundle selection of C.77. Negative
$\lambda$ and every original component label
are consequently retained.

For the new base averages one has the uniform
pointwise bounds
\[
 \begin{split}
 \min_{k\in K}\beta_k(t)&\leq B_p(t)\leq
                          \max_{k\in K}\beta_k(t),\\
 \min_{k\in K}\beta_k^{-1}(v)&\leq B_p^{-1}(v)\leq
                          \max_{k\in K}\beta_k^{-1}(v).
 \end{split}                                                   \tag{C.78.17}
\]
The first is positivity. For the second, an
input below all $\beta_k^{-1}(v)$ has every
$\beta_k$-value below $v$, and an input above
all of them has every value above $v$.
The older base prescriptions and inverses
have the same bounds, since they lie between
the original map and the identity, both in $K$.

Now let retained interval points $t_n$ tend
to $a\in A$, allowing the interval and all
parameters to vary. Compactness of $K$ and
its inverse family, all fixing the base point
$a$, makes both bounds (C.78.17) tend to $a$.
Choose a finite endpoint $a_n$ of the source
interval between $t_n$ and $a$; then $a_n\to a$.
Every label in (C.78.15) has
$u_k(a_n)=r_nu_{\tau_n}(a_n)$: for an abelian
strip this is its vector constraint; for the
new strip it follows from its scalar constraint
and the proportionality of both endpoint
characters; for a zero-character strip both
sides vanish. Uniform continuity over the
one compact family $M$ gives
$u_k(t_n)-u_k(a_n)\to0$, uniformly in its
allowed labels. Convexification preserves
the estimate. This proves the specified
fractional endpoint displacement, including
varying intervals, just as in (C.76.21).
The inverse base bounds in (C.78.17) permit
the identical argument at the inverse input,
with the negative fiber displacement. Thus
both families glue continuously to every
line over $A$.

\emph{The original circle and the exact remaining region.}
Each selected actual fiber value is between
two $M$-values at the same real target level.
The original source-plane C.62 triangle and
inverse bracketing give
\[
       K_{r,\tau}(x)\in M M^{-1}M(x),\qquad
       K_{r,\tau}^{-1}(x)\in M^{-1}M M^{-1}(x).
                                                               \tag{C.78.18}
\]
On lines over $A$, the prescribed scalar
fraction is attained by an original prefix
value and its inverse, so these bounds hold
there too. Both displayed word families are
compact families of the original closed-disk
homeomorphisms and fix the circle pointwise.
They therefore send a sequence tending to
an original boundary point uniformly to that
same point. This proves joint continuity
and inverse continuity at the original
circle, with arbitrarily varying intervals
and unbounded principal coordinates.

The resulting piecewise bijections preserve
all the retained pieces, have the displayed
continuous inverses, and are actual
homeomorphisms of the closed set $Z_1$.
Their compact-parameter evaluations give
continuity in its actual homeomorphism
topology. The restriction subgroup and
homomorphism are precisely those of (C.76.4)
with this named closed subset. On the subgroup
preserving the older subset $Z_{\mathrm{ab}}$,
further restriction is a continuous homomorphism;
every constructed map preserves that subset.
This supplies the exact comparison with the
earlier restricted loop, without asserting
identity of independently chosen fillings.

If a strip remains excluded, its base image
is nonabelian and its endpoint vector image
has dimension two. It must therefore have
two finite endpoints, and their joint image
is all of $\mathbb R^2$, proving (C.78.2).
The scalar-character comparison plane used
in (C.78.11) does not retain both independent
characters of such a strip. No vertical
interval containment in the corresponding
three-dimensional comparison space has
been proved here, and no filling through
that excluded region is inferred.
The other remaining cases of C.12 are
unchanged.
\end{proof}

\begin{claim}{C.79}{Proved}{Every relative disk loop for a one-dimensional proper operator action fills}
Let $O=\operatorname{int}\overline D$ be the interior
of the actual closed disk. Let
$Q\leq\operatorname{Homeo}(O)$ be closed, locally
compact, free and proper, with $\dim Q=1$.
Use the actual relative centralizer
\[
 \mathscr C_Q=
 \{h\in\operatorname{Homeo}(\overline D):
       h|_{\partial\overline D}=\id,
       \ hq=qh\text{ on }O\text{ for every }q\in Q\}
                                                               \tag{C.79.1}
\]
with its compact-open topology on maps and inverses.
Every continuous based loop $f_\tau$ in this group
has a continuous based contraction in the same group.
In particular
\[
                         \pi_1(\mathscr C_Q,1)=0.
                                                               \tag{C.79.2}
\]
The construction uses the original disk boundary
and every original operator. Its map and inverse
families are jointly continuous there, including
arbitrary accumulation of principal strips.
It does not require a specified Fuchsian disk
comparison, and does not replace the original
compactification by another one.

The two-endpoint step retains one common increasing
base map and one compact global source family in
both endpoint selections. Their fiber displacements
are blended in that same original real fiber.
This gives an actual filling on the whole disk;
no extension through an omitted strip is left in
the stated dimension-one disk assertion.
\end{claim}
\passage{Section 23, folios 140-141, scans 279-281; Part C(b), complete actual closed-disk relative loop filling for one-dimensional free proper interior operator groups}


\begin{proof}
\emph{The exact strips and the common preliminary family.}
Use the actual principal coordinate and inverse
of C.75. Write
\[
        f_\tau(s,t)=(s+u_\tau(t),\beta_\tau(t)),
        \qquad A=\operatorname{Fix}\langle\beta_\tau\rangle.
                                                               \tag{C.79.3}
\]
Retain the original group $G=\langle f_\tau\rangle$
and its closure $L$ in the actual relative group
(C.79.1). The complete compatible metric, all
component operators, and the jointly continuous
principal-coordinate functions on compact source
families are exactly those in C.76--C.78.
There are countably many complementary intervals.

C.78 fills all their types except bounded intervals
$I=(a,b)$ with nonabelian effective base image and
actual joint character image
\[
                  c_I=(c_a,c_b):G\longrightarrow\mathbb R^2,
                  \qquad c_I(G)=\mathbb R^2.
                                                               \tag{C.79.4}
\]
We construct all strip fillings together, using
one compact family, so that their accumulation
estimates remain uniform. Begin with precisely
the preliminary compact $K$ and $K_n$ of C.78:
$K_0=P^5$, where $P$ is the original loop and its
inverse, followed by fifteenfold short-increment
families whose infinite products converge in $L$.
This construction is independent of the strip.

\emph{A full-vector average gives the one common base.}
On a strip (C.79.4), replace (C.78.4) by the
full two-character fibers
\[
       F_n(r,\tau)=
          \{k\in K_n:c_I(k)=r c_I(f_\tau)\},
       \qquad
       F_\infty(r,\tau)=
          \{k\in K:c_I(k)=r c_I(f_\tau)\}.
                                                               \tag{C.79.5}
\]
The planar additive identity (C.37.3) gives
nonempty fibers already in $P^3\subset P^5$.
For nearby parameters the required correction
is the actual vector
$r'c_I(f_{\tau'})-r c_I(f_\tau)$.
C.76's six-short-increment neighborhood covers
this vector when it is sufficiently small; it
is a subfamily of the fifteenfold increments
used here. The same increment corrects every
label in the given fiber. Thus the Hilbert
convex-average proof (C.78.5)--(C.78.8) applies
to (C.79.5) without change: its lower-limit
argument needs only a common arbitrarily small
character correction, not a scalar target.

Choose its continuous positive average $q_p$,
$p=(r,\tau)$, supported in the precise sense
$q_p\in\mathcal P(F_\infty(p))$, with the same
four exact parameter faces. Put
\[
               B_p(t)=\Lambda_{q_p}(k\mapsto\beta_k(t)).
                                                               \tag{C.79.6}
\]
The compact positive-gap and inverse-envelope
arguments (C.78.9)--(C.78.10) make this a jointly
continuous family of increasing homeomorphisms
of $I$, with jointly continuous inverses.
It has the original base map at $r=1$ and the
identity on the other based faces. Crucially,
this is the \emph{same} $B_p$ for both endpoints.

\emph{Two actual scalar comparison planes.}
For $e=a,b$ retain the actual source action
\[
        \begin{split}
        T_e(k)(z,t)&=(z+c_e(k),\beta_k(t)),\\
        T_e(k)^{-1}(z,t)&=(z-c_e(k),\beta_k^{-1}(t)).
        \end{split}
                                                               \tag{C.79.7}
\]
In each of these planes use $z$ as the real
target, whose fibers are the proper copies
of $I$. Every map has an increasing target
translation and an increasing fiber map.
The minimum and maximum labels for the
positive average (C.79.6) belong to
$F_\infty(p)$ and hence have \emph{both}
endpoint characters prescribed. Apply the
C.62 triangle separately in the two planes,
along a path in the same compact $K$ joining
those labels. It follows that the single
compact family
\[
                         M_0=K K^{-1}K
                                                               \tag{C.79.8}
\]
realizes $(r c_a(f_\tau),B_p(t))$ under $T_a$
at $(0,t)$, and separately realizes
$(r c_b(f_\tau),B_p(t))$ under $T_b$ there.
The realizing labels need not be equal.
No realization retaining both characters
at once is claimed at this step.

Each scalar action is transitive. Indeed,
C.77 proves that the actual source derived
subgroup has transitive base image and
vanishes under every endpoint character;
each character in (C.79.4) has full real
image. Thus the precise based path-group
orbit maps of both scalar actions are open,
with the actual kernels and stabilizers
described in (C.77.4) and (C.77.12).

Add these two comparison spaces per interval
to the countable collection in C.78. Apply
its proved compact local pair-transport
construction (C.78.13) to the whole collection.
The second infinite-product construction
(C.78.14) then yields \emph{one} compact
path-connected global family $M\subset L$
and nested compact families $M_n$ starting
with (C.79.8). This includes the comparison
spaces for all the older strip types as well.
All maps and inverses are original closed-disk
maps fixing the original circle.

\emph{The two fiber selections and their exact blend.}
For $e=a,b$, define
\[
 A_n^e(p,t)=\operatorname{conv}
 \{u_k(t):k\in M_n,
       \ c_e(k)=r c_e(f_\tau),\ \beta_k(t)=B_p(t)\}.
                                                               \tag{C.79.9}
\]
They are nonempty compact intervals by
(C.79.8). The local pair transports in their
respective scalar comparison planes prove
the interlaced lower-limit property exactly
as in (C.78.16). The four relative face
values are represented by the identity
and the original $f_\tau$. The same relative
interval selection therefore supplies two
continuous displacements $U_a(p,t)$ and
$U_b(p,t)$, each in its limiting compact
interval, with those four faces exactly.
Their actual base map is the common $B_p$.

Choose the explicit cutoff
\[
 \begin{split}
 \theta_I(t)&=\max\left(0,
          \min\left(1,3\frac{t-a}{b-a}-1\right)\right),\\
 U(p,t)&=(1-\theta_I(t))U_a(p,t)+\theta_I(t)U_b(p,t).
 \end{split}
                                                               \tag{C.79.10}
\]
It is zero on a neighborhood of $a$ and
one on a neighborhood of $b$. Define
\[
 \begin{split}
 H_p(s,t)&=(s+U(p,t),B_p(t)),\\
 H_p^{-1}(s,v)&=
        (s-U(p,B_p^{-1}(v)),B_p^{-1}(v)).
 \end{split}                                                   \tag{C.79.11}
\]
These are exact inverse homeomorphisms,
jointly continuous on the open strip.
Both selections have the same four faces,
so their blend does too.

Every value of $U_a$ or $U_b$ lies between
actual $M$-label displacements at that same
base level. Their convex blend lies in the
convex hull of the union of those labels.
The minimum and maximum of the compact
set of \emph{all} $M$-labels at that level
therefore bracket $U$. Applying the original
source-plane C.62 triangle and its inverse
bracketing gives directly
\[
           H_p(x)\in M M^{-1}M(x),\qquad
           H_p^{-1}(x)\in M^{-1}M M^{-1}(x).
                                                               \tag{C.79.12}
\]
There is no new word-family bound inferred
from an unrelated coordinate interpolation:
both selected displacements and their blend
have the identical actual real base level.
The older strip constructions use the same
$M$ and have the same two bounds by C.78.

\emph{All component operators.}
The original component generator is
$q(s,t)=(\lambda s+\kappa,t+d)$, with
$d>0$ and $\lambda\ne0$, when it exists.
For the translated interval $I+d$, the
two endpoint characters are multiplied by
$\lambda$, so the fibers (C.79.5) are the
same actual subsets of $K$. Transport the
positive average by that identical label
choice. Then $B'_p(t+d)=B_p(t)+d$, exactly
as in C.78. Transport both fiber selections
by $U'_{e+d}(p,t+d)=\lambda U_e(p,t)$.
The explicit cutoff satisfies
$\theta_{I+d}(t+d)=\theta_I(t)$, and hence
the blended displacement has the same
multiplier equation. This verifies actual
commutation with $q$, including negative
$\lambda$; the original constant $\kappa$
cancels from the equation and remains in
the operator. No nonzero component power
stabilizes a bounded interval, so transport
from interval representatives is independent
of the chosen operator. The older interval
types, including $I=\mathbb R$ with its
possible component stabilizer, retain their
exact C.78 line-bundle and transport equations.

\emph{The two finite-endpoint sequence alternatives.}
On every line over $a\in A$, prescribe
$H_p(s,a)=(s+r u_\tau(a),a)$.
Consider a sequence of source points with
$t_n\to a\in A$, allowing all parameters
and all interval indices to vary. The older
strip types have the joint limit already
proved in C.78, with the same compact $M$.
For the new bounded strips, every subsequence
has a further subsequence of one of these
two forms.

First, one fixed interval $I=(a_0,b_0)$
occurs throughout. The limiting point of
$A$ must be $a_0$ or $b_0$.
Near $a_0$ the cutoff is identically zero,
so $U=U_a$ and every defining label has
$u_k(a_0)=r u_\tau(a_0)$.
Near $b_0$ it is identically one, and the
same statement uses $U_b$ and $b_0$.
Uniform continuity over the compact actual
$M$ proves the required limit of $U$.

Second, the bounded intervals
$I_n=(a_n,b_n)$ are pairwise distinct.
Then both $a_n$ and $b_n$ tend to $a$.
For example, if along a subsequence
$b_n\geq a+2\epsilon$ while $t_n\to a$,
then for all sufficiently large $n$ the
fixed point $a+\epsilon$ lies in $I_n$,
contradicting disjointness of distinct
components. The analogous argument with
$a-\epsilon$ handles the other endpoints.
Thus uniform continuity of the one compact
family $M$ gives, respectively,
\[
 \begin{split}
 U_a(p_n,t_n)-r_nu_{\tau_n}(a_n)&\longrightarrow0,\\
 U_b(p_n,t_n)-r_nu_{\tau_n}(b_n)&\longrightarrow0.
 \end{split}                                                   \tag{C.79.13}
\]
For the first estimate use the labels in
the limiting interval of $A_n^a$, and for
the second use those for $A_n^b$; they
need not be the same labels. Each estimate
persists under its convexification.
The original compact loop makes both
endpoint terms approach the prescribed
value at $a$, uniformly in the parameters.
Any convex combination in (C.79.10) has
the same limit. This proves the desired
limit even when the cutoff values vary
without converging. The two subsequence
alternatives establish the assertion for
the original arbitrary sequence.

The base average and its inverse satisfy
the compact envelopes (C.78.17), as do
the older prescriptions. Therefore their
values tend to $a$ uniformly whenever
the source or inverse input tends to $a$.
For the inverse, put $v_n=B_{p_n}^{-1}(t_n)$.
Then $v_n\to a$ and it stays in the same
interval $I_n$. Apply the same two
alternatives at $v_n$ in (C.79.11).
This gives the negative prescribed fiber
displacement and proves inverse continuity
at the limiting line. Along $A$ itself,
both limits follow from the original
continuous function $r u_\tau(a)$.

\emph{The actual whole-disk filling.}
Use (C.79.11) on every new strip, the
simultaneously constructed C.78 maps on
every other strip, the prescribed map
on each line over $A$, and the identity
on the original boundary circle. These
form a bijection of the actual closed
disk with the displayed piecewise inverse,
since every base interval and every line
over $A$ is preserved. Interior joint
continuity of maps and inverses is exactly
the strip and endpoint comparison just
proved.

On a line over $A$, the scalar fraction
$r u_\tau(a)$ is attained by an original
prefix value; its negative is attained
by the corresponding inverse. Thus the
compact bounds (C.79.12) hold there too.
Both global word families in (C.79.12)
are compact families of actual closed-disk
homeomorphisms fixing the original circle.
As a point approaches any original circle
point, all their evaluations approach that
same point uniformly. This proves joint
boundary continuity of the pasted maps
and their inverses, including arbitrarily
varying strips and unbounded principal
coordinates.

Hence the family consists of actual
homeomorphisms of $\overline D$ and is
continuous in the topology of (C.79.1).
Its operator equations were checked for
every actual component and identity-component
operator. Its four parameter faces are
$H_{0,\tau}=1$, $H_{1,\tau}=f_\tau$,
and $H_{r,0}=H_{r,1}=1$.
This is the asserted based contraction,
proving (C.79.2).

The proof is a loop statement for the
specified actual disk group. It asserts
neither contractibility of that entire
group nor a homotopy lift through C.73's
prime-end injection for a different
original domain. Those exact comparison
and broader C.12 requirements remain
as previously stated.
\end{proof}

\begin{claim}{C.80}{Proved}{Every based sphere family fills in the one-dimensional proper disk centralizer}
Let $O=\operatorname{int}\overline D$ be the actual
open disk, and let $Q\leq\operatorname{Homeo}(O)$
be closed, locally compact, free and proper, with
$\dim Q=1$. Retain the actual group and topology
\[
 \mathscr C_Q=
 \{h\in\operatorname{Homeo}(\overline D):
    h|_{\partial\overline D}=\id,
    \ hq=qh\text{ on }O\text{ for every }q\in Q\},
                                                               \tag{C.80.1}
\]
using compact-open convergence of maps and inverses.
For every $n\geq1$ and every continuous based map
$f:(S^n,*)\to(\mathscr C_Q,1)$, there is a continuous
map $H:[0,1]\times S^n\to\mathscr C_Q$ satisfying
\[
 \begin{gathered}
 H_{0,\xi}=1,\qquad H_{1,\xi}=f_\xi,\qquad H_{r,*}=1.\\
 \text{Consequently }\quad
 \pi_n(\mathscr C_Q,1)=0\quad(n\geq1).
 \end{gathered}
                                                               \tag{C.80.2}
\]
Both evaluation families $H_{r,\xi}(x)$ and
$H_{r,\xi}^{-1}(x)$ are jointly continuous on the
original closed disk. Every original operator is
retained. This proves the stated based homotopy
groups, without asserting contractibility of the
whole group or describing its components.
\end{claim}
\passage{Section 23, folios 140-141, scans 279-281; Part C(b), all positive-degree based homotopy groups of the actual closed-disk relative centralizer for one-dimensional free proper interior operator groups}


\begin{proof}
\emph{The actual source group and the sphere parameter.}
Put $\Sigma=S^n$ and
$P=f(\Sigma)\cup f(\Sigma)^{-1}$. This is a compact
path-connected family containing $1$. Use
$G=\langle P\rangle$ and its closure $L$ in the
actual group (C.80.1). The latter is a closed
subgroup of the Polish closed-disk homeomorphism
group: the boundary equations and all commutator
equations on $O$ are closed under convergence of
maps and inverses. Thus $L$ has the complete
compatible group topology used in C.76--C.79.

C.33's actual principal coordinate, with its
inverse, has $Q^0$ acting by real translations
in the first coordinate. Hence every $k\in L$
has the form
\[
 \begin{split}
 k(s,t)&=(s+u_k(t),\beta_k(t)),\\
 k^{-1}(s,t)&=
    (s-u_k(\beta_k^{-1}(t)),\beta_k^{-1}(t)).
 \end{split}                                                   \tag{C.80.3}
\]
The base maps are increasing. This holds first
for $f_\xi$, since a path in $\Sigma$ from $*$
to $\xi$ gives an actual based path of maps,
and then for $G$ and its closure. Every finite
word in $P$ has a based path within the same
finite power of $P$: choose a parameter path
for each letter and multiply the resulting
paths, with the indicated inverses. A jointly
continuous choice of those paths is unnecessary.

Let $A=\{t:\beta_{f_\xi}(t)=t\text{ for all }\xi\}$.
It is closed. Every complementary component
$I$ is a transitive base interval for $G$.
Indeed a based path moving a point and its
inverse give a neighborhood in its base orbit;
the orbit is connected because the increasing
union $\bigcup_m P^m(t)$ is connected. A finite
endpoint of that interval orbit is fixed by
all its increasing base maps. This proves
the stated component equality, exactly as
in C.75, for this actual sphere-generated $G$.
Every $k\in L$ fixes $A$ pointwise on the base
and preserves each $I$.

The endpoint characters are the actual evaluations
$c_a(k)=u_k(a)$, $a\in A$. They are continuous
characters by the composition formula (C.75.6).
All other source facts needed from C.75--C.77
have the same precise hypotheses here. For an
abelian effective base image, its closure is
abelian, transitive and free; the verified
evaluation homeomorphism identifies the actual
image with real translations, as in C.77.
This gives an increasing $b_I:I\to\mathbb R$
and the continuous character
\[
 d_I(k)=b_I(\beta_k(t_0))-b_I(t_0),\qquad
 b_I\beta_k(t)=b_I(t)+d_I(k).
                                                               \tag{C.80.4}
\]
For a nonabelian base image, $[G,G]$ is path
connected: each commutator and each finite
product of commutators has a based path from
the finite-word paths above. Its nontrivial
base image is normal. Its common fixed set
is $G$-invariant and therefore empty by base
transitivity. Its interval orbits are open
by the same path argument. They partition
the connected interval $I$, so it is transitive
there. All endpoint characters vanish on
$[G,G]$. Thus every scalar character/base
comparison used below is exactly the transitive
action in (C.77.3), with its stated kernel.

\emph{Short increments with the complete original topology.}
Give $\Sigma$ its round geodesic metric $d$.
For $0<\delta<\pi$ put
\[
 J_\delta=
 \{f_\xi f_\eta^{-1}:d(\xi,\eta)\leq\delta\},
 \qquad E_\delta=J_\delta^{15}.
                                                               \tag{C.80.5}
\]
The pair space in this formula is compact and
path connected. The unique shortest geodesic
from $\xi$ to $\eta$ contracts its first entry
to the second, continuously on this pair space
and without increasing their distance. The
diagonal is a copy of the path-connected sphere.
Consequently $J_\delta$ is compact, symmetric,
path connected, and contains $1$. Uniform
continuity of $f$ and of multiplication and
inversion on its compact image proves
$J_\delta\to\{1\}$ in the actual topology as
$\delta\downarrow0$.

For any endpoint-character vector, or the
vector consisting of $d_I$ and all finite
endpoint characters, there are at most three
coordinates. Its actual image $V$ is a real
vector subspace. To check this before closure,
apply the five-letter additive path identity
(C.50.3) to a finite-word based path ending
at any specified $g\in G$. This puts every
fraction of its vector in the actual image.
Addition, negatives and integer multiples
then give the whole real span. Applying the
same identity to a path from $*$ to $\xi$
within $\Sigma$ also gives
\[
 r a(f_\xi)\in a(P^5),\qquad
 a(E_\delta)\text{ contains a neighborhood of }0
                          \text{ in }V.
                                                               \tag{C.80.6}
\]
For the second assertion, subdivide a parameter
path into pieces of length less than $\delta$.
Telescoping shows that $a(J_\delta)$ spans $V$.
Choose at most three basis vectors in this
compact symmetric path-connected image. The
same five-letter identity supplies the full
segment from minus to plus each basis vector;
their sum contains a neighborhood of zero.
This uses at most fifteen factors. Every such
character extends continuously to $L$, and
closedness of $V$ gives $a(L)=V$.

Choose the successive $\delta_m$ in (C.80.5)
as in C.77--C.78, so multiplication by the
next compact increment changes every preceding
compact product by less than $2^{-m}$ in a
complete compatible metric. Starting with
$K_0=P^5$, the resulting nested families
$K_m$ lie in one compact path-connected limit
family $K\subset L$. The product parameter
space is a product of compact path-connected
spaces, and uniform Cauchy convergence gives
its continuous map into $L$. These statements
include convergence of the inverse maps.
One such $K$ serves every strip.

\emph{The exact relative parameter conditions.}
Use the compact metric parameter space
\[
 T=[0,1]\times\Sigma,\qquad
 D=(\{0,1\}\times\Sigma)\cup([0,1]\times\{*\}).
                                                               \tag{C.80.7}
\]
On $D$ the prescribed source labels are $1$
at $r=0$, $f_\xi$ at $r=1$, and $1$ at $\xi=*$.
They agree on all intersections because $f_*=1$.
They are continuous on $D$ and belong to $K_0$.

For every nonabelian strip with nonzero endpoint
vector $c_I$, form
\[
 F_m(p)=\{k\in K_m:c_I(k)=r c_I(f_\xi)\},
 \qquad p=(r,\xi)\in T.
                                                               \tag{C.80.8}
\]
These compact fibers are nonempty by (C.80.6).
Use exactly C.78's embedding $\chi:K\to\ell^2$
and compact convex sets $\mathcal P(F)$, with
their continuous positive averages $\Lambda_q$.
The actual character correction from $p$ to
$p'$ is $r'c_I(f_{\xi'})-r c_I(f_\xi)$. By
(C.80.6), an arbitrarily small next increment
realizes it and corrects every label in the
old fiber at once. The finite-convex-combination
argument (C.78.8) gives the interlaced lower-limit
inclusion for these Hilbert convex sets.

The relative selection proof of C.63 and C.78
applies on $T$, with $D$ as its closed relative
subset. Explicitly, replace the allowed sets
on $D$ by the continuous singleton labels
$\chi(1)$ or $\chi(f_\xi)$ just specified.
At a point outside $D$ a neighborhood avoids
$D$; at a point of $D$ the old lower-limit
inclusion and continuity of these singleton
values give the relative inclusion. The
locally finite partitions of unity and
summable successive approximations used in
C.63 require a metrizable paracompact domain,
not a rectangular parameter boundary. Both
$T$ and $T\times I$ have that property.
Their limits lie in the closed limiting
compact-family convex sets, and their errors
to the prescribed singletons tend to zero.
We therefore obtain continuous positive
averages $q_p\in\mathcal P(F_\infty(p))$ with
\[
 q_{0,\xi}=\chi(1),\qquad q_{1,\xi}=\chi(f_\xi),
                       \qquad q_{r,*}=\chi(1).
                                                               \tag{C.80.9}
\]

\emph{All base prescriptions and their compact realizations.}
For an abelian base strip use
$B_p(t)=b_I^{-1}(b_I(t)+r d_I(f_\xi))$.
For a nonabelian base strip with nonzero
endpoint vector use
$B_p(t)=\Lambda_{q_p}(k\mapsto\beta_k(t))$.
In the rank-two case the average uses the
full vector in (C.80.8), so both scalar
endpoint selections below have the same base.
For a nonabelian strip with zero endpoint
vector, use instead the formula
\[
                  B_p(t)=(1-r)t+r\beta_{f_\xi}(t).
                                                               \tag{C.80.10}
\]
This last choice removes any dependence on
an interval ordering of the sphere parameters.
It is an increasing homeomorphism of $I$:
strict inequalities are preserved by its
convex combination, and both summands tend
to the same respective endpoints of $I$.
Its inverse is jointly continuous by strict
monotone bracketing. For each fixed $p,t$,
a parameter path from $*$ to $\xi$ shows by
the intermediate value theorem that an
original label in $f(\Sigma)\subset P$
realizes $B_p(t)$ at $t$. No continuous
choice of this path or label is used.

The other two prescriptions and their inverses
are jointly continuous by C.78's compact strict
gap and inverse-envelope arguments. All three
have exactly the relative conditions (C.80.7).
They and their inverses lie between the
corresponding compact $K$-evaluation envelopes.

Use the same countable comparison spaces as
in C.78--C.79: the joint character vector
space for an abelian strip; $I$ for a zero
endpoint vector; and each required scalar
plane with the actual action
\[
 \begin{split}
 T_e(k)(z,t)&=(z+c_e(k),\beta_k(t)),\\
 T_e(k)^{-1}(z,t)&=(z-c_e(k),\beta_k^{-1}(t)).
 \end{split}
                                                               \tag{C.80.11}
\]
For a rank-one vector choose a nonzero endpoint
character; the other endpoint character, when
present, is its actual real multiple. For a
rank-two vector use both planes, separately.
Their desired targets are exactly
$(r c_e(f_\xi),B_p(t))$ at source $(0,t)$.
Two extremal labels for the average have
the same full endpoint vector. A path in
$K$ joins them. The C.62 triangle in either
scalar plane, with real target $z$, realizes
the entire intervening fiber interval in
the single actual family $M_0=K K^{-1}K$.
The vector targets and (C.80.10) already
have labels in $P^5$ and $P$, respectively.
Thus $M_0$ is a common initial realizing
family for every listed comparison target.

All these actions of the actual Polish based
path group $\mathscr P(L)$ are transitive.
For the scalar planes this follows from the
transitive derived base image and the full
nonzero character image proved above; the
other cases use the actual vector image or
base orbit. The orbit map is endpoint action
on the named source point. Its stabilizer
is exactly the paths whose endpoint fixes
that point; the action kernel retains the
zero character and identity-base equations.
The verified Polish orbit theorem makes
these precise orbit maps open.

C.78's compact local pair-transport construction
therefore supplies, for every small identity
neighborhood, a compact family of actual maps
that realizes arbitrarily small pair transports
near each fixed comparison point. It is built
from compact preimages under these path-group
orbit maps and includes all countably many
spaces at once. Its proof uses no parameter
ordering. The second compact-product construction
gives one compact path-connected family $M\subset L$,
with nested $M_m$ starting at $M_0$, exactly
as in (C.78.13)--(C.78.14).

\emph{Relative fiber selection and the common-base splice.}
For each target equality just listed, take
the convex interval of $u_k(t)$ over labels
$k\in M_m$ satisfying that equality. These
intervals are nonempty. Small pair transports
correct an allowed label at $(p,t)$ to the
precise nearby target at $(p',t')$ and converge
to the identity. This is exactly the lower-limit
proof (C.78.16). Replace the intervals on
$D\times I$ by the values $0$, $u_{f_\xi}(t)$,
and $0$, in the order prescribed in (C.80.7).
The same relative approximation proof gives
continuous selections with these values exactly.

On a rank-two interval $I=(a,b)$ perform the
two scalar selections using the same $B_p$
and the same $M$. Denote them by $U_a,U_b$,
and use precisely the cutoff and blend of
(C.79.10):
\[
 \begin{split}
 \theta_I(t)&=\max\left(0,\min\left(1,
                     3\frac{t-a}{b-a}-1\right)\right),\\
 U&=(1-\theta_I)U_a+\theta_I U_b.
 \end{split}
                                                               \tag{C.80.12}
\]
All other intervals have their single selected
$U$. Define the maps and their actual inverses by
\[
 \begin{split}
 H_p(s,t)&=(s+U(p,t),B_p(t)),\\
 H_p^{-1}(s,v)&=
    (s-U(p,B_p^{-1}(v)),B_p^{-1}(v)).
 \end{split}
                                                               \tag{C.80.13}
\]
Every defining label in a fiber selection has
that identical actual base level. The rank-two
blend stays between the extrema of their union
at that level. The original source-plane C.62
triangle and its inverse bracketing therefore
give, uniformly over all parameters and strips,
\[
 H_p(x)\in M M^{-1}M(x),\qquad
 H_p^{-1}(x)\in M^{-1}M M^{-1}(x).
                                                               \tag{C.80.14}
\]
The required connecting paths lie in the
compact path-connected family $M$; no continuous
choice of them enters these set containments.

\emph{Every original component operator and every limit.}
Retain C.75's actual component generator
$q(s,t)=(\lambda s+\kappa,t+d)$ when it exists,
with $d>0$ and $\lambda\ne0$. Choose constructions
on interval representatives and transport them.
Endpoint characters multiply by $\lambda$;
the full-vector fibers are the same actual
label sets after that multiplication. Transport
their averages by those same labels, and their
fiber selections by $U'(p,t+d)=\lambda U(p,t)$.
All base prescriptions satisfy
$B'_p(t+d)=B_p(t)+d$. In particular (C.80.10)
has this identity by its displayed affine formula.
The cutoff in (C.80.12) has it as well.
For $I=\mathbb R$ stabilized by a component
translation, the actual line-bundle selection
of C.76--C.78 retains the equation with
$\lambda$, using its positive norm
$|\lambda|^{-t/d}|z|$. The relative approximation
works in that same bundle over $T$ times the
quotient base. Thus every operator equation,
including negative $\lambda$ and the unchanged
constant $\kappa$, holds exactly.

On each $a\in A$ put
$H_p(s,a)=(s+r u_{f_\xi}(a),a)$.
The finite-line limit proofs of C.78--C.79
are uniform for the new compact parameter
space $T$. In detail, compactness of $M$
and $f(\Sigma)$ gives uniform continuity
of their principal displacements on bounded
base intervals. Labels constrained at a
finite endpoint have precisely the specified
fraction there. For rank two, a fixed interval
approaching an endpoint uses the corresponding
selection because the cutoff is identically
zero or one nearby. For pairwise distinct
bounded intervals approaching a finite line,
both endpoints approach that line, by the
disjoint-interval argument of C.79. Each scalar
selection therefore has the same limiting
prescribed displacement; so does their blend.
Every arbitrary sequence has a subsequence
of one of those two forms. The compact inverse
base envelopes first give the same limiting
input line for (C.80.13), and the identical
argument then proves its inverse limit.
These assertions allow $r$ and $\xi$ to vary.

On a line over $A$, the continuous image
$\{u_{f_\eta}(a):\eta\in\Sigma\}$ is an interval
containing $0$ and $u_{f_\xi}(a)$. It therefore
contains $r u_{f_\xi}(a)$; the corresponding
inverse label realizes the negative displacement.
Thus the compact bounds (C.80.14) hold on
these lines too. Set $H_p=1$ on the original
boundary circle. Both word families there
are compact families of original closed-disk
maps fixing that circle. All their evaluations
at a point tending to a given circle point
converge uniformly to that same point. This
proves both required boundary limits, even
through varying strips and unbounded principal
coordinates.

The piecewise formulas preserve each base
interval and line and are inverse bijections.
The preceding limits make both maps jointly
continuous on $T\times\overline D$. Because
$\overline D$ is compact, this is continuity
$T\to\mathscr C_Q$ in its specified map-and-inverse
topology. The relative selections give exactly
the three equations in (C.80.2). Collapsing
$\{0\}\times S^n$ gives the cone on $S^n$,
homeomorphic to the closed $(n+1)$-disk, and
the basepoint ray remains the identity.
This is the asserted based filling.

The proof uses path connectedness and compactness
of $S^n$ for $n\geq1$. It makes no assertion
about $\pi_0$ or a single contraction of all
of $\mathscr C_Q$. It also supplies no homotopy
lift through the prime-end injection of C.73
for a different original compactification;
the remaining general C.12 extension question
retains that precise scope.
\end{proof}

\begin{claim}{C.81}{Proved}{Exact metric, compact-support, and finite-configuration comparisons for discrete operators}
Let $O=\operatorname{int}\overline D$ be the actual
open disk, let $d$ be a compatible metric on its
closed disk, and let $Q\leq\operatorname{Homeo}(O)$
be closed, discrete, free and proper. Retain the
actual relative centralizer
\[
 \mathscr C_Q=
 \{h\in\operatorname{Homeo}(\overline D):
 h|_{\partial\overline D}=\id,
 \ hq=qh\text{ on }O\text{ for every }q\in Q\}
                                                               \tag{C.81.1}
\]
with compact-open convergence of maps and inverses.
Assume that an actual subgroup $G\leq\mathscr C_Q$
is transitive on $O$. Write $N=O/Q$ and retain
the original quotient covering $\pi:O\to N$.
No extension of the operators $q$ to the circle
is assumed. The following comparisons hold.

The metric
\[
 \begin{split}
 d_Q(x,y)&=\sup_{q\in Q}d(qx,qy),\\
 \rho(\pi x,\pi y)&=\inf_{q\in Q}d_Q(x,qy).
 \end{split}
                                                               \tag{C.81.2}
\]
gives the original topology on $O$ and its
quotient topology on $N$. The specified covering
is locally isometric for these metrics and each
fiber is uniformly discrete. Restriction embeds
$\mathscr C_Q$ as a closed subgroup of the uniform
$d_Q$-homeomorphisms commuting with $Q$, and is
an isometry for the sum of the uniform map and
inverse distances and the corresponding actual
closed-disk distances.

Every compact parameter family of equivariant
interior homeomorphisms and their inverses,
supported over one fixed compact subset of $N$,
extends jointly by the identity to the original
circle. In particular, every based sphere family
in $\mathscr C_Q$ whose descents are supported
in the interior of one compact connected
subsurface $N_0\subset N$ with nonempty boundary
has a based filling in $\mathscr C_Q$, retaining
that quotient support and every original label.
Here sphere dimension is any $n\geq1$.

For finite $A\subset O$ with distinct quotient
labels put
$\mathscr C_Q(A)=\{h\in\mathscr C_Q:h|_A=\id\}$.
For $x\notin QA$, the actual evaluation map is
a locally trivial principal bundle
\[
 \mathscr C_Q(A\cup\{x\})\longrightarrow
 \mathscr C_Q(A)\xrightarrow{\ h\mapsto h(x)\ }O\setminus QA.
                                                               \tag{C.81.3}
\]
Its source and fiber retain their original
closed-disk topologies. Every inclusion
$\mathscr C_Q(A)\hookrightarrow\mathscr C_Q$
induces an isomorphism on $\pi_n$ at $1$ for
every $n\geq1$.

When $N$ is noncompact, descent from the subgroup
of actual lifts of compactly supported quotient
isotopies is not continuous from its induced
$\mathscr C_Q$ topology to the Whitney topology
on the quotient homeomorphism group. The proof
gives an explicit convergent sequence and one
positive continuous majorant excluding all its
descents.

These are finite-configuration and fixed compact
quotient-support comparisons. They do not provide
a coherent infinite-configuration homotopy or
a controlled fragmentation of every original
family into compact quotient supports.
\end{claim}
\passage{Section 23, folios 140-141, scans 279-281; Part C(b), exact discrete-operator topology, fixed quotient-support sphere fillings, and finite-configuration continuation for the remaining level-zero question}


\begin{proof}
\emph{The invariant metric and the actual open evaluation map.}
Let $L$ be the closure of $G$ in (C.81.1).
It is a Polish group, remains transitive on $O$,
and commutes with every $q$. Closedness of the
commutator equations uses only the interior:
if $h_j\to h$ with inverses, then $h_jx\to hx\in O$,
so continuity of $q$ gives $qh_jx\to qhx$.
The verified Polish orbit theorem in C.30 makes
$L\to O$, $g\mapsto gx$, open.

The first formula (C.81.2) is finite because
$d$ is bounded on the compact disk. The identity
term separates points, and the triangle inequality
holds term by term. Reindexing by right
multiplication proves $Q$-invariance. Also
$d(x,y)\leq d_Q(x,y)$.
For the reverse topology comparison, take
$x_j\to x$ in $O$. Openness of evaluation
supplies $g_j\to1$ in $L$ with $g_jx=x_j$:
choose a decreasing identity-neighborhood basis
$U_m$, let $m(j)\to\infty$ so that
$x_j\in U_{m(j)}x$, and choose one such $g_j$.
No section has been assumed. Commutation gives
\[
 d_Q(x_j,x)=\sup_{q\in Q}d(g_jqx,qx)
       \leq\sup_{z\in\overline D}d(g_jz,z)
       \longrightarrow0.
                                                               \tag{C.81.4}
\]
This proves equality of the two interior
topologies, and in particular separability
of the displayed invariant metric space.

\emph{The specified quotient is a metric covering.}
Properness, discreteness and freeness give an
open neighborhood $U$ of any $x$ whose distinct
$Q$-translates are disjoint. Choose $r>0$ with
$B_{d_Q}(x,r)\subset U$. Then
$d_Q(x,qx)\geq r$ for $q\ne1$, and invariance
makes the entire fiber $Qx$ uniformly $r$-discrete.

The formula for $\rho$ is independent of both
representatives by invariance and reindexing.
Symmetry and the triangle inequality follow
from those of $d_Q$. A zero value gives
$q_jy\to x$ for some $q_j$; the closed-orbit
property of the proper action then gives
$x\in Qy$. Thus it is a metric. The equality
\[
 \pi B_{d_Q}(x,\epsilon)=B_\rho(\pi x,\epsilon)
                                                               \tag{C.81.5}
\]
shows that it gives the quotient topology,
and $\rho(\pi x,\pi y)\leq d_Q(x,y)$.
For $y,z\in B_{d_Q}(x,r/5)$ and $q\ne1$,
\[
 d_Q(y,qz)\geq r-2r/5=3r/5,
                    \qquad d_Q(y,z)<2r/5.
                                                               \tag{C.81.6}
\]
Consequently $\rho(\pi y,\pi z)=d_Q(y,z)$.
The inverse image of $\pi B_{d_Q}(x,r/5)$
is precisely the disjoint union of the translates
of that ball, each mapped isometrically onto
this quotient neighborhood. These are all
three metric-cover conditions in
\cite[Definition~2.1]{C_YagasakiUniform2014},
proved here for the actual map $\pi$.
The separation constant may depend on the
fiber; no compact quotient is inferred.

\emph{The exact isometry of the relative source group.}
Let $\mathscr U_Q$ denote the group of uniform
homeomorphisms of $(O,d_Q)$ commuting with $Q$,
with the uniform topology on maps and inverses.
An $h\in\mathscr C_Q$ belongs to this group
after restriction. Indeed its modulus of
uniform continuity on the original compact
disk applies to every pair $qx,qy$; commutation
then bounds $d(hqx,hqy)$ uniformly in $q$.
This is uniform $d_Q$-continuity of $h$.
Apply the identical argument to $h^{-1}$.

For $h,k\in\mathscr C_Q$ the exact formula is
\[
 \begin{split}
 \sup_{x\in O}d_Q(hx,kx)
   &=\sup_{x\in O,\ q\in Q}d(hqx,kqx)\\
   &=\sup_{z\in\overline D}d(hz,kz).
 \end{split}                                                   \tag{C.81.7}
\]
The last equality uses $q=1$ and density of
$O$. Repeat it with both inverses. Thus
restriction is an isometry for the stated
sum metrics, and its inverse on its image
is continuous. The source metric is complete:
simultaneous uniform limits of maps and
inverses on the compact disk are continuous
inverse maps, and all boundary and commutator
equations are closed. Its isometric image
is therefore closed in $\mathscr U_Q$.

This does not give an extension for arbitrary
members of $\mathscr U_Q$. For the explicit
case $Q=\{1\}$ and the standard round disk,
$d_Q=d$ and a nonidentity rotation is a
uniform interior homeomorphism. Its boundary
value is the same rotation, so it has no
preimage under this particular restriction
map. The named relative image condition
must consequently be retained in any use
of a uniform-homeomorphism deformation.

\emph{Compact cores and one common original source bound.}
Fix $x_0\in O$. C.64's proved compact-covering
lemma for the actual open evaluation map
gives, for every compact $K\subset O$, a
compact family $S\subset L$ with $K\subset S(x_0)$.
One covers $K$ by finitely many evaluation
neighborhoods and takes the union of their
compact preimages. For compact $K,K'$ choose
such $S,T$, respectively. For $z=sx_0\in K$
and $z'=tx_0\in K'$, commutation gives
\[
 \begin{split}
       qz'&=tqx_0=ts^{-1}(qz),\\
       qK'&\subset TS^{-1}(qz)\quad(q\in Q,\ z\in K).
 \end{split}
                                                               \tag{C.81.8}
\]
Every member of the one compact family
$TS^{-1}$ is an original circle-fixed map.
Thus $z_j\in K$ and $q_jz_j\to y$ on the
original circle imply uniform convergence
of the whole sets $q_jK'$ to $y$. No boundary
extension or convergence of $q_j$ is used.

Now take a compact parameter family $h_p$ of
$Q$-equivariant interior maps and inverses,
identity outside $\pi^{-1}(K_0)$ for a fixed
compact $K_0\subset N$. Choose compact
$C\subset O$ with $\pi C=K_0$, using finitely
many local sections of the covering. Include
in $C'$ all of $C$, $h_p(C)$ and $h_p^{-1}(C)$.
Joint interior continuity and compactness make
$C'$ compact. Since $\pi^{-1}(K_0)=QC$,
(C.81.8) supplies one compact original
family $P=TS^{-1}\cup\{1\}$ with
\[
 \begin{gathered}
                h_p(x)\in P(x),\qquad h_p^{-1}(x)\in P(x)\\
                       (p\text{ any parameter},\ x\in O).
 \end{gathered}
                                                               \tag{C.81.9}
\]
Inside $QC$ use equivariance and (C.81.8);
outside it use the identity. Compact evaluation
of this boundary-fixed $P$ proves both joint
circle limits. The pasted maps and inverses
are inverse homeomorphisms of the actual disk,
continuous in the actual topology of (C.81.1).

\emph{A based sphere with fixed compact quotient support.}
Let $f:(S^n,*)\to(\mathscr C_Q,1)$, $n\geq1$,
have descents $b_\xi$ identity outside
$\operatorname{Int}N_0$, where $N_0$ is as
in the statement. Descent and its inverse
family are continuous by the open covering
quotient and its local sections. The restrictions
to $N_0$ lie in the identity component of
$\operatorname{Homeo}(N_0\mathrel{\mathrm{rel}}\partial N_0)$:
the sphere is path connected and its basepoint
is the identity.

That component is contractible by
\cite[Lemma~2.1(ii)]{A_Yagasaki2000}.
The nonempty entire boundary is fixed, so
all its compact-surface exceptions are excluded.
A contraction can be normalized to fix the
identity throughout by replacing its value
$C_r(h)$ by $C_r(h)C_r(1)^{-1}$, with the
parameter directed from $1$ to $h$.
Applied to $b_\xi$ and extended by the identity
off $N_0$, it gives a continuous family
\[
 B_{0,\xi}=1,\qquad B_{1,\xi}=b_\xi,
                         \qquad B_{r,*}=1
                                                               \tag{C.81.10}
\]
of homeomorphisms of $N$ supported in $N_0$.

Lift the homotopy $B_{r,\xi}\pi$ from its
specified initial lift $H_{0,\xi}=\id_O$.
The covering homotopy property gives a joint
lift on $[0,1]\times S^n\times O$; it does
not require the sphere parameter to be simply
connected. Lift the inverse homotopy as well.
The two composites lift the identity and
start at the identity, hence are the identity
by uniqueness of lifts. For each original
$q$, the maps $H_{r,\xi}q$ and $qH_{r,\xi}$
lift the same family with identical initial
value $q$. They agree. The basepoint family
in (C.81.10) lifts to the constant identity.

At $r=1$, both $H_{1,\xi}$ and $f_\xi$ lift
the same $b_\xi$. Their relative deck label
is locally constant in $\xi$ and equals $1$
at $*$. Connectedness of the sphere therefore
gives $H_{1,\xi}=f_\xi$ for every parameter.
This proves agreement with the original
family upstairs, rather than only after
descent. Points above $N\setminus\operatorname{Int}N_0$
have constant base tracks and initial identity
lifts, so they stay fixed. Formula (C.81.9)
applies to this compact parameter family and
its inverse. The resulting joint original-circle
extension is the asserted based filling in
$\mathscr C_Q$ with every deck label retained.

\emph{Continuous local sections for actual orbit evaluation.}
For finite $A$ as in (C.81.3), its $Q$-saturation
is closed and locally finite in $O$. Commutation
identifies $\mathscr C_Q(A)$ with the subgroup
fixing that saturation pointwise. Choose
$x\notin QA$ and a small quotient disk $U$
about $\pi x$, with compact closure disjoint
from $\pi A$, that is evenly covered.
A local coordinate bump translation moves
$\pi x$ to each sufficiently nearby $y$,
using a fixed compact subdisk as support.
One explicit local formula is
$z\mapsto z+\chi(z)v$, where $\chi=1$ near
the center, $\chi=0$ outside that subdisk,
and $|v|\operatorname{Lip}(\chi)<1/2$.
The two-sided Lipschitz estimate makes it
injective with continuously varying inverse;
properness and invariance of domain make
it onto. The parameter $av$, $0\leq a\leq1$,
gives its based supported isotopy.

Lift that isotopy from the identity, retaining
the sheet containing $x$. It fixes $QA$,
commutes with every $q$, and its endpoint
maps $x$ to the corresponding local lift
of $y$. The compact quotient-support proof
above gives its continuity with inverses
in the original closed-disk topology.
Denote this actual local section by $\sigma_y$,
now with $y$ denoting the upstairs point,
and choose $\sigma_x=1$.

The evaluation map is onto even from its
identity path component. Given any other
$x'\in O\setminus QA$, choose a path from
$x$ to $x'$ in that connected punctured plane.
Cover its compact projected path by finitely
many quotient coordinate disks avoiding $\pi A$.
Their short supported point-moving isotopies
compose to move the point along that path.
The upstairs path is the chosen one by
uniqueness of the original covering lift.
All supports lie in one compact quotient
set, so their lifted family extends by
(C.81.9) and fixes $QA$. Its endpoint sends
$x$ to $x'$ in the asserted original group.

Set $F=\mathscr C_Q(A\cup\{x\})$. The local
bundle chart and its inverse are exactly
\[
 \begin{split}
 h&\longmapsto\bigl(h(x),\sigma_{h(x)}^{-1}h\bigr),\\
 (y,k)&\longmapsto\sigma_y k\qquad(k\in F).
 \end{split}                                                   \tag{C.81.11}
\]
They are continuous inverse maps in the
specified topology. Translation supplies
the other charts. Thus (C.81.3) is a locally
trivial principal bundle. Its metrizable
manifold base makes it numerable, so the
usual based homotopy lifting and exact
sequence apply to these actual maps.

\emph{The exact positive-degree effect of fixing finitely many orbits.}
For the first orbit the base in (C.81.3)
is the contractible plane $O$. Its homotopy
exact sequence gives an isomorphism on
every positive homotopy group after fixing
that first orbit. Suppose $|Q|\geq2$ and
$A$ is now nonempty. The surface
$M=O\setminus QA$ has at least two punctures.
It is aspherical: a compact sphere image
lies in a compact planar subsurface with
nonempty boundary; that subsurface retracts
onto a finite graph and has no homotopy
in degrees at least two.

Moreover $\operatorname{Homeo}(M)_0$ is
contractible by \cite[Theorem~1.1]{A_Yagasaki2000}.
Its exceptions are explicitly excluded here.
The surface is oriented, and forgetting all
but two punctures gives a surjection onto
a free group of rank two: the two small
meridians and their connecting paths can
be chosen to avoid the remaining locally
finite punctures. Thus it is neither the
plane nor an annulus, and cannot be the
nonorientable exception.

Every loop in $\mathscr C_Q(A)$ restricts
continuously, with inverses, to a loop in
$\operatorname{Homeo}(M)_0$. Its evaluation
at $x$ factors through that restriction,
and hence has zero class in $\pi_1(M,x)$.
Asphericity and the exact sequence of
(C.81.3) now give
\[
 \pi_n(\mathscr C_Q(A\cup\{x\}),1)
       \xrightarrow{\ \cong\ }
 \pi_n(\mathscr C_Q(A),1)
                              \qquad(n\geq1).
                                                               \tag{C.81.12}
\]
For $n=1$, injectivity uses $\pi_2(M)=0$
and surjectivity uses the just proved zero
evaluation homomorphism. The higher degrees
use both adjacent vanishing groups of $M$.
Induction proves the asserted finite-set
comparison for $|Q|\geq2$. For $Q=\{1\}$,
\cite[Lemma~2.1(ii)]{A_Yagasaki2000} applies
directly on the compact disk with its
boundary and the finite fixed set. The
identity components have zero positive
homotopy groups on both sides, giving
the same inclusion isomorphism.

This proves that any sphere class has a
representative fixing any specified finite
number of orbit labels, through an actual
based homotopy in the original relative
group. No uniform choice or limit over
an infinite configuration has been used.

\emph{The explicit comparison with Whitney topology.}
Descent $h\mapsto\bar h$ from $\mathscr C_Q$
is continuous into the compact-open and
inverse topology of $\operatorname{Homeo}(N)$:
local covering sections prove joint continuity
on each compact quotient set, also for
the descended inverse maps.
Assume now that $N$ is noncompact. Choose
points $y_j$ escaping every compact set
and pairwise disjoint, locally finite
coordinate disks $V_j$ around them.
Choose lifts $x_j$. Shrink each disk so
that its lift $U_j$ at $x_j$ has compact
closure, disjoint $Q$-translates, and is
contained in a $d_Q$-ball of radius less
than $1/(4j)$. Then every $qU_j$ has
original $d$-diameter less than $1/j$.

Choose a disk homeomorphism $b_j$ supported
in $V_j$, isotopic to the identity there,
with $b_j(y_j)\ne y_j$. Lift its supported
isotopy from the identity. The endpoint
$h_j$ preserves each lifted disk $qU_j$
and is the identity elsewhere. It has
the actual circle-fixed extension already
proved. Consequently
\[
 \begin{split}
 \sup_{z\in\overline D}d(h_jz,z)&<1/j,\\
 \sup_{z\in\overline D}d(h_j^{-1}z,z)&<1/j.
 \end{split}
                                                               \tag{C.81.13}
\]
Thus $h_j\to1$ in $\mathscr C_Q$ and,
by (C.81.7), in its exact uniform $d_Q$
topology as well.

Put $a_j=\rho(b_j(y_j),y_j)>0$. There is
a continuous positive $\epsilon$ on $N$
with $\epsilon(y_j)<a_j/2$: on each
locally finite disk use a bump function
interpolating between that positive chosen
value at $y_j$ and the value $1$ outside
the disk. These locally finite formulas
give a positive continuous global function.
The one Whitney neighborhood
\[
 \{b:\rho(b(y),y)<\epsilon(y)\text{ for every }y\in N\}
                                                               \tag{C.81.14}
\]
contains none of the $b_j$. Hence the
specified descent homomorphism, even on
the subgroup of actual lifts of compactly
supported quotient isotopies with its
induced $\mathscr C_Q$ topology, is not
continuous to this Whitney topology.

The definition and fragmentation results
in \cite[p.~395 and \S2, Proposition~2.9,
Corollary~2.10]{C_LechMichalikRybicki2018}
use that majorant topology and compactly
supported isotopies. Formula (C.81.13)
and the single excluded neighborhood
(C.81.14) prove the exact comparison
obstruction to their direct use for an
arbitrary original continuous family.
Likewise the metric-cover deformation
in \cite[Theorem~1.1]{C_YagasakiUniform2014}
has a compact quotient hypothesis and
does not itself establish preservation
of our named relative image. The uniform
controlled fragmentation and the coherent
infinite-configuration limit remain the
stated general discrete-case requirements.
\end{proof}

\begin{claim}{C.82}{Proved}{A null lifted triangulation and its actual skeleton stabilizer}
Retain the actual disk, discrete operator group,
transitive relative subgroup, quotient covering and
metrics of C.81. There is a locally finite
triangulation $\mathcal T$ of $N=O/Q$ whose full
lift $\widetilde{\mathcal T}$ through the original
covering $\pi:O\to N$ has the following property:
for every $\epsilon>0$, only finitely many lifted
closed vertex stars have $d$-diameter at least
$\epsilon$. In particular, this holds for the
lifted closed triangles. Every such star is
mapped injectively by $\pi$.

Write $\Lambda=|\widetilde{\mathcal T}^{(1)}|$
for the actual lifted one-skeleton and set
\[
 \mathscr S_{Q,\Lambda}
   =\{h\in\mathscr C_Q:h|_\Lambda=\id\}.
                                                               \tag{C.82.1}
\]
Choose one closed lifted triangle $D_\sigma$
for each two-simplex $\sigma$ of $\mathcal T$.
Restriction and $Q$-equivariant pasting are
mutually inverse topological group isomorphisms
\[
 \mathscr S_{Q,\Lambda}
    \ \cong\ 
 \prod_{\sigma\in\mathcal T^{(2)}}
       \operatorname{Homeo}(D_\sigma\mathbin{\mathrm{rel}}
                                      \partial D_\sigma).
                                                               \tag{C.82.2}
\]
Here $\mathcal T^{(2)}$ indexes the two-simplices,
the product has the product topology, and the
left side retains the original closed-disk
topology of maps and inverses. Coordinatewise
Alexander deformation gives a continuous
contraction of this entire actual subgroup to
$1$, fixing $1$ throughout.

The conclusion concerns the pointwise skeleton
stabilizer. A deformation of every original
based sphere family in $\mathscr C_Q$ into that
stabilizer, with joint original-boundary control,
is not supplied by this statement or by the
finite-label comparisons of C.81.
\end{claim}
\passage{Section 23, folios 140-141, scans 279-281; Part C(b), exact null lifted triangulation and contractible skeleton stabilizer in the remaining discrete-operator level-zero question}


\begin{proof}
\emph{The quotient charts and the chosen subdivision.}
The quotient $N$ is a second-countable surface,
possibly nonorientable. Indeed the original
free proper discrete action gives the actual
covering charts, and the open quotient of the
countable base of $O$ is a countable base of $N$.
Rad\'o's original triangulation theorem,
\cite[\S3, Hilfssatz 2, pp.~111--114]{C_RadoTriangulation},
gives a locally finite triangulation of this
surface. Its proof explicitly establishes
local finiteness before subdividing the finite
polygons into triangles. We retain the
homeomorphism from its geometric realization
onto the actual $N$.

Choose a continuous proper function
$a:N\to[0,\infty)$. To recall its construction,
take compact sets $K_j\subset\operatorname{int}K_{j+1}$
exhausting $N$, and continuous functions
$\varphi_j$ equal to zero on $K_j$ and one
outside $\operatorname{int}K_{j+1}$, with values
in $[0,1]$. The locally finite sum
$a=\sum_j\varphi_j$ is proper: outside $K_{m+1}$
its first $m$ summands equal one. In the compact
case the constant function zero suffices.

For each $y\in N$ choose a relatively compact
open coordinate disk $V_y$ containing $y$ so
small that every sheet above $V_y$ is mapped
isometrically onto $V_y$ for the metrics
$d_Q,\rho$ of C.81, and
\[
 \begin{gathered}
 |a(z)-a(y)|<1\quad(z\in V_y),\\
 \operatorname{diam}_{\rho}(V_y)<\frac1{2+a(y)}.
 \end{gathered}
                                                               \tag{C.82.3}
\]
Such disks are obtained by shrinking the
actual isometric covering neighborhoods of
C.81. No coordinate on the original circle
is changed in this step.

Apply Haver's subdivision lemma with $k=0$
to the initial triangulation, its actual map
to $N$, and the cover $\{V_y\}$
\cite[Lemma 1, p.~458]{C_Haver1975}.
It gives a locally finite subdivision
$\mathcal T$ for which every closed vertex
star is contained in one $V_y$. The notation
$N(v,\mathcal T)$ in that lemma is the union
of the closed simplices meeting the vertex,
so it is precisely the closed star used here.

Lift this actual triangulation through $\pi$.
On each covering chart the lifted simplices
and all their incidences are the pullbacks
under its inverse section. On overlapping
charts these pullbacks agree by covering
uniqueness. They therefore give a triangulation
$\widetilde{\mathcal T}$ of $O$, with the
original $Q$-action by simplicial deck maps.
Every lifted closed vertex star lies in one
sheet over a $V_y$, so $\pi$ restricts to a
homeomorphism on that star. If its quotient
vertex is $v$, then (C.82.3) gives
\[
 \operatorname{diam}_{d_Q}
       \bigl(\overline{\operatorname{st}}(\widetilde v)
                                             \bigr)
       <\frac1{1+a(v)}.
                                                               \tag{C.82.4}
\]
The lift is locally finite, directly from
the covering charts and local finiteness
downstairs. Alternatively, a compact subset
of $O$ projects to a compact set meeting
finitely many quotient stars; properness
allows only finitely many translates of
their compact chosen lifts to meet it.

\emph{Nullity in the original closed disk.}
We use the precise compact-core conclusion
of C.81. For a compact set $C\subset O$,
choose a compact family $T$ of actual
boundary-fixed commuting homeomorphisms
with $C\subset T(x_0)$. Such a family exists
by the compact-covering lemma applied to
the open evaluation map of the closure of
the transitive subgroup. For $z,z'\in C$
choose $s,t\in T$ with $z=sx_0$, $z'=tx_0$.
For every original operator $q\in Q$,
\[
           qz'=t s^{-1}(qz),
       \qquad qC\subset TT^{-1}(qz).
                                                               \tag{C.82.5}
\]
The compact family $TT^{-1}$ consists of
actual maps fixing every point of the
original circle. Hence
\[
 z_j\in C,\quad q_jz_j\longrightarrow b\in\partial\overline D
       \quad\Longrightarrow\quad
            q_jC\longrightarrow\{b\}
             \text{ uniformly in }d.
                                                               \tag{C.82.6}
\]
For completeness, failure of uniformity
would give $h_j\in TT^{-1}$ with
$d(h_jq_jz_j,b)$ bounded away from zero.
A convergent subsequence of $h_j$ in the
actual map topology contradicts $h_jb=b$.
No extension of $q_j$ across the circle
appears in this argument.

Suppose infinitely many distinct lifted
closed stars have diameter at least
$\epsilon>0$. Choose a point $x_j$ in each.
Compactness of $\overline D$ gives a
convergent subsequence. Its limit cannot
lie in $O$, because the lifted stars are
locally finite there. Thus $x_j\to b$ on
the original circle.

Let $v_j$ be their quotient vertices.
An unbounded subsequence of $a(v_j)$
has a further subsequence tending to
infinity. Equation (C.82.4), together
with $d\leq d_Q$, contradicts the lower
diameter bound on that subsequence.
Otherwise the $v_j$ lie in one compact
sublevel of $a$. Local finiteness gives
only finitely many such vertices.
After taking a subsequence, the stars
are therefore $q_jC$ for one chosen
compact lifted star $C$.
Writing $x_j=q_jz_j$ with $z_j\in C$,
(C.82.6) gives the same contradiction.
This proves the asserted finite-large-star
property.

Every closed triangle is contained in
each of its three closed vertex stars.
Only finitely many triangles belong to
any one star. The finite-large-star
property consequently implies the
finite-large-triangle property. It also
gives the exact boundary form that will
be used below: for any lifted triangles
$E_j$ and points $x_j\in E_j$,
\[
 x_j\longrightarrow b\in\partial\overline D
       \quad\Longrightarrow\quad
       \sup_{z\in E_j}d(z,b)\longrightarrow0.
                                                               \tag{C.82.7}
\]
Indeed each fixed finite list of compact
triangles is eventually avoided, and all
others have arbitrarily small diameter.

\emph{The actual product and both inverse maps.}
Choose one lift $D_\sigma$ for each
quotient triangle. A deck transformation
preserving $D_\sigma$ is the identity:
$\pi$ is injective there, so that
transformation fixes its points, and
the original action is free.

For a tuple $(h_\sigma)$ of the factors
in (C.82.2), define $\mathcal P(h_\sigma)$
on $\overline D$ by
\[
 \begin{cases}
   q h_\sigma q^{-1}&\text{on }qD_\sigma,
                           \quad q\in Q,\\
   \id&\text{on }\Lambda\cup\partial\overline D.
 \end{cases}
                                                               \tag{C.82.8}
\]
The first prescription is unambiguous
on every open face. On common edges
and vertices all prescriptions are
the identity. Local finiteness gives
continuity on $O$ by finite pasting.
Equation (C.82.7) gives continuity at
the original circle: the image of a
point in a triangle stays in that same
triangle. Replacing every $h_\sigma$
by its inverse gives a continuous map
by exactly the same proof, and the two
composites are the identity face by
face and on the skeleton and circle.
Thus (C.82.8) is an actual closed-disk
homeomorphism. Reindexing its formula
by left multiplication in $Q$ proves
commutation with every original operator.

Conversely, an element $h$ of (C.82.1)
preserves each closed lifted triangle.
Its Jordan boundary is fixed pointwise.
Of the two complementary components
in $O$, only its interior has compact
closure in $O$. An interior homeomorphism
preserves this property, so it cannot
exchange the two components. Restriction
therefore defines
\[
 \mathcal R(h)=(h|_{D_\sigma})_\sigma.
                                                               \tag{C.82.9}
\]
The commutator equations imply the
formula (C.82.8) on every other lift.
It follows that $\mathcal P\mathcal R$
and $\mathcal R\mathcal P$ are the
identity. Both maps respect the group
operations. This proves the asserted
group isomorphism with its exact maps.

Restriction $\mathcal R$ is continuous
for maps and inverses, since each
$D_\sigma$ is compact and invariant.
For continuity of $\mathcal P$, fix
$\epsilon>0$. Only finitely many actual
lifted triangles have diameter at least
$\epsilon/4$. On every other triangle,
two pasted maps, and separately their
inverses, take each point into the same
set of diameter less than $\epsilon/4$.
On the finitely many large triangles,
uniform convergence of the corresponding
finitely many factors gives uniform
convergence after their specified
conjugations by $q$. This uses continuity
of that fixed $q$ on a compact triangle,
and requires no boundary extension of
$q$. The finitely many factor neighborhoods
therefore control both uniform distances
on all of $\overline D$. This is exactly
continuity from the product topology
to the original map-and-inverse topology.
It proves (C.82.2) as a topological
group isomorphism.

\emph{A contraction in the actual topology.}
Choose a homeomorphism from the standard
closed unit disk to each $D_\sigma$.
In that coordinate, for a boundary-fixed
disk homeomorphism $h$ put
\[
 A_r(h)(z)=
 \begin{cases}
      r h(z/r),&|z|\leq r,\quad r>0,\\
      z,&|z|\geq r,
 \end{cases}
 \qquad A_0(h)=\id.
                                                               \tag{C.82.10}
\]
The two formulas agree at $|z|=r$,
and $A_r(h)^{-1}=A_r(h^{-1})$.
For positive $r$, evaluation is jointly
continuous, including across that seam;
compactness of the standard disk gives
the required compact-open continuity.
At $r=0$ the map and inverse displacements
are each at most $2r$, uniformly in $h$.
This proves joint continuity for the
full map-and-inverse group topology.
Also $A_1(h)=h$ and $A_r(\id)=\id$.

Take the product of these deformations
and transport it through the mutually
inverse maps $\mathcal P,\mathcal R$.
A map into a product is continuous
exactly when its coordinate maps are
continuous. Thus this gives a continuous
map $[0,1]\times\mathscr S_{Q,\Lambda}
\to\mathscr S_{Q,\Lambda}$ from the
constant identity map to the identity
map, fixing the group identity throughout.
All original operator labels, the
pointwise skeleton condition, and the
original circle condition remain exact.

The construction proves an actual
infinite product comparison because
the nullity estimate supplies uniform
control outside finitely many triangles.
It does not assert an infinite-configuration
lifting theorem for restriction from
$\mathscr C_Q$. In particular, C.81's
finite-orbit evaluation bundles have
not been passed to an inverse limit.
The deformation of a general compact
parameter family to this skeleton
stabilizer, with controlled vertex and
edge motions, remains the next step
in the discrete-operator part of C.12.
\end{proof}

\begin{claim}{C.83}{Proved}{An actual equivariant infinite-evaluation path without a closed-disk lift}
There is an actual closed disk and a closed,
discrete, free proper interior operator group
$Q\cong\mathbb Z$ satisfying every hypothesis of
C.81, including an actual transitive
boundary-fixed subgroup of its centralizer,
with the following property.

For a specified ordered, locally finite,
$Q$-invariant vertex set $V$, let
$\mathscr E_V$ be the actual image of
simultaneous labelled evaluation
\[
 \operatorname{ev}_V:\mathscr C_Q\longrightarrow\mathscr E_V.
                                                               \tag{C.83.1}
\]
Give $\mathscr E_V$ the subspace topology of
ordered coordinate convergence together with
vague convergence of counting measures in the
specified interior plane coordinate.
The map (C.83.1) is continuous and onto.
There is an explicit based null-homotopic loop in
$\mathscr E_V$ which has no continuous lift
to $\mathscr C_Q$ starting at $1$.
Every point of that loop is the evaluation
of a displayed actual boundary-fixed,
$Q$-commuting homeomorphism. All original
$Q$-labels are retained.

The vertex set can be taken to be the full
lift of a quotient triangulation having the
null-star property of C.82. Thus, for these
exact source and target topologies, infinite
evaluation has neither the path-lifting
property nor the Serre or Hurewicz fibration
property. The proof does not assert a
nontrivial homotopy class in the actual
centralizer.
\end{claim}
\passage{Section 23, folios 140-141, scans 279-281; Part C(b), exact infinite-configuration topology and lifting obstruction in the remaining discrete-operator family-to-skeleton step}


\begin{proof}
\emph{The actual disk, operators and transitive source.}
Let $\overline D$ be the closed Euclidean
unit disk, let $O=\operatorname{int}\overline D$,
and use its original Euclidean distance $d$.
The specified interior coordinate and inverse are
\[
 \begin{gathered}
 j:\mathbb R^2\longrightarrow O,\qquad
       j(z)=\frac{z}{1+\|z\|},\\
 j^{-1}:O\longrightarrow\mathbb R^2,\qquad
       j^{-1}(w)=\frac{w}{1-\|w\|}.
 \end{gathered}                                                  \tag{C.83.2}
\]
For $a\in\mathbb R^2$ write $T_a(z)=z+a$.
Set $q=jT_{(1,0)}j^{-1}$ and $Q=\langle q\rangle$,
with exactly these powers as its labels.
The conjugacy (C.83.2) proves freeness and
properness: two compact subsets of the plane
meet under only finitely many integer horizontal
translations. It also proves that $Q$ is a
closed discrete subgroup of $\operatorname{Homeo}(O)$.
Indeed, evaluation at $j(0)$ of any convergent
sequence of its elements bounds the integer
translations in the plane, and a convergent
subsequence of integers is eventually constant.

Each $jT_a j^{-1}$ extends by the identity
on the original circle. We will use the
following explicit estimate repeatedly.
For $z,v\in\mathbb R^2$,
\[
 \|j(z+v)-j(z)\|
       \leq\frac{2\|v\|}{1+\|z\|}.
                                                               \tag{C.83.3}
\]
To prove it, put $r=\|z\|$ and $r'=\|z+v\|$
and write the difference as
$v/(1+r)+(z+v)((1+r')^{-1}-(1+r)^{-1})$.
The reverse triangle inequality bounds the
second summand by $\|v\|/(1+r)$.
The same estimate applies to inverse
translations. Uniform bounds for $a$ on
compact parameter sets therefore give joint
continuity of the extended maps and inverses.

Let $G$ consist of these extended translations.
Their displayed composition law makes $G$
an actual subgroup of $\mathscr C_Q$,
and $G$ is transitive on $O$.
Here $\mathscr C_Q$ is exactly the group in
C.81.1, with metric
\[
 D(h,k)=
 \sup_{w\in\overline D}\|h(w)-k(w)\|
 +\sup_{w\in\overline D}\|h^{-1}(w)-k^{-1}(w)\|.
                                                               \tag{C.83.4}
\]
This is its original compact-open topology
on maps and inverses. In particular, no
interior-only topology is substituted.

\emph{The complete label set and target topology.}
In the plane put
$v_{k,l}=(k/4,l)$ for $(k,l)\in\mathbb Z^2$,
and let $V=\{j(v_{k,l})\}$.
Enumerate $\mathbb Z^2$ by increasing
$\max\{|k|,|l|\}$, using lexicographic order
within each finite shell. This fixes one
literal enumeration by the positive integers.
The operator $q$ acts on its labels by
$(k,l)\mapsto(k+4,l)$.

An ordered configuration is a sequence of
distinct points $z_i\in\mathbb R^2$ for which
every compact planar set contains only
finitely many $z_i$. Write
$\mu_z=\sum_i\delta_{z_i}$ and use
\[
 \begin{split}
 d_{\rm prod}(z,w)
    &=\sum_{i\geq1}2^{-i}\min\{1,\|z_i-w_i\|\},\\
 d_\Sigma(z,w)
    &=d_{\rm prod}(z,w)+d_V(\mu_z,\mu_w),
 \end{split}                                                    \tag{C.83.5}
\]
where $d_V$ metrizes vague convergence:
$\mu_m\to\mu$ means
$\int\varphi\,d\mu_m\to\int\varphi\,d\mu$
for every continuous compactly supported
planar function $\varphi$.
This is the sum topology specified in
\cite[Definition 2.6 and equation (2), pp.~5--6]
      {C_TehConfigurations2026}.
Only that definition is used here.

The actual evaluation is
\[
 \begin{gathered}
 \bigl(\operatorname{ev}_V(h)\bigr)_{k,l}
       =j^{-1}h j(v_{k,l}),\qquad
 \mathscr E_V=\operatorname{ev}_V(\mathscr C_Q),\\
 z_{k+4,l}=z_{k,l}+(1,0)
                \quad(z\in\mathscr E_V).
 \end{gathered}                                                  \tag{C.83.6}
\]
The last equation follows from the original
commutation relation with $q$, not from
discarding its labels. A homeomorphism of
the plane sends this grid to a locally
finite set, because its inverse sends
compact sets to compact sets. Thus every
evaluation in (C.83.6) belongs to the stated
configuration space.

We verify continuity of this exact map.
Suppose $h_m\to h$ in (C.83.4).
Each labelled coordinate converges.
For a compact planar set $K$, uniform
convergence of the inverse maps implies that
all $h_m^{-1}(j(K))$, for sufficiently large
$m$, lie in one compact subset of $O$:
take a small closed neighborhood of the
compact set $h^{-1}(j(K))$ still inside $O$.
Only finitely many vertices of $V$ lie
in that neighborhood. Hence the evaluation
of any fixed compactly supported test
function is eventually a fixed finite
sum of continuous coordinate evaluations.
This proves vague convergence and then
continuity for (C.83.5).

The grid also has the required triangulation
interpretation. Divide its rectangles of
width $1/4$ and height $1$ by the same
diagonal, and transport this plane
triangulation through $j$.
Horizontal translation by $1$ preserves it.
Each closed vertex star has horizontal
width $1/2<1$, so its projection to
$O/Q$ is injective. The planar stars have
one common diameter bound. Formula
(C.83.3) shows that their images have
diameters tending to zero as their
vertices escape in the plane; only
finitely many vertices lie in any fixed
compact set. This is precisely the
finite-large-star property used in C.82.

\emph{The actual up-and-back paths.}
Define the continuous tent function
$b(u)=\max\{0,1-4|u|\}$.
For $n\geq1$ and $s\in[0,1]$ put
\[
 \begin{split}
 F_{n,s}(x,y)&=(x+sn\,b(y-n),y),\\
 F_{n,s}^{-1}(x,y)&=(x-sn\,b(y-n),y).
 \end{split}                                                    \tag{C.83.7}
\]
These formulas prove both inverse identities
and commutation with every $T_{(m,0)}$.
Their displacement is at most $n$.
Consequently (C.83.3) extends
$H_{n,s}=jF_{n,s}j^{-1}$ by the identity
to the original circle, with a jointly
continuous inverse. For each fixed $n$,
$s\mapsto H_{n,s}$ is continuous in
(C.83.4). To check the uniform assertion,
use (C.83.3) outside one large planar disk
and ordinary uniform continuity on its
compact closure. The identical argument
works for the displayed inverse.
Thus $H_{n,s}\in\mathscr C_Q$ and
$H_{n,0}=1$.

Only row $l=n$ of the vertex set moves:
\[
 \bigl(\operatorname{ev}_V(H_{n,s})\bigr)_{k,l}
  =
  \begin{cases}
     (k/4+sn,n),&l=n,\\
     (k/4,l),&l\ne n.
  \end{cases}                                                    \tag{C.83.8}
\]
These equations retain every label in
(C.83.6), including labels in the same
operator orbit.

Let $I_n=[2^{-n-1},2^{-n}]$ and write
\[
 \begin{gathered}
 \lambda_n(t)=2^{n+1}t-1,\qquad
 s_n(t)=2\min\{\lambda_n(t),1-\lambda_n(t)\},
       \quad t\in I_n,\\
 t_n=3\cdot2^{-n-2}.
 \end{gathered}                                                  \tag{C.83.9}
\]
The amplitude is zero at both endpoints
of $I_n$ and one at $t_n$.
Define
\[
 c(t)=
 \begin{cases}
  \operatorname{ev}_V(H_{n,s_n(t)}),&t\in I_n,\ n\geq1,\\
  \operatorname{ev}_V(1),&t=0\text{ or }1/2\leq t\leq1.
 \end{cases}                                                    \tag{C.83.10}
\]
The formulas agree at all overlapping
interval endpoints.

We check both components of the topology
at $t=0$. Any finite collection of fixed
labels avoids the moving row once $n$
is sufficiently large. Its coordinates
then equal their original values for
every $s\in[0,1]$; the remaining product
metric terms are bounded by the tail of
$\sum_i2^{-i}$.
For a compact planar set $K$, its
$y$-coordinates are bounded. For all
sufficiently large $n$, the moving row
and all its images are disjoint from
$K$, since (C.83.7) preserves $y$.
Every test function supported in $K$
therefore has exactly its original
counting-measure evaluation, uniformly
in $s$. This proves vague convergence.
Equivalently, for a metric defined by
a weighted countable test family, first
fix finitely many supports and then
bound the remaining summands by their
weight tail. It proves continuity in
$d_\Sigma$ at zero. Continuity on each
$I_n$ and at its other endpoints has
already been proved. Thus (C.83.10) is
a continuous based loop in the actual
subspace $\mathscr E_V$.
Its based null-homotopy in this same
subspace is the exact map
\[
 C(r,t)=
 \begin{cases}
  \operatorname{ev}_V(H_{n,(1-r)s_n(t)}),
                           &t\in I_n,\ n\geq1,\\
  \operatorname{ev}_V(1),&t=0\text{ or }1/2\leq t\leq1,
 \end{cases}
 \qquad 0\leq r\leq1.
                                                               \tag{C.83.11}
\]
The preceding finite-coordinate and
test-support arguments are uniform in
$r$, since every displayed amplitude
still belongs to $[0,1]$. Together with
continuity on each parameter rectangle
and equality at its interval endpoints,
they prove joint continuity of $C$.
Its faces are $C(0,t)=c(t)$,
$C(1,t)=\operatorname{ev}_V(1)$ and
$C(r,0)=C(r,1)=\operatorname{ev}_V(1)$.

\emph{The lower bound for every proposed lift.}
Suppose a continuous map
$\widetilde c:[0,1]\to\mathscr C_Q$
satisfied
$\widetilde c(0)=1$ and
$\operatorname{ev}_V(\widetilde c(t))=c(t)$.
At the peak time $t_n\to0$, the literal
label $(0,n)$ forces
\[
 \begin{gathered}
 \widetilde c(t_n)j(0,n)=j(n,n),\\
 \sup_{w\in\overline D}
       \|\widetilde c(t_n)w-w\|
 \ \geq\ \|j(n,n)-j(0,n)\|
 \ \geq\ \frac{n}{1+\sqrt2\,n}
 \ >\ \frac13 .
 \end{gathered}                                                  \tag{C.83.12}
\]
The penultimate inequality uses the first
coordinate alone. The last follows from
$\sqrt2<2$ and $n/(1+2n)\geq1/3$.
It holds for every homeomorphism having
the required labelled evaluation,
regardless of its behavior between
vertices. It contradicts continuity at
zero in (C.83.4).

This proves the claimed failure of
path lifting. The lifting axiom of
either named fibration, applied to a
one-point initial space, would supply
precisely that lift, so it also fails.
No contraction of an original compact
sphere family has been disproved:
the loop just constructed lives in
the explicitly stated evaluation-image
topology and has the displayed based
null-homotopy there. In particular this calculation
does not decide the remaining general
centralizer homotopy question in C.12.
\end{proof}

\begin{claim}{C.84}{Proved}{The exact radial translation topology and a contractible monotone subspace}
Retain the actual radial coordinate
$j(z)=z/(1+\|z\|)$, its inverse, the original
closed unit disk, and the operator
$q=jT_{(1,0)}j^{-1}$ of C.83.
Conjugation $h\mapsto j^{-1}hj$ is a
topological group isomorphism from the
actual relative centralizer $\mathscr C_Q$
onto the following explicitly topologized
group $\mathscr P$.
Its elements are plane homeomorphisms
$F=(X_F,Y_F)$ commuting with $T_{(1,0)}$
and satisfying
\[
 \lim_{R\to\infty}
  \sup_{\substack{0\leq x\leq1\\ |y|\geq R}}
   \frac{\|F(x,y)-(x,y)\|}{|y|}=0.
                                                               \tag{C.84.1}
\]
Every inverse satisfies the same condition.
The exact topology on $\mathscr P$ is
given by the finite metric
\[
 \begin{gathered}
 \|F-H\|_*=
  \sup_{\substack{0\leq x\leq1\\ y\in\mathbb R}}
     \frac{\|F(x,y)-H(x,y)\|}{1+|y|},\\
 \delta_*(F,H)=\|F-H\|_*+\|F^{-1}-H^{-1}\|_*.
 \end{gathered}                                                  \tag{C.84.2}
\]
In (C.84.1) the limit is uniform over
every compact family in the actual group.

Let $\mathscr M\subset\mathscr P$ consist
of those $F$ for which, at every $x$,
$y\mapsto Y_F(x,y)$ is an increasing
homeomorphism of $\mathbb R$.
The proof gives an explicit continuous
contraction of this entire subspace,
fixing the identity, through the same
subspace. Conjugating back through $j$
gives the contraction in the original
closed-disk topology, preserving its
circle and every original operator label.
The inclusion $\mathscr M\hookrightarrow\mathscr P$
is proper, as witnessed by an explicit
periodic core rotation below.
\end{claim}
\passage{Section 23, folios 140-141, scans 279-281; Part C(b), exact actual topology and controlled contracting subspace in the radial discrete-cyclic operator model}


\begin{proof}
\emph{The exact set correspondence and its compact-family bound.}
For a commuting plane homeomorphism write
\[
 F(x,y)=(x+u_F(x,y),\,y+w_F(x,y)).
                                                               \tag{C.84.3}
\]
Both displacement functions are periodic
of period one in $x$. We first prove
necessity of (C.84.1) uniformly over a
compact family of actual disk maps.

Take any sequence $F_i$ from such a
family and points $(x_i,y_i)$ with
$0\leq x_i\leq1$ and $|y_i|\to\infty$.
Choose integers $m_i^+,m_i^-$ such that
\[
 x_i+m_i^+=y_i+e_i^+,\qquad
 x_i+m_i^-=-y_i+e_i^-,\qquad
 |e_i^\pm|\leq\tfrac12.
                                                               \tag{C.84.4}
\]
Pass to a subsequence on which the sign
of $y_i$ is constant. The two radial
input sequences approach the corresponding
two diagonal points of the original
circle. Compactness of the actual family,
whose members fix every circle point,
forces the output sequences to approach
those same points. To justify the uniform
assertion directly, any alleged failure
has a subsequence of disk maps converging
uniformly; its limiting map fixes the
specified boundary point.

Periodicity makes the two output second
coordinates equal to $Y_i=Y_{F_i}(x_i,y_i)$.
The limiting directions have nonzero
second coordinate. Thus the following
ratios are defined eventually and satisfy
\[
 \frac{X_{F_i}(x_i,y_i)+m_i^+}{Y_i}\longrightarrow1,
 \qquad
 \frac{X_{F_i}(x_i,y_i)+m_i^-}{Y_i}\longrightarrow-1.
                                                               \tag{C.84.5}
\]
Subtracting and using
$(m_i^+-m_i^-)/y_i\to2$ gives $Y_i/y_i\to1$.
Either ratio then gives
$u_{F_i}(x_i,y_i)/y_i\to0$.
Consequently the normalized norm in
(C.84.1) tends to zero along every such
sequence. This proves the required
uniform compact-family limit.
Apply the same proof to the compact
family of actual inverse maps.

Conversely, suppose a commuting plane
homeomorphism satisfies (C.84.1).
For any $\epsilon>0$ its displacement
is at most $\epsilon|y|$ outside one
horizontal band, uniformly in all $x$
by periodicity. Within that band the
displacement has a finite bound, by
continuity on its compact quotient.
It follows that
\[
 \frac{\|F(z)-z\|}{1+\|z\|}\longrightarrow0
                   \quad(\|z\|\to\infty).
                                                               \tag{C.84.6}
\]
The exact radial estimate (C.83.3)
therefore proves that $jFj^{-1}$
extends continuously by the identity
to the original circle. This extension
is a continuous bijection of the
compact disk onto itself, so its
inverse is continuous. It is an
actual member of $\mathscr C_Q$.
Applying necessity to that inverse
also proves its version of (C.84.1).
Restriction and extension by the
specified identity are mutually inverse,
and conjugation preserves multiplication
and every original power of $q$.
This proves the set and group
correspondence without assuming a
different compactification.

\emph{Both inverse topology maps.}
The expressions in (C.84.2) are finite:
on compact $y$-bands they are bounded
by periodicity, and on the tails this
follows from (C.84.1).
The displayed formula is a metric.

Suppose actual disk maps converge,
with their inverses, to an actual map.
The sequence together with its limit
is a compact family, so the uniform
version of (C.84.1) bounds both weighted
tails in (C.84.2). On each compact
rectangle $[0,1]\times[-R,R]$,
conjugation through $j$ and $j^{-1}$
gives uniform convergence in plane
coordinates. The images remain in
one compact interior set, where
$j^{-1}$ is uniformly continuous.
Combining this compact-rectangle
convergence with the tail bound proves
convergence of $\|F_i-F\|_*$.
The same proof for inverses gives
convergence in $\delta_*$.

For the other direction fix
$H\in\mathscr P$ and set
\[
 c_H=\sup_{(x,y)\in\mathbb R^2}
       \frac{1+|y|}{1+\|H(x,y)\|}<\infty.
                                                               \tag{C.84.7}
\]
The finiteness follows from
$|Y_H(x,y)|\geq|y|/2$ on sufficiently
far tails, and boundedness of the
numerator on the remaining band.
Apply (C.83.3) at the base point
$H(x,y)$ to obtain
\[
 \sup_{z\in\mathbb R^2}\|jF(z)-jH(z)\|
             \leq2c_H\|F-H\|_*.
                                                               \tag{C.84.8}
\]
The periodicity of $F-H$ permits
exactly the supremum in (C.84.2).
The inverse version of (C.84.8)
proves convergence in the original
disk map-and-inverse metric.
These two metrizable-space arguments
prove the asserted topological group
isomorphism.

\emph{Exact factorization on the monotone subspace.}
For $F\in\mathscr M$ define
\[
 \begin{gathered}
 V_F(x,y)=(x,Y_F(x,y)),\qquad
 H_F=F V_F^{-1},\\
 H_F(x,y)=(h_{F,y}(x),y),\qquad F=H_F V_F.
 \end{gathered}                                                  \tag{C.84.9}
\]
The inverse of $V_F$ is continuous.
Indeed, for a converging sequence of
$x$-coordinates and target heights,
two fixed height values bracketing
the limiting inverse bracket all
nearby inverses. Every convergent
subsequence between those brackets
must solve the same strictly monotone
equation. Shrinking the brackets
proves convergence of the inverse.
This gives an actual plane homeomorphism.
Its commutation with $T_{(1,0)}$
is literal periodicity, and its
displacement satisfies (C.84.1).
Thus $V_F\in\mathscr P$, and so
$H_F\in\mathscr P$ as well.

The formula for $H_F$ fixes the height.
For each height its first coordinate
is a real-line homeomorphism and satisfies
$h_{F,y}(x+1)=h_{F,y}(x)+1$.
It is therefore increasing.
For $0\leq r\leq1$ put
\[
 \begin{split}
 V_{F,r}(x,y)
    &=(x,(1-r)y+rY_F(x,y)),\\
 H_{F,r}(x,y)
    &=((1-r)x+rh_{F,y}(x),y),\\
 A_r(F)&=H_{F,r}V_{F,r}.
 \end{split}                                                    \tag{C.84.10}
\]
Every displayed one-dimensional
convex combination is strictly
increasing and proper, and its
inverse depends continuously on
the other coordinates by the
same bracketing argument.
The maps retain the translation
equations exactly. Their
displacements are scalar multiples
of the respective original fiber
displacements, so (C.84.1) holds
uniformly for $r\in[0,1]$.
They are actual elements of
$\mathscr P$.

The second coordinate of $A_r(F)$
is the increasing function
$(1-r)y+rY_F(x,y)$.
Thus $A_r(F)\in\mathscr M$.
The faces satisfy
$A_0(F)=1$, $A_1(F)=F$, and
$A_r(1)=1$.

\emph{Full family continuity, including inverses.}
The map $F\mapsto V_F$ is continuous
in weighted forward distance, by
its formula and (C.84.2).
Equation (C.84.8) converts this
to uniform convergence of the
actual disk maps.
For homeomorphisms of this compact
disk, uniform forward convergence
to a homeomorphism implies uniform
inverse convergence. Explicitly,
if $\omega_{B^{-1}}$ is a modulus
of continuity of the fixed inverse,
\[
 \sup\|B_i^{-1}-B^{-1}\|
     \leq\omega_{B^{-1}}
                     \bigl(\sup\|B_i-B\|\bigr).
                                                               \tag{C.84.11}
\]
This follows by applying $B^{-1}$
to $B(B_i^{-1}z)$ and $z$.
Consequently $F\mapsto V_F$ is
continuous in the full actual
group topology. The actual group
operations make $F\mapsto H_F$
continuous as well.

On either fiber subgroup the
weighted forward difference of
the interpolations at $(r,F)$
and $(s,H)$ is bounded by
$r\|F-H\|_*+|r-s|\|H-1\|_*$.
Equations (C.84.8) and (C.84.11)
therefore give joint continuity
of each factor in the original
map-and-inverse topology.
Their product is jointly
continuous. This proves the
contraction of the entire
subspace, rather than only an
individual pathwise construction.

\emph{The explicit scope witness.}
For each integer $k$, take the
disk of radius $1/10$ centered
at $(k+1/2,0)$. These disks are
pairwise disjoint. In polar
coordinates about each center
use the identical homeomorphism
\[
 (\rho,\theta)\longmapsto(\rho,\theta+\alpha(\rho)),
 \qquad
 \alpha(\rho)=
 \begin{cases}
  \pi,&0\leq\rho\leq1/20,\\
  \pi(2-20\rho),&1/20\leq\rho\leq1/10.
 \end{cases}                                                    \tag{C.84.12}
\]
Use the identity elsewhere.
The inverse subtracts the same
angle, and the boundary values
agree. Local finiteness proves
that this is a plane homeomorphism.
It commutes with the original
unit translation and its
displacement is at most $1/5$.
Hence it belongs to $\mathscr P$.
On the vertical line through
any core center it sends the
heights $-1/40,1/40$ to
$1/40,-1/40$, respectively.
It consequently lies outside
$\mathscr M$.

The contraction (C.84.10) retains
every original operator and
boundary label on its stated
subspace. C.85--C.87 below
construct the continuous cutoff
for arbitrary original families
and a contraction on the whole
group $\mathscr P$.
\end{proof}

\begin{claim}{C.85}{Proved}{A parametric cutoff and all positive homotopy groups in the radial cyclic model}
Retain the exact coordinate $j$, original disk,
operator $q=jT_{(1,0)}j^{-1}$, and topological
group $\mathscr P$ of C.84.
There is a continuous map
\[
 \begin{gathered}
 \mathcal C:\mathscr P\times[1,\infty]
                 \longrightarrow\mathscr P,\qquad
 \mathcal C_R(F)=\mathcal C(F,R),\\
 \mathcal C_\infty(F)=F,\qquad
 \mathcal C_R(\id)=\id.
 \end{gathered}                                                   \tag{C.85.1}
\]
For every finite $R$, $\mathcal C_R(F)$ agrees
with $F$ on $|y|\leq R$ and is the identity
outside a finite horizontal band. The band
can be chosen uniformly for $F$ in any
compact family and $R$ in any compact
finite interval. Moreover,
$\mathcal C_R(F)\to F$ as $R\to\infty$
uniformly on compact families in the
original closed-disk map-and-inverse topology.
Consequently every based sphere family in
the actual relative centralizer
$\mathscr C_Q$ has a based filling there:
\[
                 \pi_n(\mathscr C_Q,1)=0
                         \qquad(n\geq1).                         \tag{C.85.2}
\]
Every original boundary point and every
original integer-translation label is
retained. This statement concerns the
specified radial cyclic action.
\end{claim}
\passage{Section 23, folios 140-141, scans 279-281; Part C(b), continuous actual cutoff and all positive homotopy groups for the specified radial discrete-cyclic model}


\begin{proof}
\emph{The actual upper component and its marked conformal lift.}
Write $z=x+iy$ and put
\[
 \begin{gathered}
 E(z)=e^{2\pi iz},\qquad
 m=m_R(F)=\max_{0\leq x\leq1}|F(x,R)-(x,R)|,\\
 M=1+m,\qquad B=2(M+1)=2m+4,\qquad L=100B.
 \end{gathered}                                                   \tag{C.85.3}
\]
All these quantities refer to the actual
plane coordinates of C.84.
The proper line $\Gamma=F(\mathbb R\times\{R\})$
is periodic under $T_{(1,0)}$ and lies in
$R-m\leq y\leq R+m$.
Its upper complementary component is
$\Omega=F(\{y>R\})$. Indeed, the sublinear
displacement condition makes the image of
the positive end approach the positive end,
which identifies this component exactly.

The curve
$\gamma(x)=E(F(x,R))$, $x\in\mathbb R/\mathbb Z$,
is a Jordan curve. Equality of two values
means that their plane images differ by an
integer translation; injectivity of $F$
and its exact translation relation give
equality of the source labels modulo one.
Its winding about zero is one, since the
real part of $F(x+1,R)$ exceeds that of
$F(x,R)$ by one. Its bounded Jordan domain
$D$ therefore contains zero.
The preimage of $\gamma$ under $E$ is
exactly $\Gamma$. Points sufficiently high
above $\Gamma$ map into $D$, so
$E^{-1}(D\setminus\{0\})=\Omega$.

Let $f:\overline{\mathbb D}\to\overline D$
be the conformal Jordan-domain map normalized
by
\[
                  f(0)=0,\qquad f(1)=E(F(0,R)).                  \tag{C.85.4}
\]
Existence and uniqueness follow from the
Riemann map and its Jordan boundary
extension, followed by the unique rotation
of the source circle fixing the displayed
boundary value. The function
$g(w)=f(w)/w$, with $g(0)=f'(0)$,
is holomorphic and nonzero in the disk
and is continuous and nonzero on its
closure. Choose its continuous logarithm
by the exact equation
$\log g(1)=2\pi iF(0,R)$.
Define, for $\operatorname{Im}z\geq0$,
\[
 \begin{gathered}
 \phi(z)=z+\frac{1}{2\pi i}\log g(E(z)),\\
 d(z)=\phi(z)-z-iR.
 \end{gathered}                                                   \tag{C.85.5}
\]
The displayed logarithm exists because
$g$ is nonzero on the simply connected
closed disk; its holomorphic logarithm
in the interior extends continuously
to the boundary. Formula (C.85.5) fixes
the entire lift, including its integer
label: $\phi(0)=F(0,R)$ and
$\phi(z+1)=\phi(z)+1$.

It is a conformal bijection from the
upper half-plane onto $\Omega$.
For injectivity, equality of its values
first gives equality of $E(z)$ by the
injectivity of $f$, and then the period
relation removes the resulting integer.
For surjectivity, given $w\in\Omega$,
lift $f^{-1}(E(w))$ to the upper half-plane
and change that lift by the unique integer
making its $\phi$-image equal to $w$.
The same formulas give its continuous
boundary extension and inverse.

There is an increasing real homeomorphism
$\beta$ satisfying
\[
 \begin{gathered}
 \phi(x)=F(\beta(x),R),\qquad
 \beta(0)=0,\qquad \beta(x+1)=\beta(x)+1,\\
 \alpha=\beta^{-1},\qquad
 |\beta(x)-x|\leq1,\qquad |\alpha(x)-x|\leq1.
 \end{gathered}                                                   \tag{C.85.6}
\]
These statements follow by applying the
actual inverse of $F$ to the boundary
map. Its increasing orientation follows
from the period-one relation for a real
homeomorphism. The two bounds follow on
$[0,1]$ from the fixed endpoints and then
on all of $\mathbb R$ by periodicity.

The function $d$ in (C.85.5) is the
pullback of a holomorphic function on
the disk, continuous on its closure;
in particular the end $y=\infty$ is
its removable point $w=0$.
On the real boundary,
$|d(x)|\leq m+1=M$ by (C.85.6).
The maximum principle and the Cauchy
estimate on a radius-$y/2$ disk give
\[
             |d(z)|\leq B,\qquad
             |d'(z)|\leq\frac{2B}{y}\quad(y>0).                 \tag{C.85.7}
\]
Thus the bound includes the entire
original half-plane, without a new
boundary compactification assumption.

\emph{A cutoff which remains a homeomorphism.}
Fix the following $C^1$ function:
\[
 \chi(s)=
 \begin{cases}
  0,&s\leq1,\\
  3(s-1)^2-2(s-1)^3,&1\leq s\leq2,\\
  1,&s\geq2.
 \end{cases}
 \qquad 0\leq\chi\leq1,\quad |\chi'|\leq2.
                                                                    \tag{C.85.8}
\]
Put
\[
            \widetilde\phi(z)
             =z+iR+(1-\chi(y/L))d(z).                            \tag{C.85.9}
\]
For $0<y<L$ this is the original conformal
map $\phi$. For $y\geq L$, the operator
norm of the real derivative of its
perturbation of $z+iR$ is at most
\[
         |d'(z)|+\frac{2|d(z)|}{L}
                    \leq\frac{4B}{L}=\frac1{25}<1.              \tag{C.85.10}
\]
Hence its derivative is invertible there
and preserves orientation: the straight
segment from the identity derivative to
this derivative consists of invertible
matrices. The cutoff has matching first
derivatives at $L$ and $2L$.
It follows that $\widetilde\phi$ is a
local homeomorphism everywhere in the
open half-plane.

Its image stays in $\Omega$. Below $L$
this follows from $\widetilde\phi=\phi$;
above $L$ its height is at least
$R+L-B>R+m$, above the entire boundary
curve. It is proper as a map to $\Omega$.
Indeed, its distance from $z+iR$ is
bounded by $B$. An unbounded source
sequence therefore has unbounded image.
A bounded source sequence escaping
compact subsets of the open half-plane
has a subsequence approaching its real
boundary, whose image approaches
$\Gamma$ and cannot remain in a compact
subset of $\Omega$.

A proper local homeomorphism has an
open and closed image here and is a
covering onto that image. To recall
the latter assertion directly, a fiber
is compact and discrete, hence finite;
choose disjoint local inverse charts
at its points. Properness excludes
additional preimages over a sufficiently
small target neighborhood, giving an
evenly covered neighborhood.
Connectedness of $\Omega$ gives
surjectivity. A target of height
greater than $R+2L+B$ has exactly one
preimage: every preimage has $y>2L$,
where (C.85.9) is exactly $z+iR$.
The covering therefore has one sheet.
Its boundary extension is a
homeomorphism of the corresponding
closed half-planes as well. It is a
continuous bijection, and the same
bounded-distance estimate makes it
proper, including on the real boundary.
Its inverse is consequently continuous.

Set $\eta(y)=\min\{y,1\}$ for $y\geq0$ and
\[
 \begin{gathered}
 A(x,y)=\big((1-\eta(y))\alpha(x)+\eta(y)x,\ y\big),\\
 E^+_{F,R}(x,R+y)=\widetilde\phi(A(x,y)).
 \end{gathered}                                                   \tag{C.85.11}
\]
Each first-coordinate function of $A$
is increasing, proper, and period one.
Its inverse is continuous by the
two-sided monotone bracketing argument
of C.84. Thus $A$ is a homeomorphism
of the closed half-plane, commutes with
the original unit translation, and
has displacement at most one.
The inverse of $E^+_{F,R}$ is explicitly
the inverse of $\widetilde\phi$, followed
by $A^{-1}$ and by $(x,y)\mapsto(x,R+y)$.
It is a homeomorphism onto
$\overline\Omega$, and
\[
 \begin{gathered}
 E^+_{F,R}(x,R)=F(x,R),\\
 E^+_{F,R}(x,y)=(x,y)\quad(y\geq R+2L),\\
 |E^+_{F,R}(x,y)-(x,y)|\leq B+1=2m+5.
 \end{gathered}                                                   \tag{C.85.12}
\]
The identity above $R+2L$ uses $2L>1$,
so both cutoffs have already become
the identity there.

Let $J(x,y)=(x,-y)$. This exact reflection
commutes with $T_{(1,0)}$, and
$F\mapsto JFJ$ preserves $\mathscr P$
and its weighted metric. Define
$E^-_{F,R}=J E^+_{JFJ,R}J$ on $y\leq-R$.
The complete cutoff is
\[
 \mathcal C_R(F)(x,y)=
 \begin{cases}
   E^-_{F,R}(x,y),&y\leq-R,\\
   F(x,y),&-R\leq y\leq R,\\
   E^+_{F,R}(x,y),&y\geq R.
 \end{cases}                                                     \tag{C.85.13}
\]
The shared values on both boundary
lines agree. The images are respectively
the lower component, the middle strip
image under $F$, and the upper component
of the two actual image lines.
They cover the plane with disjoint
interiors. The inverses just described,
together with $F^{-1}$ on the middle
image, therefore paste to the actual
inverse of (C.85.13).
Every piece commutes with the original
unit translation. Outside a finite
band the map is the identity; on that
band its displacement is bounded by
periodicity. It thus belongs to
$\mathscr P$ by C.84. For $F=\id$,
the normalized conformal map is
$f(w)=e^{-2\pi R}w$, so $d=0$,
$\alpha=\id$, and every piece in
(C.85.13) is exactly the identity.

\emph{Finite-parameter continuity, with the chosen labels.}
Here is the precise conformal parameter
input. Rad\'o's original theorem
\cite[Section~1, p.~182; proof pp.~183--186]{C_Rado1923Variable}
states that normalized Riemann maps
with value zero and positive derivative
there converge uniformly on the entire
closed disk exactly when their enclosing
Jordan curves converge in Fr\'echet
distance. We use that theorem for the
specific Jordan curves $\gamma$ above.
The remaining normalization and topology
comparisons are as follows.

Take $F_i\to F$ in $\mathscr P$ and
finite $R_i\to R$. Their maps and
inverses converge uniformly on compact
plane sets by C.84. Hence
$E(F_i(x,R_i))\to E(F(x,R))$
uniformly for $0\leq x\leq1$; this
is an explicit Fr\'echet correspondence.
Rad\'o's theorem gives uniform
convergence of the positive-derivative
maps on the closed disk. Their boundary
preimages of $E(F_i(0,R_i))$ converge
as well: any convergent subsequence of
those circle points must have image
$E(F(0,R))$, whose limiting preimage
is unique. Rotating by these precise
circle points proves uniform convergence
$f_i\to f$ with normalization (C.85.4).

The functions $g_i=f_i/w$ also converge
uniformly on the closed disk. Indeed,
the difference quotient $(f_i-f)/w$
has a removable value at zero and its
maximum modulus is bounded by
$\|f_i-f\|_\infty$ on the unit circle.
Since $g$ is nonzero on the compact
disk, $g_i/g\to1$ uniformly. The small
logarithm of this ratio tends uniformly
to zero. It equals $\log g_i-\log g$
with our actual branches: this equality
holds at $w=1$ for large $i$ by
$F_i(0,R_i)\to F(0,R)$, and hence
everywhere by connectedness.
Consequently $d_i\to d$ uniformly on
the entire closed half-plane.

The quantities $m_i,B_i,L_i$ converge.
The functions $\chi(y/L_i)$ converge
uniformly for $y\geq0$ to $\chi(y/L)$:
they agree above one common finite
height, and convergence is uniform on
the remaining compact interval.
Thus $\widetilde\phi_i\to\widetilde\phi$
uniformly on the closed half-plane.
On $[0,1]$ the exact boundary inverses
\[
       \beta_i(x)=\operatorname{pr}_1
                    F_i^{-1}(\phi_i(x))                          \tag{C.85.14}
\]
converge uniformly to $\beta$.
Monotone bracketing and periodicity then
give uniform convergence of
$\alpha_i=\beta_i^{-1}$ to $\alpha$
on $\mathbb R$ in their displacement
coordinates. The maps $A_i$ converge
uniformly to $A$ on the closed
half-plane. Each fixed
$\widetilde\phi$ is uniformly continuous
there: its perturbation is periodic
in $x$, continuous on one compact
vertical strip modulo that period,
and zero above $2L$.
It follows that
$\widetilde\phi_i A_i\to\widetilde\phi A$
uniformly.

The two reflected versions have the
same properties. When the interfaces
$\pm R_i$ move, the formulas agree
exactly on each interface, so their
pasting is jointly continuous in the
parameters and the source point.
All cutoffs in this convergent sequence
are the identity outside one common
horizontal band. Uniform convergence
on its compact period rectangle gives
weighted forward convergence in C.84.
The radial estimate gives uniform
convergence of the original disk maps.
Their inverses converge uniformly too:
for homeomorphisms $a_i,a$ of the
compact disk,
\[
 \|a_i^{-1}-a^{-1}\|_\infty
       \leq\omega_{a^{-1}}
                     (\|a_i-a\|_\infty),                        \tag{C.85.15}
\]
where $\omega_{a^{-1}}$ is a modulus
of continuity of the fixed inverse.
C.84 now gives convergence in the
full actual group topology. This
proves finite-$R$ continuity, including
both inverse topology maps.

\emph{The infinite interface and compact families.}
For a compact $K\subset\mathscr P$ put
\[
 a_K(R)=\sup_{\substack{F\in K,\ 0\leq x\leq1\\|y|\geq R}}
                  \frac{|F(x,y)-(x,y)|}{|y|}.
                                                                    \tag{C.85.16}
\]
C.84 proves $a_K(R)\to0$.
The two boundary displacement maxima
at $\pm R$ are at most $Ra_K(R)$.
On the middle band (C.85.13) equals
$F$. On either tail (C.85.12) and its
reflection therefore give
\[
 \sup_{F\in K}\|\mathcal C_R(F)-F\|_*
          \leq 3a_K(R)+\frac{5}{1+R}
                   \longrightarrow0.                            \tag{C.85.17}
\]
The constants $c_F$ from (C.84.7) are
bounded uniformly for $F\in K$:
the common tail bound gives
$|Y_F(x,y)|\geq|y|/2$ off one band,
and the numerator there is bounded
on the remaining band. Formula
(C.84.8) thus turns (C.85.17) into
uniform forward convergence of the
original disk maps, uniformly over
$K$. The compact family of original
inverses has a common modulus of
continuity. Applying (C.85.15) with
that common modulus proves the same
uniform assertion for inverses.

This also proves joint continuity
at $R=\infty$. For a sequence
$F_i\to F$, $R_i\to\infty$, the set
$\{F,F_1,F_2,\ldots\}$ is compact,
so the estimate just proved gives
$\mathcal C_{R_i}(F_i)\to F$ in the
actual topology. Together with the
finite-parameter argument this proves
(C.85.1). In particular compact
quotient-support elements are dense,
with the displayed continuous cutoff
maps providing the parameter comparison.

\emph{The based sphere and its exact fixed-support filling.}
Let $a:(S^n,*)\to(\mathscr C_Q,1)$,
$n\geq1$, and write $F_\xi=j^{-1}a_\xi j$.
For $0<s\leq1$ use
$j\mathcal C_{1/s}(F_\xi)j^{-1}$,
and at $s=0$ use $a_\xi$.
Formula (C.85.1) proves that this is
a continuous based homotopy in the
original source group. Its last
family $a'_\xi=j\mathcal C_1(F_\xi)j^{-1}$
is supported, in the quotient
$N=\mathbb R^2/\langle T_{(1,0)}\rangle$,
in the interior of one fixed compact
annulus $N_0$. Such an annulus exists
because the two values of $L$ in
(C.85.3) are bounded on the compact
sphere family. Choose its two boundary
circles strictly beyond that common
support band.

Apply the exact compact-support
sphere filling of C.81 to the
descended family $b_\xi$ on $N_0$.
Its hypotheses hold here: the actual
extended plane translations form a
transitive boundary-fixed group
commuting with the original $Q$, as
in C.83; the quotient support is the
specified compact annulus with its
entire boundary fixed. In the notation
of (C.81.10), this gives
$B_{0,\xi}=1$, $B_{1,\xi}=b_\xi$,
$B_{r,*}=1$. Lift $B_{r,\xi}\pi$ from
the identity, where
$\pi:\mathbb R^2\to N$ is the original
covering. Lifting the inverse family
proves the inverse equations; uniqueness
of lifts proves commutation with every
original integer translation.
At $r=1$ the difference from
$\mathcal C_1(F_\xi)$ is an integer
deck label, locally constant on $S^n$
and zero at $*$. It is zero everywhere.
The whole lifted family and its inverse
fix the complement of the same band.
C.81's original-boundary extension,
or the bounded periodic displacement
estimate here, gives their joint
closed-disk continuity.

Writing this exact lifted disk family
as $H_{r,\xi}$, a complete based
null-homotopy is
\[
 K_{s,\xi}=
 \begin{cases}
   a_\xi,&s=0,\\
   j\mathcal C_{1/(2s)}(F_\xi)j^{-1},&0<s\leq\tfrac12,\\
   H_{2-2s,\xi},&\tfrac12\leq s\leq1.
 \end{cases}                                                     \tag{C.85.18}
\]
The two middle values agree with the
same original family $a'_\xi$.
The basepoint stays the identity,
$K_{0,\xi}=a_\xi$, and $K_{1,\xi}=1$.
This proves (C.85.2) in the original
topological group, with both maps
and inverses controlled throughout.

The conclusion includes all compact
based sphere families for this exact
radial cyclic action. The construction
uses its literal period-one coordinate,
its radial boundary topology, and the
sublinear comparison of C.84. The
general discrete-operator deformation
and general C.12 remain open.
\end{proof}

\begin{claim}{C.86}{Proved}{The exact spiral centralizer comparison and its based sphere fillings}
Retain the actual radial coordinate $j$ of C.83--C.85
and the actual spiral coordinate $p$ of C.72:
\[
 j(z)=\frac{z}{1+|z|},\qquad
 p(z)=\frac{z}{1+|z|}e^{i\log(1+|z|)},\qquad
 \mathcal R(z)=iz.
                                                               \tag{C.86.1}
\]
Let $\mathscr C_{\mathrm{rad}}$ be the relative
closed-disk centralizer of $jT_1j^{-1}$, and let
$\mathscr C_{\mathrm{sp}}$ be the relative
closed-disk centralizer of the original C.72
operator $pT_i p^{-1}$. All these operators have
their specified identity extensions on the circle.
There is an actual topological group isomorphism
\[
 \begin{gathered}
 \Theta:\mathscr C_{\mathrm{rad}}
                \longrightarrow\mathscr C_{\mathrm{sp}},\\
 \Theta(h)|_{\mathbb D}
       =p\mathcal Rj^{-1}h j\mathcal R^{-1}p^{-1},
       \qquad \Theta(h)|_{S^1}=\id.
 \end{gathered}                                                   \tag{C.86.2}
\]
Its inverse is given by the inverse interior
conjugation and the same identity boundary
extension. The proof establishes both topology
maps in the original map-and-inverse topology.
In particular,
\[
             \pi_n(\mathscr C_{\mathrm{sp}},1)=0
                          \qquad(n\geq1).                        \tag{C.86.3}
\]
The interior coordinate
$p\mathcal Rj^{-1}$ in (C.86.2) has no
continuous extension at the boundary point
$1$. The group comparison and the stated
fillings are established by the exact
sublinearity and topology calculations below.
\end{claim}
\passage{Section 23, folios 140-141, scans 279-281; Part C(b), exact actual-relative-centralizer comparison and based sphere fillings for the specified spiral discrete-cyclic model}


\begin{proof}
\emph{A global spiral displacement estimate.}
For all $z,v\in\mathbb C$ one has
\[
              |p(z+v)-p(z)|
                  \leq\frac{4|v|}{1+|z|}.                       \tag{C.86.4}
\]
Put $r=|z|$, $r'=|z+v|$, and $M=|v|$.
For $M\leq(1+r)/2$, the reverse triangle
inequality gives
$\min\{1+r,1+r'\}\geq(1+r)/2$.
The mean value theorem for the logarithm
and (C.83.3) therefore give
\[
 \begin{split}
 |\log(1+r')-\log(1+r)|&\leq\frac{2M}{1+r},\\
 |p(z+v)-p(z)|
  &\leq |j(z+v)-j(z)|
       +|j(z)|\,|e^{i\log(1+r')}-e^{i\log(1+r)}|\\
  &\leq\frac{4M}{1+r}.
 \end{split}                                                     \tag{C.86.5}
\]
For $M>(1+r)/2$, both spiral values lie
in the unit disk, so their distance
is at most $2<4M/(1+r)$. This proves
the global estimate, including at zero.

We first compare the horizontal operator
$pT_1p^{-1}$ with $jT_1j^{-1}$.
The final rotation will restore the
literal vertical label of C.72.

\emph{Necessity, uniformly on compact actual families.}
Let $K$ be a compact family of original
closed-disk homeomorphisms fixing the
circle and commuting with $pT_1p^{-1}$.
Write their actual plane conjugates as
$F=p^{-1}hp$.
Each $F$ commutes with the unit horizontal
translation. Compactness and boundary
fixation give the exact uniform limit
\[
 \sup_{h\in K,\ |z|\geq R}|pF(z)-p(z)|
                              \longrightarrow0.
                                                               \tag{C.86.6}
\]
Indeed, $|p(z)|\to1$. Any failing
sequence has a subsequence on which
$p(z)$ approaches one circle point
and $h$ converges uniformly to a map
fixing that point. Its displayed
difference then tends to zero.

Take arbitrary $F_i$ from these conjugates
and $(x_i,y_i)$ with $0\leq x_i\leq1$
and $|y_i|\to\infty$. Put
\[
 U_i=\frac{X_{F_i}(x_i,y_i)-x_i}{y_i},\qquad
 V_i=\frac{Y_{F_i}(x_i,y_i)}{y_i}.
                                                               \tag{C.86.7}
\]
For any fixed real $c$ choose the actual
integer $m_i(c)$ so that
\[
 x_i+m_i(c)=cy_i+e_i(c),\qquad |e_i(c)|\leq\tfrac12.
                                                               \tag{C.86.8}
\]
The input and output at this integer
translate are exactly
\[
 \begin{gathered}
 z_i(c)=y_i(c+i)+e_i(c),\\
 F_i(z_i(c))=y_i(c+U_i+iV_i)+e_i(c).
 \end{gathered}                                                   \tag{C.86.9}
\]
Here $i$ in the complex expressions is
the imaginary unit; subscripts retain
the sequence index.

First $(U_i,V_i)$ is bounded. Suppose a
subsequence has norm tending to infinity.
The two outputs for $c=0$ and $c=1$
differ by the real integer
$m_i(1)-m_i(0)$, of absolute value at
most $|y_i|+1$. The norm of the output
at zero, divided by $|y_i|$, tends to
infinity. Applying (C.86.4) to these
two outputs shows that their spiral
images have distance tending to zero.
In contrast, the distance of the two
spiral input images tends to a positive
number. To see its exact value, define
\[
       \Psi(w)=\frac{w}{|w|}e^{i\log|w|}
                         \quad(w\ne0).
                                                               \tag{C.86.10}
\]
After removal of the common unit factor
$\operatorname{sign}(y_i)e^{i\log|y_i|}$,
the two input phases tend to $\Psi(i)$
and $\Psi(1+i)$. Their angle difference
is
$\delta=\pi/4-(\log2)/2$, with
$0<\delta<\pi/4$; the first inequality
follows from $\log2<1<\pi/2$.
Their distance is $2\sin(\delta/2)>0$.
Formula (C.86.6) would force the input
and output distances to have the same
limit, a contradiction. This proves
boundedness without presupposing a
normalized nonzero output limit.

Pass to a subsequence with
$(U_i,V_i)\to(U,V)$. For every real
$c$ with $c+U+iV\ne0$, formula (C.86.9)
and removal of the same common unit
factor give, by (C.86.6),
\[
                     \Psi(c+U+iV)=\Psi(c+i).
                                                               \tag{C.86.11}
\]
This is equality of circle-valued
functions. No equality of real
arguments, or choice of an integer
multiple of $2\pi$, has been imposed.
On every open interval where the
left argument is nonzero, differentiate
the circle-valued identity and divide
by its nonzero value. The exact
logarithmic derivative is
\[
 \frac{d}{dc}\Psi(c+U+iV)\big/\Psi(c+U+iV)
       =i\frac{c+U-V}{(c+U)^2+V^2}.
                                                               \tag{C.86.12}
\]
Thus
\[
 (c+U-V)(c^2+1)
       =(c-1)((c+U)^2+V^2).                                    \tag{C.86.13}
\]
The identity holds as a polynomial
because it holds on an open interval.
This includes the possible limit
$V=0$: only the single value $c=-U$
was omitted, and no conclusion at
that zero was needed.
Comparing the coefficient of $c^2$
gives $U+V=1$; comparing constants
then gives
$U-V+U^2+V^2=2U^2=0$.
Consequently $U=0$ and $V=1$.

Every convergent subsequence of the
bounded pair has this same limit.
Hence
\[
 \frac{|F_i(x_i,y_i)-(x_i,y_i)|}{|y_i|}
           =\sqrt{U_i^2+(V_i-1)^2}\longrightarrow0.
                                                               \tag{C.86.14}
\]
Since the original sequences were
arbitrary, this proves precisely the
uniform compact-family sublinearity
condition (C.84.1) for the spiral
relative centralizer. Apply the same
argument to the compact inverse family
to retain its inverse condition.

\emph{Sufficiency and the two inverse topology maps.}
Conversely, let $F\in\mathscr P$ as
defined in C.84. Periodicity and its
sublinearity give
$|F(z)-z|/(1+|z|)\to0$ at infinity,
as proved in (C.84.6). Formula (C.86.4)
shows that $pFp^{-1}$ extends by the
identity to the original circle.
The extension is a continuous bijection
of the compact disk and thus a
homeomorphism. It commutes with the
specified horizontal spiral operator;
its inverse is the extension of
$pF^{-1}p^{-1}$. We have proved an
exact group bijection
\[
 \begin{gathered}
 I_p:\mathscr P\longrightarrow
        C_{\mathcal D_\partial}(pT_1p^{-1}),\\
 I_p(F)|_{\mathbb D}=pFp^{-1},\qquad
 I_p(F)|_{S^1}=\id,
 \end{gathered}                                                   \tag{C.86.15}
\]
whose inverse is restriction followed
by $p^{-1}(\cdot)p$.

For its forward continuity fix $H\in\mathscr P$.
The finite constant $c_H$ of (C.84.7)
and (C.86.4), applied at $H(z)$, give
\[
 \|I_p(F)-I_p(H)\|_\infty
          \leq4c_H\|F-H\|_*.
                                                               \tag{C.86.16}
\]
The estimate uses the exact periodic
displacement and the original source
height in the definition of $\|\cdot\|_*$.
Applying it to $F^{-1},H^{-1}$ proves
continuity for the original inverses
as well.

For the inverse topology map, a
convergent sequence of actual disk
maps together with its limit is
compact. The necessity proof above
gives uniform sublinear tails for
the corresponding plane maps and
their inverses. On each compact
period rectangle, conjugation by
$p,p^{-1}$ gives uniform convergence:
the disk images stay in one compact
interior set, where $p^{-1}$ is
uniformly continuous. Combining that
convergence with the common tail
bounds gives convergence of the
weighted forward differences and
of the weighted inverse differences.
Thus the inverse of (C.86.15) is
continuous in $\delta_*$.
These metrizable-space arguments prove
the asserted topological group
isomorphism with the actual topology.

\emph{The original vertical label, the nonextendible coordinate, and the filling.}
Both $j$ and $p$ commute with the
unit complex rotation $\mathcal R$.
Its closed-disk extension $\overline{\mathcal R}(w)=iw$
conjugates the horizontal spiral
operator to the original vertical
operator of C.72. The exact C.84
isomorphism $I_j$ and (C.86.15)
therefore give
\[
 \Theta(h)=\overline{\mathcal R}\,
               I_p(I_j^{-1}(h))\,
                    \overline{\mathcal R}^{-1}.
                                                               \tag{C.86.17}
\]
This is (C.86.2) on the interior and
the identity on the boundary.
Its inverse is
$I_j I_p^{-1}(\overline{\mathcal R}^{-1}(\cdot)
\overline{\mathcal R})$.
The two inverse identities and
both continuous topology maps follow
from the specified component maps;
every original integer operator
label is preserved by the displayed
rotation and conjugations.

For completeness, the interior
coordinate itself does not extend
at the indicated boundary point.
Put $r_k=e^{2\pi k}-1$ and
$s_k=e^{(2k+1)\pi}-1$.
Both $j(r_k)$ and $j(s_k)$ tend to
$1$, whereas
\[
 \begin{gathered}
 p\mathcal Rj^{-1}(j(r_k))
        =i\frac{r_k}{1+r_k}\longrightarrow i,\\
 p\mathcal Rj^{-1}(j(s_k))
        =-i\frac{s_k}{1+s_k}\longrightarrow-i.
 \end{gathered}                                                   \tag{C.86.18}
\]
Thus the group isomorphism above
is proved by (C.86.6)--(C.86.17),
with its own actual boundary
continuity, rather than by extending
that coordinate.

Given a based sphere family in
$\mathscr C_{\mathrm{sp}}$, apply
the continuous inverse $\Theta^{-1}$,
then the exact C.85 filling, and
finally $\Theta$ at every homotopy
parameter. This is a continuous
based filling in the original
spiral relative centralizer, for
both maps and inverses, and proves
(C.86.3). It covers the entire
relative centralizer of C.72,
including families beyond its
explicitly contracted example.
The general discrete-operator case
and general C.12 remain open.
\end{proof}

\begin{claim}{C.87}{Proved}{Global contractions of the radial and spiral relative centralizers}
The entire topological groups $\mathscr P$ of
C.84, $\mathscr C_{\mathrm{rad}}$ of C.85,
and $\mathscr C_{\mathrm{sp}}$ of C.86
are contractible by continuous contractions
fixing the identity throughout.
The proof constructs one contraction on
the whole group from the parametric C.85
cutoff, a continuous support height, an
exact compact-annulus lift comparison,
and the compact boundary-relative
identity-component contraction.
Every original disk boundary point and
every specified integer operator label
is retained, including in the inverse
topology. The compact-annulus comparison
identifies its identity component exactly
with the kernel of the displayed integer
boundary-shift homomorphism.
\end{claim}
\passage{Section 23, folios 140-141, scans 279-281; Part C(b), global actual contractions of the two specified discrete-cyclic relative centralizers}


\begin{proof}
\emph{The exact compact-annulus comparison.}
Put $A=(\mathbb R/\mathbb Z)\times[-1,1]$
and let $\mathscr H$ be its
orientation-preserving homeomorphism
group fixing both boundary circles
pointwise. Each $b\in\mathscr H$
induces the identity on the fundamental
group, since the lower boundary is
fixed. It has a unique lift
$\widetilde b:\mathbb R\times[-1,1]
\to\mathbb R\times[-1,1]$ fixing
$(0,-1)$. Its restriction to the
whole lower boundary is the identity,
and uniqueness of lifts gives
$\widetilde bT_1=T_1\widetilde b$.
On the upper boundary it has the form
\[
 \widetilde b(x,-1)=(x,-1),\qquad
 \widetilde b(x,1)=(x+\tau(b),1),\qquad
                         \tau(b)\in\mathbb Z.                   \tag{C.87.1}
\]
The integer is independent of $x$
by connectedness. The canonical lift
of a composite is the composite of
the canonical lifts, so $\tau$ is
a homomorphism. The maps
\[
     b_k([x],y)=\bigl([x+k(y+1)/2],y\bigr)
                         \quad(k\in\mathbb Z)                  \tag{C.87.2}
\]
fix both boundary circles and have
$\tau(b_k)=k$; they give a literal
homomorphic section.

The canonical lift varies continuously
in the compact-open map-and-inverse
topology, uniformly on one period
rectangle. To verify this directly,
let $b_i\to b$. For large $i$, the
two circle coordinates of $b_i(z)$
and $b(z)$ differ by a unique small
real value $\delta_i(z)\in(-1/4,1/4)$,
continuous in $z\in A$ and tending
uniformly to zero. The difference
of their lifted first coordinates
minus $\delta_i$ is a continuous
integer on the connected covering
strip. It is zero on the lower
boundary and hence zero everywhere.
Their second coordinates converge
uniformly as well. Apply the same
argument to inverses. In particular
$\tau(b_i)=\tau(b)$ eventually, since
$\delta_i=0$ on the upper boundary.
Thus $\tau$ is continuous with
discrete target.

We now prove its exact kernel.
Suppose $\tau(b)=0$ and let
$\gamma(y)=([0],y)$ be the spanning
arc. Its image has a lift
$\widetilde b(0,y)$ with endpoints
$(0,-1)$ and $(0,1)$.
The straight homotopy in the convex
strip from this lifted arc to
$(0,y)$ fixes these two endpoints;
its interior stays in the interior
of the strip. Its projection proves
that $b\gamma$ and $\gamma$ are
homotopic with their endpoints
fixed. Epstein's actual relative
arc theorem
\cite[Theorem~3.1, pp.~89--90; Appendix~A1--A2 and their proofs, pp.~105--106]{C_Epstein1966}
gives an ambient isotopy fixed
pointwise on both boundary circles
carrying $b\gamma$ to $\gamma$.
The Appendix supplies the stated
topological arc case: A2's proof
first makes the arc piecewise linear
near its endpoints and then applies
the construction in the proof of A1.

Any remaining parametrization of
$\gamma$ is an increasing interval
homeomorphism $\rho$ fixing its
endpoints. In a rectangular
neighborhood with coordinates
$(u,y)$ and $\gamma=\{u=0\}$,
choose $0\leq\zeta\leq1$ with
$\zeta(0)=1$ and $\zeta=0$ on
the two side edges. The formula
\[
       (u,y)\longmapsto
       \bigl(u,(1-s\zeta(u))y+s\zeta(u)\rho^{-1}(y)\bigr)
                              \quad(0\leq s\leq1)               \tag{C.87.3}
\]
is an isotopy of that rectangle,
fixed on its entire boundary,
and has a continuous inverse by
monotone bracketing. Extend it by
the identity. After this correction
the composed annulus map fixes
$\gamma$ pointwise. Cutting along
$\gamma$ produces a disk map
fixing its boundary pointwise.
The disk Alexander isotopy descends
to an isotopy of the annulus
relative to its boundary and
$\gamma$. Thus $b$ belongs to the
identity path component of
$\mathscr H$. Conversely the
continuous integer $\tau$ vanishes
on the identity component. We have
proved the exact equality
\[
                         \ker\tau=\mathscr H_0.                 \tag{C.87.4}
\]
The arc isotopies in this argument
establish membership in that
component; no parameter choice
of those isotopies is used below.

Let
\[
 \mathscr E=\{F\in\mathscr P:
                    F(x,y)=(x,y)\text{ for }|y|\geq1\}.
                                                               \tag{C.87.5}
\]
Restriction and descent identify
$\mathscr E$ with $\mathscr H_0$.
The actual inverse map
$\Lambda:\mathscr H_0\to\mathscr E$
takes the lower-fixed lift and
extends it by the identity outside
the strip. Equation (C.87.4) fixes
the upper line as well, so this
pasting is an actual plane
homeomorphism. Its inverse is
$\Lambda(b^{-1})$ and its
translation labels are unchanged.
Its displacement is bounded and
periodic, so it lies in
$\mathscr P$.
These are inverse topological
group maps. Restriction and
descent are continuous on the
compact period rectangle by
C.84. The canonical lift
continuity just proved, together
with its identity extension,
gives continuity of $\Lambda$
in both weighted terms of
(C.84.2).

The compact boundary-relative
identity component is contractible
by
\cite[Lemma~2.1(ii), pp.~2--3]{A_Yagasaki2000}.
Here the whole nonempty boundary
of the compact annulus is fixed
pointwise, as required there.
Choose one continuous contraction
$C_s$ of $\mathscr H_0$ with
$C_0(b)=b$ and $C_1$ constant.
Replacing it by
$C_s(b)C_s(1)^{-1}$ fixes the
identity throughout and makes
the last map the constant
identity. Transport it through
$\Lambda$ and denote the resulting
single continuous contraction
on the whole of $\mathscr E$ by
\[
 \widehat\Gamma:[0,1]\times\mathscr E\to\mathscr E,
 \quad
 \widehat\Gamma_0(G)=G,\quad
 \widehat\Gamma_1(G)=1,\quad
 \widehat\Gamma_s(1)=1.
                                                               \tag{C.87.6}
\]

\emph{The actual vertical normalizers.}
For $w>0$ put $D_w(x,y)=(x,wy)$
and
\[
 \begin{gathered}
 \mathcal N_w(F)=D_wFD_w^{-1},\\
 \mathcal N_w(F)(x,y)
       =\bigl(X_F(x,y/w),\ wY_F(x,y/w)\bigr).
 \end{gathered}                                                   \tag{C.87.7}
\]
The exact inverse parameter is
$w^{-1}$ and $D_w$ commutes with
the original $T_1$.
Sublinearity in (C.84.1) is
preserved by the displayed
coordinate formula. For fixed
$w$ the weighted metric satisfies
\[
 \delta_*(\mathcal N_w(F),\mathcal N_w(H))
                  \leq\max\{w,w^{-1}\}\,\delta_*(F,H).
                                                               \tag{C.87.8}
\]
Indeed the Euclidean operator
norm of $D_w$ is $\max\{1,w\}$,
and the change of source height
has ratio
$(1+|y|/w)/(1+|y|)\leq\max\{1,w^{-1}\}$.
Apply the same estimate to the
exact inverse conjugations of
$F^{-1}$ and $H^{-1}$.

For joint continuity in $w$
retain the original closed disk.
For $0\leq\rho\leq1$ and $|u|=1$
the actual normalizer is
\[
 \overline D_w(\rho u)
    =\frac{\rho D_wu}{1-\rho+\rho|D_wu|},\qquad
 \overline D_w(0)=0,\qquad
 \overline D_w^{-1}=\overline D_{1/w}.
                                                               \tag{C.87.9}
\]
For $\rho<1$ it is exactly
$jD_wj^{-1}$, and on the circle
it is $D_wu/|D_wu|$.
On any compact positive interval
of $w$, the denominator is
bounded away from zero. The
formula is jointly continuous,
also at $\rho=0$, and so is the
displayed inverse. It gives
joint continuity into the
closed-disk homeomorphism group.
It commutes with the specified
operator $jT_1j^{-1}$. Conjugation
by it preserves pointwise
boundary fixation. Transporting
that conjugation through the
exact C.84 isomorphism proves
joint continuity of (C.87.7)
as a map $(0,\infty)\times
\mathscr P\to\mathscr P$,
with both original inverse
topology maps included.

\emph{A single contraction on the entire group.}
Let $L_+(F)=L(F,1)$ be the
positive cutoff height of
(C.85.3), and let
$L_-(F)=L(JFJ,1)$ use the
reflection $J$ of C.85.
Both are continuous functions
of $F$. Put
\[
 \begin{gathered}
 W(F)=3+2\max\{L_+(F),L_-(F)\},\\
 G(F)=\mathcal N_{1/W(F)}(\mathcal C_1(F)).
 \end{gathered}                                                   \tag{C.87.10}
\]
The cutoff is the identity beyond
heights $1+2L_+(F)$ and
$-1-2L_-(F)$. Since each of their
absolute values is at least two
less than $W(F)$, the map
$G(F)$ is the identity near
both boundary lines of the
fixed strip and outside it.
Thus $G(F)\in\mathscr E$.
It depends continuously on
$F$ by C.85 and (C.87.7).
The exact original lift fixes
both boundary lines, so the
annulus component used in
(C.87.6) is indeed the
zero-$\tau$ component just
proved, for every $F$.

Define on the whole group
\[
 \mathcal H_s(F)=
 \begin{cases}
  F,&s=0,\\
  \mathcal C_{1/(2s)}(F),&0<s\leq\tfrac12,\\
  \mathcal N_{W(F)}
    \bigl(\widehat\Gamma_{2s-1}(G(F))\bigr),
       &\tfrac12\leq s\leq1.
 \end{cases}                                                     \tag{C.87.11}
\]
C.85 gives joint continuity
of the first part, including
at zero. At $s=1/2$ the second
part equals
$\mathcal N_{W(F)}G(F)=\mathcal C_1(F)$,
so the two parts agree exactly.
Every component map in the
last expression is continuous
on its full stated domain.
The function $W$ is locally
bounded and positive; in
particular the last part near
any fixed $F$ has support in
one common horizontal band,
and its inverse has that
same support. The normalizer
comparison already proves
the required original
map-and-inverse continuity.
Pasting gives a single
continuous map
$[0,1]\times\mathscr P\to\mathscr P$.
Its faces are
\[
 \mathcal H_0(F)=F,\qquad
 \mathcal H_1(F)=1,\qquad
 \mathcal H_s(1)=1.
                                                               \tag{C.87.12}
\]
For the last equality use
$\mathcal C_R(1)=1$ and
$\widehat\Gamma_s(1)=1$;
no normalization of $W(1)$
is required.

Conjugate (C.87.11) through
the C.84 topological group
isomorphism to obtain the
contraction on
$\mathscr C_{\mathrm{rad}}$.
Conjugate that contraction
through the two actual
topology maps $\Theta$ and
$\Theta^{-1}$ of C.86 to
obtain the contraction on
$\mathscr C_{\mathrm{sp}}$.
These formulas act on the
whole respective groups,
fix the identity, and
preserve the original
boundary and integer
operator labels throughout.
The general discrete-operator
case and general C.12 remain
open.
\end{proof}

\begin{claim}{C.88}{Proved}{The exact vanishing-displacement image and its compact-family topology}
Retain all the actual hypotheses and maps of
C.81: the original closed disk $\overline D$,
its interior $O$, a compatible metric $d$
on $\overline D$, a closed discrete free
proper operator group $Q$ on $O$, and a
transitive subgroup of the original
boundary-fixed centralizer $\mathscr C_Q$.
Write $\pi:O\to N=O/Q$ and retain
$d_Q(x,y)=\sup_{q\in Q}d(qx,qy)$.
For an interior homeomorphism $h$ commuting
with every actual $q$, put
\[
 \delta_h(\pi x)=d_Q(hx,x),\qquad
 \mathscr V_Q=
 \{h\in C_{\operatorname{Homeo}(O)}(Q):
                         \delta_h\in C_0(N)\}.                 \tag{C.88.1}
\]
Here $C_0(N)$ means continuous real-valued
functions whose absolute values are uniformly
small off compact subsets of $N$.
Then restriction and identity extension are
mutually inverse topological group maps
\[
 \begin{aligned}
 R:\mathscr C_Q&\longrightarrow\mathscr V_Q,
                       &R(a)&=a|_O,\\
 E:\mathscr V_Q&\longrightarrow\mathscr C_Q,
                       &E(h)|_O&=h,\quad
                         E(h)|_{\partial\overline D}=\id .
 \end{aligned}                                                   \tag{C.88.2}
\]
The topology on $\mathscr V_Q$ is the sum
of the uniform $d_Q$ distances of maps and
inverses. These maps are isometries for the
corresponding original closed-disk distances.
Every member of $\mathscr V_Q$ and its inverse
is automatically uniformly $d_Q$-continuous;
no uniform-continuity assumption is added
to (C.88.1).

There is also an exact compact-parameter
criterion. Let $P$ be a compact metric space
and let $h_p$ be actual equivariant interior
homeomorphisms such that $(p,x)\mapsto h_p(x)$
is continuous on $P\times O$. Their identity
extensions form a continuous $P$-family of
original closed-disk homeomorphisms and
inverses precisely when
\[
 \forall\epsilon>0\ \exists K\subset N
 \text{ compact}\quad
 \sup_{\substack{p\in P\\ \pi x\notin K}}
                    d_Q(h_p x,x)<\epsilon.                    \tag{C.88.3}
\]
In that case the inverse family has the
same uniform vanishing property. Thus
(C.88.3), with the actual interior maps
retained, expresses the original boundary
and compact-family requirements exactly.
It supplies no cutoff deformation by itself.
\end{claim}
\passage{Section 23, folios 140-141, scans 279-281; Part C(b), exact actual-boundary image and compact-parameter criterion for the remaining discrete-operator cutoff}


\begin{proof}
\emph{The actual displacement functions and algebra.}
C.81 proves that $d_Q$ induces the original
topology of $O$, that $\pi$ is the specified
covering onto the locally compact surface
$N$, and that $d_Q$ is $Q$-invariant.
Consequently $x\mapsto d_Q(hx,x)$ is
continuous. Commutation and invariance give
\[
 d_Q(hqx,qx)=d_Q(qhx,qx)=d_Q(hx,x).
                                                                  \tag{C.88.4}
\]
The quotient map therefore gives exactly
the continuous function in (C.88.1).
The map $h$ descends to the homeomorphism
$\bar h(\pi x)=\pi(hx)$, whose displayed
inverse is $\overline{h^{-1}}$.
Symmetry and the triangle inequality give
\[
 \begin{split}
 \delta_{h^{-1}}&=\delta_h\circ\bar h^{-1},\\
 \delta_{hk}&\leq\delta_h\circ\bar k+\delta_k .
 \end{split}                                                   \tag{C.88.5}
\]
A homeomorphism of $N$ and its inverse
carry compact sets to compact sets.
Thus its pullback preserves $C_0(N)$.
Equations (C.88.5) prove that
$\mathscr V_Q$ is a group, before using
the boundary extension.

\emph{Necessity, including the common tail for compact families.}
Let $a_p$, $p\in P$, be a continuous compact
family in $\mathscr C_Q$. Joint evaluation
on $P\times\overline D$ is continuous.
Since every $a_p$ fixes every boundary
point, for any $\epsilon>0$ there is a
compact set $C\subset O$ such that
\[
          d(a_p z,z)<\epsilon
                \quad(p\in P,\ z\in O\setminus C).             \tag{C.88.6}
\]
For completeness, failure would give
$p_j\to p$ after taking a subsequence
and $z_j$ escaping all interior compact
sets. A further subsequence converges
to $z\in\partial\overline D$, and joint
continuity gives $a_{p_j}z_j\to z$,
contradicting the displayed failure.

The set $K=\pi C$ is compact. If
$\pi x\notin K$, then $qx\notin C$
for every $q\in Q$. The original
commutator equations imply
\[
 \begin{split}
 d_Q(a_p x,x)
   &=\sup_{q\in Q}d(a_p qx,qx)\\
   &\leq\epsilon
       \qquad(p\in P,\ \pi x\notin K).
 \end{split}                                                   \tag{C.88.7}
\]
Starting with $\epsilon/2$ gives the
strict bound in (C.88.3).
The same proof applies to the compact
inverse family. Taking $P$ to be a
point proves that $R$ maps into
$\mathscr V_Q$.

\emph{One common control for compact quotient cores.}
We recall the exact C.81 comparison
needed for the reverse direction.
Let $C,C'\subset O$ be compact.
For a fixed $x_0\in O$, openness of
evaluation for the actual transitive
Polish source, and its compact-covering
lemma, give compact families $S,T$
of original boundary-fixed maps with
$C\subset S(x_0)$ and $C'\subset T(x_0)$.
For $c=sx_0\in C$ and $c'=tx_0\in C'$,
commutation gives
\[
          qc'=tqx_0=ts^{-1}(qc).
                                                                  \tag{C.88.8}
\]
All choices on the right lie in the
single compact actual family $TS^{-1}$.
It follows that, for $c_j\in C$ and
$q_jc_j\to z\in\partial\overline D$,
\[
                   q_j C'\longrightarrow\{z\}
                  \quad\hbox{uniformly}.                       \tag{C.88.9}
\]
Indeed joint evaluation of the compact
family $TS^{-1}$ fixes $z$ throughout.
No limiting map or original-circle
extension of the operators $q_j$
has been used.

\emph{Sufficiency and all parameter and inverse maps.}
Suppose the interior family in the
statement satisfies (C.88.3), and define
$E(h_p)$ by the literal formula (C.88.2).
Interior continuity is given. To prove
joint continuity at the original circle,
take $p_j\to p$ and $x_j\to z$ with
$z\in\partial\overline D$.
Terms already on the circle are fixed.
If the desired convergence fails for
the remaining interior terms, retain
a sequence on which it fails by a
fixed positive distance.

The surface $N$ has a compact exhaustion.
After passage to a subsequence, either
$\pi x_j$ leaves every compact subset
of $N$, or all $\pi x_j$ belong to
one compact $K\subset N$. This
dichotomy follows by taking infinitely
many terms in one exhaustion set when
there is such a set; otherwise the
sequence itself leaves every set.
In the escaping case, (C.88.3) gives
\[
             d(h_{p_j}x_j,x_j)
                  \leq d_Q(h_{p_j}x_j,x_j)\longrightarrow0,
                                                                  \tag{C.88.10}
\]
which contradicts failure.

In the compact-core case choose a
compact $C\subset O$ with $\pi C=K$.
Finitely many local covering sections
over relatively compact neighborhoods
of points of $K$ give such a $C$.
Write $x_j=q_jc_j$, $c_j\in C$.
The joint interior continuity and
compactness of $P\times C$ make
\[
                 C'=\{h_p(c):p\in P,\ c\in C\}
                                                                  \tag{C.88.11}
\]
compact in $O$. By (C.88.9),
$q_j C'\to\{z\}$ uniformly.
But $h_{p_j}x_j=q_jh_{p_j}c_j$
belongs to $q_j C'$, again a
contradiction. This proves joint
boundary continuity in every case.

Each $E(h_p)$ is bijective: its
interior restriction is a homeomorphism
onto $O$ and its boundary restriction
is the identity. The actual map
\[
 \mathcal E:P\times\overline D
       \longrightarrow P\times\overline D,\qquad
             (p,x)\longmapsto(p,E(h_p)x)                         \tag{C.88.12}
\]
is consequently a continuous bijection
between compact Hausdorff spaces.
Its inverse is continuous and has
the exact formula
$(p,y)\mapsto(p,E(h_p)^{-1}y)$.
Thus both the family and its inverses
are continuous in the original
closed-disk topology. The original
operator equations hold on $O$
unchanged. Applying the proved
necessity to this inverse family
gives its uniform $C_0$ tails.
This proof even obtains joint inverse
continuity without assuming it for
the initial interior family.

With $P$ a point, this establishes
that $E$ maps every member of
$\mathscr V_Q$ into $\mathscr C_Q$.
The two formulas (C.88.2) compose
to the identity on their exact
domains. Extension commutes with
composition and inversion, since
these identities hold on $O$ and
both sides fix the original circle.

\emph{The complete topology comparison.}
A modulus of uniform continuity
of $E(h)$ on $\overline D$ applies
to each pair $qx,qy$. Commutation
then proves uniform continuity of
$h$ for $d_Q$. Apply the same
argument to $E(h)^{-1}$.
For $h,k\in\mathscr V_Q$, the
actual equality is
\[
 \begin{split}
 \sup_{x\in O}d_Q(hx,kx)
  &=\sup_{\substack{x\in O\\ q\in Q}}
               d(E(h)qx,E(k)qx)\\
  &=\sup_{z\in\overline D}d(E(h)z,E(k)z).
 \end{split}                                                   \tag{C.88.13}
\]
The last equality uses $q=1$ and
density of $O$. Repeating it for
both inverses proves the asserted
isometry and both continuous maps.
In particular the displacement map
to $C_0(N)$ satisfies
\[
 \|\delta_h-\delta_k\|_\infty
       \leq\sup_{x\in O}d_Q(hx,kx).                             \tag{C.88.14}
\]
All image and topology requirements
are therefore expressed using the
specified covering and the original
operator labels.

The remaining general cutoff problem
is to deform an arbitrary original
based sphere family to a family with
one compact quotient support while
retaining joint interior continuity,
all equivariance equations and the
uniform tail condition (C.88.3).
C.81 fills the resulting supported
family. C.88 provides the exact
boundary criterion for such a
deformation; it does not construct
that deformation or replace its
uniform tail control by convergence
on compact subsets of $N$.
\end{proof}

\begin{claim}{C.89}{Proved}{A global section for the actual compact skeleton image}
Retain the actual hypotheses of C.81 and
the null lifted triangulation of C.82.
Write $\Lambda$ for its full one-skeleton
and put
\[
                    X=\Lambda\cup\partial\overline D .
                                                                  \tag{C.89.1}
\]
This is a compact connected subset of the
original closed disk. Define $\mathscr B$
to be the following space of embeddings
$b:X\to\overline D$, with uniform convergence
on this actual compact domain:
$b$ fixes the original circle, maps
$\Lambda$ into $O$, and satisfies
$b(qx)=q b(x)$ for every $x\in\Lambda$
and every original $q\in Q$.
For each lifted closed triangle $\Delta$,
its boundary image is positively oriented
relative to the orientation of $\partial\Delta$,
and the Jordan disk $D_b(\Delta)$ bounded
by this image satisfies
\[
                 \operatorname{int}D_b(\Delta)\cap b(X)
                               =\varnothing .                  \tag{C.89.2}
\]
Then $\mathscr B$ is exactly the image
of the original restriction map
$\mathscr C_Q\to\operatorname{Emb}(X,\overline D)$.
That map has a global continuous section
$\mathfrak s:\mathscr B\to\mathscr C_Q$
with $\mathfrak s(\id_X)=\id$.
Its construction retains all three labelled
vertices of each representative triangle
and every original deck operator, including
orientation-reversing operators.

If $\mathscr K=\mathscr S_{Q,\Lambda}$
is the pointwise stabilizer in C.82,
the maps
\[
 \begin{aligned}
 \Phi:\mathscr B\times\mathscr K&\longrightarrow\mathscr C_Q,
       &\Phi(b,k)&=\mathfrak s(b)k,\\
 \Psi:\mathscr C_Q&\longrightarrow\mathscr B\times\mathscr K,
       &\Psi(h)&=\bigl(h|_X,\mathfrak s(h|_X)^{-1}h\bigr)
 \end{aligned}                                                   \tag{C.89.3}
\]
are mutually inverse homeomorphisms of
spaces, for the stated skeleton topology
and the original map-and-inverse group
topologies. The C.82 contraction of
$\mathscr K$ supplies a strong deformation
retraction of $\mathscr C_Q$ onto
$\mathfrak s(\mathscr B)$.
Thus the remaining based-family deformation
can be expressed exactly on this compact
equivariant skeleton image. No contraction
of $\mathscr B$ is asserted here.
\end{claim}
\passage{Section 23, folios 140-141, scans 279-281; Part C(b), global section and compact-skeleton image topology for the remaining discrete-operator family deformation}


\begin{proof}
\emph{The compact skeleton and its face labels.}
The lifted triangulation is locally finite
on $O$, so $\Lambda$ is closed there.
Adding the original circle makes $X$
closed in $\overline D$, hence compact.
Its one-skeleton is connected: any two
vertices are joined by a finite edge
path in the connected triangulated
surface, and each edge belongs to that
same connected set.
Every original boundary point is in
the closure of $\Lambda$. Indeed choose
interior points tending to it and the
closed triangles containing them.
If these triangles did not leave all
finite sets of triangles, a subsequence
would lie in one compact interior
triangle. Otherwise C.82 nullity makes
their diameters tend to zero, so choosing
one point on each boundary gives the
same original boundary limit.
It follows that $X=\overline\Lambda$
is connected.

Fix once and for all an oriented disk
coordinate
$\iota:\overline D\to\overline{\mathbb D}$.
Use $d_E(x,y)=|\iota(x)-\iota(y)|$ in the
conformal constructions below. This
coordinate does not change the original
operators: each final formula is
transported back by $\iota^{-1}$.
Uniform continuity of $\iota$ and
$\iota^{-1}$ on their compact domains
gives moduli
\[
              d_E\leq\omega_1(d),\qquad
                    d\leq\omega_2(d_E),                       \tag{C.89.4}
\]
so uniform convergence in these two
metrics, including the inverse metrics
on the original group, is the same.
C.82 nullity also holds for $d_E$.

Choose one lifted closed triangle
$\Delta_\sigma$ for each quotient
two-simplex $\sigma$.
Choose its three vertices
$(v_{\sigma,0},v_{\sigma,1},v_{\sigma,2})$
in positive cyclic order in the
coordinate $\iota$.
The stabilizer of any lifted triangle
under $Q$ is trivial. Its projection
is one-to-one, so an operator preserving
the triangle fixes the unique lift
there of each labelled quotient
vertex; freeness then makes the
operator the identity.
Every lifted triangle consequently has
a unique expression $q\Delta_\sigma$.

\emph{The intrinsic face conditions and the extensions.}
For $b\in\mathscr B$, the curve
$b(\partial\Delta)$ is a Jordan curve
contained in $O$.
Its bounded Jordan disk in the
$\iota$ coordinate lies in $O$:
the disk is contained in the convex
hull of its compact boundary, which
is contained in $\mathbb D$.

The interiors of the disks
$D_b(\Delta)$ are pairwise disjoint.
To see this, the connected interior
of one such disk avoids the boundary
of every other by (C.89.2).
If two interiors meet, this connected
interior therefore lies inside the
other disk. The first boundary is
then contained in that closed disk.
Condition (C.89.2) forces it to lie
entirely in the second boundary.
A Jordan curve contained in another
Jordan curve is the whole curve:
after deleting a missing point,
a proper such subset would embed a
circle in a real interval, which
is impossible by the two local
directions at every point.
Thus the two boundaries coincide.
Injectivity of $b$ and the simplicial
triangulation force the original
triangles to coincide. This proves
the claimed disjointness.

Set $\zeta_j=e^{2\pi i j/3}$.
For each representative triangle fix
the conformal Jordan-disk map
\[
 f_\sigma:\overline{\mathbb D}
                 \longrightarrow\iota(\Delta_\sigma),
       \qquad f_\sigma(\zeta_j)=\iota(v_{\sigma,j}).
                                                                  \tag{C.89.5}
\]
For $b\in\mathscr B$, let
$f_{\sigma,b}$ be the conformal map
onto $\iota(D_b(\Delta_\sigma))$
with
$f_{\sigma,b}(\zeta_j)=\iota(b(v_{\sigma,j}))$.
The positive cyclic order supplies
the unique three-point normalization.
All these maps include their Jordan
boundary homeomorphisms.
Define the exact circle homeomorphism
\[
 \eta_{\sigma,b}
   =f_{\sigma,b}^{-1}\circ\iota\circ b
                  \circ\iota^{-1}\circ f_\sigma
                       \bigm|_{\mathbb S^1}.                   \tag{C.89.6}
\]
It fixes each $\zeta_j$.
Its radial extension is
$\widehat\eta_{\sigma,b}(r\zeta)
 =r\eta_{\sigma,b}(\zeta)$,
with value zero at zero. Its inverse
is the radial extension of
$\eta_{\sigma,b}^{-1}$.
The required map on the representative
triangle and its inverse are
\[
 \begin{aligned}
 F_{\sigma,b}
   &=\iota^{-1} f_{\sigma,b}
                 \widehat\eta_{\sigma,b} f_\sigma^{-1}\iota,\\
 F_{\sigma,b}^{-1}
   &=\iota^{-1} f_\sigma
           \widehat\eta_{\sigma,b}^{-1}f_{\sigma,b}^{-1}\iota .
 \end{aligned}                                                   \tag{C.89.7}
\]
It agrees with $b$ on the entire
boundary, including all labelled
vertices, not only at those vertices.

For the unique expression
$\Delta=q\Delta_\sigma$, define
\[
       \mathfrak s(b)|_{\Delta}
                  =q F_{\sigma,b}q^{-1},\qquad
       \mathfrak s(b)|_{\partial\overline D}=\id .
                                                                  \tag{C.89.8}
\]
Boundary compatibility follows from
the actual equation $bq=qb$ on
$\Lambda$. A plane homeomorphism
$q$ maps the compact Jordan disk
bounded by a curve to the compact
Jordan disk bounded by its image.
Consequently the image in (C.89.8)
is exactly $D_b(q\Delta_\sigma)$.
No conformality of $q$ is used.
An orientation-reversing $q$ occurs
on both sides of the conjugation,
so the extension is still orientation
preserving. The unique triangle
labels give, for every $q_0\in Q$,
$\mathfrak s(b)q_0=q_0\mathfrak s(b)$
on every triangle and on $\Lambda$.

\emph{Pasting on the original closed disk.}
The formula is continuous in the
interior by local finiteness.
For a sequence approaching the
original circle through infinitely
many triangles, the source triangle
diameters tend to zero.
Uniform continuity of $b$ on $X$
then makes their boundary-image
diameters tend to zero, with the
correct original boundary limit.
The convex-hull observation gives
the exact diameter equality
\[
 \operatorname{diam}_{d_E}D_b(\Delta)
             =\operatorname{diam}_{d_E}b(\partial\Delta).
                                                                  \tag{C.89.9}
\]
Thus the entire face images have
that same limit. This proves
continuity of (C.89.8) on the
original circle.

The disjoint-face result and
injectivity on $X$ make this a
continuous injection of
$\overline D$ into itself.
It is surjective. Its interior
image is open by invariance of
domain and closed relative to $O$:
the whole compact image is closed
and its boundary image is precisely
the original circle.
The interior image is nonempty
and $O$ is connected.
The compact-to-Hausdorff bijection
is a homeomorphism, with inverse
given by the inverses in (C.89.7),
their original $Q$-transports, and
the identity on the circle.
It therefore belongs to
$\mathscr C_Q$ and restricts to $b$.

Conversely, any $h\in\mathscr C_Q$
restricts to an embedding satisfying
every condition defining $\mathscr B$.
The map fixes the circle and hence
preserves the disk orientation.
It carries every original triangle
onto the Jordan disk bounded by
the image of its boundary, with
interior disjoint from $h(X)$.
This proves the asserted exact
image equality.
When $b=\id_X$, uniqueness of
the three-point normalization gives
$f_{\sigma,b}=f_\sigma$ and
$\eta_{\sigma,b}=\id$.
Every formula (C.89.8) is then
the identity, proving the required
normalization of the section.

\emph{The full parameter topology of the section.}
Suppose $b_n\to b$ uniformly on $X$.
For each representative triangle
the parametrized Jordan boundaries
converge uniformly, hence in
Fr\'echet distance.
Here is the exact use of Rad\'o's
verified parameter theorem
\cite[\S1, p.~182; proof, pp.~183--186]{C_Rado1923Variable}.
Choose one point inside the limiting
Jordan domain. It remains inside
the nearby domains by the winding
number of the uniformly convergent
curves about that point.
Normalize the Riemann maps there
with positive derivative.
Rad\'o gives uniform convergence
on the closed disk.
The preimages of each of the
three marked boundary vertices
also converge: any subsequential
limit maps to the specified limiting
vertex, whose preimage is unique.
The disk automorphism carrying
the three fixed $\zeta_j$ to these
three preimages varies continuously.
This follows, for example, from
the cross-ratio formulas for the
unique M\"obius map of distinct
cyclic triples.
Composing with this automorphism
gives exactly $f_{\sigma,b_n}$.
Therefore
$f_{\sigma,b_n}\to f_{\sigma,b}$
uniformly on the closed disk.
The normalization is unique, so
the temporary interior point has
not introduced any extra choice
into (C.89.7).

Equation (C.89.6) now gives uniform
convergence of the circle maps
$\eta_{\sigma,b_n}$.
One may verify the inverse
composition without presuming a
common target domain: for any
convergent circle input sequence,
compactness of the possible
preimages and uniform convergence
of the Riemann maps force every
preimage limit to be the unique
one in (C.89.6).
The same argument with varying
inputs proves uniform convergence.
The radial extensions converge
uniformly as well.
Formula (C.89.7) consequently
converges uniformly on each fixed
representative triangle.
For any other fixed triangle,
(C.89.8) gives the same conclusion:
the fixed map $q$ is uniformly
continuous on a compact interior
neighborhood of the convergent
image faces.

It remains to control all
triangles at once. The compact
family $\{b,b_1,b_2,\ldots\}$
is equicontinuous on $X$.
By C.82 nullity, for any
$\epsilon>0$ all but finitely
many source triangles have both
$b$-boundary and $b_n$-boundary
diameter less than $\epsilon$,
uniformly in $n$.
For a boundary point $a_\Delta$
of any such triangle and
$x\in\Delta$, (C.89.9) gives
\[
 \begin{split}
 d_E\bigl(\mathfrak s(b_n)x,\mathfrak s(b)x\bigr)
   &\leq
       \operatorname{diam}_{d_E}b_n(\partial\Delta)
       +d_E(b_n a_\Delta,b a_\Delta)\\
   &\hspace{1.2em}
       +\operatorname{diam}_{d_E}b(\partial\Delta)\\
   &<2\epsilon+\sup_{a\in X}d_E(b_n a,b a).
 \end{split}                                                   \tag{C.89.10}
\]
On the finitely many remaining
faces the proved individual
convergence applies. This proves
uniform convergence of the entire
original disk maps.
Their inverses converge uniformly
too: for homeomorphisms $u_n,u$
of a compact metric space,
\[
 \sup_y d_E(u_n^{-1}y,u^{-1}y)
       \leq
 \omega_{u^{-1}}\!
       \left(\sup_x d_E(u_nx,ux)\right).                        \tag{C.89.11}
\]
This follows by writing $y=u_nx$
and applying the modulus of
$u^{-1}$ to $u_nx,ux$.
Together with (C.89.4), these
estimates prove continuity into
the actual map-and-inverse group.
Both domain and target are
metrizable, so the sequential
argument proves continuity of
the global section.

\emph{The product and the actual deformation.}
Restriction is continuous for
the stated uniform topology
on $X$. The second coordinate
of $\Psi(h)$ fixes $X$
pointwise and hence lies in
$\mathscr K$.
Composition and inversion in
the actual closed-disk group
are continuous. Thus both
maps in (C.89.3) are continuous;
their displayed formulas
compose to the identity.

Let $A_t:\mathscr K\to\mathscr K$
be the single C.82 contraction,
with $A_0(k)=k$, $A_1(k)=1$
and $A_t(1)=1$.
The exact strong deformation
retraction is
\[
 \mathcal R_t(h)=
   \mathfrak s(h|_X)\,
   A_t\!\left(\mathfrak s(h|_X)^{-1}h\right).
                                                                  \tag{C.89.12}
\]
Its initial value is $h$,
its final value is
$\mathfrak s(h|_X)$, and it
fixes that entire section
image throughout.
Every factor retains the
original boundary values and
each individual $Q$ label.
This supplies the claimed
global family extension and
product comparison.
Constructing a controlled
deformation of arbitrary
based families in $\mathscr B$
to $\id_X$, or to one compact
quotient support for C.81,
remains the outstanding
general-to-skeleton step.
\end{proof}

\begin{claim}{C.90}{Proved}{Cocompact discrete operators in the original boundary topology}
Retain all the actual disk, operator, transitivity,
covering and topology hypotheses of C.81.
For a compact quotient $N=O/Q$, every based
sphere family in the original relative centralizer
has a based filling in that same group:
\[
          \pi_n(\mathscr C_Q,1)=0\qquad(n\geq1).
                                                               \tag{C.90.1}
\]
The filling fixes the original circle pointwise,
commutes with each individual original $q\in Q$,
and is continuous for both maps and inverses.
The quotient may be nonorientable; in particular
the torus and Klein-bottle cases are included.
No boundary extension of an operator or of an
interior conjugacy is required.

More precisely, let $p\in N$, let $x\in O$
be a specified lift of $p$, and write
\[
 \mathscr C_Q(x)=\{h\in\mathscr C_Q:h(x)=x\},\qquad
 \mathscr H_p=\operatorname{Homeo}(N\mathbin{\mathrm{rel}}\{p\})^0.
                                                               \tag{C.90.2}
\]
The superscript on $\mathscr H_p$ denotes its
identity component with its compact-open group
topology. Every based sphere family in
$\mathscr C_Q(x)$ descends into $\mathscr H_p$.
A single identity-fixed contraction of
$\mathscr H_p$, lifted through the original
covering from the identity, fills that sphere
with its original endpoint maps retained.
The inclusion $\mathscr C_Q(x)\hookrightarrow
\mathscr C_Q$ is the positive-homotopy
isomorphism of C.81, and transfers this filling
to (C.90.1).
\end{claim}
\passage{Section 23, folios 140-141, scans 279-281; Part C(b), zero-level kernel and original-boundary discrete-operator fillings for compact quotient}


\begin{proof}
\emph{The precise pointed quotient group.}
The original quotient is a connected surface
without boundary, since the free proper action
gives covering neighborhoods modeled on open
disks. It is compact by the present data.
The covering $\pi:O\to N$ has contractible
simply connected source $O$, so
$\pi_2(N)\cong\pi_2(O)=0$.
In particular $N$ is neither the sphere nor
the projective plane: their universal covers
are the sphere and their second homotopy
groups are $\mathbb Z$.

Use the quotient triangulation of C.82.
It is finite by compactness and local
finiteness. A subdivision in a coordinate
disk can make the specified point $p$
a vertex, so that $\{p\}$ is a compact
subpolyhedron of this PL surface.
Apply the verified compact-surface result
\cite[Lemma~2.1(i), p.~2; proof, p.~3]{A_Yagasaki2000}
to the exact pair $(N,\{p\})$.
The exceptions with boundary cannot occur.
Among its closed-surface exceptions, the
sphere and projective-plane pairs have just
been excluded, while the torus and
Klein-bottle exceptions have empty fixed
set. Our fixed set is the one specified
point. Thus the cited statement proves
contractibility of $\mathscr H_p$ for every
quotient in the present assertion, including
the nonorientable cases.

Choose a contraction
$C_t:\mathscr H_p\to\mathscr H_p$ with
$C_0(b)=b$ and $C_1(b)=c$ independent of $b$.
Its endpoint need not initially be the
identity. The normalized formula
\[
           J_r(b)=C_{1-r}(b)\,C_{1-r}(1)^{-1}
                        \qquad(0\leq r\leq1)
                                                               \tag{C.90.3}
\]
is continuous in this actual topological
group and has
$J_0(b)=1$, $J_1(b)=b$, and $J_r(1)=1$.
Every factor fixes $p$.

\emph{Descent of the actual pointed sphere.}
Let $a:(S^n,*)\to(\mathscr C_Q(x),1)$,
where $n\geq1$. Commutation with each
deck operator defines its exact descent
\[
 b_\xi(\pi z)=\pi(a_\xi z),\qquad
 b_\xi^{-1}(\pi z)=\pi(a_\xi^{-1}z).
                                                               \tag{C.90.4}
\]
These formulas are independent of the
chosen lift $z$. They are mutually
inverse, and $b_\xi(p)=p$.
Local sections of the original covering
prove joint continuity in $\xi$ and the
downstairs point for both formulas.
Since $N$ is compact, this is continuity
of $\xi\mapsto b_\xi$ in the stated
map-and-inverse topology. The path
connected sphere has $b_*=1$, so its
entire image lies in $\mathscr H_p$.
We make no assertion about descents of
other components of $\mathscr C_Q(x)$.

Put $B_{r,\xi}=J_r(b_\xi)$ and retain
the actual homotopy
\[
 (r,\xi,z)\longmapsto B_{r,\xi}(\pi z)
          \quad\text{from }[0,1]\times S^n\times O\text{ to }N.
                                                               \tag{C.90.5}
\]
The covering homotopy property lifts it
uniquely from $H_{0,\xi}=\id_O$.
This supplies a jointly continuous
family $H_{r,\xi}:O\to O$.
Apply the same covering property to
$B_{r,\xi}^{-1}\pi$, from the identity,
to obtain $K_{r,\xi}$.
The two composites $H_{r,\xi}K_{r,\xi}$
and $K_{r,\xi}H_{r,\xi}$ lift the
constant homotopy $\pi$ and start at
the identity. Uniqueness of the lifted
homotopy makes both composites the
identity. Thus $H_{r,\xi}$ is a
homeomorphism with jointly continuous
inverse $K_{r,\xi}$.

For every original $q\in Q$, the maps
$H_{r,\xi}q$ and $qH_{r,\xi}$ lift
the same homotopy $B_{r,\xi}\pi$
and start at the same map $q$.
They agree by uniqueness. This retains
each individual deck label.
The track $H_{r,\xi}(x)$ lies over
$B_{r,\xi}(p)=p$ and starts at $x$;
discreteness of that fiber gives
$H_{r,\xi}(x)=x$ throughout.
For the sphere basepoint, the downstairs
homotopy is constantly the identity,
so $H_{r,*}=\id_O$.
At $r=1$, both $H_{1,\xi}$ and
$a_\xi|_O$ lift $b_\xi\pi$ and
take $x$ to $x$. Uniqueness of a lift
with specified value on the connected
space $O$ proves exact endpoint agreement:
\[
 H_{0,\xi}=1,\qquad H_{1,\xi}=a_\xi|_O,
       \qquad H_{r,*}=1,\qquad H_{r,\xi}(x)=x.
                                                               \tag{C.90.6}
\]

\emph{Extension in the original group topology.}
Use C.88 with the compact metric parameter
space $P=[0,1]\times S^n$ and the actual
interior family $H_{r,\xi}$ just constructed.
Its required uniform quotient tail holds
with $K=N$: the set of $z$ with
$\pi z\notin K$ is empty.
Thus the explicit identity extensions
\[
 \overline H_{r,\xi}(z)=
 \begin{cases}
 H_{r,\xi}(z),&z\in O,\\
 z,&z\in\partial\overline D
 \end{cases}
                                                               \tag{C.90.7}
\]
form a continuous family in the original
$\mathscr C_Q$, including its inverse
topology. The inverse family is the
identity extension of $K_{r,\xi}$.
Equations (C.90.6) now hold on the
entire original closed disk, so this
is a based filling of $a$ in the
pointed original group itself.
No change of compactification has
occurred at any step.

\emph{Return to an arbitrary sphere.}
C.81 gives the actual locally trivial
evaluation bundle
\[
 \mathscr C_Q(x)\longrightarrow\mathscr C_Q
       \xrightarrow{\ h\mapsto h(x)\ }O.
                                                               \tag{C.90.8}
\]
Its base is contractible. Its exact
based homotopy sequence therefore
makes the displayed inclusion an
isomorphism on $\pi_n$ for every
$n\geq1$, precisely the finite-orbit
comparison already proved in C.81.
For an arbitrary based sphere $f$
in $\mathscr C_Q$, choose its
representative $a$ in that stabilizer
and a based homotopy $U_t$ from $f$
to $a$ supplied by this isomorphism.
Both families use the original group
topology. The complete contraction of
the given sphere is
\[
 F_{t,\xi}=
 \begin{cases}
 U_{2t,\xi},&0\leq t\leq\tfrac12,\\
 \overline H_{2-2t,\xi},&\tfrac12\leq t\leq1.
 \end{cases}
                                                               \tag{C.90.9}
\]
The two middle values are the same
original map $a_\xi$. The entire
homotopy is based, begins at $f$,
ends at the identity, and retains
the claimed boundary and operator
conditions for maps and inverses.
This proves (C.90.1).

Compactness is used here exactly in
the C.88 tail condition with $K=N$
and in the pointed compact-surface
contraction. For a noncompact
quotient, that choice of $K$ is
unavailable; the controlled quotient
tail for the lifted family remains
part of the unresolved general
discrete-operator deformation.
\end{proof}

\begin{claim}{C.91}{Proved}{A contraction of the whole cocompact original centralizer}
Retain the actual hypotheses of C.81 and the compact quotient
$N=O/Q$ of C.90.  Thus $O=\operatorname{int}\overline D$ is
the original open disk, $Q$ is closed, discrete, free and proper,
and an actual subgroup of the boundary-fixed centralizer
$\mathscr C_Q$ is transitive on $O$.
Fix $x\in O$, put $p=\pi x$, and write
$\mathscr F=\mathscr C_Q(x)$ and
$\mathscr H_p=\operatorname{Homeo}(N\mathbin{\mathrm{rel}}\{p\})^0$.
Give these groups their stated compact-open map-and-inverse
topologies.  There are exact topological group isomorphisms
\[
 \begin{aligned}
 D:\mathscr F&\longrightarrow\mathscr H_p,
       &D(h)(\pi z)&=\pi(hz),\\
 L:\mathscr H_p&\longrightarrow\mathscr F,
       &L(b)|_{\partial\overline D}&=\id.
 \end{aligned}                                                    \tag{C.91.1}
\]
The interior restriction of $L(b)$ is exactly the unique
lift of $b\pi$ that fixes $x$.
In particular this identifies the entire actual point stabilizer,
not only its identity component.  The proof establishes that every
one of its descents lies in the indicated pointed component.
The evaluation map admits a continuous section
$s:O\to\mathscr C_Q$ with $s(y)x=y$ and $s(x)=1$.
Consequently the mutually inverse homeomorphisms
\[
 \begin{aligned}
 \Phi:O\times\mathscr F&\longrightarrow\mathscr C_Q,
       &\Phi(y,k)&=s(y)k,\\
 \Psi:\mathscr C_Q&\longrightarrow O\times\mathscr F,
       &\Psi(h)&=(hx,s(hx)^{-1}h)
 \end{aligned}                                                    \tag{C.91.2}
\]
give an identity-fixed contraction of the whole original
$\mathscr C_Q$.  All maps in the contraction fix the original
circle and commute with every individual original $q\in Q$.
The assertion includes the torus, Klein bottle, and all the
other compact quotient surfaces permitted by these data.

\end{claim}
\passage{Section 23, folios 140-141, scans 279-281; Part C(b), original-boundary discrete-operator zero-level kernel and its full cocompact contraction}


\begin{proof}
\emph{A specified plane coordinate used only for homotopies and cells.}
Choose a disk coordinate
$\theta:\overline D\to\{u\in\mathbb C:|u|\leq1\}$ and put
$a=\theta(x)$.  Define
\[
 \begin{aligned}
 m_a(u)&=\frac{u-a}{1-\overline a u},
       &m_a^{-1}(v)&=\frac{v+a}{1+\overline a v},\\
 \psi(z)&=\frac{m_a(\theta(z))}{1-|m_a(\theta(z))|},
       &\psi^{-1}(w)&=
          \theta^{-1}m_a^{-1}\!\left(\frac{w}{1+|w|}\right).
 \end{aligned}                                                    \tag{C.91.3}
\]
These displayed inverse formulas give a homeomorphism
$\psi:O\to\mathbb R^2$ with $\psi(x)=0$.
They are a coordinate on the specified interior; all operator
equations below continue to use the original $Q$ and the
original quotient covering.  In particular the coordinate
does not replace the original boundary topology of the group.

\emph{Every actual pointed descent is pointed homotopic to the identity.}
For $h\in\mathscr F$, the formulas
\[
             b(\pi z)=\pi(hz),\qquad
             b^{-1}(\pi z)=\pi(h^{-1}z)                         \tag{C.91.4}
\]
are well-defined, mutually inverse continuous maps, because
$h$ and $h^{-1}$ commute with every $q$.  Both fix $p$.
For clarity, the induced fundamental-group comparison retains
the original labels.  A loop $\gamma$ at $p$ has a unique lift
$\widetilde\gamma$ starting at $x$, ending at some $qx$.
Its image $h\widetilde\gamma$ lifts $b\gamma$, starts at $x$,
and ends at $h(qx)=qh(x)=qx$.  Lifts with the same endpoints
are homotopic rel endpoints in the simply connected $O$.
Projecting such a homotopy shows
$[b\gamma]=[\gamma]$ in the actual $\pi_1(N,p)$.

Here is a direct based homotopy, so no classification of
pointed mapping classes is being presumed from that computation.
Take a finite triangulation of $N$ in which $p$ is a vertex,
as in C.90, and lift it through $\pi$.
Every lifted simplex maps homeomorphically onto its quotient
simplex.  A deck map preserving such a lift fixes a point
of it by injectivity of that restriction and hence is the
identity.  Choose one lift of every quotient simplex,
with the chosen lift of $p$ equal to $x$.
We construct a $Q$-equivariant homotopy
\[
 W:O\times[0,1]\longrightarrow O,
 \qquad W_0=h|_O,\quad W_1=\id_O,\quad W_t(qx)=qx.            \tag{C.91.5}
\]
On a chosen vertex $v$, use the path
$\psi^{-1}((1-t)\psi(hv)+t\psi(v))$.
For $v=x$ this is constantly $x$.
Define the paths on every translate by the literal formula
$W_t(qv)=qW_t(v)$.

For a chosen lifted edge $e$, the paths on its two endpoints,
together with $h|_e$ at time zero and $\id|_e$ at time one,
give a continuous map on the boundary of the square
$e\times[0,1]$.  This extends across the square.
An explicit extension for any boundary map $f:S^j\to O$
is obtained in cone coordinates by
\[
       (u,r)\longmapsto\psi^{-1}(r\psi(f(u))),
       \qquad 0\leq r\leq1,                                  \tag{C.91.6}
\]
with the cone vertex sent to $x$; a fixed ball-to-prism
homeomorphism supplies the square coordinates here.
Continuity at the cone vertex follows because $f(S^j)$
is compact and $\psi(f(S^j))$ is bounded.
Transport this extension to every translate of $e$ by $q$.
It agrees with the already defined vertex paths.

For a chosen lifted triangle, the two end faces and the
three edge homotopies define a continuous map on the
boundary of its triangular prism.  Apply (C.91.6) with
$j=2$ and a fixed ball-to-prism homeomorphism, and transport
the extension to all its $Q$-translates.
The boundary prescriptions agree on every common face.
The lifted triangulation is locally finite, so the closed-cell
pasting lemma on $O\times[0,1]$ proves continuity of $W$.
Each equivariance equation holds on every cell and hence
on all of $O$.  The vertex prescription proves the final
fixed-orbit equation in (C.91.5).

It follows that
\[
             B_t(\pi z)=\pi W_t(z)                            \tag{C.91.7}
\]
is a well-defined continuous homotopy from $b$ to $\id_N$
fixing $p$ at every time.  Continuity follows, for example,
by using local sections of the original covering.
This homotopy is proper: its domain $N\times[0,1]$ is
compact and the inverse image of a compact subset of
the Hausdorff surface $N$ is closed in that domain.
The surface $N$ has empty boundary, so the boundary-pair
condition in Epstein's theorem is satisfied as well.

Apply exactly \cite[Theorem~6.3, p.~99; proof, pp.~100--101]
{C_Epstein1966}.  This is the theorem for a proper homotopy
of the boundary pair; its last sentence retains a basepoint
fixed by the homotopy.  Its proof explicitly treats a
non-PL homotopy on p.~101 by a basepoint-preserving isotopy
and PL approximation, using the Appendix; the stated
topological conclusion therefore applies to our actual $b$.
We obtain an isotopy from $b$ to $\id_N$ fixing $p$.
On the compact $N$, the isotopy is continuous in the
compact-open topology of maps, and its inverse family is
continuous as well: the level map on $N\times[0,1]$ is
a continuous bijection of a compact space onto a Hausdorff
space.  Thus $b\in\mathscr H_p$ in the stated group topology.
This proves the asserted component membership for every
$h\in\mathscr F$, without first restricting $h$ to a component.

\emph{The exact pointed compact-surface contraction and all exceptions.}
As in C.90, $N$ is a closed connected PL surface with
universal cover $O$.  The covering isomorphism on second
homotopy groups gives $\pi_2(N)=0$, whereas $S^2$ and
$\mathbb{RP}^2$ have second homotopy group $\mathbb Z$.
These two closed-surface cases are therefore excluded
by the actual covering.
For the pair $(N,\{p\})$, Yagasaki's
\cite[Lemma~2.1(i), p.~2; proof, p.~3]{A_Yagasaki2000}
proves contractibility of $\mathscr H_p$.
The exceptional disk, annulus, and M\"obius-band pairs
have boundary and cannot be our $N$.
The exceptional torus and Klein-bottle pairs have empty
fixed set, whereas our fixed set is $\{p\}$.
The remaining listed closed exceptions belong to the
already excluded sphere and projective plane.
No orientability assumption occurs in this application.
Likewise Epstein's Theorem~6.3, used above with the
proper based homotopy (C.91.7), includes the closed
nonorientable surfaces.  We have not used the distinct
nonproper-homotopy statement of his Theorem~6.4(b)
or its open or closed M\"obius-band qualification.

Write $C_t$ for a contraction of $\mathscr H_p$,
with $C_0(b)=b$ and constant endpoint $C_1(b)=c$.
The exact normalization already used in C.90 is
\[
 J_r(b)=C_{1-r}(b)C_{1-r}(1)^{-1},
 \qquad J_0(b)=1,\quad J_1(b)=b,\quad J_r(1)=1.                \tag{C.91.8}
\]
Every value and every factor fixes $p$.

\emph{The inverse lift and the two actual topologies.}
For $b\in\mathscr H_p$, lift the homotopy
$J_r(b)\pi:O\to N$ from the identity lift at $r=0$.
Let $H_r(b)$ be the resulting interior family.
Lift $J_r(b)^{-1}\pi$ from the identity to obtain $K_r(b)$.
The two compositions $H_r(b)K_r(b)$ and $K_r(b)H_r(b)$
lift the constant homotopy $\pi$ and start at the identity;
uniqueness makes both compositions the identity.
For every literal $q\in Q$, the two maps
$H_r(b)q$ and $qH_r(b)$ lift the same homotopy and
start at the same map $q$, so they agree.
The track of $x$ lies in the discrete fiber over $p$
and starts at $x$, hence remains $x$.

C.88 applies to this compact parameter interval with
$K=N$: the required tail is empty.
It supplies the continuous original-circle identity
extensions of $H_r(b)$ and $K_r(b)$, which remain inverse.
Define $L(b)$ as the extension of $H_1(b)$.
It is therefore an actual member of $\mathscr F$.
It descends to $b$.  Any two lifts of $b\pi$ that take
$x$ to $x$ agree on the connected $O$; their boundary
extensions agree too.  This proves
\[
              D L(b)=b,\qquad L D(h)=h.                       \tag{C.91.9}
\]
The same uniqueness gives $L(bc)=L(b)L(c)$, so both maps
in (C.91.1) are group maps with the displayed inverse.

The descent topology has a direct metric estimate.
Use the actual metrics $d,d_Q,\rho$ of C.81 and set
\[
 \begin{aligned}
 d_{\mathscr C}(h,k)
   &=\sup_{z\in\overline D}d(hz,kz)
      +\sup_{z\in\overline D}d(h^{-1}z,k^{-1}z),\\
 d_{\mathscr H}(b,c)
   &=\sup_{u\in N}\rho(bu,cu)
      +\sup_{u\in N}\rho(b^{-1}u,c^{-1}u).
 \end{aligned}
\]
They induce the asserted map-and-inverse topologies
because their underlying spaces are compact.
The quotient inequality and the exact C.81 isometry give
\[
                 d_{\mathscr H}(Dh,Dk)
                       \leq d_{\mathscr C}(h,k).              \tag{C.91.10}
\]
Thus $D$ is continuous in the original group topology.

For the inverse, take any sequence $b_n\to b$ in
$\mathscr H_p$ and use the compact metric parameter space
$P=\{0\}\cup\{1/n:n\geq1\}$ with $b_0=b$.
The family $J_r(b_u)$ and its inverse are continuous
on $[0,1]\times P\times N$.
Lift them jointly from the identity through the original
covering, exactly as above.  The resulting interior
families have every original commutation equation and
fix $x$.  C.88, now with compact parameter
$[0,1]\times P$ and again $K=N$, supplies joint
continuity of their identity extensions and inverses
on the original closed disk.  At $r=1$ this gives
$L(b_n)\to L(b)$ in $d_{\mathscr C}$.
Both groups in (C.91.1) are metrizable subspaces of
the indicated compact-domain homeomorphism groups,
so this sequential argument proves continuity of $L$.
It proves the entire asserted topological group
isomorphism, rather than only a set-theoretic lifting rule.

In particular the single formula
\[
 A_t(k)=L\bigl(J_{1-t}(Dk)\bigr),
 \quad A_0(k)=k,\quad A_1(k)=1,\quad A_t(1)=1                \tag{C.91.11}
\]
is a continuous identity-fixed contraction of the whole
point stabilizer $\mathscr F$ in its original topology.

\emph{A global section with explicit local formulas.}
C.81 supplies actual continuous local sections of
$h\mapsto hx$ over an open cover $\{U_i\}$ of $O$.
Write a chosen one on $U_i$ as
$\sigma_i:U_i\to\mathscr C_Q$, so that
$\sigma_i(y)x=y$; these maps already have the original
map-and-inverse topology by C.81.11.
Start with the standard locally finite unit-grid
triangulation of $\mathbb R^2$, with $0$ a vertex.
Apply exactly the $k=0$ subdivision lemma used in C.82,
\cite[Lemma~1, p.~458]{C_Haver1975}, to this complex,
the actual map $\psi^{-1}$ into $O$, and the cover
$\{U_i\}$.  It gives a locally finite subdivision
whose closed vertex stars map into members of that
cover.  Each closed simplex is in the closed star
of any one of its vertices, so its image lies in
some $U_i$.  Pull the subdivided triangles back to
$O$ by $\psi^{-1}$.  The original vertex $0$ remains
a vertex under subdivision, and its image is $x$.

At every vertex $v$, choose once an actual
$s(v)\in\mathscr C_Q$ taking $x$ to $v$,
using the surjectivity in C.81.  Choose $s(x)=1$.
We extend this section over the edges and then over
the two-simplices.  Suppose its values on the boundary
of a closed simplex $\tau$ have been defined.
Choose $i(\tau)$ with $\tau\subset U_{i(\tau)}$ and put
\[
 a_\tau(u)=\sigma_{i(\tau)}(u)^{-1}s(u)\in\mathscr F
                         \qquad(u\in\partial\tau).            \tag{C.91.12}
\]
In the affine coordinates of the plane simplex, let
$c_\tau$ be its barycenter.  Every noncentral point
has a unique expression $y=(1-r)c_\tau+ru$ with
$u\in\partial\tau$ and $0<r\leq1$; this notation
is transported to $O$ by $\psi^{-1}$.
Define the extension by the exact formulas
\[
 \begin{aligned}
 s(y)&=\sigma_{i(\tau)}(y)
                   A_{1-r}\bigl(a_\tau(u)\bigr)
           &&(0<r\leq1),\\
 s(c_\tau)&=\sigma_{i(\tau)}(c_\tau).
 \end{aligned}                                                    \tag{C.91.13}
\]
The second factor always fixes $x$, so $s(y)x=y$.
At $r=1$ its value is exactly the already prescribed
boundary value.  At the center, continuity follows
from $A_1=1$ uniformly on the compact set
$a_\tau(\partial\tau)$: the restriction of the
continuous map $A$ to
$[0,1]\times a_\tau(\partial\tau)$ is uniformly
continuous in any compatible metric on its image.
This checks the possible polar-coordinate singularity
without imposing any uniform condition at infinity.
The edge extensions fit at their endpoints, and
the face extensions retain every edge value.
Local finiteness and the closed-cell pasting lemma
therefore give one global continuous
$s:O\to\mathscr C_Q$ in the actual group topology.
The original value $s(x)=1$ is retained.

Now the maps in (C.91.2) are defined on the displayed
spaces: $s(hx)^{-1}h$ fixes $x$.
Their two composites are identities by multiplication
and $s(y)x=y$.  Evaluation, multiplication, inversion,
and $s$ are continuous in the original topology,
so these are mutually inverse homeomorphisms.
This also proves that the constructed section is a
section of the actual C.81 bundle, with its original
fiber and topology.

\emph{The contraction of the whole original group.}
In the specified coordinate (C.91.3), put
\[
             R_t(y)=\psi^{-1}((1-t)\psi(y)).                   \tag{C.91.14}
\]
This is a continuous contraction of $O$ to $x$,
fixing $x$ throughout.  For $h\in\mathscr C_Q$ let
$y=hx$ and $k=s(y)^{-1}h$.  The required global
contraction and its pointwise inverse map are
\[
 \begin{aligned}
 \mathcal C_t(h)&=s(R_t(y))A_t(k),\\
 \mathcal C_t(h)^{-1}&=A_t(k)^{-1}s(R_t(y))^{-1}.
 \end{aligned}                                                    \tag{C.91.15}
\]
All factors belong to the original $\mathscr C_Q$,
so every value fixes its original circle pointwise
and commutes with each original $q$.
The formulas and (C.91.2), (C.91.11), (C.91.14) prove
joint continuity in the map-and-inverse topology.
At time zero the first formula is $s(y)k=h$;
at time one it is $s(x)\cdot1=1$.
For $h=1$, we have $y=x$, $k=1$, so
$\mathcal C_t(1)=1$ for every $t$.
This proves the claimed contraction of the entire group.
\end{proof}

\begin{claim}{C.92}{Proved}{Arbitrary pointed descents and the exact original lift topology}
Retain all the actual data of C.81: the original closed disk
$\overline D$, its interior $O$, the compatible closed-disk metric $d$,
the closed discrete free proper group $Q$ of operators on $O$,
the boundary-fixed centralizer $\mathscr C_Q$ with a transitive actual
subgroup, and the specified quotient covering $\pi:O\to N=O/Q$.
The quotient $N$ is allowed to be noncompact throughout this note.
Fix $x\in O$, put $p=\pi x$, and set $\mathscr F=\mathscr C_Q(x)$.
Write $\mathscr I_p$ for the identity path component of
$\operatorname{Homeo}(N\mathbin{\mathrm{rel}}\{p\})$ in the compact-open
topology of maps and inverses. Thus a path in this definition has
jointly continuous maps and inverses on $[0,1]\times N$.

Every member of the entire original $\mathscr F$ descends into
$\mathscr I_p$. Below we prove this assertion without assuming
component membership upstairs. We also identify its exact image
inside $\mathscr I_p$, with inverse maps and both relevant topologies.
The original compact-quotient contraction C.91 remains unchanged.


More precisely, let $\lambda(b)$ be the unique lift of $b\pi$
fixing $x$ for $b\in\mathscr I_p$. The image of $D$ is exactly
the subgroup $\mathscr B_p$ defined by
$\delta_b(\pi z)=d_Q(\lambda(b)z,z)\in C_0(N)$.
Descent and extension of this lift are inverse group maps and
are isometric for the original sum metric and the lifted sum
metric in (C.92.15). On a compact parameter space their exact
continuity criterion is the uniform tail (C.92.17).
The explicit noncompact example (C.92.18)--(C.92.21) gives a
sequence in this exact image converging in the quotient
compact-open topology whose pointed lifts stay more than
$2/5$ from the identity in the original metric.

\end{claim}
\passage{Section 23, folios 140-141, scans 279-281; Part C(b), exact original-boundary pointed descent, lift topology and the remaining noncompact level-zero kernel}


\begin{proof}
\emph{Literal descent and the original orbit labels.}
For $h\in\mathscr F$, define
\[
 D(h)(\pi z)=\pi(hz),\qquad
 D(h)^{-1}(\pi z)=\pi(h^{-1}z).
                                                               \tag{C.92.1}
\]
Commutation with every actual $q$ makes these formulas well-defined.
The quotient property makes them continuous, and substitution proves
they are inverse maps. They fix $p$. Composition gives
$D(hk)=D(h)D(k)$. Two such maps with equal descents are lifts of the
same map on the connected $O$, agree at $x$, and hence agree everywhere;
their boundary values are both the identity. Thus $D$ is injective.

For a loop $\gamma$ at $p$, let its lift starting at $x$ end at $qx$.
The image under $h$ of that lift starts at $x$ and ends at
$h(qx)=qh(x)=qx$. It lifts $D(h)\gamma$.
Since $O$ is simply connected, the two lifted paths with these
endpoints are homotopic rel endpoints. Their projections show
$D(h)_*[\gamma]=[\gamma]$ in the specified $\pi_1(N,p)$.
In particular every original orbit label is retained.

\emph{A based homotopy on the locally finite lifted triangulation.}
Here is a direct homotopy, including its maps, rather than an
inference of pointed isotopy from the preceding loop calculation.
Choose a disk coordinate $\theta:\overline D\to\{|u|\leq1\}\subset
\mathbb C$, put $a=\theta(x)$, and use the mutually inverse maps
\[
 \begin{aligned}
 m_a(u)&=(u-a)/(1-\overline a u),
 &m_a^{-1}(v)&=(v+a)/(1+\overline a v),\\
 \psi(z)&=\frac{m_a(\theta z)}{1-|m_a(\theta z)|},
 &\psi^{-1}(w)&=\theta^{-1}m_a^{-1}\left(\frac{w}{1+|w|}\right).
 \end{aligned}                                                  \tag{C.92.2}
\]
They give $\psi:O\to\mathbb R^2$ with $\psi(x)=0$.
The original covering and every operator $q$ remain the ones in
(C.92.1).

Use C.82's locally finite triangulation of the quotient surface,
refined locally so that $p$ is a vertex, and lift it by $\pi$.
Every lifted closed simplex maps homeomorphically to its quotient
simplex. The lifted family is locally finite, by the local covering
charts. Any deck operator preserving a lifted simplex fixes it
pointwise, by injectivity of the projection on that simplex, and is
therefore the identity by freeness. The lifts of a given quotient
simplex form one $Q$-orbit: match one point in two lifts by the unique
operator in its fiber and then use uniqueness of the lift on the
connected simplex. Choose one lift of each simplex and choose $x$
as the lift of the vertex $p$.

For a chosen vertex $v$ define
$W_t(v)=\psi^{-1}((1-t)\psi(hv)+t\psi(v))$, and define
$W_t(qv)=qW_t(v)$ on all translates. The path at $x$ is constant.
For a chosen edge, its endpoint paths and its prescribed maps
$h$ and the identity at times zero and one give a map $f$ on the
boundary of its square prism. For a chosen triangle the already
defined edge homotopies and the two end faces likewise give a map
on the boundary of its triangular prism.

For completeness, use the following explicit extension on either
convex coordinate prism $P$. Choose an interior point $c$. Every
point other than $c$ is uniquely $c+r(u-c)$ with $u\in\partial P$
and $0<r\leq1$. Define
\[
 \operatorname{Cone}_f(c+r(u-c))
       =\psi^{-1}\bigl(r\psi(f(u))\bigr),
 \qquad \operatorname{Cone}_f(c)=x.                            \tag{C.92.3}
\]
These radial coordinates are continuous away from $c$: on each
ray the boundary parameter is the first intersection with one
of the finitely many facet hyperplanes of the convex prism.
Equivalently, for a unit vector $v$ and inequalities
$a_i\cdot y\leq b_i$ defining $P$, that parameter is
$\min_{a_i\cdot v>0}(b_i-a_i\cdot c)/(a_i\cdot v)$.
This is positive and continuous on the unit sphere.
Continuity of (C.92.3) at $c$ follows from boundedness of
$\psi(f(\partial P))$. The value on $\partial P$ is exactly $f$.

Apply (C.92.3) first to each chosen edge prism and then to each
chosen triangle prism, using the coordinates of the triangulation,
and transport each extension by the literal operator $q$.
The boundary prescriptions agree on common faces. Local finiteness
gives continuity by closed-cell pasting on $O\times[0,1]$.
Thus we have actual maps
\[
 \begin{gathered}
 W:O\times[0,1]\to O,\qquad W_0=h|_O,\quad W_1=\id_O,\\
 W_t(qz)=qW_t(z),\qquad W_t(qx)=qx,\\
 B_t(\pi z)=\pi W_t(z),\qquad
 B_0=D(h),\quad B_1=\id_N,\quad B_t(p)=p .
 \end{gathered}                                                \tag{C.92.4}
\]
The last family is continuous by the local sections of $\pi$.
This is a based homotopy for every original $h\in\mathscr F$.
The argument uses local finiteness and does not require the
triangulation to be finite.

\emph{The orientation maps, including a nonorientable quotient.}
We verify every orientation qualification in the source theorem.
The connecting map
\[
 \partial:H_2(\overline D,\partial\overline D;\mathbb Z)
             \longrightarrow H_1(\partial\overline D;\mathbb Z)
                                                               \tag{C.92.5}
\]
is an isomorphism by the exact sequence of the contractible disk.
Naturality and $h|_{\partial\overline D}=\id$ give
$\partial h_*=\partial$, so $h_*$ is the identity on the left group.

Let $j_x$ be the map from
$H_2(\overline D,\partial\overline D)$ to
$H_2(\overline D,\overline D\setminus\{x\})$ induced by inclusion
of pairs. It is an isomorphism. Indeed, in the coordinate
$\kappa=m_a\theta$, the formula
\[
 R_s(z)=\kappa^{-1}
       \left(\left(1-s+\frac{s}{|\kappa(z)|}\right)\kappa(z)\right),
 \quad z\ne x,\quad 0\leq s\leq1,                             \tag{C.92.6}
\]
deformation retracts $\overline D\setminus\{x\}$ onto the boundary
while fixing the boundary. The natural diagram of the two exact
sequences then proves the assertion about $j_x$.
Let
$e_x:H_2(O,O\setminus\{x\})\to
H_2(\overline D,\overline D\setminus\{x\})$
be the excision isomorphism. Naturality for $h$, which fixes $x$,
gives, with $\eta_x=e_x^{-1}j_x$,
\[
 h_{x*}=\eta_x\,h_*\,\eta_x^{-1}
       =\id\quad\hbox{on }H_2(O,O\setminus\{x\};\mathbb Z).
                                                               \tag{C.92.7}
\]

Choose one specified evenly covered neighborhood $V_0$ of $p$
and its sheet $U_0$ containing $x$. Write $i_U,i_V$ for the
excision maps from the local homology groups of $U_0,V_0$
to those of $O,N$, respectively. The precise orientation
comparison at $x$ and $p$ is the isomorphism
\[
 \omega_x=i_V\circ(\pi|_{U_0})_*\circ i_U^{-1}:
 H_2(O,O\setminus\{x\})
        \longrightarrow H_2(N,N\setminus\{p\}).
                                                               \tag{C.92.8}
\]
Choose a smaller open neighborhood
$U\subset U_0\cap h^{-1}(U_0)$ of $x$.
On $\pi U$ the literal local formula is
\[
 D(h)|_{\pi U}
     =(\pi|_{U_0})\circ h|_U\circ(\pi|_U)^{-1}.
                                                               \tag{C.92.9}
\]
Excision and naturality applied to this formula prove
$D(h)_{p*}=\omega_xh_{x*}\omega_x^{-1}=\id$.
Thus the local orientation at the fixed point is preserved.
The use of a sheet in (C.92.8) is essential to the definition:
the global map $\pi$ sends every other point $qx$ in this fiber
to $p$, so the required local comparison is the displayed
excision composite.

On a connected orientable $N$, the sign of a homeomorphism on
the continuously chosen local fundamental classes is locally
constant and hence constant. Its sign at $p$ is $+1$ by
(C.92.7)--(C.92.9), so $D(h)$ preserves orientation globally.
This verifies the plane and open-cylinder cases.
On an open M\"obius band the same local-homology identity
verifies preservation of orientation near the basepoint.
Each point of $N$ has an open-disk covering chart, so
$\partial N=\varnothing$; the closed-cylinder and closed
M\"obius-band cases in the source statement therefore do
not occur as this $N$. No orientation of a nonorientable
$N$, and no orientation assumption on the original $q$,
has entered the local calculation.

\emph{The exact Epstein result and the topology of the resulting path.}
Use \cite[Theorem~6.4(b), p.~101; proof, pp.~101--103]{C_Epstein1966}.
Its based statement is the following: a self-homeomorphism of a
surface whose boundary components are compact, based homotopic
to the identity, is isotopic to the identity fixing that basepoint,
with orientation preservation required for the plane and the open
or closed cylinder, and local orientation preservation at the
basepoint required for the open or closed M\"obius band.
Equations (C.92.4) and (C.92.7)--(C.92.9) verify these hypotheses;
the boundary-component requirement holds because $\partial N$
is empty.

The exceptional M\"obius cases are proved in Lemma~6.5 on p.~102.
For the general case the proof on p.~103 concatenates isotopies
fixing successively larger compact submanifolds and explicitly
constructs a level-preserving homeomorphism of $N\times[0,1]$.
The topological version is included in the statement. Appendix
Theorem~A4 on p.~107 supplies its basepoint-fixed passage to a
PL homeomorphism; the corresponding topological proof paragraph
on p.~101 records the relative approximation step.
The hypotheses of Theorem~6.4(b) are met by the displayed
based homotopy itself; a properness assertion for (C.92.4)
on a noncompact $N$ is not used.

There is a small topology point in using this source:
Definition~1.4 on p.~84 defines an isotopy of embeddings.
The proof used here supplies homeomorphisms throughout.
For the general case this is explicit on p.~103.
The plane and open-cylinder steps in Theorems~5.4 and 5.7,
pp.~96--97, use Theorem~5.1, p.~95: its level maps for
$t<1$ are conjugates of the displayed homeomorphism, and
at $t=1$ are the identity. The compact disk step is the
explicit level-preserving homeomorphism of Theorem~5.2,
pp.~95--96. The open-M\"obius proof in Lemma~6.5 uses a
compact M\"obius core and the same half-cylinder construction
on its complement. In the cylinder step our based homotopy
also retains the generator of its fundamental group, as
proved after (C.92.1). After the based core-circle adjustment
in the source proof, the map fixes the core pointwise;
local orientation preservation makes it preserve its two
sides, so the two half-cylinder constructions apply to
the actual complementary components.
The topological approximation from A4 is an ambient
isotopy and likewise consists of homeomorphisms.

Write the resulting level-preserving homeomorphism as
$(y,t)\mapsto(b_t(y),t)$, where $b_0=D(h)$, $b_1=\id_N$,
and $b_t(p)=p$. Its inverse is exactly
$(y,t)\mapsto(b_t^{-1}(y),t)$ and is continuous.
This inverse continuity can also be checked without a
compactness assumption on $N$: extend the homeomorphism
family constantly for $t<0$ and $t>1$. Its level map is
a continuous bijection of the boundaryless $3$-manifold
$N\times\mathbb R$. Invariance of domain makes it open,
so its inverse is continuous; restrict back to $[0,1]$.
The compact-open exponential law on the locally compact
surface $N$ therefore gives a continuous path of maps and
inverses. Consequently $D(h)\in\mathscr I_p$ for every
$h\in\mathscr F$. This also places each such descent in the
pointed identity component; no equality of all components
or all possible descents has been presumed.

\emph{The exact pointed lift and its compact-open continuity.}
For a pointed homeomorphism $b$ of $N$, define the actual path
and its specified lift by
\[
 \beta_{b,z}(t)=b\!\left(\pi\psi^{-1}(t\psi(z))\right),
 \qquad
 \lambda(b)(z)=\widetilde\beta_{b,z}(1),\quad
 \widetilde\beta_{b,z}(0)=x.
                                                               \tag{C.92.10}
\]
The covering homotopy property applied with parameter $(b,z)$
gives joint continuity of the lifted paths. Indeed evaluation
of $b$ is continuous for the compact-open topology because
$N$ is locally compact, and the displayed formula is jointly
continuous before lifting. It gives
$\pi\lambda(b)=b\pi$ and $\lambda(b)(x)=x$.
Uniqueness of lifts on the connected $O$ makes it the unique
lift with that value. The same uniqueness proves the exact
identities
\[
 \lambda(bc)=\lambda(b)\lambda(c),\qquad
 \lambda(b^{-1})=\lambda(b)^{-1},\qquad
 \lambda(\id_N)=\id_O.
                                                               \tag{C.92.11}
\]
For the inverse identity, the two composites lift $\pi$
and fix $x$, and so are the identity. Thus these are actual
homeomorphisms. Applying the joint argument in (C.92.10) to
$b^{-1}$ also proves joint inverse continuity.
The exponential law on $O$ proves that $\lambda$ is continuous
from the pointed compact-open map-and-inverse topology to
the compact-open map-and-inverse topology on $O$.

For $b\in\mathscr I_p$, take a pointed path $b_t$ from the
identity to $b$. The continuous path $\lambda(b_t)(qx)$
lies in the discrete fiber $Qx$ and starts at $qx$.
It is therefore constantly $qx$. Now
$\lambda(b)q$ and $q\lambda(b)$ lift the same map $b\pi$
and agree at $x$, so uniqueness gives
\[
                  \lambda(b)q=q\lambda(b)
                         \quad\text{for every original }q.
                                                               \tag{C.92.12}
\]
This proves the individual operator equations without choosing
a deck representative modulo any subgroup.

Descent is continuous in the other direction as well.
The map $(h,z)\mapsto\pi(hz)$ on $\mathscr F\times O$
is continuous and is constant on the fibers of
$\id_{\mathscr F}\times\pi$. This product map is an open
quotient map, since $\pi$ is open. It therefore induces
the continuous joint map $(h,y)\mapsto D(h)y$.
Repeat the argument with $h^{-1}$ and use the exponential
law on $N$. This proves continuity of
$D:\mathscr F\to\mathscr I_p$ in the quotient compact-open
map-and-inverse topology.

\emph{The exact original image and its isometric topology.}
Retain $d_Q(z,w)=\sup_{q\in Q}d(qz,qw)$ from C.81 and put
\[
 \begin{split}
 \delta_b(\pi z)&=d_Q(\lambda(b)z,z),\\
 \mathscr B_p&=\{b\in\mathscr I_p:\delta_b\in C_0(N)\},\\
 L(b)|_O&=\lambda(b),\qquad
 L(b)|_{\partial\overline D}=\id\quad(b\in\mathscr B_p).
 \end{split}                                                   \tag{C.92.13}
\]
The displacement is well-defined and continuous by (C.92.12)
and $Q$-invariance of $d_Q$. C.88 gives its original boundary
extension exactly when this displacement belongs to $C_0(N)$.
For $h\in\mathscr F$, uniqueness gives
$\lambda(D(h))=h|_O$, so C.88 puts $D(h)$ in $\mathscr B_p$.
Conversely C.88 extends every $\lambda(b)$ in (C.92.13) to
an element of $\mathscr F$, and its descent is $b$. Thus
\[
 D:\mathscr F\longrightarrow\mathscr B_p,\qquad
 L:\mathscr B_p\longrightarrow\mathscr F,\qquad
 LD=\id_{\mathscr F},\quad DL=\id_{\mathscr B_p}
                                                               \tag{C.92.14}
\]
are the exact inverse group maps.
For example the subgroup assertion also follows directly from
$\delta_{bc}\leq\delta_b\circ c+\delta_c$ and
$\delta_{b^{-1}}=\delta_b\circ b^{-1}$: a homeomorphism and its
inverse carry compact sets to compact sets, so these pullbacks
preserve $C_0(N)$.

Define a metric on this exact image by the actual lifts:
\[
 d^\uparrow(b,c)=
 \sup_{z\in O}d_Q(\lambda(b)z,\lambda(c)z)
 +\sup_{z\in O}d_Q(\lambda(b)^{-1}z,\lambda(c)^{-1}z).
                                                               \tag{C.92.15}
\]
It is finite since $d$ is bounded on the compact original disk.
It is a metric by injectivity of the lift. Equivariance,
reindexing, and density of the interior give
\[
 \begin{split}
 \sup_{z\in O}d_Q(\lambda(b)z,\lambda(c)z)
 &=\sup_{z\in O,\ q\in Q}d(\lambda(b)qz,\lambda(c)qz)\\
 &=\sup_{w\in\overline D}d(L(b)w,L(c)w).
 \end{split}                                                   \tag{C.92.16}
\]
The same identity for the inverse maps proves that (C.92.14)
are isometric topological group isomorphisms for the original
closed-disk sum metric and $d^\uparrow$.
In particular the identity from this lifted metric on
$\mathscr B_p$ to its quotient compact-open topology is
continuous, by the descent argument above.

For a compact metric parameter space $P$ and a continuous
map $u\mapsto b_u$ into $\mathscr B_p$ with its quotient
compact-open topology, (C.92.10) gives joint interior continuity
of the actual pointed lifts and their inverses. C.88 now gives
the exact criterion for $u\mapsto L(b_u)$ to be continuous
in the original topology:
\[
 \forall\epsilon>0\ \exists K\subset N\text{ compact}\quad
 \sup_{\substack{u\in P\\y\notin K}}\delta_{b_u}(y)<\epsilon .
                                                               \tag{C.92.17}
\]
Its inverse family then satisfies the same uniform tail
condition. Thus the compact-family requirement is retained
on the unique lifts in (C.92.10), with every deck label fixed.

When $N$ is compact, taking $K=N$ makes (C.92.17) automatic,
and $C_0(N)=C(N)$ in the definition used here.
Then $\mathscr B_p=\mathscr I_p$.
For a convergent sequence $b_n\to b$ in the quotient topology,
apply (C.92.17) to $\{0\}\cup\{1/n:n\geq1\}$ and the lifts from
(C.92.10). This gives $L(b_n)\to L(b)$ in the original metric.
Both topologies are metrizable, so this proves continuity
of the inverse comparison in the compact case.
C.91's pointed-component contraction makes that entire
component path connected, giving
$\mathscr I_p=\operatorname{Homeo}(N\mathbin{\mathrm{rel}}\{p\})^0$
there. These are precisely C.91's original maps and compact
hypothesis.

\emph{An exact noncompact sequence, including its deck labels.}
The following actual C.81 example records why the last compact
argument cannot be replaced just by quotient compact-open
convergence. In this example $\overline D$ is the closed unit disk with its
Euclidean metric $d$. Use the original radial disk coordinate
\[
 j:\mathbb R^2\to\{|w|<1\},\quad j(z)=\frac{z}{1+|z|},
 \qquad j^{-1}(w)=\frac{w}{1-|w|},
                                                               \tag{C.92.18}
\]
and take $Q=\langle jTj^{-1}\rangle$ with
$T(s,t)=(s+1,t)$. The quotient coordinate is
$\pi(j(s,t))=([s],t)$ on $N=(\mathbb R/\mathbb Z)\times\mathbb R$.
The actual extended translations $j(z\mapsto z+a)j^{-1}$
are transitive, fix the original circle, and commute with
every displayed operator. The action is closed, discrete,
free and proper by the explicit plane coordinate, so this
is an instance of the retained hypotheses.

For an integer $R\geq2$ put
$\beta(v)=\max(0,1-4|v|)$ and
\[
 F_{R,t}(s,y)=(s+tR\beta(y-R),y),\qquad
 F_{R,t}^{-1}(s,y)=(s-tR\beta(y-R),y),
 \quad 0\leq t\leq1 .
                                                               \tag{C.92.19}
\]
They commute with every $T^n$, fix $(0,0)$, and have quotient
support in
$(\mathbb R/\mathbb Z)\times[R-\tfrac14,R+\tfrac14]$.
Both displacements in the plane are bounded by $R$.
For $|w|\leq R$ and $|z|>R$, direct subtraction gives
\[
       |j(z+w)-j(z)|\leq\frac{2R}{1+|z|-R}.
                                                               \tag{C.92.20}
\]
Thus $jF_{R,t}j^{-1}$ and its inverse extend jointly over
$[0,1]$ by the identity on the original circle.
Let $b_R$ be the descent of this endpoint.
Then $b_R\in\mathscr B_p$, with $x=j(0,0)$.
Its unique pointed lift is exactly $jF_{R,1}j^{-1}$.
Both $b_R$ and $b_R^{-1}$ are eventually the identity
on any specified compact subset of $N$, so $b_R\to1$
in the quotient compact-open map-and-inverse topology.

At $y_R=([0],R)$ the quotient value is $b_R(y_R)=y_R$,
because the horizontal translation there is the integer $R$.
The actual pointed lift, however, sends
$j(0,R)$ to $j(R,R)$, and therefore
\[
 \begin{gathered}
 \rho(b_R(y_R),y_R)=0,\qquad
 \delta_{b_R}(y_R)\geq
       |j(R,R)-j(0,R)|\geq\frac{R}{1+\sqrt2R}>\frac25,\\
 d^\uparrow(b_R,1)>\frac25 .
 \end{gathered}                                                \tag{C.92.21}
\]
For the strict scalar bound, put $R=1+s$, $s\geq0$.
Then $5R-2>0$ and
$(5R-2)^2-8R^2=1+14s+17s^2>0$.
The point-fixing lift in (C.92.10) cannot be replaced by a
different deck translate: a nonidentity translate sends $x$
to a different point of its fiber.
This is a sequence on the exact image $\mathscr B_p$ for
which the inverse compact-open comparison fails to be continuous.
It also fails the compact-family tail (C.92.17) on the
convergent-sequence parameter space, since the $y_R$ escape
all compact sets.

The proved comparison is consequently the original-image
description (C.92.13)--(C.92.17), together with pointed isotopy
of every actual descent. A deformation of a noncompact
quotient family still requires (C.92.17) for the actual
lift and its inverse throughout the deformation.
No such deformation is constructed in this note.
C.91 continues to assert the whole-group contraction under
its existing compact-quotient hypothesis.

\end{proof}

\section{D. Profinite extensions and the limits of the realization problem}

We write $P=\widehat\pi_0$, where $\pi_0=\pi_1(U,p)$ for a connected orientable surface of finite type with $\chi(U)<0$. All maps between profinite groups in this section are continuous. A lift of an outer action means a homomorphism inducing the prescribed outer homomorphism, with its labels in the operator group retained. For a finite group $G$ acting faithfully on $U$, $E$ denotes the group of lifts to the universal cover, not a mapping class group. The completion map $E\to\widehat E$ is injective: $E$ is a finitely generated linear hyperbolic orbifold group and is residually finite. We use $E^+$ for its orientation-preserving subgroup. The passages answered here are \S31, folios 198--199, scans 391--394; the inserted $G=G^+$ on scan 391 is material.

\begin{claim}{D.1}{Proved}{Exact completion of the lift extension}
The completion of $1\to\pi_0\to E\xrightarrow qG\to1$ is the exact sequence
\[
1\longrightarrow P\longrightarrow\widehat E\xrightarrow{\widehat q}G\longrightarrow1.
\]
It is the extension attached to the completed outer action. An invariant discretization therefore supplies a canonical comparison of extensions, with identity maps on $P$ and $G$.
\end{claim}
\passage{31; 198--199 verso; 391--394}

\begin{proof}
Every subgroup of finite index in $\pi_0$ contains a subgroup normal and of finite index in $E$: take its core in $E$, using $[E:\pi_0]<\infty$. Thus $\pi_0$ has its full profinite topology induced from $E$, and its closure is $P$. The kernel of $\widehat q$ is that closure, since $\pi_0=q^{-1}(1)$ is open. Conjugation in the completed extension is the continuous extension of conjugation in the dense subgroup $E$. The uniqueness of a center-free kernel extension, proved in Part A, identifies this extension with the pullback of $\Aut(P)\to\Out(P)$. The identification fixes the named kernel and quotient, so the assertion of canonicity has no unmentioned conjugator.
\end{proof}

\begin{claim}{D.2}{Proved}{Finite subgroups of completed surface orbifold groups, including reflections}
Every finite subgroup of $\widehat E$ is conjugate to a finite subgroup of $E$. Two finite subgroups of $E$ conjugate in $\widehat E$ are conjugate in $E$. These finite subgroups are cyclic rotation groups or dihedral groups, with a single reflection included as the dihedral case of order two.
\end{claim}
\passage{31; 198--199 verso; 391--394}

\begin{proof}
We specify the two inputs and the reduction. Bridson--Conder--Reid, Theorem 5.1 \cite{R-BCR}, applies to the finitely generated Fuchsian group $E^+$: every finite subgroup of $\widehat{E^+}$ is conjugate into $E^+$. Boggi--Zalesskii, Theorem 6.3 \cite{R-BZ}, gives conjugacy and normalizer comparison for finite $p$-subgroups of $E$:
\[
N_{\widehat E}(A)=\overline{N_E(A)}.
\]
Here its hypotheses hold because $E$ has a finite-index surface group (a free group in the noncompact case), hence is virtually compact special; its good-group and finiteness hypotheses are precisely those in Theorem 6.3. We use only this prime-power theorem, not an assertion obtained by commuting arbitrary intersections with closure.

Put $K^+=K\cap\widehat{E^+}$ for a finite $K\leq\widehat E$. The first input makes $K^+$ cyclic. When $K^+=1$, the orientation character embeds $K$ in $C_2$, and the prime-power input handles $K$. Otherwise choose the unique subgroup $A$ of order $p$ in $K^+$, for one prime $p\mid |K^+|$. It is characteristic in $K^+$ and normal in $K$. Conjugate $A$ into $E^+$. A nontrivial elliptic rotation of the hyperbolic plane has one fixed point $x$. Its normalizer in the discrete group $E$ stabilizes $x$, and conversely every element of $E_x$ normalizes the unique subgroup of order $p$ in its cyclic rotation subgroup. Thus $N_E(A)=E_x$ is finite cyclic or dihedral. The displayed comparison gives $N_{\widehat E}(A)=E_x$, so $K\leq E_x\leq E$.

For injectivity, suppose $K,L\leq E$ and $tKt^{-1}=L$, with $t\in\widehat E$. The orientation character gives $tK^+t^{-1}=L^+$. In the case $K^+=1$ apply prime-power conjugacy separation. Otherwise choose $A\leq K^+$ as above and $B=tAt^{-1}\leq L^+$. Prime-power conjugacy separation supplies $e\in E$ with $eAe^{-1}=B$. Now $e^{-1}t$ normalizes $A$, hence belongs to $N_E(A)\subset E$. Therefore $t\in E$, after this adjustment. This proves the required conjugacy assertion. Finally finite discrete hyperbolic-plane isometry groups stabilize a point and embed in $O(2)$; the listed cyclic and dihedral possibilities follow.
\end{proof}

\begin{claim}{D.3}{Proved}{The comparison of full and partial section classes}
For every subgroup $H\leq G$, completion induces a bijection
\[
\operatorname{Sec}(E_H\to H)/\pi_0
\ \longrightarrow\
\operatorname{Sec}(\widehat E_H\to H)/P,
\qquad E_H=q^{-1}(H).
\]
Both sets are finite. In particular the first conjecture on scan 393 holds for geometrically realized actions, with orientation-reversing operators allowed.
\end{claim}
\passage{31; 198--199 verso; 391--394; 18; 99--101; 197--203}

\begin{proof}
We prove the change from total-group conjugacy to kernel conjugacy, since omission of this step changes the question. A profinite section has a finite image $K$ mapping isomorphically to $H$. Apply D.2 to $E_H$, and choose $t\in\widehat E_H$ with $tKt^{-1}\leq E_H$. Choose $e\in E_H$ with $q(e)=\widehat q(t)$. Then $e^{-1}t\in P$, and conjugation by this element puts $K$ in $E_H$. The restriction of $q$ identifies the resulting subgroup with $H$, so it defines a discrete section with the correct labels.

Conversely, let discrete section images $K,L$ be conjugate by an element of $P$. D.2, applied in $E_H$, gives an $e\in E_H$ with $eKe^{-1}=L$. Let $l\in L$ be the unique element satisfying $q(l)=q(e)$. Then $l^{-1}e\in\pi_0$ still conjugates $K$ to $L$. Because the projection is injective on each section image, this equality of subgroups is equality of the labelled section homomorphisms after conjugation. This proves injectivity.

There are finitely many $E_H$-conjugacy classes of finite subgroups: they occur in the finite list of cone stabilizers, corner stabilizers, and mirror components of a finite-type quotient orbifold, and each finite stabilizer has finitely many subgroups. Each such class splits into at most $[E_H:\pi_0]=|H|$ kernel-conjugacy classes. The asserted finiteness follows.
\end{proof}

\begin{claim}{D.4}{Proved}{The exact profinite fixed subgroup of a section}
Let $s:H\to E_H$ be a nontrivial discrete section. Its fixed subgroup $P^{s(H)}=C_P(s(H))$ is trivial when $s(H)$ contains a nontrivial rotation. For a reflection $s(H)=\langle t\rangle$, it is $1$ for a noncompact real component and is the closure of the primitive oval subgroup, isomorphic to $\widehat\Z$, for a compact real component. The same assertions hold for every profinite section by D.3.
\end{claim}
\passage{31; 199; 393; 33; 222 verso--225; 440--445}

\begin{proof}
For a rotation choose $A$ of prime order inside $s(H)^+$. D.2's proof gives $N_{\widehat E}(A)=E_x$, a finite group. Thus $C_P(s(H))\leq E_x\cap P=1$, because $P$ is torsion-free. For a reflection the prime-power centralizer comparison gives
\[
C_{\widehat E}(t)=\overline{C_E(t)},\qquad
C_P(t)=\overline{C_{\pi_0}(t)}.
\]
The second equality uses that $P$ is open and $P\cap E=\pi_0$: the intersection of the closure of a subgroup with an open subgroup is the closure of their intersection. On the universal cover, $C_{\pi_0}(t)$ is the group of deck transformations preserving and commuting with reflection in its fixed line. It acts freely on that line. Its quotient is the corresponding real component of $U$. The deck group is trivial for a line and infinite cyclic for a circle; the circle is an essential primitive simple curve. Its induced profinite topology is full, so its closure is $\widehat\Z$.

This proves, rather than assumes, the correction to the expected answer in Part D(a): a dihedral section with a nontrivial rotation has fixed subgroup $1$. The value $\widehat\Z$ belongs to a reflection in an oval. It is not the fixed subgroup of a larger dihedral stabilizer containing that reflection.
\end{proof}

\begin{claim}{D.5}{Refuted}{Unique maximal sections require a rotation}
A finite subgroup of $\widehat E$ containing a nontrivial rotation is contained in a unique maximal finite subgroup. The assertion for every nontrivial finite subgroup is false, already for type $(0,3)$ with reflections.
\end{claim}
\passage{18; 99--101; 197--203; 31; 198--199; 391--393}

\begin{proof}
For the positive assertion conjugate the subgroup into $E$ using D.2. Its rotation part has a subgroup $A$ of prime order. Every finite overgroup has cyclic orientation-preserving part, so normalizes this same $A$. All finite overgroups therefore lie in the finite group $N_{\widehat E}(A)=E_x$, which is itself the unique maximal one.

For the counterexample use the $(2,3,\infty)$ reflection group
\[
W=\langle s_0,s_1,s_2\mid s_i^2=1,
(s_0s_1)^2=(s_1s_2)^3=1\rangle
       = V_4*_{C_2}S_3.
\]
Map it onto $S_3\times C_2$, with $r^3=a^2=\tau^2=1$, $ara=r^{-1}$ and $\tau$ central, by
$s_0\mapsto\tau$, $s_1\mapsto a\tau$, $s_2\mapsto ra\tau$.
The vertex groups inject. The kernel is an orientable torsion-free group of index $12$, Euler characteristic $-1$, with three cusps; hence it is $\pi_{0,3}$. The two maximal finite subgroups
$\langle s_0,s_1\rangle$ and $\langle s_1,s_2\rangle$ contain the same reflection $s_1$ and have distinct images in the quotient. They remain distinct and are not kernel-conjugate after completion. Their generated subgroup is $W$, so no finite overgroup contains both. Part B supplies the finite certificate and the geometric kernel computation.
\end{proof}

\begin{claim}{D.6}{Proved}{A well-defined transfer of the geometric rotation invariant}
For a realized $E$, its rotation invariant on finite-order elements transfers uniquely to the torsion of $\widehat E$ by conjugating each torsion element into $E$. This assertion uses the orientation module, so that conjugation by an orientation-reversing element acts by $-1$ on its values. Its values embed canonically in $\widehat T\otimes_{\Z}\Q/\Z$, where $\widehat T=T\otimes_{\Z}\widehat\Z$.
\end{claim}
\passage{17; 77--79 verso; 155--160; 31; 198--199; 391--393}

\begin{proof}
Use D.2 for existence. Explicitly, for torsion $h\in\widehat E$ choose $t\in\widehat E$ with $u=tht^{-1}\in E$ and set
\[
 \mu_{\widehat E}(h)=\varepsilon(t)\mu_E(u),
\]
where $\varepsilon:\widehat E\to\{\pm1\}$ is the orientation character, acting on $\widehat T$. Suppose $u,v\in E$ and $v=tut^{-1}$ in $\widehat E$. For a rotation apply the proof of D.2 to an order-$p$ subgroup of $\langle u\rangle$: first align it by discrete conjugation, then put the remaining conjugator in the finite discrete stabilizer. Thus the elements themselves, and not just their cyclic subgroups, are discretely conjugate, with the same orientation character of the conjugator. For reflections apply prime-power element conjugacy separation, or D.2 because an order-two subgroup has a unique nonidentity element. This establishes independence with the indicated action on $T$. A rotation of order $n$ has a value in $T/nT\subset T\otimes\Q/\Z$; the identification $T/nT=\widehat T/n\widehat T$ defines the claimed inclusion.
\end{proof}

% Replacement for D.7.  The bibliography entry bWilkes2017Seifert is supplied
% in d_mu_references.bib.  This fragment has no global preamble.
\begin{claim}{D.7}{Proved}{Equivariance of the profinite rotation invariant}
Let $E$ be the lift group of a finite action on an orientable hyperbolic
surface of finite type, and put $P=\widehat\pi\triangleleft\widehat E$.
Every automorphism $\alpha$ of $\widehat E$ preserving $P$ and its indexed
peripheral structure satisfies
\[
 \mu_{\widehat E}(\alpha(u))=\alpha_T\bigl(\mu_{\widehat E}(u)\bigr)
 \qquad (u\in\widehat E\text{ of finite order}).
\]
Here $\alpha_T$ is the induced automorphism of the orientation line
$\widehat T$ defined in A.2, and $\mu_{\widehat E}$ is the invariant of
D.6.  The assertion holds for compact and noncompact surfaces, and for
actions containing orientation-reversing elements.
\end{claim}
\passage{17; 77--79 verso; 155--160; 31; 198--199; 391--393}

\begin{proof}
Choose an orientation, so that $\alpha_T$ is multiplication by some
$\lambda\in\widehat\Z^\times$.  Conjugation by an element of
$\widehat E$ acts on $\widehat T$ by its geometric orientation sign.
Naturality of this action and commutativity of the automorphism group of a
rank-one module show that $\alpha$ preserves that sign.  Thus $\alpha$
restricts to an automorphism of $\widehat B$, where $B=E^+$ is the
orientable Fuchsian subgroup.  A nonidentity negative torsion element is a
reflection and has invariant $1/2$; multiplication by a profinite unit
fixes this value.  It remains to treat torsion in $\widehat B$.

We first make explicit the behavior of orientation under finite covers.
For a degree-$d$ connected finite cover of oriented surfaces, the map on
orientation homology takes the fundamental class upstairs to $d$ times
the fundamental class downstairs.  For punctured surfaces this follows
directly in A.2's kernel model: the peripheral factors above any one
puncture map by their positive local degrees, whose sum is $d$.
For closed surfaces it is the degree formula for fundamental classes.
The same formulas hold over $\widehat\Z$, either by this peripheral
matrix or by the natural finite-coefficient cohomology comparison.
Consequently, if $\alpha$ carries an open surface subgroup $V$ to $V'$,
the multipliers on their orientation lines and on $P$ satisfy
$d\lambda_V=d\lambda$.  Multiplication by a positive integer is
injective on $\widehat\Z=\prod_p\Z_p$, so $\lambda_V=\lambda$.
This proves the required compatibility with further finite covers; it
does not divide by $d$ in a finite coefficient ring.

\emph{The noncompact case.}
Let $L$ be a peripheral subgroup of $\pi$, and let $K$ be the primitive
cusp subgroup of $B$ containing it.  Write $K=\langle w\rangle$ and
$L=\langle w^n\rangle$, with compatible positive generators.
We have
\[
 N_{\widehat B}(\overline L)=\overline K\simeq\widehat\Z.
\]
Here is the reduction from Part A, rather than an assumption that arbitrary
normalizers commute with completion.  For an element of the normalizer,
choose a representative in $B$ of its coset modulo $P$.  That coset
stabilizes the profinite peripheral class.  Distinct discrete peripheral
classes remain distinct in $P$, so the same coset stabilizes the discrete
peripheral class.  Adjust its representative by an element of $\pi$ to
put it in $N_B(L)=K$.  The remaining factor belongs to
$N_P(\overline L)=\overline L$, as proved in Part A.  This proves the
displayed equality.

The loop condition on $\alpha|_P$ sends a positive generator of
$\overline L$ to a $P$-conjugate of the $\lambda$th power of the
corresponding generator of another peripheral subgroup $\overline L'$.
After making this conjugation, $\alpha$ sends $\overline K$ to
$\overline K'$ and has the form $w\mapsto(w')^a$.
The indices
$[\overline K:\overline L]=[\overline K':\overline L']=n$
are equal because $\alpha$ preserves $P$.  Thus $na=n\lambda$ in
$\widehat\Z$, and hence $a=\lambda$.  Every primitive orbifold cusp
therefore has precisely the same multiplier $\lambda$.

Write the signature of $B$ as $(h;m_1,\ldots,m_r;s)$ with $s\geq1$.
Choose cone generators $c_j$, of orders $m_j$, and cusp generators $w_i$
in a common orientation.  The abelianized presentation is
\[
 \widehat B^{\mathrm{ab}}
 =\left(\widehat\Z^{2h}\oplus
       \bigoplus_{j=1}^r\Z/m_j\,c_j\oplus
       \bigoplus_{i=1}^s\widehat\Z\,w_i\right)
       \big/\left\langle\sum_jc_j+\sum_iw_i\right\rangle.
\]
D.2 applied to $B$ makes $\alpha$ permute the maximal cone subgroups.
Hence, up to conjugacy in $\widehat B$,
$\alpha(c_j)=c_{\tau(j)}^{k_j}$, where
$m_{\tau(j)}=m_j$ and $k_j$ is a unit modulo $m_j$.
The preceding paragraph gives
$\alpha(w_i)\sim w_{\sigma(i)}^\lambda$.
Apply $\alpha$ to the displayed relation and subtract $\lambda$ times
that relation.  This yields
\[
 \sum_j(k_j-\lambda)c_{\tau(j)}=0.
\]
Eliminating $w_s$ exhibits the cone factors as an independent direct sum,
so $k_j\equiv\lambda\pmod{m_j}$ for every $j$.
The local formula of B.12 now proves the assertion on every cone
subgroup and, by conjugation in $\widehat B$, on all its torsion.

\emph{The compact case.}
First suppose that there are at least two cone points.  We use Wilkes's
cohomological calculation with its exact scope:
for a closed orientable hyperbolic orbifold group $B$ with cone
generators $c_j$ of orders $m_j$, an automorphism of $\widehat B$
sending each $c_j$ to a conjugate of $c_j^{k_j}$ acts on
$H^2(\widehat B;\Z/n)$, for every $n$, by one scalar
$\kappa\in\widehat\Z$ satisfying
$\kappa\equiv k_j\pmod{m_j}$.
This is \cite[Proposition 5.7]{bWilkes2017Seifert}, numbered Proposition
6.6 in arXiv:1512.05587v1.  Its proof concerns arbitrary automorphisms
of the indicated form; no lifting to a Seifert group is used in that
calculation.  In particular the scalar is independent of $n$.

For transparency, the part of the calculation that supplies both facts
is as follows.  After taking coinvariants, the relation cells $y_0,y_j$
have boundaries
$\partial y_0=\sum_jx_j$ and $\partial y_j=m_jx_j$.
The automorphism sends $x_j$ to $k_jx_j$, $y_j$ to $k_jy_j$, and
$y_0$ to $\kappa y_0+\sum_j\rho_jy_j$.
Commutation with the boundary gives
$\kappa+m_j\rho_j=k_j$ in $\widehat\Z$.
The dual coboundaries are $y_0^*+m_jy_j^*$, which shows that the
induced action is multiplication by $\kappa$ for every finite
coefficient group.  The completed resolution and the naturality of this
calculation are established in the cited proposition and its use of
Proposition 2.9 there.

The cone permutation of our $\alpha$ is realized by an
orientation-preserving homeomorphism of the underlying surface carrying
the marked cone points to the correspondingly ordered points.  It
induces an orbifold automorphism $\beta$ with the same permutation and
power $1$ on each cone.  Put $\psi=\beta^{-1}\alpha$ and apply the
preceding result to $\psi$.  Let
\[
 \eta_n\in H^2(\widehat B;\Z/n)
\]
be the image of the oriented fundamental cohomology class of the
underlying closed surface under the canonical map from the orbifold
classifying space to its underlying space.  Equivalently, in the
displayed resolution it is $[y_0^*]$, up to the common orientation sign.
The finite-coefficient class is transported to continuous cohomology by
goodness of Fuchsian groups, as in Wilkes's calculation.
Since $\beta$ preserves orientation, $\beta^*\eta_n=\eta_n$.
Thus $\alpha^*\eta_n=\psi^*\beta^*\eta_n=\kappa\eta_n$.
This argument does not require $\beta$ to preserve the chosen subgroup
$P$.

Set $N=[B:\pi]$.  Restriction to the surface subgroup gives
\[
 \operatorname{res}_{P}(\eta_n)=N\omega_n,
 \qquad\omega_n\in H^2(P;\Z/n)
\]
with $\omega_n$ the positive surface class: the underlying map from
the covering surface to the orbifold's underlying surface is an
orientation-preserving branched map of degree $N$.  Naturality of
restriction and A.2's duality yield
\[
 N\kappa\,\omega_n=N\lambda\,\omega_n
       \quad\text{for every }n.
\]
Indeed $\alpha_*[U]=\lambda[U]$ implies
$\langle\alpha^*\omega_n,[U]\rangle
=\langle\omega_n,\alpha_*[U]\rangle=\lambda$.
Taking all finite coefficients gives $N(\kappa-\lambda)=0$ in
$\widehat\Z$, hence $\kappa=\lambda$.  The local congruences are
therefore $k_j\equiv\lambda\pmod{m_j}$, exactly as required.

We remove the single-cone case from the resolution input explicitly.
When $r=1$, the underlying genus $h$ is positive and
$B^{\mathrm{ab}}=\Z^{2h}$, because the one cone generator is zero in
abelianization.  The characteristic subgroup
\[
 B_1=\ker\bigl(B\longrightarrow B^{\mathrm{ab}}/2B^{\mathrm{ab}}\bigr)
\]
has index $2^{2h}$ and contains every conjugate of the cone group.
Its orbifold has $2^{2h}$ cone points of the same order.  The subgroup
$\widehat{B_1}$ and $P\cap\widehat{B_1}$ are $\alpha$-invariant.
The argument for at least two cones applies there, and the finite-cover
calculation at the beginning identifies its surface multiplier with
$\lambda$.  For $r=0$ there is no nonidentity torsion in $\widehat B$,
by D.2.  Thus no exceptional compact signature remains.

Finally, choose positive local cone generators, for which
$\mu(c_j)=1/m_j$ in the chosen orientation line.  The proved congruence
says $\mu(\alpha(c_j))=\lambda/m_j$; its powers have the same
property.  Every positive torsion element is conjugate in $\widehat B$
to one of these powers by D.2.  Together with the reflection case this
proves the stated equation on all torsion of $\widehat E$.
\end{proof}

\noindent\emph{Why the earlier rank-one argument was insufficient.}
Integral relative orbifold cohomology can have a torsion summand:
for $C_3*C_3$ with cusp $\langle ab\rangle$ it has the presentation
$\langle u,v,w\mid3u=w=3v\rangle\cong\Z\oplus\Z/3$.
The proof above uses the cusp relation in the noncompact case and the
action on the full finite-coefficient cohomology in the compact case.
It therefore supplies the missing local-power congruences without
choosing an unproved canonical generator through that torsion summand.

\subsection{Completion of geometric mapping class groups and full loop automorphisms}

\begin{claim}{D.8}{Proved}{All finite subgroups in the geometric completion in the three requested types}
Let $M=\Gamma_{g,\nu}$ include puncture permutations and orientation reversal. For $(g,\nu)=(0,3),(0,4),(1,1)$, completion induces a bijection on conjugacy classes of finite subgroups of $M$ and $\widehat M$. Every finite subgroup of $\widehat M$ is solvable. For every finite $K\leq M$,
\[
 N_{\widehat M}(K)=\overline{N_M(K)},\qquad
 C_{\widehat M}(K)=\overline{C_M(K)}.
\]
These assertions concern $\widehat M$ with its given embedding in $\Out_{\rm lac}(P)$.
\end{claim}
\passage{30; 194 verso; 384; 30; 196 verso; 388}

\begin{proof}
For the three-punctured sphere $M=S_3\times C_2$ is finite. For the punctured torus $M=\mathrm{GL}_2(\Z)$; its principal congruence subgroup of level $3$ is a torsion-free free group of finite rank. For the four-punctured sphere the subgroup fixing each puncture and preserving orientation is a free group of rank two and has finite index. Thus all three $M$ are finitely generated virtually free groups.

Bridson--Piwek, Proposition 5.10 \cite{R-BP}, states for every finitely generated virtually free group that finite subgroup conjugacy classes are preserved bijectively by profinite completion. Its proof uses the finite graph of finite groups for $M$: completion gives the profinite fundamental group of the same graph, a finite subgroup fixes a vertex of its profinite Bass--Serre tree, and subgroup conjugacy separability distinguishes the resulting discrete subgroups. This applies without an orientation restriction and proves the first assertion. The normalizer equality is Ribes--Zalesskii, Theorem 2.6 \cite{R-RZ}, for finitely generated subgroups of finitely generated virtually free groups. To deduce the centralizer equality, the homomorphism $N_{\widehat M}(K)\to\Aut(K)$ has finite discrete target. Its kernel is open; the intersection of a dense subgroup with this open kernel is dense in the kernel. This is exactly $C_M(K)$.

For solvability, a finite subgroup of $M_{0,4}$ embeds in $S_4\times C_2$ by the torsion-free pure oriented kernel; one in $\mathrm{GL}_2(\Z)$ embeds in $\mathrm{GL}_2(\F_3)$ by the level-$3$ kernel. The former group is solvable. The latter has normal series with factors inside $C_2$, $A_4$, and the central $C_2$, since $\mathrm{PGL}_2(\F_3)=S_4$; it too is solvable. The finite case $S_3\times C_2$ is immediate. Transfer by the proved conjugacy bijection gives solvability in $\widehat M$.
\end{proof}

\begin{claim}{D.9}{Proved}{The nonsolvable discrete possibilities in genus at most two}
A nonsolvable finite subgroup of the extended mapping class group of a punctured orientable surface of genus at most two is realized on a sphere and has group $A_5$ or $A_5\times C_2$. Its puncture set is invariant under the corresponding icosahedral action. In particular it has at least twelve punctures. This classification does not assert the absence of additional nonsolvable finite subgroups in completions for all larger puncture numbers.
\end{claim}
\passage{30; 194 verso; 384}

\begin{proof}
Realize the finite subgroup by Part C. On the sphere its orientation-preserving subgroup is cyclic, dihedral, $A_4$, $S_4$, or $A_5$, by the finite rotation-group classification. The first four are solvable and an extension of one of them by an orientation group of order at most two is solvable. The orthogonal normalizer of the icosahedral rotation group is $A_5\times\{\pm I\}$: a normalizing orthogonal transformation preserves its rotation axes and its icosahedral cellulation, and composition with an icosahedral rotation leaves either $I$ or $-I$. A nonsolvable full group is consequently $A_5$ or $A_5\times C_2$. The smallest nontrivial orbit has twelve vertices; an invariant puncture set has at least twelve points.

On a torus the kernel of the action on $H_1$ consists of a finite translation group, which is abelian. The orientation-preserving linear quotient is cyclic of order $1,2,3,4$, or $6$, and the extended linear quotient is cyclic or dihedral. These groups and their extensions by the translation group are solvable.

In genus two the hyperelliptic involution is central in the full automorphism group of the invariant Riemann surface. Adjoin it when it is absent. The quotient is a sphere with six branch points. A nonsolvable orientation-preserving quotient would contain $A_5$, whose smallest orbit has twelve points, contradicting preservation of the six branch points. The quotient is therefore solvable; so is its central extension by $C_2$ and its extension by orientation reversal. Punctures can only reduce the finite group. This proves the asserted list and its restriction to the discrete setting.
\end{proof}

\begin{claim}{D.10}{Proved}{Finite torsion in the odd-prime normalized pro-$p$ centralizer}
Let $p$ be odd, let $F^{(p)}$ be the free pro-$p$ group on $x,y$, and define $GT_{0,p}$ by the normalized automorphisms
\[
 F_{\lambda,f}(x)=x^\lambda,\qquad
 F_{\lambda,f}(y)=f^{-1}y^\lambda f,
 \quad \lambda\in\Z_p^\times,\quad f\in(F^{(p)})',
\]
satisfying the two-cycle and three-cycle equations. Every finite subgroup of $GT_{0,p}$ is cyclic of order dividing $p-1$. There is a subgroup of order $p-1$, and every finite subgroup is conjugate in $GT_{0,p}$ to a subgroup of it. For $p\geq5$ this supplies finite loop-preserving outer automorphism groups which are not conjugate to geometric groups of the three-punctured sphere.
\end{claim}
\passage{25 bis; 157 verso--160; 312--317; 30; 194 verso; 384}

\begin{proof}
The normalized description is faithful: a conjugator between two normalized automorphisms first belongs to $C(x)=\langle x\rangle$; the normalization $f\in F'$ and the $y$ equation remove its exponent. The same observation identifies these automorphisms with the pure outer centralizer of the three puncture permutations. The equations are closed conditions, composition is the usual substitution law, and their inverses satisfy the same equations. Hence this is a profinite group. The Galois action on $\mathbf P^1_{\overline\Q}-\{0,1,\infty\}$ gives elements of this group with multiplier the $p$-adic cyclotomic character. This character surjects onto $\Z_p^\times$: the finite cyclotomic extensions have Galois groups $(\Z/p^n)^\times$, and their inverse limit is $\Z_p^\times$.

The kernel of $\Aut(F^{(p)})\to\Aut(F^{(p)}/\Phi(F^{(p)}))$ is pro-$p$; the exact statement and its application are also given in Herfort--Zalesskii, proof of Theorem 6.3 \cite{R-HZ}. For the normalized group the action on this Frattini quotient is the scalar $\lambda\bmod p$. Thus the kernel of
\[
 GT_{0,p}\longrightarrow\F_p^\times
\]
is pro-$p$, and the map is surjective. The profinite Schur--Zassenhaus theorem supplies a complement $C_{p-1}$ and conjugacy of complements in the inverse image of any subgroup of $\F_p^\times$. One can see the passage from the finite theorem directly: apply the finite theorem to every quotient by an open normal subgroup of the pro-$p$ kernel, then take the inverse limit of the nonempty finite sets of compatible complements and conjugators. The finite intersection property follows by passing to a common finite quotient.

We check the necessary absence of torsion in the kernel of the multiplier. A free pro-$p$ group has its injective Magnus map into
$\Z_p\langle\!\langle X,Y\rangle\!\rangle^\times$, with $x\mapsto1+X$, $y\mapsto1+Y$. The map is injective because reduction modulo $p$ is the usual faithful free pro-$p$ Magnus map; it extends continuously from the free group for the $(p,X,Y)$-adic topology on the power-series algebra. The ordinary homogeneous degree is used only in the following leading-term calculation, not as the topology for that extension. An automorphism acting identically on $F^{(p)}_{\rm ab}=\Z_p^2$ acts identically on the degree-one quotient of this power-series algebra. For a nonidentity such automorphism choose the least degree $d\geq2$ of a nonzero correction to $X$ or $Y$. The degree-$d$ correction to its $n$th iterate is $n$ times that correction: substitution of the higher terms can change only degrees above $d$. Each homogeneous coefficient module is a torsion-free $\Z_p$-module, so no positive iterate is the identity. Consequently the multiplier is injective on a finite subgroup.

The torsion subgroup of $\Z_p^\times$ is cyclic of order $p-1$. Therefore every finite subgroup maps injectively to $\F_p^\times$. It is a complement in the inverse image of its image, so the preceding conjugacy theorem proves the full assertion. Finally the geometric pure group of type $(0,3)$ consists of $1$ and complex conjugation, with multipliers $1$ and $-1$. A generator with multiplier a primitive $(p-1)$st root cannot be conjugated into it; the multiplier is unchanged by conjugation even in the full loop-preserving outer group. The first example has order four at $p=5$.
\end{proof}

\begin{claim}{D.11}{Proved}{Executed finite checks and their exact scope}
The program \texttt{code/root/finite\_checks.py} verifies the two-cycle and three-cycle equations in every Heisenberg group $H(\F_p)$ for $p=5,7,11$ and every $\lambda\in\F_p^\times$, using
\[
 f(a,b)=[a,b]^{(\lambda^2-1)/24},\qquad
 m=(\lambda-1)/2.
\]
In particular it verifies an order-four normalized automorphism at $(p,\lambda\bmod p)=(5,2)$ and an order-three one at $(7,2)$. The multipliers of the corresponding infinite finite-order subgroups are Teichm\"uller lifts of these residues; the integer $2\in\Z_5^\times$ itself has infinite order. D.10, not the finite checks alone, proves the existence of compatible infinite pro-$p$ lifts.
\end{claim}
\passage{25 bis; 159 verso--160; 316--317; 30; 194 verso; 384}

\begin{proof}
The program defines all $p^3$ elements as triples and uses
$(a,b,c)(a',b',c')=(a+a',b+b',c+c'+ab')$.
It enumerates inverses and powers and evaluates each side of both equations in that group. Its output records every multiplier and both truth values. Since all commutators are central, the associated automorphism is
$(a,b,c)\mapsto(\lambda a,\lambda b,\lambda^2c)$; the program checks multiplication preservation and computes its order. Division by $24$ is only used at the three displayed primes. No assertion of this formula at $2$ or $3$ is made. At $p=3$, D.10 gives only order-two prime-to-$p$ torsion. At $p=2$, the Magnus argument still proves that a finite subgroup injects into $\{\pm1\}$, but the complement argument does not settle conjugacy of involutions in $GT_{0,2}$.
\end{proof}

\begin{claim}{D.12}{Open}{Torsion in the full profinite $GT_0$ and the full low-type groups}
D.29 refutes the finite geometric-recognition assertion for $(g,\nu)=(0,4)$ by an actual involution with a complete preserved peripheral system. We do not settle whether every nonidentity torsion element of the full profinite $GT_0$ is conjugate to $(-1,1)$, or the full finite-subgroup recognition questions for $(g,\nu)=(0,3),(1,1)$. The geometric-completion assertion in all three low types is proved in D.8, and the odd-prime analogue is decided in D.10. D.32 proves that every finite subgroup of the full pure three-cusp group, including $GT_0$, has order at most two; it gives the faithful map $G\to S_3\times C_2$ for the full three-cusp group. D.34, D.36, D.41 and D.42 identify the global orientation character of every such finite group: every involution has its global $2$-adic sign, and every order-three element has multiplier $1$. D.35 identifies the exact remaining third-cusp conjugator equation, proves full generation from it, and identifies its conjugator sets in the pure group and in $GT_0$. D.38 gives a complete characteristic $A_5^{19}$ quotient test, including proved exclusions of finite candidates from full lifts. D.44 solves the coboundary and the first symmetry equation simultaneously in the full group; the third-cusp equation and hence existence of the conjugator for each actual full action remain unresolved. D.33 proves that full homology detects every finite subgroup in type $(1,1)$ and that its order divides $24$. D.43 proves that every involution in this full punctured-torus group has its global orientation sign. D.45 identifies the exact characteristic binary-icosahedral refinement and its full finite solution counts; the same cusp-triangle survivors pass. D.46 proves that every cyclic order-four outer action on the punctured-torus group has an actual order-four lift, using the three full section classes of its positive square. D.47 checks the exact conjugator coordinates in the binary refinement: precisely 10,125,000 of the 324,000,000 coordinate pairs pass, with the count unchanged by finite invertibility and both symmetry relations. D.48 proves that a full punctured-torus order-three action has exactly two section classes and a torsion-free cusp normalizer whose cube exponent is a 3-adic unit. D.49 computes the standard pure involution centralizer and the exact two-element fibres of its conjugacy orbit map. D.50 classifies the finite two-subgroups of the full punctured-torus group as exactly $1,C_2,C_4,C_2\times C_2,D_8$, exhibits all five, and proves the global multiplier character on every four-group or dihedral group of order eight. D.51 identifies the second torsion-free normal index-three kernel of every order-three torus extension, with its exact type-$(1,3)$ intersection cover and three labelled peripheral normalizers. D.52 gives the twelve possible abstract finite-subgroup types, realizes nine of them, and proves the reflection multiplier for every subgroup of type $S_3,D_{12},D_{24}$. D.53 proves the unique full labelled order-four section class and its exact normalizer. D.54 excludes every nonzero prime-to-three invariant character for a cubic torus action by an explicit peripheral norm comparison in a cyclic cover. D.55 supplies the corresponding square comparison, proves that every positive involution acts as $-I$ at every prime, and proves faithfulness of every primary homology map on every finite subgroup. The finite torus abstract types are consequently exactly the nine realized types $1,C_2,C_3,C_4,C_2\times C_2,C_6,S_3,D_8,D_{12}$; the three extra types in the earlier allowed list are excluded. D.56 computes the prime-section permutations of every cyclic-six action and proves that its full extension has no order-six subgroup or finite-order representative of the outer generator. D.57 proves actual labelled lifts for every four-group and order-eight dihedral action, with the exact three-or-one and one full section counts. D.58 determines the exact negative-involution section counts and their full fixed-subgroup mod-two cohomology, and proves that the two internal cusp-reflection classes fuse under full conjugacy. D.59 identifies the negative cusp section as the class with zero full fixed-subgroup mod-two first cohomology, proves that its entire fixed group has zero ambient abelian image, and proves that every actual cusp-reflection transporter generates the full fixed homology line. Its odd cyclic covers have exactly the displayed cusp cycles and full section counts, and the other fixed group, when present, surjects onto that homology line. D.60 identifies the quotient by the closed normal closure of the cusp reflection with the exact full homology coinvariants, and the quotient by all torsion with $\widehat\Z$. It constructs a section in $F$, and in the two-class case a section inside the other fixed group, giving an actual embedded $\widehat\Z\times C_2$. D.61 proves that the peripheral commutator belongs to the closed derived subgroup of no proper closed subgroup of the full rank-two free group. It consequently gives exact negative-conjugator parameter spaces with no additional generation condition, retaining both cusp-reflection indices and all full unit multipliers. Existence of these parameters, pure-involution orbit exhaustion, full fixed-group recognition, the remaining scalar cubic torus multipliers, and geometric conjugacy remain unresolved; D.70 below settles the quartic multiplier.

D.62 gives the exact full negative-conjugator transporter cosets. D.63--D.65 compute the full standard negative centralizers and their homological characters. D.66 computes the exact scalar-cover norm preimage and primary abelianized-extension splittings. D.67 gives the standard Klein-group and dihedral normalizers and all finite overgroups of the specified two-section negative involution. D.68 identifies the unique invariant double cover of every actual quartic action, its integral homology maps, and the full invariant dihedral quotient. D.69 determines the exact rotation image of the positive-square fixed subgroup in that finite quotient, its two labelled section classes, its fixed-kernel first cohomology and the full integral positive-square action on the nontrivial deck summand. D.70 uses the actual closed filling and its pro-two quotient to exclude every negative odd quartic coordinate by a degree-three restriction contradiction. Hence every order-four full torus element has global multiplier one, and every finite two-subgroup has its multiplier in the two global signs. D.71 constructs the actual cubic Frobenius cover, proves the full integral regular module on its nontrivial deck summand, and determines the image and first mod-three cohomology of the relevant section fixed subgroup. D.72 gives an explicit surjective quasi-isomorphism from the actual normalized pants cochain algebra to the three-section Boolean graded algebra; its acyclic kernel proves the stated strong Massey vanishing and exact finite unitriangular representation image. The arbitrary negative conjugator equations, scalar cubic multiplier and full low-type recognition remain unresolved.
D.73 proves the exact compatible $q$-adic cubic Frobenius quotient and its full fixed-group splitting, the integral module and section constraints at every $q$-power level, the limiting fixed-kernel character vanishing, and an actual two-torsion-generator obstruction. D.74 proves that all three full pants section fixed subgroups are trivial. Its actual finite-coefficient relative duality, full Brown maps and coset-module witness give self-centralizing sections and real projectivity. The full coboundary $h$ in D.35 is consequently unique, including D.44's first symmetry. The exact third-cusp equation with this $h$, arbitrary negative torus conjugacy, and the scalar cubic realization or exclusion remain unresolved.
D.75 proves full triviality of every negative torus cusp-reflection fixed subgroup, uniqueness of the actual cusp-fusion transporter, and both singleton transporter formulas for every full unit. In the one-section-class case its extension is real projective. Its exact remaining equation is the conjugacy (D.75.5); the other fixed group in the two-section-class case retains its full homology surjection without a procyclic recognition conclusion.
D.76 identifies the actual pants section augmentation module with $\mathbb F_2[[E]]^2$, using all finite coefficient modules and compatible finite-quotient isomorphisms. Its specified incidence map has cokernel equal to the augmentation kernel of $\mathbb F_2[[E/H]]$, where $H$ is the closed subgroup generated by the three specified section representatives. Thus its invertibility records exactly $H=E$. Four explicit involutions generate $E$, with the duplicated class and the resulting three-involution subgroup both written out. Generation by the specified triple and the original third-cusp equation remain unresolved.
D.77 constructs a complete three-class involution triple in the standard full pants extension whose two kernel generators give automorphisms of every maximal pro-$p$ quotient, while an actual $A_5$ quotient sends the triple to $D_{10}$ and its kernel to $C_5$. Its specified incidence cokernel maps onto a five-dimensional augmentation module. The same quotient proves failure of conjugacy to every unit power of the original second cusp. That retained cusp condition prevents this example from deciding the distinguished triple of D.76 or the equation (D.74.40).
D.78 proves that every open subgroup of the other negative torus fixed group has mod-two first cohomology of dimension one. Its specified dyadic procyclic subgroup is a Sylow subgroup, and the full fixed group has the actual split coordinates $B=K\rtimes\widehat{\mathbb Z}$ with a pro-odd kernel $K$. The finite-coefficient inflation and section-restriction maps retain $Q^K\hookrightarrow Q$. The original finite-coset invariant map detects every nonidentity element of $K$; its surjectivity and full procyclic recognition remain unresolved.
\end{claim}
\passage{25 bis; 157 verso--160; 312--317; 30; 194 verso; 384}

Here is the exact obstruction reached, rather than an extrapolation from computation. The homomorphism from a free profinite group to the product of its maximal pro-$p$ quotients does not detect all its finite quotients. Our program constructs $A_5$ from a $5$-cycle and a $3$-cycle, enumerates its sixty elements, and verifies that its commutator subgroup has sixty elements. Thus $A_5$ is a nontrivial two-generated finite quotient invisible to every nilpotent quotient. The pro-$p$ complements constructed in D.10 do not give a compatible solution of the defining equations in such quotients.

We also checked the $f=1$ attack: the three-cycle relation becomes the power-product relation of Lochak--Schneps, Lemma 9 \cite{R-LS}; their finite modular quotients force the global alternatives $\lambda=1$ or $-1$. A choice of independent signs in the factors of $\widehat\Z$ is not a solution of this relation. The finite nilpotent search therefore supplies no full-profinite counterexample. D.32 excludes every element of odd prime order in the full pure three-cusp group and every element of finite order greater than two. The smallest unresolved full-$GT_0$ case is therefore an additional involution up to conjugacy; D.34 proves that its multiplier is globally $-1$, supplies its three section classes and a commutator coboundary, and leaves the conjugator inside the specified loop group to be determined. D.23 supplies the exact order-bounded compatibility criterion over all isolated finite shadows. D.24 executes nonnilpotent refinement tests and proves that specific order-three shadows cannot lift with order three; D.25 computes the entire dihedral inverse limit and the exact residual multiplier-one norm equation. D.28 proves the additional prime-order section and fixed-subgroup constraint, including its exact maximal pro-$p$ comparison. The finite-cover argument of D.32 now excludes all odd-order full lifts, beyond those bounded computations; the involution conjugacy problem retains its exact isolated-refinement formulation. D.29 uses an arbitrary Galois pants isomorphism and its inverse on two different pieces; this produces a finite involution in type $(0,4)$ without requiring that pants isomorphism itself have finite order. For genus at most two and arbitrary puncture number, the solvable comparison is supplied by the checked argument in F.7. D.31 now proves that every nonsolvable finite subgroup in these completions is abstractly $A_5$ or $A_5\times C_2$ on a punctured sphere, and gives exactly the geometric positive puncture-orbit sizes. Its proof also excludes every nonsolvable finite subgroup in the genus-one and genus-two completions. What remains is geometric conjugacy of the actual $A_5$ embeddings and the corresponding centralizer comparison of F.7.

The later tests also close two specific attempted constructions of a lower-type counterexample. D.39 constructs an actual rank-two Haran involution but proves, using its unique section class and the trivial centralizer of that section, that it preserves neither hyperbolic loop marking under any full profinite change of basis. D.37 gives the exact four-cusp descent condition for a negative transposition. D.40 excludes the required additional deck symmetry for the actual D.29 example by the unequal homomorphism counts $10200$ and $4920$ in $\operatorname{PSL}_2(\F_{11})$; the argument also excludes deck maps moving the original separator. These are proved obstructions for the specified constructions. They do not provide the conjugator still required in D.35.

\subsection{Descent along a finite covering}

Put $A(P)=\Aut_{\rm lac}(P)$ and let a discretization mean the image of a loop-preserving isomorphism $\widehat{\pi_0}\to P$ restricted to $\pi_0$. Thus it is an actual dense subgroup with its peripheral classes. Define $\operatorname{Discr}(P)$ using actual subgroups, and define $\operatorname{Discr}^{\rm out}(P)$ by quotienting this set by conjugation by $P$. This convention is necessary for an outer action: an outer automorphism acts canonically on the second set. Write $J=\Aut_{\rm lac}(\pi_0)$, $B=\Out_{\rm lac}(P)$ and $M=\Out_{\rm lac}(\pi_0)$ for their embedded images.

\begin{claim}{D.13}{Proved}{Restriction and normal descent with the peripheral step included}
For an open subgroup $V\leq P$, restriction of the corresponding finite covering defines an injective map
\[
 \operatorname{Discr}(P)\longrightarrow\operatorname{Discr}(V),
 \qquad\Delta\longmapsto\Delta\cap V.
\]
For normal $V$, its image consists exactly of the lattices $\Lambda\subset V$ whose outer deck action $P/V\to\Out_{\rm lac}(V)$ is induced by orientation-preserving loop automorphisms of $\Lambda$. The cyclic peripheral subgroups of the descended lattice have closures equal to the specified peripheral subgroups of $P$.
\end{claim}
\passage{28; 178--179; 352--354}

\begin{proof}
We prove the self-normalizer input, including the class-preserving step. Paris \cite[Theorems 1.7 and 3.2]{R-Paris} proves conjugacy $p$-separability for closed orientable surface groups and free groups. Thus two elements of $\Lambda$ conjugate in $\widehat\Lambda$ are already conjugate in $\Lambda$.

For a free group $F_r=\langle x_1,\ldots,x_r\rangle$, let $\varphi$ preserve every conjugacy class. After an inner adjustment we have $\varphi(x_1)=x_1$. Write $\varphi(x_i)=w_ix_iw_i^{-1}$. Choose a reduced representative $w$ for the double coset $\langle x_1\rangle w_i\langle x_i\rangle$, so that $w_i=x_1^awx_i^b$ and $w$ begins with neither $x_1$ nor $x_1^{-1}$ and ends with neither $x_i$ nor $x_i^{-1}$. The word $x_1wx_iw^{-1}$ is cyclically reduced with length $2+2|w|$. Since $\varphi(x_1x_i)$ is conjugate to $x_1x_i$, we obtain $w=1$. Hence
$\varphi(x_i)=x_1^{n_i}x_ix_1^{-n_i}$. For distinct $i,j\geq2$, conjugating $\varphi(x_ix_j)$ by $x_1^{-n_i}$ gives
$x_i x_1^{n_j-n_i}x_jx_1^{n_i-n_j}$; cyclic length is two exactly when $n_i=n_j$. Thus all $n_i$ agree, and $\varphi$ is inner. The rank-two case has only one exponent and gives the same conclusion.

For a closed orientable surface group of genus $g\geq2$, Dehn--Nielsen--Baer realizes the outer class of a class-preserving automorphism by a mapping class $f$. It fixes every oriented closed-curve class and acts identically on $H_1$, so preserves orientation. Choose a finite filling Alexander system of simple curves. The Alexander method (\cite[\S2.3]{A_FarbMargalit2012}) makes the subgroup fixing all those curve classes finite: its possible permutations of intersection points and edges form a finite group; the kernel fixes the resulting graph and is isotopic to the identity on its complementary discs. Therefore $f$ has finite order. Realize the cyclic group generated by $f$ using Part C. For an action of order $d$, rational cochains of the quotient are the invariant cochains upstairs, by an equivariant triangulation and averaging over the finite stabilizers. The quotient's first Betti number is consequently $2g$, because the action on $H_1$ is trivial. Its underlying genus is $g$. Riemann--Hurwitz now gives
\[
 2-2g=d\left(2-2g-\sum_j(1-1/m_j)\right).
\]
For $d>1$ the right side is strictly smaller than $2-2g<0$, a contradiction. Thus $f=1$ and the automorphism is inner.

Now let $h\in N_{\widehat\Lambda}(\Lambda)$. Its conjugation automorphism is class-preserving by the first paragraph, hence is conjugation by an element of $\Lambda$ by the two arguments just given. Dividing $h$ by that element centralizes the dense lattice and therefore all of $\widehat\Lambda$. Part A gives trivial centre, so $h\in\Lambda$. This proves $N_{\widehat\Lambda}(\Lambda)=\Lambda$ with the exact comparisons needed in LM3 B.8.

For normal $V$, a descended lattice $\Delta$ normalizes $\Lambda=\Delta\cap V$. Conversely, an element of $N_P(\Lambda)$ has a coset in the finite group $P/V$. The intersection $N_P(\Lambda)\cap V$ is $\Lambda$ by the just-proved input. Given that the deck action is induced on $\Lambda$, form the pullback discrete extension $E_\Lambda$ of $P/V$ by $\Lambda$ inside $\Aut_{\rm lac}(\Lambda)$. We construct its map into $P$ explicitly. An element of the pullback is a pair $(g,\varphi)$, with $g\in P/V$, $\varphi\in\Aut_{\rm lac}(\Lambda)$, and prescribed outer class. Extend $\varphi$ continuously to $\widehat\varphi\in\Aut(V)$ and choose a lift $p_g\in P$ of $g$. Equality of outer classes gives a $v\in V$ with $\widehat\varphi=\Inn(vp_g)|_V$. This $v$ is unique, since $C_P(V)=1$ by Part A. Send $(g,\varphi)$ to $vp_g$. Changing $p_g$ changes $v$ inversely and leaves the product unchanged; multiplication is respected because both products induce the same conjugation on $V$ and have the same quotient label. This defines an injective homomorphism: an element in its kernel has $g=1$ and induces the identity on the dense subgroup $\Lambda$. Its image is exactly $N_P(\Lambda)$, by restricting a normalizing conjugation to $\Lambda$ and applying the same uniqueness. The discrete group is torsion-free, as a subgroup of $P$. The finite-extension realization theorem of Part C gives a surface lattice $\Delta=E_\Lambda$; its orientation is positive because the deck action on $V$ has multiplier $+1$. Its completion is $P$ by D.1.

We supply the peripheral argument left incomplete in LM3 B.10. Let $C\leq\Lambda$ be a primitive peripheral subgroup, and let $Q=C_P(\overline C)$ be the maximal procyclic peripheral subgroup in $P$ containing $\overline C$. Put $D=C_\Delta(C)$. The quotient surface supplied by Part C makes $D$ its primitive cusp subgroup. We have $\overline D\leq Q$ and
\[
 \overline D\cap V=\overline C=Q\cap V.
\]
The image of $D$ in $P/V$ is the stabilizer of the $\Lambda$-conjugacy class of $C$. To see this, adjust a representative of a stabilizing coset by an element of $\Lambda$ to normalize $C$; orientation preservation changes normalization into centralization. Conversely a centralizing element fixes the class. The identical argument in the profinite group identifies the image of $Q$ with the stabilizer of the $V$-conjugacy class of $\overline C$. The map from discrete cusp classes to their profinite classes is bijective and equivariant by peripheral noncommensurability and the finite covering construction of Part A. These two stabilizers are therefore equal. The two closed subgroups $\overline D\leq Q$ have the same intersection with the open subgroup $V$ and the same image in $P/V$, so they are equal. This proves the promised loop identification.

In particular a descended lattice is forced to equal $N_P(\Lambda)$, proving injectivity for normal $V$. For arbitrary $V$, take $W=\bigcap_{p\in P}pVp^{-1}$, an open normal subgroup, and use the actual intersection map $\Lambda\mapsto\Lambda\cap W$. Equality after restriction to $V$ implies equality after restriction to $W$, hence equality of the original lattices by the normal case. This proves injectivity without assuming $V$ normal.
\end{proof}

\begin{claim}{D.14}{Proved}{The image for a nonnormal finite covering}
Let $V\leq P$ be open, $W=\operatorname{core}_P(V)$ and $\Lambda\in\operatorname{Discr}(V)$. Set $\Lambda_W=\Lambda\cap W$. The image condition for $\Lambda$ is that the full deck action of $P/W$ on $W$ is induced by orientation-preserving loop automorphisms of $\Lambda_W$. This condition is independent of a choice of coset representatives. The unique extension is $N_P(\Lambda_W)$, and its intersection with $V$ is $\Lambda$.
\end{claim}
\passage{28; 178--179; 352--354}

\begin{proof}
Necessity follows by restriction of an actual lattice in $P$. For sufficiency apply the normal descent statement D.13 to $W\triangleleft P$, obtaining $\Delta=N_P(\Lambda_W)$. Both $\Delta\cap V$ and $\Lambda$ are discretizations of $V$ with intersection $\Lambda_W$ with the normal open subgroup $W$ of $V$. The injectivity already proved inside $V$ gives $\Delta\cap V=\Lambda$. The deck action is an outer action, so changing a representative multiplies its automorphism by an inner automorphism of $W$; its membership in the embedded geometric outer subgroup is unaffected. This proves the asserted independence and the exact nonnormal image description.
\end{proof}

\begin{claim}{D.15}{Proved}{The explicit nonextending discretization in scan 354}
Let $U=\mathbf P^1-\{0,1,\infty\}$ and
$U'=\mathbf P^1-\{0,\infty,1,-1\}$, with the double cover
$z\mapsto z^2$. Put $P=\widehat{\pi_1(U)}$, let $V\triangleleft P$ be its index-two subgroup, and let $\Lambda=\pi_1(U')\subset V$ be the geometric discretization. There is an explicit continuous loop automorphism $\gamma$ of $V$ for which $\gamma(\Lambda)$ is a discretization of $V$ that does not extend to $P$, although it has the same prediscretization as $\Lambda$.
\end{claim}
\passage{28; 178--179; 352--354}

\begin{proof}
The deck involution $s:z\mapsto-z$ fixes $0,\infty$ and exchanges $1,-1$. Let $a$ be a simple curve separating $\{0,1\}$ from $\{\infty,-1\}$ and let $b=s(a)$; write $T_a,T_b$ for positive Dehn twists. They give a free basis of the pure oriented mapping class group $M^{!+}_{0,4}\simeq F_2$: in the ordered moduli space $\mathcal M_{0,4}\simeq\mathbf P^1-\{0,1,\infty\}$ they are monodromies around two distinct boundary points. One can verify the completed embedding directly from A.12: filling the fourth puncture identifies this pure group with the point-pushing group $\pi_{0,3}$, and the peripheral-normalization map sends its completion into $\Aut(\widehat\pi_{0,3})$ as $\Inn(\widehat\pi_{0,3})$. The latter is isomorphic to $\widehat\pi_{0,3}$ by triviality of its centre. The composite is injective, so the completion embeds already in $\Out_{\rm lac}(V)$. Thus the pure group's closure is its free profinite completion, with abelianization $\widehat\Z^2$ on the indicated twists.

Choose the fully specified profinite integer
\[
 \alpha_2=0\in\Z_2,\qquad
 \alpha_\ell=1\in\Z_\ell\quad(\ell\ {\rm odd}),
 \qquad \alpha\in\widehat\Z.
\]
It is not an integer: its $2$-adic coordinate would force that integer to be zero, contrary to its $3$-adic coordinate. Let $u=T_a^\alpha$ in the closure of $M^{!+}_{0,4}$. Since $s$ preserves orientation,
$sT_as^{-1}=T_b$, and the continuous power maps give
\[
 u^{-1}su=s\,T_b^{-\alpha}T_a^\alpha .
\]
The pure factor on the right has abelianization $(\alpha,-\alpha)$. Every element of the discrete pure group has integral abelianization. Moreover
$M_{0,4}\cap\overline{M^{!+}_{0,4}}=M^{!+}_{0,4}$: the puncture permutation and orientation sign are detected by finite quotients. Consequently $u^{-1}su$ does not belong to the discrete extended mapping class group.

Choose any continuous automorphism $\gamma\in\Aut_{\rm lac}(V)$ representing $u$. Its ambiguity is conjugation by an element of $V$, which does not change membership in the outer geometric subgroup. The geometric outer subgroup associated to $\gamma(\Lambda)$ is $uM_{0,4}u^{-1}$. The calculation says that the deck involution is absent from this subgroup. The necessary direction of D.13 therefore proves that $\gamma(\Lambda)$ does not extend to $P$. Since $u$ lies in the closure of the pure geometric group, its coset modulo the closed geometric stabilizer is the identity, proving equality of the prediscretizations.

The code in \texttt{code/root/finite\_checks.py} verifies compatible Chinese-remainder residues of $\alpha$ and the puncture partitions interchanged by $s$. Nonintegrality and the infinite completed-group conclusion follow from the coordinate argument above. This spells out and verifies the proposed element in LM3's review, \S4.2.
\end{proof}
\noindent\emph{Manuscript:} \S28, folios 178--179, scans 352--354, especially \emph{Un tel $\gamma$ existe toujours}. We have proved the indicated index-two example; the word \emph{toujours} is not used to assert the construction for every finite covering.


\begin{claim}{D.16}{Proved}{A precise fixed-set formula for every actual action}
Let $A$ be a group, $J\leq A$, and $G\leq A$ a finite subgroup. Under the left action of $G$ on $A/J$,
\[
 (A/J)^G=\{aJ:a^{-1}Ga\leq J\}.
\]
It is the disjoint union of the $C_A(G)$-orbits
\[
 \coprod_{[a]} C_A(G)\big/\bigl(aC_J(a^{-1}Ga)a^{-1}\bigr),
\]
where $[a]$ ranges over $J$-conjugacy classes of the labelled embeddings $g\mapsto a^{-1}ga$ occurring in $J$. This formula applies to actual discretizations with $(A,J)=(A(P),\Aut_{\rm lac}(\pi_0))$, to outer discretizations with $(A,J)=(B,M)$, and to prediscretizations with $(A,J)=(B,\overline M)$.
\end{claim}
\passage{28; 178--181; 352--358}

\begin{proof}
The fixed-coset condition is $gaJ=aJ$ for every $g$, exactly the displayed inclusion. Two fixed cosets $aJ,bJ$ are in the same centralizer orbit exactly when there are $c\in C_A(G)$ and $j\in J$ with $b=caj$. Substitution gives $b^{-1}gb=j^{-1}a^{-1}gaj$ for every labelled $g$. Conversely this equality makes $b j^{-1}a^{-1}$ centralize $G$. The stabilizer of $aJ$ is $C_A(G)\cap aJa^{-1}=aC_J(a^{-1}Ga)a^{-1}$, proving the orbit formula.

The homogeneous-space identifications use the extension of each discrete loop automorphism to completion and its injectivity. The stabilizer of an actual lattice in $A(P)$ is exactly $J$. Its inner-conjugacy quotient is $B/M$. Since the closure of $\Inn(\pi_0)$ is $\Inn(P)$, the closed based stabilizer contains all inner automorphisms and its quotient is $\overline M$; hence the prediscretization set is $B/\overline M$. These identifications also show explicitly why a bare outer element is not assigned an action on the set of actual subgroups without choosing a lift.
\end{proof}

\subsection{Torsion-free extensions and recognition}

\begin{claim}{D.17}{Proved}{Torsion-free extensions with an invariant discretization}
Let $1\to P\to\mathcal E\to G\to1$ be a profinite extension with $G$ finite, $P=\widehat\pi_0$ hyperbolic, and the outer action preserving the loop structure and a discretization. A torsion-free $\mathcal E$ is the profinite completion of a torsion-free discrete surface group, with the descended peripheral system. The quotient surface may be nonorientable.
\end{claim}
\passage{28; 178 verso; 353; 24; 147 verso; 293}

\begin{proof}
First the outer action is faithful. Its kernel is the image of $C_{\mathcal E}(P)$: an element acting inner can be divided by its unique inner element of $P$. The centralizer meets $P$ trivially, so embeds in the finite quotient $G$ and is finite. Torsion-freeness makes it trivial. Construct the discrete extension $E$ of the stabilized lattice with this faithful outer action, using centre-freeness. By D.1, its completion is the given $\mathcal E$ with the identity on kernel and quotient. The discrete extension embeds in its completion because it is the finite extension realized by Part C and hence residually finite. Torsion-freeness of $\mathcal E$ therefore implies torsion-freeness of $E$. In a realization, a nontrivial point stabilizer gives a nontrivial finite-order element of the lift group. Thus $G$ acts freely on the surface associated to $\pi_0$, and its quotient is a surface. Its fundamental group is $E$ by lifting paths. The primitive cusp group in $E$ is the centralizer of the covering cusp group; D.13 proves equality of their closures with the specified procyclic peripherals. Orientation-reversing free deck transformations yield the nonorientable quotient and its orientation character. This is the exact scope of the descent argument.
\end{proof}

\begin{claim}{D.18}{Proved}{The unmarked pro-$p$ extension statement for punctured surfaces}
A finitely generated torsion-free pro-$p$ group $\mathcal E$ containing an open free pro-$p$ group $P$ of rank $r$ is free pro-$p$. Its rank is
\[
 1+\frac{r-1}{[\mathcal E:P]}.
\]
In particular the displayed quantity is an integer. This proves the underlying-group assertion for punctured pro-$p$ surfaces, without identifying the peripheral system. More generally, every torsion-free profinite group with an open normal pro-$p$ subgroup is itself pro-$p$. Thus D.19--D.20 apply to arbitrary finite-quotient extensions of their pro-$p$ kernels; the quotient is forced to be a $p$-group.
\end{claim}
\passage{28; 178 verso; 353; 24; 147 verso; 293}

\begin{proof}
Herfort--Zalesskii, Theorem 1.1 \cite{R-HZ}, writes such a finitely generated virtually free pro-$p$ group as the fundamental pro-$p$ group of a finite graph of finite $p$-groups, with vertex groups embedded in $\mathcal E$. Torsion-freeness makes all vertex and edge groups trivial. The fundamental pro-$p$ group is therefore the free pro-$p$ group on the cycles of the graph. The pro-$p$ Schreier rank formula gives $r-1=[\mathcal E:P](\operatorname{rank}(\mathcal E)-1)$, proving the formula. The theorem retains no designated peripheral subgroups in its graph-of-groups conclusion; their identification is a further assertion, not supplied by this calculation.

We prove the assertion about a general finite quotient. Write $G=\mathcal E/P$ and fix a prime $q\ne p$. For every open normal $U\triangleleft\mathcal E$ contained in $P$, the kernel of $\mathcal E/U\to G$ is a finite $p$-group. A Sylow $q$-subgroup of $\mathcal E/U$ therefore has order $|G|_q$ and maps isomorphically to a Sylow $q$-subgroup of $G$. Fix one such subgroup $Q\leq G$. The finite nonempty sets of Sylow $q$-subgroups above $Q$ admit compatible choices under quotient maps: an image is again a Sylow subgroup, and any specified one lifts by taking a Sylow subgroup of its inverse image. Compactness gives an inverse-limit subgroup of $\mathcal E$ isomorphic to $Q$, since all these transition maps on the chosen subgroups are isomorphisms. Torsion-freeness forces $Q=1$. Hence every prime divisor of $|G|$ is $p$, and every finite quotient of $\mathcal E$ is a $p$-group. This is exactly the assertion that $\mathcal E$ is pro-$p$.
\end{proof}

\begin{claim}{D.19}{Proved}{Recognition of the peripheral pro-$p$ extension}
Let $p$ be any prime, let $P$ be the pro-$p$ completion of an orientable hyperbolic surface group of type $(g,\nu)$ with $\nu>0$, and let
\[
 1\longrightarrow P\longrightarrow\mathcal E\longrightarrow G\longrightarrow1
\]
be an exact sequence of pro-$p$ groups with $G$ finite. Assume that $\mathcal E$ is torsion-free and its outer action preserves the $P$-conjugacy classes of the peripheral subgroups. Choose one $L_j$ from each $\mathcal E$-orbit of these classes. Then
\[
 \mathcal N_j=N_{\mathcal E}(L_j)=C_{\mathcal E}(L_j)\simeq\Z_p,
\]
and the pair $(\mathcal E;\{\mathcal N_j\})$ is continuously isomorphic, preserving the indexed peripheral conjugacy classes, to the pro-$p$ completion of a compact surface with its boundary circles. Its interior is orientable for odd $p$. At $p=2$ the interior may be nonorientable; its orientation character is the induced action on the orientation line of $P$.
\end{claim}
\passage{28; 178 verso; 353; 24; 147 verso; 293}

\begin{proof}
Each $\mathcal N_j$ is closed, and
$\mathcal N_j\cap P=N_P(L_j)=L_j$ by Part A. Thus $\mathcal N_j$ is a torsion-free pro-$p$ group with an open subgroup $\Z_p$. It is finitely generated, by lifting generators of its finite quotient and adjoining a generator of $L_j$. Apply D.18 to $\mathcal N_j$: the free pro-$p$ rank formula gives rank one. Consequently $\mathcal N_j\simeq\Z_p$, and it centralizes $L_j$.

We identify the restricted peripheral family, including its multiplicities. The restriction of $(\mathcal E;\{\mathcal N_j\})$ to $P$ contains the subgroups
\[
 P\cap e\mathcal N_je^{-1}=eL_je^{-1},
 \qquad e\in P\backslash\mathcal E/\mathcal N_j .
\]
The map from this double-coset set to the orbit of the $P$-conjugacy class of $L_j$ is surjective by definition. It is injective: an equality
$e'L_je'^{-1}=a eL_je^{-1}a^{-1}$ with $a\in P$ is equivalent to
$e'^{-1}ae\in N_{\mathcal E}(L_j)=\mathcal N_j$, exactly equality of the double cosets. Thus the restricted family consists of all original peripheral classes, each once.

Wilkes, \cite[Proposition 5.11]{R-WilkesRelative}, proves finite-index ascent and descent of Poincar\'e duality pairs at $p$ for a $p$-torsion-free profinite ambient group, with compatible dualizing modules. The pair $(P;\{L_i\})$ is orientable of dimension two: this follows from its standard free pro-$p$ surface presentation and the converse direction of \cite[Theorem 3.3]{R-Wilkes}. We have checked the exact restricted family and torsion hypothesis. Hence $(\mathcal E;\{\mathcal N_j\})$ is a Poincar\'e duality pair of dimension two.

Its dualizing module has underlying group $\Q_p/\Z_p$, with action through a continuous character $\chi:\mathcal E\to\Z_p^\times$ whose restriction to $P$ is trivial. Naturality of duality under conjugation identifies its character with the action on the surface orientation line, using the inverse convention on the dual cohomology line. The image of $\chi$ is a finite $p$-subgroup because it factors through $G$: it is trivial at odd $p$ and is trivial or $\{\pm1\}$ at $p=2$. Thus inversion of the character does not affect the answer. Its restriction to each $\mathcal N_j$ is trivial, by the boundary-duality restriction in \cite[Proposition 5.9]{R-WilkesRelative} and the trivial duality character of $\Z_p$.

D.18 makes $\mathcal E$ free pro-$p$. Write the peripheral count as $b+1$ and the rank as $n+b$. The condition $n\geq0$ in the classification follows from the relative cohomology sequence: the cokernel of
\[
 H^1(\mathcal E;\F_p)\longrightarrow
       \bigoplus_{j=0}^{b}H^1(\mathcal N_j;\F_p)
\]
has dimension one, since $H^2(\mathcal E;\F_p)=0$ and the relative top group has dimension one. Hence $b\leq\operatorname{rank}(\mathcal E)$.
The classification in \cite[Theorem 3.3]{R-Wilkes} gives a free basis
$x_1,\ldots,x_n,s_1,\ldots,s_b$, with each $s_i$ generating a conjugate of $\mathcal N_i$, and a generator $s_0$ of $\mathcal N_0$ in one of these forms:
\[
\begin{array}{ll}
 \chi=1:&
 s_0=[x_1,x_2]_{-}\cdots[x_{n-1},x_n]_{-}\,s_1\cdots s_b,
 \\[2pt]
 \chi\ne1,\ n\ {\rm odd}:&
 s_0=x_1^2[x_2,x_3]_{-}\cdots[x_{n-1},x_n]_{-}\,s_1\cdots s_b,
 \\[2pt]
 \chi\ne1,\ n\ {\rm even}:&
 s_0=x_1^2[x_1,x_2]_{-}[x_3,x_4]_{-}\cdots
 [x_{n-1},x_n]_{-}\,s_1\cdots s_b .
\end{array}
\]
Here $[x,y]_{-}=x^{-1}y^{-1}xy$ matches that source's convention. The first line has $n$ even. The other two have $p=2$; their parameter $2^f$ is zero, since $\chi$ has no infinite principal-unit part.

These are surface words, with explicit changes of basis. Inverting each handle generator changes the commutator convention in the first line. The second line gives one crosscap and orientable handles. For the third, put
\[
 a=x_1x_2^{-1},\qquad b'=x_2,\qquad
 x_1=ab',\quad x_2=b'.
\]
Then $x_1^2[x_1,x_2]_{-}=a^2(b')^2$, which gives two crosscaps and orientable handles. Replacing $s_0$ by its inverse gives the usual boundary relation and preserves its subgroup. Sending the discrete surface generators to these free pro-$p$ generators induces the asserted isomorphism of completed pairs, with exactly the indexed peripheral conjugacies already specified. The free rank formula gives
$\chi(\mathcal E)=\chi(P)/[\mathcal E:P]<0$, so its interior is hyperbolic. This proves the peripheral statement at every prime, including $2$ and $3$.
\end{proof}

\begin{claim}{D.20}{Proved}{Recognition of the closed pro-$p$ extension}
Let $P$ be the pro-$p$ completion of a closed orientable surface group of genus $g\geq2$. A torsion-free pro-$p$ group $\mathcal E$ containing $P$ as an open normal subgroup is the pro-$p$ completion of a closed surface group. The surface is orientable at odd $p$. At $p=2$ a nonorientable surface is allowed, and its orientation character is the action on the orientation line of $P$.
\end{claim}
\passage{28; 178 verso; 353; 24; 147 verso; 293}

\begin{proof}
We give the absolute finite-index duality calculation. The orientable surface presentation makes $P$ a Poincar\'e duality pro-$p$ group of dimension two, with
\[
 H^i(P,\Z_p[[P]])=
 \begin{cases}\Z_p,&i=2,\\0,&i\ne2,\end{cases}
\]
and trivial right $P$-action on the top group. This is the trivial-character case of the Demushkin classification in \cite[Theorem 1.9]{R-Wilkes}.

Serre's theorem \cite[\S1, Theorem]{R-SerreCD} equates the $p$-cohomological dimensions of an open subgroup and a profinite group without elements of order $p$. Therefore $\operatorname{cd}_p(\mathcal E)=2$. Wilkes \cite[Proposition 1.26]{R-WilkesRelative} gives finite-index ascent and descent of type $p$-$FP_\infty$ for profinite groups. The surface resolution gives this property for $P$, and hence for $\mathcal E$.

At finite index, induction and coinduction of compact modules agree, and continuous Shapiro gives
\[
\begin{split}
 H^i(\mathcal E,\Z_p[[\mathcal E]])
 &=
 H^i\!\left(\mathcal E,
        \operatorname{coind}_{P}^{\mathcal E}\Z_p[[P]]\right)\\
 &\simeq H^i(P,\Z_p[[P]]).
\end{split}
\]
The first coefficient isomorphism and compact Shapiro are \cite[Propositions 1.16 and 1.21, and the remark after 1.21]{R-WilkesRelative}. We make their module compatibility explicit. Put $\Lambda_H=\Z_p[[H]]$ and let $\operatorname{pr}_P:\Lambda_{\mathcal E}\to\Lambda_P$ retain the $P$-coset component. The coefficient isomorphism is
\[
 \Lambda_{\mathcal E}\longrightarrow\operatorname{Coind}_P^{\mathcal E}\Lambda_P,
 \qquad \lambda\longmapsto \bigl(x\mapsto\operatorname{pr}_P(x\lambda)\bigr).
\]
Here coinduction consists of left-$P$-linear functions and its left $\mathcal E$-action is $(e f)(x)=f(xe)$. For representatives $t$ of the finite set $P\backslash\mathcal E$, the inverse is $f\mapsto\sum_t t^{-1}f(t)$. These continuous maps are left-$\mathcal E$-linear and right-$P$-linear. Evaluation at $1$ in Shapiro is natural for the right-$P$ coefficient endomorphisms. Consequently the top right duality character restricts to that of $P$. The cohomology criterion for Poincar\'e duality now makes $\mathcal E$ a duality pro-$p$ group of dimension two with dualizing module $\Q_p/\Z_p$. Its orientation character $\chi:\mathcal E\to\Z_p^\times$ has finite $p$-group image, trivial at odd $p$ and contained in $\{\pm1\}$ at $p=2$.

Its one-dimensional degree-two cohomology and nondegenerate degree-one cup product make it an infinite Demushkin group. The classification \cite[Theorem 1.9]{R-Wilkes} gives, for these characters, exactly the three words of D.19 with the boundary letters omitted. The same invertible free-basis change converts the even-rank nonorientable word to the surface form. Hence an actual closed surface presentation induces a continuous isomorphism of its pro-$p$ completion with $\mathcal E$. Finite-index multiplicativity for the Euler characteristic of the finite two-dimensional resolution gives
\[
 \chi(\mathcal E)=\frac{2-2g}{[\mathcal E:P]}<0.
\]
The resulting surface is hyperbolic, and conjugation-naturality of the dualizing module gives the stated orientation.
\end{proof}

\begin{claim}{D.21}{Refuted}{A full profinite normalizer shortcut}
The assertion that every torsion-free profinite group with an open subgroup $\widehat\Z$ is itself $\widehat\Z$ is false. The pro-$p$ version used in D.19 is proved in D.18.
\end{claim}
\passage{28; 178 verso; 353}

\begin{proof}
Define
\[
 N=(\Z_3\rtimes\Z_2)\times\prod_{\ell\ne2,3}\Z_\ell,
 \qquad (a,b)(a',b')=(a+(-1)^{\,b\bmod2}a',\,b+b')
\]
on the first two factors, and use addition on the remaining factors.
This is a profinite group with the open normal subgroup
\[
 L=(\Z_3\times2\Z_2)\times\prod_{\ell\ne2,3}\Z_\ell
       \simeq\widehat\Z
\]
of index two. A finite-order element projects trivially to the torsion-free quotient
$\Z_2\times\prod_{\ell\ne2,3}\Z_\ell$, and then lies in the torsion-free kernel $\Z_3$; it is therefore the identity. Yet conjugation by $(0,1)$ sends $(1,0)$ to $(-1,0)$, so $N$ is nonabelian. The code in \texttt{code/root/finite\_checks.py} checks its quotient
$C_3\rtimes C_4$, the index-two cyclic subgroup of order six, and this noncommutation. The proof in the inverse limit establishes torsion-freeness; no finite quotient is claimed to be torsion-free.

The multiplier on $L$ is $-1$ in its $3$-adic coordinate and $+1$ in every other coordinate. Thus passing separately to pro-$p$ groups loses precisely this interaction between two different primes. D.29 uses this $N$ as an explicit retract in its full peripheral counterexample. Its cusp normalizers are all procyclic, so that example does not depend on making $N$ itself a cusp normalizer.
\end{proof}

\begin{claim}{D.22}{Refuted}{Full profinite recognition without an invariant discretization}
The assertion that every torsion-free finite extension of a full
profinite hyperbolic surface group with preserved loop structure
is a completed surface pair is false.  D.29 constructs a
counterexample with kernel of type $(0,4)$ and quotient $C_2$,
including the exact peripheral comparisons and a finite quotient
excluding every free profinite presentation. D.30 constructs a
counterexample with closed genus-two kernel and quotient $C_2$;
its degree-two cohomology at two odd primes excludes every
surface completion. Recognition with an
invariant discretization holds by D.17.  The separate pro-$p$
recognition theorems, with their precise group and peripheral
hypotheses, are proved in D.18--D.20.
\end{claim}
\passage{28; 178 verso; 353; 24; 147 verso; 293}

\begin{proof}
The basis-derived four-cusp presentation, faithful outer involution,
continuous twisted extension and torsion-free retraction are given
in D.29.  Its order-eighteen quotient needs three generators, whereas
an index-two extension of a free rank-three group would have rank two
if it were free.  The exact finite-cover comparison excludes closed
surface completions.  Thus the counterexample satisfies the full
peripheral hypothesis and fails the claimed conclusion.

For the closed case, D.30 gives the extension by its continuous
twisted multiplication and by an equivalent profinite presentation.
Its torsion-free maximal pro-$2$ quotient detects the parity of
every element, proving that every torsion element lies in its
torsion-free surface kernel. The two restriction isomorphisms give
$H^2(\mathcal E;\mathbf F_3)=\mathbf F_3$ and
$H^2(\mathcal E;\mathbf F_5)=0$. The cohomology of every orientable,
nonorientable, or punctured surface completion is computed in
D.30 and excludes this pair. This settles the closed hyperbolic
assertion in its smallest genus as well.

This also identifies the failed step in the earlier recognition
attempt.  Separate pro-$p$ comparisons do not produce a single
full profinite surface presentation.  The mixed-prime multiplier
in D.29 is an actual finite-order loop automorphism, and all its
cusp normalizers are procyclic.  The preserved global surface
relation therefore does not restore the proposed conclusion.
The first possible punctured nonabelian rank is three: D.27 excludes
proper torsion-free finite extensions of $\widehat F_2$.
The explicit free-basis comparison in D.29 settles the formerly
undecided peripheral marking on D.26's rank-three extension.

We checked Collas--Philip--Yamaguchi \cite{R-CPY}, Definition 2.1,
Theorems 1.3 and 2.4, against the supplied source.  Their profinite
F-groups already have the stated orbifold presentation.
Theorem 2.4 constructs a torsion-free open normal subgroup, and
Theorem 1.3 bounds the derived length of a finite quotient in its
nonperfect case.  The maps go from a prescribed presentation to
an open subgroup; they do not assert the reverse reconstruction
from an arbitrary finite extension.  D.29 refutes that reverse
assertion under the present full profinite hypotheses.  The
positive conclusions of D.17--D.20 retain their proved maps and
hypotheses throughout.
\end{proof}


\subsection{A full profinite lifting criterion and nonnilpotent finite tests}

\begin{claim}{D.23}{Proved}{Order-preserving lifting and conjugacy in the full shadow system}
Let $\mathcal I$ be the cofiltered poset of isolated finite-index normal
subgroups $N\unlhd B_3$ contained in the pure braid group $PB_3$, and put
$Q_N=GT(N)$.  The finite groups and their reduction homomorphisms
$R_{N,H}:Q_N\to Q_H$ for $N\subseteq H$ are those of
\cite[Proposition 3.14, Theorem 5.2 and Proposition 5.5]{A-DGShadows}.
We use the continuous isomorphism
\[
 \Psi:GT_0\xrightarrow{\ \sim\ }\varprojlim_{N\in\mathcal I}Q_N,
 \qquad (\lambda,f)\longmapsto
 \bigl([ (\lambda-1)/2,f]_N\bigr)_N.
\]
For a positive integer $d$, write $T_d(Q)=\{q\in Q:q^d=1\}$.
For each $H\in\mathcal I$, the image of $T_d(GT_0)$ in $Q_H$ is exactly
\[
 \bigcap_{\substack{N\in\mathcal I\\N\subseteq H}}
 R_{N,H}\bigl(T_d(Q_N)\bigr).                         \tag{D.23a}
\]
An element of order exactly $d$ in this intersection has a lift of
order exactly $d$. Every element of order exactly $d$ in $GT_0$ has
an image of order exactly $d$ at some isolated level.
For $a,b\in GT_0$, conjugacy in $GT_0$ is equivalent to conjugacy
of $\Psi(a)_N$ and $\Psi(b)_N$ in every $Q_N$.
\end{claim}
\passage{\S{}31; folios 198--199 verso; scans 391--394; \S{}30; folio 194 verso; scan 384; Part D(b) of the task: full GT\_0 torsion}

\begin{proof}
The cited theorem identifies the two profinite groups, including their
topologies. Its Proposition 5.1 makes the coordinates cofinal in the
finite quotients of $\widehat\Z$ and $\widehat F_2$. The map between
the normalization in that paper and ours is exactly
$(\widehat m,\widehat f)\mapsto(2\widehat m+1,\widehat f)$;
Proposition 5.5 checks the equations and the multiplication under this
map. Thus its inverse limit concerns the full $GT_0$ of Part E.

Fix $q\in Q_H$ in the right side of (D.23a). For every isolated
$N\subseteq H$ define the finite nonempty set
\[
 S_N=\{v\in Q_N:v^d=1,\ R_{N,H}(v)=q\}.
\]
Every reduction $R_{N,M}$ restricts to a function $S_N\to S_M$.
The inverse limit of these sets is nonempty. To verify the exact
compactness step, any finite list of coordinates and compatibility
conditions is dominated by one isolated subgroup
$L\subseteq H\cap N_1\cap\cdots\cap N_t$; Proposition 3.14 supplies
an isolated refinement of the intersection. Choose an element of
$S_L$ and project it to the specified coordinates. This satisfies
all the finitely many conditions. The closed conditions in the
compact product $\prod_NS_N$ consequently have the finite
intersection property. A compatible element maps under
$\Psi^{-1}$ to $g\in GT_0$ with $g^d=1$ and $g_H=q$.
The reverse inclusion follows by taking the coordinates of such a
$g$. The levels below $H$ are cofinal, so the constructed family is
a full inverse-limit point.

When $q$ has order $d$, the order of $g$ both divides $d$ and is a
multiple of the order of its image $q$, hence equals $d$.
Conversely, for $g$ of order $d$, choose, for each prime $p\mid d$,
a finite coordinate detecting $g^{d/p}\ne1$. A common isolated
refinement detects all these powers. The order of its image divides
$d$ and divides none of the proper divisors $d/p$, so it is $d$.

For the conjugacy assertion, the sets
\[
 C_N(a,b)=\{v\in Q_N:v\Psi(a)_Nv^{-1}=\Psi(b)_N\}
\]
are finite nonempty sets with restriction maps. The same finite
intersection argument gives a compatible conjugator, which
$\Psi^{-1}$ sends to a conjugator in $GT_0$. Necessity follows by
projection. The construction does not assume surjectivity of a
reduction map or that every element of $Q_N$ comes from $GT_0$.
\end{proof}

\begin{claim}{D.24}{Proved}{Dihedral tests with an order bound and the full dyadic quotient}
The program \texttt{code/part\_a/gt\_shadow\_refinement.py} executes
the finite order-two and order-three tests below. Each entry counts
nonidentity elements with order dividing $d$ at the lower level,
and then those admitting a lift with order dividing $d$ at the
upper level.
\begin{center}
\begin{tabular}{rrrrr}
\toprule
$n$&$q$&$d$&Lower elements&Elements admitting such a lift\\
\midrule
6&18&2&7&7\\
6&18&3&2&0\\
18&54&3&8&0\\
54&162&3&8&0\\
14&98&3&14&14\\
4&8&2&3&1\\
8&16&2&7&2\\
16&32&2&11&4\\
32&64&2&19&8\\
\bottomrule
\end{tabular}
\end{center}
The levels are isolated dihedral shadows $Q_{K^{(n)}}$. Their
subsystem with $n=2^\alpha$, $\alpha\geq2$, has inverse limit
\[
 H_2=\{(k,(-1)^a5^b)\in
 \Z_2\rtimes\Z_2^\times:
 a\in\{0,1\},\ b\in\Z_2,\ k\equiv b\pmod2\}.
                                                               \tag{D.24a}
\]
Every finite subgroup of $H_2$ is trivial or cyclic of order two,
and all its order-two subgroups are conjugate to
$\langle(0,-1)\rangle$. There is a continuous surjection
$GT_0\to H_2$. The family of all dihedral levels is not cofinal
in $\mathcal I$; an explicit $A_5$ quotient below proves this.
\end{claim}
\passage{\S{}31; folios 198--199 verso; scans 391--394; \S{}30; folio 194 verso; scan 384; Part D(b) of the task: finite-quotient torsion program}

\begin{proof}
Write $D_n=\langle r,s\mid r^n=s^2=1,\ srs^{-1}=r^{-1}\rangle$.
The homomorphism
\[
 \psi_n:PB_3=\langle x,y\rangle\times\langle c\rangle
       \longrightarrow D_n^3,
 \quad x\longmapsto(r,s,s),\quad
 y\longmapsto(rs,r,rs),\quad c\longmapsto1
\]
defines $K^{(n)}=\ker\psi_n$. For even $n\geq4$, the explicit
finite model of \cite[Theorem 4.3 and Proposition 4.5]{A-BDHPShadows}
is
\[
 Q_{K^{(n)}}\simeq
 \left\{(k,u)\in(\Z/(n/2))\rtimes(\Z/(2n))^\times:
 \begin{array}{l}
 k\equiv\nu(u)\pmod2\text{ when }4\mid n
 \end{array}\right\},                                  \tag{D.24b}
\]
where $\nu(u)=(u-1)/4$ for $u\equiv1\pmod4$ and
$\nu(u)=(u+1)/4$ for $u\equiv3\pmod4$.
Multiplication is $(k,u)(l,v)=(k+ul,uv)$, with the indicated
moduli. For $n\mid q$, reduction is
$(k\bmod(q/2),u\bmod2q)\mapsto(k\bmod(n/2),u\bmod2n)$.

The code constructs the actual subgroup
$G_n=\langle(r,s,s),(rs,r,rs)\rangle\leq D_n^3$ for
$n=4,6,8,12,14$. It computes its commutator subgroup and exhausts
all pairs $(m,f)$ with $0\leq m<n$, $\gcd(2m+1,n)=1$ and
$f\in G_n'$. It tests the two cycle equations
\[
 \theta(f)f=1,\qquad
 \tau^2(y^mf)\tau(y^mf)y^mf=1,
\]
where $\theta(x)=y$, $\theta(y)=x$, $\tau(x)=y$ and
$\tau(y)=(xy)^{-1}$. It checks descent of these maps, and of
$x\mapsto x^{2m+1}$, $y\mapsto f^{-1}y^{2m+1}f$, on every
labelled edge of the finite Cayley graph, and checks that the
resulting maps are bijective. The accepted pairs are exactly
those in (D.24b), of respective orders $4,12,16,24,84$.
This checks the reduction to the quoted finite model on the
tested groups, with no reliance on floating point arithmetic.
Projection of $G_6,G_{12},G_{14}$ onto $D_3,D_3,D_7$
exhibits finite quotients in which the odd-order rotation and
order-two reflection do not commute; their nilpotence would
force these coprime Sylow subgroups to commute. These tests
therefore include nonnilpotent finite quotients of $\widehat F_2$.

For the refinement counts, the code enumerates every element of
(D.24b), computes its $d$th power using the affine multiplication,
reduces every upper root, and counts the exact fibres. For example,
at $n=6$ the group is
$\operatorname{Aff}(\F_3)\times C_2$, and its two order-three
elements are the nonzero translations. At $q=18$, an element of
order dividing three has multiplier $u\equiv1,4,7\pmod9$ and
translation divisible by $3$. Both its affine coordinates reduce
to the identity modulo $3$. Thus neither order-three element
at $n=6$ has an order-dividing-three lift, although the unbounded
reduction map is surjective. At $n=14$ the multipliers $2,4$
modulo $7$ have order three. Each lifts to its order-three
root modulo $49$, and each of the seven lifts of a translation
coordinate is allowed because $1+u+u^2=0$ there. This gives
seven order-three lifts of each of the fourteen lower elements,
as the exhaustive output records. The same program gives the
remaining rows and the full table in its accompanying CSV.

The compatible affine coordinate reductions identify the dyadic
inverse limit with (D.24a): the residue of $b$ modulo $2$ is the
condition in (D.24b), and the inverse limits of the two coordinate
rings are $\Z_2$ and $\Z_2^\times$. The surjection from $GT_0$
is \cite[Theorem 5.7 and Proposition 5.8]{A-BDHPShadows}; more
precisely, that theorem proves surjectivity already after
restriction along $G_{\Q}\to GT_0$.

The translation kernel $\Z_2$ is torsion-free, so a finite
subgroup of $H_2$ injects into the finite-order subgroup of
$\Z_2^\times$. This subgroup is $\{1,-1\}$: the parametrization
$\Z_2^\times=\{\pm1\}\times5^{\Z_2}$ has torsion-free second
factor, as follows from
$v_2(5^{2^j}-1)=j+2$ and passage to the inverse limit.
Every nonidentity finite-order element of $H_2$ is consequently
$(k,-1)$ with $k\in2\Z_2$. Put $t=k/2$ and choose
$v=5^{t\bmod2}$. Then $(t,v)\in H_2$, and direct multiplication
gives
\[
 (t,v)(0,-1)(t,v)^{-1}=(2t,-1)=(k,-1).
\]
This proves the full finite-subgroup and conjugacy assertions
for this infinite quotient, rather than extrapolating the
finite dyadic counts.

Finally every $G_n$ is metabelian: the commutator subgroup of
$D_n^3$ lies in its abelian rotation subgroup, and $G_n'$ is
a subgroup of that rotation subgroup. Define the epimorphism
$\eta:F_2\to A_5$ by
$x\mapsto(1\ 2\ 3\ 4\ 5)$ and $y\mapsto(1\ 2\ 3)$.
The code enumerates its image and derived subgroup; both have
order $60$, proving directly that this image is a nontrivial
perfect group. Proposition 5.1 of \cite{A-DGShadows} gives an
isolated $N$ with $N\cap F_2\subseteq\ker\eta$.
An inclusion $K^{(n)}\subseteq N$ would give the surjection
\[
 G_n=F_2/(K^{(n)}\cap F_2)
   \longrightarrow F_2/\ker\eta=A_5.
\]
Its image would be metabelian, contradicting the computed
nontrivial perfect derived subgroup. Hence no such inclusion
exists. This proves the asserted failure of cofinality by an
explicit quotient and the exact factoring map.
\end{proof}

These tests give explicit finite obstructions for D.12.
D.32 proves, by its full fixed-subgroup and finite-cover argument,
that $GT_0$ has no nonidentity element of odd order and that every
finite subgroup has order at most two, detected by the dyadic
multiplier. D.34 further proves that a nonidentity involution
has the global multiplier $-1$.
The fourteen order-three shadows at $K^{(14)}$ survive the
executed test through $K^{(98)}$. Combining D.32 with D.23
shows that each fails the order-three lifting condition at some
isolated refinement; a common refinement excludes all fourteen.
This existence conclusion does not identify that refinement or
alter the recorded finite counts. D.25 explains why the entire
dihedral subsystem alone cannot supply the exclusion.
The remaining full $GT_0$ question is conjugacy of its
involutions. D.35 gives the exact third-cusp equation and the
finite conjugator criterion, and D.44 supplies a coboundary
satisfying its first symmetry. D.42 also proves the global
multiplier restriction for every finite subgroup of the full
three-cusp loop-outer group.

\begin{claim}{D.25}{Proved}{All-prime dihedral compatibility and the multiplier-one obstruction}
The inverse limit $D$ over all the dihedral levels of D.24 has the
following explicit isomorphism of profinite groups:
\[
 D\xrightarrow{\ \sim\ }
 \left(\prod_{p\text{ odd}}\operatorname{Aff}(\Z_p)\right)
 \times H_2.                                                   \tag{D.25a}
\]
There is a canonical continuous homomorphism $GT_0\to D$.
The group $D$ contains an element of every prime order. In
particular, imposing order three throughout this entire
dihedral system still permits a multiplier whose $7$-adic
coordinate is a primitive cube root and whose other prime
coordinates are $1$.

Put $F=\widehat F_2=\langle x,y\rangle$ and define the continuous
epimorphism and its closed kernel by
\[
 \kappa:F\longrightarrow\prod_pF_2^{(p)},\qquad
 R=\ker\kappa.
\]
Write $\lambda:GT_0\to\widehat\Z^\times$ for the multiplier.
For every integer $d\geq2$, the elements $a\in\ker\lambda$
with $a^d=1$ correspond exactly to the following elements
$f\in R$. The continuous endomorphism
\[
 \alpha_f:F\to F,\qquad x\mapsto x,\quad y\mapsto f^{-1}yf
\]
must be an automorphism and satisfy
\[
 \theta(f)f=1,\qquad
 \tau^2(f)\tau(f)f=1,\qquad
 \prod_{j=0}^{d-1}\alpha_f^{j}(f)=1,                         \tag{D.25b}
\]
where the product is in increasing order of $j$,
$\theta(x)=y$, $\theta(y)=x$, $\tau(x)=y$, and
$\tau(y)=(xy)^{-1}$. Every such finite-order element has identity
image in $D$. The kernel $R$ itself maps onto the explicit
perfect group $A_5$ used in D.24.
The finite-order solutions of (D.25b) in the stated
automorphism domain are exactly $f=1$, by D.32's subsequent
full fixed-subgroup argument.  Thus the nontrivial prime-order
points of the dihedral inverse limit exhibited here do not
lift with that order to $GT_0$ when the prime is odd.
\end{claim}
\passage{\S{}25 bis; folios 157 verso--160; scans 312--317; \S{}30; folio 194 verso; scan 384; \S{}31; folios 198--199 verso; scans 391--394}

\begin{proof}
It suffices to use the even levels $n=2^\alpha n_0$ with
$\alpha\geq1$ and $n_0$ odd; every odd level agrees with its
double. Apply the Chinese remainder isomorphism separately to
the translation coordinate modulo $n/2$ and to the multiplier
coordinate modulo $2n$ in (D.24b). The odd factors give
$\operatorname{Aff}(\Z/n_0)$, and the powers of two give
the finite groups defining $H_2$. These coordinate maps commute
with every reduction $n\mid q$. Conversely, specified odd
coordinates and a specified point of $H_2$ reconstruct every
finite pair $(k,u)$ by the Chinese remainder theorem. They
satisfy the parity condition because that condition involves
only the power-of-two coordinates. The resulting inverse
maps are continuous, being continuous at every finite
coordinate. This proves (D.25a), including its topology.
Restriction of the full inverse-limit point in D.23 defines
the indicated map from $GT_0$; surjectivity onto $D$ is not
used or asserted.

Let $q$ be a prime, and choose a prime divisor $p$ of $2^q-1$.
The residue of $2$ in $\F_p^\times$ has order $q$: its order
divides $q$, and order one would give $p\mid2-1$. Hence
$q\mid p-1$, in particular $p$ is odd and $p\ne q$.
The element $2\bmod p$ lifts to a unique $q$th root
$\zeta\in\Z_p^\times$. Explicitly, after a root $a$ modulo
$p^j$ is chosen, substitute $a+p^jt$, $0\leq t<p$,
in $X^q-1$ modulo $p^{j+1}$. Its coefficient of $t$ is
$p^jqa^{q-1}$, and $qa^{q-1}$ is a unit modulo $p$;
there is exactly one permitted value of $t$.
The compatible limit has order exactly $q$ because its
residue already does. Put $(0,\zeta)$ in the $p$-factor
of (D.25a), and put the identity in all the other factors.
This gives the asserted element of order $q$ in $D$.

For $q=3$, take $p=7$. The executed program
\texttt{code/part\_a/gt\_shadow\_all\_prime.py} constructs ten
successive roots of $X^2+X+1$ starting with $2\bmod7$, and
uses the exact Chinese remainder formulas to construct
compatible shadow coordinates at sixteen levels through
$n=2744$. At $n=14$ its affine pair is $(0,9\bmod28)$;
at $n=98$ it is $(0,177\bmod196)$. Their orders are three
and the latter reduces to the former. The preceding
root-lifting proof constructs the infinite point, while
the program verifies the displayed finite coordinates.
Thus the equations tested over all dihedral levels admit
prime-dependent torsion multipliers. They do not force
the global alternatives $1$ and $-1$.

We next prove the statement about $R$. Each projection
$F\to F_2^{(p)}$ is the canonical maximal pro-$p$ quotient
map. Their product is onto: in any finite collection of
finite $p$-group quotients at distinct primes, its image
surjects to each factor. Its order is consequently divisible
by every factor order, hence by their product, since these
orders are pairwise coprime. The image is the full product.
These finite products form the finite quotients of
$\prod_pF_2^{(p)}$, so the compact image of $F$ is the whole
product. Also $R\subset F'$, since the abelianization
$F^{\mathrm{ab}}=\widehat\Z^2$ is detected by the
abelianizations of all these pro-$p$ quotients.

Let $a=(1,f)\in GT_0$ have finite order dividing $d$.
It induces a finite-order automorphism of each free pro-$p$
quotient with trivial action on its abelianization. The
Magnus leading-term argument in D.10, which also applies
at $p=2$, makes this automorphism the identity. Its normalized
formula then gives
$f_p\in C_{F_2^{(p)}}(y)\cap(F_2^{(p)})'=1$ by A.7 and
the primitive abelianization coordinate of $y$.
Thus $f\in R$.

The two first equations of (D.25b) are precisely the defining
cycle equations at multiplier one; the required invertibility
is stated explicitly. To obtain the last equation, put
$w_d=f\alpha_f(f)\cdots\alpha_f^{d-1}(f)$.
Induction on $d$ gives
\[
 \alpha_f^d(x)=x,\qquad
 \alpha_f^d(y)=w_d^{-1}yw_d.
\]
Therefore $\alpha_f^d=1$ is equivalent to
$w_d\in C_F(y)$. But $w_d\in F'$ and
$C_F(y)\cap F'=\overline{\langle y\rangle}\cap F'=1$,
again by A.7 and abelianization. This proves the necessity
of the last equation and its sufficiency. The normalized
identification of $GT_0$ with these automorphisms gives
the asserted bijection. The kernel of the multiplier on
$D$ is a subgroup of the additive translation group
$\prod_p\Z_p$, which is torsion-free. Consequently every
finite-order multiplier-one element maps to the identity
of $D$.

Finally let $\eta:F\to A_5$ be the epimorphism from D.24.
The quotient $A_5/\eta(R)$ is a finite quotient of
$F/R=\prod_pF_2^{(p)}$ and is therefore nilpotent.
It is also perfect, since the executed permutation
calculation in D.24 proves $A_5'=A_5$. A finite nilpotent
perfect group is trivial: its lower central series would
otherwise be constantly equal to the nontrivial group
and could not terminate. Hence $\eta(R)=A_5$.
This identifies a nontrivial part of the residual group in the
exact equation (D.25b). The full three-cusp argument in D.32
proves that this equation, together with its stated cycle and
invertibility requirements, has only the identity solution.
The surjection $R\to A_5$ and the compatible points of $D$
remain the exact quotient calculations proved above.
\end{proof}

\begin{claim}{D.26}{Refuted}{Unmarked recognition at index two and the peripheral obstruction}
The assertion that a torsion-free profinite group with an open
free profinite subgroup of rank three is free is false, already
at index two. We give the rank-three instance of
\cite[Proposition, pp.~350--352]{A-Haran1985} as an explicit
extension
\[
 1\longrightarrow H=\widehat F(l,a,b)
   \longrightarrow E\longrightarrow C_2\longrightarrow1.
                                                               \tag{D.26a}
\]
It has no section. Its outer involution preserves no loop
structure of type $(1,2)$ on $H$. Every preserved type-$(0,4)$
structure would have a double-transposition cusp permutation,
multiplier $-\epsilon$, and canonical character
$\chi(l)=1$, $\chi(a)=\chi(b)=-1$ sending all four cusp
generators to $-1$. Two explicit type-$(0,4)$ candidates fail,
with finite $S_3$ and $S_4$ certificates. D.29 supplies a
different complete type-$(0,4)$ marking, satisfying exactly
these necessary conditions, and thereby turns this extension
into a counterexample to the loop-preserving assertion D(d).
\end{claim}
\passage{\S{}28; folios 178--179; scans 352--354; Part D(d) of the task: full profinite recognition}

\begin{proof}
Choose $\epsilon,\kappa\in\widehat\Z$ by their prime
coordinates:
\[
 \epsilon_3=-1,\quad\epsilon_p=1\ (p\ne3),\qquad
 \kappa_2=1,\quad\kappa_p=0\ (p\ne2).
\]
Thus $\epsilon$ is a unit, $\epsilon^2=1$ and
$\epsilon\kappa=\kappa$. Define the continuous endomorphism
of the free profinite group $H$ by
\[
 \alpha(l)=l^\epsilon,\qquad
 \alpha(a)=b,\qquad
 \alpha(b)=l^\kappa a l^{-\kappa},\qquad c=l^\kappa.
                                                               \tag{D.26b}
\]
Substitution on the three free generators gives
$\alpha^2=\operatorname{Inn}(c)$ and $\alpha(c)=c$.
In particular $\operatorname{Inn}(c^{-1})\alpha$ is an
inverse, so $\alpha$ is an automorphism.

Define $E$ on the compact space $H\times\{0,1\}$ by
\[
 (h,i)(h',j)=
 \bigl(h\alpha^i(h')c^{\lfloor(i+j)/2\rfloor},
             i+j\bmod2\bigr).
                                                               \tag{D.26c}
\]
The two identities just checked prove associativity by
the eight triples of quotient labels. When the first two
labels are $1$, the required interchange is
$\alpha^2(h'')c=ch''$; the triple $(1,1,1)$ additionally
uses $\alpha(c)=c$. The other triples follow directly.
The identity and inverses follow by the same formulas.
All operations are continuous, so this is a profinite group.
Writing $t=(1,1)$ gives $tht^{-1}=\alpha(h)$ and $t^2=c$.
The projection to the second coordinate gives (D.26a).

The closed subgroup $N=\overline{\langle l,t\rangle}$
has the explicit isomorphism
\[
 N\simeq (\Z_3\rtimes\Z_2)
          \times\prod_{p\ne2,3}\Z_p,
\]
where the parity of the $\Z_2$ coordinate acts by inversion
on $\Z_3$. Namely, the normal form $l^u t^i$ maps to
$(u_3,2u_2+i,(u_p)_{p\ne2,3})$.
This is a continuous bijection preserving the displayed
relations and multiplication; its inverse reads the parity
$i$ and then all coordinates of $u$.
The projection of $N$ to
$\Z_2\times\prod_{p\ne2,3}\Z_p$ has torsion-free kernel
$\Z_3$ and torsion-free target, hence $N$ is torsion-free.

There is a continuous retraction $E\to N$ given on the
generators by $l\mapsto l$, $t\mapsto t$, $a,b\mapsto1$;
(D.26b) verifies its compatibility with the group law.
An element of finite order maps trivially under this
retraction and therefore has trivial image in $C_2$.
It lies in $H$, which is torsion-free by A.7, so it is
the identity. Thus $E$ is torsion-free and the extension
has no section. The induced outer action is faithful:
its nonidentity element acts nontrivially on $H^{\mathrm{ab}}$.

For comparison with Haran's construction, the homomorphism
$N\amalg\widehat\Z\to E$ from the free profinite product,
identity on $N$ and sending the other generator to $a$,
is onto since $b=tat^{-1}$. In the opposite direction send
$l,a,b,t$ to $l,a,tat^{-1},t$ in that free product.
Equations (D.26b) hold there, using the two relations in
$N$, so (D.26c) gives a continuous homomorphism in the
opposite direction. The composites fix every generator.
These are mutually inverse isomorphisms.

To prove that $E$ is not free, form the finite group
\[
 Q=(\F_3^2)\rtimes C_2,
 \qquad (v,i)(w,j)=(v+(-1)^iw,i+j).
\]
The formulas $l\mapsto(e_1,0)$, $a\mapsto(e_2,0)$,
$t\mapsto(0,1)$, and $b\mapsto(-e_2,0)$ give an
epimorphism $E\to Q$: $\epsilon\equiv-1\pmod3$ and
$\kappa\equiv0\pmod3$ verify all relations.
This group requires three generators. Two elements of
$\F_3^2$ do not generate its $C_2$ quotient. When one
of two proposed generators lies outside $\F_3^2$, it is
an involution. Multiplying the other by it when necessary
puts that other generator in $\F_3^2$. Conjugation then
only inverts this one vector, so the generated subgroup
has order at most six. Three generators, as displayed,
do generate $Q$.
Were $E$ a free profinite group of rank $r$, its open
index-two subgroup would have rank $2r-1$ by the finite
coset Schreier rewriting. Since $H$ has rank three,
this would give $r=2$, contradicting its quotient $Q$.
The program \texttt{code/part\_a/haran\_rank3.py}
checks all $18^3$ associativity instances and all $18^2$
ordered pairs in $Q$. Their generated subgroup orders
are $1,2,3,6,9$, with multiplicities $1,27,32,216,48$;
none is $18$.

We now check the peripheral condition rather than
discarding it. On the basis $[l],[a],[b]$ of
$H^{\mathrm{ab}}=\widehat\Z^3$, the outer involution has
matrix and determinant
\[
 M=\begin{pmatrix}\epsilon&0&0\\0&0&1\\0&1&0\end{pmatrix},
 \qquad\det M=-\epsilon.
                                                               \tag{D.26d}
\]
This determinant has value $+1$ at the prime $3$ and
$-1$ at every other prime, and is invariant under a
change of basis. For a preserved type-$(1,2)$ peripheral
structure, let $W$ be the primitive rank-one cusp
submodule in $H^{\mathrm{ab}}$. Write $\lambda$ for the
action on the orientation line and $\delta\in\{1,-1\}$
for the sign of the permutation of the two punctures.
The action on $W$ is multiplication by $\lambda\delta$.
The quotient $H^{\mathrm{ab}}/W$ is the first homology
of the filled torus. The orientation comparison of A.2
identifies its determinant with $\lambda$.
The square of the outer action is the identity, so
$\lambda^2=1$. Multiplicativity of determinant in the
exact sequence of these free modules gives
\[
 \det M=(\lambda\delta)\lambda=\delta.
\]
This contradicts (D.26d). Thus no such type-$(1,2)$
peripheral structure is preserved.

For type $(0,4)$ this determinant argument permits a
prime-dependent orientation multiplier. We test the
specific free basis $l,a,b^\epsilon$, which is a basis
because the power substitution on $b$ has inverse
$b\mapsto b^\epsilon$. Set
\[
 P=lab^\epsilon,
 \quad L_1=\overline{\langle l\rangle},\quad
 L_2=\overline{\langle a\rangle},\quad
 L_3=\overline{\langle b^\epsilon\rangle},\quad
 L_4=\overline{\langle P^{-1}\rangle}.
\]
This is a full type-$(0,4)$ peripheral system. The first
three classes transform as required with multiplier
$\epsilon$ and permutation $(2\ 3)$. The fourth would
require $\alpha(P)$ to be conjugate to $P^\epsilon$.
The actual words are
\[
 \alpha(P)=l^\epsilon b l^\kappa a^\epsilon l^{-\kappa},
 \qquad P^\epsilon=(lab^\epsilon)^\epsilon.
\]
Map the free group $H$ continuously onto $S_3$ by
\[
 l\longmapsto(1\ 2),\qquad
 a\longmapsto(2\ 3),\qquad
 b\longmapsto(1\ 2\ 3).
\]
At its exponent six, $\epsilon\equiv5$ and
$\kappa\equiv3$. Direct multiplication gives
$P\mapsto1$, $\alpha(P)\mapsto(1\ 2\ 3)$,
and $P^\epsilon\mapsto1$. These images are not conjugate,
so the fourth loop class fails. The program checks these
values and all $6^3$ possible assignments, finding
twenty-four nonconjugacy certificates.

There is a basis-independent restriction on any remaining
type-$(0,4)$ structure. Write $\sigma\in S_4$ for its cusp
permutation and $\lambda\in\widehat\Z^\times$ for its
multiplier. The square of the outer action is the identity,
so $\sigma^2=1$ and $\lambda^2=1$. Its action on
$H^{\mathrm{ab}}$ is $\lambda$ times the action of
$\sigma$ on $\widehat\Z^4/\widehat\Z(1,1,1,1)$.
Comparing traces with (D.26d) gives
\[
 \epsilon=\lambda\bigl(|\operatorname{Fix}(\sigma)|-1\bigr).
                                                               \tag{D.26e}
\]
The identity permutation would give $3\lambda=\epsilon$,
which is impossible modulo $3$. A single transposition
forces $\lambda=\epsilon$; a double transposition forces
$\lambda=-\epsilon$.

For each such peripheral structure there is a canonical
character $\chi:H\to\{1,-1\}$ which is $-1$ on all
four cusp generators. The product relation makes it
well-defined, and a free peripheral basis makes it
nonzero. It is preserved by $\alpha$, since $\lambda$
is odd and the cusp classes are permuted. In the basis
$l,a,b$ the three nonzero invariant characters are exactly
\[
 \chi_1=(-1,1,1),\qquad
 \chi_2=(1,-1,-1),\qquad
 \chi_3=(-1,-1,-1).                               \tag{D.26f}
\]
We compute the action on continuous cohomology with
coefficients $\Q_p(\chi)$ for an odd prime $p$; equivalently
one may use the inverse limit of the finite coefficient
modules and then tensor with $\Q_p$. A crossed
homomorphism $u$ is determined freely by
$(L,A,B)=(u(l),u(a),u(b))$. Put
$\eta=\chi(l)$ and $\zeta=\chi(a)=\chi(b)$, and define
\[
 e=\begin{cases}\epsilon_p&\eta=1,\\1&\eta=-1,\end{cases}
 \qquad
 k=\begin{cases}0&\eta=1,\\1&\eta=-1.\end{cases}
\]
Here $u(l^\epsilon)=eL$ and $u(l^\kappa)=kL$.
For $\eta=1$ this is multiplication by the $p$-coordinate
of the exponent; for $\eta=-1$ the geometric sum is
$(1-(-1)^v)/2=1$ for either odd exponent
$v=\epsilon,\kappa$. Substitution in (D.26b) gives
\[
 \alpha^*u\quad\longleftrightarrow\quad
 A_{\chi,p}\begin{pmatrix}L\\A\\B\end{pmatrix},
 \qquad
 A_{\chi,p}=
 \begin{pmatrix}
 e&0&0\\0&0&1\\(1-\zeta)k&\eta&0
 \end{pmatrix}.                                 \tag{D.26g}
\]
The coboundary line is spanned by
$(\eta-1,\zeta-1,\zeta-1)$, and this nonzero vector
is fixed by the displayed matrix. Thus
\[
 \det\bigl(\alpha^*|H^1(H,\Q_p(\chi_i))\bigr)
 =\begin{cases}
 1&i=1,3,\\-\epsilon_p&i=2.
 \end{cases}                                    \tag{D.26h}
\]

We identify the determinant required by the peripheral
structure through actual restriction and filling maps.
Let $K=\ker\chi$. Under the specified marking it is the
completion of the double cover of the four-punctured
sphere in which every cusp has nontrivial monodromy.
The cover has four punctures and genus one, since its
Euler characteristic is $-4$. The deck involution fixes
each of the four peripheral homology classes. Their
span is the three-dimensional invariant subspace of
$H_1(K,\Q_p)$: restriction and averaging identify the
invariant cohomology with $H^1(H,\Q_p)$, of dimension
three. Filling the four punctures therefore identifies
the anti-invariant part with the first homology of the
filled torus. The restriction of crossed homomorphisms
induces the precise comparison
\[
 H^1(H,\Q_p(\chi))
   \xrightarrow{\ \sim\ }H^1(K,\Q_p)^-
   \xleftarrow{\ \sim\ }
          H^1\bigl(\widehat{\pi_1(T^2)},\Q_p\bigr).
                                                               \tag{D.26i}
\]
The first isomorphism follows by averaging over the
order-two deck group, which is invertible in $\Q_p$;
the second is the dual of the filling quotient on
the anti-invariant part. Both commute with the induced
action of $\alpha$. Upstairs cusp generators are squares
of the downstairs generators. A conjugator in $H$ can
be moved into $K$ by multiplying it by the corresponding
downstairs cusp generator when necessary; that generator
centralizes its own square. Thus $\alpha$ has the same
multiplier $\lambda$ on this covering pair. The transfer
of the relative orientation class, followed by filling,
is the comparison of A.2; it identifies the determinant
on the filled torus with $\lambda_p$. Consequently the
determinant in (D.26h) must be $\lambda_p$.

At $p=3$ all three determinants in (D.26h) are $1$.
This excludes the single-transposition case, which by
(D.26e) would have $\lambda_3=\epsilon_3=-1$.
For the remaining double transposition,
$\lambda=-\epsilon$; at $p=5$ its value is $-1$.
Formula (D.26h) then excludes $\chi_1$ and $\chi_3$,
leaving exactly $\chi_2$. This proves the necessary
conditions stated at the beginning for every marking.

One can satisfy the first two cusp conditions in this
remaining case by the explicit free basis
\[
 x=a,\qquad y=b^{-\epsilon},\qquad z=b^\epsilon l,
 \qquad w=(xyz)^{-1}=(al)^{-1}.                    \tag{D.26j}
\]
The inverse substitution is $a=x$, $b=y^{-\epsilon}$,
$l=yz$, proving that this is a free basis.
It has canonical character $\chi_2$, and
$\alpha(x)$ is $y^{-\epsilon}$ while $\alpha(y)$
is conjugate to $x^{-\epsilon}$. For permutation
$(x\ y)(z\ w)$ and multiplier $-\epsilon$, the next
condition would require
\[
 \alpha(z)\sim w^{-\epsilon},\quad\hbox{equivalently}\quad
                 a^\epsilon l^\epsilon\sim(al)^\epsilon.
                                                               \tag{D.26k}
\]
Indeed $\alpha(z)=c a^\epsilon c^{-1}l^\epsilon$
is conjugate to $a^\epsilon l^\epsilon$ because
$c=l^\kappa$ commutes with $l$. In $S_4$ take
$a=(2\ 3\ 4)$, $l=(1\ 2\ 3\ 4)$ and $b=1$.
Its exponent is twelve, with $\epsilon\equiv5$.
The left member of (D.26k) is the transposition $(1\ 4)$,
whereas the right member is the four-cycle
$(1\ 3\ 2\ 4)$. They cannot be conjugate.
The supplied program checks all $24^2$ pairs in $S_4$,
finding $96$ such nonconjugacy witnesses. It also checks
all six matrices (D.26g) over $\F_3$ and $\F_5$, their
fixed coboundary lines and the determinants on their
two-dimensional cohomology quotients.

These calculations reduce the permissible type-$(0,4)$
markings to double transpositions with the specified
multiplier and character. The second candidate (D.26j)
is excluded, but D.29 constructs another marking by
gluing a genuine three-cusp automorphism to its inverse.
The exact basis change $b'=\tau(a)$ in that construction
identifies its extension with (D.26a). Thus an invariant
marking does exist; the finite failures above concern
the two displayed candidates and the proved exclusions
of the other permutation and character types.
\end{proof}

\begin{claim}{D.27}{Proved}{The full-profinite rank divisibility and prime-order lifting test}
For a torsion-free profinite extension
\[
 1\longrightarrow\widehat F_r\longrightarrow E
   \longrightarrow G\longrightarrow1,
 \qquad 2\le r<\infty,\quad |G|<\infty,
\]
the integer $|G|$ divides $r-1$. In particular a torsion-free
finite extension of $\widehat F_2$ equals its prescribed kernel.
Every homomorphism $C_p\to\Out(\widehat F_2)$, for any prime $p$,
has a homomorphic lift to $\Aut(\widehat F_2)$; a faithful one has
a lift of order exactly $p$. This last assertion holds in the full
profinite category, without an invariant discretization and without
first passing to a pro-$p$ quotient.
\end{claim}
\passage{\S{}28; folios 178--179; scans 352--354; \S{}31; folios 198--199 verso; scans 391--394}

\begin{proof}
We give the normal-subgroup version of Jarden's index argument
\cite[Theorems 1.2 and 4.1, pp.~496--498]{R-Jarden1982}, with the
quotient comparisons written out. Put $F=\widehat F_r$. For every
prime $q$, $\operatorname{cd}_q(F)=1$. The torsion-free finite-index
dimension theorem of Serre \cite[\S1]{R-SerreCD}, whose torsion
hypothesis applies to $E$ itself, gives $\operatorname{cd}_q(E)=1$.
The same is true of every open subgroup of $E$. Thus these groups
are projective profinite groups. Here we use the precise
projectivity criterion recalled in \cite[\S2]{R-Jarden1982}:
a profinite group is projective exactly when its cohomological
dimension at every prime is at most one; equivalently, all its
finite embedding problems have weak solutions. This is the
criterion for the full profinite category, including finite
embedding problems whose kernels are nonabelian.

Fix a prime divisor $p$ of $|G|$, and let $p^a$ be its exact
$p$-part. Choose a Sylow $p$-subgroup $G_p\le G$ and write
$E_p$ for its inverse image in $E$. Define $O^p(J)$, for a
profinite group $J$, as the kernel of its maximal pro-$p$ quotient.
The characteristic subgroup $O^p(F)$ of $F$ is normal in $E_p$.
There are exact continuous quotient maps
\[
 1\longrightarrow F/O^p(F)\longrightarrow E_p/O^p(F)
 \longrightarrow G_p\longrightarrow1.
\]
The middle group is pro-$p$, since both its kernel and quotient
are pro-$p$. Therefore $O^p(E_p)\le O^p(F)$. Conversely the
restriction to $F$ of the maximal pro-$p$ quotient of $E_p$ is a
pro-$p$ quotient of $F$, so $O^p(F)\le O^p(E_p)$. This proves
\[
 O^p(E_p)=O^p(F),\qquad
 [E_p^{(p)}:F^{(p)}]=p^a.                         \tag{D.27a}
\]

The quotient $E_p^{(p)}$ is projective as a pro-$p$ group. Indeed a
finite embedding problem with $p$-group target and covering group
pulls back to one for $E_p$; its projective solution factors through
$E_p^{(p)}$ because its target is a $p$-group. A projective pro-$p$
group is free pro-$p$. It has finite rank $d_p$ because $E_p$ is
finitely generated. The subgroup $F^{(p)}$ in (D.27a) is free
pro-$p$ of rank $r$, with its full induced topology, as the exact
quotient construction proves. Finite-coset Schreier rewriting
therefore gives
\[
                         r-1=p^a(d_p-1).          \tag{D.27b}
\]
Consequently each exact prime-power divisor $p^a$ of $|G|$ divides
$r-1$. These divisors are pairwise coprime, proving the first
assertion. At $r=2$, this gives $|G|=1$.

Now let $\rho:C_p\to\Out(\widehat F_2)$ be a prescribed continuous
outer action, with the finite group given its discrete topology.
Form the pullback
\[
 E_\rho=\{(a,g)\in\Aut(\widehat F_2)\times C_p:
                         [a]=\rho(g)\}.
\]
This is a profinite group. Centre-freeness identifies its kernel
over $C_p$ with $\Inn(\widehat F_2)=\widehat F_2$ by the injective
map $x\mapsto(\Inn(x),1)$. If $E_\rho$ were torsion-free, the
already proved rank-two assertion would force its quotient to be
trivial, contradicting $|C_p|=p$. Hence $E_\rho$ contains a
nonidentity torsion element $v$. Its finite cyclic subgroup
intersects the torsion-free kernel trivially, so projection embeds
it in $C_p$. Its order is therefore exactly $p$, and it projects
isomorphically onto $C_p$. The inverse of that projection is a
section with the prescribed labels. Composing with the first
projection $E_\rho\to\Aut(\widehat F_2)$ gives the asserted lift.
For faithful $\rho$, the lift is faithful because its outer
projection is. No choice of conjugacy class of this section is
claimed; the argument proves existence.

Applied to a prime-order element of $GT_0$, this gives a genuine
finite-order loop-preserving automorphism representing its outer
class. It does not impose an unproved restriction on its multiplier,
or identify its fixed subgroup. In particular it neither solves
the residual norm equation D.25 nor proves conjugacy to the
geometric involution. It does prove that the nontrivial
torsion-free recognition test of D.22 first reaches free rank
three, where (D.27b) permits index two and D.26 supplies an
explicit unmarked example.
\end{proof}

\begin{claim}{D.28}{Proved}{Prime-order sections and the full fixed subgroup}
Let $p$ be a prime and let $\rho:C_p\to GT_0$ be a continuous
homomorphism whose multiplier has $p$-adic component $1$.
This condition is automatic for odd $p$. Let $\alpha$ be the
normalized automorphism of $F=\widehat F(x,y)$ associated to
the chosen generator of $C_p$, and put
\[
 E=F\rtimes_\alpha C_p,\qquad S=\{1\}\times C_p,
 \qquad C=C_F(S).
\]
There is exactly one $F$-conjugacy class of sections of
$E\to C_p$. Each of the three peripheral normalizers contains
a section. The inclusion $C\hookrightarrow F$ induces an
isomorphism of maximal pro-$p$ quotients
\[
                         C^{(p)}\xrightarrow{\sim}F^{(p)}.
                                                               \tag{D.28a}
\]
This assertion concerns the maximal pro-$p$ quotient of the
actual closed fixed subgroup. It does not identify that
subgroup with $F$. An explicit order-three automorphism below
has a proper fixed subgroup and exactly the isomorphism
(D.28a), while its third peripheral class fails in a finite
dihedral quotient.
\end{claim}
\passage{\S{}31; folios 198--199verso; scans 391--394; \S{}30; folio 194verso; scan 384; Part D(a) of the task: actual full GT\_0 torsion}

\begin{proof}
The normalized map $GT_0\to\Aut(F)$ is a homomorphism:
for the pair $(\lambda,f)$ it sends
$x\mapsto x^\lambda$, $y\mapsto f^{-1}y^\lambda f$, and
composition gives precisely the product $f\alpha(g)$ in
the second coordinate. Consequently $\alpha^p=1$.
On the maximal pro-$p$ quotient its multiplier is $1$;
the Magnus argument of D.10 and D.25 makes this finite-order
normalized automorphism the identity. Thus the quotient of
the semidirect product is the explicit map
\[
 \pi:E\longrightarrow F^{(p)}\times C_p,
                  \qquad(h,s^j)\longmapsto(\bar h,s^j).
                                                               \tag{D.28b}
\]
For odd $p$, a finite-order unit in $\Z_p^\times$ of order
dividing $p$ is $1$: reduction modulo $p$ eliminates the
prime-to-$p$ part, and the logarithm on $1+p\Z_p$ eliminates
the remaining torsion. This explains the automatic condition.

Write $c_1=x$, $c_2=y$, $c_3=(xy)^{-1}$, and
$L_i=\overline{\langle c_i\rangle}$. Preservation of the
three peripheral classes gives an exact sequence
\[
 1\longrightarrow L_i\longrightarrow N_E(L_i)
                \longrightarrow C_p\longrightarrow1;
                                                               \tag{D.28c}
\]
its kernel is $L_i$ by A.7. Choose $a_i$ over the specified
generator of $C_p$. Then
\[
 a_i c_i a_i^{-1}=c_i^\lambda,\qquad
 a_i^p=c_i^{m_i},\qquad (\lambda-1)m_i=0.
                                                               \tag{D.28d}
\]
The first component of $\pi(a_i)$ centralizes $\bar c_i$,
since $\lambda_p=1$. The centralizer of this primitive
element in the free pro-$p$ group is
$\overline{\langle\bar c_i\rangle}$. Hence
$\pi(a_i)=(\bar c_i^{v_i},s)$ for $v_i\in\Z_p$, and
comparison of $p$th powers yields
\[
                            (m_i)_p=pv_i.
                                                               \tag{D.28e}
\]
Let $N_\lambda=1+\lambda+\cdots+\lambda^{p-1}$.
In the $q$-adic coordinate where $\lambda_q\ne1$,
the equation $\lambda_q^p=1$ and the integral-domain
property give $(N_\lambda)_q=0$ and $(m_i)_q=0$.
Where $\lambda_q=1$, the norm is $p$; division by $p$
is possible for $q\ne p$, and (D.28e) supplies it at
$q=p$. Define $u_i\in\widehat\Z$ coordinatewise by
$(u_i)_q=-(m_i)_q/p$ in the latter coordinates and zero
in the former. This gives
$N_\lambda u_i=-m_i$. The element
$b_i=a_i c_i^{u_i}$ has order $p$, by (D.28d), and
projects to the prescribed generator. Thus every
normalizer in (D.28c) contains a section $S_i$.

Let $e_p\in\widehat\Z$ be the idempotent with
$p$-coordinate $1$ and all other coordinates $0$.
Each $S_i$ centralizes $c_i^{e_p}$. Its image in
$F^{\rm ab}\otimes\F_p$ is the nonzero primitive vector
$[c_i]$. It follows that
\[
       \dim_{\F_p}H^1(C_F(S_i),\F_p)\ge1.
                                                               \tag{D.28f}
\]
This inference uses restriction of a linear functional
on $F^{\rm ab}/p$ that is nonzero on the displayed
element; it does not infer first homology from the mere
presence of a pro-$p$ subgroup.

We now use the finite-cohomology form of the profinite
Brown theorem, quoted as Theorem~2 and applied in
\cite[\S1.1, equations (1)--(6)]{A-KochloukovaZalesskii2007}.
There is no subgroup $C_p\times C_p$ in $E$, because
every finite subgroup injects into $E/F=C_p$.
The Lyndon--Hochschild--Serre spectral sequence and
$\operatorname{cd}_p(F)=1$ give, for every $n>1$,
\[
 \dim_{\F_p}H^n(E,\F_p)\le
             1+\dim_{\F_p}H^1(F,\F_p)=3.
\]
For each subgroup $T$ of order $p$ in $E$, projection
to $C_p$ shows that its normalizer equals its centralizer:
a normalizing element must induce the identity on the
generator of the quotient. Dividing by the unique element
of $T$ with the same quotient gives
\[
                 N_E(T)=C_F(T)\times T.
\]
The Brown restriction map is a dense embedding into the
product of the cohomology of these normalizers. Its
finite-dimensional source makes it an isomorphism and
makes the number of factors finite. Each closed subgroup
$C_F(T)$ has $p$-cohomological dimension at most one.
The direct-product cohomology calculation therefore gives
the exact finite sum and bound
\[
 \sum_{[T]}
    \left(1+\dim_{\F_p}H^1(C_F(T),\F_p)\right)
       =\dim_{\F_p}H^n(E,\F_p)\le3,
                                                               \tag{D.28g}
\]
where $[T]$ runs through $E$-conjugacy classes of order-$p$
subgroups. These are also their $F$-conjugacy classes:
for any such $T$, the decomposition $E=FT$ gives the
same conjugacy orbit under $E$ and $F$.

By (D.28f), two distinct cusp-section classes would
contribute at least four to (D.28g). Thus all three
$S_i$ are conjugate. Conjugating $S_1$ and $S_2$ to a
single representative places conjugates of
$x^{e_p}$ and $y^{e_p}$ in its fixed subgroup.
Their images in $F^{\rm ab}/p$ are still the independent
vectors $[x],[y]$. This gives two independent restricted
linear functionals, so that common factor contributes
at least three. Bound (D.28g) is exhausted. There are
no other classes, and its fixed subgroup has
$\dim H^1(-,\F_p)=2$. Every order-$p$ subgroup projects
isomorphically onto $C_p$ and determines one labelled
section; hence there is exactly one conjugacy class of
sections. In particular these conclusions apply to $S$.

The image of $C$ in $F^{(p)}$ is closed and surjects
on its Frattini quotient by the two independent cusp
vectors, so it is all of $F^{(p)}$. Since $C$ is a
closed subgroup of a free profinite group, it is
projective. Pulling a finite pro-$p$ embedding problem
back to $C$ shows that $C^{(p)}$ is projective pro-$p$,
and hence free pro-$p$. Its rank equals
$\dim H^1(C,\F_p)=2$. The resulting surjection from
a rank-two free pro-$p$ group onto $F^{(p)}$ is an
isomorphism: after identifying free bases it is a
surjective endomorphism of a finitely generated profinite
group, which is injective. This proves (D.28a) in the
claimed direction.

Here is an explicit limitation of that last comparison.
Choose $u\in\widehat\Z^\times$ with $u_q=1$ for
$q\ne7$, and $u_7$ the order-three root with
$u_7\equiv2\pmod7$. The Hensel construction is given
in D.25. The substitutions
\[
                    \beta(x)=x^u,\qquad\beta(y)=y^u
                                                               \tag{D.28h}
\]
have inverse given by exponent $u^{-1}$ and have order
exactly three. The fixed subgroup $C_\beta$ is proper,
since $\beta(x)\ne x$ in the cyclic quotient of order
seven. But it contains $x^{e_3}$ and $y^{e_3}$, so its
map to $F^{(3)}$ is onto. The same Brown bound, now
using this one section and these two independent vectors,
gives $\dim H^1(C_\beta,\F_3)=2$. The preceding
projectivity argument proves
$C_\beta^{(3)}\cong F^{(3)}$, despite $C_\beta\ne F$.

The precise missing cusp condition is detected by
\[
 D_7=\langle r,s\mid r^7=s^2=1, srs=r^{-1}\rangle,
                  \qquad x\mapsto s,\quad y\mapsto sr.
\]
Here $xy\mapsto r$, $u\equiv9\pmod{14}$, and the
two involutions are unchanged under exponent $u$.
Thus $x^u y^u\mapsto r$ whereas $(xy)^u\mapsto r^2$.
The conjugacy class of $r$ consists of $r,r^{-1}$,
so these are not conjugate. Consequently (D.28h) is
not an element of the three-cusp loop-outer group, and
in particular not a $GT_0$ torsion construction. The
program \texttt{code/part\_a/gt\_torsion\_sections.py}
executes this finite certificate and the coordinate
norm-solvability tests used above. D.32 uses the third cusp
in a finite-cover rank argument to strengthen (D.28a) to
$C=F$ under the three-cusp hypotheses, and thereby excludes
the nontrivial prime-order actions considered here.
\end{proof}

\begin{claim}{D.29}{Refuted}{A preserved four-puncture structure on the torsion-free Haran extension}
The full profinite recognition assertion in D(d) is false, including
its peripheral hypothesis.  There is an exact sequence of profinite
groups
\[
 1\longrightarrow H=\widehat\pi_{0,4}
 \longrightarrow E\longrightarrow C_2\longrightarrow1
                                                        \tag{D.29.1}
\]
in which $E$ is torsion-free and its faithful outer action preserves
a complete type-$(0,4)$ loop structure, while $E$ is the completion
of no finite-type surface group, orientable or nonorientable.
The extension is explicitly isomorphic to D.26's group.
Its outer involution has a mixed-prime orientation multiplier and
is conjugate to no subgroup of the geometric completion inside
$\Out_{\rm lac}(\widehat\pi_{0,4})$.
For this involution both the discretization and prediscretization
fixed sets of F.1 are empty.
\end{claim}
\passage{28; 178--179; 352--354; 30; 194--196; 383--387; 31; 198--199 verso; 391--394}

\begin{proof}
Define units and an exponent by their prime coordinates:
\[
 \epsilon_3=-1,\qquad\epsilon_q=1\ (q\ne3),\qquad
 \kappa_2=1,\quad\kappa_q=0\ (q\ne2),\qquad
 \lambda=-\epsilon.
                                                        \tag{D.29.2}
\]
Thus $\epsilon^2=\lambda^2=1$ and
$\epsilon\kappa=\kappa$.  In particular $\lambda_3=1$
and $\lambda_q=-1$ for every other prime.

We first obtain the required pants isomorphism, without any
finite-order requirement on it.  The compatible assignments
$\zeta_n\mapsto\zeta_n^{\lambda\bmod n}$ define an automorphism
of the cyclotomic extension $\Q^{\rm cyc}/\Q$.
An automorphism of an algebraic extension extends to its algebraic
closure: extend its embedding along algebraic generators, or take
a maximal extension of the embedding and adjoin a root of a missing
minimal polynomial to prove maximality gives all of
$\overline\Q$.  We therefore choose an element
$\sigma\in G_{\Q}$ with cyclotomic character $\lambda$.
The Galois action on the three-punctured projective line commutes
outerly with its rational $S_3$ action and preserves each cusp;
on its local inertia it has exponent $\lambda$.
Harbater--Schneps's normalized description
\cite[\S0.2, Main Theorem (a), Theorem A]{R-HS}
therefore supplies an element $(\lambda,f)\in GT_0$, with
$f\in\widehat F(x,y)'$, whose automorphism is
\[
 x\longmapsto x^\lambda,\qquad
 y\longmapsto f(x,y)^{-1}y^\lambda f(x,y).
                                                        \tag{D.29.3}
\]
Its third cusp $(xy)^{-1}$ goes to a conjugate of its
$\lambda$th power.  This is the only property of $f$ needed
below, besides invertibility of the displayed automorphism.
Every choice of the extension $\sigma$ has these properties,
and every one gives the construction and conclusions that follow.

Let $H=\widehat F(l,a,b)$ and define its closed free factors
\[
 P_1=\widehat F(l,a),\qquad P_2=\widehat F(l,b),\qquad
 L=\overline{\langle l\rangle}.
\]
Their shared-factor universal property is
$H=P_1\mathbin{*_{L}}P_2$: homomorphisms out of the two
free factors agreeing on $l$ are exactly arbitrary continuous
assignments of $l,a,b$.  Give $P_1$ the ordered pants basis
$(l^{-1},a)$ and $P_2$ the ordered pants basis $(l,b)$.
Formula (D.29.3), between these bases, defines a continuous
isomorphism $\tau:P_1\to P_2$ with
\[
 \begin{split}
 \tau(l)&=l^\epsilon,\\
 \tau(a)&=f(l,b)^{-1}b^\lambda f(l,b),\\
 \tau(a^{-1}l)&\sim_{P_2}(b^{-1}l^{-1})^\lambda.
 \end{split}                                           \tag{D.29.4}
\]
The symbol $\sim_{P_2}$ means conjugacy by an element of
$P_2$; in particular it specifies a conjugacy in $H$.
Because $\lambda^{-1}=\lambda$, the inverse isomorphism has
the same peripheral multiplier.

Put $c=l^\kappa$ and prescribe
\[
 \alpha|_{P_1}=\tau,\qquad
 \alpha|_{P_2}=\operatorname{Inn}(c)\tau^{-1}.
                                                        \tag{D.29.5}
\]
Both maps send the shared generator $l$ to $l^\epsilon$:
$\tau^{-1}(l)=l^\epsilon$ and $c$ commutes with $l$.
They consequently extend to a continuous endomorphism of
$H$.  Direct substitution on the two free factors gives
\[
                  \alpha(c)=c,\qquad
                  \alpha^2=\operatorname{Inn}(c).      \tag{D.29.6}
\]
For example, on $P_1$ the second iterate is
$\operatorname{Inn}(c)\tau^{-1}\tau$; on $P_2$ it is
$\operatorname{Inn}(\alpha(c))\tau\tau^{-1}$.
These identities also prove invertibility, with inverse
$\operatorname{Inn}(c^{-1})\alpha$.

Here is the full peripheral structure, with its relation:
\[
 \begin{split}
 (u_1,u_2,u_3,u_4)&=(a,\ a^{-1}l,\ b,\ b^{-1}l^{-1}),\\
                         u_1u_2u_3u_4&=1.
 \end{split}                                           \tag{D.29.7}
\]
The first three entries form a free profinite basis, since
$a=u_1$, $l=u_1u_2$, and $b=u_3$ are its continuous
inverse substitutions.  Thus (D.29.7) is precisely a
type-$(0,4)$ presentation, not just four procyclic subgroups
whose generators happen to satisfy a relation.
Equations (D.29.4)--(D.29.5) give
\[
 \alpha(u_i)\sim_H u_{s(i)}^\lambda,
                  \qquad s=(1\ 3)(2\ 4).              \tag{D.29.8}
\]
For the first two entries this is (D.29.4), and for the last
two it is the inverse comparison, followed by conjugation
by $c$.  Hence the action preserves every required loop
source, its conjugacy ambiguity, and the common orientation
multiplier.  It permutes distinct cusp classes nontrivially,
so its outer class is nonidentity.  Equation (D.29.6) gives
that class exact order two.

On the profinite space $H\times\{0,1\}$ define
\[
 (h,i)(k,j)=
   \bigl(h\alpha^i(k)c^{ij},\ i+j\bmod2\bigr).
                                                        \tag{D.29.9}
\]
Associativity follows by the two carry identities in
(D.29.6); checking the binary indices gives precisely
$\alpha^2(k)c=ck$ and $\alpha(c)=c$.
The element $t=(1,1)$ satisfies $t^2=c$ and
$tht^{-1}=\alpha(h)$.  The inclusion $h\mapsto(h,0)$
and parity projection give (D.29.1), with the profinite
product topology and continuous operations.

We make the comparison with D.26 explicit.  Set
$b'=\tau(a)$.  The pair $(l^\epsilon,b')$ is a free basis
of $P_2$, and replacing its first member by its
$\epsilon$th power gives the free basis $(l,b')$.
Thus $(l,a,b')$ is a free basis of $H$, and
\[
 \alpha(l)=l^\epsilon,\qquad
 \alpha(a)=b',\qquad
 \alpha(b')=c a c^{-1}.                               \tag{D.29.10}
\]
The map from D.26's named free basis to $(l,a,b')$, with
$t\mapsto t$, is a continuous isomorphism of its entire
twisted extension onto $E$.  It is an isomorphism of
extensions with identity map on $C_2$.  Formula (D.29.7)
therefore supplies the previously missing preserved
peripheral system on that exact group.  In particular
D.26's necessary double-transposition and double-cover
character conditions are attained.

For clarity, we give the torsion and finite-quotient
arguments also in the present basis.  The retraction
$r:H\to\widehat\Z$ sends $l$ to $1$ and $a,b$ to $0$.
Since both $\tau(a)$ and $\tau^{-1}(b)$ lie in the
normal closures of the respective nonshared generators,
$r\alpha=\epsilon r$.  It extends to a retraction of
$E$ onto the group $N=\widehat\Z\times\{0,1\}$ with law
(D.29.9), action $\epsilon$, and square parameter
$\kappa$.  An element $(z,1)$ in $N$ has square
$((1+\epsilon)z+\kappa,0)$, whose $2$-coordinate is
$2z_2+1$ and is never zero.  Since its open subgroup
$\widehat\Z$ is torsion-free, every torsion element
outside it would have order two; hence $N$ is torsion-free.
A torsion element of $E$ has trivial image in $N$, so
has even parity and lies in the torsion-free $H$.
It is therefore the identity.

Let $V=\F_3^2$, with basis $v_1,v_2$, and send
\[
         l\longmapsto v_1,\qquad
         a\longmapsto v_2,\qquad
         b\longmapsto -v_2.
                                                        \tag{D.29.11}
\]
Here $\lambda_3=1$ and $\epsilon_3=-1$; conjugators
vanish in the abelian target.  The induced action of
$\alpha$ is $-1$ on $V$, and $c$ has image zero because
$\kappa_3=0$.  Sending $t$ to the involution consequently
extends (D.29.11) to a continuous epimorphism
\[
              E\longrightarrow V\rtimes_{-1}C_2.      \tag{D.29.12}
\]
The target has order eighteen and requires three
generators.  Two positive elements do not generate the
quotient $C_2$; one positive and one negative element
generate at most one line of $V$; two negative elements
reduce to the preceding case by taking their product.
The three displayed images do generate it.
The independent program
\path{haran_quotient_independent_audit.py} checks all
$18^3$ associativity cases, all $18^2$ ordered pairs,
all ordered triples, and the actual matrix identity in
(D.29.11).  The ordered-pair subgroup-size histogram is
\[
 \begin{array}{c|rrrrr}
 \text{generated subgroup order}&1&2&3&6&9\\\hline
 \text{ordered pairs}&1&27&32&216&48.
 \end{array}                                          \tag{D.29.13}
\]
These calculations agree with the independently executed
D.26 certificate under (D.29.10).

A free profinite $E$ of finite rank $d$ would have an
index-two free subgroup of rank $1+2(d-1)$.  The actual
rank of $H$ is three, forcing $d=2$, contrary to
(D.29.12).  Thus $E$ is not free.  This excludes all
punctured finite-type surface completions.
For a putative isomorphism from a closed surface
completion to $E$, the inverse image of $H$ is an open
index-two subgroup.  Its intersection with the dense
surface group has index two, and completion of this
intersection is exactly that open subgroup by the
finite-index topology argument of D.1.  It is the
completion of the fundamental group of the corresponding
closed double cover.  Every aspherical closed surface
completion has cohomological $p$-dimension two, whereas
$\widehat F_3$ has dimension one.  The spherical and
projective-plane groups and their completions are finite
and cannot give $H$.  This excludes every closed case
as well.

Finally A.2 supplies the continuous orientation
homomorphism
\[
 \operatorname{or}:\Out_{\rm lac}(H)
                       \longrightarrow\widehat\Z^\times.
\]
It is unchanged by conjugation.  Every discrete geometric
class has value $1$ or the global unit $-1$, and the
closure of the geometric group still has values in that
closed two-element subgroup.  The value of $[\alpha]$
is the mixed unit $\lambda$ in (D.29.2), which belongs
to neither value.  No conjugate of $\langle[\alpha]\rangle$
is therefore contained in the geometric completion.
The transporter formulas of D.16 and F.1 now give both
empty fixed sets claimed above.

The cusp normalizers themselves cause no exception to
this argument.  Since $s$ fixes none of the four cusp
classes, their normalizers in $E$ map trivially to $C_2$;
each equals its normalizer in $H$, namely its own
procyclic loop group by A.7.  The failure is present
even with these procyclic normalizers retained.
\end{proof}

\begin{claim}{D.30}{Refuted}{A closed genus-two counterexample to full profinite recognition}
The assertion that every torsion-free finite extension of a full
profinite closed hyperbolic surface group is a completed surface group
is false.  There is an exact extension
\[
 1\longrightarrow H=\widehat\pi_1(\Sigma_2)
 \longrightarrow E\longrightarrow C_2\longrightarrow1
                                                        \tag{D.30.1}
\]
with $E$ torsion-free and with a faithful outer action, such that
\[
 H^2(E;\mathbb F_3)=\mathbb F_3,\qquad
 H^2(E;\mathbb F_5)=0.                                  \tag{D.30.2}
\]
No profinite completion of a discrete surface group has these two
cohomology groups.  The exact maximal pro-$p$ quotients are
\[
 E^{(2)}\cong\widehat{\pi_1(N_3)}^{(2)},\qquad
 E^{(3)}\cong\mathbb Z_3^2,\qquad
 E^{(p)}\cong\widehat F_2^{(p)}\quad(p\ne2,3),           \tag{D.30.3}
\]
where $N_3$ is the closed nonorientable surface of genus three.
Thus the counterexample also gives the precise maps through which
its separate pro-$p$ surface presentations arise.
Its outer involution is conjugate to no element of the extended
geometric completion. Its fixed discretization and
prediscretization sets are both empty.
\end{claim}
\passage{28; 178 verso; 353; 24; 147 verso; 293; Part D(d) of the task}

\begin{proof}
We use the commutator convention $[x,y]=xyx^{-1}y^{-1}$.
Set
\[
 \epsilon_3=-1,\quad\epsilon_p=1\ (p\ne3),\qquad
 \kappa_2=1,\quad\kappa_p=0\ (p\ne2),\qquad
 \lambda=-\epsilon
                                                        \tag{D.30.4}
\]
in $\widehat{\mathbb Z}=\prod_p\mathbb Z_p$.  These elements
satisfy $\epsilon^2=\lambda^2=1$ and
$\epsilon\kappa=\kappa$.

\emph{An actual punctured-torus isomorphism.}
Take the elliptic curve $Y^2=4X^3-4X$ over $\mathbb Q$,
with its origin removed and its rational Weierstrass tangential
base point.  Nakamura's full profinite construction
\cite[\S1, equation (1.2), and \S2.5, p.~424]{C_Nakamura2013}
gives an actual homomorphism
\[
 G_{\mathbb Q}\longrightarrow\operatorname{Aut}
 \langle x,y,z\mid[x,y]z=1\rangle^{\widehat{\ }},
 \qquad \sigma(z)=z^{\chi(\sigma)}.
                                                        \tag{D.30.5}
\]
Here $\chi$ is the full cyclotomic character.  The local
coordinate at the rational puncture and its Kummer covers identify
its inertia with $\widehat{\mathbb Z}(1)$, which explains the exact
power equation.  Every finite cyclotomic extension has Galois group
$(\mathbb Z/n)^\times$; the resulting surjections and compactness
show that $\chi:G_{\mathbb Q}\to\widehat{\mathbb Z}^{\times}$
is onto.  Choose $\sigma$ with $\chi(\sigma)=\lambda$.
Renaming the target generators gives a continuous isomorphism
\[
 \tau:F_1=\widehat F(x,y)\xrightarrow{\sim}
       F_2=\widehat F(u,v),
 \qquad \tau([x,y])=[u,v]^\lambda.                     \tag{D.30.6}
\]
This uses a single Galois element of prescribed multiplier.
The construction below proves its outer involution directly from
$\tau$ and $\tau^{-1}$.

\emph{The automorphism of the closed surface group.}
Use its defining profinite presentation
\[
 H=\langle x,y,u,v\mid[x,y][u,v]=1\rangle^{\widehat{\ }},
 \qquad l=[x,y]^{-1}=[u,v],\qquad c=l^\kappa.           \tag{D.30.7}
\]
Define $\alpha$ on the first pair by the two words of $\tau$,
and on the second pair by the two words of
$\operatorname{Inn}(c)\tau^{-1}$.
Equation (D.30.6) and $\lambda^{-1}=\lambda$ give
\[
 \alpha([x,y])=l^\lambda,\qquad
 \alpha([u,v])=c\,[x,y]^\lambda c^{-1}=l^{-\lambda}.
\]
Their product is one, so this substitution defines a continuous
endomorphism of the presented group $H$.  It gives
$\alpha(l)=l^\epsilon$ and $\alpha(c)=c$.
Substitution on each of the four generators, using the actual inverse
$\tau^{-1}$, then gives
\[
                     \alpha^2=\operatorname{Inn}(c).
                                                        \tag{D.30.8}
\]
Consequently $\operatorname{Inn}(c^{-1})\alpha$ is its inverse.
These identities use the universal property of the displayed
presentation, so no assertion about efficiency of a profinite
amalgam is required.

Define $E$ on the compact space $H\times\{0,1\}$ by
\[
 (h,i)(k,j)=
 \bigl(h\alpha^i(k)c^{\lfloor(i+j)/2\rfloor},
                           i+j\bmod2\bigr).            \tag{D.30.9}
\]
For quotient labels $(1,1,0)$ associativity is precisely
$\alpha^2(k)c=ck$; for $(1,1,1)$ it also uses
$\alpha(c)=c$.  The other six label triples follow by direct
substitution.  The identity is $(1,0)$; the inverse of $(h,0)$ is
$(h^{-1},0)$, and that of $(h,1)$ is
$(\alpha^{-1}(h^{-1}c^{-1}),1)$.
All operations are continuous.  Thus $E$ is a profinite group
and its projection gives (D.30.1).  Write $t=(1,1)$, so that
$tht^{-1}=\alpha(h)$ and $t^2=c$.

\emph{An exact presentation and torsion-freeness.}
Set $u'=\tau(x)$ and $v'=\tau(y)$.  This is a free basis
of the formal second punctured-torus group, with
\[
 [u',v']=[u,v]^\lambda=l^\lambda,\qquad
       u'=txt^{-1},\quad v'=tyt^{-1}.
\]
The relations for $t$ on $u',v'$ follow from $t^2=c$.
Eliminating these two generators therefore gives the exact
profinite presentation
\[
 E\cong\langle x,y,t\mid
          tlt^{-1}=l^\epsilon,\quad t^2=l^\kappa
                 \rangle^{\widehat{\ }},
                 \qquad l=[x,y]^{-1}.                 \tag{D.30.10}
\]
For completeness, the homomorphism to (D.30.9) sends $x,y,t$
to the named elements and is onto because their conjugates
$u',v'$ generate the second factor.  In the reverse direction,
recover $u,v$ by applying the words of $\tau^{-1}$ to
$txt^{-1},tyt^{-1}$.  Their commutator is
$[txt^{-1},tyt^{-1}]^\lambda=l$, by (D.30.10).
Thus the surface relator holds, and the conjugation relations on
both pairs follow from $t^2=c$ and the defining words of
$\alpha$.  The two maps fix all named generators and are inverse.

The universal property of the maximal pro-$2$ quotient substitutes
$\epsilon_2=\kappa_2=1$ into (D.30.10).  Its first relation follows
from the second, giving
\[
 E^{(2)}\cong
       \langle x,y,t\mid[x,y]t^2=1\rangle_{\mathrm{pro}\text{-}2}.
                                                        \tag{D.30.11}
\]
This is the rank-three Demushkin presentation with $q=2$ and
$f=\infty$: cyclically move $t^2$ to the beginning and replace
$x,y$ by their inverses to match the commutator convention of
\cite[\S1.4 and Theorem~1.9]{R-Wilkes}.  The converse in that
theorem proves the Demushkin assertion.  Its abelianization is
$\mathbb Z_2^2\times C_2$, so it is infinite.  The same source's
preceding cohomological-dimension statement gives
$\operatorname{cd}_2(E^{(2)})=2$.  A subgroup of order two has
infinite $2$-cohomological dimension, so $E^{(2)}$ is torsion-free.

There is also an explicit free-basis identification with the
three-crosscap word.  The substitution
\[
 x\mapsto a d^{-1}b^{-1}a^{-1},\qquad
 y\mapsto a b d a b a^{-1},\qquad t\mapsto abd
                                                        \tag{D.30.12}
\]
sends $[x,y]t^2$ to $a^2b^2d^2$.  Its inverse is
\[
 a\mapsto xt,\qquad
 b\mapsto t^{-1}x^{-1}t^{-1}yxt,\qquad
 d\mapsto t^{-1}x^{-1}y^{-1}t^2.
\]
Free reduction checks the relator and both composites.  The maps
extend to mutually inverse pro-$2$ isomorphisms, proving the first
comparison in (D.30.3).

The parity quotient $E\to C_2$ factors through $E^{(2)}$ and
sends $t$ to its nonidentity element and $x,y$ to the identity.
Every finite-order element of $E$ maps trivially to the torsion-free
$E^{(2)}$, hence has trivial parity and lies in $H$.
The latter is torsion-free by A.7.  This proves torsion-freeness of
$E$ itself and, in particular, the absence of a section of (D.30.1).

For an odd prime $p$, (D.30.10) gives $t^2=1$ in $E^{(p)}$,
so $t=1$.  At $p=3$ the remaining relation is $l=l^{-1}$,
which forces $l=1$, and the quotient is
$\langle x,y\mid[x,y]=1\rangle_{\mathrm{pro}\text{-}3}
=\mathbb Z_3^2$.  At every other odd prime the remaining relation
is the identity $l=l$.  The universal properties in both directions
give the free pro-$p$ group on $x,y$.  This proves all of (D.30.3).

\emph{The orientation line and the obstruction to every surface.}
Let $A\in\operatorname{GL}_2(\widehat{\mathbb Z})$ be the
abelianization matrix of $\tau$.  In the finite Heisenberg group
modulo $n$, the commutator of two elements is the central generator
raised to the determinant of their two abelianization columns.
Applying (D.30.6) there gives $\det A=\lambda\pmod n$ for every
$n$, hence $\det A=\lambda$ in $\widehat{\mathbb Z}$.
On the basis $([x],[y],[u],[v])$ of $H^{\mathrm{ab}}$, inner
conjugation disappears and the matrix of $\alpha$ is
\[
 M=\begin{pmatrix}0&A^{-1}\\A&0\end{pmatrix}.
\]
For $J=\left(\begin{smallmatrix}0&1\\-1&0\end{smallmatrix}\right)$,
the identity $A^tJA=(\det A)J$ gives
\[
 M^t\operatorname{diag}(J,J)M
       =\lambda\operatorname{diag}(J,J),                \tag{D.30.13}
\]
because $\lambda^{-1}=\lambda$.  The companion identity
$M\operatorname{diag}(J,J)M^t=\lambda\operatorname{diag}(J,J)$
also holds by $AJA^t=(\det A)J$.  Since the matrix of $\alpha^*$
on degree-one cohomology is $M^t$, this last identity gives the
multiplier of its cup pairing with all dual maps specified.
The completed surface cup-product comparison of A.2 is natural
for automorphisms.  Its nondegenerate degree-one pairing therefore
shows that $\alpha^*$ acts by $\lambda_p$ on the one-dimensional
$H^2(H;\mathbb F_p)$.  The matrix $M$ exchanges two nonzero direct
summands, so $\alpha$ is not inner.  Equation (D.30.8) thus gives
a faithful outer $C_2$ action.

The orientation-line character is a continuous homomorphism
\[
 \operatorname{Out}(H)\longrightarrow
                    \widehat{\mathbb Z}^{\times}.
\]
Its restriction to the extended geometric mapping class group
has image the two diagonal elements $+1,-1$, and continuity gives
the same bound for its closure. Conjugation in $\operatorname{Out}(H)$
preserves this character because its target is abelian. Our element
has value $\lambda$, with $\lambda_3=1$ and $\lambda_5=-1$, so no
conjugate lies in that closure. D.16's exact transporter formulas
therefore give empty fixed discretization and prediscretization
sets for this actual involution.

For odd $p$, averaging over $C_2$ makes all its higher cohomology
with $\mathbb F_p$-coefficients vanish.  The continuous
Hochschild--Serre spectral sequence of (D.30.1) consequently gives
\[
 H^2(E;\mathbb F_p)\xrightarrow[\mathrm{restriction}]{\sim}
                      H^2(H;\mathbb F_p)^{C_2}.
                                                        \tag{D.30.14}
\]
At $p=3$ the scalar is $\lambda_3=1$, while at $p=5$ it is
$\lambda_5=-1$.  This proves (D.30.2).

Finally, the completion of a closed orientable surface group of
positive genus has one-dimensional $H^2(-;\mathbb F_p)$ at every
prime, by the surface comparison of A.2 (including its torus
calculation).  A closed nonorientable surface of genus at least
two has an orientable double cover.  Exact completion at index
two and the same restriction map (D.30.14) compute its degree-two
cohomology at odd primes as the invariants of the orientation
reversal, hence zero.  The sphere group is trivial and the
projective-plane group is $C_2$, so both have zero degree-two
cohomology at odd primes.  A surface with boundary or a noncompact
surface has a free fundamental group, whose profinite completion
has cohomological dimension at most one.  These cases exhaust the
surface groups and exclude the pair of values in (D.30.2).

The independently executed program
\path{code/part_b/closed_extension_d30_audit.py}
checks the free-word isomorphism in (D.30.12), the intersection
matrices and the finite cochain computation of the two invariant
lines.  The profinite construction, presentation comparisons and
torsion-freeness are proved above; the finite computations verify
the stated identities.
\end{proof}

\begin{claim}{D.31}{Proved}{Abstract finite groups in the geometric completions of genus at most two}
Let $M_{g,n}^+=\Gamma^+_{g,n}$ allow permutations of the $n$
punctures, and let $M_{g,n}^{\pm}$ also allow orientation reversal.
Every finite nonsolvable subgroup of $\widehat{M_{0,n}^+}$ is
isomorphic to $A_5$. Its permutation action on the punctures has
\[
       n=12a+20b+30c+60m,\qquad
       a,b,c\in\{0,1\},\quad m\geq0,                 \tag{D.31.1}
\]
where the summands are the actual orbit sizes. In particular, it
has at least twelve punctures. Every finite nonsolvable subgroup
of $\widehat{M_{0,n}^{\pm}}$ is isomorphic to $A_5$ or
$A_5\times C_2$. Every finite subgroup of
$\widehat{M_{g,n}^{\pm}}$ for $g=1,2$ is solvable.

These statements determine the abstract groups and, in the positive
sphere case, their puncture-permutation representations. The
surjectivity of the map from geometric $A_5$ conjugacy classes to
completed $A_5$ conjugacy classes, and the completed $A_5$
centralizer comparison of F.7, remain to be proved.
\end{claim}
\passage{\S{}30; folios 194--196; scans 383--387; \S{}31; folios 198--199 verso; scans 391--394; Part D(b): finite subgroups of the geometric completion and full completion comparison}

\begin{proof}
For $n=0,1,2$ the positive sphere mapping class groups are
respectively $1,1,C_2$; including orientation reversal gives
groups of order at most four. Their completions are the same
finite solvable groups. The sphere arguments below therefore
take $n\geq3$.

We give the completed exact sequences first, so that every use of
a subgroup or quotient has its specified completion map.

\emph{A completion and configuration-group preparation.}
For the forgetting sequences used below, the kernel $N$ is finitely
generated and the quotient $Q$ is good. These properties imply
that $\widehat N\to\widehat E$ is injective. Here is the required
argument. Take a characteristic finite-index subgroup $N_0$ of $N$;
these are cofinal because $N$ is finitely generated. In $E/N_0$ the
kernel $B=N/N_0$ is finite. Its centralizer has finite index and fits
into the central extension
\[
 1\longrightarrow Z(B)\longrightarrow C_{E/N_0}(B)
       \longrightarrow Q_0\longrightarrow1,
\]
where $Q_0$ has finite index in $Q$. Goodness passes to $Q_0$.
The class in $H^2(Q_0,Z(B))$ is therefore represented by a continuous
cohomology class of $\widehat Q_0$, hence is inflated from a finite
quotient. Restriction to the kernel of that finite quotient kills
it. The central extension splits over this finite-index subgroup;
a chosen complement meets $B$ trivially and has finite index in
$E/N_0$. Its core is a finite-index normal subgroup still meeting
$B$ trivially. Pulling this core back to $E$ gives a finite quotient
whose restriction to $N$ has kernel exactly $N_0$. This proves the
full induced topology and thus
\[
 1\longrightarrow\widehat N\longrightarrow\widehat E
       \longrightarrow\widehat Q\longrightarrow1.     \tag{D.31.2}
\]
The last equality follows directly from the finite quotients of
$E/N$.

The braid groups on the plane, on finitely punctured surfaces,
and on the closed surfaces of genus one or two used here are
good and have finite cohomological dimension.
We recall how the ingredients apply here. The ordered
configuration-space fibrations of Fadell--Neuwirth
\cite[Theorems~1 and~3]{A-FadellNeuwirth1962} forget one point and
have fibre the surface with finitely many points removed.
The base surface and the fibres are aspherical; hence their
configuration spaces are aspherical by the homotopy exact sequence.
The fibres' groups are finitely generated free groups once a point
has been removed. Starting with the good surface group, the
preceding finite-kernel argument and the Lyndon--Hochschild--Serre
comparison spectral sequences prove goodness inductively.
The required cohomology of each fibre with finite coefficients is
finite. For the plane the initial group is trivial; for the closed
surfaces its goodness is the usual surface-group goodness also
used in \cite{R-BZ}. The unordered group has the ordered one as a
finite-index subgroup, so goodness follows by the finite-index
comparison and Shapiro's lemma. The configuration spaces are
finite-dimensional aspherical manifolds, which give finite
cohomological dimension. Goodness therefore makes each completed
braid group torsion-free: its cohomology with every finite
$p$-primary coefficient module vanishes above that dimension,
whereas a subgroup $C_p$ has infinite $p$-cohomological dimension.
This also follows directly from the prime-subgroup comparison in
\cite[Theorem~A]{R-BZ}.

\emph{Finite subgroups of a completed puncture stabilizer are cyclic.}
For $n\geq4$, capping the boundary of a disc by a once-punctured
disc gives the exact sequence
\[
 1\longrightarrow\langle\Delta^2\rangle\cong\Z
 \longrightarrow B_{n-1}\longrightarrow
 (M_{0,n}^+)_{p}\longrightarrow1.                      \tag{D.31.3}
\]
Here the distinguished puncture is $p$, and the kernel is the
central full twist; this is
\cite[Lemma~2.2]{A-BigelowBudney2001}, whose proof specifies the
capping map. Braid abelianization sends each Artin generator to
$1$ and $\Delta^2$ to $(n-1)(n-2)$. Thus the composite
\[
 \widehat\Z\longrightarrow\widehat B_{n-1}
       \longrightarrow\widehat\Z
\]
is multiplication by this nonzero integer and is injective.
Consequently completion of (D.31.3) is the exact central extension
\[
 1\longrightarrow\widehat\Z\longrightarrow
 \widehat B_{n-1}\longrightarrow
 \widehat{(M_{0,n}^+)_{p}}\longrightarrow1.             \tag{D.31.4}
\]
The last group identifies with the stabilizer of $p$ in
$\widehat{M_{0,n}^+}$: its discrete source has finite index, and
stabilizing $p$ is the inverse image of a subgroup in the finite
puncture-permutation quotient.

Let $K$ be a finite subgroup of that completed stabilizer and let
$E_K$ be its inverse image in $\widehat B_{n-1}$. The subgroup
$A=\widehat\Z$ is central and open in $E_K$, of index $d=|K|$.
Continuous transfer over its finitely many cosets gives
\[
             \operatorname{Ver}:E_K\longrightarrow A,
             \qquad \operatorname{Ver}(a)=d a
             \quad(a\in A).                           \tag{D.31.5}
\]
Its kernel intersects $A$ trivially, so injects into the finite
group $E_K/A$. It is finite, and is trivial because $E_K$ is a
subgroup of the torsion-free completed braid group. Transfer is
therefore a continuous injection into $\widehat\Z$. Its compact
image is a closed procyclic subgroup. Hence $E_K$, and then its
finite quotient $K$, are procyclic and cyclic respectively.
For $n=3$ the same capping map is
$1\to2\Z\to B_2=\Z\to C_2\to1$, and the conclusion is immediate.
For $n=1,2$ the positive point stabilizer is trivial.

It follows in particular that a finite subgroup of a completed
positive sphere mapping class group preserving a specified set
of $1\leq b\leq3$ punctures is solvable. The kernel of its permutation
map to $S_b$ fixes a puncture and is cyclic, while $S_b$ is
solvable. This assertion also applies to every finite-index
decorated subgroup through its actual inclusion into the completed
sphere group.

\emph{The prime-power inputs and elimination of a solvable radical.}
The good-group prime-power theorem
\cite[Theorem~A]{R-BZ} conjugates each finite $p$-subgroup of
$\widehat{M_{0,n}^+}$ to a geometric one and gives its completed
normalizer. The checked prime-power topology argument in F.7
identifies that normalizer with the full completion of its
discrete normalizer. Thus, for a nontrivial normal $p$-subgroup
$A$ of a finite $G\leq\widehat{M_{0,n}^+}$, conjugate $A$ into
$M_{0,n}^+$ and use the exact covering-normalizer sequence
\[
 1\longrightarrow A\longrightarrow N_{M_{0,n}^+}(A)
       \longrightarrow Q_A\longrightarrow1.           \tag{D.31.6}
\]
The group $Q_A$ is a finite-index decorated subgroup of the
mapping class group of the quotient sphere. Its specified branch
set has two points for a cyclic rotation group, and three for a
noncyclic rotation group. This branch set includes the images of
ramification points that were already punctures upstairs; the
ramification indices are retained in the covering decoration.
Completion of (D.31.6) and the normalizer comparison give an
actual injection
\[
                 G/A\lhook\joinrel\longrightarrow\widehat Q_A.
                                                               \tag{D.31.7}
\]
It preserves that set of two or three branch points. The puncture
stabilizer argument makes $G/A$ solvable, and then $G$ is solvable.

For a nontrivial finite solvable normal subgroup $R$ of $G$, the
last nonzero derived term of $R$ is abelian. Its nontrivial
$p$-primary component $A$ is characteristic in $R$, hence normal
in $G$. The preceding argument again makes $G$ solvable.
Consequently a nonsolvable finite $G$ has trivial solvable radical.
Only the prime-power comparison was used in this step.

The same prime-power theorem and the finite rotation-group list
give cyclic odd Sylow subgroups of $G$, and cyclic or dihedral
Sylow $2$-subgroups, including $C_2^2$ in the latter case.
There is also the precise normalizer constraint
\[
 \operatorname{im}\bigl(N_G(C_p)\longrightarrow\Aut(C_p)\bigr)
                      \subset\{1,-1\}.               \tag{D.31.8}
\]
Indeed normalize a geometric rotation. A representative of a
discrete normalizer either fixes its two poles or exchanges them.
The oriented local rotation numbers are opposite, so the induced
automorphism is respectively $1$ or $-1$. The normalizer comparison
preserves this image because $\Aut(C_p)$ is finite. The same
argument with the local orientation sign proves (D.31.8) in the
extended geometric completion as well. Equivalently one can apply
the prime-power normalizer theorem to the good index-two extended
mapping class group; the discrete representatives still induce
only these two exponents.

We record the elementary finite-group reduction needed below.
A finite group all of whose Sylow subgroups are cyclic is solvable.
For the least prime $p$ dividing its order, its cyclic Sylow group
$P$ has $N(P)/C(P)$ of order both prime to $p$ and dividing
$p^{a-1}(p-1)$, where $|P|=p^a$. Minimality of $p$ forces this
quotient to be trivial. Transfer to $P$ has restriction
$x\mapsto x^{[G:P]}$ and is onto. To check this formula without
a fusion assumption, any two elements of the abelian Sylow group
$P$ conjugate in $G$ are conjugate by $N(P)$: the two Sylow
subgroups containing the conjugate element are conjugate inside
its centralizer. Thus they are equal when $P\leq Z(N(P))$.
The coset-cycle formula for transfer now gives the stated exponent.
Its kernel is a normal $p$-complement. Induction on the group order
proves solvability. It follows that every odd-order subgroup here
is solvable, and that a nonsolvable $G$ has a noncyclic dihedral
Sylow $2$-subgroup.

\emph{The verified Gorenstein--Walter application.}
Let $S$ be such a Sylow $2$-subgroup and $z\in Z(S)$ an involution.
The quotient of $C_G(z)$ by $\langle z\rangle$ in (D.31.7)
preserves the two branch points of the rotation quotient.
Its pointwise branch stabilizer is cyclic. Its inverse image
$C_0\leq C_G(z)$ has index at most two and is a central extension
of a finite cyclic group by $\langle z\rangle$, hence is abelian:
a lift of a quotient generator and the central kernel generate it.
The odd-order part of $C_0$ is characteristic in $C_0$, normal in
$C_G(z)$, and cyclic because all odd Sylow subgroups are cyclic.
It is a normal abelian $2$-complement of $C_G(z)$.

We can therefore apply the original, precisely matching theorem
\cite[Theorem~I, p.~553]{A-GorensteinWalter1962}. It states that
$G$ has a normal odd-order subgroup $K$ such that $G/K$ is one
of
\[
       \operatorname{PSL}_2(q),\quad
       \operatorname{PGL}_2(q)\quad(q\text{ odd}),
       \qquad A_7,\qquad\text{a Sylow }2\text{-subgroup}.
                                                               \tag{D.31.9}
\]
The source explicitly includes $C_2^2$ among dihedral groups.
All its hypotheses have just been verified. The odd group $K$
is solvable by the preceding transfer argument, so the triviality
of the solvable radical gives $K=1$.

The group $A_7$ contains the commuting independent cycles
$(1\ 2\ 3)$ and $(4\ 5\ 6)$, contradicting cyclic odd Sylow
subgroups. A $2$-group is solvable. Write $q=p^f$ in either
matrix case. The unipotent subgroup
$u(t)=\left(\begin{smallmatrix}1&t\\0&1\end{smallmatrix}\right)$
for $t\in\F_q$ is elementary abelian of order $p^f$, so $f=1$.
In $\operatorname{PSL}_2(p)$ the diagonal matrices
$\operatorname{diag}(a,a^{-1})$ act on this $C_p$ by
$u(t)\mapsto u(a^2t)$, supplying $(p-1)/2$ distinct automorphisms.
In $\operatorname{PGL}_2(p)$ the matrices
$\operatorname{diag}(a,1)$ supply all $p-1$ automorphisms.
Constraint (D.31.8) leaves $p\leq5$ in the first case and
$p\leq3$ in the second. At $p=3$, the faithful projective-line
action embeds these groups into $S_4$, so they are solvable.
The remaining group is $\operatorname{PSL}_2(5)\cong A_5$.

For this last isomorphism the exact map is conjugation on the
five Sylow $2$-subgroups. The supplied matrix enumeration finds
fifteen involutions, partitions them into the five four-groups,
and verifies that the conjugation image has order sixty and
consists of even permutations. With the recorded ordering of
the five subgroups, the matrices
\[
 \begin{pmatrix}1&1\\0&1\end{pmatrix},\qquad
 \begin{pmatrix}0&-1\\1&0\end{pmatrix}
\]
map respectively to $(1\ 2\ 4\ 3\ 5)$ and
$(2\ 4)(3\ 5)$. Thus the map is injective and its image is exactly
$A_5$, with its sources and targets fixed.

\emph{The puncture representation.}
The pure positive sphere group is the ordered configuration
group of $n-3$ points on the three-punctured sphere: fix the
first three marked points and use their contractible identity
homeomorphism component in the evaluation fibration. It is good
of finite cohomological dimension by the preceding configuration
argument, so its completion is torsion-free. This is exactly
the kernel of the completed map to $S_n$, by the finite-index
completion sequence. Thus every finite subgroup acts faithfully
on the puncture set.

For a finite $A_5$ subgroup, each puncture stabilizer is cyclic by
(D.31.4)--(D.31.5). The cyclic subgroups of $A_5$ have orders
$1,2,3,5$, so its transitive puncture orbits are precisely
$A_5/1$, $A_5/C_2$, $A_5/C_3$, and $A_5/C_5$.
The last three have sizes $30,20,12$. For each $d=2,3,5$,
an element generating $C_d$ fixes
$|N_{A_5}(C_d)/C_d|=2$ points in the corresponding orbit and
none in the other three types. Prime-power conjugacy makes its
puncture permutation that of a rotation of the sphere, which
has at most two fixed punctures. Each exceptional orbit occurs
at most once. This proves (D.31.1), including its actual
permutation representation, and gives the lower bound twelve.

\emph{The extended sphere group.}
Let $G\leq\widehat{M_{0,n}^{\pm}}$ be finite and nonsolvable.
Its positive subgroup $G^+$ has index at most two and is
nonsolvable, so is $A_5$. The centralizer $C_G(A_5)$ intersects
$A_5$ trivially and injects into $G/G^+$, so has order at most
two. With a nontrivial centralizer, its generator is negative
and $G=A_5\times C_2$. With a trivial centralizer, conjugation
embeds $G$ into $\Aut(A_5)\cong S_5$ with image containing the
inner $A_5$, leaving $A_5$ or $S_5$. The automorphism isomorphism
is induced by the action on the five subgroups $A_4\leq A_5$.
An automorphism fixing all five makes $g$ and its image induce
the same permutation of them for every $g\in A_5$; the faithful
natural action of $A_5$ then makes that automorphism the identity.
Conjugation by $S_5$ is faithful as well and supplies one hundred
twenty automorphisms, proving the isomorphism. In $S_5$, the normalizer of
a $5$-cycle acts on it by all four nonzero exponents modulo five.
This violates the extended version of (D.31.8). Therefore only
$A_5$ and $A_5\times C_2$ occur as abstract nonsolvable groups.

\emph{Forgetting punctures in genus two.}
Evaluation on the unordered puncture set, with isotopy extension,
gives the discrete forgetting sequence
\[
 1\longrightarrow B_n(S_2)\longrightarrow M_{2,n}^+
        \longrightarrow M_{2,0}^+\longrightarrow1.      \tag{D.31.10}
\]
The identity homeomorphism component of $S_2$ is contractible,
as used in Part C, so there is no preceding evaluation-loop
kernel. The quotient is good, the braid kernel is finitely
generated, and (D.31.2) completes this sequence exactly.
The completed kernel is torsion-free by the configuration
argument. Every finite subgroup of $\widehat{M_{2,n}^+}$
therefore injects under forgetting into $\widehat{M_{2,0}^+}$.
The hyperelliptic quotient is the exact sequence
\[
 1\longrightarrow C_2\longrightarrow M_{2,0}^+
        \longrightarrow M_{0,6}^+\longrightarrow1;
                                                               \tag{D.31.11}
\]
see \cite[Propositions~3.2--3.3]{A-BigelowBudney2001}.
The central involution is detected on homology modulo three,
so it injects into the completion. Finite-kernel completion
then makes (D.31.11) exact with hats. Every finite subgroup
has solvable image in $\widehat{M_{0,6}^+}$ by (D.31.1),
and kernel contained in $C_2$. It is solvable. The same holds
in the extended group because the positive subgroup has
index at most two. The case $n=0$ uses (D.31.11) directly.

\emph{Forgetting punctures in genus one.}
For $n\geq1$, evaluation gives
\[
 1\longrightarrow N_n\longrightarrow M_{1,n}^+
       \longrightarrow\operatorname{SL}_2(\Z)
       \longrightarrow1,\qquad
 N_n=B_n(T^2)/\Z^2.                                    \tag{D.31.12}
\]
Here $\Z^2$ is the image of the translation loops in the
identity homeomorphism component, identified by evaluation
as in C.9. It is central in the braid group: applying a loop
of translations to every configuration of a braid gives the
square homotopy interchanging the two loop products. Its injection
can be checked on the ordered configuration space by evaluating
the first strand. To verify its full topology in the unordered
braid group, sum the homology classes of all strand paths:
\[
 B_n(T^2)\longrightarrow H_1(T^2,\Z)=\Z^2.
\]
Endpoints are permuted, so the sum is a cycle and gives a
homomorphism. On diagonal translation loops it is multiplication
by $n$. The resulting completed injection proves the exact
central extension
\[
 1\longrightarrow\widehat\Z^2\longrightarrow
 \widehat{B_n(T^2)}\longrightarrow\widehat N_n
       \longrightarrow1.                              \tag{D.31.13}
\]
The completed braid group is torsion-free by the configuration
argument. For completeness the ordered torus configuration
has the explicit product coordinates
\[
 (x_1,\ldots,x_n)\longmapsto
 (x_1,x_2-x_1,\ldots,x_n-x_1)
 \in T^2\times\operatorname{Conf}_{n-1}(T^2\setminus\{0\}),
\]
so its goodness also follows directly from that of $\Z^2$
and the iterated punctured-surface configuration groups.

For a finite $K\leq\widehat N_n$, pull it back in (D.31.13).
Transfer to the central open $\widehat\Z^2$ has restriction
$a\mapsto |K|a$. Its kernel again injects into $K$ and is
finite, so torsion-freeness makes transfer injective.
The pullback, and hence $K$, is abelian. Since
$\operatorname{SL}_2(\Z)$ is good and $N_n$ is finitely
generated, (D.31.2) also completes (D.31.12) exactly.
The image of every finite subgroup lies in a finite subgroup
of $\widehat{\operatorname{SL}_2(\Z)}$, which is cyclic
of order $1,2,3,4$, or $6$ by D.8's virtually free comparison.
The finite kernel is abelian by (D.31.13), proving solvability.
For orientation reversal the positive subgroup has index at
most two, which preserves this conclusion. For $n=0$ use
D.8 directly with $\operatorname{GL}_2(\Z)$.

The executed program
\path{code/part_a/sphere_completion_restriction.py}
checks the projective-matrix automizer calculations at
$p=3,5,7,11,13$, constructs the displayed
$\operatorname{PSL}_2(5)\to A_5$ isomorphism, compares
the $C_5$ normalizer actions in $A_5$ and $S_5$, and
checks the exceptional coset fixed-point counts and the
$C_3^2$ subgroup of $A_7$. The general reductions are
the completed exact sequences and transfer arguments above.
\end{proof}

\begin{claim}{D.32}{Proved}{The three-cusp rank test and finite pure loop groups}
Let $F=\widehat F(x,y)$, with the three named peripheral generators
$c_1=x$, $c_2=y$, and $c_3=(xy)^{-1}$.  Let $P$ be the subgroup of
$\Out(F)$ preserving each of their three procyclic conjugacy classes.
Its common peripheral multiplier is denoted by $\lambda$.
Every finite subgroup of $P$ maps injectively, by $\lambda_2$, to
$\{1,-1\}\leq\Z_2^\times$.  In particular it has order at most two,
and a nonidentity element of finite order has order two and
$\lambda_2=-1$.  These statements hold for $GT_0\leq P$ as well.

For the full three-cusp loop-outer group, allowing puncture permutations,
every finite subgroup $G$ has the explicit faithful representation
\[
 G\longrightarrow S_3\times C_2,\qquad
       g\longmapsto\bigl(\sigma_g,\lambda_2(g)\bigr).
                                                        \tag{D.32a}
\]
This is an assertion about the displayed homomorphism and abstract
finite groups.  Conjugacy of its original action to the geometric
action with this image is not asserted here. D.34 determines the
global multiplier of a surviving involution; its geometric conjugacy
remains to be proved.
\end{claim}
\passage{30; folios 194--196; scans 383--387; 31; folios 198--199 verso; scans 391--394; Part D(b) of the task}

\begin{proof}
We first prove the closed-subgroup fact which uses all three cusps.
For a prime $p$, let $e_p\in\widehat\Z$ have $p$-coordinate one and
all other coordinates zero.  Suppose a closed subgroup $C\leq F$
has $\dim_{\F_p}H^1(C,\F_p)=2$ and contains, for each $i=1,2,3$,
a conjugate
\[
                    g_i c_i^{e_p\Z_p}g_i^{-1}.             \tag{D.32b}
\]
We prove $C=F$.

The comparison with open subgroups is the actual restriction map
\[
 \underset{U\text{ open},\ C\leq U\leq F}{\operatorname{colim}}
       H^1(U,\F_p)\ \xrightarrow{\ \sim\ } H^1(C,\F_p),
                                                        \tag{D.32c}
\]
where the transition maps are restrictions as $U$ decreases.
Indeed a continuous character of $C$ kills $C\cap V$ for some open
normal $V\triangleleft F$.  It extends to $CV$ by the quotient
$CV/V=C/(C\cap V)$.  Conversely a character on $U$ which vanishes
on $C$ vanishes after restriction to its open kernel, which contains
$C$.  This proves both surjectivity and injectivity of (D.32c).

If $C$ is open, finite-coset Schreier rewriting identifies it with
a free profinite group of rank $1+[F:C]$.  Its given first-cohomology
dimension therefore forces $[F:C]=1$.
Suppose instead that $C$ has infinite index.  The $p$-valuations
of the finite indices $[F:U]$, for open $U$ containing $C$, must
be unbounded.  To prove this, assume they were bounded and choose
$U_0$ attaining their maximum.  For $C\leq V\leq U_0$, every
further restriction between such open subgroups has prime-to-$p$
index and is injective on $H^1(-,\F_p)$: a nonzero character has
an index-$p$ kernel and cannot vanish on a subgroup of prime-to-$p$
index.  Thus (D.32c) contains the first cohomology of every such
$V$.  The indices $[F:V]$ are unbounded, since their intersection
is $C$ and $[F:C]=\infty$.  Their Schreier dimensions
$1+[F:V]$ are therefore unbounded, a contradiction.

Choose an open $U\geq C$ for which $p\mid [F:U]$.  Put
$\Delta=F(x,y)$, the named dense discrete subgroup of $F$, and
$U_0=U\cap\Delta$.  The finite covering of the pair of pants
defined by $U_0\leq\Delta$ has degree $n=[F:U]$, and its
completed fundamental group is exactly $U$: finite-index subgroups
of a discrete group carry the full induced profinite topology.
For each $g_i$ choose $h_i\in\Delta$ in its left $U$-coset and
write $g_i=u_i h_i$, with $u_i\in U$.  The subgroup
\[
 U\cap g_i\overline{\langle c_i\rangle}g_i^{-1}
  =u_i\overline{\langle h_i c_i^{d_i}h_i^{-1}\rangle}u_i^{-1}
                                                        \tag{D.32d}
\]
is consequently the completion, up to the displayed inner
conjugation, of a particular boundary subgroup in this finite
cover.  Here $d_i$ is its boundary covering degree.  Inclusion
(D.32b) proves $p\nmid d_i$.

The three selected boundaries lie over the three different original
boundaries, so they are distinct.  There is at least one more
boundary.  Otherwise each original boundary would have just one
lift, and its degree would be $n$, contrary to $p\nmid d_i$ and
$p\mid n$.  If the cover has genus $h$ and $b\geq4$ boundary
components, its presentation is
\[
 \left\langle a_1,b_1,\ldots,a_h,b_h,d_1',\ldots,d_b'
       \ \middle|\ \prod_{j=1}^h[a_j,b_j]\,
                         d_1'\cdots d_b'=1\right\rangle.
\]
Thus the only relation between its boundary homology classes over
$\F_p$ is their sum.  The three selected classes are independent,
and evaluation on them gives a surjection
$H^1(U,\F_p)\to\F_p^3$.

This remains a surjection after restriction of these characters to
$C$ and evaluation on the three elements in (D.32b).
To check the comparison explicitly, the element
$g_i c_i^{e_p}g_i^{-1}$ is the $e_p/d_i$ power of the boundary
generator in (D.32d).  Inner conjugation does not change a
character, and its evaluation is therefore multiplied by the
nonzero scalar $d_i^{-1}\in\F_p$.  We obtain three independent
characters of $C$, contradicting its dimension two.  This proves
the closed-subgroup assertion.

We now apply it to finite loop actions, retaining the precise
normalization map.  If an outer class fixes the three peripheral
classes, its exponents on $c_1,c_2,c_3$ are equal: in
$F^{\rm ab}=\widehat\Z^2$, the relation
$[c_1]+[c_2]+[c_3]=0$ gives equality of the two coordinates.
Conjugate a representative $\beta$ to have
$\beta(x)=x^\lambda$, and write
$\beta(y)=k^{-1}y^\lambda k$.  Write $[k]=(a,b)$ in the
named abelian basis.  Conjugating this representative by $x^a$
gives the normalized pair
\[
 \alpha(x)=x^\lambda,\qquad
 \alpha(y)=f^{-1}y^\lambda f,
       \qquad f=y^{-b}kx^{-a}\in F'.                     \tag{D.32e}
\]
Multiplication on the left of the conjugator by $y^{-b}$ changes
neither image of $y$ nor the automorphism.  Two normalized
representatives differ by conjugation by an element centralizing
$x^\lambda$, hence by $x^v$ by A.7.  Their two conjugators then
differ by left multiplication by a power of $y$ and right
multiplication by $x^{-v}$.  Since both are in $F'$, their two
abelian coordinates force both powers to be zero.  Normalization
is therefore unique.  A product of normalized pairs has second
coordinate $f\alpha(g)\in F'$, so (D.32e) is a homomorphic
section into $\Aut(F)$.  Its image is precisely the Belyi group
defined in \cite[p.~3120, equation (2)]{R-HS}; the argument has
not identified that entire group with its subgroup $GT_0$.

Let an element of $P$ have prime order $p$ with $\lambda_p=1$.
Its normalized automorphism has order $p$.  The proof of D.28
uses exactly the properties just established: finite order of the
normalized automorphism, the common multiplier, and preservation
of the three cusps.  Neither of the extra $GT_0$ symmetry equations
is used in its norm computation or Brown restriction argument.
In particular that proof gives a single conjugacy class of
sections in $F\rtimes C_p$ and the actual fixed subgroup
$C=C_F(C_p)$ with $\dim H^1(C,\F_p)=2$.
Each of the three peripheral sections centralizes the whole
$c_i^{e_p\Z_p}$; conjugating those sections to the specified
$C_p$ places the three groups (D.32b) in this same $C$.
The closed-subgroup assertion gives $C=F$, so the normalized
automorphism is the identity.  This contradicts its prime order.

For odd $p$, an element of order dividing $p$ in
$\Z_p^\times$ equals one, by reduction modulo $p$ and the
torsion-free group $1+p\Z_p$.  Thus no odd prime occurs in
the order of a finite subgroup of $P$.  The finite image of
$\lambda_2$ lies in $\{1,-1\}$, since $1+4\Z_2$ is
torsion-free.  Its kernel in any finite subgroup has no element
of order two by the preceding argument, and no element of odd
prime order either.  Cauchy's theorem makes that kernel trivial.
This proves the first assertion.

For a class permuting cusps, compose it with the inverse of the
corresponding positively oriented geometric permutation.  The
preceding abelian calculation gives its common multiplier too,
and composition on cusp generators proves that both $\sigma$
and $\lambda$ are homomorphisms.  On a finite subgroup, the
kernel of (D.32a) is a finite subgroup of $P$ on which
$\lambda_2=1$, hence is trivial.  All subgroups of the target
already have their geometric actions on the three-punctured
sphere: use an equilateral triple on the equator, the positive
rotational $S_3$, and reflection in that equator for the
commuting $C_2$.  This proves the claimed abstract group bound
and its attainment, while keeping the action-conjugacy question
separate through the actual map (D.32a).

The program \texttt{code/part\_c/three\_cusp\_rank\_d32.py}
enumerates the finite covers used for a validation through degree
six.  It computes the Schreier graph, every boundary cycle,
the unique boundary-homology relation, and explicit dual
characters for the selected boundaries, including the
$d_i^{-1}$ restriction scalars.  The proof above applies to
arbitrary degrees and closed subgroups; no bounded enumeration
is used in place of it.
\end{proof}

\paragraph{Source check for the fixed-second-coordinate assertion.}
Lochak--Schneps, Lemma~9 on p.~583 of
\cite{R-LS}, asserts that a fixed $f$ occurs only with the two
multipliers $\lambda$ and $-\lambda$.  The displayed class-two
identity $\lambda^2=24r+1$ determines the squares of the prime
coordinates.  The second proof in that source adds the required
full-group equation, whose elementary verification is as follows.

Comparing $(\lambda,f)$ and $(\mu,f)$ by the exact group law
gives $(\lambda^{-1}\mu,1)$.  Set
$n=(\lambda^{-1}\mu-1)/2\in\widehat\Z$.
The three-cycle relation is
\[
                       x^n(xy)^{-n}y^n=1.                 \tag{LS.1}
\]
In the class-two Heisenberg quotients its central coordinate
gives $n(n+1)=0$.  Thus in each integral domain $\Z_p$ one
has $n_p=0$ or $n_p=-1$.  These choices are linked by actual
finite quotients: for every odd prime $q$, evaluate (LS.1) in
\[
 D_q=\langle r,s\mid r^q=s^2=1,\ srs=r^{-1}\rangle,
                         \qquad x=s,\quad y=r.
\]
Both $x$ and $xy$ have order two.  If $n_2=0$, (LS.1) is
$r^{n_q}=1$, forcing $n_q=0$.  If $n_2=-1$, it is
$r^{1+n_q}=1$, forcing $n_q=-1$.  Hence $n$ is globally
zero or globally minus one, proving $\mu=\lambda$ or
$\mu=-\lambda$.  Both occur, because multiplication on the
right by $(-1,1)$ leaves the second coordinate unchanged.

This verifies the source's fixed-$f$ conclusion by exact
quotient morphisms.  Powers of a general torsion element
have second coordinates $f\alpha(f)\cdots\alpha^{j-1}(f)$;
the argument does not identify those coordinates with each
other.  D.32 supplies the separate argument restricting
finite subgroups, and leaves the stated involution questions
unchanged.

\begin{claim}{D.33}{Proved}{Finite punctured-torus actions are detected by their primary homology}
Put $F=\widehat F(x,y)$ and $c=[x,y]$, and let
$T=\Out_{\rm lac}(F;\overline{\langle c\rangle})$ be the full
once-punctured-torus loop-outer group.  For every prime $p$, the
homomorphism induced by the named abelianization,
\[
              A_p:T\longrightarrow\GL_2(\Z_p),            \tag{D.33a}
\]
is injective on every finite $p$-subgroup of $T$.
The full abelianization homomorphism
$T\to\GL_2(\widehat\Z)$ is therefore injective on every finite
subgroup.  Every such subgroup has order dividing $24$.

The fixed-subgroup calculation underlying these assertions is also
explicit.  For an automorphism $\alpha$ of $\widehat F_r$ with
$\alpha^p=1$ whose induced automorphism on the maximal pro-$p$
quotient is the identity, put $E=\widehat F_r\rtimes_\alpha C_p$.
There is one kernel-conjugacy class of sections of $E\to C_p$,
and for its fixed subgroup $C$ the actual restriction map is an
isomorphism
\[
 H^1(\widehat F_r,\F_p)\ \xrightarrow{\ \sim\ }H^1(C,\F_p).
                                                        \tag{D.33b}
\]
For the punctured-torus loop structure and $r=2$, its cusp forces
$C=F$ and hence $\alpha=1$.
\end{claim}
\passage{30; folios 194--196; scans 383--387; 31; folios 198--199 verso; scans 391--394; Part D(b) of the task}

\begin{proof}
We first prove (D.33b), keeping the multiplicative information in
the cohomology comparison.  Write $P=\widehat F_r^{(p)}$ and
$S=\{1\}\times C_p\leq E$.  The quotient homomorphism is
\[
 \pi:E\longrightarrow P\times C_p,\qquad
              (h,s^i)\longmapsto(\bar h,s^i).              \tag{D.33c}
\]
Restriction along $\widehat F_r\to P$ is an isomorphism on
mod-$p$ cohomology: in degree one both sides are the same
$r$-dimensional character group, and higher cohomology vanishes
because both groups have $p$-cohomological dimension one.
The Hochschild--Serre comparison for (D.33c) is consequently an
isomorphism on its second pages and gives a ring isomorphism
\[
 \pi^*:H^*(P\times C_p,\F_p)\xrightarrow{\sim}H^*(E,\F_p).
                                                        \tag{D.33d}
\]
All maps here are induced by the specified quotient, not by an
identification of the two groups.  In particular
$\dim H^n(E,\F_p)=r+1$ for every $n>1$.

Let $u\in H^2(C_p,\F_p)$ be the nonzero periodic class, and
inflate it to either side of (D.33d).  The K\"unneth calculation
for the direct product gives, for every $m\geq1$, an injective map
\[
 H^1(P,\F_p)\longrightarrow H^{2m+1}(E,\F_p),
                         \qquad\chi\longmapsto\pi^*(\chi u^m).
                                                        \tag{D.33e}
\]
For a subgroup $Q\leq E$ of order $p$, its projection to
$C_p$ is an isomorphism, and
$N_E(Q)=C_{\widehat F_r}(Q)\times Q$, exactly as in D.28.
The first coordinate of $\pi(Q)$ is trivial because $P$ is
torsion-free.  Thus every character $\pi^*\chi$ in (D.33e)
vanishes on $Q$.  Its restriction to this normalizer lies
entirely in $H^1(C_{\widehat F_r}(Q),\F_p)$, and multiplication
by $u^m$ places it in
\[
 H^1(C_{\widehat F_r}(Q),\F_p)\otimes H^{2m}(Q,\F_p).
                                                        \tag{D.33f}
\]

Apply the profinite Brown restriction theorem in the precise
form verified in \cite[Theorem 2 and \S1.1]{A-KochloukovaZalesskii2007}.
There is no $C_p\times C_p$ in $E$, since a finite subgroup
injects into $E/\widehat F_r=C_p$.  Its open subgroup
$\widehat F_r$ has $p$-cohomological dimension one.  The
finite-dimensional source makes the theorem's dense restriction
embedding an isomorphism onto a finite product.  If $k$ is
the number of conjugacy classes and
$d_Q=\dim H^1(C_{\widehat F_r}(Q),\F_p)$, its dimension formula is
\[
                   k+\sum_{[Q]}d_Q=r+1.                  \tag{D.33g}
\]
The injection (D.33e), followed by this same restriction map,
lands in the summands (D.33f).  Therefore
$\sum d_Q\geq r$.  Since $k\geq1$, (D.33g) forces
$k=1$ and $d_Q=r$.  For the specified $S$, (D.33e)--(D.33f)
then say exactly that restriction
$H^1(\widehat F_r,\F_p)\to H^1(C_{\widehat F_r}(S),\F_p)$
is injective; both spaces have dimension $r$, proving (D.33b).
As in D.28, conjugacy under $E$ equals conjugacy under the kernel,
because $E=\widehat F_r Q$ for any section $Q$.

Now take $r=2$ and assume that $\alpha$ preserves the conjugacy
class of $L=\overline{\langle c\rangle}$, where $c=[x,y]$.
Write its peripheral multiplier as $\lambda$.  In a class-two
Heisenberg quotient, the image of the commutator under an
abelianization matrix $A$ is the $\det(A)$ power of the
commutator.  The finite Heisenberg quotients for all moduli
therefore give the exact identity
\[
                        \lambda=\det(\alpha^{\rm ab}).   \tag{D.33h}
\]
In the present calculation $\alpha$ is the identity on $P$,
so $\lambda_p=1$.

The cusp normalizer has the exact sequence
\[
 1\longrightarrow L\longrightarrow N_E(L)
                        \longrightarrow C_p\longrightarrow1.
                                                        \tag{D.33i}
\]
The kernel is $L$ by A.7.  For a lift $a$ of the chosen
generator, write
$aca^{-1}=c^\lambda$ and $a^p=c^v$, so $(\lambda-1)v=0$.
The first coordinate of $\pi(a)$ centralizes the image $\bar c$
in $P$ and belongs to $\overline{\langle\bar c\rangle}$,
by the pro-$p$ peripheral centralizer assertion of A.7.
Thus it is $\bar c^w$ and $v_p=pw$.
This uses the whole cusp subgroup.  Its primitive pro-$p$
parameter can also be checked directly: in the mod-$p$ Magnus
expansion, $[x,y]-1$ starts with $XY-YX$.  A $p$th power
starts in degree divisible by $p$.  At $p=2$, the possible
degree-two square is $(aX+bY)^2$; its $X^2$ and $Y^2$
coefficients prevent it from equaling $XY+YX$.  Hence the
commutator is not a $p$th power.

Since $\alpha^p=1$, its multiplier satisfies $\lambda^p=1$.
Set $N_\lambda=1+\lambda+\cdots+\lambda^{p-1}$.
At a prime $q$ with $\lambda_q\ne1$, both
$(N_\lambda)_q$ and $v_q$ are zero.  Where $\lambda_q=1$
the norm is $p$; $v_q$ is divisible by $p$, using $v_p=pw$
at $q=p$ and invertibility of $p$ at all other primes.
Consequently there is $z\in\widehat\Z$ with
$N_\lambda z=-v$.  The element $ac^z$ has order $p$
and centralizes $c^{e_p\Z_p}$.
The unique section-conjugacy class just proved places a
conjugate $g c^{e_p\Z_p}g^{-1}$ inside
$C=C_F(S)$.

Every character of $C$ vanishes on $g c^{e_p}g^{-1}$:
by (D.33b) it is the restriction of a character of $F$,
and $c$ is a commutator.  We prove that this forces $C=F$.
Its mod-$p$ first-cohomology dimension is two.  If $C$ is
open, Schreier gives $1+[F:C]=2$.  If $C$ is proper and
not open, the restriction-colimit proof of D.32 shows that
its index has infinite $p$-part.  Choose an open $U\geq C$
with $p\mid n=[F:U]$, and take the degree-$n$ covering of
the punctured torus corresponding to $U\cap F(x,y)$.
The chosen cusp has a boundary lift of degree $d$ prime
to $p$, since its subgroup contains the full chosen
$c^{e_p\Z_p}$.  The comparison with a discrete boundary
generator is exactly (D.32d), with a dense discrete
representative for the profinite conjugator $g$.
There must be at least two boundary lifts: a unique lift
would have degree $n$, divisible by $p$.
The selected boundary has nonzero homology in this cover,
since the sole relation among the at least two boundary
classes is their sum.  A character of $U$ detecting this
boundary restricts to a character of $C$ detecting
$g c^{e_p}g^{-1}$; the value is multiplied by
$d^{-1}\bmod p$, as computed in D.32.  This is the required
contradiction.  Hence $C=F$ and $\alpha=1$.

To apply this to (D.33a), let an outer element have order
$p$ and $A_p=I$.  D.27 supplies an actual order-$p$
automorphism representing it.  Its induced pro-$p$
automorphism has identity abelianization.  The Magnus
argument in D.10 and D.25 makes this finite-order
automorphism the identity: a first nonzero homogeneous
deviation in the $\Z_p$ Magnus algebra would be multiplied
by $p$ under the $p$th iterate, and the homogeneous module
has no $p$-torsion.  Thus (D.33c) applies to this lift
directly; there is no subsequent conjugation which could
change its order.  The cusp argument proves that the lift,
and hence its outer class, is the identity.
The kernel of (D.33a) on a finite $p$-subgroup cannot
contain an order-$p$ element, so it is trivial.  Likewise,
the kernel of full abelianization on any finite subgroup
would contain an element of some prime order and is trivial.

We finish with the numerical bound on these finite groups.
For $p>3$, a matrix of order $p$ over $\Q_p$ has a primitive
$p$th root among its eigenvalues.  The polynomial
$\Phi_p(t)$ is irreducible of degree $p-1$ over $\Q_p$,
as seen by Eisenstein after $t\mapsto t+1$.  It cannot be
the factor of a degree-two characteristic polynomial.
Thus no such prime divides the order of a finite subgroup
of $T$.  A finite $3$-subgroup injects by $A_3$ into
$\GL_2(\Z_3)$ and then injects modulo $3$: the kernel
$1+3M_2(\Z_3)$ is torsion-free by the matrix logarithm.
Since $|\GL_2(\F_3)|=48$, its order divides $3$.

For completeness, a finite $2$-subgroup $K$ of
$\GL_2(\Q_2)$ has order at most eight.  The same
Eisenstein argument for $\Phi_{2^a}$ excludes order
$2^a$ for $a\geq3$.  Every matrix of order four has
minimal polynomial $t^2+1$, so its square is $-I$.
Consequently the projective image of $K$ has exponent
two.  Extend the matrices to the algebraic closure
$\overline{\Q}_2$; the scalar kernel is unchanged.
An exponent-two subgroup of $\PGL_2(\overline{\Q}_2)$
has order at most four.  To prove this
last bound, conjugate a nonidentity element to
$z\mapsto-z$.  Every commuting transformation preserves
its pair of fixed points, so is $z\mapsto az$ or
$z\mapsto a/z$.  The scaling subgroup has at most two
elements under the exponent-two condition and has index
at most two.  The scalar kernel of $K$ is contained in
$\{I,-I\}$, the finite $2$-torsion of $\Q_2^\times$.
Thus $|K|\leq8$.  Apply this to a Sylow $2$-subgroup
through $A_2$.  The Sylow bounds and the exclusion of
other primes give $|G|\mid 8\cdot3=24$.

The program \texttt{code/part\_c/torus\_homology\_d33.py}
checks the Magnus degree-two coefficients, the periodic
cohomology multiplication used in (D.33e), and the
one-cusp restriction argument on every connected finite
cover through degree six.  The matrices and boundary
characters are obtained from actual Schreier walks.
The bound and comparison maps above are proved for the
full profinite groups, independently of the bounded checks.
\end{proof}

\begin{claim}{D.34}{Proved}{Global signs and the cusp cycle of a pure involution}
Every nonidentity finite-order element of the full pure
three-cusp loop-outer group $P$ of D.32 is an involution
whose common multiplier is the global element
$-1\in\widehat\Z^\times$.  In particular this holds in
$GT_0$, and the multiplier-one kernel in either group is
torsion-free.

For its normalized lift
\[
 \alpha(x)=x^{-1},\qquad\alpha(y)=f^{-1}y^{-1}f,
                      \qquad f\in F',
                                                        \tag{D.34a}
\]
there are precisely three $F$-conjugacy classes of sections
of $F\rtimes_\alpha C_2\to C_2$.  Each of their fixed
subgroups has zero first cohomology with $\F_2$ coefficients.
The two section classes in each cusp normalizer are the
endpoints of one edge of a triangle.  Moreover there is
$h\in F'$ with
\[
                         f=h^{-1}\alpha(h).               \tag{D.34b}
\]
Thus $\alpha$ inverts both $x$ and $hyh^{-1}$.
The endomorphism $x\mapsto x$, $y\mapsto hyh^{-1}$ induces
an automorphism of every maximal pro-$p$ quotient.  Its
surjectivity on the full $F$, and conjugacy of the original
involution to complex conjugation in the pure loop group,
are not conclusions of this statement.
\end{claim}
\passage{30; folios 194--196; scans 383--387; 31; folios 198--199 verso; scans 391--394; Part D(b) of the task}

\begin{proof}
We first prove a genus-zero incidence calculation for an
arbitrary number $n\geq3$ of cusps.  Let
\[
 F_n=\widehat{\langle c_1,\ldots,c_n\mid
                              c_1\cdots c_n=1\rangle}
\]
and let $\beta$ be an actual involution preserving each
peripheral conjugacy class, with common multiplier $\mu$
and $\mu_2=-1$.  Put $E_n=F_n\rtimes_\beta C_2$.
The equality $\beta^2=1$ gives $\mu^2=1$; hence each
$\mu_q$ is $1$ or $-1$.

For each cusp $L_i=\overline{\langle c_i\rangle}$,
A.7 gives an exact sequence
\[
 1\longrightarrow L_i\longrightarrow N_{E_n}(L_i)
                          \longrightarrow C_2\longrightarrow1.
                                                        \tag{D.34c}
\]
Choose a lift $a_i$ of the nonidentity element and write
$a_i c_i a_i^{-1}=c_i^\mu$ and $a_i^2=c_i^{v_i}$.
Then $(\mu-1)v_i=0$.  In every coordinate with
$\mu_q=-1$ this forces $(v_i)_q=0$, including $q=2$.
In a coordinate with $\mu_q=1$ the prime $q$ is odd,
so division by two is possible.  Thus
$(1+\mu)u_i=-v_i$ has a solution in $\widehat\Z$,
and $b_i=a_i c_i^{u_i}$ is an involution.

The set of its further order-two translates by $L_i$,
modulo conjugation by $L_i$, is exactly
\[
 \frac{\ker(1+\mu:\widehat\Z\to\widehat\Z)}
      {(1-\mu)\widehat\Z}\ \simeq\ \F_2.               \tag{D.34d}
\]
At an odd prime with $\mu_q=-1$, multiplication by two
is onto, and a coordinate with $\mu_q=1$ contributes
zero to the kernel.  The remaining coordinate is
$\Z_2/2\Z_2$.  Representatives for the two classes are
$b_i$ and $c_i^{e_2}b_i$.

Since $\beta$ is pure and the cusp classes span homology,
its action on $V=H_1(F_n,\F_2)$ is the identity.  Therefore
the actual mod-two abelianization of the semidirect product is
\[
                  H_1(E_n,\F_2)=V\oplus\F_2.             \tag{D.34e}
\]
The two involutions just exhibited have images differing by
the nonzero vector $[c_i]$ in $V$.  They cannot be conjugate
in $E_n$ or in $F_n$.  Thus the two local classes in
(D.34d) give two distinct global section classes.

Build a finite graph with one edge for each of the $n$
cusps and with these two global classes as its endpoints.
This graph uses some of the global section classes; we
now bound all of them.  The Brown calculation of D.28,
with free rank $n-1$ and prime two, gives
\[
 \sum_{[Q]}
   \bigl(1+\dim_{\F_2}H^1(C_{F_n}(Q),\F_2)\bigr)
       =\dim_{\F_2}H^k(E_n,\F_2)\leq n\qquad(k>1).
                                                        \tag{D.34f}
\]
Every order-two subgroup is a section, and its total-group
and kernel conjugacy classes agree, because $E_n=F_nQ$.
Hence there are at most $n$ vertices available for the graph.
An edge has no loop by (D.34e).

Assign to every vertex its image $(v,1)$ in (D.34e).
The sum of the edge vectors around a graph cycle is zero,
since its successive vertex differences cancel.  But the
$n$ vectors $[c_1],\ldots,[c_n]$ have just one nonzero
linear relation over $\F_2$, namely
$[c_1]+\cdots+[c_n]=0$.  Consequently any simple cycle
in the graph uses all $n$ edges.  The graph cannot be a
forest: a forest with $n$ edges needs at least $n+1$
vertices, more than (D.34f) permits.  Thus it is exactly
a simple $n$-cycle.  It has $n$ vertices, so there are
no additional section classes outside this graph, and
(D.34f) gives
\[
            H^1(C_{F_n}(Q),\F_2)=0
                   \quad\text{for every section }Q.       \tag{D.34g}
\]
This construction is invariant under any automorphism
of the extension preserving its kernel and peripheral
classes, with any permutation of those classes: such
an automorphism sends the normalizer and its two classes
in (D.34c)--(D.34d) to those of the image cusp.

Return to a nontrivial finite-order element of $P$.
D.32 gives its order two, its normalized actual lift
$\alpha$, and $\lambda_2=-1$.  We show that
$\lambda_q=-1$ for every odd prime $q$.  The only other
possibility is $\lambda_q=1$.  Suppose this occurs and
use the exact cyclic quotient
\[
 \chi:F=\widehat F(x,y)\longrightarrow C_q,
          \qquad\chi(x)=1,\quad\chi(y)=0,
          \qquad N=\ker\chi.                            \tag{D.34h}
\]
The common multiplier gives
$\chi\alpha=\lambda\chi=\chi$, since its reduction
modulo $q$ is one.  In particular $N$ is invariant.
The corresponding degree-$q$ covering of the pair of
pants has one boundary over $x$, one over
$z=(xy)^{-1}$, and $q$ over $y$, with representative
generators
\[
 x^q,\qquad z^q,\qquad x^i yx^{-i}\quad(0\leq i<q).
                                                        \tag{D.34i}
\]
The two first degrees are $q$ and all others are one.
Its Euler characteristic is $-q$ and it has $q+2$
boundaries, so its genus is zero.  Finite-index
completion identifies its full profinite group with
the specified subgroup $N$.

The restriction of $\alpha$ is pure on this entire
peripheral structure, with the same multiplier
$\lambda$.  The first two boundaries are the unique
lifts of their original boundaries.  For the other
$q$ boundaries use the actual normalized formula
\[
 \alpha(x^i yx^{-i})
        =x^{\lambda i}f^{-1}y^\lambda f x^{-\lambda i}.
                                                        \tag{D.34j}
\]
Here $\chi(f)=0$, because $f\in F'$, and
$\lambda i\equiv i\pmod q$.  The conjugator in
(D.34j) therefore belongs to the same $N$-coset
as $x^i$: explicitly
$k_i=x^{\lambda i}f^{-1}x^{-i}\in N$ satisfies
$\alpha(x^i yx^{-i})=k_i(x^i yx^{-i})^\lambda k_i^{-1}$.
This is the required conjugacy inside $N$.  The action on
$H_1(N,\Z_2)$ is consequently the scalar $-1$,
so the restriction remains a nontrivial involution.

Put $E=F\rtimes_\alpha C_2$ and
$E_N=N\rtimes_\alpha C_2$.  The quotient maps give
\[
                     E/N=C_q\times C_2.                 \tag{D.34k}
\]
In particular $E_N$ is normal in $E$.  Conjugation by
$x$ defines an automorphism of $E_N$ preserving its
map to $C_2$, and thus acts on its $N$-conjugacy
classes of sections.  Its $q$th power is conjugation
by $x^q\in N$, so this action has order dividing
$q$ on those classes and on the cusp incidence graph.
On the cusp edges (D.34i), it fixes the $x^q$ and
$z^q$ classes and cycles the $q$ classes over $y$.
The second fixed class is the unique boundary over
$z$, or follows directly from its cyclic coset
decomposition.  Thus its graph action has order
exactly $q$.

By the preceding calculation the graph is a cycle
with $q+2$ edges.  Its automorphism group has order
$2(q+2)$: an image of a vertex and a choice of which
neighbor is its successor determine the automorphism.
The odd prime $q$ does not divide $2(q+2)$.
This contradicts the displayed order-$q$ graph action.
Therefore $\lambda_q=-1$ for every odd $q$, proving
the global multiplier assertion.  This argument used
the whole pure group $P$, without an additional
$GT_0$ symmetry assumption.

At $n=3$ the incidence graph is the stated triangle.
Write $s$ for the chosen section in $E=F\rtimes C_2$.
The normalized group law gives $f\alpha(f)=1$,
so $fs$ is also an involution and it normalizes $y$
by inversion.  The three involutions
\[
                     s,\quad x^{e_2}s,\quad y^{e_2}fs
                                                        \tag{D.34l}
\]
have the distinct images $(0,0,1)$, $(1,0,1)$,
and $(0,1,1)$ in (D.34e).  They represent all three
classes.  Since $f\in F'$, the image of $fs$ is
$(0,0,1)$, so it is conjugate to $s$ by an element
of $F$.  Choose $h$ with $fs=h^{-1}sh$.  This is
exactly (D.34b).  In abelianization it says
$0=(-2)[h]$ in the torsion-free group
$\widehat\Z^2$, hence $h\in F'$.
Substitution gives
\[
 \alpha(hyh^{-1})
   =(hf)(f^{-1}y^{-1}f)(f^{-1}h^{-1})
   =hy^{-1}h^{-1}.
\]
The two displayed inverted generators have the same
abelian images as $x,y$.  On each maximal pro-$p$
quotient their images therefore generate the Frattini
quotient and hence the whole group.  The induced
endomorphism of that finitely generated free pro-$p$
group is surjective and thus an automorphism.
No full-profinite surjectivity is inferred from those
separate comparisons.

The program \texttt{code/part\_c/pure\_involution\_d34.py}
checks the normalizer orbit calculation, the unique
binary cusp relation, and the exact cyclic-cover
boundary and deck permutations for several odd
primes.  The exclusion of all odd primes is the
unrestricted divisibility argument above.
\end{proof}

\begin{claim}{D.35}{Proved}{One exact conjugator equation and its finite obstructions}
Let $P$ be the full pure three-cusp group of D.32, let
$F=\widehat F(x,y)$, and put
$\delta(x)=x^{-1}$, $\delta(y)=y^{-1}$.
For a nonidentity involution $\alpha\in P$, use its normalized
lift from D.34,
\[
 \alpha(x)=x^{-1},\qquad
 \alpha(y)=f^{-1}y^{-1}f,\qquad f\in F'.
\]
Define the compact solution space
\[
\begin{split}
 \mathcal S_\alpha=\{(\mu,h,u)\in
       \widehat\Z^\times\times F'\times F:\quad&
       h^{-1}\alpha(h)=f,\\
       &x^\mu h y^\mu h^{-1}
                    =u(xy)^\mu u^{-1}\}.
\end{split}                                                   \tag{D.35a}
\]
The assignment
\[
 (\mu,h,u)\longmapsto
 \left[\beta_{\mu,h}\right],\qquad
 \beta_{\mu,h}(x)=x^\mu,\quad
 \beta_{\mu,h}(y)=h y^\mu h^{-1},                              \tag{D.35b}
\]
has image exactly
\[
 \operatorname{Conj}_{P}(\delta,\alpha)
       =\{\beta\in P:\ \beta\delta\beta^{-1}=\alpha\}.
                                                                  \tag{D.35c}
\]
Every map (D.35b) is an automorphism of the full profinite
group.  The fibers of the assignment are exactly
$(\mu,h,u)\mapsto(\mu,h,u(xy)^a)$ for $a\in\widehat\Z$.
It therefore induces a homeomorphism from the corresponding
quotient of $\mathcal S_\alpha$ to (D.35c).

For $\alpha\in GT_0$, the inclusion $GT_0\hookrightarrow P$
gives the equality of actual conjugator sets
\[
       \operatorname{Conj}_{GT_0}(\delta,\alpha)
             =\operatorname{Conj}_{P}(\delta,\alpha).             \tag{D.35d}
\]
The defining Harbater--Schneps relations for the conjugating
pair follow from the proof below; they are not discarded.
Finally, the existence of a solution of (D.35a) is equivalent
to the nonemptiness of its explicit finite solution sets
(D.35h) at every characteristic finite quotient of $F$.
\end{claim}
\passage{\S{}31; folios 198-199verso; scans 391-394; \S{}30; folio 194verso; scan 384; Part D(b): actual full GT\_0 torsion and conjugacy}

\begin{proof}
For $(\mu,h,u)\in\mathcal S_\alpha$, the free universal
property first gives a continuous endomorphism
$\beta_{\mu,h}:F\to F$ with the displayed generator images.
Write $H=\beta_{\mu,h}(F)$, a closed subgroup because $F$
is compact.  It is generated by two elements, so
$\dim H^1(H,\F_2)\leq2$.  The restrictions of the two
standard characters of $F$ to these generators are the
two coordinate characters: $\mu\bmod2=1$ and conjugation
does not change abelianization.  These restrictions are
independent, proving
\[
                         \dim H^1(H,\F_2)=2.                   \tag{D.35e}
\]
The unit exponent $\mu$ shows that $H$ contains
$\overline{\langle x\rangle}$ and
$h\overline{\langle y\rangle}h^{-1}$.
The second equation in (D.35a) puts
$u\overline{\langle xy\rangle}u^{-1}$ in $H$ as well.
In particular $H$ contains the three full cusp
$\Z_2$ subgroups required by D.32.  Its closed-subgroup
rank test, applied to (D.35e), gives $H=F$.
A surjective endomorphism of a finitely generated
profinite group is injective, so $\beta_{\mu,h}$ is
an automorphism.  Its images of the three cusps,
including the exact displayed equation for the third,
place its outer class in $P$.

The first equation in (D.35a) is $\alpha(h)=hf$.
Substitution on the two free generators gives
\[
 \alpha\beta_{\mu,h}(x)=x^{-\mu}
                          =\beta_{\mu,h}\delta(x),\qquad
 \alpha\beta_{\mu,h}(y)=h y^{-\mu}h^{-1}
                          =\beta_{\mu,h}\delta(y).
                                                                  \tag{D.35f}
\]
These are equalities of actual automorphisms, proving
that the image of (D.35b) lies in (D.35c).

Conversely, let $\beta\in P$ satisfy (D.35c), and take its
unique normalized representative from D.32.
Write it as $\beta_{\mu,h}$ with $h\in F'$.
The normalization is a homomorphic section
$P\to\Aut(F)$, so the outer conjugacy equality lifts
to the actual equality $\alpha\beta_{\mu,h}
=\beta_{\mu,h}\delta$.
Its value on $y$ implies that
$h^{-1}\alpha(h)f^{-1}$ centralizes $y^{-\mu}$.
By A.7 it is a power of $y$.  It also lies in $F'$,
so its abelianization forces this power to be zero.
Thus $h^{-1}\alpha(h)=f$.
The third-cusp condition on $\beta$ gives a conjugator
$u$ in the second equation of (D.35a).
This proves that the image is exactly (D.35c).

The normalized pair $(\mu,h^{-1})$ is unique.
For this fixed pair the two third-cusp conjugators
$u,u'$ differ by right multiplication by the
centralizer of $(xy)^\mu$, which is
$\overline{\langle xy\rangle}$ by A.7.
These are precisely the asserted fibers.
The assignment is continuous in the profinite
automorphism topology, and its domain is compact.
Its quotient is therefore homeomorphic to its
Hausdorff image.

For an alternative form using D.34, fix
$h_0\in F'$ with $h_0^{-1}\alpha(h_0)=f$, and put
$C=C_F(\alpha)$.
All solutions of the first equation of (D.35a)
are exactly $h=ch_0$ with $c\in C$:
$\alpha(h)=hf$ and $\alpha(h_0)=h_0f$ give
$\alpha(hh_0^{-1})=hh_0^{-1}$, and the reverse
substitution proves the converse.
Since $\alpha$ acts by $-1$ on $F^{\rm ab}$,
every $c\in C$ has $2[c]=0$, hence $[c]=0$.
Thus $C\subset F'$ and this substitution respects
the stated domain of (D.35a).  The remaining exact
equation is
\[
 x^\mu c h_0 y^\mu h_0^{-1}c^{-1}
                         =u(xy)^\mu u^{-1}.
                                                                  \tag{D.35g}
\]

We prove (D.35d) by identifying the relevant centralizer,
rather than assuming that a conjugator in $P$ lies in
$GT_0$.  Let $P^1=\ker(\lambda:P\to\widehat\Z^\times)$.
First,
$C_{P^1}(\delta)=1$.
Indeed a normalized element $(1,k)\in P^1$ commuting
with $\delta$ has $k=\delta(k)$, by uniqueness of the
normalized representative and its group law.
The following explicit mutually inverse homomorphisms
identify the semidirect product:
\[
 \begin{aligned}
 C_2\amalg C_2\amalg C_2
    &\xrightarrow{\ \sim\ }F\rtimes_\delta\langle s\rangle,
       & (s_0,s_1,s_2)&\longmapsto(s,xs,ys),\\
 (x,y,s)&\longmapsto(s_1s_0,s_2s_0,s_0).
 \end{aligned}
\]
The orders and conjugation relations follow by direct
substitution.  A free profinite factor has trivial
intersection with a conjugate by an element outside
that factor; the original full-profinite statement is
recalled in \cite[Theorem 5 of the inspected author
version]{A-HerfortZalesskii2008}.
Consequently the centralizer of $s_0$ in this free
product is $\langle s_0\rangle$.
Its intersection with $F$ is trivial.
Hence $C_F(\delta)=1$, so $k=1$ and
$C_{P^1}(\delta)=1$ as claimed.

Let $\Sigma\cong S_3$ be the geometric cusp-permutation
subgroup of the full loop-outer group, generated by
the outer classes
$\theta(x)=y,\ \theta(y)=x$ and
$\omega(x)=y,\ \omega(y)=(xy)^{-1}$.
The exact identification
$GT_0=C_P(\Sigma)$ follows from
\cite[Propositions 3--4 and Theorem 2]{R-HS}:
commutation with $\theta$ is relation (I), and,
after (I), commutation with $\omega$ is relation (II).
Both $\alpha$ and $\delta$ centralize $\Sigma$.
For any $\beta\in\operatorname{Conj}_P(\delta,\alpha)$
and $\sigma\in\Sigma$, form the actual element of
the outer group
$r_\sigma=\beta^{-1}\sigma\beta\sigma^{-1}$.
The pure subgroup $P$ is normal in the full
loop-outer group, and the common multiplier is a
homomorphism there.  Therefore $r_\sigma\in P^1$.
Moreover
\[
 \begin{split}
 r_\sigma\delta r_\sigma^{-1}
  &=\beta^{-1}\sigma\beta
        (\sigma^{-1}\delta\sigma)
        \beta^{-1}\sigma^{-1}\beta\\
  &=\beta^{-1}\sigma\alpha\sigma^{-1}\beta
    =\beta^{-1}\alpha\beta=\delta.
 \end{split}
\]
The proved equality $C_{P^1}(\delta)=1$ gives
$r_\sigma=1$.  Thus $\beta$ centralizes $\Sigma$,
so its normalized pair satisfies both (I) and (II)
and lies in $GT_0$.  This proves (D.35d).

Here is the exact finite form of the remaining
existence problem.  For each open characteristic
subgroup $N\triangleleft F$, put $Q_N=F/N$ and
$m_N=\exp(Q_N)$.
The automorphism $\alpha$ induces $\alpha_N$.
Use bars for the images of $x,y,f$ in $Q_N$ and define
\[
\begin{split}
 \mathcal S_{\alpha,N}=\{(a,b,v)\in
       (\Z/m_N)^\times\times Q_N'\times Q_N:\quad&
       b^{-1}\alpha_N(b)=\bar f,\\
       &\bar x^a b\bar y^a b^{-1}
              =v(\bar x\bar y)^a v^{-1}\}.
\end{split}                                                       \tag{D.35h}
\]
These are finite sets, obtained entirely from the
specified finite quotient and its induced
automorphism.  Reduction sends
$\mathcal S_\alpha$ to every
$\mathcal S_{\alpha,N}$.
Conversely, a point of a nonempty
$\mathcal S_{\alpha,N}$ lifts to
$\widehat\Z^\times\times F'\times F$:
units lift by the Chinese remainder theorem,
and the image of $F'$ in $Q_N$ is exactly $Q_N'$.
The lifts satisfying the two equations modulo $N$
form a nonempty closed subset of this compact space.
Every finite collection of such conditions is
implied by the one at a common characteristic
refinement.  Their finite intersection property
therefore gives a point satisfying both equations
in $F$ itself.  Open characteristic subgroups are
cofinal for a finitely generated profinite group,
so equality at all these quotients is equality in
$F$.  This proves the asserted finite criterion.

The need for the full third-cusp equation can be
checked by a short executed finite certificate.
Put $d=[x,y]=xyx^{-1}y^{-1}$ and take the actual
commutator word
\[
                 w=y d^{-1}y^{-1}x^{-1}d^{-1}x\in F'.
\]
Consider the endomorphism
$\rho(x)=x$, $\rho(y)=wyw^{-1}$.
Its abelianization is the identity, so the Frattini
argument makes its induced endomorphism on every
maximal pro-$p$ quotient an automorphism.
In the finite quotient
\[
 F\twoheadrightarrow A_5,\qquad
             x\longmapsto(1\,2\,3),\quad
             y\longmapsto(3\,4\,5),
\]
the program
\path{code/part_a/pure_involution_conjugator.py}
enumerates the generated group of order $60$ and
computes
\[
 \bar w=(1\,4\,3\,2\,5),\qquad
 \overline{wyw^{-1}}=(1\,2\,3).
\]
The image of $\rho(F)$ in this quotient therefore
has order $3$, proving $\rho(F)\ne F$.
Moreover $\overline{xy}$ has order $5$, whereas
$\overline{xwyw^{-1}}$ has order $3$; the
third-cusp equation fails in this same quotient.
This word is a certificate concerning the
generation test.  It is not asserted to solve
the coboundary equation for a nonstandard
involution.  No nonemptiness or emptiness
assertion for the spaces (D.35a) of arbitrary
involutions is inferred from this finite example.
\end{proof}

\begin{claim}{D.36}{Proved}{The global multiplier of a positive transposition}
Let $B$ be the full three-cusp loop-outer group, and let
$g\in B$ be an involution whose $2$-adic multiplier is
$\lambda_2(g)=1$.  Its puncture permutation is a
transposition, and its common multiplier is the global
element $1\in\widehat\Z^\times$.

Consequently a finite subgroup of $B$ whose map (D.32a)
has image $S_3\times C_2$ has the usual global orientation
character.  The same is true when that image is a
$V_4$ containing the pure involution, or is the positively
oriented $S_3\times\{1\}$.  These statements identify
the actual multiplier homomorphism; they do not assert
conjugacy of those finite actions.
\end{claim}
\passage{\S{}30; folios194--196; scans383--387; \S{}31; folios198--199verso; scans391--394; Part D(b) of the task}

\begin{proof}
We first record the genus-zero extension of the actual
automorphism calculation already proved in D.32--D.33.
Let $F_n=\widehat\pi_{0,n}$ with $n\geq3$, and let
$\beta$ be a pure automorphism with $\beta^p=1$ and
$p$-adic peripheral multiplier one.  Then $\beta=1$.
Indeed $F_n$ is free of rank $r=n-1$, and its cusp
classes span its abelianization.  Its induced pro-$p$
automorphism has identity abelianization and finite
order, so the $\Z_p$ Magnus argument of D.33 makes
it the identity.  The ring calculation (D.33c)--(D.33g)
gives one conjugacy class of sections of
$F_n\rtimes_\beta C_p\to C_p$, and its actual fixed
subgroup $C$ has
\[
                       \dim H^1(C,\F_p)=n-1.               \tag{D.36a}
\]
The norm calculation of D.28 applies to every cusp:
its image is a primitive element of the free pro-$p$
group, the first coordinate of its normalizer lift
centralizes that element, and the lift's $p$th power
has $p$-coordinate divisible by $p$.  Coordinatewise
division of the norm therefore supplies a section in
each peripheral normalizer.  Uniqueness of the
section class places conjugates of all $n$ full
peripheral $\Z_p$ subgroups inside the same $C$.

The finite-cover rank proof of D.32 now applies with
all $n$ cusps.  If $C$ is open, its first-cohomology
dimension is $1+[F_n:C](n-2)$, which equals $n-1$
only at index one.  If it has infinite index and
finite $p$-part of that index, the same exact
restriction colimit and eventual prime-to-$p$
injections give unbounded first-cohomology dimension.
Otherwise choose an open $U\geq C$ of $p$-divisible
index.  The $n$ selected boundary lifts all have
degrees prime to $p$ and lie over different original
boundaries.  They cannot be all the boundaries:
a single lift over any original boundary would
have degree $[F_n:U]$.  With at least $n+1$ boundaries,
the $n$ selected boundary homology classes are
independent.  Their dual characters restrict to
$C$ with the nonzero multipliers $d_i^{-1}\bmod p$
computed in (D.32d).  Thus $\dim H^1(C,\F_p)\geq n$,
contradicting (D.36a).  Hence $C=F_n$ and $\beta=1$.
This proof concerns an actual finite-order automorphism;
it does not presume a lift for arbitrary finite outer
actions in larger rank.

Return to $g$.  D.32 excludes a nonidentity pure
involution with $\lambda_2=1$, so its permutation
must be a transposition.  Relabel the original cusp
generators as $x,y,z=(xy)^{-1}$ so that $x$ is fixed
and $y,z$ are exchanged.  D.27 provides an actual
order-two lift $\alpha$ of this outer action.
Its peripheral multiplier satisfies $\lambda^2=1$,
so all its prime coordinates are $1$ or $-1$.

Suppose $\lambda_q=-1$ for an odd prime $q$, and use
the exact cyclic quotient
\[
 \chi:F\longrightarrow C_q,\qquad
       \chi(x)=0,\quad\chi(y)=1,\quad\chi(z)=-1,
                   \qquad N=\ker\chi.                    \tag{D.36b}
\]
The puncture permutation and the multiplier give
$\chi\alpha=\chi$: on $y$ the exponent $-1$ at
$q$ is multiplied by $\chi(z)=-1$, and similarly
for $z$.  Thus $N$ is invariant and the induced
action on its cyclic deck quotient is the identity.
The degree-$q$ cover is a sphere with $q+2$ cusps.
There are $q$ degree-one cusps over $x$, and one
degree-$q$ cusp over each of $y,z$.

Here are the exact covering and filling generators.
Finite-coset Schreier rewriting with coset representatives
$1,y,\ldots,y^{q-1}$ gives the free basis of $N$
\[
 a_i=y^i x y^{-i}\quad(0\leq i<q),\qquad b=y^q.
\]
The remaining boundary generator is
\[
 d=z^q=b^{-1}a_{q-1}^{-1}\cdots a_0^{-1},
                     \qquad a_0\cdots a_{q-1}bd=1.        \tag{D.36c}
\]
For example the latter relation is the actual free-word
identity
$(xy)^q=a_0\cdots a_{q-1}y^q$.
Completing the discrete Schreier comparison identifies
this presented free profinite group with the specified
open subgroup $N$.

The lift $\alpha$ fixes the conjugacy class of every
$a_i$ inside $N$.  To see this without choosing an
unproved normalization, write
$\alpha(x)=h^{-1}x^\lambda h$ for some $h\in F$.
Since $\chi\alpha(y)=1$, its action on the indices
of these cusp classes is the translation
$i\mapsto i-\chi(h)$ in $C_q$.  Its square is the
identity, so $2\chi(h)=0$ and $\chi(h)=0$.
Explicitly
\[
 k_i=\alpha(y)^i h^{-1}y^{-i}\in N,\qquad
                \alpha(a_i)=k_i a_i^\lambda k_i^{-1}.     \tag{D.36d}
\]
The two unique boundary classes over $y,z$ are
exchanged, with multiplier $\lambda$.

Fill exactly these two cusps.  The map is the quotient
by the closed normal subgroup of $N$ generated by $b,d$:
\[
 N\longrightarrow
 K=N/\overline{\langle\!\langle b,d\rangle\!\rangle}_N
   =\widehat{\langle a_0,\ldots,a_{q-1}\mid
                             a_0\cdots a_{q-1}=1\rangle}
   \cong\widehat\pi_{0,q}.                               \tag{D.36e}
\]
Both presentation maps are explicit: kill $b$ and
send the remaining $a_i$ to the named peripheral
generators; the inverse identifies $a_{q-1}$ with
$(a_0\cdots a_{q-2})^{-1}$.  These substitutions
verify both universal properties and compositions.
Since $\alpha$ exchanges the normal generators'
peripheral classes, the closed normal subgroup
in (D.36e) is invariant.  Consequently $\alpha$
induces an actual automorphism $\bar\alpha$ of
$K$ with $\bar\alpha^2=1$.  By (D.36d) it is
pure, with the same multiplier $\lambda$.

The proved genus-zero calculation at $p=2$ applies:
$q\geq3$ and $\lambda_2=1$, so $\bar\alpha=1$.
But its $q$-adic multiplier on each surviving
nontrivial cusp is $-1$, which is not the
identity on the cusp's $\Z_q$ subgroup.  Equivalently,
it acts as $-I$ on the nonzero $q$-adic abelianization
of the genus-zero group $K$.  This contradiction
excludes every such $q$.  Together with
$\lambda_2=1$ it proves $\lambda=1$ globally.

For the finite-group consequences, use the faithful
map (D.32a).  In a $V_4$ with pure kernel, the
pure involution has multiplier $-1$ by D.34,
and the positive transposition has multiplier
$1$ by the present proof.  These generate it.
The positively oriented $S_3$ is generated by
its positive transpositions.  Finally
$S_3\times C_2$ is generated by those
transpositions and its pure involution.
Multiplicativity of the common multiplier proves
each asserted character identity.

The program \texttt{code/part\_c/positive\_transposition\_d36.py}
checks (D.36c) as an exact free-word identity,
the finite cyclic cusp permutations and their
involutive translations, and the two presentation
substitutions in (D.36e) for several odd primes.
The proof above supplies the all-prime assertion.
\end{proof}

\begin{claim}{D.37}{Proved}{Negative transpositions and the specified four-cusp deck action}
Every involution of the full three-cusp loop-outer group whose
puncture permutation is a transposition and whose multiplier has
$\lambda_2=-1$ gives a faithful $V_4$ action on a specified
four-cusp double cover.  One generator is the standard free positive
deck involution; another has double-transposition permutation and
multiplier $\lambda$, and its center-free extension is torsion-free.
Conversely every such $V_4$ action containing the specified standard
deck action descends to a negative transposition on the three-cusp
group.  The constructions are inverse as marked extensions.

For the double-transposition permutation of D.29, neither of the
commuting single transpositions preserves any separating two-pants
partition whose pants the double transposition exchanges.  Thus an
additional deck action in this construction necessarily moves the
separator.  This is an obstruction to preserving that decomposition,
and asserts no obstruction to an action moving it.
\end{claim}
\passage{Part D(b) and D(d) of the task; \S{}24, folio 147 verso, scan 293; \S{}30, folios 194--196, scans 383--387; \S{}31, folios 198--199 verso, scans 391--394}

\begin{proof}
We specify both the markings and all the extension maps.  Put
\[
 F=\widehat F(x,y),\qquad z=(xy)^{-1},\qquad
 \chi:F\longrightarrow C_2,\quad (\chi(x),\chi(y))=(0,1),
 \qquad N=\ker\chi.
                                                        \tag{D.37a}
\]
Schreier rewriting identifies $N$ with the full free profinite
group on
\[
 A_0=x,\qquad A_1=yxy^{-1},\qquad B=y^2,
 \qquad D=z^2=B^{-1}A_1^{-1}A_0^{-1}.
                                                        \tag{D.37b}
\]
Its four peripheral generators satisfy $A_0A_1BD=1$.
These are the actual degree-one lifts of $x$ and the degree-two
lifts of $y,z$; completion of the discrete index-two Schreier
isomorphism gives precisely the specified inclusion $N\subset F$.
Let $\delta\in\Out(N)$ be conjugation by $y$.  Its named
representative is
\[
 A_0\longmapsto A_1,\qquad
 A_1\longmapsto B A_0 B^{-1},\qquad B\longmapsto B,
 \qquad \delta^2=\operatorname{Inn}(B).
                                                        \tag{D.37c}
\]
Here the final equality concerns the named automorphism;
its outer class has order two, multiplier $1$, and permutation
$(A_0\ A_1)$.  Its center-free extension is the explicitly
identified group $F$, hence is torsion-free.  In the converse
below, ``the specified deck action'' means this action transported
by a given peripheral isomorphism.  It does not mean an arbitrary
positive outer involution with a torsion-free extension.

Let $g$ be a negative transposition, labeled to fix $x$ and
exchange $y,z$.  Its center-free extension $E$ contains $F$
as an open subgroup of index two.  The peripheral normalizer
surjects onto $\langle g\rangle$, and A.7 gives its kernel
$\overline{\langle x\rangle}=\widehat\Z$.
For a lift $v$ in this normalizer write
\[
 v x v^{-1}=x^\lambda,\qquad v^2=x^a,
 \qquad (\lambda-1)a=0,\qquad \lambda^2=1.
                                                        \tag{D.37d}
\]
At every prime where $\lambda_p=-1$ the last commutation
identity gives $a_p=0$.  At every prime where $\lambda_p=1$
the prime is odd, since $\lambda_2=-1$, so $2$ is invertible.
Choose $b_p=-a_p/2$ in these latter coordinates and $b_p=0$
in the former.  Then $s=x^b v$ has order two.  Thus
\[
 E=F\rtimes_\alpha\langle s\rangle,\qquad
 \alpha^2=1,\qquad \alpha(x)=x^\lambda.
                                                        \tag{D.37e}
\]
This constructs the required finite-order lift inside the
peripheral normalizer, without a further choice of inner
normalization.  The multiplier and the cusp permutation give
$\chi\alpha=\chi$.  Therefore $N$ is normal in $E$, and
$E/N=C_2\times C_2$, with basis the images of $y,s$.

There is exactly one $F$-conjugacy class of order-two subgroups
of $E$.  Indeed the action on $H^1(F,\F_2)$ is the two-dimensional
nontrivial unipotent module $J_2$, isomorphic to the regular
$\F_2[C_2]$ module.  Consequently
$H^i(C_2,H^1(F,\F_2))=0$ for $i>0$.
The Hochschild--Serre spectral sequence, with the split
inflation on its bottom row, gives
\[
                 \dim H^n(E,\F_2)=1\qquad(n\geq3).
                                                        \tag{D.37f}
\]
The full profinite Brown restriction theorem used in D.28
identifies this high-degree cohomology with the product over
the order-two conjugacy classes.  Each class contributes
$1+\dim H^1(C_F(T),\F_2)$, since its normalizer is
$C_F(T)\times C_2$; see
\cite[Theorem 2 and \S1.1]{A-KochloukovaZalesskii2007}.
The section in (D.37e) exists, so (D.37f) forces exactly one
class and zero first mod-two cohomology for its centralizer.
Conjugacy by $E$ agrees here with conjugacy by $F$: for any
section $T$, one has $E=FT$, and the $T$ factor centralizes
that subgroup itself.
Every finite-order element of $E$ injects into $E/F$,
because $F$ is torsion-free, and consequently has order two.
All nonidentity torsion therefore maps to the one line
$\langle\bar s\rangle$ in the abelian group $E/N$.

Conjugation gives a faithful action of $E/N$ on $N$ in
$\Out(N)$.  To check the kernel, first
$C_F(N)=1$: an element centralizing both $x$ and $y^2$
lies in $\langle x\rangle\cap\langle y\rangle=1$ by A.7.
For $e\in C_E(N)$ and $f\in F$, the commutator $[e,f]$
lies in $F\cap C_E(N)$ and hence is trivial.  Thus $e$
centralizes $F$; faithfulness of the original nonidentity
outer action and $Z(F)=1$ give $e=1$.

Write $a$ for the outer restriction of $\alpha$ to $N$.
It fixes the classes of $A_0,A_1$ and exchanges $B,D$.
For the second fixed class use the explicit conjugator
$\alpha(y)y^{-1}\in N$:
\[
 \alpha(A_1)
  =(\alpha(y)y^{-1}) A_1^\lambda
                         (\alpha(y)y^{-1})^{-1}.
                                                        \tag{D.37g}
\]
The unique lifts of the other two cusps are exchanged inside
$N$.  Thus $\delta$ and $a$ commute in $\Out(N)$, and
\[
 \tau=\delta a,\qquad
 \sigma(\tau)=(A_0\ A_1)(B\ D),\qquad
 \operatorname{or}(\tau)=\lambda.
                                                        \tag{D.37h}
\]
The preimage $E_\tau$ of the diagonal line
$\langle\bar y+\bar s\rangle$ in $E/N$ is the center-free
extension of this outer involution.  It is torsion-free:
all nonidentity torsion has image $\bar s$, and its
intersection with $N$ is already torsion-free.
For an explicit representative set
\[
 \beta=\operatorname{Inn}(y)\alpha|_N,\qquad
 c=y\alpha(y)\in N;\qquad
 \beta^2=\operatorname{Inn}(c),\quad \beta(c)=c.
                                                        \tag{D.37i}
\]
Then $E_\tau$ is exactly the twisted extension generated by
$N$ and $ys$, whose square is $c$; the multiplication rule
is the rule of (D.29.9) with these named $\beta,c$.

Conversely let a marked four-cusp group carry the specified
$\delta$ and an involution $\tau$ with the permutation and
multiplier in (D.37h), satisfying $[\delta,\tau]=1$.
The distinct permutations make
$\Gamma=\langle\delta,\tau\rangle\cong V_4$.
Take its center-free extension
\[
 K=\{(\varphi,\gamma)\in\Aut(N)\times\Gamma:
                         [\varphi]=\gamma\}.
                                                        \tag{D.37j}
\]
Here $N$ embeds by its inner automorphisms.  The preimage
$K_\delta$ of $\langle\delta\rangle$ is identified with
the exact group $F$ in (D.37a) by
\[
 f\longmapsto
       (\operatorname{Inn}(f)|_N,\delta^{\chi(f)}).
                                                        \tag{D.37k}
\]
Injectivity follows from $C_F(N)=1$; surjectivity follows
because an automorphism in either of these two outer
classes differs from $1$ or $\operatorname{Inn}(y)|_N$
by an inner automorphism from $N$.
The quotient $K/F=C_2$ therefore gives an outer action on
this very $F$.  It is faithful, since any element of $K$
centralizing $F$ centralizes its inner copy of $N$, and
an automorphism of the center-free $N$ with that property
is the identity.

We check descent of the peripheral powers, not only of the
four unordered procyclic subgroups.  The two classes
$A_0,A_1$ become the one $F$-class of $x$.
The other classes are $B=y^2$ and $D=z^2$.
For a representative of $\tau$, the induced automorphism
$\psi$ of $F$ satisfies
$\psi(y)^2=h z^{2\lambda}h^{-1}$ for some $h\in N$.
A.7 implies
$C_F(hz^{2\lambda}h^{-1})=h\langle z\rangle h^{-1}$,
so $\psi(y)=hz^\mu h^{-1}$ and $2\mu=2\lambda$.
Multiplication by two is injective on $\widehat\Z$;
hence $\mu=\lambda$.  The same proof applies to $z$,
and the $x$ exponent is already given by $A_0$.
Thus the descended action fixes $x$, exchanges $y,z$,
and has exactly the original $\lambda$, in particular
$\lambda_2=-1$.

For completeness its normalizer construction (D.37d)--(D.37e)
and (D.37f) also recover the torsion-free assertion upstairs,
without making that assertion an extra converse hypothesis.
An order-two lift in the normalizer of $\langle x\rangle$
fixes $A_0$, so its image in $\Gamma$ is $\delta\tau$,
not $\tau$.  All torsion has this same image by (D.37f).
The preimage of $\langle\tau\rangle$ is therefore torsion-free.
The conjugation maps (D.37j)--(D.37k) identify the recovered
extension with the original one on every element; conversely
restriction to $N$ recovers the given $\Gamma$ and both
named generators.  This proves the asserted inverse
correspondence of marked extensions.

We finish by making the obstacle for the D.29 gluing explicit.
Write its ordered cusps as $1,2,3,4$, so that
$\sigma(\tau)=(1\ 3)(2\ 4)$.  The unordered two-pants
partitions and their images are
\[
\begin{array}{c|c|c}
 \text{partition}&\tau\text{ on its two blocks}
                     &(1\ 3)\text{ or }(2\ 4)\\\hline
 12\mid34&\text{exchanges}&14\mid23\\
 13\mid24&\text{preserves each}&13\mid24\\
 14\mid23&\text{exchanges}&12\mid34.
\end{array}                                            \tag{D.37l}
\]
The only single transpositions commuting with
$(1\ 3)(2\ 4)$ are $(1\ 3)$ and $(2\ 4)$:
a commuting permutation must send its two two-cycles to
two-cycles, and a transposition can do this only by
exchanging the members of one of them.  Hence a deck
involution of the required permutation cannot act on the
same two-vertex pants graph as the inverse-pair gluing,
under any relabeling or choice of its exchanged partition.

One can nevertheless write the required geometric deck
map in the D.29 free basis.  For
$(u_1,u_2,u_3,u_4)=(a,a^{-1}l,b,b^{-1}l^{-1})$, the
change of four-cusp marking
\[
 A_0=a,\quad A_1=(a^{-1}l)b(a^{-1}l)^{-1},\quad
 B=a^{-1}l,\quad D=b^{-1}l^{-1}
                                                        \tag{D.37m}
\]
satisfies $A_0A_1BD=1$.  Its inverse substitutions are
$a=A_0$, $l=A_0B$, $b=B^{-1}A_1B$.
Transporting (D.37c) through these maps gives
\[
 \Delta(l)=a^{-1}lb,\qquad
 \Delta(a)=a^{-1}lbl^{-1}a,\qquad
 \Delta(b)=a,\qquad
 \Delta^2=\operatorname{Inn}(a^{-1}l).
                                                        \tag{D.37n}
\]
It has multiplier $1$ and permutation $(1\ 3)$.
The existence question for D.29 is the existence of a
deck action of this marked geometric kind commuting
outerly with its actual $\alpha$; (D.37l) alone does
not decide it.

In particular failure of commutation for one selected
deck representative cannot settle the question.  In the
geometric specialization $\lambda=-1$, $\epsilon=1$,
$\kappa=1$, $f=1$, D.29 gives
\[
 \alpha(l)=l,\quad\alpha(a)=b^{-1},\quad
                         \alpha(b)=la^{-1}l^{-1}.
\]
For (D.37n), the words
$\alpha\Delta(l)=bl^2a^{-1}l^{-1}$ and
$\Delta\alpha(l)=a^{-1}lb$ are not conjugate even
in the full profinite group: the finite quotient
$l\mapsto(1\ 2\ 3)$, $a,b\mapsto(2\ 3)$ in $S_3$
sends the first to $1$ and the second to a nonidentity
three-cycle.  These two automorphisms therefore do not
commute outerly.  Yet another geometric deck does commute:
\[
 \begin{aligned}
 \Delta_0(l)&=bla^{-1},\qquad \Delta_0(a)=b,\qquad
                  \Delta_0(b)=blal^{-1}b^{-1},\\
 \Delta_0^2&=\operatorname{Inn}(bl),\qquad
       \alpha\Delta_0=\operatorname{Inn}(b^{-1})\Delta_0\alpha.
 \end{aligned}
                                                        \tag{D.37o}
\]
Its origin is also explicit.  In the geometric extension
$E=\widehat F(a,t)$, put $l=t^2$ and
$b=ta^{-1}t^{-1}$.  The involution
$\eta(a)=a^{-1}$, $\eta(t)=a^{-1}t$ preserves the
index-two kernel.  Restricting
$\operatorname{Inn}(t)\eta$ to that kernel gives
exactly (D.37o).  The formulas prove its peripheral
permutation and multiplier as well as the two displayed
composition identities.  Thus this geometric
specialization supplies an explicit example of the
separator being moved.
To verify directly that $\Delta_0$ is a deck of the
specified geometric kind, adjoin its lift $r$ with
$r^2=bl$.  The extension presentation eliminates to
$\widehat F(a,r)$ by
$b=rar^{-1}$ and $l=b^{-1}r^2=ra^{-1}r$.
The peripheral generators downstairs are
$a,a^{-1}r,r^{-1}$, whose product is one.
Their lifts are the two classes of $a,b$, the square
$(a^{-1}r)^2=a^{-1}l$, and the square $r^{-2}$,
conjugate to $b^{-1}l^{-1}$.  These substitutions identify
its four-cusp inclusion with (D.37a)--(D.37b).

The executable \texttt{code/part\_c/negative\_transposition\_d37.py}
verifies the free-word covering and inverse substitutions,
the deck squares, all the equations in (D.37o), the
two-dimensional mod-two cohomology calculation, and the
complete permutation calculation (D.37l) under all
relabelings.  The extension and torsion arguments above
prove the statements for the full profinite groups.
\end{proof}

\begin{claim}{D.38}{Proved}{A nonnilpotent finite obstruction on the universal \(A_5\) quotient}
Let $R$ be the intersection of the kernels of all continuous
epimorphisms $\widehat F(x,y)\twoheadrightarrow A_5$.
This is an open characteristic subgroup, and the quotient
specified by the nineteen generating-pair representatives in
the accompanying program is
\[
                 Q=F/R\xrightarrow{\ \sim\ } A_5^{19}.
                                                               \tag{D.38a}
\]
Write $X,Y,Z=(XY)^{-1}$ for the images of the three cusps,
and $\theta,\omega$ for the automorphisms of $Q$ induced by
$x\leftrightarrow y$ and $x\mapsto y,\ y\mapsto(xy)^{-1}$.
The image of every full $GT_0$ involution belongs to the
following explicitly computed finite set:
\[
\begin{split}
 \mathcal T=\{t\in Q:\quad&
 t^2=1,\quad tXt^{-1}=X^{-1},\\
 &\omega(t)=Y^{-1}\theta(t)\}.
\end{split}                                                     \tag{D.38b}
\]
Its normalized coordinates are
$\alpha_Q=\operatorname{Inn}(t)$ and
$f_Q=\theta(t)t^{-1}$.
Every tuple in $\mathcal T$ satisfies the finite
normalized relations (I), (II), and
$f_Q\alpha_Q(f_Q)=1$, and defines an actual automorphism
of $Q$ of order two.

The executed exact enumeration gives
\[
\begin{array}{l|r}
\text{finite set or test}&\text{cardinality}\\ \hline
\mathcal T&10\,125\,000\\
\text{tuples surviving the full cusp-triangle obstruction}
       &5\,062\,500\\
\text{tuples conjugate to standard inversion in the finite
pure model (D.38f)}&2\,531\,250\\
\text{triangle survivors outside that finite conjugacy set}
       &2\,531\,250 .
\end{array}                                                     \tag{D.38c}
\]
Thus $5\,062\,500$ explicit finite tuples satisfying the
finite involution equations have no full pure-involution
lift.  The last row is disjoint from the image of the
standard full conjugacy class.  No tuple in that last
row is asserted to have a full lift.
\end{claim}
\passage{\S{}31; folios 198-199verso; scans 391-394; Part D(b): actual full GT\_0 torsion and conjugacy}

\begin{proof}
The program enumerates all $60^2$ ordered pairs in $A_5$,
tests generation by closure under multiplication, and
obtains $2\,280$ generating pairs.
Simultaneous conjugation by all $120$ elements of $S_5$
has nineteen free orbits on this set.
These are all $\Aut(A_5)$-orbits: an automorphism
permutes the five Sylow-$2$ subgroups, and equivariance
with the faithful inner $A_5$ action makes that
permutation representation of $\Aut(A_5)$ injective;
conjugation by $S_5$ already supplies its $120$ elements,
as checked in D.31.
Two epimorphisms to $A_5$ have the same kernel exactly
when they differ by an automorphism of $A_5$.
The nineteen kernels are therefore distinct.
The joint quotient map is a subdirect product of
nineteen copies of the nonabelian simple group $A_5$.
The elementary subdirect-product argument shows that
a proper subdirect product of copies of one simple
group identifies two coordinates by an automorphism:
induct on the number of factors and apply the
normal-subgroup description of a direct product of
nonabelian simple groups to the last projection.
Such an identification would equate two of the
nineteen kernels.  This proves surjectivity of the
specified map to $A_5^{19}$ and hence (D.38a).
The kernel is characteristic because every automorphism
of $F$ permutes all epimorphisms to $A_5$.

For each representative pair $(X_i,Y_i)$ the program
also records the ordered triple of conjugacy classes
of $X_i,Y_i,Z_i$ in $A_5$, allowing the simultaneous
exchange of the two classes of elements of order five
by odd conjugation in $S_5$.
All nineteen resulting triples are different.
Every triple contains an element of order five,
and has at least two elements of odd order.
These assertions are checked by complete finite
enumeration, with representatives and component
permutations retained in the output.

By D.34 a full pure involution has multiplier $-1$.
Inversion preserves each conjugacy class of $A_5$:
this is immediate for involutions, and reversal of
a $3$-cycle or $5$-cycle is achieved by an even
permutation.
An automorphism of $A_5^{19}$ permutes its minimal
normal simple factors and acts on each factor by
an element of $\Aut(A_5)=S_5$.
Preservation of the three cusp classes therefore
forces its factor permutation to preserve the
recorded class triple.  The nineteen distinct
triples force the permutation to be the identity.
Since each triple contains an order-five entry,
the component automorphisms must have even parity;
odd conjugation exchanges its two order-five
classes.  Thus the induced automorphism of $Q$
is $\operatorname{Inn}(t)$ for a unique $t\in Q$.
The group $Q$ has trivial center, so its
involutivity gives $t^2=1$.

The exact Harbater--Schneps automorphism identities
from \cite[Propositions 3--4]{R-HS}, with
$m=(-1-1)/2=-1$, descend to $Q$:
\[
 \theta\alpha_Q\theta^{-1}\alpha_Q^{-1}
      =\operatorname{Inn}(f_Q),\qquad
 \omega\alpha_Q\omega^{-1}\alpha_Q^{-1}
      =\operatorname{Inn}(Y^{-1}f_Q).
\]
Center-freeness yields
$f_Q=\theta(t)t^{-1}$ and
$\omega(t)t^{-1}=Y^{-1}f_Q$.
Together with the normalized equation on $X$ these
give precisely (D.38b).

Conversely take $t$ in (D.38b) and define
$f=\theta(t)t^{-1}$.
Then $\theta(t)$ inverts $Y$, so
\[
 tYt^{-1}=f^{-1}Y^{-1}f,\qquad
 \theta(f)f=1,\qquad f\operatorname{Inn}(t)(f)=1.
\]
The second relation in (D.38b) involving $\omega$
can be rewritten $f=Y\omega(t)t^{-1}$.
Apply $\omega$ and $\omega^2$ and use $\omega^3=1$:
\[
 \omega(f)=Z\omega^2(t)\omega(t)^{-1},\qquad
 \omega^2(f)=Xt\omega^2(t)^{-1}.
\]
Substitution now proves the complete finite
three-cycle relation by cancellation,
\[
       \omega^2(f)Z^{-1}\omega(f)Y^{-1}fX^{-1}=1.
                                                               \tag{D.38d}
\]
The tuple $X$ has components of odd order, so
$tXt^{-1}=X^{-1}$ excludes $t=1$ in $Q$.
Thus $\operatorname{Inn}(t)$ has exact order two.
This verifies the assertions for every member
of the finite set, independently of any lifting.

We next impose an obstruction from the full
group.  For each component the extension quotient
is explicitly
\[
 A_5\rtimes_{\operatorname{Inn}(t_i)}C_2
       \xrightarrow{\ \sim\ }A_5\times C_2,\qquad
                (g,s^j)\longmapsto(g t_i^j,s^j).
                                                               \tag{D.38e}
\]
The three vertices of the full cusp triangle in
D.34 are represented by $s$, $xs$, and $yf s$.
Their first coordinates under (D.38e) are
$t_i$, $X_i t_i$, and $Y_i f_i t_i$.
None can be the identity in $A_5$.
Indeed that would make the image of the
corresponding section central in $A_5\times C_2$.
This section class is incident to two different
cusp edges of the full triangle.  At either
edge choose its representative in the cusp
normalizer.  Its image is still central, because
centrality is unchanged by conjugation.
It must therefore both act trivially on the
cusp image and act by inversion there.
Both incident cusp images would have order
at most two.  No cusp image is trivial for
a generating $A_5$ pair.  Two involutions
generate a dihedral group, while any two
entries of the generating triple generate
$A_5$.  This is a contradiction.
Possible fusion of distinct section classes
in the finite quotient does not change this
argument.

For all components together, this obstruction
is exactly
$t_i\notin\{1,X_i\}$ for every $i$.
The third vertex condition follows from these:
$f_i t_i=\theta(t)_i$ and $\theta(X)=Y$,
so $Y_i\theta(t)_i=1$ is another occurrence
of $t_j=X_j$ under the component action of
$\theta$.

Here are the complete finite counts, with the
factorization used to execute them.
The relation $\omega(t)=Y^{-1}\theta(t)$
joins components in twelve independent
sets, using indices $0,\ldots,18$:
\[
 \{0,2\},\{1\},\{3,5\},\{4\},\{6,7\},\{8\},
 \{9,17\},\{10\},\{11,14\},\{12,18\},\{13,15\},\{16\}.
\]
Exhaustive propagation of each initial component
value through its specified $S_5$ conjugation
maps gives respectively
\[
                 3,3,3,3,4,2,5,5,5,5,5,5
\]
assignments.  The equations on a component set
involve no other set, so their Cartesian product
has exactly $10\,125\,000$ elements.
The triangle obstruction changes only the
count for $\{6,7\}$, from four to two.
Its surviving cardinality is therefore
$5\,062\,500$.
All forty-eight component assignments are
executed and checked, including the original
equations (D.38d) and the finite norm; the
Cartesian-product count is exact.

For precision, the finite conjugacy model in
(D.38c) consists of the following actual
automorphisms of $Q$ with a common unit exponent:
\[
\begin{split}
 \mathcal P_Q=\{(\mu,\gamma)\in
      (\Z/30)^\times\times\Aut(Q):\quad&
 \gamma(X)=X^\mu,\\
 &\gamma(Y)\sim_QY^\mu,\quad
 \gamma(Z)\sim_QZ^\mu\}.
\end{split}                                                     \tag{D.38f}
\]
Every normalized full pure automorphism
maps to this set, since $R$ is characteristic.
The recorded class triples again imply that
an element of (D.38f) acts on the $i$th
factor by conjugation by some $g_i\in S_5$,
without permuting factors.
Its parity is even for $\mu\bmod5\in\{1,4\}$
and odd for $\mu\bmod5\in\{2,3\}$.
The remaining condition is exactly
$g_iX_i g_i^{-1}=X_i^\mu$; with that parity,
the $Y_i$ and $Z_i$ class conditions follow.

The standard inversion induces
$\operatorname{Inn}(d)$ on $Q$, where the
program supplies each simultaneous inverter
$d_i$ of $(X_i,Y_i)$.
The conjugacy condition is the finite equation
$t_i=g_i d_i g_i^{-1}$ for every $i$.
The program tests all $120$ choices of $g_i$
for each of the eight units modulo thirty.
There are exactly two distinct compatible
patterns, each with component counts
\[
                  3,3,3,3,1,1,5,5,5,5,5,5.
\]
They are disjoint, so their total cardinality
is $2\cdot1\,265\,625=2\,531\,250$.
This proves the last two rows of (D.38c).
It proves an upper constraint on the image
of the known full conjugacy class, not
surjectivity from that class onto the
finite model.

The complete implementation is
\path{code/part_a/a5_involution_quotient.py};
its output retains the nineteen pairs,
the $\theta,\omega$ component maps, the
standard inversion tuple, every component
count, and an explicit rejected tuple.
No family of lifts through all further
characteristic quotients has been supplied
for a surviving tuple.
\end{proof}

\begin{claim}{D.39}{Proved}{A rank-two Haran involution admitting neither hyperbolic loop type}
Let $\epsilon,\kappa\in\widehat\Z$ have the coordinates
\[
 \epsilon_3=-1,\quad\epsilon_q=1\ (q\ne3),\qquad
 \kappa_2=1,\quad\kappa_q=0\ (q\ne2).
\]
There is an actual involution of the rank-two free
profinite group $H=\widehat F(a,b)$ given by
\[
 \alpha(a)=b^{-1}a^\epsilon b,\qquad
 \alpha(b)=b^{-1}a^\kappa.                                  \tag{D.39a}
\]
Its split extension has the explicit identification
\[
 H\rtimes_\alpha C_2\ \xrightarrow{\ \sim\ }\ N\amalg C_2,
 \qquad
 N=(\Z_3\rtimes\Z_2)\times\prod_{q\ne2,3}\Z_q,               \tag{D.39b}
\]
where parity in $\Z_2$ acts by inversion on $\Z_3$.
Its actual fixed subgroup $C_H(\alpha)$ is trivial.
The outer involution preserves no three-cusp
marking and no once-punctured-torus marking of $H$,
under any full profinite change of free basis.

For the natural once-punctured-torus marking
$\overline{\langle[a,b]\rangle}$, the peripheral
condition also fails: a quotient of order $21$
separates $\alpha([a,b])$ from every conjugate
of $[a,b]^{-\epsilon}$.
\end{claim}
\passage{\S{}31; folios 198-199verso; scans 391-394; Part D: full rank-two loop actions and precise peripheral conditions}

\begin{proof}
The relations $\epsilon^2=1$ and
$\epsilon\kappa=\kappa$ give, by substitution
on the specified free generators,
\[
 \begin{split}
 \alpha^2(a)
  &=(a^{-\kappa}b)(b^{-1}ab)(b^{-1}a^\kappa)=a,\\
 \alpha^2(b)
  &=(a^{-\kappa}b)(b^{-1}a^{\epsilon\kappa}b)=b.
 \end{split}
\]
Thus the universal continuous endomorphism
defined by (D.39a) is an automorphism with
square one.
On full abelianization its matrix is
\[
       A=\begin{pmatrix}\epsilon&\kappa\\0&-1\end{pmatrix},
       \qquad \det A=-\epsilon,\qquad
       \operatorname{tr}A=\epsilon-1.                       \tag{D.39c}
\]
It is nonidentity on abelianization, so both
the automorphism and its outer class have
exact order two.

For the exact identification (D.39b), use
the normal form for $N$ verified in D.26:
\[
 N=\overline{\langle l,t\rangle},\qquad
 tlt^{-1}=l^\epsilon,\qquad t^2=l^\kappa.
\]
The map $l^u t^i\mapsto
(u_3,2u_2+i,(u_q)_{q\ne2,3})$, $i\in\{0,1\}$,
is its displayed isomorphism with the group
in (D.39b).  The parity character sends
$l$ to zero and $t$ to one.
In $N\amalg\langle r\mid r^2=1\rangle$
extend this character by $r\mapsto1$.
The substitutions
\[
             a\longmapsto l,\qquad
             b\longmapsto tr,\qquad
             s\longmapsto r
\]
define a homomorphism from
$H\rtimes_\alpha\langle s\rangle$ to that
free profinite product.  To check its
relations, note that $r=(tr)^{-1}t$,
so
\[
 rlr=(tr)^{-1}l^\epsilon(tr),\qquad
 r(tr)r=(tr)^{-1}l^\kappa.
\]
In the reverse direction send
$l\mapsto a$, $t\mapsto bs$, and $r\mapsto s$.
Here
$(bs)^2=b\alpha(b)=a^\kappa$ and
$(bs)a(bs)^{-1}=b\alpha(a)b^{-1}=a^\epsilon$.
The normal-form multiplication of $N$
therefore gives its continuous homomorphism
to the semidirect product, and the free
product property extends it over $r$.
The two composites fix every displayed
generator, proving the isomorphism.
In particular the parity kernel is exactly
the free profinite group on the stated
generators $a,b$, and the extension has
the actual section $\langle r\rangle$.

The free factor $\langle r\rangle$ is
malnormal by the full-profinite
free-factor intersection property
\cite[Theorem 5 of the inspected author
version]{A-HerfortZalesskii2008}.
Its centralizer in $N\amalg C_2$ is
therefore precisely $\langle r\rangle$.
Intersecting with the parity kernel
proves $C_H(\alpha)=1$.

In fact every involution $u\in H\rtimes_\alpha C_2$
has $C_H(u)=1$.  To prove this with the specified
index-two kernel, reduce (D.39c) modulo two.
The resulting matrix is
$\left(\begin{smallmatrix}1&1\\0&1\end{smallmatrix}\right)$.
Thus its dual action on $H^1(H,\F_2)$ is the
regular two-dimensional $\F_2[C_2]$ module $J_2$.
Its positive-degree $C_2$ cohomology vanishes:
both differentials of the periodic resolution are
$1+\alpha$, and their image equals their kernel,
both of dimension one.
Since $H$ has cohomological dimension one and
the extension is split, the Hochschild--Serre
spectral sequence gives
\[
       \dim H^n(H\rtimes_\alpha C_2,\F_2)=1
                        \qquad(n\geq3).                 \tag{D.39g}
\]
The full profinite Brown restriction theorem
\cite[Theorem 2 and \S1.1]{A-KochloukovaZalesskii2007},
applied exactly as in (D.37f), expresses this
cohomology as the product over involution classes.
Every class contributes
$1+\dim H^1(C_H(u),\F_2)$, because its normalizer
is $C_H(u)\times C_2$.
The existing section $\langle s\rangle$ therefore
accounts for the unique conjugacy class.
Conjugacy by the extension agrees with conjugacy
by $H$: for any involution $u$ the decomposition
$E=H\langle u\rangle$ absorbs its second factor,
which centralizes $u$.
Every involution is consequently $H$-conjugate
to $s$, and the already proved equality
$C_H(s)=1$ gives
\[
                  C_H(u)=1
           \quad\text{for every involution }u\in E.     \tag{D.39h}
\]

There is a basis-independent obstruction
to a three-cusp structure.  An involution
on three cusps has either the identity
permutation or a transposition.
In the former case its homology is a
scalar matrix $\lambda I$ and has trace
$2\lambda$, whose $2$-coordinate is
nonzero.  This contradicts the zero
$2$-coordinate of the trace in (D.39c).
For a transposition the reduced
permutation module has trace zero, and
multiplication by a common unit keeps
that trace zero.  This contradicts the
$3$-coordinate $-2$ of (D.39c).
Trace is unchanged by conjugating the
homology matrix under a change of free
basis.  Hence neither peripheral
permutation can occur.

There is also a basis-independent obstruction
to a once-punctured-torus marking.
Take any full profinite free basis $a',b'$ of
$H$ and put $c'=[a',b']$,
$L=\overline{\langle c'\rangle}$.
Preservation of the conjugacy class of $L$
would give the multiplier
$\nu=\det A=-\epsilon$.
The exact comparison is the one in
(D.33h): in every finite class-two
Heisenberg quotient the commutator
exponent is the determinant of the
abelianization matrix, so equality for
all moduli identifies these elements
of $\widehat\Z^\times$.
Changing the free basis conjugates the
matrix and leaves its determinant unchanged.

The normalizer $N_E(L)$ surjects onto
$E/H=C_2$, and its kernel is $L$ by A.7
transported through this free-basis isomorphism.
For a lift $v$ of the nonidentity element,
write
\[
  vc'v^{-1}=(c')^\nu,\qquad
  v^2=(c')^d,
  \qquad(\nu-1)d=0.                                    \tag{D.39i}
\]
Here $\nu_3=1$ and $\nu_q=-1$ for every
$q\ne3$.  Therefore $d_q=0$ for $q\ne3$,
including $q=2$, since multiplication by two
is injective on $\Z_q$.
Choose $j\in\widehat\Z$ with
$j_3=-d_3/2$ and $j_q=0$ otherwise.
The lift $u=(c')^jv$ then satisfies
\[
 u^2=(c')^{(1+\nu)j+d}=1,
 \qquad
 u(c')^{e_3w}u^{-1}=(c')^{e_3w}
                       \quad(w\in\Z_3),                 \tag{D.39j}
\]
where $e_3$ is the idempotent supported at
the prime three.
The lift $u$ has exact order two because
its image in $E/H$ is nonidentity.
The peripheral injection
$\widehat\Z\xrightarrow{\sim}L$ makes
$(c')^{e_3\Z_3}$ a nontrivial subgroup of
$C_H(u)$, contradicting (D.39h).
This excludes every once-punctured-torus
marking, including those obtained from
arbitrary full profinite free bases.

For an additional finite certificate in
the original basis, put
$c=[a,b]=aba^{-1}b^{-1}$.

Use the finite group
\[
 Q=\F_7\rtimes C_3,\qquad
 (v,i)(w,j)=(v+2^i w,\ i+j),
                                                               \tag{D.39d}
\]
where the second coordinate is modulo
three and $2$ has order three in
$\F_7^\times$.
The assignment
\[
                    a\mapsto(0,1),\qquad b\mapsto(1,1)
                                                               \tag{D.39e}
\]
is an epimorphism: $a,b$ have order
three, their commutator is $(1,0)$,
and these elements generate the
normal $\F_7$ and the $C_3$ quotient.
Since $\epsilon\equiv-1\pmod3$
and $\kappa\equiv0\pmod3$, (D.39a)
has the following exact images:
\[
 \begin{array}{c|ccccc}
 \text{word}&c&\alpha(a)&\alpha(b)&\alpha(c)&c^{-\epsilon}\\ \hline
 \text{image}&(1,0)&(5,2)&(3,2)&(1,0)&(6,0).
 \end{array}                                                   \tag{D.39f}
\]
The last entry uses
$(-\epsilon)_7=-1$.
Conjugacy on the normal $\F_7$ in
(D.39d) is multiplication by a
power of two.  The two relevant
orbits are exactly
\[
             \{(1,0),(2,0),(4,0)\},\qquad
             \{(3,0),(5,0),(6,0)\}.
\]
They are disjoint.  Thus the images
of $\alpha(c)$ and $c^{-\epsilon}$
are not conjugate, proving failure
of the natural torus cusp.

The program
\path{code/part_a/haran_torus_test.py}
executes this order-$21$ quotient
and an independent matrix check
inside $\operatorname{PSL}_2(\F_7)$.
For the latter it enumerates all
$168$ elements and uses
\[
 a=\begin{pmatrix}0&1\\6&1\end{pmatrix},
 \qquad b=\begin{pmatrix}0&2\\3&1\end{pmatrix}
 \quad\bmod\{\pm I\}.
\]
Both have order three,
$c=\left(\begin{smallmatrix}1&1\\3&4\end{smallmatrix}\right)$
has order seven, and
$\alpha(c)=
\left(\begin{smallmatrix}0&3\\2&2\end{smallmatrix}\right)$.
There are seven conjugators from
$c$ to this image and none from
$c^{-1}$.  The matrix calculation
agrees with (D.39f).
The finite calculation detects the original
cusp explicitly.  The argument
(D.39g)--(D.39j) supplies the exclusion
for every alternative marking.
The program also verifies the regular
mod-two module used in (D.39g).
\end{proof}

\begin{claim}{D.40}{Proved}{A finite homomorphism-count obstruction to the extra deck symmetry}
The mixed-multiplier four-cusp involution constructed in D.29
commutes with no single-transposition involution of multiplier $1$ in its
full loop-outer group. In particular it admits none of the
specified geometric deck actions required for the descent of D.37,
including actions moving the two-pants separator.

More precisely, for the exact profinite group
\[
 E=\overline{\langle a,l,t\mid
       tlt^{-1}=l^\epsilon,\ t^2=l^\kappa\rangle},
 \qquad p=a^{-1}l,
                                                        \tag{D.40a}
\]
with $\epsilon,\kappa$ as in D.29 and $\lambda=-\epsilon$,
there is no continuous automorphism sending the conjugacy
classes of $a,p$ to those of $a^\lambda,p$, respectively.
There is also none sending those classes to $a,p^\lambda$.
The obstruction is the exact inequality
\[
 \begin{aligned}
 \#\{\phi:E\to Q:\phi(a),\phi(p)\in C\}&=10200,\\
 \#\{\phi:E\to Q:\phi(a)\in C^{-1},\phi(p)\in C\}&=4920,
 \end{aligned}
                                                        \tag{D.40b}
\]
where all homomorphisms are continuous,
$Q=\operatorname{PSL}_2(\F_{11})$, and $C$ is the
conjugacy class of
$\left(\begin{smallmatrix}0&1\\10&2\end{smallmatrix}\right)$
modulo scalar matrices $\{\pm I\}$.
\end{claim}
\passage{Section 24; folio 147 verso; scan 293. Section 30; folios 194-196; scans 383-387. Section 31; folios 198-199 verso; scans 391-394. Part D(b) and D(d)}

\begin{proof}
We first identify the group and its peripheral marking.
In the basis $(l,a,b')$ of (D.29.10), adjoining the outer
lift $t$ gives the relations
\[
 tlt^{-1}=l^\epsilon,\quad tat^{-1}=b',\quad
 tb't^{-1}=l^\kappa a l^{-\kappa},\quad t^2=l^\kappa.
\]
Eliminating $b'$ gives (D.40a): the third relation follows
from the second and fourth. Conversely $b'=tat^{-1}$
defines the inverse homomorphism from (D.40a) to this
presentation. Both composites fix the named generators.
Thus this is an isomorphism of profinite groups, and
\[
                         E=N_{\rm bad}\amalg\widehat\Z(a)
                                                        \tag{D.40c}
\]
with the first factor generated by $l,t$.
The index-two subgroup $H$ is the kernel of the character
sending $t$ to $1$ and $a,l$ to $0$ in $C_2$.

The four $H$-peripheral classes of D.29 are represented by
\[
 a,\quad p=a^{-1}l,\quad b,\quad b^{-1}l^{-1}.
\]
Since $\alpha$ interchanges the first and third classes,
and the second and fourth, with multiplier $\lambda$,
and conjugation by $t$ is $\alpha$, these are equally
the classes of
\[
 a,\quad p,\quad (tat^{-1})^\lambda,
                         \quad(tpt^{-1})^\lambda.       \tag{D.40d}
\]
Here $\lambda^2=1$, and the original cusp conjugators
belong to $H$; thus (D.40d) records actual conjugacy
classes under the specified inclusion into $E$.
This description is independent of the Galois conjugator
$f$ used to construct the two pants isomorphisms.

For a finite group $Q$ define the central function
\[
 F_Q(L)=\#\{T\in Q:T L T^{-1}=L^\epsilon,
                                      \ T^2=L^\kappa\}.
                                                        \tag{D.40e}
\]
Every assignment $a\mapsto A$, $p\mapsto P$ forces
$l\mapsto AP$. The presentation (D.40a) then gives
exactly $F_Q(AP)$ continuous homomorphisms to $Q$.
For two conjugacy classes $C,D$ put
\[
             K_Q(C,D)=\sum_{A\in C,\ P\in D}F_Q(AP).
                                                        \tag{D.40f}
\]
Conjugating $L$ conjugates all its possible $T$, so
$F_Q$ is central. Since $PA=A^{-1}(AP)A$, this also proves
$K_Q(C,D)=K_Q(D,C)$.

Precomposition by a continuous automorphism of $E$ is
a bijection on its continuous homomorphisms to $Q$.
An automorphism with the first pair of class images in
the claim would therefore give
$K_Q(C,D)=K_Q(C^\lambda,D)$ for every pair $C,D$.
The inverse map has the same power effect because
$\lambda^2=1$, so this is equality of the two full sets,
without a choice of conjugators in the target.
The second proposed class effect would similarly give
$K_Q(C,D)=K_Q(C,D^\lambda)$.

We compute these integers in one completely specified
finite group. Enumerate the determinant-one matrices
over $\F_{11}$, and identify each matrix with its
negative, using their lexicographically smaller
representative. Matrix multiplication followed by this
identification constructs all $660$ elements of $Q$.
The class $C$ has size $60$ and order $11$.
Its inverse class is distinct and has the same size.
Since $\lambda_{11}=-1$, we have $C^\lambda=C^{-1}$.
The following table gives the full product counts needed
in (D.40f). The representative of each nonidentity
column is $\left(\begin{smallmatrix}0&1\\10&r\end{smallmatrix}\right)$;
the last column is the identity. We write
$N_{C,D}(L)=\#\{(A,P)\in C\times D:AP\in[L]\}$.
\[
\begin{array}{c|rrrrrrrr}
 r&0&1&2&3&4&5&9&I\\\hline
 |L|&2&3&11&5&5&6&11&1\\
 \text{class size}&55&110&60&132&132&110&60&1\\
 F_Q(L)&0&6&1&1&1&0&1&56\\
 N_{C,C}(L)
   &0&1320&120&660&660&0&840&0\\
 N_{C^{-1},C}(L)
   &660&0&120&660&660&1320&120&60
\end{array}                                               \tag{D.40g}
\]
Multiplying either of the last two rows by the
$F_Q$ row and summing gives respectively
\[
 7920+120+660+660+840=10200,
 \qquad 120+660+660+120+3360=4920.
\]
These unequal integers prove the first nonexistence
assertion. Symmetry of $K_Q$ proves the second.

All entries of (D.40g) are verified by the finite
enumeration in \path{code/root/haran_pair_counts.py}.
It constructs the multiplication, inverse, power and
conjugacy tables directly from the displayed matrices,
computes the residues of $\epsilon,\kappa$ modulo each
element order, and tests both equations in (D.40e) for
every pair $(L,T)$. It then counts each product in the
two displayed class pairs. Two independent computations,
\path{code/part_a/haran_pair_count_audit.py} and
\path{code/part_c/deck_obstruction_psl11_independent.py},
instead generate the faithful twelve-point permutation
action by $z\mapsto z+1$ and $z\mapsto-1/z$ on
$\mathbf P^1(\F_{11})$. The first enumerates allowed
$(L,T)$ before setting $P=A^{-1}L$; the second also
enumerates the triples directly. Both reproduce (D.40b).
The attached outputs include the full integer data.

Finally suppose a positive single-transposition
involution $\delta$ commuted with the D.29 involution
$\tau$. Their distinct peripheral permutations make
$\langle\delta,\tau\rangle\simeq V_4$.
Its center-free pullback extension inside $\Aut(H)$
contains $E$ as the normal preimage of $\langle\tau\rangle$,
using the canonical conjugation identification of
center-free extensions. Conjugation by any lift of
$\delta$ therefore defines a continuous automorphism
$\eta$ of this exact $E$.
By D.37's permutation computation, $\delta$ is either
$(1\ 3)$ or $(2\ 4)$ in the labeling (D.40d).
Its positive multiplier is $1$.
In the first case $\eta$ sends $a$ to an $E$-conjugate
of $a^\lambda$ and $p$ to an $E$-conjugate of $p$;
in the second case the roles are reversed.
Both have just been excluded. The argument imposes
no preservation of the original pants separator and
no order condition on the chosen lift of $\delta$.
It proves the stated obstruction for this actual
four-cusp construction; the general negative
three-cusp conjugacy question remains in D.12.
\end{proof}

\begin{claim}{D.41}{Proved}{The global multiplier of a negative transposition}
An involution of the full three-cusp loop-outer group whose
puncture permutation is a transposition and whose $2$-adic
multiplier is $-1$ has the global multiplier
$-1\in\widehat\Z^\times$.
Together with D.34 and D.36, this identifies the multiplier
of every involution in that full group with its global
$2$-adic sign.  Consequently every finite subgroup generated
by involutions has the usual global orientation character.
\end{claim}
\passage{Part D(b) of the task; \S{}24, folio 147 verso, scan 293; \S{}30, folios 194--196, scans 383--387; \S{}31, folios 198--199 verso, scans 391--394}

\begin{proof}
Label the original peripheral generators as
$x,y,z=(xy)^{-1}$ so that the involution fixes $x$ and
exchanges $y,z$.  Write its common multiplier as $\lambda$.
It satisfies $\lambda^2=1$, so every coordinate $\lambda_p$
is $1$ or $-1$.  The peripheral-normalizer construction
(D.37d)--(D.37e) supplies an actual order-two automorphism
$\alpha$ of $F=\widehat F(x,y)$ representing it and satisfying
$\alpha(x)=x^\lambda$.

Suppose a positive odd-prime coordinate exists, and let
$q$ be the least odd prime with $\lambda_q=1$.  Use the
exact cyclic quotient
\[
 \chi:F\longrightarrow C_q,\qquad
       \chi(x)=-2,\quad\chi(y)=\chi(z)=1,\qquad N=\ker\chi .
                                                        \tag{D.41a}
\]
The relation $xyz=1$ is respected.  Since $\lambda=1$
modulo $q$ and $\alpha$ exchanges the last two peripheral
classes, $\chi\alpha=\chi$.  The cover corresponding to
$N$ is the completion of the specified connected cyclic
cover of the pair of pants.  Each of its three original
cusps has one lift of degree $q$: all three displayed
monodromies are nonzero in $C_q$.  Its Euler characteristic
is $-q$, so its genus is
\[
                             g=(q-1)/2.
                                                        \tag{D.41b}
\]

Fill the two cusp lifts represented by $y^q,z^q$.
The quotient and its deck extension are explicit:
\[
\begin{split}
 R&=F/\overline{\langle\!\langle y^q,z^q\rangle\!\rangle}_F
       =\langle y,z\mid y^q=z^q=1\rangle_{\rm prof}
       =C_q\amalg C_q,\\
 K&=\ker(R\longrightarrow C_q,\ y,z\longmapsto1)\\
  &=N/\overline{\langle\!\langle y^q,z^q
                                      \rangle\!\rangle}_N
       \cong\widehat\pi_{g,1}.
\end{split}                                             \tag{D.41c}
\]
Here $x=z^{-1}y^{-1}$, and the one remaining cusp is
$c=x^q$.  To verify the equality of normal subgroups
used in (D.41c), note $F=N\langle y\rangle=N\langle z\rangle$.
Conjugating $y^q$ or $z^q$ by any element of $F$ therefore
gives an $N$-conjugate of the same element.  Thus its
closed $F$-normal closure is its closed $N$-normal closure.
The completion comparison is also exact: start with the
discrete cyclic surface cover, fill those two boundaries,
and complete the displayed quotient presentation.
Finite-index subgroup topology identifies its completed
kernel with the specified kernel in $C_q\amalg C_q$.
The inverse presentation substitutions are the original
$y,z$ and $x=z^{-1}y^{-1}$.

The automorphism $\alpha$ preserves this normal subgroup,
because it exchanges the two filled peripheral classes.
It therefore induces an actual automorphism of $R$,
preserves $K$, and induces an involution $\beta$ of $K$
with
\[
                           \beta(c)=c^\lambda.
                                                        \tag{D.41d}
\]
Let $T$ be conjugation by $y$ on $K$.  In $R$, $y^q=1$,
so $T^q=1$ as an actual automorphism.  The equality
$\chi\alpha=\chi$ gives $\alpha(y)y^{-1}\in K$; hence
$\beta T\beta^{-1}$ and $T$ differ by an inner
automorphism of $K$.  In particular their maps on
abelianization commute.

We compute that abelianization as a deck module, without
requiring a comparison of traces.  Reidemeister--Schreier
rewriting with representatives $1,y,\ldots,y^{q-1}$ gives
\[
 u_i=y^i z y^{-(i+1)}\quad(i\in\Z/q),\qquad
 u_0u_1\cdots u_{q-1}=1,\qquad T(u_i)=u_{i+1}.
                                                        \tag{D.41e}
\]
Exponents of $y$ here are interpreted in $R$, where
$y^q=1$.  The first $q-1$ elements form a free profinite
basis; eliminate the last by
$u_{q-1}=(u_0\cdots u_{q-2})^{-1}$.
These substitutions give both maps in the Schreier
presentation.  The remaining boundary also has the exact
word
\[
 c=x^q=\prod_{j=0}^{q-1}u_{\,q-1-2j}^{-1},
                       \qquad\text{indices modulo }q.
                                                        \tag{D.41f}
\]
Indeed $z=u_0y$ gives
$x=z^{-1}y^{-1}=y^{-1}u_0^{-1}y^{-1}
 =u_{q-1}^{-1}y^{-2}$.
Multiplying its $q$ successive translates by $y^{-2}$
gives exactly (D.41f), because $y^{2q}=1$.

For any prime $p$ put $V=H_1(K,\Q_p)$ and let $A$ be
the map induced by $\beta$.  From (D.41e) we obtain
the exact module isomorphism
\[
 V\cong
 \Q_p[C_q]\big/\Q_p(1+T+\cdots+T^{q-1})
 \cong\Q_p[X]/\Phi_q(X),\qquad
              A^2=1,\quad AT=TA.
                                                        \tag{D.41g}
\]
Every irreducible cyclotomic factor occurs with
multiplicity one.

The one-boundary surface structure in (D.41c) supplies
a nondegenerate alternating pairing, and its relation
to (D.41d) is algebraic.  In a surface free basis the
boundary word is the product of the $g$ commutators,
up to inversion and conjugation.  Its degree-two image
in the class-two free pro-$p$ quotient is a nondegenerate
bivector $\omega\in\bigwedge^2 V$.
Equation (D.41d) gives
$(\bigwedge^2 A)\omega=\lambda_p\omega$.
The deck map preserves its peripheral exponent, so
$(\bigwedge^2 T)\omega=\omega$.
If $\Omega$ is the nonsingular skew matrix of $\omega$,
these are precisely
\[
 A\Omega A^{\mathsf t}=\lambda_p\Omega,\qquad
 T\Omega T^{\mathsf t}=\Omega.
\]
Consequently the alternating form with matrix
$\Omega^{-1}$ satisfies
\[
 \langle Av,Aw\rangle=\lambda_p\langle v,w\rangle,
 \qquad
 \langle Tv,Tw\rangle=\langle v,w\rangle.
                                                        \tag{D.41h}
\]
This proves the multiplier comparison also at $p=2$:
one may first take the integral commutator calculation
in each finite class-two quotient and then tensor with
$\Q_2$.  It uses no geometric realization of $\beta$.

Now let $p\ne q$ be a prime with
$d=\operatorname{ord}_q(p)$ even.  The irreducible
factors of $\Phi_q$ over $\Q_p$ have root sets
the Frobenius orbits
\[
                  \{\zeta^{j p^k}:0\leq k<d\},
                     \qquad j\in(\Z/q)^\times.
                                                        \tag{D.41i}
\]
For completeness these roots lie in the unramified
extension with residue field $\F_{p^d}$: the primitive
$q$th roots in that field lift uniquely since $p\ne q$,
and Frobenius lifts their map $z\mapsto z^p$.
The least such degree is $d$, proving the factor
description.  Since $d$ is even, $p^{d/2}=-1$ modulo
$q$, and each of these root sets is closed under
inversion.

Write $V=\bigoplus V_i$ for those irreducible summands.
By (D.41g), the commutant of $T$ is the product of their
fields $\Q_p(\zeta^j)$.  Therefore $A|_{V_i}$ is
multiplication by an element of a field whose square
is one, and so is the scalar $1$ or $-1$ on that entire
summand.  Each $V_i$ is nondegenerate for the form in
(D.41h).  To see this explicitly after a splitting-field
extension, a $T$-eigenvector with eigenvalue $\zeta^j$
can pair only with the eigenvalue $\zeta^{-j}$, since
the pairing is $T$-invariant.  Both belong to $V_i$ by
(D.41i); different summands are orthogonal.  Global
nondegeneracy thus restricts to nondegeneracy on $V_i$.
Multiplication by either scalar $1$ or $-1$ preserves
that restricted form.  Comparing with (D.41h) gives
\[
 \lambda_p=1
       \quad\text{whenever }p\ne q
                    \text{ and }\operatorname{ord}_q(p)
                    \text{ is even}.
                                                        \tag{D.41j}
\]

There is a prime divisor $p$ of $q-1$ having this even
order.  Otherwise every prime factor of $q-1$ would
lie in the odd-order subgroup of the cyclic group
$(\Z/q)^\times$.  Their product, with multiplicities,
would also have odd order.  But that product is
$q-1=-1$ modulo $q$, of order two.  This contradiction
proves the assertion.  Such a prime satisfies $p<q$.
For $p=2$ the hypothesis gives $\lambda_p=-1$; for an
odd $p<q$ the minimal choice of $q$ gives the same
value.  Both contradict (D.41j).

Thus no positive odd-prime coordinate exists, and
$\lambda=-1$ globally.  D.34 handles a pure
involution, while D.36 handles an involution with
positive $2$-adic multiplier.  All involutions are
therefore accounted for.  Multiplicativity proves
the last assertion for every subgroup generated by
them.  This conclusion determines these actual
multiplier homomorphisms; it does not assert geometric
conjugacy of the finite actions or settle the
multiplier of an isolated order-three subgroup.

The program \texttt{code/part\_c/negative\_multiplier\_d41.py}
checks the exact Schreier boundary word, the induced
cyclotomic deck matrices and the nondegenerate
degree-two boundary form in examples, and checks
the prime-factor and Frobenius-orbit calculation.
The displayed argument proves the assertion for
every odd prime, rather than only for the executed
range.
\end{proof}

\begin{claim}{D.42}{Proved}{Order-three pants actions have global multiplier one}
Every element of order three in the full three-cusp
loop-outer group has orientation multiplier
$1\in\widehat\Z^\times$.
Consequently every finite subgroup of that full
group has multiplier image contained in the two
global units $\{1,-1\}$, and its multiplier is
determined by its $2$-adic coordinate.

The closed-surface calculation used in the proof is
also unconditional.  For $g\geq2$ and every prime $p$,
the homomorphism
\[
 \Aut(\widehat\pi_{g,0})
       \longrightarrow\GL(H_1(\widehat\pi_{g,0},\Z_p))
                                                        \tag{D.42a}
\]
is injective on every finite $p$-subgroup.
Here the source consists of actual continuous
automorphisms.  The order-three outer action in the
first assertion is supplied with an actual lift by
the exact rank-two lifting result D.27.
\end{claim}
\passage{\S{}30; folios 194--196; scans 383--387; \S{}31; folios 198--199 verso; scans 391--394; Part D: full loop-outer actions of type (0,3)}

\begin{proof}
We first prove (D.42a), keeping the two completion
maps and the fixed subgroup explicit.

\emph{The pro-$p$ automorphism calculation.}
Let $P=\pi_{g,0}^{(p)}$ be the pro-$p$ completion.
The closed lower central terms $\gamma_nP$
satisfy
\[
 \bigcap_n\gamma_nP=1,
 \qquad P/\gamma_nP\text{ is torsion-free for every }n.
                                                        \tag{D.42b}
\]
These are precisely
\cite[Propositions 2.14--2.15, pp.~1515--1516]{A-Morales2024},
applied to the orientable surface presentation.
The second proposition checks the primitive degree-two
initial form of its single defining relation; the
first passes the torsion-free nilpotent quotients to
the pro-$p$ completion by the stated canonical quotient
maps.  In particular every
$\gamma_nP/\gamma_{n+1}P$, as a subgroup of
$P/\gamma_{n+1}P$, is torsion-free.

Put
$K_n=\ker(\Aut(P)\to\Aut(P/\gamma_{n+1}P))$ for $n\geq1$.
For $b\in K_n$, the formula
\[
 d_n(b)(\bar z)=b(z)z^{-1}\pmod{\gamma_{n+2}P}
                                                        \tag{D.42c}
\]
defines a continuous homomorphism from $P^{\rm ab}$
to $\gamma_{n+1}P/\gamma_{n+2}P$.
Indeed this last quotient is central modulo
$\gamma_{n+2}P$, which makes the formula additive
in $z$ and trivial on commutators.
An automorphism acting trivially on $P^{\rm ab}$
acts trivially on every lower central graded quotient:
these quotients are generated by the iterated
commutators of its abelianization generators.
Consequently $d_n$ is a group homomorphism on $K_n$,
and its kernel is exactly $K_{n+1}$.
Its target is torsion-free by (D.42b).
A finite-order element of $K_1$ therefore lies
successively in every $K_n$ and is the identity.
Thus the kernel of
$\Aut(P)\to\GL(H_1(P,\Z_p))$ is torsion-free.

\emph{The full fixed subgroup.}
Write $\Pi=\widehat\pi_{g,0}$ and let
$a\in\Aut(\Pi)$ satisfy $a^p=1$ and act trivially
on $H_1(\Pi,\Z_p)$.
Its induced automorphism on $P=\Pi^{(p)}$ is the
identity by the preceding calculation.
There is therefore a specified quotient homomorphism
\[
 E=\Pi\rtimes_a C_p\longrightarrow P\times C_p,
             \qquad(h,s^i)\longmapsto(\bar h,s^i).
                                                        \tag{D.42d}
\]
The natural map $\Pi\to P$ induces an isomorphism
on mod-$p$ cohomology, including products.
Both groups have cohomological $p$-dimension two,
first cohomology $\F_p^{2g}$, and second cohomology
$\F_p$ with the standard nondegenerate surface
cup pairing.  The map is the identity on the
generator characters, so the pairing shows that
it is also an isomorphism in degree two; higher
degrees vanish.  These are the surface duality
comparisons of A.2 and A.7, with the full and
pro-$p$ formations respectively.
Hochschild--Serre comparison for the named map
(D.42d) now gives a ring isomorphism
\[
 H^*(P\times C_p,\F_p)\xrightarrow{\sim}H^*(E,\F_p).
                                                        \tag{D.42e}
\]
In particular $\dim H^n(E,\F_p)=2g+2$ for $n>2$.

Choose a nonzero surface class
$\eta\in H^2(P,\F_p)$ and a nonzero periodic
class $u\in H^2(C_p,\F_p)$.
For every $m\geq1$, the class
$\eta u^m\in H^{2m+2}(E,\F_p)$ is nonzero
under (D.42e).
Every subgroup $S\leq E$ of order $p$ projects
isomorphically to $C_p$, because $\Pi$ is
torsion-free.  Its first coordinate in $P\times C_p$
is trivial because $P$ is torsion-free.
Its normalizer is exactly
\[
               N_E(S)=C_\Pi(S)\times S.                 \tag{D.42f}
\]
An element normalizing $S$ fixes each of its elements
by projection to the abelian quotient $C_p$;
splitting off the uniquely determined $S$ component
proves (D.42f).

Apply the full profinite Brown restriction theorem
\cite[Theorem 2 and \S1.1]{A-KochloukovaZalesskii2007}
with virtual cohomological dimension two.
Its hypotheses hold: $\Pi$ is an open subgroup
of cohomological $p$-dimension two, and there is
no $C_p\times C_p$ in $E$, since every finite
subgroup injects into $E/\Pi=C_p$.
The finite-dimensional source makes its dense
restriction embedding an isomorphism onto a
finite product of the normalizer cohomologies.
Restriction of the nonzero class $\eta u^m$
is nonzero in at least one factor.
For that factor put $C=C_\Pi(S)$.
Because the first coordinate of $S$ in (D.42d)
is trivial, its restriction is exactly
\[
        (\eta|_C)\otimes(u^m|_S)
       \in H^2(C,\F_p)\otimes H^{2m}(S,\F_p).
                                                        \tag{D.42g}
\]
Thus $\eta|_C\ne0$.
The restriction
$H^1(\Pi,\F_p)\to H^1(C,\F_p)$ is injective:
for any nonzero $v$, nondegeneracy supplies $w$
with $v\smile w$ a nonzero multiple of $\eta$,
and its restriction cannot vanish.
Hence $\dim H^1(C,\F_p)\geq2g$ and
$\dim H^2(C,\F_p)\geq1$.
For every normalizer (D.42f), the K\"unneth formula
in degree $n>2$ contributes
\[
  1+\dim H^1(C_\Pi(S),\F_p)
                +\dim H^2(C_\Pi(S),\F_p).              \tag{D.42h}
\]
This holds also for odd $p$: the exterior generator
of $H^*(C_p,\F_p)$ and its degree-two polynomial
generator give one dimension in each nonnegative
degree.  At $p=2$ the degree-one polynomial
generator has the same dimension consequence.
The one factor in (D.42g) already accounts for
all $2g+2$ dimensions.  There is exactly one
conjugacy class of sections, and for its centralizer
\[
 \dim H^1(C,\F_p)=2g,\qquad H^2(C,\F_p)=\F_p.           \tag{D.42i}
\]
Conjugacy under $E$ equals conjugacy under $\Pi$,
since $E=\Pi S$ for each section $S$ and the
$S$ factor centralizes that section.
Consequently these conclusions apply to
$C=C_\Pi(a)$ for the specified section in (D.42d).

We now prove directly that this closed subgroup
equals $\Pi$.
Continuous cohomology has the restriction comparison
\[
 H^j(C,\F_p)=
  \mathop{\rm colim}_{C\leq U\leq_o\Pi}H^j(U,\F_p).
                                                        \tag{D.42j}
\]
For an explicit cochain comparison, choose an
open normal subgroup $V$ of $\Pi$ so small that
the given continuous finite-valued cochain on $C$
factors through $C/(C\cap V)$ in every variable.
The isomorphism $CV/V=C/(C\cap V)$ inflates it to
the open subgroup $CV$.  Cocycle equalities descend
to that finite quotient; for a coboundary choose
a common $V$ also for its primitive cochain.
This proves (D.42j) with the displayed restriction maps.
Each $U$ is the completion of its finite covering
surface, of genus
$g_U=1+[\Pi:U](g-1)$.
For $V\leq U$ the restriction on $H^2$ multiplies
the fundamental class by $[U:V]$.
A nonzero class surviving in (D.42j) in degree two
therefore has a representative at some $U$ for
which every smaller open $V\geq C$ has
$p\nmid[U:V]$.
For these subgroups the restriction on $H^1$
is injective: corestriction composed with restriction
is multiplication by this prime-to-$p$ index.
If $C$ had infinite index, the indices
$[\Pi:V]$ would be unbounded, making
\[
 \dim H^1(V,\F_p)=2+2[\Pi:V](g-1)
\]
unbounded inside the fixed finite-dimensional
$H^1(C,\F_p)$, contrary to (D.42i).
Thus $C$ is open.  The same genus formula and
(D.42i) give $[\Pi:C]=1$.
It follows that $a=1$.
The kernel of (D.42a) on a finite $p$-subgroup
cannot contain an element of order $p$, proving
its asserted injectivity.

\emph{The order-three cyclic cover.}
Let $g$ have order three in the full pants group,
with multiplier $\lambda$.
By D.32 its cusp permutation is a three-cycle.
Label $x,y,z=(xy)^{-1}$ so that it sends their
classes successively to $y,z,x$, with the common
power $\lambda$.
D.27 gives an actual order-three representative
$\alpha\in\Aut(\widehat F(x,y))$.
The equality $\lambda^3=1$ implies $\lambda_3=1$:
an order-three unit of $\Z_3$ reduces to one,
and for $v\in\Z_3^\times$, $r\geq1$, the
valuation of $(1+3^rv)^3-1$ is $r+1$.

Suppose a coordinate $\lambda_q$ were nonidentity.
Then $q\ne3$, and its reduction
$\zeta\in\F_q^\times$ has exact order three;
in particular $q\equiv1\pmod3$ and $q\geq7$.
Define the surjective character
\[
 \chi:F\longrightarrow C_q=\F_q,\qquad
 \chi(x)=\zeta,\quad\chi(y)=1,\quad\chi(z)=\zeta^2.
                                                        \tag{D.42k}
\]
The relation is respected because
$1+\zeta+\zeta^2=0$.
The three scalar identities
$\zeta\cdot1=\zeta$,
$\zeta\cdot\zeta^2=1$, and
$\zeta\cdot\zeta=\zeta^2$
show exactly that $\chi\alpha=\chi$;
all inner conjugators disappear under $\chi$.
Let $N=\ker\chi$.
The specified degree-$q$ cover has one cusp
above each of the original three, since all
three values in (D.42k) are nonzero.
Its type is $(h,3)$ with
$h=(q-1)/2\geq3$, by
$2-2h-3=-q$.
Fill these three cusps and put
\[
 \Pi_h=N/\overline{\langle\!\langle x^q,y^q,z^q
                         \rangle\!\rangle}_N
          =\widehat\pi_{h,0}.                            \tag{D.42l}
\]
Completion of the finite-cover presentation
gives this exact quotient map.
The peripheral powers under $\alpha$ preserve
the indicated closed normal subgroup, so it
induces an actual automorphism $b$ of $\Pi_h$
with $b^3=1$.
Conjugation by $y$ gives the deck automorphism
$T$ of this quotient; $T^q=1$, because $y^q$
has been killed.
The relation $\chi\alpha=\chi$ says that
$bT$ and $Tb$ differ by an inner automorphism.
Their actions on homology therefore commute.

Here is an explicit calculation of that deck
module.  First fill only $y^q,z^q$ upstairs.
The group downstairs is
$D=C_q(y)\amalg C_q(z)$, with its map to
$C_q$ given by $y\mapsto1$, $z\mapsto r=\zeta^2$.
Its kernel $K$ has Schreier generators
\[
 v_i=y^iz y^{-i-r}\quad(i\in\Z/q),\qquad
           v_0v_rv_{2r}\cdots v_{(q-1)r}=1.              \tag{D.42m}
\]
The displayed relation contains every generator
once and eliminates one; these are the completed
discrete Schreier presentation and maps.
Conjugation by $y$ sends $v_i$ to $v_{i+1}$.
The remaining cusp $x^q=(z^{-1}y^{-1})^q$
rewrites as
\[
             \prod_{j=0}^{q-1}
                    v_{-r-j(r+1)}^{-1}.                \tag{D.42n}
\]
Since $r+1\ne0$, it also uses every index once.
Its abelianization is the negative of the
already imposed sum relation.  Filling it changes
no first homology.  It follows, for the actual
map $K\to\Pi_h$, that
\[
 H_1(\Pi_h,\Q_3)
   =\Q_3[C_q]/\Q_3(1+T+\cdots+T^{q-1})
   \simeq\Q_3[X]/\Phi_q(X).                             \tag{D.42o}
\]
This is the regular augmentation module, with
each irreducible factor occurring once.

Because $q\ne3$, reduction of $\Phi_q$ modulo
three is square-free.  Hensel lifting its
irreducible factors expresses the last algebra
in (D.42o) as a product of finite unramified
extensions of $\Q_3$.
None contains a nonidentity cube root of one:
reduction to the residue field, whose multiplicative
group has order $3^d-1$, first gives residue one;
the same valuation calculation used for $\Z_3$
then applies, since three remains a uniformizer.
The centralizer of $T$ on its cyclic module
(D.42o) is exactly that product algebra.
Thus the commuting operator $b_*$, with
$b_*^3=1$, equals one in every factor.
The torsion-free lattice
$H_1(\Pi_h,\Z_3)$ embeds in (D.42o), so $b_*$
is also the identity on that lattice.
The first part of this proof, applied with $p=3$,
now gives $b=1$ on the full $\Pi_h$.

Finally the quotient filling comparison in A.2
identifies the relative orientation line of
$(N;\langle x^q\rangle,\langle y^q\rangle,
\langle z^q\rangle)$ with the closed orientation
line of $\Pi_h$.
The three cusp maps have the common power
$\lambda$, so $b$ acts on the latter line by
$\lambda$ as well.
Its being the identity forces
$\lambda_q=1$, contradicting the chosen coordinate.
There is no such $q$, proving $\lambda=1$.

For the finite-subgroup assertion, D.32 embeds
every finite subgroup in $S_3\times C_2$.
Its elements therefore have orders $1,2,3$, or $6$.
D.34, D.36, and D.41 give the global sign for
every involution; the preceding proof handles
order three.  An element $v$ of order six is
the product of its commuting powers $v^3$ and
$v^4$, of orders two and three respectively.
Multiplicativity gives the assertion for every
element, with its prescribed $2$-adic sign.

The program \path{code/part_a/pants_order_three.py}
executes the character invariance, both Schreier
index cycles, the cyclic deck matrix and its
cyclotomic polynomial, and the characteristic-three
factorization checks for both primitive cube roots
at every prime $q\equiv1\pmod3$ through $97$.
The arguments above prove the result for every
prime and for the full profinite groups.
\end{proof}

\begin{claim}{D.43}{Proved}{Global orientation of punctured-torus involutions}
Every involution in the full once-punctured-torus loop-outer
group has the global orientation multiplier $1$ or $-1$,
with the same sign as its $2$-adic determinant.
Consequently the orientation character of every finite
subgroup generated by involutions has values in the global
two-element subgroup $\{1,-1\}\subset\widehat\Z^\times$.
\end{claim}
\passage{Part D(b) of the task; \S{}24, folio 147 verso, scan 293; \S{}30, folios 194--196, scans 383--387; \S{}31, folios 198--199 verso, scans 391--394}

\begin{proof}
Write $F=\widehat F(x,y)$ and $c=[x,y]$ for the torus
cusp.  D.27 supplies an actual order-two representative
$\alpha$ of the given outer involution.  Its peripheral
multiplier is
\[
 \lambda=\det(\alpha|_{F^{\rm ab}})\in\widehat\Z^\times,
                    \qquad \lambda^2=1,
                                                        \tag{D.43a}
\]
by the exact Heisenberg comparison of D.33.
Thus every $\lambda_p$ is $1$ or $-1$.

\emph{The characteristic Heisenberg cover.}
For an odd prime $q$, take the characteristic quotient
\[
 F\longrightarrow G_q,\qquad
 G_q=\langle X,Y,Z\mid X^q=Y^q=Z^q=1,\
              [X,Y]=Z,\ [X,Z]=[Y,Z]=1\rangle,
                                                        \tag{D.43b}
\]
with $x\mapsto X$, $y\mapsto Y$, $c\mapsto Z$.
Collection of these words identifies $G_q$ with the
order-$q^3$ Heisenberg group.  Explicitly it is
$\F_q^2\times\F_q$, with multiplication
\[
 (v,z)(w,t)=(v+w,z+t+\tfrac12\omega(v,w)),
 \qquad \omega((a,b),(a',b'))=ab'-ba'.
                                                        \tag{D.43c}
\]
Take $X=((1,0),0)$, $Y=((0,1),0)$ and $Z=((0,0),1)$.
The kernel $N_q$ is the closed subgroup generated by
the laws $\gamma_3F$ and all $q$th powers, so it is
characteristic.  More precisely it is their closed
normal closure: the exponent-$q$, class-two quotient
has the normal forms $X^aY^bZ^d$ in (D.43b), and
(D.43c) verifies that all $q^3$ forms are distinct.
This proves the characteristic-quotient identification.

The corresponding discrete cover of a once-punctured
torus has $q^2$ cusps, one for each coset of
$\langle Z\rangle$ in $G_q$, each of degree $q$.
Fill all these cusps.  Completing the finite-cover
presentation gives the exact quotient
\[
 N_q\longrightarrow
 \Pi_q=N_q/
   \overline{\langle\!\langle h c^q h^{-1}:
                       h\in F\rangle\!\rangle}_{N_q}
       \cong\widehat\pi_{g_q,0},
 \qquad 2g_q=2+q^3-q^2 .
                                                        \tag{D.43d}
\]
Representatives for the $q^2$ cosets suffice in the
normal closure.  The genus formula follows by multiplying
the punctured torus's Euler characteristic by $q^3$
and then adding its $q^2$ filled discs.
The original deck group induces its actual homology
representation on this closed surface.
Since $\alpha$ preserves $N_q$ and permutes all the
filled peripheral classes, it induces an actual
involution $\beta_q$ of $\Pi_q$.

For every prime $p$, put $V_p=H_1(\Pi_q,\Q_p)$,
let $\rho$ be its $G_q$ representation, and let
$A$ be the map induced by $\beta_q$.  If $\theta$
is the automorphism induced by $\alpha$ on $G_q$,
the precise deck comparison is
\[
 A^2=1,\qquad
 A\rho(h)A^{-1}=\rho(\theta(h)),\qquad
                    \theta(Z)=Z^{\lambda\bmod q}.
                                                        \tag{D.43e}
\]
Indeed representatives in $F$ for the two sides differ
only by an element of $N_q$, whose conjugation is inner
on $\Pi_q$ and hence trivial on its first homology.

The nondegenerate alternating surface pairing satisfies
\[
 \langle Av,Aw\rangle=\lambda_p\langle v,w\rangle,
       \qquad
 \langle\rho(h)v,\rho(h)w\rangle=\langle v,w\rangle .
                                                        \tag{D.43f}
\]
Here is an exact verification of the first multiplier.
The involution permutes the odd number $q^2$ of lifted
cusps, so one cusp class is fixed.  Fill every other
cusp first.  This is an invariant closed normal
quotient, and gives a surface group of genus $g_q$
with one cusp, which $\alpha$ preserves with power
$\lambda$.  Filling its last cusp does not change
first homology.  The degree-two boundary bivector
calculation in (D.41h) therefore proves the first
identity in (D.43f) for precisely $A$ on $V_p$.
The second identity comes from the orientation-preserving
geometric deck maps of the specified finite surface cover.
Thus no realization of $\alpha$ as a homeomorphism
has been used.

We compute the whole deck module by finite cover chains.
The cover before filling retracts onto the regular
cover of a two-loop rose, so its first homology has
representation class $1+\Q_p[G_q]$.
The relative second homology of the attached discs
is $\Q_p[G_q/\langle Z\rangle]$.
The exact homology sequence of this filling, including
the invariant orientation class in degree two, gives
\[
 [V_p]=
       2[1]+[\Q_p[G_q]]
                 -[\Q_p[G_q/\langle Z\rangle]].
                                                        \tag{D.43g}
\]
All these maps are deck-equivariant.  They compute
the homology of the completed surface groups through
the finite-index and surface-homology comparisons
already used in A.2 and D.42.

Over a characteristic-zero field containing a primitive
$q$th root $\zeta$, let $V_j$ be the subspace where $Z$
acts as $\zeta^j$, for $j\in\F_q^\times$.
The last permutation module in (D.43g) is trivial on
$Z$, and the regular representation has each of its
$q$ central eigenvalues with multiplicity $q^2$.
Semisimplicity therefore gives
\[
                         \dim V_j=q^2 .
                                                        \tag{D.43h}
\]
The $G_q$-invariant alternating form pairs $V_j$
only with $V_{-j}$, and gives a nondegenerate
pairing between them.

\emph{One negative odd-prime sign determines every sign.}
Suppose $\lambda_q=-1$ at an odd prime $q$.
For any prime $p$, extend scalars to $\Q_p(\zeta)$.
By (D.43e), $A$ exchanges $V_j$ and $V_{-j}$.
Choose bases of these two spaces and write $P$ for
their nonsingular pairing matrix, $B$ for
$A:V_j\to V_{-j}$ and $C$ for $A:V_{-j}\to V_j$.
Then $CB=BC=1$, and alternation in (D.43f) gives
\[
                  -B^{\mathsf t}P^{\mathsf t}C
                                      =\lambda_p P.
                                                        \tag{D.43i}
\]
The dimension $q^2$ is odd.  Taking determinants
in (D.43i) yields $-1=\lambda_p^{q^2}$, hence
$\lambda_p=-1$ by (D.43a).
This works also at $p=q$ and $p=2$, since the
calculation is over characteristic-zero fields.
Thus
\[
       \lambda_q=-1\text{ for an odd }q
                  \quad\Longrightarrow\quad\lambda=-1.
                                                        \tag{D.43j}
\]

\emph{The positive $3$-adic case.}
It remains to prove that $\lambda_3=1$ forces
$\lambda_2=1$.  The matrix of $\alpha$ on
$H_1(F,\Z_3)$ has square one and determinant one.
Its two eigenvalues over $\Q_3$ are consequently
both $1$ or both $-1$, and it is diagonalizable,
so the matrix is $I$ or $-I$.

In the group $G_3$ of (D.43c), the resulting $\theta$
has the form
\[
                \theta(v,z)=(Mv,z+\ell(v)),
                      \qquad M=I\text{ or }-I,
                                                        \tag{D.43k}
\]
with $\ell$ a linear functional.
For $M=I$, $\theta^2=1$ implies $2\ell=0$ and
$\ell=0$.
For $M=-I$, conjugation of $\theta$ by the inner
automorphism of $(u,0)$ replaces $\ell(v)$ by
$\ell(v)-2\omega(u,v)$.  Choose $u$ with
$\omega(u,v)=\ell(v)/2$.  Lifting $(u,0)$ to $F$
and conjugating the actual automorphism $\alpha$
by that inner automorphism preserves its order two
and its peripheral multiplier.  It gives exactly
\[
             \theta=1\quad\hbox{or}\quad
             \theta=\iota,\qquad\iota(v,z)=(-v,z).
                                                        \tag{D.43l}
\]
This is conjugation of an actual automorphism,
not multiplication by an arbitrary inner lift.

Work now over the unramified quadratic field
$L=\Q_2(\zeta_3)$, with
$\sigma(\zeta_3)=\zeta_3^{-1}$.
For $j=1,2$ the irreducible $G_3$ module $W_j$
with central character $\zeta_3^j$ has basis
$e_r$, $r\in\F_3$, and matrices
\[
 X e_r=e_{r+1},\qquad
 Y e_r=\zeta_3^{-jr}e_r,\qquad
 Z e_r=\zeta_3^j e_r .
                                                        \tag{D.43m}
\]
These matrices satisfy (D.43b).
The distinct eigenvalues of $Y$ and the cyclic
action of $X$ show irreducibility.  The diagonal
projectors for $Y$, followed by powers of $X$,
span the full matrix algebra.  Hence the central
block of $L[G_3]$ is $M_3(L)$.
Together with (D.43g)--(D.43h), this gives
\[
                       V_j=W_j\otimes_L U_j,
                             \qquad\dim U_j=3.
                                                        \tag{D.43n}
\]

The matrix $J=I$ implements $\theta=1$ on $W_j$.
For $\theta=\iota$, take $J e_r=e_{-r}$.
This rational permutation matrix implements
$X\mapsto X^{-1}$, $Y\mapsto Y^{-1}$, fixes $Z$,
and has $J^2=I$ and $J^{\mathsf t}J=I$.
Thus (D.43e) and the full matrix-algebra commutant
give
\[
                A|_{V_j}=J\otimes B_j,\qquad B_j^2=I.
                                                        \tag{D.43o}
\]
Choose a multiplicity basis for $U_1$ and its
$\sigma$-image for $U_2$.
Since the original representation and $A$ are
defined over $\Q_2$, we then have
$B_2=\sigma(B_1)$.
The canonical invariant pairing between $W_1$
and $W_2$ is the ordinary dot product in (D.43m).
Consequently the nondegenerate cross-pairing in
(D.43f) has the form $I_3\otimes S$ for an
invertible $3$-by-$3$ matrix $S$ on the multiplicity
spaces.  Equation (D.43f) becomes
\[
                 B_1^{\mathsf t}S\,\sigma(B_1)
                                      =\lambda_2 S.
                                                        \tag{D.43p}
\]
Because $B_1^2=I$, its determinant is $1$ or $-1$,
both fixed by $\sigma$.
Taking determinants in (D.43p) therefore gives
$1=\lambda_2^3$.  Together with $\lambda_2^2=1$,
this proves $\lambda_2=1$.

If $\lambda_2=-1$, the last conclusion forces
$\lambda_3=-1$, so (D.43j) gives $\lambda=-1$.
If $\lambda_2=1$, (D.43j) excludes every negative
odd-prime coordinate, and $\lambda=1$.
This proves the claim.  Multiplicativity proves
the stated consequence for subgroups generated
by involutions; no conclusion about their
conjugacy is inferred.

The program \texttt{code/part\_c/torus\_involution\_d43.py}
checks the characteristic Heisenberg presentations,
the finite-cover chain ranks and central-character
dimensions, and every $G_3$ involution with determinant
one in (D.43k), including the exact inner-conjugacy
formula in (D.43l).  It also verifies the
Schr\"odinger matrices and their rational implementer.
The field and determinant arguments above prove
the assertion for the full profinite groups.
\end{proof}

\begin{claim}{D.44}{Proved}{A first-symmetry choice for the involution coboundary}
Let $\alpha=(-1,f)\in GT_0$ be an involution in its
normalized actual action on $F=\widehat F(x,y)$:
\[
 \alpha(x)=x^{-1},\qquad
 \alpha(y)=f^{-1}y^{-1}f,\qquad f\in F'.
\]
Write $\theta(x)=y$, $\theta(y)=x$.
There exists $h\in F'$ satisfying both exact equations
\[
               f=h^{-1}\alpha(h),\qquad h\theta(h)=1.
                                                        \tag{D.44a}
\]
Thus the coboundary in D.35 can be chosen with the
first Harbater--Schneps symmetry.  For every unit
$\mu\in\widehat\Z^\times$, the associated continuous
endomorphism
$\rho_{\mu,h}(x)=x^\mu$,
$\rho_{\mu,h}(y)=hy^\mu h^{-1}$ satisfies
$\alpha\rho_{\mu,h}=\rho_{\mu,h}\delta$, where
$\delta$ inverts $x,y$.
Its full peripheral-conjugator condition is precisely
D.35's remaining equation
$x^\mu hy^\mu h^{-1}\sim_F(xy)^\mu$.

At the characteristic finite quotient $A_5^{19}$
of D.38, this first-symmetry condition is soluble
for every one of the $5\,062\,500$ tuples that
pass the triangle test.  This finite equation
therefore gives no further exclusion among them.
\end{claim}
\passage{\S{}30; folios 194--196; scans 383--387; \S{}31; folios 198--199 verso; scans 391--394}

\begin{proof}
The normalized involution and the first
Harbater--Schneps equation give
\[
                  f\alpha(f)=1,\qquad\theta(f)f=1.
                                                        \tag{D.44b}
\]
Checking the two free generators gives the
identity of actual automorphisms
\[
                 \theta\alpha\theta
                        =\operatorname{Inn}(f)\alpha.
                                                        \tag{D.44c}
\]
Indeed both sides send $x$ to $fx^{-1}f^{-1}$
and $y$ to $y^{-1}$.

Form $E=F\rtimes_\alpha\langle s\mid s^2=1\rangle$.
Define an automorphism $\Theta$ of this exact
extension by
\[
                      \Theta|_F=\theta,\qquad
                      \Theta(s)=fs.                    \tag{D.44d}
\]
The first equation of (D.44b) preserves $s^2=1$;
(D.44c) preserves the conjugation relation;
and the second equation of (D.44b) gives
$\Theta^2(s)=\theta(f)fs=s$.
Thus $\Theta^2=1$ on all of $E$.
Adjoin its named involution $\tau$ and put
\[
 G=E\rtimes_\Theta\langle\tau\rangle,\qquad
                         G/F=C_2\times C_2.
                                                        \tag{D.44e}
\]
The two quotient coordinates are represented
by $s,\tau$.

The preimage of the $\tau$ line is the normal
subgroup
$G_\theta=F\rtimes_\theta\langle\tau\rangle$.
It has the explicit free-product identification
\[
 G_\theta\xrightarrow{\sim}
                  C_2(\tau)\amalg\widehat\Z(v),
       \qquad x\mapsto v,\quad y\mapsto\tau v\tau.
                                                        \tag{D.44f}
\]
The inverse sends $v$ to $x$ and $\tau$ to
the named involution; both composites fix
the indicated generators.
The free-factor intersection theorem
\cite[Theorem 5 of the inspected author version]{A-HerfortZalesskii2008}
therefore gives $C_F(\tau)=1$.

There is exactly one $F$-conjugacy class of
involutions in $G_\theta$.
For a direct verification, its mod-two
$H^1(F,\F_2)$ module is the regular
two-dimensional $C_2$ module, since $\theta$
exchanges its two basis characters.
Its positive-degree $C_2$ cohomology vanishes.
The split Hochschild--Serre calculation
therefore gives
$\dim H^n(G_\theta,\F_2)=1$ for $n\geq3$.
Brown restriction
\cite[Theorem 2 and \S1.1]{A-KochloukovaZalesskii2007}
gives one involution class, by exactly the
dimension argument of (D.37f).
Conjugacy by $G_\theta$ equals conjugacy by
$F$, since $G_\theta=F\langle u\rangle$
for every involution $u$ and its second
factor centralizes that element.

Consider the closed normalizer
$V=N_G(\langle\tau\rangle)$.
For every $g\in G$, the conjugate $g\tau g^{-1}$
belongs to $G_\theta$, and its image in
$G/F$ is still the $\tau$ coordinate.
The preceding uniqueness gives an $F$-conjugator
taking it back to $\tau$.
It follows that $V\to G/F$ is onto.
Its kernel is $C_F(\tau)=1$, so
\[
                   V\xrightarrow{\sim}C_2\times C_2.
                                                        \tag{D.44g}
\]
In particular $V$ is a finite subgroup,
contains $\tau$, and has a unique involution
$u$ mapping to the $s$ coordinate.
That involution belongs to $E$ and commutes
with $\tau$.

The three $F$-conjugacy classes of involutions
in $E$ were determined in D.34.  They are
represented by
\[
                            s,\quad xs,\quad yfs.
                                                        \tag{D.44h}
\]
Their distinct mod-two abelianization vectors
also identify the classes.
Formula (D.44d) sends $s$ to $fs$, which
belongs to the class of $s$, and gives the
exact transposition
\[
                      \Theta(xs)=yfs,\qquad
                      \Theta(yfs)=xs.
                                                        \tag{D.44i}
\]
Thus the class of $s$ is the unique fixed
class under $\Theta$.
Since $u$ commutes with $\tau$, its class
is fixed, so $u=ksk^{-1}$ for some $k\in F$.
Conjugate the finite subgroup (D.44g)
by $k^{-1}$.
It now contains $s$ and
\[
                k^{-1}\tau k=h\tau,\qquad
                              h=k^{-1}\theta(k).
                                                        \tag{D.44j}
\]
The square of this second involution gives
$h\theta(h)=1$.
Its commutation with $s$ gives
$hf=\alpha(h)$, which is precisely the
first equation of (D.44a).
On abelianization, this equation reads
$-2h^{\rm ab}=f^{\rm ab}=0$.
Multiplication by two is injective on
$\widehat\Z^2$, hence $h\in F'$.
Finally substitution in $\alpha(hy^\mu h^{-1})$
using $\alpha(h)=hf$ proves
$\alpha\rho_{\mu,h}(y)=\rho_{\mu,h}(y^{-1})$; the
$x$ identity is immediate.

We check the finite scope explicitly.
Write $Q=A_5^{19}$, and use D.38's component
description of $\theta$ and its involution tuple
$t$.  The triangle test makes every component
$t_i$ a nonidentity involution of $A_5$.
The equations to solve in this quotient are
\[
 h^{-1}t h=\theta(t),\qquad h\theta(h)=1,
       \qquad f=\theta(t)t^{-1}.
                                                        \tag{D.44k}
\]
They separate over the one-element and
two-element orbits of $\theta$ on the factors.

For a two-element orbit, write
$\theta(v)_i=g v_jg^{-1}$ and
$\theta(v)_j=g^{-1}v_i g$, with $g\in S_5$.
All nonidentity involutions of $A_5$ are
conjugate within $A_5$.
Choose $h_i\in A_5$ satisfying
$h_i^{-1}t_i h_i=g t_jg^{-1}$, and set
$h_j=g^{-1}h_i^{-1}g$.
These substitutions give both equations
of (D.44k) on that orbit.

For a fixed factor, $\theta$ is conjugation
by an involution $g\in S_5$, allowing $g=1$.
Indeed its square is the identity and
$C_{S_5}(A_5)=1$.
Choose an element $v$ with $v^2=1$, commuting
with $t_i$ and having the same parity as $g$.
When $g$ is even choose $v=1$.
When $g$ is odd, take either transposition
in the disjoint-transposition expression
for $t_i$.
Then $h_i=vg^{-1}\in A_5$ satisfies both
equations in (D.44k).
Thus every triangle survivor has a finite
first-symmetry solution, including the
$2\,531\,250$ survivors failing D.38's
finite peripheral-conjugacy test.
No lifting of these finite solutions is
asserted.

The program
\path{code/part_a/pure_involution_first_symmetry.py}
executes the local construction for all
$390$ fixed-factor cases and all
$27\,000$ two-factor cases.
It checks both equations in (D.44k),
which cover every factor orbit in D.38.
\end{proof}

\begin{claim}{D.45}{Proved}{A characteristic binary-icosahedral refinement}
Let $F=\widehat F(x,y)$ and let $R$ be D.38's characteristic
kernel with $F/R=A_5^{19}$.  Intersecting the kernels of all
epimorphisms $F\twoheadrightarrow S=\mathrm{SL}_2(\F_5)$ gives
an open characteristic subgroup $\widetilde R\subset R$ and
the following exact quotient and reduction map:
\[
 \widetilde Q=F/\widetilde R
       \xrightarrow{\ \sim\ }S^{19}\times C_2^2,
 \qquad
 \pi:\widetilde Q\longrightarrow A_5^{19},\quad
       ((g_i),a)\longmapsto(g_i\bmod\{\pm I\})_i.
                                                        \tag{D.45a}
\]
For the $10\,125\,000$ finite tuples $t$ of D.38, impose at
this specified refinement the normalized multiplier $-1$
action, both Harbater--Schneps cycle equations, and the
exact norm $f\alpha(f)=1$.
Exactly $5\,062\,500$ base tuples admit such a refined
$f$-coordinate.  They are precisely D.38's cusp-triangle
survivors, and each has $64$ admissible $f$-coordinates.
Thus the total number of refined coordinate pairs is
$324\,000\,000$.
For a fixed $t$, these $64$ coordinates define the same
actual automorphism of $\widetilde Q$, of order two.
In particular this named characteristic refinement
detects exactly D.38's first excluded half.  All
$2\,531\,250$ triangle survivors outside D.38's finite
pure conjugacy set still pass this refinement.
\end{claim}
\passage{\S{}30; folios 194--196; scans 383--387; \S{}31; folios 198--199 verso; scans 391--394; Part D of the task: full GT\_0 torsion and involution conjugacy}

\begin{proof}
All matrix calculations below are in $\F_5$.
The program enumerates the $120$ determinant-one matrices
and verifies that their commutator subgroup has order $120$.
Their center is $\{I,-I\}$, and their projective quotient is
identified explicitly with D.38's $A_5$ by a generating pair
of orders two and three with product of order five.
The resulting bijection is checked against all $60^2$
products.  Conjugation by the $120$ projective invertible
matrices gives all automorphisms of $S$.
Indeed $\Aut(S)\to\Aut(A_5)$ is injective: an element of
its kernel has the form $g\mapsto c(g)g$, where
$c:S\to\{\pm I\}$ is a homomorphism, and perfection makes
$c=1$.  D.38 gives $|\Aut(A_5)|=120$, and the enumerated
projective conjugations already give $120$ distinct
automorphisms of $S$.

Choose lifts $X_i,Y_i\in S$ of the nineteen generating
pairs in D.38, and put $Z_i=(X_iY_i)^{-1}$.
Every pair of lifts generates $S$: its image generates
$A_5$, so a proper image would have index two in $S$,
contrary to perfection.  The program also checks all
$4\cdot19=76$ choices directly.
Two maps to $A_5$ with a given kernel differ by an
automorphism, and every automorphism lifts uniquely to
$S$.  Over each of the nineteen $A_5$ kernels there
are exactly four $S$ kernels, from the independent
sign choices for the lifts of $x$ and $y$.
If two such lifts had the same kernel, the induced
automorphism of $S$ would act trivially on $A_5$;
the proved injectivity would identify the lifts.
Consequently these are exactly the $76$ epimorphism
kernels defining $\widetilde R$.

The joint map
\[
 x\longmapsto((X_i),e_x),\qquad
 y\longmapsto((Y_i),e_y),
        \qquad e_x,e_y\text{ the basis of }C_2^2,
                                                        \tag{D.45b}
\]
is onto.  To see this without a rank inference, its
projection $D$ to $S^{19}$ maps onto $A_5^{19}$ by
D.38.  Thus $D Z(S^{19})=S^{19}$; centrality gives
$[S^{19},S^{19}]=[D,D]\leq D$, so perfection gives
$D=S^{19}$.  The commutator subgroup of the image
of (D.45b) therefore contains $S^{19}\times\{0\}$,
and its projection to $C_2^2$ is onto.  This proves
surjectivity of (D.45b).
Every other sign lift of the $i$th pair is obtained by
composing (D.45b) with
\[
 ((g_j),a)\longmapsto g_i(-I)^{\ell(a)},
                         \qquad \ell\in\Hom(C_2^2,C_2).
\]
The intersection of these four kernels, for all $i$,
is trivial: the zero character first makes each
$g_i=I$, and all characters then make $a=0$.
This proves (D.45a), including the specified kernel.
It is characteristic because every automorphism of
$F$ permutes its epimorphisms to $S$.

Here are the complete finite equations used in the
refinement.  The named word substitutions $\theta,
\omega$ of D.38 descend to $\widetilde Q$.
On its characteristic perfect subgroup $S^{19}$ they
permute the factors by D.38's recorded maps and apply
the unique lifted $S$ automorphisms just described.
Write $\widetilde t_i\in S$ for any lift of $t_i$ and
put
\[
 f_{0,i}=\theta(\widetilde t)_i\widetilde t_i^{-1},
 \qquad
 f_i=(-I)^{b_i}f_{0,i},\qquad b_i\in\F_2.
                                                        \tag{D.45c}
\]
These are all the $2^{19}$ lifts of
$f_Q=\theta(t)t^{-1}$ into
$\widetilde Q'=S^{19}$.
For every such lift, set
\[
 \begin{split}
 c_{x,i}&=X_i^{-1}
       (\widetilde t_i X_i\widetilde t_i^{-1})^{-1},\\
 c_{y,i}&=(f_i^{-1}Y_i^{-1}f_i)
       (\widetilde t_i Y_i\widetilde t_i^{-1})^{-1}.
 \end{split}                                             \tag{D.45d}
\]
D.38's normalization implies $c_{x,i},c_{y,i}\in
\{I,-I\}$.  The actual automorphism induced by the
normalized pair is exactly
\[
 \alpha((g_i),(a_x,a_y))=
 \bigl((\widetilde t_i g_i\widetilde t_i^{-1}
           c_{x,i}^{a_x}c_{y,i}^{a_y})_i,
                      (a_x,a_y)\bigr).
                                                        \tag{D.45e}
\]
This formula is a continuous automorphism; its
restriction to the first factor is an automorphism,
and its two correction characters are central.
It sends the two generators (D.45b) to the required
$X^{-1}$ and $f^{-1}Y^{-1}f$.
Since $\widetilde t_i^2$ is central and both correction
characters have exponent two, its square is the
identity.  Its reduction is the nonidentity order-two
automorphism of D.38, so its order is exactly two.
Multiplying $f_i$ by a central sign changes neither
(D.45d) nor (D.45e).  Conversely these images of the
two generators determine the automorphism uniquely.

The program tests the following three exact equations
in $S^{19}$ for all choices (D.45c):
\[
 \theta(f)f=1,\qquad
 \omega^2(f)Z^{-1}\omega(f)Y^{-1}fX^{-1}=1,
 \qquad f\alpha(f)=1.
                                                        \tag{D.45f}
\]
On $f$, which has zero $C_2^2$ coordinate, formula
(D.45e) is conjugation by $\widetilde t$.
Thus every operation in (D.45f) is a specified
matrix multiplication or a specified factor map;
the norm is tested as an equation, not inferred
from the order of the finite automorphism.

The factor maps for $\theta,\omega$ have six joint
orbits.  Exhaustion within these orbits gives the
following complete counts, with D.38's zero-based
factor indices:
\[
\begin{array}{c|r|l}
 \text{factor orbit}&\text{base tuples}&
       \text{numbers of admissible central lifts}\\ \hline
 \{0,2,10\}&15&2\text{ for each of }15\\
 \{1,9,17\}&15&2\text{ for each of }15\\
 \{3,5,6,7,11,14\}&60&0\text{ for }30;\quad4\text{ for }30\\
 \{4,12,18\}&15&2\text{ for each of }15\\
 \{8,13,15\}&10&2\text{ for each of }10\\
 \{16\}&5&1\text{ for each of }5.
\end{array}                                               \tag{D.45g}
\]
The rejected thirty assignments in the six-factor
orbit are exactly its triangle failures; the other
five orbits have no triangle failure.
The equations on one joint orbit involve no other
orbit, so their solution sets have the displayed
Cartesian-product decomposition.
The surviving count is
$15\cdot15\cdot30\cdot15\cdot10\cdot5=5\,062\,500$,
and the number of central lifts per survivor is
$2\cdot2\cdot4\cdot2\cdot2\cdot1=64$.
This proves every asserted count.
The executed implementation is
\path{code/part_a/a5_schur_refinement.py}; its output
records each orbit, the lift histogram, and sample
coordinates.  These computations impose a necessary
finite refinement on full involutions.  They assert
no full lift of a surviving finite coordinate pair.
\end{proof}

\begin{claim}{D.46}{Proved}{Three positive involution sections and order-four lifting}
Let $T$ be the full once-punctured-torus loop-outer group
of $F=\widehat F(x,y)$.
For every involution $b\in T$ whose $2$-adic orientation
is positive, its centre-free extension
\[
             1\longrightarrow F\longrightarrow E_b
                  \longrightarrow\langle b\rangle
                  \longrightarrow1
\]
has exactly three $F$-conjugacy classes of order-two
sections.  For every such section $S$,
\[
                         H^1(C_F(S),\F_2)=0.             \tag{D.46a}
\]
Every cyclic order-four subgroup of $T$ has a homomorphic
lift to $\operatorname{Aut}(F)$, with its prescribed
outer action.  Thus every element of order four in
$T$ has an actual order-four representative.
\end{claim}
\passage{Part D(b) of the task; \S{}24, folio 147 verso, scan 293; \S{}30, folios 194--196, scans 383--387; \S{}31, folios 198--199 verso, scans 391--394}

\begin{proof}
\emph{The homology of a positive involution.}
D.27 gives an actual order-two representative $\beta$
of $b$, and identifies $E_b$ with
$F\rtimes_\beta\langle s\rangle$.
The map of D.33 on the finite $2$-group $\langle b\rangle$
is faithful.  The matrix of $\beta$ on
$H_1(F,\Z_2)$ is therefore nonidentity.
Its square is the identity and its determinant is $1$.
Over $\Q_2$ it is diagonalizable; the determinant
condition makes its two eigenvalues equal.
It follows that
\[
                    \beta^{\rm ab}_2=-I.               \tag{D.46b}
\]
This uses the punctured-torus statement D.33 to exclude
the identity matrix; no analogous exclusion for an
arbitrary unmarked rank-two action is assumed.

\emph{The full cohomology algebra and its Bockstein.}
Let $u:E_b\to\F_2$ be the quotient character, with
$u(s)=1$.  Choose a basis $\chi_1,\chi_2$ of
$H^1(F,\F_2)$ and extend these characters to $E_b$
by declaring $\chi_i(s)=0$.
These extensions exist because (D.46b) is the identity
modulo $2$.
The Hochschild--Serre spectral sequence has precisely
the two rows
\[
 E_2^{a,0}=\F_2 u^a,\qquad
 E_2^{a,1}=H^1(F,\F_2)\,u^a.                            \tag{D.46c}
\]
The transgression on $E_2^{0,1}$ is zero since the two
characters have just been extended to $E_b$.
Its other possible differentials are the products of
this transgression with powers of $u$, and hence are
zero as well.  There are no further nonzero rows.
Consequently, for every $n\geq1$,
\[
 H^n(E_b,\F_2)
   =\langle u^n,u^{n-1}\chi_1,u^{n-1}\chi_2\rangle_{\F_2},
 \qquad \dim H^n(E_b,\F_2)=3.                           \tag{D.46d}
\]
Independence follows from the two associated graded
pieces in (D.46c).
Multiplication by $u$ is an isomorphism in every
degree at least one: it is an isomorphism on those
associated graded pieces.

For each $i$, lift $\chi_i|_F$ to a continuous
homomorphism $\varphi_i:F\to\Z/4$.
Equation (D.46b) gives
$\varphi_i(\beta(h))=-\varphi_i(h)$.
The actual continuous one-cochain on $E_b$
\[
             \widetilde\chi_i(hs^\epsilon)=\varphi_i(h)
                      \quad(\epsilon=0,1)
\]
therefore satisfies
\[
 \begin{split}
 \delta\widetilde\chi_i(hs^\epsilon,ks^\eta)
  &=\varphi_i(h)+\varphi_i(k)
             -\varphi_i(h\beta^\epsilon(k))\\
  &=2\epsilon\,\varphi_i(k)
       \quad\text{in }\Z/4.
 \end{split}                                           \tag{D.46e}
\]
For the coefficient sequence
$0\to\F_2\xrightarrow{\;2\;}\Z/4\to\F_2\to0$,
this says that the Bockstein of $\chi_i$ is
$u\chi_i$.
It is also $\chi_i^2$.
For completeness, the lift of any $\F_2$ character
which takes its values in $\{0,1\}\subset\Z/4$
has coboundary, divided by $2$, equal to the cup
square of that character; comparison of two lifts
changes this cocycle only by a coboundary.
Thus
\[
                            \chi_i^2=u\chi_i.           \tag{D.46f}
\]

Invert the central class $u$ and take degree zero:
\[
             R=\bigl(H^*(E_b,\F_2)[u^{-1}]\bigr)^0.
\]
It has dimension three, with basis
$1,z_1=\chi_1/u,z_2=\chi_2/u$ by (D.46d).
Equation (D.46f) makes both $z_i$ idempotent.
In particular $R$ is a quotient of the algebra
\[
 \F_2[Z_1,Z_2]/(Z_1^2-Z_1,Z_2^2-Z_2)
           \cong\prod_{(a_1,a_2)\in\F_2^2}\F_2,
 \quad P\longmapsto(P(a_1,a_2))_{(a_1,a_2)}.             \tag{D.46g}
\]
The displayed evaluation map is an isomorphism:
the four products of $Z_i$ and $1-Z_i$ are its four
coordinate idempotents.
Every ideal of a finite product of fields is a product
of ideals in those fields.  Every quotient is therefore
reduced, so $R$ has no nonzero nilpotent element.

\emph{Brown restriction and the three full sections.}
Every finite subgroup of $E_b$ injects into
$E_b/F=C_2$, since $F$ is torsion-free.
Thus the nontrivial elementary abelian $2$-subgroups
are exactly its order-two sections, with no inclusions
between distinct ones.
For such a section $S$ put $C_S=C_F(S)$.
Its normalizer is the exact direct product
\[
                         N_{E_b}(S)=C_S\times S.        \tag{D.46h}
\]
Indeed every normalizer element centralizes $S$, since
$\operatorname{Aut}(C_2)=1$; multiplying it by its
unique image in $S$ leaves an element of $C_S$.

The full-profinite Brown restriction theorem used in
D.28 and D.33 applies because $F$ is open in $E_b$
and has $2$-cohomological dimension one
\cite[Theorem~2 and \S1.1]{A-KochloukovaZalesskii2007}.
In sufficiently high degrees it identifies cohomology
with the product of the cohomologies of (D.46h);
the identification is the actual restriction map and
preserves cup products.
Each factor contributes a copy of $\F_2$ in every
positive degree, so (D.46d) in particular makes the
number of factors finite.
Since $C_S$ is a closed subgroup of $F$, it has
$\operatorname{cd}_2(C_S)\leq1$.
After the same localization, restriction identifies
$R$ as a ring with
\[
       \prod_{[S]}
         \bigl(\F_2\oplus H^1(C_S,\F_2)\,u^{-1}\bigr).
                                                        \tag{D.46i}
\]
The second summand in each factor is an ideal with
square zero, since $H^2(C_S,\F_2)=0$.
Here $u|_{C_S}=0$ and $u|_S$ is the nonzero
character, so the localization in (D.46i) is exactly
the localization already used in (D.46g).
Reducedness of $R$ forces every second summand to
vanish.  Its dimension three then proves that there
are exactly three section classes, and proves
(D.46a).

These are also precisely the $F$-conjugacy classes.
To see this directly, let $e\in E_b$ and choose
$s'\in S$ with the same image as $e$ in $E_b/F$.
Then $eS e^{-1}=(es'^{-1})S(es'^{-1})^{-1}$,
and $es'^{-1}\in F$.

\emph{An invariant section in the order-four extension.}
Now let $a\in T$ have order four, and form its
centre-free extension
\[
 E=\{(\phi,j)\in\operatorname{Aut}(F)\times C_4:
                                      [\phi]=a^j\}.
                                                        \tag{D.46j}
\]
Its kernel is identified with $F$ by
$h\mapsto(\operatorname{Inn}(h),0)$.
Let $E_2$ be the inverse image of $\langle a^2\rangle$.
The orientation of $a^2$ at $2$ is positive:
the only roots of unity of order dividing four
in $\Z_2^\times$ are $1,-1$, whose squares are $1$.
Equivalently, D.33 embeds $\langle a\rangle$ in
$\operatorname{GL}_2(\Z_2)$, where its square is $-I$.
The first part of the proof thus applies to $E_2$.

Conjugation by $E$ permutes the three $F$-conjugacy
classes of sections of $E_2$.
Every element of $E_2$ acts trivially on this class
set, by the preceding paragraph, so the permutation
action factors through $E/E_2=C_2$.
An action of $C_2$ on a set of three elements has a
fixed element.
Choose a section $S$ in that class.
If $e\in E$ projects to the specified generator of
$C_4$, there is $f\in F$ with
$eSe^{-1}=fSf^{-1}$.
Hence $f^{-1}e\in N_E(S)$ still projects to that
generator.
Put $N=N_E(S)$ and $C=C_F(S)$.
We have proved the exact sequences
\[
 \begin{gathered}
 1\longrightarrow C\longrightarrow N
              \longrightarrow C_4\longrightarrow1,
       \qquad S\subseteq Z(N),\\
 1\longrightarrow C\longrightarrow J:=N/S
              \longrightarrow C_2\longrightarrow1.
 \end{gathered}                                         \tag{D.46k}
\]
The embedding of $C$ in the second line is injective
because $C\cap S=1$.

\emph{Projectivity supplies the order-four section.}
Equation (D.46a) implies
\[
                              C^{(2)}=1:                \tag{D.46l}
\]
every nontrivial pro-$2$ group has a nontrivial finite
$2$-group quotient, which in turn has a quotient
$C_2$.
In particular every homomorphism from $J$ to a finite
$2$-group kills $C$ and factors through the last
$C_2$ in (D.46k).  That quotient is nontrivial, so
\[
                              J^{(2)}=C_2.              \tag{D.46m}
\]

We show that $J$ has a nonidentity torsion element.
If it were torsion-free, its open subgroup $C$ has
$\operatorname{cd}_p(C)\leq1$ at every prime $p$,
as a closed subgroup of the free profinite group $F$.
Serre's torsion-free finite-index dimension theorem,
in exactly the form used in D.27, would give
$\operatorname{cd}_p(J)\leq1$ at every prime.
Thus $J$ would be projective as a full profinite
group.
Its maximal pro-$2$ quotient would then be a
projective, hence free, pro-$2$ group: pull a finite
$2$-group embedding problem back to $J$, solve it
by projectivity, and factor the solution through
$J^{(2)}$.
This contradicts (D.46m), since $C_2$ is not a free
pro-$2$ group.

Let $j\in J$ be a nonidentity torsion element.
Its cyclic subgroup injects into $J/C=C_2$, because
$C$ is torsion-free.
Consequently $j$ has order two and lies outside $C$.
Choose $n\in N$ over $j$.
Then $n^2\in S$, whereas its image in $C_4$ is an
odd element and has order four.
It follows that $n^2$ is the nonidentity element
of $S$, and $n$ has order exactly four.
After replacing $n$ by $n^{-1}$ if necessary, its
projection is the specified generator of $C_4$.
Thus $C_4\to E$, sending that generator to $n$,
is a section of (D.46j).
Projection to $\operatorname{Aut}(F)$ gives the
claimed homomorphic lift with its original labels.

The program \texttt{code/part\_c/torus\_order4\_lift\_d46.py}
checks the finite cochain formula (D.46e), its
comparison with the cup square, and the explicit
Boolean-algebra evaluation in (D.46g).
The dimension and normalizer arguments above prove
the statement for the full profinite extensions.
No claim of conjugacy to a geometric order-four
action is made.  D.43 applied to $a^2$ does also
give $\lambda(a)^2=1$ globally, so the remaining
orientation issue for $a$ concerns mixed signs,
rather than fourth roots of unity.
\end{proof}

\begin{claim}{D.47}{Proved}{Conjugator coordinates at the binary refinement}
Use the specified quotient
$\widetilde Q=S^{19}\times C_2^2$, $S=\mathrm{SL}_2(\F_5)$,
of D.45, and its $324\,000\,000$ refined pairs $(t,f)$.
Every one of these pairs satisfies D.44's finite
first-symmetry coboundary condition
\[
       h^{-1}\alpha(h)=f,\qquad h\theta(h)=1,
                         \qquad h\in S^{19}.             \tag{D.47a}
\]
Exactly $10\,125\,000$ of the pairs satisfy the two
finite conjugator equations of D.35 at this quotient.
The count stays the same after requiring both actual
invertibility of the finite substitution and both
Harbater--Schneps symmetries for the conjugating pair.
There are $5\,062\,500$ pairs with admissible unit
parameters congruent to a square modulo five, and
$5\,062\,500$ with parameters congruent to a nonsquare.
These two sets are disjoint.

Thus $313\,875\,000$ refined coordinate pairs lie
outside the image of the standard full $GT_0$
involution conjugacy class.  For each of D.38's
$2\,531\,250$ conjugator-compatible base tuples,
exactly four of its $64$ refined $f$-coordinates
pass this test; for each of the other
$2\,531\,250$ base tuples, none does.
These are statements about the exact normalized
coordinates: D.45's $64$ coordinates for a fixed
$t$ define the same finite automorphism.
No surviving pair is asserted to have a full
involution or conjugator lift.
\end{claim}
\passage{\S{}30; folios 194--196; scans 383--387; \S{}31; folios 198--199 verso; scans 391--394; Part D of the task: full GT\_0 involution conjugacy}

\begin{proof}
Write $X_i,Y_i,Z_i,\widetilde t_i$ and $\alpha$ for
the matrices and actual automorphism fixed in D.45.
Since $S$ has exponent dividing $60$, the power
parameter in D.35 factors through
$(\Z/60)^\times$.  For the second Harbater--Schneps
symmetry we retain the finer, explicit parameter map
\[
 a\in(\Z/120)^\times\longmapsto
       \left(a\bmod60,\ m=(a-1)/2\bmod60\right).
                                                        \tag{D.47b}
\]
It is well-defined, and every such $a$ is the
reduction of a unit in $\widehat\Z$.
Changing the representative of $a$ by $120$ changes
$m$ by $60$, so all subsequent equations are
independent of the choice of integer representative.

The image of $F'$ is exactly $S^{19}$ by D.45.
For fixed $(t,f)$ the first conjugator equation
therefore separates into the nineteen exact equations
\[
 h_i^{-1}\widetilde t_i h_i\widetilde t_i^{-1}=f_i,
                                  \qquad h_i\in S.       \tag{D.47c}
\]
Each has exactly four solutions.  The program
verifies this by enumeration of all $120$ matrices.
It also follows from the order-four conjugacy class
in $S$: $\widetilde t_i$ has square $-I$, and the
norm equation in D.45 gives
$(f_i\widetilde t_i)^2=\widetilde t_i^2=-I$.
The thirty matrices with square $-I$ form one
conjugacy class with centralizer of order four,
as the same matrix enumeration verifies.

For each parameter $a\bmod60$ and each of these
four matrices, the second equation of D.35 is
tested by the exact matrix conjugacy condition
\[
      X_i^a h_iY_i^a h_i^{-1}
                       \sim_S (X_iY_i)^a.               \tag{D.47d}
\]
Indeed the $C_2^2$ coordinates of both sides are
$e_x+e_y$, since $a$ is odd, and that direct factor
is central.  Conjugacy in $\widetilde Q$ is therefore
exactly these nineteen independent matrix conjugacies.
The program computes every conjugacy class in $S$
by conjugation by all $120$ matrices; (D.47d) is
tested against those complete classes.

The finite invertibility test is explicit as well.
Reduce the ordered pair
$(X_i^a,h_iY_i^a h_i^{-1})$ to $A_5$ and test whether
it belongs to the $i$th recorded $\Aut(A_5)$ orbit
of generating pairs in D.38.
This is sufficient: its unique component
automorphism lifts to an automorphism $\gamma_i$
of $S$, and the two central differences from
$\gamma_i(X_i),\gamma_i(Y_i)$ are corrected by
characters of $C_2^2$, exactly as in (D.45d)--(D.45e).
The resulting automorphism of $\widetilde Q$ has
the prescribed two generator images.
It is necessary for a finite pure automorphism
with these images: its reduction preserves the
three specified power-conjugacy classes, and
D.38's distinct passports force it to preserve
each simple factor.  Its pair must therefore be
in that recorded orbit.  This tests the actual
finite substitution, with its characteristic
quotient map retained.

For completeness the program also imposes both
symmetries on the conjugating pair, not just its
three peripheral images.
Put $k=h^{-1}$ and use $m$ from (D.47b).
The equations tested are
\[
       h\theta(h)=1,\qquad
       \omega^2(k)Z^m\omega(k)Y^m kX^m=1.
                                                        \tag{D.47e}
\]
All $h,k,f$ coordinates lie in the perfect factor
$S^{19}$; the actions of $\theta,\omega$ there
are the specified lifted factor maps of D.45.
The $C_2^2$ coordinate of the second equation is
$m(e_x+e_y+e_z)=0$.
Thus these are exact equations at the indicated
finite quotient, including the half-exponent
required by relation (II).

The first equation of (D.47e) is exhausted without
assuming independent choices where factors are
exchanged.  If
$\theta(v)_i=G_i v_jG_i^{-1}$, a choice of $h_i$
forces
\[
                       h_j=G_i^{-1}h_i^{-1}G_i.         \tag{D.47f}
\]
For a fixed factor it must satisfy the same equation
with $j=i$.  The program enumerates one of the
four candidates from (D.47c) for each orbit of
$\theta$, uses (D.47f), and checks membership in
the corresponding candidate sets.
It verifies that (D.47a) is soluble for all the
refined pairs.  For the full conjugator test it
then checks (D.47d), finite invertibility, and the
second equation of (D.47e) for every remaining
choice and all $32$ units modulo $120$.

Here is the entire factorization of the result.
Let $U_+$ and $U_-$ denote the subsets of
$(\Z/120)^\times$ whose residues modulo five are
respectively squares and nonsquares; each has
$16$ elements.
For each of D.45's six joint factor orbits, the
table gives the number of refined pairs having
the specified set of admissible parameters.
\[
\begin{array}{c|r|r|r|r}
 \text{factor orbit}&U_+&U_-&U_+\cup U_-&\varnothing\\ \hline
 \{0,2,10\}&15&15&0&0\\
 \{1,9,17\}&15&15&0&0\\
 \{3,5,6,7,11,14\}&30&30&0&60\\
 \{4,12,18\}&15&15&0&0\\
 \{8,13,15\}&10&10&0&0\\
 \{16\}&0&0&5&0.
\end{array}                                               \tag{D.47g}
\]
The admissible global parameter set is the
intersection of the six local sets, because
the parameter is common while the factor
equations have the exact orbit decomposition
of D.45.
It is nonempty precisely when the first five
nonempty choices all select the same one of
$U_+,U_-$.
Each choice gives
$15\cdot15\cdot30\cdot15\cdot10\cdot5
 =5\,062\,500$ pairs.
Their sum is $10\,125\,000$, and subtraction
from D.45's $324\,000\,000$ gives the stated
$313\,875\,000$ failures.
The program records the same table when only
(D.47c)--(D.47d) are imposed, and when finite
invertibility and the first symmetry are added.
Thus the asserted equality of these counts is
itself checked at every factor orbit.
It also records the counts for each base tuple
before multiplying over the six orbits.  Their
Cartesian product gives the exact histogram
\[
 \begin{array}{c|r}
 \text{admissible }f\text{-coordinates per base tuple}
          &\text{number of base tuples}\\ \hline
 0&2\,531\,250\\
 4&2\,531\,250.
 \end{array}                                             \tag{D.47h}
\]

Every full conjugator in D.35 reduces to solutions
of (D.47c)--(D.47d), and D.35's equality of the
$P$ and $GT_0$ conjugator sets also supplies
(D.47e).
Consequently failure of the displayed finite
test excludes that exact coordinate pair from
the image of the standard full conjugacy class.
The executed implementation is
\path{code/part_a/pure_involution_binary_conjugacy.py}.
Its output retains all unit sets, the six local
tables, the per-base histogram, and the separate
first-symmetry counts.  The proof uses only this
finite refinement and makes no assertion that
one of its surviving pairs lifts through all
characteristic refinements of $F$.
\end{proof}

\begin{claim}{D.48}{Proved}{The two sections of a punctured-torus order-three action}
Let $a$ have order three in the full once-punctured-torus
loop-outer group of $F=\widehat F(x,y)$, and let $E_a$
be its centre-free extension.
There are exactly two $F$-conjugacy classes of order-three
sections in $E_a$.  Each such section $S$ satisfies
\[
                         H^1(C_F(S),\F_3)=0.             \tag{D.48a}
\]
The normalizer in $E_a$ of the torus cusp
$L=\overline{\langle[x,y]\rangle}$ is torsion-free.
If $v$ is any element of this normalizer projecting
to the prescribed generator $a$, and
\[
             v[x,y]v^{-1}=[x,y]^\lambda,\qquad
                            v^3=[x,y]^\kappa ,
\]
then its $3$-adic coordinate satisfies
$\kappa_3\in\Z_3^\times$.
\end{claim}
\passage{Part D(b) of the task; \S{}24, folio 147 verso, scan 293; \S{}30, folios 194--196, scans 383--387; \S{}31, folios 198--199 verso, scans 391--394}

\begin{proof}
\emph{The actual action and its abelian quotient.}
D.27 gives an actual order-three representative
$\alpha$ of $a$, so
$E_a=F\rtimes_\alpha\langle s\rangle$.
Let $A$ be its matrix on $H_1(F,\Z_3)$.
It has order exactly three by D.33.
Over $\Q_3$ the polynomial $X^2+X+1$ is irreducible:
a primitive cube root would give a square root of
$-3$, whose $3$-adic valuation is odd.
Since $A^3=I$ and $A\ne I$, it follows that
\[
             A^2+A+I=0,\qquad
             \det(A-I)=3,\qquad
             \operatorname{coker}(A-I)\cong C_3.
                                                        \tag{D.48b}
\]
The last equality follows by Smith reduction over
$\Z_3$, since the determinant has valuation one.
Abelianizing the actual semidirect product gives
the exact maximal pro-$3$ abelian quotient
\[
                     E_a^{\rm ab,(3)}\cong C_3^2.        \tag{D.48c}
\]
One factor is (D.48b), the other is the specified
section quotient.

The reduction of $A$ modulo $3$ is a nontrivial
unipotent Jordan block.  Indeed, if $A=I+3B$,
then $A^3=I$ implies
$B(I+3B+3B^2)=0$ after division by $9$.
The second factor is invertible, giving $B=0$,
contrary to $A\ne I$.
The induced action on
$M=H^1(F,\F_3)$ is therefore also one nontrivial
Jordan block of dimension two.
Writing $D$ for the action matrix on this dual
space, the exact cyclic resolution has successive
operators
\[
              D-I,\quad 1+D+D^2=0,\quad D-I,\quad0,\ldots,
 \qquad \operatorname{rank}(D-I)=1.
                                                        \tag{D.48d}
\]
Consequently $H^j(C_3,M)$ has dimension one for
every $j\geq0$.

\emph{Bockstein injectivity and periodicity.}
Let $u\in H^1(E_a,\F_3)$ be the quotient character,
and let $\chi$ be a nonzero invariant character of
$F$, extended to $E_a$ by $\chi(s)=0$.
They form a basis of $H^1(E_a,\F_3)$ by (D.48c).
For the coefficient sequence
\[
              0\longrightarrow\F_3
                 \xrightarrow{\;3\;}\Z/9
                 \longrightarrow\F_3\longrightarrow0
\]
write $\mathfrak b$ for the Bockstein.
Every homomorphism $E_a\to\Z/9$ factors through
(D.48c), so its values are in $3(\Z/9)$ and its
reduction modulo $3$ is zero.
The coefficient long exact sequence therefore makes
\[
     \mathfrak b:H^1(E_a,\F_3)\longrightarrow
                              H^2(E_a,\F_3)
                  \quad\hbox{injective}.                \tag{D.48e}
\]
Put $v=\mathfrak b(u)$ and $w=\mathfrak b(\chi)$.
The notation $v$ in this cohomology calculation
denotes a class, not the normalizer element in the
statement.
The class $v$ is the inflation of the usual degree-two
period class of $C_3$; it is nonzero on the section
$\langle s\rangle$.
The classes $v,w$ are linearly independent by
(D.48e).

The Hochschild--Serre spectral sequence has only
rows zero and one.  By (D.48d), at each positive
total degree its $E_2$ page has two
one-dimensional terms.  Thus its degree-two
abutment has dimension at most two.
The two independent Bockstein classes make its
dimension exactly two.
The only possible differential is
$d_2:E_2^{j,1}\to E_2^{j+2,0}$.
At $j=0$ it vanishes because $\chi$ extends to
$E_a$, and at $j=1$ it vanishes because the
one-dimensional term $E_2^{1,1}$ must survive
to obtain dimension two in degree two.

Multiplication by the period class $v$ identifies
each cyclic-resolution term with the term two
places later.
On the row with coefficient $M$ this also holds
starting at $j=0$, since the norm in (D.48d) is
zero, so $H^0(C_3,M)\to H^2(C_3,M)$ is the identity
on $\ker(D-I)$.
On the row with trivial coefficients it is the
usual period isomorphism.
The multiplicative spectral sequence therefore
makes every remaining $d_2$ a multiple of one of
the two already vanishing differentials.
It collapses, and multiplication by $v$ is an
isomorphism on its two associated graded pieces
at every positive total degree.  Hence
\[
 \begin{gathered}
        \dim H^n(E_a,\F_3)=2\quad(n\geq1),\\
        v:H^n(E_a,\F_3)\xrightarrow{\ \sim\ }
                             H^{n+2}(E_a,\F_3)
                       \quad(n\geq1),\\
                         wv^m\ne0\quad(m\geq0).
 \end{gathered}                                         \tag{D.48f}
\]

\emph{Exactly two full section classes.}
Every finite subgroup of $E_a$ injects into its
quotient $C_3$, because $F$ is torsion-free.
For an order-three section $S$ the normalizer is
the exact direct product
\[
                         N_{E_a}(S)=C_F(S)\times S.      \tag{D.48g}
\]
Conjugation by a normalizer element induces the
identity on $S$: projection to the abelian quotient
$C_3$ identifies every element of $S$ with its
label, so that label cannot change under
conjugation.
Multiplying by the section element with the same
quotient label then leaves precisely $C_F(S)$.

The full-profinite Brown restriction theorem
\cite[Theorem~2 and \S1.1]{A-KochloukovaZalesskii2007},
applied with the open subgroup $F$ of
$3$-cohomological dimension one, identifies
sufficiently high cohomology with the product
over (D.48g).
Since $C_F(S)$ is a closed subgroup of $F$,
its $3$-cohomological dimension is at most one.
Writing $k$ for the number of section classes,
the dimensions in (D.48f) therefore give
\[
                  k+\sum_{[S]}\dim H^1(C_F(S),\F_3)=2.
                                                        \tag{D.48h}
\]
In particular the class set is finite.

It cannot consist of a single class.
Otherwise take that class to be the chosen
$S=\langle s\rangle$.
On its normalizer (D.48g), the character $\chi$
is pulled back from $C_F(S)$, since $\chi(s)=0$.
Its Bockstein therefore restricts to zero:
\[
       w|_{N_{E_a}(S)}
        =\operatorname{pr}_{C_F(S)}^*
               \mathfrak b(\chi|_{C_F(S)})=0,
                                                        \tag{D.48i}
\]
because $H^2(C_F(S),\F_3)=0$.
Every nonzero class $wv^m$ from (D.48f) would
then have zero Brown restriction, a contradiction
for large $m$.
Thus $k\geq2$; (D.48h) makes $k=2$ and makes
both summands zero, proving (D.48a).
As in D.46, $E_a$-conjugacy equals $F$-conjugacy:
for any section $S$ choose its unique element
with the same quotient label as a given
conjugator and multiply the conjugator by its
inverse.

\emph{The peripheral normalizer.}
We spell out the boundary consequence of the
vanishing in (D.48a).
A closed subgroup $C\leq F$ with
$H^1(C,\F_3)=0$ cannot contain a conjugate of
the full $3$-part
$\langle[x,y]^{e_3}\rangle\cong\Z_3$
of the torus cusp.
First its supernatural index in $F$ must have
infinite $3$-part.  Otherwise choose an open
$U_0\geq C$ whose index has maximal $3$-valuation.
For every open $V$ with $C\leq V\leq U_0$,
the index $[U_0:V]$ is prime to $3$.
Transfer makes
$H^1(U_0,\F_3)\to H^1(V,\F_3)$ injective.
The restriction-colimit comparison of D.32
would then inject the nonzero group
$H^1(U_0,\F_3)$ into $H^1(C,\F_3)$,
a contradiction.

Choose an open $U\geq C$ with $3\mid[F:U]$.
Its corresponding finite cover of the punctured
torus has a lift of the selected cusp of degree
prime to $3$, because its completed boundary
subgroup contains the selected full $\Z_3$.
The exact discrete boundary comparison uses
the dense discrete representative for the
profinite conjugator, as in D.32 and D.33.
There are at least two boundary lifts: if there
were only one, its degree would be $[F:U]$.
The selected boundary class is nonzero in
$H_1(U,\F_3)$, since the only relation among
the boundary classes is their sum.
A boundary-dual character restricts nontrivially
to the selected $\Z_3$ in $C$; the scalar
comparison is the inverse of its covering
degree modulo $3$.
This again contradicts $H^1(C,\F_3)=0$ and proves
the assertion.

For $L=\overline{\langle[x,y]\rangle}$, the
peripheral normalizer has the exact sequence
\[
           1\longrightarrow L\longrightarrow N_{E_a}(L)
                    \longrightarrow C_3\longrightarrow1.
                                                        \tag{D.48j}
\]
Surjectivity follows from preservation of the cusp
class, and its kernel is $L$ by A.7.
The multiplier satisfies $\lambda^3=1$, and
$\lambda_3=1$, since $\Z_3^\times$ has no
nontrivial root of unity of order three.
An order-three section in (D.48j) would therefore
centralize the full cusp $\Z_3$.
Its fixed subgroup in $F$ would contradict the
boundary consequence of (D.48a).
Every torsion subgroup of (D.48j) injects into
$C_3$, so this proves that the normalizer is
torsion-free.

Finally take the element called $v$ in the
statement, with $v^3=[x,y]^\kappa$.
Commuting $v$ with its own cube gives
$(\lambda-1)\kappa=0$.
Where $\lambda_p\ne1$, the coordinate $\kappa_p$
and the norm $1+\lambda_p+\lambda_p^2$ are zero.
Where $\lambda_p=1$, that norm is $3$.
If $\kappa_3$ were divisible by $3$, the exact
norm equation
\[
                (1+\lambda+\lambda^2)z=-\kappa
                               \quad\text{in }\widehat\Z
                                                        \tag{D.48k}
\]
would consequently have a solution: divide by
$3$ at the positive coordinates, and use zero
at the other coordinates.
The element $v[x,y]^z$ would then have order
three and project to $a$, contradicting the
torsion-free conclusion.
Thus $\kappa_3$ is a unit, as asserted.

The program
\texttt{code/part\_c/torus\_order3\_sections\_d48.py}
checks the integral and mod-$3$ matrix calculations,
the abelian character and Bockstein input in its
finite semidirect quotient, and the cyclic
resolution operators.
The full restriction and finite-cover arguments
above establish the section and normalizer
assertions.  They do not assert global triviality
of the order-three multiplier or geometric
conjugacy.
\end{proof}

\begin{claim}{D.49}{Proved}{The standard centralizer and the two-element conjugator fibres}
Let $P$ be the full pure three-cusp group of D.32 and
$\delta(x)=x^{-1}$, $\delta(y)=y^{-1}$.
Its full centralizer and normalizer are
\[
             C_P(\delta)=N_P(\langle\delta\rangle)
                         =\langle\delta\rangle.         \tag{D.49a}
\]
Define the actual conjugacy orbits
\[
 \mathcal O_P=\{\beta\delta\beta^{-1}:\beta\in P\},
 \qquad
 \mathcal O_{GT_0}=\{\beta\delta\beta^{-1}:\beta\in GT_0\}.
\]
Every fibre of $P\to\mathcal O_P$ is exactly the
right coset $\{\beta,\beta\delta\}$, and
\[
                    \mathcal O_{GT_0}
                            =\mathcal O_P\cap GT_0.     \tag{D.49b}
\]
For $\alpha\in\mathcal O_P$, its two conjugators
have the exact normalized formulas
\[
 \begin{array}{ll}
  x\mapsto x^\mu,& y\mapsto h y^\mu h^{-1},\\
  x\mapsto x^{-\mu},& y\mapsto h y^{-\mu}h^{-1},
 \end{array}
       \qquad \mu\in\widehat\Z^\times,\quad h\in F'.   \tag{D.49c}
\]
Here $h$ is unique, and the unordered pair
$\{\mu,-\mu\}$ is unique.
In D.35's solution space, a choice $(\mu,h,u)$
therefore gives all the solutions over $\alpha$ as
\[
 (\mu,h,u(xy)^a),\qquad
 (-\mu,h,x^{-\mu}u(xy)^a),\qquad a\in\widehat\Z.
                                                        \tag{D.49d}
\]
These assertions concern the specified orbit map;
they do not assert that it contains every full
pure or $GT_0$ involution.
\end{claim}
\passage{\S{}30; folios 194--196; scans 383--387; \S{}31; folios 198--199 verso; scans 391--394; Part D of the task: full GT\_0 involution conjugacy}

\begin{proof}
Lochak--Schneps prove that complex conjugation is
self-centralizing in the usual $GT$
\cite[Proposition 4(ii), with Lemma 9]{R-LS}.
Their argument uses the first two defining relations
and applies to $GT_0$ as well.  We give the passage
from the full group $P$ explicitly, together with
the normalized fibre formulas.

Let $\beta\in C_P(\delta)$ have normalized pair
$(\mu,k)$, so
$\beta(x)=x^\mu$ and $\beta(y)=k^{-1}y^\mu k$,
with $k\in F'$.
D.32's normalization is a homomorphic section
into $\Aut(F)$, so the commutation is an actual
automorphism equality.
Its value on $y$ makes $k\delta(k)^{-1}$
centralize $y^\mu$.
The peripheral centralizer is $\langle y\rangle$
by A.7, while $k\delta(k)^{-1}\in F'$.
Abelianization therefore gives $k=\delta(k)$.
The exact free-product identification of D.35,
$F\rtimes_\delta C_2\cong C_2\amalg C_2\amalg C_2$,
gives $C_F(\delta)=1$.  Hence $k=1$.

We also need the precise inclusion in $GT_0$.
For each geometric cusp permutation $\sigma\in\Sigma$,
form $r_\sigma=\beta^{-1}\sigma\beta\sigma^{-1}$.
It lies in $P^1=\ker\lambda$, and commutes with
$\delta$, because both $\beta$ and $\sigma$ do.
D.35 proved $C_{P^1}(\delta)=1$, so $r_\sigma=1$.
The exact identity $GT_0=C_P(\Sigma)$ of D.35
therefore puts this $\beta$ in $GT_0$.

Its second Harbater--Schneps relation, with
$k=1$ and $m=(\mu-1)/2\in\widehat\Z$, is
\[
                (xy)^{-m}y^m x^m=1,
       \qquad\text{equivalently}\quad
                         (xy)^m=y^m x^m.                \tag{D.49e}
\]
For every positive integer $n$ use the explicit
finite Heisenberg quotient with coordinates
$(a,b,c)\in(\Z/n)^3$ and multiplication
\[
 (a,b,c)(a',b',c')=(a+a',b+b',c+c'+ab'),
       \quad x=(1,0,0),\quad y=(0,1,0).
\]
The two sides of (D.49e) have coordinates
\[
 (m,m,m(m+1)/2),\qquad(m,m,0),
                                                        \tag{D.49f}
\]
where the binomial expression is the continuous
integer-valued function on $\widehat\Z$.
These identities first follow for nonnegative
integers by multiplication, and then for all
$m\in\widehat\Z$ by continuity; modulo $n$ they
depend only on $m\bmod2n$.
Equality in every such quotient gives
$m(m+1)/2=0$ in $\widehat\Z$, and hence
\[
                           \mu^2=1.                     \tag{D.49g}
\]
Since $k=1$, the actual square of $\beta$ fixes
both $x$ and $y$ by (D.49g).  Thus $\beta^2=1$.
D.34 gives the global multiplier $-1$ for a
nonidentity pure involution; the identity has
multiplier $1$.
The normalized pair is therefore $(1,1)$ or
$(-1,1)$, proving $C_P(\delta)=\langle\delta\rangle$.
A normalizer of its order-two subgroup must
centralize its sole nonidentity element, proving
the normalizer equality in (D.49a).

Two elements $\beta,\beta'$ have the same image
under the orbit map exactly when
$\beta^{-1}\beta'\in C_P(\delta)$.
This proves its two-element fibres.
For an orbit element lying in $GT_0$, D.35's
proved equality of conjugator sets places both
conjugators in $GT_0$; the reverse inclusion is
immediate.  This proves (D.49b).

Write one conjugator in D.35's convention as
$\beta_{\mu,h}$.  Right composition with $\delta$
gives exactly $\beta_{-\mu,h}$, so (D.49c) follows
with the same $h$.
They are distinct because $\mu\ne-\mu$ in
$\widehat\Z^\times$.
Moreover actual normalized conjugation identifies
$C_F(\alpha)$ with $\beta(C_F(\delta))=1$.
The description of all coboundaries in (D.35g)
then makes $h$ the unique solution of
$h^{-1}\alpha(h)=f$ in $F'$.
The two-element fibre proves uniqueness of the
unordered multiplier pair.

Finally put $A=x^\mu$, $B=h y^\mu h^{-1}$ and
suppose $AB=u(xy)^\mu u^{-1}$.
The exact word identity
\[
           A^{-1}B^{-1}=A^{-1}(AB)^{-1}A
\]
shows that $x^{-\mu}u$ is a third-cusp conjugator
for the opposite multiplier.
D.35's centralizer calculation gives all other
third-cusp conjugators by right multiplication
by $(xy)^a$.  This proves (D.49d).
The program
\path{code/part_a/standard_involution_centralizer.py}
checks (D.49f) over sixteen finite moduli and
checks the displayed word identity.
The full conclusion (D.49a) uses the proved
global-sign theorem D.34, beyond those finite
checks.
\end{proof}

\begin{claim}{D.50}{Proved}{The finite two-subgroups of the full punctured-torus group}
Let
$T=\Out_{\rm lac}(\widehat F(x,y);
\overline{\langle[x,y]\rangle})$ be the full
once-punctured-torus loop-outer group.
The possible abstract isomorphism types of its finite
$2$-subgroups are exactly
\[
                  1,\quad C_2,\quad C_4,\quad
                         C_2\times C_2,\quad D_8,
                                                        \tag{D.50a}
\]
where $D_8$ has order eight. All five types occur.
For a subgroup of type $C_2\times C_2$ or $D_8$,
the full multiplier takes its values in the globally
constant signs $\{1,-1\}\subset\widehat\Z^\times$.
In $C_2\times C_2$ there is one positive involution
and two negative involutions. In $D_8$ the rotation
subgroup has multiplier $1$ and the four reflections
have multiplier $-1$.

For an arbitrary cyclic subgroup of order four,
a generator $a$ satisfies the already proved
restrictions
\[
           \lambda(a)^2=1,\qquad \lambda(a)_2=1,
           \qquad \lambda(a^2)=1.                       \tag{D.50b}
\]
The abstract classification (D.50a) does not establish
conjugacy of arbitrary such subgroups, or replace
(D.50b) by a global multiplier determination.
\end{claim}
\passage{Section 30, folios 194-196, scans 383-387; Section 31, folios 198-199 verso, scans 391-394; Part D(b), full loop-outer finite actions in type (1,1)}

\begin{proof}
D.33 makes the actual homology homomorphism
\[
                     A_2:T\longrightarrow\GL_2(\Z_2)
                                                        \tag{D.50c}
\]
injective on every finite $2$-subgroup $K$ and proves
$|K|\leq8$. We classify its image using that same
specified injection.

No matrix in $\GL_2(\Q_2)$ has order eight:
an eigenvalue of order eight has minimal polynomial
$X^4+1$ over $\Q_2$, which is irreducible by
Eisenstein after $X\mapsto X+1$.
The same argument excludes higher powers of two.
Every matrix $R$ of order four has minimal
polynomial $X^2+1$, since this polynomial is
irreducible over $\Q_2$ and $R^2\ne I$.
Consequently
\[
                         R^2=-I.                       \tag{D.50d}
\]
An exponent-two subgroup is abelian, because
$(ab)^2=a^2=b^2=1$ implies $ab=ba$.
Its matrices are commuting diagonalizable
involutions over $\Q_2$, and simultaneous
diagonalization embeds it in
$\{\operatorname{diag}(\epsilon_1,\epsilon_2):
\epsilon_i=\pm1\}$. Its order is at most four.
Thus an image of order eight contains an $R$
of order four.

The matrix centralizer of $R$ is the multiplicative
group of the field $\Q_2[R]\cong\Q_2(i)$.
For completeness, choose a nonzero vector $v$:
the vectors $v,Rv$ are independent, since an
eigenvalue in $\Q_2$ would contradict the
irreducibility of $X^2+1$.
A commuting endomorphism is determined by its
value on $v$ and is therefore $aI+bR$.
Every element of the finite $2$-subgroup in this
field has order at most four by the preceding
matrix argument. It is consequently one of
the at most four roots of $X^4-1$.
Since $\langle R\rangle$ already has four elements,
\[
                       C_K(R)=\langle R\rangle
                                                        \tag{D.50e}
\]
for a subgroup of order eight.

Choose $S\in K\setminus\langle R\rangle$.
The index-two subgroup $\langle R\rangle$ is normal;
by (D.50e) conjugation by $S$ induces its nonidentity
automorphism. Hence
\[
          SRS^{-1}=R^{-1},\qquad
                    S^2\in\{I,R^2\}.                  \tag{D.50f}
\]
The second inclusion follows because $S^2\in
\langle R\rangle$ and $S$ fixes its own square,
whereas it inverts $R$.
The alternative $S^2=I$ gives the dihedral
presentation of order eight. We exclude the
alternative $S^2=R^2=-I$ explicitly.

In the basis $(v,Rv)$ used above one has
\[
 R=\begin{pmatrix}0&-1\\1&0\end{pmatrix}.
\]
The equation $SR=-RS$ forces
\[
             S=\begin{pmatrix}a&b\\b&-a\end{pmatrix},
             \qquad S^2=(a^2+b^2)I.
                                                        \tag{D.50g}
\]
Thus the second alternative would require
$a^2+b^2=-1$ in $\Q_2$.
Here is the exact valuation obstruction.
Put $m=\min(v_2(a),v_2(b))$ and write
$a=2^m u$, $b=2^m v$, with $u,v\in\Z_2$
and at least one a unit. If both are units,
$u^2+v^2\equiv2\pmod8$, so this sum has
valuation one. If exactly one is a unit,
the sum has valuation zero.
On the other hand the equation says
\[
                      u^2+v^2=-2^{-2m};
                                                        \tag{D.50h}
\]
the valuation of its right side is even.
Both-unit values are therefore impossible.
The other case forces $m=0$; reduction modulo
four then gives $1=-1$, also impossible.
Zero entries cause no exception to this
argument. This excludes the quaternion
alternative and leaves only $D_8$ at order
eight. Orders four, two, and one give the
remaining groups in (D.50a).

We specify actual full-group realizations.
On $\widehat F(x,y)$ define
\[
 \begin{array}{c|cc}
          &x&y\\ \hline
 r        &y&x^{-1}\\
 s        &y&x .
 \end{array}                                            \tag{D.50i}
\]
The free universal property gives continuous
endomorphisms, and the displayed relations
$r^4=s^2=1$, $srs=r^{-1}$, verified on
$x,y$, give their inverses and make them
automorphisms.
For $c=[x,y]$ direct multiplication gives
\[
                   r(c)=x^{-1}cx,\qquad
                   s(c)=c^{-1}.                        \tag{D.50j}
\]
Thus their outer classes lie in $T$.
Their homology matrices are
\[
             R=\begin{pmatrix}0&-1\\1&0\end{pmatrix},
       \qquad S=\begin{pmatrix}0&1\\1&0\end{pmatrix},
                                                        \tag{D.50k}
\]
and the eight matrices $R^iS^j$
($0\leq i<4$, $0\leq j<2$) are distinct.
These maps therefore exhibit an embedded
$D_8$ in $T$. The subgroups
$\langle r\rangle$, $\langle r^2,s\rangle$,
$\langle s\rangle$, and the identity realize
$C_4$, $C_2\times C_2$, $C_2$, and $1$.

We finish with the asserted multiplier data.
D.43 makes every involution's multiplier a
global sign. A positive nonidentity involution
has determinant $1$ in (D.50c). Its order-two
matrix is diagonalizable in characteristic
zero, and faithfulness excludes $I$; its matrix
is consequently $-I$.
There is at most one such involution in any
finite $2$-subgroup.
For $C_2\times C_2$, simultaneous diagonalization
identifies its full image with the four sign
matrices above, giving exactly one positive
and two negative involutions.
For $D_8$, the square of a rotation is
positive by (D.50d), so every other involution
is negative. The reflections generate $D_8$,
and a rotation is a product of two reflections.
The homomorphic multiplier is therefore
globally $1$ on the rotations and globally
$-1$ on the reflections.

Finally let $a$ have order four in an arbitrary
cyclic subgroup. Its matrix in (D.50c) has
minimal polynomial $X^2+1$ and determinant one.
Thus $\lambda(a)_2=1$ and $\lambda(a^2)_2=1$.
D.43 applied to the involution $a^2$ gives
$\lambda(a^2)=1$ globally. The multiplier
homomorphism then gives (D.50b).
No assertion of conjugacy follows from this
matrix classification.

The finite certificate
\path{code/part_c/torus_two_subgroups_d50.py}
verifies the full-word maps (D.50i)--(D.50j),
their eight distinct matrices, and the unit
square computations modulo eight used in
(D.50h). The passage from that finite
calculation to all of $\Q_2$ is the valuation
argument just given.
\end{proof}

\begin{claim}{D.51}{Proved}{The second index-three kernel and its three peripheral normalizers}
Let $a$ have order three in the full punctured-torus
loop-outer group, put $F=\widehat F(x,y)$,
$c=[x,y]$, and let $E$ be its centre-free extension.
Besides $F$, there is exactly one other torsion-free
normal subgroup $U$ of index three in $E$.
Write $V=F\cap U$. Then
\[
 \begin{gathered}
 [F:V]=[U:V]=3,\qquad E/V\cong C_3^2,\\
 V=\widehat{\pi_1(S_{1,3})},\qquad
 U\ \hbox{is projective},\qquad
 U^{(3)}\cong\widehat F^{(3)}_2 .
 \end{gathered}                                         \tag{D.51a}
\]
The cusp normalizer $N=N_E(\langle c\rangle)$
is contained in $U$. The three cusp subgroups
$L_0,L_1,L_2$ of the displayed degree-three
cover have normalizers $N_i=N_U(L_i)$ with
the exact sequences
\[
              1\longrightarrow L_i
                 \longrightarrow N_i
                 \longrightarrow U/V\longrightarrow1.
                                                        \tag{D.51b}
\]
They are three distinct $U$-conjugacy classes,
cyclically permuted by $E/U$. All have the
given multiplier $\lambda$ under the same
generator label in $U/V$.

For a prime $q\ne3$ at which $\lambda_q$ is
a nonidentity cube root of unity, the
$U/V$-character multiplicities of
$H_1(V,\Q_q)$, ordered as $(1,\lambda_q,\lambda_q^2)$,
are exactly one of
\[
                         (1,3,0),\qquad (0,2,2).
                                                        \tag{D.51c}
\]
Correspondingly $U^{(q)}$ is free pro-$q$ of rank
one or zero. These are constraints on the actual
subgroup $U$; they do not identify it with a full
free profinite group.
\end{claim}
\passage{Section 30, folios 194-196, scans 383-387; Section 31, folios 198-199 verso, scans 391-394; Part D(b), full loop-outer finite actions in type (1,1)}

\begin{proof}
\emph{The section and cusp lines in the actual abelian quotient.}
Use D.27 to choose an actual order-three representative
$\alpha$ and write $E=F\rtimes_\alpha\langle s\rangle$.
Let $u:E\to\F_3$ be its quotient character, with $u(s)=1$.
D.48 gives exactly two full section classes,
$S_0=\langle s\rangle$ and $S_1$, and gives
$E^{\rm ab,(3)}\cong C_3^2$.

Let $M$ be the reduction of the homology matrix of
$\alpha$ modulo three. By D.48 it is a nonidentity
unipotent matrix with $(M-I)^2=0$.
On $\F_3^2$ write
$\omega((a,b),(a',b'))=ab'-ba'$.
Choose $w$ with $e=(M-I)w\ne0$.
Then $(e,w)$ is a basis, $Me=e$, $Mw=w+e$,
and $d=\omega(e,w)\ne0$.
The linear functional
\[
                  \chi(r)=\omega(e,r)/d
                                                        \tag{D.51d}
\]
is $M$-invariant. It gives a character of $F$,
which we extend to $E$ by $\chi(s)=0$.
Thus $(u,\chi):E\to\F_3^2$ is the actual
abelian quotient above.

The generator of $S_1$ labelled by $u=1$
has $\chi$-value $b_1\ne0$.
Here is the required separation, using the maps
already proved in D.48 rather than an assumption
about section labels.
The Bockstein of $\chi$ has nonzero high periodic
multiples. On a section normalizer its restriction
is the Bockstein of $\chi$ restricted to the section:
the fixed-group first cohomology vanishes by D.48.
If $\chi$ vanished on both section classes,
all those multiples would have zero full Brown
restriction. This contradicts D.48f.
It follows that $b_1\ne0$.

To locate the cusp line, use the characteristic
class-two exponent-three quotient
\[
 Q=\F_3^2\times\F_3,\qquad
 (r,z)(r',z')=(r+r',z+z'+\tfrac12\omega(r,r')) .
                                                        \tag{D.51e}
\]
The images of $x,y$ are the two standard vectors
with central coordinate zero, and $c$ maps to
$Z=(0,1)$.
Since $\lambda_3=1$, the induced automorphism is
\[
               \theta(r,z)=(Mr,z+\ell(r))
                                                        \tag{D.51f}
\]
for a linear functional $\ell$.
For an element $(r,z)s$ of $Q\rtimes_\theta C_3$
with quotient label one, direct multiplication gives
\[
       ((r,z)s)^3
          =Z^{\,d\chi(r)^2+\ell(e)\chi(r)} .
                                                        \tag{D.51g}
\]
To check the formula, write $r=be+tw$, so
$\chi(r)=t$. The successive vector parts are
$r,r+te,r+2te$, whose sum is zero.
The sum of the three central correction terms
is $\ell((2I+M)r)=t\ell(e)$.
The three alternating cross terms, multiplied
by $1/2$, sum to $dt^2$ in $\F_3$.
This proves (D.51g) for every element, including
the arbitrary central coordinate $z$.

The image of the actual generator of $S_1$
has cube one. Equation (D.51g) and $b_1\ne0$
therefore give $\ell(e)=-db_1$.
Thus its cube polynomial is
\[
                           db(b-b_1).
                                                        \tag{D.51h}
\]
Choose $v\in N$ with $u(v)=1$, as in D.48.
It satisfies $v^3=c^\kappa$ with
$\kappa_3$ a unit. Its image in (D.51e)
has nonzero cube exponent. The two roots of
(D.51h) are $0$ and $b_1$; the only other
element of $\F_3$ is $-b_1$. Consequently
\[
             (u,\chi)(S_0)=\langle(1,0)\rangle,\quad
             (u,\chi)(S_1)=\langle(1,b_1)\rangle,\quad
             (u,\chi)(N)=\langle(1,-b_1)\rangle .
                                                        \tag{D.51i}
\]
The last equality uses that $N$ is generated
topologically by $v$ and $\langle c\rangle$,
and both characters kill $c$.

\emph{The full kernels and peripheral labels.}
Define the continuous characters and subgroups
\[
           \psi=\chi+b_1u,\qquad
           U=\ker\psi,\qquad V=\ker u\cap\ker\chi .
                                                        \tag{D.51j}
\]
These are actual subgroups of $E$.
The character $\psi$ is nonzero on both
section classes in (D.51i), and it kills $N$.
Every nonidentity finite subgroup of $E$
is one of the order-three section groups,
up to conjugacy. Hence $U$ is torsion-free.
The four lines in $E^{\rm ab,(3)}=\F_3^2$
give all normal index-three subgroups of $E$.
Two contain respectively $S_0$ and $S_1$.
The other two are exactly $\ker u=F$ and
$\ker\psi=U$, proving the uniqueness assertion.
Also $E=U\rtimes S_0$, since $\psi|_{S_0}$
is an isomorphism onto $\F_3$.

The restriction of $\chi$ to $F$ is nonzero
and kills its commutator cusp.
Choose a word $f$ in the dense discrete
group $F(x,y)$ with $\chi(f)=1$.
The degree-three covering of the punctured
torus corresponding to $\ker(\chi|_{F(x,y)})$
has three boundary lifts of degree one.
Its Euler characteristic is $-3$; thus its
type is $(1,3)$, and its fundamental group
has rank four.
The induced profinite subgroup is exactly $V$,
with cusp representatives
\[
          L_i=f^i\langle c\rangle f^{-i},
          \qquad
          N_i=f^i N f^{-i}\quad(0\leq i<3).
                                                        \tag{D.51k}
\]
The full peripheral normalizer assertion A.7
gives $N_E(L_i)=N_i$.
Since $N\subset U$ and $U$ is normal in $E$,
all $N_i$ lie in $U$. Their kernels under
$U\to U/V$ are exactly $L_i$, because
$N_i\cap F=L_i$; their images are all of
$U/V$, because $u(v)=1$.
This proves (D.51b), with the common generator
label supplied by $u$.

For the class assertions, the actual quotient
$(u,\chi)$ identifies $E/V$ with $\F_3^2$,
and the image of every $N_i$ is the same
line $U/V$.
The three $V$-conjugacy classes of the cusp
lifts are therefore the three cosets of this
line in $E/V$. The subgroup $U/V$ fixes each
class, so they remain distinct under $U$.
The quotient $E/U$ acts regularly on these
three cosets, proving its asserted cyclic
permutation. This also proves that conjugation
by a generator of $U/V$ fixes each cusp
class with multiplier $\lambda$.

The open subgroup $V$ of $U$ is free profinite
of rank four and $U$ is torsion-free.
Serre's finite-index cohomological-dimension
theorem, used in D.27, makes $U$ projective.
Its maximal pro-three quotient is projective
as a pro-three group, by the finite embedding
problem characterization, and is therefore
free pro-three.
The pro-three topology on $V$ is exactly the
one induced from $U$: taking the intersection
of the three conjugates of a normal
three-power-index subgroup of $V$ produces
one normal in $U$, still of three-power index.
Thus $V^{(3)}$ embeds as an open subgroup
of index three in $U^{(3)}$.
Schreier's formula now reads
\[
                 4=1+3(\operatorname{rank}U^{(3)}-1),
\]
and gives rank two. This proves (D.51a).

\emph{The exact character restriction at another prime.}
The three oriented boundary classes of $V$
have their sole relation given by their sum.
They form a two-dimensional submodule of
$H_1(V,\Q_q)$ on which $U/V$ acts by
$\lambda_q$.
Filling the three boundaries gives the compact
degree-three cover of the original torus.
Its map on homology is the inclusion of a
rank-two sublattice of index three, hence an
isomorphism over $\Q_q$, and over $\Z_q$
when $q\ne3$.
Conjugation by $v$ projects to the original
outer action $a$ on the base; its inner
factor acts trivially on that base homology.
Therefore, if $A_q$ denotes the original
matrix, the exact boundary sequence and
averaging over $C_3$ give
\[
       H_1(V,\Q_q)\cong
                  \lambda_q^{\oplus2}
                    \oplus(\Q_q^2,A_q)
                 \quad\hbox{as }U/V\hbox{-modules}.
                                                        \tag{D.51l}
\]
This comparison is induced by the stated
boundary inclusion and covering homology map.

At the specified prime, $A_q^3=I$ and
$\det A_q=\lambda_q$.
The three cube roots $1,\lambda_q,\lambda_q^2$
are distinct in $\Q_q$. Diagonalization gives
exactly the two eigenvalue multisets
\[
                 \{1,\lambda_q\},\qquad
                 \{\lambda_q^2,\lambda_q^2\}.
                                                        \tag{D.51m}
\]
Adding the two boundary characters in (D.51l)
proves (D.51c).
For $q\ne3$, the same decomposition holds
over $\Z_q$: the averaging idempotents have
unit denominator three, and the cube-root
differences are units.
Hochschild--Serre for $V\triangleleft U$ then
identifies $H^1(U,\F_q)$ with the invariant
subspace of $H^1(V,\F_q)$, of dimension one
or zero respectively.
The projectivity of $U$ makes its maximal
pro-$q$ quotient free pro-$q$, with this
number of generators. This proves the last
assertion.

The program
\path{code/part_c/torus_order3_sibling_d51.py}
checks (D.51g) for every central-linear
automorphism above every nonidentity unipotent
matrix over $\F_3$, the resulting section and
cusp lines, and the finite eigenvalue patterns
in (D.51m). These finite calculations validate
the displayed formulas. The full subgroup,
topology, and projectivity conclusions are
proved by the maps above.
\end{proof}

\paragraph{Exact remaining recognition problem.}
The construction supplies a specified torsion-free
projective group $U$ with an open full free subgroup
of rank four and three specified normalizers.
It supplies no identification of $U$ with
$\widehat F_2$.
At the primes in (D.51c), its maximal pro-$q$
quotient has rank zero or one, whereas the
maximal pro-three quotient has rank two.
The existence or exclusion of such a full
$U$ with these three normalizers remains the
step needed for this approach to determine
every torus order-three multiplier.

\begin{claim}{D.52}{Proved}{An abstract allowed list for every finite full torus subgroup}
Let $T$ be the full once-punctured-torus loop-outer
group of D.50. Every finite subgroup of $T$ has
one of the following twelve abstract types:
\[
\begin{gathered}
1,\ C_2,\ C_3,\ C_4,\ C_2\times C_2,\ C_6,\ S_3,\ D_8,\\
C_{12},\ D_{12},\ \operatorname{Dic}_{12},\ D_{24}.
\end{gathered}                                             \tag{D.52a}
\]
Here $D_{2n}$ has order $2n$, and
$\operatorname{Dic}_{12}=C_3\rtimes C_4$, with
a generator of $C_4$ inverting $C_3$.
The nine types obtained by omitting
$C_{12},\operatorname{Dic}_{12},D_{24}$ all occur.
Occurrence of those three remaining types is
not determined by this allowed list.
The subsequent full covering arguments in D.54--D.55
exclude all three.

Every subgroup of type $S_3,D_{12}$ or $D_{24}$
has the global usual reflection multiplier:
rotations have multiplier $1$ and reflections
have multiplier $-1$.
In a subgroup of type $\operatorname{Dic}_{12}$,
the subgroup of order three has multiplier $1$,
and every generator $b$ of a cyclic subgroup
of order four satisfies
\[
       \lambda(b)^2=1,\qquad
       \lambda(b)_2=1,\qquad\lambda(b)_3=-1.              \tag{D.52b}
\]
For a subgroup of type $C_{12}$, its order-four
generator instead has $\lambda(b)_2=\lambda(b)_3=1$,
and still has $\lambda(b)^2=1$.
These restrictions retain the unresolved cyclic
order-three and order-four multiplier questions.
\end{claim}
\passage{Section 30, folios 194-196, scans 383-387; Section 31, folios 198-199 verso, scans 391-394; Part D(b), full loop-outer finite actions in type (1,1)}

\begin{proof}
\emph{The two primary representations and their actual kernels.}
For a finite subgroup $G\leq T$, D.33 gives
\[
 |G|\mid24,\qquad
 A_p:G\longrightarrow\GL_2(\Z_p),\qquad
 \det A_p(g)=\lambda(g)_p.                              \tag{D.52c}
\]
It makes $A_p$ injective on every finite
$p$-subgroup. Put $K=\ker(A_2|_G)$.
An element of order two in $K$ would contradict
that injection. Cauchy's theorem and (D.52c)
therefore give
\[
                      K=1\quad\hbox{or}\quad K\cong C_3.
                                                               \tag{D.52d}
\]
We classify the actual image $H=A_2(G)$ and
then the extension by this specified kernel.
We do not identify $G$ with a representation
at one prime when the kernel is nontrivial.

\emph{Finite images over $\Q_2$.}
The two-subgroups of $H$ have the types in D.50.
Indeed a Sylow two-subgroup of the preimage
of any such subgroup maps isomorphically onto
it, since the kernel has odd order.
The matrix arguments in D.50 also explicitly
exclude the quaternion group of order eight.

Suppose $3$ divides $|H|$.
Its number of Sylow three-subgroups is either
one or four. We exclude four.
Conjugation gives a transitive homomorphism
$H\to S_4$. Its kernel has order at most six.
If that kernel had order divisible by three,
its unique Sylow three-subgroup would be
characteristic in the kernel and normal in $H$,
contrary to the assumed four Sylow subgroups.
Thus the transitive image has order divisible
by both four and three, and is $A_4$ or $S_4$.
The kernel consequently has order one or two.

There is no faithful two-dimensional
representation of $A_4$ over $\Q_2$.
Its normal Klein four subgroup would be
simultaneously diagonalizable and, by
faithfulness, would be the full group
of diagonal sign matrices. One of its
three involutions would therefore be $-I$.
Conjugation fixes this scalar, whereas the
order-three elements of $A_4$ cyclically
permute all three involutions.
This is a contradiction.
It excludes the cases with trivial kernel,
whose images are $A_4$ or $S_4$.
In the remaining case $|H|=24$ and the image
is $A_4$ with a kernel of order two.
The preimage $P$ of its Klein four subgroup
is a normal subgroup of order eight, hence
is $D_8$ by D.50. But
$|\operatorname{Aut}(D_8)|=8$: a generator
of its rotation subgroup has two possible
images, and a chosen reflection has four.
Every element of order three therefore
centralizes $P$. Multiplication by a Sylow
three-subgroup gives $H=P\times C_3$, which
has only one Sylow three-subgroup.
This contradiction proves that $H$ has a
normal subgroup $P_3$ of order three.

Let $R$ generate $P_3$.
Its minimal polynomial is
$X^2+X+1$, irreducible over $\Q_2$ by
reduction modulo two.
The basis $(v,Rv)$ for any nonzero vector $v$
identifies the matrix centralizer with
\[
 L_2^\times,\qquad L_2=\Q_2(\zeta_3).
                                                               \tag{D.52e}
\]
Under this identification, a commuting
endomorphism is multiplication by $a+b\zeta_3$.
The full normalizer consists exactly of these
multiplication maps and their compositions
with $\sigma(\zeta_3)=\zeta_3^{-1}$:
conjugation must send $R$ to $R$ or $R^{-1}$,
and division by $\sigma$ in the second case
gives a commuting map.

The field $L_2$ contains no element of order
four. To check this directly, the equation
$(a+b\zeta_3)^2=-1$ has coefficient equations
\[
                   a^2-b^2=-1,\qquad b(2a-b)=0.
                                                               \tag{D.52f}
\]
They require either $a^2=-1$ or $a^2=1/3$
in $\Q_2$. Both are impossible: these are
units congruent respectively to seven and
three modulo eight, whereas an odd square
is congruent to one.
Since $|H|\mid24$, it follows that
$C=H\cap L_2^\times$ is $C_3$ or $C_6$.
Here its two-part is contained in
$\{1,-1\}$, and its three-part is $P_3$,
so no general classification of local
roots of unity is required.

If $H=C$, these are the two possible groups.
Otherwise choose $B=a\sigma\in H\setminus C$.
Then
\[
 B^2=N_{L_2/\Q_2}(a)\,I
        \in C\cap\Q_2^\times=\{I,-I\}\cap C.             \tag{D.52g}
\]
The last scalar restriction follows because
a finite-order scalar of order dividing
twenty-four in $\Q_2$ is a sign: there are
no primitive cube or fourth roots in $\Q_2$.
For $C=C_3$, equation (D.52g) gives $S_3$.
For $C=C_6$, its two values give respectively
$D_{12}$ and $\operatorname{Dic}_{12}$.
Thus the images with order divisible by
three have the complete allowed list
\[
                  C_3,\ C_6,\ S_3,\ D_{12},
                             \operatorname{Dic}_{12}.  \tag{D.52h}
\]

\emph{The three-adic normalizer and its determinant.}
We will need a statement about any normal
subgroup $P_3\cong C_3$ of the actual $G$.
The map $A_3$ is injective on $P_3$.
Its generator has minimal polynomial
$X^2+X+1$ over $\Q_3$, irreducible since
its discriminant $-3$ has odd valuation.
The same centralizer and normalizer
calculation identifies
\[
 C_{\GL_2(\Q_3)}(A_3(P_3))=L_3^\times,\qquad
 N_{\GL_2(\Q_3)}(A_3(P_3))=L_3^\times\rtimes\langle\sigma\rangle,
 \quad L_3=\Q_3(\zeta_3).                              \tag{D.52i}
\]
This field also has no element of order four:
(D.52f) would require $a^2=-1$, impossible
modulo three, or $a^2=1/3$, impossible by
valuation. Thus every finite-order element
in its centralizer with order dividing
twenty-four lies in $\mu_6$ and has norm one.

There is also no element of norm $-1$ in $L_3$.
Here is an explicit calculation, including
possible nonintegral coordinates.
The norm of $a+b\zeta_3$ is
$a^2-ab+b^2$. Put
$m=\min(v_3(a),v_3(b))$ and write
$a=3^m u,b=3^m v$, where $u,v\in\Z_3$
and at least one is a unit.
When $u+v\not\equiv0\pmod3$, the norm of
$u+v\zeta_3$ has valuation zero and residue
$(u+v)^2=1$ modulo three.
When $u+v\equiv0\pmod3$, both are units;
writing $u=-v+3k$ gives
\[
                  u^2-uv+v^2
                       =3(v^2-3vk+3k^2),               \tag{D.52j}
\]
of valuation one.
The norm $-1$ would have valuation zero.
The second case cannot give that valuation,
since $2m+1\ne0$.
The first forces $m=0$ and residue one,
whereas $-1$ has residue two.
This proves the asserted norm obstruction.

An element of the finite image $A_3(G)$
outside the centralizer has the form
$B=a\sigma$, and $B^2=N(a)I$ is a
finite-order scalar, hence a sign.
The norm obstruction forces $B^2=I$.
In the basis $(1,\zeta_3)$, the determinant
of $\sigma$ is $-1$, so
$\det B=-N(a)=-1$.
Together with the centralizer calculation,
this proves the exact character identity
\[
 \lambda(g)_3=
 \begin{cases}
  1,&g\text{ centralizes }P_3,\\
 -1,&g\text{ inverts }P_3.
 \end{cases}                                            \tag{D.52k}
\]
The argument uses faithfulness only on
$P_3$, not on all of $G$.

\emph{Reconstructing the actual finite subgroup.}
If $K=1$ in (D.52d), the injection into
$H$ gives precisely D.50 and (D.52h).
Now let $K=P_3\cong C_3$.
A Sylow two-subgroup $Q$ of $G$ has order
$|G|/3$ and maps isomorphically onto $H$.
The actual multiplication map
\[
                  P_3\rtimes Q\longrightarrow G,
                         \qquad (a,q)\longmapsto aq,    \tag{D.52l}
\]
is an isomorphism, with action given by
conjugation. D.50 lists the five possibilities
for $Q$.
For $Q=1,C_2,C_4$, their homomorphisms
to $\operatorname{Aut}(C_3)=C_2$ give
$C_3$; $C_6$ or $S_3$; and $C_{12}$ or
$\operatorname{Dic}_{12}$, respectively.

For $Q=C_2\times C_2$ or $D_8$, D.50
already determines its global multiplier.
Equation (D.52k) therefore makes its
action on $P_3$ exactly the reflection
character of D.50.
Write $a$ for a generator of $P_3$.
For $Q=C_2\times C_2$, write $z$ for its
positive involution and $s$ for a negative
one. The element $az$ has order six,
is inverted by $s$, and generates $a,z$
through $(az)^4=a$ and $(az)^3=z$.
This gives an explicit $D_{12}$ presentation.
For $Q=D_8=\langle r,s\rangle$, where
$r$ is its order-four rotation, the
element $ar$ has order twelve, is inverted
by $s$, and satisfies
\[
                       (ar)^4=a,\qquad (ar)^9=r.
                                                               \tag{D.52m}
\]
It gives the explicit $D_{24}$ presentation.
These actual maps finish the allowed list.

\emph{Multipliers and realized cases.}
In every group just obtained that inverts
its normal $C_3$, the multiplier on that
$C_3$ is one: a conjugation-invariant
homomorphism into the abelian group
$\widehat\Z^\times$ takes an inverted
element to its own inverse, and a value
of order dividing three with square one
is one.
In $S_3,D_{12},D_{24}$ each reflection
inverts $C_3$ and hence has three-adic
multiplier $-1$ by (D.52k).
D.43 makes its full multiplier $-1$.
Every rotation is a product of two
reflections, proving the stated global
reflection character.
For $\operatorname{Dic}_{12}$, an
order-four generator inverts $C_3$,
so (D.52k) gives its negative
three-coordinate. D.50 gives its positive
two-coordinate and global square one.
For $C_{12}$ the same generator centralizes
$C_3$, giving instead the positive
three-coordinate. This proves all the
claimed multiplier restrictions.

D.50 realizes its five types by explicit
full profinite free-group automorphisms.
The other four realized types have the
following simultaneous exact surface model.
On $\R^2/\Z^2$ the integral matrices
\[
       R=\begin{pmatrix}0&-1\\1&1\end{pmatrix},
       \qquad
       S=\begin{pmatrix}0&1\\1&0\end{pmatrix}            \tag{D.52n}
\]
fix the origin and satisfy
$R^6=S^2=I$, $SRS=R^{-1}$.
Their twelve matrices $R^iS^j$ are distinct.
Removing the origin therefore gives an
actual $D_{12}$ action on a once-punctured
torus. Passing to its fundamental group,
then to its full profinite completion and
outer automorphisms, gives a homomorphism
to $T$. A small peripheral circle is
carried to an ellipse of local degree
$\det R=1$ or $\det S=-1$, so this
homomorphism preserves the stated full
cusp with the indicated multiplier.
The induced homology matrices are exactly
$R,S$ in $\GL_2(\widehat\Z)$; their
distinctness proves that the homomorphism
is injective.
The subgroups
$\langle R^2\rangle,\langle R\rangle,
\langle R^2,S\rangle$, and
$\langle R,S\rangle$ realize
$C_3,C_6,S_3,D_{12}$.

The program
\path{code/part_c/torus_finite_types_d52.py}
checks the matrix relations, the local
norm and square residues, the eight
automorphisms of $D_8$, and the explicit
finite semidirect-product identifications
in (D.52l)--(D.52m).
The prime-wise kernel and valuation
arguments above prove the full allowed
list. No assertion of occurrence or
nonoccurrence of its three remaining
types is inferred from those finite checks.
\end{proof}

\begin{claim}{D.53}{Proved}{The unique labelled order-four section class}
Let $a$ have order four in the full punctured-torus
loop-outer group, and let
\[
          1\longrightarrow F\longrightarrow E
             \xrightarrow{\;\pi\;}C_4=\langle a\rangle
             \longrightarrow1,\qquad F=\widehat F_2,
\]
be its centre-free extension. Its homomorphic
sections form exactly one $F$-conjugacy class,
with the generator label $a$ retained.
For an actual lift $b\in E$ with
$b^4=1$ and $\pi(b)=a$, put $B=\langle b\rangle$.
Then
\[
       N_E(B)=C_F(b)\times B,\qquad
                         H^1(C_F(b),\F_2)=0.            \tag{D.53a}
\]
Conjugation by a lift of $a$ fixes exactly one
of the three full section classes for $a^2$
and interchanges the other two.
Consequently $E$ has two conjugacy classes
of subgroups of order two; exactly one of
them is contained in its class of cyclic
subgroups of order four.
\end{claim}
\passage{Section 30, folios 194-196, scans 383-387; Section 31, folios 198-199 verso, scans 391-394; Part D(b), full loop-outer finite actions in type (1,1)}

\begin{proof}
D.46 supplies the actual lift $b$ with its
prescribed label. Put $T=\langle b^2\rangle$
and
\[
               E_2=\pi^{-1}(\langle a^2\rangle)
                      =F\rtimes T.
                                                               \tag{D.53b}
\]
The action of $b^2$ on $H_1(F,\Z_2)$ is
$-I$ by D.46. Thus
\[
                      E_2^{\rm ab,(2)}\cong C_2^3.
                                                               \tag{D.53c}
\]
Use coordinates $(v,u)\in\F_2^2\times\F_2$,
where $u$ is the quotient character and
the first two coordinates are the extended
characters $\chi_1,\chi_2$ of D.46.
The three full section classes of $E_2$
give three distinct points of the affine
plane $u=1$.
Indeed, the degree-zero localized algebra
in D.46g--i has dimension three, is generated
by $\chi_1/u,\chi_2/u$, and restriction to
the three section normalizers identifies
it with $\F_2^3$.
Each restriction is evaluation at
$(\chi_1(t),\chi_2(t))$ for the nonidentity
element $t$ of that section.
If two such points agreed, the two
restriction coordinates would agree on
both generators and hence on the entire
algebra, contradicting the stated product
isomorphism.

Let $R=A_2(a)\in\GL_2(\Z_2)$.
D.50 gives $R^2=-I$.
Moreover $R\not\equiv I\pmod2$.
For otherwise $R=I+2D$ with an integral
matrix $D$, so $R^2=-I$ would give
$I+2D+2D^2=0$, an impossibility modulo two.
Conjugation by $b$ on (D.53c) fixes $u$
and induces the actual affine map
\[
               v\longmapsto\overline R v+d
                        \quad\text{on }u=1             \tag{D.53d}
\]
for a vector $d\in\F_2^2$.
The linear part follows from the
restriction of conjugation to the
specified subgroup $F$; the translation
is the image of a chosen point of $u=1$.
Conjugation by an element of $E_2$
acts trivially on the three
$F$-conjugacy classes, as shown in D.46.
Thus (D.53d) permutes their three points
with order dividing two.
Three distinct points of $\F_2^2$ affinely
span it. Fixing all three would therefore
make this affine map the identity,
contrary to $\overline R\ne I$.
Its permutation is consequently a
transposition, with exactly one fixed class.
The class of $T=\langle b^2\rangle$ is
that fixed class.

Put $C=C_F(T)$ and $N=N_E(T)$.
Since $b$ normalizes $T$, the image of $N$
under $\pi$ is all of $C_4$.
Its kernel is exactly $C$, because an
element of $F$ normalizing $T$ centralizes
its nonidentity element.
The chosen section gives the actual
semidirect-product identification
\[
          N=C\rtimes\langle b\rangle,\qquad
          T\subseteq Z(N),\qquad
          J:=N/T=C\rtimes\langle\overline b\rangle,
                    \quad|\overline b|=2.              \tag{D.53e}
\]
In the last expression the embedding of
$C$ is injective since $C\cap T=1$.
Conjugation by $b$ on $C$ has square the
identity because $b^2$ centralizes $C$.

As a closed subgroup of the full free
profinite group $F$, the group $C$ is
projective and has $\operatorname{cd}_2(C)\leq1$.
D.46 gives $H^1(C,\F_2)=0$.
All its positive-degree mod-two cohomology
therefore vanishes.
Hochschild--Serre for the last extension
in (D.53e) identifies inflation as
\[
               H^*(C_2,\F_2)\xrightarrow{\;\cong\;}
                                  H^*(J,\F_2).         \tag{D.53f}
\]
This uses the displayed full group $C$,
not an identification of $C$ with its
maximal pro-two quotient.

We spell out the full Brown hypotheses
for this new normalizer quotient.
The subgroup $C$ is open in $J$ and has
finite $2$-cohomological dimension.
It is torsion-free. Every finite subgroup
of $J$ consequently injects into $J/C=C_2$;
in particular $J$ has no subgroup $C_2^2$.
Theorem~2 of
\cite[\S1.1]{A-KochloukovaZalesskii2007}
therefore supplies, in degrees $n>1$,
the dense restriction embedding into the
product over the conjugacy classes of
order-two subgroups of $J$.
By (D.53f) its source has dimension one.
Each product factor is nonzero: for a
section $S$ of $J\to C_2$, its normalizer is
\[
                        C_C(S)\times S,
                                                               \tag{D.53g}
\]
and the restriction from the $S$ factor
already supplies one dimension.
Thus the product has finitely many
factors and the finite dense image is
the entire product.
Since $\operatorname{cd}_2(C_C(S))\leq1$,
the dimensions give the exact equality
\[
             1=\sum_{[S]}
                       \bigl(1+\dim H^1(C_C(S),\F_2)\bigr).
                                                               \tag{D.53h}
\]
There is a section $\langle\overline b\rangle$.
It follows that it is the unique
conjugacy class and that
$H^1(C_C(\overline b),\F_2)=0$.
These are also $C$-conjugacy classes:
an element of $J$ can be multiplied by
the appropriate element of a given
section to leave a conjugator in $C$.

Now let $B'$ be any subgroup of order
four in $E$. Its intersection with $F$
is trivial, so it is cyclic and $\pi|_{B'}$
is an isomorphism onto $C_4$.
Choose its unique generator $b'$ with
$\pi(b')=a$. Its square generates an
order-two section $T'$ of $E_2$.
Conjugation by $b'$ has the same action
on the section classes of $E_2$ as
conjugation by $b$, because
$b'b^{-1}\in F$.
The class of $T'$ is fixed, and is
therefore the unique class of $T$.
An $F$-conjugation puts $T'$ equal to $T$,
and then $B'\subseteq N_E(T)=N$.
Its image in $J=N/T$ has order two.
By (D.53h), a further conjugation by
an element of $C$ puts that image equal
to $\langle\overline b\rangle$.

The exact inverse image of this subgroup
under $N\to N/T$ has order four,
because its kernel is the specified
subgroup $T$ of order two.
It contains $\langle b\rangle$ and
hence equals $\langle b\rangle$.
Thus the conjugated $B'$ equals $B$.
All conjugations used lie in $F$ and
preserve $\pi$, so its labelled generator
$b'$ is sent to $b$, proving uniqueness
of the homomorphic section with labels.

The fixed group $C_C(\overline b)$ in
(D.53h) is exactly $C_F(b)$.
For an element of $C$ fixed by
$\overline b$, its commutator with $b$
lies in $T$ and also in $C$, and
$T\cap C=1$.
The converse is immediate, and an
element commuting with $b$ already
centralizes $T$.
This proves the cohomology statement
in (D.53a).
Finally, an element normalizing $B$
must fix its generator $b$: conjugation
preserves its label in the abelian
quotient $C_4$, and $\pi|_B$ is injective.
Multiplication by the matching power
of $b$ then leaves an element of
$C_F(b)$, proving the direct-product
normalizer in (D.53a).

The two orbits of the transposition
in (D.53d) are exactly the two
$E$-conjugacy classes of order-two
subgroups. Only its fixed class can
be the square of a subgroup of order
four, and $B$ realizes that containment.
This proves the final assertion.

The finite program
\path{code/part_c/torus_order4_sections_d53.py}
checks the mod-four matrix obstruction
and every invariant three-point affine
configuration above a nonidentity
order-two matrix in $\GL_2(\F_2)$.
The normalizer and restriction arguments
prove the full section result.
No conjugacy of the outer action $a$
to a prescribed geometric action is
asserted here.
\end{proof}

\begin{claim}{D.54}{Proved}{No fixed character for a cubic torus action}
Let $a$ have order three in the full
once-punctured-torus loop-outer group.
For every prime $q\ne3$, it fixes no nonzero
continuous character
\[
                         \chi:F\longrightarrow C_q,
                      \qquad F=\widehat F(x,y).         \tag{D.54a}
\]
Consequently the reduction of $A_q(a)$ modulo
$q$ has no eigenvalue one.
For every prime $p$, the endomorphism
$A_p(a)-I$ of $\Z_p^2$ is injective.
The full fixed group of every actual
order-three representative is therefore
contained in the full commutator subgroup
$F'$.

In particular, $A_2$ is injective on every
finite subgroup of the full torus group.
The allowed list of D.52 loses
$C_{12}$ and $D_{24}$.
At a nontrivial cubic multiplier
coordinate $q\ne3$, the remaining
homology matrix is exactly
\[
                         A_q(a)=\lambda(a)_q^2 I.
                                                               \tag{D.54b}
\]
This statement does not exclude the
remaining scalar cubic multiplier case.
\end{claim}
\passage{Section 30, folios 194-196, scans 383-387; Section 31, folios 198-199 verso, scans 391-394; Part D(b), full loop-outer finite actions in type (1,1)}

\begin{proof}
Use D.27 to choose an actual representative
$\alpha$ of order three and put
$E=F\rtimes_\alpha\langle s\rangle$,
where $s$ has quotient label $a$.
Write $c=[x,y]$ and $L=\langle c\rangle$.
D.48 supplies a labelled cusp-normalizer
element $v$ with
\[
       v\in N_E(L),\qquad v^3=c^\kappa,\qquad
       vcv^{-1}=c^\lambda,\qquad
                         \kappa_3\in\Z_3^\times.        \tag{D.54c}
\]
It also gives $\lambda_3=1$ and
\[
                  I+A_3(a)+A_3(a)^2=0
                         \quad\text{on }\Z_3^2.         \tag{D.54d}
\]

Suppose $\chi$ in (D.54a) were nonzero
and invariant. Its extension
\[
       \widetilde\chi:E\longrightarrow C_q,\qquad
              \widetilde\chi(hs^j)=\chi(h)
                                                               \tag{D.54e}
\]
is an actual continuous homomorphism.
Put
\[
               N=\ker\chi,\qquad
               E_N=\ker\widetilde\chi=N\rtimes\langle s\rangle .
                                                               \tag{D.54f}
\]
Both have their indicated index $q$
in $F$ and $E$. Since $\chi(c)=0$,
(D.54c) gives
$3\widetilde\chi(v)=0$. The prime $q$
differs from three, so $v\in E_N$.
Thus $v=ns$ for a specified $n\in N$.

The restriction of $\chi$ to the dense
discrete free group is onto $C_q$.
Its covering surface is the degree-$q$
cover of the punctured torus with
boundary monodromy zero.
It has $q$ boundary lifts of degree one,
Euler characteristic $-q$, and hence
type $(1,q)$. Its full fundamental-group
completion is exactly $N$, by the
finite-index completion map.
Choose a discrete word $f$ with
$\chi(f)=1$. Its boundary representatives
are
\[
                c_i=f^i c f^{-i}\quad(0\leq i<q).
                                                               \tag{D.54g}
\]
Their homology classes have exactly the
single relation $\sum_i[c_i]=0$.
In particular each $[c_i]$ is nonzero
in $H_1(N,\F_3)$, because $q\geq2$.
Equivalently, omitting any other boundary
from the usual genus-one presentation
puts $c_i$ into a free basis.

The action on every boundary class has
multiplier $\lambda$ and fixes its label.
This follows from the actual $v\in E_N$:
the element
$n_i=(vf^iv^{-1})f^{-i}$ belongs to $N$,
and
$vc_iv^{-1}=n_i c_i^\lambda n_i^{-1}$.
Conjugation by $s=n^{-1}v$ differs from
this by an inner automorphism of $N$.
Thus $\alpha$ fixes all $q$ boundary
homology classes over $\Z_3$.

We give the required splitting of
the entire homology module, including
its comparison with the base torus.
Let $D=F/N=C_q$ be the deck group and set
\[
          M=H_1(N,\Z_3),\qquad
                       e=\frac1q\sum_{d\in D}d .
                                                               \tag{D.54h}
\]
The action of $\alpha$ on $M$ commutes
with the deck action: it induces the
identity on $D$, and two lifts of a
deck transformation differ by an inner
automorphism of $N$ on this homology.
Consequently $e$ commutes with $\alpha$.

Filling the $q$ boundaries gives the
compact degree-$q$ torus cover.
Explicitly its lattice is the kernel
$\Lambda\leq\Z^2$ of the character
$\Z^2\to C_q$ induced by $\chi$,
and its map to the base torus is
$\R^2/\Lambda\to\R^2/\Z^2$.
The deck maps are translations and
act trivially on the filled homology.
The boundary submodule
\[
              B_\partial
                =\Z_3[D]/\Z_3\Bigl(\sum_{d\in D}d\Bigr)
                                                               \tag{D.54i}
\]
has no deck-invariant vector, because
$q$ is a unit in $\Z_3$.
The exact boundary sequence and (D.54h)
therefore give
\[
             (1-e)M=B_\partial,\qquad
             eM\cong H_1(\R^2/\Lambda,\Z_3)
                         \cong H_1(\R^2/\Z^2,\Z_3).
                                                               \tag{D.54j}
\]
The last isomorphism is induced by the
displayed covering: its lattice inclusion
has index $q$, a $3$-adic unit.
These maps intertwine the induced
automorphisms, so $\alpha|_{eM}$ is
exactly $A_3(a)$.
On $(1-e)M$ it is the identity by
(D.54g). Hence
\[
                  I+\alpha+\alpha^2=0
                       \quad\text{on }H_1(N,\F_3):
                                                               \tag{D.54k}
\]
on the first summand this is (D.54d),
and on the boundary summand the norm
is multiplication by three.

The cube of the actual element $ns$
now gives the contradiction:
\[
 \begin{split}
  [v^3]
      &=[n\,\alpha(n)\,\alpha^2(n)]\\
      &=(I+\alpha+\alpha^2)[n]=0
                       \quad\text{in }H_1(N,\F_3),\\
  [c^\kappa]&=(\kappa_3\bmod3)[c]\ne0 .
 \end{split}                                             \tag{D.54l}
\]
Equations (D.54c) and (D.54l) concern
the same two elements of $N$.
Thus no invariant character (D.54a)
exists.

For $q\ne3$, a matrix and its transpose
have the same rank, so a one-eigenvector
would supply a nonzero invariant
character. It has just been excluded.
In particular $A_q(a)-I$ is invertible
over $\Z_q$.
At $q=3$ its determinant is three,
by D.48, so it is still injective.
If $h$ is fixed by an actual representative,
its image in every $\Z_p^2$ is therefore
zero. The common kernel of the actual
abelianization map
$F\to\widehat\Z^2=\prod_p\Z_p^2$
is $F'$, proving the fixed-group assertion.

At a nontrivial cubic multiplier
coordinate the two eigenvalue multisets
were exactly
$\{1,\lambda_q\}$ and
$\{\lambda_q^2,\lambda_q^2\}$ in D.51.
The first is now excluded.
A finite-order matrix in characteristic
zero is diagonalizable, so the second
gives the actual scalar matrix
(D.54b).

Finally let $G$ be any finite torus
subgroup. D.33 gives $|G|\mid24$.
A nontrivial kernel of $A_2|_G$ would
contain a prime-order element.
Its prime can only be two or three.
D.33 excludes an element of order two,
and (D.54a), with $q=2$, excludes
an element of order three acting
trivially on two-adic homology.
Thus $A_2|_G$ is injective.
The faithful-image part (D.52h), together
with D.50, now excludes $C_{12}$ and
$D_{24}$ from that allowed list.

The program
\path{code/part_c/torus_cover_norm_d54_d55.py}
builds the cyclic covering homology,
its individual boundary vectors and
the deck-average splitting for several
primes $q\ne3$, and verifies the forced
norm contradiction over $\F_3$.
The displayed maps prove it for every
such prime.
\end{proof}

\begin{claim}{D.55}{Proved}{Full primary faithfulness and the nine finite torus types}
Every positive involution $b$ in the full
once-punctured-torus loop-outer group has
\[
                      A_p(b)=-I
                           \quad\text{for every prime }p.
                                                               \tag{D.55a}
\]
Consequently, for each prime $p$, the actual
homology map $A_p$ is injective on every
finite subgroup of that group.
The possible finite abstract subgroup
types are exactly the following nine,
all realized:
\[
       1,\ C_2,\ C_3,\ C_4,\ C_2\times C_2,\ C_6,\
                                      S_3,\ D_8,\ D_{12}.
                                                               \tag{D.55b}
\]
The groups $D_8,D_{12}$ have orders eight
and twelve respectively.

For an arbitrary order-four element $a$,
\[
                A_p(a)^2=-I\quad\text{for every }p.
                                                               \tag{D.55c}
\]
At a negative multiplier coordinate $q$
it is necessarily a scalar matrix
$i_q I$ or $-i_q I$, where $i_q^2=-1$
in $\Q_q$; thus $q\equiv1\pmod4$.
Together with (D.54b), this retains
the precise scalar cubic and quartic
multiplier cases. The abstract
classification (D.55b) does not
settle those scalar cases or the
conjugacy of arbitrary finite subgroups.
\end{claim}
\passage{Section 30, folios 194-196, scans 383-387; Section 31, folios 198-199 verso, scans 391-394; Part D(b), full loop-outer finite actions in type (1,1)}

\begin{proof}
\emph{The positive involution's full cusp normalizer.}
Let $b$ be any positive involution.
D.43 gives its global multiplier $1$.
D.27 gives an actual representative
$\beta$ of order two, and put
$E_b=F\rtimes_\beta\langle s\rangle$.
The first half of D.46 applies to
every such positive involution; it
does not require $b$ to have first
been exhibited as an order-four square.
It gives
\[
        A_2(b)=-I,\qquad
        H^1(C_F(S),\F_2)=0
          \quad\text{for every order-two section }S.
                                                               \tag{D.55d}
\]

We recall the prime-independent part
of the closed-subgroup argument in
D.48. If $C\leq F$ is closed and
$H^1(C,\F_p)=0$, it cannot contain
the full subgroup $c^{e_p\Z_p}$ of
the torus cusp, where $c=[x,y]$.
If the $p$-part of $[F:C]$ were
finite, choose an open neighborhood
$W\geq C$ with maximal such part.
All subsequent indices are prime
to $p$, so transfer makes the
restriction maps on first cohomology
injective. Their direct limit is
$H^1(C,\F_p)$, contradicting the
nonzero first cohomology of the open
full free group $W$.
If the $p$-part is infinite, choose
an open $W\geq C$ of index divisible
by $p$. Its selected cusp lift has
degree prime to $p$, because it
contains the full cusp $\Z_p$.
There must be at least two boundary
lifts: a unique lift would have
degree $[F:W]$.
A boundary character of this finite
cover then restricts nontrivially
to that cusp $\Z_p$ and hence to
$C$, again a contradiction.
These are the actual open-subgroup
restriction and finite-cover maps
used in D.48, with an arbitrary
prime $p$ in place of three.

Apply this with $p=2$ to (D.55d).
An involution in $N_{E_b}(\langle c\rangle)$
would be a section and would centralize
the entire cusp, since the global
multiplier is $1$.
Its full fixed group would therefore
contain the cusp $\Z_2$, which has
just been excluded.
Thus this cusp normalizer is
torsion-free.
Choose its labelled element $v$:
\[
              vcv^{-1}=c,\qquad v^2=c^\kappa.
                                                               \tag{D.55e}
\]
If $\kappa_2$ were even, the equation
$2z=-\kappa$ would be soluble in
$\widehat\Z$: at two by that
divisibility, and at every odd prime
because two is a unit.
Then $(vc^z)^2=1$, and its nontrivial
quotient label would make it an
involution in the cusp normalizer.
Consequently
\[
                             \kappa_2\in\Z_2^\times.
                                                               \tag{D.55f}
\]
The use of the global multiplier
$1$ in (D.55e) is essential to
this exact norm equation.

\emph{The primitive boundary rules out every odd fixed character.}
Let $q$ be odd and suppose a nonzero
$\beta$-invariant character
$\chi:F\to C_q$ existed.
Extend it by $\widetilde\chi(s)=0$
to $E_b$, and put
$N=\ker\chi$ and
$E_N=\ker\widetilde\chi=N\rtimes\langle s\rangle$.
Equation (D.55e) gives
$2\widetilde\chi(v)=0$, so $v\in E_N$.
Write $v=ns$ with $n\in N$.

The exact cover and its maps are
those of (D.54g)--(D.54j), now over
$\Z_2$. It has type $(1,q)$.
The deck-average idempotent
$e=q^{-1}\sum_{C_q}d$ commutes
with $\beta$, identifies $eH_1(N,\Z_2)$
with the base torus homology by the
index-$q$ lattice inclusion, and
has boundary augmentation as its
complement.
On the first summand $\beta=-I$
by (D.55d). On the second it is
the identity: the actual $v\in E_N$
fixes each cusp label and centralizes
its cusp generator up to $N$-conjugacy,
exactly as in (D.54g).
Hence
\[
                  I+\beta=0
                       \quad\text{on }H_1(N,\F_2).
                                                               \tag{D.55g}
\]
The same two elements in (D.55e)
then have incompatible homology:
\[
             [v^2]=[n\,\beta(n)]=0,\qquad
             [c^\kappa]=(\kappa_2\bmod2)[c]\ne0.
                                                               \tag{D.55h}
\]
The second class is nonzero because
$c$ is a primitive boundary in
the cover with $q\geq3$ cusps.
This excludes the invariant character.

At every odd prime $q$ the matrix
$A_q(b)$ squares to $I$ and has
determinant $1$.
Its characteristic-zero eigenvalues
are both $1$ or both $-1$.
The first possibility makes it $I$
and supplies a nonzero invariant
character modulo $q$, contrary
to (D.55h).
The second makes it $-I$.
Together with (D.55d), this proves
(D.55a).

\emph{Faithfulness at every prime and the exact abstract list.}
Fix a prime $p$ and a finite
subgroup $G$ of the full torus group.
By D.33 its order divides twenty-four.
A nontrivial kernel of $A_p|_G$
would contain an element of order
two or three.
A positive involution has image
$-I$ by (D.55a).
A negative involution has determinant
$-1$ at every prime by D.43, so
also cannot lie in the kernel.
An order-three element has nonidentity
image at $p=3$ by D.33, and at
every other prime by D.54's
fixed-character obstruction.
The kernel is therefore trivial,
as claimed.

D.54 already removed $C_{12}$
and $D_{24}$ from D.52.
The only additional type in that
allowed list was $\operatorname{Dic}_{12}$.
Its normal subgroup of order three
would be faithfully represented over
$\Q_3$, and an order-four element
inverting it would also have order
four by the just-proved faithfulness.
But the exact field-normalizer
calculation (D.52i)--(D.52k) proves
that every finite-order matrix
inverting this $C_3$ has square $I$:
the alternative would require norm
$-1$ in $\Q_3(\zeta_3)$.
This contradiction excludes
$\operatorname{Dic}_{12}$.
The nine remaining types all have
the explicit realizations of
D.50 and D.52, proving (D.55b).

For an arbitrary order-four $a$,
its square is a positive involution,
so (D.55a) gives (D.55c).
The full multiplier still has square
one by D.50.
At a negative coordinate the two
eigenvalues of $A_q(a)$ are roots
of $X^2+1$ whose product is $-1$,
so they are equal.
The matrix is diagonalizable in
characteristic zero, hence scalar
$i_q I$ or $-i_q I$.
The two-coordinate is positive by
D.50. For an odd $q$ a square
root of $-1$ in $\Q_q$ reduces to
one in $\F_q$, which requires
$q\equiv1\pmod4$.
This proves the stated quartic
restriction without excluding its
remaining scalar case.

The program
\path{code/part_c/torus_cover_norm_d54_d55.py}
also checks the square norm and
primitive-boundary comparison over
$\F_2$ in explicit cyclic odd-degree
covering homology modules.
The full subgroup, covering and
normalizer arguments above give
the all-prime conclusion.
\end{proof}

\begin{claim}{D.56}{Proved}{The prime sections of a cyclic-six torus action}
Let $a$ have order six in the full once-punctured-torus
loop-outer group of $F=\widehat F(x,y)$, and let
\[
 E=\{(\phi,j)\in\operatorname{Aut}(F)\times C_6:
                                      [\phi]=a^j\}
                                                               \tag{D.56a}
\]
be its centre-free extension, with kernel identified
with $F$ by the inner-automorphism map.
There are exactly three $F$-conjugacy classes of
order-two sections over $\langle a^3\rangle$ and
exactly two $F$-conjugacy classes of order-three
sections over $\langle a^2\rangle$.
Conjugation by an element with quotient label $a^2$
cycles the three former classes; conjugation by an
element with quotient label $a^3$ exchanges the two
latter classes.

Consequently $E$ has exactly one conjugacy class
of subgroups of order two and one of order three,
and these are its only nontrivial finite subgroup
types.  In particular the full extension (D.56a)
does not split, and $a$ has no finite-order
representative in $\operatorname{Aut}(F)$.
For each prime section $S$ of order $p=2$ or $3$,
\[
                 N_E(S)=C_F(S)\times S,
                 \qquad H^1(C_F(S),\F_p)=0.             \tag{D.56b}
\]
The assertions concern these specified extension
and section maps; they do not identify $E$ with
a free profinite product.
\end{claim}
\passage{Section 30, folios 194-196, scans 383-387; Section 31, folios 198-199 verso, scans 391-394; Part D(b), finite actions and actual lifts in type (1,1)}

\begin{proof}
Put
\[
 E_2=\pi^{-1}(\langle a^3\rangle),\qquad
 E_3=\pi^{-1}(\langle a^2\rangle),
 \quad \pi:E\longrightarrow\langle a\rangle .          \tag{D.56c}
\]
The involution $a^3$ is positive.
Indeed D.55 embeds $\langle a\rangle$ faithfully
in $\operatorname{GL}_2(\Q_2)$.
An order-six matrix there has its order-three
part in the irreducible quadratic field
$\Q_2(\zeta_3)$; its commuting involution is
the nonidentity scalar $-I$.
For completeness, the centralizer is that field
because the cubic part has irreducible polynomial
$X^2+X+1$, and its only elements squaring to one
are $1,-1$. Faithfulness excludes the former
for $a^3$.
Thus $a^3$ has determinant one at two, and D.43
gives its global positive sign.
D.27 supplies the actual prime-order sections
of both groups in (D.56c).

\emph{The three involution classes and their affine coordinates.}
Choose an order-two section generator $t$ in
$E_2$ and identify $E_2=F\rtimes\langle t\rangle$.
By D.46 and D.55, its action on two-adic
homology is $-I$, and its maximal elementary
abelian two-quotient has coordinates
\[
 (u,\chi_1,\chi_2):E_2\longrightarrow\F_2^3,
 \quad u(t)=1,\quad\chi_1(t)=\chi_2(t)=0,                \tag{D.56d}
\]
where $\chi_1|_F,\chi_2|_F$ form a basis of
$H^1(F,\F_2)$.
Every section generator has $u=1$, so its
coordinate is a point of the affine plane
$\F_2^2$.
The three $F$-conjugacy classes of D.46 give
three distinct such points.
One can read the distinctness directly from
that proof: Brown restriction identifies the
localized degree-zero ring with $\F_2^3$,
and this ring is generated by $\chi_1/u$ and
$\chi_2/u$. Equal values on two section
classes would prevent these two elements
from generating the corresponding two
coordinate factors.

Conjugation by an element $e\in E$ with
$\pi(e)=a^2$ preserves $u$ and acts on the
plane (D.56d) by
\[
                           r\longmapsto Rr+d,           \tag{D.56e}
\]
where the linear part is the reduction of
$A_2(a^2)$ in the chosen homology coordinates.
D.54, at $q=2$, says that $R-I$ is invertible.
Also $R^3=I$, so $R$ is a nonidentity cubic
matrix over $\F_2$ and $I+R+R^2=0$.
The affine map (D.56e) has one fixed point,
and its other three points form a single
orbit of length three: translate its unique
fixed point to zero and use the corresponding
linear statement.
Its invariant three-point set of section
classes must be that orbit.
This proves the asserted cyclic permutation.
The action on the class set factors through
$E/E_2=C_3$, because conjugation by $E_2$
is already $F$-conjugation on its sections,
as proved in D.46.

\emph{The two cubic classes and their affine coordinates.}
Choose an order-three section generator $s$
in $E_3$ with label $a^2$.
Use the actual abelian quotient of D.48 and
the separation proved in D.51:
\[
 (u_3,\chi):E_3\longrightarrow\F_3^2,
 \qquad u_3(s)=1,\quad\chi(s)=0.                        \tag{D.56f}
\]
The two $F$-conjugacy classes of section
generators with prescribed label $u_3=1$
have $\chi$-values $0$ and $b$, with
$b\ne0$.
For an element $e\in E$ with label $a^3$,
conjugation preserves $u_3$ because $E/F$
is abelian.  On the restriction of $\chi$
to $F$ it acts by $-1$, since D.55 gives
$A_3(a^3)=-I$.
The difference between the conjugated
character and $-\chi$ vanishes on $F$,
so is a scalar multiple of $u_3$.
On the line $u_3=1$ the action consequently
has the exact form
\[
                              r\longmapsto-r+d'.        \tag{D.56g}
\]
Preserving the two-point set $\{0,b\}$
forces $d'=b$: the alternative of fixing
both would give $2b=0$ in $\F_3$.
Thus this action exchanges the two classes.
Conjugation by $E_3$ is $F$-conjugation on
its sections, by D.48, so the action factors
through $E/E_3=C_2$.

\emph{The full finite subgroups and the failure of lifting.}
Every finite subgroup of $E$ injects into
$E/F=C_6$, because $F$ is torsion-free.
It is therefore cyclic of order one, two,
three or six.  A subgroup of order six
would normalize its order-two subgroup.
Its element with quotient label $a^2$
would then fix that subgroup's
$F$-conjugacy class, contradicting
(D.56e).  Thus no subgroup of order six
exists.
The two affine actions show respectively
that all order-two subgroups are conjugate
in $E$ and all order-three subgroups are
conjugate in $E$.

For a chosen section $S$ of order two,
the stabilizer of its $F$-conjugacy class
in $E/F$ is exactly $\langle a^3\rangle$;
for order three it is exactly
$\langle a^2\rangle$.
The image of $N_E(S)$ equals this
stabilizer: a conjugator whose class is
fixed can be multiplied by an element
of $F$ to normalize $S$ itself.
Thus the full normalizer is already
its normalizer in $E_p$.
Equations (D.46h) and (D.48g) identify
it as $C_F(S)\times S$ and give the
cohomology assertion in (D.56b).

The first projection of (D.56a) is
injective and identifies $E$ with the
inverse image of $\langle a\rangle$
in $\operatorname{Aut}(F)$.
Indeed the outer class of the first
coordinate determines its label $j$.
A homomorphic lift of $\langle a\rangle$
would give a subgroup of order six in
this exact inverse image, which has
been excluded.
More generally, a finite-order
representative $\phi$ of $a$ would
satisfy $\phi^6\in\operatorname{Inn}(F)$.
The latter group is isomorphic to the
torsion-free group $F$, so
$\phi^6=1$, making $\phi$ a forbidden
element of order six.
This proves the labelled nonlifting
assertion while preserving the actual
prime-order lifts in (D.56c).

The program
\path{code/part_a/torus_order6_sections_d56.py}
checks every affine cubic action on
$\F_2^2$ with nonidentity cubic linear
part, both nonzero cubic section
separations in (D.56f), and the
corresponding quotient stabilizers.
The extension and cohomological
arguments above establish that these
are the actions on the full section
class sets.
\end{proof}

\begin{claim}{D.57}{Proved}{Full lifts of the noncyclic torus two-groups}
Let $Q$ be a subgroup of the full
once-punctured-torus loop-outer group,
isomorphic to $C_2\times C_2$ or to
$D_8$ of order eight.  Put
\[
 E_Q=\{(\phi,q)\in\operatorname{Aut}(F)\times Q:
                        [\phi]=q\},\qquad F=\widehat F(x,y).
                                                               \tag{D.57a}
\]
The projection $E_Q\to Q$ has a homomorphic
section with exactly its specified outer labels.
For $Q=D_8$ there is exactly one
$F$-conjugacy class of such sections.

For $Q=C_2\times C_2$, let $z$ be its
positive involution and let $r$ be either
of its two negative involutions.  There
are exactly three section classes when
\[
                         A_2(r)\bmod2=I,                \tag{D.57b}
\]
and exactly one otherwise.  This choice
does not depend on which negative
involution is called $r$.
For the image $H$ of every full section,
\[
                  N_{E_Q}(H)=C_F(H)\times H,
                  \qquad H^1(C_F(H),\F_2)=0.            \tag{D.57c}
\]
Together with D.27 and D.46, this supplies
actual lifts for every finite two-subgroup
in the list of D.50.  It is a statement
about labelled automorphism lifts, not
a geometric conjugacy classification.
\end{claim}
\passage{Section 30, folios 194-196, scans 383-387; Section 31, folios 198-199 verso, scans 391-394; Part D(b), finite two-groups and actual labelled lifts in type (1,1)}

\begin{proof}
\emph{The invariant prime-section classes for a four-group.}
Suppose first that $Q=C_2\times C_2$.
By D.50 and D.55 its faithful dyadic
representation has exactly one nonidentity
positive element $z$, with $A_2(z)=-I$;
the two other involutions are negative.
In the inverse image $E_z$ of
$\langle z\rangle$, D.46 supplies exactly
three $F$-conjugacy classes of sections
$T$ of order two, each with
\[
                         H^1(C_F(T),\F_2)=0.             \tag{D.57d}
\]
As in (D.56d), their elementary abelian
coordinates are three distinct points
of the affine plane $\F_2^2$.

Conjugation by an element of $E_Q$
with label $r$ induces an affine map
\[
                           v\longmapsto Mv+d,
                    \qquad M=A_2(r)\bmod2.              \tag{D.57e}
\]
Its square is the identity on the plane:
the square of the conjugating element
lies in $F$, and inner automorphisms
act trivially on the elementary abelian
quotient of $E_z$.
Its three-point set of section classes
is invariant.

If $M=I$, a nonzero translation would
have only two-element orbits, so could
not preserve a three-point set.
Thus $d=0$, and all three classes are
fixed. If $M\ne I$, it is the unique
two-dimensional unipotent type over
$\F_2$, with
\[
          \ker(I+M)=\operatorname{im}(I+M),
                         \quad\dim\ker(I+M)=1.
\]
The square identity gives $(I+M)d=0$,
so (D.57e) has two fixed points and
one orbit of size two.
An invariant three-point subset must
contain that orbit and one fixed
point. Thus exactly one section
class is fixed in this case.
Since $A_2(zr)=-A_2(r)$, the reductions
of the two negative involutions are
equal, proving the independence in
(D.57b).

Choose any fixed class and one of
its section groups $T$.
Its normalizer $N=N_{E_Q}(T)$ maps
onto $Q$: adjust a conjugator with
label $r$ by an element of $F$ to
normalize $T$, and use $T$ itself
for the label $z$.
Its kernel is $C=C_F(T)$, since an
element of $F$ normalizing the
order-two group $T$ centralizes it.
Moreover $T$ is central in $N$.
The quotient $J=N/T$ therefore has
the exact sequence
\[
                  1\longrightarrow C\longrightarrow J
                     \longrightarrow Q/\langle z\rangle
                     \longrightarrow1,
                  \qquad Q/\langle z\rangle=C_2.        \tag{D.57f}
\]
The inclusion of $C$ in this quotient
is injective, because $C\cap T=1$.

\emph{The invariant order-four class for a dihedral group.}
Now suppose that $Q=D_8$ and let
$R\cong C_4$ be its rotation subgroup.
By D.53 the inverse image $E_R$ has
exactly one $F$-conjugacy class of
labelled order-four sections $B$,
and
\[
                         H^1(C_F(B),\F_2)=0.             \tag{D.57g}
\]
Every element of $E_Q$ preserves this
class, since $R$ is normal in $Q$.
The normalizer $N=N_{E_Q}(B)$ thus
maps onto $Q$, by the same adjustment
of a conjugator by an element of $F$.
Its kernel is $C=C_F(B)$: an element
of $F$ normalizing $B$ fixes the label
of every element of $B$, and therefore
centralizes $B$.
The subgroup $B$ is normal in $N$,
and the quotient $J=N/B$ has
\[
                  1\longrightarrow C\longrightarrow J
                        \longrightarrow Q/R\longrightarrow1,
                              \qquad Q/R=C_2.           \tag{D.57h}
\]
Again $C$ embeds, because $C\cap B=1$.
Conjugation on $B$ in this construction
is the actual conjugation on its
specified rotation labels in $D_8$.

\emph{Torsion and exact finite preimages in the two constructed quotients.}
In each of (D.57f) and (D.57h), the
group $C$ is a closed subgroup of
$F$, and hence is torsion-free with
$\operatorname{cd}_p(C)\leq1$ at every
prime. Equations (D.57d) and (D.57g)
give $C^{(2)}=1$.
Every map from $J$ to a finite two-group
therefore kills $C$, and the displayed
nontrivial quotient gives
\[
                              J^{(2)}=C_2.              \tag{D.57i}
\]
The same full cohomological-dimension
argument used in D.27 and D.46 proves
that $J$ has torsion. Indeed, a
torsion-free $J$ would be projective,
since $C$ is open and has the stated
cohomological dimensions. Its maximal
pro-two quotient would then be a free
pro-two group, contrary to (D.57i).
Every finite subgroup of $J$ injects
into $J/C=C_2$. Consequently $J$ has
an involution $j$ outside $C$.

Take the full inverse image $H$ of
$\langle j\rangle$ in $N$.
In the four-group construction its
kernel is $T$, so $|H|=4$; in the
dihedral construction its kernel is
$B$, so $|H|=8$.
In either case $H\cap F=1$.
To check this in the actual maps,
an element of $H\cap F$ has image
in $\langle j\rangle\cap C=1$ in
$J$, and then lies in $T\cap F=1$
or in $B\cap F=1$ respectively.
Thus the projection $H\to Q$ is
injective. Its order equals $|Q|$,
so it is an isomorphism with the
prescribed labels.
Its inverse is the required section.
This exact injection also fixes the
isomorphism type of the finite lift;
no alternative central extension
is substituted for $Q$.

\emph{Uniqueness above each invariant class.}
The section just constructed identifies
$N=C\rtimes H$ and
\[
                   J=C\rtimes\langle j\rangle .        \tag{D.57j}
\]
The Hochschild--Serre sequence, with
$H^1(C,\F_2)=0$ and
$\operatorname{cd}_2(C)\leq1$, gives
the actual inflation isomorphism
\[
                   H^*(J,\F_2)\cong H^*(C_2,\F_2).
                                                               \tag{D.57k}
\]
The full Brown restriction theorem
\cite[Theorem~2 and \S1.1]{A-KochloukovaZalesskii2007}
applies: $C$ is open of finite
two-cohomological dimension and no
finite subgroup of $J$ can be a
four-group, since it injects into
$C_2$.
In high degree the source of its
dense restriction map is the finite
one-dimensional space in (D.57k).
Restriction to the displayed section
in (D.57j) is already nonzero in
every degree, so this map is injective.
Its image is therefore closed as
well as dense, so the restriction
is an isomorphism onto the product
of the section-normalizer cohomologies.
For a section $S\leq J$ its normalizer
is $C_C(S)\times S$, by the quotient
labels. Each such factor contributes
\[
                          1+\dim H^1(C_C(S),\F_2)
\]
in every sufficiently high degree.
There is consequently exactly one
conjugacy class of sections in $J$,
and its fixed group has zero first
cohomology over $\F_2$.
These are also exactly the
$C$-conjugacy classes: multiply a
given conjugator by the section
element with its quotient label.

Full sections of $E_Q$ containing
the chosen $T$ or $B$ lie in $N$,
and are exactly the full inverse
images of the order-two sections
of $J$. This follows from their
orders and the same kernel calculation
used above.
They therefore form one $C$-conjugacy
class. Every full four-group section
has a positive subgroup $T$ in one
of the invariant classes of
(D.57e); distinct such classes
cannot be identified by $F$-conjugacy.
The numbers of full section classes
are thus exactly three and one
as asserted. Every full dihedral
section has a rotation subgroup
in the unique class of D.53, so
the dihedral count is exactly one.

The fixed group $C_C(j)$ in the
Brown calculation is precisely
$C_F(H)$ in the original extension.
For example, commuting with $j$
modulo $T$ or $B$ means that a
commutator lies in that finite
subgroup. The commutator also lies
in $F$, so is trivial. The elements
of $C$ already commute with $T$
or $B$. This proves the cohomology
part of (D.57c).
Finally an element normalizing $H$
can be multiplied by the inverse
of the unique element of $H$ with
its outer label. The result lies
in $F$ and normalizes $H$, and
therefore centralizes each element
of $H$ by preservation of labels.
This gives the exact direct product
in (D.57c).

The program
\path{code/part_a/torus_two_group_lifts_d57.py}
checks all affine involutions of
$\F_2^2$ together with their invariant
three-point sets and verifies the
three-versus-one fixed-class count.
The full projectivity, Brown and
finite-preimage arguments above
provide the labelled lifts and
their uniqueness statements.
\end{proof}

\begin{claim}{D.58}{Proved}{Negative torus involution sections and cusp fusion}
Let $b$ be a negative involution of the full
once-punctured-torus loop-outer group of
$F=\widehat F(x,y)$, put $c=[x,y]$, and let
$E_b$ be its centre-free extension.
Write $A_2$ for its action on $H_1(F,\Z_2)$.
The $F$-conjugacy classes of order-two sections and
their full fixed-subgroup cohomology are exactly
\[
\begin{array}{c|c|c}
 A_2\bmod2&
 \text{number of section classes}&
 \{\dim_{\F_2}H^1(C_F(S),\F_2)\}_{[S]}\\ \hline
 I&2&\{1,0\}\\
 \ne I&1&\{0\}.
\end{array}                                             \tag{D.58a}
\]
For every such section,
$N_{E_b}(S)=C_F(S)\times S$.
The cusp normalizer is the full profinite dihedral
group
\[
 N_{E_b}(\overline{\langle c\rangle})
       =\overline{\langle c\rangle}
                       \rtimes\langle v\rangle,\qquad
 v^2=1,\quad vcv^{-1}=c^{-1}.                            \tag{D.58b}
\]
Its two internal conjugacy classes of order-two
sections become the same $F$-conjugacy class in $E_b$.
In particular, for every choice of the cusp reflection
$v$, there exists $g\in F$ with
\[
                  gvg^{-1}=cv,\qquad
                  c=[g,v]=g\,(vgv^{-1})^{-1}.            \tag{D.58c}
\]
The dimensions in (D.58a) describe continuous cohomology
of the full closed fixed subgroups.
\end{claim}
\passage{Section 30, folios 194-196, scans 383-387; Section 31, folios 198-199 verso, scans 391-394; Part D(b), full loop-outer finite actions in type (1,1)}

\begin{proof}
\emph{The actual extension and the two integral module cases.}
D.27 supplies an actual order-two representative
$\alpha$ and the exact identification
$E_b=F\rtimes_\alpha\langle s\rangle$.
D.43 gives its global multiplier $\lambda=-1$,
so $A_2^2=I$ and $\det A_2=-1$.
Let $u:E_b\to\F_2$ be the quotient character,
$u(s)=1$.

First suppose $A_2\equiv I\pmod2$.
On the free $\Z_2$-module $H^1(F,\Z_2)$, the induced
involution $D$ also satisfies $D\equiv I\pmod2$.
The endomorphism
\[
                         P=(I-D)/2                     \tag{D.58d}
\]
is integral and satisfies $P^2=P$.
Its image and kernel are direct summands of rank one,
because $D$ has the two eigenvalues $1,-1$ over $\Q_2$.
Thus choose continuous characters
$\varphi_+,\varphi_-:F\to\Z_2$ which form a basis and
satisfy
\[
        \varphi_+\alpha=\varphi_+,\qquad
        \varphi_-\alpha=-\varphi_- .                    \tag{D.58e}
\]
This is a decomposition of the specified cohomology
module. All subsequent maps use these actual
characters; no free generating pair is changed.

Their mod-$2$ reductions extend to characters
$a,b:E_b\to\F_2$ by $a(s)=b(s)=0$.
The Hochschild--Serre spectral sequence has the two
rows $\F_2u^i$ and $\F_2^2u^i$.
The degree-one characters extend, so their
transgression vanishes.
All other possible differentials are its products
with powers of $u$ and also vanish.
The associated graded pieces give, for every $n\geq1$,
\[
 H^n(E_b,\F_2)
     =\langle u^n,u^{n-1}a,u^{n-1}b\rangle_{\F_2}.
                                                        \tag{D.58f}
\]
Multiplication by $u$ is an isomorphism in degrees
at least one, since it is an isomorphism on both
associated graded pieces.

\emph{The two Bocksteins and the actual localized ring.}
Reduce (D.58e) modulo four and use the continuous
one-cochains
\[
 \widetilde a(hs^\epsilon)=\varphi_+(h)\bmod4,\qquad
 \widetilde b(hs^\epsilon)=\varphi_-(h)\bmod4.
\]
The first is a homomorphism, while the second has
coboundary
\[
 \delta\widetilde b(hs^\epsilon,ks^\eta)
        =2\epsilon\varphi_-(k)\pmod4.                   \tag{D.58g}
\]
For the coefficient sequence
$0\to\F_2\xrightarrow{2}\Z/4\to\F_2\to0$,
the Bockstein of a degree-one character is its
cup square. This follows directly by taking its
$\{0,1\}$-valued lift; changing the lift changes
the divided coboundary by a coboundary.
Thus (D.58g) gives
\[
                         a^2=0,\qquad b^2=ub.           \tag{D.58h}
\]

Write $ab=\epsilon ua+\delta ub+\gamma u^2$ using
(D.58f), with coefficients in $\F_2$.
Restriction to the chosen section $\langle s\rangle$
gives $\gamma=0$.
Multiplication by $a$ and (D.58h) then give
\[
           0=a^2b
             =\delta\epsilon\,u^2a+\delta^2u^2b .
\]
The two displayed degree-three classes are independent
by (D.58f), so $\delta=0$.
Consequently the degree-zero localization
$R=(H^*(E_b,\F_2)[u^{-1}])^0$ has the exact presentation
\[
 R=\F_2[A,B]\big/
       (A^2,\ B^2-B,\ AB-\epsilon A),\qquad
 A=a/u,\quad B=b/u,\quad \dim_{\F_2}R=3.                 \tag{D.58i}
\]
For either value of $\epsilon$ the map
\[
 \begin{split}
 R&\longrightarrow\F_2[\nu]/(\nu^2)\times\F_2,\\
 1&\longmapsto(1,1),\\
 A&\longmapsto(\nu,0),\\
 B&\longmapsto(\epsilon,1-\epsilon)
 \end{split}                                           \tag{D.58j}
\]
is a ring isomorphism.
It preserves the stated relations and sends the
three-dimensional basis $1,A,B$ to a basis.
In particular $R$ has exactly two primitive
idempotents, and its nilpotent ideal has dimension one.

\emph{The nonidentity action modulo two.}
Now suppose $A_2\bmod2\ne I$.
The dual two-dimensional module
$M=H^1(F,\F_2)$ is the regular $\F_2[C_2]$-module:
$D-I=D+I$ has rank one, with image equal to kernel.
The periodic cyclic resolution therefore gives
\[
       \dim M^{C_2}=1,\qquad H^i(C_2,M)=0\quad(i>0).
\]
Choose the nonzero invariant character on $F$ and
extend it to $\chi:E_b\to\F_2$ with $\chi(s)=0$.
This again makes the only possible transgression zero.
The two-row spectral sequence now gives
\[
 H^1(E_b,\F_2)=\langle u,\chi\rangle,\qquad
 H^n(E_b,\F_2)=\F_2u^n\quad(n\geq2).                    \tag{D.58k}
\]
Both $u\chi$ and $\chi^2$ restrict to zero on
$\langle s\rangle$, whereas $u^2$ restricts nontrivially.
Thus $u\chi=\chi^2=0$.
In this case the degree-zero localization is
\[
             (H^*(E_b,\F_2)[u^{-1}])^0=\F_2.             \tag{D.58l}
\]

\emph{Full Brown restriction and the section labels.}
Every finite subgroup of $E_b$ injects into $E_b/F=C_2$,
because $F$ is torsion-free.
In particular $E_b$ has no subgroup $C_2\times C_2$.
For a section $S$, conjugation in its normalizer
induces the identity on $S$.
Multiplying a normalizer element by the section
element with its quotient label leaves an element
of $C_S=C_F(S)$. This proves the exact direct product
\[
                         N_{E_b}(S)=C_S\times S.        \tag{D.58m}
\]
The closed subgroup $C_S\leq F$ has
$\operatorname{cd}_2(C_S)\leq1$.

The full-profinite Brown restriction theorem applies
with the open subgroup $F$ of cohomological
$2$-dimension one
\cite[Theorem~2 and \S1.1]{A-KochloukovaZalesskii2007}.
For $n>1$ its restriction map to the product of the
cohomologies of (D.58m) is a dense embedding.
The source is finite by (D.58f) or (D.58k);
therefore that map is an isomorphism onto a finite
product, as in the cited proof.
These are the actual restriction maps and preserve
cup products. Since $u$ restricts trivially on
$C_S$ and nontrivially on $S$, localization gives
the ring isomorphism
\[
 R\ \xrightarrow{\ \operatorname{res}\ }\
   \prod_{[S]}
      \left(\F_2\oplus H^1(C_S,\F_2)u^{-1}\right).
                                                        \tag{D.58n}
\]
The second summand in each factor has square zero,
because $H^2(C_S,\F_2)=0$.
Each factor is local, with that summand as its
nilpotent ideal and with residue field $\F_2$.

If $A_2\bmod2\ne I$, (D.58l) and (D.58n) give
one factor and a zero second summand.
If $A_2\bmod2=I$, (D.58j) gives two factors,
one with a one-dimensional second summand and one
with a zero second summand.
This proves all the numbers in (D.58a).
The classes in this comparison are also precisely
the $F$-conjugacy classes:
for a conjugator $e\in E_b$, multiply it by the
inverse of the unique element of $S$ having the
same quotient label. This replaces $e$ by an
element of $F$ without changing $eSe^{-1}$.

In the identity-modulo-two case the two classes
are distinguished by the value of $b$ on their
nonidentity section element.
Indeed the two residue-field evaluations of
(D.58j) give $B=0$ and $B=1$.
Restriction to a section is exactly that
residue-field evaluation in (D.58n).
Also $a$ vanishes on every section element, since
$a^2=0$ whereas the square of the nonzero
character of $C_2$ is nonzero.

\emph{The exact cusp reflections and their fusion.}
Put $L=\overline{\langle c\rangle}\cong\widehat\Z$.
Preservation of the cusp class and A.7 give
\[
       1\longrightarrow L\longrightarrow N_{E_b}(L)
                     \longrightarrow C_2\longrightarrow1.
                                                        \tag{D.58o}
\]
Choose $v$ over its nonidentity label.
The global multiplier $-1$ gives $vcv^{-1}=c^{-1}$.
Write $v^2=c^\kappa$.
Commutation of $v$ with its own square implies
$c^\kappa=c^{-\kappa}$, hence $2\kappa=0$ in
$\widehat\Z$ and $\kappa=0$.
This proves (D.58b), with the multiplication and
specified label of its finite factor retained.

Every section of this dihedral group is
$\langle c^t v\rangle$, $t\in\widehat\Z$.
Conjugation by $c^z$ changes $t$ to $t+2z$,
and conjugation by $v$ changes it to $-t$.
Its two internal classes are therefore represented
by $v$ and $cv$.
These elements have the same image under every
abelian character of $E_b$, since $c=[x,y]\in[F,F]$.
When there is only one full section class they
are already $F$-conjugate.
When there are two, their equal $b$-values and
the preceding residue-field distinction again
put them in the same $F$-class.
An $F$-conjugator sends the unique nonidentity
element to the unique nonidentity element, giving
the first equation of (D.58c).
Multiplication by $v^{-1}$ gives its other
two exact equations.

The executable
\texttt{code/part\_c/torus\_negative\_involution\_d58.py}
checks the mod-four one-cochains, the cyclic
resolution modules, the two explicit localized
algebra isomorphisms, and geometric word examples
of the cusp equation.
The full-group conclusions above follow from
restriction and the actual normalizer maps.
They assert no full triviality or procyclicity
of the fixed subgroups from their first
cohomology dimensions.
\end{proof}

\begin{claim}{D.59}{Proved}{The negative cusp section and its fixed homology line}
Let $b$ be a negative involution in the full
once-punctured-torus loop-outer group, let
$F=\widehat F(x,y)$ and $c=[x,y]$, and choose a
cusp reflection $v$ in its centre-free extension $E$.
Thus
\[
 E=F\rtimes\langle v\rangle,\qquad
 v^2=1,\qquad vcv^{-1}=c^{-1}.
\]
Put $\alpha=\operatorname{Inn}(v)|_F$,
$C=C_F(v)$, and
\[
 M=H_1(F,\widehat\Z),\qquad
 M^+=\ker(\alpha_*-I:M\to M).
\]
Then the cusp section is specifically the class with
\[
             H^1(C,\F_2)=0,\qquad C\subseteq\overline{[F,F]}.
                                                        \tag{D.59a}
\]
For every actual transporter $g$ supplied by D.58,
\[
                    gvg^{-1}=cv,
                                                        \tag{D.59b}
\]
its homology class is a generator of the full fixed
line:
\[
 \widehat\Z\xrightarrow{\ \cong\ }M^+,\qquad
                         t\longmapsto t[g].
                                                        \tag{D.59c}
\]
This line is a direct summand of $M$.

More precisely, let $m\geq3$ be odd and let
$\chi:F\to C_m$ be any surjective continuous
$\alpha$-invariant character. Extend it to
$\widetilde\chi:E\to C_m$ by
$\widetilde\chi(v)=0$, and put
$F_m=\ker\chi$, $E_m=\ker\widetilde\chi$.
Then
\[
 C\subset F_m,\qquad
 \chi(g)\text{ generates }C_m,\qquad
 E_m=F_m\rtimes\langle v\rangle .
                                                        \tag{D.59d}
\]
The $m$ cusp-dihedral edges form one cycle on the
$m$ $F_m$-section classes above the original cusp
class. Their fixed groups are conjugates of the
same full group $C$.
The complete section counts in this cover are
\[
\begin{array}{c|c|c}
 \alpha_*\bmod2&
 \text{cusp classes in }E_m&
 \text{other classes in }E_m\\ \hline
 I&m\text{, each with }H^1(-,\F_2)=0&
 1\text{, with }\dim H^1(-,\F_2)=1\\
 \ne I&m\text{, each with }H^1(-,\F_2)=0&0 .
\end{array}                                             \tag{D.59e}
\]
In the first row, if $D=C_F(w)$ is the fixed group
of the other full section class, its actual
homology-image map is surjective:
\[
                         D\longrightarrow M^+ .
                                                        \tag{D.59f}
\]
These statements specify full subgroup cohomology
and the indicated ambient homology maps. They make
no identification of $C$ with the trivial group,
of $D$ with a procyclic group, or of $g$ with a member
of a full free profinite basis.
\end{claim}
\passage{Section 30, folios 194-196, scans 383-387; Section 31, folios 198-199 verso, scans 391-394; Part D(b), full loop-outer finite actions in type (1,1)}


\begin{proof}
D.43 gives the global multiplier $-1$.
D.58 gives the stated actual reflection, the
transporter (D.59b), and exactly one or two full
section classes according to the two rows above.
Their normalizers are the actual direct products
$C_F(S)\times S$.
Write $A_2$ for the action on $H_1(F,\Z_2)$.

\emph{The odd-degree cover and its boundary injection.}
The displayed extension of $\chi$ is a homomorphism
because $\chi\alpha=\chi$.
Its kernel is precisely the split extension in
(D.59d). Every order-two element of $E$ lies in
that kernel, since a homomorphism to the odd cyclic
group kills such an element.

The restriction of $\chi$ to the dense discrete
free group is onto. Its degree-$m$ covering of
the punctured torus has $m$ degree-one boundary
lifts, because $\chi(c)=0$.
Filling those boundaries gives the corresponding
unramified degree-$m$ cover of the compact torus.
Consequently its type is $(1,m)$, its completed
group is exactly $F_m$, and
$F_m\cong\widehat F_{m+1}$.
Choose a discrete word $f$ with $\chi(f)=1$, and put
\[
 c_i=f^icf^{-i},\qquad
 v_i=f^ivf^{-i},\qquad 0\leq i<m .
                                                        \tag{D.59g}
\]
Each $v_i$ belongs to $E_m$, has order two, and
inverts $c_i$.

Here is the exact homology comparison required
for the section argument.
Let $\Delta=F/F_m=C_m$ be the deck group and let
$P=H_1(F_m,\Z_2)$. The boundary sequence is
\[
 0\longrightarrow B\longrightarrow P
       \longrightarrow H_1(\overline S_m,\Z_2)
       \longrightarrow0,\qquad
 B=\Z_2^m/\Z_2(1,\ldots,1).
                                                        \tag{D.59h}
\]
The image of the $i$th standard generator of $B$
is $[c_i]$. Thus these classes have their sole
relation $\sum_i[c_i]=0$.
Deck transformations act trivially on the compact
torus homology. Since $m$ is odd, the actual
idempotent
\[
                    e=\frac1m\sum_{\delta\in\Delta}\delta
                                                        \tag{D.59i}
\]
is defined on $P$. It annihilates $B$ and induces
the identity on the last term of (D.59h).
It follows that
\[
             P=eP\oplus(1-e)P,\qquad (1-e)P=B .
                                                        \tag{D.59j}
\]
The covering map identifies $eP$ with the base
torus homology: on integral homology its image
is the index-$m$ covering lattice, and that
inclusion is an isomorphism over $\Z_2$.

The automorphism $\alpha$ induces the identity
on $\Delta$. Its action on $P$ therefore commutes
with the deck action; a discrepancy between two
chosen lifts is an inner automorphism of $F_m$
and has trivial action on $P$.
It preserves every individual boundary class
with multiplier $-1$. Explicitly,
$\alpha(f^i)f^{-i}\in F_m$ conjugates $c_i^{-1}$
to $\alpha(c_i)$. Thus $\alpha=-I$ on $B$.
Its action on $eP$ is the actual base matrix $A_2$,
through the covering homology map just specified.

Reduction modulo two makes $\alpha$ the identity
on $B/2B$. The abelianization of the split extension
gives the exact comparison
\[
 H_1(E_m,\F_2)
   \cong \bigl(H_1(F_m,\F_2)\bigr)_{\alpha}
                    \oplus\F_2[v].
                                                        \tag{D.59k}
\]
The direct summand $B/2B$ survives injectively
in these coinvariants by (D.59j).
In particular the images
$d_i=[c_i]\in H_1(E_m,\F_2)$ still have precisely
their one total-sum relation. Their span has
dimension $m-1$.

\emph{The actual section graph and its transporter labels.}
Let $h=|\chi(C)|$.
The $F_m$-conjugacy classes lying above the
original full cusp-section class are exactly
the vertices
\[
                              C_m/\chi(C).
                                                        \tag{D.59l}
\]
Indeed $f^ivf^{-i}$ and $f^jvf^{-j}$ are
$F_m$-conjugate precisely when $i-j\in\chi(C)$:
a conjugator $n\in F_m$ gives
$f^{-j}nf^i\in C$, and the converse follows by
reversing that formula.
All original order-two sections lie in $E_m$,
so this is the full class splitting, not a
selection of section images in a finite quotient.

For each $i$, take the edge with endpoints
$\langle v_i\rangle$ and $\langle c_i v_i\rangle$.
Both endpoints are actual order-two sections.
Equation (D.59b) identifies the second endpoint
with the conjugate by $f^ig$ of $\langle v\rangle$.
Thus, with $\delta=\chi(g)$, this edge has the
exact vertex labels
\[
              i+\chi(C)\quad\text{and}\quad
                         i+\delta+\chi(C).
                                                        \tag{D.59m}
\]
There are $m$ labelled edges and $m/h$ vertices.
Assign to a vertex the homology class in
$H_1(E_m,\F_2)$ of its nonidentity section element.
This is independent of the conjugate representative.
The sum of the two endpoint values of the $i$th
edge is exactly $d_i$, by the group equality
$c_i v_i=c_i\cdot v_i$.

For a finite graph with $r$ vertices and $k$
connected components, its edge differences span
a subspace of dimension at most $r-k$ in any
vector space of vertex values.
One proves this by choosing one vertex in each
component and expressing every edge difference
in the differences from those chosen vertices.
Apply this to the actual values above. By (D.59k),
\[
                         m-1\leq m/h-k .
                                                        \tag{D.59n}
\]
Since $h\geq1$ and $k\geq1$, equality is forced:
$h=1$ and $k=1$.
The edges now join $i$ to $i+\delta$ in $C_m$.
Connectedness says exactly that $\delta$ generates
$C_m$. Since $m$ is odd and at least three,
these edges form a single $m$-cycle.
This proves all assertions in (D.59d) and the
cusp part of the graph assertion.
Because $C\subset F_m$, the fixed group of each
such section is the corresponding conjugate of
the same full group $C$.

\emph{Which full section is the cusp section.}
We compute the high-degree cohomology of $E_m$
using the actual module (D.59j).
Over $\F_2$ its dual consists of $m-1$ trivial
boundary summands and the dual base module.
When $A_2\bmod2=I$ this last module is two
trivial summands; otherwise it is the regular
two-dimensional $\F_2[C_2]$-module.
The two-row Hochschild--Serre sequence has no
nonzero differential into the quotient row:
the actual section makes inflation from $C_2$
injective in every degree. There are no other
possible differentials.
The periodic cyclic resolution therefore gives,
for every $n\geq2$,
\[
 \dim_{\F_2}H^n(E_m,\F_2)=
 \begin{cases}
   m+2,&A_2\bmod2=I,\\
   m,&A_2\bmod2\ne I.
 \end{cases}                                           \tag{D.59o}
\]

The full-profinite Brown restriction theorem
applies with the open subgroup $F_m$ of
$2$-cohomological dimension one
\cite[Theorem~2 and \S1.1]{A-KochloukovaZalesskii2007}.
Every nontrivial finite subgroup has order two.
For a section $T$, its normalizer in $E_m$ is
$C_{F_m}(T)\times T$, with
$\operatorname{cd}_2C_{F_m}(T)\leq1$.
The finite source in (D.59o) makes the dense
Brown restriction an isomorphism onto its
finite product, so each section class contributes
\[
                     1+\dim H^1(C_{F_m}(T),\F_2)
                                                        \tag{D.59p}
\]
to (D.59o).
Conjugacy by $E_m$ and by $F_m$ agrees: multiply
an $E_m$-conjugator by the element of its source
section with the same quotient label.

D.58 already gives $H^1(C,\F_2)=0$ when
$A_2\bmod2\ne I$.
In the other case its dimension is zero or one.
Choose an invariant surjective character to
$C_3$: the integral involution on $\Z_3^2$ has
distinct eigenvalues $1,-1$, so projection onto
its fixed direct summand supplies this character.
The graph just proved gives three cusp classes
with the same group $C$. A dimension-one $H^1(C,\F_2)$
would make them contribute six to (D.59o),
whose value is five. Thus the dimension is zero
also in this case, proving the first assertion
of (D.59a).

If $A_2\bmod2\ne I$, D.58 has no other full
section class, so the $m$ cusp classes exhaust
$E_m$. If $A_2\bmod2=I$, let $D=C_F(w)$ be the
fixed group of the other full class.
As observed above, $w\in E_m$.
The same exact class-splitting formula gives
$m/|\chi(D)|$ classes, each with fixed group
conjugate to $D\cap F_m$. Subtracting the $m$
cusp contributions from (D.59o) yields
\[
 \frac{m}{|\chi(D)|}
       \bigl(1+\dim H^1(D\cap F_m,\F_2)\bigr)=2 .
                                                        \tag{D.59q}
\]
The first factor is odd. It is therefore one,
and the cohomology dimension is one.
This proves (D.59e), as well as
$\chi(D)=C_m$ for every such odd character.

\emph{The full ambient homology images.}
Every element of $C$ has homology in $M^+$.
At an odd prime $p$, the two idempotents
$(I\pm A_p)/2$ split $\Z_p^2$ into its rank-one
fixed and negative summands.
For every $r\geq1$, projection onto a chosen
generator of the fixed summand followed by
reduction modulo $p^r$ is an actual invariant
surjective character $F\to C_{p^r}$.
Equation (D.59d) makes the image of $C$ vanish
under every one of these maps, and hence in the
entire fixed $\Z_p$-line.
At the prime two, a nonzero compact image of
$C\to M^+_2\cong\Z_2$ would be $2^r\Z_2$ for
some $r$, and its quotient by twice itself would
give a nonzero continuous character $C\to\F_2$.
The already proved $H^1(C,\F_2)=0$ excludes that.
Thus $C$ has zero image in every primary factor
of $M$, which is exactly
$C\subseteq\ker(F\to M)=\overline{[F,F]}$.

From $c=g\alpha(g)^{-1}$, the homology of $g$
belongs to $M^+$. At every odd prime the same
characters and the generator assertion of
(D.59d) say that $[g]$ generates the fixed
$\Z_p$-line. We check the dyadic factor directly.

In the maximal pro-two group, let $D_2,D_3$
be the second and third Zassenhaus terms.
The quotient $D_2/D_3$ is the three-dimensional
$\F_2$-space with basis
\[
                        X^2,\quad Y^2,\quad Z=[X,Y].
\]
For completeness this is the central subgroup of
the finite characteristic quotient with
coordinates $(a,b,k)\in(\Z/4)^2\times\F_2$ and
product
\[
 (a,b,k)(a',b',k')
     =(a+a',\,b+b',\,k+k'+ba').
                                                        \tag{D.59r}
\]
The generators are $(1,0,0),(0,1,0)$ and
their commutator is $(0,0,1)$.
The subgroup with $a,b$ even has the stated basis.
Equivalently its kernel is the closed subgroup
generated by fourth powers, squares of commutators,
and triple commutators, so the quotient is
characteristic.

The induced action on this degree-two space depends
only on $\overline A=A_2\bmod2$. A vector
$aX+bY$ has square
$aX^2+bY^2+abZ$, and $\overline A$ fixes $Z$.
If $\overline A=I$, the degree-two action is
the identity. Otherwise an actual linear basis
change over $\F_2$ gives
$X\mapsto X$, $Y\mapsto Y+X$.
It preserves $Z$ since every invertible
two-dimensional matrix over $\F_2$ has determinant
one. In these coordinates
\[
 \operatorname{im}(I+\overline A|_{D_2/D_3})
                        =\langle X^2+Z\rangle .
                                                        \tag{D.59s}
\]
In either case $Z$ is not in that image.
If $[g]$ vanished modulo two, its pro-two image
would lie in $D_2$, and the exact equation
$c=g\alpha(g)^{-1}$ would give
$Z=(I+\overline A)[g]_2$ in $D_2/D_3$.
This contradicts (D.59s).
Thus $[g]\bmod2\ne0$.

The kernel of $A_2-I$ on $\Z_2^2$ is a saturated
rank-one submodule: if $2z$ lies in it, then
$2(A_2-I)z=0$ in a torsion-free module.
It is therefore generated by a primitive vector,
and the nonzero reduction of $[g]$ makes its
coefficient in that generator a unit.
Together with the odd-primary conclusions, this
proves (D.59c).
For each prime choose a complementary basis vector
to $[g]_p$. Taking their product gives an actual
decomposition $M=\widehat\Z[g]\oplus\widehat\Z e$
and a continuous homology retraction
$M\to\widehat\Z$ sending $[g]$ to one.

Finally suppose $A_2\bmod2=I$.
Equation (D.59q) makes the ambient image of $D$
surject onto every odd-primary fixed line.
At two, use the actual positive character
$\varphi_+$ and its reduction $a$ in D.58e--D.58j.
The nonzero nilpotent class $a/u$ restricts under
the full Brown map to the one-dimensional
$H^1(D,\F_2)$ factor; the cusp factor has no such
summand by (D.59a). Hence $a|_D\ne0$.
The integral idempotent decomposition of D.58
makes $\varphi_+$ primitive on the fixed dyadic
line. Since the ambient image of $D$ lies in
that line, its image contains a dyadic unit
and is all of the line.
The total image is a compact subgroup of
$M^+\cong\widehat\Z$ surjecting onto every primary
factor. Its image in each finite cyclic quotient
is therefore the full quotient; compactness
proves surjectivity onto $\widehat\Z$.
This establishes (D.59f).

The executable
\path{code/part_c/torus_cusp_cycle_d59.py}
checks the exact finite incidence ranks, boundary
relations, cyclic-module dimensions and the
degree-two quotient calculation. The full group
and topology comparisons, including the Brown
restriction and all ambient homology maps, are
proved above.
\end{proof}

\begin{claim}{D.60}{Proved}{The exact torsion quotient of a negative torus extension}
Use D.59's notation
\[
 E=F\rtimes\langle v\rangle,\quad F=\widehat F(x,y),
 \quad c=[x,y],\quad v^2=1,\quad
 \alpha=\operatorname{Inn}(v)|_F,\quad M=F^{\rm ab},
\]
and choose an actual transporter $g$ with
$gvg^{-1}=cv$.
For a subset of a profinite group, double angle
brackets below mean its closed normal closure in
the specified ambient group. There are exact
continuous quotient identifications
\[
 \begin{split}
 E/\langle\!\langle v\rangle\!\rangle_E
 &\ \xrightarrow{\ \cong\ }\ M/(\alpha_*-I)M,\\
 E/\langle\!\langle\operatorname{Tor}(E)\rangle\!\rangle_E
 &\ \xrightarrow{\ \cong\ }\ \widehat\Z ,
 \end{split}                                             \tag{D.60a}
\]
where $\operatorname{Tor}(E)$ is the set of
finite-order elements. More explicitly,
\[
\begin{array}{c|c|c}
 A_2\bmod2&
 E/\langle\!\langle v\rangle\!\rangle_E&
 \chi(g)\\ \hline
 I&\widehat\Z\times C_2&1\\
 \ne I&\widehat\Z&2 .
\end{array}                                              \tag{D.60b}
\]
Here $\chi:E\to\widehat\Z$ is the second quotient
in (D.60a), with the normalization specified in
the proof. Its restriction to $F$ is onto.
It has a continuous section whose image lies
in $F$, so
\[
 E=\langle\!\langle\operatorname{Tor}(E)\rangle\!\rangle_E
                         \rtimes s(\widehat\Z).          \tag{D.60c}
\]
In the first row, let $w$ represent the other
full section class. One can choose that section
$s$ inside $C_F(w)$, and the resulting map
\[
 \widehat\Z\times C_2\longrightarrow E,\qquad
            (t,\epsilon)\longmapsto s(t)w^\epsilon       \tag{D.60d}
\]
is a continuous embedding. These statements specify
the quotient, splitting, and embedded subgroup;
they give no further free-product presentation
of $E$.
\end{claim}
\passage{Section 30, folios 194-196, scans 383-387; Section 31, folios 198-199 verso, scans 391-394; Part D(b), full loop-outer finite actions in type (1,1)}


\begin{proof}
\emph{The two quotient maps before any module decomposition.}
Write $R=\langle\!\langle v\rangle\!\rangle_E$ and
$Q=E/R$. The map $F\to Q$ is onto, since $E$
is generated by $F$ and $v$.
D.58's actual transporter gives
\[
                  c=[g,v]=g\alpha(g)^{-1}.               \tag{D.60e}
\]
Consequently $c$ dies in $Q$.
The profinite quotient
$F/\langle\!\langle[x,y]\rangle\!\rangle_F$
is exactly $\widehat\Z^2=M$: its two generating
images commute, while the abelianization map
to $\widehat\Z^2$ gives the inverse universal map.
Thus the specified map $F\to Q$ factors through
$M$.
Since $v$ dies, $\alpha(f)$ and $f$ have the same
image for every $f\in F$. This supplies a
continuous surjection
\[
                 M/(\alpha_*-I)M\longrightarrow Q.       \tag{D.60f}
\]
The image $(\alpha_*-I)M$ is compact and hence
closed.

In the other direction, the formula
\[
 E\longrightarrow M/(\alpha_*-I)M,\qquad
           f v^\epsilon\longmapsto [f]\bmod(\alpha_*-I)M
                                                               \tag{D.60g}
\]
is a continuous homomorphism. Indeed the product
in the specified semidirect extension is
$f\alpha^\epsilon(h)v^{\epsilon+\eta}$, and
$\alpha_*$ acts trivially on this quotient.
Formula (D.60g) kills $v$, so it factors through
$Q$. Its composite with (D.60f) is the identity
on the images of $F$; the reverse composite has
the same property. Those images generate both
quotients, proving the first isomorphism in
(D.60a).
Equivalently, all of its full-group relations
can be retained as
\[
 Q=F/\langle\!\langle
       c,\ \alpha(f)f^{-1}\ (f\in F)
                         \rangle\!\rangle_F .
                                                               \tag{D.60h}
\]
The relation $c$ is already implied by the
displayed relation with $f=g$, through (D.60e).

\emph{Explicit coordinates for both integral module cases.}
Put $R_0=\widehat\Z$.
D.59 proves that $e_+=[g]$ generates the entire
fixed line in $M$ and that this line is a direct
summand. Choose a complementary basis vector
$e$ in the \emph{homology module}, so
$M=R_0 e_+\oplus R_0 e$.
No free generating pair of $F$ is changed.
Since the global determinant of $\alpha_*$ is
$-1$ by D.43, its matrix in these coordinates is
\[
 A=\begin{pmatrix}1&a\\0&-1\end{pmatrix},\qquad
  (A-I)M=R_0(a,-2)                                      \tag{D.60i}
\]
for a specified $a\in R_0$.
The two cases of (D.60b) are exactly whether
$a\in2R_0$ or $a_2\in\Z_2^\times$.

In the first case write $a=2b$.
The following map is onto:
\[
 R_0^2\longrightarrow R_0\times C_2,\qquad
                   (r,t)\longmapsto(r+bt,\ t\bmod2).
                                                               \tag{D.60j}
\]
Its kernel is exactly $R_0(a,-2)$.
Indeed $t=2z$ and $r+bt=0$ are equivalent to
$(r,t)=(-az,2z)$.
Conversely each multiple of $(a,-2)$ has zero
image. The map is onto by choosing $t=0$ or $1$
for the prescribed second coordinate and then
choosing $r$ for the first.
It therefore gives the first-row quotient
homeomorphism, since a continuous bijection
between compact Hausdorff spaces is a
homeomorphism.

In the second case the map
\[
 R_0^2\longrightarrow R_0,\qquad
                         (r,t)\longmapsto2r+at          \tag{D.60k}
\]
is onto.
At every odd prime its coefficient $2$ is a
unit; at the prime two its coefficient $a_2$
is a unit. These coordinatewise solutions
give an element of $R_0^2$ for every prescribed
value. Its kernel is again exactly
$R_0(a,-2)$:
if $2r+at=0$, the dyadic unit $a_2$ forces
$t_2\in2\Z_2$, hence $t\in2R_0$.
Write $t=-2z$. The equation becomes
$2(r-az)=0$, and $R_0$ has no additive
two-torsion, so $r=az$.
The converse follows by substitution.
This gives the second-row quotient
homeomorphism.

For $[f]=r e_++t e$, define $\chi(fv^\epsilon)$
to be $r+bt$ in the first case and $2r+at$
in the second. Formulas (D.60g), (D.60j), and
(D.60k) prove that these are continuous
homomorphisms. They give exactly the values
$\chi(g)=1$ and $2$, respectively, and
$\chi|_F$ is onto in both cases.

\emph{The second section and the precise finite factor.}
The free profinite group $F$ is torsion-free.
Every nontrivial finite-order element of $E$
therefore has order two and generates a section:
its finite cyclic subgroup injects into
$E/F=C_2$.
D.58 gives one full $F$-conjugacy class in the
second row and exactly two in the first.

In the first row, the dyadic coordinate
character $[f]=r e_++t e\mapsto t_2$ is a
primitive anti-invariant character
$\varphi_-:F\to\Z_2$, since (D.60i) sends
$t$ to $-t$.
Its mod-two extension
\[
 \beta:E\longrightarrow\F_2,\qquad
                    \beta(fv^\epsilon)=t_2\bmod2        \tag{D.60l}
\]
is a homomorphism and vanishes on $v$.
It is exactly the D.58 section-distinguishing
character normalized to vanish on this cusp
reflection.
To make that normalization explicit, if D.58
was written using another initial section $s$,
its character $b$ has the same restriction
to $F$ as (D.60l), and
\[
                         \beta=b+b(v)u,                 \tag{D.60m}
\]
where $u:E\to C_2$ is the original quotient
character. Primitive choices of the
anti-invariant integral character differ by
a dyadic unit and have the same mod-two
restriction. D.58 proves that the two section
classes have opposite $b$-values.
Equation (D.60m) thus gives
$\beta(v)=0$ and $\beta(w)=1$ for every
representative $w$ of the other class.

The image of $w$ in the quotient (D.60j)
has order dividing two.
The additive group $R_0$ is torsion-free,
so this image has first coordinate zero.
Equation (D.60l) is its second coordinate
and has value one. Thus that image is
exactly $(0,1)$.
The closed normal subgroup generated by
all torsion in the first row is consequently
the inverse image of $\{0\}\times C_2$ under
$E\to Q$.
In the second row all torsion is conjugate
to $v$, so its closed normal closure is
already $R$.
This proves in both cases
\[
 \ker\chi=
       \langle\!\langle\operatorname{Tor}(E)\rangle\!\rangle_E
                                                               \tag{D.60n}
\]
and proves the second isomorphism in (D.60a).
It also identifies the exact relation between
the two quotient maps, including the entire
finite factor in the first row.

\emph{Actual splittings and the embedded commuting factor.}
Choose $h\in F$ with $\chi(h)=1$, possible
because $\chi|_F$ is onto.
The continuous profinite power map
\[
                 s:R_0\longrightarrow E,\qquad t\longmapsto h^t
                                                               \tag{D.60o}
\]
satisfies $\chi(s(t))=t$.
It is therefore injective, and compactness
identifies its image with $R_0$ topologically.
Every $z\in E$ has the unique factorization
\[
                  z=\bigl(zs(\chi(z))^{-1}\bigr)s(\chi(z)),
 \qquad zs(\chi(z))^{-1}\in\ker\chi .
\]
This proves (D.60c), with the given maps.

In the first row, D.59 proves that
$D=C_F(w)\to M^+$ is onto.
Choose $h\in D$ with $[h]=[g]$.
Then $\chi(h)=1$, so the same construction
gives a section $s$ inside $D$.
The formula (D.60d) is a homomorphism
because $h$ commutes with $w$.
If its value is the identity, application
of $\chi$ gives $t=0$, and then
$w^\epsilon=1$ gives $\epsilon=0$.
It is a continuous embedding by compactness.
This establishes every stated splitting and
subgroup without identifying their generated
subgroup with the whole extension.

The executable
\texttt{code/part\_c/torus\_torsion\_quotient\_d60.py}
checks the displayed quotient maps and their
exact kernels in finite coefficient rings.
The full quotient identifications follow from
the two universal maps and the actual
transporter relation above.
\end{proof}

\begin{claim}{D.61}{Proved}{Full generation and exact negative conjugator equations}
Put $F=\widehat F(x,y)$ and $c=[x,y]$.
Every closed subgroup $J\leq F$ satisfying
$c\in\overline{[J,J]}$ is $F$ itself.
Consequently every continuous endomorphism of $F$
carrying $c$ to a conjugate of $c^\mu$, with
$\mu\in\widehat\Z^\times$, is an automorphism.

Let $b$ be any negative full once-punctured-torus
outer involution, with the specified extension
\[
 E=F\rtimes\langle v\rangle,\quad
 v^2=1,\quad vcv^{-1}=c^{-1}.
\]
Let $A_2$ be its action on $H_1(F,\Z_2)$.
Set $v_i=c^iv$ and
$\alpha_i=\operatorname{Inn}(v_i)|_F$, for $i=0,1$.
On a separate free profinite group
$F_0=\widehat F(X,Y)$, with $c_0=[X,Y]$, use the
two standard actual negative involutions
\[
 \begin{array}{c|cc}
      &X&Y\\ \hline
 \beta_0&YXY^{-1}&Y^{-1}\\
 \beta_1&Y&X .
 \end{array}                                             \tag{D.61a}
\]
For $j=0,1$, let $\mathcal T_j(b)$ be the space
of peripheral outer isomorphisms $[\phi]:F_0\to F$
conjugating $[\beta_j]$ to $b$.
The peripheral requirement is
$\phi(c_0)\sim_F c^\mu$ for a full unit $\mu$.
Identifying $F_0$ with $F$ by $X\mapsto x$,
$Y\mapsto y$ identifies this with the usual
outer conjugator space.

For $A_2\bmod2=I$, define
\[
 \begin{split}
 \mathcal S_{0,i}=\{(\mu,h,w):\
 &\mu\in\widehat\Z^\times,\ h\in F,\quad
 w^2=1,\ w\in E\setminus F,\\
 &w\text{ belongs to the other full section class},\\
 &[h,w]=1,\quad [h,v_i]=c^\mu\}.
 \end{split}                                             \tag{D.61b}
\]
Here \emph{other} means the class distinct from the
common cusp-reflection class of D.58.
There is a continuous bijection, with continuous
inverse,
\[
 \mathcal S_{0,0}\sqcup\mathcal S_{0,1}
       \xrightarrow{\ \cong\ }\mathcal T_0(b),\qquad
 (\mu,h,w)\longmapsto
             [\,X\mapsto h,\ Y\mapsto v_iw\,].            \tag{D.61c}
\]
For $A_2\bmod2\ne I$, define
\[
 \mathcal S_{1,i}
   =\{(\mu,f)\in\widehat\Z^\times\times F:
                         [f,\alpha_i(f)]=c^\mu\}.
                                                               \tag{D.61d}
\]
The corresponding continuous bijection is
\[
 \mathcal S_{1,0}\sqcup\mathcal S_{1,1}
       \xrightarrow{\ \cong\ }\mathcal T_1(b),\qquad
 (\mu,f)\longmapsto
             [\,X\mapsto f,\ Y\mapsto\alpha_i(f)\,].
                                                               \tag{D.61e}
\]
Equivalently, writing $z=fv_i$, the last equation
is
\[
                z\in E\setminus F,\qquad
                         [z^2,v_i]=c^\mu.               \tag{D.61f}
\]
The two basis images are then $zv_i,v_iz$.
Thus no separate generation condition occurs in
either parameter space. The equations and their
outer conjugator spaces are identified exactly;
nonemptiness for an arbitrary $b$ remains unproved.
\end{claim}
\passage{Section 30, folios 194-196, scans 383-387; Section 31, folios 198-199 verso, scans 391-394; Part D(b), full loop-outer finite actions in type (1,1)}


\begin{proof}
\emph{The boundary test for a full closed subgroup.}
Suppose $J$ were proper. A proper open subgroup
$U\leq F$ contains $J$: separate a point outside
the compact closed subgroup $J$ in a finite
quotient and take the inverse image of the image
of $J$. In particular $c\in U$.
Write $\Pi=F(x,y)$ for the dense discrete
punctured-torus group. Its subgroup
$\Pi_U=\Pi\cap U$ has the same finite index
$d=[F:U]>1$, and its full profinite completion
is exactly $U$.

In the associated connected covering surface,
the boundary lift through the base point has
degree one because $c\in\Pi_U$.
The degrees of all boundary lifts sum to $d$.
There must therefore be a second boundary
component. In the homology of an orientable
surface with at least two boundary components,
the boundary classes have only their total-sum
relation. Each individual boundary class is
nonzero over every $\F_p$.
The chosen class is the class of $c$, with its
specified base path. Its image in
$H_1(U,\F_p)=H_1(\Pi_U,\F_p)$ is nonzero.
On the other hand
$c\in\overline{[J,J]}\subseteq\overline{[U,U]}$,
so that image is zero. This contradiction proves
$J=F$.

Now let $\rho:F\to F$ be an endomorphism with
$\rho(c)=q c^\mu q^{-1}$ for a full unit $\mu$.
Its image is compact and closed. In the conjugate
image $J=q^{-1}\rho(F)q$, the element $c^\mu$
belongs to $\overline{[J,J]}$.
Taking its $\mu^{-1}$ power puts $c$ in the same
closed commutator subgroup. The result just proved
makes $\rho$ onto.
A finitely generated profinite group is Hopfian:
the inverse-image operation of a surjective
endomorphism permutes the finite set of open
normal subgroups of each given index, so its
kernel lies in every such subgroup. Their
intersection is trivial. Thus $\rho$ is an
automorphism.
The same argument applies to a homomorphism
$F_0\to F$ after the explicit basis identification
in the claim.

\emph{The exact two standard extensions.}
The substitutions in (D.61a) square to the
identity and send $c_0$ to $c_0^{-1}$.
Let $E_j=F_0\rtimes_{\beta_j}\langle v_0\rangle$.
For $j=0$ put $w_0=Y^{-1}v_0$.
There are mutually inverse continuous maps
\[
 \begin{aligned}
 (\widehat\Z\langle X\rangle\times\langle w_0\rangle)
                   \amalg\langle v_0\rangle
       &\longrightarrow E_0,\\
 (X,w_0,v_0)&\longmapsto(X,Y^{-1}v_0,v_0),\\
 E_0&\longrightarrow
  (\widehat\Z\langle X\rangle\times C_2)\amalg C_2,\\
 (X,Y,v_0)&\longmapsto(X,v_0w_0,v_0).
 \end{aligned}                                             \tag{D.61g}
\]
Both displayed finite generators have order
two. The equations $w_0^2=1$ and $[X,w_0]=1$
follow directly from the action $\beta_0$.
Conversely $Y=v_0w_0$ gives
$v_0Yv_0=Y^{-1}$ and
$v_0Xv_0=YXY^{-1}$.
Substitution on the displayed generators proves
that the two maps are inverse.
The quotient label sends $X$ to zero and
$w_0,v_0$ to one in $C_2$.
The two finite factors represent distinct
section classes: the character taking $w_0$
to one and $X,v_0$ to zero separates them.

For $j=1$ put $z_0=Xv_0$. The exact maps are
\[
 \begin{aligned}
 \langle v_0\rangle\amalg\widehat\Z\langle z_0\rangle
       &\longrightarrow E_1,\\
 (v_0,z_0)&\longmapsto(v_0,Xv_0),\\
 E_1&\longrightarrow C_2\amalg\widehat\Z,\\
 (X,Y,v_0)&\longmapsto(z_0v_0,v_0z_0,v_0).
 \end{aligned}                                             \tag{D.61h}
\]
Here $v_0$ swaps the two displayed images of
$X,Y$, so all semidirect relations hold.
The two composites fix the generators.
The quotient label is one on both $v_0$ and
$z_0$. These maps specify the full profinite
products through their universal properties.

\emph{Normalization of every peripheral outer conjugator.}
Take $[\phi]\in\mathcal T_j(b)$.
After an inner change in $F$, choose a
representative satisfying
\[
                         \phi(c_0)=c^\mu .              \tag{D.61i}
\]
The unit $\mu$ is unique: a conjugator between
two full unit powers of $c$ normalizes
$L=\overline{\langle c\rangle}$; the full
peripheral normalizer inside $F$ is $L$ itself,
which centralizes $c$.

The centre-free extension construction gives
the unique isomorphism
$\Phi:E_j\to E$ restricting to $\phi$ and
inducing the identity on $C_2$.
Explicitly, conjugation by $\phi$ carries
the full inverse image of
$\langle[\beta_j]\rangle$ in
$\operatorname{Aut}(F_0)$ to the full inverse
image of $\langle b\rangle$ in
$\operatorname{Aut}(F)$.
The embeddings of these extensions are faithful:
an element centralizing their free kernel has
trivial quotient label and then belongs to
the trivial centre of that kernel.
This gives both the extension map and its
inverse.

Put $t=\Phi(v_0)$. It has order two,
has nonzero quotient label, and inverts
$c^\mu$ by (D.61i).
As $\mu$ is a unit, it also inverts $c$.
D.58's actual cusp normalizer therefore gives
the unique expression
\[
                           t=c^zv,\qquad z\in\widehat\Z.
                                                               \tag{D.61j}
\]
Let $i\in\{0,1\}$ be the residue of $z$ modulo
two and put $r=(i-z)/2\in\widehat\Z$.
Replacing $\Phi$ by $\operatorname{Inn}(c^r)\Phi$
retains (D.61i) and sends (D.61j) to
$c^{z+2r}v=v_i$.
Thus every outer class has a representative
with the exact two normalization equations
\[
                    \phi(c_0)=c^\mu,\qquad
                              \Phi(v_0)=v_i.            \tag{D.61k}
\]

This representative and the index $i$ are
unique within the outer class.
Indeed two representatives satisfying the
first equation differ by an inner conjugation
by an element of $C_F(c^\mu)=L$.
Conjugation by $c^r$ changes the reflection
exponent by $2r$. Equality of two indices in
$\{0,1\}$ modulo two first makes the indices
equal, and then $2r=0$ gives $r=0$.
All peripheral centralizers and normalizers
used here are the full specified subgroups.

\emph{The two-class parameter map and its inverse.}
Suppose first that $A_2\bmod2=I$.
For a tuple in (D.61b), put $k=v_iw\in F$.
Commutation of $h$ with $w$ gives the exact
identity
\[
                       [h,k]=[h,v_i]=c^\mu.             \tag{D.61l}
\]
The homomorphism $X\mapsto h$, $Y\mapsto k$
is consequently an automorphism $F_0\to F$
by the boundary test.
Also
\[
 \alpha_i(k)=k^{-1},\qquad
                         \alpha_i(h)=khk^{-1},           \tag{D.61m}
\]
the second equation following by inserting
$w h w=h$.
Equations (D.61m) prove
$\phi\beta_0=\alpha_i\phi$.
Since $\alpha_i$ represents $b$, this is
the asserted peripheral outer conjugator.
Its actual extension sends $v_0$ to $v_i$
and $w_0$ to $w$.

Conversely, normalize any member of
$\mathcal T_0(b)$ by (D.61k).
Put $h=\phi(X)$ and $w=\Phi(w_0)$.
The source relations in (D.61g) give
$w^2=1$, $w\notin F$, and $[h,w]=1$.
That isomorphism carries the two distinct
source section classes to the two distinct
target classes, so $w$ is in the other class.
Since $\phi(Y)=v_iw$, equation (D.61i)
gives (D.61l), hence (D.61b).
These two constructions recover the exact
tuple and the normalized representative.
This proves bijectivity in (D.61c).

\emph{The one-class parameter map and its inverse.}
Now let $A_2\bmod2\ne I$.
A tuple in (D.61d) defines
$\phi(X)=f$ and $\phi(Y)=\alpha_i(f)$.
Its commutator is $c^\mu$, so the boundary
test again proves full invertibility.
The equality $\alpha_i^2=1$ makes
$\phi\beta_1=\alpha_i\phi$.
This gives the indicated outer conjugator
and an extension sending $v_0$ to $v_i$.
Conversely, for the normalized representative
of any member of $\mathcal T_1(b)$, the same
intertwining equation forces
$\phi(Y)=\alpha_i(\phi(X))$.
Taking $f=\phi(X)$ and using (D.61i) yields
exactly (D.61d).

Finally $z=fv_i$ gives
$z^2=f\alpha_i(f)$ and
\[
 [z^2,v_i]
   =f\alpha_i(f)\bigl(\alpha_i(f)f\bigr)^{-1}
   =[f,\alpha_i(f)] .                                   \tag{D.61n}
\]
Conversely $f=zv_i$ and
$\alpha_i(f)=v_iz$.
This proves the alternative parameterization
and both basis-image formulas in (D.61f).

\emph{Topology and the remaining existence obligation.}
Each parameter space is closed in the indicated
compact product: the full other section class
in (D.61b) is the compact image of its conjugation
map, and all equations are continuous.
The maps to automorphisms are continuous on
generator images, and hence so are their maps
to the peripheral outer isomorphism space.
The latter space is Hausdorff and closed in
the profinite outer group, after the explicit
basis identification. The continuous
bijections just proved are therefore
homeomorphisms.

The two integral homology actions in (D.61a)
have respectively identity and nonidentity
reduction modulo two. Conjugacy preserves
this distinction, so only the indicated
standard action can occur in each row.
The formulas supply full actual generators
for every tuple in the displayed spaces.
They establish no tuple for an arbitrary
negative full outer involution.
The remaining problem is the existence of
solutions of those precise equations,
with their retained peripheral element and
reflection index.
The executable
\texttt{code/part\_c/torus\_negative\_parameters\_d61.py}
checks the exact standard free-word and
free-product identities and the boundary
test on all labelled connected covers of
degrees two through five.
\end{proof}

\begin{claim}{D.62}{Proved}{Primitive homology centralizers and the exact cusp-transporter root test}
Let $F=\widehat F(x,y)$ and let $g\in F$ have primitive
abelianization: there is a continuous homomorphism
$\rho:F\to\widehat\Z$ with $\rho(g)=1$.
Then, for every nonzero integer $n$,
\[
 C_F(g^n)=\overline{\langle g\rangle}\cong\widehat\Z,
 \qquad
 N_F(\overline{\langle g\rangle})
                  =\overline{\langle g\rangle}.
                                                        \tag{D.62a}
\]
This conclusion uses homology primitivity; it does not
require a free basis containing $g$.

Now use the exact negative torus extension of
D.58--D.61:
\[
 \begin{gathered}
 E=F\rtimes\langle v\rangle,\qquad v^2=1,\qquad
 \alpha=\operatorname{Inn}(v)|_F,\qquad
 c=[x,y],\qquad \alpha(c)=c^{-1},\\
 C=C_F(v),\qquad gvg^{-1}=cv.
 \end{gathered}                                         \tag{D.62b}
\]
D.59 supplies the required primitive abelianization of
this $g$ and the exact inclusion $C\subseteq\overline{F'}$.
For every $g'\in gC$, the following two conditions
are equivalent:
\[
 \begin{split}
 &\alpha(g')\text{ is conjugate to }g'\text{ in }F;\\
 &C_E(g')\text{ contains an element outside }F.
 \end{split}                                           \tag{D.62c}
\]
Their consequences and converse maps are exact:

\begin{enumerate}
\item
In the two-section-class case of D.58,
they are equivalent to the existence of an involution
$w\in E\setminus F$ centralizing $g'$.
That involution is unique, belongs to the other
section class, and
\[
 C_E(g')=\overline{\langle g'\rangle}\times\langle w\rangle.
                                                        \tag{D.62d}
\]
The maps
\[
 X\longmapsto g',\qquad Y\longmapsto vw,\qquad
 v_0\longmapsto v
                                                        \tag{D.62e}
\]
give a full isomorphism from the two-class standard
extension in D.61 to $E$, preserving the quotient labels
and carrying $[X,Y]$ to $c$.

\item
In the one-section-class case of D.58,
they are equivalent to the existence of $z\in E\setminus F$
with $z^2=g'$.
This root is unique, and
\[
 C_E(g')=\overline{\langle z\rangle}\cong\widehat\Z,
 \qquad
 F\cap C_E(g')=\overline{\langle z^2\rangle}
                         =\overline{\langle g'\rangle}.
                                                        \tag{D.62f}
\]
The maps
\[
 X\longmapsto zv,\qquad Y\longmapsto vz,\qquad
 v_0\longmapsto v
                                                        \tag{D.62g}
\]
give the corresponding full isomorphism from the
one-class standard extension, with the same exact
peripheral and quotient labels.
\end{enumerate}

Thus the multiplier-one, $v_0\mapsto v$ portion of the
D.61 conjugator problem is exactly the existence of
a member $g'\in gC$ satisfying (D.62c).
The explicit continuation (D.62q)--(D.62w) treats every
full unit multiplier and both actual reflection labels,
and gives the same criterion for the entire D.61
conjugator space. Neither equivalence asserts that
its required conjugacy occurs.
\end{claim}
\passage{Section 30, folios 194-196, scans 383-387; Section 31, folios 198-199 verso, scans 391-394; Part D(b), full loop-outer finite actions in type (1,1)}


\begin{proof}
\emph{The primitive-homology centralizer.}
The equation $\rho(g)=1$ first proves that
$a\mapsto g^a$ embeds $\widehat\Z$ as the indicated
closed procyclic subgroup. Put $D=C_F(g^n)$.
We show $\ker(\rho|_D)=1$.

Suppose $k\ne1$ belongs to this kernel.
The procyclic group generated by $k$ has a nontrivial
$p$-primary component for some prime $p$.
Since $F$ is torsion-free, that component is infinite
procyclic, isomorphic to $\Z_p$.
Let $k_p$ be a topological generator of it.
The $p$-primary component of $g^n$ also generates
a copy of $\Z_p$: under $\rho$ it is $n\Z_p$
in the $p$ coordinate, and $n\ne0$ in $\Z_p$.
These two procyclic groups commute, because $k$
centralizes $g^n$. Their intersection is trivial:
$\rho$ is zero on the first and injective on the
second. Multiplication consequently gives a continuous
injection
\[
                 \Z_p\times\Z_p\longrightarrow F.
                                                        \tag{D.62h}
\]
Its image is closed by compactness. This contradicts
$\operatorname{cd}_p(F)=1$, since a closed subgroup
has no larger $p$-cohomological dimension and
$\operatorname{cd}_p(\Z_p^2)=2$.
Thus $\rho|_D$ is injective.

Every $h\in D$ satisfies
$h g^{-\rho(h)}\in\ker(\rho|_D)$, because $g\in D$.
Hence $h=g^{\rho(h)}$, proving the centralizer equality.
For a normalizing element $h$, write
$hgh^{-1}=g^\lambda$ with
$\lambda\in\widehat\Z^\times$.
Applying $\rho$ gives $1=\lambda$.
The normalizer equality now follows from the
centralizer equality with $n=1$.
No change from homology primitivity to free-basis
primitivity has been made.

\emph{The two exact labels of a central extension.}
Fix $g'\in gC$. Then
\[
 [g',v]=c,\qquad
 \alpha(g')=c^{-1}g',\qquad
 [g']_{\rm ab}=[g]_{\rm ab}.
                                                        \tag{D.62i}
\]
The first equality uses the defining commutation of
$C$ with $v$, and the last uses $C\subseteq\overline{F'}$.
In particular (D.62a) applies to $g'$.
An outside element has the unique form $hv$, $h\in F$,
and
\[
 (hv)g'(hv)^{-1}=h\alpha(g')h^{-1}.
                                                        \tag{D.62j}
\]
This is the actual comparison proving (D.62c), in
both directions.

Suppose its conditions hold and choose
$t\in C_E(g')\setminus F$.
The kernel of the parity map on $C_E(g')$ is exactly
$\overline{\langle g'\rangle}$ by (D.62a).
Thus $t^2=(g')^\kappa$ for a unique
$\kappa\in\widehat\Z$, and $t$ commutes with $g'$.
The entire centralizer has the presentation and
coordinate maps
\[
 \begin{split}
 P_\kappa&=
 \langle G,T\mid [G,T]=1,\ T^2=G^\kappa\rangle_{\rm prof}
          \xrightarrow{\ \sim\ } C_E(g'),\\
 G&\longmapsto g',\qquad T\longmapsto t,\\
 (a,\epsilon)&\longmapsto (g')^a t^\epsilon
       \quad(a\in\widehat\Z,\ \epsilon\in\{0,1\}).
 \end{split}                                           \tag{D.62k}
\]
Surjectivity is the two-coset description.
For injectivity, a relation with $\epsilon=1$
has nonzero parity, whereas a relation with
$\epsilon=0$ is detected by the primitive character
of $g'$. The coordinate multiplication is
\[
 (a,\epsilon)(b,\delta)=
 \bigl(a+b+\epsilon\delta\kappa,\,
                       \epsilon+\delta\bmod2\bigr).
                                                        \tag{D.62l}
\]
These maps also give the specified compact topology.

An outside element $(g')^a t$ is an involution
precisely when $2a+\kappa=0$ in $\widehat\Z$.
This equation has a solution exactly when
$\kappa_2\in2\Z_2$; at every odd prime division
by two is unique. In that case
\[
       w=(g')^{-\kappa/2}t,\qquad
       (a,\epsilon)\longmapsto(g')^a w^\epsilon
\]
is the exact isomorphism
$\widehat\Z\times C_2\to C_E(g')$.
Its outside involution is unique.

If $\kappa_2$ is odd, choose $b\in\widehat\Z$ with
\[
 b_2=0,\qquad
 b_p=(1-\kappa_p)/2\quad(p\text{ odd}),
 \qquad \kappa'=\kappa+2b.
                                                        \tag{D.62m}
\]
The element $\kappa'$ is a full unit: its dyadic
coordinate is $\kappa_2$ and every odd coordinate
is one. Set
\[
 t'=(g')^b t,\qquad
 \nu=(\kappa')^{-1},\qquad z=(t')^\nu.
                                                        \tag{D.62n}
\]
Then $(t')^2=(g')^{\kappa'}$, $z^2=g'$, and $z$
has outside parity because $\nu$ is odd at two.
Moreover $t'=z^{\kappa'}$, so $z$ generates the
entire centralizer. If $z^a=1$, then
$(g')^a=(z^2)^a=1$, whence $a=0$.
This proves the isomorphism with $\widehat\Z$.
Multiplication by two is injective there, proving
uniqueness of the specified square root.
The displayed odd-prime corrections retain the
full profinite group, not only its pro-two quotient.

\emph{Which central extension occurs.}
In the two-class case, the explicit torsion quotient
$\chi:E\to\widehat\Z$ of D.60 has
$\chi(g)=1$ and vanishes on $C$.
Thus $\chi(g')=1$.
Applying it to $t^2=(g')^\kappa$ gives
$\kappa=2\chi(t)$, so $\kappa_2$ is even.
The preceding split case gives the involution $w$.
It cannot be $F$-conjugate to $v$:
otherwise $g'\in C_F(w)$ would lie in a conjugate
of $C\subseteq\overline{F'}$, contradicting its
primitive abelianization. By D.58 it therefore
belongs to the other actual section class.

In the one-class case, the same contradiction
excludes every outside involution centralizing $g'$,
because all such involutions are $F$-conjugate to $v$.
Thus $\kappa_2$ must be odd.
Equations (D.62m)--(D.62n) give precisely the root
and centralizer in (D.62f). Conversely either an
outside centralizing involution or an outside
square root is an outside centralizer element,
so these implications are equivalences.

\emph{Full basis and extension maps.}
For the two-class case, put $X'=g'$ and $Y'=vw$.
Both lie in $F$, and $[g',w]=1$ gives
\[
 [X',Y']=[g',v]=c,\qquad
 \alpha(X')=Y'X'(Y')^{-1},\qquad
 \alpha(Y')=(Y')^{-1}.
                                                        \tag{D.62o}
\]
For the one-class case put $X'=zv$ and $Y'=vz$.
Again both lie in $F$, and
\[
 [X',Y']=[z^2,v]=c,\qquad
 \alpha(X')=Y',\qquad \alpha(Y')=X'.
                                                        \tag{D.62p}
\]
In both cases the full boundary test of D.61 proves
that $X\mapsto X'$, $Y\mapsto Y'$ is an automorphism
of the full rank-two free profinite group.
The displayed intertwining identities extend it
to the asserted isomorphism with $v_0\mapsto v$;
its inverse on the kernel is the inverse free-group
automorphism, and on an element $fv^\epsilon$ it is
that inverse applied to $f$, followed by
$v_0^\epsilon$. This proves the exact extension
inverse and retains both parity labels.

Conversely, an isomorphism in this normalized
portion of D.61 sends $X$ in the two-class standard
model, or $XY$ in the one-class standard model,
to an element $g'$ with $[g',v]=c$.
The set of all such elements is precisely $gC$:
equality of the two conjugations of $v$ says that
$g^{-1}g'$ centralizes $v$.
The image of the standard other reflection in the
first model, or of its element $Xv_0$ in the second,
supplies the centralizer element in (D.62c).
This proves both directions of the final assertion.

The centralizer calculation and root formula give
no proof that $\alpha(g')$ is $F$-conjugate to $g'$
for any member of this coset. The continuation below
retains every unit multiplier and the other reflection
label through explicit additional transporter maps.
\end{proof}

\paragraph{Every unit multiplier and both actual reflection labels.}
The preceding criterion extends to the full parameter spaces
of D.61 by explicit transporter maps. No peripheral marking
is rescaled. Retain $c,v,g,C$ from (D.62b), and put
\[
 \begin{gathered}
 v_i=c^iv,\qquad \alpha_i=\operatorname{Inn}(v_i)|_F,
 \qquad C_i=C_F(v_i),\qquad i\in\{0,1\},\\
 C_0=C,\qquad C_1=gCg^{-1}.
 \end{gathered}                                         \tag{D.62q}
\]
For each actual unit $\mu\in\widehat\Z^\times$, define
\[
 g_{\mu,0}=c^{(\mu-1)/2}g,\qquad
 g_{\mu,1}=c^{(\mu+1)/2}g^{-1}.
                                                        \tag{D.62r}
\]
Both exponents are elements of $\widehat\Z$:
the dyadic coordinate of $\mu$ is odd, and division
by two is available at every odd prime.
The conjugation identities
\[
 c^avc^{-a}=c^{2a}v,\qquad
 gvg^{-1}=cv,\qquad g^{-1}(cv)g=v
\]
give, with no change of the named cusp,
\[
 \begin{gathered}
 g_{\mu,i}v_i g_{\mu,i}^{-1}=c^\mu v_i,\\
 \{a\in F:[a,v_i]=c^\mu\}=g_{\mu,i}C_i,\\
 [a]_{\rm ab}=(-1)^i[g]_{\rm ab}
           \quad(a\in g_{\mu,i}C_i).
 \end{gathered}                                         \tag{D.62s}
\]
The coset equality follows by dividing two transporters.
The homology equality follows from
$c\in\overline{F'}$ and
$C_i\subseteq\overline{F'}$, the latter being (D.59a)
and its displayed conjugate. In particular every one
of these transporters is homologically primitive;
its exact homology is independent of the unit $\mu$.

For a specified $a\in g_{\mu,i}C_i$ we now have
\[
 \alpha_i(a)=c^{-\mu}a,\qquad
 \alpha_i(a)\sim_F a
       \ \Longleftrightarrow\
                 C_E(a)\not\subset F.                  \tag{D.62t}
\]
The equivalence is the exact map $h\mapsto hv_i$
of (D.62j). In the two-class case, D.60 gives
$\chi(a)=(-1)^i$. Thus $t^2=a^\kappa$ for an
outside centralizer element implies
$2\chi(t)=(-1)^i\kappa$.
The even-dyadic argument (D.62k)--(D.62l) supplies
the unique outside commuting involution $w$.
It is in the other full section class because
$a$ has nonzero primitive homology whereas
every cusp-class fixed group lies in
$\overline{F'}$. The exact basis and extension maps are
\[
 X\longmapsto a,\qquad Y\longmapsto v_iw,\qquad
 v_0\longmapsto v_i,\qquad [X,Y]\longmapsto c^\mu.
                                                        \tag{D.62u}
\]
In the one-class case, that same fixed-group
argument excludes every outside involution in
$C_E(a)$. The odd-dyadic construction
(D.62m)--(D.62n) gives the unique $z\in E\setminus F$
with $z^2=a$, and the exact maps are
\[
 X\longmapsto zv_i,\qquad Y\longmapsto v_i z,\qquad
 v_0\longmapsto v_i,\qquad [X,Y]\longmapsto c^\mu.
                                                        \tag{D.62v}
\]
In both cases D.61's boundary test proves full
invertibility because $\mu$ is a full unit.
The inverse extension maps are the inverse kernel
map together with the specified label-$1$ reflection,
as in (D.62o)--(D.62p).

Conversely, every normalized conjugator in D.61
has some specified $\mu$ and $i$. In the two-class
model its image of $X$, and in the one-class model
its image of $XY$, belongs to the precise coset
in (D.62s). The image of the standard other
reflection or of $Xv_0$ supplies the outside
centralizer element, respectively. Therefore
\[
 \mathcal T_j(b)\ne\varnothing
 \quad\Longleftrightarrow\quad
 \left\{\begin{array}{l}
 \text{there are }\mu\in\widehat\Z^\times,\ i\in\{0,1\},
       \ a\in g_{\mu,i}C_i\\
 \text{with }c^{-\mu}a\sim_F a .
 \end{array}\right.
                                                        \tag{D.62w}
\]
Here $\mathcal T_j(b)$ is the exact peripheral
conjugator space defined in D.61, with its standard
model chosen by the section count. Equations
(D.62u)--(D.62v) and the converse give the actual
maps proving this equivalence. They retain the
unit exponent $c^\mu$ and both reflection labels
throughout. The existence in the right-hand set
remains unresolved; none of these identities
asserts it for an arbitrary full negative action.

The executable
\path{code/part_c/torus_unit_transporters_d62.py}
checks both reflection-label identities in the full
free-product word normal form for thirty-two odd
integer exponents. These identities hold before
imposing the boundary condition; the full unit
hypothesis is retained where D.61 proves invertibility.
The full centralizer and root statements are proved
above by their exact subgroup and coefficient maps.

\begin{claim}{D.63}{Proved}{The exact homology characters of a negative torus centralizer}
Let $\Omega$ be the full once-punctured-torus peripheral
outer group, and let $b\in\Omega$ be a negative involution.
Keep the actual extension and peripheral labels of D.59--D.62:
\[
 \begin{gathered}
 F=\widehat F(x,y),\quad c=[x,y],\quad
 E=F\rtimes\langle v\rangle,\quad
 \alpha=\operatorname{Inn}(v)|_F,\quad
 v^2=1,\quad \alpha(c)=c^{-1},\\
 M=F_{\rm ab},\quad A=\alpha_*,\quad
 C=C_F(v),\quad gvg^{-1}=cv,\quad
 M^+=\ker(A-I)=\widehat\Z[g].
 \end{gathered}                                               \tag{D.63a}
\]
In particular, D.59 says that $[g]$ is primitive and
$C\subseteq\overline{F'}$.
Put $Z_b=C_\Omega(b)$.
There are actual continuous characters
\[
 \sigma:Z_b\longrightarrow\{1,-1\},\qquad
 \lambda:Z_b\longrightarrow\widehat\Z^\times
                                                               \tag{D.63b}
\]
such that $\lambda$ is the peripheral multiplier,
$\beta_*[g]=\sigma(\beta)[g]$, and the full homology
matrix of every $\beta\in Z_b$ is exactly
\[
 \rho(\beta)=
 \frac{\sigma(\beta)}2
 \bigl((1+\lambda(\beta))I+(1-\lambda(\beta))A\bigr)
       \quad\text{in }\operatorname{GL}_2(\widehat\Z).
                                                               \tag{D.63c}
\]
Both half-coefficients in this formula belong to
$\widehat\Z$: a full unit is odd at two, and two
is invertible at every odd prime.
In particular,
\[
          \lambda(\beta)=1
       \quad\Longrightarrow\quad
          \rho(\beta)=\sigma(\beta)I.                          \tag{D.63d}
\]

Write
\[
 \begin{split}
 \Xi_b&=\{(\sigma(\beta),\lambda(\beta)):\beta\in Z_b\},\\
 I_b&=Z_b\cap\ker\bigl(\Omega\longrightarrow
                          \operatorname{GL}_2(\widehat\Z)\bigr).
 \end{split}
\]
The exact character sequence is
\[
 1\longrightarrow I_b\longrightarrow Z_b
       \xrightarrow{\,(\sigma,\lambda)\,}\Xi_b
       \longrightarrow1.                                    \tag{D.63e}
\]
Its kernel has the following description by actual
automorphisms, with inverse maps given in the proof:
\[
 I_b\ \cong\
 J_b:=
 \{\psi\in\operatorname{Aut}(F):
          \psi(c)=c,\ \psi\alpha=\alpha\psi\}.                 \tag{D.63f}
\]
Every member of $J_b$ automatically acts trivially on $M$.
The group $I_b$ is torsion-free.
This description does not assert that $I_b$ is trivial,
or that $\Xi_b$ is the whole product in (D.63b).

Every finite subgroup $K\leq Z_b$ has the exact
faithful homology embedding
\[
 K\ \lhook\joinrel\longrightarrow\
       \{I,-I,A_2,-A_2\}\cong C_2\times C_2,
       \qquad \beta\longmapsto \rho_2(\beta).                 \tag{D.63g}
\]
Here $A_2$ is the original dyadic involution matrix.
Thus its possible abstract types are $1,C_2,C_2^2$;
a finite subgroup containing $b$ has order two or four.
This bound makes no assertion that a four-element
subgroup occurs for every $b$.
\end{claim}
\passage{Section 30, folios 194-196, scans 383-387; Section 31, folios 198-199 verso, scans 391-394; Part D(b), full loop-outer finite actions in type (1,1)}


\begin{proof}
\emph{The extension map for a centralizing outer class.}
The map
\[
 E\longrightarrow\operatorname{Aut}(F),\qquad
       f v^\epsilon\longmapsto
                  \operatorname{Inn}(f)\alpha^\epsilon
                                                               \tag{D.63h}
\]
identifies $E$ with the full inverse image of
$\langle b\rangle$ under
$\operatorname{Aut}(F)\to\operatorname{Out}(F)$.
Indeed an element of that inverse image is
$\operatorname{Inn}(f)\alpha^\epsilon$ by its outer label.
The map is injective because $F$ is centre-free and
$b$ has order two. It is a homeomorphism onto its image
by compactness. This also gives $C_E(F)=1$.

Let $\beta\in Z_b$, and choose an actual representative
$\psi\in\operatorname{Aut}(F)$.
Conjugation by $\psi$ preserves the inverse image in
(D.63h), since its outer class centralizes $b$.
Transport through (D.63h) gives an actual automorphism
$\Psi:E\to E$, with
$\Psi|_F=\psi$ and the identity map on $E/F=C_2$.
This constructs the required extension directly from
the centralizer; no conjugator from a standard
extension has been presumed.

Write the peripheral equation as
$\psi(c)=u c^\mu u^{-1}$, with
$\mu=\lambda(\beta)\in\widehat\Z^\times$.
Conjugate $\Psi$ by $u^{-1}$ to obtain
$\Psi(c)=c^\mu$.
The element $t=\Psi(v)$ has order two, has outside
label, and inverts $c^\mu$.
Taking its $\mu^{-1}$ power shows that it inverts $c$.
The full cusp normalizer from D.58 is
$\langle c\rangle\rtimes\langle v\rangle$.
Consequently
\[
                      t=c^zv
                 \quad(z\in\widehat\Z)
\]
with a unique exponent.
Let $i\in\{0,1\}$ be the reduction of $z$ modulo two.
Conjugating the entire extension map by
$c^{(i-z)/2}$ now gives
\[
                \Psi(c)=c^\mu,\qquad
                \Psi(v)=v_i:=c^iv.                           \tag{D.63i}
\]
The half-exponent is defined at every prime, with
even numerator at two.
All these conjugations lie in $F$, so the outer
class and its homology action have been preserved.

For completeness, (D.63i) determines the actual
representative uniquely once its outer class is fixed.
An inner comparison preserving $c^\mu$ belongs to
$C_F(c^\mu)=\langle c\rangle$, by the full peripheral
centralizer calculation. Write it as $c^a$.
It changes the exponent of $v_i$ by $2a$.
Two exponents in $\{0,1\}$ can therefore agree only
with the same $i$, and then $2a=0$ gives $a=0$.
The multiplier itself is unique by the same full
peripheral normalizer calculation.

\emph{The actual transporter determines the fixed-line sign.}
Apply $\Psi$ to $gvg^{-1}=cv$.
The resulting equation is
\[
             \Psi(g)v_i\Psi(g)^{-1}=c^\mu v_i.
                                                               \tag{D.63j}
\]
D.62 gives the full solution cosets, retaining the
specified generator $c$:
\[
 \begin{array}{c|c|c}
 i&\Psi(g)\text{ belongs to}&[\Psi(g)]\text{ in }M\\ \hline
 0&c^{(\mu-1)/2}g\,C&[g]\\
 1&c^{(\mu+1)/2}g^{-1}\,gCg^{-1}&-[g].
 \end{array}                                                  \tag{D.63k}
\]
The last column uses $C\subseteq\overline{F'}$.
Thus the action on the fixed line is specifically
$\sigma(\beta)=(-1)^i$.
Because $[g]$ generates the full direct summand
$M^+$, its actual action defines a continuous
character; composition multiplies these two signs.

\emph{The full matrix and the exact kernel.}
The peripheral multiplier is also
$\det\rho(\beta)$.
One can verify the equality at each prime $p$ by
mapping $F$ to the group of upper unitriangular
three-by-three matrices over $\Z_p$, with $x,y$
the two elementary generators.
The commutator coordinate of their images under
$\psi$ is $\det\rho_p(\beta)$.
The equation $\psi(c)=c^\mu$ makes it $\mu_p$;
a peripheral conjugator would have the same effect
because that commutator coordinate is central.
This proves the equality in every $\Z_p$ and hence
in $\widehat\Z$.

The matrices $\rho_p(\beta)$ and $A_p$ commute,
since their outer classes commute.
On $M\otimes_{\widehat\Z}\Q_p$, the matrix $A_p$
has the two distinct eigenvalues $1,-1$.
Each eigenspace has dimension one, since its
determinant is the negative multiplier $-1$.
The fixed-space eigenvalue of $\rho_p(\beta)$
is $\sigma$, and its determinant is $\mu_p$.
Its other eigenvalue is consequently $\sigma\mu_p$.
This proves (D.63c) over $\Q_p$.
The two half-coefficients are integral as explained
in the statement, and $\Z_p^2\to\Q_p^2$ is injective.
The same matrix equality therefore holds over
$\Z_p$ for every $p$, proving the displayed full
formula and (D.63d).

More explicitly, extend $[g]$ to a basis of $M$.
In that basis there is an $a\in\widehat\Z$ with
\[
 A=\begin{pmatrix}1&a\\0&-1\end{pmatrix},
 \qquad
 \rho(\beta)=
 \sigma\begin{pmatrix}
           1&a(1-\mu)/2\\0&\mu
         \end{pmatrix}.                                    \tag{D.63l}
\]
This is the same formula on the actual integral
lattice, including its possibly nonsplit dyadic
eigenspaces. It does not replace that lattice by
a direct sum of the two integral eigenlattices.

The formula proves that the kernel of
$(\sigma,\lambda)$ is exactly $I_b$.
Its image is the stated compact group $\Xi_b$,
so (D.63e) is an exact sequence of profinite groups.
In particular all commutators in $Z_b$ lie in $I_b$.
On $Z_b^1=\ker\lambda$ it restricts to the exact
sequence
\[
 1\longrightarrow I_b\longrightarrow Z_b^1
       \xrightarrow{\,\sigma\,}\sigma(Z_b^1)
       \longrightarrow1,\qquad
                 \sigma(Z_b^1)\leq\{1,-1\}.                 \tag{D.63m}
\]
No larger image is inferred.

For $\beta\in I_b$, the unique normalized
representative (D.63i) has $\mu=1$ and $i=0$.
Its extension fixes $v$ and $c$, so its restriction
belongs to $J_b$.
Conversely, a member $\psi\in J_b$ extends by
$v\mapsto v$ to $E$, and its outer class centralizes
$b$. It is already normalized with $\mu=1,i=0$.
Equations (D.63c) and (D.63k) give $\psi_*=I$,
so that outer class belongs to $I_b$.
These two maps are inverse.
Composition of two representatives fixing both
$c$ and $v$ again fixes both, so they are group
homomorphisms. They are homeomorphisms because
$J_b$ is a closed subgroup of the profinite group
$\operatorname{Aut}(F)$ and the outer projection is
a continuous bijection onto $I_b$.
Equivalently, an inner automorphism in $J_b$ would
be conjugation by $c^a$ and would require
$c^a=\alpha(c^a)=c^{-a}$, giving $a=0$.
This also verifies directly that no inner ambiguity
remains in (D.63f).

D.55 proves faithfulness of every primary homology
representation on every finite subgroup of $\Omega$.
A nonidentity torsion element of $I_b$ would generate
a finite subgroup killed by that representation.
This proves that $I_b$ is torsion-free.
It does not remove infinite-order elements of the
actual group (D.63f).

\emph{Finite centralizers.}
For $\beta$ in a finite $K\leq Z_b$, its dyadic
multiplier is a finite-order element of
$\Z_2^\times$, hence is $1$ or $-1$.
Indeed $\Z_2^\times=\{1,-1\}\times(1+4\Z_2)$,
and the latter factor is torsion-free by the
injective $2$-adic logarithm.
Formula (D.63c) at two therefore puts its matrix
in $\{I,-I,A_2,-A_2\}$.
These four matrices are distinct and form a
Klein four group: $A_2^2=I$ and
$\det A_2=-1$, which excludes $A_2=\pm I$.
D.55 makes the displayed map injective on $K$.
This proves (D.63g) and all the stated finite bounds.
The proof leaves exactly the kernel in (D.63f)
and the actual image $\Xi_b$ to be determined.
\end{proof}

\begin{claim}{D.64}{Proved}{Standard negative torus centralizers and their actual kernels}
Let $b_0,b_1$ be the two standard negative outer
involutions on $F_0=\widehat F(X,Y)$:
\[
 \begin{array}{c|cc}
       &X&Y\\ \hline
 \alpha_0&YXY^{-1}&Y^{-1}\\
 \alpha_1&Y&X .
 \end{array}                                                  \tag{D.64a}
\]
Both have the named cusp $c_0=[X,Y]$ and
$\alpha_j(c_0)=c_0^{-1}$.
For $b_0$, put $\zeta_0(X)=X^{-1}$,
$\zeta_0(Y)=Y$.
Its full peripheral outer centralizer and normalizer
are exactly
\[
 C_\Omega(b_0)=N_\Omega(\langle b_0\rangle)
      =\langle b_0,[\zeta_0]\rangle\cong C_2^2.               \tag{D.64b}
\]
Thus every fibre of the actual orbit map
$\theta\mapsto\theta b_0\theta^{-1}$ is precisely
a right coset of this specified four-element group.
Conjugation by $\theta$ gives the exact corresponding
centralizer on that orbit.

For both $j=0,1$, the IA kernel $I_{b_j}$ of
D.63 is trivial. Consequently the full centralizer
of $b_1$ has the faithful character embedding
\[
 C_\Omega(b_1)\ \lhook\joinrel\longrightarrow\
          \{1,-1\}\times\widehat\Z^\times,\qquad
                         \beta\longmapsto(\sigma,\lambda).
                                                               \tag{D.64c}
\]
In particular it is abelian.
For the actual automorphism
$\zeta_1(X)=X^{-1}$, $\zeta_1(Y)=Y^{-1}$,
\[
 \begin{split}
 C_\Omega(b_1)\cap\ker\lambda&=\langle[\zeta_1]\rangle
                                                  \cong C_2,\\
 \{\beta\in C_\Omega(b_1):\lambda(\beta)\in\{1,-1\}\}
                   &=\langle b_1,[\zeta_1]\rangle\cong C_2^2.
 \end{split}                                                  \tag{D.64d}
\]
This is also the entire torsion subgroup of this
abelian centralizer. The remaining multiplier
calculation for (D.64c) is supplied by D.65 below.
\end{claim}
\passage{Section 30, folios 194-196, scans 383-387; Section 31, folios 198-199 verso, scans 391-394; Part D(b), full loop-outer finite actions in type (1,1)}


\begin{proof}
\emph{The exact standard extensions.}
D.61 gives the mutually inverse full presentation maps
\[
 \begin{aligned}
 E_0&=(\widehat\Z\langle X\rangle\times C_2\langle w\rangle)
                                      \amalg C_2\langle v\rangle,\\
 Y&=vw,\qquad w=Y^{-1}v,\\
 E_1&=C_2\langle v\rangle\amalg\widehat\Z\langle z\rangle,\\
 X&=zv,\qquad Y=vz,\qquad z=Xv .
 \end{aligned}                                                \tag{D.64e}
\]
The parity homomorphisms have
$\epsilon(X)=0,\epsilon(v)=\epsilon(w)=1$ in the
first row and $\epsilon(v)=\epsilon(z)=1$ in the
second. Their kernels are the actual $F_0$ in
(D.64a), and conjugation by $v$ gives $\alpha_j$.
The full free-factor intersection theorem used
in D.35 gives
\[
 C_{E_0}(v)=\langle v\rangle,\qquad
 C_{E_0}(w)=\langle X\rangle\times\langle w\rangle,\qquad
 C_{E_1}(v)=\langle v\rangle.
                                                               \tag{D.64f}
\]
For example a centralizer element outside the
specified factor would put the nontrivial involution
in the intersection of that factor with its conjugate.
The intersection theorem excludes that possibility;
the indicated factor is already abelian where needed
\cite[Theorem 5 of the inspected author version]{A-HerfortZalesskii2008}.
In particular $C_{F_0}(v)=1$ in both rows.

The actual fusion transporters are
\[
          g_0=X,\qquad g_1=z^2=XY,\qquad
          g_jvg_j^{-1}=c_0v.
                                                               \tag{D.64g}
\]
The identities follow by substitution in (D.64e).
Both have primitive homology, and
$c_0=[g_j,v]$ in their respective extensions.

\emph{Vanishing of the two actual IA kernels.}
Use D.63f to represent an IA centralizer element
by its unique actual automorphism fixing $c_0$
and commuting with $\alpha_j$.
Its extension fixes $v$.
The exact multiplier-one transporter formula in
D.62, together with $C_{F_0}(v)=1$, makes it fix
$g_j$ as well.

For $E_0$, D.62 says that $w$ is the unique
outside involution centralizing $g_0=X$.
The extension must fix it.
It fixes all three generators $X,w,v$ of $E_0$,
and is therefore the identity.
For $E_1$, D.62 says that $z$ is the unique
outside square root of $g_1=z^2$.
The extension must fix it, and hence fixes the
two generators $z,v$.
This proves $I_{b_0}=I_{b_1}=1$ using actual
extension generators and their uniqueness, rather
than generation inferred from homology.
The faithful map (D.64c) and abelianness follow
from D.63e.

\emph{All multipliers in the two-class standard model.}
Let $\beta\in C_\Omega(b_0)$.
Normalize its actual extension by D.63:
\[
        \Psi(c_0)=c_0^\mu,\qquad
        \Psi(v)=c_0^iv,\qquad i\in\{0,1\}.
\]
There is no fixed-subgroup ambiguity in its
transporter image, since the two groups $C_i$
in D.62 are trivial by (D.64f).
Thus, with $a=\Psi(X)$,
\[
       a=
       \begin{cases}
         c_0^{(\mu-1)/2}X,&i=0,\\
         c_0^{(\mu+1)/2}X^{-1},&i=1.
       \end{cases}                                          \tag{D.64h}
\]
The involution $t=\Psi(w)$ belongs to the other
full $F_0$-section class: $\Psi$ permutes those
two classes, and already preserves the cusp
class by its normalized image of $v$.
Therefore $t=hwh^{-1}$ for some actual $h\in F_0$.
Since $a$ centralizes $t$, (D.64f) gives
\[
 a\in h\bigl(C_{E_0}(w)\cap F_0\bigr)h^{-1}
       =h\langle X\rangle h^{-1}.
 \]
Write $a=hX^u h^{-1}$.
Its primitive homology in D.62 is $(-1)^i[X]$,
so $u=(-1)^i$ in the full $\widehat\Z$.
This is the exact conjugacy
\[
                       a\sim_{F_0}X^{(-1)^i}.                \tag{D.64i}
\]

For every integer $n\geq2$, use the explicit
finite dihedral quotient
\[
 \begin{split}
 q_n:E_0&\longrightarrow
       D_{2n}=\langle r,s\mid r^n=s^2=1,\ srs=r^{-1}\rangle,\\
       X&\longmapsto r,\qquad v\longmapsto s,\qquad
       w\longmapsto1 .
 \end{split}                                                  \tag{D.64j}
\]
It respects the defining relations in (D.64e)
and is surjective.
The exact peripheral image is $q_n(c_0)=r^2$.
Both lines of (D.64h) consequently give
$q_n(a)=r^\mu$.
By (D.64i), this rotation is conjugate to
$r^{(-1)^i}$.
The conjugacy class of either $r$ or $r^{-1}$
in this finite group is $\{r,r^{-1}\}$.
Hence
\[
                  \mu\bmod n\in\{1,-1\}
                       \quad\text{for every }n\geq2.
                                                               \tag{D.64k}
\]
These simultaneous congruences force one global
sign $\mu=1$ or $\mu=-1$ in $\widehat\Z$.
For a direct compactness proof, the sets
$\{\varepsilon\in\{1,-1\}:\mu\equiv\varepsilon
\pmod n\}$ are nonempty and have the finite
intersection property, since the condition for
a least common multiple implies every specified
finite set of conditions.
Their intersection supplies the one sign at
all moduli.

Together with $I_{b_0}=1$, D.63 now bounds its
full centralizer by four elements.
All four occur through actual extension maps.
The map fixing $v,w$ and sending $X$ to $X^{-1}$
is an order-two automorphism of $E_0$.
On the kernel it is precisely $\zeta_0$ in
the statement.
It commutes with $\operatorname{Inn}(v)$ and
\[
                 \zeta_0(c_0)=X^{-1}c_0^{-1}X,
\]
so it is peripheral of multiplier $-1$.
The two commuting involutions have respective
homology matrices
$\operatorname{diag}(1,-1)$ and
$\operatorname{diag}(-1,1)$.
They give four distinct outer classes.
This proves (D.64b).
A normalizer of an order-two subgroup fixes its
sole nonidentity element and is its centralizer.
Finally equality of the two orbit images of
$\theta,\theta'$ is exactly the statement that
$\theta^{-1}\theta'$ belongs to this computed
centralizer, proving the stated actual orbit fibres.

\emph{The exact multiplier-one conclusion in the one-class model.}
The automorphisms $\alpha_1$ and $\zeta_1$ commute
on both free generators and have order two.
Moreover
\[
        \zeta_1(c_0)=(XY)^{-1}c_0(XY),
        \qquad (\zeta_1)_*=-I .
                                                               \tag{D.64l}
\]
Thus $\zeta_1$ has multiplier one and fixed-line
sign $-1$, while $b_1$ has multiplier $-1$
and fixed-line sign $1$.
They provide the four indicated elements.
The injective characters in (D.64c) allow at
most two elements with multiplier one.
Multiplication by $b_1$ gives exactly the two
with multiplier $-1$.
This proves (D.64d).
Every torsion element of this abelian centralizer
generates a finite subgroup; together with the
displayed four-element subgroup it still generates
a finite group.
D.63g bounds that group by order four.
Consequently the displayed four elements are
exactly all its torsion.
This argument establishes the stated character
embedding and its torsion subgroup. D.65 below
computes its remaining multiplier image through
the specified cone-filling quotients.

The program
\path{code/part_c/standard_negative_centralizers_d64.py}
checks the exact free-group identities in
(D.64a), (D.64g), and (D.64l), and checks both
dihedral transporter formulas over a finite range.
The global conclusion (D.64k) is the proved
all-modulus statement above.
\end{proof}

\begin{claim}{D.65}{Proved}{A uniform square obstruction and the remaining standard centralizer}
For each integer $n\geq2$, put
\[
 G_n=\langle z,v\mid v^2=1,\ [z^2,v]^n=1\rangle,\qquad
 g=z^2,\quad c=[g,v].
                                                               \tag{D.65a}
\]
These are discrete groups; $\widehat G_n$ denotes their
full profinite completions.
For an integer $r$ with $\gcd(2r+1,n)=1$, the exact
square criterion is
\[
 c^rg\text{ is a square in }\widehat G_n
       \quad\Longleftrightarrow\quad
                  r\equiv0\text{ or }-1\pmod n.              \tag{D.65b}
\]
For either allowed residue the square root already
lies in the displayed discrete group.

Consequently the full peripheral outer centralizer
of the standard one-section negative torus involution
$b_1=[\alpha_1]$, $\alpha_1(X)=Y,\alpha_1(Y)=X$, is
\[
 C_\Omega(b_1)=N_\Omega(\langle b_1\rangle)
       =\langle b_1,[\zeta_1]\rangle\cong C_2^2,\qquad
       \zeta_1(X)=X^{-1},\quad\zeta_1(Y)=Y^{-1}.
                                                               \tag{D.65c}
\]
In particular all its multiplier values are the two
global units $1,-1$. Together with D.64 this computes
both standard negative centralizers, and every fibre
of each actual standard conjugacy-orbit map has exactly
four elements, with the explicit right-coset description.
It does not assert that the two standard orbits exhaust
all negative involutions.
\end{claim}
\passage{Section 30, folios 194-196, scans 383-387; Section 31, folios 198-199 verso, scans 391-394; Part D(b), full loop-outer finite actions in type (1,1)}


\begin{proof}
\emph{Two exact discrete graph-of-groups presentations.}
Let
\[
 \begin{split}
 D_n&=\langle c,v\mid c^n=v^2=1,\ vcv=c^{-1}\rangle,
           \qquad w=cv,\\
 H_n&=\langle D_n,g\mid gvg^{-1}=w\rangle.
 \end{split}                                                  \tag{D.65d}
\]
The first group is the finite dihedral group of order
$2n$. Its two specified order-two subgroups are
$\langle v\rangle,\langle w\rangle$.
The second line is their HNN extension, with the
actual edge map $v\mapsto w$.
The HNN normal form embeds $D_n$ in $H_n$.
There is an actual exponent character
\[
        \ell:H_n\longrightarrow\Z,\qquad
                   \ell(g)=1,\quad\ell(c)=\ell(v)=0.          \tag{D.65e}
\]
In particular $\langle g\rangle$ is infinite cyclic.

Eliminating $c$ gives the mutually inverse presentation
maps
\[
 \begin{aligned}
 H_n&\longrightarrow
       \langle G,V\mid V^2=1,\ [G,V]^n=1\rangle,
       &(g,v,c)&\longmapsto(G,V,[G,V]),\\
 (G,V)&\longmapsto(g,v).
 \end{aligned}                                                \tag{D.65f}
\]
Indeed $gvg^{-1}=cv$ gives $[g,v]=c$.
Conversely $V[G,V]V=[G,V]^{-1}$ follows from $V^2=1$,
so the defining dihedral relation is retained.
Now form the discrete amalgam with its stated
injective edge maps:
\[
       P_n=H_n*_{\langle g\rangle=\langle z^2\rangle}
                                                 \langle z\rangle .
                                                               \tag{D.65g}
\]
Its inverse presentation maps with (D.65a) send
\[
  (g,v,c,z)\longmapsto(z^2,v,[z^2,v],z),
        \qquad (z,v)\longmapsto(z,v).
                                                               \tag{D.65h}
\]
All defining relators have just been verified in
(D.65f), and each composite fixes its generators.
Thus $P_n=G_n$ through these explicit maps.
In particular both amalgam factors embed, and $c$
has exactly order $n$.
The following global character will be used after
completion:
\[
       \vartheta:G_n\longrightarrow\Z,\qquad
                        \vartheta(z)=1,\quad\vartheta(v)=0.
                                                               \tag{D.65i}
\]
Its restriction to $H_n$ is $2\ell$.

\emph{Length-one conjugacy and its full discrete centralizer.}
In the Bass--Serre tree of (D.65d), the vertex group
is $D_n$ and the edge groups are the specified
order-two subgroups.
Both $g$ and $a_r=c^rg$ have cyclically reduced stable
letter length one. Their successive positive and
negative powers are reduced, so they translate their
axes by one edge, and both axes contain the vertex
stabilized by $D_n$.

Suppose $hgh^{-1}=a_r$ in $H_n$.
The element $h$ carries the axis of $g$ to the axis
of $a_r$. Since $a_r$ translates by one edge, some
power of $a_r$ carries the vertex $hD_n$ to $D_n$.
Multiplying $h$ by that power therefore gives a
conjugator $d\in D_n$ with the same equation.
Rearrange it as
\[
                  g d^{-1}g^{-1}=d^{-1}c^r\in D_n.
                                                               \tag{D.65j}
\]
Britton's normal form forces
$d^{-1}\in\langle v\rangle$.
There are exactly two possibilities.
For $d=1$ the equation gives $r=0\bmod n$.
For $d=v$ the exact relation is
\[
                   vgv=c^{-1}g,
                                                               \tag{D.65k}
\]
and gives $r=-1\bmod n$.
Conversely those two displayed conjugators work.
Thus
\[
           c^rg\sim_{H_n}g
          \quad\Longleftrightarrow\quad
                      r\equiv0,-1\pmod n.                   \tag{D.65l}
\]

Assume now that $r$ is neither residue.
The character (D.65e) shows that $a_r$ cannot be
$H_n$-conjugate to any element of the amalgam edge
group $\langle g\rangle$: such a conjugacy would
have to be to $g$ itself, which (D.65l) excludes.
In the tree of (D.65g), $a_r$ fixes the $H_n$ vertex
and fixes no adjacent edge. If it fixed another
vertex, the connecting path would give an adjacent
fixed edge. Hence this is its unique fixed vertex.
Every element centralizing $a_r$ must preserve that
vertex, so
\[
                        C_{G_n}(a_r)=C_{H_n}(a_r).
                                                               \tag{D.65m}
\]
This is an argument in the discrete tree; it assumes
no fixed-path assertion for a profinite tree.

For $h\in C_{H_n}(a_r)$, commuting with a nontrivial
translation preserves and orients its axis.
Again its vertices form a single orbit under the
length-one translation $a_r$.
Thus $h=a_r^m k$ for an integer $m$ and a
$k\in D_n$ centralizing $a_r$.
The same Britton reduction as in (D.65j) forces
$k\in\langle v\rangle$.
Its nonidentity possibility is excluded by the
exact dihedral calculation
\[
              a_r v a_r^{-1}=c^{\,2r+1}v\ne v,
                                                               \tag{D.65n}
\]
where the inequality uses $\gcd(2r+1,n)=1$ and
$n\geq2$.
We have proved the complete discrete equality
\[
                         C_{G_n}(a_r)=\langle a_r\rangle .
                                                               \tag{D.65o}
\]

\emph{The actual closed surface cover.}
The homomorphism
$\epsilon:G_n\to C_2$ sending both $z$ and $v$
to the nontrivial element has kernel
\[
           \Pi_n=\langle X,Y\mid [X,Y]^n=1\rangle.
                                                               \tag{D.65p}
\]
Here the exact isomorphism
$G_n\cong\Pi_n\rtimes\langle v\rangle$ has
\[
         X=zv,\quad Y=vz,\qquad
         z=Xv,\qquad vXv=Y,\quad vYv=X.
                                                               \tag{D.65q}
\]
Substitution gives $[z^2,v]=[X,Y]$.
Conversely swapping $X,Y$ preserves the relation
in (D.65p), because it inverts their commutator.
The indicated maps between the presentation and
this semidirect product therefore invert each other,
proving both the kernel and its inclusion.

Let $Q_n$ be the finite Heisenberg group
\[
 Q_n=(\Z/n)^3,\qquad
 (a,b,t)(a',b',t')=(a+a',b+b',t+t'+ab').
\]
The map
\[
  \Pi_n\longrightarrow Q_n,\qquad
    X\longmapsto(1,0,0),\quad
    Y\longmapsto(0,1,0),\quad
    [X,Y]\longmapsto(0,0,1)                                  \tag{D.65r}
\]
is onto and has a commutator of exactly order $n$.
To identify its kernel, start with the actual
degree-$n^3$ regular cover of a once-punctured torus
given by the same homomorphism from its free group.
The boundary permutation is multiplication by that
central element of order $n$.
It has exactly $n^2$ cycles, each of length $n$,
so the covering surface has $n^2$ boundary components,
each mapping with degree $n$.
Fill each by a disk.
Van Kampen kills exactly the normal closure of
the corresponding conjugates of $[X,Y]^n$.
These are exactly the Reidemeister--Schreier relators
for the kernel of (D.65r): the original normal
closure is generated by its conjugates over the
$n^3$ cosets, and the $n$ cosets in a boundary
cycle give conjugates of the same boundary relator.
Thus the resulting closed surface group is precisely
that kernel, with the inclusion supplied by the
covering construction.
Its Euler characteristic and genus are
\[
         \chi=-n^3+n^2,\qquad
                    h_n=1+\frac{n^2(n-1)}2 .
                                                               \tag{D.65s}
\]
It has index $2n^3$ in $G_n$.
This proves directly that $G_n$ is virtually a
closed surface group.

\emph{The full completion bridge.}
Martino's theorem states that every group having
a surface subgroup of finite index is conjugacy
separable \cite[Theorem 3.7 of the inspected author
version]{C-Martino2007}.
Every finite-index subgroup of $G_n$ has the same
property, by intersection with the surface group
just constructed.
Consequently $G_n$ is hereditarily conjugacy
separable, and in particular residually finite.

We give the precise centralizer consequence and
its maps. For any such group $L$ and $a\in L$,
\[
           C_{\widehat L}(a)=\overline{C_L(a)}
                         \quad\text{inside }\widehat L.       \tag{D.65t}
\]
Let $x$ belong to the left-hand side.
For a finite-index normal subgroup $N\lhd L$,
choose $h\in L$ with $x=h\widehat n$,
$\widehat n\in\overline N$.
The commuting equation gives
$h^{-1}ah=\widehat n a\widehat n^{-1}$.
Its image modulo $N$ is therefore the image of $a$,
so the actual element $h^{-1}ah$ belongs to
$L_N=N\langle a\rangle$.
Both actual elements lie in the finite-index
subgroup $L_N$, and they are
conjugate in its full completion, since
$\overline N\subseteq\widehat{L_N}$.
Conjugacy separability of $L_N$ gives a conjugator
there. Write it as $n a^k$, with $n\in N$.
It follows that
$h^{-1}ah=n a n^{-1}$, so $hn\in C_L(a)$ and
$x\in C_L(a)\overline N$.
Intersecting these neighborhoods over $N$ puts
$x$ in the right-hand side of (D.65t).
The reverse inclusion follows from the closed
commuting equation.
Finite-index subgroups carry their full induced
profinite topologies, which is the only completion
inclusion used in this proof.

Apply (D.65t) to (D.65o).
The continuous extension of (D.65i) satisfies
$\widehat\vartheta(a_r)=2$.
Hence the map
\[
       \widehat\Z\longrightarrow
             \overline{\langle a_r\rangle}
                  \subseteq\widehat G_n,\qquad t\longmapsto a_r^t
                                                               \tag{D.65u}
\]
is injective: its composite with
$\widehat\vartheta$ is multiplication by two
on $\widehat\Z$.
It is onto the indicated procyclic closure by
definition.
Thus (D.65t) is the exact full procyclic
centralizer $\langle a_r\rangle\cong\widehat\Z$,
with no missing prime coordinates.
If $q^2=a_r$ in $\widehat G_n$, then $q$
centralizes $a_r$ and equals $a_r^t$.
The equation would give $2t=1$ in $\widehat\Z$,
which is impossible modulo two.
This excludes a square for the two disallowed
residue conditions above.
For $r=0$ the element is $g=z^2$.
For $r=-1$, (D.65k) gives
$c^{-1}g=vgv=(vzv)^2$.
Since $c$ has order $n$, these formulas prove
both directions of (D.65b).

\emph{Full units and both original reflection labels.}
Return to the full standard extension
$E_1=C_2\langle v\rangle\amalg\widehat\Z\langle z\rangle$
in D.64, with $g=z^2$ and $c_0=[g,v]$.
For each $n$, the generator maps give an actual
continuous epimorphism
\[
           q_n:E_1\longrightarrow\widehat G_n,\qquad
                              z\longmapsto z,\quad v\longmapsto v.
                                                               \tag{D.65v}
\]
Their dense discrete image generates $G_n$, so
compactness makes the map onto its full completion.
Let $\beta\in C_\Omega(b_1)$ have the unique
normalized extension from D.63.
D.64 gives $C_{F_0}(v_i)=1$, so its image
$a=\Psi(g)=\Psi(z)^2$ is exactly
\[
 a=
 \begin{cases}
   c_0^{(\mu-1)/2}g,&i=0,\\
   c_0^{(\mu+1)/2}g^{-1},&i=1,
 \end{cases}
 \qquad \mu=\lambda(\beta)\in\widehat\Z^\times .
                                                               \tag{D.65w}
\]
First take $i=0$ and put $r=(\mu-1)/2\in\widehat\Z$.
Choose an integer representative $r_n$ of $r\bmod n$.
Its image in (D.65v) is $c^{r_n}g$ and is a square.
The integer $2r_n+1$ is prime to $n$, because it
is the reduction of the full unit $\mu$.
Criterion (D.65b) gives
$r_n\equiv0,-1\pmod n$.
Equivalently,
\[
                         \mu\equiv1\text{ or }-1\pmod{2n}
                             \quad\text{for every }n\geq2.
                                                               \tag{D.65x}
\]
The finite-intersection argument of D.64k now
forces a single global sign $\mu=1$ or $\mu=-1$.

For $i=1$, inversion and the specified conjugation
give the exact comparison
\[
        g^{-1}a^{-1}g
            =c_0^{-(\mu+1)/2}g
            =c_0^{((-\mu)-1)/2}g .
                                                               \tag{D.65y}
\]
This element is the square of
$g^{-1}\Psi(z)^{-1}g$.
It is therefore precisely the already treated
first-label expression for the full unit $-\mu$,
and again $\mu\in\{1,-1\}$.
No reflection label or named cusp has been rescaled.

D.64 proves that the character pair is injective
for $b_1$, and explicitly realizes all four pairs
with these two multipliers.
This gives (D.65c).
The normalizer and actual orbit-fibre conclusions
follow exactly as in D.64.
The proof restricts the standard centralizers;
the separate existence of a conjugator for an
arbitrary negative involution remains the full
equation in D.62.

The program
\path{code/part_c/standard_oneclass_square_d65.py}
checks the defining word maps, the finite
dihedral identities in the axis calculation,
and the exact Heisenberg covering data.
The square obstruction in the full completion
uses the proved discrete centralizer and the
specified hereditary-conjugacy-separability bridge,
beyond those finite checks.
\end{proof}

\begin{claim}{D.66}{Proved}{The exact norm in a scalar cyclic torus cover}
Let $a$ be an actual element of order $d=3$ or $4$ in the
full once-punctured-torus loop-outer group, and put $p=3$
or $2$, respectively. At every scalar coordinate retained
by D.54--D.55, write
\[
 A_q(a)=uI,\qquad
 \begin{cases}
  u=\lambda_q^2,\quad \lambda_q\ne1,&d=3,\\
  u^2=-1,\quad \lambda_q=-1,&d=4.
 \end{cases}                                                   \tag{D.66a}
\]
Thus $q\equiv1\pmod d$, and the reduction $\bar u\in\F_q^\times$
has exact order $d$. For every nonzero character
$\chi:F=\widehat F(x,y)\to C_q$, its kernel $N$ is the full
completion of the actual degree-$q$ cover of type $(1,q)$.
The complete primary homology comparison is
\[
 H_1(N,\Z_p)
   \cong (\Z_p^2,A_p(a))\ \oplus\
          \Z_p[C_d]^{(q-1)/d}.                              \tag{D.66b}
\]
Here the action is the one induced by a labelled cusp-normalizer
element above $a$; the two displayed summands are the actual
deck-average handle summand and boundary summand.

Label the boundary classes by $b_i$, $i\in\F_q$, with the
distinguished original cusp labelled $0$. The boundary module
has the exact relation $\sum_i b_i=0$, and the action is
$b_i\mapsto b_{\bar u i}$. If $R$ contains one representative
of each nonzero $\langle\bar u\rangle$-orbit, then
\[
 w=-\sum_{r\in R}b_r,\qquad
 (1+T+\cdots+T^{d-1})w=b_0,                              \tag{D.66c}
\]
where $T$ denotes this action. In particular the distinguished
unit cusp-power class is an integral norm, at every primary
precision. More fully, for $M=H_1(N,\Z_p)$,
\[
 H^2(C_d,M)=0,\qquad H^1(C_d,M)\cong C_p.                \tag{D.66d}
\]
The actual cyclic extension over this cover, after quotienting
its kernel to $M$, has the explicit splitting given below.
These statements identify exactly why the invariant-character
norm contradiction of D.54--D.55 does not extend to the scalar
coordinates (D.66a). They neither construct an action having
such a coordinate nor decide the remaining scalar multiplier
or conjugacy questions.
\end{claim}
\passage{Section 30, folios 194-196, scans 383-387; Section 31, folios 198-199 verso, scans 391-394; Part D(b), full loop-outer finite actions in type (1,1)}

\begin{proof}
\emph{The actual extension and its finite quotient.}
Choose an actual representative $\alpha$ of order $d$,
using D.27 for $d=3$ and D.46 for $d=4$, and set
\[
 E=F\rtimes_\alpha\langle s\rangle,\quad c=[x,y],\quad
 L=\langle c\rangle .
\]
The chosen $s$ has quotient label $a$. Choose a cusp-normalizer
element $v$ with the same quotient label. The exact relations are
\[
 vcv^{-1}=c^\lambda,\qquad v^d=c^\kappa,\qquad
             \lambda_p=1,\quad \kappa_p\in\Z_p^\times.       \tag{D.66e}
\]
For $d=3$ these are D.48. For $d=4$, the element $v^2$
lies over the positive involution $a^2$. Its square exponent
is a two-adic unit by D.55's full cusp-normalizer calculation.
Also $\lambda_2=1$: the order-four matrix on $\Z_2^2$ has
square $-I$, characteristic polynomial $X^2+1$, and determinant
one. The peripheral multiplier equals that determinant.

Let $\bar u$ be as in (D.66a), and write $C_q$ additively.
Both conjugation by $s$ and by $v$ satisfy
\[
             \chi(vhv^{-1})=\bar u\,\chi(h)
                         \qquad(h\in F).                    \tag{D.66f}
\]
Indeed they induce the same homology map, and every character
factors through the actual abelianization of $F$.
Thus $N=\ker\chi$ is normal in $E$. Since $c\in N$,
the class of $v$ in $E/N$ has exact order $d$.
Choose a discrete word $f$ with $\chi(f)=1$. The two mutually
inverse finite quotient maps are
\[
 \begin{aligned}
 C_q\rtimes_{\bar u} C_d&\longrightarrow E/N,
             &(i,j)&\longmapsto f^i v^jN,\\
 E/N&\longrightarrow C_q\rtimes_{\bar u} C_d,
             &hv^jN&\longmapsto(\chi(h),j),
                    \quad 0\leq j<d.
 \end{aligned}                                               \tag{D.66g}
\]
The multiplication on the left is
$(i,j)(k,l)=(i+\bar u^j k,j+l)$; the relation $v^d\in N$
and (D.66f) verify both maps and their inverse identities.

Put $E_N=N\langle v\rangle$, the inverse image of
$\{0\}\rtimes C_d$ in this exact quotient. It has index
$q$ in $E$, with kernel $N$ and quotient $C_d$.
For completeness it has an actual order-$d$ section.
Write the image of the original $s$ as $(j,1)$ and choose
the integer residue
\[
                    k=\frac{j}{\bar u-1}\quad\text{in }\F_q.
\]
The denominator is nonzero. Conjugation by $(k,0)$ sends
$(j,1)$ to $(j+(1-\bar u)k,1)=(0,1)$. Consequently
$s_0=f^k s f^{-k}$ belongs to $E_N$, has exact order $d$,
and gives $E_N=N\rtimes\langle s_0\rangle$ with the
specified quotient label. No character extension
$E\to C_q$ with trivial action is being substituted for
(D.66g).

\emph{The actual cover and the boundary labels.}
The restriction of $\chi$ to the dense discrete free group
is onto. Its punctured-torus cover has degree $q$ and
trivial cusp monodromy, hence $q$ degree-one boundary
lifts. Euler characteristic $-q$ gives genus one.
The finite-index completion map identifies its full
fundamental-group completion with precisely $N$.
The chosen boundary representatives are
\[
           c_i=f^i c f^{-i}\in N\qquad(0\leq i<q).
                                                               \tag{D.66h}
\]
They have distinct cover labels. In homology their classes
$b_i$ have exactly the single relation $\sum_i b_i=0$.
For $j$ the representative of $\bar u i$ modulo $q$, put
$n_i=(vf^iv^{-1})f^{-j}\in N$, by (D.66f). Then the
actual conjugacy equation is
\[
                 vc_iv^{-1}=n_i c_j^\lambda n_i^{-1}.     \tag{D.66i}
\]
In $p$-primary homology $\lambda_p=1$, so its action is
exactly $T b_i=b_{\bar u i}$. The label zero is fixed,
and every nonzero label has orbit length $d$.

\emph{Deck averaging and both inverse boundary-module maps.}
Let $D=F/N=C_q$ act on $M=H_1(N,\Z_p)$ by conjugation
of deck lifts, and set
\[
                  e=\frac1q\sum_{h\in D}h.                \tag{D.66j}
\]
The prime $q$ differs from $p$. Although $T$ need not
commute with individual deck transformations, it permutes
them by $h\mapsto h^{\bar u}$; hence it commutes with $e$.
Filling all $q$ boundaries gives the actual compact torus
cover associated to the lattice
$\Lambda=\ker(\Z^2\to C_q)$, with covering map
$\R^2/\Lambda\to\R^2/\Z^2$. Its deck maps are translations,
so they act trivially on the filled homology. The boundary
module is
\[
 B_\partial=\Z_p\{b_i:i\in\F_q\}/\Z_p\Bigl(\sum_i b_i\Bigr).
                                                               \tag{D.66k}
\]
It has zero deck average, since $q$ is a unit. The boundary
sequence therefore gives the exact identifications
\[
 (1-e)M=B_\partial,\qquad
 eM\xrightarrow{\ \cong\ }H_1(\R^2/\Lambda,\Z_p)
       \xrightarrow{\ \cong\ }H_1(\R^2/\Z^2,\Z_p).       \tag{D.66l}
\]
The second map is induced by the index-$q$ lattice inclusion;
its determinant is a $p$-adic unit. Both maps intertwine
the induced homology actions, so $T|_{eM}=A_p(a)$ under
these specified identifications.

Write $m=(q-1)/d$. With $R$ as in the statement, define
\[
 \begin{aligned}
 \Phi:\bigoplus_{r\in R}\Z_p[C_d]e_r&\longrightarrow B_\partial,
                    &T^j e_r&\longmapsto b_{\bar u^j r},\\
 \Psi:B_\partial&\longrightarrow
                         \bigoplus_{r\in R}\Z_p[C_d]e_r,
                    &b_{\bar u^j r}&\longmapsto T^j e_r,\\
                    &&b_0&\longmapsto
                         -\sum_{r\in R}\sum_{j=0}^{d-1}T^j e_r.
 \end{aligned}                                               \tag{D.66m}
\]
The displayed image of the total boundary relation is zero.
It also makes $\Psi$ equivariant at the fixed label zero.
Every nonzero label has a unique expression $\bar u^j r$,
so these maps are mutually inverse on all their generators.
This proves the free cyclic-module assertion in (D.66b).
Applying $\Phi$ to the negative sum of the $e_r$ proves
(D.66c) integrally.

\emph{The exact cyclic cohomology and abelianized extension.}
On the handle summand $H=eM$, D.48 and D.55 give respectively
\[
 \begin{array}{c|c|c|c}
 d&p&\text{identity on }H&\det(T-I)|_H\\ \hline
 3&3&T^2+T+I=0&3\\
 4&2&T^2=-I&2
 \end{array}                                                   \tag{D.66n}
\]
Thus the cyclic norm $\mathcal N=1+T+\cdots+T^{d-1}$ is
zero on $H$, while $H^T=0$. The cokernel of $T-I$ on
$H\cong\Z_p^2$ is $C_p$: its determinant has valuation one,
so its two invariant factors have valuations zero and one.
On each free summand $\Z_p[C_d]$, the invariants are
exactly the multiples of the sum of its $d$ basis elements,
which are exactly its norm image. The kernel of the norm
is the augmentation ideal, the image of $T-I$.
The periodic cyclic cochain complex consequently gives
\[
 M^T/\mathcal N M=0,\qquad
 \ker(\mathcal N)/(T-I)M\cong C_p,
\]
which is precisely (D.66d).

Here is an explicit splitting retaining the cusp exponent.
Let $K=\ker(N\to M)$ and $P=E_N/K$. This is an actual
quotient, since $K$ is characteristic in $N$. Write the
kernel $M$ additively, and still denote the image of $v$
by $v$. Its relations are
\[
           vmv^{-1}=Tm,\qquad v^d=\kappa_p b_0.
                                                               \tag{D.66o}
\]
They are complete: every element of $P$ has a unique
normal form $m v^j$, $0\leq j<d$, and these relations
give its multiplication. By (D.66c), the element
\[
                    t=(-\kappa_p w)v
\]
has $t^d=0$ in the additive kernel, and has exact order
$d$ because of its quotient label. The inverse
isomorphisms are
\[
 \begin{aligned}
 M\rtimes_T C_d&\longrightarrow P,&(m,j)&\longmapsto m t^j,\\
 P&\longrightarrow M\rtimes_T C_d,&m&\longmapsto(m,0),\\
 &&v&\longmapsto(\kappa_p w,1).
 \end{aligned}                                               \tag{D.66p}
\]
The norm identity verifies the image of $v^d$, and
the two composites fix the indicated kernel and generator.
Reducing these maps modulo every $p^n$ gives the same
splitting at that exact precision. All module solutions
of $\mathcal N z=-\kappa_p b_0$ form
$-\kappa_p w+\ker\mathcal N$.

The actual section $s_0$ constructed above has one such
vector relative to $v$, but this calculation does not
assert that every vector in that coset lifts to a
finite-order element before the quotient by $K$.
In the invariant-character proof all boundary labels
were fixed and the primary norm vanished modulo $p$.
Here the explicit free boundary summands (D.66m) give
the precise integral preimage for the same distinguished
cusp class. Thus there is no remaining norm obstruction
from this cover module, without any assertion that a
scalar-coordinate full action exists.

The executable certificate
\texttt{code/part\_c/torus\_scalar\_cover\_norm\_d66.py}
checks the finite quotient conjugation, the two inverse
boundary-module maps, the deck conjugacy and averaging,
the integral norm preimage, and the handle matrices
for the first admissible cubic and quartic primes.
The arbitrary-prime assertions follow from the explicit
maps above, not from this finite enumeration.
\end{proof}

\begin{claim}{D.67}{Proved}{The standard negative normalizers and their finite overgroups}
Retain the two standard negative involutions $b_0,b_1$
and the actual automorphisms $\alpha_j,\zeta_j$ of D.64.
Put
\[
 R(X)=Y,\quad R(Y)=X^{-1},\quad r=[R],\quad
 \delta=r^2=[\zeta_1],\quad
 V_j=C_\Omega(b_j),\quad D=\langle r,b_0\rangle .          \tag{D.67a}
\]
Then $D$ is the specified geometric dihedral group of order
eight, and
\[
 \begin{gathered}
 b_1=rb_0,\qquad
 V_0=\langle\delta,b_0\rangle,\qquad
 V_1=\langle\delta,b_1\rangle,\qquad
 V_0\cap V_1=\langle\delta\rangle,\\
 N_\Omega(V_0)=N_\Omega(V_1)=D,\qquad
 N_\Omega(D)=D,\qquad C_\Omega(D)=\langle\delta\rangle.
 \end{gathered}                                               \tag{D.67b}
\]
The complete list of finite subgroups of $\Omega$ containing
this specified $b_0$ is
\[
                     \langle b_0\rangle,\qquad V_0,\qquad D.
                                                               \tag{D.67c}
\]
Thus $D$ is its unique maximal finite overgroup. Conjugation
by any $\theta\in\Omega$ transports this exact list and all
the normalizer and centralizer equalities to the corresponding
specified conjugate groups. The statement does not assert
that every negative involution is conjugate to $b_0$ or $b_1$.
\end{claim}
\passage{Section 30, folios 194-196, scans 383-387; Section 31, folios 198-199 verso, scans 391-394; Part D(b), full loop-outer finite actions in type (1,1)}

\begin{proof}
\emph{Actual word maps for the dihedral group.}
Let $B(X)=X$, $B(Y)=Y^{-1}$. Its relation with the
cusp-normalized representative is
\[
                         \alpha_0=\operatorname{Inn}(Y)B.
                                                               \tag{D.67d}
\]
All of $R,B,\alpha_j,\zeta_j$ are actual automorphisms of
the free group and extend to its full completion. The
following identities are identities of automorphisms,
checked on the two free generators:
\[
 \begin{aligned}
 R^4&=1,& R^2&=\zeta_1,& BRB&=R^{-1},\\
 RB&=\alpha_1,& R\alpha_0&=\operatorname{Inn}(X^{-1})\alpha_1,\\
 R\alpha_0R^{-1}&=\operatorname{Inn}(X^{-1})\zeta_0,\\
 R\zeta_0R^{-1}&=\operatorname{Inn}(Y^{-1})\alpha_0,\\
 R\alpha_1R^{-1}&=\zeta_1\alpha_1,&
 \alpha_0\zeta_0&=\operatorname{Inn}(Y)\zeta_1.
 \end{aligned}                                               \tag{D.67e}
\]
Their peripheral maps, with $c_0=[X,Y]$, include
\[
 R(c_0)=X^{-1}c_0X,\qquad
 B(c_0)=Y^{-1}c_0^{-1}Y.                                  \tag{D.67f}
\]
Consequently the indicated outer classes belong to the
full peripheral group. The homology matrices of $R$ and
$B$ are respectively
\[
                 \begin{pmatrix}0&-1\\1&0\end{pmatrix},
                 \qquad\begin{pmatrix}1&0\\0&-1\end{pmatrix}.
                                                               \tag{D.67g}
\]
Their eight matrices $R_*^k B_*^\epsilon$, $0\leq k<4$,
$0\leq\epsilon<2$, are distinct. The actual relations
(D.67e) give at most eight elements, while these matrices
give at least eight distinct outer classes. Hence
$D\cong D_8$ with exactly these labels.
The linear maps $(u,v)\mapsto(-v,u)$ and
$(u,v)\mapsto(u,-v)$ on the torus, punctured at the origin,
realize these homology matrices. The discrete
once-punctured-torus identification
$\operatorname{Out}(F(X,Y))\cong\operatorname{GL}_2(\Z)$
proved in A.13 identifies their induced outer classes with
the displayed word maps. This supplies the asserted actual
geometric realization.

By D.64--D.65 the full centralizers are exactly
$V_0=\langle b_0,[\zeta_0]\rangle$ and
$V_1=\langle b_1,\delta\rangle$, both of order four.
Equations (D.67d)--(D.67e) give
$[\zeta_0]=\delta b_0$ and $b_1=rb_0$, proving the
first line of (D.67b) and identifying their intersection.
They are the two specified index-two Klein groups in $D$.

\emph{Their full normalizers.}
The unique nonidentity multiplier-one element of each
$V_j$ is $\delta$; its two other nonidentity elements have
multiplier $-1$. The multiplier homomorphism on $\Omega$
is conjugacy invariant. Thus conjugation by a member of
$N_\Omega(V_j)$ fixes $\delta$ and either fixes or exchanges
the other two elements. Its image in $\operatorname{Aut}(V_j)$
therefore has order at most two. Its kernel is
$C_\Omega(V_j)$, and the exact inclusions
\[
 V_j\leq C_\Omega(V_j)\leq C_\Omega(b_j)=V_j
\]
make this kernel precisely $V_j$. It follows that the full
normalizer has at most eight elements. The element $r$
realizes the nontrivial exchange in each case by (D.67e),
and the group $D$ already normalizes both $V_j$.
This proves both full normalizer equalities.

The mod-two homology map on $D$ has exact kernel $V_0$:
$B_*\equiv I$ and $R_*\equiv\begin{psmallmatrix}0&1\\1&0\end{psmallmatrix}$
over $\F_2$. Every element of $N_\Omega(D)$ preserves this
kernel, since its actual homology map conjugates the
mod-two matrices and preserves the identity matrix.
It therefore belongs to $N_\Omega(V_0)=D$.
Thus $N_\Omega(D)=D$. The full centralizer is now its
ordinary group center, which is $\langle\delta\rangle$.

\emph{Every finite overgroup of the specified two-class reflection.}
Let $K\leq\Omega$ be finite and contain $b_0$.
D.55 gives the complete nine-type list
$1,C_2,C_3,C_4,C_2^2,C_6,S_3,D_8,D_{12}$.
The first and third types cannot contain an involution.
The unique involution of a cyclic order-four group is
positive by D.55; that of a cyclic order-six group is
positive by D.56. Hence neither can be $b_0$.

We exclude $S_3$ and $D_{12}$ by the exact mod-two
representation. Every element $h$ of order three has
mod-two homology matrix of exact order three by D.54.
In $S_3$ every involution inverts its normal $C_3$.
In $D_{12}$ the central involution is the positive
involution of a cyclic-six subgroup, so a negative
involution is a reflection and also inverts the
normal $C_3$. Thus either type would give
\[
                    b_0 h b_0^{-1}=h^{-1},\qquad |h|=3.
                                                               \tag{D.67h}
\]
But $A_2(b_0)\bmod2=I$. Applying the actual homology
map and reducing modulo two would make the order-three
matrix equal to its inverse, a contradiction.

For $K\cong C_2^2$, every element commutes with $b_0$,
so $K\leq V_0$, and equality follows from their orders.
For $K\cong D_8$, its central involution is the square
of an order-four element, hence is positive. It
commutes with $b_0$, so it is the unique positive
nonidentity element $\delta$ of $V_0$.
The group $\langle\delta,b_0\rangle=V_0$ has index two
in $K$ and is normal. Consequently
$K\leq N_\Omega(V_0)=D$, and equality again follows
from the orders. The remaining cyclic-two group is
exactly $\langle b_0\rangle$. This proves (D.67c).

Finally $K\mapsto\theta K\theta^{-1}$ and its inverse
are exact bijections between finite overgroups of $b_0$
and of $\theta b_0\theta^{-1}$. The same conjugation
identifies each full normalizer and centralizer.
This gives the asserted transport without any claim
that these specified orbits exhaust all negative actions.

The executable certificate
\path{code/part_c/standard_negative_overgroups_d67.py}
checks the actual word identities, all eight signed
permutation automorphisms and their multiplication,
both labelled Klein subgroups, and the mod-two
order-three inversion obstruction. The full upper
bounds use the proved centralizers and nine-type list,
not a finite enumeration of the ambient profinite group.
\end{proof}

\begin{claim}{D.68}{Proved}{The exact dihedral quotient of a quartic torus action}
Let $a$ have order four in the full once-punctured-torus
loop-outer group. Choose an actual order-four representative
$\alpha$ as in D.46, and write
\[
 F=\widehat F(X,Y),\qquad c=[X,Y],\qquad
 E=F\rtimes_\alpha\langle s\rangle .
\]
There is exactly one nonzero character $\chi:F\to C_2$
fixed by $a$. Extend it to $\widetilde\chi:E\to C_2$
by $\widetilde\chi(s)=0$. Every cusp-normalizer element
$v$ with quotient label $a$ has
\[
                   \widetilde\chi(v)=1.                    \tag{D.68a}
\]

There is an explicitly specified discrete free basis $x,y$
with $[x,y]=c$, $\chi(x)=1$, and $\chi(y)=0$.
For this basis the closed normal subgroup
\[
                K=\langle\!\langle x^2,(xy)^2
                           \rangle\!\rangle_F              \tag{D.68b}
\]
is $\alpha$-invariant. The actual quotient and its inverse
presentation maps identify
\[
 \begin{split}
 F/K&\cong\mathcal D=\widehat\Z\rtimes C_2
       =\langle r,t\mid t^2=1, trt=r^{-1}\rangle,\\
 x&\longmapsto t,\qquad y\longmapsto r,
                 \qquad c\longmapsto r^{-2}.
 \end{split}                                               \tag{D.68c}
\]
Here $r$ generates the full $\widehat\Z$ rotation subgroup.
The induced automorphism is exactly
\[
 \overline\alpha(r)=r^{-\lambda},\qquad
 \overline\alpha(t)=r^j t,\qquad
 j_2\in\Z_2^\times,\qquad (1-\lambda)j=0.                \tag{D.68d}
\]
It has order two. For every odd prime $q$ the further
specified quotient $\mathcal D\to D_{2q}$ sends $r$ to
a rotation of order $q$ and $t$ to a reflection.
On that quotient the induced action is the identity
exactly at the coordinates $\lambda_q=-1$; at
$\lambda_q=1$ it is an involution inverting the rotation.
In particular each remaining scalar quartic coordinate
has an actual nonabelian quotient $F\twoheadrightarrow D_{2q}$
fixed pointwise by this representative.

The proof also gives the full rank-three homology matrix
of the actual index-two cover, with all cusp labels retained.
These conclusions constrain every actual quartic action;
they do not construct a negative multiplier coordinate
or settle its remaining exclusion or conjugacy problem.
\end{claim}
\passage{Section 30, folios 194-196, scans 383-387; Section 31, folios 198-199 verso, scans 391-394; Part D(b), full loop-outer finite actions in type (1,1)}

\begin{proof}
\emph{The specified mod-two character and a cusp-preserving basis.}
By D.55 the positive square of $a$ acts as $-I$ on
$H_1(F,\Z_p)$ for every prime. Hence
\[
 A_p(a)^2=-I,\qquad \lambda^2=1,\qquad \lambda_2=1.       \tag{D.68e}
\]
The last assertion also follows directly from the
characteristic polynomial $Z^2+1$ over $\Q_2$ and
the determinant interpretation of the cusp multiplier.
The matrix $A_2(a)$ cannot be identity modulo two:
otherwise $A_2(a)=I+2B$ would have square identity
modulo four, contradicting its square $-I$.
Its reduction is therefore a nonidentity unipotent
matrix over $\F_2$. Its fixed dual line has exactly
one nonzero vector. This is the stated character $\chi$.
Conjugation on $F$ does not change a character, so
this character is defined by the outer action as well.

Here are the basis maps and their actual inverses,
according to the values of $\chi$ on $X,Y$:
\[
\begin{array}{c|cc|cc}
 (\chi(X),\chi(Y))&x&y&X&Y\\ \hline
 (1,0)&X&Y&x&y\\
 (0,1)&XYX^{-1}&X^{-1}&y^{-1}&yxy^{-1}\\
 (1,1)&X&YX^{-1}&x&yx.
\end{array}                                                \tag{D.68f}
\]
Free reduction verifies both composites and $[x,y]=[X,Y]$
in every row. The same substitutions give mutually
inverse maps on full free profinite groups. No peripheral
conjugation or change of multiplier has been introduced.

Put $N=\ker\chi$. Its actual discrete double cover of
the punctured torus has type $(1,2)$. Schreier rewriting
with transversal $1,x$ gives its full free profinite basis
\[
 u=x^2,\qquad B=y,\qquad C=xyx^{-1},\qquad c=CB^{-1}.
                                                               \tag{D.68g}
\]
The second based peripheral generator is $xcx^{-1}$;
its homology class is the negative of that of $c$.
Write $M=H_1(N,\widehat\Z)$ in the ordered basis
$[u],[B],[C]$ and put $b=[C]-[B]$.
The actual deck action $J$, induced by conjugation
by $x$, is
\[
 J[u]=[u],\qquad J[B]=[C],\qquad J[C]=[B].                \tag{D.68h}
\]
Before abelianization the last identity is
$xCx^{-1}=uBu^{-1}$. Homology of the actual inclusion
$N\hookrightarrow F$ is the map
\[
 I:M\longrightarrow H_1(F,\widehat\Z),\qquad
 I=\begin{pmatrix}2&0&0\\0&1&1\end{pmatrix},\qquad
 \ker I=\widehat\Z b.                                    \tag{D.68i}
\]
It factors through filling both cusps of the covered
torus and then the index-two closed-torus lattice
inclusion. Its image is precisely that lattice.
Thus after
tensoring each primary coordinate with $\Q_p$, the
$+1$ eigenspace of $J$ is the handle space and the
$-1$ eigenspace is precisely $\Q_p b$.

\emph{The cusp root has the nonzero character label.}
Choose $v$ above $a$ in the normalizer of the actual
cusp subgroup. The cusp relations, including their
primary exponent, are
\[
 vcv^{-1}=c^\lambda,\qquad v^4=c^\kappa,\qquad
                         \kappa_2\in\Z_2^\times.          \tag{D.68j}
\]
The unit assertion is D.55 applied to $v^2$, which
lies over the positive involution $a^2$.
The invariant character extends to $E$ by sending $s$
to zero. Since $\chi$ is invariant, $N$ is normal in
$E$. Let $T$ be conjugation by $v$ on $M$.
The element $vxv^{-1}$ has the same nonzero character
as $x$, so its action on $M$ is also $J$.
Consequently $T$ commutes with $J$.

On the boundary line $T$ is multiplication by $\lambda$.
On the filled handle quotient it induces the actual
matrix $A(a)$. Equations (D.68e) and (D.68i), first
over $\Q_p$ and then in the torsion-free integral
module for each $p$, prove
\[
                   T^2=-J.                                \tag{D.68k}
\]
At the two-primary coordinate $Tb=b$, and therefore
\[
 \mathcal N_4(T)=(I+T)(I+T^2)
           =(I+T)(I-J)=2(I-J),\qquad
 \mathcal N_4(T)M_2=2\Z_2b.                               \tag{D.68l}
\]
If $\widetilde\chi(v)$ were zero, write $v=ns$ with
$n\in N$. Conjugation by $v$ and by $s$ then agree
on $M$. The actual identity
$(ns)^4=n\alpha(n)\alpha^2(n)\alpha^3(n)$ would give
$\kappa_2b=\mathcal N_4(T)[n]$, contrary to
(D.68l) and the unit exponent in (D.68j).
This proves (D.68a).

Let $S$ denote the action of $\alpha$ on $M$.
Writing $v=f_0s$, the proved label means $\chi(f_0)=1$,
so conjugation by $f_0$ on $M$ is exactly $J$.
Thus the actual actions satisfy
\[
            T=JS,\qquad S=JT,\qquad
            S^2=-J,\qquad Sb=-\lambda b.                  \tag{D.68m}
\]
In particular the primary norm of $S$ at two is zero.

The entire matrix is explicit. Write the base homology
matrix in the basis $[x],[y]$ as
$A(a)=\left(\begin{smallmatrix}a_{11}&a_{12}\\
a_{21}&a_{22}\end{smallmatrix}\right)$. Then
\[
 T=\begin{pmatrix}
 a_{11}&a_{12}/2&a_{12}/2\\
 a_{21}&(a_{22}+\lambda)/2&(a_{22}-\lambda)/2\\
 a_{21}&(a_{22}-\lambda)/2&(a_{22}+\lambda)/2
 \end{pmatrix},\qquad S=JT.                               \tag{D.68n}
\]
Every half coefficient belongs to $\widehat\Z$:
the invariant mod-two character makes $a_{12}$ even,
and $a_{22}$ and $\lambda$ are odd. To verify the
matrix without choosing an integral splitting, use
$TJ=JT$, $Tb=\lambda b$, and the actual equality
$IT=A(a)I$. These equations determine the displayed
entries uniquely.

\emph{The full dihedral quotient and both presentation maps.}
Define the actual character
\[
 \ell:N\longrightarrow\widehat\Z,\qquad
       \ell(u)=0,\quad\ell(B)=1,\quad\ell(C)=-1.
                                                               \tag{D.68o}
\]
It is surjective, and $\ell\circ\operatorname{Inn}(x)=-\ell$.
Thus its kernel is normal in $F$. Equation (D.68n)
gives, as an equality of characters of $N$,
\[
                         \ell\alpha=-\lambda\ell.         \tag{D.68p}
\]
Since $-\lambda$ is a unit, $\ker\ell$ is
$\alpha$-invariant as a subgroup, in both directions.

The homomorphism $F\to\mathcal D$ sending $x$ to $t$
and $y$ to $r$ restricts on $N$ to $n\mapsto r^{\ell(n)}$.
It is onto, and its kernel is $\ker\ell$. Indeed an
element outside $N$ maps to the nontrivial reflection
coset, while the rotation parameter detects every
element of $N/\ker\ell$.
This is also precisely the quotient (D.68b): in that
quotient $x^2=1$ and $(xy)^2=1$, so
$xyx=y^{-1}$. Conversely the images $x=t,y=r$ in
$\mathcal D$ satisfy these two relators.
The inverse homomorphism sends $r$ to $yK$ and $t$
to $xK$. Both composites fix the named generators.
The rotation subgroup is full $\widehat\Z$, also
directly witnessed by the surjective character $\ell$.
This proves all assertions about the actual kernel
and maps in (D.68b)--(D.68c).

Equation (D.68p) specifies the image of $r$ under the
induced automorphism. Preservation of $\chi$ specifies
the reflection coset of the image of $t$. Hence there
is a unique $j\in\widehat\Z$ with the two formulas
in (D.68d). Their square sends $r$ to $r$ and $t$ to
$r^{(1-\lambda)j}t$. Since $\alpha^4=1$, its fourth
power gives $2(1-\lambda)j=0$ in $\widehat\Z$.
Multiplication by two is injective, proving the last
equation of (D.68d) and $\overline\alpha^2=1$.

For the retained cusp label one can compute $j$ more
precisely. The quotient extends to the actual map
\[
 E\longrightarrow\mathcal D\rtimes_{\overline\alpha}C_4,
 \qquad f\longmapsto\pi(f),\quad s\longmapsto s.
                                                               \tag{D.68q}
\]
Write $\pi(f_0)=r^m t$, using $\chi(f_0)=1$.
In this quotient the exact square is
\[
 \pi(v)^2=r^{(1+\lambda)m-j}s^2,\qquad
 \pi(v)^4=r^{2((1+\lambda)m-j)}=r^{-2\kappa}.
\]
Injectivity of multiplication by two in the rotation
parameter gives
\[
                         j=(1+\lambda)m+\kappa.            \tag{D.68r}
\]
At two, this reads $j_2=2m_2+\kappa_2$, so $j_2$ is
odd. In particular the induced action is nonidentity
and has exact order two. The formula retains the
original exponent and does not rescale the cusp.

Finally reduce the full rotation parameter modulo an
odd prime $q$. If $\lambda_q=-1$, equation
$(1-\lambda)j=0$ gives $j_q=0$ already in $\Z_q$.
Both displayed generators of $D_{2q}$ are fixed.
If $\lambda_q=1$, the rotation is inverted and is
nontrivial because $q$ is odd, so the induced action
has exact order two. The cusp has image $r^{-2}$,
which still generates the order-$q$ rotation subgroup.
Thus these are actual peripheral finite quotients,
not quotients inferred only from a dimension count.

The executed certificate
\path{code/part_c/torus_quartic_dihedral_d68.py}
checks the discrete basis maps, cusp and deck words,
the double-cover matrices and integral norm identity,
the standard quartic representative, and the exact
dihedral formulas at finite primary coordinates.
It does not assert that a permitted coordinate
model lifts to a full torus automorphism.
\end{proof}

\begin{claim}{D.69}{Proved}{A negative quartic coordinate forces a full rotation image}
Use the actual quartic representative, adapted basis, and dihedral
quotient of D.68. For every odd prime $q$ at which $\lambda_q=-1$,
write
\[
 \rho:F\twoheadrightarrow Q=D_{2q},\qquad
 N=\ker\rho,\qquad \beta=\alpha^2,\qquad
 \sigma=s^2,\qquad C=C_F(\sigma),\qquad C_0=C\cap N.
                                                               \tag{D.69a}
\]
The specified rotations form $R=\langle\rho(y)\rangle\cong C_q$.
Then restriction of the actual quotient map gives the exact
sequence and dimension
\[
 1\longrightarrow C_0\longrightarrow C
       \xrightarrow{\ \rho|_C\ }R\longrightarrow1,\qquad
                  \dim_{\F_2}H^1(C_0,\F_2)=q-1.             \tag{D.69b}
\]
The cover defining $N$ is the full completion of an actual
surface of type $(q,2)$.
There are exactly two $N$-conjugacy classes of order-two
sections in $N\rtimes\langle\sigma\rangle$; representatives are
$\langle\sigma\rangle$ and
$\langle x\sigma x^{-1}\rangle$.
Their exact normalizers are
\[
 C_0\times\langle\sigma\rangle,\qquad
 x\bigl(C_0\times\langle\sigma\rangle\bigr)x^{-1}.
                                                               \tag{D.69c}
\]
They are also the two order-two subgroup classes in
$N\rtimes\langle s\rangle$.
Let $V_2=H^1(N,\Z_2)$, let
$e=q^{-1}\sum_{g\in R}g$ be the actual deck-averaging
operator, and put $W_2=(1-e)V_2$.
On this rank-$2(q-1)$ integral summand the positive
square acts identically:
\[
                         \beta^*|_{W_2}=I.                 \tag{D.69y}
\]
The statement determines the indicated quotient of the
centralizer of the positive square. It does not determine
the image of $C_F(s)$ or the full structure of either
centralizer, and it does not exclude or realize a negative
quartic coordinate.
\end{claim}
\passage{Section 30, folios 194-196, scans 383-387; Section 31, folios 198-199 verso, scans 391-394; Part D(b), full loop-outer finite actions in type (1,1)}

\begin{proof}
\emph{The actual covers and the two cyclic extensions.}
By D.68 the homomorphism $\rho$ has
\[
 \rho(x)=t,\qquad \rho(y)=r,\qquad
 \rho(c)=r^{-2},\qquad
 \rho\alpha=\rho,\qquad
 Q=\langle r,t\mid r^q=t^2=1,\ trt=r^{-1}\rangle .
                                                               \tag{D.69d}
\]
The last equality $\rho\alpha=\rho$ is the proved identity
action at this multiplier coordinate. Thus $N$ is normal
in $E=F\rtimes\langle s\rangle$, and the actual quotient,
with inverse coordinate maps, is
\[
 \begin{aligned}
 E/N&\longrightarrow Q\times C_4,
        &fs^jN&\longmapsto(\rho(f),s^j),\\
 Q\times C_4&\longrightarrow E/N,
        &(z,s^j)&\longmapsto fs^jN\quad(\rho(f)=z).
 \end{aligned}                                               \tag{D.69e}
\]
The second map is independent of the chosen lift $f$.
The multiplication rule follows from $\rho\alpha=\rho$,
so these maps are inverse homomorphisms.
Let
\[
 E_N=N\rtimes\langle s\rangle,\qquad
 E_{2,N}=N\rtimes\langle\sigma\rangle,\qquad
 F_0=\ker\chi=\rho^{-1}(R).
                                                               \tag{D.69f}
\]
These are the inverse images of the named factors under
(D.69e), with their actual generator labels.

The original peripheral monodromy $r^{-2}$ has order $q$.
Its right action on the $2q$ cosets has two cycles, both
of length $q$. Thus the finite covering surface has two
punctures. Its Euler characteristic is $-2q$, so its
genus is $q$. The inclusion of its discrete fundamental
group has finite index in the original free group and
therefore has its full induced profinite topology.
This identifies its full completion with $N$ and gives
$\operatorname{rank}(N)=2q+1$.
The intermediate group $F_0$ is the specified type-$(1,2)$
double cover with Schreier basis $u=x^2,B=y,C=xyx^{-1}$
from D.68. Its quotient $F_0/N=R$ sends these generators
to $1,r,r^{-1}$, respectively.

\emph{The exact integral deck module.}
Put $\Lambda=\Z_2[R]$ and $\nu=1+r+\cdots+r^{q-1}$.
The actual degree-$q$ graph cover of the three-petal rose
for $F_0$ has chain boundary
\[
 \partial:\Lambda^3\longrightarrow\Lambda,\qquad
       (a,b,d)\longmapsto(r-1)b+(r^{-1}-1)d .
                                                               \tag{D.69g}
\]
The two explicit inverse maps on its cycle module are
\[
 \begin{aligned}
 \Lambda^2\oplus\Z_2&\longrightarrow\ker\partial,
       &(a,d,k)&\longmapsto(a,r^{-1}d+k\nu,d),\\
 \ker\partial&\longrightarrow\Lambda^2\oplus\Z_2,
       &(a,b,d)&\longmapsto(a,d,k),
       \quad b-r^{-1}d=k\nu .
 \end{aligned}                                               \tag{D.69h}
\]
The scalar $k$ exists uniquely because the annihilator of
$r-1$ in $\Lambda$ consists exactly of constant-coefficient
multiples of $\nu$. This follows directly by comparing
the $q$ cyclic coefficients. Hence
\[
       H_1(N,\Z_2)\cong\Z_2\oplus\Lambda^2
                       \quad\text{as actual deck modules}. \tag{D.69i}
\]
Dualizing gives the same multiplicities in
$V_2=H^1(N,\Z_2)=\operatorname{Hom}_{\rm cont}(N,\Z_2)$.
Let $e=q^{-1}\sum_{g\in R}g$ and
\[
                V=H^1(N,\F_2),\qquad
                W_2=(1-e)V_2,\qquad W=W_2/2W_2 .
                                                               \tag{D.69j}
\]
Thus $eV_2$ has rank three and $W_2$ has rank $2(q-1)$.
Restriction and transfer for the index-$q$ inclusion
$N\subset F_0$, with $q$ invertible in $\Z_2$, identify
$eV_2$ with the actual $H^1(F_0,\Z_2)$. These maps are
equivariant for the original $\alpha$.

Let $S$ act on $V_2$ by pullback under $\alpha^{-1}$, and
write $B=S^2$ for the pullback under $\beta$.
Because $\alpha$ induces identity on $Q$, $S$ commutes
with every deck operator from $R$: the two actual
conjugations differ only by conjugation by an element
of $N$, which acts trivially on this homology.
Choose the finite unramified extension $\mathcal O$ of
$\Z_2$ containing the $q$th roots of unity. Its explicit
character projectors are
$q^{-1}\sum_{g\in R}\psi(g)^{-1}g$.
Every nontrivial character summand of
$W_2\otimes\mathcal O$ is free of rank two by (D.69i).
On its residue field $S^4=1$ makes $S$ unipotent.
A two-dimensional unipotent matrix in characteristic
two has square identity. Consequently
\[
                          B|_W=I.                          \tag{D.69k}
\]
On the deck-trivial rank-three summand, let
$S_{\rm hom}$ be the matrix induced by $\alpha$
on the retained homology basis of $F_0$.
D.68 gives $S_{\rm hom}^2=-J$.
In the dual basis the present pullback convention is
$S=(S_{\rm hom}^{-1})^{\mathsf T}$; therefore
$B=S^2=-J^{\mathsf T}$, since $J^2=I$.
The reduction of $J^{\mathsf T}-I$ has rank one.
Hence $S$ has one Jordan block of size three there,
and $B$ has blocks of sizes two and one.
On all of $V$, therefore,
\[
                   \operatorname{rank}(B-I)=1.             \tag{D.69l}
\]

\emph{Which full section classes occur in these covers.}
D.46 gives three $F$-conjugacy classes of order-two
sections over $\beta$, with zero mod-two first
cohomology for each of their full fixed groups.
The proof of D.53 gives their three distinct affine
labels in $H_1(F,\F_2)$. Use $\sigma=s^2$ as label zero.
Conjugation by $s$ now has zero translation part,
since it fixes $\sigma$ itself. Its linear part is
$A_2(a)$ modulo two, a nonidentity transvection.
Its unique fixed line is $\ker\chi$.
D.53 proves that its action fixes exactly the class
of $\sigma$ and interchanges the other two classes.
Their two labels consequently have $\chi$-value one.
An element $f\sigma$ in either class has
$\chi(f)=1$, so its image under (D.69e) has nontrivial
$Q$ coordinate. Neither class meets $E_{2,N}$.
Every finite subgroup of $E_N$ injects into its
quotient $C_4$, since $N$ is torsion-free.
Its order-two subgroups lie in $E_{2,N}$.
Thus the only global class meeting either group is
the class of $\langle\sigma\rangle$.

Put $C=C_F(\sigma)$. Its mod-two first cohomology
vanishes by D.46, so $\chi|_C=0$.
The image $\rho(C)$ is contained in $R$ and is
therefore either $1$ or $R$.
The exact class parametrization is
\[
 Q/\rho(C)\longrightarrow
 \{\text{$N$-conjugacy classes of sections in $E_{2,N}$}\},
 \qquad
 \rho(f)\rho(C)\longmapsto
          [\,f\langle\sigma\rangle f^{-1}\,].
                                                               \tag{D.69m}
\]
Indeed two such subgroups are $N$-conjugate precisely
when their conjugators differ by an element of $N$
and an element of $N_F(\langle\sigma\rangle)=C$.
Every conjugate belongs to $E_{2,N}$ because its
image in $Q\times C_4$ is the same central subgroup.
These are also its $E_{2,N}$-conjugacy classes:
multiply any conjugator by the matching element of
the section to leave a conjugator in $N$.
Conjugation by $s$ on (D.69m) sends $\rho(f)$ to
$\rho(\alpha(f))=\rho(f)$ and fixes $\sigma$.
Hence (D.69m) also parametrizes the $E_N$ classes.

\emph{The trivial-image alternative and its Brown algebra.}
Suppose $\rho(C)=1$.
Then $C\subset N$ and (D.69m) has exactly $2q$ classes.
The normalizer in $E_{2,N}$ of the representative
$f\langle\sigma\rangle f^{-1}$ is exactly
\[
                       f(C\times\langle\sigma\rangle)f^{-1}.
                                                               \tag{D.69n}
\]
The product map sends $(h,\sigma^\epsilon)$ to
$fh\sigma^\epsilon f^{-1}$; its inverse reads the
unique $F$ coordinate and the specified quotient
coordinate after conjugation by $f^{-1}$.
As a closed subgroup of the full free profinite
group $F$, $C$ has $\operatorname{cd}_2(C)\leq1$.
Since $H^1(C,\F_2)=0$, all its positive-degree
mod-two cohomology vanishes.

The full Brown restriction theorem
\cite[Theorem~2 and \S1.1]{A-KochloukovaZalesskii2007}
applies to $E_{2,N}$: its open subgroup $N$ has
$2$-cohomological dimension one, every finite subgroup
injects into $C_2$, and the only nontrivial elementary
abelian two-subgroups are the order-two sections,
with no inclusions between distinct ones.
The two-row Hochschild--Serre spectral sequence
has finite modules in every degree, since $N$
is finitely generated and $C_2$ is finite.
Thus $H^n(E_{2,N},\F_2)$ is finite in every degree;
the theorem's dense restriction embedding is
therefore onto its Hausdorff target.
The actual restriction map is a ring isomorphism
in sufficiently high degrees onto the product of
the cohomologies of (D.69n).
There are exactly $2q$ factors, each one-dimensional
in each positive degree.
For the quotient character $u:E_{2,N}\to\F_2$,
localizing that restriction and taking degree zero gives
\[
 \mathcal R=
 \bigl(H^*(E_{2,N},\F_2)[u^{-1}]\bigr)^0
                         \cong\F_2^{\,2q}.                \tag{D.69o}
\]
Here $u$ is zero on each fixed group and nonzero
on its order-two section. This is the actual
localization of the displayed restriction maps.

We need the following precise injection into this algebra.
Every $\theta\in W$ is invariant under $\beta$ by
(D.69k), and extends to a character of $E_{2,N}$
by declaring $\theta(\sigma)=0$.
The Hochschild--Serre spectral sequence has only
rows zero and one. The actual section
$\langle\sigma\rangle\to E_{2,N}$ makes inflation
from $H^*(C_2,\F_2)$ injective in every degree.
Thus every bottom-row class survives, and all
possible differentials, whose targets are in that
row, are zero. On $W$ the row-one module is
$W[u]$, since $\beta$ acts trivially modulo two.
The associated graded class of $u^m\theta$ is the
nonzero element $\theta u^m$ whenever $\theta\ne0$.
It follows that
\[
                W\longrightarrow\mathcal R,\qquad
                         \theta\longmapsto\theta/u
                    \quad\text{is injective}.             \tag{D.69p}
\]

\emph{The exact mod-four character cochain.}
Since $B=I$ on $W_2/2W_2$, the integral operator
$(I-B)/2$ is defined on $W_2$. Let $D$ be its
reduction modulo two.
Lift $\theta\in W$ to a continuous character
$\varphi:N\to\Z/4$ from the integral summand $W_2$.
The cochain
$\widetilde\theta(n\sigma^\epsilon)=\varphi(n)$
has exact coboundary
\[
 \begin{aligned}
 \delta\widetilde\theta
       (n\sigma^\epsilon,m\sigma^\eta)
 &=\varphi(n)+\varphi(m)
                   -\varphi(n\beta^\epsilon(m))\\
 &=\epsilon\bigl(\varphi(m)-\varphi(\beta(m))\bigr)
       =2\epsilon(D\theta)(m)\quad\text{in }\Z/4.
 \end{aligned}                                               \tag{D.69q}
\]
Changing the lift by twice a character does not
change $D\theta$, because $\beta$ is identity
modulo two on $W$.
Thus its Bockstein, which equals its cup square,
is the actual identity
\[
                           \theta^2=uD\theta .
                                                               \tag{D.69r}
\]
Every element of the product algebra (D.69o) is
idempotent. Applying this to $\theta/u$ and using
(D.69r) and the injection (D.69p) gives $D\theta=\theta$.
Consequently
\[
                         B|_{W_2}\equiv -I\pmod4.          \tag{D.69s}
\]

\emph{The incompatible quartic section count.}
Extend (D.69s) to the unramified ring $\mathcal O$
used above. On a nontrivial character summand,
if the two-dimensional reduction of $S$ were
identity, write $S=I+2L$ over $\mathcal O$.
Then $S^2\equiv I\pmod4$, contrary to
$S^2=B\equiv-I\pmod4$.
Hence every nontrivial character summand of
$V\otimes\overline{\F}_2$ has one Jordan block
of size two for $S$. The deck-trivial summand
has one block of size three, as already proved.
There are $q-1$ nontrivial characters.
In the periodic $C_4$ resolution the two maps
are $S-I$ and
$1+S+S^2+S^3=(S-I)^3$.
Each of these blocks contributes exactly one
dimension in every positive cohomological degree.
Descent of dimension from the splitting field gives
\[
                         \dim H^j(C_4,V)=q\quad(j>0).
                                                               \tag{D.69t}
\]
The two-row spectral sequence for $E_N$ therefore
gives the upper bound
\[
                         \dim H^n(E_N,\F_2)\leq q+1
                                      \quad(n\geq2).       \tag{D.69u}
\]

On the other hand, the exact full normalizer in
$E$ of $\langle\sigma\rangle$ is
$C\rtimes\langle s\rangle$, as in D.53.
For each representative in (D.69m), its conjugate
lies entirely in $E_N$ under the present
assumption $\rho(C)=1$: (D.69e) sends its fixed
group to identity and its conjugated section
generator to the named $C_4$ factor.
Thus its $E_N$ normalizer is exactly
\[
                           f(C\rtimes\langle s\rangle)f^{-1}.
                                                               \tag{D.69v}
\]
The map $(h,s^j)\mapsto fhs^jf^{-1}$ and its
coordinate inverse give this semidirect-product
identification with labels.
All positive-degree cohomology of $C$ is zero,
so inflation identifies the cohomology of each
normalizer with $H^*(C_4,\F_2)$.
Brown applies here as well: $N$ is the open
dimension-one subgroup, finite subgroups inject
into $C_4$, and there is no elementary abelian
subgroup of rank two. Its elementary two-subgroup
classes are exactly the $2q$ classes in (D.69m);
conjugation by $s$ does not fuse them.
The same two-row finite-module spectral sequence
makes $H^n(E_N,\F_2)$ finite, so this dense
restriction embedding is again onto.
Restriction therefore gives
$\dim H^n(E_N,\F_2)=2q$ in high degrees.
This contradicts (D.69u), since $q>1$.
The alternative $\rho(C)=1$ is impossible.
We have proved the surjectivity in (D.69b).

\emph{The exact remaining fixed-group cohomology.}
Now $\rho(C)=R$. Formula (D.69m) has exactly two
cosets, represented by $1$ and $t=\rho(x)$.
Intersecting the full normalizers with $E_{2,N}$
gives precisely (D.69c), since its fixed kernel is
$C_0=C\cap N$. Both are conjugate by the original
element $x$ and have the same first-cohomology dimension.

The split spectral sequence for $E_{2,N}$ still
collapses by the actual section argument above.
By (D.69l) the operator $B-I$ on its
$(2q+1)$-dimensional module $V$ has rank one.
For $C_2$ both periodic differentials are $B-I$,
so
\[
 \dim H^j(C_2,V)=(2q+1)-2=2q-1\quad(j>0),\qquad
 \dim H^n(E_{2,N},\F_2)=2q\quad(n\geq2).
                                                               \tag{D.69w}
\]
Brown restriction to the two normalizers (D.69c)
then gives the exact equality
\[
                    2q=2\bigl(1+\dim H^1(C_0,\F_2)\bigr).
                                                               \tag{D.69x}
\]
In particular this first cohomology is finite-dimensional,
and its dimension is $q-1$, as claimed.
No assertion that $C_0$ is itself a full free profinite
group, or that $C_F(s)$ has the same finite image,
is used.

\emph{The full integral action on the nontrivial deck summand.}
We now use the proved image $\rho(C)=R$, not the
trivial-image alternative discarded above.
Choose an actual $h\in C$ with $\rho(h)=r$.
Conjugation by $h$ preserves $E_{2,N}$ and fixes
$\sigma$ itself. Its $q$th power is conjugation
by $h^q\in C_0\subset E_{2,N}$, so the resulting
cohomology action factors through $R$.
On $V_2$ it is exactly the deck action: replacing
$h$ by any other lift of $r$ changes conjugation
by an inner automorphism of $N$.
The extension of $\theta\in W$ which is zero on
$\sigma$ is consequently $R$-equivariant, and
the quotient character $u$ is fixed.

By the same finite-dimensional Brown restriction
argument, the actual localized degree-zero ring is now
\[
 \mathcal R\cong
 \bigl(\F_2\oplus H^1(C_0,\F_2)u^{-1}\bigr)
 \times
 \bigl(\F_2\oplus H^1(xC_0x^{-1},\F_2)u^{-1}\bigr).
                                                               \tag{D.69z}
\]
In each factor the product of two degree-one fixed-group
classes vanishes, because its full fixed group has
$2$-cohomological dimension at most one.
Thus the two displayed augmentation ideals are square
zero; their product $\mathcal I$ is the nilradical of
$\mathcal R$, and $\mathcal R/\mathcal I=\F_2^2$.
The action of $R$ fixes both section classes, which
are the two cosets in $Q/R$. Its action on this
residue algebra is therefore identity.

The injection (D.69p) still holds, since its split
spectral-sequence proof used no assumption on
$\rho(C)$. It is now the $R$-equivariant map
$W\to\mathcal R$, $\theta\mapsto\theta/u$.
The composite with $\mathcal R/\mathcal I$ vanishes:
deck averaging is zero on $W$ and identity on that
residue algebra, and $q$ is invertible in $\F_2$.
Hence its image is contained in $\mathcal I$.
Both vector spaces have dimension $2(q-1)$ by
(D.69j) and (D.69x), so it identifies $W$ with
that square-zero ideal.
In particular $(\theta/u)^2=0$ for every $\theta\in W$.
The exact cochain identity (D.69r) and the same
injection give $D\theta=0$ for every $\theta$.
Thus $B|_{W_2}\equiv I\pmod4$.

Finally $B^2=I$ on the free $\Z_2$-module $W_2$.
Writing $B=I+4L$ gives
$L(I+2L)=0$ after cancelling the nonzero scalar
$8$ in its torsion-free endomorphism module.
The matrix $I+2L$ is invertible over $\Z_2$,
because its reduction is identity.
Therefore $L=0$ and $B=I$ on $W_2$.
This is precisely (D.69y), since $\beta^{-1}=\beta$.
Together with the preceding dual-basis calculation,
it determines the integral square as
$-J^{\mathsf T}$ on $eV_2$ and identity on $W_2$.

The executable certificate
\path{code/part_c/torus_quartic_bad_prime_d69.py}
checks the exact finite covering cycles and module maps,
the two section-label possibilities, the cyclic
resolution dimensions, and the mod-four cochain identity.
The full normalizer and Brown arguments above supply
the stated centralizer conclusion.
\end{proof}

\begin{claim}{D.70}{Proved}{The quartic torus multiplier is globally one}
Every element of order four in the full loop-outer group
of the once-punctured torus has multiplier
$1\in\widehat\Z^\times$.
Consequently every finite two-subgroup of that full group
has multiplier image in the two global units $\{1,-1\}$,
and the multiplier is determined by its $2$-adic coordinate.
The conclusion concerns the full group, without a
geometric-membership assumption.
It determines the quartic multiplier; the conjugacy of
an arbitrary quartic action to a specified geometric
action is not asserted.
\end{claim}
\passage{Section 30, folios 194-196, scans 383-387; Section 31, folios 198-199 verso, scans 391-394; Part D(b), full loop-outer finite actions in type (1,1)}

\begin{proof}
\emph{An exact degree-three restriction obstruction.}
We first give the cohomological calculation used below.
Let $P$ be the pro-two completion of an orientable closed
surface group of genus $g\geq2$, and let $b$ be an actual
involution of $P$ whose action on $H^1(P,\F_2)$ is
identity and whose action on $H^1(P,\Z_2)$ is nonidentity.
Write
\[
 G=P\rtimes_b\langle\tau\rangle,\qquad
 \tau^2=1,\qquad
 u:G\longrightarrow\F_2,\quad u(h\tau^\epsilon)=\epsilon .
                                                               \tag{D.70a}
\]
Then any two $b$-fixed continuous characters
$\theta,\psi:P\to\Z_2$ have zero reduced cup product
in $H^2(P,\F_2)$.

Here are all maps and hypotheses for this assertion.
The group $P$ is a torsion-free Demushkin pro-two group:
its cohomological dimension is two, its first cohomology
has dimension $2g$, and its second cohomology is
$\F_2$ with the nondegenerate surface cup pairing.
The integral character module is $\Z_2^{2g}$.
These are the surface duality comparisons of A.2 and A.7,
used with both completions in D.42.

Every $\xi\in H^1(P,\F_2)$ has the actual character
extension
\[
 \widetilde\xi:G\longrightarrow\F_2,\qquad
                  \widetilde\xi(h\tau^\epsilon)=\xi(h),
                                                               \tag{D.70b}
\]
because $b$ fixes $\xi$.
The action on $H^2(P,\F_2)$ is also identity, since this
space is one-dimensional over $\F_2$.
The multiplicative Hochschild--Serre spectral sequence
therefore has its actual second page
\[
 E_2^{i,j}
   =H^i(C_2,\F_2)\otimes H^j(P,\F_2)
   =\F_2[u]\otimes H^j(P,\F_2),\qquad 0\leq j\leq2 .
                                                               \tag{D.70c}
\]
The quotient class $u$ is permanent: the section
$C_2\to G$, $\tau^\epsilon\mapsto\tau^\epsilon$,
makes inflation injective.
All classes in $E_2^{0,1}$ are permanent by (D.70b).
The surface class in $E_2^{0,2}$ is the cup product of
two of them, so it is permanent as well.
These classes generate the entire second-page algebra.
Thus every differential vanishes, and multiplication by
$u$ on its row-two term gives the nonzero class
$u\eta\in E_\infty^{1,2}$ for each nonzero surface
class $\eta$.

The fixed integral characters have the stronger actual
extensions
\[
 \widetilde\theta:G\longrightarrow\Z_2,\qquad
 \widetilde\psi:G\longrightarrow\Z_2,\qquad
 \widetilde\theta(h\tau^\epsilon)=\theta(h),\quad
 \widetilde\psi(h\tau^\epsilon)=\psi(h).
                                                               \tag{D.70d}
\]
Every finite subgroup of $G$ injects into $G/P=C_2$,
because $P$ is torsion-free. For each order-two
subgroup $T=\langle z\rangle$, the two characters
in (D.70d) vanish on $T$: twice their values on
$z$ is zero in the torsion-free group $\Z_2$.
There is the exact direct-product normalizer
\[
 \begin{aligned}
 C_P(T)\times T&\longrightarrow N_G(T),
                    &(h,z^\epsilon)&\longmapsto hz^\epsilon,\\
 N_G(T)&\longrightarrow C_P(T)\times T,
                    &w&\longmapsto(wz^{-\epsilon},z^\epsilon),
                    \qquad \epsilon=u(w).
 \end{aligned}                                                \tag{D.70e}
\]
Indeed every normalizer element fixes the unique
nonidentity element of $T$, and its quotient label
determines $\epsilon$ uniquely.
These are inverse continuous homomorphisms.

Write $z=n\tau$. Its actual action on $P$ is
$\operatorname{Inn}(n)b$, with square identity.
It has the same integral homology action as $b$,
so it is nonidentity by the stated hypothesis.
The hypotheses of
\cite[Corollary~2 and its proof in \S2]{A-KochloukovaZalesskii2007}
therefore hold for this particular action:
$P$ is a pro-two Poincare duality group of dimension
two and type $\mathrm{FP}_\infty$, and the acting
subgroup has order two in $\operatorname{Aut}(P)$.
The corollary proves that the full fixed subgroup
$C_P(T)$ is a free pro-two group.
In particular
\[
                          H^2(C_P(T),\F_2)=0.                \tag{D.70f}
\]
For clarity, the infinite-index step in that proof
uses its Lemma~1 for an actual nontrivial finite
two-group of automorphisms of a torsion-free pro-two
group; it then applies the dimension bound for an
infinite-index subgroup of a Demushkin group.
It is applicable separately to every section $T$.

Let bars denote the reductions of (D.70d) modulo two.
Under (D.70e), both restricted characters are pulled
back from $C_P(T)$, since they vanish on $T$.
Their cup product is zero by (D.70f).
It follows that
\[
 \operatorname{res}_{N_G(T)}
           \bigl(u\smile\overline{\widetilde\theta}
                    \smile\overline{\widetilde\psi}\bigr)=0
                 \quad\text{for every }T.                  \tag{D.70g}
\]
The full Brown restriction theorem
\cite[Theorem~2]{A-KochloukovaZalesskii2007}
is injective in degree three here: $P$ is open
of cohomological dimension two, and $G$ has no
$C_2\times C_2$ because its finite subgroups
inject into $C_2$.
Thus the class in (D.70g) is zero in $H^3(G,\F_2)$.
If $\bar\theta\smile\bar\psi=\eta\ne0$ on $P$,
its associated graded class under (D.70c) is
the nonzero element $u\eta\in E_\infty^{1,2}$,
a contradiction. This proves the asserted
zero-cup property of the integral fixed characters.
Only the injectivity part of Brown restriction
is needed in this calculation.

\emph{The actual full cusp filling and its pro-two quotient.}
Now take an order-four full torus action and the
actual lift $\alpha$ of D.46.
D.43 and D.50 give
$\lambda^2=1$ and $\lambda_2=1$.
Suppose an odd coordinate satisfies $\lambda_q=-1$.
D.68 gives its actual dihedral quotient, and D.69
gives the actual type-$(q,2)$ cover
\[
 N=\ker\bigl(\rho:F\twoheadrightarrow D_{2q}\bigr),
 \qquad \beta=\alpha^2,\qquad
 R=\langle\rho(y)\rangle\cong C_q .
                                                               \tag{D.70h}
\]
In particular $q\equiv1\pmod4$, so $q\geq5$.
Let $\Gamma_N$ be the corresponding discrete
finite-index subgroup of the original punctured-torus
group, with its original covering surface.
Use the two actual based peripheral generators
$d_0=c^q$ and $d_1=xc^qx^{-1}$ of this cover.
The second base path is the lift of the original
word $x$ to the other peripheral orbit.
Their full procyclic subgroups in $N$ are the
two covering cusps.
Let
\[
 K=\overline{\langle\!\langle d_0,d_1\rangle\!\rangle}_N,
 \qquad
 \Pi=N/K,\qquad P=\Pi^{(2)},\qquad
 p:N\twoheadrightarrow\Pi\twoheadrightarrow P .
                                                               \tag{D.70i}
\]
These are the specified quotient maps.
Filling the two punctures of the discrete cover
gives the closed genus-$q$ surface group
$\Gamma_N/\langle\!\langle d_0,d_1\rangle\!\rangle$.
Its full completion is exactly $\Pi$:
a map from either group to a finite group is
precisely a map from $\Gamma_N$ killing $d_0,d_1$.
The finite-index subgroup $\Gamma_N$ has its full
induced profinite topology in $F$, as in D.69,
so this identifies that universal finite quotient
with the displayed $N/K$.
Restricting the same finite targets to two-groups
gives the displayed maximal pro-two quotient $P$.
Consequently $P$ is the genus-$q$ pro-two closed
surface group used in the preceding calculation.

The preservation of $K$ has explicit conjugators
in $N$. Retain the cusp root $v=f_0s$ of D.68,
so $\alpha(c)=f_0^{-1}c^\lambda f_0$.
For $g\in\{1,x\}$ choose $g'\in\{1,x\}$
with $\rho(g')R=\rho(\alpha(g)f_0^{-1})R$.
Since $\rho(c)=r^{-2}$ generates $R$, there is
an integer $k$ for which
$\rho(\alpha(g)f_0^{-1})=\rho(g'c^k)$.
The actual element
$n=\alpha(g)f_0^{-1}c^{-k}(g')^{-1}$ then belongs to $N$,
and direct multiplication gives
\[
 \alpha(gc^qg^{-1})
       =n\bigl(g'c^q(g')^{-1}\bigr)^\lambda n^{-1}.
\]
The inverse automorphism gives the reverse inclusion.
Thus $\alpha(K)=K$ with the original two labels
and their full unit powers retained.
The maximal pro-two kernel of $\Pi$ is
characteristic, so $\alpha$, and hence $\beta$,
induce actual automorphisms through both arrows
in (D.70i). Denote the latter induced automorphism
of $P$ by $b$; it satisfies $b^2=1$.
No section is changed in this passage.

\emph{The dual inclusion and the closed integral deck splitting.}
The actual pullback of characters gives an injection
\[
 p^*:L=H^1(P,\Z_2)\hookrightarrow
                     V_2=H^1(N,\Z_2).
                                                               \tag{D.70j}
\]
Its image is exactly the characters which vanish
on $d_0,d_1$.
One direction follows from (D.70i).
Conversely such a character kills $K$, factors
through $\Pi$, and then through $P$ because its
target $\Z_2$ is pro-two.
This proves both directions of the assertion
for the exact displayed maps.
The same argument holds with target $\F_2$,
and reduction identifies $L/2L$ with $H^1(P,\F_2)$.

The rotation deck action on the punctured cover
extends over each of its two filled punctures.
It preserves orientation and gives the actual
action of $R$ on $L$ and its surface cup form.
As in D.69 put
\[
 e=\frac1q\sum_{r'\in R}r',\qquad W_2=(1-e)V_2 .
                                                               \tag{D.70k}
\]
Both cusp restriction functionals on $V_2$
are $R$-invariant. Indeed $R$ fixes each of the
two cusp labels, and its oriented deck action
changes a based peripheral generator only by
conjugation in $N$.
Consequently each such functional vanishes on
$W_2$, so (D.70j) identifies $W_2$ with
$(1-e)L$.

On the deck-trivial part, restriction and transfer
identify $eV_2$ with $H^1(F_0,\Z_2)$, where
$F_0=\rho^{-1}(R)$ is the type-$(1,2)$ double cover
of D.68. Its retained basis is
$u=x^2,B=y,C=xyx^{-1}$.
The original boundary is $c=CB^{-1}$,
and its two covering boundary homology classes
are opposite.
Restriction to $N$ sends each peripheral evaluation
to $q$ times the corresponding evaluation on $F_0$.
Since $q$ is a unit in $\Z_2$, the condition in
(D.70j) on this part is exactly
\[
                              \xi(B)=\xi(C).
                                                               \tag{D.70l}
\]
Thus $eL$ has rank two, and the actual direct
sum of integral character modules is
\[
                  L=eL\oplus W_2,\qquad
             \operatorname{rank}(W_2)=2q-2 .
                                                               \tag{D.70m}
\]
D.68 and the exact dual convention in D.69 give
$\beta^*=-J^{\mathsf T}$ on $eV_2$, with $J$
fixing $u$ and interchanging $B,C$.
On (D.70l), $J^{\mathsf T}$ is identity.
D.69 gives $\beta^*=I$ on $W_2$.
Equivariance of (D.70j) therefore gives the
full integral action
\[
                  b^*|_{eL}=-I,\qquad b^*|_{W_2}=I.
                                                               \tag{D.70n}
\]
It follows that $b$ has order exactly two,
its integral homology action is nonidentity,
and its action on $L/2L$ is identity.
All three assertions follow from the actual
rank-two summand in (D.70n); none is inferred
from the order of the original outer class alone.

\emph{The unimodular fixed cup pair and the contradiction.}
Choose the oriented surface generator of
$H^2(P,\Z_2)\cong\Z_2$.
The cup pairing
\[
               \langle\, ,\,\rangle:L\times L\to\Z_2
                                                               \tag{D.70o}
\]
is alternating and unimodular.
Every deck operator $r'$ preserves this generator,
so
$\langle r'\xi,\eta\rangle
 =\langle\xi,(r')^{-1}\eta\rangle$.
Summing over $R$ proves that the actual integral
idempotent $e$ is self-adjoint:
\[
                         \langle e\xi,\eta\rangle
                                  =\langle\xi,e\eta\rangle .
                                                               \tag{D.70p}
\]
Hence (D.70m) is an orthogonal direct sum.
Each restricted form is unimodular: in bases
of the two integral direct summands, the whole
Gram matrix is block diagonal with unit
determinant, so each block determinant is a
unit in $\Z_2$.
In particular the rank-$2q-2$ form on $W_2$
is nonzero and unimodular.
Choose $\theta,\psi\in W_2$ with
$\langle\theta,\psi\rangle$ a unit.
Their reductions have nonzero cup product
in $H^2(P,\F_2)$, while (D.70n) makes both
characters fixed by the actual $b$.
The degree-three restriction obstruction
(D.70a)--(D.70g) applies to exactly this $P,b$,
and says that their reduced cup product is zero.
This contradiction excludes the chosen
coordinate $\lambda_q=-1$.

There is no negative odd coordinate.
Together with $\lambda^2=1$ and $\lambda_2=1$,
this proves $\lambda=1$ in the full product
$\widehat\Z^\times=\prod_p\Z_p^\times$.
The possible finite two-subgroups are exactly
$1,C_2,C_4,C_2\times C_2,D_8$ by D.50.
D.43 gives the global sign for their involutions;
the present result gives multiplier one for
their order-four elements. These statements
cover every element of each of the listed
groups and prove the final assertion.
\end{proof}

\emph{Scope.}
This result closes the quartic multiplier branch
left in D.50 and D.68. It uses the exact
rotation-cover conclusion of D.69 and the actual
closed filling (D.70i).
It supplies no conjugacy to a geometric quartic
representative and no recognition of either full
fixed subgroup in D.69.
The remaining scalar order-three branch is
unaffected by this degree-three calculation.

\begin{claim}{D.71}{Proved}{The Frobenius cover of a scalar cubic torus coordinate}
Let $a$ be an order-three element of the full
once-punctured-torus loop-outer group, with actual
order-three representative $\alpha$ and labelled
extension $E=F\rtimes_\alpha\langle s\rangle$.
At every prime $q\ne3$ with
$\lambda_q\ne1$, there is the following actual
finite quotient, after the explicit based peripheral
change of generators in the proof:
\[
 \begin{gathered}
 Q=\langle r,t\mid r^q=t^3=1,\ trt^{-1}=r^\ell\rangle,
 \qquad
 \ell=\lambda_q\bmod q\in\F_q^\times,\\
 \rho:F\twoheadrightarrow Q,\qquad
 \rho(x)=t,\quad\rho(y)=r,\quad
 \rho\alpha=\rho .
 \end{gathered}                                                \tag{D.71a}
\]
Here $\ell$ has order three and $|Q|=3q$.
Put $R=\langle r\rangle$, $N=\ker\rho$,
$C=C_F(s)$, and $C_0=C\cap N$.
The actual cover defining $N$ has type
$((3q-1)/2,3)$, and
\[
 1\longrightarrow C_0\longrightarrow C
      \xrightarrow{\ \rho|_C\ }R\longrightarrow1,
 \qquad
                       H^1(C_0,\F_3)=0 .                    \tag{D.71b}
\]
There are exactly three $N$-conjugacy classes of
order-three sections in $N\rtimes\langle s\rangle$.
With their original quotient-generator labels,
their representatives and exact normalizers are
\[
 \begin{gathered}
 T_i=x^i\langle s\rangle x^{-i},\\
 N_{N\rtimes\langle s\rangle}(T_i)
       =x^i(C_0\times\langle s\rangle)x^{-i}
                      \qquad(0\leq i<3).
 \end{gathered}                                                \tag{D.71c}
\]
On the integral character module
$V_3=H^1(N,\Z_3)$, retain the actual rotation
projector and its complement
\[
 e=\frac1q\sum_{g\in R}g,\qquad W_3=(1-e)V_3 .
                                                               \tag{D.71d}
\]
For the action $S$ given by pullback under
$\alpha^{-1}$, the module $W_3$ is free of rank
$q-1$ over $\Z_3[\langle S\rangle]\cong\Z_3[C_3]$.
The isomorphism is supplied by lifting cyclic
vectors for the actual action.
The averaged rank-four module $eV_3$ reduces to
two Jordan blocks of size two modulo three.
These are exact quotient, lattice, and fixed-group
constraints. No exclusion or realization of the
scalar cubic coordinate is asserted.
\end{claim}
\passage{Section 30, folios 194-196, scans 383-387; Section 31, folios 198-199 verso, scans 391-394; Part D(b), full loop-outer finite actions in type (1,1)}

\begin{proof}
\emph{The actual index-three cover and its based generators.}
Use D.48 and D.51 for the two section classes of
$E$, with $\langle s\rangle$ the chosen one.
There is an invariant character $\chi:F\to\F_3$,
extended by $\chi(s)=0$, which can be scaled
uniquely so that
\[
 \chi(v)=1,\qquad
 \chi(s_1)=-1,\qquad
 vcv^{-1}=c^\lambda,\qquad v^3=c^\kappa ,
 \quad c=[X,Y],
                                                               \tag{D.71e}
\]
where $v$ is a cusp-normalizer element labelled by
$s$, and $s_1$ is the labelled generator of the
other full section class.
Indeed D.51 gives opposite nonzero values for
the cusp line and the other section line.
Also $\kappa_3$ is a unit by D.48.

Here is an explicit peripheral change of basis
with $\chi(x)=1,\chi(y)=0$, without rescaling
the now specified character.
For any current basis $(X,Y)$ the transformation
\[
 Q_0(X,Y)=(XYX^{-1},X^{-1}),\qquad
 Q_0^{-1}(x,y)=(y^{-1},yxy^{-1})
                                                               \tag{D.71f}
\]
has the displayed inverse and preserves
$[X,Y]$ exactly. Its action on the two character
values is $(a,b)\mapsto(b,-a)$.
Among its four iterates, at least one has first
character value $1$, since $(a,b)\ne(0,0)$.
Choose the first such iterate $(x,Y')$,
and let $d\in\{0,1,-1\}$ represent $\chi(Y')$.
Set $y=Y'x^{-d}$.
This shear has inverse $Y'=yx^d$ and
$[x,Y'x^{-d}]=[x,Y']$.
These actual discrete free-group automorphisms,
with the stated inverses, therefore supply
the required basis in the full completion
and keep $c=[x,y]$ exactly.

Let $F_0=\ker\chi$.
Its actual type-$(1,3)$ covering has the
Schreier basis and peripheral representatives
\[
 \begin{gathered}
 u=x^3,\qquad B_i=x^iyx^{-i}\quad(0\leq i<3),\\
 c_i=x^icx^{-i}=B_{i+1}B_i^{-1},\qquad
 B_3=uB_0u^{-1}.
 \end{gathered}                                                \tag{D.71g}
\]
These identities are actual words, including
the final value $i=2$.
On $M=H_1(F_0,\widehat\Z)$, deck conjugation by
$x$ gives $J$, which fixes $u$ and cyclically
permutes the $B_i$.
The inclusion homology map is the actual map
\[
 I:M\longrightarrow H_1(F,\widehat\Z),\qquad
 I(u)=3[x],\quad I(B_i)=[y].
                                                               \tag{D.71h}
\]
Its kernel is the rank-two boundary augmentation,
generated by the three differences $B_{i+1}-B_i$
with their sole sum relation.
Filling the three punctures gives the actual
degree-three closed torus cover; (D.71h) factors
through its integral lattice inclusion.

Let $S_{\rm hom}$ be induced by $\alpha$ on $M$.
Since $\chi\alpha=\chi$, it commutes with $J$:
the conjugating elements $\alpha(x)$ and $x$
differ by an element of $F_0$, whose inner
action on $M$ is trivial.
Write $v=f_0s$. Equation (D.71e) gives
$\chi(f_0)=1$, so the homology action $T$ of
conjugation by $v$ is
\[
                              T=JS_{\rm hom}.                \tag{D.71i}
\]
It fixes each of the three peripheral classes
with multiplier $\lambda$.
To check its labels, conjugation by $v$ preserves
$\chi$ on $F$, so $vx^iv^{-1}$ differs from
$x^i$ by an element of $F_0$, while
$vcv^{-1}=c^\lambda$.
Thus
\[
       S_{\rm hom}|_{\ker I}=\lambda J^{-1}.
                                                               \tag{D.71j}
\]

There is also a fully specified matrix for this
comparison. If $A=(a_{ij})$ is the original
homology matrix in the basis $(x,y)$ and
$\boldsymbol1=(1,1,1)^{\mathsf T}$, then
\[
 S_{\rm hom}=
 \begin{pmatrix}
 a_{11}&(a_{12}/3)\boldsymbol1^{\mathsf T}\\
 a_{21}\boldsymbol1&
       \lambda J_3^{-1}
          +((a_{22}-\lambda)/3)\boldsymbol1\boldsymbol1^{\mathsf T}
 \end{pmatrix},
                                                               \tag{D.71k}
\]
where $J_3$ is the cycle on $(B_0,B_1,B_2)$.
The numerator divisibilities hold at $3$:
$\chi\alpha=\chi$ gives $a_{11}\equiv1$ and
$a_{12}\equiv0$ modulo three; the determinant
is $\lambda$, with $\lambda_3=1$, so
$a_{22}\equiv1\equiv\lambda$ there.
At every other prime division by three is
invertible. To prove the formula, the $u$-column
is determined by $J$-invariance and (D.71h);
the remaining three columns are determined by
(D.71j), commutation with $J$, and
$IS_{\rm hom}=AI$.
These determinations can first be read over each
$\Q_p$; the stated integral divisibilities
and torsion-free lattices give (D.71k) in
the full module.

\emph{The invariant character and the actual Frobenius quotient.}
Fix the specified bad prime $q$.
The unit $\lambda_q$ is a primitive cube root.
It reduces to one modulo $q$ only for the
identity cube root, because $q\ne3$ and the
three roots have distinct reductions.
Thus $\ell=\lambda_q\bmod q$ has order three
and $q\equiv1\pmod3$.
Define the continuous character
\[
 \ell_q:F_0\longrightarrow\Z_q,\qquad
           \ell_q(u)=0,\quad
           \ell_q(B_i)=\lambda_q^i .
                                                               \tag{D.71l}
\]
Its values sum to zero, and
$\ell_qJ=\lambda_q\ell_q$.
Over $\Z_q$, the row therefore annihilates the
deck-trivial summand and is supported on the
boundary augmentation. Equation (D.71j), or directly
(D.71k), gives the exact equation
\[
                             \ell_q\alpha=\ell_q.             \tag{D.71m}
\]
Put $N=\ker(\ell_q\bmod q)$ inside $F_0$.
It is normal in $F$: conjugation by $x$ scales
the character by the unit $\ell$, and
conjugation by $F_0$ does not change a character.
It is also $\alpha$-invariant in both directions
by (D.71m).
Its quotient in $F_0$ has order $q$, so
$[F:N]=3q$.

In $F/N$, the images of $x,y$ satisfy
$x^3=u=1$, $y^q=1$, and
$xyx^{-1}=y^\ell$.
The latter equality is the equality of their
images under $\ell_q\bmod q$ in $F_0/N$.
The presented group in (D.71a) has the explicit
multiplication
\[
 (a,i)(b,j)=(a+\ell^ib,\ i+j)
         \quad\text{on }\F_q\times\F_3,
 \qquad r=(1,0),\quad t=(0,1).
                                                               \tag{D.71n}
\]
It has order $3q$.
The inverse coordinate homomorphisms are
$r\mapsto yN,t\mapsto xN$ and
$xN\mapsto t,yN\mapsto r$.
The relations and the already computed index
make them inverse, proving exactly (D.71a)'s
group and quotient, rather than just a
surjection to a smaller group.

The induced automorphism on $Q$ fixes $r$
by (D.71m) and preserves the quotient label
in $Q/R$ by $\chi\alpha=\chi$.
It consequently sends $t$ to $r^jt$ for
some $j\in\F_q$.
Its third iterate sends $t$ to $r^{3j}t$.
Since the original $\alpha$ has order three
and $q\ne3$, $j=0$.
Thus $\rho\alpha=\rho$.
The actual extension quotient and inverse are
\[
 \begin{aligned}
 E/N&\longrightarrow Q\times C_3,
               &fs^iN&\longmapsto(\rho(f),s^i),\\
 Q\times C_3&\longrightarrow E/N,
               &(z,s^i)&\longmapsto fs^iN
                         \quad(\rho(f)=z).
 \end{aligned}                                                \tag{D.71o}
\]
Independence of the lift and multiplication
follow from $\rho\alpha=\rho$.

The cusp has image
$\rho(c)=[t,r]=r^{\ell-1}$ of order $q$.
Its action on the $3q$ cosets has three cycles
of length $q$.
The corresponding discrete cover therefore
has three punctures and Euler characteristic
$-3q$, hence genus $(3q-1)/2$.
Its finite-index subgroup carries its full
induced profinite topology, giving exactly
the stated full surface completion $N$.

\emph{The integral rotation module, including its inverse maps.}
Write $\Lambda_3=\Z_3[R]$ and
$\nu=\sum_{i=0}^{q-1}r^i$.
Let $a,b\in\{1,\ldots,q-1\}$ be the integer
representatives of $\ell,\ell^2$, and put
$P_k=1+r+\cdots+r^{k-1}$.
The actual degree-$q$ graph cover of the
four-generator rose for $F_0$ has boundary
\[
 \partial:\Lambda_3^4\longrightarrow\Lambda_3,\qquad
 (A,B,D,E)\longmapsto
       (r-1)B+(r^a-1)D+(r^b-1)E .
                                                               \tag{D.71p}
\]
Its exact cycle-coordinate maps are
\[
 \begin{aligned}
 \Lambda_3^3\oplus\Z_3&\longrightarrow\ker\partial,
 &(A,D,E,k)&\longmapsto
          (A,-P_aD-P_bE+k\nu,D,E),\\
 \ker\partial&\longrightarrow\Lambda_3^3\oplus\Z_3,
 &(A,B,D,E)&\longmapsto(A,D,E,k),\\[-2pt]
 &&&\hspace{-15mm}B+P_aD+P_bE=k\nu .
 \end{aligned}                                                \tag{D.71q}
\]
The scalar $k$ exists uniquely: the annihilator
of $r-1$ consists exactly of constant-coefficient
multiples of $\nu$, as comparison of the cyclic
coefficients proves.
Thus these are inverse integral module maps and
\[
                  H_1(N,\Z_3)\cong
                             \Z_3\oplus\Lambda_3^3 .
                                                               \tag{D.71r}
\]
Dualizing gives the same deck multiplicities
in $V_3$.
Restriction and transfer for the actual
index-$q$ subgroup $N\subset F_0$ identify
$eV_3$ with $H^1(F_0,\Z_3)$.
The maps are equivariant under the original
$\alpha$; the remaining module $W_3$
has rank $3(q-1)$.

On $H_1(F_0,\Q_3)$, the boundary part of
$S_{\rm hom}$ is $J^{-1}$ by (D.71j).
Its two eigenvalues are the primitive cube
roots. The filled rank-two quotient has
the original nonidentity cubic matrix
$A_3$ by the lattice map (D.71h).
D.48 proves its polynomial is $X^2+X+1$.
Consequently the same polynomial annihilates
the whole rank-four module:
first over $\Q_3$, by the two commuting
$J$-summands, and then over $\Z_3$ by
torsion-freeness.
For the dual convention
$S=(S_{\rm hom}^{-1})^{\mathsf T}$ this gives
\[
                        I+S+S^2=0\quad\text{on }eV_3.
                                                               \tag{D.71s}
\]
The ring $\Z_3[z]/(z^2+z+1)$ is a discrete
valuation ring: with $\pi=z-1$ its defining
polynomial is $\pi^2+3\pi+3$, and
$(3)=(\pi^2)$.
The rank-four $\Z_3$-lattice in (D.71s)
is a torsion-free module over this ring,
so is free of rank two.
Its reduction modulo three is therefore
exactly two Jordan blocks of size two.

\emph{The integral fixed-cup obstruction at three.}
We use the following calculation with the
actual pro-three closed-surface group.
Let $P$ be such a group of genus at least two
and let $b$ be an actual automorphism of
order three whose integral homology action
is nonidentity.
Any two $b$-fixed characters
$\theta,\psi:P\to\Z_3$ have cup pairing
divisible by three.

Indeed put $G=P\rtimes_b\langle\tau\rangle$.
Both characters extend by
$\widetilde\theta(h\tau^i)=\theta(h)$ and
$\widetilde\psi(h\tau^i)=\psi(h)$.
They kill every subgroup of order three,
since $\Z_3$ is torsion-free.
Such a subgroup $T=\langle n\tau\rangle$
has the exact normalizer
$C_P(T)\times T$, by splitting off its
unique element with a prescribed quotient
label, as in (D.70e).
Its actual action on $P$ is a nonidentity
inner twist of $b$, because inner twists
do not change the given nonidentity
integral homology matrix.
The Demushkin conclusion of
\cite[Corollary~2 and \S2]{A-KochloukovaZalesskii2007}
therefore makes $C_P(T)$ free pro-three.
The restricted product of the two character
classes is zero on every normalizer.

If their reduced cup were a nonzero surface
class $\eta$, the actual product of the
extensions would lift $\eta$ to $H^2(G,\F_3)$.
It is consequently a permanent nonzero
class in the highest row $E_2^{0,2}$
of the Hochschild--Serre spectral sequence.
The action on $H^2(P,\F_3)$ is identity,
since an order-three scalar in $\F_3^\times$
is identity.
Let $z\in H^2(C_3,\F_3)$ be its period
generator. Its inflation is permanent by
the actual section, so $z\eta$ is a
nonzero permanent element of $E_2^{2,2}$.
There is no incoming differential to
this highest row, and its outgoing
differentials vanish by multiplicativity.
The class
$z\overline{\widetilde\theta}
       \overline{\widetilde\psi}$ is therefore
nonzero in $H^4(G,\F_3)$.
Its restrictions to all normalizers are zero.
This contradicts Brown injectivity in
degree four
\cite[Theorem~2]{A-KochloukovaZalesskii2007}:
$P$ is open of cohomological dimension two,
and finite subgroups of $G$ inject into
$C_3$, so there is no $C_3\times C_3$.
This proves the divisibility assertion.
It requires no assumption that $b$ is
identity on mod-three homology.

\emph{The actual filled character module and its paired lattices.}
Fill the three actual cusps
$d_i=x^ic^qx^{-i}$ of $N$ and put
\[
 P=\left(N/
       \overline{\langle\!\langle d_0,d_1,d_2\rangle\!\rangle}
                   \right)^{(3)},\qquad
 p:N\twoheadrightarrow P .
                                                               \tag{D.71t}
\]
The same finite-target universal property
as in (D.70i) identifies the quotient
before maximal pro-three completion with
the full completion of the closed covering
surface. Thus $P$ is the genus-$(3q-1)/2$
closed-surface pro-three group.
The normal closure is $\alpha$-invariant
with the original labels: for
$g=x^i$, write
$\alpha(g)f_0^{-1}=n x^j c^k$
with $n\in N$, choosing $j\in\{0,1,2\}$
by its coset in $Q/R$ and $k$ by
the rotation coordinate.
This is possible because $\rho(c)$
generates $R$.
Then
$\alpha(gc^qg^{-1})=n d_j^\lambda n^{-1}$.
The inverse proves equality of the
closed normal subgroups.
Hence $\alpha$ induces an actual
automorphism of $P$.

Pullback gives the exact dual injection
$H^1(P,\Z_3)\hookrightarrow V_3$ whose
image is the characters vanishing on
all $d_i$: such a character factors
through the full quotient and then
through its maximal pro-three quotient,
and the converse follows from (D.71t).
The three cusp evaluations are
$R$-invariant, since the rotations
fix their labels and their positive
generators up to conjugation.
Thus $W_3$ lies in this closed character
module. Its complement there is the
averaged rank-two closed-torus module.
The original covering homology map identifies
the action on the filled rank-two homology
over $\Q_3$ with the nonidentity matrix $A_3$.
On its dual the pullback convention is
$(A_3^{-1})^{\mathsf T}$ after that rational
identification. It is nonidentity as well.
Thus the induced automorphism of $P$ has
order exactly three and is nonidentity
on integral homology.

The rotation action extends over the
three filled punctures and preserves
the integral surface orientation.
The odd projector $e$ is therefore
self-adjoint for the unimodular
alternating cup form, exactly as in
(D.70p), with $q$ a unit in $\Z_3$.
Consequently its complementary closed
summand $W_3$ has a unimodular
alternating form.
The automorphism $S$ preserves this
form: its action on the integral
surface class has order dividing three
in $\Z_3^\times$, which has no
nonidentity element of order three.

Choose a finite unramified extension
$\mathcal O/\Z_3$ containing the $q$th
roots of unity.
For a character $\psi:R\to\mathcal O^\times$,
the actual projector is
$e_\psi=q^{-1}\sum_{g\in R}\psi(g)^{-1}g$.
Each nontrivial character lattice
\[
                 V_\psi=e_\psi(W_3\otimes\mathcal O)
                                                               \tag{D.71u}
\]
has rank three by (D.71r).
The automorphism $S$ commutes with
every rotation operator because
$\rho\alpha=\rho$, so preserves
each such lattice.
The cup form pairs $V_\psi$
perfectly with $V_{\psi^{-1}}$
and pairs it trivially with every
other character lattice.
The inverse character is distinct,
since $q$ is odd.

Let $K_3$ be the fraction field of
$\mathcal O$.
The polynomial $X^2+X+1$ is irreducible
over $K_3$: a square root of $-3$
would have half-integral valuation
in this unramified extension.
Since $S^3=1$, its fixed subspace
on the rank-three $K_3$-space has
dimension one or three.
The fixed-cup divisibility just proved
holds after extension to $\mathcal O$.
Indeed $\mathcal O$ is finite free
over $\Z_3$, so kernels of $S-I$
commute with this flat scalar
extension, and the bilinear divisibility
passes to their scalar extensions.

Fixed dimension three is impossible.
It would make $S=I$ on $V_\psi$,
and preservation of the perfect paired
form would give $S=I$ on
$V_{\psi^{-1}}$.
The perfect integral pairing would
then give two fixed characters with
unit cup pairing, contradicting that
divisibility.
Thus both fixed lattices have rank one.
Choose primitive generators
$u_\psi$ and $u_{\psi^{-1}}$.
They are primitive in the full lattices:
each fixed kernel is saturated, since
the image of $S-I$ is torsion-free.
Semisimplicity over $K_3$ and preservation
of the pairing make the pairing of the
two fixed lines nonzero.

Unimodularity supplies
$w\in V_\psi$ with
$\langle w,u_{\psi^{-1}}\rangle=1$.
Writing the actual norm as
$\mathcal N=I+S+S^2$, we have
$\mathcal Nw=b\,u_\psi$ for some
$b\in\mathcal O$, and
\[
       b\,\langle u_\psi,u_{\psi^{-1}}\rangle
          =\langle\mathcal Nw,u_{\psi^{-1}}\rangle
          =3 .
                                                               \tag{D.71v}
\]
Fixed-cup divisibility says the
displayed pairing lies in $3\mathcal O$.
Since $\mathcal O$ is unramified,
(D.71v) forces that pairing to be
three times a unit and forces $b$
to be a unit.
The norm therefore maps $V_\psi$
onto its primitive fixed line.

Modulo three, the norm is
$(S-I)^2$, and its rank on this
three-dimensional block is one.
As $(S-I)^3=0$, the block is exactly
one Jordan block of size three.
This also supplies an actual integral
cyclic basis: lift a vector whose
three cyclic images form a basis
modulo three. The columns
$w,Sw,S^2w$ then have unit determinant,
so the map
$\mathcal O[C_3]\to V_\psi$ sending
the group generator to $S$ and $1$
to $w$ is an isomorphism.

For descent, no unspecified module
form is used.
After the splitting residue-field
extension, the norm on $W_3/3W_3$
has rank $q-1$ in dimension $3(q-1)$.
Ranks are unchanged by this field
extension. Since every Jordan block
has size at most three and each
size-three block contributes one
to the norm rank, all blocks over
$\F_3$ are of size three.
Choose a basis of the actual
$\F_3[C_3]$-module by cyclic
vectors and lift them to $W_3$.
The resulting explicit action map
\[
 \Z_3[C_3]^{q-1}\longrightarrow W_3,\qquad
             (f_i(S))_i\longmapsto\sum_i f_i(S)w_i
                                                               \tag{D.71w}
\]
reduces to an isomorphism modulo three.
It is a map between free $\Z_3$-modules
of the same finite rank, so its matrix
has unit determinant and is invertible.
Its inverse is $C_3$-equivariant.
This proves the integral freeness
in the statement.

\emph{The three full section classes and the fixed-kernel map.}
Put $E_N=N\rtimes\langle s\rangle$.
D.48 gives $H^1(C,\F_3)=0$,
so $\chi(C)=0$ and $\rho(C)\subset R$.
Its image is either $1$ or $R$.
The other global section class has
nonzero $\chi$-value by (D.71e),
so none of its conjugates meets $E_N$.
Every conjugate of $\langle s\rangle$
does meet $E_N$, because its image
in (D.71o) is the same central
order-three subgroup.
There is the exact class bijection
\[
 Q/\rho(C)\longrightarrow
 \{\text{$N$-conjugacy classes of sections in }E_N\},
 \qquad
 \rho(f)\rho(C)\longmapsto[f\langle s\rangle f^{-1}].
                                                               \tag{D.71x}
\]
Two conjugators give the same class
precisely when their ratio differs by
an element of $N$ and an element of
$N_F(\langle s\rangle)=C$.
Conjugacy under $E_N$ gives the same
classes, by removing the matching
section label from a conjugator.

Let $V=H^1(N,\F_3)$.
On $eV$ the action is two Jordan
blocks of size two, while (D.71w)
makes the complementary module free
over $\F_3[C_3]$.
The periodic cyclic resolution has
maps $S-I$ and
$I+S+S^2=(S-I)^2$.
Each size-two block contributes one
dimension in every positive degree;
the free blocks contribute zero.
Thus
\[
          \dim H^j(C_3,V)=2\quad(j>0),\qquad
          \dim H^n(E_N,\F_3)\leq3\quad(n\geq2).
                                                               \tag{D.71y}
\]
The last bound is the actual two-row
Hochschild--Serre sequence, since $N$
has cohomological dimension one.

Brown restriction applies to $E_N$:
$N$ is open of dimension one, and
every finite subgroup injects into
$C_3$, so there is no elementary
abelian subgroup of rank two.
The finite-dimensional source in
(D.71y) makes the dense restriction
embedding onto its Hausdorff target.
If $\rho(C)=1$, (D.71x) supplies
$3q$ normalizer factors, already
contradicting (D.71y).
Therefore $\rho(C)=R$, proving the
surjectivity in (D.71b).
The classes are exactly its three
cosets, represented by $1,t,t^2$.
Intersecting the full direct-product
normalizers with $E_N$ gives
exactly (D.71c).
The maps
$(h,s^j)\mapsto x^ihs^jx^{-i}$
and their inverses, which read
the fixed-group coordinate and the
specified quotient label after
conjugation by $x^{-i}$, prove
these identifications as full
topological groups.

Each fixed kernel $x^iC_0x^{-i}$
is a closed subgroup of the free
profinite group $N$ and has
$3$-cohomological dimension at most
one. Kunneth and Brown therefore give,
in high degrees,
\[
 3\bigl(1+\dim H^1(C_0,\F_3)\bigr)
                  =\dim H^n(E_N,\F_3)\leq3.
                                                               \tag{D.71z}
\]
All dimensions are nonnegative,
so $H^1(C_0,\F_3)=0$ and the
high-degree dimension is exactly
three. This completes the proof
of the fixed-kernel and section
assertions.

The executable certificate
\path{code/part_c/torus_cubic_frobenius_d71.py}
checks all eight based character
adaptations and their word inverses,
the integral covering matrices,
the Frobenius quotients and cusp
cycles for the first admissible
primes, the graph-cycle maps with
their inverses, and every rank-three
order-three matrix over $\F_3$.
The integral fixed-cup argument
and the full Brown normalizers
above supply the conclusions
beyond those finite identities.
\end{proof}

\begin{claim}{D.72}{Proved}{The full mod-two cochain comparison for a pants involution}
Let $\alpha$ be the actual normalized representative of any
nonidentity finite-order element of the full pure three-cusp
loop-outer group. Retain all the data of D.34:
\[
 F=\widehat F(x,y),\qquad
 \alpha(x)=x^{-1},\quad \alpha(y)=f^{-1}y^{-1}f,
 \quad f\in F',\qquad E=F\rtimes_\alpha\langle s\rangle.
                                                               \tag{D.72a}
\]
Write $e_2\in\widehat\Z$ for the element whose $2$-adic
coordinate is one and whose other primary coordinates are zero.
The three actual section generators and their label maps are
\[
 s_0=s,\qquad s_1=x^{e_2}s,\qquad s_2=y^{e_2}fs,
 \qquad \iota_i:C_2\longrightarrow E,\quad\tau\longmapsto s_i.
                                                               \tag{D.72b}
\]
Let $\mathcal C^\bullet(E)$ be the differential graded
$\F_2$-algebra of normalized continuous group cochains with
trivial coefficients, with its ordinary cup multiplication.
Put
\[
 \mathcal A=\F_2[t_0,t_1,t_2]/(t_it_j:i\ne j),
              \qquad |t_i|=1,\qquad d_{\mathcal A}=0.
                                                               \tag{D.72c}
\]
The following is an actual surjective quasi-isomorphism of
unital differential graded algebras:
\[
 \begin{split}
 q:\mathcal C^\bullet(E)&\longrightarrow\mathcal A,\\
 q_0(a)&=a,\\
 q_n(c)&=\sum_{i=0}^2 c(s_i,\ldots,s_i)t_i^n
                                      \quad(n>0).
 \end{split}                                                   \tag{D.72d}
\]
Thus these are the complete mod-two cohomology ring and a
specified formality comparison, without an assumption that
the full centralizers of the sections are trivial.

Every sequence of at least three positive-degree cohomology classes whose
neighboring cup products are zero has a defined higher
Massey product containing zero. There is also the following
exact finite-matrix image calculation. For $n\geq1$, let
$U_{n+1}(\F_2)$ be the group of upper unitriangular matrices,
and define
\[
 \begin{gathered}
 \mathcal X_n=\bigl\{(\chi_1,\ldots,\chi_n)\in H^1(E,\F_2)^n:\
                    \chi_j\chi_{j+1}=0\ (1\leq j<n)\bigr\}.
 \end{gathered}
\]
The map reading the first superdiagonal gives a surjection
\[
 \operatorname{Hom}_{\rm cont}(E,U_{n+1}(\F_2))
                  \longrightarrow\mathcal X_n.
                                                               \tag{D.72e}
\]
It has exactly the displayed image. A representing
homomorphism can be chosen with the section values
\[
 s_i\longmapsto
       I+\sum_{j=1}^n\chi_j(s_i)E_{j,j+1}
                                  \quad(0\leq i<3).
                                                               \tag{D.72f}
\]
Writing $F_0=0,F_1=1,F_{k+2}=F_{k+1}+F_k$ for the
ordinary Fibonacci numbers, the number of possible
superdiagonal tuples is
\[
                              |\mathcal X_n|=F_{n+2}^{\,3}.
                                                               \tag{D.72g}
\]
No equality of full profinite groups, or solution of the
third-cusp conjugator equation, is inferred from these
cochain and matrix comparisons.
\end{claim}
\passage{Section 31, folios 198-199 verso, scans 391-394; Part D(b), full pure three-cusp involution recognition and the exact remaining conjugator equation}

\begin{proof}
\emph{The three literal restrictions and their normalizers.}
D.34 proves that the elements in (D.72b) have order two
and represent precisely the three $F$-conjugacy classes
of sections. The chosen generator of each quotient is
retained. Put $C_i=C_F(s_i)$ and $N_i=N_E(\langle s_i\rangle)$.
There are inverse continuous homomorphisms
\[
 C_i\times C_2\longrightarrow N_i,
       \quad(h,\tau^\epsilon)\longmapsto hs_i^\epsilon,
 \qquad
 w\longmapsto(ws_i^{-\epsilon},\tau^\epsilon),
                                                               \tag{D.72h}
\]
where $\epsilon$ is the actual image of $w$ in $E/F=C_2$.
A normalizer fixes the unique nonidentity element of its
order-two subgroup, so these formulas have the stated
domains and are inverse. Each $C_i$ is a closed subgroup
of the free profinite group $F$; consequently it has
$2$-cohomological dimension at most one. D.34 gives
$H^1(C_i,\F_2)=0$. Hence all its positive-degree
cohomology with these coefficients vanishes.

The Hochschild--Serre sequence for
$C_i\times C_2\to C_2$ has just its row zero. Inflation
and restriction to its displayed section therefore give
inverse isomorphisms
\[
           H^k(N_i,\F_2)\ \mathop{\rightleftarrows}^\sim\
                         H^k(C_2,\F_2)\qquad(k\geq0).
                                                               \tag{D.72i}
\]
The Brown restriction theorem
\cite[Theorem~2 and the finite-dimensional argument in
\S1.1]{A-KochloukovaZalesskii2007} applies to $E$:
$F$ is open of cohomological dimension one, and every
finite subgroup injects into $E/F$, so $E$ has no
$C_2\times C_2$. The two-row spectral sequence for
$F\to E\to C_2$ bounds $\dim H^k(E,\F_2)$ by three
for $k>1$. Its dense restriction embedding is therefore
onto its Hausdorff target. With the three representatives
and (D.72i), it gives the actual isomorphisms
\[
 (\iota_0^*,\iota_1^*,\iota_2^*):H^k(E,\F_2)
                   \xrightarrow{\ \sim\ }\F_2^3
                                      \qquad(k>1).
                                                               \tag{D.72j}
\]

\emph{Degree one and cochains in every degree.}
Let $u:E\to\F_2$ be the quotient character. Let
$\xi_x,\xi_y:E\to\F_2$ read the two mod-two
abelianization coordinates on $F$, and be zero on $s$.
These are actual characters because $\alpha$ is
minus identity on the full abelianization and therefore
identity modulo two. Every character of $E$ is
uniquely determined by its values on $x,y,s$, and
each triple of values is permitted by the semidirect
relations. Thus $u,\xi_x,\xi_y$ form a basis, with
the exact restriction table
\[
 \begin{array}{c|ccc}
       &s_0&s_1&s_2\\ \hline
 u     &1&1&1\\
 \xi_x &0&1&0\\
 \xi_y &0&0&1.
 \end{array}                                                   \tag{D.72k}
\]
Its inverse basis is
$\eta_0=u+\xi_x+\xi_y,\eta_1=\xi_x,\eta_2=\xi_y$;
each $\eta_i$ restricts to one on $s_i$ and zero
on the other two generators. Consequently (D.72j)
also holds in degree one. Degree zero is the common
scalar field, as in (D.72d).

For an order-two group with generator $\tau$, a
normalized cochain in degree $k>0$ is determined by
its value on $(\tau,\ldots,\tau)$. The differential
is zero: all internal terms have an identity argument,
and the two endpoint terms cancel in characteristic
two. The cup product multiplies the two specified
values. Its normalized cochain algebra is therefore
literally $\F_2[t]$, with $t$ in degree one and zero
differential. The three restrictions agree in degree
zero and give exactly (D.72d), including its cup
multiplication and its unit. This proves that $q$
is a differential graded algebra homomorphism.
The cohomology calculations above prove that it is
a quasi-isomorphism.

Surjectivity is also literal on cochains. Define the
finite quotient character
\[
 \rho=(u,\xi_x,\xi_y):E\longrightarrow\F_2^3,
 \qquad
 v_0=(1,0,0),\quad v_1=(1,1,0),\quad v_2=(1,0,1).
                                                               \tag{D.72l}
\]
For $k>0$, the continuous normalized cochain
\[
 e_{i,k}(g_1,\ldots,g_k)=
   \begin{cases}
       1,&\rho(g_1)=\cdots=\rho(g_k)=v_i,\\
       0,&\text{otherwise}
   \end{cases}
                                                               \tag{D.72m}
\]
satisfies $q(e_{i,k})=t_i^k$. The values $v_i$
are distinct and nonzero, so this formula is
normalized and gives all three basis vectors.
It factors through the stated finite quotient;
no extension of a function from an unspecified
subset is needed. Therefore
$\mathcal K^\bullet=\ker q$ is an acyclic differential
graded ideal, by the cohomology exact sequence of
this degreewise surjection.

The target is the graded algebra associated with
the eight-element Boolean ring $\F_2^3$ in
\cite[Definition~6.1]{R-PalQuick2025}. The displayed
restriction is a direct formality proof for these
particular extensions; it does not apply the
real-projective hypothesis to $E$.

\emph{All defining systems and their continuous lifts.}
Take $n\geq3$ positive-degree classes $a_1,\ldots,a_n$, with
$a_j\in H^{d_j}(E,\F_2)$, whose neighboring products
are zero. Their images in $\mathcal A$ have the
same property, by the proved cohomology isomorphism.
Choose cocycles $c_{j,j+1}$ representing them.
Their $q$-values are exactly the specified images,
because the target differential is zero.

We now choose cochains $c_{i,j}$ for every gap
$j-i\geq2$, including $(i,j)=(1,n+1)$, with
\[
 |c_{i,j}|=\sum_{k=i}^{j-1}d_k-(j-i-1),\qquad
 q(c_{i,j})=0,\qquad
 dc_{i,j}=\sum_{i<k<j}c_{i,k}\smile c_{k,j}.
                                                               \tag{D.72n}
\]
Suppose all smaller gaps have been chosen.
The right side is a cocycle: applying $d$ and
using the equations for smaller gaps makes every
triple product appear twice, so the terms cancel.
Its $q$-image is zero. At gap two this is exactly
the assumed neighboring product. At a larger gap,
each summand has a factor from a previously chosen
higher entry with $q$-value zero. The cocycle is
therefore in $\mathcal K^\bullet$. Acyclicity gives
the required $c_{i,j}$ inside that kernel. All its
degrees are positive because every $d_j\geq1$.
All cochains in this construction are continuous,
since the complex and its kernel consist of
continuous cochains throughout.

Omitting the final entry gives a defining system
for the higher Massey product. Including it in
(D.72n) exhibits its value as a coboundary. Thus
the product is defined and contains zero, with
every degree and every actual lifting complex
specified.

\emph{The exact upper unitriangular image.}
For a tuple in $\mathcal X_n$, apply the same
construction with $d_j=1$ and $c_{j,j+1}=\chi_j$.
The gap recursion also applies for $n=2$: its single
higher entry is obtained from acyclicity of
$\mathcal K^\bullet$ in exactly the displayed degree.
All entries are now normalized one-cochains.
Set
\[
   M(g)_{i,i}=1,\qquad M(g)_{i,j}=c_{i,j}(g) (i<j),
                   \qquad M(g)_{i,j}=0 (i>j).
                                                               \tag{D.72o}
\]
For a one-cochain the equation in (D.72n) reads
\[
 c_{i,j}(gh)=c_{i,j}(g)+c_{i,j}(h)
                     +\sum_{i<k<j}c_{i,k}(g)c_{k,j}(h).
\]
These are exactly all entries of $M(gh)=M(g)M(h)$.
The finite matrix-valued map is continuous and
has the required superdiagonal. Its higher
entries vanish on every $s_i$ because they lie
in $\ker q$, proving (D.72f). Conversely any
continuous unitriangular representation has
\(d c_{j,j+2}=\chi_j\smile\chi_{j+1}\), so its
superdiagonal belongs to $\mathcal X_n$.
This proves both directions of the image statement.
For $n=1$ the map is just the character itself.

Finally use the isomorphism (D.72k) to record
each $\chi_j$ by its three section values.
The neighboring cup product vanishes exactly
when those two three-bit vectors have disjoint
supports, by (D.72j). Each of the three coordinates
therefore gives a binary string of length $n$
with no two consecutive ones. Their number
$b_n$ satisfies $b_0=1,b_1=2$ and
$b_n=b_{n-1}+b_{n-2}$ by its final digit.
Hence $b_n=F_{n+2}$, and the three independent
coordinates give (D.72g).

The same three restriction maps for the standard
involution $\alpha_0(x)=x^{-1},\alpha_0(y)=y^{-1}$
give the explicit zigzag
\[
 \mathcal C^\bullet(E)\xrightarrow{\ q\ }\mathcal A
       \xleftarrow{\ q_{\rm std}\ }
       \mathcal C^\bullet(F\rtimes_{\alpha_0}C_2).
                                                               \tag{D.72p}
\]
Both arrows are the specified quasi-isomorphisms
of differential graded algebras. This is the
exact comparison established here. It leaves
the third-cusp equation of D.35 and the full
section-centralizer equations unchanged.
\end{proof}

\begin{claim}{D.73}{Proved}{The full \(q\)-adic cubic quotient and its section kernels}
Retain the actual action, based generators, character and cusp labels
of D.71 at a prime \(q\ne3\) with \(\lambda_q\ne1\).
Write \(\omega=\lambda_q\in\Z_q^\times\), so that
\(\omega^2+\omega+1=0\) and \(\omega-1\) is a unit.
There is an actual quotient
\[
 \begin{gathered}
 Q_\infty=\Z_q\rtimes_\omega C_3,\qquad
 (a,i)(b,j)=(a+\omega^i b,i+j),\\
 \rho_\infty:F\twoheadrightarrow Q_\infty,\qquad
 \rho_\infty(x)=t=(0,1),\quad
 \rho_\infty(y)=r=(1,0),\quad
 \rho_\infty\alpha=\rho_\infty .
 \end{gathered}                                                 \tag{D.73a}
\]
The labelled extension quotient is
\(E/\ker\rho_\infty=Q_\infty\times\langle s\rangle\).
For \(n\geq1\), let \(\rho_n\) be reduction modulo \(q^n\),
and put
\[
 \begin{gathered}
 N_n=\ker\rho_n,\qquad N_\infty=\ker\rho_\infty,\qquad
 C=C_F(s),\\
 C_n=C\cap N_n,\qquad D=C\cap N_\infty .
 \end{gathered}                                                 \tag{D.73b}
\]
The finite-level cover has type
\(((3q^n-1)/2,3)\), and the following are exact:
\[
 \begin{gathered}
 \rho_\infty(C)=\Z_q,\qquad
 H^1(C_n,\F_3)=0\quad(n\geq1),\qquad H^1(D,\F_3)=0,\\
 1\longrightarrow D\longrightarrow C
       \xrightarrow{\ \rho_\infty|_C\ }\Z_q\longrightarrow1 .
 \end{gathered}                                                 \tag{D.73c}
\]
The last sequence splits by a specified procyclic-power construction
in the proof. In particular \(E\) contains the actual subgroup
\(\Z_q\times\langle s\rangle\) supplied by that section.

For both \(n\geq1\) and \(n=\infty\), there are exactly three
\(N_n\)-conjugacy classes of sections in
\(E_n=N_n\rtimes\langle s\rangle\). They have representatives
\[
 T_i=x^i\langle s\rangle x^{-i}\quad(0\leq i<3),
                                                               \tag{D.73d}
\]
and their full normalizers are, respectively,
\(x^i(C_n\times\langle s\rangle)x^{-i}\) and
\(x^i(D\times\langle s\rangle)x^{-i}\).
At each finite level, for the actual rotation projector \(e_n\)
on \(V_n=H^1(N_n,\Z_3)\), the pullback action of
\(\alpha^{-1}\) gives
\[
 (1-e_n)V_n\cong\Z_3[C_3]^{\,q^n-1};
 \qquad e_nV_n\bmod3\cong J_2\oplus J_2 .
                                                               \tag{D.73e}
\]
The first isomorphism is obtained by lifting cyclic vectors,
with the actual action retained.

There is also an exact finite obstruction to two-section generation:
the closed image of every continuous homomorphism
\[
                         C_3\amalg C_3\longrightarrow E
                                                               \tag{D.73f}
\]
is proper. This is proved in the actual quotient
\(Q_1\times C_3\), with both full section classes labelled below.
These statements retain the remaining cubic-coordinate question:
they do not assert existence or exclusion of that coordinate,
or a presentation of the full extension by its section subgroups.
\end{claim}
\passage{Section 30, folios 194-196, scans 383-387; Section 31, folios 198-199 verso, scans 391-394; Part D(b), full loop-outer finite actions in type (1,1)}

\begin{proof}
\emph{The quotient before reduction at \(q\).}
The based character change in D.71 keeps
\(c=[x,y]\) exactly and gives
\(\chi(x)=1,\chi(y)=0\). Its index-three cover has
\[
 F_0=\ker\chi,\qquad
 u=x^3,\quad B_i=x^iyx^{-i}\ (0\leq i<3),\qquad
 B_3=uB_0u^{-1}.
                                                               \tag{D.73g}
\]
Use the continuous additive character already constructed
integrally in (D.71l)--(D.71m):
\[
 \ell_q:F_0\longrightarrow\Z_q,\qquad
 \ell_q(u)=0,\quad\ell_q(B_i)=\omega^i,\qquad
                    \ell_q\alpha=\ell_q .
                                                               \tag{D.73h}
\]
In particular its image is all of \(\Z_q\), since it contains \(1\).
Conjugation by \(x\) multiplies it by \(\omega\);
conjugation by \(F_0\) fixes it. Thus its kernel is a
closed normal subgroup of \(F\), preserved in both directions
by the actual automorphism \(\alpha\).

The universal homomorphism from the full free profinite group
which sends \(x\) to \(t\) and \(y\) to \(r\) in (D.73a)
is onto: the closure of the powers of \(r\) is the entire
normal \(\Z_q\), and \(t\) gives the complement.
Its restriction to \(F_0\), followed by the rotation coordinate,
has exactly the four values (D.73h), so its kernel is
\(\ker\ell_q\). The inverse quotient map is
\[
 Q_\infty\longrightarrow F/\ker\ell_q,\qquad
 (a,i)\longmapsto y^{e_q a}x^i\ker\ell_q ,
                                                               \tag{D.73i}
\]
where \(e_q a\in\widehat\Z\) has \(q\)-coordinate \(a\)
and every other primary coordinate zero.
Here \(y^{e_q a}\ker\ell_q\) is the rotation \(a\),
\(x^3\in\ker\ell_q\), and conjugation by \(x\) is
multiplication by \(\omega\). These verify multiplication
and both inverse compositions, including the labels.

The induced automorphism fixes \(r\) and sends \(t\)
to \(r^j t\) for some \(j\in\Z_q\). Since it fixes every
rotation, its actual cube sends \(t\) to \(r^{3j}t\).
The equality \(\alpha^3=1\) and invertibility of \(3\)
in \(\Z_q\) give \(j=0\). This proves
\(\rho_\infty\alpha=\rho_\infty\).
The extension comparison and its inverse are consequently
\[
 \begin{aligned}
 E/N_\infty&\longrightarrow Q_\infty\times C_3,
       &fs^jN_\infty&\longmapsto(\rho_\infty(f),s^j),\\
 Q_\infty\times C_3&\longrightarrow E/N_\infty,
       &(z,s^j)&\longmapsto fs^jN_\infty
                       \quad(\rho_\infty(f)=z).
 \end{aligned}                                                 \tag{D.73j}
\]
The choice of \(f\) changes neither coset, and the maps
are inverse continuous homomorphisms.

\emph{The fixed subgroup and its explicit procyclic section.}
D.48 gives \(H^1(C,\F_3)=0\). Therefore
\(\chi(C)=0\), so \(\rho_\infty(C)\) lies in the rotation
\(\Z_q\). It is closed, since it is a continuous image
of the compact group \(C\). D.71 proves that its
reduction modulo \(q\) is all of \(C_q\).
Every nonzero closed subgroup of the additive group
\(\Z_q\) is \(q^a\Z_q\): choose an element with least
valuation, whose integer powers are dense in that
subgroup. Reduction modulo \(q\) thus forces \(a=0\).
This proves the surjection in (D.73c).

Choose \(h\in C\) with \(\rho_\infty(h)=r\).
For \(a\in\Z_q\) define
\[
                 \iota(a)=h^{e_q a}.
                                                               \tag{D.73k}
\]
Continuous powers in the actual procyclic closure of
\(\langle h\rangle\) show that this is a continuous
homomorphism. Its composite with \(\rho_\infty\) is
\(a\mapsto r^a\), so it is injective and splits the
displayed surjection. In particular the domain is the
full \(\Z_q\), rather than a finite cyclic quotient.
With the conjugation action given by this actual
section there are inverse maps
\[
 \begin{aligned}
 D\rtimes_\iota\Z_q&\longrightarrow C,
                   &(d,a)&\longmapsto d\iota(a),\\
 C&\longrightarrow D\rtimes_\iota\Z_q,
                   &h'&\longmapsto
          \bigl(h'\iota(\rho_\infty(h'))^{-1},
                        \rho_\infty(h')\bigr).
 \end{aligned}                                                 \tag{D.73l}
\]
In these expressions the value of \(\rho_\infty\) on \(C\)
means its additive rotation coordinate.
Every \(\iota(a)\) commutes with \(s\).
Thus \((a,s^j)\mapsto\iota(a)s^j\) gives the asserted
embedded direct product, with inverse given by the
rotation and original quotient-generator labels.

\emph{Every finite power cover, with its integral lattice.}
Put \(m=q^n\), and let \(\omega_n\) denote the reduction
of \(\omega\) modulo \(m\). The quotient \(Q_n\) has
order \(3m\), with exact multiplication obtained by
reducing (D.73a); the actual quotient maps between levels
are reduction of the first coordinate.
The cusp \(c=[x,y]\) has image \(r^{\omega_n-1}\),
which generates the rotation \(C_m\), since
\(\omega-1\) is a \(q\)-adic unit.
It has three cycles of length \(m\) on the \(3m\)
covering sheets. Thus \(N_n\) has three cusps and
Euler characteristic \(-3m\), giving the stated type
\(((3m-1)/2,3)\). A set of actual cusp representatives is
\[
                          d_i=x^i c^m x^{-i}\quad(0\leq i<3).
                                                               \tag{D.73m}
\]

Here the graph-cycle computation in D.71 has the following
literal version for \(m\). Set
\(\Lambda_m=\Z_3[C_m]\), \(\nu_m=1+r+\cdots+r^{m-1}\),
and let \(a,b\in\{1,\ldots,m-1\}\) represent
\(\omega_n,\omega_n^2\). The four Schreier generators
\(u,B_0,B_1,B_2\) have rotation images \(0,1,a,b\).
The graph boundary and its inverse cycle coordinates are
\[
 \begin{aligned}
 \partial(A,B,D',E')&=(r-1)B+(r^a-1)D'+(r^b-1)E',\\
 (A,D',E',k)&\longmapsto
    (A,-P_aD'-P_bE'+k\nu_m,D',E'),\\
 P_j&=1+r+\cdots+r^{j-1}.
 \end{aligned}                                                 \tag{D.73n}
\]
The inverse reads \(A,D',E'\) and the unique \(k\) for
which \(B+P_aD'+P_bE'=k\nu_m\).
Indeed the annihilator of \(r-1\) in \(\Lambda_m\)
is exactly \(\Z_3\nu_m\), by comparing adjacent
coefficients. These are actual inverse
\(\Lambda_m\)-module maps, so
\[
 H_1(N_n,\Z_3)\cong\Z_3\oplus\Lambda_m^3 .
                                                               \tag{D.73o}
\]
Duality gives the same rotation multiplicities for \(V_n\).
The projector \(e_n=m^{-1}\sum_{g\in C_m}g\) is integral,
and restriction and transfer identify \(e_nV_n\) with
the actual rank-four module \(H^1(F_0,\Z_3)\).
They are equivariant for \(\alpha\), since the actual
quotient action is the identity. Its two Jordan-two
blocks are exactly those computed in (D.71s).

For completeness, the paired-lattice argument applies at
every level with the following precise quotient maps.
Fill the three actual cusps (D.73m), then take the
maximal pro-three quotient:
\[
 P_n=\bigl(N_n/
       \overline{\langle\!\langle d_0,d_1,d_2\rangle\!\rangle}
                                                  \bigr)^{(3)} .
                                                               \tag{D.73p}
\]
It is the closed orientable genus-\((3m-1)/2\)
surface pro-three group by the same finite-target
filling universal property used in D.71.
Write the cusp normalizer element as \(v=f_0s\);
then \(\alpha(c)=f_0^{-1}c^\lambda f_0\).
For \(g=x^i\), the quotient \(Q_n/C_m\) selects a unique
\(j\in\{0,1,2\}\); since \(\rho_n(c)\) generates \(C_m\),
there is \(k\in\Z\) and \(h\in N_n\) such that
\[
 \alpha(g)f_0^{-1}=h x^j c^k,\qquad
 \alpha(g c^m g^{-1})=h d_j^\lambda h^{-1}.
                                                               \tag{D.73q}
\]
The inverse automorphism gives the opposite inclusion
of the normal closures. Thus the actual action descends
to (D.73p).
The injection \(H^1(P_n,\Z_3)\to V_n\) has as its image
exactly the characters killing the three \(d_i\).
All cusp evaluations are rotation-invariant, so
\(W_n=(1-e_n)V_n\) lies in this image.
The averaged filled summand has rank two, with the
nonidentity integral action dual to the base torus
matrix \(A_3\); this is the same restriction/transfer
comparison and cusp vanishing test as in D.71,
with \(m\) a unit in \(\Z_3\).
The rotation deck transformations extend positively
over the filled cusps. Hence \(e_n\) is self-adjoint
for the unimodular alternating cup form, and its
orthogonal summand \(W_n\) is integrally unimodular.

Choose a finite unramified extension \(\mathcal O/\Z_3\)
containing all \(m\)-th roots of unity. The idempotents
\[
       e_\psi=\frac1m\sum_{g\in C_m}\psi(g)^{-1}g
                                                               \tag{D.73r}
\]
split \(W_n\otimes\mathcal O\) into rank-three lattices
for all nontrivial \(\psi\).
Since \(m\) is odd, \(\psi\ne\psi^{-1}\); the cup form
pairs precisely those two lattices, perfectly.
The action \(S\), pullback under \(\alpha^{-1}\),
commutes with all these projectors.
Its action on \(H^2(P_n,\Z_3)\) is trivial:
it has order dividing three, whereas \(\Z_3^\times\)
has no nontrivial element of order three.
Thus \(S\) preserves the cup form used here.

The integral fixed-cup lemma proved in D.71 applies
to the actual action on \(P_n\).
Its nonidentity hypothesis, including every inner-twisted
section action, follows from the retained rank-two
averaged action just specified.
Thus pairings of two \(S\)-fixed integral characters
are divisible by three, and the same is true after
flat extension to \(\mathcal O\).
The polynomial \(X^2+X+1\) remains irreducible over
the fraction field of this unramified extension.
On each rank-three character lattice the fixed rank
is therefore one or three. Fixed rank three on
one lattice would give the identity there and on
its perfectly paired inverse lattice, contradicting
that divisibility and the perfect form.
The fixed rank is consequently one on both.
Choose primitive fixed generators \(u,u'\).
For \(w\) with \(\langle w,u'\rangle=1\), write
\[
   (1+S+S^2)w=b u,\qquad
                  b\langle u,u'\rangle=3 .
                                                               \tag{D.73s}
\]
The divisibility and unramified valuation show that
\(\langle u,u'\rangle\) is three times a unit and
\(b\) is a unit. Thus the norm has rank one modulo
three on this rank-three lattice, giving one Jordan
block \(J_3\).
After scalar extension all nontrivial character
summands have this property. Hence over \(\F_3\)
the norm has rank \(m-1\) on a space of dimension
\(3(m-1)\). Every Jordan block has size at most three;
the rank forces every one to have size three.
Lift cyclic vectors for these actual blocks.
The resulting homomorphism
\(\Z_3[C_3]^{m-1}\to W_n\) reduces to an isomorphism
modulo three, hence has unit determinant and an
equivariant inverse. This proves (D.73e) at every
finite level, without a descent or generation
assumption about the full groups.

\emph{The exact section classes and their fixed kernels.}
The other full section generator \(s_1\) has
\(\chi(s_1)=-1\), whereas \(\chi(s)=0\), as in (D.71e).
Consequently none of its \(F\)-conjugates belongs
to \(E_n\), for finite \(n\) or for \(n=\infty\).
Every \(F\)-conjugate of \(s\) does belong, because
the image of \(s\) in (D.73j) is central.
The criterion for two of these sections to be
\(N_n\)-conjugate gives the exact bijections
\[
 \begin{aligned}
 Q_n/\rho_n(C)&\longrightarrow
        \{\text{\(N_n\)-classes of sections in \(E_n\)}\},\\
 \rho_n(f)\rho_n(C)&\longmapsto
                         [f\langle s\rangle f^{-1}] .
 \end{aligned}                                                 \tag{D.73t}
\]
The identical maps apply to \(Q_\infty\).
Indeed a conjugacy between the labelled generators
holds precisely when the corresponding elements
of \(F\) differ by a left \(N_n\)-factor and a
right \(C\)-factor. An \(E_n\)-conjugator may be
multiplied by an element of the source section to
remove its quotient label, so it gives the same
classes as \(N_n\)-conjugacy.
The proved \(\rho_\infty(C)=\Z_q\) makes every coset
space in (D.73t) the three cosets represented by
\(1,t,t^2\).
Any normalizer fixes the labelled generator,
because its conjugation action on \(E/F=C_3\) is
trivial. Removing its quotient power gives exactly
the normalizers stated after (D.73d).

At finite level, the LHS sequence for
\(N_n\to E_n\to C_3\), together with (D.73e), bounds
\[
                    \dim H^j(E_n,\F_3)\leq3\qquad(j\geq2).
                                                               \tag{D.73u}
\]
In fact the two \(J_2\) blocks each contribute one
to positive cyclic-group cohomology, while the
regular blocks contribute zero.
The subgroup \(N_n\) is open with cohomological
dimension one, and every finite subgroup of \(E_n\)
injects into \(C_3\). The full profinite Brown
restriction theorem \cite[Theorem~2 and \S1.1]{A-KochloukovaZalesskii2007}
therefore applies, with no
\(C_3\times C_3\) subgroup. The displayed bound
makes its dense restriction embedding onto.
Each \(C_n\) is closed in \(F\), of
\(3\)-cohomological dimension at most one.
The three actual direct-product normalizers give
\[
        3\bigl(1+\dim H^1(C_n,\F_3)\bigr)
                            =\dim H^j(E_n,\F_3)\leq3
                                  \qquad(j\geq2).
                                                               \tag{D.73v}
\]
This proves the finite-level vanishing.

It also proves the vanishing for \(D\), as follows
directly on continuous characters.
Let \(\xi:D\to\F_3\) be continuous.
Choose an open normal subgroup \(U\) of \(C\) with
\(D\cap U\subseteq\ker\xi\).
Then \(\xi\) descends to \(D/(D\cap U)\cong DU/U\)
and inflates to a character of \(DU\).
The quotient \(C/D=\Z_q\) identifies the open
subgroup \(DU\) with the preimage of some
\(q^n\Z_q\), including possibly \(n=0\).
Thus \(DU=C_n\) when \(n\geq1\), or \(DU=C\)
when \(n=0\). All these character groups vanish,
by (D.73v) and D.48. Hence \(\xi=0\).
The normalizer and class statements at \(n=\infty\)
were already proved by the actual coset maps;
no finite-type surface assertion for \(N_\infty\)
is made.

\emph{The finite image of any two torsion generators.}
Reduce (D.73j) modulo \(q\), writing
\(Q_1=\langle r,t\rangle\).
Every nonidentity finite-order element of \(E\)
has order three, since \(F\) is torsion-free and
its image in \(E/F\) is injective.
Replace it by its inverse if needed to give the
specified quotient-generator label \(s\).
The two full section classes then have images
\[
                  s,\qquad r^a t^{-1}s\quad(a\in\F_q).
                                                               \tag{D.73w}
\]
For the first class this follows from centrality
of \(s\) in the quotient; for the other it follows
from \(\chi(s_1)=-1\), and conjugation by rotations
gives exactly the displayed possibilities because
\(\omega_1-1\) is a unit.

Two elements from the first class generate at most
\(\langle s\rangle\).
One element from each class generates the order-nine
abelian subgroup
\(\langle s,r^a t^{-1}s\rangle\), which omits the
order-\(q\) rotation.
Two from the second class lie in the kernel of
the explicit homomorphism
\[
 Q_1\times C_3\longrightarrow C_3,\qquad
                       (r^a t^i,s^j)\longmapsto i+j .
                                                               \tag{D.73x}
\]
This is again a proper subgroup. The conclusions
also hold if one image is the identity.
Thus every subgroup generated by two finite-order
elements has proper image in this particular
finite quotient. A homomorphism from the full
free profinite product in (D.73f) has precisely
such generator images, proving its proper-image
assertion. No converse generation statement is used.

The exact producer
\path{code/part_c/torus_cubic_qadic_d73.py}
checks the compatible prime-power roots, quotient
coordinate maps, three cusp cycles, graph-module
inverse formulas, procyclic primary splitting
coordinates, and all two-section image types in
the specified finite cases.
It passes 155,189 exact checks over fourteen
prime-power cases and 3,252 section-image pairs.
The infinite quotient, fixed-kernel vanishing
and integral lattice conclusions are the actual
maps and cohomology arguments above.
\end{proof}

\begin{claim}{D.74}{Proved}{Full section centralizers from twisted relative duality}
Retain the actual normalized involution, full profinite group,
and labelled extension of D.34:
\[
 \begin{gathered}
 F=\widehat F(x,y),\qquad z=(xy)^{-1},\qquad
 E=F\rtimes_\alpha\langle s\rangle,\\
 \alpha(x)=x^{-1},\qquad \alpha(y)=f^{-1}y^{-1}f,
 \qquad f\in F'.
 \end{gathered}
 \tag{D.74.1}
\]
Thus $\alpha^2=1$, and its multiplier at every prime is
the global element $-1$.  Put
\[
 s_0=s,\qquad s_1=x^{e_2}s,\qquad s_2=y^{e_2}fs,
 \qquad S_i=\langle s_i\rangle,\qquad C_i=C_F(s_i).
 \tag{D.74.2}
\]
Then
\[
                 C_i=1,\qquad N_E(S_i)=S_i
                       \quad(0\leq i<3).
 \tag{D.74.3}
\]
These equalities concern the full profinite groups.  The
extension $E$ is real projective in Haran's sense, with the
actual finite embedding-problem formulation given below.
There is exactly one $h\in F$ satisfying
$f=h^{-1}\alpha(h)$; it lies in $F'$.
For $\alpha\in GT_0$, this same unique element satisfies
$h\theta(h)=1$, where $\theta$ exchanges $x$ and $y$.

The proof gives an exact finite-coefficient boundary
restriction calculation.  D.35's remaining peripheral
conjugator problem is the displayed equation (D.74.40)
with this unique $h$.  Its solution and a free-product
presentation of $E$ are not assertions of this result.
\end{claim}
\passage{Section 30, folios 194-196, scans 383-387; Section 31, folios 198-199 verso, scans 391-394; Section 24, folio 147 verso, scan 293; Part D(b), full loop-outer finite actions in type (0,3)}

\begin{proof}
\emph{The three actual peripheral extensions.}
Let $L_1=\overline{\langle x\rangle}$,
$L_2=\overline{\langle y\rangle}$, and
$L_3=\overline{\langle z\rangle}$, and set
$D_j=N_E(L_j)$.  A.7 and D.34 give the exact maps
\[
 \begin{split}
 1\longrightarrow L_j\longrightarrow D_j
       &\xrightarrow{\ u\ }C_2\longrightarrow1,\\
 \widehat{\mathbb Z}\rtimes_{-1}C_2
       &\longrightarrow D_j,\qquad
       (a,\tau^\epsilon)\longmapsto c_j^a b_j^\epsilon ,
 \end{split}
 \tag{D.74.4}
\]
where $(c_1,c_2,c_3)=(x,y,z)$ and $b_j$ is an actual
involution in $D_j$ with $u(b_j)=1$.
The inverse of the second map reads $\epsilon=u(d)$ and
then the unique exponent of $d b_j^{-\epsilon}\in L_j$.
Here and throughout, $u:E\to C_2=\mathbb F_2$ is the
specified quotient map.

The two $D_j$-conjugacy classes of subgroups of order two
are represented by $\langle b_j\rangle$ and
$\langle c_j^{e_2}b_j\rangle$.  Each has normalizer
equal to itself in $D_j$: an element of $L_j$ commuting
with either involution has exponent $a$ with $2a=0$,
and hence is the identity.
Every finite subgroup of $E$ injects into $E/F$, because
$F$ is torsion-free.  D.34 identifies all three
$E$-conjugacy classes of nontrivial finite subgroups as
$S_0,S_1,S_2$ and proves that the six occurrences just
listed form the endpoints of a triangle.  Consequently
each $S_i$ occurs exactly twice among those six classes.

The full normalizers are already specified without any
assumption on $C_i$.  There are inverse continuous maps
\[
 \begin{aligned}
 C_i\times S_i&\longrightarrow N_E(S_i),
                         &(h,s_i^\epsilon)&\longmapsto hs_i^\epsilon,\\
 N_E(S_i)&\longrightarrow C_i\times S_i,
                         &w&\longmapsto(ws_i^{-\epsilon},s_i^\epsilon),
                         \quad \epsilon=u(w).
 \end{aligned}
 \tag{D.74.5}
\]
A normalizer fixes the unique nonidentity element of
$S_i$, which proves that the second formula has its
stated target.  These are the normalizer maps that will
occur in every restriction below.

\emph{One actual coefficient module for the relative theory.}
Write $k=\mathbb F_2$.  Form the compact permutation
$k[[E]]$-module and its augmentation kernel
\[
 P=\bigoplus_{j=1}^3 k[[E/D_j]],\qquad
 0\longrightarrow\Delta\longrightarrow P
             \xrightarrow{\ \varepsilon\ } k\longrightarrow0 .
 \tag{D.74.6}
\]
There is an actual $F$-equivariant homeomorphism
\[
             F/L_j\longrightarrow E/D_j,\qquad
                        aL_j\longmapsto aD_j .
 \tag{D.74.7}
\]
It is injective since $F\cap D_j=L_j$ and surjective
since $E=FD_j$; compactness gives its inverse continuity.
Thus $\Delta|_F$ is precisely the augmentation kernel
for the three specified full peripheral subgroups of
$F$.  The additional $E$-action is retained, rather
than replacing it by an action fixing the $L_j$
pointwise.  In the coordinates of (D.74.7), it is
\[
       e\cdot(aL_j)=e a b_j^{-u(e)}L_j .
 \tag{D.74.8}
\]
The displayed representative lies in $F$.  Changing
$a$ by $a\ell$, $\ell\in L_j$, gives the same coset
because $b_j$ normalizes $L_j$.  This is the action
transported from $E/D_j$, so it obeys the group law.

For a finite continuous $E$-module $M$ over $k$, put
\[
 \mathscr H_M=\operatorname{Hom}_{k,\mathrm{cont}}(\Delta,M),
 \qquad (e\varphi)(d)=e\varphi(e^{-1}d).
 \tag{D.74.9}
\]
This is a discrete continuous $E$-module.  A continuous
map from the compact module $\Delta$ to a finite
discrete module factors through a finite quotient;
the compact-open topology on this Hom space is
discrete, and each such map has an open stabilizer.
In particular ordinary continuous group cohomology
with coefficients $\mathscr H_M$ is defined.

The relative groups used here are, in every positive
degree, the following specified ordinary groups:
\[
 R^n(E,M):=H^n(E,\{D_j\};M)
       =H^{n-1}(E,\mathscr H_M)\qquad(n\geq1).
 \tag{D.74.10}
\]
The corresponding formula for $F$ uses precisely
$\Delta|_F$ and the $L_j$ by (D.74.7).
For clarity, applying continuous Hom to (D.74.6) gives
the exact sequence of discrete $E$-modules
\[
 0\longrightarrow M
 \longrightarrow\operatorname{Hom}_{k,\mathrm{cont}}(P,M)
 \longrightarrow\mathscr H_M\longrightarrow0 .
 \tag{D.74.11}
\]
Exactness follows also directly because (D.74.6)
splits as a sequence of compact $k$-vector spaces.
Shapiro on the three summands gives the actual long
exact sequence
\[
 \cdots\longrightarrow H^n(E,M)
 \xrightarrow{\ r\ }\bigoplus_{j=1}^3 H^n(D_j,M)
 \xrightarrow{\ \delta\ }R^{n+1}(E,M)
 \xrightarrow{\ j\ }H^{n+1}(E,M)\longrightarrow\cdots .
 \tag{D.74.12}
\]
Here $r$ consists of literal subgroup restrictions.
These are the augmentation and Hom constructions of
\cite[Definition~2.1, Propositions~2.3--2.4 and~2.10]
{C-WilkesRelativeAccepted2018}.  In particular the spectral sequence
below will be the ordinary extension spectral sequence
for the actual coefficient module (D.74.9).

\emph{The pairing and its equivariance.}
Let $M^\vee=\operatorname{Hom}_k(M,k)$, with
$(e\ell)(m)=\ell(e^{-1}m)$.  There is the actual
$E$-equivariant coefficient pairing
\[
 M\otimes_k\mathscr H_{M^\vee}\longrightarrow\mathscr H_k,
 \qquad
 m\otimes\varphi\longmapsto
                \bigl[d\longmapsto\varphi(d)(m)\bigr].
 \tag{D.74.13}
\]
Its cup products identify, under (D.74.10), with the
ordinary absolute-relative cup products.  This
identification can be made on the completed bar
resolution: tensor that resolution over $k$ with
$\Delta$, with diagonal group action, and use the
continuous adjunction
\[
 \operatorname{Hom}_F(P_\bullet(F)\widehat\otimes_k\Delta,M)
 \simeq
 \operatorname{Hom}_F(P_\bullet(F),\mathscr H_M).
 \tag{D.74.14}
\]
The augmentation makes the tensor complex a resolution
of $\Delta$; its terms are projective, since the diagonal
action on a free group-module factor can be untwisted
and compact $k$-vector spaces are projective over $k$.
The Alexander--Whitney diagonal followed by (D.74.13)
is the stated cup product.  This also proves compatibility
with the coefficient sequence (D.74.11) and its connecting
maps.  Equivalently these compatibilities are the
natural cup constructions in
\cite[Section~3, especially Propositions~3.1 and~3.5]
{C-WilkesRelativeAccepted2018}.

The full profinite pair $(F,\{L_j\})$ is the completion
of the ordinary pair of pants with its three boundary
subgroups.  Its mod-two duality is the actual cup
duality
\[
 \begin{split}
 H^0(F,M)\otimes H^2(F,\{L_j\};M^\vee)&\longrightarrow k,\\
 H^1(F,M)\otimes H^1(F,\{L_j\};M^\vee)&\longrightarrow k,
 \end{split}
 \tag{D.74.15}
\]
and both pairings are perfect.  One direct comparison
is restriction along the dense inclusions
$F(x,y)\to F$ and $\langle c_j\rangle\to L_j$.
For finite coefficient modules these induce the usual
cohomology isomorphisms for a free group and for an
infinite cyclic group.  The relative long exact
sequences give the relative isomorphisms, with the
same cup products by bar restriction.  Thus the
perfect pairings are those of the compact surface
with its boundary, including its original peripheral
maps.  The completed surface-pair statement and its
cup formulation are also
\cite[Proposition~6.16 and Proposition~5.7]
{C-WilkesRelativeAccepted2018}.

All these pairings are $E/F$-equivariant.  Indeed
(D.74.13) is $E$-equivariant, and the target
\[
 H^1(F,\mathscr H_k)
       =H^2(F,\{L_j\};k)=k
 \tag{D.74.16}
\]
has trivial $E/F$-action.  This last assertion is
literal: the target is one-dimensional over $k$,
so its automorphism group is trivial.  Its
fundamental coclass is the boundary-sum coclass;
the conjugations in (D.74.8) act by inversion on
each peripheral exponent and hence act trivially
on every mod-two boundary character.  No geometric
representative for $\alpha$ is needed for this
equivariance.

Put $A=H^0(F,M)$ and $B=H^1(F,M)$.  They are finite
$k[C_2]$-modules.  Equations (D.74.10) and (D.74.15)
give the following \emph{equivariant} identifications:
\[
 H^0(F,\mathscr H_{M^\vee})=B^\vee,\qquad
 H^1(F,\mathscr H_{M^\vee})=A^\vee .
 \tag{D.74.17}
\]
Every higher row is zero, since $\operatorname{cd}_2(F)=1$.

\emph{The two transgressions have equal rank.}
For a finite $k[C_2]$-module $V$, write
$T=1+\sigma$ and
\[
                 \overline V=\ker T/\operatorname{im}T .
 \tag{D.74.18}
\]
The normalized positive-degree cyclic cochain complex
has value group $V$ in every degree and differential
$T$.  Thus $H^a(C_2,V)=\overline V$ for $a\geq1$.
Multiplication by the standard class
$u\in H^1(C_2,k)$ is the identity under these
identifications.  The pairing
\[
              \overline V\otimes\overline{V^\vee}
                         \longrightarrow k
 \tag{D.74.19}
\]
induced by evaluation is perfect.  To check this,
the annihilator of $\ker T$ is $\operatorname{im}T^\vee$,
and the annihilator of $\operatorname{im}T$ is
$\ker T^\vee$.  The radicals on the kernels are
therefore precisely the images being divided out.
On invariant representatives, the cyclic cup formula
is evaluation itself, since every factor $\sigma^a$
acts trivially on such a representative.

Apply the ordinary Hochschild--Serre spectral sequence
to $F\to E\to C_2$, first with coefficient $M$ and
then with coefficient $\mathscr H_{M^\vee}$.
Its $E_2$ terms in the two cases are
\[
 H^a(C_2,H^b(F,M)),\qquad
 H^a(C_2,H^b(F,\mathscr H_{M^\vee})).
 \tag{D.74.20}
\]
Both have exactly the two possibly nonzero rows
$b=0,1$.  In positive columns their rows are,
respectively,
\[
 \begin{array}{c|cc}
       &b=0&b=1\\ \hline
 \text{absolute}&\overline A&\overline B\\
 \text{shifted relative}&\overline{B^\vee}&\overline{A^\vee}.
 \end{array}
 \tag{D.74.21}
\]
The coefficient pairing (D.74.13), on the same filtered
bar resolutions, supplies a product of these spectral
sequences into that for $\mathscr H_k$.  The
Alexander--Whitney product preserves the quotient
filtration and obeys the cochain Leibniz identity;
therefore its $d_2$ obeys that identity as well.
Its products on the fibre rows are exactly
(D.74.15)--(D.74.17).

Only $d_2$ can be nonzero.  In any sufficiently
positive columns denote the two maps by
\[
 d:\overline B\longrightarrow\overline A,\qquad
 d':\overline{A^\vee}\longrightarrow\overline{B^\vee}.
 \tag{D.74.22}
\]
They do not depend on a further column shift, under
the cyclic identifications, because they commute
with multiplication by $u$.  Take an absolute
row-one element $b$ and a shifted-relative row-one
element $a^\vee$, both in positive columns.
Their product on the fibre row two is zero:
$H^2(F,\mathscr H_k)=0$.  The $d_2$ Leibniz
identity in the target fibre row one is consequently
\[
           \langle d b,a^\vee\rangle+
                    \langle b,d'a^\vee\rangle=0 .
 \tag{D.74.23}
\]
By (D.74.19) these are perfect pairings in the
positive quotient columns.  Hence $d'$ is the
transpose of $d$ for the specified pairings, and
they have a common rank $r$.

There are no later differentials.  It follows that
the two stable total dimensions are
\[
 \begin{split}
 \dim H^n(E,M)
       &=\dim\overline A+\dim\overline B-2r,\\
 \dim H^{n-1}(E,\mathscr H_{M^\vee})
       &=\dim\overline{B^\vee}+
                    \dim\overline{A^\vee}-2r
        =\dim H^n(E,M)
 \end{split}
 \tag{D.74.24}
\]
for all sufficiently large $n$.  The degree shift
in the second line is accounted for by the explicitly
proved one-periodicity (D.74.18) of the positive
quotient columns.  More precisely,
multiplication by the actual inflated $u\in H^1(E,k)$
is an isomorphism in sufficiently high degree on
each group in (D.74.12): on the two-row spectral
sequences it is an isomorphism on both surviving
positive-column rows, and hence on their finite
filtrations.  The same proof applies to each $D_j$,
using its open subgroup $L_j$.

\emph{The rank of the actual boundary restriction.}
Suppose now that $M$ has a specified invariant
perfect bilinear pairing with itself, so that
$M\simeq M^\vee$ as $E$-modules.  This will be the
permutation pairing used below.
For precision, take degree-zero localizations by
the actual class $u$:
\[
 \begin{split}
 \mathcal A_M&=H^\bullet(E,M)[u^{-1}]_0,\\
 \mathcal B_M&=\left(\bigoplus_jH^\bullet(D_j,M)\right)
                                                    [u^{-1}]_0,\\
 \mathcal R_M&=R^\bullet(E,M)[u^{-1}]_0 .
 \end{split}
 \tag{D.74.25}
\]
These are finite vector spaces; each is an actual
sufficiently high cohomology group under the
proved multiplication-by-$u$ maps.  All maps
in (D.74.12) commute with these multiplications.
Localization, or equivalently the direct limit
of that exact sequence under those maps, gives
the exact cyclic sequence
\[
 \mathcal A_M\xrightarrow{\ \bar r\ }\mathcal B_M
       \xrightarrow{\ u^{-1}\bar\delta\ }\mathcal R_M
       \xrightarrow{\ \bar j\ }\mathcal A_M .
 \tag{D.74.26}
\]
For example the middle arrow sends a class
$b/u^n$ to $\delta(b)/u^{n+1}$.
Equation (D.74.24) gives
$\dim\mathcal R_M=\dim\mathcal A_M$.
If the ranks of the three arrows in (D.74.26) are
$a,b,c$, exactness says
\[
 \dim\mathcal A_M=a+c,\qquad
 \dim\mathcal B_M=a+b,\qquad
 \dim\mathcal R_M=b+c .
 \tag{D.74.27}
\]
Consequently $a=b$, and therefore
\[
             \operatorname{rank}\bar r
                         =\tfrac12\dim\mathcal B_M .
 \tag{D.74.28}
\]
This is a rank equality for the actual restriction
map with the original boundary labels.

\emph{Brown restriction and the two occurrences of each section.}
The full finite-coefficient Brown restriction theorem
\cite[Theorem~2 and Section~1.1]
{A-KochloukovaZalesskii2007} applies to $E$:
$F$ is an open subgroup of $2$-cohomological
dimension one, and $E$ has no $C_2\times C_2$.
The two-row sequence (D.74.20) shows that
$H^n(E,M)$ is finite for finite $M$.  Thus the
theorem's dense restriction embedding is onto
its Hausdorff target.  The three full normalizers
are exactly those in (D.74.5).  It follows, for
$n>1$, that
\[
 H^n(E,M)\xrightarrow{\ \sim\ }
                       \bigoplus_{i=0}^2H^n(N_E(S_i),M).
 \tag{D.74.29}
\]
The same theorem applies to each $D_j$, whose
open subgroup $L_j$ has dimension one and whose
two section normalizers are the finite sections
themselves.  Its cohomology is finite by its
two-row cyclic spectral sequence.  Consequently
\[
 H^n(D_j,M)\xrightarrow{\ \sim\ }
           H^n(T_{j,0},M)\oplus H^n(T_{j,1},M),
                    \qquad n>1 ,
 \tag{D.74.30}
\]
where $T_{j,0}=\langle b_j\rangle$ and
$T_{j,1}=\langle c_j^{e_2}b_j\rangle$.

Choose, for each of these six occurrences, an actual
$a_{j,\epsilon}\in F$ with
\[
 T_{j,\epsilon}=a_{j,\epsilon}S_{i(j,\epsilon)}
                                    a_{j,\epsilon}^{-1}.
 \tag{D.74.31}
\]
Such an $F$-conjugator exists by D.34, which identifies
$E$- and $F$-conjugacy for sections.  Its coefficient
transport on cochains is the literal map
\[
 b(t_1,\ldots,t_n)\longmapsto
 a_{j,\epsilon}^{-1}
 b(a_{j,\epsilon}t_1a_{j,\epsilon}^{-1},\ldots,
   a_{j,\epsilon}t_na_{j,\epsilon}^{-1}).
 \tag{D.74.32}
\]
Canonical inner conjugation with this coefficient
transport acts identically on group cohomology.
For example, on the homogeneous bar resolution it
is induced by simultaneous right multiplication
of all vertices by $a_{j,\epsilon}^{-1}$.
The bar prism between that map and the identity
has terms
\[
 \sum_{\ell=0}^{n}
 (g_0,\ldots,g_\ell,
   g_\ell a_{j,\epsilon}^{-1},\ldots,
   g_n a_{j,\epsilon}^{-1});
 \tag{D.74.33}
\]
its boundary identity is the required cochain
homotopy after applying equivariant Hom.

Thus (D.74.29)--(D.74.32) identify $\bar r$ in
(D.74.26) with three duplicated restriction maps.
More explicitly, let
\[
 U_i=H^\bullet(N_E(S_i),M)[u^{-1}]_0,\qquad
 V_i=H^\bullet(S_i,M)[u^{-1}]_0,\qquad
 r_i:U_i\longrightarrow V_i
 \tag{D.74.34}
\]
be induced by the actual inclusion $S_i\subset N_E(S_i)$.
The triangle incidence and the coefficient transports
give
\[
 \mathcal A_M\simeq\bigoplus_iU_i,\qquad
 \mathcal B_M\simeq\bigoplus_i(V_i\oplus V_i),
 \qquad
 \bar r((w_i)_i)=((r_iw_i,r_iw_i))_i .
 \tag{D.74.35}
\]
There is no unidentified automorphism of a coefficient
module in this display; its two copies use exactly
(D.74.32).  In particular
\[
 \operatorname{rank}\bar r=\sum_i\operatorname{rank}r_i,
 \qquad
 \dim\mathcal B_M=2\sum_i\dim V_i .
 \tag{D.74.36}
\]
Together with (D.74.28) these equalities force every
$r_i$ to be surjective.

\emph{A finite permutation coefficient detects every nontrivial commuter.}
Fix $i$ and suppose $1\ne c\in C_i$.
Choose an open normal subgroup $N$ of $E$ with
$N\subset F$ and $c\notin N$.  This is possible
by separating $c$ in a finite continuous quotient
and then intersecting its kernel with $F$.
Write
\[
 J=E/N,\qquad Q=F/N,\qquad \bar S_i=S_iN/N,\qquad
                    M=k[J/\bar S_i].
 \tag{D.74.37}
\]
The quotient still has the specified map to $C_2$,
so $\bar S_i$ has order two,
$Q\cap\bar S_i=1$, and $J=Q\bar S_i$.
The basis pairing
$\langle\delta_a,\delta_b\rangle=1$ for $a=b$
and zero otherwise is a perfect $E$-invariant
self-duality of $M$.  Therefore all of
(D.74.25)--(D.74.36) apply to this exact coefficient
module.

For $n\geq1$, the cyclic calculation (D.74.18)
identifies $H^n(S_i,M)$ with the vector space
on those cosets in $J/\bar S_i$ fixed by $S_i$.
Indeed a two-point orbit gives the regular
$k[C_2]$-module, whose invariant vector is
also its norm image, so it contributes zero.
A singleton orbit contributes its basis vector
in every positive degree.  This identification
commutes with multiplication by $u$ and hence
also describes $V_i$.

The coset $\bar S_i$ is fixed by $S_i$.
The element $\bar c\in Q$ is nonidentity and
centralizes $\bar S_i$, so $\bar c\bar S_i$
is another fixed coset, different from
$\bar S_i$.  Consequently the nonzero class
of $\delta_{\bar S_i}$ in $V_i$ is moved by
$c$, to the different class of
$\delta_{\bar c\bar S_i}$.

On the other hand the image of the restriction
$r_i:U_i\to V_i$ is contained in $C_i$-invariants.
For $c\in C_i$, the cochain transport
$b(g_1,\ldots,g_n)\mapsto
c\,b(c^{-1}g_1c,\ldots,c^{-1}g_nc)$
is inner on $N_E(S_i)$ and acts identically on that
group's cohomology, by (D.74.33) with $a=c^{-1}$.
On $S_i$ its group arguments are unchanged, so it
acts there precisely by $c$ on $M$.
Naturality of restriction
therefore gives the asserted invariance
also after localization.
The class of $\delta_{\bar S_i}$ is not invariant,
contradicting the surjectivity of $r_i$ proved
in (D.74.36).  This proves $C_i=1$.  Equation
(D.74.5) then gives $N_E(S_i)=S_i$, as claimed.

\emph{The exact real-projectivity consequence.}
Every involution in $E$ is conjugate to one of the
$s_i$, so its full centralizer is its own order-two
subgroup.  Also $F$ is an actual open projective
subgroup of index two.  The set of involutions is
closed: it is $\{t\in E:t^2=1\}\cap(E\setminus F)$,
with $E\setminus F$ clopen.  These verify all the
hypotheses of Haran's characterization
\cite[Proposition~2.2]{C-HaranArtinSchreier1993}.
Thus $E$ is real projective.  In the formulation of
\cite[Definition~1.1 and Theorem~2.1]{R-PalQuick2025},
this has the following precise meaning.  For finite
$A,B$, a surjection $\pi:B\to A$, and a continuous
$\varphi:E\to A$, assume that each involution $t$
with $\varphi(t)\ne1$ has an involution $b_t\in B$
with $\pi(b_t)=\varphi(t)$.  There is then a
continuous homomorphism
\[
                 \gamma:E\longrightarrow B,
                         \qquad \pi\gamma=\varphi .
 \tag{D.74.38}
\]
This statement keeps the given finite quotient maps;
it does not require or conclude surjectivity of
$\gamma$.

\emph{Uniqueness of the actual coboundary.}
D.34 supplies $h\in F'$ with
$f=h^{-1}\alpha(h)$.  If $h'$ is any other
solution in $F$, then
\[
 \alpha(h'h^{-1})=h'f(hf)^{-1}=h'h^{-1}.
 \tag{D.74.39}
\]
The element on the right belongs to $C_F(s)=1$,
so $h'=h$.  For $\alpha\in GT_0$, D.44 supplies
a solution with $h\theta(h)=1$.  Uniqueness shows
that this is the same element, with no remaining
fixed-subgroup adjustment.

Keeping the exact original cusps and every unit
exponent, D.35's conjugator set is now parametrized
by the solutions $(\mu,u)$ of
\[
 \begin{gathered}
 \mu\in\widehat{\mathbb Z}^{\times},\qquad u\in F,\\
             x^\mu h y^\mu h^{-1}
                         =u(xy)^\mu u^{-1}.
 \end{gathered}
 \tag{D.74.40}
\]
Its associated maps are the actual maps
$x\mapsto x^\mu$, $y\mapsto hy^\mu h^{-1}$
of D.35.  The closed-subgroup test there proves
that every solution gives a full automorphism;
its converse gives every peripheral conjugator.
The proof above has removed the $C_F(s)$ variable
from that equation.  It has not established the
existence of a pair $(\mu,u)$ satisfying it.
\end{proof}

\paragraph{Scope and proof record.}
The finite coefficient used at the last step is the
actual coset permutation module (D.74.37); no
pro-$2$ quotient is substituted for $F$ or $E$.
The two transgressions in (D.74.22) are not assumed to
vanish.  Their equality of ranks is proved by the
specified cup pairing, and the relative-to-absolute
dimension comparison is then passed through the
actual relative long exact sequence.  The argument
uses D.34's three global section classes and its
labelled triangle, including the fact that each
class occurs twice.  The resulting full self-centralizing property and
real-projectivity are proved above.  Recognition and
the existence in (D.74.40) require their own subsequent maps.

\begin{claim}{D.75}{Proved}{The full negative torus cusp centralizer and unique transporters}
Let $b$ be any negative involution of the full
once-punctured-torus loop-outer group, with the actual
cusp-normalized extension of D.58--D.59:
\[
 \begin{gathered}
 F=\widehat F(x,y),\qquad c=[x,y],\qquad
 E=F\rtimes\langle v\rangle,\\
 v^2=1,\qquad vcv^{-1}=c^{-1},\qquad
 \alpha=\operatorname{Inn}(v)|_F .
 \end{gathered}                                             \tag{D.75.1}
\]
Then
\[
                    C_F(v)=1,\qquad N_E(\langle v\rangle)
                                      =\langle v\rangle .
                                                               \tag{D.75.2}
\]
Every cusp reflection $c^t v$, $t\in\widehat{\mathbb Z}$,
has the corresponding full trivial fixed subgroup in $F$.
In particular there is a unique $g\in F$ with
\[
                            gvg^{-1}=cv .                   \tag{D.75.3}
\]
Its homology is the primitive generator of the full fixed
line specified in D.59.

For every full unit $\mu\in\widehat{\mathbb Z}^{\times}$
and both literal reflection labels $v_i=c^i v$, $i=0,1$,
the entire transporter set is a singleton:
\[
 \begin{gathered}
 \{a\in F:[a,v_i]=c^\mu\}=\{a_{\mu,i}\},\\
 a_{\mu,0}=c^{(\mu-1)/2}g,\qquad
 a_{\mu,1}=c^{(\mu+1)/2}g^{-1}.
 \end{gathered}                                             \tag{D.75.4}
\]
If $\alpha_*\bmod2\ne I$, the extension has one section
class and is real projective.  If $\alpha_*\bmod2=I$,
the other section class retains the fixed group and
surjective homology map of D.59; no full procyclic
identification of that group is asserted here.

With the standard model selected by these section counts,
the complete D.61 peripheral-conjugator existence problem
is exactly
\[
 \text{there exist }\mu\in\widehat{\mathbb Z}^{\times},
 \ i\in\{0,1\}\text{ with }\quad
                     c^{-\mu}a_{\mu,i}\sim_F a_{\mu,i}.
                                                               \tag{D.75.5}
\]
The maps proving this equivalence are retained below.
Existence in (D.75.5), and recognition of an arbitrary
negative action by one of the standard models, remain
unproved.
\end{claim}
\passage{Section 30, folios 194-196, scans 383-387; Section 31, folios 198-199 verso, scans 391-394; Section 24, folio 147 verso, scan 293; Part D(b), full loop-outer finite actions in type (1,1)}

\begin{proof}
\emph{The one actual cusp module.}
Let $u:E\to C_2=\mathbb F_2$ be the quotient character,
$u(v)=1$, and set
$L=\overline{\langle c\rangle}$ and $D=N_E(L)$.
D.58 gives the actual inverse homomorphisms
\[
 \begin{aligned}
 \widehat{\mathbb Z}\rtimes_{-1}C_2
        &\longrightarrow D,
        &(t,\tau^\epsilon)&\longmapsto c^t v^\epsilon,\\
 D&\longrightarrow\widehat{\mathbb Z}\rtimes_{-1}C_2,
        &d&\longmapsto(t,u(d)),
        \qquad d v^{-u(d)}=c^t .
 \end{aligned}                                             \tag{D.75.6}
\]
Its two internal section classes have representatives
$S=\langle v\rangle$ and $S'=\langle cv\rangle$.
Their normalizers in $D$ are $S$ and $S'$:
an element $c^a$ commuting with a reflection has
$2a=0$, so $a=0$.  D.58 supplies an actual
$g\in F$ with $S'=gSg^{-1}$ and (D.75.3).
We will prove its uniqueness after proving (D.75.2).

Put $k=\mathbb F_2$, and form the actual compact module
\[
 0\longrightarrow\Delta
 \longrightarrow k[[E/D]]
            \xrightarrow{\ \varepsilon\ }k\longrightarrow0 .
                                                               \tag{D.75.7}
\]
The explicit restriction map
\[
            F/L\longrightarrow E/D,\qquad fL\longmapsto fD
                                                               \tag{D.75.8}
\]
is a homeomorphism: $E=FD$ and $F\cap D=L$
give bijectivity, and both spaces are compact Hausdorff.
In these coordinates the full $E$-action is
\[
                   e\cdot(fL)=efv^{-u(e)}L .              \tag{D.75.9}
\]
The representative is in $F$, and its coset is
independent of the representative of $fL$, because
$v$ normalizes $L$.  Thus $\Delta|_F$ is exactly
the augmentation module of the full completed
once-punctured-torus pair $(F,L)$, with the original
commutator cusp and the transported $E$-action.

For a finite continuous $E$-module $M$ over $k$, define
\[
 \mathscr H_M=\operatorname{Hom}_{k,\mathrm{cont}}(\Delta,M),
 \qquad (e\varphi)(d)=e\varphi(e^{-1}d).
                                                               \tag{D.75.10}
\]
This is the same actual discrete coefficient
construction as (D.74.9), now for (D.75.7).
The augmentation sequence, its continuous Hom sequence,
and Shapiro give
\[
 \begin{gathered}
 R^n(E,M)=H^n(E,D;M)=H^{n-1}(E,\mathscr H_M)
                       \quad(n\geq1),\\
 \cdots\longrightarrow H^n(E,M)\xrightarrow{\ r\ }H^n(D,M)
 \longrightarrow R^{n+1}(E,M)
 \longrightarrow H^{n+1}(E,M)\longrightarrow\cdots .
 \end{gathered}                                             \tag{D.75.11}
\]
All maps here are the actual coefficient maps and
subgroup restrictions.  No other section normalizer
has been introduced into this relative pair.

\emph{The exact duality and half-rank comparison.}
The pair $(F,L)$ in (D.75.8) is the full completion
of a compact oriented genus-one surface with its
one boundary.  Apply the
surface-pair duality of
\cite[Propositions~5.7 and~6.16 of the accepted version]
{C-WilkesRelativeAccepted2018} to this specified pair.
Equivalently the dense discrete subgroup and its
cyclic boundary give the absolute and relative
cohomology comparisons, with their actual cup products,
as in (D.74.15).  For
$A=H^0(F,M)$ and $B=H^1(F,M)$, they give
\[
 H^0(F,\mathscr H_{M^\vee})=B^\vee,\qquad
 H^1(F,\mathscr H_{M^\vee})=A^\vee .
                                                               \tag{D.75.12}
\]
The identifications are equivariant for the actual
$E/F$ action.  The coefficient evaluation map is
\[
 M\otimes\mathscr H_{M^\vee}\longrightarrow\mathscr H_k,
 \qquad m\otimes\varphi\longmapsto
                 [d\longmapsto\varphi(d)(m)] .
                                                               \tag{D.75.13}
\]
Its target relative orientation group
$H^1(F,\mathscr H_k)=H^2(F,L;k)=k$ has trivial action.
In boundary coordinates this also follows from
$v c v^{-1}=c^{-1}$, which fixes the mod-two
boundary character.  All higher fibre rows vanish
because $\operatorname{cd}_2(F)=1$.

The two ordinary Hochschild--Serre sequences for
$M$ and $\mathscr H_{M^\vee}$ therefore have exactly
the rows in (D.74.21), with these $A$ and $B$.
Their products are induced by (D.75.13).
The positive-column cyclic complex has differential
$1+\sigma$, with the perfect kernel/image pairing
and period-one $u$ map of (D.74.18)--(D.74.19).
The product of the two fibre-row-one classes lies
in $H^2(F,\mathscr H_k)=0$.  Thus the same explicit
$d_2$ derivation identity gives
\[
                    \langle db,a^\vee\rangle
                          =\langle b,d'a^\vee\rangle .
                                                               \tag{D.75.14}
\]
These transgressions have equal rank, without
either one being assumed zero.

For a finite self-dual $M$, the high-degree
absolute and relative dimensions are consequently
equal.  Multiplication by the actual $u$ is an
isomorphism in sufficiently high degree on all
terms of (D.75.11), by the same two-row calculation.
Taking degree-zero localizations, write
\[
 \mathcal A=H^\bullet(E,M)[u^{-1}]_0,\quad
 \mathcal B=H^\bullet(D,M)[u^{-1}]_0,\quad
 \mathcal R=R^\bullet(E,M)[u^{-1}]_0 .
                                                               \tag{D.75.15}
\]
They are finite-dimensional, and the original
long exact sequence becomes
$\mathcal A\to\mathcal B\to\mathcal R\to\mathcal A$,
with the connecting map divided by $u$.
Its dimensions satisfy $\dim\mathcal A=\dim\mathcal R$.
The rank identities (D.74.27), applied to these
actual three maps, give
\[
                     \operatorname{rank}\bar r
                                  =\tfrac12\dim\mathcal B .
                                                               \tag{D.75.16}
\]
This use of D.74's coefficient argument has
replaced its three-cusp augmentation module
by (D.75.7); it has not replaced the torus
peripheral structure or its involution matrix.

\emph{Both incidence transports and the other section factor.}
Let $C=C_F(v)$.  Before proving it trivial,
the exact normalizer is
$N_E(S)=C\times S$, via multiplication and the
inverse $e\mapsto(ev^{-u(e)},v^{u(e)})$.
In the identity-modulo-two case let $T$ be the
other full section class of D.58 and write
$N_E(T)=C_F(T)\times T$ with the same labelled
coordinate maps.  In the other case there is
no such factor.

The finite-coefficient Brown theorem applies to
$E$ and to $D$: their open subgroups $F$ and
$L$ have cohomological dimension one, and every
finite subgroup injects into $C_2$.  Their high
cohomology is finite by the two-row sequences
with finite coefficients.  Thus the actual
dense restriction embeddings of
\cite[Theorem~2 and Section~1.1]
{A-KochloukovaZalesskii2007} are isomorphisms.
Put
\[
 U=H^\bullet(N_E(S),M)[u^{-1}]_0,\qquad
 V=H^\bullet(S,M)[u^{-1}]_0 .
                                                               \tag{D.75.17}
\]
Brown identifies $\mathcal A$ with $U$ together
with the other normalizer factor when present.
It identifies $\mathcal B$ with the two section
factors for $S,S'$.

Transport the first local factor by the identity,
and the second by the actual $g$ of (D.75.3).
The second cochain transport is
\[
 b(t_1,\ldots,t_n)\longmapsto
             g^{-1}b(gt_1g^{-1},\ldots,gt_ng^{-1}).
                                                               \tag{D.75.18}
\]
As proved by the explicit bar prism (D.74.33),
inner conjugation with exactly this coefficient
transport acts identically on global cohomology.
Consequently, if $r_S:U\to V$ is literal
restriction, the boundary map is exactly
\[
        \bar r:U\oplus U_{\rm other}\longrightarrow V\oplus V,
                 \qquad (a,a_{\rm other})
                             \longmapsto(r_Sa,r_Sa).
                                                               \tag{D.75.19}
\]
Here $U_{\rm other}=0$ in the one-class case.
The additional global section factor has no
boundary occurrence and its map in this
display is zero.  Both local sections occur
in the cusp class, by D.58's actual $F$-fusion.
Thus (D.75.16) says
$\operatorname{rank}r_S=\dim V$, proving
surjectivity of $r_S$ for every finite
self-dual coefficient module.

Suppose $1\ne k_0\in C$.  Choose an open normal
$N\triangleleft E$, with $N\subset F$ and
$k_0\notin N$, and use the exact finite module
\[
              J=E/N,\qquad \bar S=SN/N,\qquad
                             M=k[J/\bar S].
                                                               \tag{D.75.20}
\]
Its basis pairing is invariant and perfect.
In positive $S$-cohomology the singleton
$S$-fixed cosets are a basis; two-point orbits
have equal invariant and norm subspaces and
contribute zero.  The cosets $\bar S$ and
$\bar k_0\bar S$ are distinct and both fixed.
They are distinct because $\bar k_0\in F/N$
is nonidentity and $(F/N)\cap\bar S=1$.
The class of $\delta_{\bar S}$ is therefore
moved by $k_0$.

Every class in the image of $r_S$ is
$C$-invariant: the cochain map
$b\mapsto k_0 b(k_0^{-1}(\,\cdot\,)k_0)$ is
inner on $N_E(S)$, whereas on $S$ it is exactly
the coefficient action of $k_0$.
This contradicts the proved surjectivity.
Therefore $C=1$, proving (D.75.2).
All cusp reflections are $F$-conjugate to $v$:
inside $D$ each is conjugate to $v$ or $cv$,
and these two fuse by $g$.  Their fixed
subgroups in $F$ are therefore all trivial.
Dividing any two solutions of (D.75.3) gives
an element of $C$, so that solution is unique.

\emph{Every unit and both literal cusp reflections.}
The two half-exponents in (D.75.4) belong
to $\widehat{\mathbb Z}$, because a full unit
is odd at two and division by two is defined
at every other prime.  The identities
\[
 c^a v c^{-a}=c^{2a}v,\qquad
                 gvg^{-1}=cv,\qquad g^{-1}(cv)g=v
                                                               \tag{D.75.21}
\]
show directly that
$a_{\mu,i}v_i a_{\mu,i}^{-1}=c^\mu v_i$.
Dividing two such transporters gives an
element of $C_F(v_i)=1$.  This proves
the full singleton assertion (D.75.4).
Since $c$ is a commutator, D.59's primitive
homology statement gives
\[
 [a_{\mu,i}]_{\rm ab}=(-1)^i[g]_{\rm ab},\qquad
          \alpha_i(a_{\mu,i})=c^{-\mu}a_{\mu,i},
      \quad \alpha_i=\operatorname{Inn}(v_i)|_F .
                                                               \tag{D.75.22}
\]
In particular D.62's full primitive-homology
centralizer theorem applies to every one of
these elements, without a free-basis
primitivity assumption.

For completeness retain the maps relating
(D.75.5) to the peripheral conjugator.
Fix $(\mu,i)$ and write $a=a_{\mu,i}$.
Every outside element is uniquely $hv_i$,
$h\in F$, and
\[
             (hv_i)a(hv_i)^{-1}
                         =h(c^{-\mu}a)h^{-1}.
                                                               \tag{D.75.23}
\]
Thus the required $F$-conjugacy is exactly
the existence of an outside centralizer
element.  Its kernel under $u$ is the full
group $\overline{\langle a\rangle}\cong
\widehat{\mathbb Z}$ by D.62.
Choose an outside centralizer element $t$;
then $t^2=a^\kappa$, and the complete
centralizer has the actual coordinates
$a^r t^\epsilon$, as in (D.62k)--(D.62l).

In the two-class case the D.60 character
$\chi:E\to\widehat{\mathbb Z}$ has
$\chi(a)=(-1)^i$.  Hence
$2\chi(t)=(-1)^i\kappa$ and $\kappa_2$ is even.
There is the unique outside commuting involution
\[
                         w=a^{-\kappa/2}t .
                                                               \tag{D.75.24}
\]
It belongs to the other class, because a
cusp-class involution has trivial fixed
subgroup in $F$ and $a\ne1$.

In the one-class case every outside involution
has trivial fixed subgroup in $F$, so no such
involution centralizes $a\ne1$.  Thus
$\kappa_2$ is odd.  The full root is obtained
by the exact primewise correction
\[
 \begin{gathered}
 b_2=0,\qquad b_p=(1-\kappa_p)/2\quad(p\ne2),\\
 \kappa'=\kappa+2b\in\widehat{\mathbb Z}^{\times},
 \qquad z=(a^b t)^{(\kappa')^{-1}},
 \qquad z^2=a .
 \end{gathered}                                             \tag{D.75.25}
\]
The centralizer is the full procyclic group
$\overline{\langle z\rangle}$; multiplication
by two in this group is injective, so this
root is unique.  These are the full-coordinate
maps of D.62, rather than a root only in a
maximal pro-two quotient.

The resulting extension maps from the
specified D.61 standard models are
\[
 \begin{array}{c|ccc}
       &X&Y&v_0\\ \hline
 \text{two classes}&a&v_iw&v_i\\
 \text{one class}&zv_i&v_i z&v_i .
 \end{array}
 \qquad [X,Y]\longmapsto c^\mu .
                                                               \tag{D.75.26}
\]
The commutator identity in the first row
uses $[a,w]=1$; in the second it uses
$z^2=a$ and (D.75.4).  Conjugation by
$v_i$ has exactly the standard D.61 action
on the two displayed basis images.
D.61's actual boundary test makes the kernel
map a full isomorphism.  The inverse
extension map is its inverse on $F$ and
sends the specified $v_i$ back to $v_0$.
Conversely every D.61 normalized
peripheral conjugator has a unit $\mu$
and one of these labels $i$.  Its image
of $X$ in the two-class model, or of $XY$
in the one-class model, is the unique
element in (D.75.4).  The image of its
other reflection or of $Xv_0$, respectively,
gives the outside centralizer element in
(D.75.23).  This proves (D.75.5) in both
directions with all cusp and quotient labels.

\emph{The one-class consequence and the other fixed group.}
If $\alpha_*\bmod2\ne I$, D.58 gives only
the cusp section class.  Every involution
is therefore self-centralizing in $E$.
The actual open subgroup $F$ is projective
of index two, and the involutions form
the closed subset
$\{t:t^2=1\}\cap(E\setminus F)$.
Haran's Proposition~2.2
\cite{C-HaranArtinSchreier1993} applies.
The resulting real-projectivity has
exactly the finite embedding-problem maps
spelled out in (D.74.38).

If $\alpha_*\bmod2=I$, D.58 gives the other
full section class $T=\langle w_1\rangle$.
D.59 retains the exact assertions
\[
 \begin{gathered}
 N_E(T)=C_F(w_1)\times T,\qquad
             \dim_k H^1(C_F(w_1),k)=1,\\
 C_F(w_1)\longrightarrow
       \ker(\alpha_*-I:H_1(F,\widehat{\mathbb Z})
                            \to H_1(F,\widehat{\mathbb Z}))
                  \quad\text{is onto}.
 \end{gathered}                                             \tag{D.75.27}
\]
This factor is precisely the unused factor
in (D.75.19).  The proof does not identify
it with a procyclic group.  In either case,
the remaining recognition question is the
exact conjugacy in (D.75.5); its existence
has not been inferred from the uniqueness
of the transporter.
\end{proof}

\begin{claim}{D.76}{Proved}{The full section augmentation module and its incidence map}
Retain the actual pants extension $E=F\rtimes_\alpha\langle s\rangle$
of D.74, with $F=\widehat F(x,y)$, and put
$k=\mathbb F_2$, $R=k[[E]]$, and $R_F=k[[F]]$.
Choose one actual subgroup $T_i=\langle t_i\rangle$ in each
of its three section classes, $0\leq i<3$.
The compact left $R$-module
\[
 \Delta_T=\ker\left(
       \bigoplus_{i=0}^2 k[[E/T_i]]
                    \xrightarrow{\ \varepsilon\ }k\right)
                                                               \tag{D.76.1}
\]
is projective and has a full continuous module isomorphism
\[
                         R^2\xrightarrow{\ \cong\ }\Delta_T .
                                                               \tag{D.76.2}
\]
The proof constructs this isomorphism from compatible
isomorphisms at every finite group-ring quotient.
For every finite continuous $E$-module $M$ over $k$,
\[
             \dim_k\operatorname{Hom}_E(\Delta_T,M)
                                      =2\dim_k M.             \tag{D.76.3}
\]

There is a separate, specified incidence map.
Let $e_i$ denote the identity coset in the $i$th summand
of (D.76.1), and put
$H=\overline{\langle T_0,T_1,T_2\rangle}$.
Then
\[
 \begin{aligned}
 d_T:R^2&\longrightarrow\Delta_T,
       &(a,b)&\longmapsto a(e_1-e_0)+b(e_2-e_0),\\
 \operatorname{coker}d_T
       &\xrightarrow{\ \cong\ }
       \ker\bigl(k[[E/H]]\xrightarrow{\ \varepsilon\ }k\bigr).
 \end{aligned}                                                 \tag{D.76.4}
\]
Under any isomorphism (D.76.2), $d_T$ is a square
module endomorphism. It is invertible exactly when
these three actual involutions generate $E$.
This criterion retains the chosen subgroup maps;
the existence of (D.76.2) does not assert invertibility
of the specified $d_T$.

In the normalized D.74 coordinates, four actual
involutions already generate $E$:
\[
                        s,\quad xs,\quad fs,\quad yfs.
                                                               \tag{D.76.5}
\]
The first and third belong to the same section class,
with the exact unique transporter
$fs=h^{-1}sh$ from D.74.
For the three actual involutions
\[
                   s,\qquad xs,\qquad hyh^{-1}s,             \tag{D.76.6}
\]
the subgroup $H\cap F$ in (D.76.4) is exactly
$\overline{\langle x,hyh^{-1}\rangle}$.
Generation by (D.76.6), and D.74's remaining third-cusp
equation, are not asserted by this module calculation.
\end{claim}
\passage{Section 30, folios 194-196, scans 383-387; Section 31, folios 198-199 verso, scans 391-394; Section 24, folio 147 verso, scan 293; Part D(b), the full loop-outer finite-action and actual peripheral-recognition problem for type (0,3)}


\begin{proof}
\emph{All actual coset and restriction maps.}
Every $T_i$ maps isomorphically to $E/F=C_2$, so
$E=FT_i$ and $F\cap T_i=1$. Hence
\[
        F\longrightarrow E/T_i,\qquad a\longmapsto aT_i
                                                               \tag{D.76.7}
\]
is an $F$-equivariant homeomorphism: it is a continuous
bijection between compact Hausdorff spaces.
It identifies the restriction of the $i$th permutation
module with the actual regular module $R_F$.
If $T_i=a_iS_i a_i^{-1}$ for a representative in D.74,
the full $E$-equivariant coset comparison is
$gT_i\mapsto ga_iS_i$, with inverse
$gS_i\mapsto ga_i^{-1}T_i$.
Thus all changes of the three representatives retain
explicit maps.

Write $I_F=\ker(R_F\to k)$.
In the coordinates (D.76.7), the mutually inverse
$R_F$-linear maps
\[
 \begin{aligned}
 I_F\oplus R_F\oplus R_F&\longrightarrow\Delta_T,
                &(a,b,c)&\longmapsto(a+b+c,b,c),\\
 \Delta_T&\longrightarrow I_F\oplus R_F\oplus R_F,
                &(a_0,a_1,a_2)&\longmapsto
                             (a_0+a_1+a_2,a_1,a_2)
 \end{aligned}                                                \tag{D.76.8}
\]
give the actual restricted module.
The augmentation ideal of the full free group is
free on $x-1,y-1$:
\[
       R_F^2\longrightarrow I_F,\qquad
                  (a,b)\longmapsto a(x-1)+b(y-1).
                                                               \tag{D.76.9}
\]
Here is a direct verification in the completed category.
For every finite $F$-module $M$, a continuous linear
map $I_F\to M$ is exactly a continuous crossed
homomorphism $F\to M$, by $g\mapsto g-1$ and
$d(gh)=d(g)+g\,d(h)$.
The images of $x,y$ of such a crossed homomorphism
are arbitrary elements of $M$: use the free profinite
universal property on the two lifts in the finite
semidirect product $M\rtimes\operatorname{im}(F\to\operatorname{Aut}M)$.
This also proves uniqueness.
Consequently (D.76.9) induces a bijection on continuous
Hom into every finite module. Applying those bijections
to the finite quotients of its source supplies compatible
inverse maps; finite quotients separate points.
Thus (D.76.9) is the asserted full continuous isomorphism.
In particular
\[
                         \Delta_T|_F\cong R_F^4.             \tag{D.76.10}
\]
It follows that $\Delta_T$ is finitely generated over
$R$ as well. No finiteness statement for a hypothetical
discrete group has been used.

\emph{A direct degree-one bridge for every coefficient module.}
Let $u:E\to C_2=k$ be the original quotient character,
and let $M$ be any finite continuous left $E$-module
over $k$. Give
$W_M=M\otimes_k k[E/F]$ the diagonal action.
The exact coefficient sequence is
\[
 \begin{gathered}
 0\longrightarrow M\xrightarrow{\ j\ }W_M
                    \xrightarrow{\ \epsilon\ }M
                       \longrightarrow0,\\
 j(m)=m\otimes(F+sF),\qquad
 \epsilon(m\otimes gF)=m.
 \end{gathered}                                               \tag{D.76.11}
\]
The actual induction comparison and its inverse are
\[
 \begin{aligned}
 R\widehat\otimes_{R_F}M&\longrightarrow W_M,
              &g\otimes m&\longmapsto gm\otimes gF,\\
 W_M&\longrightarrow R\widehat\otimes_{R_F}M,
              &m\otimes gF&\longmapsto g\otimes g^{-1}m .
 \end{aligned}                                                \tag{D.76.12}
\]
Changing $g$ to $gf$ with $f\in F$ leaves the second
formula unchanged by the tensor relation. Both composites
are the identity. Since the subgroup has finite index,
induction is also coinduction. Explicitly, $g\otimes m$
corresponds to the function on $E$ whose value at $a$ is
$(ag)m$ when $ag\in F$, and zero otherwise; the coinduced
action is right translation of the argument.
Shapiro therefore gives
$H^2(E,W_M)=H^2(F,M)=0$.

The connecting map from (D.76.11) is precisely
cup product by the named $u$. To check it without a
choice of quotient generator, lift a $1$-cocycle
$a:E\to M$ by $\widetilde a(g)=a(g)\otimes F$.
Its coboundary at $(g,h)$ is
\[
 j\bigl(u(g)\,g\,a(h)\bigr).
 \qquad\text{Thus}\qquad
 \partial[a]=[u\smile a].
                                                               \tag{D.76.13}
\]
The coefficient long exact sequence proves that
$u\smile-:H^1(E,M)\to H^2(E,M)$ is onto.

D.74 gives the actual equalities $N_E(T_i)=T_i$.
The finite-coefficient Brown theorem applies in degree
two: the open group $F$ has $\operatorname{cd}_2(F)=1$,
all finite subgroups inject into $E/F$, and the three
$T_i$ are exactly its full section classes.
The two-row Hochschild--Serre sequence with this finite
module makes $H^2(E,M)$ finite, so the dense embedding of
\cite[Theorem~2 and Section~1.1]{A-KochloukovaZalesskii2007}
is onto. Hence
\[
       H^2(E,M)\xrightarrow{\ \cong\ }
                            \bigoplus_{i=0}^2H^2(T_i,M).
                                                               \tag{D.76.14}
\]
On every actual $T_i$, $u$ is its nonzero quotient
character. The cyclic complex has differential $1+t_i$
in each positive degree. Cup product by $u$ identifies
$H^1(T_i,M)$ with $H^2(T_i,M)$.
Given three degree-one local classes, use (D.76.14)
on their cups with $u$, then use surjectivity in
(D.76.13). Restriction and cup commute, and local
periodicity gives exactly the original three classes.
Thus
\[
 H^1(E,M)\longrightarrow\bigoplus_iH^1(T_i,M)
       \text{ is onto},\qquad
 H^2(E,\{T_i\};M)=0.                                        \tag{D.76.15}
\]

\emph{Projectivity with compact module maps.}
For this specified augmentation kernel, the continuous
Hom construction and the relative long exact sequence
are the same ones used in D.74, now with $T_i$ in place
of the peripheral normalizers. The augmentation sequence
splits as compact $k$-vector spaces. Consequently
\[
 \operatorname{Ext}^1_R(\Delta_T,M)
   =H^1\bigl(E,\operatorname{Hom}_{k,\mathrm{cont}}(\Delta_T,M)\bigr)
   =H^2(E,\{T_i\};M)=0
                                                               \tag{D.76.16}
\]
for every finite $M$.
These are the compact/discrete Hom and relative
augmentation identifications of
\cite[Section~1 and Definition~2.1]{C-WilkesRelativeAccepted2018}.
They may also be read directly as follows. A compact
module extension of $\Delta_T$ by finite $M$ has a
continuous $k$-linear splitting; the difference of its
translates is a continuous cocycle with values in the
displayed Hom module. Vanishing of its class is
exactly an equivariant splitting.

For completeness, finite-coefficient vanishing here
does give projectivity in the compact module category.
Choose a surjection $p:R^m\to\Delta_T$, possible by
(D.76.10), and let $K=\ker p$.
For every open $R$-submodule $U$ of $R^m$ there is an
exact finite module sequence
\[
 0\longrightarrow (K+U)/U
 \longrightarrow R^m/U
 \longrightarrow\Delta_T/p(U)\longrightarrow0 .
                                                               \tag{D.76.17}
\]
Its pullback along $\Delta_T\to\Delta_T/p(U)$ splits
by (D.76.16). It gives a continuous $R$-linear map
$\ell_U:\Delta_T\to R^m/U$ lifting that quotient map.
The possible maps are a nonempty finite set, since
$\Delta_T$ is finitely generated and the target is finite.
Reduction for a finer $U$ carries such a lift to a
lift at every coarser $U$.
Any finitely many compatibility requirements are
realized at a common finer submodule. Compactness
of the product of these finite sets supplies compatible
$\ell_U$ for all $U$. Their inverse limit is an actual
continuous $R$-linear section of $p$.
Therefore $\Delta_T$ is a direct summand of $R^m$,
and is projective.

\emph{The full coefficient count and finite free quotients.}
Apply $\operatorname{Hom}_R(\Delta_T,-)$ to (D.76.11).
Projectivity gives an exact short sequence whose first
and last spaces are $\operatorname{Hom}_R(\Delta_T,M)$.
Frobenius reciprocity and (D.76.10) identify its middle
space with
$\operatorname{Hom}_{R_F}(R_F^4,M)$.
The reciprocity maps, using the coinduced functions in
(D.76.12), are evaluation at $1$ and its inverse
\[
        q\longmapsto
            \bigl[d\longmapsto(a\longmapsto q(ad))\bigr].
                                                               \tag{D.76.18}
\]
These are inverse continuous equivariant maps.
All three Hom spaces are finite. Their dimensions give
$2\dim\operatorname{Hom}_R(\Delta_T,M)=4\dim M$,
which is (D.76.3).

For an open normal subgroup $N\triangleleft E$, write
\[
 R_N=k[E/N],\qquad
 P_N=R_N\widehat\otimes_R\Delta_T
     =\Delta_T/\overline{\langle(n-1)d:n\in N,\ d\in\Delta_T\rangle}.
                                                               \tag{D.76.19}
\]
This is a finite projective left $R_N$-module:
base change of the direct-summand description above
is a direct summand of the finite free module $R_N^m$.
For every simple $R_N$-module $M$, (D.76.3) says
\[
       \dim_k\operatorname{Hom}_{R_N}(P_N,M)
         =2\dim_k M
         =\dim_k\operatorname{Hom}_{R_N}(R_N^2,M).
                                                               \tag{D.76.20}
\]
Let $J_N$ be the Jacobson radical of the finite algebra
$R_N$. Both Hom spaces factor through the respective
semisimple tops. For every simple $M$, its multiplicity
in such a top is the dimension of its Hom space divided
by $\dim_k\operatorname{End}_{R_N}(M)$.
Thus (D.76.20) identifies the two semisimple tops.
Lift that identification to a map $R_N^2\to P_N$.
Its cokernel has zero top, so is zero by Nakayama's
lemma, since $J_N$ is nilpotent.
Projectivity splits this surjection. The complementary
summand again has zero top and hence is zero.
This proves
$P_N\cong R_N^2$ as actual finite modules.

Let $\mathcal I_N=\operatorname{Iso}_{R_N}(R_N^2,P_N)$.
These are nonempty finite sets.
For $N'\subseteq N$, the canonical base-change identity
$R_N\otimes_{R_{N'}}P_{N'}=P_N$ sends an isomorphism
at $N'$ to one at $N$.
Every finite list of compatibility equations has a
solution obtained at a common finer open normal subgroup.
The finite intersection property therefore gives a
compatible family $(\phi_N)_N$.
The kernels in (D.76.19) form a neighbourhood basis
for $\Delta_T$: every finite module quotient has an
open normal subgroup acting trivially, and its kernel
contains the corresponding coinvariant kernel.
Consequently
\[
   \Delta_T=\varprojlim_NP_N,\qquad
   R^2=\varprojlim_NR_N^2,\qquad
   \phi=\varprojlim_N\phi_N
                                                               \tag{D.76.21}
\]
give the isomorphism (D.76.2), with inverse
$\varprojlim_N\phi_N^{-1}$.
This argument uses neither surjectivity of the maps
between the finite isomorphism sets nor an assumed
preferred basis.

\emph{The specified incidence map and its exact cokernel.}
Put $P_T=\bigoplus_i k[[E/T_i]]$.
The map
\[
        P_T\longrightarrow k[[E/H]],\qquad
                           gT_i\longmapsto gH
                                                               \tag{D.76.22}
\]
is well defined, continuous and equivariant.
It kills the image of $d_T$.
In $P_T/\operatorname{im}d_T$, the three identity
cosets are equal. Each $t_i$ fixes the common class,
because it fixes $e_i$ and $e_i=e_0$ there.
Continuity makes their closed generated group $H$
fix that class. Thus
$gH\mapsto[g e_0]$ is a well-defined continuous inverse
to the map induced by (D.76.22).
Checking the generators $gT_i$ and $gH$ proves both
compositions are the identity.
All images are closed, since the source modules are
compact. The augmentations commute with these maps.
Taking their kernels gives exactly the cokernel
isomorphism in (D.76.4).

The augmentation kernel of $k[[E/H]]$ vanishes exactly
when $E/H$ is a singleton. If two distinct cosets
exist, some finite clopen quotient separates them;
the difference of their two basis elements survives
in its finite permutation module.
Therefore $d_T$ is onto exactly when $H=E$.
Choose $\phi$ in (D.76.21) and set
$A_T=\phi^{-1}d_T:R^2\to R^2$.
If this map is onto, its reduction at every $N$
is an onto endomorphism of the finite module $R_N^2$,
hence is invertible. Taking the compatible inverse
limits makes $A_T$ invertible. The converse is immediate.
This proves the stated exact incidence criterion.
It supplies no reason to replace $d_T$ by the unrelated
isomorphism $\phi$.

\emph{The literal generators and the remaining removal target.}
The normalized action gives
$f\alpha(f)=1$ and
$\alpha(y)=f^{-1}y^{-1}f$.
Direct substitution shows that every element in
(D.76.5) squares to one.
They generate the original generators through
\[
            x=(xs)s,\qquad y=(yfs)(fs),\qquad s=s.
                                                               \tag{D.76.23}
\]
D.74's unique equation $h^{-1}\alpha(h)=f$ gives
$fs=h^{-1}sh$.
The other two involutions have respectively the
second and third section classes of D.34: replacing
$x^{e_2}s$ by $xs$, and $y^{e_2}fs$ by $yfs$, is
conjugation by $x^{(1-e_2)/2}$ and $y^{(1-e_2)/2}$.
These full half-exponents exist because $1-e_2$
has zero dyadic component.
Thus the four generators have exactly the duplicated
class stated in the claim.

Set $Y=hyh^{-1}$. The coboundary equation gives
$\alpha(Y)=Y^{-1}$, and
$h(yfs)h^{-1}=Ys$.
Hence the three elements in (D.76.6) represent the
three actual classes.
The subgroup $K=\overline{\langle x,Y\rangle}$ is
$\alpha$-invariant, and
\[
       \overline{\langle s,xs,Ys\rangle}
                          =K\langle s\rangle,\qquad
       K\cap\langle s\rangle=1,\qquad
       H\cap F=K.
                                                               \tag{D.76.24}
\]
These follow from $x=(xs)s$, $Y=(Ys)s$ and the
literal quotient label $u(s)=1$.
Equations (D.76.4) and (D.76.24) specify the exact
three-generator removal question.
Even its solution would require the original
third-cusp equation to give the peripheral conjugator
of D.35. Neither that equation nor a new group
presentation is inferred from the free module (D.76.2).
\end{proof}

\begin{claim}{D.77}{Proved}{A proper three-section subgroup with every primary image full}
There is an explicit choice of one representative of each
of the three section classes in the standard pants extension
for which the specified incidence map of D.76 is not onto,
although the chosen orientation-kernel generators have full
image in every maximal pro-$p$ quotient.

In the full profinite coproduct
\[
 E_0=C_2(a)\amalg C_2(b)\amalg C_2(c),\qquad
 F_0=\ker\bigl(E_0\longrightarrow C_2,\ a,b,c\longmapsto1\bigr),
                                                               \tag{D.77.1}
\]
put
\[
 x=ba,\quad y=ca,\quad
 \delta=\operatorname{Inn}(a)|_{F_0},\quad
 g=x[x,y]x^{-1},\quad
 Y=g\,y\,\delta(g)^{-1}.
                                                               \tag{D.77.2}
\]
Then $F_0=\widehat F(x,y)$, $\delta$ inverts $x,y$, and
\[
 \begin{gathered}
 T_0=\langle a\rangle,\qquad T_1=\langle b\rangle,\qquad
 T_2=\langle gcg^{-1}\rangle,\\
 H=\overline{\langle T_0,T_1,T_2\rangle}
       =K\rtimes\langle a\rangle,\qquad
 K=\overline{\langle x,Y\rangle}=H\cap F_0 .
 \end{gathered}                                                \tag{D.77.3}
\]
The pair $x,Y$ induces the identity on
$F_0^{\rm ab}=\widehat\Z^2$ and generates $F_0^{(p)}$
for every prime $p$, but $K$ is a proper subgroup of $F_0$
and $H$ is a proper subgroup of $E_0$.
The corresponding section augmentation module is the free
module of rank two in D.76. Its specified incidence cokernel
has a quotient of dimension five over $\mathbb F_2$.

The same finite quotient proving properness proves that
$Y$ is not conjugate to any full unit power of the original
$y$. Thus this example retains the three section labels
and all primary generation tests, and explicitly fails the
additional original-cusp alignment in D.76's distinguished
triple.
\end{claim}
\passage{Section 30, folios 194-196, scans 383-387; Section 31, folios 198-199 verso, scans 391-394; Section 24, folio 147 verso, scan 293; Part D(b), the full loop-outer finite-action and actual peripheral-recognition problem for type (0,3)}

\begin{proof}
\emph{The full extension and inverse maps.}
Let $F=\widehat F(x,y)$ and let $\delta$ be the actual
automorphism inverting its named generators.
The mutually inverse full profinite homomorphisms are
\[
 \begin{aligned}
 E_0&\longrightarrow F\rtimes_\delta C_2(a),
       &(a,b,c)&\longmapsto(a,xa,ya),\\
 F\rtimes_\delta C_2(a)&\longrightarrow E_0,
       &(x,y,a)&\longmapsto(ba,ca,a).
 \end{aligned}                                                \tag{D.77.4}
\]
The three images in the first row square to one.
The second row preserves
$a(ba)a=(ba)^{-1}$ and $a(ca)a=(ca)^{-1}$.
Both composites fix every indicated generator.
Their quotient characters to $C_2$ agree on those
generators, so their restriction identifies the actual
kernel in (D.77.1) with the full $F$, including its topology.

The subgroups $\langle a\rangle,\langle b\rangle,\langle c\rangle$
are the three section classes for this standard inversion
extension, as in D.34. Their images in the actual
abelian quotient
\[
 E_0\longrightarrow C_2^3,\qquad
       a\longmapsto e_0,\quad b\longmapsto e_1,\quad
       c\longmapsto e_2
                                                               \tag{D.77.5}
\]
are distinct and unchanged by conjugation.
Replacing $\langle c\rangle$ by its conjugate $T_2$
therefore retains exactly one representative of each class.

Since $c=ya$, direct multiplication gives
\[
 gcg^{-1}=gya g^{-1}
        =g y\delta(g)^{-1}a=Ya,\qquad
 \delta(Y)=Y^{-1}.
                                                               \tag{D.77.6}
\]
Thus $K$ is $\delta$-invariant. The identities
$x=ba$ and $Y=(gcg^{-1})a$ show that the group generated
by the three $T_i$ is exactly $K\langle a\rangle$.
The quotient label of $a$ is nonzero, whereas $K\leq F_0$;
hence the intersection and semidirect-product statements
in (D.77.3) follow.

\emph{The actual all-prime images.}
The word $g=x[x,y]x^{-1}$ lies in $F_0'$.
Taking full abelianizations in (D.77.2) gives
\[
                     x^{\rm ab}=(1,0),\qquad
                     Y^{\rm ab}=(0,1).
                                                               \tag{D.77.7}
\]
For each prime $p$, the continuous map
$x\mapsto x,\ y\mapsto Y$ induces an endomorphism
of the free pro-$p$ group $F_0^{(p)}$.
Its image maps onto the elementary abelian quotient
$F_0^{(p)}/\Phi(F_0^{(p)})$ by (D.77.7).
To verify full generation here, a proper closed image
would have proper image in some finite $p$-group quotient.
A maximal subgroup of that finite $p$-group has index $p$,
so a nonzero continuous character to $\mathbb F_p$
would vanish on both specified generators.
That contradicts (D.77.7).
The induced endomorphism is therefore onto.

It is also injective. For completeness, a finitely
generated pro-$p$ group has only finitely many
homomorphisms to the finitely many $p$-groups of any
bounded order. The intersections of their kernels
are open fully invariant subgroups and separate points.
An onto endomorphism induces a bijection on each
of these finite quotients. Its kernel is therefore
contained in their trivial intersection.
This proves the asserted automorphism on every
maximal pro-$p$ quotient by the actual specified maps.
Equation (D.77.6) also gives
$\delta\rho=\rho\delta$ for the full endomorphism
$\rho(x)=x,\ \rho(y)=Y$.

\emph{The finite quotient detecting both proper inclusions.}
Use permutations of $\{1,2,3,4,5\}$ with multiplication
$(uv)(i)=u(v(i))$. Define the full continuous quotient
$q:E_0\to A_5$ on the three order-two factors by
\[
 \begin{aligned}
 q(a)&=(1\,2)(3\,4),&
 q(b)&=(1\,3)(2\,5),&
 q(c)&=(1\,4)(2\,3).
 \end{aligned}                                                \tag{D.77.8}
\]
The coproduct universal property gives this homomorphism.
The exact finite multiplication and generation
calculation gives
\[
 \begin{gathered}
 q(x)=(1\,5\,2\,3\,4),\qquad
 q(y)=(1\,3)(2\,4),\qquad q(g)=(1\,3\,2),\\
 q(gcg^{-1})=q(a),\qquad q(Y)=1,\\
 \langle q(a),q(b),q(c)\rangle
       =\langle q(x),q(y)\rangle=A_5 .
 \end{gathered}                                                \tag{D.77.9}
\]
These equalities are also exhaustively verified by
the finite producer below: closure under multiplication
lists exactly sixty elements in each of the last two
generated groups.
In particular the restriction $q|_{F_0}$ is onto.

The images of the chosen subgroups are instead
\[
 q(H)=\langle q(a),q(b)\rangle=D_{10},\qquad
 q(K)=\langle q(x)\rangle=C_5 .
                                                               \tag{D.77.10}
\]
Indeed $q(x)=q(b)q(a)$ has order five and
$q(a)$ conjugates it to its inverse. The five rotations
and five reflections are distinct, proving the first
order in (D.77.10); the second is immediate.
Both are proper subgroups of the specified image $A_5$.
This proves $H\ne E_0$ and $K\ne F_0$ in the full
profinite groups.

\emph{The specified incidence cokernel and cusp comparison.}
Let $k=\mathbb F_2$, $R=k[[E_0]]$, and use the actual
$T_i$ from (D.77.3). D.76 applies to give
$\Delta_T\cong R^2$ and
\[
 \begin{aligned}
 d_T:R^2&\longrightarrow\Delta_T,\qquad
 (r_1,r_2)\longmapsto r_1(e_1-e_0)+r_2(e_2-e_0),\\
 \operatorname{coker}d_T
      &\cong\ker\bigl(k[[E_0/H]]\to k\bigr).
 \end{aligned}                                                \tag{D.77.11}
\]
There is a specified continuous surjection of coset spaces
\[
 E_0/H\longrightarrow A_5/D_{10},\qquad
                         zH\longmapsto q(z)D_{10}.
                                                               \tag{D.77.12}
\]
It is well defined by (D.77.10) and is onto by (D.77.9).
The induced permutation-module surjection respects
augmentation. Its restriction maps the kernel in
(D.77.11) onto
$\ker(k[A_5/D_{10}]\to k)$, which has dimension
$6-1=5$. This is an explicit nonzero finite quotient
of the actual incidence cokernel.

Finally $q(y)$ has order two, while $q(Y)=1$.
For every $\mu\in\widehat\Z^\times$, its residue modulo
two is one, so $q(y^\mu)=q(y)$ still has order two.
A conjugacy $Y=w y^\mu w^{-1}$ in $F_0$ would induce
a conjugacy between elements of orders one and two
under $q|_{F_0}$, which is impossible.
This proves the claimed failure of original-y alignment
through the same full homomorphism.
D.76's distinguished element $hyh^{-1}$ retains that
alignment by its defining equation; this example does
not solve its remaining generation or third-cusp problem.

The producer
\path{code/part_c/pants_incidence_limit.py}
executes twenty exact checks: the free-word expression
for $g$, its two zero abelian coordinates,
$\delta(Y)=Y^{-1}$, the identity abelianization
of $x,Y$, and all finite permutations, element orders
and generated image orders in (D.77.8)--(D.77.10).
The preceding Frattini argument supplies the
all-prime assertion; no finite list of primes is substituted.
\end{proof}

\begin{claim}{D.78}{Proved}{The other negative fixed subgroup and its actual pro-odd kernel}
Retain the two-section-class situation of D.75:
\[
 F=\widehat F(x,y),\quad c=[x,y],\quad
 E=F\rtimes\langle v\rangle,\quad vcv^{-1}=c^{-1},\quad
 gvg^{-1}=cv.
\]
The transporter $g$ is the unique original D.75 transporter. Put
$S=\langle v\rangle$, let $T=\langle w\rangle$ be an actual representative
of the other section class, and put
\[
 B=C_F(w),\qquad N_E(T)=B\times T,\qquad k=\mathbb F_2.
\]
D.59 gives $\dim_kH^1(B,k)=1$. Retain the actual D.60 quotient
$\chi:E\to\widehat{\mathbb Z}$, normalized by $\chi(g)=1$.
Its restriction to $B$ is onto. Choose $a\in B$ with
$[a]_{\mathrm{ab}}=[g]_{\mathrm{ab}}$, as in D.60, so $\chi(a)=1$.
Define the actual closed groups
\[
 A=\overline{\langle a\rangle},\qquad
 K=\ker(\chi|_B),\qquad
 O=\ker(B\xrightarrow{\chi}\widehat{\mathbb Z}\to\mathbb Z_2).
\]

Every open subgroup $B_0$ of the full $B$ satisfies
\[
                         \dim_k H^1(B_0,k)=1.                 \tag{D.78.1}
\]
The actual subgroup $A_2=\overline{\langle a^{e_2}\rangle}$ is a
Sylow pro-$2$ subgroup of $B$, and the restriction of $\chi_2$ to it
is an isomorphism onto $\mathbb Z_2$. Both $O$ and $K$ are pro-odd.
In particular the full original fixed group has the exact inverse coordinates
\[
 \begin{aligned}
 K\rtimes\widehat{\mathbb Z}&\longrightarrow B,
       &(k,t)&\longmapsto k a^t,\\
 B&\longrightarrow K\rtimes\widehat{\mathbb Z},
       &b&\longmapsto(ba^{-\chi(b)},\chi(b)).
 \end{aligned}                                                \tag{D.78.2}
\]
The action is conjugation by the specified $a^t$.
For every finite continuous $2$-primary $B$-module $Q$, the actual inflation map is
an isomorphism
\[
       H^n(\widehat{\mathbb Z},Q^K)
                       \xrightarrow{\ \sim\ }H^n(B,Q).       \tag{D.78.3}
\]
Under $A\xrightarrow{\chi}\widehat{\mathbb Z}$, restriction to $A$
is precisely the coefficient-inclusion map
\[
 H^n(\widehat{\mathbb Z},Q^K)
                \longrightarrow H^n(\widehat{\mathbb Z},Q).   \tag{D.78.4}
\]
The action on $Q$ in the right-hand group is explicitly
$t\cdot q=a^t q$; on $Q^K$ it is the quotient action of $B/K$.
These actions agree on the indicated submodule.
For every finite continuous $E$-set $X$, put $Q_X=k[X^T]$ with its
actual $B$-action. The proof supplies perfect finite-coefficient pairings
and the exact comparison $\mathcal U_T(k[X])\cong
H^0(B,Q_X)\oplus H^1(B,Q_X)$. The full recognition $B=A$ is equivalent
to surjectivity for every such $X$ of the specified inclusion
$Q_X^B\to Q_X^A$. Every nonidentity element of $K$ has the explicit
finite-coset coefficient witness given below.
Triviality of $K$ and that invariant-surjectivity assertion remain
unproved. The original cusp $c$, transporter $g$, and D.75.5 remain
unchanged.

\end{claim}
\passage{Section 30, folios 194-196, scans 383-387; Section 31, folios 198-199 verso, scans 391-394; Section 24, folio 147 verso, scan 293; Part D(b), full loop-outer finite actions in type (1,1)}

\begin{proof}
\emph{The actual perfect localized duality.}
Use the original cusp normalizer $D=N_E(\overline{\langle c\rangle})$
and the exact augmentation kernel and coefficient modules of D.75.7--13.
For a finite continuous $E$-module $M$, write
\[
 \begin{aligned}
 \mathcal A_M&=H^\bullet(E,M)[u^{-1}]_0,\\
 \mathcal B_M&=H^\bullet(D,M)[u^{-1}]_0,\\
 \mathcal R_M&=H^\bullet(E,D;M)[u^{-1}]_0,
 \end{aligned}
\]
where $u:E\to C_2=k$ is the original quotient character. The maps
\[
 \mathcal A_M\xrightarrow{r}\mathcal B_M
       \xrightarrow{u^{-1}\delta}\mathcal R_M
       \xrightarrow{j}\mathcal A_M                         \tag{D.78.5}
\]
are the actual localized restriction, connecting and relative-to-absolute
maps. They are linear over the localized coefficient ring $\mathcal A_k$.

The D.75.14 transgression calculation gives more than equality of ranks.
Put $U_M=H^0(F,M)$ and $V_M=H^1(F,M)$. In positive quotient columns the
absolute rows are $\overline U_M,\overline V_M$, whereas the shifted
relative rows for $M^\vee$ are
$\overline{V_M^\vee},\overline{U_M^\vee}$, with the kernel/image cyclic
pairings of D.74.19. The only possible differentials are transpose maps
\[
 d:\overline V_M\to\overline U_M,
 \qquad d':\overline{U_M^\vee}\to\overline{V_M^\vee}.
\]
Consequently the pairings on the two surviving pairs of rows are perfect:
\[
 \operatorname{coker}d\ \times\ \ker d'\longrightarrow k,
 \qquad
 \ker d\ \times\ \operatorname{coker}d'\longrightarrow k.   \tag{D.78.6}
\]
These are the actual pairings induced by evaluation and the fibre
surface-pair fundamental class, not chosen vector-space identifications.
For sufficiently large absolute degree $n$ and relative degree $m$, the
surviving filtration sequences are exactly
\[
 \begin{gathered}
 0\longrightarrow\operatorname{coker}d\longrightarrow H^n(E,M)
      \longrightarrow\ker d\longrightarrow0,\\
 0\longrightarrow\operatorname{coker}d'\longrightarrow H^m(E,D;M^\vee)
      \longrightarrow\ker d'\longrightarrow0.
 \end{gathered}
\]
The first subspace has quotient-column filtration $n$, and the second
has filtration $m-1$ in the shifted relative complex.

Here is the required trace on the target. For trivial coefficient $k$,
the absolute transgression is zero: the actual section $v$ makes
inflation from $C_2$ injective in every degree. Its transpose in the
shifted relative sequence is therefore zero as well. The top fibre row
$H^1(F,\mathscr H_k)=H^2(F,\langle c\rangle;k)=k$ survives.
For sufficiently large $n$ its quotient edge map supplies
\[
 H^n(E,D;k)\longrightarrow
 E_\infty^{n-2,1}=H^{n-2}(C_2,k)=k u^{n-2}\longrightarrow k,
                                                               \tag{D.78.7}
\]
where the last map reads the coefficient of the specified power of $u$.
These maps commute with a further multiplication by $u$. Thus they give
a specified trace $\operatorname{tr}:\mathcal R_k\to k$.

Cup product followed by this trace gives
\[
 \mathcal A_M\times\mathcal R_{M^\vee}\longrightarrow k.       \tag{D.78.8}
\]
Its associated two-row pairings are exactly (D.78.6), in opposite filtration
positions. The product of the two filtration subspaces just displayed has
quotient filtration at least $n+m-1$, whereas the trace reads column
$n+m-2$; that product therefore has trace zero. An absolute vector with
nonzero image in $\ker d$ pairs nontrivially with the relative subspace
$\operatorname{coker}d'$. A nonzero absolute vector in
$\operatorname{coker}d$ pairs nontrivially with a relative vector mapping
to $\ker d'$. These statements follow from the two perfect pairings in
(D.78.6), and the same argument applies with the sides reversed. This proves
nondegeneracy on the full finite filtered spaces, not only on their
associated graded spaces. Passing through the specified $u$ maps proves
that (D.78.8) is perfect.
The construction is natural in the coefficient maps and balanced over
$\mathcal A_k$, by the original cup-product and connecting-map identities.

\emph{The other Brown factor and its exact relative lift.}
D.75 proves $C_F(v)=1$. Brown restriction therefore gives the actual
coefficient decomposition
\[
 \mathcal A_M\xrightarrow{\ \sim\ }
 \mathcal V_S(M)\oplus\mathcal U_T(M),\qquad
 \begin{cases}
 \mathcal V_S(M)=H^\bullet(S,M)[u^{-1}]_0,\\
 \mathcal U_T(M)=H^\bullet(B\times T,M)[u^{-1}]_0.
 \end{cases}                                                  \tag{D.78.9}
\]
The two cusp factors are transported by $1$ and the original $g$, with
exactly the coefficient transport of D.75.18. The restriction in (D.78.5)
is consequently
\[
 \mathcal V_S(M)\oplus\mathcal U_T(M)
            \longrightarrow\mathcal V_S(M)\oplus\mathcal V_S(M),
              \qquad(z,t)\longmapsto(z,z).                    \tag{D.78.10}
\]
In particular the other factor has zero boundary restriction for every
finite module, including every finite permutation module.

For coefficient $k$, let $e\in\mathcal A_k$ be the Brown inverse of
$(0,1)$, the identity in the other normalizer factor and zero in the cusp
factor. It is a specified idempotent. Its restriction to both cusp factors
is zero, so $e\mathcal B_M=0$. Applying this idempotent to (D.78.5) gives
an actual isomorphism
\[
            j:e\mathcal R_M\xrightarrow{\ \sim\ }
                         e\mathcal A_M=\mathcal U_T(M).       \tag{D.78.11}
\]
Its inverse is unambiguous: lift an element through $j$, then multiply
its lift by $e$. Two lifts differ by a connecting image, which $e$ kills.

The pairing (D.78.8) splits orthogonally under $e$ and $1-e$; its restriction
to $e\mathcal A_M\times e\mathcal R_{M^\vee}$ is perfect. Transporting it
through (D.78.11) gives a perfect pairing
\[
 \mathcal U_T(M)\times\mathcal U_T(M^\vee)\longrightarrow k,
 \qquad(z,z')\longmapsto\tau\bigl(\operatorname{ev}_*(z\smile z')\bigr),
                                                               \tag{D.78.12}
\]
where $\tau=\operatorname{tr}\,j^{-1}:\mathcal U_T(k)\to k$ uses
(D.78.11). The formula follows because $j$ commutes with cup multiplication
by an absolute class; it therefore retains the actual normalizer cup
product and the actual coefficient evaluation.

\emph{Arbitrary finite permutation coefficients.}
Let $X$ be any finite continuous $E$-set and $M=k[X]$, with its invariant
basis pairing. The $T$-fixed points give a specified $B\times T$-submodule
\[
                         Q_X=k[X^T]\subseteq M.              \tag{D.78.13}
\]
Projection onto these basis elements is an inverse on this summand.
The complementary summand consists of the two-point $T$-orbits and is
free as a $k[T]$-module. It is $B$-stable because $B$ centralizes $T$.
Its positive $T$-cohomology vanishes. The extension spectral sequence
for $T\to B\times T\to B$, with $\operatorname{cd}_2B\leq1$, shows
that this complementary summand contributes no cohomology in degree
at least two. Thus inclusion of $Q_X$ induces the exact localized
coefficient comparison.

The action of $T$ on $Q_X$ is trivial. The direct-product resolution,
or the two-row complex with zero $T$ differential, gives the actual
isomorphism
\[
 \begin{aligned}
 H^0(B,Q_X)\oplus H^1(B,Q_X)&\xrightarrow{\ \sim\ }\mathcal U_T(M),\\
 (q,z)&\longmapsto q+u^{-1}\operatorname{inf}_B(z).
 \end{aligned}                                                \tag{D.78.14}
\]
At a positive degree $n\geq2$, this formula means
$q u^n+\operatorname{inf}_B(z)u^{n-1}$, followed by the coefficient
inclusion (D.78.13). Projection to the fixed-point summand and the direct
product cochain decomposition supply its inverse.
Restriction to the literal subgroup $T$ is exactly
\[
 H^0(B,Q_X)\oplus H^1(B,Q_X)\longrightarrow Q_X,
                         \qquad(q,z)\longmapsto q.           \tag{D.78.15}
\]
Thus its image is $Q_X^B$, with precisely that subgroup action retained.

The $H^1(B,Q_X)$ summand in (D.78.14) is totally isotropic for (D.78.12):
the product of two such summands is a coefficient image of
$H^2(B,Q_X\otimes Q_X)=0$. Since (D.78.12) is perfect on the finite
self-dual module $M$, it follows that
\[
                 \dim H^1(B,Q_X)\leq\dim H^0(B,Q_X).         \tag{D.78.16}
\]
Decompose the actual finite $B$-set $X^T$ into its orbits, with
stabilizers $B_1,\ldots,B_r$. The usual finite-coset Shapiro maps give
\[
 H^0(B,Q_X)=k^r,\qquad
 H^1(B,Q_X)=\bigoplus_{i=1}^rH^1(B_i,k).                       \tag{D.78.17}
\]
Every $B_i$ is open in $B$, and $\chi(B_i)$ has finite index in
$\widehat{\mathbb Z}$ because $\chi(B)=\widehat{\mathbb Z}$.
Write this image as $d_i\widehat{\mathbb Z}$, with $d_i$ a positive
integer. The actual map
\[
       B_i\xrightarrow{\chi}d_i\widehat{\mathbb Z}
        \xrightarrow{\div d_i}\widehat{\mathbb Z}
          \longrightarrow\mathbb F_2                        \tag{D.78.18}
\]
is a nonzero continuous character. Division here is the continuous
inverse of multiplication by $d_i$ onto its indicated image.
Thus each summand in (D.78.17) has dimension at least one. Together with
(D.78.16) this proves equality and dimension exactly one for each stabilizer.

Every open subgroup $B_0\leq B$ occurs as such a stabilizer. Choose an
open normal $N\triangleleft E$ contained in $F$ with
$N\cap B\subseteq B_0$. The group $U=B_0TN$ is open in $E$, since
$B_0T$ is a subgroup and $N$ is normal. The quotient label $u$ and
$N\subset F$ show $U\cap B=B_0$. In the actual finite $E$-set
$X=E/U$, the identity coset is $T$-fixed and its $B$-stabilizer is $B_0$.
This proves (D.78.1), with no passage to a maximal pro-$2$ quotient.

The two summands in (D.78.14) now have equal dimension. The mixed part
of the perfect pairing (D.78.12) is therefore perfect as well. Indeed
a vector of $H^1$ pairing trivially with $H^0$ pairs trivially with
both summands and must vanish, and the dimensions agree. Hence the
actual cup pairings, with the degree-correct trace
$\tau_1(z)=\tau(u^{-1}\operatorname{inf}_Bz)$, are
\[
 H^0(B,Q_X)\times H^1(B,Q_X^\vee)
     \longrightarrow H^1(B,k)\xrightarrow{\tau_1}k            \tag{D.78.19}
\]
and their reversed versions are perfect. The fixed-point basis pairing
identifies $Q_X^\vee$ with $Q_X$. On the one-dimensional $H^1(B,k)$,
$\tau_1$ is nonzero: otherwise that square-zero summand would lie in the
radical of (D.78.12) for $M=k$. Thus $\tau_1$ takes the original nonzero class
$\chi|_B\bmod2$ to one. This records the original finite-coefficient duality, including the
actual stabilizers and coefficient actions.

\emph{The full pro-odd kernel consequence.}
The map $t\mapsto a^t$ is an actual continuous section of $\chi|_B$.
In particular $A_2=\overline{\langle a^{e_2}\rangle}$ is a copy of
$\mathbb Z_2$ mapping isomorphically onto the dyadic quotient. Choose a
Sylow pro-$2$ subgroup $P$ of $B$ containing $A_2$.
There cannot be two linearly independent continuous characters
$P\to k$. Their joint map factors through $P/(P\cap N_0)$ for some
open normal $N_0\triangleleft B$. It therefore extends to the open
subgroup $PN_0$ by the quotient isomorphism
$PN_0/N_0=P/(P\cap N_0)$, contradicting (D.78.1).

It follows that $P$ is procyclic. One elementary verification is that
every nontrivial finite $2$-group quotient has at most one independent
character to $C_2$, hence a unique maximal subgroup; any element outside
that maximal subgroup generates the finite group. Thus every finite
quotient is cyclic, and compatible generators give a procyclic inverse
limit. Since $P$ maps onto $\mathbb Z_2$, it is $\mathbb Z_2$, and that
surjection is an isomorphism. Its subgroup $A_2$ already maps onto the
same quotient, so $P=A_2$.

Every Sylow pro-$2$ subgroup of $B$ is conjugate to $A_2$.
The homomorphism $\chi_2$ is invariant under conjugation and is injective
on all these subgroups. Hence $O=\ker\chi_2$ has trivial Sylow pro-$2$
subgroup and all its finite quotients have odd order. It is pro-odd.
The actual subgroup $K=\ker(\chi|_B)$ lies in $O$, so is pro-odd too.
The multiplication and inverse formulas (D.78.2) follow directly from
the chosen section $a^t$ and the original character $\chi$.

\emph{The surviving coefficient map and the original scope.}
For any finite continuous $2$-primary $B$-module $Q$, average over the
finite odd image of $K$ on $Q$. Its odd order is invertible on $Q$. The averaging map
\[
       e_K:Q\longrightarrow Q^K,
       \qquad q\longmapsto\frac1{|\overline K|}
                               \sum_{h\in\overline K}hq       \tag{D.78.20}
\]
is independent of the chosen finite image through which the action
factors. It commutes with $B$, since $K$ is normal. It makes the
invariants functor exact; higher $2$-primary cohomology of $K$ vanishes.
The extension spectral sequence for the actual sequence
$K\to B\xrightarrow{\chi}\widehat{\mathbb Z}$ therefore gives (D.78.3).
Restriction to the specified section $A$ commutes with inflation and
is induced by the inclusion $Q^K\hookrightarrow Q$, proving (D.78.4).
Evaluation, coefficient maps and cup products commute with this comparison.

For $Q=Q_X=k[X^T]$, the invariant space is concretely
$k[K\backslash X^T]$, mapped to $Q_X$ by its orbit sums. Every such orbit
has odd size, so the original basis pairing restricts to the ordinary
basis pairing on these sums. The action of the quotient is the action
of the specified powers of $a$ on the $K$-orbits. This is the actual
finite-coefficient comparison; it retains the action of $K$ through
its orbit partition.

There is an exact finite-permutation detection of the surviving kernel.
The equality $B=A$ is equivalent to surjectivity, for every finite
continuous $E$-set $X$, of the actual restriction map
\[
       Q_X^B\longrightarrow Q_X^A.                            \tag{D.78.21}
\]
The forward implication follows from equality of the subgroups.
For the reverse implication, take a nonidentity $k_0\in K$.
The compact subgroup $AT$ is closed, and $k_0\notin AT$:
its $u$-label is zero, so membership would put it in $A$, where $\chi$
is injective. Choose an open normal $N\triangleleft E$ contained in
$F$ such that $k_0\notin ATN$. For the finite set $X=E/(ATN)$,
the identity coset is fixed by both $A$ and $T$, whereas $k_0$ moves it.
The corresponding basis vector in $Q_X$ lies in $Q_X^A$ and not in
$Q_X^B$. Thus (D.78.21) fails for that actual finite coefficient.

The cusp restriction (D.78.10) vanishes on the entire other factor.
The perfect duality proves (D.78.1)--(D.78.3), but has supplied no
surjectivity for (D.78.21). This is the retained coefficient map that
still detects the full pro-odd kernel. The coefficient comparison retains the actual
$K$-orbit partition and supplies no replacement of $Q_X^K$ by $Q_X$. The original cusp $c$, unique transporter $g$, and
D.75 conjugacy equation are unchanged throughout.

\end{proof}

\section{E. Special groups and two operator groups}

\subsection{E(a). Special groups from conjugators and from framed diagrams}
\label{sec:ea}
We give the constructions of \S25, folios 149--154, scans 296--306,
with their precise hypotheses. Let $\pi$ be a discrete hyperbolic
surface group or its full profinite completion, and let $I$ be its finite
peripheral index set. Choose representatives $L_i$ and coherent
parametrizations $c_i:T\to L_i$, where $T=\Z$ or $\widehat\Z$.
The self-centralizing and self-normalizing properties needed below
are the results of Part A. Put
$\mathcal E=\Aut_{\rm lac}(\pi)$ and $\Gamma=\Out_{\rm lac}(\pi)$.
Every $\alpha\in\mathcal E$ has a permutation $\sigma_\alpha$ of $I$
and an automorphism $\chi_\alpha$ of $T$.

\begin{claim}{E.1}{Proved}{The conjugator construction and the product kernel}
Define $S\mathcal E$ to consist of pairs $(\alpha,(g_i)_{i\in I})$
with $\alpha\in\mathcal E$, $g_i\in\pi$, and
\[
 \alpha(c_i(t))=g_i^{-1}c_{\sigma_\alpha(i)}(\chi_\alpha t)g_i
 \quad(t\in T).
\]
The multiplication is
\[
 (\beta,(h_i))(\alpha,(g_i))
 =\bigl(\beta\alpha,
       (h_{\sigma_\alpha(i)}\,\beta(g_i))_{i\in I}\bigr).
\]
There are exact sequences and a canonically embedded normal subgroup
\[
 1\to T^I\to S\mathcal E\to\mathcal E\to1,
 \qquad
 \iota:\pi\hookrightarrow S\mathcal E,
 \quad a\mapsto(\operatorname{Int}(a),(a^{-1})_i).
\]
The groups $T^I$ and $\iota(\pi)$ commute and intersect trivially.
Consequently
\[
 S\Gamma=S\mathcal E/\iota(\pi),\qquad
 1\to T^I\times\pi\to S\mathcal E\to\Gamma\to1,
 \qquad 1\to T^I\to S\Gamma\to\Gamma\to1.
\]
The superscript $!$ takes the inverse image of the kernel of the
permutation action on $I$. All maps in the profinite case are continuous
maps of profinite groups, and all kernels and quotients are closed.
\end{claim}
\passage{\S{}25; folios 149--151; scans 296--300}

\begin{proof}
Apply $\beta$ to the defining equation for $\alpha$. The two
conjugating factors multiply as
$\beta(g_i)^{-1}h_{\sigma_\alpha(i)}^{-1}$, which is the inverse of the
factor displayed in the product. Associativity follows either by
substitution in that formula or by associativity of composition of the
specified conjugations. The identity is $(\id,(1)_i)$, and the inverse
has automorphism $\alpha^{-1}$ and $j$th conjugator
$\alpha^{-1}(g_{\sigma_\alpha^{-1}(j)}^{-1})$.

Surjectivity to $\mathcal E$ is the definition of preserving the loop
classes. In its kernel every $g_i$ centralizes $L_i$, hence belongs to
$L_i$; this identifies the kernel with $T^I$. Substituting
$(\operatorname{Int}(a),(a^{-1})_i)$ proves that $\iota$ is a
homomorphism and that conjugation by $(\alpha,(g_i))$ takes $\iota(a)$
to $\iota(\alpha(a))$. Its injectivity follows from $Z(\pi)=1$, including
when $I$ is empty. To check commutation with the kernel, the two products
with $(\id,(\ell_i))$, $\ell_i\in L_i$, both have conjugators
$\ell_i a^{-1}$. Their intersection is trivial because projection to
$\mathcal E$ sends $\iota(a)$ to $\operatorname{Int}(a)$. Finally an
element lying over the identity in $\Gamma$ has an inner automorphism
component and can be multiplied by a unique $\iota(a)^{-1}$ to put that
component equal to the identity. This proves that the whole kernel is
the displayed direct product, not merely an abstract extension of its
two factors.

For a finitely generated profinite $\pi$, its automorphism group has the
usual profinite topology obtained from characteristic open quotients.
The defining equations of $S\mathcal E$ are closed equations in
$\mathcal E\times\pi^I$. Thus it is profinite, the multiplication and
all maps above are continuous, and the compact subgroups giving the
quotients are closed. The same algebraic proof gives the exact
continuous sequences.
\end{proof}

\begin{claim}{E.2}{Proved}{Agreement with autoequivalences of the framed diagram}
The group $S\Gamma$ of E.1 is the group of isomorphism classes of
autoequivalences of
\[
 B_I\longrightarrow\coprod_{i\in I}BL_i\longrightarrow B\pi
\]
with the common orientation action retained. Adding a fixed base object
in $B\pi$ gives $S\mathcal E$. Thus the two constructions on scans
296--300 agree, with the comparison maps and inner ambiguities exactly
as in E.1.
\end{claim}
\passage{\S{}25; folios 150--151; scans 298--300}

\begin{proof}
An autoequivalence gives a permutation of the selected objects, the
specified automorphism $\chi$ on their cyclic automorphism groups,
and an automorphism $\alpha$ at the generic base object. The natural
isomorphism making the second square commute is, at its $i$th object,
an element $g_i\in\pi$. Its naturality equation is exactly the defining
equation of E.1. Pasting the squares gives the product there, including
the index $\sigma_\alpha(i)$.

Initially the first square also contains an isomorphism at each selected
peripheral object. A natural transformation at that object changes this
isomorphism by the corresponding element of $L_i$. Use it to make the
first-square isomorphism the identity. After this gauge choice there
is no remaining peripheral-object natural-transformation ambiguity.
A natural transformation at the generic object is an element $a\in\pi$
and changes the representative by $\iota(a)$, with the convention of
E.1. These are precisely the equivalences being divided out in
$S\Gamma$. Fixing the generic base object and its comparison
isomorphism removes that last ambiguity and gives $S\mathcal E$.
Conversely every pair in E.1 defines the three functors and their two
commutativity squares. The two constructions undo each other on the
chosen representatives and on the specified natural transformations,
so the group isomorphisms are canonical relative to the initial
peripheral representatives. Changing those representatives transports
the diagrams by their explicit conjugating isomorphisms.
\end{proof}

\begin{claim}{E.3}{Proved}{The puncture-stabilizer pushout is the normalizer group}
Let $\Gamma_i\subset\Gamma$ stabilize $i$. Restrict the extension
$S\Gamma\to\Gamma$ to $\Gamma_i$ and push its kernel out by
$\operatorname{pr}_i:T^I\to T$. The resulting extension is canonically
\[
 1\to L_i\xrightarrow{\operatorname{Int}}
 \{\alpha\in\Aut_{\rm lac}(\pi):\alpha(L_i)=L_i\}
 \longrightarrow\Gamma_i\to1.
\]
\end{claim}
\passage{\S{}25; folios 153--154; scans 305--306}

\begin{proof}
For a representative $(\alpha,(g_j))$ with $\sigma_\alpha(i)=i$, put
\[
 F_i[(\alpha,(g_j))]=\operatorname{Int}(g_i)\alpha.
\]
The defining equation shows that this automorphism normalizes the
actual subgroup $L_i$. Replacing the pair by its generic-basepoint
change replaces $\alpha$ by $\operatorname{Int}(a)\alpha$ and $g_i$ by
$g_i a^{-1}$; the displayed value is unchanged. Hence the map descends
to $S\Gamma_i$. The multiplication formula gives
\[
 \operatorname{Int}(h_i\beta(g_i))\beta\alpha
 =\operatorname{Int}(h_i)\beta\operatorname{Int}(g_i)\alpha,
\]
so it is a homomorphism. An automorphism normalizing $L_i$ has a
representative with $g_i=1$, proving surjectivity. On the kernel $T^I$
the map is $\ell_i\mapsto\operatorname{Int}(\ell_i)$; its kernel is
exactly the other factors because $Z(\pi)=1$. This is the universal
pushout claimed. The normalizer equality in Part A gives the last
kernel $L_i$ in the displayed sequence. This is the compatibility of
scans 305--306, with its actual group of automorphisms rather than an
unspecified outer class.
\end{proof}

\begin{claim}{E.4}{Refuted}{The product-kernel claim requires the rigidity hypotheses}
The formula $\ker(S\mathcal E\to\Gamma)=T^I\times\pi$ does not hold
unchanged for all abelian types. For the unpunctured torus,
$I=\varnothing$, $S\mathcal E=\operatorname{GL}_2(\Z)=\Gamma$ and the
kernel is trivial, whereas $\pi=\Z^2$. For the annulus the inner
basepoint subgroup intersects the peripheral kernel, and the special
outer kernel is a quotient of $T^2$ by one diagonal copy of $T$.
\end{claim}
\passage{\S{}25; folios 149--153; scans 296--304}

\begin{proof}
For the torus, all inner automorphisms of $\Z^2$ are the identity and
there are no conjugators indexed by punctures. Thus the pair
construction contains only the automorphism itself, giving the stated
identity map and its trivial kernel.

For the annulus, both peripheral images are the whole group
$\pi=\Z$. The conjugator kernel is $\pi^2$, but the generic-basepoint
change $a$ is $(\id,(a^{-1},a^{-1}))$ and lies in that same kernel.
Dividing by these changes gives $\Z^2/\Delta\Z\cong\Z$, with the
antidiagonal convention after orienting both boundary loops. Keeping
the two end labels and the independent orientation data supplies the
appropriate extended quotient; their omission would create an
additional error. For the plane the peripheral homomorphism
$\Z\to1$ is not injective, so the normalizer calculation used in E.1
is unavailable. Its correct fixed-boundary positive special group is
trivial, as C.20 computes. These explicit maps identify the failed
hypotheses in each case.
\end{proof}

\subsection{E(b). Two commuting operator groups}

We retain the names $\theta$ and $\omega$ for the standard order-two and order-three lifts on $\widehat F_2$; Grothendieck's $\sigma_0$ and $\rho$ denote the corresponding choices on scans 316--317, \S25 bis, folios 159--160. A change of base path conjugates a lift by $\widehat F_2$. We keep this ambiguity in the statements below.

\begin{claim}{E.20}{Proved}{The two lifting conditions cut out exactly the multiplier-$1$ subgroup modulo $3$}
In $B=\Out_{\rm lac}(\widehat\pi_{0,3})$ let $S_3$ be the geometric positive subgroup. Its centralizer is $GT_0$, with the normalized description $(\lambda,f)$ of Harbater--Schneps. Every element of $GT_0$ has an automorphism representative commuting with the chosen order-two lift $\theta$. The elements with an automorphism representative commuting with the chosen order-three lift $\omega$ are exactly
\[
 GT_0^{(3)}=\{(\lambda,f)\in GT_0:\lambda\equiv1\pmod3\}.
\]
This subgroup has index two. The two representatives need not be the same representative; this is the two-condition formulation on scans 316--317.
\end{claim}
\passage{25 bis; 157 verso--160; 312--317}

\begin{proof}
Harbater--Schneps, Main Theorem (a) and Theorem A(i) \cite{R-HS}, identify the centralizer and state precisely
\[
 [\theta,F]=\Inn(f),\qquad
 [\omega,F]=\Inn(y^mf),\qquad m=(\lambda-1)/2,
\]
for $F(x)=x^\lambda$, $F(y)=f^{-1}y^\lambda f$. Here $[a,b]=aba^{-1}b^{-1}$. A representative $\Inn(h)F$ commutes with $\theta$ precisely when
\[
 \theta(h)f=h,
\]
since $F_2$ has trivial centre. It commutes with $\omega$ precisely when $\omega(h)y^mf=h$. These are the two exact coboundary equations. We determine their solvability without imposing the pentagon equation.

The positive lift extension of the $S_3$ action on the three-punctured sphere is
\[
 E=C_2*C_3=\langle s,r\mid s^2=r^3=1\rangle,
 \qquad E\longrightarrow S_3,
\]
whose kernel is the peripheral free group $F_2$. Its cusp generator can be chosen as $e=sr$; $e^2$ is a primitive cusp in that kernel. Complete this exact sequence by D.1. An outer centralizer element, represented by $F$, acts on $\widehat E$ by conjugating its image in $\Aut(\widehat F_2)$ and induces the identity on $S_3$. It sends $e^2$ to a $\widehat F_2$-conjugate of $e^{2\lambda}$. After that conjugation it sends $e$ into $C_{\widehat E}(e^{2\lambda})=\overline{\langle e\rangle}$. To prove this equality, put $L=\overline{\langle e^2\rangle}$. A centralizing element preserves its $\widehat F_2$-conjugacy class, so its image in $S_3$ stabilizes the associated puncture. That stabilizer is $\langle q(e)\rangle$ of order two. Multiplication by $e$ when necessary puts the element in $\widehat F_2$, where Part A gives $C_{\widehat F_2}(e^{2\lambda})=L$. This proves the equality and identifies the whole centralizer with $\widehat\Z$. Doubling is injective in that group, so the image of $e$ is $e^\lambda$.

Consequently this automorphism acts on $\widehat E_{\rm ab}=C_6$ by multiplication by $\lambda$. Every order-two element is conjugate to $s$, and every order-three subgroup is conjugate to $\langle r\rangle$, by D.2 or the free profinite product theorem. A lift of the given transposition is $F_2$-conjugate to $s$: a total-group conjugator projects to $C_{S_3}(s)=\langle s\rangle$, and multiplication by $s$ adjusts it into the kernel. Thus there is one kernel-conjugacy class for the order-two lift. It is fixed by every centralizer element, proving solvability of the first equation.

There are exactly two kernel-conjugacy classes of lifts of the specified three-cycle. They are represented by $r$ and $s r^{-1}s^{-1}$. Indeed a total-group conjugator has either an even or an odd image in $S_3$; adjusting by $\langle r\rangle$ or by $s\langle r\rangle$ gives these two representatives. They are distinct because their images in the order-three quotient of $E_{\rm ab}$ are inverse nonidentity elements. The action of $(\lambda,f)$ preserves the class of $r$ exactly when $\lambda\equiv1\pmod3$, by the computed action on $C_6$. When it does, the same adjustment by $r$ puts its conjugator in the kernel. This is precisely solvability of the second coboundary equation. The argument treats elements, so it retains the label of the three-cycle.

For comparison, Lochak--Schneps, Proposition 3 and Theorem 2 \cite{R-LS}, compute the same two nonabelian cohomology sets: one class for $\theta$, two for $\omega$, with the second class selected by $\lambda\equiv-1\pmod3$. The first two parts of their proof use only the two-cycle and three-cycle equations, so their computation also applies to $GT_0$. Finally $(-1,1)\in GT_0$ exchanges the two classes. The homomorphism to $(\Z/3)^\times$ is therefore surjective and the asserted index is exactly two. The residue calculation is independently recorded in the executed program.
\end{proof}

\begin{claim}{E.21}{Proved}{The arithmetic subgroup and the two fixed points}
The intersection of the usual Galois image $G_{\Q}\to GT_0$ with $GT_0^{(3)}$ is the image of $G_{\Q(\sqrt{-3})}$. The order-three fixed points in $U=\mathbf P^1-\{0,1,\infty\}$ are $e^{\pm i\pi/3}$, and the order-two lift has one fixed point in $U$. Thus the requirement to fix the chosen order-three section is a proper finite-index restriction, even on the Galois subgroup.
\end{claim}
\passage{25 bis; 157 verso--160; 312--317}

\begin{proof}
The multiplier of the Galois action is the cyclotomic character, by its action on the tame inertia at $0$. Its reduction modulo $3$ is the action on $\mu_3$ and has fixed field $\Q(\zeta_3)=\Q(\sqrt{-3})$. E.20 identifies precisely its kernel. The three-cycle $z\mapsto1/(1-z)$ has fixed-point equation $z^2-z+1=0$, and the transposition $z\mapsto1-z$ has fixed-point equation $z=1/2$. These points lie in $U$. Complex conjugation exchanges the former pair and fixes the latter point, which also proves properness directly.

On scan 312 the relevant subgroup is obtained by fixing a finite set of critical data; it is therefore of finite index. An assertion that every element of $G_\Q$ lifts commuting with a specified one of the order-three lifts would contradict the displayed computation. We record the scan's wording as editorial evidence, not as a settled reading.
\end{proof}

\begin{claim}{E.22}{Proved}{Normalizer, stabilizer, and the forced lift}
Let $1\to\pi\to E\xrightarrow q\mathcal G\to1$ have centre-free kernel, and let $K\leq E$ be a finite subgroup on which $q$ is injective. Assume $C_\pi(K)=1$, as holds for a nontrivial rotation stabilizer in the anabelian geometric situation and its completion. Put $G_Q=q(K)$ and let $Q$ denote the $\pi$-conjugacy class of $K$. The restriction
\[
 q:N_E(K)\longrightarrow\mathcal G
\]
is injective and has image the stabilizer $\mathcal G_Q$ of $Q$. Its inverse on that image is a canonical homomorphism. In a two-group situation $\mathcal G=\Gamma G$ with $\Gamma$ centralizing $G$, its restriction gives the forced lift of $\Gamma_Q=\Gamma\cap\mathcal G_Q$, and this lift centralizes $K$.
\end{claim}
\passage{25 bis; 155--157 verso; 307--312}

\begin{proof}
An element of $N_E(K)\cap\pi$ acts trivially on $q(K)$, and the restriction of $q$ to $K$ is injective. It therefore centralizes $K$, so the kernel is $C_\pi(K)=1$. A quotient element stabilizes $Q$ precisely when one, and hence every, lift carries $K$ to a $\pi$-conjugate. Multiplying that lift by the appropriate element of $\pi$ produces an element of $N_E(K)$ above it. This proves the image statement and uniqueness of the lift. The inverse is a homomorphism because it is the inverse of an injective homomorphism onto its image. For $\gamma\in\Gamma_Q$, conjugation of its lift on $K$ projects to conjugation by $\gamma$ on $G_Q$, which is trivial. Injectivity on $K$ then makes this conjugation trivial. In the profinite case all these maps are continuous, and the inverse is continuous by compactness of the closed normalizer and the Hausdorff property of the quotient.

For the centralizer, the precise image is $\mathcal G_Q\cap C_{\mathcal G}(G_Q)$, rather than all of $\mathcal G_Q$. When $G_Q$ is nonabelian dihedral, its own elements already distinguish these images. For a reflection in an oval the kernel is $\Z$ or $\widehat\Z$ by Part B and D.4, so the stated injectivity is replaced by an exact sequence with that kernel. These computations identify exactly the rigidity used in the manuscript's argument.
\end{proof}

\begin{claim}{E.23}{Proved}{Explicit discrete centralizers and pointwise kernels}
For the three actions in the question, the full extended geometric groups, their centralizers, their normalizers, and their pointwise kernels on the exceptional points are as follows. Here $D_{2n}$ has order $2n$.
\[
\begin{array}{c|c|c|c|c}
 (g,\nu),G&M&C_M(G)&N_M(G)&C_M(G)^{\natural}\\ \hline
 (0,3),S_3&S_3\times C_2&C_2&S_3\times C_2&1\\
 (1,1),C_4&\mathrm{GL}_2(\Z)&C_4&D_8&C_2\\
 (1,1),C_6&\mathrm{GL}_2(\Z)&C_6&D_{12}&1
\end{array}
\]
The last column fixes every exceptional point individually, not merely its orbit under $G$.
\end{claim}
\passage{25 bis; 155--160; 307--317}

\begin{proof}
For type $(0,3)$, the geometric group is $S_3$ times complex conjugation. The exceptional points are the three transposition fixed points $-1,1/2,2$ and the two order-three fixed points of E.21. Complex conjugation fixes the first three and interchanges the last two. The centralizer of $S_3$ is therefore $C_2$ and its pointwise kernel is trivial; its normalizer is the whole product.

Use the matrices
\[
 R_4=\begin{pmatrix}0&-1\\1&0\end{pmatrix},\qquad
 R_6=\begin{pmatrix}0&-1\\1&1\end{pmatrix}
\]
on $\R^2/\Z^2$, punctured at $0$. A rational matrix commuting with either one is $aI+bR_n$. For an integral commuting matrix $a,b\in\Z$. Its determinant is $a^2+b^2$ for $R_4$ and $a^2+ab+b^2$ for $R_6$. Equating this nonnegative norm to $1$ gives four and six units respectively, proving the centralizer assertions. Conjugation in the normalizer embeds its quotient by the centralizer in $\Aut(C_n)=C_2$ for $n=4,6$. Reflection in the square or hexagonal lattice realizes inversion, proving the two normalizer assertions.

For $C_4$ the exceptional set outside the puncture is exactly the three nonzero two-torsion points. The matrix $-I$ fixes all three, while $R_4$ interchanges $(1/2,0)$ and $(0,1/2)$. Thus the kernel is $\langle-I\rangle$. For $C_6$ the exceptional set consists of the same three two-torsion points and $(1/3,1/3),(2/3,2/3)$. The three two-torsion points force a power of $R_6$ to be $I$ or $-I$, and the latter exchanges the two order-three points. Hence the kernel is trivial. These lists are exhaustive: each nonidentity power has its fixed points given by $(R_n^j-I)v\in\Z^2$; their union is the displayed finite set. The program enumerates these solutions, the kernels, and the matrix centralizers and normalizers. The norm calculation proves exhaustiveness beyond its bounded integer search.
\end{proof}

\begin{claim}{E.24}{Refuted}{Recovery of an arbitrary commuting group from its pointwise subgroup}
For an arbitrary specified $\Gamma\leq C_{\Gamma'}(G)$, its pointwise subgroup is
\[
 \Gamma^{\natural}=\Gamma\cap\Gamma'^{\natural},
\]
and it need not equal $C_{\Gamma'^{\natural}}(G)$. Nor is $\Gamma$ determined as a normalizer by $\Gamma^{\natural}$. For the maximal choice $\Gamma=C_{\Gamma'}(G)$ the proposed equality holds, by the displayed intersection identity.
\end{claim}
\passage{25 bis; 157 verso; 312}

\begin{proof}
Both identities follow by testing separately the two defining conditions: membership in $\Gamma$ and pointwise stabilization of each critical datum. To refute recovery, take the first row of E.23. The choices $\Gamma=1$ and $\Gamma=C_2$ have the same pointwise subgroup $1$. Its normalizer inside the centralizer $C_2$ is $C_2$, so the procedure cannot recover the first choice. For the proposed equality without the maximal-choice hypothesis, use the second row and take $\Gamma=1$: its pointwise subgroup is $1$, whereas the centralizer pointwise kernel is $C_2$. All groups and their kernels are included in the finite check. Thus the failure persists even in the specifically requested discrete cases and has a complete intersection formula as its replacement.
\end{proof}

\section{F. Statements usable in the later sections}

We use $\Gamma_{g,\nu}$ for the extended geometric group permitting puncture permutations and $\Gamma_{g,\nu}^{!+}$ for the pure oriented group, following \S19, folio 102, scans 204--205. In each formula below the chosen variant is stated explicitly. The full profinite loop-outer group is $B=\Out_{\rm lac}(P)$; the closure $H=\overline M$ of a chosen geometric subgroup $M$ is a subgroup of $B$. We do not identify that closure with the abstract profinite completion without the congruence assertion required for that identification.

\begin{claim}{F.1}{Proved}{LM4-F.a: fixed discretizations and prediscretizations}
\label{LM4-F.a}
For $M=\Gamma_{g,\nu}$ and $H=\overline M\leq B$, the sets of outer discretizations and prediscretizations are $B/M$ and $B/H$. For a given finite $G\leq B$ their fixed sets are exactly D.16's transporter sets and centralizer-orbit decompositions. In the three types of D.8, every fixed prediscretization has a fixed outer discretization above it. The same lifting conclusion holds for a finite solvable $G$ and genus at most two for the orientation-preserving geometric group covered by Boggi's Theorem C. Over a fixed coset represented by $b$, the fibre is
\[
 (H/M)^{b^{-1}Gb},
\]
with its full labelled-embedding decomposition as in D.16.
\end{claim}
\passage{28; 178--181; 352--358}

\begin{proof}
The identifications and the fixed-set computation are proved in D.16, including their inner-automorphism ambiguities. A fixed coset $bH$ has $K=b^{-1}Gb\leq H$. Its fibre consists of $bhM$, $h\in H$; this coset is fixed precisely when $h^{-1}Kh\leq M$, proving the formula. In the three low types D.8 conjugates every finite $K\leq H$ into $M$, giving a point in that fibre. In genus at most two and for solvable $K$, Boggi's Theorem C \cite{R-Boggi} gives the same conjugation into the oriented geometric group. The congruence subgroup property in these genera identifies its completion with $H$, so the theorem applies to the displayed subgroup. We retain the labelled embeddings in the fibre decomposition because subgroup conjugacy alone does not fix an automorphism of $K$.

The actual based version has $A(P)/J$ in place of $B/M$, as D.16 proves. A chosen genuine lift of $G$ acts on it; an outer subgroup alone acts on its quotient by $P$. These are the exact maps through which the fixed-set assertions of \S28, folios 178--181, scans 352--358 are used. For a finite group in the full $B$ not known to stabilize a prediscretization, the same formula still determines the fixed set as a transporter; its nonemptiness is exactly the unresolved recognition question D.12, not an additional conclusion of this lemma.
\end{proof}

\begin{claim}{F.2}{Refuted}{LM4-F.b, finiteness and equality of the two centralizers}
\label{LM4-F.b}
The centralizer of a finite subgroup $G\leq\Gamma^+_{g,\nu}$ need not be finite, either discretely or in the geometric completion. In genus at most two, for finite solvable $G$, the normalizer comparison and the open-kernel argument below give
\[
 C_{\widehat{\Gamma^+_{g,\nu}}}(G)
       =\overline{C_{\Gamma^+_{g,\nu}}(G)}.
\]
The corresponding equality with the full loop-outer group in place of the geometric completion is false already for $(g,\nu)=(0,3)$ and $G=S_3$.
\end{claim}
\passage{30; 194--196; 383--387}

\begin{proof}
On a closed genus-two surface let $G$ be generated by the hyperelliptic involution. This involution is central in $\Gamma_2^+$, so its centralizer is the entire group. A Dehn twist about a nonseparating curve has infinite order: its action on integral homology sends a transverse vector $b$ to $b+a$, and its $n$th power sends it to $b+na$. This also gives infinite order in the congruence image, since some reduction modulo an integer detects every nonzero $n$. Thus the completed centralizer is infinite as well. For finite solvable $G$, the normalizer comparison, with the source induction and its topology input audited in F.7, gives $N_{\widehat M}(G)=\overline{N_M(G)}$ for $M=\Gamma^+_{g,\nu}$. The continuous conjugation homomorphism from this profinite normalizer to the finite group $\Aut(G)$ has open kernel $C_{\widehat M}(G)$. Intersecting the dense subgroup $N_M(G)$ with that open kernel gives a dense subgroup of the kernel: every open subset of the kernel is open in the normalizer and meets $N_M(G)$. The intersection is exactly $C_M(G)$, proving the displayed centralizer equality. The map of the completion to the loop-outer group is injective in these genera by the congruence subgroup property. The cited sources state the equality for every finite $G$; F.7 records the precise missing intersection step for their nonsolvable assertion. The passage from any such equality to finiteness would still require the discrete centralizer to be finite, contradicted by the hyperelliptic example.

For the smallest hyperbolic type, $\Gamma_{0,3}^+=S_3$, and $C_{S_3}(S_3)=1$. Harbater--Schneps identify $C_B(S_3)=GT_0$. This group is infinite: its multiplier contains the cyclotomic image $\widehat\Z^\times$ of $G_\Q$. These are actual homomorphisms on the same group $\widehat F_2$, with the geometric $S_3$ supplied by rational M\"obius transformations, so the commutation with the Galois image is the naturality of the fundamental-group action. The two centralizers are therefore $1$ and an infinite group. With orientation reversal included the geometric centralizer becomes $C_2$, and the same failure remains.

The statement denoted (c) on scans 383--385, \S30, folios 194--195, concerns the centralizer of an arithmetic germ $\Gamma^{\natural}$. It cannot be replaced merely by a finite $G$: scan 385 itself discusses infinite centralizers of finite groups. The counterexamples above compute the result of the requested replacement explicitly and explain which datum that replacement removes.
\end{proof}

\begin{claim}{F.3}{Proved}{LM4-F.b, the intrinsic open-image statement}
Let $M\leq B$ be an oriented geometric subgroup, $H=\overline M$, $N=N_B(H)$, and let $i:G\hookrightarrow H$ be a fixed finite embedding. Conjugation defines an action of $N/H$ on the set of $H$-conjugacy classes of embeddings $G\hookrightarrow H$. The image of the continuous homomorphism
\[
 C_N(i(G))\longrightarrow N/H
\]
is exactly the stabilizer of $[i]$, and its kernel is $C_H(i(G))$. When the orbit of $[i]$ is finite, this image is open. The orbit is finite for solvable finite $G$ in genus at most two and for every finite $G$ in the three low types of D.8.
\end{claim}
\passage{30; 195 verso--196; 386--387}

\begin{proof}
An element $nH$ stabilizes $[i]$ exactly when there is $h\in H$ with
\[
 ni(g)n^{-1}=hi(g)h^{-1}\quad(g\in G).
\]
This is equivalent to $h^{-1}n\in C_N(i(G))$. It proves the image assertion, and the kernel is obtained by intersecting the centralizer with $H$. Because $H$ is compact, the stabilizer is the image of a compact closed transporter and is closed in $N/H$. A closed subgroup of finite index in a profinite group is open: its finitely many closed cosets have open complements, so each coset is open.

For the claimed cases, finite subgroup conjugacy classes in $H$ come from the finite list of geometric classes, by D.8 or Boggi's solvable statement. There are finitely many geometric classes for each fixed finite group: a finite action has an orbifold quotient of bounded signature by the Euler characteristic formula; there are finitely many epimorphisms from each finitely generated orbifold group to the fixed finite group. Passing from subgroup classes to labelled-embedding classes adds at most $|\Aut(G)|$ choices per subgroup class. Thus the orbit is finite and the image is open. F.6 constructs the arithmetic quotient directly from universal monodromy and proves the open-image assertion in that quotient for every finite geometric subgroup and every hyperbolic type. The inclusion of that quotient into $N/H$ is retained there as the exact comparison map.
\end{proof}

\begin{claim}{F.4}{Proved}{LM4-F.c: the exact canonical lift of a section stabilizer}
\label{LM4-F.c}
Let $\pi$ be a centre-free discrete surface group or its full profinite completion, and put
\[
 A=\Aut_{\rm lac}(\pi),\qquad B=\Out_{\rm lac}(\pi).
\]
Let $G\leq B$ be finite and let $s:G\to A$ be a specified lift with $C_\pi(s(G))=1$. The natural map
\[
 C_A(s(G))\longrightarrow C_B(G)
\]
is injective. Its image is the stabilizer of the $\pi$-conjugacy class $[s]$ in the set of lifts of the labelled inclusion $G\hookrightarrow B$. Its inverse is the canonical homomorphism lifting that stabilizer to $A$ and centralizing $s(G)$. In the profinite case it is continuous. For geometrically realized actions with a nontrivial rotation the section-class set is finite, so this stabilizer has finite index and is open in the profinite case.
\end{claim}
\passage{31; 197 verso--199; 390--393}

\begin{proof}
The kernel is the inner automorphisms centralizing $s(G)$, identified, by centre-freeness, with $C_\pi(s(G))$. For $b\in C_B(G)$ choose an automorphism representative $v$. The conjugate homomorphism $g\mapsto vs(g)v^{-1}$ still lifts the labelled inclusion, because $b$ centralizes each $g$. It represents $[s]$ exactly when some $p\in\pi$ makes $\Inn(p)^{-1}v$ commute with every $s(g)$. This proves that the image equals the claimed stabilizer. The kernel computation proves uniqueness. The inverse is a homomorphism by uniqueness, since products of commuting lifts still commute with $s(G)$.

For a finitely generated profinite $\pi$, $A$ and $B$ are profinite, and the centralizer is closed. The map is a continuous injection from a compact space into a Hausdorff space, so its inverse on its image is continuous. In the discrete geometric situation Part B proves finiteness of the lift classes; D.3 proves their equality with the profinite classes. The orbit of a class is therefore finite. Its stabilizer is closed, as the image of the compact centralizer, hence open. This proves all topology assertions as well as the group assertion.

For the notation of scans 390--393, \S31, folios 197--199, the target $\mathfrak S(\pi_0)^G$ is the centralizer of the chosen genuine lift inside the loop automorphism group. The source is the stabilizer of the corresponding section class inside $\mathcal T(\pi_0)^G=C_{\Out_{\rm lac}(\pi_0)}(G)$. A point $P$ determines that section class after a base path; changing the path conjugates the resulting lift by $\pi_0$, exactly as displayed. Thus the canonicity belongs to the chosen section, and the geometric point supplies its conjugacy class.
\end{proof}

\begin{claim}{F.5}{Refuted}{The unrestricted uniqueness assertion for reflection ovals}
For a reflection fixing an oval, the lifting map in F.4 has kernel $\Z$ in the discrete case and $\widehat\Z$ after completion. Consequently it is not injective and its lifts are not unique without further data.
\end{claim}
\passage{31; 197 verso--199; 390--393; 33; 222 verso--225; 440--445}

\begin{proof}
Part B and D.4 identify $C_\pi(s(G))$ with the primitive oval subgroup and its closure. Every element of that subgroup induces an inner automorphism centralizing the chosen reflection lift and maps to the identity outer class. The kernel computation in F.4 therefore gives the displayed groups. A nonidentity primitive oval element is an explicit second lift of the identity outer element. The corrected general statement is the exact sequence with this kernel and with image the same section stabilizer, proved by the same representative-and-conjugator calculation. No splitting of that sequence is supplied by the point class alone.
\end{proof}

\begin{claim}{F.6}{Proved}{LM4-F.b: open arithmetic image for every finite geometric subgroup}
Let $2g-2+\nu>0$, let $M=\Gamma^+_{g,\nu}$ permit permutations
of the $\nu$ punctures, and retain
\[
 P=\widehat\pi_{g,\nu},\qquad
 B=\Out_{\rm lac}(P),\qquad H=\overline M\leq B.
\]
Write $A_{\rm ar}\leq B$ for the image of universal \emph{arithmetic}
\'etale monodromy of the moduli stack of smooth genus-$g$ curves
with an unordered reduced divisor of degree $\nu$, over $\Q$.
Then $H\triangleleft A_{\rm ar}$, and put
$\Gamma_{\rm ar}=A_{\rm ar}/H$.
For every finite subgroup $G\leq M$, the continuous homomorphism
\[
 C_{A_{\rm ar}}(G)\longrightarrow\Gamma_{\rm ar}
                                                        \tag{F.6.1}
\]
has open image.  Its kernel is $C_H(G)$.
In the notation $N=N_B(H)$ and $q:N\to N/H$, this is exactly
the assertion that
\[
 q(C_N(G))\cap\Gamma_{\rm ar}
       =q(C_{A_{\rm ar}}(G))
       \quad\hbox{is open in }\Gamma_{\rm ar}.           \tag{F.6.2}
\]
The same assertion holds with ordered punctures and the pure
oriented group throughout.  No identification of $H$ with the
abstract completion of $M$ is used.
\end{claim}
\passage{30; 195--196; 385--387}

\begin{proof}
We first type the arithmetic group and its comparison maps.
Denote the indicated moduli stack over $\Q$ by $\mathscr M$,
and choose a geometric base point and an orientation-preserving
topological marking of its complex fiber.  Its universal
punctured curve supplies the outer monodromy homomorphism
\[
 R:\Pi_{\rm ar}=\pi_1^{\rm et}(\mathscr M)
           \longrightarrow B.
\]
The geometric comparison and the arithmetic exact sequence are
\[
 1\longrightarrow\widehat M
  \longrightarrow\Pi_{\rm ar}
  \longrightarrow G_{\Q}\longrightarrow1,
 \qquad R(\widehat M)=H.                               \tag{F.6.3}
\]
These are the constructions in \cite[\S\S7.1--7.2,
equation (29) and the universal \'etale monodromy
representation]{R-BoggiAut2026}.  For the geometric comparison,
finite \'etale covers analytify to finite topological covers;
the inverse comparison is compatible on an \'etale atlas and
its double and triple overlaps, so descends to the stack.
The marked complex moduli stack has contractible Teichm\"uller
space as its covering, with group $M$, which identifies its
topological fundamental group.  The universal curve then
identifies its monodromy with the completed geometric action
on $P$, rather than introducing a separate representation.
All these maps preserve peripheral conjugacy classes, since
the removed divisor is carried to itself in the universal
family.  Their base-point ambiguities are inner automorphisms
of $P$, already removed in $B$; a change of topological marking
conjugates the displayed subgroups by the corresponding element
of $M$.

Set $A_{\rm ar}=R(\Pi_{\rm ar})$.  It is compact and hence
closed in $B$.  Normality in (F.6.3) gives
$H\triangleleft A_{\rm ar}\leq N$.  The quotient map induces
a continuous surjection
\[
 \overline R:G_{\Q}\longrightarrow\Gamma_{\rm ar}.     \tag{F.6.4}
\]
The proof uses this quotient of $G_{\Q}$ as defined.  Its kernel
and the kernel of $\widehat M\to H$ need not be computed.

We next construct a number-field curve with the specified
\emph{labelled topological} $G$-action.  Finite realization
C.5 supplies a compact Riemann surface $X$, a reduced marked
divisor $D$, and a faithful holomorphic action of $G$ preserving
$D$ whose marked mapping classes are the given inclusion
$G\hookrightarrow M$.  Choose projective equations for $X$,
equations for $D$, and formulas for the finitely many
automorphisms and their inverses.  Their finitely many
coefficients lie in a finitely generated $\Q$-subalgebra
$T\subset\C$.  Adjoin the coefficients needed for their
identities and invert finitely many nonzero elements.  This
gives a proper smooth family of curves $\mathcal X\to V$,
a finite \'etale reduced divisor $\mathcal D\to V$, and
the same labelled $G$-action over a nonempty smooth open
subscheme $V\subset\operatorname{Spec}T$.  Smoothness and
disjointness are open conditions.  Faithfulness can be kept
on this open set as well: for each nonidentity $h\in G$,
the equalizer of $h$ and the identity is a proper closed
subset of the generic fiber, and this property persists
after shrinking.  The original inclusion $T\hookrightarrow\C$
defines a point of every such nonempty open, since it has
generic underlying point.

Base change to $\overline\Q$ and take the connected component
whose complex base change contains that point.  After a
finite extension this component and the family are defined
over a number field.  It is a smooth connected variety; its
complex points form a connected, locally path-connected
complex manifold and hence a path-connected space.  It
contains an $\overline\Q$-point, by taking a closed point
in any nonempty affine open.  Choose a piecewise smooth
complex path from the original fiber to this fiber.

For precision, parallel transport along this path preserves
the labelled action.  A smooth horizontal splitting of the
relative tangent sequence can be chosen to agree with the
tangent splitting on $\mathcal D$: locally the divisor is
a disjoint union of sections, and a partition of unity
extends their prescribed splittings.  Average this splitting
over the finite group $G$.  Averaging retains its right-inverse
property and its prescribed values on the divisor, and
makes it $G$-equivariant.  The horizontal lift along the
compact path has complete transport, because the family
is proper and its fibers are compact.  Its transport is an
orientation-preserving diffeomorphism of pairs commuting
with each labelled element of $G$.  Thus the endpoint curve
has exactly the required topological action.  This constructs
a deformation and its comparison diffeomorphism; it does not
assert that the original complex curve is isomorphic to a
curve over $\overline\Q$.

All coefficients of the endpoint curve, divisor, and the
finitely many labelled automorphisms lie in one finite
extension $K/\Q$.  Let $x_K:\operatorname{Spec}K\to\mathscr M$
be the resulting moduli point.  Choose the geometric path
to $x_K$ supplied by the transport just constructed and the
finite-cover comparison.  The morphism $x_K$ then supplies
a continuous homomorphism
\[
 s_K:G_K\longrightarrow\Pi_{\rm ar}
\]
whose composite with $\Pi_{\rm ar}\to G_{\Q}$ is the
inclusion $G_K\hookrightarrow G_{\Q}$.  Under the transported
marking, $R\circ s_K$ is the Galois outer action on the
punctured endpoint curve.

Every $h\in G$ is an automorphism over $K$ of that curve and
its marked divisor.  Naturality of the fiber functor on finite
\'etale covers gives
\[
 (R\circ s_K)(\sigma)\,h
       =h\,(R\circ s_K)(\sigma)
 \quad\hbox{in }B\qquad(\sigma\in G_K).                \tag{F.6.5}
\]
Indeed, pullback by $h$ commutes with scalar conjugation by
$\sigma$; choosing paths for the individual base points can
alter the induced automorphisms only by inner automorphisms
of $P$.  Passing to $B$ gives precisely the equality written.
The equivariant transport aligns these $h$ with the original
labelled inclusion, so (F.6.5) places $R(s_K(G_K))$ in
$C_{A_{\rm ar}}(G)$ itself.

The image of (F.6.1) therefore contains
$\overline R(G_K)$.  The latter has index at most
$[K:\Q]$ in $\Gamma_{\rm ar}$ by the surjectivity of
(F.6.4).  It is compact and hence closed; a closed finite-index
subgroup of a profinite group is open, as proved in F.3.
Every subgroup containing it is a union of its open cosets.
This proves openness in (F.6.1), with kernel $C_H(G)$ by
intersection with $H$.

Finally $q^{-1}(\Gamma_{\rm ar})=A_{\rm ar}$, since
$H\subset A_{\rm ar}$ and $\Gamma_{\rm ar}=q(A_{\rm ar})$.
An element of $C_N(G)$ whose quotient lies in $\Gamma_{\rm ar}$
therefore already belongs to $C_{A_{\rm ar}}(G)$.
This proves both inclusions in (F.6.2).  Ordering the divisor
and using the pure moduli stack changes none of the maps or
arguments when $G$ belongs to the pure oriented group.

For \S30, folios 195--196, scans 385--387, formula (F.6.2)
is the arithmetic open-image assertion with its exact
quotient made explicit.  F.3 supplies the centralizer image
in the entire quotient $N/H$; F.6 supplies its openness
after intersection with the arithmetic image for every
finite geometric $G$.  The inclusion
$\Gamma_{\rm ar}\hookrightarrow N/H$ is the comparison
map between the two statements.  Its image has not been
identified with all of $N/H$ here, so no larger openness
assertion is obtained by suppressing that inclusion.
\end{proof}

\begin{claim}{F.7}{Open}{The nonsolvable centralizer comparison and its checked finite image}
For an oriented mapping class group $M$ of genus at most two,
the normalizer and centralizer comparisons for finite solvable
subgroups have the complete proof below. For a geometric
$A_5\le M=\Gamma^+_{0,\nu}$, the equality
\[
                 C_{\widehat M}(A_5)
                    =\overline{C_M(A_5)}                 \tag{F.7.1}
\]
is stated in \cite[Theorem~C(iii)]{R-Boggi} and
\cite[Theorem~5.1(iii)]{R-BZ}. The displayed proofs contain
an additional closure--intersection equality for which we
have not verified a justification. We do not assert that
(F.7.1) is false.

There is nevertheless equality of their images in the
finite puncture-permutation quotient. Write the invariant
puncture set as a union of some of the exceptional orbits
of sizes $12,20,30$ and $m$ free orbits of size $60$.
Both images fix the exceptional punctures individually
and induce the full group $A_5^m\rtimes S_m$ on the
free orbits. In particular, with only exceptional orbits
the unresolved part of (F.7.1) lies in the pure profinite
mapping class group; its discrete right side is trivial.

F.12--F.14 identify the residual twelve-puncture $A_5$ centralizer with an exact common-involution equalizer in the pure eight-point completion and prove injectivity after the specified ordinary-point forgetting. F.15 proves the full two-point pushing comparison and removes every cocycle in the actual intersection of the two forgetting kernels. Its correcting elements form the explicitly described fixed-subgroup cosets. The original source image has trivial intersection with that kernel; membership of a corrected element in that image remains the full equalizer equation. F.16 constructs a full closed procyclic subgroup of the correcting centralizer that every proper labelled forgetting kills and whose intersection with the source image is trivial. The exact membership test using all proper forgetting images, even with the common-involution centralization equation, therefore has a strictly larger solution group than the original source image. F.17 identifies the sixteen additional punctures required by the literal rational fixed-point drilling and proves its nonfactorization through the original source, using primitive full procyclic winding detectors. F.18 proves that every original residual element has a unique full branch-ordered centralizing lift. Its membership in the rational drilling image is the distinguished fibre of the explicit continuous coset map $\mathfrak o:\mathcal E\to\mathcal J/K^+$. Whether that fibre is the whole residual, and whether the residual is trivial, remain unresolved.
F.19 identifies the joint forgetting kernel $D^+$ as the closed ambient normal closure of eight actual cross-collision meridians and proves its exact one-class cocycle comparison with a unique corrector. The canonical correction therefore lies in $D^+$, and the source-membership obstruction has the precise target $D^+/K^+$. The actual drilled image and the full two-track pullback have common explicit splittings over the original rank-eleven source, with kernels $K^+$ and $D^+$ respectively. Equality of those two subgroups and triviality of the original residual remain unresolved.
F.20 proves the actual rank-three cross-divisor coset coordinate, its eleven based source actions and all original transports. Exact configuration computations and finite group calculations prove transitivity on the $82$ invariant $A_5$ quotients, all $3,321$ distinct paired quotients, and six noncentral quotients of order $60\cdot2^{121}$. The specified continuous finite wreath maps consequently have the same full images on $K^+$ and $D^+$, for every coefficient prime. The full source-membership equation remains unresolved.
F.21 proves transitivity on all eight full cross-divisor coset spaces and constructs lifts of their entire divisor groups in the actual selected-meridian centralizers. Consequently the closed normal closure of $K^+$ inside $D^+$ is $D^+$, and $K^+$ maps onto the full abelianization and maximal pronilpotent quotient of $D^+$. The exact equality $K^+=D^+$ and the original residual source-membership equation remain unresolved.
F.22 gives the actual two-point kernel as a semidirect product of two explicit compactly generated profinite kernels, retains its full Artin action, and identifies the interaction kernel killed by both one-point forgettings. The source map has image exactly $K^+$; its joint one-point image is the product of its two actual marginal images. Fullness of those marginals and surjectivity of the displayed source map remain unresolved.
F.23 gives an explicit based source-sign symmetry carrying the two actual marginal inclusions isomorphically to one another. It specifies all twenty-seven original source images by exact nine-strand words, with the original basepoint paths and inverse fiber coordinates. For every continuous homomorphism from either named ambient rank-seven group to $S_3$ or $A_4$, the marginal source and marginal target have equal images. The full marginal equalities, the interaction source image and the original residual membership remain unresolved.
F.24 proves injectivity of the actual restriction on continuous rank-one crossed homomorphisms and first cohomology for every finite-field character nontrivial on the marginal target. Sixty-four actual source rows have rank seven on four explicit Laurent charts over every characteristic. The signed comparison retains the original source maps. Every scalar affine representation has the same fixed-point set on the marginal source and target, and its two marginal images agree whenever the source image meets the translation subgroup trivially. The remaining nonzero translation intersections, full marginal images and interaction source image remain unresolved.
F.25 proves equality of those affine images for every representation whose source translation intersection is linear over the ambient scalar field. This includes every representation with equal ambient and marginal scalar fields, and every representation into $\operatorname{Aff}_1(\mathbb F_p)$ for every prime $p$. The chosen scalar lift, its translation conjugation, and the separating coefficient maps are explicit. Unequal scalar fields without that linearity property, the full marginal equalities and the interaction source image remain unresolved.
F.26 gives an exact equivalent trace-cocycle kernel criterion for a proper affine marginal image, with the actual conjugation and coefficient homomorphism. Its tensor and Frobenius coefficient maps, including their inverses, identify the remaining rank target as $6m+1$ on $7m$ columns. Complete finite verification and an independently written replay prove equality of the actual marginal images for every scalar field of order nine, even inside a larger original coefficient field. General attainment of the trace rank, the full marginal equalities and the interaction source image remain unresolved.
F.27 proves the exact trace rank and affine marginal image equality over every finite coefficient field on the actual two-parameter character locus $(1,s,t,1,1,1,(st)^{-1})$. Its integer Laurent commutators force an essential source line, and actual character-one seed conjugates give an invertible Frobenius Vandermonde comparison. F.28 extends this to the three-parameter locus $(r,s,t,1,1,1,(rst)^{-1})$: an explicit two-dimensional invariant inclusion and four quotient coordinates separate all minimal source submodules across distinct Frobenius factors. The individual restriction theorem then proves summed restriction injective and rank $6m+1$. Both results retain the exact based sign map and include larger original coefficient fields. The general scalar-character locus, full marginal equalities, interaction source image and original residual membership remain unresolved.
F.29 proves the larger exact four-parameter locus
$(r,s,t,u,1,1,(rstu)^{-1})$ over every finite coefficient field.
Four actual principal-coordinate charts retain all original source words.
Three charts have complete comparisons of their minimal invariant
subspaces across Frobenius factors; the exceptional overlap has a
forced-lift residual with coordinate $-2$ in its proved odd
characteristic. The fourth chart has an essential line detected by
actual character-one seed conjugates and an explicit Vandermonde
inverse. Summed restriction has rank $6m+1$, and the exact trace
and based sign maps give both marginal image equalities on this locus.
The general character locus and the original full marginal and
interaction comparisons remain unresolved.
\end{claim}
\passage{\S{}30; folios 194--196; scans 383--387; Part D(b) and Part F(b): nonsolvable completion centralizers}

\begin{proof}[Source audit and proved comparisons]
The supplied preprint snapshot of \cite{R-Boggi}
has since appeared in \emph{S\~ao Paulo Journal of
Mathematical Sciences}, volume~20, article~35,
published 23 September 2026, DOI
\texttt{10.1007/s40863-026-00558-y}.
The publisher's full text was checked as well:
Proposition~2.4 retains the same displayed
closure--intersection step, and Lemma~3.2
still invokes that proposition.
The audit below applies to both inspected versions;
it is a proof-step audit and makes no claim based
on the paper's publication status.

For every finite subgroup $J\le M$, there are canonical
continuous maps
\[
 \widehat{C_M(J)}\longrightarrow
       \overline{C_M(J)}\lhook\joinrel\longrightarrow
       C_{\widehat M}(J).                              \tag{F.7.2}
\]
The first is onto by density and compactness. Its
injectivity states that the topology induced from $M$
on the discrete centralizer is its full profinite
topology. Surjectivity of the second arrow is the
ambient centralizer comparison. These are two specific
properties of the displayed maps; injectivity of the
first arrow does not supply surjectivity of the second.

In \cite[Proposition~2.4]{R-Boggi}, choose the displayed
generators $x_1,\ldots,x_k$ of the finite group, each of
prime-power order. The prime-power centralizer theorem
does prove
\[
 C_{\widehat M}(J)
   =\bigcap_i C_{\widehat M}(x_i)
   =\bigcap_i\overline{C_M(x_i)}.                      \tag{F.7.3}
\]
The source then identifies this last intersection with
$\overline{\bigcap_i C_M(x_i)}$. The same final equality
occurs in the proof of
\cite[Theorem~5.1(iii)]{R-BZ}, there indexed by all
$x\in J$. Neither displayed argument supplies a
closure--intersection hypothesis. This exact step is
what is left unverified here. The later virtually
compact special application in \cite[\S6.2]{R-BZ}
has a separate relatively quasiconvex product-separability
input; no such input has been proved here for the
punctured-sphere mapping class group.

\emph{A noncircular solvable proof.}
We first record the prime-power inputs used in the
induction. Goodness and finite virtual cohomological
type of mapping class groups of genus at most two,
and their finite-index subgroups, permit the direct
application of \cite[Theorem~A]{R-BZ}. Thus for a
finite $p$-subgroup $A$ we have
\[
 N_{\widehat M}(A)=\overline{N_M(A)},\qquad
 C_{\widehat M}(A)=\overline{C_M(A)}.                   \tag{F.7.4}
\]
The proof of \cite[Theorem~A and Corollary~B]{R-Boggi}
also supplies the full induced profinite topology on
$C_M(A)$ and hence on $N_M(A)$ in this prime-power case.
Here is why this use avoids the preceding gap. Its
free-action argument is the restriction-injectivity
argument on the finite covering surface group. Its
ramification-removal argument uses Proposition~2.4
only for the specified group $A$, now a $p$-group,
in the good configuration-group extension. For that
application the second equality in (F.7.4), directly
from \cite[Theorem~A]{R-BZ}, replaces Proposition~2.4's
intersection argument. Thus the two required inputs
are justified without using any nonsolvable assertion.
The passage from centralizer to normalizer topology
uses its finite-index kernel under conjugation into
$\Aut(A)$; this kernel is open in the induced topology,
as is detected by any finite quotient in which $A$
injects.

We prove the solvable normalizer comparison, together
with its full induced topology, by induction on the
order of the finite solvable group, simultaneously for
all finite-index subgroups of these mapping class
groups. Finite-index passage preserves the comparison:
intersect the completed normalizer with the corresponding
open subgroup and use density in an open set.

For a nontrivial finite solvable $J$, take a nontrivial
$p$-primary subgroup $A$ of the last nonzero term of
its derived series. That term is abelian, so its
$p$-primary subgroup is characteristic in $J$.
Consequently every normalizer of $J$ normalizes $A$.
The discrete covering normalizer sequence is
\[
 1\longrightarrow A\longrightarrow N_M(A)
                 \longrightarrow Q\longrightarrow1,   \tag{F.7.5}
\]
where $Q$ is a finite-index subgroup of the mapping
class group of the quotient surface with its branch
and puncture decorations. Its genus is again at most
two. The exact quotient map is induced by descent of
the commuting covering diagram; the decorations
record the original punctures, the removed branch
points and the covering subgroup. This is the discrete
normalizer sequence in \cite[\S2.1]{R-Boggi}.

By the prime-power topology input, the completed
normalizer in (F.7.4) is $\widehat{N_M(A)}$.
Completing (F.7.5) therefore gives its quotient by
the finite embedded kernel $A$ as $\widehat Q$.
The subgroup $J/A\le Q$ is smaller and solvable.
The induction hypothesis applies to its normalizer
in $\widehat Q$. The inverse image of that normalizer
under the quotient in (F.7.5), on either side, is
exactly the normalizer of $J$: $J$ is the full inverse
image of $J/A$ because it contains $A$. Pulling back
the dense discrete normalizer proves
\[
                  N_{\widehat M}(J)=\overline{N_M(J)}.
                                                               \tag{F.7.6}
\]
It also proves the topology assertion. Indeed the
completion of this inverse image maps to
$\widehat{N_M(A)}$; on the quotient by $A$ it is
injective by induction, and on $A$ it is the specified
inclusion. Its kernel is therefore trivial.
This completes the induction.

Conjugation gives a continuous homomorphism
\[
       N_{\widehat M}(J)\longrightarrow\Aut(J)
\]
to a finite group. Its kernel is the ambient centralizer
and is open. The dense subgroup $N_M(J)$ from (F.7.6)
meets that open kernel densely, and its intersection
there is exactly $C_M(J)$. This proves
\[
                  C_{\widehat M}(J)=\overline{C_M(J)}
\]
for every finite solvable $J$, without taking an
unjustified intersection of closures.

\emph{The finite image for $A_5$.}
Realize the geometric $A_5$ as the icosahedral rotation
group. Its exceptional point stabilizers have orders
$2,3,5$, and the corresponding transitive orbits are
$A_5/C_2$, $A_5/C_3$, $A_5/C_5$. Every permutation
commuting with this transitive action is right
multiplication by an element of $N_{A_5}(C_d)/C_d$.
This quotient has order two in each of the three
cases; its nonidentity permutation exchanges the
two poles of each relevant rotation axis.

An orientation-preserving homeomorphism commuting
with an odd-order nonidentity rotation fixes its
two poles individually. The local oriented rotation
numbers at the two poles are opposite, and an
orientation-preserving commuting homeomorphism
preserves this local rotation number. This excludes
the nonidentity permutation on the $12$- and
$20$-point orbits for the discrete cyclic centralizers.
Their completed centralizers have the same finite
permutation images by (F.7.4). The full $A_5$
centralizer lies in these cyclic centralizers, so
it fixes these exceptional punctures individually.

For an order-two stabilizer choose a tetrahedral
subgroup $A_4\le A_5$ containing it. An
orientation-preserving homeomorphism commuting
with $A_4$ fixes every ramification point of the
$A_4$ quotient. To see this, its quotient fixes the
branch point of order two. It also fixes each of the
two branch points of order three, rather than
exchanging them, by the preceding local rotation
argument. Isotope the quotient map to the identity
relative to these three branch points and lift the
isotopy. The lift stays in the centralizer; its
endpoint is a deck transformation in $Z(A_4)=1$.
The ramification set is finite and invariant along
the isotopy, so every one of its points is fixed.
The usual discrete covering-centralizer comparison
provides a commuting representative for a centralizing
mapping class. Since $A_4$ is solvable, the proved
centralizer comparison transfers this fixed-puncture
property to $C_{\widehat M}(A_4)$. It therefore holds
for $C_{\widehat M}(A_5)$, excluding the last possible
nonidentity permutation on the $30$-point orbit.

On $m$ free orbits the finite commutant is
$A_5^m\rtimes S_m$. Every element of this commutant
is already realized by the discrete centralizer.
For a direct construction, remove the three branch
points from the quotient sphere. Its unramified
$A_5$-cover has surjective monodromy onto $A_5$.
Move one of the $m$ quotient punctures around a
loop, keeping all the others fixed, and lift the
isotopy starting at the identity. The lift commutes
with $A_5$ and acts on that one free orbit by its
specified monodromy, while fixing the other orbits
and all ramification points. These loops realize
each factor $A_5$. Half-twists interchanging the
quotient punctures realize $S_m$ modulo these factors.
Thus the displayed wreath product is the exact
finite permutation image on both sides of (F.7.1).

The program
\path{code/part_a/a5_centralizer_permutations.py}
constructs $A_5$ as the sixty even permutations of
five letters. It checks its order-two and order-three
generating pair, the three coset actions, normalizer
orders $4,6,10$, and commutants of order two. The
fixed-pole constraints leave only the identity in
each exceptional block; the order-two block uses
the six ramification points for the specified $A_4$.

For completeness, the residual intersection question
has an exact finite-quotient formulation. Take
$u=(1\ 2)(3\ 4)$ and $v=(1\ 3\ 5)$ in $A_5$, which
generate it, and put $D_u=C_M(u)$, $D_v=C_M(v)$.
For open normal subgroups $V\subset U$ of $M$, let
$D_u(V),D_v(V)$ denote their images in $M/V$.
Compactness of coset fibres proves
\[
 \operatorname{im}\bigl(C_{\widehat M}(A_5)\to M/U\bigr)
 =\bigcap_{V\subset U}
       \operatorname{im}\bigl(D_u(V)\cap D_v(V)\to M/U\bigr).
                                                               \tag{F.7.7}
\]
Indeed choose a lift $t\in\widehat M$ of a point
in the right side. For each $V\subset U$, the set
\[
 (\overline{D_u}\,\overline V)\cap
 (\overline{D_v}\,\overline V)\cap t\overline U
\]
is nonempty and compact. These sets are directed
under refinement, so their finite intersections
are nonempty. Compactness produces a single element
in their intersection, which is
$\overline{D_u}\cap\overline{D_v}\cap t\overline U$.
Conversely any such element supplies every finite
coordinate. Equations (F.7.3) and (F.7.7) identify
the remaining assertion precisely: its right side
must equal the image of $C_M(A_5)$ for every $U$.
The finite permutation computation proves that
equality at the puncture-permutation level. It does
not establish it at every refinement. With only
exceptional orbits the discrete centralizer is
trivial, since the quotient sphere has just the
three differently ordered branch points and
$Z(A_5)=1$; the still unproved assertion is then
that the pure completed centralizer is trivial.

\emph{The minimal twelve-puncture case.}
There is a further exact localization of this remaining group.
Write $P=P\Gamma_{0,12}$ and choose $A=A_4\leq A_5$.
The action of $A$ on the twelve vertices is free and transitive.
Its quotient has one free puncture and three branch points of
orders $2,3,3$. The labelled commuting homeomorphism group on
the filled sphere is identified, by descent and the unique
commuting lift, with the identity homeomorphism component of
the quotient fixing those three branch points. It is
contractible. Evaluating the free quotient puncture consequently
identifies the mapping-class centralizer with the free group
$\pi_1(S^2\setminus\{3\text{ branch points}\})=F_2$.
Its permutation on the free fibre is the surjective monodromy
onto $A_4$. The pure centralizer is therefore the exact kernel
\[
 1\longrightarrow C_P(A)=F_{13}\longrightarrow F_2
                    \longrightarrow A_4\longrightarrow1.
                                                               \tag{F.7.8}
\]
Indeed it is the fundamental group of
$X=S^2\setminus R_A$, where the $A_4$ ramification set $R_A$
has $6+4+4=14$ points; its rank is also
$1+12(2-1)=13$ by Schreier. Label the vertices by $g\in A_4$
using a chosen vertex $z_0$. The inclusion is the configuration
map $z\mapsto(gz)_{g\in A_4}$, followed by the quotient by
M\"obius transformations. The proved solvable comparison and
full topology identify its completion with $C_{\widehat P}(A)$.
In particular the unknown group $C_{\widehat M}(A_5)$ lies in
this specified embedded $\widehat F_{13}$.

Each of the six subgroups $C_5\leq A_5$ has two fixed vertices
and two free five-cycles. Retain either one fixed vertex and
one five-cycle, giving twenty-four six-point sets, or retain
both fixed vertices and one five-cycle, giving twelve
seven-point sets. For each retained set $I$, the continuous
forgetting map $r_I$ is defined on $\widehat P$ and is equivariant
for the specified $C_5$. In either target the quotient has
exactly its two branch points and one free puncture. Its
discrete $C_5$ centralizer is therefore exactly the deck group
$C_5$: the quotient mapping class group fixing these three
decorated points is trivial. Prime-power comparison gives the
same completed centralizer. Its intersection with the pure
group is trivial, since the deck group permutes the five-cycle.
It follows that
\[
 C_{\widehat M}(A_5)\subset
 C_{\widehat P}(A)\cap
       \bigcap_{I\text{ among these }36\text{ sets}}\ker(r_I).
                                                               \tag{F.7.9}
\]

The six-point maps in (F.7.9) have a proved common kernel even
on the discrete $F_{13}$. The antipodal permutation on the
twelve vertices commutes with $A_5$. In the regular $A_4$
labelling it is right multiplication by an involution $h\in A_4$.
Every retained six-point set contains exactly one member of
each antipodal pair. Collisions of two entries $g_i z,g_j z$
at a ramification point with stabilizer $H\leq A_4$ occur
precisely when $g_j^{-1}g_i\in H$; these are the right cosets
$gH$. Thus there is no collision among the retained entries
at either fixed point of $h$. Each six-point configuration map
extends across both of these deleted points of $X$. Their
nontrivial peripheral meridians in $F_{13}$ map to the identity
under all twenty-four maps. Completion retains these meridians
as nontrivial primitive procyclic subgroups. Consequently the
joint six-point map on the embedded $\widehat F_{13}$ is
noninjective by these explicit elements.

Every seven-point set retains a collision at each of the
fourteen ramification points: seven choices from six pairs
force a pair, and seven choices from four triples force a
triple collision. The program verifies all twelve sets and
the preceding common invisible axis for the six-point sets,
using right cosets throughout. It also computes the induced
maps on $H_1(F_{13},\F_{101})$. A meridian at an $A_4$
ramification point maps to the sum of the pair generators
within its collision clusters. In the target pure sphere
group the pair generators satisfy the vertex-sum and total-sum
relations. The two fixed points of each cyclic stabilizer give
the same homology column. Each seven-point map and their joint
map have rank six, with seven-dimensional kernel; the joint
six-point map has rank five. These are exact calculations of
the indicated abelian quotients. They neither produce a
nontrivial $A_5$ centralizer element nor prove injectivity of
the joint seven-point map on $\widehat F_{13}$.

F.8 proves that all twelve seven-point kernels in (F.7.9)
are the same, both discretely and after completion. It further
proves that the discrete kernel contains no nonidentity reducible
mapping class and that the completed kernel meets every conjugate
of every full source cusp subgroup trivially. The exact four-point
characters detecting those cusps and all eighty-four finite
subset--inertia checks are supplied there. F.9 subsequently
constructs an explicit nontrivial common kernel element with
a full procyclic closure, and proves that all twenty-four
six-point restrictions factor through these seven-point maps.
The thirty-six-fold kernel is therefore nontrivial. The full
centralizer comparison and geometric conjugacy remain unproved;
no converse to the inclusion (F.7.9) is inferred.

\medskip\noindent\textbf{The exact common-$C_3$ comparison.}

The following maps give an exact model for the residual group in (F.7.1).
The fixed intersection displayed below is the remaining unresolved group.
All permutations use the five-letter identification of $A_5$ already
fixed in F.7--F.10.  Put
\[
 A=\operatorname{Stab}_{A_5}(5),\qquad
 r=(1\ 2\ 3),\qquad K=\langle r\rangle,
 \qquad\beta=(1\ 2)(4\ 5).
\]
Then $A\cong A_4$, $\beta K\beta^{-1}=K$ by inversion,
and $\langle A,\beta\rangle=A_5$.

\medskip\noindent\textbf{The actual quotient homomorphism.}
Let $P=P\Gamma_{0,12}$ at the icosahedral configuration.
The discrete branched-cover sequence for the $K$ action is
\[
 1\longrightarrow K\longrightarrow N(K)
       \longrightarrow Q\longrightarrow1,
\]
where $Q$ is the finite-index liftable subgroup of the
mapping class group of the quotient sphere with two decorated
branch points and four free punctures.  The centralizer
$C_P(K)$ has finite index in $N(K)$: both the centralizing
condition and the pure puncture condition have finite index.
Its intersection with $K$ is trivial, since nonidentity deck
maps permute the twelve punctures.  Its image $J$ in $Q$ is
therefore isomorphic to $C_P(K)$ and has finite index in $Q$.
Every element of $J$ fixes the four free quotient punctures.
It also fixes the two branch points separately: exchanging
them conjugates the order-three local deck generator to its
inverse, and hence cannot centralize it.  Thus
\[
             C_P(K)\xrightarrow{\ \sim\ }J
                         \leq P\Gamma_{0,6},
               \qquad [P\Gamma_{0,6}:J]<\infty.          \tag{F.7.10}
\]
The prime-power comparison and full induced topology proved
in F.7 identify $C_{\widehat P}(K)$ with
$\widehat{C_P(K)}$.  Completing the finite-index embedding
(F.7.10) identifies it with the specified open subgroup
$\widehat J\leq\widehat{P\Gamma}_{0,6}$.  This gives an
injective continuous homomorphism on the actual completed
centralizer; no completed intersection identity is used.

By F.7.8 the orbit map identifies $C_P(A)$ with the specified
$F_{13}=\pi_1(X)$, where
\[
 X=\mathbf P^1\setminus R_A,\qquad
 R_A=\{\infty\}\cup\{z:z(z^4-1)(z^8+14z^4+1)=0\}.
\]
Since $K\leq A$, composition with (F.7.10) is the injection
$j:F_{13}\hookrightarrow P\Gamma_{0,6}$ described next.
The solvable comparison and full topology for $A$ and $K$
show that its continuous extension
\[
                \widehat j:\widehat F_{13}
                       \hookrightarrow\widehat{P\Gamma}_{0,6}
                                                               \tag{F.7.11}
\]
is injective.  This is a different, explicitly specified map
from the thirty-six forgetful maps whose common kernel was
computed in F.9.

\medskip\noindent\textbf{The six labelled quotient points.}
In the F.9 tetrahedral coordinate,
\[
 r(z)=\frac{z+i}{z-i},\qquad
 q(z)=z+r(z)+r^2(z)
     =\frac{z^3+3iz+1-i}{(z-1)(z-i)}.
\]
The degree-three quotient $q$ is invariant under $r$.
Its branch points upstairs are
\[
 p_\pm=\frac{(1+i)(1\pm\sqrt3)}2,
\]
and their images are $3p_\pm$.  In the twelve-point
labelling of F.8 the four free $K$-orbits, in order, are
\[
 (0,8,1),\quad(11,3,5),\quad(2,6,9),\quad(7,10,4).
                                                               \tag{F.7.12}
\]
Their representatives have coordinates $z,1/z,-1/z,-z$.
Consequently the unnormalized labelled quotient map is
\[
 z\longmapsto
   (3p_+,3p_-,q(z),q(1/z),q(-1/z),q(-z)).                \tag{F.7.13}
\]
Its entries are distinct on $X$: any equality between two
free quotient points would identify their points in the
$A$ orbit, and a collision with a branch point would give a
nontrivial stabilizer in $A$.

Put $U(z)=(z-p_+)/(z-p_-)$.
Direct substitution gives
\[
       \frac{q(z)-3p_+}{q(z)-3p_-}=U(z)^3.
\]
Normalizing the first, second, and third points in (F.7.13)
to $0,\infty,1$ gives the actual map
\[
 \Sigma:X\longrightarrow M_{0,6},\qquad
 z\longmapsto(s_1(z),s_2(z),s_3(z)),
\]
where
\[
 s_1=\left(\frac{U(1/z)}{U(z)}\right)^3,
 \qquad
 s_2=\left(\frac{U(-1/z)}{U(z)}\right)^3,
 \qquad
 s_3=\left(\frac{U(-z)}{U(z)}\right)^3.                 \tag{F.7.14}
\]
The homomorphism $j$ in (F.7.11) is precisely $\Sigma_*$,
by descent of the same labelled orbit path.

\medskip\noindent\textbf{The second $A_4$ and the exact based involution.}
The involution $\beta$ exchanges the two branch points and
permutes the four free orbits in (F.7.12) by $(0\ 1)(2\ 3)$.
On the normalized moduli coordinates its relabelling is
therefore
\[
 B(s_1,s_2,s_3)=(s_1,s_1/s_3,s_1/s_2).                 \tag{F.7.15}
\]
For example, after swapping the branch points the new
normalizing coordinate is $u\mapsto s_1/u$, since the
first free point has become the old point $s_1$.
This derives (F.7.15) from the actual labelled quotient,
rather than selecting a conjugate linear action.

Use $z_*$ and the path
$\gamma:z_*\to5+z_*\to b=5+4i$ from (F.10.12).
The point $\Sigma(z_*)$ is fixed by $B$.  Let
$\tau=\Sigma\circ\gamma$, with the indicated direction.
For a chronological loop $\ell$ based at $\Sigma(b)$, the
transported involution is exactly
\[
 \beta_b(\ell)=
 \tau^{-1}*B(\tau)*B(\ell)*B(\tau)^{-1}*\tau.          \tag{F.7.16}
\]
The path endpoints in this expression match in order;
$B^2=1$ and $B\Sigma(z_*)=\Sigma(z_*)$ show directly that
$\beta_b^2=1$.  It induces the conjugation action of the
same transported finite mapping class $\beta$ used in F.10
on the injective quotient image (F.7.10).

Since $A$ and $\beta$ generate $A_5$, the exact residual
comparison is
\[
 \widehat j\bigl(C_{\widehat P}(A_5)\bigr)
   =\widehat j(\widehat F_{13})\cap
       \operatorname{Fix}\!\left(
          \beta_b:\widehat{P\Gamma}_{0,6}
                    \longrightarrow\widehat{P\Gamma}_{0,6}
                          \right).                    \tag{F.7.17}
\]
Indeed centralizing $A_5$ is exactly centralizing $A$ and
$\beta$.  For elements centralizing $A$, equation (F.7.11)
is injective; also $\beta$ normalizes $K$, so both the
element and its $\beta$ conjugate lie in the domain of
that injection.  Equality of their quotient images is
therefore equivalent to their equality upstairs.

\medskip\noindent\textbf{Executed exact checks and their scope.}
The program \path{code/part_a/second_a4_c3_attack.py}
uses exact arithmetic in $\Q(i,\sqrt5,\sqrt{\varphi+2})$
to check the $60$ M\"obius matrices, $180$ generator
products, and $720$ vertex-label images.  It verifies all
four orbits in (F.7.12), all representative functions,
the quotient identity, and (F.7.14).  It also checks
the fixed-locus calculation
\[
 \Sigma(z)\in\operatorname{Fix}(B)
 \quad\Longleftrightarrow\quad s_1(z)=s_2(z)s_3(z),
\]
whose numerator, up to a nonzero scalar and after
cancelling the displayed rational denominators, is
\[
 z(z^2+1)(z^8+16z^6+46z^4+16z^2+1).                    \tag{F.7.18}
\]
It verifies exactly that $z_*$ satisfies the relevant
quartic factor.  Finite intersection of these geometric
loci does not by itself calculate the intersection of
their completed fundamental groups in (F.7.17).

The auxiliary numerical braid discovery in
\path{work/part_a/f9_c3_johnson_probe.py} currently gives
no exclusion at the first nonzero outer Johnson degree.
It is explicitly exploratory: the crossings there have
not been certified for the actual based map (F.7.16),
and a higher-degree comparison must retain that path.
No conclusion about (F.7.17) uses that calculation.

\medskip\noindent\textbf{The modular orbifold and the exact $GT_0^1$ image.}
There is a further exact description of the ambient peripheral
centralizer in the minimal twelve-puncture case.  It also gives
five specified three-cusp tests for the remaining geometric
intersection.  The intersection itself remains unresolved.

Let
\[
 \Delta=\langle a,b\mid a^2=b^3=1\rangle,
 \qquad E=\widehat\Delta=C_2\amalg C_3,
 \qquad d=ab,\quad e=ba.                                \tag{F.7.19}
\]
The isomorphism with the full profinite coproduct follows from
the same universal maps on the two finite factors.
Use the epimorphisms
\[
 \begin{array}{c|cc|c}
     &a&b&\text{kernel}\\ \hline
 q_3:E\to S_3 &(2\ 3)&(1\ 2\ 3)&K\\
 q_5:E\to A_5 &(1\ 2)(3\ 4)&(1\ 3\ 5)&H .
 \end{array}                                             \tag{F.7.20}
\]
The discrete modular orbifold has one order-two point,
one order-three point, and one puncture, whose oriented
loop is $d$.  In the first quotient $d$ has order two,
and in the second it has order five.  The corresponding
torsion-free covers have the following exact data:
\[
 \begin{array}{c|c|c|c|c}
  &\text{degree}&\text{cusp width}&\text{number of cusps}
        &\text{surface type}\\ \hline
  K&6&2&3&(0,3)\\
  H&60&5&12&(0,12).
 \end{array}                                             \tag{F.7.21}
\]
Indeed the quotient maps inject both finite orbifold
stabilizers, so the covers are surfaces.  Their Euler
characteristics are the degree times $-1/6$; the cusp
counts are the indices of $\langle q_i(d)\rangle$ in
the finite quotient.  The resulting genera are zero.
Completion of these finite-index inclusions retains their
full profinite topology.  Thus $K\cong\widehat F_2$,
$H\cong\widehat F_{11}$, and the second extension is
the actual centre-free orbifold extension of the
icosahedral twelve-puncture action.

Here is the precise identification with the specified
action in F.7.  Start with that oriented sphere $S$
and its twelve-vertex puncture set $V$.  Choose a
basepoint $s\in S\setminus V$ outside every nontrivial
stabilizer and use its image in the quotient orbifold.
Choose ordinary positive orbifold loops $u_2,u_3$
around the order-two and order-three points, with
whiskers so that the positive puncture loop is
$u_c=(u_2u_3)^{-1}$.  Set
\[
             a=u_3^{-1}u_2u_3,\qquad b=u_3^{-1};
                   \qquad ab=u_c.
\]
Path lifting gives the monodromy map to the actual
deck group $A_5$.  Its generating pair can be labelled
as the pair $q_5(a),q_5(b)$ in (F.7.20), by the exact
single-orbit calculation below.  This relabels the
deck group and does not change its kernel or the
given subgroup of the sphere mapping class group.
Based loop lifting now identifies the discrete
kernel with $\pi_1(S\setminus V,s)$, and its completion
identifies $H$ with the same specified profinite
surface group.  The chosen lift of the cusp whisker
selects a vertex $v_0\in V$ stabilized by $q_5(d)$.
The actual identification of the twelve labels is
\[
        A_5/\langle q_5(d)\rangle\longrightarrow V,
                   \qquad g\langle q_5(d)\rangle
                                    \longmapsto g v_0.
\]
Under it the cusp $[t d^5t^{-1}]_H$ is exactly the
puncture $q_5(t)v_0$.  In particular the identification
of $H$ carries its pure geometric mapping class
subgroup to the one used in F.7.
Changing the base path changes this identification
by an inner surface-group map; changing the chosen
lift changes it by the specified deck action.
An equivalent choice of the oriented orbifold
presentation gives the corresponding based covering
isomorphism.  These maps conjugate the peripheral
outer group and its geometric subgroup together,
and carry the displayed $A_5$ to the same geometric
action.  They therefore carry the intersection
below to the identical residual centralizer under
that covering identification.

For precision, $x=d^2$ and $y=e^2$ are a full free basis
of $K$.  Schreier rewriting with transversal
$1,a,b,ab,ba,b^{-1}$ gives the following kernel words
$t_g s t_{q_3(gs)}^{-1}$:
\[
 \begin{array}{c|cc}
 t_g&s=a&s=b\\ \hline
 1&1&1\\ a&1&1\\ b&1&1\\
 ab&x&y^{-1}\\ ba&1&y\\ b^{-1}&x^{-1}&1 .
 \end{array}                                             \tag{F.7.22}
\]
Rewriting $a^2$ and $b^3$ at these six representatives
gives only freely trivial words in $x,y$; these twelve
relator checks are included in the certificate.
The resulting Schreier presentation is exactly
$\langle x,y\mid\ \rangle$, giving the stated basis
and then its full profinite basis.  Conjugation by $a$ sends
$(x,y)$ to $(y,x)$, while conjugation by $b$ sends it
to $(y,x^{-1}y^{-1})$.  Writing
$\omega(x)=y$, $\omega(y)=(xy)^{-1}$, the latter
automorphism is exactly $\operatorname{Inn}(y)\omega$.
Therefore the outer
deck action on $K$ is exactly the geometric $S_3$
used in D.35, with these peripheral generators.

\emph{Characteristic kernels and the finite normalization.}
There are six generating order-two/order-three pairs in
$S_3$, forming one orbit under its six inner automorphisms.
There are $120$ such generating pairs in $A_5$, forming
one orbit under conjugation by $S_5$.
For the latter count, fix $(1\ 2)(3\ 4)$: the generating
three-cycles use the fifth letter and one letter from
each of its two pairs, giving eight choices.
Their products with the involution have order five.
The remaining choices generate proper subgroups.
The finite certificate below verifies all $300$ pairs,
as well as the transitive $S_5$ action.
Consequently each epimorphism in (F.7.20) has the unique
kernel of its indicated quotient type.  In particular
both $K$ and $H$ are characteristic in $E$.

We also need the exact quotient-automorphism statement.
The five $V_4$ subgroups of $A_5$ form an intrinsic set.
Its conjugation action is faithful.  An automorphism
acting trivially on these five subgroups acts trivially
on that faithful permutation representation, and hence
on $A_5$.  Thus $\operatorname{Aut}(A_5)$ embeds in $S_5$;
all $120$ conjugations by $S_5$ already give automorphisms,
so this is the full automorphism group.  An odd
conjugation exchanges the two $A_5$ conjugacy classes
of five-cycles.  It follows that an automorphism
preserving the conjugacy class of $q_5(d)$ is inner.
The certificate verifies the five $V_4$ groups,
faithfulness, and both twelve-element five-cycle classes.

\emph{The full peripheral centralizers.}
For $T=K,H$, let $\mathcal P_T^1$ denote the subgroup
of $\operatorname{Out}(T)$ preserving each individual
oriented peripheral element up to conjugacy with
multiplier one.  Put
\[
 \Omega_E^1=\{[\eta]\in\operatorname{Out}(E):
                         \eta(d)\sim_E d\}.
\]
There are canonical continuous isomorphisms, with
the normalizations described below,
\[
 C_{\mathcal P_H^1}(A_5)
       \xrightarrow{\ \sim\ }\Omega_E^1
       \xrightarrow{\ \sim\ }
 C_{\mathcal P_K^1}(S_3)=GT_0^1,
 \qquad GT_0^1=\ker(\lambda:GT_0\to\widehat\Z^\times).
                                                               \tag{F.7.23}
\]
We prove the maps rather than identify the two actions
only by their surface types.

First, for $m=2$ and $T=K$, or for $m=5$ and $T=H$,
\[
                         C_E(d^m)=\overline{\langle d\rangle}
                                      \cong\widehat\Z . 
                                                               \tag{F.7.24}
\]
The element $d^m$ is a primitive cusp of the free group
$T$, so $C_T(d^m)=\overline{\langle d^m\rangle}$.
This follows, for example, by writing $T$ as the
free profinite product of that primitive cyclic factor
and its free complement and using factor malnormality
as in D.35.  The quotient action on the cusp classes
is the actual action on
$Q/\langle q(d)\rangle$, where $Q=S_3$ or $A_5$.
An element centralizing $d^m$ therefore has quotient
in $\langle q(d)\rangle$.  Multiplying it by a suitable
integral power of $d$ puts it in $T$ without changing
centralization.  This proves the equality in (F.7.24).
The procyclic group is a full $\widehat\Z$ because
$d^m$ already generates such a subgroup: a relation
$d^u=1$ would imply $(d^m)^u=1$ and hence $u=0$.

The centralizer $C_E(T)$ is trivial.  Its image in
$Q$ would act trivially in $\operatorname{Out}(T)$,
but that action is faithful on the cusp permutation
set.  Thus this image is trivial, and the remaining
element lies in $Z(T)=1$.
Consequently conjugation embeds $E$ in
$\operatorname{Aut}(T)$ as the full inverse image
of the specified $Q\leq\operatorname{Out}(T)$.

Let $[\phi]\in C_{\mathcal P_T^1}(Q)$.
Conjugation by an actual representative $\phi$ on
$\operatorname{Aut}(T)$ preserves that inverse image
and gives an automorphism $\eta$ of $E$ inducing
the identity on $Q$ and restricting to $\phi$ on $T$.
Changing $\phi$ by an inner automorphism of $T$
changes $\eta$ by the same inner automorphism of $E$.
Also $\eta(d)^m$ is conjugate in $T$ to $d^m$.
After that conjugation, (F.7.24) places $\eta(d)$
in $\overline{\langle d\rangle}$; the equation
$m u=m$ in $\widehat\Z$ forces its exponent to be
one.  Thus $[\eta]\in\Omega_E^1$.

Conversely, choose $[\eta]\in\Omega_E^1$.
Its induced automorphism of $Q$ is inner: for
$S_3$ this follows from its full automorphism group;
for $A_5$ it follows from preservation of the
five-cycle class proved above.
If that quotient automorphism is $\operatorname{Inn}(q_0)$,
choose $u\in E$ with $q(u)=q_0^{-1}$ and replace
$\eta$ by $\eta_0=\operatorname{Inn}(u)\eta$.
The new automorphism induces the identity on $Q$.
Different choices of $u$ differ by an element of $T$,
so the outer class of its restriction to $T$ is
well defined.  Changing the representative of
$[\eta]$ has the same inner effect after normalization;
the centre of $Q$ is trivial, so there is no further
quotient ambiguity.

Write $\eta_0(d)=rdr^{-1}$.  Because the quotient
action is now the identity, $q(r)$ centralizes
$q(d)$.  In both finite groups that centralizer is
exactly $\langle q(d)\rangle$.  Hence $r=t d^j$
for $t\in T$, and $\eta_0(d)=tdt^{-1}$.
For every $s\in E$, the elements $\eta_0(s)$ and
$s$ have the same quotient.  Therefore the restriction
preserves each $T$-cusp class
$[s d^m s^{-1}]_T$ with exponent one.
It belongs to $\mathcal P_T^1$ and centralizes the
outer $Q$ action.

These constructions are inverse homomorphisms.
For injectivity, a normalized automorphism whose
restriction to $T$ is inner can first be adjusted
to fix $T$ pointwise.  Then
$\eta(s)s^{-1}\in C_E(T)=1$ for every $s\in E$,
so it is the identity.  In the other direction,
an inner automorphism of $E$ arising from a
centralizing outer class of $T$ has quotient in
$Z(Q)=1$, hence its conjugator lies in $T$.
Continuity follows from conjugation on the specified
closed subgroup of $\operatorname{Aut}(T)$ and
normalization by the finite quotient.  The maps
are homeomorphisms because all the groups here
are compact Hausdorff.
Finally the equality with $GT_0^1$ is precisely the
Harbater--Schneps centralizer identification used
in D.35, now with the exact basis and outer actions
in (F.7.22)\cite[Propositions 3--4 and Theorem 2]{R-HS}.

Let $\iota:GT_0^1\hookrightarrow\mathcal P_H^1$
denote the inverse of the left-to-right isomorphism
in (F.7.23).  Genus-zero congruence identifies the
completed pure geometric group with its actual
image in $\mathcal P_H^1$.  Thus the remaining group
has the exact description
\[
             C_{\widehat{P\Gamma}_{0,12}}(A_5)
                 =\widehat{P\Gamma}_{0,12}
                                \cap\iota(GT_0^1)
                 \quad\text{inside }\mathcal P_H^1.     \tag{F.7.25}
\]
This proves the relationship with the full $GT_0$
problem; it does not calculate this intersection.

\emph{Explicit substitutions and three-cusp quotient maps.}
Write $\gamma=(1,f)\in GT_0^1$ in its normalized
representative on $K$, so
$x\mapsto x$, $y\mapsto f^{-1}yf$ and $f\in K'$.
Its actual extension to (F.7.19) has the particularly
simple formulas
\[
          \eta_f(a)=f^{-1}a,\qquad
          \eta_f(b)=f^{-1}b,\qquad \eta_f(d)=d .         \tag{F.7.26}
\]
Indeed the first symmetry is $\theta(f)=f^{-1}$.
Direct conjugation on $K$ gives
$\gamma\theta=\operatorname{Inn}(f^{-1})\theta\gamma$,
so $C_E(K)=1$ forces the stated image of $a$.
Since $\gamma(x)=x$, (F.7.24) with $m=2$ forces
$\eta_f(d)=d$; the displayed image of $b=a^{-1}d$
then follows.  These formulas describe the actual
automorphism furnished by (F.7.23), including its
relations and its invertibility.

The quotient automorphism on $A_5$ fixes $q_5(d)$
and is inner.  There is consequently a unique
$j\in\{0,1,2,3,4\}$ for which it is
$\operatorname{Inn}(q_5(d)^j)$; equivalently
\[
             q_5(f)=q_5(e)^j q_5(d)^{-j}.
\]
The pure normalized representative on $H$ is
\[
       \eta_f^H=
          \left.\operatorname{Inn}(d^{-j})\eta_f\right|_H.
                                                               \tag{F.7.27}
\]
For example, with $c_0=d^5$ and $c_1=a d^5a^{-1}=e^5$,
put $h_1=d^{-j}f^{-1}e^j\in H$.  Then
\[
                 c_0\longmapsto c_0,
        \qquad c_1\longmapsto h_1c_1h_1^{-1}.           \tag{F.7.28}
\]
This uses the finite quotient to ensure $h_1\in H$;
it is not a formal substitution of a word outside
the source of a filling homomorphism.

Choose the cusp labels as the twelve left cosets
$A_5/\langle q_5(d)\rangle$, in the breadth-first
order specified by the certificate.  The corresponding
representatives $t_i\in\Delta$ and peripheral words
$c_i=t_i d^5t_i^{-1}$ are recorded there; $t_0=1$
and $t_1=a$.
For a three-element label set $J$ define the actual
continuous filling homomorphism
\[
 q_J:H\longrightarrow H_J
     =H/\langle\!\langle c_i:i\notin J\rangle\!\rangle_H
                         \cong\widehat F_2.            \tag{F.7.29}
\]
The isomorphism type follows from filling precisely
those nine punctures of the specified sphere;
completion preserves the displayed quotient
presentation.  Since (F.7.27) is pure peripheral,
it descends to $H_J$.  Thus there is a continuous
homomorphism
\[
 \mathcal D_J:GT_0^1\longrightarrow\mathcal P_{H_J}^1,
 \qquad\mathcal D_J(\gamma)=
 \bigl[q_J(h)\longmapsto q_J(\eta_f^H(h))\bigr].       \tag{F.7.30}
\]
The brackets on the right specify the outer
automorphism induced by the displayed substitution.
Any change of representative in (F.7.27) changes
it only by an inner automorphism of $H_J$.
Formula (F.7.26), the finite determination of $j$,
and (F.7.29) define each map without assuming that
any two arbitrarily based retained cusp words
form a full free basis.

The pure geometric mapping class group of the
three-punctured sphere is trivial.  The actual
forgetting homomorphism on pure mapping classes,
its action on the filled fundamental group, and
continuity on completions therefore give
\[
 \iota^{-1}\!\left(
       C_{\widehat{P\Gamma}_{0,12}}(A_5)\right)
               \subseteq\bigcap_{|J|=3}\ker\mathcal D_J.
                                                               \tag{F.7.31}
\]
For $g\in A_5$, a lift in $E$ induces the actual
quotient isomorphism $H_J\to H_{gJ}$ by conjugation;
different lifts differ by an inner automorphism.
Because $\iota(\gamma)$ centralizes $A_5$ in
$\operatorname{Out}(H)$, these isomorphisms intertwine
$\mathcal D_J(\gamma)$ and $\mathcal D_{gJ}(\gamma)$.
Their kernels are consequently equal within each
actual $A_5$ orbit of three-element subsets.

The $220$ triples form exactly five orbits, with
the following representatives in the recorded coset
labelling:
\[
 \begin{array}{c|ccccc}
 J&\{0,1,2\}&\{0,1,4\}&\{0,1,8\}&\{0,1,10\}&\{0,4,9\}\\
 |A_5J|&20&60&60&60&20 .
 \end{array}                                             \tag{F.7.32}
\]
The certificate supplies a labelled $A_5$ transporter,
as a word in $a,b$, for every triple.  Hence the
intersection of $220$ kernels in (F.7.31) equals the
intersection of these five specified kernels.
No injectivity or triviality of that intersection
has been proved here.

The standard-library executable
\path{code/part_a/modular_a5_gt0_f7.py} has been run.
It checks the six and $120$ epimorphism pairs and
their single postcomposition orbits, both faithful
cusp actions, both cyclic cusp centralizers, the
five-$V_4$ automorphism argument, the exact twelve
Schreier rewritings in (F.7.22), all twelve peripheral
words, and all $220$ subset transporters in (F.7.32).
These finite calculations certify the finite inputs
and the displayed words.  The centre-free extension,
normalization, completion, and filling maps are
proved above.


\medskip\noindent\textbf{Exact correspondences for the two orbit-twenty detectors.}
The two orbit-twenty tests in (F.7.32) admit the following
explicit subgroup and filling description. It gives
continuous word substitutions and a common section of
the surface-group quotient maps. Faithfulness of the
outer-automorphism detectors remains unresolved.

Use the precise $E,q_5,H,K,a,b,d,e$ of (F.7.19)--(F.7.22).
Put $T=\langle q_5(b)\rangle$ and $B=q_5^{-1}(T)$.
The index-twenty orbifold cover has two order-three
points and four width-five cusps, and no order-two
points. This is the exact coset action: $a$ has ten
two-cycles, $b$ has two fixed cosets and six three-cycles,
and $d$ has four five-cycles. Its signature is
$(0;3,3;4\text{ cusps})$.

The following words give a specified full profinite
presentation:
\[
 \begin{split}
 z&=aba b^{-1},\qquad \tau=zaz^{-1},\\
 r&=b,\qquad s=\tau b\tau^{-1},\\
 u&=(ab^{-1})^4a,\qquad v=(ab)u(ba),\\
 w&=(aba b^{-1}a)(ba)^3(b^{-1}aba).
 \end{split}
 \qquad
 B=\langle r,s,u,v,w\mid r^3=s^3=1\rangle_{\rm prof}.
                                                               \tag{F.7.33}
\]
Thus $B=C_3\amalg C_3\amalg\widehat F_3$ with the
displayed inclusion into $E$. The exact Schreier
certificate uses twenty right cosets $Tg$ and the
prefix-closed breadth-first transversal in $a,b$.
Its forty directed generator edges have nineteen
tree edges, leaving twenty-one Schreier generators.
Rewriting $a^2,b^3$ at the twenty vertices gives
forty relators. Sixteen Tietze eliminations, each
solving a relator in a generator occurring exactly
once, leave the five displayed generators and
precisely the two cube relators.
The complete transversal, original generator words,
substitutions, and eliminations are recorded.
Reduction in $C_2*C_3$ checks the compact words in
(F.7.33). These are inverse presentation maps;
completion of the finite-index discrete inclusion
retains its full profinite topology.

The four based cusp loops and their actual
$T$-orbits of $H$ punctures are
\[
 \begin{array}{c|c|c}
 i&\kappa_i\text{ in }B&J_i\\ \hline
 0&u^{-1}r&\{0,1,2\}\\
 1&v^{-1}&\{3,4,6\}\\
 2&wu&\{5,7,8\}\\
 3&vw^{-1}s&\{9,10,11\}
 \end{array}                                                \tag{F.7.34}
\]
Each is the rewritten word $t_i d^5t_i^{-1}$
for its recorded transversal representative.
The first and fourth triples lie in the
$A_5$-orbit of $\{0,1,2\}$; the middle two lie
in the orbit of $\{0,4,9\}$.

Fill the other three orbifold cusps and put
$B_i=B/\langle\!\langle\kappa_j:j\ne i\rangle\!\rangle_B$.
To
$A=\langle R,S\mid R^3=S^3=1\rangle_{\rm prof}=C_3\amalg C_3$
send $r$ to $R$, $s$ to $S$, and use the indicated row:
\[
 \begin{array}{c|ccc|c}
 i&u&v&w&\kappa_i\\ \hline
 0&S^{-1}&1&S&SR\\
 1&R&S^{-1}R^{-1}&R^{-1}&RS\\
 2&R&1&S&SR\\
 3&R&1&R^{-1}&RS
 \end{array}                                                \tag{F.7.35}
\]
Substitution kills the three required cusp relators.
Conversely those three relators solve for $u,v,w$
as the entries of the row. Inclusion of $r,s$ in
the quotient gives the inverse map from $A$.
This proves the exact quotient isomorphism and
identifies the retained positive cusp as the
displayed product without a power ambiguity.

Let $E^+=\ker(\epsilon:E\to C_2)$, where
$\epsilon(a)=1,\epsilon(b)=0$.
The two-coset Schreier calculation with transversal
$1,a$ gives
\[
 A\xrightarrow{\ \sim\ }E^+,\qquad
                  R\longmapsto b,\quad S\longmapsto aba.
                                                               \tag{F.7.36}
\]
Write $p_i:B\to E^+$ for the composite of the row
(F.7.35) and (F.7.36).
The quotient $B/H=T$ sends
\[
 (r,s,u,v,w)\longmapsto
       (\bar b,\bar b^{-1},\bar b,1,\bar b^{-1}),
\]
whereas $E^+/K=C_3$ sends $b$ to its generator
and $aba$ to the inverse. Identify these two cyclic
groups by their named generators.
All four rows intertwine the specified quotient maps,
and hence $p_i(H)=K$.

The restriction to $H$ is precisely the filling
of the nine punctures outside $J_i$.
Indeed the closed $B$-normal closure of the other
three cusp words is contained in $H$.
Conjugation by $B/H=T$ supplies the three $H$ cusp
classes in each omitted orbit, and further
$H$ conjugation gives their closed normal closure.
Thus
\[
 \begin{split}
 \langle\!\langle\kappa_j:j\ne i\rangle\!\rangle_B
 &=\langle\!\langle c_\ell:\ell\notin J_i
                                      \rangle\!\rangle_H,\\
 1\longrightarrow H_{J_i}&\longrightarrow B_i
                 \longrightarrow C_3\longrightarrow1,
       \qquad H_{J_i}\xrightarrow[\ p_i\ ]{\ \sim\ }K .
 \end{split}                                                \tag{F.7.37}
\]
Each omitted cusp has width one for $H\leq B$:
$\kappa_j\in H$ and its three lifts have precisely
the labels in (F.7.34). No peripheral power is
inserted into this normal-closure identity.

There is a common actual section. Define
\[
 \iota_0:E\longrightarrow E,\qquad
                a\longmapsto\tau,\quad b\longmapsto b .
                                                               \tag{F.7.38}
\]
The orders give a homomorphism. On
$E^+=\langle b,aba\rangle$ it maps the two free
profinite factors to $r,s$, whose inclusion in
$B$ is injective by its retraction onto those factors.
Moreover $q_5(\tau)$ is an involution inverting
$q_5(b)$, so $q_5\iota_0(E)$ is the order-six
normalizer of $T$. Its quotient by $T$ is
exactly $\epsilon$. Thus $\ker\iota_0$ is contained
in $E^+$, where it is trivial, proving that
(F.7.38) is a continuous embedding.

The kernel of
$E\to q_5\iota_0(E)\cong S_3$ is the characteristic
$K$ proved in (F.7.20). Hence $\iota_0(K)\leq H$.
Every $p_i$ is the identity on the displayed two
factors after (F.7.36), so
\[
 p_i\iota_0|_{E^+}=1_{E^+},\qquad
             p_i|_H\circ\iota_0|_K=1_K.                 \tag{F.7.39}
\]
The program also rewrites $\iota_0(x),\iota_0(y)$
through all four quotient rows and verifies that
their images are exactly $x,y$.

For $\gamma=(1,f)\in GT_0^1$, the actual normalized
automorphism (F.7.27) extends to an automorphism of
$E$ inducing the identity on $A_5$. It preserves
$B$ and each filled normal subgroup.
Using (F.7.37) for the target identification, the
actual detector is
\[
 \mathcal D_{J_i}(\gamma)
        =\bigl[p_i\,\eta_f^H\,\iota_0|_K\bigr]
                          \quad\hbox{in }\mathcal P_K^1.
                                                               \tag{F.7.40}
\]
Here $\eta_f^H$ is applied to $\iota_0(K)\leq H$.
Its preservation of that embedded subgroup is
not assumed. The subsequent quotient map is
part of the composition. The composition is
an automorphism because it is the actual descended
automorphism under (F.7.37), evaluated on the
section (F.7.39).
The two pairs of equal detector kernels are
the actual orbit transports of (F.7.31)--(F.7.32).

The four source quotient maps themselves have
a nontrivial common kernel:
\[
 [\kappa_0,\kappa_1]
           =u^{-1}r v^{-1}r^{-1}uv\ne1.                 \tag{F.7.41}
\]
Every row kills this word.
Its nontriviality in the full completion is
detected by $B\to S_3$ sending $r,s,w$ to $1$,
$u$ to $(1\ 2)$, and $v$ to $(2\ 3)$:
its image has order three.
This source word gives an inner automorphism
of $E$, whose class in $\operatorname{Out}(E)$
is the identity. It therefore supplies no
nonidentity $GT_0^1$ element in the kernel of
(F.7.40). Faithfulness of an outer detector,
or of the five-detector family, remains open.

The executed program
\path{code/part_a/modular_c3_correspondence_f7.py}
records the full Schreier presentation, its
sixteen eliminations, the four three-step filling
eliminations, compact word identities, exact
finite quotient labels, eight common-section
word checks, and the finite detector in (F.7.41).


\medskip\noindent\textbf{An exact normalized detector and its naturality equation.}
The first correspondence in (F.7.35) gives a particularly
short formula for the orbit-twenty detector.  All maps
in this paragraph are the based maps already fixed in
(F.7.19)--(F.7.40).  Thus $x=d^2$, $y=e^2$ is the
specified free basis of $K$, and $p_0:B\to E^+$ is the
row retaining $J_0=\{0,1,2\}$.

For $\gamma=(1,f)\in GT_0^1$, let $j\in\{0,\ldots,4\}$
be the actual finite correction in (F.7.27), and put
\[
 \begin{split}
 W_f&=p_0(d^{-j}f^{-1}e^j)\in K,\\
 A_f&=\ell_x(W_f),\qquad B_f=\ell_y(W_f),\\
 F_0(f)&=y^{B_f}W_f^{-1}x^{A_f}\in K'.
 \end{split}                                                   \tag{F.7.42}
\]
Here $\ell_x,\ell_y:K\to\widehat\Z$ are the
coordinates of the displayed abelianization.
The source of $p_0$ in this formula is legitimate:
$d^{-j}f^{-1}e^j=h_1$ lies in $H\leq B$ by
(F.7.27)--(F.7.28).
The exact normalized representative of the detector is
\[
 \mathcal D_{J_0}(\gamma)
       =\bigl[(x\mapsto x,\quad
                   y\mapsto F_0(f)^{-1}yF_0(f))\bigr]
                    \quad\hbox{in }\mathcal P_K^1.              \tag{F.7.43}
\]
In particular this is a formula in the actual target
pure peripheral group; membership of its normalized
coordinate in $GT_0^1$ is not used.

To prove the formula, use the two actual $H$ cusp
words $c_0=d^5$ and $c_1=e^5=b c_0b^{-1}$.
The row in (F.7.35) and the identification (F.7.36)
give
\[
       p_0(c_0)=(aba)b=x,\qquad
       p_0(c_1)=bxb^{-1}=y .
\]
These images are the previously proved full basis,
so (F.7.28) shows that the actual descended
automorphism is
$\psi_f(x)=x$, $\psi_f(y)=W_fyW_f^{-1}$.
This agrees with $p_0\eta_f^H\iota_0|_K$ because
$p_0\iota_0|_K=1_K$ and $\psi_fp_0=p_0\eta_f^H$
on $H$.  Composing $\psi_f$ with
$\operatorname{Inn}(x^{-A_f})$ gives the representative
in (F.7.43), since
$F_0(f)^{-1}yF_0(f)=
x^{-A_f}W_fyW_f^{-1}x^{A_f}$.
The abelianization of $F_0(f)$ is zero.
Uniqueness is the normalized-pair uniqueness:
two representatives fixing $x$ differ by conjugation
in $\overline{\langle x\rangle}$, and equality of the
images of $y$ then leaves a factor in
$\overline{\langle y\rangle}$; the two zero
abelian coordinates remove both factors.

Consequently the kernel and the equality comparison
with the original normalized action have the exact
descriptions
\[
 \begin{split}
 \ker\mathcal D_{J_0}
   &=\{\,(1,f)\in GT_0^1:
                     W_f=x^{A_f}y^{B_f}\,\},\\
 \{\gamma\in GT_0^1:\mathcal D_{J_0}(\gamma)=\gamma
                                      \text{ in }\mathcal P_K^1\}
   &=\{\,(1,f)\in GT_0^1:
                     y^{B_f}W_f^{-1}x^{A_f}=f\,\}.
 \end{split}                                                   \tag{F.7.44}
\]
For the first equality one can also argue directly.
An inner automorphism fixing $x$ is conjugation
by $x^t$.  Its equality with $\psi_f$ is exactly
$x^{-t}W_f\in C_K(y)=\overline{\langle y\rangle}$,
which forces $t=A_f$ and the remaining exponent
to be $B_f$.
No assertion that the second set in (F.7.44)
is all of $GT_0^1$ is made.

The normalized formula is independent of the lift
of the finite correction.  Replacing $j$ by
$j+5m$, with $m\in\widehat\Z$, gives the exact
identities
\[
 h_1\longmapsto c_0^{-m}h_1c_1^m,\qquad
 W_f\longmapsto x^{-m}W_fy^m,\qquad
 (A_f,B_f)\longmapsto(A_f-m,B_f+m).
                                                               \tag{F.7.45}
\]
Substitution in (F.7.42) leaves $F_0(f)$ unchanged.
This also proves explicitly that this normalization
absorbs the permissible $H$-inner change.

For completeness the exact naturality defect can
be written before any filling.  Put
$\eta_j=\operatorname{Inn}(d^{-j})\eta_f$ and
\[
 L_f=x f^{-1}b,\qquad
 T_f=L_f f^{-1}a L_f^{-1}.
\]
Since $z=xb$, these are $\eta_f(z)$ and
$\eta_f(\tau)$, respectively.  The actual two
basis values are
\[
 \begin{split}
 \eta_j\iota_0(x)
     &=d^{-j}(T_f f^{-1}b)^2d^j,\\
 \eta_j\iota_0(y)
     &=d^{-j}(f^{-1}bT_f)^2d^j .
 \end{split}                                                   \tag{F.7.46}
\]
Define two elements of $H$ by
\[
 \begin{split}
 \Delta_x(f)&=\eta_j\iota_0(x)\,\iota_0(x)^{-1},\\
 \Delta_y(f)&=\eta_j\iota_0(y)
              \bigl(\iota_0(f)^{-1}
                       \iota_0(y)\iota_0(f)\bigr)^{-1}.
 \end{split}
\]
For every one of the four exact filling maps, the
actual representative of (F.7.40) satisfies
\[
 \begin{split}
 (p_i\eta_f^H\iota_0)(x)&=p_i(\Delta_x(f))\,x,\\
 (p_i\eta_f^H\iota_0)(y)&=
                   p_i(\Delta_y(f))\,f^{-1}yf .
 \end{split}                                                   \tag{F.7.47}
\]
For $i=0$ the first defect is killed, by the
$c_0$ computation above.  The second comparison,
including its permitted inner normalization, is
exactly (F.7.44).  A section of the quotient map
supplies the evaluations in (F.7.47); it does not
supply an additional intertwining identity.

There is also an explicit finite certificate for
the size of that section's image.  The same
$\iota_0$ maps $K$ into $K$, with exact word values
\[
 \begin{split}
 \iota_0(x)&=xyx^{-1}y^{-1}x^{-1}yx,\\
 \iota_0(y)&=yx^{-1}y^{-1}xy^{-1}x^{-1}.
 \end{split}                                                   \tag{F.7.48}
\]
Reduction in $C_2*C_3$ verifies both identities.
Their abelianization matrix is
$\left(\begin{smallmatrix}0&-1\\1&-1\end{smallmatrix}\right)$,
so the induced endomorphism of each free pro-$p$
quotient is an automorphism by its Frattini
criterion and the Hopf property.  Nevertheless
$q_5\iota_0(K)=1$, whereas $q_5(K)=A_5$:
the images of $x=d^2$ and $y=e^2$ generate $A_5$,
as the finite certificate verifies.
Thus $\iota_0|_K$ is an injective, proper full
profinite self-embedding with the displayed
finite quotient obstruction to surjectivity.
Its common-section property is fully compatible
with this exact comparison.

The executable
\path{code/part_a/modular_detector_defect_f7.py}
checks the two based cusp images, all four common
section basis images, the words (F.7.48), their
$S_3$ and $A_5$ images, and the normalization,
kernel, and lift-independence identities on all
$4\,373$ reduced words of length at most seven.
These word checks verify the displayed formulas;
they make no faithfulness assertion about
(F.7.43) or the full five-detector family.


\medskip\noindent\textbf{One cone conjugator and all five exact filling equations.}
The kernel equation (F.7.44) has an exact formulation
using a single element of $H$.  The full profinite
generation assertion needed in this formulation
follows from the actual degree-one cusp of a finite
orbifold cover.

First, every closed subgroup $L\leq E$ satisfying
\[
              d\in L,\qquad d^6\in\overline{[L,L]}
                                                               \tag{F.7.49}
\]
is $E$.  Indeed a proper such $L$ is contained in a
proper open subgroup $U$: separate a point outside
$L$ in a finite quotient and take the inverse image
of the image of $L$.  The subgroup
$U\cap\Delta$ has the same finite index, and its full
completion is $U$.  In the corresponding connected
orbifold cover, $d\in U\cap\Delta$ gives a width-one
cusp through the base point.  A proper cover has
at least one more cusp, since the remaining cosets
cannot lie in that length-one cycle of $d$.
Its cone orders divide two or three.
For any prime $p>3$, the cone classes vanish in
$H_1(U,\F_p)$; the cusp classes have precisely their
total-sum relation, together with independent
handle classes.  Each cusp class is therefore
nonzero.  The class of $d^6$ is six times the
specified width-one cusp and is nonzero, contradicting
(F.7.49).

It follows that every continuous endomorphism
$\rho:E\to E$ satisfying
$\rho(d)=t d^\mu t^{-1}$ for a full unit
$\mu\in\widehat\Z^\times$ is an automorphism.
The compact image $L=t^{-1}\rho(E)t$ contains
$d^\mu$ and hence its $\mu^{-1}$ power $d$.
Since $E^{\rm ab}=C_2\times C_3$, the element
$d^6$ belongs to the closed commutator subgroup
of $E$.  Thus $d^{6\mu}\in\overline{[L,L]}$,
and its $\mu^{-1}$ power gives (F.7.49).
Surjectivity follows, and finite generation of $E$
gives Hopficity.  Both the conjugation and the
unit-root operation here are performed on the
specified image subgroup; no pro-$p$ quotient
replaces $E$.

Define the following compact solution space:
\[
 \mathscr S=
 \{\,g\in H:(d g b^{-1}g^{-1})^2=1\,\},\qquad
 \eta_g(d)=d,\quad
 \eta_g(b)=g b g^{-1},\quad
 \eta_g(a)=d g b^{-1}g^{-1}.
                                                               \tag{F.7.50}
\]
The equation gives the $a^2$ relation, and the
$b^3$ relation holds by conjugation.  Hence the
displayed substitutions define an endomorphism
of $E$, and the preceding full boundary test makes
it an automorphism.  Its action on $A_5$ is the
identity because $g\in H$.  Consequently its outer
class is the element of $GT_0^1$ specified by
(F.7.23).

Every element of that $GT_0^1$ has this form.
For $\gamma=(1,f)$, use the actual
$\eta_j=\operatorname{Inn}(d^{-j})\eta_f$ in
(F.7.27).  The full free profinite product torsion
theorem makes $\eta_j(b)$ conjugate to $b$ or
$b^{-1}$; these are the only order-three factor
elements in $C_2\amalg C_3$
\cite[Theorem A and Section 3]{A-HerfortRibes1985}.
Fixing $d$ forces the identity on
$E^{\rm ab}=C_6$, which excludes $b^{-1}$.
Choose $u$ with $\eta_j(b)=ubu^{-1}$.
Since the quotient action is the identity,
$q_5(u)\in C_{A_5}(q_5(b))=\langle q_5(b)\rangle$.
Write $q_5(u)=q_5(b)^r$ and replace $u$ by
$g_f=ub^{-r}\in H$.  The same source's full
free-factor centralizer theorem gives
$C_E(b)=\langle b\rangle$
\cite[Theorem B and Section 4]{A-HerfortRibes1985}.
As $H\cap\langle b\rangle=1$, this $g_f$ is unique.
The conjugation map $H\to E$, $g\mapsto gbg^{-1}$,
is an embedding of a compact space and thus has
a continuous inverse on its image.
We have the exact identities
\[
 \eta_j=\eta_{g_f},\qquad
 h_1=g_f b g_f^{-1}b^{-1},\qquad
 f=b d^j g_f b^{-1}g_f^{-1}d^{-j}.
                                                               \tag{F.7.51}
\]
The last formula retains the particular finite
$j$ selected for the original normalized pair;
it follows from $\eta_f=\operatorname{Inn}(d^j)\eta_j$.

Two solutions $g,h\in\mathscr S$ have the same
outer class exactly when
\[
                   h=d^{5t}g\quad(t\in\widehat\Z).
 \qquad
 \mathscr S/\bigl(g\sim d^{5t}g\bigr)
                         \xrightarrow{\ \sim\ }GT_0^1.
                                                               \tag{F.7.52}
\]
To prove the assertion, an inner comparison of
$\eta_g,\eta_h$ must centralize $d$, so its
conjugator lies in $\overline{\langle d\rangle}$
by (F.7.24).  Both quotient actions are the
identity; since $Z(A_5)=1$, that conjugator is
$d^{5t}\in H$.  Comparing the two images of $b$
then gives $h=d^{5t}g$, because the possible
remaining factor lies in
$H\cap C_E(b)=1$.  Conversely left multiplication
by $d^{5t}$ changes $\eta_g$ by precisely
$\operatorname{Inn}(d^{5t})$ and preserves
(F.7.50).
This proves the bijection, and compactness gives
the stated homeomorphism of the solution-orbit
space with the compact group on the right.

For the first detector the exponent condition
sharpens to one procyclic condition:
\[
 \begin{split}
 \ker\mathcal D_{J_0}
   &=\{\, (1,f):W_f=x^{A_f}y^{-A_f}\,\}\\
   &=\{\, (1,f):p_0(g_f)\in
                            \overline{\langle x\rangle}\,\},\\
 \mathscr S\cap\ker(p_0|_H)
   &\xrightarrow{\ \sim\ }\ker\mathcal D_{J_0}.
 \end{split}                                                   \tag{F.7.53}
\]
For the proof let $\varphi_g$ be the actual
automorphism of $E^+$ descended from $\eta_g$
through $p_0$.  It sends $b$ to
$p_0(g)b p_0(g)^{-1}$, and restricts to the
surface automorphism in (F.7.43).
When its outer restriction is the identity,
that restriction is $\operatorname{Inn}(x^{A_f})$,
by (F.7.44).  Since $C_{E^+}(K)=1$, two
automorphisms of $E^+$ with this same restriction
are equal: their quotient fixes $K$ pointwise,
and its difference on each element centralizes
$K$.  It follows that
$\varphi_g=\operatorname{Inn}(x^{A_f})$.
Comparing the images of $b$ gives
$p_0(g)=x^{A_f}$, since
$K\cap C_{E^+}(b)=1$.
Therefore
$W_f=x^{A_f}b x^{-A_f}b^{-1}=x^{A_f}y^{-A_f}$.
Conversely this expression is exactly the inner
kernel equation (F.7.44).
Finally (F.7.52) replaces a solution with
$p_0(g)=x^t$ uniquely by $d^{-5t}g$ with
$p_0(d^{-5t}g)=1$.  This gives the last
homeomorphism in (F.7.53).
The group law on this slice is explicitly
$g\star h=\eta_g(h)g$: here $\eta_g$ descends
to the identity through $p_0$, so the product
still lies in the slice.

The other four representative fillings can be
written in the same actual coordinates.
An independent index-sixty Schreier calculation
gives $H$ with $61$ initial generators and
$120$ rewritten relators.  Fifty single-occurrence
Tietze eliminations leave eleven generators and
no relators.  For each of the five triples in
(F.7.32), the nine omitted cusp relators give
nine more such eliminations, leaving a full
free basis $P,Q$ with the following values:
\[
 \begin{array}{c|c|ccc|cc}
 J&m&c_0&c_m&c_{\rm third}&P\text{ in }x,y&Q\text{ in }x,y\\ \hline
 \{0,1,2\}&1&QP&Q^{-1}&P^{-1}&yx&y^{-1}\\
 \{0,1,4\}&1&Q&P^{-1}&Q^{-1}P&y^{-1}&x\\
 \{0,1,8\}&1&QP&Q^{-1}&P^{-1}&yx&y^{-1}\\
 \{0,1,10\}&1&Q&P^{-1}&Q^{-1}P&y^{-1}&x\\
 \{0,4,9\}&4&P&P^{-1}Q&Q^{-1}&x&xy
 \end{array}                                                   \tag{F.7.54}
\]
The last two columns are inverse substitutions:
they identify $c_0,c_m$ with the specified full
basis $x,y$ of $K$.
All omitted cusp words go to one.
Thus the resulting maps
$\pi_J:H\to K$ are precisely the actual fillings
(F.7.29) followed by these based quotient
isomorphisms.  The complete original generators,
their words in $a,b$, all relators, all
eliminations, and both inverse basis maps
are recorded in the certificate.
For $J_0$, this is exactly $p_0|_H$,
as both identifications send $c_0,c_1$ to $x,y$.

For $g\in\mathscr S$ put $h(g)=g b g^{-1}b^{-1}\in H$.
Use the original $y=e^2\in E$ in the following
source word:
\[
 \begin{array}{c|c}
 J& V_J(g)\\ \hline
 \{0,1,2\},\{0,1,4\},\{0,1,8\},\{0,1,10\}
       &\pi_J(h(g))\\
 \{0,4,9\}&
          \pi_J\bigl(h(g)y h(g)^{-1}y^{-1}\bigr).
 \end{array}                                                   \tag{F.7.55}
\]
The whole word in the last row lies in $H$,
because $H$ is normal; the map $\pi_J$ is
applied to that word and is not evaluated
separately on $y$ outside its source.
To check the row, the exact label-four
whisker is $(ba)^2=y$, so $c_4=yc_0y^{-1}$.
The substitutions (F.7.50) fix $d$ and $x$,
and give $\eta_g(y)=h(g)y h(g)^{-1}$.
Consequently the conjugator for $c_4$ is
$h(g)y h(g)^{-1}y^{-1}$.
For the first four rows the conjugator for
$c_1=e^5$ is $h(g)$, as in (F.7.51).
Every descended map therefore fixes $x$ and
sends $y$ to $V_J(g)yV_J(g)^{-1}$.
If $A_J=\ell_x(V_J(g))$ and
$B_J=\ell_y(V_J(g))$, all five exact kernel
equations are
\[
 \begin{split}
 \mathcal D_J([\eta_g])&=
       [\,x\mapsto x,\ 
          y\mapsto F_J(g)^{-1}yF_J(g)\,],\\
 F_J(g)&=y^{B_J}V_J(g)^{-1}x^{A_J},\\
 \mathcal D_J([\eta_g])=1
             &\quad\Longleftrightarrow\quad
                   V_J(g)=x^{A_J}y^{B_J}.
 \end{split}                                                   \tag{F.7.56}
\]
Here $[\eta_g]$ is read as its exact $GT_0^1$
class under (F.7.23).
This supplies the same single full profinite
$g$ in every row, together with the
quadratic equation (F.7.50).

The four $C_3$-stable triples of (F.7.34) admit
one further common word.  With the same
$\tau$ from (F.7.33), define
\[
 R(g)=g^{-1}\eta_g(\tau)g\tau^{-1}\in H.
 \qquad
 \mathcal D_{J_i}([\eta_g])=1
       \quad\Longleftrightarrow\quad p_i(R(g))=1
       \quad(0\leq i\leq3).
                                                               \tag{F.7.57}
\]
Indeed $r=b$ is conjugated by $g$, while
$s=\tau b\tau^{-1}$ is conjugated by
$v(g)=\eta_g(\tau)g\tau^{-1}$.
After filling, $r,s$ generate $E^+$.
An inner surface action extends uniquely across
$E^+$ and its conjugator must be $p_i(g)$,
by $C_{E^+}(K)=1$ and
$K\cap C_{E^+}(b)=1$.
Its agreement on $s$ is then exactly
$p_i(v(g))=p_i(g)$, because the other finite
factor also has trivial centralizer intersection
with $K$.  This is (F.7.57), in both directions.
Replacing $g$ by $d^{5t}g$ and $\eta_g$ by its
specified inner change leaves $R(g)$ unchanged,
so this common-word comparison retains the
outer normalization.

The actual geometric residual group in (F.7.25)
therefore maps into the simultaneous solution
set (F.7.50), (F.7.53), and (F.7.56), with the
common-word constraints (F.7.57).
Triviality of that solution set has not been
proved.  The nontrivial common source kernel
in (F.7.41) supplies no solution to the additional
quadratic equation.  Conversely the exact
section and the primewise images in (F.7.48)
supply no step replacing these full quotient
equations.

The executed standard-library program
\path{code/part_a/modular_five_fillings_f7.py}
records all $95$ Tietze eliminations, all five
inverse full-basis maps, every original cusp
word, the order-three finite centralizer,
the five $j$-recovery word identities in
(F.7.51), and the exact label-four conjugator.
The orbifold generation test, the full torsion
and centralizer inputs, and the simultaneous
solution-space identifications are proved above.


This audit concerns the comparison requested for
\S30, folios 194--196, scans 383--387. F.6's
arithmetic open-image proof applies to $A_5$ without
using (F.7.1), so that proved arithmetic assertion
does not depend on the unresolved intersection.
\end{proof}

\begin{claim}{F.8}{Proved}{The seven-point kernel and its peripheral restrictions}
In the regular $A=A_4$ labelling of (F.7.8), all twelve
seven-point subsets in (F.7.9) form one free orbit for
left multiplication by $A$.  Consequently their kernels
on the specified $F_{13}$ are equal; their kernels on
its profinite completion are equal as well.
The completed kernel meets every conjugate of every
full source cusp subgroup $\widehat\Z$ trivially.
The discrete kernel contains no nonidentity reducible
mapping class under the specified index-twelve
inclusion $F_{13}\le P\Gamma_{0,4}$.
\end{claim}
\passage{\S{}30; folios 194--196; scans 383--387; Part F: nonsolvable finite centralizer comparison}

\begin{proof}

Here is the actual comparison of maps.  Regard the target
as $M_{0,I}$ with its seven labels $I\subset A$, and write
\[
 f_I:X\longrightarrow M_{0,I},\qquad
             z\longmapsto[(gz)_{g\in I}].
\]
For $h\in A$, let
$r_h:M_{0,I}\xrightarrow{\sim}M_{0,hI}$ be the relabelling
which assigns to the label $hg$ the former coordinate
with label $g$.  The configuration defining $f_{hI}(z)$
has coordinate $h(gz)$ at that label.  Applying the same
M\"obius transformation $h^{-1}$ to every coordinate
therefore gives the exact identity of maps
\[
                         f_{hI}=r_h\circ f_I.
\]
This identity also identifies the base configurations;
it induces an isomorphism of the target fundamental
groups intertwining the two homomorphisms from
$\pi_1(X)$.  Functoriality of profinite completion gives
the same commuting diagram for their completions.
Since $r_h$ and its completion are isomorphisms, the
two kernel equalities follow.

The exact label certificate is
\path{code/part_a/seven_point_same_kernel.py}.
It constructs the twelve subsets in the action
$A_5/C_5$ and gives all twelve left multipliers and
the corresponding seven-label permutations.  With
the certificate's zero-based vertex numbering, one
subset is $I=\{0,1,2,3,4,7,10\}$.  Its intersections
with the three cosets of the normal subgroup
$V_4\le A_4$ are respectively
\[
 \begin{array}{c|c}
 \text{coset}&I\text{ retained in that coset}\\ \hline
 \{0,2,7,11\}&\{0,2,7\}\\
 \{1,4,5,9\}&\{1,4\}\\
 \{3,6,8,10\}&\{3,10\}.
 \end{array}
\]
The program checks that none of the twelve subsets
contains a full $V_4$ coset.  Thus the four-label
$V_4$-orbit cross-ratio is not one of the forgetful
projections of any of these seven-point maps.

The intersection of the twelve kernels in (F.7.9)
is therefore the kernel of any one of these named
maps.  F.9 below proves it nontrivial by an exact braid certificate.
This discrete kernel contains no nonidentity reducible
mapping class in the source subgroup of
$P\Gamma_{0,4}$.  To prove this, use the exact
$A_4$ covering
\[
 X=S^2\setminus R_A\longrightarrow
 S^2\setminus\{b_2,b_3,b'_3\}
\]
and its index-twelve injection on fundamental groups
in (F.7.8).  On the four-punctured quotient sphere,
every essential curve separates the moving puncture
$q$ and one branch point $b$ from the other two
branch points.  Let $e=2$ or $3$ be the order at $b$.
The preimage of the disk containing $q,b$ consists
of $12/e$ disks, each containing $e$ of the twelve
points above $q$.  These assertions hold for every
such curve: its monodromy subgroup is the conjugate
of the inertia subgroup at $b$, so its connected
components have degree $e$, and each is a disk
branched only at the one point above $b$.

The lift of a nonzero twist power belonging to the
index-twelve subgroup is a product of twists about
these boundary curves.  More precisely, its exponent
is divisible by $e$, and each lifted twist has exponent
equal to the original exponent divided by $e$.  After
forgetting five points, at least one of these boundaries
remains essential: seven retained points distributed
among six pairs or four triples force a disk containing
at least two retained points, and its complement
contains at least four.  The surviving twists all have
the same nonzero sign; parallel essential components
add their exponents.  Independence of twists about
pairwise disjoint nonisotopic essential curves makes
the resulting mapping class nonidentity.  Finally,
each nonidentity reducible pure mapping class of a
four-punctured sphere is a nonzero power of one such
twist: its invariant multicurve has just one component,
and cutting along it gives two three-punctured spheres.
This proves the asserted exclusion of reducible
elements from the discrete kernel.

There is also an injective detector for each full
profinite cusp.  Let $c$ be a primitive peripheral
generator at any one of the fourteen punctures of
$X$.  Its image in $P\Gamma_{0,4}$ is a conjugate
of $T^e$, where $T$ is a twist as above.  Thus its
lift is the indicated product of twists of exponent
one.  Choose two retained labels in an essential
lifted disk and two outside that disk.  Forgetting
all other labels gives a specified map to
$P\Gamma_{0,4}$.  The inside disk contributes one
twist $U$.  Among the other disks, another twist
survives exactly when both outside labels lie in
one disk; it then gives the same isotopy class and
the same sign.  Hence the image of $c$ is $U^d$
with $d\in\{1,2\}$, up to conjugation.

There is a character
$\ell:P\Gamma_{0,4}\to\Z$ with $\ell(U)=1$.
For example the three peripheral twists in the
free group $P\Gamma_{0,4}=F_2$ have abelianization
vectors $(1,0),(0,1),(-1,-1)$, and an essential
twist is conjugate to one of these under the pure
group.  Choose the corresponding integral linear
functional.  Completing this character and the
two forgetting maps gives a continuous homomorphism
\[
 \widehat F_{13}\longrightarrow
 \widehat{P\Gamma}_{0,7}\longrightarrow
 \widehat{P\Gamma}_{0,4}\longrightarrow\widehat\Z
\]
whose restriction to $\overline{\langle c\rangle}$
is $a\mapsto da$.  Multiplication by $1$ or $2$
is injective on $\widehat\Z$.  The cusp is a full
embedded $\widehat\Z$ in $\widehat F_{13}$, since
its primitive generator is a member of a free
basis after choosing the omitted sphere peripheral.
The same final character detects every conjugate
$g c^a g^{-1}$ with $g\in\widehat F_{13}$, because
its target is abelian.  This proves the asserted
trivial intersection of the completed kernel with
all conjugate cusp groups.  The program gives the
four retained labels and $d$ in each of the
$12\cdot7=84$ subset--inertia-partition checks;
the two fixed points of each inertia subgroup
have the same partition.

The exclusion just proved confines nonidentity discrete
kernel elements to pseudo-Anosov mapping classes. F.9 supplies
an explicit such element by an exact braid certificate.
No assertion about a Nielsen--Thurston type for profinite
elements is used.

The preceding rank-six calculation is a calculation
of its induced abelian map, and does not identify
the kernel of the fundamental-group homomorphism.
The separate file
\path{code/part_a/seven_point_braid_explore.py}
contains a numerical crossing exploration; it is
explicitly exploratory and is not used in this
map comparison. F.9 replaces the discovery computation
by a separate exact rational certificate of a kernel relation.
\end{proof}

\begin{claim}{F.9}{Proved}{An explicit common kernel element for the seven-point maps}
Every one of the twelve seven-point maps in F.7--F.8 has
a nontrivial kernel, both on the specified discrete
$F_{13}$ and on its full profinite completion.  With the
source coordinate, basepoint, and meridians specified
below, the same element
\[
 w=m_{-1}m_i m_1^{-1}m_i^{-1}                         \tag{F.9.1}
\]
lies in all twelve kernels.  A source winding character
sends $w$ to $1$, so its procyclic closure is a full
$\widehat\Z$ contained in their common completed kernel.
The joint kernel of all thirty-six six- and seven-point
maps in (F.7.9) is exactly this same kernel.
This refutes injectivity of these particular restriction
maps; the residual $A_5$ centralizer comparison in F.7
remains unresolved.
\end{claim}
\passage{Section 30, folios 194-196, scans 383-387; Part F, nonsolvable A5 centralizer comparison on the twelve-punctured sphere}

\begin{proof}
\textbf{The specified source and target.}
Use the faithful tetrahedral action generated by
\[
 z\longmapsto-z,\qquad z\longmapsto\frac{z+i}{z-i}.
\]
Its fourteen ramification points are infinity and the
zeros of
\[
 D(z)=z(z^4-1)(z^8+14z^4+1).
\]
Thus $X=\mathbf P^1\setminus R_{A_4}$ has
$\pi_1(X,b)=F_{13}$ for $b=5+4i$.
The seven labels $I=\{0,1,2,3,4,7,10\}$ of F.8,
in that order, have the following coordinate functions:
\[
 \begin{split}
 (m_0(z),\ldots,m_6(z))=
 \left(z,\frac{i(z+1)}{z-1},-\frac1z,
 -\frac{z-i}{z+i},\frac{i(z-1)}{z+1},-z,
 \frac{z-i}{z+i}\right).
 \end{split}                                                   \tag{F.9.2}
\]
Here the subscripts $0,\ldots,6$ denote positions in
the displayed seven-set, not its original labels.
The exact finite part of the certificate below verifies
all $144$ matrix products in this $A_4$, its correspondence
with the labels in F.8, and the ramification polynomial
$D$ by multiplying the fixed-point polynomials of its
seven nontrivial cyclic subgroups.

For $z\in X$, let
\[
 L_z(t)=\frac{(t-m_0(z))(m_1(z)-m_6(z))}
              {(t-m_6(z))(m_1(z)-m_0(z))}.
\]
The seven orbit points are distinct, so this is the
unique M\"obius transformation carrying the first,
second, and seventh points to $0,1,\infty$.
Define the ordered configuration map
\[
 \Phi:X\longrightarrow\operatorname{Conf}_6(\mathbf C),
 \qquad z\longmapsto(R_0(z),\ldots,R_5(z)),
 \quad R_j(z)=L_z(m_j(z)).                              \tag{F.9.3}
\]
The rational functions extend across any removable
singularities in their displayed expressions; their
only poles are among $R_{A_4}$.  In particular
$R_0=0$ and $R_1=1$.
Appending the seventh point $\infty$ and taking the
labelled M\"obius quotient gives a map
\[
 q:\operatorname{Conf}_6(\mathbf C)\longrightarrow M_{0,I}
 \quad\hbox{with}\quad q\circ\Phi=f_I.                 \tag{F.9.4}
\]
These are actual continuous maps, with the displayed
base configurations.  The induced map of fundamental
groups therefore factors through the pure braid group
$P_6\le B_6$.  The inclusion here is induced by the
cover from ordered to unordered configurations and is
injective by path lifting.  We shall prove
$\Phi_*(w)=1$ in $B_6$, which proves $f_{I,*}(w)=1$
by this factorization.

\textbf{The three based source meridians.}
Put $r=27/1000$.  For $c=-1,i,1$, set
\[
 t_c=(\operatorname{Re}c-r)+4i,
 \quad b_c=c-r-ri,
 \quad d_c=c+r-ri,
 \quad e_c=c+r+ri,
 \quad a_c=c-r+ri.
\]
Let $m_c$ be the loop obtained by traversing the
straight polygonal segments with successive vertices
\[
 b,\ t_c,\ b_c,\ d_c,\ e_c,\ a_c,\ b_c,\ t_c,\ b.
                                                               \tag{F.9.5}
\]
The square is counterclockwise; its stem is followed
by its reverse.  Each square contains the single
ramification point $c$.  This also follows directly
from the above polynomial: the other eight finite
points have real and imaginary coordinates with
absolute values $(\sqrt3-1)/2$ or $(\sqrt3+1)/2$,
whereas the remaining finite points are
$0,\pm1,\pm i$.
The exact segment check below also verifies that every
stem and every square edge avoids all zeros of $D$.
Source loop products in (F.9.1) are chronological
concatenation.

\textbf{Exact extraction of the three braids.}
Multiply every coordinate in (F.9.3) by
$\eta=1+(173/1000)i$ and project to its real part.
This multiplication is an orientation-preserving
homeomorphism of the plane.  At the basepoint the
six coordinate positions, in increasing projected
order, are
\[
                      (0,5,3,1,4,2).
\]
Identify this chamber with the standard real ordered
base configuration by a path within the chamber.
The chamber of any fixed strict real ordering is
convex, so the identification is unique up to
homotopy.  Read crossings chronologically.  Write
$j$ for $\sigma_j$, and $-j$ for $\sigma_j^{-1}$;
the positive sign means that the strand on the left
immediately before the crossing has larger imaginary
coordinate.  This fixes both the numbering and the
crossing convention.

The executed certificate
\path{code/part_a/seven_point_kernel_f9.py}
uses only rational and Gaussian-rational arithmetic.
For a segment $z=p+(q-p)t$, $0\le t\le1$, and each
pair of normalized coordinates, write their projected
complex difference before taking real parts as
$A(t)/B(t)$ with $A,B\in\mathbf Q(i)[t]$.
It verifies that $A$, $B$, every coordinate denominator,
and $D(p+(q-p)t)$ have no complex zero on the real
segment: the real and imaginary parts have a greatest
common divisor with no real root there.
The possible crossings are precisely the roots in
$(0,1)$ of
\[
       P(t)=\operatorname{Re}(A(t)\overline{B(t)}).
                                                               \tag{F.9.6}
\]
Clearing rational denominators preserves its roots.
Sturm isolation over $\mathbf Q$ enumerates all of
them.  Their rational isolating intervals are pairwise
disjoint, their multiplicities are one, and
$\operatorname{Im}(A\overline B)$ has a certified
constant nonzero sign on each interval.  No endpoint
is a crossing.  The two strands are checked to be
adjacent in the current order at every crossing,
and the final order is checked against direct exact
evaluation at the endpoint.

There are $18$ different segments up to reversal,
$270$ pair checks, and $64$ isolated crossings in
these segments; the largest crossing-polynomial
degree is $12$.  Reversed segments give the inverse
words.  Between successive crossings the path lies
in a convex ordering chamber.  At a crossing only
the named adjacent pair changes order, with its
certified nonzero imaginary separation.  Straightening
these intervals and crossing neighborhoods gives the
usual elementary braid pieces.  Thus the recorded
words represent the actual monodromies, not just a
linear representation of them.  This is also the
projection-isotopy argument of Artin's Theorem~3
\cite{A-Artin1947}.

The freely reduced words for $m_{-1},m_i,m_1$ are,
respectively, $A,B,C$ below; a line break only separates
consecutive parts of the same word:
\[
\begin{aligned}
 A={}&(1,4,5,2,1,3,2,3,2,1,2,-5,-4,-5,\\
     &\hspace{9mm}-5,4,5,-2,-1,-2,-3,-2,-3,-1,-2,-5,-4,-1),\\
 B={}&(1,4,5,2,3,1,4,2,3,4,-3,1,2,3,4,5,\\
     &\hspace{9mm}-2,5,4,3,2,1,-4,-3,-2,-4,-1,-3,-2,-5,-4,-1),\\
 C={}&(1,4,5,2,3,4,1,2,3,4,5,-1,3,2,-5,-4,\\
     &\hspace{9mm}-4,5,-2,-3,1,-5,-4,-3,-2,-1,-4,-3,-2,-5,-4,-1).
\end{aligned}                                                   \tag{F.9.7}
\]

\textbf{The braid identity and the source detector.}
Use the faithful Artin substitution on
$F_6=\langle X_1,\ldots,X_6\rangle$:
\[
 \sigma_j(X_j)=X_jX_{j+1}X_j^{-1},\qquad
 \sigma_j(X_{j+1})=X_j,
 \qquad \sigma_j(X_k)=X_k\ (k\ne j,j+1).              \tag{F.9.8}
\]
Its inverse is used for a negative letter.  Products
of substitutions are composed from the displayed
word order; the certificate implements composition
by substitution and cancellation of adjacent inverse
letters.  Artin's Theorem~14, with the inverse
generator formula on p.~113, proves faithfulness
\cite{A-Artin1947}.  Simultaneously reversing all
generator signs is an automorphism of the braid
group, and reversing the stacking convention is an
anti-automorphism; neither changes the identity test.

For the exact words (F.9.7), substitution gives
\[
       \operatorname{Artin}(AB)=\operatorname{Artin}(BC),
 \qquad
       \operatorname{Artin}(ABC^{-1}B^{-1})(X_j)=X_j
       \quad(1\le j\le6).                              \tag{F.9.9}
\]
This is an integer free-word calculation.  The output
records all six freely reduced common images for the
first equality and all six single-letter images for
the second.  Faithfulness gives
$AB=BC$ in $B_6$.  By the certified identification of
the three words with the three source monodromies,
(F.9.1) lies in $\ker\Phi_*$, and hence in
$\ker f_{I,*}$ by (F.9.4).

The map $X\to\mathbf C^*$, $z\mapsto z+1$, gives
the integral winding character
\[
 \ell:\pi_1(X,b)\longrightarrow\Z,
 \qquad
 \ell(m_{-1})=1,\quad \ell(m_i)=\ell(m_1)=0.
\]
Therefore $\ell(w)=1$ and $w\ne1$.
Completing this exact character gives
$\widehat\ell:\widehat F_{13}\to\widehat\Z$ with
$\widehat\ell(w)=1$; this proves nontriviality in the
full completion without a detection assumption.
The continuous power map $\widehat\Z\to\widehat F_{13}$,
$a\mapsto w^a$, has composite $a\mapsto a$ with
$\widehat\ell$, so it is injective.  Its image lies
in the closed completed kernel.

Finally F.8 gives the actual map identities
$f_{hI}=r_h\circ f_I$ for the twelve left translates
$hI$, with $r_h$ a labelled target isomorphism.
They carry this same source word, and all its
profinite powers, into the identity in every target.

The twenty-four six-point conditions in (F.7.9)
follow from these twelve conditions by specified
forgetting maps.  For one of the six subgroups
$C_5\le A_5$, let $p_+,p_-$ be its fixed poles and
let $O$ be either of its two five-cycles.  Put
\[
 K=O\cup\{p_+,p_-\},\qquad J_\pm=O\cup\{p_\pm\}.
\]
Forgetting the other pole gives the labelled morphism
$d_\pm:M_{0,K}\to M_{0,J_\pm}$ and the exact identity
$f_{J_\pm}=d_\pm\circ f_K$ on the same based source.
The twelve choices of $(C_5,O)$ give all twelve
seven-sets, and their two pole choices give all
twenty-four six-sets.  The separate executed certificate
\path{code/part_a/seven_point_six_factorization_f9.py}
lists all twenty-four containments, the omitted pole,
and the six positions retained in each seven-tuple.
Functoriality of fundamental groups and of full
profinite completion proves on both sources that
\[
 \bigcap_{I\text{ among all }36}\ker f_{I,*}
       =\bigcap_{K\text{ among the }12\text{ seven-sets}}
                                      \ker f_{K,*}.     \tag{F.9.10}
\]
For the completed equality the starred maps in
(F.9.10) mean their completed homomorphisms.
Consequently (F.9.1) and its full procyclic closure
also satisfy every six-point condition.

The common-kernel injectivity route in F.7 therefore
fails.  The necessary seven-point restrictions of
an element of the residual $A_5$ centralizer still
vanish, but the converse implication has not been
proved.  No centralizer element is constructed here.

The discovery computation used numerical crossings
and a finite Burau image to find the short candidate.
Neither computation is an input to this proof.
The file
\path{code/part_a/seven_point_kernel_f9_output.json}
contains the deterministic exact certificate:
the coordinate functions, Gaussian-rational path
vertices, every crossing polynomial, every rational
root interval and sign, the three words, and the
Artin identities.  The companion stdout file records
the successful execution.
\end{proof}

\begin{claim}{F.10}{Proved}{Discrete procyclic directions in the residual centralizer}
Let $A_5\leq\Gamma^+_{0,12}$ be the geometric
icosahedral group acting on its twelve-vertex orbit,
and put $P=P\Gamma_{0,12}$.  For every nonidentity
$x\in P$, its procyclic closure is a full
$\widehat\Z$, and
\[
 \overline{\langle x\rangle}\cap
              C_{\widehat P}(A_5)=1.                  \tag{F.10.1}
\]
In particular, the full procyclic subgroup of the
common forgetting kernel constructed in F.9 meets
the actual residual $A_5$ centralizer trivially.
For that particular word $w$, it already has
trivial intersection with the centralizer of the
specified involution $\beta=(1\ 2)(4\ 5)$ in $A_5$.
\end{claim}
\passage{Section 30, folios 194-196, scans 383-387; Part F, nonsolvable A5 centralizer comparison in the minimal twelve-puncture case}

\begin{proof}
\textbf{An explicit separating family of integral representations.}
We first construct the finite quotients used in the
proof.  For $n\geq4$, normalizing the last marked
point to infinity and the first two to $0,1$
identifies
\[
 M_{0,n}=\operatorname{Conf}_{n-3}
                              (\mathbf C\setminus\{0,1\}).
\]
The exact homeomorphism
\[
\begin{split}
 \operatorname{Conf}_{n-1}(\mathbf C)
   &\longrightarrow\mathbf C\times\mathbf C^*\times M_{0,n},\\
 (z_0,z_1,z_2,\ldots,z_{n-2})
   &\longmapsto\left(z_0,z_1-z_0,
        \left(\frac{z_j-z_0}{z_1-z_0}\right)_{j=2}^{n-2}\right)
\end{split}                                                   \tag{F.10.2}
\]
has inverse $(a,b,(u_j))\mapsto(a,a+b,(a+bu_j))$.
The section fixing the first two factors at $(0,1)$
therefore gives an injective homomorphism
\[
          P\Gamma_{0,n}\longrightarrow P_{n-1},
 \qquad P_{n-1}\cong\Z\times P\Gamma_{0,n}.             \tag{F.10.3}
\]
Here $P_{n-1}$ is the ordered plane braid group.
All basepoints in (F.10.2) are the displayed normalized
configurations.  Thus (F.10.3) is the actual direct-factor
map, rather than an inference from the ranks of the groups.

Put $r=n-1$ and $F_r=\langle t_1,\ldots,t_r\rangle$.
Artin's faithful braid action gives
$P_r\hookrightarrow\operatorname{Aut}(F_r)$
\cite[Theorem~14]{A-Artin1947}.  We use the homomorphism
convention
\[
 \sigma_i(t_i)=t_it_{i+1}t_i^{-1},\qquad
 \sigma_i(t_{i+1})=t_i,
\]
fixing the other generators, with automorphism products
composed in that order.  Equivalently, for chronological
point-motion concatenation, use the inverse of the right
monodromy action on the punctured-plane group.  This
conversion makes $A(bc)=A(b)\circ A(c)$ and preserves
faithfulness.  Every pure braid fixes $F_r^{\rm ab}$.

Let
\[
 \mu:F_r\longrightarrow
          \Z\langle\!\langle X_1,\ldots,X_r\rangle\!\rangle^*,
 \qquad \mu(t_i)=1+X_i
\]
be the integral Magnus expansion, with
$(1+X_i)^{-1}=1-X_i+X_i^2-\cdots$.
Here is the injectivity needed below.  Write a nonidentity
reduced word in maximal syllables as
$t_{i_1}^{a_1}\cdots t_{i_s}^{a_s}$, where every
$a_j\ne0$ and successive $i_j$ differ.  The coefficient
of the monomial $X_{i_1}\cdots X_{i_s}$ is
$a_1\cdots a_s\ne0$.  Indeed that monomial has $s$
successive runs of distinct variables.  Omitting a
syllable contribution leaves at most $s-1$ runs,
whereas using all syllables and total degree $s$
forces their linear terms.  Thus $\mu$ is injective.

For each $d\geq1$ let
\[
 T_{r,d}=\Z\langle X_1,\ldots,X_r\rangle/
                   (\text{monomials of degree }>d).
\]
An automorphism $\alpha$ of $F_r$ acts on this finite
free $\Z$-module by the substitution
\[
             X_i\longmapsto\mu(\alpha(t_i))-1
                                      \pmod{\deg>d}.    \tag{F.10.4}
\]
The substitution identities for group words show
directly that these are homomorphisms of automorphism
groups.  For an automorphism fixing $F_r^{\rm ab}$,
the right side is $X_i$ plus terms of degree at least
two.  Order the monomial basis by decreasing degree.
The matrix in (F.10.4) is then upper unitriangular,
and all its entries are integers.  Combining
(F.10.3)--(F.10.4) gives actual homomorphisms
\[
 \rho_d:P\Gamma_{0,n}\longrightarrow
                      \operatorname{UT}_{N_d}(\Z),
 \qquad N_d=1+r+\cdots+r^d.                             \tag{F.10.5}
\]
This family separates elements.  Two different pure
mapping classes give different automorphisms of $F_r$
by (F.10.3) and Artin faithfulness.  Some generator has
different images; injectivity of $\mu$ makes their
expansions different in some finite degree.  That
degree gives a representation (F.10.5) which separates
the two classes.  In particular, no passage from a
discrete lower-central quotient to a completed one
is assumed here: the detecting maps are integral
matrices and their reductions in finite matrix groups.

\textbf{The actual finite-group action on the matrices.}
Return to $n=12$ and put $G=A_5$.  Its intersection
with $P$ is trivial, as is seen from its faithful
puncture permutation.  The subgroup $PG$ of
$\Gamma^+_{0,12}$ is the actual semidirect product
$P\rtimes G$ and has finite index.  Thus its profinite
completion is the corresponding open subgroup of
$\widehat\Gamma^+_{0,12}$, with the specified copies
of $\widehat P$ and $G$.

The discrete centralizer $C_P(G)$ is trivial.
Indeed the quotient sphere has exactly the three
branch/puncture decorations of orders $2,3,5$.
The discrete covering normalizer sequence of F.7
has trivial quotient mapping class group, so its
centralizer is $Z(A_5)=1$.
For $x\ne1$ choose $a\in G$ with
$axa^{-1}\ne x$, and then choose $d$ in (F.10.5)
which separates these two elements.  Write $\rho=\rho_d$.
On a direct sum of $|G|$ copies of $\Z^{N_d}$,
indexed by $g\in G$, define
\[
       R(h)=\operatorname{diag}_{g\in G}
                              \rho(g^{-1}hg),\qquad h\in P.
                                                               \tag{F.10.6}
\]
Let $T_a$ carry the block with index $g$ to the block
with index $ag$.  Then
\[
 T_aT_b=T_{ab},\qquad
 T_aR(h)T_a^{-1}=R(aha^{-1}).                            \tag{F.10.7}
\]
For example the block with index $g$ on the left of
the second identity is
$\rho((a^{-1}g)^{-1}h(a^{-1}g))$
$=\rho(g^{-1}aha^{-1}g)$, which is its right-hand
block.  Consequently
\[
       h a\longmapsto R(h)T_a                          \tag{F.10.8}
\]
is a homomorphism from the actual group $P\rtimes G$
to $\operatorname{GL}_{|G|N_d}(\Z)$.
Its reductions modulo every positive integer are
finite quotients of this group.

\textbf{Every prime coordinate is detected.}
Put $W=R(x)$.  It is unipotent.  By the choice of
$d$, $T_aWT_a^{-1}\ne W$, already on the block
with index $1$.  The finite polynomial logarithm
\[
        L=\log W=\sum_{j\geq1}
                      \frac{(-1)^{j+1}}j(W-I)^j
\]
is a rational nilpotent matrix; the sum terminates.
The exponential and logarithm polynomials are
mutually inverse on nilpotent and unipotent matrices
over a field of characteristic zero.  Hence
\[
                       T_aLT_a^{-1}-L\ne0.             \tag{F.10.9}
\]
In particular this is the same nonzero rational
matrix after embedding $\Q$ in $\Q_p$, for every
prime $p$.

The reductions of (F.10.8) modulo $p^k$, for all
$k$, give a continuous homomorphism from
$\widehat{P\rtimes G}$ to
$\operatorname{GL}_{|G|N_d}(\Z_p)$.
The image of the procyclic power $x^u$, for
$u\in\widehat\Z$, is
\[
                          \exp(u_p L),                 \tag{F.10.10}
\]
where $u_p$ is its $\Z_p$ coordinate.
To check this equality, first take integral $u$;
the identity $W^u=\exp(uL)$ follows from the finite
polynomial identities.  For arbitrary $u_p$, pass
to limits of integers in $\Z_p$.  The matrix
polynomial is continuous, and the powers of $W$
modulo $p^k$ lie in an upper unitriangular finite
$p$-group, so these powers depend only on $u_p$.
The limit lies in $\operatorname{GL}(\Z_p)$ because
the integral matrices form a closed subset there.
This proves (F.10.10) through the specified finite
quotients.

Now an element $x^u$ centralizing $G$ commutes with
$a$.  Apply the continuous representation and then
the polynomial logarithm to (F.10.10).  It gives
\[
                   u_p\bigl(T_aLT_a^{-1}-L\bigr)=0
                      \quad\hbox{in }M_{|G|N_d}(\Q_p).
                                                               \tag{F.10.11}
\]
By (F.10.9), $u_p=0$ for every prime $p$.
Thus $u=0$ in $\widehat\Z$ and (F.10.1) follows.
The same argument without $T_a$, using any
$\rho_d(x)\ne I$, proves that
$u\mapsto x^u$ is injective.  Its image is exactly
the procyclic closure, which is therefore a full
$\widehat\Z$ as claimed.

\textbf{The particular F.9 direction.}
The following transport fixes the basepoint convention.
Put $\varphi=(1+\sqrt5)/2$, $\rho=\sqrt{\varphi+2}$,
and $z_*=-i(\rho+\varphi)$.  In the tetrahedral coordinate
of F.9 the orbit of $z_*$ is the specified icosahedral
twelve-point configuration.  The extra involution has
M\"obius matrix
\[
 \begin{pmatrix}
  1&\varphi-1-i\varphi\\
  \varphi-1+i\varphi&-1
 \end{pmatrix}.                                         \tag{F.10.12}
\]
This matrix represents $(1\ 2)(4\ 5)$ under the
fixed five-letter identification of $A_5$, and its
twelve-vertex permutation is checked in the certificate
below.  Use the
polygonal source path
$\gamma:z_*\longrightarrow5+z_*\longrightarrow b=5+4i$.
It misses $R_A$: the finite points of $R_A$ have real
and imaginary coordinates of absolute value less than
$3/2$, the first segment has imaginary part less than
$-3$, and the second has real part $5$.
Transport along the full labelled orbit path
$(g\gamma)_{g\in A_4}$ identifies the marked mapping
class groups at these two configurations.  Equivariant
isotopy extension keeps the specified $A_4$ action
fixed.  The symbol $\beta$ at the F.9 basepoint denotes
the transport of (F.10.12) by this path.  The exact
matrix and vertex-label certificate
\path{code/part_a/second_a4_c3_attack.py}
checks all $180$ generator products and all $720$
vertex images in the number field
$\Q(i,\sqrt5,\rho)$.

The F.9 source homomorphism is the full twelve-point
orbit map of F.7.8, before the stated forgetting.
For every source path, its configurations remain
invariant under the same fixed M\"obius $A_4$.
Equivariant isotopy extension along the path, or
equivalently the point-pushing lift of the quotient
cover in F.7.8, places its monodromy in the actual
discrete $C_P(A_4)$.  That map identifies the source
$F_{13}$ with this discrete centralizer.  Thus the
word $w$ is a nonidentity element of $C_P(A_4)$,
not merely an element whose seven-point images
are compatible with $A_4$.

In the fixed five-letter description, $A_4$ fixes
the fifth letter and
$\beta=(1\ 2)(4\ 5)$ carries it to the subgroup
fixing the fourth.  These subgroups generate $A_5$;
equivalently $\langle A_4,\beta\rangle=A_5$.
Therefore $\beta w\beta^{-1}\ne w$, since the
opposite equality would put $w$ in the just proved
trivial discrete $A_5$ centralizer.  We may take
$a=\beta$ throughout (F.10.6)--(F.10.11).
This proves the last assertion of the claim, on
the whole full profinite procyclic closure.

\textbf{Executed formula certificate.}
The exact-integer program
\path{code/part_a/procyclic_magnus_f10.py}
checks the maximal-syllable coefficient on all $484$
nonempty reduced words of length at most five in two
generators.  It verifies the substitution-composition
identities and upper unitriangular form for two actual
pure Artin automorphisms on the degree-three tensor
algebra of rank $40$, and checks eleven integral
power/logarithm identities.  For the actual action of
$A_5$ on $\operatorname{Inn}(F_5)$ by permutation of
the five free generators, it checks all $180$
displayed block-conjugation identities and $60$
block-product identities.  Finally it checks finite
power reductions modulo $p^2$ for $p=2,3,5,7,11$.
This is a certificate for the algebraic formulas;
it does not compute a separating tensor degree for
the particular F.9 word.  Existence of that degree
is proved above by faithful Artin action and the
integral coefficient argument.

The conclusion restricts the actual residual
centralizer on every displayed discrete procyclic
direction.  It does not assert that an arbitrary
element of that closed residual group lies in one
of these subgroups.  No such covering assertion is
used, and the full equality in F.7 remains unresolved.
\end{proof}

\begin{claim}{F.11}{Proved}{The rank-six cyclic centralizer and an exact common kernel}
Let $P=P\Gamma_{0,12}$ for the specified geometric
icosahedral action in F.7.  Fix $D=C_5\leq A_5$ and
$N=N_{A_5}(D)\cong D_{10}$, where $D_{10}$ has order ten.
The explicit configuration and marking below give
\[
 \begin{split}
 C_P(D)&=C_P(N)\cong\pi_1(X_5)=F_6,\\
 C_{\widehat P}(D)&=C_{\widehat P}(N)
                  \cong\widehat{\pi_1(X_5)}=\widehat F_6,
 \qquad
 X_5=\mathbf P^1\setminus\{0,\infty,\mu_5\}.
 \end{split}                                                   \tag{F.11.1}
\]
These are the actual discrete and completed inclusions
constructed below.  Of the twelve seven-point restrictions
in F.7.9, two are trivial on this subgroup; the other ten
have exactly the same kernel.  Their joint kernel, including
all twenty-four six-point restrictions, is the kernel of
the specified map
\[
 f:X_5\longrightarrow M_{0,7},\qquad
 t\longmapsto
 (1,\zeta^2,\zeta^4,t,t\zeta,t\zeta^4,\infty),
 \qquad \zeta=e^{2\pi i/5}.                                  \tag{F.11.2}
\]
The equalities of kernels hold for the discrete groups
and their full profinite completions.

At $b=2+2i$, the precisely based meridians below give
\[
          w=[m_{\zeta^2},\,m_\zeta m_{\zeta^3}m_0]
                    \in\ker(f_*),\qquad w\ne1.               \tag{F.11.3}
\]
Its procyclic closure is a full $\widehat\Z$ in this
common completed kernel.  After the specified basepoint
transport and inclusion into $\widehat P$, that
procyclic subgroup has trivial intersection with
$C_{\widehat P}(A_5)$, by F.10.  Thus the restrictions
again have a concrete common kernel, while the residual
centralizer in F.7 remains unresolved.
\end{claim}
\passage{Section 30, folios 194-196, scans 383-387; Part F, nonsolvable A5 centralizer comparison on the twelve-punctured sphere}


\begin{proof}
\emph{The actual action, labels, and basepoint transport.}
Use the same five-letter generators and based coset
labels as F.7.19--57:
\[
 a=(1\ 2)(3\ 4),\qquad b=(1\ 3\ 5),\qquad d=ab,\qquad
                       D=\langle d\rangle .
\]
The indices $0,\ldots,11$ refer to the twelve cosets
of $D$, in that previously specified order.  Put
\[
 \theta=2\sin(2\pi/5),\quad
 \theta^4-5\theta^2+5=0,\quad 19/10<\theta<2,\qquad
 k=1-\theta^2,\qquad t_0=k\zeta^2 .
                                                               \tag{F.11.4}
\]
The exact marked configuration is
\[
 \begin{array}{c|ccccc}
 j&0&1&2&3&4\\ \hline
 \text{label of }\zeta^j&1&3&5&6&2\\
 \text{label of }t_0\zeta^j&4&7&9&11&8
 \end{array}
 \qquad 0\longleftrightarrow0,\quad
                      10\longleftrightarrow\infty .           \tag{F.11.5}
\]
Here are actual M\"obius transformations realizing the
named generators and the normalizer:
\[
 \begin{aligned}
 a(z)&=\frac{1-z}{1-z/k},&
 b(z)&=a(\zeta z),\\
 d(z)&=\zeta z,&
 \tau(z)&=\frac{t_0}{z},\\
 \tau&=(1\ 5)(3\ 4)\in A_5,&
 \tau d\tau^{-1}&=d^{-1}.
 \end{aligned}                                                 \tag{F.11.6}
\]
The finite certificate
\path{code/part_a/c5_centralizer_labels_f11.py}
checks these formulas in
$\mathbf Q(\theta,i)$, the $720$ images of the
twelve points under the sixty matrices, and all
$120$ products by the two generators.  The sixty
distinct matrices therefore form exactly the given
$A_5$, with precisely its stated coset action.
In particular this supplies an actual equivariant
identification of the marked action, not just an
isomorphism of abstract groups.  It is the explicit
M\"obius realization of the based covering choice
in F.7; changing that initial marking conjugates
all the inclusions and all the named forgetting
maps by the same marking isomorphism.

For every $t\in X_5$ replace $t_0$ in (F.11.5) by $t$.
This gives a map
\[
 j:X_5\longrightarrow M_{0,12},\qquad
                   t\longmapsto(0,\infty,\mu_5,t\mu_5)
                                                               \tag{F.11.7}
\]
with the individual labels in (F.11.5).
The functions $d(z)=\zeta z$ and $\tau_t(z)=t/z$
give a labelled $D_{10}$ action on this entire
family: $\tau_t$ exchanges the two poles and
sends the first orbit point with index $j$ to
the second with index $-j$, and conversely.
They satisfy the group relations on every fibre.
Monodromy therefore centralizes the actual
$D_{10}$ at $t_0$.
Equivalently, each of the ten corresponding
label permutations fixes the map to labelled
moduli (F.11.7), through these specified M\"obius
transformations, and hence fixes every induced
based loop.

The straight segment $\gamma$ from $t_0$ to
$b=2+2i$ lies in $X_5$: the exact field calculation
gives $\operatorname{Re}t_0>2$, so the whole
segment has real part at least two.
The element of $C_P(N)$ associated to (F.11.3)
is explicitly
$j_*(\gamma w\gamma^{-1})$.
All loop products here mean chronological
concatenation.  This also fixes the basepoint
transport for the seven- and six-point target maps.

\emph{The exact discrete centralizer and its full topology.}
The quotient by $d$ is the actual cover
$z\mapsto z^5$; its four punctures are
$0,\infty,1,t_0^5$.  An equivalent based group
description is
\[
 L=\pi_1\bigl(\mathbf P^1\setminus\{0,\infty,1,t_0^5\}\bigr)
       =F_3,\qquad
 U=\ker(\chi:L\longrightarrow C_5)=\pi_1(S_{0,12}),
                                                               \tag{F.11.8}
\]
where the positive meridians at $0,\infty$ have
$\chi$-values $1,-1$ and the other two have value
zero.  The inclusion and the quotient map in
(F.11.8) are induced by this covering.

We spell out why descent accounts for the entire
mapping-class centralizer.  Conjugation embeds
$L$ in $\operatorname{Aut}(U)$: an element of the
free group $L$ centralizing the nonabelian subgroup
$U$ is trivial.  Its image in $\operatorname{Out}(U)$
is the specified deck group $D$.
For an actual automorphism $\alpha$ representing
a pure mapping class centralizing $D$, conjugation
by $\alpha$ preserves this embedded $L$.
Indeed its effect on every deck outer class is
the identity, so its effect on the conjugation
automorphism of an $l\in L$ differs from that
automorphism by conjugation by an element of $U$.
The embedding of $L$ makes the resulting element
unique.  This defines an automorphism
$\alpha_L$ of $L$, restricting to $\alpha$ on $U$
and inducing the identity on $C_5$.

The two ramified cusp meridians in $U$ are fifth
powers of the corresponding meridians of $L$;
the other cusp meridians are their specified
unramified lifts.  Since $\alpha$ fixes each
upstairs peripheral class with exponent one,
the uniqueness of roots in a discrete free group
shows that $\alpha_L$ fixes each of the four
downstairs peripheral classes with exponent one.
The discrete Dehn--Nielsen--Baer identification
used in A.12 therefore gives a pure mapping class
on the four-punctured quotient.
Changing $\alpha$ by an inner automorphism of $U$
does not change that downstairs class.
The kernel of this descent consists of deck
classes, since an inner automorphism of $L$
restricts to precisely such a class on $U$.
A deck class fixing the upstairs labels is the
identity, because $D$ acts freely on each
five-cycle.  This proves injectivity of descent
on the pure centralizer.

Conversely, every pure class on the quotient
preserves $\chi$ and has five commuting lifts.
Choose the unique lift fixing the labelled
point $1$ of the first five-cycle; it then fixes
that entire cycle.  Identify
$P\Gamma_{0,4}$ with
$\pi_1(\mathbf P^1\setminus\{0,1,\infty\},t_0^5)$
by pushing its fourth point while the first
three are fixed, as in A.12.
The chosen lift fixes the second five-cycle
exactly when this loop lifts to a closed loop
under $t\mapsto t^5$.
Its image is consequently the kernel of winding
about zero modulo five, namely
$\pi_1(X_5,t_0)$.
The lift is exactly the family (F.11.7).
This proves
$j_*\pi_1(X_5)=C_P(D)$ and identifies its source
as the index-five subgroup $F_6$ of this specified
$F_2$.
Since its monodromy also centralizes $\tau$,
it proves $C_P(D)=C_P(N)$.

There is a particularly direct check of the full
induced topology.  Retain the eight labels
\[
              I_8=\{0,10\}\cup\{1,3,5,6,2\}\cup\{4\}.
                                                               \tag{F.11.9}
\]
After this forgetting map, (F.11.7) has the seven
fixed points $0,\infty,\mu_5$ and the single
moving point $t$.  Filling that moving point
is precisely the normalization map of A.12.
On its point-pushing subgroup the resulting
automorphism of $\pi_1(X_5)$ is inner conjugation
by the source loop.  After completion the exact
composite is therefore
\[
 \widehat{\pi_1(X_5)}
 \xrightarrow{\widehat j_*}\widehat P
 \longrightarrow\widehat{P\Gamma}_{0,8}
 \longrightarrow\operatorname{Aut}(\widehat{\pi_1(X_5)}),
               \qquad u\longmapsto\operatorname{Inn}(u).
                                                               \tag{F.11.10}
\]
The final extension is continuous on each finite
characteristic quotient.  Its displayed restriction
is injective because the free profinite group has
trivial centre.  Thus $\widehat j_*$ is injective
and its image is exactly the closure of the
discrete centralizer.  The prime-order centralizer
comparison already justified in F.7 identifies
that closure with $C_{\widehat P}(D)$.
Taking closures of $C_P(D)=C_P(N)$ gives both
inclusions needed for the second equality in
(F.11.1).

\emph{All thirty-six restrictions and the actual normalizer transports.}
For $D$ itself the two seven-sets are
\[
 \begin{split}
 I_A&=\{0,1,2,3,5,6,10\},\\
 I_B&=\{0,4,7,8,9,10,11\}.
 \end{split}                                                   \tag{F.11.11}
\]
Their maps to labelled moduli are constant:
the first configuration is already fixed, and
the second becomes fixed under $z\mapsto z/t$.
For the other ten sets choose
$I=\{1,2,4,5,7,8,10\}$ as representative.
In the ordered target labels
$(1,5,2,4,7,8,10)$ its map is exactly (F.11.2).
The following are the actual $D_{10}$ transports:
\[
\begin{array}{c|l|l}
 k& d^k I&d^k\tau I\\ \hline
 0&\{1,2,4,5,7,8,10\}&\{0,1,2,3,4,7,11\}\\
 1&\{1,3,4,6,7,9,10\}&\{0,1,3,5,7,8,9\}\\
 2&\{2,3,5,7,9,10,11\}&\{0,3,4,5,6,9,11\}\\
 3&\{1,5,6,8,9,10,11\}&\{0,2,5,6,7,8,11\}\\
 4&\{2,3,4,6,8,10,11\}&\{0,1,2,4,6,8,9\}.
\end{array}                                                     \tag{F.11.12}
\]
The label certificate enumerates the six order-five
subgroups, verifies this orbit, and records each
transporter as a five-letter permutation, a word
in $a,b$, and its permutation of the twelve labels.
For $n=d^k\tau^\epsilon$, the actual homeomorphism
in (F.11.6) gives the target isomorphism
$n_I:P\Gamma_{0,I}\to P\Gamma_{0,nI}$.
Naturality of forgetting gives
\[
        r_{nI}\circ\operatorname{conj}(n)
                           =n_I\circ r_I .
                                                               \tag{F.11.13}
\]
Conjugation by $n$ is the identity on the subgroup
(F.11.1), so these ten restrictions have equal
kernels there.  The same identity extends to
the full completions.

Each of the twenty-four six-sets is obtained by
removing one of the two poles from its specified
seven-set.  The certificate lists all twelve
pairs of these literal subset inclusions.
Forgetting that pole gives the exact factorization
of the corresponding six-point map through the
seven-point map.  Hence adding all twenty-four
constraints does not change the joint kernel.
Together with (F.11.11)--(F.11.13), this proves
the asserted joint-kernel equality for both
the discrete and completed groups.

\emph{Precisely based loops and exact monodromy.}
Put $r=31/1000$.  For
$q\in\{\zeta^2,\zeta,\zeta^3,0\}$,
define $m_q$ by the six successive straight
segments through
\[
 b,\quad q+r+ri,\quad q-r+ri,\quad
 q-r-ri,\quad q+r-ri,\quad q+r+ri,\quad b.
                                                               \tag{F.11.14}
\]
The square is counterclockwise, with its
whisker traversed in reverse at the end.
The certificate checks that every segment
avoids $0,\mu_5$, and that each square contains
exactly its named puncture.

Remove the last, fixed point $\infty$ from
(F.11.2).  This gives an actual ordered
configuration map to $\operatorname{Conf}_6(\mathbf C)$.
Appending $\infty$ and taking labelled M\"obius
classes recovers $f$.
Its loop map therefore factors through
$P_6\leq B_6$, with the injective inclusion
given by the ordered-to-unordered covering.
Multiply all six finite coordinates by
$1+(371/1000)i$, and project to their real parts.
The increasing base order of the six positions,
numbered $0,\ldots,5$, is
$(4,1,2,0,3,5)$.
Use the convention of F.9: a positive
$\sigma_j$ has its left strand above its
right strand just before the crossing,
and crossings are read chronologically.
Identifying the base chamber with the real
ordered configuration fixes the braid
generators up to its unique chamber isotopy.

The program
\path{code/part_a/c5_seven_kernel_exact.py}
uses fractions in the exact field
$\mathbf Q[\theta]/(\theta^4-5\theta^2+5)$,
with
\[
             \zeta=\frac{\theta^2-3+i\theta}{2}.
                                                               \tag{F.11.15}
\]
The polynomial is strictly increasing on
$(19/10,2)$ and has opposite endpoint signs.
Rational bisection isolates its specified
root; exact interval evaluation certifies
each nonzero field sign.
For every one of the $24$ source segments
and every one of the $15$ coordinate pairs,
the projected difference is affine in
the segment parameter.  Its endpoint signs
enumerate all crossings, and its single
possible root is computed exactly in the field.
Every crossing has nonzero imaginary
separation.  Distinct crossings have distinct
parameters, the strands are adjacent in
the current order, and the final ordering
agrees with direct endpoint evaluation.
The exact source-segment tests also exclude
every configuration collision.

The output contains $360$ pair checks,
$58$ simple crossings, and $2530$ certified
algebraic signs, with a rational isolating
interval obtained by $128$ bisections.
Thus the words describe the actual paths:
between crossings the ordering chambers
are convex; each crossing has exactly the
certified adjacent pair and sign.
This is the projection-isotopy argument
of Artin's Theorem~3 \cite{A-Artin1947}.
Writing signed integers for braid generators,
the words for $m_{\zeta^2}$ and for
$m_\zeta m_{\zeta^3}m_0$, after free reduction,
are respectively
\[
\begin{split}
 A={}&(-4,-3,-5,-1,-4,-2,-1,-1,2,4,1,5,3,4),\\
 B={}&(-4,-3,-5,-1,-4,-2,-3,-2,-1,-4,\\
     &\hspace{9mm}-1,-4,-3,-2,-3,4,-5,-1,3,4).
\end{split}                                                     \tag{F.11.16}
\]
The complete words for all four individual
meridians, of lengths $14,8,20,16$, and
every crossing are retained in the output.

Use the faithful Artin substitutions of F.9.8,
with the stated inverse substitution for
each negative generator.  Integer free-word
substitution and cancellation give
\[
       \operatorname{Artin}(AB)
           =\operatorname{Artin}(BA),\qquad
       \operatorname{Artin}(ABA^{-1}B^{-1})(X_i)=X_i
                    \quad(1\leq i\leq6).                       \tag{F.11.17}
\]
The output records the six common images
and the six single-letter final images.
Artin's faithfulness theorem, Theorem~14
and formulas (14)--(15) of the original
article \cite{A-Artin1947}, proves
$[A,B]=1$ in $B_6$.
The actual configuration factorization
therefore proves $f_*(w)=1$.

\emph{Nontriviality and the full procyclic subgroup.}
Fill all source punctures except
$\zeta^2,\zeta,\infty$.
This is an actual map to a twice-punctured
plane.  The loops $m_{\zeta^3}$ and $m_0$
become trivial; the other two loops have
winding vectors $(1,0)$ and $(0,1)$.
In the integral Heisenberg quotient with
product
\[
 (a,b,t)(a',b',t')
       =(a+a',b+b',t+t'+ab'),
\]
their images consequently have the form
$(1,0,s)$ and $(0,1,t)$.
Their commutator is exactly $(0,0,1)$,
for every $s,t\in\Z$.
Thus the image of $w$ is nontrivial.
For each prime $p$, reducing these integer
matrices modulo every $p^n$ extends this
homomorphism continuously to the full
source completion and then to
$\operatorname{UT}_3(\Z_p)$.
The central coordinate of $w^u$ is $u_p$.
All prime coordinates are therefore
detected, so
$u\mapsto w^u$ embeds $\widehat\Z$.

The already proved full embedding
(F.11.10) preserves this conclusion for
$j_*(\gamma w\gamma^{-1})$ in $\widehat P$.
Its whole procyclic closure lies in all
thirty-six kernels by continuity.
Finally this is a nontrivial discrete
pure mapping class, so the general
procyclic exclusion of F.10 applies and
gives the asserted trivial intersection
with $C_{\widehat P}(A_5)$.
No element of that residual centralizer
is produced, and no equality between it
and this joint kernel is asserted.
\end{proof}

\begin{claim}{F.12}{Proved}{The common involution quotient and its winding equations}
Use the actual action and source in F.11.  Put
\[
 \tau=(1\ 5)(3\ 4),\qquad
 \beta=(1\ 3)(4\ 5),\qquad
 Y=\mathbf P^1\setminus\{0,\infty,\pm\mu_5\},
               \qquad t=s^2 .
                                                               \tag{F.12.1}
\]
The map $f$ in F.11 induces a split injection on
abelianizations, including their full profinite
completions.  Consequently its completed kernel
lies in the closed commutator subgroup of the
specified $\widehat F_6$.

There is an explicit, fully injective map
\[
 \widehat\phi:\widehat{\pi_1(Y,s_0)}
                     =\widehat F_{11}
                   \longrightarrow\widehat{P\Gamma}_{0,8},
 \qquad
 s_0=i(\theta^2-2)\zeta,\quad s_0^2=t_0,
                                                               \tag{F.12.2}
\]
induced by the configuration
\[
 \phi(s)=(0,\infty,1,v_1(s),\ldots,v_5(s)),
 \qquad
 v_{j+1}(s)=
       \left(\frac{\zeta^j-s}{\zeta^j+s}\right)^2
                         \quad(0\leq j\leq4).
                                                               \tag{F.12.3}
\]
The actual quotient action of $\beta$ induces
the following based involution on this normalized
eight-point moduli space:
\[
 R(v_1,v_2,v_3,v_4,v_5)
       =(v_2/v_3,\ v_2,\ v_2/v_1,\ v_2/v_5,\ v_2/v_4).
                                                               \tag{F.12.4}
\]
The residual centralizer in F.7 is exactly the
image, under the specified covering inclusion,
of the full equalizer
\[
 \mathcal E=
 \{u\in\widehat{\pi_1(Y,s_0)}:
                    \widehat\phi(u)=R_*\widehat\phi(u)\}.
                                                               \tag{F.12.5}
\]
In particular this identifies all the groups and
maps in the remaining fixed-intersection problem.

Write $e_j^\pm$ for the positively oriented source
meridian classes at $\pm\zeta^j$.  The abelian
image of every $u\in\mathcal E$ has the form
\[
        [u]=\sum_{j=0}^4 a_j(e_j^+-e_j^-),\qquad
        a_0-a_1+a_2=0,\quad -a_1+a_3+a_4=0
                  \quad\hbox{in }\widehat\Z.
                                                               \tag{F.12.6}
\]
These are exactly the restrictions of the twenty
target winding equalities to that anti-invariant
source submodule.  Their kernel is the explicitly
specified free rank-three $\widehat\Z$-module
in the proof; no equality with the actual image
of $\mathcal E$ is asserted.

For the element $w$ in F.11, the lift to this
double-cover source has
$[w]=e_1^- -e_1^+$.
Its first target winding defect is exactly four.
Hence its entire procyclic closure meets
$\mathcal E$ trivially, by this specific
abelian quotient computation.
Triviality of the full equalizer (F.12.5)
remains unresolved.
\end{claim}
\passage{Section 30, folios 194-196, scans 383-387; Part F, nonsolvable A5 centralizer comparison on the twelve-punctured sphere}


\begin{proof}
\emph{Six exact characters for the seven-point map.}
On the target of F.11.2, use the last labelled
point as infinity and normalize the first two
finite points to $0,1$.  Denote the resulting
six finite coordinates by $z_0,\ldots,z_5$,
so $z_0=0,z_1=1$.  Their pullbacks under $f$ are
\[
 (z_0,\ldots,z_5)=
 \frac{(1,\zeta^2,\zeta^4,t,t\zeta,t\zeta^4)-1}
                         {\zeta^2-1}.
                                                               \tag{F.12.7}
\]
For distinct $i,j$, the function $z_i-z_j$
is a globally defined nonzero function on
this labelled moduli space, so its winding
defines a homomorphism to $\Z$.
Take the six pairs, in this order,
\[
            (3,4),\ (3,0),\ (5,0),\
                         (3,1),\ (5,1),\ (4,0).
                                                               \tag{F.12.8}
\]
The six pullbacks are nonzero constant multiples
of
\[
             t,\quad t-1,\quad t-\zeta,\quad
                       t-\zeta^2,\quad t-\zeta^3,\quad t-\zeta^4.
                                                               \tag{F.12.9}
\]
Thus the induced homomorphism
$\ell:P\Gamma_{0,7}\to\Z^6$ satisfies
$\ell f_*=\operatorname{ab}$ on
$\pi_1(X_5)$, in its ordinary finite-puncture
meridian basis.  Completion gives the same
identity with $\widehat\Z^6$.
This proves the asserted split injection and
the closed-commutator containment.
In particular the residual centralizer, which
belongs to this kernel by F.7.9 and F.11,
lies in the index-two subgroup corresponding
to $t=s^2$.
The source is exactly $Y$ in (F.12.1);
finite-index subgroups have their full
induced profinite topology.

\emph{The actual common-involution quotient.}
For the family in F.11, the involution
$\tau_t(z)=t/z$ has fixed points $s,-s$.
On $Y$, label them in that order.
The actual quotient map, including these
branch labels, is
\[
                  Q_s(z)=\left(\frac{z-s}{z+s}\right)^2 .
                                                               \tag{F.12.10}
\]
It sends the two branch points to $0,\infty$.
The original pole pair $\{0,\infty\}$ goes to
$1$.  Its other five fibres on the original
twelve punctures are
\[
    \{\zeta^j,s^2\zeta^{-j}\},\qquad 0\leq j\leq4,
\]
and their images are precisely (F.12.3).
No two of the eight target points coincide
for $s\in Y$.
Indeed $v_{j+1}$ is $0$ or infinity only at
$s=\zeta^j$ or $-\zeta^j$, and equals $1$
only at $s=0,\infty$.
For $i\ne j$, the other collision equation
is $s^2=\zeta^{i+j}$, whose two solutions
again lie in $\pm\mu_5$.

Here is the full completion comparison for
this quotient.  The discrete normalizer
sequence for $\langle\tau\rangle$ is the
specified sequence F.7.5:
\[
 1\longrightarrow\langle\tau\rangle
 \longrightarrow N_M(\langle\tau\rangle)
 \longrightarrow Q\longrightarrow1 .
                                                               \tag{F.12.11}
\]
Here $Q$ is the group of the quotient sphere
with its six punctures and two branch points,
with their types retained.  It has finite
index in the mapping class group with these
decorations.  The prime-order comparison and
the full topology in F.7.4--5 identify its
completed source with the actual completed
normalizer.  Consequently completion of
(F.12.11) has exactly the displayed finite
deck kernel and has target $\widehat Q$.
The deck involution permutes the twelve
original punctures.  Its intersection with
the completed pure centralizer is therefore
trivial, as detected by the finite permutation
quotient.  Descent is injective on
$C_{\widehat P}(\tau)$.
The finite-index inclusion of $Q$ supplies
its full embedding in the decorated
eight-point mapping class completion.

The map from $\widehat{\pi_1(Y)}$ into
$\widehat P$ is the index-two covering
inclusion followed by F.11.10.
It is injective and its image centralizes
$\tau$.  On $Y$ the two branch labels do
not exchange, so the quotient map lands
in $\widehat{P\Gamma}_{0,8}$.
Its composite is exactly
$\widehat\phi$ in (F.12.2), by (F.12.10).
Both components are injective.  This
proves the full injectivity in (F.12.2),
with all branch labels retained.

\emph{The exact second involution and based naturality.}
In the twelve-label table of F.11, the six
$\tau$-orbits, in the order used for the
quotient punctures $1,v_1,\ldots,v_5$, are
\[
 \{0,10\},\quad\{1,4\},\quad\{3,8\},\quad
              \{5,11\},\quad\{6,9\},\quad\{2,7\}.
                                                               \tag{F.12.12}
\]
The element $\beta=(1\ 3)(4\ 5)$ commutes
with $\tau$, is outside $D_{10}$, and
together with $D_{10}$ generates $A_5$.
Its permutation on these six orbits is
\[
              p=(0\ 2)(1\ 3)(4\ 5).
                                                               \tag{F.12.13}
\]
It exchanges the two branch points.
These facts are checked in the same actual
M\"obius action as F.11, not merely in a
separate permutation representation.
Specifically the word
$b^{-1}aba b^{-1}ab$ has the named action
of $\beta$, sends $s_0$ to $-s_0$ and
$-s_0$ to $s_0$, and in the coordinate
$(z-s_0)/(z+s_0)$ has the form $c/z$.
The certificate verifies $c^2=v_2(s_0)$.
Thus its actual quotient representative is
$V\mapsto v_2(s_0)/V$.

For general normalized quotient coordinates
write $v_0=1$.
After applying the label permutation and
renormalizing the branch values to
$0,\infty$ and puncture zero to $1$,
the new coordinate at puncture $i$ is
$v_{p(0)}/v_{p(i)}$.
This is precisely (F.12.4).
At $s_0$ the actual quotient representative
just specified proves
$R(\phi(s_0))=\phi(s_0)$.
The based action $R_*$ is therefore fixed
without an unspecified base path, and
agrees with conjugation by that quotient
mapping class.
The source basepoint comparison with F.11
uses the unique lift, starting at $s_0$,
of the already specified straight segment
from $t_0$ to $2+2i$.

Let $\jmath:\widehat{\pi_1(Y)}\to\widehat P$
denote the actual covering inclusion above.
For the residual centralizer, (F.12.7)--(F.12.9)
show that its elements have unique preimages
under $\jmath$.
The descent of each such element commutes
with the quotient of $\beta$, proving that
its preimage lies in $\mathcal E$.
Conversely, for $u\in\mathcal E$, both
$\jmath(u)$ and
$\beta\jmath(u)\beta^{-1}$ lie in
$C_{\widehat P}(\tau)$.
Their descents agree by (F.12.5).
The injectivity proved from (F.12.11)
therefore gives
$\beta\jmath(u)\beta^{-1}=\jmath(u)$.
The element $\jmath(u)$ already centralizes
$D_{10}$, so it centralizes their generated
group $A_5$.
This proves both directions of (F.12.5);
the resulting isomorphism is continuous
with continuous inverse on these compact
subgroups.

\emph{All winding equations on the source homology.}
Use the finite source meridians at
$0,\zeta^j,-\zeta^j$ as a basis of
$H_1(Y;\widehat\Z)$, omitting the infinity
meridian.  The map $s\mapsto s^2$ sends
the zero meridian to twice the zero
meridian downstairs, and sends each
$e_j^\pm$ to the downstairs meridian at
$\zeta^{2j}$.
The downstairs abelianization of an element
of $\mathcal E$ vanishes by the first part
of the proof.  Since multiplication by
two on $\widehat\Z$ is injective, its
zero coefficient is zero and its two
coefficients above each $\zeta^{2j}$ sum
to zero.  This gives the first expression
in (F.12.6).

For clarity the remaining calculation is
entirely integral.  Take the twenty winding
characters on the normalized eight-point
space:
\[
 L_i=\operatorname{wind}(v_i),\quad
 M_i=\operatorname{wind}(v_i-1)\quad(1\leq i\leq5),
 \qquad
 N_{ij}=\operatorname{wind}(v_i-v_j)\quad(1\leq i<j\leq5).
                                                               \tag{F.12.14}
\]
These give its usual integral first-cohomology
basis.  One can see the basis directly by
forgetting the last variable in turn: the
fibre is the plane with the earlier points
and $0,1$ removed, its meridians are detected
by the newly added differences, and pure
monodromy fixes their abelian classes.
The successive ranks are $2,3,4,5,6$.

The exact functions in (F.12.3) satisfy
\[
 \begin{split}
 v_{j+1}-1
     &=\frac{-4\zeta^j s}{(\zeta^j+s)^2},\\
 v_{i+1}-v_{j+1}
     &=\frac{4s(\zeta^i-\zeta^j)(\zeta^{i+j}-s^2)}
             {(\zeta^i+s)^2(\zeta^j+s)^2}
                         \qquad(i\ne j).
 \end{split}                                                   \tag{F.12.15}
\]
Their zero and pole orders at the twelve
source punctures determine the pullbacks.
On an anti-invariant class
$\sum a_j(e_j^+-e_j^-)$ the resulting values
are
\[
       L_{j+1}=4a_j,\qquad M_{j+1}=2a_j,\qquad
                   N_{i+1,j+1}=2(a_i+a_j).
                                                               \tag{F.12.16}
\]
The extra zeros of the last expression,
at the two points with $s^2=\zeta^{i+j}$,
cancel on an anti-invariant class.

Put $A_0=0$ and $A_{j+1}=a_j$.
For $v_0=1$, the winding of any difference
$v_i-v_j$ is uniformly $2(A_i+A_j)$,
and that of $v_i$ is $4A_i$.
Under $R$, these functions are transformed
using $v_i\mapsto v_2/v_{p(i)}$.
The twenty equalities of winding values
are consequently equivalent to
\[
                   A_i+A_{p(i)}=A_2
                             \qquad(0\leq i\leq5).            \tag{F.12.17}
\]
Here necessity already follows from
$M_1$ and $M_4$, whose defects are
$2(a_0-a_1+a_2)$ and
$2(-a_1+a_3+a_4)$.
Multiplication by two is injective.
For sufficiency, substitute (F.12.17)
in the winding of $v_2/v_{p(i)}$,
$(v_2-v_{p(i)})/v_{p(i)}$, and
$v_2(v_{p(j)}-v_{p(i)})/(v_{p(i)}v_{p(j)})$;
these are exactly the three types in
(F.12.14).
The only two nontrivial relations in
(F.12.17) are (F.12.6).
Their kernel has the following integral
basis, which also gives its full
$\widehat\Z$-basis:
\[
          (-1,0,1,0,0),\qquad
          (1,1,0,1,0),\qquad
          (1,1,0,0,1).
                                                               \tag{F.12.18}
\]
These equations are necessary for
membership in $\mathcal E$; their kernel
is only the module defined by the
displayed winding equations.

\emph{The concrete detector for the element of F.11.}
At $t_0$ the chosen root is
$s_0=i(\theta^2-2)\zeta$.
The lift of the straight path to $b=2+2i$
ends at the negative square root of $b$:
throughout that path the real part of
$t$ is positive, and $s_0$ is the root
with negative real part.
The whisker and the small square for
$m_{\zeta^2}$ lie in the upper half-plane.
Continuation on this sheet therefore
lifts that meridian to a positive
meridian at $-\zeta$.
Its lift on the other sheet is a meridian
at $+\zeta$.

The second commutator factor
$m_\zeta m_{\zeta^3}m_0$ has winding one
about zero, and hence exchanges the
two sheets.  In the lifted commutator,
its path and its reverse cancel in
abelianization.  The two remaining
meridians give exactly
\[
                         [w]=e_1^- -e_1^+ .
                                                               \tag{F.12.19}
\]
By (F.12.4), $R^*L_1=L_2-L_3$.
The winding difference for this class
is therefore
\[
       (L_1-R^*L_1)\widehat\phi(w)
               =4(a_0-a_1+a_2)=4.
                                                               \tag{F.12.20}
\]
For a full profinite power $w^u$, the
continuous extension of this character
has value $4u\in\widehat\Z$.
It vanishes only for $u=0$.
This is an explicit family of finite
cyclic detectors for the procyclic
exclusion, and uses the actual common
involution quotient.

The executed certificate
\path{code/part_a/c5_common_reflection_f12.py}
verifies the common finite labels, the
M\"obius involution and its branch
normalization, the six-character splitting,
the twenty-by-twelve integral divisor
matrix, the exact involution on the
twenty winding characters, and the
restriction to the two rows in (F.12.6).
It also records the winding defect four.
Its larger unreduced matrices have source
rank ten and fixed-difference rank six;
no rank calculation is used to infer
triviality of the full equalizer.
\end{proof}

\begin{claim}{F.13}{Proved}{The exact homology kernel and the retraction obstruction}
Retain the specified injection
$\widehat\phi$ of F.12.2.  Write $e_0,e_\infty,e_j^+,e_j^-$
for the twelve positive source cusp classes, with the
single relation
\[
 e_0+e_\infty+\sum_{j=0}^4(e_j^++e_j^-)=0.
                                                               \tag{F.13.1}
\]
For $A=\Z$, $\widehat\Z$, or $\Z_p$ at any prime $p$,
the actual first-homology map induced by $\phi$ has exact
kernel
\[
 \ker\phi_{*,1}
   =A\,k,\qquad
 k=e_0-e_\infty
   =2e_0+\sum_{j=0}^4(e_j^++e_j^-).
                                                               \tag{F.13.2}
\]
Here the second expression uses the genuine basis
omitting $e_\infty$.  In particular $k$ is primitive
over $\Z$ and has nonzero reduction modulo every prime.

The discrete map, its full profinite completion, and its
maps to maximal pro-$p$ completions have no homomorphic
left inverse on the stated domains, continuous where
profinite groups are involved.  The same obstruction
holds after every further target-forgetting map.
No assertion about a finite-index target subgroup, the
injectivity of a forgotten map, or the full equalizer
F.12.5 is made.
\end{claim}
\passage{Section 30, folios 194-196, scans 383-387; Section 31, folios 198-199 verso, scans 391-394; Part F(b), exact residual centralizer and the specified common-involution quotient}

\begin{proof}
To prove the kernel equality, write an arbitrary class
in that basis as
\[
 z=c e_0+\sum_{j=0}^4(p_j e_j^++m_j e_j^-).
\]
The twenty target winding characters are the actual
integral basis of F.12.14.  Their pullbacks from
F.12.3 and F.12.15 give
\[
 \begin{aligned}
 L_{j+1}(\phi_*z)&=2(p_j-m_j),\\
 M_{j+1}(\phi_*z)&=c-2m_j,\\
 N_{i+1,j+1}(\phi_*z)
   &=c+p_k+m_k-2m_i-2m_j,
      &2k&\equiv i+j\pmod5\quad(i\ne j).
 \end{aligned}                                                \tag{F.13.3}
\]
Multiplication by two is injective in each of the
specified coefficient rings.  Vanishing of the first
two rows therefore gives $p_j=m_j=t$ for all $j$ and
$c=2t$.  The third row then vanishes identically.
This proves (F.13.2), without identifying an actual
group kernel with its abelian image.

For completeness, at a specified primary coordinate
$p$, the vector $k$ has ten entries equal to one in
the displayed $\Z_p$-basis.  It is consequently
nonzero over $\Z_p$ and modulo $p$.  In the full
$\widehat\Z$-module, multiplying $k$ by the idempotent
whose $p$-coordinate is one and whose other primary
coordinates are zero gives a nonzero kernel vector
supported exactly at that primary coordinate.

There is no homomorphic left inverse of the discrete
map $\phi_*:\pi_1(Y)\to P\Gamma_{0,8}$ and no continuous
homomorphic left inverse of its full profinite map.
Indeed such a left inverse would induce a left inverse
on abelianizations, contradicting (F.13.2).
More precisely, for every prime $p$ there is no
continuous homomorphism
\[
 r_p:\widehat{P\Gamma}_{0,8}
                  \longrightarrow\pi_1(Y)^{(p)}
 \quad\hbox{with}\quad
 r_p\widehat\phi=\pi_p,
                                                               \tag{F.13.4}
\]
where $\pi_p$ is the actual maximal pro-$p$ quotient
map.  Apply abelianization modulo $p$ to this equality:
its left side kills $k\bmod p$, whereas its right side
is the identity on the source's $\F_p$-homology.
The same argument rules out a left inverse of the
induced map between maximal pro-$p$ completions.

Every homomorphism obtained by first applying
$\widehat\phi$ and then forgetting any target labels
also kills $k$ on first homology, by functoriality.
Thus none of those composites has a homomorphic
left inverse.  This does not decide which such
composites are injective, nor does it decide the
equalizer F.12.5.  The obstruction is stated for
these displayed domains; passing to a finite-index
target subgroup would change its first-homology
map and is not covered by this argument.

The separate exact certificate
\path{code/part_a/common_reflection_retraction.py}
constructs all twenty divisor rows from (F.13.3),
checks the source puncture-sum relation and the
primitive kernel vector, computes rank ten over
$\Q$, and checks the explicit nonzero kernel vector
in the listed finite primary reductions.  Its
finite reductions verify the displayed formulas;
the uniform statement for every prime follows
from the coefficient-one argument above.
\end{proof}

\begin{claim}{F.14}{Proved}{An injective seven-point comparison and its exact remaining defect}
Retain the based map $\phi$ and the involution $R$ of F.12.
Label its eight target points
\[
             0,\ \infty,\ v_0=1,\ v_1,\ldots,v_5 .
                                                               \tag{F.14.1}
\]
Let $q_{v_0}$ forget the ordinary point $v_0$, retaining
both branch labels $0,\infty$.  The actual composite
\[
 \widehat\psi=(q_{v_0})_*\widehat\phi:
       \widehat{\pi_1(Y,s_0)}
                     \longrightarrow\widehat{P\Gamma}_{0,7}
                                                               \tag{F.14.2}
\]
is injective.  With $v_1$ renormalized to one, its four
remaining coordinates and those of the second composite
$\psi_R=q_{v_0}R\phi$ are exactly
\[
 \begin{split}
 \psi(s)&=(v_2/v_1,\ v_3/v_1,\ v_4/v_1,\ v_5/v_1),\\
 \psi_R(s)&=(v_3,\ v_3/v_1,\ v_3/v_5,\ v_3/v_4).
 \end{split}                                                   \tag{F.14.3}
\]
Both maps have the same specified target basepoint.
Their second coordinates agree as actual functions.

Forgetting either branch point $0$ or $\infty$ from
(F.14.1) instead kills each of the two explicitly
whiskered source meridians at $s=0,\infty$ below.
Each meridian generates a full closed $\widehat\Z$
in the source.  Both branch-forgetting composites
factor through the specified filling
\[
 \widehat{\pi_1(Y)}
       \longrightarrow
 \widehat{\pi_1\bigl(\mathbf P^1\setminus\{\pm\mu_5\}\bigr)} .
                                                               \tag{F.14.4}
\]

There is also an exact description of the information
retained by (F.14.3).  Put $P_8=\widehat{P\Gamma}_{0,8}$,
write $K_i$ for its actual kernel of forgetting $v_i$,
and define
\[
 \begin{split}
 \mathcal E_7&=\{u:\widehat\psi(u)=\widehat\psi_R(u)\},\\
 B&=K_0\cap K_2,\\
 d(u)&=\widehat\phi(u)^{-1}R_*\widehat\phi(u)
                         \qquad(u\in\mathcal E_7).
 \end{split}                                                   \tag{F.14.5}
\]
Then $d$ takes values in $B$, satisfies
\[
 \begin{split}
 R_*d(u)&=d(u)^{-1},\\
 d(uv)&=\widehat\phi(v)^{-1}d(u)\widehat\phi(v)d(v),\\
 \mathcal E&=\{u\in\mathcal E_7:d(u)=1\},
 \end{split}                                                   \tag{F.14.6}
\]
where $\mathcal E$ is the full residual equalizer F.12.5.
In the actual point-pushing identification $K_0=\widehat F_6$,
$B$ is precisely the closed normal subgroup generated by
the meridian about $v_2$.  These statements retain both
distinct forgetting composites; they do not cancel
$q_{v_0}$ from the equalizer.
\end{claim}
\passage{Section 30, folios 194-196, scans 383-387; Section 31, folios 198-199 verso, scans 391-394; Part F(b), residual nonsolvable centralizer and exact quotient comparisons}


\begin{proof}
\emph{Filling the two order-five poles.}
Let $S_{12}$ be the sphere punctured at the labelled
configuration
\[
                  0,\infty,\mu_5,t_0\mu_5,
\]
and let $S_{10}$ be obtained by filling the first two
punctures.  The named rotation $d(z)=\zeta z$ acts on
both surfaces.  The two fixed poles have become
unmarked ramification points on $S_{10}$.
After removal of those ramification points, both
covering problems are literally the same unramified
cover
\[
 S_{12}\longrightarrow
       X=\mathbf P^1\setminus\{0,\infty,1,t_0^5\},
                           \qquad z\longmapsto z^5 .
                                                               \tag{F.14.7}
\]
Its four meridians have the same monodromy values
$1,-1,0,0$ in $C_5$ as F.11.8.

Write $N_{12}$ and $N_{10}$ for the discrete full
normalizers of this specified $C_5$ in the two
mapping class groups.  The actual covering normalizer
sequences of F.7.5 form the commuting diagram
\[
\begin{array}{ccccccccc}
1&\longrightarrow&C_5&\longrightarrow&N_{12}
 &\longrightarrow&Q&\longrightarrow&1\\
 &&\Vert&&\downarrow\mathrm{fill}&&\Vert\\
1&\longrightarrow&C_5&\longrightarrow&N_{10}
 &\longrightarrow&Q&\longrightarrow&1 .
\end{array}                                                     \tag{F.14.8}
\]
Here $Q$ is the same subgroup of the mapping class
group of $X$: its elements preserve the covering
kernel and the two branch labels as a set, and
preserve the two free-puncture labels as a set.
The nonzero and zero monodromy values already
distinguish these two types.  The extra description
of a branch point as an original puncture on
$S_{12}$ therefore imposes no new permutation
condition.  On $S_{10}$ those same labels are retained
as branch decorations.  Descent of a filled
representative is the same homeomorphism of $X$,
which proves that both right vertical maps are
literally the identity after these identifications.
The kernel map is the named deck action itself.
The exact-sequence argument proves that the middle
filling map is an isomorphism.

Both finite groups in (F.14.8) are the same order-five
group.  The prime-order normalizer comparison and
full induced topology established in F.7.4--5 identify
their completed normalizers with their actual
closures in the ambient completions.  Thus completion
of the middle isomorphism is the actual injective
filling map on the completed normalizer.  In
particular it is injective on the source subgroup
from F.11 and on its index-two subgroup
$\widehat{\pi_1(Y)}$.
The ambient filling map is defined on the stabilizer
of the named pole pair; the normalizer preserves
that pair, and the source is pure, so this restriction
uses the stated ambient maps.

\emph{The commuting quotient square.}
The source still centralizes the actual involution
$\tau_t(z)=t/z$ after the poles are filled.
On $S_{10}$ this involution permutes the remaining
punctures in five pairs.  It has the same two
unmarked fixed points $s,-s$, labelled over $Y$.
Its decorated quotient has these two branch
points and five ordinary punctures.
The quotient map is still the literal function
$Q_s(z)$ in F.12.10.  The removed pole pair has
image $v_0=1$ and every remaining pair has the
same image $v_j$.

The completed involution-normalizer quotient again
has exactly its named finite deck kernel, by
F.7.4--5.  That deck involution permutes the ten
punctures nontrivially.  It has trivial intersection
with the completed pure subgroup, detected by
the finite permutation map.  Hence descent is
injective on the pure completed centralizer of
$\tau$ in the ten-puncture group.
The branch-labelled source lands in the pure
seven-puncture target.  Write $C^{\rm br}_{12}$
and $C^{\rm br}_{10}$ for the subgroups of the
two pure completed $\tau$-centralizers whose
descents also fix its two branch labels
individually.  Filling the pole pair preserves
this condition.
We have the actual commuting square
\[
\begin{array}{ccc}
 \widehat{\pi_1(Y)}
   &\longrightarrow&C^{\rm br}_{12}\\
 \big\downarrow\widehat\psi
   &&\big\downarrow\mathrm{fill}_{\{0,10\}}\\
 \widehat{P\Gamma}_{0,7}
   &\longleftarrow&C^{\rm br}_{10}.
\end{array}                                                     \tag{F.14.9}
\]
The upper arrow is the F.11--F.12 source inclusion;
the lower arrow is the just specified decorated
descent.  On that source the right arrow is
injective by (F.14.8), and the lower arrow is
injective on the pure subgroup.  Their composite
is precisely $(q_{v_0})_*\widehat\phi$, by the
pointwise formula for $Q_s$.  This proves
(F.14.2).

\emph{Both seven-point maps and their common coordinate.}
After forgetting $v_0$, the three labels $0,\infty,v_1$
give the unique normalization fixing $0,\infty$
and sending $v_1$ to one.  This gives the first
line of (F.14.3).
Substitution of F.12.4 followed by this same
normalization gives its second line.
Since $R(\phi(s_0))=\phi(s_0)$ with the exact
based representative already proved in F.12,
no additional target base path is required.
For clarity their common coordinate is
\[
 \frac{v_3(s)}{v_1(s)}
   =\left(
       \frac{(\zeta^2-s)(1+s)}
            {(\zeta^2+s)(1-s)}
     \right)^2 .
                                                               \tag{F.14.10}
\]
Thus forgetting all but its three normalization
labels and this coordinate gives equal composites
to $M_{0,4}$ as actual maps of spaces.
No injectivity of that four-point composite is used.

\emph{The branch-forgetting maps and based meridians.}
Forget the branch point $0$ instead, retaining
$\infty,v_0=1,v_1,\ldots,v_5$.
Normalize $v_0$ to zero and $v_1$ to one, keeping
infinity fixed.  Its four coordinates are
\[
 u_j(s)=\frac{v_{j+1}(s)-1}{v_1(s)-1}
       =\frac{\zeta^j(1+s)^2}{(\zeta^j+s)^2},
                \qquad 1\le j\le4.
                                                               \tag{F.14.11}
\]
They extend over $s=0$ and $s=\infty$, with respective
values $\zeta^{-j}$ and $\zeta^j$.  At either end
these four values are pairwise distinct and avoid
$0,1,\infty$.  Thus this is an extension into the
actual seven-point configuration space, not merely
a finite limit in its compactification.
The other branch-forgetting map has the same
property: the actual target coordinate inversion
$V\mapsto1/V$ exchanges the two branch labels and
satisfies $1/v_j(s)=v_j(-s)$.  Apply the same
normalization and extension with $-s$.
This proves the space-level factorization through
filling the two indicated source punctures and
hence (F.14.4) on full completions.

Here are explicit based choices.  Put
$\epsilon=1/10$ and use the polygonal whiskers
\[
 \begin{split}
 \omega_0&=(s_0,\,-2+2i,\,\epsilon+2i,\,\epsilon),\\
 \omega_\infty&=(s_0,\,-2+2i,\,10+2i,\,10).
 \end{split}                                                   \tag{F.14.12}
\]
From the exact expression
$s_0=\theta-\theta^3/2+i/2$ and $\theta>19/10$,
its real part is less than $-1$.
The first segments therefore lie to the left of
the unit circle; the horizontal segments have
imaginary part two.  The final segment of
$\omega_\infty$ lies outside that circle.
The final segment of $\omega_0$ has real part
$1/10$, different from zero and from every real
part of a point of $\pm\mu_5$.  Indeed their
nonzero real parts have absolute value at least
$(\sqrt5-1)/4>1/10$.
Both whiskers lie in $Y$.

Use the local circles
$s=\epsilon e^{2\pi i t}$ and $s=10e^{-2\pi i t}$,
respectively, and return along the reverse
whisker.  These give positive based source
meridians $m_0,m_\infty$.
The extended maps in (F.14.11) carry each local
circle to a loop filled by its actual disk
at the added source point; composing with the
image whiskers gives based null-homotopies.
The same holds for the other branch-forgetting
map.
In source abelianization $m_0$ is $e_0$ and
$m_\infty$ is $e_\infty=-e_0-\sum(e_j^++e_j^-)$.
A finite-puncture winding character takes each
of these meridians to $1$ or $-1$.
Its full continuous extension consequently
identifies the closure of each cyclic subgroup
with a full $\widehat\Z$ and detects every nonzero
profinite exponent.  Thus these are genuine
nontrivial full procyclic subgroups of both
branch-forgetting kernels.

\emph{The exact point-pushing defect.}
In the target labels of (F.14.1), the permutation
underlying $R$ exchanges $v_0$ with $v_2$.
Naturality of forgetting with this actual
labelled map gives
\[
                    R_*(K_0)=K_2.                            \tag{F.14.13}
\]
For $u\in\mathcal E_7$, the defining equality in
(F.14.5) says exactly that $d(u)\in K_0$.
Since $R_*^2=1$, direct multiplication gives
$R_*d(u)=d(u)^{-1}$.
Equation (F.14.13) then also puts $d(u)$ in
$K_2$.  The second identity in (F.14.6) follows
by expanding the definition on the product $uv$.
Both $K_i$ are normal, so the displayed conjugation
preserves $B$.  The third identity is the defining
full equality F.12.5, with both implications now
retained inside the explicitly specified
seven-point equalizer.

Finally the actual Birman point-pushing diagrams
identify $K_0$ with the completion of the
fundamental group of the sphere punctured at
the other seven labelled points, based at
$v_0$.  Forgetting $v_2$ restricts on this kernel
to the ordinary puncture-filling map
$\widehat F_6\to\widehat F_5$.
The point-pushing subgroup in the seven-point
target after forgetting $v_2$ is embedded.  The
actual completed Birman square is
\[
\begin{array}{ccccc}
 K_0=\widehat F_6&\longrightarrow&P_8
    &\longrightarrow&\widehat{P\Gamma}_{0,7}\\
 \downarrow\mathrm{fill}_{v_2}
    &&\downarrow\mathrm{forget}_{v_2}
    &&\downarrow\mathrm{forget}_{v_2}\\
 K'_0=\widehat F_5&\longrightarrow&
       \widehat{P\Gamma}_{0,7}
    &\longrightarrow&\widehat{P\Gamma}_{0,6}.
\end{array}                                                     \tag{F.14.14}
\]
The horizontal quotient maps forget $v_0$;
the lower middle group still retains $v_0$.
Consequently the kernel of that
restriction is exactly $K_0\cap K_2$.
The map of free groups kills the named peripheral
meridian about $v_2$ and keeps the other five
finite-puncture generators; the same finite
presentation and the universal property of
completion identify its full kernel with the
closed normal subgroup generated by that
meridian.

These full Birman maps can also be obtained
directly from normalized configuration spaces.
Forgetting one variable of an ordered configuration
in $\mathbf C\setminus\{0,1\}$ is a fibration
with the stated punctured-plane fibre, and has
a continuous section putting the forgotten
variable on the positive real axis beyond
$2+\sum|z_j|$.
The base and fibre are aspherical by induction
through these fibrations, giving the actual
discrete exact sequences.
The split group extension induces the full
profinite topology on its finitely generated
free kernel: use its characteristic finite-index
subgroups and the finite action on each resulting
finite quotient in the semidirect product.
Completion is then exact, and the forgetting
squares are the completions of the displayed
topological squares.
This justifies the kernel comparison on the
full completions used above.

The executed certificate
\path{code/part_a/common_reflection_forgetting_f14.py}
checks the two normalized composite formulas,
their common coordinate, the two endpoint
extensions, the exact source basepoint formula,
and the fourteen-by-twenty matrix of the
ordinary-point forgetting map on winding
characters.  The injectivity in (F.14.2)
uses the actual normalizer and covering maps,
not these homology checks.
\end{proof}

\begin{claim}{F.15}{Proved}{The two-point kernel and its full involution cocycles}
Retain the actual based configuration of F.12--14 and put
$c_0=v_2(s_0)$, $x_0=v_0=1$, $y_0=c_0$.  On the sphere
with the other six labels retained, set
\[
\begin{split}
 M&=\mathbf P^1\setminus
       \{0,\infty,v_1(s_0),v_3(s_0),v_4(s_0),v_5(s_0)\},\\
 r(z)&=c_0/z,\qquad
 X=\operatorname{Conf}_2(M),\qquad
 \mathcal R(x,y)=(r(y),r(x)).
\end{split}                                                     \tag{F.15.1}
\]
Thus $\mathcal R(x_0,y_0)=(x_0,y_0)$.  Write $a_+,a_-$
for the two roots of $z^2=c_0$, which are unmarked points of
$M$.  Write $\mathsf r$ for the actual eight-puncture
mapping class of $r$, so conjugation by $\mathsf r$ induces
the based automorphism $R_*$ of F.12.

There are the following actual completed homomorphisms:
\[
\begin{split}
 J&=\pi_1(X,(x_0,y_0)),\\
 L_0&=\pi_1(M,x_0),\qquad L_2=\pi_1(M,y_0),\\
 1\longrightarrow B\longrightarrow\widehat J
   \xrightarrow{\ p\ }\widehat L_0\times\widehat L_2
     \longrightarrow1 .
\end{split}                                                     \tag{F.15.2}
\]
Here $\widehat J$ is the kernel in
$P_8=\widehat{P\Gamma}_{0,8}$ of forgetting the pair
$\{v_0,v_2\}$; $p$ records the two labelled point tracks,
and its kernel is exactly $B=K_0\cap K_2$ of F.14.
Let $\rho=r_*:\widehat L_0\longrightarrow\widehat L_2$.
The induced involution and its fixed subgroup are
\[
 R_*(u,v)=(\rho^{-1}(v),\rho(u)),\qquad
 (\widehat L_0\times\widehat L_2)^{R_*}
        =\{(u,\rho(u)):u\in\widehat L_0\}.                       \tag{F.15.3}
\]

The actual fixed-configuration map
\[
 i:M\setminus\{a_+,a_-\}\longrightarrow X,\qquad
                         z\longmapsto(z,r(z))                 \tag{F.15.4}
\]
induces a full profinite embedding, and
\[
\begin{split}
 C_{\widehat J}(\mathsf r)
   &=\widehat i_*\,
        \widehat{\pi_1(M\setminus\{a_+,a_-\},x_0)}
       \ \cong\widehat F_7,\\
 p\bigl(C_{\widehat J}(\mathsf r)\bigr)
   &=(\widehat L_0\times\widehat L_2)^{R_*},\\
 C_B(\mathsf r)
   &\cong\ker\bigl(
       \widehat{\pi_1(M\setminus\{a_+,a_-\},x_0)}
           \xrightarrow{\ \mathrm{fill}_{a_+,a_-}\ }
                  \widehat{\pi_1(M,x_0)}\bigr).
\end{split}                                                     \tag{F.15.5}
\]
In particular the last kernel is the closed normal subgroup
generated by the two based branch meridians in this
specified $\widehat F_7$.

Every full continuous $R_*$-cocycle in $B$ is a
$B$-coboundary.  More precisely, there is a homeomorphism
of compact spaces
\[
\begin{split}
 C_B(\mathsf r)\backslash B
    &\xrightarrow{\ \sim\ }
        Z^1(R_*,B):=\{\delta\in B:\delta R_*(\delta)=1\},\\
 C_B(\mathsf r)b&\longmapsto b^{-1}R_*(b).
\end{split}                                                     \tag{F.15.6}
\]
Consequently the quotient map $\pi:P_8\to P_8/B$ gives the
exact sequence
\[
 1\longrightarrow C_B(\mathsf r)
  \longrightarrow C_{P_8}(\mathsf r)
  \xrightarrow{\ \pi\ }(P_8/B)^{R_*}
  \longrightarrow1 .                                          \tag{F.15.7}
\]

For the exact equalizer $\mathcal E_7$ of F.14, set
\[
 \mathcal B_u=\{b\in B:b^{-1}R_*(b)=d(u)\}.
                                                               \tag{F.15.8}
\]
This is a nonempty left $C_B(\mathsf r)$-coset for every
$u\in\mathcal E_7$.  The centralizing representatives of
its actual $B$-coset are exactly
\[
 \bigl(\widehat\phi(u)B\bigr)\cap C_{P_8}(\mathsf r)
       =\widehat\phi(u)b^{-1}C_B(\mathsf r)
                 \qquad(b\in\mathcal B_u).                    \tag{F.15.9}
\]
The source embedding has trivial intersection with $B$.
It follows that the intersection of (F.15.9) with
$\widehat\phi(\widehat{\pi_1(Y)})$ is either empty or the
singleton $\{\widehat\phi(u)\}$, and is the latter
exactly for $u\in\mathcal E$.  These formulas determine
the correction and its ambiguity; triviality of the
full residual group $\mathcal E$ is still unresolved.
\end{claim}
\passage{Section 30, folios 194-196, scans 383-387; Section 31, folios 198-199 verso, scans 391-394; Part F(b), the exact residual centralizer and its common-involution quotient}


\begin{proof}
\emph{The labelled configuration and the two tracks.}
The actual quotient representative in F.12 is
$z\mapsto v_2(s_0)/z$.  Its permutation on the eight
labels is
\[
       (0\ \infty)(v_0\ v_2)(v_1\ v_3)(v_4\ v_5).
                                                               \tag{F.15.10}
\]
It has no fixed marked point.  Its two fixed points are
therefore precisely the unmarked $a_+,a_-$ of (F.15.1).
Deleting the invariant pair $\{v_0,v_2\}$ retains an
actual homeomorphism of the indicated six-puncture
sphere, with no additional base path.
Under point pushing, conjugation by its mapping class
acts on the ordered two-point fibre by $\mathcal R$.
This proves both the specified action and (F.15.3).

The discrete Birman fibration identifies $J$ with the
kernel of forgetting those two labels.  Projecting $X$
to one coordinate has fibre a sphere with seven
punctures, and gives a split extension
\[
              1\longrightarrow F_6
                \longrightarrow J
                \longrightarrow F_5
                \longrightarrow1 .                            \tag{F.15.11}
\]
One continuous section can be seen in a coordinate with
one retained puncture at infinity: put the second point
on the positive real axis beyond
$2+|x|+\sum_{z\in P_6\setminus\{\infty\}}|z|$.
A path in the fibre adjusts its basepoint to the
specified one.  The maps in (F.15.11), including that
adjustment, remain the ordinary labelled forgetful maps.

The finitely generated free kernel in (F.15.11) has its
full induced profinite topology.  Indeed, its
characteristic finite-index subgroups are cofinal among
its finite-index subgroups.  For any such subgroup, the
quotient of the free kernel is finite and the action of
the free quotient on it has finite image; the associated
finite semidirect product supplies the required finite
quotient of $J$.  Completing (F.15.11) is therefore
exact.
Map this completed sequence into the two successive
completed Birman sequences of F.14.14.  On its free
kernel and its free quotient the maps are the already
specified embeddings $\widehat F_6$ and $\widehat F_5$.
An element in the kernel of $\widehat J\to P_8$ must
first map trivially to $\widehat F_5$, then trivially to
$\widehat F_6$.  The map is injective.
Conversely, an element in the kernel of forgetting both
points can first be matched in the $\widehat F_5$
quotient by an element of $\widehat J$; the remaining
difference lies in the embedded $\widehat F_6$.
This proves the asserted equality with the actual
completed two-point kernel.

Each coordinate projection of $J$ records the track of
that named point.  The fibre of the first projection
maps onto $L_2$ by filling $x_0$, so the product track
map is onto $L_0\times L_2$.  The same is true after
completion.  The completion of a direct product is the
product of the full completions: any homomorphism to a
finite group restricts to homomorphisms on its two
commuting factors, and their continuous extensions
multiply to its unique continuous extension.
The kernel of the product track map is the simultaneous
kernel of the two point-forgetting restrictions.
Inside the embedded $\widehat F_6$ it is precisely
the kernel of filling the other named point.
The last Birman square in F.14 identifies it with
$K_0\cap K_2$, including its closed normal
meridian description.  This proves (F.15.2) without a
closure--intersection assumption.

\emph{The discrete fixed subgroup from the actual quotient cover.}
Consider the specified double cover obtained by deleting
the two fixed points of $r$:
\[
 M\setminus\{a_+,a_-\}\longrightarrow
   (M/\langle r\rangle)\setminus\{\overline a_+,
                                      \overline a_-\}.
                                                               \tag{F.15.12}
\]
The target is a sphere with three ordinary punctures
and two branch points removed.  Its group is $F_4$.
The monodromy takes the two branch meridians to the
nonidentity element of $C_2$ and takes the three
ordinary puncture meridians to the identity.
Its kernel is the actual $\pi_1$ of the source of
(F.15.12), of index two and rank seven.

Use the discrete covering-normalizer sequence F.7.5
for $r$ before and after forgetting the pair
$\{x_0,y_0\}$.  Downstairs this forgets its one ordinary
orbit point.  The kernel of this downstairs
forgetting map is the point-pushing $F_4$ just
specified.  A point-pushing isotopy lifts from the
identity through the double cover, fixing all retained
punctures individually.  Its endpoint fixes both
forgotten labels individually precisely when its
monodromy in (F.15.12) is trivial.
The finite deck kernel has trivial intersection with
the pure upstairs group, because $r$ exchanges the
retained punctures in three pairs.
Conversely, an element of $J$ that centralizes $r$
descends through this normalizer sequence.  After
forgetting the orbit point its descent is the identity,
since its upstairs filling is the identity and the
quotient map has only the displayed deck kernel.
It is therefore exactly one of these monodromy-zero
point pushes.  This proves
\[
 C_J(\mathsf r)
      =i_*\pi_1(M\setminus\{a_+,a_-\},x_0)
      \cong F_7.                                               \tag{F.15.13}
\]
The lift of its downstairs point track is the path
$(z,r(z))$.  Thus the equality uses the actual map
(F.15.4), not merely an abstract isomorphism of free
groups.  Projection to the first track is exactly
filling the two fixed points, and projection to the
second is its composite with $\rho$.

This $F_7$ also has its full induced profinite topology.
The downstairs point-pushing $F_4$ has its full topology
in the decorated quotient mapping class group by the
completed Birman argument of F.14.  Its index-two
monodromy kernel inherits its full topology.
The pure centralizer is a finite-index subgroup of the
normalizer of the specified order-two group, and its
descent is injective because the deck element is
nonpure.  Its image is finite index in the decorated
quotient, so this restriction is an isomorphism onto
that subgroup.  The prime-order normalizer topology in
F.7.4--5 consequently transports the full topology of
that $F_7$ to its actual subgroup in the eight-puncture
mapping class group.  It remains full in its subgroup
$J$ as well.  Hence the completion of (F.15.13)
is an embedding of this specified $\widehat F_7$ in
$\widehat J$.

\emph{One discrete conjugacy class of labelled sections.}
Let
\[
                       E=J\rtimes\langle\mathsf r\rangle
                                                               \tag{F.15.14}
\]
with the actual action already specified.
It embeds in the eight-puncture mapping class group:
the nontrivial coset is detected by the finite label
permutation (F.15.10).
Every order-two section in $E$ has that same permutation
and fills to the specified mapping class of $r$ on $M$.

We first justify pure conjugacy in the eight-puncture
group.  The cyclic case of Nielsen realization, in
the punctured-surface form of
\cite[Theorem 4 and Section IV, pp.~255--256]{C_Kerckhoff1983},
realizes such a section as an orientation-preserving
involution of the punctured sphere.  It extends over
the punctures.  With the invariant conformal structure,
filling gives a Riemann sphere, and a nontrivial
holomorphic involution has two fixed points.
Neither is marked, by (F.15.10).
Its quotient and that of $r$ are spheres with two
branch points and four specified ordinary orbit
points.  Choose an orientation-preserving quotient
homeomorphism matching those branch points and the
four named orbit pairs.  Both double covers have
monodromy one at the branch points and zero at the
ordinary points, so it lifts to an equivariant
homeomorphism.

The lift may interchange the two names inside some
of the four pairs.  Each mismatch can be corrected
independently: in the target quotient push that
ordinary orbit point once around one branch point,
holding the other orbit points and both branch points
fixed.  Lift the isotopy starting at the identity.
The endpoint commutes with $r$, interchanges exactly
the two labels in the selected orbit, and fixes the
other six labels individually.  The covering
monodromy proves these exact permutations.
Composing the required corrections yields a pure
conjugator.  Thus every section is $g\mathsf r g^{-1}$
with $g\in P\Gamma_{0,8}$.

The actual pure centralizer forgetting map is onto:
\[
 C_{P\Gamma_{0,8}}(\mathsf r)\longrightarrow
                              C_{P\Gamma_{0,6}}(r).
                                                               \tag{F.15.15}
\]
To prove this, take a class on the right.  The
discrete covering-normalizer map realizes it by a
homeomorphism commuting with the actual $r$.
Its value at $x_0$ is in the connected surface
$M\setminus\{a_+,a_-\}$.  Choose a path in that
surface from this value to $x_0$.  The quotient
path moves the one ordinary orbit point while
avoiding all branch points and retained punctures.
An isotopy moving that point, supported away from
those fixed labels, lifts from the identity to an
$r$-equivariant isotopy.  Its chosen lift carries
the image of $x_0$ to $x_0$ and the paired point to
$y_0$.  Composing its endpoint with the original
representative gives a pure eight-label centralizer
representative with the same six-label class.
This proves (F.15.15).

For a section $s=g\mathsf r g^{-1}$ in $E$, its
filled class is $r$, so the filled image of $g$
centralizes $r$.  Choose $c$ in the left side of
(F.15.15) with that same image.  Then $h=gc^{-1}$
belongs to $J$ and
$s=h\mathsf r h^{-1}$.  Consequently
\[
          \{\hbox{order-two sections of }E\}/J
                            \quad\hbox{has one element}.       \tag{F.15.16}
\]

\emph{The hypotheses for full profinite section comparison.}
We verify them for this particular $E$, rather than
assuming that goodness passes to an arbitrary
subgroup of a mapping class group.
The split sequence (F.15.11) and its exact completion
show that $J$ is residually finite: an element with
nontrivial image in $F_5$ is detected there, and an
element of the free kernel is detected by one of
the finite semidirect products used above.
The finitely generated group $J$ has cofinal
characteristic finite-index subgroups.  Taking
these to be invariant under the order-two
automorphism shows directly that the finite
extension $E$ is residually finite and that
$\widehat E=\widehat J\rtimes\langle\mathsf r\rangle$.

The free kernel and free quotient in (F.15.11)
are good.  With a finite coefficient module their
cohomology is finite and vanishes above degree one.
The discrete and profinite
Lyndon--Hochschild--Serre spectral sequences for
(F.15.11) therefore have isomorphic second pages:
the coefficients $H^b(F_6,A)$ for $b=0,1$ are
finite, so goodness of $F_5$ applies to them.
This proves goodness of $J$ and finiteness of its
cohomology with every finite module.
The same spectral-sequence comparison for
$1\to J\to E\to C_2\to1$ proves these assertions
for $E$.
The free-by-free group $J$ has a finite
two-dimensional classifying complex: take a
finite graph for the free kernel over a graph
for the free quotient, using the finitely many
monodromy homotopy equivalences on its edges.
Its universal cover is a tree of contractible
covering spaces.  In particular
$\operatorname{cd}(J)\leq2$.
The resulting finite-type resolutions and the
finite-group resolution for $C_2$ show that
cohomology of $E$ commutes with filtered unions
of coefficient modules.  A discrete
$\mathbf F_2[\![\widehat E]\!]$-module is the
union of its finite invariant submodules:
the span of each finite orbit is finite.
Thus the finite-coefficient comparison just
proved supplies the comparison for all
discrete modules required in
\cite[Theorem A]{R-BZ}.
The virtual cohomological dimension and
finite-cohomology hypotheses there are now
verified as well.

Apply that prime-order theorem to
$E\hookrightarrow\widehat E$.
It identifies the conjugacy classes of its
finite $2$-subgroups and gives the closure
of the centralizer of $\mathsf r$.
The group $\widehat J$ is torsion-free, by
its exact extension of the two torsion-free
free profinite groups in (F.15.11).
Hence every order-two subgroup of $\widehat E$
is a section.  By (F.15.16) there is one such
conjugacy class.  A conjugator in $\widehat E$
can be replaced by one in $\widehat J$:
write it as $h\mathsf r^\epsilon$ and use
that $\mathsf r^\epsilon$ centralizes
$\mathsf r$.
The centralizer comparison, intersected with
the open subgroup $\widehat J$, gives
\[
 C_{\widehat J}(\mathsf r)
                  =\overline{C_J(\mathsf r)}.
                                                               \tag{F.15.17}
\]
Combining this equality with the exact
full-topology statement following (F.15.13)
proves the first line of (F.15.5).
Its projection is precisely the graph of
the full filling epimorphism
$\widehat F_7\to\widehat F_5$ and its
composite with $\rho$.  This proves the
other two lines, including the named kernel.

\emph{Removing the fixed quotient coordinate.}
Take $\delta\in B$ with
$\delta R_*(\delta)=1$.
Then $\delta\mathsf r$ is an order-two
section of $\widehat E$.  The single
section class just proved gives $h\in\widehat J$
with
\[
                  \delta=h^{-1}R_*(h).
                                                               \tag{F.15.18}
\]
Its image $p(h)$ is fixed by (F.15.3).
By the second line of (F.15.5), choose
$c\in C_{\widehat J}(\mathsf r)$ with
$p(c)=p(h)$.  Then $b=c^{-1}h$ lies in
$B$ and direct multiplication yields
$b^{-1}R_*(b)=\delta$.
This proves surjectivity in (F.15.6).
Two such elements $b,b'$ have the same
coboundary exactly when
$R_*(b'b^{-1})=b'b^{-1}$, or equivalently
$b'\in C_B(\mathsf r)b$.  This proves the
asserted fibres.  The map is continuous,
its source is compact, and the cocycle
set is a closed subspace of the profinite
group $B$.  The bijection is therefore
a homeomorphism.

The subgroup $B$ is closed and normal in
$P_8$, and $R_*$ preserves it by F.14.13.
For any fixed element of $P_8/B$, choose a
lift $z\in P_8$.  Its defect
$z^{-1}R_*(z)$ is an $R_*$-cocycle in $B$.
Choose $b$ with the same coboundary by
(F.15.6).  Then
$R_*(zb^{-1})=zb^{-1}$, giving a
centralizing lift with the same quotient
image.  Its kernel is exactly
$C_B(\mathsf r)$.
This proves every map and the surjectivity
in (F.15.7).

\emph{The original source and its remaining membership equation.}
F.14 places $d(u)$ in the cocycle set
of (F.15.6), so (F.15.8) is nonempty
and has the specified left coset form.
For any one of its elements,
\[
 R_*(\widehat\phi(u)b^{-1})
  =\widehat\phi(u)d(u)R_*(b)^{-1}
  =\widehat\phi(u)b^{-1}.
                                                               \tag{F.15.19}
\]
Any two centralizing lifts of the same
$B$-coset differ on the right by an
element of $C_B(\mathsf r)$, proving
(F.15.9).
Finally $B\subset K_0$, and F.14 proves
that the composite of $\widehat\phi$
with forgetting $v_0$ is injective.
Thus $\widehat\phi(\widehat{\pi_1(Y)})\cap B=1$.
If two source elements have images in the
same $B$-coset, the same injectivity
makes them equal.
The only possible source representative
of (F.15.9) is consequently
$\widehat\phi(u)$ itself; it centralizes
$\mathsf r$ precisely for
$u\in\mathcal E$, by F.12.5.
No choice of the correcting element
$b$ has been asserted to remain in
the original source image.
\end{proof}

\begin{claim}{F.16}{Proved}{A full procyclic ambiguity invisible to every proper forgetting}
Use $P_8$, the actual involution $\mathsf r$, and the
source image
\[
       H=\widehat\phi\bigl(\widehat{\pi_1(Y,s_0)}\bigr)
                         \ \leq P_8
\]
of F.12--15.  There is an explicitly constructed discrete
element $w\in P\Gamma_{0,8}$ such that
\[
\begin{split}
 Z_w:=\overline{\langle w\rangle}&\cong\widehat\Z,\\
 Z_w&\subset C_B(\mathsf r),\\
 q_I(Z_w)&=1
       \quad\text{for every proper labelled forgetting }q_I,\\
 Z_w\cap H&=1 .
\end{split}                                                     \tag{F.16.1}
\]
Here $q_I$ retains any proper subset $I$ of the eight
labels, with the usual induced homomorphism on full
pure mapping class completions.  In particular all
eight individual one-point forgettings kill $Z_w$.

The exact image-membership test ruled out by this
construction is the following pullback subgroup:
\[
 \mathcal T=
       \bigcap_{I\ {\rm proper}}
                     q_I^{-1}\bigl(q_I(H)\bigr).
                                                               \tag{F.16.2}
\]
Both inclusions
\[
 H\ \subsetneqq\ \mathcal T,\qquad
 H\cap C_{P_8}(\mathsf r)\
       \subsetneqq\
          \mathcal T\cap C_{P_8}(\mathsf r)
                                                               \tag{F.16.3}
\]
are strict, witnessed by every nonidentity element of
$Z_w$.  The same strictness holds when the intersection
uses only the eight individual forgettings, or any
subfamily of proper forgettings.
For each $h\in H$, the elements $hz$ with $z\in Z_w$
have exactly the same images as $h$ under all these
maps, and $h$ is their only member in $H$.
Within each nonempty correction coset of F.15.9,
right multiplication by $Z_w$ likewise preserves
centralization and every proper forgetting image.
The original source-membership equation and the
full residual centralizer remain unresolved.
\end{claim}
\passage{Section 30, folios 194-196, scans 383-387; Section 31, folios 198-199 verso, scans 391-394; Part F(b), exact residual centralizer and the common-involution source-membership problem}


\begin{proof}
\emph{The actual six-point quotient and its labelled point-pushing group.}
Keep the two fixed points $a_+,a_-$ from F.15,
with $a_-=-a_+$ and $a_+^2=c_0$.
The specified quotient of $r(z)=c_0/z$ is
\[
    q(z)=\left(\frac{z-a_+}{z+a_+}\right)^2.
                                                               \tag{F.16.4}
\]
Indeed the coordinate $(z-a_+)/(z+a_+)$ changes
sign under $r$.  Label the two quotient branch
points $A=q(a_+)=0$, $D=q(a_-)=\infty$.
The four ordinary quotient labels and their
literal fibres on the eight marked points are
\[
\begin{array}{c|c|c}
 i&o_i&\text{marked fibre}\\ \hline
 0&1&\{0,\infty\}\\
 1&q(1)&\{v_0,v_2\}\\
 2&q(v_1(s_0))&\{v_1,v_3\}\\
 3&q(v_4(s_0))&\{v_4,v_5\}.
\end{array}                                                     \tag{F.16.5}
\]
These six quotient points are distinct: the four
ordinary values are the four different free
$r$-orbits, and none is a branch orbit.
This is a quotient of the actual configuration,
with the same choice of $a_+$ as F.15.

In the pure mapping class group of
$(\mathbf P^1;A,D,o_0,o_1,o_2,o_3)$,
forget the branch label $A$.  The actual Birman
kernel and point-pushing map are
\[
 \begin{split}
 L&=\pi_1\bigl(\mathbf P^1\setminus
                    \{D,o_0,o_1,o_2,o_3\},A\bigr)\cong F_4,\\
 \operatorname{push}_A:L&\hookrightarrow P\Gamma_{0,6}.
 \end{split}                                                    \tag{F.16.6}
\]
Choose a fixed geometric system of four disjoint
whiskers from $A$ to small circles about the four
$o_i$, with their interiors avoiding all five
other labels.  Orient the circles positively.
Their based loops $m_0,m_1,m_2,m_3$ are the
specified free basis of $L$; the omitted
$D$-meridian is the inverse of their ordered
product.  The based point-pushing map is induced
by moving $A$ along these loops and fixing
the other five labels.  Its inverse on the
Birman kernel is the track of $A$ in an isotopy
after that label has been forgotten.  The
aspherical Birman fibration in F.14 makes this
inverse independent of the isotopy, and gives
both the displayed injection and the full
induced profinite topology.

Let $\gamma_i$ be the boundary of a small regular
neighbourhood of the $i$th whisker and its two
ends.  It encloses exactly $A,o_i$.
Write
$T_i=\operatorname{push}_A(m_i)$.
This is the one-turn Dehn twist about $\gamma_i$,
with its sign fixed by this point-pushing
convention.  To see the representative directly,
move $A$ once round the small circle at $o_i$
and back along its whisker; outside the
regular neighbourhood the endpoint isotopy is
the identity, and on its boundary annulus it
is precisely that one-turn twist.

\emph{The pure lift, including its annulus degree.}
The double-cover monodromy is nontrivial around
each of $A,D$, and trivial around every $o_i$.
Therefore the inverse image of $\gamma_i$ is one
connected annulus covering its downstairs
annulus with degree two.  The disk it bounds
upstairs is the double cover of the disk
enclosing $A,o_i$, branched at its one point
$A$.  Its Euler characteristic is $2-1=1$,
and its boundary is connected, so it is a disk.
It contains exactly the two marked points in
row $i$ of (F.16.5), together with the unmarked
fixed point $a_+$, and no other marked point.

Use angular coordinates in turns on the two
annuli, so the covering map is
$\theta=2\vartheta$ modulo one.
Choose the radial coordinate $t\in[0,1]$
according to the preceding twist convention.
The square of the downstairs twist and its
pure lift are exactly
\[
 \begin{split}
  T_i^2:(\theta,t)&\longmapsto(\theta+2t,t),\\
  b_i:(\vartheta,t)&\longmapsto(\vartheta+t,t).
 \end{split}                                                    \tag{F.16.7}
\]
Both boundary circles are fixed pointwise by
the second map.  It extends by the identity
off the annulus.  The deck action there is
$\vartheta\mapsto\vartheta+\tfrac12$, which
commutes with (F.16.7).
Thus $b_i$ is a pure mapping class that
centralizes the actual $r$ and descends
to $T_i^2$.
It is the unique pure lift: the only other
lift differs by $r$, whose permutation is the
four nontrivial pairs (F.16.5).

The class $b_i$ is the one-turn twist about
the boundary of the displayed two-puncture
disk.  Forgetting either of those two points
leaves a disk with one marked point.
Its twist is isotopic to the identity
relative to that point and all exterior
labels: extend the untwisting isotopy
across its disk by rotations about that
remaining point.  Consequently
\[
 b_i\in C_{P\Gamma_{0,8}}(\mathsf r)
        \ \cap\
        \ker q_{p_i}\ \cap\ \ker q_{p'_i},
 \qquad
       \{p_i,p'_i\}=q^{-1}(o_i)\cap
                     \{0,\infty,v_0,\ldots,v_5\}.
                                                               \tag{F.16.8}
\]
All kernels here refer to the actual named
point-forgetting maps.

\emph{The nontrivial word and all eight forgettings.}
Take the commutator convention
$[g,h]=ghg^{-1}h^{-1}$ and define
\[
 \begin{split}
 w&=[[[b_0,b_1],b_2],b_3],\\
 \omega&=[[[m_0^2,m_1^2],m_2^2],m_3^2]\in L .
 \end{split}                                                    \tag{F.16.9}
\]
Every $b_i$ centralizes $\mathsf r$, so $w$
does.  Forgetting any one of the eight
points kills the corresponding $b_i$
by (F.16.8); substituting the identity
for any input of the displayed nested
commutator gives the identity.
This proves all eight individual
vanishings.  Any proper forgetting factors
through one of them, which proves all
the asserted vanishings after completion
as well.
Descent through the actual quotient gives
\[
             \operatorname{desc}(w)
                       =\operatorname{push}_A(\omega).
                                                               \tag{F.16.10}
\]
In particular this is an identity in the
specified Birman kernel, with its based
inverse in (F.16.6).

Here is an integral detector for $\omega$.
Let $U_5(\widehat\Z)$ be the group of
upper unitriangular five-by-five matrices,
and let $E_{ij}$ denote its usual matrix
units.  The free basis map has a unique
continuous extension
\[
 \widehat L\longrightarrow U_5(\widehat\Z),
       \qquad m_i\longmapsto I+E_{i+1,i+2}
                              \quad(0\leq i\leq3).
                                                               \tag{F.16.11}
\]
The four squares go to $I+2E_{i+1,i+2}$.
Successive multiplication in the actual
unitriangular group gives
\[
 \begin{split}
 [I+2E_{12},I+2E_{23}]&=I+4E_{13},\\
 [I+4E_{13},I+2E_{34}]&=I+8E_{14},\\
 [I+8E_{14},I+2E_{45}]&=I+16E_{15}.
 \end{split}                                                    \tag{F.16.12}
\]
Thus $\omega$ is nontrivial already
discretely.  For every $z\in\widehat\Z$,
its full profinite power has image
\[
                         I+16zE_{15}.                          \tag{F.16.13}
\]
Multiplication by sixteen is injective
on $\widehat\Z$; this detects every
nonzero primary coordinate, including
the two-primary coordinate.

\emph{The precise completion maps for this detector.}
Let $N_r$ be the discrete normalizer of
the specified $\langle r\rangle$ in the
eight-puncture mapping class group.
Its decorated quotient sequence F.7.5
is
\[
       1\longrightarrow\langle r\rangle
         \longrightarrow N_r
         \xrightarrow{\ \operatorname{desc}\ }Q
         \longrightarrow1 ,
                                                               \tag{F.16.14}
\]
where the two branch labels and four
ordinary labels of (F.16.5) are retained.
Every pure six-point mapping class
preserves the double-cover monodromy:
it fixes the peripheral homology classes,
and that monodromy is the specified
$C_2$-valued homology character.
Thus $P\Gamma_{0,6}$ is a finite-index
subgroup of this $Q$.
The Birman kernel $L$ in (F.16.6) has
its full induced topology in
$P\Gamma_{0,6}$ by F.14, and hence in
$Q$.  In particular the actual completed
maps give embeddings
\[
           \widehat L\hookrightarrow
             \widehat{P\Gamma}_{0,6}
                    \hookrightarrow\widehat Q .
                                                               \tag{F.16.15}
\]

Set $C_r=C_{P\Gamma_{0,8}}(\mathsf r)$.
It is a finite-index subgroup of $N_r$.
Its intersection with the finite deck
kernel is trivial, since that nonidentity
element permutes the marked points.
Consequently descent identifies $C_r$
with its finite-index image $Q_0$ in $Q$.
The inverse on $Q_0$ is the unique
lift fixing all eight labels individually.
The prime-order normalizer and full-topology
comparison F.7.4--5 make the completion
of this isomorphism the actual continuous
descent on the closure of $C_r$ in $P_8$:
\[
       \overline{C_r}\cong\widehat{C_r}
             \xrightarrow[\ \operatorname{desc}\ ]{\ \sim\ }
                           \widehat{Q_0}
               \hookrightarrow\widehat Q .
                                                               \tag{F.16.16}
\]
In particular (F.16.10) persists under
these maps, and its right side lies in
the embedded $\widehat L$ of (F.16.15).
If $w^z=1$ in $P_8$, its descent is
$\omega^z=1$ in $\widehat Q$.
The embedding (F.16.15) then puts this
equality in $\widehat L$, and
(F.16.13) gives $16z=0$, hence $z=0$.
The usual continuous power map
$\widehat\Z\to\overline{\langle w\rangle}$
is therefore injective as well as onto.
This proves the full procyclic
identification in (F.16.1), with the
inverse-source and topology comparisons
needed to apply the matrix detector.

\emph{The exact obstruction to the proposed membership test.}
The eight vanishing maps include
$q_{v_0}$ and $q_{v_2}$, so
$Z_w\subset B$.  Its elements centralize
$\mathsf r$, proving $Z_w\subset C_B(\mathsf r)$.
F.14 states that $q_{v_0}|_H$ is
injective.  Since $q_{v_0}(Z_w)=1$,
we obtain $Z_w\cap H=1$.

For every proper $I$ the identity lies in
$q_I(H)$, and $q_I(Z_w)=1$.
Thus $Z_w\subset\mathcal T$, while
$H\subset\mathcal T$ by its definition.
Every nonidentity element of $Z_w$
witnesses the first strict inclusion.
It also centralizes $\mathsf r$ and
therefore witnesses the second.
The same witnesses work for every
subfamily of the displayed maps.
For $h\in H$ and $z\in Z_w$,
$q_I(hz)=q_I(h)$ for all $I$, whereas
$hz\in H$ forces $z\in H\cap Z_w=1$.
Finally $Z_w\subset C_B(\mathsf r)$
places right multiplication by $Z_w$
inside each actual centralizing
correction coset (F.15.9), while all
proper forgetting maps remain unchanged.
These statements prove precisely
(F.16.2)--(F.16.3); they supply no
nonidentity element of the original
source equalizer.

The executed certificate
\path{code/part_a/common_reflection_forgetting_ambiguity_f16.py}
records the four actual orbit pairs,
checks the degree-two annulus formulas,
all eight individual and all proper
word-forgetting consequences, the
integer matrix value (F.16.12), and
its power formulas in twenty-one
finite primary quotients.
The geometric lifting and full
completion maps are proved above.
\end{proof}

\begin{claim}{F.17}{Proved}{The actual branch drilling and its nonfactorization}
Retain $Y,s_0,\zeta$ and the labelled map $\phi$ of F.12.
Write
\[
 v_u(s)=\left(\frac{u-s}{u+s}\right)^2\quad(u\in\mu_5),
 \qquad a(s)=\frac{\zeta-s}{\zeta+s},
 \qquad v_\zeta=a^2.
                                                               \tag{F.17.1}
\]
Thus $v_{\zeta^j}=v_{j+1}$ in the earlier notation.
For $u\in\mu_5\setminus\{\zeta\}$ put
\[
 A_u=u(2\zeta-u),\qquad B_u=\frac{\zeta u^2}{2u-\zeta},
 \qquad
 \Delta=\{s:s^2=A_u\text{ or }s^2=B_u
                    \text{ for one such }u\}.
                                                               \tag{F.17.2}
\]
The set $\Delta$ consists of sixteen distinct points of $Y$.
Set
\[
 Y^\dagger=Y\setminus\Delta,\qquad
 Y_0=\mathbf P^1\setminus\{0,\infty,\zeta,-\zeta\}.
                                                               \tag{F.17.3}
\]
There is an actual ten-label configuration map
\[
 \phi^\dagger(s)=
       (0,\infty,1,(v_u(s))_{u\in\mu_5},a(s),-a(s))
             \quad(s\in Y^\dagger).
                                                               \tag{F.17.4}
\]
The two added labels are ordered by these last two formulas.
At the basepoint they are the two fixed points of the actual
involution in F.12--15.  This is a specified ordering; comparison
with the other ordering is the literal transposition of these
two new labels.

The source fillings and their completed homomorphisms are
\[
 \widehat{\pi_1(Y^\dagger,s_0)}=\widehat F_{27}
       \xrightarrow{\ q\ }\widehat{\pi_1(Y,s_0)}=\widehat F_{11}
       \xrightarrow{\ q_0\ }\widehat{\pi_1(Y_0,s_0)}=\widehat F_3 .
                                                               \tag{F.17.5}
\]
The first fills precisely $\Delta$, and the second fills
the eight points $\pm u$ with $u\ne\zeta$.
Both maps are onto.

Each based meridian $m_d$ at an added point $d\in\Delta$
is killed by $q$, but
$\widehat{\phi^\dagger}_*(m_d)$ generates a full closed
copy of $\widehat\Z$ in the ten-point pure completion.
It is the actual primitive two-point twist at the corresponding
collision in (F.17.2), with the base path induced by the
specified source whisker.
Consequently the particular map
$\widehat{\phi^\dagger}_*$ does not factor through $q$.

Forgetting the four labels $v_u$ with $u\ne\zeta$ gives
the exact map
\[
 \chi:Y_0\longrightarrow\mathcal M_{0,6},\qquad
                 \chi(s)=(0,\infty,1,a(s)^2,a(s),-a(s)).
                                                               \tag{F.17.6}
\]
Its induced full profinite map is injective and its discrete
image has the full induced topology.
The retained involution is $V\mapsto a(s)^2/V$.
Its actual branched quotient has the four labelled values
$0,\infty,1,t$, where $t=s^2/\zeta^2$.
The two maps $Y_0\to\mathcal M_{0,4}$ involved here are
specified in the proof: ordinary forgetting has parameter
$a^2$, whereas the involution quotient has parameter $t$.
Each is a degree-two unramified covering.
The first proves the full topology assertion directly.

These statements give an exact additional labelled
correspondence for the source-membership problem.
They make no injectivity assertion for the whole
ten-point map (F.17.4), and do not determine the
intersection in F.12.5.
\end{claim}
\passage{Section 30, folios 194-196, scans 383-387; Section 31, folios 198-199 verso, scans 391-394; Part F(b), the exact residual centralizer and its source-membership problem}


\begin{proof}
\emph{All collision points of the two new labels.}
The rational labels $a,-a$ are the fixed points of
$V\mapsto v_\zeta(s)/V$.
They are distinct and avoid $0,\infty,1,v_\zeta$
on $Y$: the excluded cases are $s=0,\infty,\zeta,-\zeta$.
For $u\ne\zeta$, the two exact differences have the
factorizations
\[
\begin{split}
 v_u-a
   &=\frac{2s(s^2-A_u)}{(u+s)^2(\zeta+s)},\\
 v_u+a
   &=\frac{2(\zeta-2u)(s^2-B_u)}
                {(u+s)^2(\zeta+s)}.
\end{split}                                                     \tag{F.17.7}
\]
This proves the stated collision equations, without
discarding a denominator or a source puncture.

Here is also a direct separation check.
In the specified complex embedding $|u|=|\zeta|=1$,
\[
 |A_u|=|2\zeta-u|>1,\qquad
                  B_u=\frac1{\overline{A_u}},\qquad |B_u|<1.
                                                               \tag{F.17.8}
\]
Equality in the first inequality would force $u=\zeta$.
For two distinct fifth roots $u,w$,
$A_u=A_w$ would imply $u+w=2\zeta$; the equality case
in the triangle inequality again forces $u=w=\zeta$.
Thus the four $A_u$ are distinct.  Their reciprocal
conjugates are distinct, and the two sets cannot
intersect by (F.17.8).
All eight values are nonzero.
Their square roots are therefore sixteen distinct
points.  None has absolute value one or equals
$0,\infty$, so none is an old source puncture.

The denominators in (F.17.7) are nonzero at these
sixteen points, since the excluded zeros have
absolute value one.
Every zero displayed there is simple:
the $s^2-A_u$ and $s^2-B_u$ factors have nonzero
derivatives at their roots, and the remaining
numerator factors are nonzero.
There is only one new collision at each such point.
The original eight labels are distinct on $Y$ by F.12.
This proves that (F.17.4) is defined precisely on
$Y^\dagger$ among the points of $Y$.

The basepoint $s_0$ belongs to $Y^\dagger$.
Indeed its actual involution permutes all eight
original labels in four nontrivial pairs, by F.15.10.
Neither of its two fixed points is one of those labels.
The formula in (F.17.1) squares to $v_\zeta(s_0)=v_2(s_0)$,
so it selects exactly the required ordered new pair.

\emph{The actual filling maps and local twists.}
Choose a fixed system of based geometric whiskers
from $s_0$ to small positively oriented circles at
the finite punctures of $Y^\dagger$.
The circles are disjoint, and the whisker interiors
avoid all punctures.
The corresponding meridians are a free basis, with
the infinity meridian given by their inverse ordered
product.
Filling the sixteen punctures $\Delta$ sends precisely
those meridians to the identity and retains the
meridians of $Y$ with their induced whiskers.
Filling the eight additional old punctures gives $Y_0$.
The same maps after completion are exactly (F.17.5):
the finite-target universal property identifies
each filling kernel with the closed normal closure
of its named meridians.  In particular both completed
maps are onto, and their composite fills all twenty-four
additional points of $Y^\dagger$ relative to $Y_0$.

Let $d^2=A_u$.
In a small disk about $d$, the difference $v_u-a$
has winding one by (F.17.7); every other label
difference stays nonzero there.
The local monodromy in the ordered configuration
is therefore the one-turn pure twist around the
two labels $v_u,a$.
Concretely, the center and difference of these two
coordinates give local configuration coordinates;
the center varies in a contractible disk, and the
nonzero difference runs once around zero.
Extending this representative by the identity outside
a small disk containing the two labels gives that
twist.  Its based conjugating path is the image
under $\phi^\dagger$ of the already chosen whisker.
For $d^2=B_u$ the identical argument uses $v_u+a$
and the pair $v_u,-a$.
Thus the asserted monodromy includes the actual
pair of labels, exponent, and base transport.

For a full profinite detection, normalize the original
three labels $0,\infty,1$ in the ten-point moduli space.
Write $Z_u,Z_+,Z_-$ for the coordinates of $v_u,a,-a$.
The globally nonzero functions $Z_u-Z_+$ and $Z_u-Z_-$
define their integer winding homomorphisms on its
pure mapping class group.
Their pullbacks under $\phi^\dagger$ are precisely
the two functions in (F.17.7).
At the corresponding added meridian the relevant
integer winding is one.  Completing that character
gives the literal equality
\[
 \widehat{\operatorname{wind}}_{u,\pm}
       \bigl(\widehat{\phi^\dagger}_*(m_d^{\,z})\bigr)
                         =z\quad(z\in\widehat\Z).
                                                               \tag{F.17.9}
\]
The continuous power map onto the closure of that
twist is injective by (F.17.9), as well as surjective.
This proves the full $\widehat\Z$ assertion at every
primary coordinate.  It does not require injectivity
of $\widehat{\phi^\dagger}_*$ on the entire source.

Since $q(m_d)=1$ but its ten-point image is nonidentity,
there cannot be a homomorphism $D$ with
$Dq=\widehat{\phi^\dagger}_*$.
This proves exactly the stated nonfactorization.

\emph{Both diagrams with their literal retained labels.}
Let $\operatorname{forget}_{+,-}$ forget the two added
labels, and let $\operatorname{forget}_{\mathrm{four}}$
forget precisely the four $v_u$ with $u\ne\zeta$.
The based maps on full completions satisfy
\[
\begin{split}
 \operatorname{forget}_{+,-}\,
                    \widehat{\phi^\dagger}_*
       &=\widehat\phi_*\,q,\\
 \operatorname{forget}_{\mathrm{four}}\,
                    \widehat{\phi^\dagger}_*
       &=\widehat\chi_*\,q_0q .
\end{split}                                                     \tag{F.17.10}
\]
The first target has the original eight labels.
The second target has exactly
$0,\infty,1,v_\zeta,a,-a$.
Both identities follow from the actual configuration
coordinates, with the same basepoint and whiskers.
Their common further forgetting retains
$0,\infty,1,v_\zeta$ and has parameter
\[
                          u_0(s)=a(s)^2.                        \tag{F.17.11}
\]
This common four-point target is the ordinary
forgetting target, not the branched quotient target
described next.

\emph{The six-point double cover and its full topology.}
The coordinate $a=(\zeta-s)/(\zeta+s)$ is a M\"obius
isomorphism from $Y_0$ onto
$\mathbf P^1\setminus\{0,\infty,1,-1\}$.
Squaring $a$ is the degree-two unramified covering
of $\mathbf P^1\setminus\{0,\infty,1\}$.
It is exactly the ordinary forgetting (F.17.11).
The induced fundamental-group image is the index-two
kernel of the character that is one on a meridian
at zero and zero on a meridian at one, using the
$u_0$ coordinate and its induced basepoint.

A finite-index subgroup of a discrete group has
its full induced profinite topology: the normal core
of a finite-index subgroup of it has finite index
in the whole group.
Consequently the completed map on this covering is
an embedding $\widehat F_3\hookrightarrow\widehat F_2$.
Since it is the composite of $\widehat\chi_*$
with an actual target forgetting homomorphism,
$\widehat\chi_*$ is injective.
The same composite proves the full induced topology
of the discrete image of $\chi_*$.
This proof asserts full topology for this specified
six-point map and uses no inference about the kernel
of the ten-point map.

The involution $V\mapsto a^2/V$ fixes the two labels
$a,-a$ individually and exchanges the other labels
as $(0\ \infty)(1\ a^2)$.
Its actual quotient is
\[
 Q_a(V)=\left(\frac{V-a}{V+a}\right)^2,
 \qquad
 \begin{array}{c|c}
 \text{labelled fibre}&\text{quotient value}\\ \hline
 a&0\\
 -a&\infty\\
 \{0,\infty\}&1\\
 \{1,a^2\}&t=s^2/\zeta^2 .
 \end{array}                                                    \tag{F.17.12}
\]
Direct substitution proves all four entries;
in particular
$Q_a(1)=((1-a)/(1+a))^2=s^2/\zeta^2$.
Thus the branched-quotient parameter map is the
other degree-two unramified covering
$Y_0\to\mathbf P^1\setminus\{0,\infty,1\}$.

The sheet change $s\mapsto-s$ sends $a$ to $a^{-1}$.
The M\"obius map $V\mapsto V/a^2$ identifies the
two six-point covers over the same $t$ and exchanges
exactly the labels $1,a^2$, while preserving the
two named branch labels and the labels $0,\infty$.
The deck involution exchanges both ordinary pairs.
Hence that sheet change cannot be corrected to a
pure labelling by the deck involution.
Conversely the choices of the two ordinary
sheet orderings, modulo the simultaneous deck
exchange, are exactly these two possibilities.
This proves that $Y_0$ is the actual index-two
pure-labelling cover of this four-point quotient
family, with all of its labels and the branch order
specified.

The executed certificate
\path{code/part_a/common_reflection_drilling_f17.py}
checks the eight exact factorizations over
$\mathbf Q(\zeta)$, all separations of the eight
squared collision values, the two quotient and
sheet-label formulas, and the coefficient-one
meridian detectors in finite primary reductions.
It passes 445 exact checks.
The topology and full completion comparisons are
the explicit covering and winding arguments above.
\end{proof}

\begin{claim}{F.18}{Proved}{The unique full branch lift and its exact source coset}
Keep the actual eight-point configuration $\mathcal P_8$
and involution $\mathsf r$ of F.12--17.  Put
$a_+=a(s_0)$ and $a_-=-a(s_0)$ with the literal ordering
in F.17.  Mark these two additional points and retain
the same homeomorphism as $\mathsf r^\dagger$.
Write
\[
\begin{split}
 P_{10}^{\dagger}
  &=\widehat{P\Gamma}(\mathbf P^1;\mathcal P_8,a_+,a_-),\\
 P_8&=\widehat{P\Gamma}(\mathbf P^1;\mathcal P_8),\\
 f:P_{10}^{\dagger}&\longrightarrow P_8,\qquad
 \mathcal J=\ker f,\\
 C^\dagger&=C_{P_{10}^{\dagger}}(\mathsf r^\dagger),
 \qquad C_8=C_{P_8}(\mathsf r).
\end{split}                                                     \tag{F.18.1}
\]
Here $f$ forgets exactly the two named added points.
The branch-permutation character
$\epsilon:C_8\twoheadrightarrow C_2$ is defined by
the actual decorated quotient of $\mathsf r$.
Then
\[
 f:C^\dagger\xrightarrow{\ \sim\ }C_8^+
       :=\ker\epsilon,\qquad
             C_{\mathcal J}(\mathsf r^\dagger)=1              \tag{F.18.2}
\]
are an isomorphism of profinite groups and its exact
kernel assertion.

For the action $R^\dagger$ on $\mathcal J$ by
$\mathsf r^\dagger$, the full continuous cocycles
\[
 Z^1(R^\dagger,\mathcal J)
     =\{\delta\in\mathcal J:\delta R^\dagger(\delta)=1\}
                                                               \tag{F.18.3}
\]
have precisely two $\mathcal J$-conjugacy classes.
Every order-two section has trivial centralizer in
$\mathcal J$.  In particular the map
\[
 b\longmapsto b^{-1}R^\dagger(b)
                                                               \tag{F.18.4}
\]
is a homeomorphism from $\mathcal J$ onto the
coboundary class.
More generally, for $k\in P_{10}^{\dagger}$ with
$f(k)\in C_8$, the class of
$k^{-1}R^\dagger(k)$ is the class labelled
$\epsilon(f(k))$.

The original residual group $\mathcal E$ of F.12
has a unique full centralizing branch lift
\[
 \ell:\mathcal E\longrightarrow C^\dagger,\qquad
                   f\ell=\widehat\phi|_{\mathcal E}.
                                                               \tag{F.18.5}
\]
It is continuous and injective.
This includes the proof that its branch character
is zero; no ordered drilling is assumed in advance.

Here is the exact remaining comparison to the
rational drilling of F.17.  Put
\[
\begin{split}
 L^\dagger&=\widehat{\pi_1(Y^\dagger,s_0)},\qquad
 L=\widehat{\pi_1(Y,s_0)},\\
 q:L^\dagger&\twoheadrightarrow L,\qquad
 \Phi^\dagger=\widehat{\phi^\dagger}_*,\\
 K^+&=\Phi^\dagger(\ker q)
          =\Phi^\dagger(L^\dagger)\cap\mathcal J .
\end{split}                                                     \tag{F.18.6}
\]
For $u\in\mathcal E$ and any $w\in q^{-1}(u)$,
define the actual element
\[
             b(w)=\ell(u)^{-1}\Phi^\dagger(w)\in\mathcal J .
                                                               \tag{F.18.7}
\]
It is the unique solution of
\[
 b(w)^{-1}R^\dagger(b(w))
       =\Phi^\dagger(w)^{-1}R^\dagger(\Phi^\dagger(w)).
                                                               \tag{F.18.8}
\]
The right $K^+$-orbit $b(w)K^+$ is independent of $w$
and defines a continuous map
\[
 \mathfrak o:\mathcal E\longrightarrow\mathcal J/K^+ .
                                                               \tag{F.18.9}
\]
No group structure on this coset space is assumed.
Its distinguished fibre is exactly
\[
\begin{split}
 \mathfrak o^{-1}(K^+)
   &=\{u\in\mathcal E:
                 \ell(u)\in\Phi^\dagger(L^\dagger)\}\\
   &=q\bigl((\Phi^\dagger)^{-1}(C^\dagger)\bigr).
\end{split}                                                     \tag{F.18.10}
\]
These are equalities for the specified full maps.
Whether this fibre is all of $\mathcal E$, and
whether $\mathcal E$ is trivial, remain unresolved.
\end{claim}
\passage{Section 30, folios 194-196, scans 383-387; Section 31, folios 198-199 verso, scans 391-394; Part F(b), full residual centralizer and exact source-membership comparison}


\begin{proof}
\emph{The two quotient normalizers and their actual inverse.}
On the eight-point sphere the involution is the
M\"obius map $r(V)=a_+^2/V$, with four nontrivial
marked pairs and unmarked fixed points $a_+,a_-$.
After marking them, the same map fixes each new
label and retains the four old pairs.
The decorated quotient before marking has four
ordinary orbit labels and two branch points;
after marking it has the same four orbit labels
and the two branch points ordered individually.

Use the discrete covering-normalizer sequence F.7.5
on both sides of the forgetting map.
Its kernel is the same deck group
$\langle\mathsf r\rangle=C_2$.
On either pure centralizer this kernel is trivial,
because its nonidentity element permutes the eight
old labels.  Descent therefore identifies the
discrete pure centralizer before drilling with a
finite-index subgroup $Q_0$ of the decorated quotient
mapping class group.  Its image after ordered
drilling is exactly the subgroup $Q_0^+$ preserving
the two branch labels individually.
Indeed, a representative in $Q_0^+$ has its unique
lift fixing all old labels, and that lift fixes the
two ordered ramification points as well.
Conversely every pure lift after drilling has those
two properties.  The inverse comparison is this
same unique pure lift.

The branch permutation on $Q_0$ is onto $C_2$.
To see this with the old labels retained, interchange
the two branch points along an arc avoiding the
four ordinary quotient labels.
Lift that quotient homeomorphism.
If its endpoint exchanges the two names in some
old pair, push that ordinary quotient point once
around one branch point and lift the isotopy from
the identity.  This changes exactly that pair
permutation, fixes the branch points individually,
and leaves all other pair names fixed.
Correcting each mismatch gives a pure upstairs
class still exchanging the two unmarked fixed
points.  Thus $[Q_0:Q_0^+]=2$.

The prime-order normalizer and full-topology
comparison F.7.4--5 applies to these two actual
involutions.  In particular both actual pure
centralizers are their discrete completions.
Their branch-permutation homomorphism is continuous
in the induced topology: the full topology admits
this finite $C_2$ quotient.  Its kernel is open;
the completion of $Q_0^+$ is exactly that kernel.
Completing the discrete inverse maps above gives
the first isomorphism in (F.18.2), for the
actual forgetting map $f$.  Intersecting its
kernel with $\mathcal J$ gives the second assertion.
No intersection of two unrelated closures is used.

\emph{The actual two-point kernel and its section classes.}
Write $S_8=\mathbf P^1\setminus\mathcal P_8$ and
\[
 J_{\rm d}=\pi_1(\operatorname{Conf}_2(S_8),(a_+,a_-)).
                                                               \tag{F.18.11}
\]
The action on this ordered configuration fibre is
$(z_+,z_-)\mapsto(rz_+,rz_-)$; it does not interchange
the two new labels.
The two successive Birman sequences, with the
completed comparison of F.14--15, identify
$\widehat{J_{\rm d}}$ with the actual subgroup
$\mathcal J$ in (F.18.1).
For completeness, projection to the first point
gives the split free-by-free extension
\[
              1\longrightarrow F_8
                \longrightarrow J_{\rm d}
                \longrightarrow F_7\longrightarrow1 .
                                                               \tag{F.18.12}
\]
A section is obtained by putting the second point
far out in a plane coordinate with one original
puncture at infinity, and then adjusting the
basepoint by a path in the fibre.
Characteristic finite-index subgroups of $F_8$
and the finite action of the free quotient on
their finite quotients give the full topology on
that kernel.  Completion is exact.
Mapping this sequence into the two completed
Birman sequences proves the asserted embedding
and equality with the forgetting kernel, just
as in F.15.11, with these actual ranks and labels.

Consider the discrete group
\[
            E_{\rm d}=J_{\rm d}\rtimes
                             \langle\mathsf r^\dagger\rangle .
                                                               \tag{F.18.13}
\]
It is the actual subgroup of the ten-point mapping
class group generated by the two-point kernel and
the indicated involution; the nontrivial coset is
detected by the old-label permutation.
Each order-two section has exactly those four
old pairs and fixes the two added labels.
Nielsen realization in the punctured form used
in F.15 realizes it as a sphere involution.
Its two fixed points are precisely the two
added marked points.

An ordered quotient homeomorphism matching those
two branch points and the four named ordinary
orbits lifts to a conjugator to
$\mathsf r^\dagger$.  Correct any old pair-name
mismatches by the quotient point pushes just
described.  This proves that every such section
is conjugate to $\mathsf r^\dagger$ by a pure
ten-point class $g$.
The condition that it lies in (F.18.13) is exactly
$f(g)\in C_{8,{\rm d}}$, the discrete centralizer.
The set of its $J_{\rm d}$-conjugacy classes is
therefore the actual double-coset set
\[
 J_{\rm d}\backslash
       \{g\in P\Gamma_{0,10}^{\dagger}:f(g)\in C_{8,{\rm d}}\}
       /C^\dagger_{\rm d}
       \ \xrightarrow{\ \sim\ }\
       C_{8,{\rm d}}/C_{8,{\rm d}}^+
       \ \cong C_2 .                                         \tag{F.18.14}
\]
The map sends $g$ to the branch permutation of $f(g)$.
It is onto because forgetting the two points is
onto on pure mapping class groups.
If two images have the same branch permutation,
their quotient downstairs lifts through
$C^\dagger_{\rm d}$ by the discrete isomorphism
already proved; removing that lift leaves an
element of $J_{\rm d}$.  This proves injectivity.
All their centralizers in $J_{\rm d}$ are trivial:
for the standard section this is the discrete
kernel of (F.18.2), and pure conjugation preserves
the normal subgroup $J_{\rm d}$.

Here are the full completion hypotheses for
transporting this calculation.
The group in (F.18.12) is good, by the two-row
Lyndon--Hochschild--Serre comparison for its free
kernel and free quotient, whose finite-coefficient
cohomology groups are finite.
It has a finite two-dimensional classifying model:
use finite graphs for the free kernel and quotient,
and mapping cylinders of graph homotopy equivalences
representing the finitely many free-group automorphisms.
The finite extension (F.18.13) is residually finite
as an actual subgroup of the residually finite
mapping class group.  Its completion is
$\mathcal J\rtimes\langle\mathsf r^\dagger\rangle$,
using the finite-index subgroup $J_{\rm d}$.
The finite-extension cohomology comparison gives
goodness and finite virtual cohomological type.
The same argument for all discrete two-primary
modules follows by taking the filtered union of
their finite invariant submodules; the
$\mathrm{FP}_\infty$ resolution commutes with
these unions.  These are precisely the
hypotheses used for (F.15.14).

Apply \cite[Theorem~A]{R-BZ} to this actual group
and the prime two.
It gives the two full conjugacy classes of
order-two subgroups and the closure of their
centralizers.  An $E_{\rm d}$-conjugacy between
sections is already a $J_{\rm d}$-conjugacy:
an outside conjugator can be multiplied by
the nonidentity element of the section it
conjugates to make its quotient label zero,
without changing that conjugation.
The same argument works after completion.
Thus (F.18.14) gives exactly two
$\mathcal J$-classes.
For each discrete section $T$, its centralizer
in $E_{\rm d}$ is $T$ itself, since its intersection
with $J_{\rm d}$ is trivial and the quotient is
$C_2$.  The completed centralizer equality therefore
makes the full fixed groups in $\mathcal J$
trivial as well.
The orbit map from $\mathcal J$ to either
cocycle class is consequently injective
and onto; compactness and the Hausdorff
target make it a homeomorphism.
This proves (F.18.3)--(F.18.4).
For a full $k$ as in the claim with
$\epsilon(f(k))=0$, let $c\in C^\dagger$
be the unique lift of $f(k)$ under (F.18.2).
Then $b=c^{-1}k\in\mathcal J$ and
$k^{-1}R^\dagger(k)=b^{-1}R^\dagger(b)$.
If instead $\epsilon(f(k))=1$, a coboundary
expression for this element would put
$kb^{-1}$ in $C^\dagger$, contradicting
(F.18.2) after forgetting.
Since there are exactly two full classes,
this proves the asserted class label for
every full $k$ directly.

\emph{The residual group fixes the two actual stabilizer labels.}
We verify this before defining the lift.
Let $\jmath:L\to\widehat P\Gamma_{0,12}$ be
the injective original covering map of F.12.
For $u\in\mathcal E$, both $\jmath(u)$ and
$\beta\jmath(u)\beta^{-1}$ lie in the pure
completed centralizer of $\tau$.
Their descents agree by the defining equation
of $\mathcal E$.  That descent is injective
on the pure centralizer: its only finite
kernel is $\langle\tau\rangle$, whose
nonidentity element has a nontrivial label
permutation.  Therefore
\[
                    \beta\jmath(u)\beta^{-1}=\jmath(u).
                                                               \tag{F.18.15}
\]
The element also centralizes $\tau$.
This establishes pointwise centralization
of $V_4=\langle\tau,\beta\rangle$; it is
stronger than preserving the two lifts
$\{\beta,\tau\beta\}$ as an unordered set.

The actual coordinate $W=(z-s_0)/(z+s_0)$
on the original sphere has the formulas
\[
 \tau(W)=-W,\qquad
 \beta(W)=\frac{k}{W},\qquad
 \tau\beta(W)=-\frac{k}{W},\qquad
 k=a(s_0)=\left(1-\frac{2\theta^2}{5}\right)(1+2i).
                                                               \tag{F.18.16}
\]
These follow by conjugating the original
M\"obius matrices of F.11--12; the exact
certificate below checks them in the
same coefficient field.
The quotient by $\tau$ is $V=W^2$.
The fixed points of $\beta$ consequently
give $V=k=a_+$, and those of $\tau\beta$
give $V=-k=a_-$.
Thus the literal F.17 branch ordering
is the ordering of these two distinct
stabilizer types.

A discrete class centralizing $V_4$
preserves each stabilizer subgroup
individually in the actual quotient
cover.  Its descent under $\tau$
therefore fixes these two labels
individually.
The full statement follows from the
direct prime-power comparison for the
actual two-group $V_4$: its completed
centralizer is the closure of its
discrete centralizer, with the full
topology of F.7.4--5.
The branch character is finite and
continuous on the target centralizer,
so the same zero character holds on
that entire closure.
Applying this to (F.18.15) proves
$\widehat\phi(\mathcal E)\subset C_8^+$.
The inverse in (F.18.2) now defines
the continuous injective map $\ell$
in (F.18.5), and proves its uniqueness.

\emph{The exact remaining source-image test.}
The square F.17.10 is the literal identity
\[
                       f\Phi^\dagger=\widehat\phi\,q .
                                                               \tag{F.18.17}
\]
Since $\widehat\phi$ is injective, an
element $\Phi^\dagger(w)$ lies in
$\mathcal J$ precisely when $q(w)=1$.
This proves both inclusions in the
equality defining $K^+$ in (F.18.6).
It is a closed subgroup, being the
compact image of the closed kernel
of $q$.

For $q(w)=u\in\mathcal E$, (F.18.17)
and (F.18.5) put $b(w)$ in $\mathcal J$.
Because $\ell(u)$ centralizes
$\mathsf r^\dagger$, direct multiplication
gives (F.18.8).
Equation (F.18.4) and its trivial fixed
kernel prove uniqueness of this solution.
Replacing $w$ by $wv$ with $v\in\ker q$
gives the exact formula
\[
                         b(wv)=b(w)\Phi^\dagger(v).
                                                               \tag{F.18.18}
\]
Thus $b(w)K^+$ depends only on $u$.
The map on $q^{-1}(\mathcal E)$ is
continuous; the compact surjection to
$\mathcal E$ is a quotient map.
It therefore descends continuously
to (F.18.9).

Finally,
$\ell(u)=\Phi^\dagger(w)b(w)^{-1}$.
The definition and equality in (F.18.6)
show that this element belongs to
$\Phi^\dagger(L^\dagger)$ precisely
when $b(w)\in K^+$.
In that case choose $v\in\ker q$
with $\Phi^\dagger(v)=b(w)^{-1}$;
then $\Phi^\dagger(wv)=\ell(u)$
centralizes and $q(wv)=u$.
Conversely, a source element with
centralizing ten-point image fills
to $\widehat\phi(u)$ with $u\in\mathcal E$,
and uniqueness in (F.18.2) identifies
its image with $\ell(u)$.
These are both directions of
(F.18.10).
In particular the canonical correcting
element is unique; its membership in
the specified closed subgroup $K^+$
is the remaining literal source test.

The executed certificate
\path{code/part_a/common_reflection_branch_lift_f18.py}
checks the original M\"obius conjugations,
the exact coefficient $k$, the four old
marked pairs, and the $V_4$ stabilizer
labels.  It passes 34 exact checks.
The full topological and completion
assertions are the arguments above,
not finite-quotient lifting assumptions.
\end{proof}

\begin{claim}{F.19}{Proved}{The two-track kernel and the exact remaining source subgroup}
Use the literal base configuration and ordered branch labels of
F.17--18.  Partition its ten labels as
\[
 A=\{0,\infty,1,v_\zeta\},\qquad
 U=\{v_u:u\in\mu _5\setminus\{\zeta\}\},\qquad
 Z=\{a_+,a_-\}.
                                                               \tag{F.19.1}
\]
Let \(P_4,P_8,P_6,P_{10}^{\dagger}\) be the full pure
completions on \(A,A\cup U,A\cup Z,A\cup U\cup Z\), respectively.
The actual forgetting maps form the square
\[
 \begin{array}{ccc}
 P_{10}^{\dagger}&\xrightarrow{\ g\ }&P_6\\
 {\scriptstyle f}\downarrow&&\downarrow{\scriptstyle h}\\
 P_8&\xrightarrow{\ f_0\ }&P_4 .
 \end{array}
 \qquad
 \begin{gathered}
 \mathcal Q=P_8\times_{P_4}P_6,\\
 \pi=(f,g):P_{10}^{\dagger}\twoheadrightarrow\mathcal Q,\\
 D^+=\ker\pi=\ker f\cap\ker g .
 \end{gathered}                                                 \tag{F.19.2}
\]
In particular the displayed map to this specified profinite
fibre product is onto.

For each of the eight pairs \((v_u,a_\pm)\), choose either one
of its two collision points in F.17 and retain its actual
source whisker.  Write \(\nu_{u,\pm}\) for the image of that
source meridian under \(\Phi^\dagger\).
Then
\[
 D^+=\overline{\langle\!\langle
               \nu_{u,\pm}:u\ne\zeta
                         \rangle\!\rangle}_{\,P_{10}^{\dagger}},
 \qquad \nu_{u,\pm}\in K^+ .
                                                               \tag{F.19.3}
\]
The normal closure is taken in the indicated full ambient
group.  This does not replace \(K^+\) by that normal closure.

All four label sets in (F.19.1) are invariant under the
specified involution, and its two new labels are fixed
individually.  Denote the induced actions by \(R\) on each
group where no ambiguity arises.  With
\(\mathcal J=\ker f\) and
\(\mathcal J_4=\ker(h:P_6\to P_4)\), one has
\[
 \begin{gathered}
 1\longrightarrow D^+\longrightarrow\mathcal J
       \xrightarrow{\ g\ }\mathcal J_4\longrightarrow1,\\
 C_{\mathcal J_4}(R)=1,\qquad
 H^1(C_2,D^+)=\{1\}.
 \end{gathered}                                                 \tag{F.19.4}
\]
Here the last expression denotes the continuous nonabelian
cohomology set for this actual action.
More explicitly, the map
\[
 D^+\longrightarrow
     \{\delta\in D^+:\delta R(\delta)=1\},\qquad
                b\longmapsto b^{-1}R(b)
                                                               \tag{F.19.5}
\]
is a homeomorphism.
The centralizer map
\[
          \pi:C^\dagger\xrightarrow{\ \sim\ }\mathcal Q^R
                                                               \tag{F.19.6}
\]
is an isomorphism of profinite groups.

For the original residual group \(\mathcal E\), its unique
branch lift \(\ell\) and the source maps of F.17--18,
\[
 g\ell(u)=\widehat\chi_*(q_0u),\qquad
 b(w)=\ell(u)^{-1}\Phi^\dagger(w)\in D^+
               \quad(qw=u,\ u\in\mathcal E).
                                                               \tag{F.19.7}
\]
Thus \(K^+\subset D^+\), and the exact obstruction map of F.18
has the more precise target
\[
 \mathfrak o_D:\mathcal E\longrightarrow D^+/K^+,\qquad
                 u\longmapsto b(w)K^+ .
                                                               \tag{F.19.8}
\]
It is continuous and independent of the source lift \(w\).
Its distinguished fibre is still exactly the fibre
(F.18.10).

Here are explicit subgroup maps accompanying that comparison.
Set \(H^\dagger=\Phi^\dagger(L^\dagger)\) and
\[
 \begin{aligned}
 \gamma:L&\longrightarrow\mathcal Q,&
       \gamma(u)&=(\widehat\phi_*(u),\widehat\chi_*(q_0u)),\\
 \mathcal T&=\pi^{-1}(\gamma(L)).
 \end{aligned}                                                 \tag{F.19.9}
\]
The map \(\gamma\) is an embedding.  There is a commutative
diagram of exact full profinite sequences
\[
 \begin{array}{ccccccccc}
 1&\longrightarrow&K^+&\longrightarrow&H^\dagger
     &\xrightarrow{\ q_H\ }&L&\longrightarrow&1\\
 &&\downarrow&&\downarrow&&\Vert\\
 1&\longrightarrow&D^+&\longrightarrow&\mathcal T
     &\xrightarrow{\ \pi_L\ }&L&\longrightarrow&1 .
 \end{array}                                                   \tag{F.19.10}
\]
Both rows split by the same specified map
\(j=\Phi^\dagger i\), where \(i:L\to L^\dagger\) retains
the eleven old free meridians in the source filling.
The corresponding inverse coordinates are
\[
 \begin{aligned}
 D^+\rtimes_j L&\longrightarrow\mathcal T,
       &(d,u)&\longmapsto d\,j(u),\\
 \mathcal T&\longrightarrow D^+\rtimes_j L,
       &t&\longmapsto
          \bigl(tj(\pi_Lt)^{-1},\pi_Lt\bigr).
 \end{aligned}                                                 \tag{F.19.11}
\]
Restricting these maps gives \(K^+\rtimes_jL\cong H^\dagger\).
Moreover
\[
                  \mathcal T\cap C^\dagger=\ell(\mathcal E).
                                                               \tag{F.19.12}
\]
These are assertions about the actual images and inverse maps;
they do not assert \(H^\dagger=\mathcal T\) or \(K^+=D^+\).
The source-membership equation in (F.19.8), and triviality
of \(\mathcal E\), remain unresolved.
\end{claim}
\passage{Section 30, folios 194-196, scans 383-387; Section 31, folios 198-199 verso, scans 391-394; Part F(b), the exact residual centralizer and its source-membership comparison}


\begin{proof}
\emph{The literal six-point centralizer.}
The retained involution on \(A\cup Z\) exchanges
\((0,\infty)\) and \((1,v_\zeta)\), and fixes \(a_+,a_-\).
Let \(C_6^\dagger\) denote its full pure centralizer,
and let \(C_4\) denote the full pure centralizer after
forgetting \(Z\).
There is a branch-permutation character on \(C_4\);
write \(C_4^+\) for its kernel.
The discrete decorated quotient has two ordinary orbit
labels and two branch labels.
The deck involution is nonpure before and after adding
the branch labels, so descent is injective on both
pure centralizers.
The actual inverse lift argument of F.18.2 therefore gives
\[
              h:C_6^\dagger\xrightarrow{\ \sim\ }C_4^+ .
                                                               \tag{F.19.13}
\]
For this application the prime-order completion and full
topology comparison F.7.4--5 is applied to these very
four- and six-point actions.
The branch permutation is onto: exchange the two
quotient branch points, then correct any of the two
ordinary pair-name changes by pushing that ordinary
quotient point once around a branch point.
These corrections preserve the branch permutation.
Thus this is the same exact index-two comparison as
in F.18, with two old pairs retained.

In fact
\[
             \widehat\chi_*(\widehat{\pi_1(Y_0,s_0)})
                                =C_6^\dagger .                \tag{F.19.14}
\]
To check equality with this particular map, descend by
the actual quotient \(Q_a(V)=((V-a)/(V+a))^2\).
F.17.12 identifies its four ordered quotient values
as \(0,\infty,1,t\), \(t=s^2/\zeta^2\).
There are exactly two choices of pure ordinary sheet
names modulo simultaneous deck exchange.
The sheet change \(s\mapsto-s\), followed by
\(V\mapsto V/a^2\), exchanges precisely the names in
\(\{1,a^2\}\); the deck exchanges both ordinary pairs.
Hence the degree-two map \(Y_0\to\mathcal M_{0,4}\)
is exactly the covering belonging to the discrete
pure-lift subgroup of this quotient group.
Its fundamental group is the entire image of the
discrete six-point pure centralizer under descent.
The family of involutions on the six labelled points
shows that \(\chi_*\) lands in that centralizer.
Injective descent now proves the discrete equality.
The full topology of \(\chi_*\), proved using the
ordinary \(a^2\)-cover in F.17, and the prime-order
centralizer comparison complete that equality to
(F.19.14).
Both the quotient parameter \(t\) and ordinary
forgetting parameter \(a^2\) are retained in this
argument; no identification of those two maps is used.

\emph{The actual fibre-product map and its completion.}
Normalize the three common labels \(0,\infty,1\).
The space
\[
 X=\mathcal M(A\cup U)
          \times_{\mathcal M(A)}\mathcal M(A\cup Z)
                                                               \tag{F.19.15}
\]
has coordinates \(v_\zeta\), four ordered \(U\)-coordinates
and two ordered \(Z\)-coordinates.
The \(U\)-coordinates are distinct and avoid \(A\);
the \(Z\)-coordinates are distinct and avoid \(A\).
At this stage a \(U\)-coordinate may equal a
\(Z\)-coordinate.
Both projections to the four-point base are the
actual configuration forgetting fibrations.
The base \(\mathcal M_{0,4}\) has \(\pi_2=0\).
The homotopy sequences therefore identify \(\pi_1(X)\)
with the discrete fibre product of the two pure
mapping class groups over the four-point group:
the kernel on both sides is the product of the
two configuration-fibre groups, and the maps on
those kernels and on the base group are identities.

Here is also the full-topology comparison for this
particular fibre product.
The map \(\mathcal M(A\cup Z)\to\mathcal M(A)\)
has a continuous section.  For the variable coordinate
\(u=v_\zeta\), place the two new coordinates at
\(2+|u|\) and \(3+|u|\).
They are distinct and avoid \(0,1,u,\infty\).
A path in the connected fibre identifies the
resulting basepoint with the specified \((a_+,a_-)\).
This gives a based group section.
The discrete kernel is
\[
 J_{4,\mathrm d}
   =\pi_1(\operatorname{Conf}_2(S_{0,4}),(a_+,a_-)),
 \qquad
 1\longrightarrow F_4\longrightarrow J_{4,\mathrm d}
                     \longrightarrow F_3\longrightarrow1 .
                                                               \tag{F.19.16}
\]
The latter sequence splits by the same far-point
construction in a punctured-plane coordinate.
In particular \(J_{4,\mathrm d}\) is finitely generated.

Consequently the discrete group in (F.19.15) is
\(J_{4,\mathrm d}\rtimes P\Gamma(A\cup U)\), with
the action pulled back from the displayed actual
section over \(A\).
For every finite-index subgroup of the kernel,
choose a characteristic finite-index subgroup inside
it.  The action on this finite quotient has finite
image; quotienting the acting group by its action
kernel gives a finite semidirect quotient of the
whole group.  This proves the full induced topology
on the kernel.  It also extends the action to the
completions and proves exactness of the completed
split sequence: finite homomorphisms on both factors
are separated by these compatible finite quotients.
The completed Birman comparison identifies its kernel
with the actual \(\mathcal J_4\) in \(P_6\).
The same section thus gives
\[
 \widehat{\pi_1(X)}
   \cong\mathcal J_4\rtimes P_8
   \cong P_8\times_{P_4}P_6=\mathcal Q
                                                               \tag{F.19.17}
\]
with the actual projection maps.
This proves completion for the specified fibre product.

For clarity, surjectivity in (F.19.2) can also be
seen directly on those actual full maps.
Lift a given \(P_8\) component through \(f\).
Its \(P_6\) component differs from the desired one
by an element of \(\mathcal J_4\).
The filling map
\[
 \pi_1\operatorname{Conf}_2(S_{0,8})
                  \longrightarrow
       \pi_1\operatorname{Conf}_2(S_{0,4})
\]
is onto.  Apply the two Fadell--Neuwirth sequences:
both the one-point quotient \(F_7\to F_3\) and
the fibre map \(F_8\to F_4\) are the onto puncture
fillings, so a lift on the quotient can be corrected
in the fibre.
Its completion is onto as well, and these completions
are exactly \(\mathcal J\to\mathcal J_4\) by the
completed Birman comparisons.
Correcting the chosen lift inside \(\mathcal J\)
gives the required pair.
This also proves the exact sequence in (F.19.4).

\emph{The eight actual cross-diagonal generators.}
The ten-point moduli space is the open subset of
\(X\) obtained by removing exactly
\[
       \mathcal D_{u,\pm}=\{Z_u=Z_\pm\}\qquad(u\ne\zeta).
                                                               \tag{F.19.18}
\]
These are eight smooth complex hypersurfaces.
Each has connected smooth part after removing its
intersections with the others: identify the coincident
pair, then move the remaining finitely many labelled
coordinates along paths avoiding their finite forbidden
sets.
Different hypersurfaces meet only in complex codimension
at least two; their possible intersections represent
independent pair collisions.

Every loop in \(X\) can be moved off these hypersurfaces.
A disk witnessing a nullhomotopy in \(X\) can be made
transverse to their smooth parts and disjoint from
their intersections.  Its finitely many intersection
points, joined by based arcs in the disk, express
its original boundary as a product of conjugates
of the corresponding signed normal meridians.
Conversely every such meridian dies in \(X\).
This proves that the discrete kernel is the normal
closure of one based meridian for each hypersurface.

Choose its based path to be the actual source
whisker from F.17 approaching either of the two
simple collision points for that pair.
The other labels remain distinct there.
The coefficient-one local calculation (F.17.7--9)
identifies its normal meridian exactly with
\(\nu_{u,\pm}\), with that path and a positive
one-turn exponent.
The finite-target universal property of a quotient
now identifies the kernel on full completions
with the closed normal closure of these eight
elements.  Using (F.19.17) gives exactly (F.19.3).
Every chosen source meridian is killed by \(q\),
so every \(\nu_{u,\pm}\) belongs to the specified
subgroup \(K^+\).

There is a useful exact check on what these ordinary
winding coordinates detect.
The eight global functions \(Z_u-Z_\pm\) give
a continuous homomorphism
\[
 \operatorname{wind}:P_{10}^{\dagger}\longrightarrow
                 \widehat\Z^{\,8},\qquad
 \operatorname{wind}(\nu_{u,\pm})=e_{u,\pm}.
                                                               \tag{F.19.19}
\]
At either root for a given pair, its own coordinate
has winding one and the other seven have winding zero.
The selected eight elements already show that the
images of both \(K^+\) and \(D^+\) under this
particular homomorphism are all of
\(\widehat\Z^{\,8}\): take the product of their
specified full powers to obtain any eight coordinates.
No assertion about the whole abelianization of
\(D^+\) follows from this calculation.

\emph{The two section labels under this filling.}
The action on each two-point configuration fibre
fixes the two base points individually:
\((z_+,z_-)\mapsto(rz_+,rz_-)\).
Its fixed subgroup in \(\mathcal J_4\) is the kernel
of (F.19.13), and is therefore trivial.
Its continuous section classes have the same
two labels as in F.18.  Here are the inputs for
this smaller actual group.
The discrete extension
\(J_{4,\mathrm d}\rtimes C_2\) is the subgroup of
the six-point mapping class group generated by
the point-pushing kernel and the indicated involution.
Each order-two section exchanges the two old
pairs and fixes the two new labels.
Nielsen realization and the ordered quotient
homeomorphism identify it with the standard
involution by a pure six-point conjugator.
The double-coset comparison of those conjugators is
\(C_{4,\mathrm d}/C_{4,\mathrm d}^+=C_2\).
The label is the branch permutation of the
forgotten conjugator.
All their discrete fixed subgroups in
\(J_{4,\mathrm d}\) are trivial, by the discrete
version of (F.19.13).

The split free-by-free sequence (F.19.16) gives the
actual full topology and the two-row cohomology
comparison, just as in F.18.12.
Finite graphs for \(F_4\) and \(F_3\), with mapping
cylinders of representatives of their monodromy
homotopy equivalences, give a finite classifying
complex of dimension two.
The extension is residually finite as this actual
mapping class subgroup.  The finite extension
cohomology comparison and finite virtual
cohomological-type hypotheses are therefore the
same verified ones as in F.18.13--14.
The prime-order theorem \cite[Theorem~A]{R-BZ}
transports the two conjugacy classes and their
centralizers to the full completion.
Conjugacy by the extension can be changed to
conjugacy by its configuration kernel by removing
the quotient-generator power of the conjugator.
Thus the full \(H^1(C_2,\mathcal J_4)\) has exactly
these two branch labels.

The filling \(\mathcal J\to\mathcal J_4\) carries
the two full classes of F.18 to these two classes
with the same labels.
At the discrete level this follows by forgetting
the other two ordinary quotient orbits: it preserves
both named branch points and their permutation.
The finite branch characters are continuous by
the full-topology comparison, so the identical
label statement holds in the completions.
The induced map on these two cohomology sets is
therefore a bijection.

Let \(\delta\in D^+\) satisfy \(\delta R(\delta)=1\).
Its image in \(\mathcal J_4\) is the identity
cocycle.  The just-proved bijection forces its
\(\mathcal J\)-class to be the identity class.
By F.18 there is a unique \(b\in\mathcal J\) with
\(\delta=b^{-1}R(b)\).
Its image \(g(b)\) is fixed in \(\mathcal J_4\),
so is the identity.
Hence \(b\in D^+\).
Uniqueness also follows from
\(C_{D^+}(R)\subset C_{\mathcal J}(R)=1\).
The coboundary map is continuous; compactness of
\(D^+\) and the Hausdorff cocycle space make this
continuous bijection a homeomorphism.
This proves (F.19.4--5).

For \(z\in\mathcal Q^R\), choose \(k\in\pi^{-1}(z)\).
Then \(k^{-1}R(k)\) is a cocycle in \(D^+\).
Its unique corrector \(b\in D^+\) makes
\(kb^{-1}\) fixed, and does not change its image
under \(\pi\).
The kernel of \(\pi\) on the fixed subgroup is
\(C_{D^+}(R)=1\).
This proves (F.19.6), including its actual inverse
and continuity.

\emph{The specified source image inside the pullback.}
For \(u\in\mathcal E\), the element \(g\ell(u)\)
lies in \(C_6^\dagger\).
Its ordinary four-point image is
\(f_0\widehat\phi_*(u)\).
The exact source square (F.17.10--11) identifies
that with \(h\widehat\chi_*(q_0u)\).
The latter element also belongs to \(C_6^\dagger\)
by (F.19.14).
Injectivity in (F.19.13) proves the first equality
of (F.19.7).
For \(qw=u\), the same square gives
\(g\Phi^\dagger(w)=\widehat\chi_*(q_0u)\).
Thus \(g(b(w))=1\); its \(f\)-image is already
the identity by F.18.
This proves \(b(w)\in D^+\).
For \(qv=1\), both forgetting images of
\(\Phi^\dagger(v)\) are the identity, so
\(K^+\subset D^+\).
The right-coset comparison and its independence
are the actual equality
\(b(wv)=b(w)\Phi^\dagger(v)\).
The quotient map \(q^{-1}(\mathcal E)\to\mathcal E\)
is a compact surjection, so that equality and
continuity give the asserted continuous descended
map (F.19.8).
The inclusion \(D^+/K^+\to\mathcal J/K^+\)
is a closed embedding of compact coset spaces.
Hence its distinguished fibre is precisely
the earlier fibre (F.18.10).

The same source square puts \(\gamma(L)\) in
\(\mathcal Q\), and injectivity of
\(\widehat\phi_*\) makes \(\gamma\) an embedding.
Define \(\pi_L\) by the inverse of that embedding
after applying \(\pi\) to \(\mathcal T\).
It is onto by (F.19.2).
For \(z=\Phi^\dagger(w)\in H^\dagger\), the equality
\(\pi(z)=\gamma(qw)\) shows that \(q_H(z)=qw\)
is independent of \(w\).
This is a continuous surjection: it is induced
by a continuous map on the compact source.
Its kernel is exactly \(K^+\).
These facts prove every map and kernel in
(F.19.10).

Use the free based meridians in F.17.5.
The filling \(q\) kills exactly the sixteen new
generators and sends the eleven old finite
meridians to a free basis of \(L\).
Sending that latter basis back to those eleven
specified meridians defines a continuous map
\(i:L\to L^\dagger\) with \(qi=1\).
Thus \(j=\Phi^\dagger i\) splits both rows of
(F.19.10).
Conjugation by \(j(L)\) preserves \(D^+\), and
also preserves \(K^+\), the kernel in the upper
row.  The multiplication and inverse coordinates
in (F.19.11) are consequently homomorphisms
of the stated semidirect products and are
inverse continuous maps.
This constructs the common splitting directly
from the named free meridians.

Finally an element of \(\mathcal T\cap C^\dagger\)
has projection \(\gamma(u)\) for some \(u\in L\).
Its first component being fixed is exactly the
condition \(u\in\mathcal E\), by F.12.
The second component is already fixed by
(F.19.14).
Equation (F.19.6) gives a unique fixed lift of
that pair; (F.19.7) identifies it with \(\ell(u)\).
Conversely \(\ell(u)\) belongs to \(\mathcal T\)
for every \(u\in\mathcal E\).
This proves (F.19.12).

The new rational comparison checked by
\path{code/part_a/common_reflection_cross_kernel_f19.py}
is particularly explicit.
Writing \(w_j=v_{\zeta^j}\), \(p=a_+\), \(m=a_-\),
the normalized action sends
\[
 (w_0,w_1,w_2,w_3,w_4,p,m)\longmapsto
 \left(\frac{w_1}{w_2},w_1,\frac{w_1}{w_0},
       \frac{w_1}{w_4},\frac{w_1}{w_3},
       \frac{w_1}{p},\frac{w_1}{m}\right).
                                                               \tag{F.19.20}
\]
It fixes the retained curve
\((w_1,p,m)=(a^2,a,-a)\).
For \(\sigma=(0\ 2)(3\ 4)\) and \(z=p\) or \(m\),
its cross difference is exactly
\[
 R(w_j)-R(z)
       =-\frac{w_1}{w_{\sigma(j)}z}
                      \bigl(w_{\sigma(j)}-z\bigr).
                                                               \tag{F.19.21}
\]
The coefficient is nonzero in the relevant
configuration chart.  Thus the eight cross
divisors and their labels are preserved in
the stated four pairs.
The executable checks these identities, the
involution identity and the sixteen-column
winding matrix from the simple source collisions.
It passes 44 exact checks.
The fibre-product, completion and section arguments
above establish the group assertions; the calculation
does not identify \(K^+\) with \(D^+\).
\end{proof}

\begin{claim}{F.20}{Proved}{An exact cross-divisor action and its finite wreath images}
Retain the labelled source, the two-track group, and the subgroups
$K^+\leq D^+$ of F.17--F.19. Select the cross divisor
between the old point $v_1$ and the positive new point $a$.
Its full relative path coordinate is
$\mathcal F=\widehat F(x,y,z)$, with the three specified
meridians corresponding, in order, to $0,v_\zeta,1$.
The actual source $L=\widehat{\pi_1(Y)}=\widehat F_{11}$
has the continuous action
\[
       A_g(v)=q_g^{-1}\alpha_g(v)p_g
                 \qquad(g\in L,\ v\in\mathcal F),
                                                               \tag{F.20.1}
\]
where $\alpha_g,p_g,q_g$ are the literal based braid data
specified below. The diagonal-coset map proving this
formula, its inverse, and every based source transport
are retained in the proof.

Let $\mathscr Q$ be the following explicitly enumerated
finite quotients of this same $\mathcal F$:
the $82$ source-invariant $A_5$ quotients in the enumeration
below, all $3,321$ products of two distinct such quotients,
and the six mod-two homology-cover quotients
\[
  Q_\rho=\mathcal F/
       \overline{[N_\rho,N_\rho]\,N_\rho^2},
       \qquad N_\rho=\ker(\rho:\mathcal F\to A_5),
       \qquad |Q_\rho|=60\cdot2^{121},
                                                               \tag{F.20.2}
\]
whose four peripheral images under $\rho$ are all nonidentity.
Every action (F.20.1) on every member of $\mathscr Q$ is
transitive. In each of the six last cases the image of the
source filling kernel acts by right translations through
a subgroup of index four; the remaining actual source
generators connect all four cosets.

For each $Q\in\mathscr Q$ and each prime $\ell$ there is
an explicitly specified continuous homomorphism
\[
 \chi_{Q,\ell}:D^+\longrightarrow\mathbb F_\ell[Q]
 \quad\hbox{with}\quad
 \chi_{Q,\ell}(K^+)=\chi_{Q,\ell}(D^+)=\mathbb F_\ell[Q].
                                                               \tag{F.20.3}
\]
It is obtained from the actual one-point forgetting map
and the finite wreath assignment in (F.20.15) below.
Thus these maps give the same full finite images for
the two specified subgroups. This conclusion does not
identify $K^+$ with $D^+$ or solve the remaining
source-membership equation of F.19.
\end{claim}
\passage{Section 30, folios 194-196, scans 383-387; Section 31, folios 198-199 verso, scans 391-394; Part F(b), the actual residual A5-centralizer and source-membership comparison}


\begin{proof}
\emph{The selected two-track space and its actual diagonal.}
Normalize the common labels $0,\infty,1$ as in F.19 and put
\[
 \begin{gathered}
 \zeta=e^{2\pi i/5},\qquad
 Y=\mathbb P^1\setminus\{0,\infty,\pm\zeta^j:0\leq j<5\},\\
 a(s)=\frac{\zeta-s}{\zeta+s},\quad
 B(s)=a(s)^2,\quad
 P(s)=\left(\frac{1-s}{1+s}\right)^2,\quad Q(s)=a(s).
 \end{gathered}                                                \tag{F.20.4}
\]
The letter $Q(s)$ in this display is a point coordinate;
a finite quotient $Q\in\mathscr Q$ will always be named as such.
Forget the other three old labels and the negative new label
in the F.19 two-track space $X$. The target is exactly
\[
 X_1=\{(B,U,V):B\ne0,1,\quad
                  U,V\notin\{0,1,B\}\}
       =\mathcal M_{0,5}\times_{\mathcal M_{0,4}}\mathcal M_{0,5}.
                                                               \tag{F.20.5}
\]
It allows $U=V$; that equality is its selected diagonal.
The curve map is literally $s\mapsto(B(s),P(s),Q(s))$.
The selected divisor upstairs in $X$ is its inverse image.

The comparison with the original F.17 basepoint is
also specified. Write
$s_0=i(\theta^2-2)\zeta$, as in F.12, and take the
polygonal path with successive vertices
\[
        s_0,\quad 3s_0,\quad 3s_0+6i,\quad
                         6+6i,\quad 2+2i.                    \tag{F.20.4a}
\]
It lies in $Y^\dagger$.
Indeed $|s_0|^2=\theta^2-1>1$.
For the four $A_u$ of F.17.2, the imaginary parts
of $A_u\overline{s_0^2}$ are respectively
$\theta^3/2,-\theta^3/2,
 (5\theta-3\theta^3)/2,
 (-5\theta+3\theta^3)/2$ in the order
$u=1,\zeta^2,\zeta^3,\zeta^4$; all are nonzero.
Thus its first radial segment misses the outer
collision roots. The inner collision roots have
absolute value less than one.
All later segments lie outside the disk of radius
two: $3\operatorname{Re}s_0<-2$, the middle
horizontal-direction segment has imaginary part
at least six, and the last segment has both
coordinates at least two.
Every collision root has absolute value less
than two, by $|A_u|<3$ and $|B_u|<1$.
The exact producer checks these inequalities.

Call this path $\eta$.
The based maps are $g\mapsto\eta g\eta^{-1}$
on the source and conjugation by the actual
labelled configuration path $\phi^\dagger(\eta)$
on the ten-point target, with the identical
forgettings of that path on every smaller target.
The composites commute because these are
literal paths of the displayed rational labels.
They carry the $K^+,D^+$ at $s_*=2+2i$ to the
original F.19 subgroups at $s_0$.
The homomorphisms in (F.20.3) on those original
subgroups include the inverse of this specified
basepoint change.

The forgetting projection is a configuration fibration with
connected fibers. Over its diagonal the additional three
old points and the other new point still vary in their
respective connected configuration fibers. Consequently
the image of the divisor fundamental group is the full
inverse image of the diagonal fundamental group: lift a
diagonal loop through this fibration, and correct its
endpoint within the connected fiber. The kernel of the
projection is represented by loops in that same fiber
and already lies in the divisor image.

The same comparison holds on full completions.
For clarity, write $P_j=\widehat{\pi_1(\mathcal M_{0,j})}$
with the specified label maps. The diagonal inverse image
is the fiber product $P_8\times_{P_5}P_6$; its kernel over
$P_8$ is the actual one-point-pushing $\widehat F_4$.
The discrete sequence has an actual continuous
configuration section, obtained by placing the forgotten
point on the positive real axis beyond
$2+|B|+|U|$, and then choosing the based fiber path.
The characteristic finite kernel quotients used in F.19
give its full induced topology and the explicit completed
semidirect-product inverse maps. They identify its
completion with this stated fiber product.
The identical argument for (F.20.5) gives
$\mathcal Q_1=P_5\times_{P_4}P_5$.
Thus the full coset space for the selected divisor in
F.19 is carried homeomorphically onto
$\mathcal Q_1/\operatorname{diag}(P_5)$.
No closure/intersection exchange is used.

\emph{The based source and the two reference transports.}
Use source basepoint $s_*=2+2i$, reference point $R=8+3i$,
and projection $z\mapsto\operatorname{Re}((1+371i/1000)z)$.
For each of the eleven centers
\[
       0,\ 1,\zeta,\zeta^2,\zeta^3,\zeta^4,\
          -1,-\zeta,-\zeta^2,-\zeta^3,-\zeta^4
                                                               \tag{F.20.6}
\]
take the straight whisker from $s_*$ to the northeast
corner of the axis-parallel square of half-side
$31/1000$ about that center, traverse the square
counterclockwise, and return along that same whisker.
Call these loops $g_0,\ldots,g_{10}$ in the displayed order.
The certificate verifies that the squares are pairwise
disjoint, each whisker misses all other squares, and
the whiskers have distinct directions from their common
basepoint. The basepoint lies northeast of every square,
so its own whisker first meets that square at its corner.
These squares and the embedded star give an actual
geometric free basis of $\pi_1(Y)$.

At $s_*$, transport the fourth finite point along the
straight segment from $R$ to $P(s_*)$ for the old track
and to $Q(s_*)$ for the new track, keeping $0,1,B(s_*)$
fixed. Both paths avoid those three punctures; their
exact cross-product checks are retained. Call these
paths $T_o,T_n$. A source loop gives the two based
four-strand braids
$T_o\,P(g)\,T_o^{-1}$ and $T_n\,Q(g)\,T_n^{-1}$.
The notation here retains the actual intermediate
three-strand path $B(g)$.
At the common reference configuration the projected
order is exactly $(0,B(s_*),1,R)$.
Use the same disk basepoint and straighten the four
points within their projection-order chamber. A generator
$\sigma_j$ is the literal clockwise exchange of adjacent
points: the left strand passes above the right strand.
For example, after translation and positive scaling, the
two paths
\[
 2t-1+4it(1-t),\qquad -2t+1-4it(1-t)
       \quad(0\leq t\leq1)
\]
have exactly this exchange and the producer's positive
crossing. Their squared distance from the midpoint is
$((2t-1)^2-1/2)^2+3/4$, so the pair never collides.
These geometric generators give the Artin presentation
of the actual braid group. The symbols $x,y,z,d$ below
are its ordered free coordinates; the exact identification
of $x,y,z$ with the actual point-pushing fibre, including
its inverse based loops, is given by (F.20.8)--(F.20.9).

\emph{The Artin coordinate and the diagonal-coset formula.}
Use the faithful Artin substitutions
\[
 \sigma_j:
 X_j\longmapsto X_jX_{j+1}X_j^{-1},\quad
 X_{j+1}\longmapsto X_j ,
 \qquad \sigma_j(X_k)=X_k\ (k\ne j,j+1),
                                                               \tag{F.20.7}
\]
in the word-composition convention of \cite{A-Artin1947}.
Path products are read from left to right in time, and
the substitution for a displayed word $b_1\cdots b_r$
is $\phi_{b_1}\circ\cdots\circ\phi_{b_r}$, as in the
producer. The displayed substitutions and their inverses
satisfy the Artin relations, so the named geometric
generators define this homomorphism into the free-group
automorphisms. The following fibre-coordinate comparison
determines its actual point-pushing action.

For a pure four-strand braid $\beta$, write
$\phi_\beta(d)=W_\beta dW_\beta^{-1}$ and let $w_\beta$
be the word obtained by filling $d$ in $W_\beta$.
The ambiguity in $W_\beta$ is a power of $d$ and
therefore disappears under this filling.
Write $\alpha_\beta$ for the filled automorphism of
$F(x,y,z)$. Direct substitution gives
\[
 w_{\beta\gamma}=\alpha_\beta(w_\gamma)w_\beta,\qquad
 u_\beta=w_\beta^{-1},\qquad
 u_{\beta\gamma}=u_\beta\alpha_\beta(u_\gamma).
                                                               \tag{F.20.8}
\]
To identify the actual point-pushing kernel, let
$\beta_j=\sigma_3\cdots\sigma_{j+1}\sigma_j^2
 \sigma_{j+1}^{-1}\cdots\sigma_3^{-1}$.
The first descending word carries the last strand along
the lower-half-plane whisker, the square makes a clockwise
lasso about puncture $j$, and the inverse word returns
that whisker. Deleting the last strand contracts the
old motion within its projection-order chambers. Thus
these are the three actual lower-whiskered clockwise
lassos, a free basis of the punctured-plane fibre.
Direct substitution gives their $u$-images
\[
        z^{-1}y^{-1}x^{-1}yz,\qquad
        z^{-1}y^{-1}z,\qquad z^{-1}.
                                                               \tag{F.20.9}
\]
This triangular substitution is its own inverse, by
free reduction. The inverse map sends $x,y,z$ to the
same three words with $\beta_1,\beta_2,\beta_3$
substituted for $x,y,z$. These are specified conjugates
of the inverse clockwise lassos, so they are actual
based puncture meridians. Both composites are the
identity on the named free bases. This gives the actual
fibre isomorphism and then its full profinite inverse.
Adding the last strand gives the standard section on
the three-strand braid group, with $u=1$.
The cocycle (F.20.8) therefore makes conjugation by that
section exactly the filled Artin substitution, and makes
conjugation by a general pure braid $b$ the map
$v\mapsto u_b\alpha_b(v)u_b^{-1}$.
Its restriction gives the section of $P_5\to P_4$
used in this coordinate. Filling a puncture kills its
lower-whiskered lasso; the same triangular formula
commutes with filling the corresponding variable and
retaining the order of the others. These are the actual
based filling maps in the displayed coordinates.

Here the normalized moduli groups have not been replaced
by braid groups without a map. The coordinate isomorphism
\[
 (z_0,z_1,z_B,z_R)\longmapsto
 \left(z_0,z_1-z_0,\
       \frac{z_B-z_0}{z_1-z_0},
       \frac{z_R-z_0}{z_1-z_0}\right)
                                                               \tag{F.20.10}
\]
identifies the ordered configuration space with
$\mathbb C\times\mathbb C^\times\times\mathcal M_{0,5}$;
its inverse multiplies the last two coordinates by
$z_1-z_0$ and adds $z_0$.
The analogous three-point formula has
$\mathcal M_{0,4}$ as last factor.
The original paths keep $z_0=0,z_1=1$.
The braid section commutes with forgetting to these
two labelled points, so it restricts to the kernels
of the displayed affine-coordinate projections.
The section and its retraction extend to full
completions and have the same composites.

After the based transports, both factors of
$\mathcal Q_1$ are this same $P_5$. The mutually inverse
coset maps are
\[
 \begin{aligned}
 \mathcal Q_1/\operatorname{diag}(P_5)&\longrightarrow\mathcal F,
       &(b_o,b_n)\operatorname{diag}(P_5)&\longmapsto b_nb_o^{-1},\\
 \mathcal F&\longrightarrow
                 \mathcal Q_1/\operatorname{diag}(P_5),
       &v&\longmapsto(1,v)\operatorname{diag}(P_5).
 \end{aligned}                                                \tag{F.20.11}
\]
The first value belongs to the kernel because both
factors have the same base image. Both maps are
continuous and their composites are the identity.
Left multiplication by the two source braids acts
by $v\mapsto b_nvb_o^{-1}$.
Substituting $b_o=(p_g^{-1},\alpha_g)$ and
$b_n=(q_g^{-1},\alpha_g)$ from (F.20.8) proves
(F.20.1), including its order of multiplication.

\emph{The complete generator data and exact crossing certificate.}
Set $P=xyz$, $k=Pxyx$, and $t=Pxyx^{-1}$.
Let $\tau_1,\tau_2$ be the automorphisms (F.20.7)
of $F(x,y,z)$ associated with $\sigma_1^2,\sigma_2^2$.
The complete source data are
\[
\begin{array}{c|c|c|c}
i&\alpha_{g_i}&p_{g_i}&q_{g_i}\\ \hline
0&\tau_1\tau_2^{-1}\tau_1^{-1}
  &(xy)x(xy)^{-1}P^{-1}&z^{-1}xy^{-1}x^{-1}\\
1&1&x^{-2}&1\\
2&\tau_1^{-2}&1&(xy)^{-1}\\
3&1&1&1\\
4&1&ky^{-1}k^{-1}&1\\
5&1&1&1\\
6&1&P^2&1\\
7&(\tau_1\tau_2)^2&tz^{-1}x^{-1}&t\\
8&1&1&1\\
9&1&xy^{-1}x^{-1}&1\\
10&1&1&1 .
\end{array}                                                   \tag{F.20.12}
\]
These are values in the specified basis, not
independently chosen conjugacy classes.

The executed producer
\path{code/part_a/cross_divisor_monodromy_f20.py}
works in $\mathbb Q(\theta)$, where
$\theta^4-5\theta^2+5=0$ and $19/10<\theta<2$, using
$\zeta=(\theta^2-3+i\theta)/2$.
Every sign is separated by rational interval arithmetic
at that particular real root.
For each rational configuration segment and each pair
of strands it multiplies the projected difference
numerator by the conjugate of its denominator.
Its real polynomial has exactly the projected
crossings; its imaginary polynomial gives their
heights, since the omitted denominator is positive.
All possible poles and collisions of either separate
track are among (F.20.6): in particular
$P(s)=B(s)$ has numerator
$4s(1-\zeta)(\zeta-s^2)$, whose remaining roots
are $\pm\zeta^3$.
The source-basis clearance therefore also checks that
these paths stay in the stated configuration spaces.

Exact Bernstein subdivision supplies the root count.
On a rational interval, multiplication by the positive
binomial coefficients changes its Bernstein sequence
to the coefficients of the usual fractional-linear
transform for Descartes' rule. Zero sign variation
excludes roots and one variation gives exactly one
root with multiplicity one. Intervals with larger
variation are subdivided. A subdivision point that
is an exact root is avoided by choosing another
rational fraction; a nonzero polynomial of degree
$d$ cannot vanish at $d+1$ such distinct choices.
The output retains the complete projected and height
polynomials, the terminal root-isolating intervals
and their Bernstein signs, and the counts of discarded
root-free leaves; the executed code reproduces the
subdivisions. The certificate further
separates a nonzero derivative and nonzero height
on every crossing interval.
Intervals for different pairs are disjoint.
Their exact order, initial projected order,
adjacency at each crossing, and final projected
order give the literal braid word.

There are $112$ certified forward segments,
$672$ projected-pair checks, $172$ crossings, and
$72,380$ exact number-field sign calls.
The return whisker is the literal inverse path,
so its braid is the inverse of its certified
forward braid. The two reference transports and
all eleven old and new based braids are included.
Deleting the fourth strand gives respectively the
three nontrivial base braid words
$\sigma_1^2\sigma_2^{-2}\sigma_1^{-2}$,
$\sigma_1^{-4}$, and
$(\sigma_1^2\sigma_2^2)^2$ at indices $0,2,7$.
The second producer verifies this strand deletion,
all entries of (F.20.12), and the inverse substitution
(F.20.9) by exact free reduction.
Neither producer reads a numerical discovery file.

\emph{The actual full one-point quotient and its finite wreath maps.}
Forget the negative new point in F.19's two-point
configuration kernel. The two Fadell--Neuwirth rows
and their filling maps are
\[
 \begin{array}{ccccccccc}
 1&\longrightarrow&\widehat F_8&\longrightarrow&J
                 &\longrightarrow&\widehat F_7&\longrightarrow&1\\
 &&\downarrow&&\downarrow&&\downarrow\\
 1&\longrightarrow&\widehat F_4&\longrightarrow&J_4
                 &\longrightarrow&\mathcal F&\longrightarrow&1 .
 \end{array}                                                  \tag{F.20.13}
\]
Here $J$ and $J_4$ are the actual eight- and four-puncture
two-point groups of F.19, and $D^+=\ker(J\to J_4)$.
Every vertical filling is onto. The full topologies
and exact rows are those proved in F.19 by the
actual configuration sections and finitely generated
free kernels.
Put $R_0=\ker(\widehat F_7\to\mathcal F)$.
The induced map
\[
                         D^+\longrightarrow R_0
                         \quad\hbox{is onto}.               \tag{F.20.14}
\]
Indeed lift an element of $R_0$ to $J$.
Its image in $J_4$ lies in the displayed $\widehat F_4$;
an element of $\widehat F_8$ has its inverse as image.
Multiplication by that element gives the required
lift in $D^+$ without changing its $\widehat F_7$ image.

Choose the geometric free basis of the seven-puncture
one-point group consisting of $x,y,z$ and the four old
puncture meridians $u_j$. Its filling kills the $u_j$
and retains $x,y,z$. The whisker of the selected
$u_1$ is chosen to project to the old reference
transport specified above after changing the
one-point basepoint by $T_n$; at the original
one-point basepoint its projected path is
$T_n^{-1}T_o$. The straight $T_n$ also avoids
all four other old points, as checked by their
nonzero exact cross products.
This is an explicit
allowable basis adjustment: a difference word $w$
in the retained $x,y,z$ replaces $u_1$ by
$wu_1w^{-1}$, fixes the other generators, and has
inverse $u_1\mapsto w^{-1}u_1w$.
The difference word is the class of the two stated
based paths. Thus the full filling and its splitting
are unchanged, with their exact maps retained.

For $\rho:\mathcal F\to Q$ in the stated collection
and a prime $\ell$, the finite free-group assignment is
\[
 \begin{aligned}
 \widehat F_7&\longrightarrow
                    \mathbb F_\ell[Q]\rtimes Q,\\
 x,y,z&\longmapsto(0,\rho(x)),(0,\rho(y)),(0,\rho(z)),\\
 u_1&\longmapsto(e_1,1),\qquad
 u_j\longmapsto(0,1)\quad(j\ne1).
 \end{aligned}                                                \tag{F.20.15}
\]
Here $Q$ acts by its actual left regular permutation
of the basis $e_a$, $a\in Q$.
On $R_0$ it is a homomorphism to the additive module
and is onto: conjugating $u_1$ by a retained word
with image $a$ gives $e_a$.
Composing with (F.20.14) defines $\chi_{Q,\ell}$
and proves its surjectivity on $D^+$.
This supplies continuity through an actual finite
quotient; it does not complete an infinitely generated
discrete meridian module without a topology argument.

The source filling $Y^\dagger\to Y$ kills the sixteen
specified positive collision meridians. Its image
$K^+$ is their closed normal closure inside the
drilled source image, as in F.17--F.19.
Forgetting the negative new point kills its eight
collision meridians. The assignment (F.20.15) kills
the six positive collisions with the other old labels.
The remaining two positive meridians have images
$e_{a_+},e_{a_-}$ for the two actual transported
relative path classes in $Q$.
Their exponent is one, by the simple holomorphic
collision and annular calculation of F.17.

Conjugation by a source lift transports these basis
elements by exactly (F.20.1).
One can see this either by (F.20.11) or directly:
a based old meridian is carried to its conjugate
by the old track, while the module basepoint is
carried by the new track. Their two transports
give $q_g^{-1}\alpha_g(v)p_g$ in that order.
Different lifts of the same element of $L$ differ
by an element of $R_0$, whose conjugation acts
trivially on this additive image.
Consequently, writing $\mathcal O_\pm$ for the
two actual source orbits of $a_\pm$, one has the
literal finite-image equality
\[
 \chi_{Q,\ell}(K^+)
       =\operatorname{span}_{\mathbb F_\ell}
                  \{e_a:a\in\mathcal O_+\cup\mathcal O_-\}.
                                                               \tag{F.20.16}
\]
The labels $a_\pm$ have not been moved to the
identity by an unproved normalization.
The transitivity proved next makes their particular
values unnecessary for (F.20.3).

\emph{The finite simple quotients and their paired actions.}
The executed producer
\path{code/part_a/cross_divisor_finite_f20.py}
reads only the preceding exact geometric output.
It constructs $A_5$ as the sixty even permutations
of five letters, with its exact multiplication
and inverse tables. Conjugation by $S_5$ gives
all its $120$ automorphisms: every automorphism
permutes the five Sylow four-groups, and this
action is faithful because the resulting faithful
$A_5$-action determines each inner automorphism.
The conjugacy-class sizes are checked as
$1,12,12,15,20$; no proper nontrivial union
including the identity has order dividing $60$.
This also verifies the simplicity input below.

The program enumerates all $200,160$ generating
triples and their $1,668$ automorphism orbits.
These are exactly the epimorphism kernels.
It evaluates each of the eleven specified
$\alpha_{g_i}$ on each triple and finds precisely
$82$ invariant kernels. For each one it records
the unique descending automorphism $\sigma_i$,
the two word values $\rho(p_{g_i}),\rho(q_{g_i})$,
and the sixty-point permutation
\[
       v\longmapsto
       \rho(q_{g_i})^{-1}\sigma_i(v)\rho(p_{g_i}).
                                                               \tag{F.20.17}
\]
Every one of these $82$ actions has a single orbit.

For two distinct kernels, the map to $A_5^2$
is onto. Indeed a proper subdirect subgroup of
two copies of a nonabelian simple group is the
graph of an automorphism: the two projection
kernels are normal, and simplicity leaves
only the full product or that graph.
The graph would identify the two kernels,
contrary to their enumeration.
The program applies each source generator
simultaneously to both coordinates; all
$3,321$ such actions have one orbit of size
$3,600$. This is a test of the actual joint
action, not two separate marginal counts.

\emph{The six noncentral homology-cover refinements.}
Exactly six of the invariant triples have
$\rho(x),\rho(y),\rho(z),\rho(xyz)$ all nonidentity.
For such a triple put $G=A_5$, and consider the
connected based Cayley covering graph of the
three-petalled rose. Its chain groups and
boundary over $\mathbb F_2$ are
\[
 C_1=\mathbb F_2[G]^3,\quad C_0=\mathbb F_2[G],
 \qquad \partial e_{i,g}=e_{g\rho(x_i)}-e_g .
                                                               \tag{F.20.18}
\]
For a word $w$, let $v(w)$ be its actual oriented
edge chain and $a(w)=\rho(w)$.
Concatenation gives
$v(ww')=v(w)+a(w)v(w')$.
Every chain with boundary $e_g-e_1$ is the chain
of a word ending at $g$: subtract one path to
$g$ and realize the remaining graph homology
class by based loops.
Thus (F.20.2) is exactly the finite group
\[
 \{(v,g)\in\mathbb F_2[G]^3\rtimes G:
                       \partial v=e_g-e_1\},
 \quad
 (v,g)(w,h)=(v+gw,gh).
                                                               \tag{F.20.19}
\]
Its kernel has dimension $180-59=121$.
This description does not assert that this
particular extension splits.

Each actual $\alpha_{g_i}$ preserves $N_\rho$
and its mod-two homology kernel.
The lifted action is computed by the edge
chains of its three actual generator images.
If $J_i$ is that semilinear substitution matrix,
$\sigma_i$ its permutation of $G$, and
$(v_p,a_p),(v_q,a_q)$ the two word values,
its affine action is literally
\[
 (v,a)\longmapsto
 \left(a_q^{-1}\bigl(J_i v-v_q+\sigma_i(a)v_p\bigr),
                         a_q^{-1}\sigma_i(a)a_p\right).
                                                               \tag{F.20.20}
\]
The code checks these formulas against direct
word evaluation in $660$ exact comparisons.

The filling $Y\to
\mathbb P^1\setminus\{0,\infty,\pm\zeta\}$
kills exactly $g_1,g_3,g_4,g_5,g_6,g_8,g_9,g_{10}$.
Their base automorphisms and new words in
(F.20.12) are the identity, so they act by
right translations through their old words.
Conjugating such a translation by a general
source action gives the right translation
\[
                         h\longmapsto
                  p_{g_i}^{-1}\alpha_{g_i}(h)p_{g_i}.
                                                               \tag{F.20.21}
\]
The finite subgroup $H_\rho$ generated by these
translations and their indicated conjugates is
computed exactly as follows. Choose representatives
of its projection into $G$. Schreier products
give vectors in $\ker\partial$; their span,
closed under the projection's regular action,
is precisely its kernel. Add the images of
the representatives and this kernel under
(F.20.21), and repeat until stable.
Each step is finite linear algebra and finite
group multiplication. Stability under the
given injective maps on a finite subgroup
also gives stability under their inverses.
This computes the full subgroup generated
by the named source filling kernel.

In all six cases its projection is all of
$A_5$ and its vector kernel has dimension
$119$. Hence $[Q_\rho:H_\rho]=4$.
The source action descends to its right
cosets, because
$A_g(vh)=A_g(v)\,
 (p_g^{-1}\alpha_g(h)p_g)$.
The program reduces those cosets by the
specified representatives and row-reduced
kernel basis, and finds one orbit of size
four in every case.
Right translations by $H_\rho$ are themselves
source actions, so this proves transitivity
on the entire group $Q_\rho$.

These computations prove transitivity for
every finite quotient in the statement.
Equation (F.20.16) now proves (F.20.3) for
every prime $\ell$ and for the actual
source collision labels, without requiring
a lift of an arbitrary finite candidate.
The exact maps and finite tests enlarge the
verified winding comparison of F.19.
The remaining full coset equation in
$D^+/K^+$, and the original residual
$A_5$-centralizer, retain their stated scope.
\end{proof}

\begin{claim}{F.21}{Proved}{Full cross-divisor transitivity and normal generation inside the joint kernel}
Retain the actual groups, source maps, eight selected positive
cross-collision meridians, and base whiskers of F.17--F.19.
Write
\[
 G=P_{10}^{\dagger},\qquad
 \varpi:G\twoheadrightarrow\mathcal Q,
 \qquad D^+=\ker\varpi,\qquad
 H^\dagger=\Phi^\dagger(L^\dagger),\qquad
 K^+=H^\dagger\cap D^+ .
                                                               \tag{F.21.1}
\]
For each $u\in\mu _5\setminus\{\zeta\}$ and
$\epsilon\in\{1,-1\}$, let $\mathcal P_{u,\epsilon}$
be the full closed image in $\mathcal Q$ of its actual
cross divisor $v_u=\epsilon a$, based using the source
whisker selected for $\nu_{u,\epsilon}$ in F.19.
Then the source subgroup $\gamma(L)$ has a transitive
left action on the entire compact coset space
\[
                 \mathcal Q/\mathcal P_{u,\epsilon}.
                                                               \tag{F.21.2}
\]
This holds for each of the eight named divisors, with
the original label maps and full profinite completions.
The corresponding discrete coset action is also transitive.

The selected meridian has the following actual
centralizer-image containment:
\[
 \varpi\bigl(C_G(\nu_{u,\epsilon})\bigr)
       \supseteq\mathcal P_{u,\epsilon}.
                                                               \tag{F.21.3}
\]
Write $D^+_{\mathrm{nil}}$ for the inverse limit of all
continuous finite nilpotent quotients of $D^+$.
Consequently
\[
 \begin{gathered}
 \overline{\langle\!\langle K^+\rangle\!\rangle}_{D^+}=D^+,\\
 K^+\twoheadrightarrow
       D^+_{\mathrm{ab}}
       :=D^+/\overline{[D^+,D^+]},\qquad
 K^+\twoheadrightarrow D^+_{\mathrm{nil}} .
 \end{gathered}
                                                               \tag{F.21.4}
\]
In particular every continuous finite nilpotent quotient of
$D^+$ has the same image on $K^+$ and on $D^+$.
The quotient maps and the normal closure are the displayed
actual maps and subgroups.
The remaining equality $K^+=D^+$ and the source-membership
equation of F.19 remain unresolved.
\end{claim}
\passage{Section 30, folios 194-196, scans 383-387; Section 31, folios 198-199 verso, scans 391-394; Part F(b), the actual residual A5-centralizer and source-membership comparison}


\begin{proof}
\emph{Actual source basepoints on each filled cross divisor.}
Put
\[
 \begin{gathered}
 A(s)=\frac{\zeta-s}{\zeta+s},\qquad
 C_u(s)=\frac{u-s}{u+s},\qquad B=A^2,\qquad U=C_u^2,\qquad
 V=\epsilon A,\\
 k_u=\frac{u-\zeta}{u+\zeta},\qquad
 \mathscr A=\mathbb P^1\setminus\{0,\infty,1,-1\}.
 \end{gathered}                                                \tag{F.21.5}
\]
For the selected meridian $\nu_{u,\epsilon}$, let $d$ be
its collision point in $Y$. Thus $d^2=A_u$ for
$\epsilon=1$, and $d^2=B_u$ for $\epsilon=-1$, in the
notation of F.17.2. The sixteen collision points are
distinct and lie in $Y$ by F.17. At $d$ exactly this
cross collision occurs. The filled map
$\gamma:Y\to X$ of F.19.15 is defined at $d$; its
selected two-track projection is
\[
       s\longmapsto\bigl(A(s)^2,C_u(s)^2,\epsilon A(s)\bigr).
                                                               \tag{F.21.6}
\]
Its value at $d$ is on the diagonal $U=V$.

Keep the actual source whisker $\omega$ of that meridian,
extended to its endpoint $d$. It is a path in $Y$ from
the original $s_0$ to $d$, with its open terminal part
in $Y^\dagger$. The based source isomorphism is
$g\mapsto\omega g\omega^{-1}$ from $\pi_1(Y,d)$ to
$\pi_1(Y,s_0)$. On $X$ it is conjugation by the literal
path $\gamma(\omega)$. All forgetting projections
commute with these maps, because they forget coordinates
of this same configuration path. We prove transitivity
at the basepoint $d$ and then apply these specified
isomorphisms. The cross-divisor subgroup so obtained is
exactly the one based by the chosen whisker. No change
of the selected meridian or its conjugacy coordinate
is required.

\emph{The fourfold square-root cover and its pencil.}
The common two-coordinate moduli space is
\[
 \mathcal M_{0,5}
       =\{(B,U): B\notin\{0,1\},\ U\notin\{0,1,B\}\}.
\]
Both coordinates are finite. Its actual connected
fourfold covering and an intermediate double covering are
\[
 \begin{gathered}
 W=\{(A,C)\in\mathscr A^2:C\ne A,\ C\ne-A\},\\
 W_A=\{(A,U):A\in\mathscr A,\ U\notin\{0,1,A^2\}\},\\
 W\xrightarrow{\rho}W_A\xrightarrow{\beta}\mathcal M_{0,5},\\
 \rho(A,C)=(A,C^2),\qquad\beta(A,U)=(A^2,U).
 \end{gathered}                                               \tag{F.21.7}
\]
These maps are coverings because the squared coordinates
avoid $0$ and $\infty$; their four, respectively two,
local square-root choices are distinct. The complements
are connected: their projections to $A\in\mathscr A$
have connected punctured-plane fibers, and local
trivializations arise by moving the finitely many
excluded points in disjoint small disks.

There is an everywhere-defined map
\[
 \kappa:W\longrightarrow
        \mathscr B:=\mathbb P^1\setminus\{0,1,-1\},
       \qquad \kappa(A,C)=\frac{C-A}{1-AC}.
                                                               \tag{F.21.8}
\]
The numerator and denominator cannot vanish together on
$W$: that would give $A=C$ and $A^2=1$.
The values $0,1,-1$ would force, respectively,
$C=A$, $(C-1)(A+1)=0$, or $(C+1)(1-A)=0$.
All are excluded in (F.21.7).
Let
\[
 E=\{AC=1\}\subset W,
 \qquad W^\circ=W\setminus E,
 \qquad\mathscr B^\circ=\mathscr B\setminus\{\infty\}.
                                                               \tag{F.21.9}
\]
For finite $k\in\mathscr B^\circ$, the inverse coordinate
formula is
\[
 C=\frac{A+k}{1+kA},\qquad
 A\in\mathbb P^1\setminus
 \{0,\infty,1,-1,-k,-1/k,r_+(k),r_-(k)\},
 \quad kr^2+2r+k=0.
                                                               \tag{F.21.10}
\]
The eight displayed punctures are distinct for
$k\notin\{0,1,-1,\infty\}$.
Indeed the two roots are distinct by $1-k^2\ne0$;
neither is $0$ or $\infty$; the substitutions $r=1,-1$
give $2(k+1),2(k-1)$; and the substitutions
$r=-k,-1/k$ give $k(k^2-1),(k^2-1)/k$.
The remaining six punctures are pairwise distinct
under the same exclusions.

On a small simply connected neighborhood of any such
$k$, choose the two distinct analytic roots $r_+,r_-$.
All eight puncture paths remain in pairwise disjoint
coordinate disks. Within each disk, translate its
puncture using a bump function that is one near the
path and zero near the disk boundary. These isotopies,
with the inverse isotopies, depend continuously on $k$.
Their disjoint supports produce an actual local
trivialization of the complements in (F.21.10).
Consequently $W^\circ\to\mathscr B^\circ$ is a locally
trivial bundle with the indicated eight-punctured sphere
fiber. Its base and fiber are connected, and its
fundamental-group sequence is exact.

The old source has constant pencil value:
\[
 \kappa\bigl(A(s),C_u(s)\bigr)=k_u.
                                                               \tag{F.21.11}
\]
Here $k_u\notin\{0,1,-1,\infty\}$ since $u\ne\zeta$,
$u\ne-\zeta$, and both $u,\zeta$ are nonzero.
The coordinate $s=\zeta(1-A)/(1+A)$ identifies its
entire fiber with
\[
 Y_u=\mathbb P^1\setminus
       \{0,\infty,\zeta,-\zeta,u,-u,t_u,-t_u\},
       \qquad t_u^2=\zeta u .
                                                               \tag{F.21.12}
\]
For $u=\zeta^j$, take $t_u=\zeta^h$ with
$h=3(1+j)\pmod5$. The four exponent pairs are
\[
                 (j,h)\in\{(0,3),(2,4),(3,2),(4,0)\}.
\]
The inverse coordinate in (F.21.10) becomes exactly
$C_u(s)$. Its exclusions are $s=0,\infty,\pm\zeta$
from $A\in\mathscr A$, $s=\pm u$ from $C_u\in\mathscr A$,
and $s=\pm t_u$ from $C_u=-A$; $C_u=A$ adds only
$0,\infty$. Thus $Y\to Y_u$ fills exactly four of
the old punctures, and induces a surjection on
fundamental groups. Every loop on $Y_u$ can be moved
away from those four points by a small path homotopy.
The actual image of $\pi_1(Y,d)$ in $\pi_1(W^\circ)$
is therefore the image of this whole bundle fiber.

\emph{Filling the connected multiplicity-one divisor.}
The divisor $E$ is the actual connected smooth curve
\[
 A\longmapsto(A,A^{-1}),\qquad
 A\in\mathbb P^1\setminus\{0,\infty,1,-1,i,-i\}.
                                                               \tag{F.21.13}
\]
The extra exclusions $\pm i$ express $C\ne-A$.
At $E$ the local coordinate at the target infinity is
\[
 \frac1\kappa=\frac{1-AC}{C-A},\qquad
 \left.\frac{\partial(1/\kappa)}{\partial C}\right|_{C=1/A}
       =\frac{-A}{A^{-1}-A}\ne0.
                                                               \tag{F.21.14}
\]
Thus a positive transverse meridian of $E$ maps to
a single positive meridian of infinity in the base.

Set $J^\circ=\pi_1(W^\circ,(A(d),C_u(d)))$ and
$J=\pi_1(W,(A(d),C_u(d)))$. The inclusion is onto
on fundamental groups, and its kernel is the normal
closure of the meridians of $E$. To see both statements
directly, perturb a loop in the smooth manifold $W$
away from the real-codimension-two submanifold $E$.
Perturb a null-homotopy disk relative to its boundary
to meet $E$ transversely in finitely many interior
points. Delete small disks around those points;
the resulting punctured disk expresses the original
loop as a product of conjugates of the transverse
meridians, with their intersection signs.
All positive meridians form one conjugacy class,
because $E$ is connected and its complex normal line
has its fixed orientation.

Let $N^\circ$ be the fiber image in $J^\circ$ and let
$N$ be its image in $J$. The bundle sequence and
(F.21.14) give the following quotient map with its
specified inverse identification:
\[
 \begin{aligned}
 J/N
 &\cong
 (J^\circ/N^\circ)/
       \langle\!\langle\hbox{meridian of }\infty
                                \rangle\!\rangle\\
 &\cong
 \pi_1(\mathscr B^\circ,k_u)/
       \langle\!\langle\hbox{meridian of }\infty
                                \rangle\!\rangle
 \cong\pi_1(\mathscr B,k_u).
 \end{aligned}                                                \tag{F.21.15}
\]
The maps are induced by $\kappa$ and the two actual
fillings. Their kernels are the normal subgroups
just specified, so the inverse maps are the unique
factorizations of these same quotient homomorphisms.
In particular
\[
 N=\operatorname{im}\bigl(\pi_1(Y,d)\to J\bigr)
      =\ker\bigl(\kappa_*:J\twoheadrightarrow
                                  \pi_1(\mathscr B,k_u)\bigr).
                                                               \tag{F.21.16}
\]
This is an equality in the discrete fundamental group
of the named covering. No assertion about completing
this kernel sequence is needed below.

\emph{The new section and the onto pencil map.}
Base $W_A$ at $(A(d),C_u(d)^2)=(A(d),\epsilon A(d))$
and put $T=\pi_1(W_A)$.
The cover $W\to W_A$ embeds $J$ as a normal subgroup
of index two. The map $(A,U)\mapsto A$ has the actual
based section
\[
       \sigma_\epsilon:\mathscr A\longrightarrow W_A,
                    \qquad A\longmapsto(A,\epsilon A).
                                                               \tag{F.21.17}
\]
Its image $M=\sigma_{\epsilon,*}\pi_1(\mathscr A,A(d))$
maps isomorphically to $\pi_1(\mathscr A,A(d))$ under
the projection. The kernel of that projection is
the actual one-point-pushing group
\[
 F=\pi_1\bigl(\mathbb P^1\setminus
                   \{0,\infty,1,A(d)^2\},\epsilon A(d)\bigr)
       \cong F_3 .
                                                               \tag{F.21.18}
\]
Its injection and kernel description follow from the
configuration bundle, because $\pi_2(\mathscr A)=0$.

The pullback of $W\to W_A$ by this very section is
the connected double covering
\[
 \mathscr R=\mathbb P^1\setminus\{0,\infty,1,-1,i,-i\}
       \longrightarrow\mathscr A,
              \qquad r\longmapsto\epsilon r^2,
                                                               \tag{F.21.19}
\]
and its map into $W$ is $r\mapsto(\epsilon r^2,r)$.
Its basepoint is exactly $r_0=C_u(d)$; the collision
equation gives $A(d)=\epsilon r_0^2$.
Thus its fundamental-group image is exactly $M\cap J$.
It also proves that the projection of $M$ onto $T/J$
is onto. On this pullback the pencil map is literally
\[
 f_\epsilon(r)=\kappa(\epsilon r^2,r)
               =\frac{r}{r^2+\epsilon r+1}.
                                                               \tag{F.21.20}
\]
At $r_0$ it has value $k_u$, with no intervening
basepoint path.

This map induces a surjection
$\pi_1(\mathscr R,r_0)\to\pi_1(\mathscr B,k_u)$.
Here is a direct covering proof. Its derivative is
$(1-r^2)/(r^2+\epsilon r+1)^2$, and the exact identities
\[
 \begin{aligned}
 f_\epsilon(r)-\epsilon/3
     &=\frac{-\epsilon(r-\epsilon)^2}
                       {3(r^2+\epsilon r+1)},\\
 f_\epsilon(r)+\epsilon
     &=\frac{\epsilon(r+\epsilon)^2}{r^2+\epsilon r+1},\\
 f_\epsilon(r)-\epsilon
     &=\frac{-\epsilon(r^2+1)}{r^2+\epsilon r+1}
 \end{aligned}                                                \tag{F.21.21}
\]
identify every removed preimage and both branch values.
In particular
\[
 f_\epsilon:\mathscr R\longrightarrow
       \mathscr B\setminus\{\epsilon/3\}
                                                               \tag{F.21.22}
\]
is a connected unramified degree-two covering.
The two preimages of $0$ are $0,\infty$; the preimages
of $\epsilon$ are $i,-i$; the single double preimage
of $-\epsilon$ is $-\epsilon$; and the single double
preimage of $\epsilon/3$ is $\epsilon$.
The poles of $f_\epsilon$ map to the allowed target
infinity and are ordinary unramified points.
These are precisely the six punctures of $\mathscr R$.

Write $G_0=\pi_1(\mathscr B\setminus\{\epsilon/3\},k_u)$
and let $\chi:G_0\to C_2$ be the monodromy character
of (F.21.22). Its index-two kernel is the image of
$\pi_1(\mathscr R,r_0)$.
A based meridian $c$ of $\epsilon/3$ has $\chi(c)=1$
by its local double-root expression in (F.21.21).
The filling $q:G_0\twoheadrightarrow\pi_1(\mathscr B,k_u)$
kills $c$. For every $g\in G_0$, the element
$g c^{-\chi(g)}$ belongs to $\ker\chi$ and has the same
image under $q$ as $g$. Hence $q(\ker\chi)$ is the
entire target. This proves the asserted surjectivity,
using the actual index-two cover and actual filling.
The basepoint avoids $\epsilon/3$: its fiber in
(F.21.21) consists only of $r=\epsilon$, which is
excluded from $\mathscr R$.

Combining this surjectivity with (F.21.16) proves the
exact subgroup factorizations
\[
                       J=N(M\cap J),\qquad T=JM=NM .
                                                               \tag{F.21.23}
\]
For the first, given $j\in J$, choose $m\in M\cap J$
with $\kappa_*(m)=\kappa_*(j)$ and write
$j=(jm^{-1})m$ with $jm^{-1}\in N$.
For the second, $M\to T/J$ is onto by the connected
cover (F.21.19), and $(M\cap J)M=M$.

\emph{The exact orbit map.}
At the diagonal basepoint $d$, the two copies of
$\pi_1(\mathcal M_{0,5})$ have exactly the same
base configuration. Let $P_{5,\mathrm d}$ and
$P_{4,\mathrm d}$ denote these discrete groups with
their actual forgetting homomorphism. The mutually
inverse diagonal-coset maps are the following, where
$[b_o,b_n]$ denotes the indicated right diagonal coset:
\[
 \begin{gathered}
 \mathcal Q_{\mathrm d}
     =P_{5,\mathrm d}\mathbin{\times}_{P_{4,\mathrm d}}P_{5,\mathrm d},
 \qquad
 \mathcal Z_{\mathrm d}
     =\mathcal Q_{\mathrm d}/\operatorname{diag}(P_{5,\mathrm d}),\\
 \mathcal Z_{\mathrm d}\longrightarrow F,
        \qquad[b_o,b_n]\longmapsto b_nb_o^{-1},\\
 F\longrightarrow\mathcal Z_{\mathrm d},
        \qquad v\longmapsto[1,v].
 \end{gathered}                                               \tag{F.21.24}
\]
Their composites are the identity, by right
multiplication with $(b_o^{-1},b_o^{-1})$.
The group $F$ here is the same one-point kernel
(F.21.18); passage to the cover $W_A$ does not
alter its fiber or its inclusion.

For $n\in N$, write $a(n)\in\pi_1(\mathscr A)$ for
its actual $A$-coordinate and let $m(a(n))\in M$
be its image under (F.21.17).
The actual source pair has precisely the image
\[
                   n\longmapsto(n,m(a(n)))\qquad(n\in N).
                                                               \tag{F.21.25}
\]
Indeed its old projection has image $N$ by
(F.21.16), and its new projection is
$s\mapsto(A(s)^2,\epsilon A(s))$, the displayed
section applied to the same $A$-loop.
Any element killed by the old projection is also
killed by its $A$-projection and hence by the new
projection. This proves the equality of the source
image with (F.21.25), including its maps.

Left multiplication on (F.21.24) acts on $F$ as
\[
                      v\longmapsto m(a(n))\,v\,n^{-1}.
                                                               \tag{F.21.26}
\]
For every $v\in F$, factor $v^{-1}=nm$ using
(F.21.23). Projection to $\pi_1(\mathscr A)$ gives
$a(m)=a(n)^{-1}$. The projection is an isomorphism
on $M$, so $m=m(a(n))^{-1}$ and
$v=m(a(n))n^{-1}$. This is the image of $1$ under
(F.21.26). The source action is therefore transitive
on the entire discrete coset space.

The connected-fiber forgetting map from the F.19 space
$X$ to this selected two-track space carries the
selected divisor onto its diagonal. Its divisor
image is the full inverse image of the diagonal
group: lift a diagonal loop and correct the endpoint
in the connected configuration fiber. Its kernel
already comes from that fiber in the divisor.
This is the same exact comparison proved in F.20,
now with old label $u$ and new label $\epsilon a$.
The omitted labels are still three old points and
one new point; their distinctness and avoidance
conditions give those same connected fibers.
It proves the corresponding discrete version of
(F.21.2) for each of the four rows of (F.21.12)
and both choices of sign.

On full completions, F.19--F.20 identify this precise
divisor image with the inverse image of the full
diagonal subgroup, and its full coset space with
$\widehat F_3$. The actual configuration sections
and characteristic finite quotients of the finitely
generated free kernels used there apply with these
same omitted labels. Thus no completion of an
unproved intersection is being substituted here.
The orbit of a point under the compact group
$L=\widehat{\pi_1(Y)}$ is compact and closed in
$\widehat F_3$. By the discrete transitivity just
proved it contains the dense image of $F_3$.
It is therefore all of $\widehat F_3$.
Transport by $\gamma(\omega)$ gives exactly
(F.21.2) at the original F.19 basepoint.

\emph{The tubular centralizer lift.}
Write $D_i$ for one selected cross divisor in $X$
and $D_i^{\mathrm{reg}}$ for the complement in it
of its intersections with the other seven cross
divisors. This open submanifold is connected:
$D_i$ is connected by the configuration-fiber
description above, and paths can be perturbed
away from its finite union of complex divisors.
The same perturbation makes
$\pi_1(D_i^{\mathrm{reg}})\to\pi_1(D_i)$ onto.
The point $\gamma(d)$ belongs to
$D_i^{\mathrm{reg}}$ by the distinctness of the
sixteen source collision points.

Choose a Hermitian metric on the complex normal
line to $D_i^{\mathrm{reg}}$. A sufficiently small
positive continuous radius in that line gives
a circle bundle lying in the complement of all
eight divisors, namely the ten-point moduli space.
Such a radius is obtained from a locally finite
system of tubular coordinate charts by a positive
partition-of-unity minorant of their admissible
radii. Over any compact loop the resulting
normal circles and their homotopies stay in
these tubular charts.

The image of the fiber generator is central
in the fundamental group of this circle bundle.
Indeed transport around a base loop acts on
$\pi_1(S^1)=\mathbb Z$ by $+1$, because the normal
line is complex oriented. The circle-bundle
homotopy sequence is onto on the base fundamental
group, since the fiber is connected.
Base this construction by the chosen source
whisker and its small terminal normal circle.
Its fiber generator maps to exactly
$\nu_{u,\epsilon}$: F.17.7 gives the simple
holomorphic collision, and hence positive exponent
one. Filling the divisor takes each of these
circle-bundle loops to its base loop in $D_i$.
Consequently the discrete centralizer of this
very meridian maps onto the entire discrete
divisor image in $\pi_1(X)$.

Compose the circle-bundle homomorphism with the
full completion $G$. It extends continuously
from the full completion of the circle-bundle
group. Its compact image centralizes
$\nu_{u,\epsilon}$ and its image in $\mathcal Q$
contains the closure of the whole discrete
divisor image. That closure is exactly
$\mathcal P_{u,\epsilon}$ by the specified
completion comparison above. This proves
(F.21.3). Injectivity of the circle-bundle
completion is not used.

\emph{Normal generation in the actual kernel.}
Let $N_D$ be the closed normal closure of $K^+$
inside $D^+$. Fix one selected meridian $\nu_i$.
For $g\in G$, transitivity (F.21.2) factors its
image as
\[
                 \varpi(g)=\gamma(l)p,
          \qquad l\in L,\quad p\in\mathcal P_i.
\]
Choose $h\in H^\dagger$ over $l$ using the onto
map in F.19.10, and choose
$c\in C_G(\nu_i)$ with $\varpi(c)=p$ using
(F.21.3). Then $d=g c^{-1}h^{-1}\in D^+$ and
\[
                  g=d h c,
       \qquad g\nu_i g^{-1}=d(h\nu_i h^{-1})d^{-1}\in N_D.
                                                               \tag{F.21.27}
\]
Here $\nu_i\in K^+$ and $K^+\triangleleft H^\dagger$
by the exact source row F.19.10, so the middle
conjugate belongs to $K^+$.
F.19.3 states that the closed ambient normal
closure in $G$ of these eight $\nu_i$ is $D^+$.
Equation (F.21.27) puts every such conjugate in
$N_D$, proving $D^+=N_D$.

In the quotient $D^+_{\mathrm{ab}}$, conjugation
by $D^+$ acts trivially. The image of $K^+$ is
a closed subgroup, since $K^+$ is compact.
The normal-generation equality therefore says
that this image is the entire abelianization.
For any continuous finite nilpotent quotient
$q:D^+\twoheadrightarrow P$, put $S=q(K^+)$.
The proved equality says that $S$ normally
generates $P$. Every maximal proper subgroup
$M_0$ of a finite nilpotent group is normal:
choose the first term $Z_j(P)$ of its upper
central series not contained in $M_0$.
Then $Z_{j-1}(P)\subseteq M_0$ and an element
of $Z_j(P)\setminus M_0$ normalizes $M_0$,
because its commutators with $M_0$ lie in
$Z_{j-1}(P)$. The normalizer is strictly
larger than the maximal subgroup and hence
is all of $P$. A proper $S$ would be contained
in such a normal $M_0$; its normal closure
would remain there. Therefore $S=P$.
Finally the image of compact $K^+$ in
$D^+_{\mathrm{nil}}$ is closed and maps onto
every defining finite quotient. It is dense
by the inverse-limit topology and thus equals
$D^+_{\mathrm{nil}}$. This proves all the
stated quotient surjections.

The exact symbolic certificate
\path{code/part_a/cross_divisor_geometry_f21.py}
checks (F.21.11)--(F.21.12), the coordinate
inverse in (F.21.10), the divisor derivative
(F.21.14), the section pullbacks, and all branch
and puncture identities (F.21.20)--(F.21.21).
Its finite checks supplement the bundle,
fundamental-group, and compact-orbit arguments
written here; those arguments establish the
full assertions.
\end{proof}

\begin{claim}{F.22}{Proved}{The actual two-point kernel, its compact generators, and the remaining source map}
Retain the original F.19 base configuration. Let $p=a(s_0)$,
$n=-a(s_0)$, and let $S'$ be the plane with its seven old
finite labels removed. Filling the four labels
$U=\{v_u:u\ne\zeta\}$ gives the plane $S$ with the three
old finite labels $0,1,v_\zeta$ removed. The actual
two-point groups are
\[
 J=\widehat{\pi_1(\operatorname{Conf}_2(S'),(p,n))},\qquad
 J_4=\widehat{\pi_1(\operatorname{Conf}_2(S),(p,n))},\qquad
 D^+=\ker(J\to J_4).
                                                               \tag{F.22.1}
\]
There are compatible based point-pushing sections of the
two positive-point projections, described below, and
geometric free coordinates in which
\[
 \begin{gathered}
 V_+=\widehat F(g_1,g_2,g_3,u_1,u_2,u_3,u_4),\qquad
 Q_+=\widehat F(g_1,g_2,g_3),\\
 V=\widehat F(X_1,X_2,X_3,U_1,U_2,U_3,U_4,t),\qquad
 P=\widehat F(X_1,X_2,X_3,t),\\
 R_+=\ker(V_+\to Q_+),\qquad R_1=\ker(V\to P).
 \end{gathered}                                                \tag{F.22.2}
\]
The displayed fillings kill precisely $u_j$ and $U_j$,
respectively. With the actual section $\sigma:V_+\to J$
and fiber inclusion $\iota:V\to J$, there are mutually
inverse continuous maps
\[
 \begin{aligned}
 R_1\rtimes R_+&\xrightarrow{\ \Theta\ }D^+,
       &(b,g)&\longmapsto\iota(b)\sigma(g),\\
 D^+&\longrightarrow R_1\rtimes R_+,
       &d&\longmapsto
       \bigl(\iota^{-1}(d\sigma(p_+d)^{-1}),p_+d\bigr).
 \end{aligned}                                                \tag{F.22.3}
\]
Here $p_+$ forgets the negative point; the action in this
semidirect product is the actual conjugation by $\sigma$.
Its complete free-word recipe is (F.22.11)--(F.22.12).

For a profinite space $T$, write $\mathfrak F(T)$ for the
free profinite group on continuous maps from $T$, as
constructed in the proof. This definition retains the
topology on the generator space. Put $I=\{1,2,3,4\}$.
Then the exact compact-generator identifications are
\[
 \begin{aligned}
 \mathfrak F(Q_+\times I)&\xrightarrow{\ \sim\ }R_+,
       &(q,j)&\longmapsto s_+(q)u_js_+(q)^{-1},\\
 \mathfrak F(P\times I)&\xrightarrow{\ \sim\ }R_1,
       &(v,j)&\longmapsto s_1(v)U_js_1(v)^{-1},
 \end{aligned}                                                \tag{F.22.4}
\]
where $s_+$ and $s_1$ retain the displayed free generators.

The other one-point forgetting is the actual
$p_-:J\to V_-$, with
\[
 \begin{gathered}
 V_-=\widehat F(X_1,X_2,X_3,U_1,U_2,U_3,U_4),\qquad
 Q_-=\widehat F(X_1,X_2,X_3),\\
 R_-=\ker(V_-\to Q_-),\qquad
 W=D^+\cap\ker p_+\cap\ker p_- .
 \end{gathered}                                                \tag{F.22.5}
\]
One has the exact sequence
\[
 1\longrightarrow W\longrightarrow D^+
        \xrightarrow{(p_+,p_-)}R_+\times R_-\longrightarrow1.
                                                               \tag{F.22.6}
\]
Inside the actual fiber $V$, let $s_-$ retain $Q_-$ and put
$\eta(q,e)=s_-(q)e s_-(q)^{-1}$ for
$e\in\{U_1,U_2,U_3,U_4,t\}$. Then
\[
 W=
 \overline{\langle\!\langle
       [\eta(q,U_j),\eta(q',t)]
       :q,q'\in Q_-,\ j\in I
                         \rangle\!\rangle}_{V} .
                                                               \tag{F.22.7}
\]
The literal element $[U_1,t]$ belongs to $W$ and
generates a full closed copy of $\widehat{\mathbb Z}$.

Finally let $\Delta$ be the sixteen actual source
collisions, with the F.19 source meridians and section
$j:L\to H^\dagger$. There is a specified continuous map
\[
 \Lambda:\mathfrak F(L\times\Delta)\longrightarrow D^+,
       \qquad(l,d)\longmapsto j(l)\nu_dj(l)^{-1},
       \qquad\operatorname{im}\Lambda=K^+ .
                                                               \tag{F.22.8}
\]
Its coordinates in (F.22.3) are given by the actual
two projection formulas there, applied to the displayed
element. Thus the remaining equality $K^+=D^+$ is
exactly surjectivity of this named continuous homomorphism
onto the displayed semidirect product. This statement
does not settle that surjectivity or the original
residual source-membership equation.
Write $\Delta_+$ and $\Delta_-$ for the eight collisions
with the positive and negative new point, respectively,
and put $\mathcal K_\pm=p_\pm(K^+)$. The exact joint
source image is their full direct product:
\[
 1\longrightarrow K^+\cap W\longrightarrow K^+
       \xrightarrow{(p_+,p_-)}\mathcal K_+\times\mathcal K_-
       \longrightarrow1 .
                                                              \tag{F.22.8a}
\]
Each marginal is the actual image of the corresponding
eight-collision source filling kernel, as specified
in (F.22.21). Their equality with $R_+$ and $R_-$
remains unresolved.
\end{claim}
\passage{Section 30, folios 194-196, scans 383-387; Section 31, folios 198-199 verso, scans 391-394; Part F(b), the original residual A5-centralizer and the exact source-membership problem}


\begin{proof}
\emph{Compatible sections at the original basepoint.}
All seven old finite points, $p$, and $n$ are distinct
by F.17. Choose an embedded finite tree in the plane
joining the seven old points and $p$ while avoiding
$n$. A sufficiently small regular neighborhood is a
closed disk $B$ containing those eight points in its
interior and excluding $n$. The inclusions
\[
 B^\circ\setminus\{\hbox{seven old points}\}\longrightarrow S',
 \qquad
 B^\circ\setminus\{\hbox{three retained old points}\}
                    \longrightarrow S
                                                               \tag{F.22.9}
\]
induce isomorphisms on fundamental groups based at
the original point $p$. Indeed a collar of the disk
boundary and the outside annulus retract onto that
boundary, so adjoining the outside region changes
none of the punctured-disk fundamental-group generators
or relations.

On each punctured disk the map $z\mapsto(z,n)$ is an
actual based map to its ordered two-point configuration
space. Compose its fundamental-group map with the
inverse of the relevant isomorphism (F.22.9).
This gives sections $\sigma_{\mathrm d}$ for $S'$ and
$\sigma_{4,\mathrm d}$ for $S$.
The two disk inclusions and the two maps $z\mapsto(z,n)$
commute with filling the four old punctures. Their
inverse group isomorphisms therefore commute too.
Hence these very sections form a commuting square
with the actual filling $J_{\mathrm d}\to J_{4,\mathrm d}$.
All basepoints here are $(p,n)$; the comparison paths
to the F.19 two-point coordinates are the constant
paths at that configuration.

The Fadell--Neuwirth sequences have kernels
$\pi_1(S'\setminus\{p\},n)=F_8$ and
$\pi_1(S\setminus\{p\},n)=F_4$.
Their injections follow from $\pi_2(S')=\pi_2(S)=0$.
The full-topology and completed-kernel comparisons
are exactly those established in F.19 for these
two-point configuration groups. The actual discrete
sections extend to full completions, their section
identities persist, and their filling square still
commutes. Thus
\[
 J=V\rtimes_\sigma V_+,\qquad
 J_4=P\rtimes_{\sigma_4}Q_+,
 \qquad
 (b,g)\longmapsto(\overline b,\overline g)
                                                               \tag{F.22.10}
\]
is the actual filling map in these coordinates.
The equation $\overline g=1$ puts $g$ in $R_+$,
and then $\overline b=1$ puts $b$ in $R_1$.
Multiplication and the section inverse give exactly
the two maps in (F.22.3).

\emph{Free coordinates and the retained Artin action.}
Choose a geometric chart on $B$ in which the old
retained points occupy positions $1,2,3$, the four
other old points occupy positions $4,5,6,7$, and $p$
occupies position $8$. The exterior point $n$ is
the basepoint for puncture meridians. Such a chart
is an orientation-preserving homeomorphism of the
disk with the stated finite labelled set: move
the points along polygonal arcs avoiding the other
points, using successive small disk isotopies.
The induced fundamental-group isomorphisms have
the inverse chart as their literal inverse.
Both filling maps use this same chart and its
restriction to the retained labelled set.
They therefore commute with the actual coordinate
maps, including the basepoints.

Write $X_4=U_1,\ldots,X_7=U_4,X_8=t$.
Use the literal clockwise exchanges and the Artin
substitutions calibrated in F.20.7--9 on
$\widehat F(X_1,\ldots,X_8)$.
For the positive point moving once clockwise around old
position $i$ along its lower-half-plane whisker, the
based geometric free generator $g_i$ has the braid
\[
 \beta_i=\sigma_7\cdots\sigma_{i+1}\sigma_i^2
                  \sigma_{i+1}^{-1}\cdots\sigma_7^{-1},
          \qquad1\leq i\leq7,
 \qquad u_j=g_{j+3}.
                                                               \tag{F.22.11}
\]
The descending word carries the last strand below
the intervening points, its square is the clockwise
lasso, and the inverse word returns the same whisker.
Deleting that strand gives the trivial old braid by
contracting the old motion in its projection-order
chambers. These seven lower-whiskered lassos form the
actual geometric kernel basis. The section keeps the
negative point $n$ fixed outside $B$, so the exact
fibre-coordinate comparison below makes its conjugation
on the negative-point fiber
\[
 \alpha_{g_i}=\phi_{\beta_i},
 \qquad
 \phi_{\sigma_k}(X_k)=X_kX_{k+1}X_k^{-1},
 \quad\phi_{\sigma_k}(X_{k+1})=X_k,
 \quad\phi_{\sigma_k}(X_j)=X_j\ (j\ne k,k+1).
                                                               \tag{F.22.12}
\]
The clockwise braid word is composed in the same
chronological and substitution order as F.20, and its
inverse word gives its inverse map.

Here is the precise point-pushing identification
underlying that conjugation formula. For a pure
braid on the eight punctures and the exterior ninth
point, let $d$ be the ninth Artin meridian and write
$\phi_b(d)=W_bdW_b^{-1}$. Erasing $d$ gives $w_b$;
put $u_b=w_b^{-1}$. Exactly as in F.20.8, substitution
gives $u_{bc}=u_b\phi_b(u_c)$ after erasing the ninth
strand. The section on the first eight strands has
$u=1$. Write $\gamma_i$ for the eight lower-whiskered
clockwise point-pushing lassos of the ninth point,
given by the same braid pattern with $\sigma_8$ as
its first exchange. Their values of $u$ are
\[
       \tau_i=X_8^{-1}\cdots X_{i+1}^{-1}
                    X_i^{-1}X_{i+1}\cdots X_8,
                \qquad1\leq i\leq8.
                                                               \tag{F.22.13}
\]
This triangular substitution is its own inverse:
starting with $i=8$ and substituting downward,
each adjacent inverse pair cancels and the image
of $\tau_i$ is $X_i$. The actual forward coordinate
map is $\gamma_i\mapsto\tau_i(X)$; its inverse is
$X_i\mapsto\tau_i(\gamma)$. Both composites are
the identity on the named free bases. They therefore
give the actual free-group isomorphism and its full
profinite inverse between the fiber and these Artin
coordinates. The crossed multiplication formula now
makes conjugation by the actual section exactly
(F.22.12). Erasing any of the four old labels,
or erasing $t$, commutes with (F.22.13); the
remaining triangular formula has exactly the
remaining labels in their retained order.
This proves the coordinate form of both fillings.

The plane braid groups in this calculation map
to the original normalized moduli groups through
the explicit coordinate isomorphism
\[
 (z_0,z_1,z_2,\ldots)
   \longmapsto
 \left(z_0,z_1-z_0,
           \left(\frac{z_j-z_0}{z_1-z_0}\right)_{j\geq2}\right).
                                                               \tag{F.22.14}
\]
Its inverse multiplies the last coordinates by
$z_1-z_0$ and adds $z_0$. The first two old labels
are fixed on every loop in $J$, so their two
affine coordinates are constant. This is the
actual comparison with the normalized F.19
configuration groups.

For completeness, the action extends continuously
from $F_7$ to $V_+=\widehat F_7$ on $V=\widehat F_8$.
Every open subgroup of this finitely generated
profinite $V$ contains a characteristic open
subgroup: intersect its finitely many subgroups
of at most the given index. The action on each
resulting finite characteristic quotient has
finite image and extends through $V_+$.
The compatible inverse automorphisms give the
continuous inverse action. This is also the
action already supplied by the actual section
in (F.22.10).

\emph{A free profinite group on a compact generator space.}
For a profinite space $T$, begin with the abstract
free group on its underlying set. Take its inverse
limit over finite quotients for which the map
from $T$ to that finite quotient is continuous;
denote the result by $\mathfrak F(T)$.
Products of finitely many such quotient maps again
have continuous generator maps, so these quotients
form a directed system. The canonical map
$e:T\to\mathfrak F(T)$ is continuous and its
image generates a dense subgroup. Every continuous
map from $T$ to any profinite group $H$ extends
uniquely to a continuous homomorphism
$\mathfrak F(T)\to H$: compose with every finite
quotient of $H$, extend on the abstract free
generators, and take the compatible inverse limit.
This construction proves the precise universal
property used here.

Let $Q$ be one of the finitely generated free
profinite groups in (F.22.2), and let $I$ be the
finite set of the additional generators.
Left multiplication on $Q\times I$ induces
automorphisms of $\mathfrak F(Q\times I)$.
This action is jointly continuous. Indeed for
any continuous map $Q\times I\to A$ to a finite
group, compactness gives an open normal subgroup
$N\triangleleft Q$ through which all its finitely
many coordinate functions factor. Its translates
therefore form a finite set indexed by $Q/N$.
Intersect the kernels of the corresponding
finite quotient homomorphisms of $\mathfrak F$.
The intersection is open and $Q$-invariant, and
the action on this quotient factors through a
finite quotient of $Q$. These invariant open
subgroups refine every finite quotient kernel.
They prove joint continuity and give the actual
profinite semidirect product
$\mathfrak F(Q\times I)\rtimes Q$.

Let $\widehat F(Q;I)$ denote the free profinite
group obtained by adjoining these finitely many
extra generators $v_j$ to the fixed free basis
of $Q$. Its retraction $\rho$ kills the $v_j$,
and its section $s$ retains the basis of $Q$.
There are inverse continuous homomorphisms
\[
 \begin{aligned}
 \widehat F(Q;I)&\longrightarrow
                     \mathfrak F(Q\times I)\rtimes Q,\\
 x&\longmapsto(1,x)\quad(x\text{ in the retained basis}),
       \qquad v_j\longmapsto(e(1,j),1),\\
 \mathfrak F(Q\times I)\rtimes Q
       &\longrightarrow\widehat F(Q;I),\\
 e(q,j)&\longmapsto s(q)v_js(q)^{-1},
       \qquad q\longmapsto s(q).
 \end{aligned}                                                \tag{F.22.15}
\]
The second generator map is continuous and
$Q$-equivariant, so the universal property and
the semidirect-product law give its homomorphism.
Both composites fix every retained generator,
every $v_j$, and every $e(q,j)$ by the displayed
conjugation formula. These sets generate densely,
so both composites are the identity. The
projections onto $Q$ are precisely $\rho$ and
the semidirect-product projection. Their kernels
are therefore identified by these inverse maps.
Apply this construction first to $(Q_+,u_j)$
and then to $(P,U_j)$ to prove (F.22.4).
No full completion of an infinitely generated
discrete kernel has been substituted in this
argument.

The action in (F.22.3) can now be evaluated
on every compact generator, not only on a
finite selection. The element $e(q,j)$ of
$R_+$ has image $s_+(q)u_js_+(q)^{-1}$ in
$V_+$. Its action is
$\alpha_{s_+(q)}\alpha_{u_j}\alpha_{s_+(q)}^{-1}$.
For a compact generator $e(v,k)$ of $R_1$,
apply that automorphism to
$s_1(v)U_ks_1(v)^{-1}$ in $V$, then use the
inverse of (F.22.15). The result lies in
$R_1$ because the induced action of $R_+$
on $P$ is the identity. The formula is
continuous in both $q$ and $v$ by the
actual continuous group operations and the
proved action continuity.

\emph{The interaction kernel retained by both one-point forgettings.}
On the fiber $V$, the negative-point projection
$p_-$ fills the positive puncture, so its map
$h:V\to V_-$ kills $t$ and retains every other
letter. On the section $\sigma(V_+)$ it is
the identity element: that section keeps the
negative point fixed. Hence
\[
        (p_+,p_-)(\Theta(b,g))=(g,h(b)).
                                                               \tag{F.22.16}
\]
The restriction $h:R_1\to R_-$ is onto.
For an element of $R_-$, retain its seven
letters in $V$; its old-puncture filling is
the identity, so it belongs to $R_1$ and
maps back to that element under $h$.
Together with the arbitrary $g\in R_+$,
this proves surjectivity in (F.22.6), with
kernel $W=R_1\cap\ker h$ inside $V$.

Apply (F.22.15) with retained group $Q_-$
and the five extra letters $U_1,\ldots,U_4,t$.
The kernel of $V\to Q_-$ is
\[
 \mathcal R=
 \mathfrak F\bigl((Q_-\times I)\sqcup(Q_-\times\{t\})\bigr)
       =\mathcal A\amalg\mathcal B,
                                                               \tag{F.22.17}
\]
where $\mathcal A=\mathfrak F(Q_-\times I)$,
$\mathcal B=\mathfrak F(Q_-\times\{t\})$,
and the coproduct equality follows directly
from the continuous-map universal property
on the two compact open-and-closed pieces.
The two fillings restrict to the retractions
$\mathcal R\to\mathcal B$ and
$\mathcal R\to\mathcal A$. Consequently
\[
             W=\ker(\mathcal A\amalg\mathcal B
                            \longrightarrow\mathcal A\times\mathcal B).
                                                               \tag{F.22.18}
\]
Quotienting $\mathcal R$ by the closed normal
closure of all the displayed generator
commutators makes its two factor images
commute. The multiplication map from
$\mathcal A\times\mathcal B$ to this quotient
is then a homomorphism. Its composite with
the two retractions is the identity, and
the reverse composite fixes all generators.
These are inverse homomorphisms. Thus their
kernel is exactly the stated normal closure.
It is invariant under $Q_-$, whose action
sends $(q,q')$ to $(aq,aq')$.
Its normal closure in $V=\mathcal R\rtimes Q_-$
is therefore the same subgroup. This proves
(F.22.7), with its actual ambient group.

For every integer $m\geq2$, send $U_1$ to
$I+E_{12}$, send $t$ to $I+E_{23}$, and send
the other free generators of $V$ to the
identity in the finite group
$\operatorname{UT}_3(\mathbb Z/m)$.
These assignments define a continuous
homomorphism. With $[a,b]=aba^{-1}b^{-1}$,
\[
       [U_1,t]\longmapsto I+E_{13},\qquad
       [U_1,t]^z\longmapsto I+(z\bmod m)E_{13}
                \quad(z\in\widehat{\mathbb Z}).
                                                               \tag{F.22.19}
\]
The commutator is killed by both fillings,
so it belongs to $W$. The displayed finite
images detect every nonzero
$z\in\widehat{\mathbb Z}$ through some modulus.
The exponent map to its closed cyclic
subgroup is therefore injective as well
as onto, proving the asserted full
procyclic subgroup.

\emph{The actual compact-generator source image.}
The F.19 source section retains its eleven
old free meridians in the rank-twenty-seven
source; the sixteen additional free meridians
are $m_d$, $d\in\Delta$. Apply (F.22.15)
to that very based free system. It identifies
\[
 \ker(q:L^\dagger\to L)
       \cong\mathfrak F(L\times\Delta),
       \qquad(l,d)\longmapsto i(l)m_di(l)^{-1}.
                                                               \tag{F.22.20}
\]
Composition with the actual map
$\Phi^\dagger:L^\dagger\to G$ gives exactly
(F.22.8), because $\Phi^\dagger i=j$ and
$\Phi^\dagger(m_d)=\nu_d$ with the retained
source whisker. F.19.10 identifies the image
of this source kernel with $K^+$; alternatively
its image is precisely the kernel of the
induced onto map $H^\dagger\to L$, since
the commuting source filling has that
kernel. The map in (F.22.8) is continuous
and its image is closed, since its domain
is compact.

\emph{The two marginal source images and their joint image.}
Put $Y^\pm=Y\setminus\Delta_\pm$ and
$L_\pm=\widehat{\pi_1(Y^\pm,s_0)}\cong\widehat F_{19}$.
The source fillings $q_\pm:L_\pm\to L$ retain
the eleven old basis elements and kill their
eight added meridians.
Let $F_+:G\to P_9^+$ forget the negative point,
and let $F_-:G\to P_9^-$ forget the positive point.
These are the actual full nine-point forgetting
maps; their restrictions to $J$ are $p_+$ and
$p_-$ with the displayed one-point kernel targets.
The rational source configuration after forgetting
the negative point is defined on all of $Y^+$;
the only extra singularities of that map are
the eight positive collisions. The corresponding
statement holds with the signs interchanged.
Thus the actual source maps factor as
\[
 \begin{gathered}
 F_\pm\Phi^\dagger
       =\Phi_\pm r_\pm,
       \qquad r_\pm:L^\dagger\twoheadrightarrow L_\pm,\\
 \mathcal K_\pm
       =\Phi_\pm(\ker q_\pm)
       =\operatorname{im}
        \bigl(\mathfrak F(L\times\Delta_\pm)
                                \longrightarrow R_\pm\bigr).
 \end{gathered}                                                \tag{F.22.21}
\]
Here $r_+$ fills the negative collision points,
and $r_-$ fills the positive ones.
These equalities retain their based source maps:
the signed free basis embeds $L_\pm$ back into
$L^\dagger$ by its named old and signed meridians,
so $r_\pm(\ker q)=\ker q_\pm$.
The last map sends $(l,d)$ to
$F_\pm(j(l)\nu_dj(l)^{-1})$.

For $d\in\Delta_+$, forgetting the positive
point removes the collision. The remaining
configuration map extends across $d$, so the
based source meridian has $F_-(\nu_d)=1$.
For $d\in\Delta_-$ one similarly has
$F_+(\nu_d)=1$.
For every $l\in L$ these identities give
\[
 \begin{aligned}
 F_-(j(l)\nu_dj(l)^{-1})&=1\quad(d\in\Delta_+),\\
 F_+(j(l)\nu_dj(l)^{-1})&=1\quad(d\in\Delta_-).
 \end{aligned}                                                \tag{F.22.22}
\]
The maps here are defined on all of $G$;
the calculation does not apply $p_\pm$ to
an element outside $J$.

Let $K_+^{\mathrm{src}}$ and $K_-^{\mathrm{src}}$
be the compact images of
$\mathfrak F(L\times\Delta_+)$ and
$\mathfrak F(L\times\Delta_-)$ in $K^+$
under (F.22.8). Their joint images are
$\mathcal K_+\times\{1\}$ and
$\{1\}\times\mathcal K_-$, respectively,
by (F.22.21)--(F.22.22).
Given a pair in $\mathcal K_+\times\mathcal K_-$,
lift its first coordinate to $K_+^{\mathrm{src}}$
and its second to $K_-^{\mathrm{src}}$.
Their product in $K^+$ maps to the given pair.
This proves (F.22.8a). Its kernel is
$K^+\cap W$ by (F.22.6), giving the exact
commuting diagram of actual inclusions
\[
 \begin{array}{ccccccccc}
 1&\longrightarrow&K^+\cap W&\longrightarrow&K^+
     &\longrightarrow&\mathcal K_+\times\mathcal K_-
     &\longrightarrow&1\\
 &&\downarrow&&\downarrow&&\downarrow\\
 1&\longrightarrow&W&\longrightarrow&D^+
     &\longrightarrow&R_+\times R_-
     &\longrightarrow&1 .
 \end{array}                                                  \tag{F.22.23}
\]
The first exact row concerns the proved marginal
images, with no enlargement of either marginal.

The point-pushing section $\sigma$ and the
source section $j$ have both been retained
as their actual maps. To find the coordinates
of any generator in (F.22.8), set
$z=j(l)\nu_dj(l)^{-1}$, compute the literal
one-point forgetting $g=p_+(z)$, and take
$b=\iota^{-1}(z\sigma(g)^{-1})$.
This is exactly $\Theta^{-1}(z)$, with no
identification of the two sections.
The first coordinate belongs to $R_1$ and
the second to $R_+$ by the proved inverse
maps. These formulas identify the remaining
source-membership problem in the displayed
kernel presentation.

The exact producer
\path{code/part_a/joint_kernel_presentation_f22.py}
checks all seven Artin action words and
their inverses, their pure conjugators,
their boundary product, and their two
actual filling squares. It also records
the deleted-label conjugators after the
old filling, checks that $[U_1,t]$ is
killed by both fillings, and checks the
central power formula in thirty finite
Heisenberg groups. The algebraic formula
(F.22.19), proved for every modulus,
supplies the full profinite conclusion.
The compact-generator construction and
all semidirect-product inverse maps are
proved above independently of those
finite checks.
\end{proof}

\begin{claim}{F.23}{Proved}{The based source-sign symmetry, complete source braid map, and two exact finite marginal images}
Retain the original groups, source map, basepoints, and two
one-point forgettings of F.19--F.22. In particular put
\[
 \mathcal K_\pm=p_\pm(K^+)\leq R_\pm\leq V_\pm,
       \qquad V_\pm\cong\widehat F_7,
                                                               \tag{F.23.1}
\]
with the actual geometric positive coordinates and Artin negative
coordinates of F.22. The source involution
$h(s)=\zeta^2/s$, together with the explicit path $\lambda$ below,
gives an actual based isomorphism of inclusions
\[
 \begin{array}{ccc}
 \mathcal K_-&\lhook\joinrel\longrightarrow&R_-\\
 {\scriptstyle\mathcal B_+}\downarrow\phantom{\scriptstyle\sim}
       &&\phantom{\scriptstyle\sim}\downarrow{\scriptstyle\mathcal B_+}\\
 \mathcal K_+&\lhook\joinrel\longrightarrow&R_+ .
 \end{array}                                                  \tag{F.23.2}
\]
The same map extends to $V_-\xrightarrow{\sim}V_+$.
Its formula, inverse, and all filling squares are specified below.

There is a complete exact braid description of the actual homomorphism
\[
 \Phi^\dagger:L^\dagger=\widehat{\pi_1(Y^\dagger,s_0)}
                       =\widehat F_{27}\longrightarrow G.
                                                               \tag{F.23.3}
\]
It uses the twenty-seven specified source lassos, a rational
projection perturbation, the original endpoint whiskers, and the
original F.20 path $\eta$. The exact arithmetic certificates and
twenty-seven full braid words are retained in the accompanying
files. The coordinate comparison is the clockwise-exchange and
inverse triangular comparison of F.20.7--9 and F.22.11--13.

For each sign, for each of the two groups $H=S_3,A_4$, and for
every continuous homomorphism from the indicated full point-pushing
group, the following equality holds:
\[
 \rho(\mathcal K_\pm)=\rho(R_\pm)
       \qquad\bigl(\rho\in\operatorname{Hom}_{\rm cont}(V_\pm,H),
                              \ H\in\{S_3,A_4\}\bigr).
                                                               \tag{F.23.4}
\]
This is the exact finite scope of the exhaustive calculations.
The original maps $\mathcal K_\pm\hookrightarrow R_\pm$ and
$K^+\hookrightarrow D^+$ remain the specified inclusions;
their surjectivity remains unresolved. The joint image and its
kernel remain exactly those of F.22.8a.
\end{claim}
\passage{Part F root-adjoining intersection question; Section 30, folios 194-196, scans 383-387; Section 31, folios 198-199 verso, scans 391-394; original F7/F19 residual subgroup-membership reduction}


\begin{proof}
\emph{The literal source and the sign-changing maps.}
Write
\[
 \begin{gathered}
 a(s)=\frac{\zeta-s}{\zeta+s},\qquad
 v_u(s)=\left(\frac{u-s}{u+s}\right)^2,\qquad
 u\in\{1,\zeta,\zeta^2,\zeta^3,\zeta^4\},\\
 A_u=u(2\zeta-u),\qquad
 B_u=\frac{\zeta u^2}{2u-\zeta}
       \quad(u\ne\zeta),\\
 \Delta_+=\{s:s^2=A_u,\ u\ne\zeta\},\qquad
 \Delta_-=\{s:s^2=B_u,\ u\ne\zeta\},\\
 Y^\dagger=Y\setminus(\Delta_+\cup\Delta_-).
 \end{gathered}                                                \tag{F.23.5}
\]
The complete finite configuration is
\[
 \phi^\dagger(s)=
 (0,1,v_1(s),v_\zeta(s),v_{\zeta^2}(s),
                     v_{\zeta^3}(s),v_{\zeta^4}(s),a(s),-a(s)),
                                                               \tag{F.23.6}
\]
with the original additional label at infinity. Its normalization
and its induced map to $G$ are those of F.19 and F.22.14.
All punctures and their simple collision multiplicities are
the actual ones established in F.17.

For $h(s)=\zeta^2/s$ and $u'=\zeta^2/u$, direct substitution gives
\[
 \begin{gathered}
 h^2(s)=s,\qquad a(h(s))=-a(s),\qquad
 v_u(h(s))=v_{u'}(s),\\
 \frac{\zeta^4}{A_u}=B_{u'},\qquad
 \frac{\zeta^4}{B_u}=A_{u'}.
 \end{gathered}                                                \tag{F.23.7}
\]
Thus $h$ is an actual homeomorphism of $Y^\dagger$ and $Y$,
and it interchanges the two signed sources
$Y_\pm=Y\setminus\Delta_\pm$.
Let $\Pi$ permute the old labels by $u\mapsto u'$ and exchange
the two new labels. It fixes $0,1,\infty,v_\zeta$.
The complete source equality is
\[
                  \phi^\dagger h=\Pi\phi^\dagger .            \tag{F.23.8}
\]
This equality commutes with all the original fillings and with
the two one-point forgettings, with their signs exchanged.

The based path is explicit. Put $s_1=h(s_0)$ and use the polygon
\[
                 \lambda=(s_0, s_0+6i, s_1+6i, s_1).
                                                               \tag{F.23.9}
\]
It lies in $Y^\dagger$. The exact sign-symmetry producer checks
that each of its two vertical lines has real coordinate different
from every old puncture. For either real coordinate $x$ and every
collision square $V=A_u,B_u$, it checks $x\ne0$ and
\[
                  4x^4-4\operatorname{Re}(V)x^2
                            -(\operatorname{Im}V)^2\ne0.     \tag{F.23.10}
\]
Indeed a root $x+iy$ of $z^2=V$ would have
$y=\operatorname{Im}(V)/(2x)$ and would make this expression zero.
The middle segment has imaginary part greater than two, whereas
every old and collision puncture has absolute value less than two.
These are exact inequalities in the retained real number field.

With left-to-right path products, the actual based source
automorphism and its inverse are
\[
 \begin{aligned}
 \mathcal A_{\rm src}(g)&=\lambda h(g)\lambda^{-1},\\
 \mathcal A_{\rm src}^{-1}(g)&=h(\lambda)^{-1}h(g)h(\lambda).
 \end{aligned}                                                \tag{F.23.11}
\]
Both composites cancel to $g$. The square
$\mathcal A_{\rm src}^2$ is conjugation with the literal
loop $\lambda h(\lambda)$.
These maps extend to inverse continuous maps of the full
source completions and commute with the original filling $q$.

Put $w=\phi^\dagger(\lambda)$, a path from the original target
basepoint $c_0$ to $\Pi(c_0)$. The actual target automorphism is
\[
 \mathcal A_G(z)=w\Pi(z)w^{-1},\qquad
 \mathcal A_G^{-1}(z)=\Pi^{-1}(w^{-1}zw).                       \tag{F.23.12}
\]
The literal equality (F.23.8) gives
$\mathcal A_G\Phi^\dagger=\Phi^\dagger\mathcal A_{\rm src}$.
It follows on the actual source filling kernel that
$\mathcal A_G(K^+)=K^+$; the commuting target fillings give
$\mathcal A_G(D^+)=D^+$ and preserve the corresponding old-point
kernel $J$.

Write $F_+$ and $F_-$ for the full $G$ forgettings in F.22.
Let $\Pi_+$ be the old-label permutation together with renaming
the retained negative new label as positive. Then
$F_+\Pi=\Pi_+F_-$. Set $w_+=F_+(w)$. The two inverse maps are
\[
 \begin{aligned}
 \mathcal B_+(r)&=w_+\Pi_+(r)w_+^{-1},\\
 \mathcal B_+^{-1}(r)&=\Pi_+^{-1}(w_+^{-1}rw_+),\qquad
 F_+\mathcal A_G=\mathcal B_+F_- .
 \end{aligned}                                                \tag{F.23.13}
\]
The old-point filling square identifies its restrictions
$V_-\to V_+$, and the common-three-point filling square identifies
$R_-\to R_+$. These are isomorphisms with the displayed inverse.
Finally the equality on $K^+$ and the last formula give
$\mathcal B_+(\mathcal K_-)=\mathcal K_+$, proving (F.23.2).

\emph{A full geometric source basis.}
Use $s_*=2+2i$ and axis-parallel squares of half-side
$h_0=10^{-4}$. Their twenty-seven centers are $0$, both signs of
the five old centers in the first table, and both signs of the
eight centers in the second table. Every pair $(r,t)$ below
denotes $(r+it)/10^6$.
\[
 \begin{array}{c|rr}
 \text{old label}&r&t\\\hline
 1&1000000&0\\
 \zeta&309017&951057\\
 \zeta^2&-809017&587785\\
 \zeta^3&-809017&-587785\\
 \zeta^4&309017&-951057
 \end{array}
 \qquad
 \begin{array}{cc|rr}
 u&\text{square}&r&t\\\hline
 1&A_u&882644&1077509\\
 1&B_u&454951&555393\\
 \zeta^2&A_u&80730&-1390528\\
 \zeta^2&B_u&41611&-716735\\
 \zeta^3&A_u&1260728&-1131556\\
 \zeta^3&B_u&439300&-394290\\
 \zeta^4&A_u&1685063&174411\\
 \zeta^4&B_u&587159&60773
 \end{array}                                                  \tag{F.23.14}
\]
For each square take the straight whisker from $s_*$ to its
northeast corner, traverse its boundary counterclockwise, and
return along the same whisker. Call these loops $\gamma_0,\ldots,
\gamma_{26}$ in the producer's recorded order: the eleven old
punctures first, then the sixteen signed collisions.

Here are the exact checks which make these actual free generators.
For every old center the radius-$h_0/2$ disk contains its named
old root. For a new center $c$ and its named square $V$, the checks give
\[
 |c|>\tfrac12,\qquad |c^2-V|<10^{-5},\qquad
 (h_0/2)^2+10^{-5}<h_0/2.                                    \tag{F.23.15}
\]
On the circle of radius $h_0/2$ about $c$, the linear term
$2c(z-c)$ has absolute value greater than $h_0/2$, while
$(z-c)^2+c^2-V$ has smaller absolute value. Rouch\'e's theorem
therefore gives exactly one root of $z^2-V$ in that disk.
The twenty-seven squares are pairwise disjoint by the exact
center-distance inequalities. Each whisker stays farther than
$\sqrt2h_0$ from every other square's center. Its basepoint is
northeast of its own square, and its first intersection is that
square's northeast corner. All whisker directions are distinct.
They consequently form an embedded common-base star. Cutting
the punctured plane along this star gives the usual free basis
of its fundamental group. Filling the sixteen collision punctures
retains precisely the eleven old free generators.

\emph{The exact projection perturbation and crossing certificate.}
The field is the same as in F.20:
\[
 \mathbb Q(\theta),\qquad
 \theta^4-5\theta^2+5=0,\qquad 19/10<\theta<2.
\]
Real signs are decided by rational isolation of this indicated
root and exact field arithmetic. Put $P_0=1+371i/1000$ and use
the literal offsets
\[
 \delta_j=\frac{10^{-14}j^2}{P_0},
                  \qquad 0\leq j\leq8,                       \tag{F.23.16}
\]
in the finite-label order of (F.23.6).
For each segment and each pair of strands write their original
difference as $N/D$. The producer verifies the polynomial
inequality, over the entire real parameter interval,
\[
 |P_0N|^2-\bigl(10^{-14}(i^2-j^2)\bigr)^2|D|^2>0.
                                                               \tag{F.23.17}
\]
Every denominator is nonzero on the specified segment.
This strict inequality proves that the entire interpolation
$z_j\mapsto z_j+t\delta_j$, $0\leq t\leq1$, is collision-free.
The polynomial is certified by rational Bernstein subdivision
with strictly positive coefficients on every terminal interval.
All its coefficients and terminal intervals are retained.

For the perturbed paths, real projected differences are again
exact field polynomials after multiplying by positive squared
denominators. Bernstein--Descartes subdivision isolates all their
zeros, and the certificate checks nonzero height and derivative
on each root interval. These intervals are refined until all
crossing times are disjoint. At each crossing the two labels
are adjacent in the current projection order; the height sign
then gives the clockwise Artin letter calibrated in F.20.
The final order is checked against the exact endpoint order.
The record retains the full polynomials, root intervals, height
and derivative signs, subdivision counts, and resulting words.
These operations prove the full path word by successive actual
projection-chamber comparisons.

The primary run has $137$ forward segments, $4,932$ pair checks,
and $1,412$ crossings. Its exact sign algorithm made $846,769$
field-sign calls. Return whiskers are the literal reverse paths.
The additional chart path below has five segments and ten
crossings. Thus the complete comparison uses $142$ forward
segments, $5,112$ pair checks, and $1,422$ crossings. The initial
unperturbed projection had simultaneous events from $0,a,-a$;
its rejected run is retained separately. The proved interpolation
(F.23.17) is the precise comparison used for the completed run.

\emph{The standard chart and the original endpoint whiskers.}
At $s_*$ let $b=\phi^\dagger(s_*)$. Let $r$ have the same old
coordinates as $b$ and the two new coordinates
$R_p=8+3i$, $R_n=12+6i$. The actual path $T:r\to b$ first
moves $R_p$ straight to $a(s_*)$ with $R_n$ fixed, and then
moves $R_n$ straight to $-a(s_*)$ with the positive point fixed.
Both paths and their entire offset interpolations are certified.

The standard configuration is
\[
 r_{\rm std}=(0,1,3,2,4,5,6,8+3i,12+6i).                      \tag{F.23.18}
\]
The path $G:r_{\rm std}\to r$ moves the five old nonconstant
labels one at a time along the specified straight segments,
keeping both new points and $0,1$ fixed. The certified projected
order of $r_{\rm std}$ is the F.22 order
$0,1,v_\zeta,v_1,v_{\zeta^2},v_{\zeta^3},v_{\zeta^4},p,n$.
For an explicit straightening, write these standard projected
points as $x_j+iy_j$ and take the continuous piecewise linear
function $f$ with $f(x_j)=y_j$, affine between adjacent $x_j$
and constant outside them. The oriented plane map
\[
 x+iy\longmapsto x+i(y-f(x))
                                                               \tag{F.23.19}
\]
has inverse $x+iy\mapsto x+i(y+f(x))$. Its interpolation using
$tf$ changes no projected order. This gives the actual chart
and inverse used for the clockwise generators and the lower
whiskers; the same chart is restricted for every filling.

At every endpoint $z$ retain the actual path
$H_z(t)=(z_j+t\delta_j)_j$. The basepath producer checks
the original and shifted orders of all $284$ segment endpoints,
in $10,224$ exact pair comparisons. Every $H_z$ stays in its
original projection-order chamber. Writing a subscript $1$
for an offset path, the certified two-parameter homotopies give
the following identities in the actual path groupoid:
\[
 \begin{gathered}
 \gamma=H_b\gamma_1H_b^{-1},\qquad
 TH_b=H_rT_1,\qquad GH_r=H_{r_{\rm std}}G_1,\\
 GT\gamma T^{-1}G^{-1}
 =H_{r_{\rm std}}
       (G_1T_1\gamma_1T_1^{-1}G_1^{-1})H_{r_{\rm std}}^{-1}.
 \end{gathered}                                                \tag{F.23.20}
\]
Thus the computed word is carried to the original standard
configuration by its displayed endpoint whisker. The chamber
comparison preserves precisely the named clockwise generators.
The first two coordinates on an offset loop are the constant
values $0,1+t\delta_1$. Dividing all coordinates by
$1+t\delta_1$, which is nonzero by the same pairwise certificate,
is the literal affine normalization of F.22.14. Consequently
(F.23.20) also holds in every original normalized target.

Retain the original F.20.4a path $\eta:s_0\to s_*$.
Its existing exact certificate places it in $Y^\dagger$:
the radial first segment misses the old and inner collision
roots by its radius, and misses all outer roots by the four
nonzero imaginary parts $\operatorname{Im}(A_u\overline{s_0^2})$.
The remaining segments lie outside radius two.
Let $\mathcal N$ be the exact affine normalization of F.22.14.
The full original comparison path in the normalized target is
\[
             \Xi=\mathcal N(GT\phi^\dagger(\eta)^{-1}).         \tag{F.23.21}
\]
All three paths $G,T,\phi^\dagger(\eta)$ keep their first
two coordinates literally $0,1$. Hence the normalization
$(z_j-z_0)/(z_1-z_0)$ leaves every remaining coordinate
unchanged at every point of this entire path. Thus $\Xi$
has exactly the original standard and $s_0$ endpoints.
Let $b_j$ denote the recorded standard nine-strand braid
of $GT\phi^\dagger(\gamma_j)T^{-1}G^{-1}$, interpreted through
(F.23.20). Then the complete source formula is
\[
 \Phi^\dagger(\eta\gamma_j\eta^{-1})
           =\Xi^{-1}\mathcal N_*(b_j)\Xi,
                          \qquad0\leq j\leq26.                 \tag{F.23.22}
\]
This follows by cancelling the identical based paths; it keeps
the original source basepoint and every original target label.

Extend the finite-point path $\Xi$ by successive supported local
disk isotopies. The image of the standard disk containing the
first eight points and excluding the ninth specifies the
permitted F.22 disk at the original base configuration. The
inverse isotopy and (F.23.19) give its inverse chart. Applying
the isotopy to a loop gives a square whose two side paths are
$\Xi$; that square identifies the induced fiber-coordinate map
with the groupoid conjugation in (F.23.22). The original F.22
section still has the constant comparison paths at $(p,n)$.
Forgetting labels from the same square proves all its based
filling comparisons.

The twenty-seven actual source generators freely generate
$\pi_1(Y^\dagger,s_0)$. Their recorded images therefore give
its entire homomorphism. The universal full completion maps
give the unique continuous map on $L^\dagger$ with these images.
It equals the original completed $\Phi^\dagger$, since the
two maps agree on its dense discrete source. This proves
the complete-source assertion (F.23.3).

\emph{The two retained source sections and the exact marginal maps.}
The eleven old generators give a particular section
$i_{\rm geo}:L\to L^\dagger$ of the original filling $q$.
Retain the F.19 section $i$ as well. On the actual target
coordinates put
\[
 c(\ell)=\Phi^\dagger(i_{\rm geo}(\ell)i(\ell)^{-1})\in K^+,
       \qquad j_{\rm geo}(\ell)=c(\ell)j(\ell).                \tag{F.23.23}
\]
These formulas follow because both sections have the same
$q$-image $\ell$. The exact comparison of the two pullback
semidirect coordinates is
\[
 (d,\ell)\longmapsto(d c(\ell)^{-1},\ell),
       \qquad(e,\ell)\longmapsto(e c(\ell),\ell).              \tag{F.23.24}
\]
Both maps express the same actual element, and their
composites are the identity. The original F.19 meridians
keep their original whiskers. The twenty-seven new lassos
and the old generators define the same actual source
filling kernel; restricting the same $\Phi^\dagger$ to
that kernel gives exactly $K^+$ in either coordinate system.

Delete strand $9$ from each $b_j$ to obtain its positive
eight-strand braid, delete strand $8$ to obtain its negative
one, and delete both to obtain the seven-old-strand braid.
The algorithm tracks the literal labels through every crossing.
The resulting paths and (F.23.22) are the full $G$ forgetting
maps $F_+,F_-$ before their restrictions to $J$.
The inverse triangular fiber maps of F.20 and F.22 commute
with these fillings: a deleted lower-whiskered clockwise
generator becomes trivial, and every retained one becomes
the corresponding generator in the retained order.
Equivalently, filling the same variables in
\[
 \tau_i(X)=X_n^{-1}\cdots X_{i+1}^{-1}
                         X_i^{-1}X_{i+1}\cdots X_n
                                                               \tag{F.23.25}
\]
gives exactly the triangular word in that retained order.
The additional calibration producer checks the full pure-braid
filling identity on every standard pure generator at the ranks
and deletion sets used here, as well as the Artin relations and
both triangular inverse maps. Thus its finite list of free-word
identities proves those homomorphism identities on the entire
pure braid groups.

Each of the sixteen collision words has trivial old braid and
trivial opposite-sign marginal braid, both by free reduction
of the recorded words and by the actual forgetting maps.
Its remaining linking vector is supported on its selected
old/new pair with value $-1$. This is the counterclockwise
simple holomorphic collision in the retained clockwise Artin
convention. The full words, including this sign, are retained.

\emph{Exact finite evaluation of the marginal maps.}
The following recurrence evaluates each recorded eight-strand
word in any finite group, without expanding intermediate free
words. Begin with $x_1,\ldots,x_7$, set $x_8=1$, and set all
conjugator values $c_j=1$. For a positive letter $\sigma_j$ use
\[
 (x_j,x_{j+1},c_j,c_{j+1})
  \longmapsto
 (x_jx_{j+1}x_j^{-1},\,x_j,\,x_jc_{j+1},\,c_j),                 \tag{F.23.26}
\]
and for its inverse use
\[
 (x_j,x_{j+1},c_j,c_{j+1})
  \longmapsto
 (x_{j+1},\,x_{j+1}^{-1}x_jx_{j+1},\,c_{j+1},
                                      \,x_{j+1}^{-1}c_j).     \tag{F.23.27}
\]
Induction on the displayed Artin substitutions proves these
identities for the images and their labelled conjugators at
each prefix. At a pure final word put $u=c_8^{-1}$; the final
seven $x_i$ evaluate the filled automorphism $\alpha$.
The actual conjugation on the fiber, proved by the cocycle
and inverse triangular comparison, evaluates as
\[
                         X_i\longmapsto u x_i u^{-1}.          \tag{F.23.28}
\]
For the negative marginal these are precisely the Artin
coordinates of $V_-$. For the positive marginal, a seven-tuple
of values on its actual geometric generators first undergoes
$\tau_7$; use those as the initial $x_i$. A collision seed then
has value $u$. For a source action apply $\tau_7$ also to the
final tuple in (F.23.28). This is exactly the evaluation of
the two inverse coordinate maps, since $\tau_7^2=1$.

There are eight collision seeds for either sign, and eleven
old source automorphisms. They are the computed actual
images and actions for $i_{\rm geo}$. The image of the signed
source filling kernel is the closed subgroup generated by
all their $L$-conjugates, as follows from the free source basis
and its full completed filling. Its image is exactly
$\mathcal K_\pm$ by F.22.21 and (F.23.23)--(F.23.24).

\emph{The finite candidate reduction with actual peripheral maps.}
F.21 and the actual surjections $p_\pm:D^+\to R_\pm$ give
\[
        \overline{\langle\!\langle\mathcal K_\pm
                    \rangle\!\rangle}_{R_\pm}=R_\pm.          \tag{F.23.29}
\]
Consequently in every finite image its image normally
generates the image of $R_\pm$, and has that entire image
when the latter is nilpotent, by F.21's finite argument.

The four deleted-puncture coordinates also retain their
actual meridians. For a collision over old label $u$, the
simple local parameter $a-v_u$, or $-a-v_u$, takes its
counterclockwise source circle once around that collision.
Transport its whisker by the labelled plane isotopy along
the specified source path. The inverse endpoint isotopy
carries the local circle to a based peripheral loop in the
original seven-punctured plane. Comparing its transported
whisker with the specified lower whisker gives the actual
conjugating loop between these two peripheral loops.
In the positive geometric coordinates this identifies the
seed with a conjugate of $u_j^{-1}$. In the negative
coordinates (F.23.25) identifies $U_j$ with a conjugate of
the inverse clockwise lasso; the seed is consequently a
conjugate of $U_j$. These comparisons preserve orientation
and the named puncture. In particular, in every finite
homomorphism the seed and its named deleted coordinate
have exactly the same order.

For $H=S_3$, the proper subgroups are the trivial subgroup,
three subgroups $C_2$, and the normal subgroup $C_3$.
For $H=A_4$, they are the trivial subgroup, three subgroups
$C_2$, four subgroups $C_3$, and the normal subgroup $V_4$.
The $S_3$ list follows directly from its three transpositions
and two three-cycles; a subgroup with two distinct
transpositions, or with both a transposition and a
three-cycle, is the full group. The $A_4$ list, its actual
element sets and normal closures, is also obtained by the
exhaustive finite subgroup closure algorithm in the
coefficient certificate. All these
proper subgroups are nilpotent. The proper subgroups
normally generating $S_3$ are its three $C_2$ subgroups;
those normally generating $A_4$ are its four $C_3$ subgroups.

Thus a failure of (F.23.4) requires both
$\rho(V_\pm)=H$ and $\rho(R_\pm)=H$. Its marginal image
then lies in a proper subgroup $M=C_2$ for $S_3$, or
$M=C_3$ for $A_4$. Conjugation in $H$ carries these respective
subgroups to one fixed $M$, so all homomorphisms with such
a failure occur in the finite enumeration with that fixed
subgroup. The peripheral comparison just proved makes all
four deleted-coordinate values have order one or two in
the first case, and order one or three in the second.
At least one is nonidentity, since they normally generate
the image of $R_\pm$.

Every seven-tuple in $H^7$ is the values of a unique
continuous homomorphism from the actual $\widehat F_7$.
The complete possible-failure candidate counts are therefore
\[
 \begin{aligned}
 S_3:&\quad 6^3(4^4-1)-3\,2^3(2^4-1)=54,720,\\
 A_4:&\quad 12^3(9^4-1)-4\,3^3(3^4-1)=11,327,040.
 \end{aligned}                                                \tag{F.23.30}
\]
The subtracted families are exactly the tuples contained
in a proper subgroup while having a nonidentity deleted
coordinate. They are disjoint because that coordinate
determines the particular $C_2$ or $C_3$ subgroup.

\emph{The exhaustive source-orbit test and its precise conclusion.}
For each candidate tuple evaluate the eight actual collision
words by (F.23.26)--(F.23.28), and retain the tuple precisely
when all eight values lie in the fixed $M$. Call the finite
set of retained tuples $\mathcal T_0$. Each of the eleven
actual old source automorphisms acts as a permutation of
$\operatorname{Hom}(V_\pm,H)$, by precomposition.
It preserves generation of the image. Its actual pure
labelled point-pushing action carries each deleted
peripheral generator to a conjugate of itself, so it
also preserves the candidate order restrictions and the
condition that some deleted value is nonidentity.
Record all eleven images of every element of $\mathcal T_0$.
Delete an element when one of these images is outside
$\mathcal T_0$, and repeat on the retained set.

This computes exactly the largest subset of
$\mathcal T_0$ invariant under the eleven permutations.
Indeed a subset mapped into itself by a permutation of a
finite set is also invariant under its inverse. Hence
forward invariance under the eleven maps is invariance
under their generated group. Their profinite source
completion has the same finite orbits: the continuous
action on this finite representation set has dense
discrete source image and therefore equal finite image.
The resulting invariant set consists exactly of the
candidates for which every source-conjugate collision
seed has image in $M$. Those source-conjugates generate
the actual marginal subgroup, so this is the exact
remaining failure set.

Both signs were enumerated separately. The complete
executed results are
\[
 \begin{array}{c|r|r|r|r}
 H&\text{candidates}&|\mathcal T_0|&
       \text{first deletion}&\text{survivors}\\\hline
 S_3&54,720&624&624&0\\
 A_4&11,327,040&15,552&15,552&0
 \end{array}                                                  \tag{F.23.31}
\]
The table holds for each sign. The outputs retain every
element of $\mathcal T_0$, its seven coordinates, all
eleven image keys and edges, and its deletion layer.
Multiplication is literal composition of the enumerated
permutations. The $A_4$ run streams the same exhaustive
candidate set in finite batches; it does not sample it.
The empty invariant sets prove (F.23.4) for every specified
continuous homomorphism, including the cases already
settled by the finite nilpotent image argument.

\emph{Executed certificate and retained scope.}
The complete source geometry is produced by
\texttt{full\_cross\_source\_geometry\_stage.py};
\texttt{full\_cross\_source\_basepaths\_stage.py}
retains the chart and endpoint comparisons;
\texttt{full\_cross\_marginal\_braids\_stage.py}
records all full and deleted words.
The two finite producers are
\texttt{marginal\_s3\_source\_stage.py} and
\texttt{marginal\_a4\_source\_stage.py}.
The sign-symmetry producer passes $64$ exact checks,
and \texttt{source\_artin\_calibration\_stage.py}
passes $2,473$ exact polynomial and free-word checks.
They are retained under \texttt{code/part\_a/} with their
full outputs, stdout, and hashes. Independent isolated
execution reproduced all $35$ complete-source, symmetry,
and finite-scan output files byte for byte; the calibration
output was independently reproduced as well.

All original sections, fillings, signs, and based
comparison paths enter the displayed maps. The finite
claim is the equality of the two images under every map
$V_\pm\to S_3$ or $V_\pm\to A_4$. The remaining full
marginal inclusion and the interaction subgroup
$K^+\cap W$ of F.22 retain their original source-membership
questions.
\end{proof}

\begin{claim}{F.24}{Proved}{All-field rank-one cocycles and the actual marginal affine fixed points}
Retain the actual continuous inclusions
$\mathcal K_\pm\leq R_\pm\leq V_\pm$ and the source maps of
F.19--F.23. Let $k$ be a finite field, with its discrete topology,
and let $\chi:V_\pm\to k^\times$ be a continuous character whose
restriction to $R_\pm$ is nontrivial. Give the additive group of
$k$ the left action $g\cdot b=\chi(g)b$, and denote this module
by $k_\chi$. The actual restriction maps are injective:
\[
\begin{aligned}
 Z^1_{\mathrm{cont}}(V_\pm,k_\chi)&\longrightarrow
 Z^1_{\mathrm{cont}}(\mathcal K_\pm,k_{\chi|\mathcal K_\pm}),\\
 H^1_{\mathrm{cont}}(V_\pm,k_\chi)&\longrightarrow
 H^1_{\mathrm{cont}}(\mathcal K_\pm,k_{\chi|\mathcal K_\pm}).
\end{aligned}                                                   \tag{F.24.1}
\]
Here a crossed homomorphism satisfies
$d(gh)=d(g)+\chi(g)d(h)$, and the principal crossed
homomorphisms are $d_b(g)=(\chi(g)-1)b$, for $b\in k$.

Write $\operatorname{Aff}_1(k)=k\rtimes k^\times$, acting on
$k$ by $(b,a)z=b+az$, and let
$T_k=\{(b,1):b\in k\}$ be its translation subgroup.
For every continuous homomorphism
$\rho:V_\pm\to\operatorname{Aff}_1(k)$ one has
\[
 \operatorname{Fix}_k\rho(\mathcal K_\pm)
       =\operatorname{Fix}_k\rho(R_\pm).                        \tag{F.24.2}
\]
For the representations whose scalar character is nontrivial
on $R_\pm$, both sides also equal
$\operatorname{Fix}_k\rho(V_\pm)$.
For every representation in the specified class
\[
 \left\{\rho\in\operatorname{Hom}_{\mathrm{cont}}
       (V_\pm,\operatorname{Aff}_1(k)):
           \rho(\mathcal K_\pm)\cap T_k=\{1\}\right\},
\]
the exact marginal images agree:
\[
                 \rho(\mathcal K_\pm)=\rho(R_\pm).             \tag{F.24.3}
\]
The coefficient calculation proving (F.24.1) has integer
Laurent coefficients and proves its rank assertion over every
field, including characteristics two and three. Its input is
the literal negative marginal braid data of F.23, with the
same source lassos and comparison maps. The positive assertion
is transported by the explicit F.23.13 isomorphism.
The full marginal equalities and the equality $K^+=D^+$
remain unresolved.
\end{claim}
\passage{Part F root-adjoining intersection question; Section 30, folios 194-196, scans 383-387; Section 31, folios 198-199 verso, scans 391-394; original F7/F19 residual subgroup-membership reduction}


\begin{proof}
\emph{The retained words and the exact coefficient maps.}
We first use the negative coordinates
\[
 \begin{gathered}
 V_-=\widehat F(X_1,X_2,X_3,X_4,X_5,X_6,X_7),\qquad
 X_{j+3}=U_j\quad(1\leq j\leq4),\\
 R_-=\ker\bigl(V_-\longrightarrow
                     \widehat F(X_1,X_2,X_3)\bigr).
 \end{gathered}                                                \tag{F.24.4}
\]
These are the actual F.22.5 filling coordinates. Put
$\Lambda=\mathbb Z[t_1^{\pm1},\ldots,t_7^{\pm1}]$.
Let $F_7$ be the discrete free group on these same seven
generators, with its dense completion map $F_7\to V_-$.
The universal character and crossed homomorphism are
\[
 t:F_7\longrightarrow\Lambda^\times,\quad t(X_i)=t_i,
 \qquad
 D:F_7\longrightarrow\Lambda^7,\quad D(X_i)=e_i,
                                                               \tag{F.24.5}
\]
where $\Lambda^7$ consists of row vectors and
$D(vw)=D(v)+t(v)D(w)$. These maps are well defined on
reduced words: the inverse formula is
$D(v^{-1})=-t(v)^{-1}D(v)$, which cancels an adjacent
letter and its inverse. Equivalently they are the two
coordinates of the group homomorphism
$F_7\to\Lambda^7\rtimes\langle t_1,\ldots,t_7\rangle$
specified on its free generators.

For a character $\chi$ and a continuous crossed homomorphism
$d$ as in the statement, let
\[
 \operatorname{ev}_\chi:\Lambda\longrightarrow k,
       \qquad t_i\longmapsto\chi(X_i),\qquad
 v_d=(d(X_1),\ldots,d(X_7))^{\mathsf T}.
\]
This is a ring homomorphism because each $\chi(X_i)$ is
nonzero. Induction on the reduced word gives the exact formula
\[
             d(w)=\operatorname{ev}_\chi(D(w))v_d
                      \qquad(w\in F_7).                       \tag{F.24.6}
\]
All words below are finite words in this dense free group,
interpreted in $V_-$ by its stated completion map.
In particular, the calculation does not substitute a
completion of a different source subgroup.

Let $\nu_0,\ldots,\nu_7$ be the eight negative collision
words in F.23, in the order
\[
 (u,\epsilon)=
 (1,+),(1,-),(\zeta^2,+),(\zeta^2,-),
 (\zeta^3,+),(\zeta^3,-),(\zeta^4,+),(\zeta^4,-).
                                                               \tag{F.24.7}
\]
The sign here selects the two roots of each negative
collision square $B_u$; every word in (F.24.7) belongs
to the negative marginal. The original paths $\eta$,
$G$, $T$, and $\Xi$ and their affine normalization are
exactly those retained in F.23.19--22.
Let $B_i:V_-\to V_-$, $0\leq i\leq10$, be conjugation
on this point-pushing group by the image of the $i$th
old source generator under the geometric source section
$j_{\mathrm{geo}}$ of F.23.23. Thus $B_i$ is a continuous
automorphism and
\[
 B_i(\mathcal K_-)=\mathcal K_- .                              \tag{F.24.8}
\]
Indeed the source filling kernel is normal in the source;
its image is $K^+$; projection to the negative marginal
commutes with this conjugation. These are the original
maps, with the old and geometric section comparison
$j_{\mathrm{geo}}(l)=c(l)j(l)$ from F.23.23--24.
The old source indices are $0$ for its zero puncture,
$1,\ldots,5$ for $+\zeta^j$, and $6,\ldots,10$ for
$-\zeta^j$, in each case $0\leq j\leq4$.

Every $B_i$ preserves the abelianization of $V_-$.
For the explicit proof in these coordinates, start the
decorated Artin recurrence F.23.26--28 with the first
seven characters $t_1,\ldots,t_7$ and the eighth character
one. A crossing exchanges the two characters, and each
recorded old braid is pure. Its final conjugation by $u$
does not change them. Hence $t(B_i(X_j))=t_j$, and
\[
 A_i=\bigl(D(B_i(X_j))\bigr)_{1\leq j\leq7},\qquad
 D(B_i(w))=D(w)A_i .                                          \tag{F.24.9}
\]
Both identities follow on generators and then on products
and inverses from the crossed-homomorphism identity.
Consequently the following sixty-four rows evaluate
the crossed homomorphism on actual elements of
$\mathcal K_-$:
\[
 \begin{gathered}
 a_r=D(\nu_r)\quad(0\leq r\leq7),\qquad
 a_rA_i\quad(0\leq r\leq7,\quad i\in\mathscr I),\\
                 \mathscr I=(2,9,8,10,1,0,7).
 \end{gathered}                                                \tag{F.24.10}
\]
Number the rows from zero, first the eight $a_r$, then
the eight rows for each index of $\mathscr I$ in the
displayed order.

\emph{Exact production and rank of the sixty-four rows.}
For completeness we give the full arithmetic recipe
and the terminal identities of its certificate. Encode
an element of the universal affine group by $(a,v)$,
where $a\in\mathbb Z^7$ is its character exponent and
$v\in\Lambda^7$ its crossed-homomorphism row. Then
\[
 (a,v)(b,w)=(a+b,v+t^a w),\qquad
 (a,v)^{-1}=(-a,-t^{-a}v).                                    \tag{F.24.11}
\]
Initialize the eight Artin images as $(e_j,e_j)$ for
$1\leq j\leq7$ and $(0,0)$ for the erased eighth
generator, and initialize all eight conjugators as
$(0,0)$. Apply the exact positive and inverse substitutions
F.23.26--27 in this group, using (F.24.11). Put
$u=c_8^{-1}$ at the end. The row of $u$ is the seed row;
the rows of $u x_j u^{-1}$ form $A_i$ for an old word.
This applies the proved decorated recurrence to a
specified group, so it evaluates the same actual
words as F.23.28. The separate producer records
eleven old matrices and eight seed rows. All of them
also satisfy
\[
 D(w)(t_1-1,\ldots,t_7-1)^{\mathsf T}=t(w)-1,                 \tag{F.24.12}
\]
which follows directly from (F.24.5).

Here is the exact row elimination used to check their
rank. In chart $j$, work over
$\Lambda[(t_j-1)^{-1}]$. A row normalization multiplies
by a Laurent monomial and divides by a common power of
$t_j-1$. A pivot $p$ is accepted only with the exact
coefficient identity
$p=\varepsilon t^a(t_j-1)^m$, where
$\varepsilon\in\{1,-1\}$ and $m\geq0$. With coefficient
$c$ in its column, replace a later row $r$ by $pr-c r_p$
and normalize again. Each operation has an invertible
row multiplier in the displayed ring and preserves
the row module there. Every polynomial is stored as
its integer coefficients and seven integer exponents;
the record retains the input coefficients, shifts,
division counts, pivot columns, and resulting rows.

The following are the complete pivot orders, with
zero-based row indices and one-based column indices.
The last chart first specializes $t_4=t_5=t_6=1$.
Each listed final pivot coefficient is a signed Laurent
monomial, so every pivot remains nonzero in every field
on its chart.
\[
\begin{array}{c|l|l|r}
 j&\text{pivot rows}&\text{pivot columns}&\text{row operations}\\\hline
 4&11,8,0,6,4,1&5,4,1,7,6,3&389\\
 5&11,8,3,6,4,2,34&5,4,1,7,6,2,3&444\\
 6&11,8,13,14,20,5,28&5,4,6,7,2,1,3&445\\
 7&0,2,4,14,39,6,46&4,5,6,7,2,1,3&196
\end{array}                                                     \tag{F.24.13}
\]
The first chart has one remaining column, column two.
After its six displayed pivots, the retained rows
numbered five and forty-eight are exactly
\[
 \begin{aligned}
 r_5&=(1-t_6)e_2,\\
 r_{48}&=\bigl(t_1t_2t_3t_4t_5t_7(t_6-1)-1\bigr)e_2.
 \end{aligned}
\]
They give the explicit seventh pivot
\[
           r_{48}+t_1t_2t_3t_4t_5t_7 r_5=-e_2 .              \tag{F.24.14}
\]
Thus all four charts have rank seven. The first three
charts cover $t_4\ne1$, $t_5\ne1$, and $t_6\ne1$.
The remaining stratum has those three values equal to
one and $t_7\ne1$, and is covered by the fourth chart.
This proves rank seven at every seven-tuple of nonzero
field elements for which some $t_j\ne1$, $4\leq j\leq7$.
No integer other than $1$ or $-1$ is inverted.

The accompanying exact certificate contains every row
and operation just described. Its independent verifier
imports none of the three producers: it reconstructs
the nineteen word evaluations with the group law
(F.24.11), reconstructs all sixty-four rows, replays
every row operation, proves every division by multiplying
back by $t_j-1$, and checks the triangular pivot orders
and (F.24.14). It also replays the longer scalar-ideal
certificate retained from the original search. There
are $13\,982$ exact assertions in this independent
verification. The manifest fixes the braid input,
all producers and outputs, and this verification by
size and SHA-256. These finite polynomial identities,
together with their invertible operations, establish
the asserted rank over all characteristics.

\emph{Restriction, continuity, and the sign comparison.}
The restriction $\chi|R_-$ is nontrivial precisely when
at least one of $\chi(X_4),\ldots,\chi(X_7)$ differs
from one: $R_-$ is the closed normal subgroup generated
by these four generators, and the target $k^\times$
is abelian. For a crossed homomorphism vanishing on
$\mathcal K_-$, every row in (F.24.10) annihilates $v_d$
by (F.24.6) and (F.24.8)--(F.24.9). The rank assertion
therefore gives $v_d=0$. The crossed-homomorphism rule
then gives zero on the dense group $F_7$; continuity
into the finite Hausdorff space $k$ gives zero on
$V_-$. This proves injectivity on $Z^1$.
For a class restricting to zero in $H^1$, its
representative restricts to $(\chi-1)b$ for some $b\in k$.
Subtract that same principal crossed homomorphism on
$V_-$ and apply the proved $Z^1$ injectivity. This
proves injectivity on $H^1$ with the exact coefficient
restriction in (F.24.1).

For the positive sign use the specific continuous
isomorphism $\mathcal B_+:V_-\to V_+$ of F.23.13.
It carries $R_-$ onto $R_+$ and $\mathcal K_-$ onto
$\mathcal K_+$. Pullback sends $(\chi,d)$ to
$(\chi\mathcal B_+,d\mathcal B_+)$; pullback along
$\mathcal B_+^{-1}$ is its inverse. It preserves the
crossed-homomorphism equations and the principal
rows and commutes literally with restriction.
Nontriviality on $R_+$ transfers to nontriviality on
$R_-$. These maps prove both positive assertions.

\emph{The affine representations and their fixed points.}
Write a given representation as
$\rho(g)=(d(g),\chi(g))$. Its multiplication law makes
$d$ a continuous crossed homomorphism. For a point
$z\in k$ put
\[
       d_z(g)=\rho(g)z-z=d(g)+(\chi(g)-1)z .                  \tag{F.24.15}
\]
It is another continuous crossed homomorphism for
the same character. A point fixed by
$\rho(\mathcal K_\pm)$ makes $d_z$ vanish there.
For $\chi|R_\pm$ nontrivial, (F.24.1) makes $d_z$
zero on $V_\pm$. The inclusions of fixed sets from
$\mathcal K_\pm\leq R_\pm\leq V_\pm$ now give the
three equal fixed sets claimed above.

For $\chi|R_\pm=1$, the image of $R_\pm$ lies in
the abelian group $T_k$. Apply the F.21 normal-generation
conclusion through the actual surjection
$p_\pm:D^+\to R_\pm$ of F.22.6 and then through
$\rho|R_\pm$. The normal closure of
$\rho(\mathcal K_\pm)$ in $\rho(R_\pm)$ is
$\rho(R_\pm)$; in this abelian group the closure is
the subgroup itself. Hence their images agree and
their fixed sets agree. This completes (F.24.2).

Finally put $M=\rho(\mathcal K_\pm)$ for a representation
in the class stated before (F.24.3). Its scalar
projection $M\to k^\times$ is injective because its
kernel is $M\cap T_k$. Thus $|M|$ divides $|k|-1$
and is invertible in $k$. The explicit point
\[
                     z_M=|M|^{-1}\sum_{m\in M}m(0)            \tag{F.24.16}
\]
is fixed by $M$. For $h=(b,a)\in M$, applying its
affine formula to this average gives
$h(z_M)=|M|^{-1}\sum_{m\in M}(hm)(0)=z_M$, because
left multiplication by $h$ permutes $M$.
For a nontrivial character on $R_\pm$, (F.24.2) and
its stronger assertion put all of $\rho(V_\pm)$ in
the stabilizer of $z_M$. The scalar map on that
stabilizer has the explicit inverse
$a\mapsto((1-a)z_M,a)$, so the stabilizer is abelian.
F.21 normal generation, transported through the same
two actual maps as above, gives (F.24.3). For a
trivial character on $R_\pm$, the preceding abelian
translation argument already gives (F.24.3).

The remaining finite-affine case consists exactly of
representations with $\rho(\mathcal K_\pm)\cap T_k$
nontrivial. The argument has not proved equality of
the two marginal images in that class. In particular
it has not constructed an ambient invariant quotient
that kills this translation intersection. The actual
inclusions of F.23.1 and the interaction kernel of
F.22.6--8 remain the maps whose surjectivity is under
investigation.
\end{proof}

\begin{claim}{F.25}{Proved}{The actual affine marginal images with a common scalar field}
Retain $\mathcal K_\pm\leq R_\pm\leq V_\pm$ and their
actual maps from F.19--F.24. Let $k$ be a finite field
of characteristic $p$, and let
$\rho:V_\pm\to\operatorname{Aff}_1(k)$ be continuous.
Write $\chi$ for its scalar character and put
\[
 \begin{gathered}
 M=\rho(\mathcal K_\pm),\qquad N=\rho(R_\pm),\qquad
 D=\chi(V_\pm),\qquad C=\chi(R_\pm),\\
 U=\{b\in k:(b,1)\in M\},\qquad
 k_0=\mathbb F_p(D),\qquad k_1=\mathbb F_p(C).
 \end{gathered}                                                \tag{F.25.1}
\]
The subgroup $U$ is a vector subspace of $k$ over
$k_1$. For the class of these representations in
which $U$ is also a vector subspace over $k_0$, the
actual finite images agree:
\[
                              M=N.                            \tag{F.25.2}
\]
Consequently the equality holds for every representation
in the explicitly specified class
\[
 \left\{\rho:\ C=\{1\}\ \text{or}\ k_1=k_0\right\}.
                                                               \tag{F.25.3}
\]
In particular, for every prime $p$ and every continuous
homomorphism to the full affine group of its prime field,
\[
 \rho(\mathcal K_\pm)=\rho(R_\pm)
       \qquad\bigl(\rho:V_\pm\to\operatorname{Aff}_1(\mathbb F_p)
                       \text{ continuous}\bigr).             \tag{F.25.4}
\]
All comparison maps, the choice of an actual scalar
lift in $M$, and the coefficient maps used in the proof
are specified below. The finite scalar-field class
$k_1\subsetneq k_0$ with $U$ failing to be a $k_0$-subspace
remains outside the proved equality; full marginal
surjectivity and $K^+=D^+$ remain unresolved.
\end{claim}
\passage{Part F root-adjoining intersection question; Section 30, folios 194-196, scans 383-387; Section 31, folios 198-199 verso, scans 391-394; original F7/F19 residual subgroup-membership reduction}


\begin{proof}
Fix either sign throughout. First transport the F.21
normal-generation statement through the actual continuous
surjection $p_\pm:D^+\to R_\pm$ of F.22.6. Its restriction
to $K^+$ has image $\mathcal K_\pm$. The image of the
closed normal closure is the closed normal closure of
the image: one inclusion follows by applying the map
to conjugates, and the other by taking the inverse
image of a closed normal subgroup containing the image.
The image is closed because the source is compact and
the target Hausdorff. Thus
\[
 \overline{\langle\!\langle\mathcal K_\pm
                         \rangle\!\rangle}_{R_\pm}=R_\pm,
 \qquad
 \langle\!\langle M\rangle\!\rangle_N=N,
 \qquad \chi(\mathcal K_\pm)=C.                               \tag{F.25.5}
\]
The second equality applies the actual map
$\rho|R_\pm$; the third applies its scalar projection
to the finite abelian group $k^\times$. When $C=\{1\}$,
the group $N$ is a subgroup of the abelian translation
group, so (F.25.5) already gives $M=N$.

We now take $C\ne\{1\}$. The set $U$ is an additive
subgroup of $k$, hence an $\mathbb F_p$-vector subspace.
Conjugation of a translation by an actual element of
$M$ with scalar $c$ is given by the literal formula
\[
             (b,c)(u,1)(b,c)^{-1}=(cu,1).                    \tag{F.25.6}
\]
The last equality of (F.25.5) supplies such an element
for every $c\in C$. Consequently $U$ is stable under
every scalar in $C$, hence under the ring
$\mathbb F_p[C]$. This finite subring of $k$ is a field:
multiplication by any nonzero element is an injection
of its finite set and therefore is onto, giving its
inverse in the ring. Thus $\mathbb F_p[C]=k_1$ and
$U$ is a $k_1$-subspace, as asserted.

The finite multiplicative subgroup $C$ is cyclic.
One direct verification chooses an element of maximal
order $n$: the prime-power components of element
orders in the abelian group show that every element
has order dividing $n$, since otherwise their commuting
components would give a larger order. All elements
then satisfy $X^n-1=0$ in the field, so $|C|\leq n$.
The chosen cyclic subgroup has $n$ elements, proving
$|C|=n$. Choose a generator $c\ne1$ and an actual
lift $m=(b,c)\in M$. Define
\[
 z=\frac{b}{1-c},\qquad \tau_z(w)=w+z,
 \qquad
 \rho_z(g)=\tau_z^{-1}\rho(g)\tau_z
           =\bigl(d(g)+(\chi(g)-1)z,\chi(g)\bigr),             \tag{F.25.7}
\]
where $\rho(g)=(d(g),\chi(g))$. The maps $\tau_z$ and
$\tau_z^{-1}=\tau_{-z}$ are explicit inverse affine
bijections of $k$, and their conjugations are inverse
group automorphisms of $\operatorname{Aff}_1(k)$.
The chosen lift now equals $(0,c)$. Conjugating by a
translation fixes every element of $T_k$, so the
translation intersection remains exactly $U$.

Let $M_z=\rho_z(\mathcal K_\pm)$ and $N_z=\rho_z(R_\pm)$.
Every element of $M_z$ has scalar $c^j$; multiplication
by its element $(0,c^{-j})$ leaves a translation in $U$.
Conversely $M_z$ contains all translations in $U$ and
the generator $(0,c)$. We have therefore proved the
literal equality and the associated inverse coordinates
\[
 \begin{aligned}
 M_z&=U\rtimes C,\\
 U\rtimes C&\longrightarrow M_z,&(u,a)&\longmapsto(u,a),\\
 M_z&\longrightarrow U\rtimes C,&(b,a)&\longmapsto(b,a).
 \end{aligned}                                                \tag{F.25.8}
\]
The group law is $(u,a)(v,a')=(u+av,aa')$, with
$U$ stable under $C$ by (F.25.6). Thus the complement
used here has been constructed from the actual
element $m\in M$.

Now let $U$ be a $k_0$-subspace, as in the first
assertion. The field $k$ is a finite-dimensional vector
space over $k_0$. Let
\[
 \mathscr L_U=\{\ell\in\operatorname{Hom}_{k_0}(k,k_0):
                                      \ell(U)=0\}.
\]
These exact coefficient maps satisfy
\[
                         U=\bigcap_{\ell\in\mathscr L_U}
                                           \ker\ell.          \tag{F.25.9}
\]
For a proof, extend a $k_0$-basis of $U$ to a basis of
$k$. Coordinate projections on the additional basis
vectors belong to $\mathscr L_U$ and detect every
nonzero coordinate outside $U$; the reverse inclusion
is part of the definition.

Put $d_z(g)=d(g)+(\chi(g)-1)z$, the translation
coordinate in (F.25.7). The affine multiplication law
makes it a continuous $k$-valued crossed homomorphism.
For each $\ell\in\mathscr L_U$ the composite
\[
 \ell d_z:V_\pm\longrightarrow(k_0,+),\qquad
 (\ell d_z)(gh)=(\ell d_z)(g)
                         +\chi(g)(\ell d_z)(h)               \tag{F.25.10}
\]
is an actual continuous crossed homomorphism with
the character $\chi:V_\pm\to k_0^\times$.
Indeed $D\subset k_0^\times$ and $\ell$ is $k_0$-linear,
which proves exactly the displayed multiplication
identity. All target spaces are finite and discrete,
so composition also preserves continuity. Equation
(F.25.8) gives $d_z(\mathcal K_\pm)\subset U$, and
therefore $\ell d_z$ vanishes on $\mathcal K_\pm$.
The character is nontrivial on $R_\pm$ because
$C\ne\{1\}$. Apply the proved F.24.1 restriction
injection over the finite field $k_0$. It gives
$\ell d_z=0$ on all of $V_\pm$, for every such $\ell$.
Together with (F.25.9) this proves
\[
 d_z(V_\pm)\subset U,\qquad
 N_z\subset U\rtimes C=M_z.                                  \tag{F.25.11}
\]
The inclusion $\mathcal K_\pm\leq R_\pm$ supplies the
reverse inclusion. Conjugating back by $\tau_z$ gives
the exact equality (F.25.2).

For the class $k_1=k_0$, the $k_0$-linearity of $U$
has already been proved using (F.25.6). Combining
this with the separate case $C=\{1\}$ proves
(F.25.3). For $k=\mathbb F_p$, both generated
subfields are $\mathbb F_p$ for every character,
so (F.25.4) follows for every prime, including two
and three. The proof applies to either of the actual
marginals, using its corresponding F.24 injection;
that injection already contains the explicit F.23
sign-symmetry comparison.

For the remaining class $k_1\subsetneq k_0$,
(F.25.6) supplies $k_1$-linearity of $U$. The
coefficient map (F.25.10) requires $k_0$-linearity.
The proof has established the passage precisely
for the representations specified in the statement,
and supplies no quotient of the actual ambient
module for an arbitrary such $U$.
\end{proof}

\begin{claim}{F.26}{Proved}{The exact trace obstruction and all affine marginal images with scalar field of order nine}
Retain the actual inclusions
$\mathcal K_\pm\leq R_\pm\leq V_\pm$ of F.19--F.25.
Fix either sign, and temporarily write these groups as
$K\leq R\leq V$. A proper affine marginal image
$\rho(K)\subsetneq\rho(R)$, for a continuous
$\rho:V\to\operatorname{Aff}_1(k)$ over a finite field,
has an actual continuous affine image of the following form:
\[
 \begin{gathered}
 \rho':V\longrightarrow k_0\rtimes D,\qquad
 \rho'(V)=k_0\rtimes D,\qquad \rho'(R)=k_0\rtimes C,\\
 \rho'(K)\ \leq\
       \ker\bigl(\operatorname{Tr}_{k_0/k_1}:k_0\to k_1\bigr)
                       \rtimes C,\\
 C=\chi(R)=\chi(K)\ne\{1\},\qquad
 D=\chi(V),\qquad
 k_1=\mathbb F_p(C)\subsetneq k_0=\mathbb F_p(D).
 \end{gathered}                                                \tag{F.26.1}
\]
Here $\chi$ is the original scalar character, retained
by the explicit conjugation and coefficient map constructed
below. In particular the last group in (F.26.1) is a
proper subgroup of the displayed $R$ image.

For any finite-field data $k_0,\chi,C,k_1$ as in
(F.26.1), including the equality case $k_1=k_0$, put
$q_1=|k_1|$ and $m=[k_0:k_1]$. The exact coefficient map is
\[
 \begin{aligned}
 \mathscr T_\chi:
 Z^1_{\rm cont}(V,k_{0,\chi})&\longrightarrow
 Z^1_{\rm cont}(K,k_{1,\chi|K}),\\
 d&\longmapsto\operatorname{Tr}_{k_0/k_1}(d|K).
 \end{aligned}                                                \tag{F.26.2}
\]
Its known kernel subspace is
\[
 \mathscr P_\chi=
 \{\,g\mapsto(\chi(g)-1)b:
                      b\in k_0,\ \operatorname{Tr}(b)=0\,\}.
                                                               \tag{F.26.3}
\]
Every element of $\ker\mathscr T_\chi\setminus\mathscr P_\chi$
gives a proper actual affine marginal image, and every
proper affine marginal image gives such an element for
the data in (F.26.1). Thus
$\ker\mathscr T_\chi=\mathscr P_\chi$ for all the stated
data is an exact remaining criterion for the full
scalar affine image equality. The proof gives an explicit
coefficient isomorphism identifying this criterion with
the summed restrictions for the distinct characters
$\chi^{q_1^j}$, $0\leq j<m$. Its rank target is $6m+1$
on $7m$ columns; no general attainment of that rank
is asserted here.

For the scalar field of order nine that rank is attained.
More precisely, for every finite field $k$ and every
continuous $\rho:V_\pm\to\operatorname{Aff}_1(k)$ whose
scalar values generate a subfield of order nine,
\[
                         \rho(\mathcal K_\pm)=\rho(R_\pm).
                                                               \tag{F.26.4}
\]
Consequently this equality holds for every continuous
map $V_\pm\to\operatorname{Aff}_1(\mathbb F_9)$, including
maps whose scalar field is $\mathbb F_3$. The finite
calculation proving the additional field case exhausts
exactly $7\,560$ character tuples and retains the
literal F.23 source automorphisms.
\end{claim}
\passage{Part F root-adjoining intersection question; Section 30, folios 194-196, scans 383-387; Section 31, folios 198-199 verso, scans 391-394; original F7/F19 residual subgroup-membership reduction}


\begin{proof}
\emph{Reduction through the actual affine image.}
Write $H=\rho(V)$, $N=\rho(R)$, $M=\rho(K)$ and
let $D,C$ be their scalar images as above. The F.21
normal generation, transported through the actual
$D^+\to R$ of F.22, gives
\[
 \langle\!\langle M\rangle\!\rangle_N=N,\qquad
                         \chi(K)=\chi(R)=C.                   \tag{F.26.5}
\]
These are exactly the maps and equalities proved
in F.25.5. A proper marginal requires $C\ne\{1\}$,
since $C=\{1\}$ makes $N$ abelian. Put
$k_0=\mathbb F_p(D)$ and $k_1=\mathbb F_p(C)$.
F.25 excludes equality of these two fields for
a proper marginal.

Choose a generator $c$ of $C$ and an actual lift
$(b,c)\in M$. Set $z=b/(1-c)$ and conjugate the
representation by $\tau_z^{-1}$, where
$\tau_z(w)=w+z$. As in F.25.7--8, this gives
\[
 M_z=U\rtimes C,\qquad
 U=\{u:(u,1)\in M\},\qquad (0,c)\in M_z,
                                                               \tag{F.26.6}
\]
with $U$ unchanged and $k_1$-linear. Let $A_H$ and
$A_N$ be the translation subgroups of $H_z$ and
$N_z$. Because $N_z$ is normal in $H_z$, conjugation
by every actual scalar lift in $H_z$ preserves both
subgroups and multiplies translations by its scalar.
Thus both are vector subspaces over
$\mathbb F_p[D]=k_0$. For $a\in A_H$ the exact commutator is
\[
                  [(a,1),(0,c)]=((1-c)a,1)\in N_z .
                                                               \tag{F.26.7}
\]
Since $1-c\ne0$ and $A_N$ is $k_0$-linear, this proves
$A_H=A_N=:A$. For an arbitrary $(b',d)\in H_z$ the
same multiplication gives
$[(b',d),(0,c)]=((1-c)b',1)\in A$ and hence $b'\in A$.
Multiplying by its inverse translation then gives
the corresponding pure scalar element. Consequently
\[
 H_z=A\rtimes D,\qquad N_z=A\rtimes C,\qquad
 M_z=U\rtimes C,\qquad U\subsetneq A .                          \tag{F.26.8}
\]
Here $A$ is $k_0$-linear and $U$ is $k_1$-linear.

Choose a $k_1$-linear functional $L:k\to k_1$
vanishing on $U$ and nonzero on $A$. To construct
one, extend a basis of $U$ by an element of
$A\setminus U$, then to a basis of $k$, and give
that selected element value one and the other
basis vectors value zero.

\emph{The precise trace and functional factorization.}
Write $m=[k_0:k_1]$ and $q_1=|k_1|$. The trace is
\[
 \operatorname{Tr}(x)=\sum_{j=0}^{m-1}x^{q_1^j}.
                                                               \tag{F.26.9}
\]
Its $q_1$th power equals itself, using
$x^{q_1^m}=x$, so its values lie in $k_1$.
It is $k_1$-linear and nonzero: its polynomial
has degree $q_1^{m-1}<|k_0|$ and leading coefficient
one, and therefore cannot vanish on all of $k_0$.
It is consequently onto $k_1$. Choose $\beta$
with $\operatorname{Tr}(\beta)=1$. For every
$a\ne0$,
$\operatorname{Tr}(a(\beta/a))=1$. Thus the pairing
$\operatorname{Tr}(ab)$ is nondegenerate, and
\[
 k_0\xrightarrow{\ \sim\ }
 \operatorname{Hom}_{k_1}(k_0,k_1),\qquad
 a\longmapsto(b\longmapsto\operatorname{Tr}(ab))
                                                               \tag{F.26.10}
\]
is an isomorphism: it is injective and both sides
have $k_1$-dimension $m$. This also covers characteristics
dividing $m$.

Choose a $k_0$-basis $v_1,\ldots,v_r$ of $k$.
For each $i$, (F.26.10) gives a unique $a_i\in k_0$
with $L(tv_i)=\operatorname{Tr}(a_it)$ for all $t\in k_0$.
Define
\[
 \ell:k\longrightarrow k_0,\qquad
 \ell\left(\sum_i t_iv_i\right)=\sum_i a_it_i,\qquad
                         L=\operatorname{Tr}\ell .
                                                               \tag{F.26.11}
\]
This is the required $k_0$-linear coefficient map.
Its image on $A$ is a nonzero $k_0$-subspace of
$k_0$, hence is all of $k_0$; its image on $U$
lies in $\ker\operatorname{Tr}$. The exact homomorphism
of finite affine groups is
\[
 P_\ell:k\rtimes D\longrightarrow k_0\rtimes D,\qquad
                         (b,d)\longmapsto(\ell(b),d).
                                                               \tag{F.26.12}
\]
The identity $\ell(db)=d\ell(b)$ proves its compatibility
with multiplication. All groups here are finite and
discrete, so it is continuous. Composing it with the
specified conjugation and with $\rho$ gives $\rho'$.
Equations (F.26.8)--(F.26.12) prove every image assertion
in (F.26.1), with the original scalar character retained.

\emph{The exact trace-cocycle criterion.}
For the data in (F.26.2), the restriction of $\chi$
to $K$ lies in $C\subset k_1^\times$. Thus composing
a crossed homomorphism with the $k_1$-linear trace
on $K$ proves precisely the target cocycle equation
in (F.26.2); continuity is preserved. Each principal
cocycle in (F.26.3) maps to zero. Because $C\ne\{1\}$,
the parameter map $b\mapsto(\chi-1)b$ is injective,
and this subspace has $k_1$-dimension $m-1$.

The representation $\rho'$ from (F.26.1) has a
translation cocycle $d'$ in $\ker\mathscr T_\chi$.
It is not principal: a principal affine representation
fixes a point of the field, whereas $\rho'(V)$ contains
every translation, including nonzero translations
which move every point. This gives an element outside
$\mathscr P_\chi$.

Conversely take $d\in\ker\mathscr T_\chi
\setminus\mathscr P_\chi$. The specified homomorphism
\[
 \rho_d:V\longrightarrow k_0\rtimes D,\qquad
                        g\longmapsto(d(g),\chi(g))             \tag{F.26.13}
\]
is continuous by the cocycle rule. Its translation
subgroup is stable under $D$ and hence is a
$k_0$-subspace of $k_0$. It cannot be zero. Otherwise
the scalar projection would identify $\rho_d(V)$
with the cyclic group $D$. A generator lift $(b,a)$,
$a\ne1$, fixes $b/(1-a)$, and all its powers fix that
same point. Then $d$ is principal; its vanishing
trace on $K$, together with $C\ne1$, makes the
parameter trace zero, contrary to the choice of $d$.
Thus $\rho_d(V)$ contains all translations of $k_0$.
Its $R$ image is normal and has nontrivial scalar
image $C$. Formula (F.26.7) therefore forces all
translations into that image, giving
$\rho_d(R)=k_0\rtimes C$. Its $K$ image lies in
$\ker\operatorname{Tr}\rtimes C$, by (F.26.2), and
is proper. This proves both directions of the criterion.

\emph{The coefficient splitting, with its inverse.}
The characters $\chi_j=\chi^{q_1^j}$, $0\leq j<m$,
are pairwise distinct: a nonidentity Frobenius power
fixing all their scalar values would fix their
generated field $k_0$. They all agree on $R$.
The exact coefficient isomorphism is
\[
 S:k_0\otimes_{k_1}k_0\xrightarrow{\ \sim\ }\prod_{j=0}^{m-1}k_0,
             \qquad b\otimes a\longmapsto
                       (a b^{q_1^j})_{j=0}^{m-1}.              \tag{F.26.14}
\]
Choose a multiplicative generator $\gamma$ of $k_0^\times$.
Its minimal polynomial $f$ over $k_1$ has degree $m$
and distinct roots $\gamma^{q_1^j}$. Identify the
tensor ring with $k_0[X]/f(X)$, with the first
$\gamma$ represented by $X$ and the second factor
supplying coefficients. Evaluation at these roots
is (F.26.14), and its inverse is
\[
 (a_j)_j\longmapsto
 \sum_{j=0}^{m-1} a_j
   \prod_{r\ne j}
     \frac{X-\gamma^{q_1^r}}{\gamma^{q_1^j}-\gamma^{q_1^r}}
                      \pmod{f(X)}.                            \tag{F.26.15}
\]
Every denominator is nonzero and the two composites
are the identity by evaluation and degree.

Give the first tensor factor the $\chi$ action.
Under $S$ its $j$th factor has the $\chi_j$ action.
A continuous cocycle on $V=\widehat F_7$ is specified
by its seven values: any such tuple defines a map
from the discrete free group into the finite affine
group and extends uniquely to its full completion.
Thus (F.26.14) gives the exact comparison
\[
 \begin{aligned}
 Z^1_{\rm cont}(V,k_{0,\chi})\otimes_{k_1}k_0
    &\xrightarrow{\ \sim\ }
       \bigoplus_{j=0}^{m-1}Z^1_{\rm cont}(V,k_{0,\chi_j}),\\
 d\otimes a&\longmapsto(a d^{q_1^j})_j .
 \end{aligned}                                                \tag{F.26.16}
\]
Expansion in a $k_1$-basis of $k_0$ likewise identifies
the scalar extension of the target of (F.26.2)
with $Z^1_{\rm cont}(K,k_{0,\chi|K})$. Coordinate
projections preserve cocycles and continuity, so
this target comparison holds for the actual
profinite subgroup $K$ as well. On simple tensors,
the trace formula proves that $\mathscr T_\chi\otimes1$
is exactly
\[
                  (d_j)_j\longmapsto\sum_j d_j|K.              \tag{F.26.17}
\]
The scalar extension of $\mathscr P_\chi$ is
$\{(b_j(\chi_j-1))_j:\sum_jb_j=0\}$, since
the trace becomes coordinate summation under $S$.
Extension by a field preserves kernels and detects
a nonzero vector space by a basis. This proves
the claimed equivalent kernel description and
the rank target $7m-(m-1)=6m+1$.

\emph{The exhaustive field-of-order-nine character set.}
Use the actual negative coordinates of F.24.4 and
write $\mathbb F_9=\mathbb F_3[\iota]/(\iota^2+1)$.
Every element is uniquely $a+b\iota$ with
$a,b\in\mathbb F_3$; the polynomial has no root
in $\mathbb F_3$, so this quotient is a field.
For any scalar field $k_0$ of order nine, choose
an element $\iota_0$ of order four in $k_0^\times$.
It satisfies $\iota_0^2=-1$. The explicit field map
$a+b\iota\mapsto a+b\iota_0$ is an isomorphism:
its source is a field, it preserves one, and both
fields have nine elements. Its inverse takes the
unique two coefficients in this basis. Apply that
inverse to both affine coordinates to obtain the
corresponding isomorphism of affine groups. It
preserves the original source representation and
its two subgroup images through this specified
coefficient map. Either choice of $\iota_0$ is
covered by the complete character enumeration.
The field code $a+3b$ in the programs records these
two coordinates exactly. Multiplication is
\[
 (a+b\iota)(c+d\iota)=(ac-bd)+(ad+bc)\iota,
 \qquad
 \operatorname{Tr}(a+b\iota)=2a,\qquad
 \operatorname{Tr}(\iota(a+b\iota))=b .
                                                               \tag{F.26.18}
\]
Thus multiplication of cocycles by $\iota$ is the
exact coefficient isomorphism from the imaginary-part
kernel computed below to the trace kernel, with
inverse multiplication by $-\iota$. It sends
real principal parameters to the trace-zero
principal parameters $\mathbb F_3\iota$.

F.25 covers trivial $C$ and equality of scalar
fields. The remaining scalar-field-nine data
have $k_1=\mathbb F_3$, so $C=\{1,-1\}$.
The four deleted character values are in
$\{1,-1\}$ with at least one equal to $-1$.
Conversely those values give exactly this $C$,
because the four deleted generators normally
generate $R$ and the character is abelian.
The three retained values lie in
$(\mathbb F_9^\times)^3$ and are not all in
$(\mathbb F_3^\times)^3$. These conditions give
$k_0=\mathbb F_9$, because the extension of
$\mathbb F_3$ has prime degree two. The exhaustive
character count is therefore
\[
                   (8^3-2^3)(2^4-1)=7\,560 .                  \tag{F.26.19}
\]
The first producer enumerates every such tuple.

For a fixed tuple $t$, write a cocycle-value vector
as $v=x+\iota y$, with $x,y\in\mathbb F_3^7$,
and put $\widetilde v=(x,y)^{\mathsf T}$.
For each of the sixty-four actual F.24 source rows
$a(t)=a_{\rm r}+\iota a_{\rm i}$ the exact condition
that its value have imaginary part zero is
\[
       (a_{\rm i},a_{\rm r})\widetilde v=0 .                   \tag{F.26.20}
\]
The seven-coordinate vector $t-\mathbf1$ with real
principal parameter gives a nonzero solution.
It is nonzero because some deleted value equals
$-1$. Every row arises from an actual element
of $K$, whose character lies in $\mathbb F_3$,
so this principal line is contained in the
kernel of all the rows and of every further
source conjugate.

The first exact row reduction gives the following
complete distribution:
\[
 \begin{array}{c|rrr}
 \text{rank on fourteen columns}&11&12&13\\\hline
 \text{character tuples}&8&396&7156 .
 \end{array}                                                  \tag{F.26.21}
\]
For the $7\,156$ rank-thirteen tuples, the kernel
is precisely the stated principal line.
The other $404$ tuples are retained explicitly
and passed to the second producer.

\emph{All eleven actual source actions on the remaining tuples.}
For each of these $404$ tuples, evaluate all eleven
matrices $A_i$ from F.24.9 by the same exact
Laurent evaluation. Write $A_i=A_{i,\rm r}
+\iota A_{i,\rm i}$. Its actual action on the
fourteen prime-field coordinates is
\[
 \widetilde A_i=
 \begin{pmatrix}
 A_{i,\rm r}&-A_{i,\rm i}\\
 A_{i,\rm i}& A_{i,\rm r}
 \end{pmatrix}.                                               \tag{F.26.22}
\]
Start with the eight original negative seed
imaginary-part rows. At each round replace their
row space $W$ by
\[
                 W+\sum_{i=0}^{10}W\widetilde A_i .
                                                               \tag{F.26.23}
\]
Every row lies in the $\mathbb F_3$-linear span of the
imaginary-part evaluation rows of actual source
conjugates of the seeds.
The pure source action preserves the character,
and its action on cocycles is exactly the
precomposition matrix proved in F.24.9.
The successive exact ranks are
\[
 \begin{array}{c|rrrr}
 \text{stage}&\multicolumn{4}{c}{\text{rank : number of tuples}}\\\hline
 \text{eight seeds}&5:80&6:284&7:2&8:38\\
 \text{one round}&12:50&13:354&&\\
 \text{two rounds}&13:404&&&\\
 \text{third round, stable}&13:404&&&
 \end{array}                                                  \tag{F.26.24}
\]
In particular two rounds already supply rank
thirteen for every remaining tuple. The third
round verifies invariance under all eleven
matrices. Their actions are invertible, since
they come from the actual source automorphisms.
An invariant subspace for an invertible linear
map on a finite-dimensional space is also
invariant under its inverse. Hence this stable
row space accounts for both directions of all
source generators.

All arithmetic here is in the explicitly described
finite field or in $\mathbb F_3$. Every nonzero
Laurent variable is evaluated using its powers
modulo eight. The field table, all $7\,560$
character choices, all $64$ row evaluations,
the $404$ retained tuples, the eleven actual
source matrices and the complete rank histories
are specified by the accompanying producers
and their fixed input files. Gaussian elimination
uses the exact inverses one and two in
$\mathbb F_3$. The independent verification imports
neither producer. It reconstructs these matrices
from the fixed universal source rows and actions,
using a separate multiplicative-log Laurent evaluation,
and computes the ranks by scalar incremental elimination
over $\mathbb F_3$. It verifies the entire character set,
all principal vectors, all eleven source actions,
every closure image, the stable row spaces and all
nine coefficient identities in (F.26.18). The universal
source data themselves are the audited literal source
evaluations of F.24.

It follows that every continuous cocycle whose
imaginary part vanishes on $K$ has values
$b(t-\mathbf1)$ for a unique $b\in\mathbb F_3$.
Indeed all the actual rows (F.26.20) and their
source conjugates must vanish on it, so the
proved ranks give precisely that one-dimensional
kernel. Conversely those principal cocycles
do have imaginary part zero on $K$.
Multiplication by $\iota$ in (F.26.18), and its
specified inverse, prove
$\ker\mathscr T_\chi=\mathscr P_\chi$ for all
$7\,560$ remaining field-of-order-nine characters.

The criterion proved in (F.26.1)--(F.26.13)
now excludes every proper marginal whose
scalar field has order nine, even when its
original coefficient field $k$ is larger.
For a representation into
$\operatorname{Aff}_1(\mathbb F_9)$ its scalar
field is either $\mathbb F_3$ or $\mathbb F_9$;
F.25 handles the former and the present argument
handles the latter. Finally, precomposition
with the actual continuous isomorphism
$\mathcal B_+:V_-\to V_+$ of F.23.13, and
its inverse, transports each positive
representation to the proved negative case,
carrying both $K$ and $R$ exactly as stated
there. This proves (F.26.4) for both signs.
\end{proof}

\begin{claim}{F.27}{Proved}{All finite affine marginal images on an exact two-parameter character locus}
Use the actual negative free basis $X_1,\ldots,X_7$ of
F.24 and the based comparison $\mathcal B_+:V_-\to V_+$
of F.23.13. Put $Y_i^-=X_i$ and
$Y_i^+=\mathcal B_+(X_i)$. For every finite field $k$
and continuous $\rho:V_\pm\to\operatorname{Aff}_1(k)$,
let $\chi$ be its scalar character. The character locus
\[
 \chi(Y_1^\pm)=\chi(Y_4^\pm)=\chi(Y_5^\pm)
       =\chi(Y_6^\pm)=1,\qquad
 \chi(Y_2^\pm)\chi(Y_3^\pm)\chi(Y_7^\pm)=1
                                                               \tag{F.27.1}
\]
has the exact image equality
\[
                         \rho(\mathcal K_\pm)=\rho(R_\pm).
                                                               \tag{F.27.2}
\]
There is no restriction on the characteristic, extension
degree, or size of the original coefficient field.

More precisely, write a negative character in this locus as
$(1,s,t,1,1,1,c^{-1})$, where $c=st$, and put
$k_0=\mathbb F_p(s,t)$, $k_1=\mathbb F_p(c)$.
For $c\ne1$ and $k_1\subsetneq k_0$, let
$q_1=|k_1|$, $m=[k_0:k_1]$ and
$\chi_j=\chi^{q_1^j}$, $0\leq j<m$.
Writing $K=\mathcal K_-$ and $V=V_-$, the actual map
\[
 \bigoplus_{j=0}^{m-1}H^1_{\rm cont}(V,k_{0,\chi_j})
   \longrightarrow H^1_{\rm cont}(K,k_{0,\chi|K}),\qquad
                 ([d_j])_j\longmapsto\sum_j[d_j|K]
                                                               \tag{F.27.3}
\]
is injective. At the cocycle level, the kernel of
the summed restriction is exactly
\[
 \left\{\bigl(b_j(\chi_j-1)\bigr)_j:
                  b_j\in k_0,\ \sum_jb_j=0\right\}.
                                                               \tag{F.27.4}
\]
Thus its rank is $6m+1$ on $7m$ cocycle coordinates.
The proof uses the literal eleven F.23 source
automorphisms and specified source conjugates of
one of its eight negative seeds.
\end{claim}
\passage{Part F root-adjoining intersection question; Section 30, folios 194-196, scans 383-387; Section 31, folios 198-199 verso, scans 391-394; original F7/F19 residual subgroup-membership reduction}


\begin{proof}
\emph{The actual scalar fields and all exceptional cases.}
Work first with the negative groups and write
$K\leq R\leq V$. The four deleted generators
$X_4,\ldots,X_7$ normally generate $R$.
The normal generation and actual marginal maps
of F.21--F.22, as used in F.25.5 and F.26.5, give
\[
 \begin{gathered}
 C=\chi(K)=\chi(R)=\langle c^{-1}\rangle,\qquad
 D=\chi(V)=\langle s,t\rangle,\\
 k_1=\mathbb F_p(C)=\mathbb F_p(c),\qquad
 k_0=\mathbb F_p(D)=k_1(s)=k_1(t),\qquad c=st.
 \end{gathered}                                                \tag{F.27.5}
\]
Indeed the last equalities follow from $t=c/s$ and
$s=c/t$. When $c=1$, the $R$ image consists of
translations and is abelian; the actual normal
generation therefore gives (F.27.2). For $c\ne1$
and $k_0=k_1$, F.25 gives the same equality.
It remains to treat $c\ne1$ and $k_1\subsetneq k_0$.
Here $s,t\notin k_1$, by (F.27.5), and consequently
\[
             s\,t\,(s-1)(t-1)(t+1)(c-1)\ne0.
                                                               \tag{F.27.6}
\]
The factors $s,t$ are nonzero character values.
The elements $1,-1$ belong to $k_1$, so neither
can equal $s$ or $t$. This includes characteristic
two, where $t+1=t-1$. The case $c=1$ has already
been settled, so every factor in (F.27.6) is justified.

\emph{The exact coefficient map and principal quotient.}
All universal word calculations in this proof are
performed on the discrete free group on the seven
specified generators, with its canonical dense map
to $V=\widehat F_7$. The coefficient homomorphism
from the F.24 Laurent ring is
\[
 \begin{aligned}
 \mathbb Z[t_1^{\pm1},\ldots,t_7^{\pm1}]&\longrightarrow
                  A=\mathbb Z[s^{\pm1},t^{\pm1}],\\
 (t_1,t_2,t_3,t_4,t_5,t_6,t_7)
           &\longmapsto(1,s,t,1,1,1,(st)^{-1}).
 \end{aligned}                                                \tag{F.27.7}
\]
Thus a monomial with exponent vector $(a_i)$ maps
to $s^{a_2-a_7}t^{a_3-a_7}$. The universal
derivation on free words has $D(X_i)=e_i$ and
$D(uv)=D(u)+\chi(u)D(v)$. Evaluation of its
coefficients in a finite field gives the actual
cocycle value on that word. Every seven-tuple
of generator values defines a homomorphism from
the discrete free group to the finite affine
group and extends uniquely to its full profinite
completion. Consequently the continuous cocycle
space is exactly $k_0^7$. No continuous map to
an infinite Laurent ring is being used.

Put $\delta=(1,s,t,1,1,1,c^{-1})^{\mathsf T}
-\mathbf1$. Its last entry is
$\delta_7=(1-c)/c\ne0$.
The quotient by principal cocycles has the explicit
projection and section
\[
 \begin{aligned}
 P(v)&=\bigl(v_r-\delta_r\delta_7^{-1}v_7\bigr)_{r=1}^6,
       &S(w)&=(w_1,\ldots,w_6,0)^{\mathsf T},\\
 PS&=I_6,
       &SP(v)&=v-\delta\delta_7^{-1}v_7.
 \end{aligned}                                                \tag{F.27.8}
\]
These formulas identify
$Q_\chi=H^1_{\rm cont}(V,k_{0,\chi})$ with $k_0^6$.
Universally they are defined over $A[(1-st)^{-1}]$;
all finite-field evaluations here respect that
localization by (F.27.6).

Let $B_l(v)=u_l\alpha_l(v)u_l^{-1}$, $0\leq l\leq10$,
be the actual old-source automorphisms of F.23.
They preserve $K$, including under their inverses,
and fix the character because their action on
the abelianization of $V$ is the identity.
Let $A_l$ be their F.24 precomposition matrices
after (F.27.7), so the $i$th row is $D(B_l(X_i))$.
Then $A_l\delta=\delta$ and their actual cohomology
actions are
\[
                           \overline A_l=PA_lS.
                                                               \tag{F.27.9}
\]
The source indices are exactly
$0=\texttt{old0}$,
$1,\ldots,5=\texttt{old\_1\_0},\ldots,\texttt{old\_1\_4}$
and $6,\ldots,10=\texttt{old\_-1\_0},\ldots,
\texttt{old\_-1\_4}$.

\emph{A basis with explicit inverse and its actual source action.}
For $P_0=X_1\cdots X_7$, the character is one,
and its derivative row is $(1,1,s,c,c,c,c)$.
Every pure $\alpha_l$ fixes this word. Hence
evaluation on $P_0$ induces the cohomology
covector $b=(1,1,s,c,c,c)$, with
$b\overline A_l=\chi(u_l)b$.
For each $r\in\{1,4,5,6\}$, the pure Artin
action writes $\alpha_l(X_r)=a_{lr}X_ra_{lr}^{-1}$.
Since $\chi(X_r)=1$, its actual conjugated source
action gives
$d(B_l(X_r))=\chi(u_la_{lr})d(X_r)$.
Thus $e_1^*,e_4^*,e_5^*,e_6^*$ are four more
actual quotient covectors. Their common kernel
together with $b$ is the line generated by
\[
                      v_0=(0,1,-s^{-1},0,0,0)^{\mathsf T}.
                                                               \tag{F.27.10}
\]
Use the basis matrix $J$ with columns
\[
       v_0,\quad e_2,\quad e_5-ce_2,\quad e_4-ce_2,
                       \quad e_1-e_2,\quad e_6-ce_2.
                                                               \tag{F.27.11}
\]
Its inverse is the explicit map
\[
 J^{-1}(w)=(-sw_3,\,bw,\,w_5,\,w_4,\,w_1,\,w_6)^{\mathsf T}.
                                                               \tag{F.27.12}
\]
For completeness, $J(z_0,\ldots,z_5)^{\mathsf T}$ has
coordinates
\[
 \bigl(z_4,\ z_0+z_1-cz_2-cz_3-z_4-cz_5,
          \ -s^{-1}z_0,\ z_3,\ z_2,\ z_5\bigr)^{\mathsf T}.
                                                               \tag{F.27.13}
\]
Substitution in either direction gives the identity,
so these are actual mutually inverse coefficient maps.

The five displayed quotient covectors give the
following form for the source actions in this basis:
\[
 T_l=J^{-1}\overline A_lJ
       =\begin{pmatrix}\nu_l&r_l\\0&\Lambda_l\end{pmatrix},
 \qquad \Lambda_l\text{ diagonal},\qquad
 (\nu_0,\ldots,\nu_{10})=(1,s,t,1,1,1,s,t,1,1,1).
                                                               \tag{F.27.14}
\]
The last character is computed from the literal
source words. More explicitly, for the seven-coordinate
lift $\widetilde v_0=S(v_0)$ the verified identities are
\[
 \begin{gathered}
 A_l\widetilde v_0=\nu_l\widetilde v_0+\pi_l\delta,\\
 \pi_1=-c,\qquad \pi_2=\pi_7=ct,\qquad
 \pi_6=s(c-t-1),\qquad
 \pi_l=0\quad(l\notin\{1,2,6,7\}).
 \end{gathered}                                                \tag{F.27.15}
\]
The certificate records every entry of all eleven
$A_l$, $\overline A_l$ and $T_l$, with the coefficient
maps (F.27.7)--(F.27.13).

\emph{An integral determinant forcing the essential line.}
Take the additive matrix commutators
$E_j=T_1T_j-T_jT_1$ for $j=2,4,5,6,7$.
They kill the first coordinate line and have image
in that line, since the lower diagonal blocks
in (F.27.14) commute. Their first-row entries on
the other five coordinates, in that order, form
the exact matrix
\[
 \begin{gathered}
 gH,\qquad g=(s-1)(t-1),\\[-2pt]
 H=\begin{pmatrix}
 0&0&0&-(t+1)/t&st\\
 0&s&-1/t&0&0\\
 0&-1/t&0&st+1&0\\
 (s-1)/(st^2)&s(s-1)&-(s-1)/t&s-1&0\\
 (t+1)/(st^2)&0&0&0&st
 \end{pmatrix}.
 \end{gathered}                                                \tag{F.27.16}
\]
Direct determinant expansion gives the integral
Laurent identity
\[
 \det H=\frac{(s-1)(t-1)(t+1)}{t^4},\qquad
 \det(gH)=\frac{(s-1)^6(t-1)^6(t+1)}{t^4}.
                                                               \tag{F.27.17}
\]
All its factors are nonzero after every finite-field
evaluation under consideration, by (F.27.6).
The identities have integer coefficients and only
signed Laurent monomial denominators; no prime
is removed from the calculation.

Let $L_\chi=k_0v_0\leq Q_\chi$. It is invariant
by (F.27.14). Every nonzero subspace $W\leq Q_\chi$
invariant under the eleven actual source operators
contains $L_\chi$. Indeed, in the $J$ coordinates
the common kernel of the five $E_j$ is precisely
the first coordinate line, by (F.27.17).
If $W$ is contained in this line, it equals it.
Otherwise choose a vector of $W$ with nonzero
last five coordinates. Some $E_j$ sends it to
a nonzero vector of the first coordinate line;
this vector remains in $W$. Since the line has
dimension one, it is contained in $W$.
This proves the required essential-socle statement
by explicit source endomorphisms.

\emph{Actual seed evaluations and all Frobenius factors.}
Let $d_0$ be the continuous cocycle with generator
values $\widetilde v_0$. Take the literal negative
seed $a=u_{\texttt{cross\_4\_-1\_1}}$ of index $25$.
The actual seed calculation gives
\[
                   \chi(a)=c^{-1},\qquad d_0(a)=c-1.
                                                               \tag{F.27.18}
\]
Both are computed by the same decorated braid
words as in F.23; all eight seed rows are audited.
Since $a\in K$ and $B_1$ preserves $K$, put
\[
          \kappa=B_1(a)a^{-1}\in K,
          \qquad \kappa_r=B_1^r(\kappa)\in K\quad(r\geq0).
                                                               \tag{F.27.19}
\]
Their characters are one. Using $\pi_1=-c$ in
(F.27.15) gives
$d_0(B_1(a))=(s+1)(c-1)$, hence
$d_0(\kappa)=s(c-1)$.
For every character-one element $h$, (F.27.15)
also gives $d_0(B_1(h))=s\,d_0(h)$: the principal
correction evaluates to zero there. Induction yields
\[
                         d_0(\kappa_r)=s^{r+1}(c-1).
                                                               \tag{F.27.20}
\]
In row coordinates the first two evaluations
are $D(a)(A_1-I_7)$ and $D(a)(A_1-I_7)A_1$;
the independent audit verifies their character-one
principal vanishing and the stated values.

Put $s_j=s^{q_1^j}$ and $t_j=t^{q_1^j}$.
Since $c\in k_1$, $s_jt_j=c$ for every $j$.
The field equality $k_0=k_1(s)$ proves that
$s_0,\ldots,s_{m-1}$ are distinct. Indeed an
equality would make a nonidentity power of the
$k_1$-Frobenius fix $s$, and hence its generated
field $k_0$; the powers before $m$ are distinct
automorphisms of the field of $q_1^m$ elements.

The coefficient map on cocycles and on cohomology is
\[
                  d\longmapsto d^{q_1^j},
 \qquad k_{0,\chi}\longrightarrow k_{0,\chi_j}.
                                                               \tag{F.27.21}
\]
It is additive and semilinear, preserves the
cocycle rule, carries principals to principals,
and preserves continuity because the field is finite.
Its inverse is the power $q_1^{m-j}$, with the
identity used at $j=0$; the two composites are
the identity on $k_0$. It commutes with the actual
source precomposition operators. Thus the explicit
bases, inverse maps, commutator identities and
essential lines transport to every $\chi_j$.
Write $Q_j=H^1_{\rm cont}(V,k_{0,\chi_j})$ and
$L_j=k_0[d_j]\leq Q_j$, where $d_j=d_0^{q_1^j}$.
Equation (F.27.20) becomes
\[
              d_j(\kappa_r)=a_js_j^r,
                     \qquad a_j=s_j(c-1)\ne0.
                                                               \tag{F.27.22}
\]

The characters $\chi_j$ all agree on $K$, since
$\chi(K)=C\subset k_1^\times$. Evaluation at
$\kappa_0,\ldots,\kappa_{m-1}$ therefore defines
a map on the common target cohomology space:
principals vanish on each character-one element.
On the summed restrictions of the $m$ essential
lines its matrix is
\[
                  \bigl(a_js_j^r\bigr)_{
                        0\leq r,j<m},\qquad
 \det=\left(\prod_j a_j\right)
                   \prod_{j<l}(s_l-s_j)\ne0.
                                                               \tag{F.27.23}
\]
The coefficient inverse is explicit. Write
\[
 \prod_{l\ne j}\frac{X-s_l}{s_j-s_l}
       =\sum_{r=0}^{m-1}h_{jr}X^r.
 \quad\text{Then}\quad
       z_j=a_j^{-1}\sum_{r=0}^{m-1}h_{jr}y_r
                                                               \tag{F.27.24}
\]
is the inverse to
$y_r=\sum_j a_js_j^r z_j$.
Evaluation of the displayed interpolation polynomial
at every $s_l$ proves both inverse identities.
Every denominator in (F.27.24) is nonzero by
the proved distinctness and (F.27.22).

\emph{Summed restriction, the trace kernel and the affine image.}
The map (F.27.3) commutes with all eleven source
operators: on both source and target they are
the actual maps $[d]\mapsto[d\circ B_l]$.
Consequently its kernel is an invariant subspace
of the finite-dimensional direct sum $\bigoplus_jQ_j$.
Suppose that kernel were nonzero and choose a
nonzero invariant subspace $T$ inside it of
smallest dimension. Each coordinate projection
$T\to Q_j$ has invariant kernel and image.
Minimality makes every nonzero projection
injective, with an image having no proper
nonzero invariant subspace. The essential-line
result puts $L_j$ inside that image; therefore
the image equals $L_j$. All projections of $T$
thus lie in their respective $L_j$, and
$T\leq\bigoplus_jL_j$. But (F.27.23)--(F.27.24)
make evaluation after summed restriction injective
on this entire direct sum of lines. This
contradicts $T$ lying in the kernel and proves
(F.27.3).

Zero summed restriction of cocycles now implies
that each $d_j$ is principal. Their restrictions
have the common character $\chi|K$, so their
sum is $(\chi|K-1)\sum_jb_j$. Since $C\ne1$,
this vanishes exactly when $\sum_jb_j=0$.
This proves (F.27.4), including both directions.
The seven-coordinate description of each
cocycle space gives rank $7m-(m-1)=6m+1$.

For the trace map of F.26.2, a cocycle $d$ in
its kernel yields the tuple $(d^{q_1^j})_j$,
whose summed restriction is exactly
$\operatorname{Tr}_{k_0/k_1}(d|K)=0$.
Thus $d=b(\chi-1)$ by the result just proved.
Trace vanishing on $K$, together with
$\chi(K)=C\ne1$, gives $\operatorname{Tr}(b)=0$.
Conversely these principal cocycles have zero
trace restriction. Hence the precise trace
kernel is $\mathscr P_\chi$ of F.26.3.

Every proper affine marginal image in the current
locus would, by the explicit conjugation and
coefficient homomorphism of F.26.1--12, give a
nonprincipal element of that trace kernel over
its scalar field $k_0$. Those maps retain the
original scalar character, so all the displayed
locus identities and fields remain the same.
The proved trace kernel excludes such an element.
This establishes (F.27.2) in the negative case,
for every original finite coefficient field $k$.

\emph{Certificate and transport.}
The symbolic producer applies exactly (F.27.7)
to the fixed, previously audited F.24 universal
rows. The separate verifier imports neither
that producer nor a computer algebra package.
It evaluates all eleven actual old-source words
and eight actual negative seed words by generic
affine multiplication over the integral Laurent
ring $A$. It checks the inverse basis maps,
the principal-quotient action after clearing
$\delta_7$, all source and socle identities,
the actual seed evaluations, and (F.27.17)
by its $120$ signed permutation terms.
The verifier performs $699$ recorded checks.
Thus the determinant certificate is an exact
integer identity and applies in every characteristic;
its nonzero finite-field factors are separately
proved in (F.27.6). The evidence manifest binds
both programs, their outputs and stdout, the
original braid input and all accepted dependencies.

Finally, a positive representation in (F.27.1)
precomposes with the actual continuous isomorphism
$\mathcal B_+:V_-\to V_+$ to give exactly the
negative character just treated. By F.23.13,
this isomorphism and its inverse carry both
$\mathcal K_-$ and $R_-$ onto their positive
counterparts. The equality of the two actual
images therefore transports through these
specified maps, proving (F.27.2) for both signs.
\end{proof}

\begin{claim}{F.28}{Proved}{All finite affine marginal images on the three-parameter product-one locus}
Retain the actual bases $Y_i^-=X_i$ and
$Y_i^+=\mathcal B_+(X_i)$ of F.27, with
$\mathcal B_+$ the based comparison of F.23.13.
For every finite field $k$ and continuous
$\rho:V_\pm\to\operatorname{Aff}_1(k)$ with scalar
character $\chi$, the locus
\[
 \chi(Y_4^\pm)=\chi(Y_5^\pm)=\chi(Y_6^\pm)=1,
 \qquad
 \chi(Y_1^\pm)\chi(Y_2^\pm)\chi(Y_3^\pm)
                         \chi(Y_7^\pm)=1
                                                               \tag{F.28.1}
\]
has the exact image equality
\[
                         \rho(\mathcal K_\pm)=\rho(R_\pm).
                                                               \tag{F.28.2}
\]
This includes every characteristic, every finite
scalar-field extension degree and every larger
original coefficient field.

For a negative character
$(r,s,t,1,1,1,c^{-1})$, $c=rst\ne1$, put
$k_1=\mathbb F_p(c)$ and $k_0=\mathbb F_p(r,s,t)$.
In the case $k_1\subsetneq k_0$, let
$q_1=|k_1|$, $m=[k_0:k_1]$ and
$\chi_j=\chi^{q_1^j}$, $0\leq j<m$.
The summed continuous $H^1$ restriction for these
$m$ characters is injective, and its cocycle
kernel is exactly the principal tuples with
parameter sum zero. Its rank is therefore
$6m+1$ on $7m$ cocycle coordinates.
\end{claim}
\passage{Part F root-adjoining intersection question; Section 30, folios 194-196, scans 383-387; Section 31, folios 198-199 verso, scans 391-394; original F7/F19 residual subgroup-membership reduction}


\begin{proof}
\emph{The actual fields and the already proved cases.}
Use the negative groups $K\leq R\leq V$ first.
The deleted generators $X_4,\ldots,X_7$ and the
actual normal-generation maps of F.21--F.22 give
\[
 \begin{gathered}
 C=\chi(K)=\chi(R)=\langle c^{-1}\rangle,
 \qquad D=\chi(V)=\langle r,s,t\rangle,\\
 k_1=\mathbb F_p(c),\qquad
 k_0=\mathbb F_p(r,s,t)=k_1(r,t),\qquad c=rst.
 \end{gathered}                                                \tag{F.28.3}
\]
The last equality follows from $s=c/(rt)$;
all character values are nonzero.
For $c=1$, the $R$ image is an abelian translation
group, and actual normal generation gives (F.28.2).
For $c\ne1$ and $k_0=k_1$, F.25 gives it.
For $r=1$, the character is exactly
$(1,s,t,1,1,1,(st)^{-1})$ in the unchanged actual
basis, so F.27 gives both the image equality
and the stated summed-restriction result.
At the coefficient-ring level this last comparison
is the homomorphism $r\mapsto1$, $s\mapsto s$,
$t\mapsto t$, sending $1-rst$ to $1-st$ in
the corresponding localized Laurent rings.
Thus it respects the precise maps used in F.27.

It remains to treat $c\ne1$, $r\ne1$ and
$k_1\subsetneq k_0$. These conditions imply
only the nonzero factors actually used below:
$r,s,t,c,c-1,r-1$. No condition on $s-1$,
$t-1$, $r+1$ or the characteristic is imposed.

\emph{The exact invariant inclusion and quotient maps.}
The literal coefficient homomorphism from F.24 is
\[
 \begin{aligned}
 \mathbb Z[t_1^{\pm1},\ldots,t_7^{\pm1}]&\longrightarrow
                   A=\mathbb Z[r^{\pm1},s^{\pm1},t^{\pm1}],\\
 (t_1,t_2,t_3,t_4,t_5,t_6,t_7)
          &\longmapsto(r,s,t,1,1,1,(rst)^{-1}).
 \end{aligned}                                                \tag{F.28.4}
\]
A monomial with exponent vector $(a_i)$ maps to
$r^{a_1-a_7}s^{a_2-a_7}t^{a_3-a_7}$.
These are identities on the specified discrete
free words. Their finite-field cocycle evaluations
extend uniquely to $V=\widehat F_7$ by the finite
affine homomorphisms, exactly as in F.27.7--8.

Let $\delta=(r,s,t,1,1,1,c^{-1})^{\mathsf T}
-\mathbf1$, with $\delta_7=(1-c)/c\ne0$.
The actual continuous cohomology space
$Q_\chi=H^1_{\rm cont}(V,k_{0,\chi})$ is $k_0^6$
through the projection and section
\[
 \begin{aligned}
 P(v)&=(v_i-\delta_i\delta_7^{-1}v_7)_{i=1}^6,
       &S(w)&=(w_1,\ldots,w_6,0)^{\mathsf T},\\
 PS&=I_6,
       &SP(v)&=v-\delta\delta_7^{-1}v_7.
 \end{aligned}                                                \tag{F.28.5}
\]
For the eleven actual F.23 source automorphisms
$B_l$, $0\leq l\leq10$, both directions preserve
$K$ and fix the character, as in F.27.
Let $A_l$ be their
specialized cocycle matrices. Then $A_l\delta=\delta$
and $\overline A_l=PA_lS$ are the actual source
operators on $Q_\chi$.

The boundary word $P_0=X_1\cdots X_7$ has
character one and derivative
$(1,r,rs,c,c,c,c)$. Pure Artin monodromy fixes
it, giving the quotient covector
$b=(1,r,rs,c,c,c)$.
The character-one meridians $X_4,X_5,X_6$ give
the three quotient covectors $e_4^*,e_5^*,e_6^*$,
as in F.27: the scalar of their actual conjugator
multiplies their cocycle value.
The four maps together define
\[
 \begin{aligned}
 \pi:Q_\chi&\longrightarrow k_0^4,
       &\pi(v)&=(bv,v_4,v_5,v_6)^{\mathsf T},\\
 \iota:k_0^2&\longrightarrow Q_\chi,
       &\iota(x,y)&=(-rx-rsy,x,y,0,0,0)^{\mathsf T}.
 \end{aligned}                                                \tag{F.28.6}
\]
Their kernel-image identity is
$\ker\pi=\operatorname{im}\iota=:W_\chi$.
The coordinate projection $(v_2,v_3)$ is inverse
to $\iota$ on this subspace. A coefficient section
of $\pi$ is
\[
 \sigma(z_\partial,z_4,z_5,z_6)
       =(z_\partial-c(z_4+z_5+z_6),0,0,z_4,z_5,z_6)^{\mathsf T}.
                                                               \tag{F.28.7}
\]
One has $\pi\sigma=I_4$ and
$\iota(v_2,v_3)+\sigma\pi(v)=v$.
Thus the basis matrix with columns
\[
 \iota(1,0),\quad\iota(0,1),\quad e_1,
             \quad e_4-ce_1,\quad e_5-ce_1,\quad e_6-ce_1
                                                               \tag{F.28.8}
\]
has inverse $v\mapsto(v_2,v_3,bv,v_4,v_5,v_6)$.
The displayed identities prove both inverse
coefficient maps explicitly.

In this basis the actual source matrices have form
\[
 \begin{pmatrix}G_l&E_l\\0&D_l\end{pmatrix},\qquad
 D_l=\operatorname{diag}
  (\lambda_{\partial,l},\lambda_{4,l},
                         \lambda_{5,l},\lambda_{6,l}).
                                                               \tag{F.28.9}
\]
Equivalently,
$\overline A_l\iota=\iota G_l$ and
$\pi\overline A_l=D_l\pi$.
All extension matrices $E_l$ are retained in
the certificate. These maps give the actual
invariant subspace $W_\chi$ and the four actual
one-dimensional quotient coordinates.

The literal source-word calculation gives the
following comparison table:
\[
 \begin{array}{c|c|cccc}
 l&G_l&\lambda_{\partial,l}&\lambda_{4,l}
                         &\lambda_{5,l}&\lambda_{6,l}\\\hline
 0&rI_2&1&r&r&r\\
 7&tI_2&t^{-1}&t&t&ct\\
 8&I_2&1&c&1&t^{-1}\\
 9&I_2&1&t^{-1}&c&1
 \end{array}
 \qquad \lambda_{\partial,2}=rt.
                                                               \tag{F.28.10}
\]
Every source index here is the original index
beginning at zero: $0=\texttt{old0}$,
$2=\texttt{old\_1\_1}$,
$7=\texttt{old\_-1\_1}$,
$8=\texttt{old\_-1\_2}$ and
$9=\texttt{old\_-1\_3}$.
The independent verifier checks all entries
against these exact labelled source words.

There are also explicit actual $K$ evaluations
on $W_\chi$. Let $a=u_{\texttt{cross\_4\_-1\_1}}$
be seed $25$, with $\chi(a)=c^{-1}$.
For $h_2=B_2(a)a^{-1}$ and $h_7=B_7(a)a^{-1}$,
both characters are one. A representative with
generator values $S\iota(x,y)$ satisfies
\[
 \bigl(d(h_2),d(h_7)\bigr)=((1-c)x,s(c-1)y).
                                                               \tag{F.28.11}
\]
Both evaluations therefore descend to cohomology.
Their inverse coefficient map is
$(u,v)\mapsto((1-c)^{-1}u,[s(c-1)]^{-1}v)$.
The determinant $-s(c-1)^2$ is nonzero in every
remaining case. These identities are also
verified from the actual seed and source words.

\emph{Frobenius transport and the classification needed for comparison.}
For $0\leq j<m$, put
$r_j=r^{q_1^j}$, $s_j=s^{q_1^j}$ and
$t_j=t^{q_1^j}$. Then $r_js_jt_j=c$.
The semilinear coefficient map
\[
             d\longmapsto d^{q_1^j},\qquad
             k_{0,\chi}\longrightarrow k_{0,\chi_j}
                                                               \tag{F.28.12}
\]
and its inverse, the power $q_1^{m-j}$ with the
identity at $j=0$, preserve cocycles, principals
and continuity. They commute with all actual
source precomposition maps. Consequently all
the preceding inclusions, projections, inverse
bases and comparison entries transport to
$Q_j=H^1_{\rm cont}(V,k_{0,\chi_j})$ and its
invariant subspace $W_j$ by the substitutions
$(r,s,t)\mapsto(r_j,s_j,t_j)$.
The characters agree on $K$, since
$\chi(K)=C\subset k_1^\times$.

The pairs $(r_j,t_j)$ are distinct. Equality of
two pairs would make a nonidentity $k_1$-Frobenius
power fix the generators of $k_0=k_1(r,t)$.
It would then fix every element of $k_0$;
the nonzero polynomial $X^{q_1^a}-X$, with
$0<a<m$, has degree smaller than $|k_0|$ and
cannot vanish on all those elements. Thus
\[
                  (r_j,t_j)=(r_l,t_l)\quad\Longrightarrow\quad j=l.
                                                               \tag{F.28.13}
\]
Also every $r_j\ne1$, since Frobenius is a field
automorphism and $r\ne1$.

Consider a nonzero invariant subspace $U\leq Q_j$
having no proper nonzero invariant subspace.
Its intersection with $W_j$ is invariant.
If the intersection is nonzero, $U\leq W_j$.
If it is zero, $\pi_j|U$ is injective. At least
one of its four actual coordinate projections
is nonzero. That projection has invariant
kernel, hence is injective by the defining
property of $U$. Thus $U$ is one-dimensional
and its source character is one of
$\lambda_{\partial,j},\lambda_{4,j},
\lambda_{5,j},\lambda_{6,j}$.
This argument uses the displayed coordinate
maps even when quotient characters repeat;
it requires no assumption about distinct
quotient characters within one factor.

\emph{Every comparison between different factors.}
We now prove that for $j\ne l$ there is no
source-intertwining isomorphism between such
subspaces of $Q_j$ and $Q_l$.
For a source operator acting by scalars $a$ and
$b$ on its two sides, an intertwiner $F$ satisfies
$(a-b)F=0$. A nonzero $F$ therefore forces $a=b$.
Apply this identity to the exact table (F.28.10).

A boundary-character subspace has source-$0$
scalar one. Every subspace contained in $W_l$
and every deleted-meridian-character subspace
has source-$0$ scalar $r_l\ne1$. This excludes
every comparison of a boundary character with
either of those types.

Two subspaces contained in $W_j$ and $W_l$
would force $r_j=r_l$ and $t_j=t_l$ through
sources $0$ and $7$. This contradicts (F.28.13).
For two boundary characters, source $7$ forces
$t_j=t_l$ and source $2$ then forces $r_j=r_l$,
giving the same contradiction. For the same
deleted type $4$ or $5$, sources $0$ and $7$
again determine $(r_j,t_j)$. For type $6$ the
source-$7$ value is $ct_j$, and the common
nonzero $c$ cancels, with the same result.

For two different deleted types, type $4$ versus
type $5$ is excluded by source $8$, whose scalars
are $c$ and $1$. Type $5$ versus type $6$ is
excluded by source $9$, with the same two values.
For type $4$ in factor $j$ versus type $6$ in
factor $l$, source $8$ would force
$t_l=c^{-1}$. Since $c\in k_1$, inverse Frobenius
then gives $t=c^{-1}$ and every $t_j=c^{-1}$.
Source $9$ would consequently compare $c$ with
$1$, again impossible because $c\ne1$.

Finally, a subspace of $W_j$ versus type $4$
in any factor is excluded by source $8$, acting
as $1$ and $c$. A subspace of $W_j$ versus
type $5$ is excluded by source $9$.
For type $6$ in factor $l$, source $8$ would
force $t_l=1$. Inverse Frobenius gives $t=1$
and all $t_j=1$. Source $7$ would then compare
$1$ on $W_j$ with $c$ on type $6$, a contradiction.
The inverse of an intertwining isomorphism
also intertwines, so these comparisons cover
both orientations of every pair of types.

\emph{The exact restriction kernel and affine image.}
Let
\[
 \Phi:\bigoplus_{j=0}^{m-1}Q_j
        \longrightarrow H^1_{\rm cont}(K,k_{0,\chi|K}),
        \qquad ([d_j])_j\longmapsto\sum_j[d_j|K].
                                                               \tag{F.28.14}
\]
This commutes with every actual source operator
$[d]\mapsto[d\circ B_l]$.
Each individual restriction $Q_j\to H^1(K,k_{0,\chi|K})$
is injective by F.24: its character on $R$ is
the same nontrivial $C$, and its coefficient
field is finite.

Suppose $\ker\Phi$ were nonzero, and choose a
nonzero invariant subspace $T$ inside it with
smallest dimension. Each nonzero coordinate
projection $p_j|T$ is injective and identifies
$T$ with an invariant subspace of $Q_j$ having
no proper nonzero invariant subspace. Two
nonzero projections would give the explicit
source-intertwining isomorphism
\[
                  p_l\circ(p_j|T)^{-1}:p_j(T)\longrightarrow p_l(T).
                                                               \tag{F.28.15}
\]
Its inverse is $p_j\circ(p_l|T)^{-1}$.
The comparisons just proved exclude this for
$j\ne l$. Thus $T$ has only one nonzero coordinate
projection and is contained in one $Q_j$.
On that factor $\Phi$ is the injective individual
restriction, a final contradiction. This proves
the summed $H^1$ restriction injective.

Consequently zero summed restriction of cocycles
forces every cocycle to be principal. As all
characters agree on $K$ and $C\ne1$, the exact
kernel is
\[
 \left\{\bigl(b_j(\chi_j-1)\bigr)_j:
                 b_j\in k_0,\ \sum_jb_j=0\right\}.
                                                               \tag{F.28.16}
\]
The converse inclusion follows by direct restriction.
Its dimension is $m-1$ on the $7m$ continuous
cocycle coordinates, proving rank $6m+1$.

A cocycle in the trace kernel of F.26.2 gives
the tuple $(d^{q_1^j})_j$ with zero summed
restriction. It is therefore principal, and
its parameter has trace zero because $C\ne1$.
Conversely every such principal cocycle has
zero trace restriction. The precise trace
kernel is thus $\mathscr P_\chi$ of F.26.3.
Every proper affine marginal image would give
a nonprincipal trace-kernel element by the
actual conjugation and coefficient homomorphism
of F.26.1--12. Those maps retain the original
scalar character, so they preserve (F.28.1)
and the fields (F.28.3). The proved kernel
equality excludes such an element. This
proves (F.28.2) for the remaining negative
case, completing all cases considered above.

\emph{Integral evidence and sign transport.}
The symbolic producer applies exactly (F.28.4)
to the retained F.24 universal data and records
all eleven full matrices, all eight seed rows,
all basis maps and all extension blocks.
A separately structured verifier imports
neither the producer nor an algebra package.
It evaluates all nineteen literal F.23 source
and seed braid words by generic affine
multiplication over the integral Laurent ring
$A$. It checks both basis inverses, the
principal-quotient identity after clearing
only $\delta_7$, the invariant inclusion,
each diagonal quotient coordinate, every
entry of (F.28.10), its exact source labels
and the actual seed evaluations (F.28.11).
It records $781$ checks. All coefficient
identities have integer coefficients; the
only nonmonomial localization is at $1-rst$,
whose nonvanishing was already proved.
The evidence manifest binds these programs,
outputs and inputs to the exact statement.

For a positive representation, precomposition
with the actual continuous isomorphism
$\mathcal B_+:V_-\to V_+$ sends the displayed
$Y_i^+$ character to precisely the negative
character just treated. This map and its
inverse carry $\mathcal K_-$ and $R_-$ onto
their positive counterparts by F.23.13.
Equality of the two negative images therefore
gives equality of the two original positive
images through the specified maps, proving
(F.28.2) for both signs.
\end{proof}

\begin{claim}{F.29}{Proved}{All finite affine marginal images on the four-parameter product-one locus}
Retain the actual bases $Y_i^-=X_i$ and
$Y_i^+=\mathcal B_+(X_i)$, where $\mathcal B_+$
is the based isomorphism of F.23.13.
For every finite field $k$ and continuous
$\rho:V_\pm\to\operatorname{Aff}_1(k)$ with
scalar character $\chi$, the exact locus
\[
 \chi(Y_5^\pm)=\chi(Y_6^\pm)=1,
 \qquad \prod_{i=1}^7\chi(Y_i^\pm)=1
                                                               \tag{F.29.1}
\]
has the image equality
\[
                         \rho(\mathcal K_\pm)=\rho(R_\pm).
                                                               \tag{F.29.2}
\]
Every characteristic, finite scalar-field extension
degree and larger original coefficient field is included.

Write a negative character in this locus as
$(r,s,t,u,1,1,c^{-1})$, where $c=rstu$.
On its nontrivial-scalar part
$C=\langle u,c^{-1}\rangle\ne\{1\}$, put
$k_1=\mathbb F_p(u,c)$, $k_0=\mathbb F_p(r,s,t,u)$,
$q_1=|k_1|$, $m=[k_0:k_1]$ and
$\chi_j=\chi^{q_1^j}$, $0\leq j<m$.
The summed continuous $H^1$ restriction to
$K=\mathcal K_-$ is injective. At the cocycle
level its kernel and rank are exactly
\[
 \begin{gathered}
 \left\{\bigl(b_j(\chi_j-1)\bigr)_{j=0}^{m-1}:
                 b_j\in k_0,\ \sum_jb_j=0\right\},\\
                       \operatorname{rank}=6m+1
                 \quad\hbox{on }7m\hbox{ coordinates}.
 \end{gathered}                                                \tag{F.29.3}
\]
\end{claim}
\passage{Part F root-adjoining intersection question; Section 30, folios 194-196, scans 383-387; Section 31, folios 198-199 verso, scans 391-394; original F7/F19 residual subgroup-membership reduction}


\begin{proof}
\emph{Actual fields, completed cases and the four remaining charts.}
Work first with $K=\mathcal K_-\leq R=R_-\leq V=V_-$.
F.21 and the actual marginal surjection of F.22
give $\overline{\langle\!\langle K\rangle\!\rangle}_R=R$.
Apply this equality to the continuous abelian
map $\chi|R$. Normal generation in its abelian
image gives $\chi(K)=\chi(R)$.
Since the actual deleted generators are
$X_4,X_5,X_6,X_7$, their scalars now give
\[
 \begin{gathered}
 C=\chi(K)=\chi(R)=\langle u,c^{-1}\rangle,
 \qquad D=\chi(V)=\langle r,s,t,u\rangle,\\
 k_1=\mathbb F_p(C)=\mathbb F_p(u,c),\qquad
 k_0=\mathbb F_p(D)=k_1(r,t),\qquad s=c/(rtu).
 \end{gathered}                                                \tag{F.29.4}
\]
All displayed character values are nonzero.
For $C=\{1\}$ the image of $R$ is a translation
group, hence abelian. The same normal-generation
map gives (F.29.2) directly. This case is precisely
$u=c=1$.
For $C\ne\{1\}$ and $k_0=k_1$, F.25 gives the
affine equality, while F.24 gives (F.29.3) with
$m=1$.

For $u=1$, the original character is exactly
$(r,s,t,1,1,1,(rst)^{-1})$ in the unchanged basis.
F.28 supplies the affine equality and, on its
nontrivial-scalar part, the stated restriction
kernel. The coefficient homomorphism making this
comparison is $u\mapsto1$, fixing $r,s,t$ and
sending $c$ to $rst$. It sends each original
word evaluation to that of F.28, and sends the
principal-pivot denominator $1-c$ to $1-rst$.
Thus the comparison uses the actual source maps.

We now work with $u\ne1$, $C\ne\{1\}$ and
$k_1\subsetneq k_0$. The four disjoint coordinate
charts are $c\ne1,r\ne1$; $c\ne1,r=1$;
$c=1,r\ne1$; and $c=1,r=1$.
We give the actual principal quotient, invariant
maps and source calculations in each chart.

\emph{Coefficient maps, principal quotients and Frobenius.}
The first two charts use the coefficient homomorphism
\[
 \begin{aligned}
 \mathbb Z[t_1^{\pm1},\ldots,t_7^{\pm1}]&\longrightarrow
     A_4=\mathbb Z[r^{\pm1},s^{\pm1},t^{\pm1},u^{\pm1}],\\
 (t_1,t_2,t_3,t_4,t_5,t_6,t_7)&\longmapsto
                  (r,s,t,u,1,1,(rstu)^{-1}).
 \end{aligned}                                                \tag{F.29.5}
\]
An exponent vector $(a_i)$ maps to
$r^{a_1-a_7}s^{a_2-a_7}t^{a_3-a_7}u^{a_4-a_7}$.
The chart $r=1$ further applies the homomorphism
fixing $s,t,u$ and sending $r$ to $1$.
All these are derivation identities on the
specified discrete free words. Evaluating in
a finite field gives the actual cocycle rule
$d(ab)=d(a)+\chi(a)d(b)$.
Every seven-tuple of generator values defines
a homomorphism to the finite affine group and
extends uniquely to $V=\widehat F_7$.
Thus the continuous cocycle space is $k_0^7$;
the Laurent calculations are not asserted to
define continuous cocycles with infinite coefficient ring.

For the character in a given chart put
$\delta=(\chi(X_i)-1)_{i=1}^7$.
Use pivot $p=7$ when $c\ne1$, and pivot $p=4$
when $c=1,u\ne1$. Let $I_p$ be the increasing
list of the other six indices. Define
\[
 \begin{aligned}
 P_p(v)&=(v_i-\delta_i\delta_p^{-1}v_p)_{i\in I_p},
       &S_p(w)_p&=0,\quad S_p(w)_{I_p}=w,\\
 P_pS_p&=I_6,
       &S_pP_p(v)&=v-\delta\delta_p^{-1}v_p.
 \end{aligned}                                                \tag{F.29.6}
\]
Here $\delta_7=(1-c)/c\ne0$ in the first charts;
the last charts have $u=h^{-1}$, $h=rst\ne1$,
and $\delta_4=(1-h)/h\ne0$.
These two composites identify the actual continuous
cohomology quotient $Q_\chi=H^1_{\rm cont}(V,k_{0,\chi})$
with $k_0^6$ in each displayed coordinate order.

Retain the eleven actual F.23 automorphisms
$B_l(v)=u_l\alpha_l(v)u_l^{-1}$, $0\leq l\leq10$.
They and their inverses preserve $K$ and fix
the character. Their cocycle matrices $A_l$
have rows $D(B_l(X_i))$, so $A_l\delta=\delta$.
Put $\overline A_l=P_pA_lS_p$. For each basis
$J$ below the cohomology matrix
$T_l=J^{-1}\overline A_lJ$ is certified by the
integral identity
\[
 \delta_p(A_lS_pJ-S_pJT_l)
             =\delta\,\bigl(\operatorname{row}_p(A_lS_pJ)\bigr).
                                                               \tag{F.29.7}
\]
This clears precisely the displayed principal
pivot and proves the claimed action after every
permitted field evaluation.
All source indices in this proof begin at zero:
$0=\texttt{old0}$, $1,\ldots,5$ are
$\texttt{old\_1\_0},\ldots,\texttt{old\_1\_4}$,
and $6,\ldots,10$ are
$\texttt{old\_-1\_0},\ldots,\texttt{old\_-1\_4}$.
Their positions in an enumeration beginning at
one are $l+1$; the word labels remain unchanged.

For $P_0=X_1\cdots X_7$, the product-one condition
makes $d(P_0)$ a covector on cohomology.
The pure Artin map fixes $P_0$, and direct
conjugation gives $d(B_l(P_0))=\chi(u_l)d(P_0)$.
Also each character-one meridian has
$d(B_l(X_i))=\chi(u_la_{li})d(X_i)$, where
$\alpha_l(X_i)=a_{li}X_ia_{li}^{-1}$ is its
actual pure braid conjugator.
These boundary and meridian evaluations are
the quotient covectors used below.

Put $r_j=r^{q_1^j}$, $s_j=s^{q_1^j}$,
$t_j=t^{q_1^j}$; $u,c\in k_1$ stay fixed.
The exact coefficient map
\[
 d\longmapsto d^{q_1^j},\qquad
 k_{0,\chi}\longrightarrow k_{0,\chi_j},\qquad
 Q_\chi\longrightarrow Q_j=H^1_{\rm cont}(V,k_{0,\chi_j})
                                                               \tag{F.29.8}
\]
is additive and semilinear and preserves cocycles,
principals and continuity. Its inverse is the
power $q_1^{m-j}$, with the identity at $j=0$.
It commutes with every actual source precomposition.
Thus each displayed basis, inclusion, quotient,
source identity and inverse map transports by
the indicated Frobenius substitutions.
All $\chi_j$ agree on $K$, by (F.29.4).
Moreover
\[
                     (r_j,t_j)=(r_l,t_l)\quad\Longrightarrow\quad j=l.
                                                               \tag{F.29.9}
\]
Indeed such an equality for distinct indices
would make a nonidentity $k_1$-Frobenius power
fix the generators of $k_0=k_1(r,t)$.
For $0<a<m$, the polynomial $X^{q_1^a}-X$ has
degree smaller than $|k_0|$, and cannot vanish
on the whole field. In either $r=1$ chart,
$k_0=k_1(t)$, so the $t_j$ are distinct and
lie outside $k_1$. In particular $t_j^2\ne1$.
In either $r\ne1$ chart every $r_j\ne1$.

\emph{Chart $c\ne1$, $r\ne1$: the three-dimensional core.}
Use $p=7$ and the six coordinates of
$X_1,\ldots,X_6$. Put $b=(1,r,rs,rst,c,c)$.
The boundary and the actual $X_5,X_6$ evaluations
give the explicit maps
\[
 \begin{aligned}
 \pi(v)&=(bv,v_5,v_6),\\
 \iota(x,y,z)&=(-rx-rsy-rstz,x,y,z,0,0),\\
 \sigma(z_\partial,z_5,z_6)
             &=(z_\partial-c(z_5+z_6),0,0,0,z_5,z_6).
 \end{aligned}                                                \tag{F.29.10}
\]
They satisfy $\pi\sigma=I_3$,
$\ker\pi=\operatorname{im}\iota=:W$ and
$\iota(v_2,v_3,v_4)+\sigma\pi(v)=v$.
In particular the basis and its inverse are
\[
 \begin{gathered}
 J=(\iota e_1,\iota e_2,\iota e_3,e_1,e_5-ce_1,e_6-ce_1),\\
 J^{-1}(v)=(v_2,v_3,v_4,bv,v_5,v_6).
 \end{gathered}                                                \tag{F.29.11}
\]
Here the $e_i$ inside $\iota$ are the three
domain coordinate vectors; the remaining ones
are the six quotient coordinate vectors.
The displayed decomposition proves both basis
composites directly.
In this basis every actual source matrix is
$\left(\begin{smallmatrix}G_l&E_l\\0&D_l\end{smallmatrix}\right)$,
where $G_l$ acts on $W$ and $D_l$ is diagonal
on the boundary, $X_5$ and $X_6$ coordinates.
Equivalently $\overline A_l\iota=\iota G_l$ and
$\pi\overline A_l=D_l\pi$.
The actual source-word table is
\[
 \begin{array}{c|c|ccc}
 l&G_l&\lambda_{\partial,l}&\lambda_{5,l}&\lambda_{6,l}\\\hline
 0&rI_3&1&r&r\\
 2&*&rt&rt/u&rct\\
 7&tI_3&t^{-1}&t/u&ct\\
 8&*&1&1&t^{-1}\\
 9&*&1&c&1
 \end{array}
                                                               \tag{F.29.12}
\]
The stars denote retained full core matrices;
the additional exact identities used for them are
\[
 (G_8-I_3)^2(G_8-(c/u)I_3)=0,
 \qquad (G_9-I_3)^2(G_9-(tu)^{-1}I_3)=0.
                                                               \tag{F.29.13}
\]
These matrix products are integral Laurent
identities verified from the literal source words.

We explain the invariant-subspace classification
used in this and the next two separating charts.
Take a nonzero source-invariant subspace
$U\leq Q_j$ with no proper nonzero invariant
subspace. Its intersection with the displayed
$W_j$ is invariant. A nonzero intersection
therefore gives $U\leq W_j$.
For zero intersection, $\pi_j|U$ is injective.
A nonzero coordinate projection of this map
has invariant kernel, so is injective as well.
It identifies $U$ with one of the displayed
one-dimensional quotient characters. This
uses the actual coordinate projections and
does not require distinct quotient characters
inside a single factor.

For $j\ne l$, compare these subspaces through
any source-intertwining isomorphism. Source
operators acting by scalars $a,b$ on its two
sides give $(a-b)F=0$, so $F\ne0$ forces $a=b$.
In the present chart source $0$ separates a
boundary type from every other type, because
its scalars are $1$ and $r_l\ne1$.
Two core subspaces force $r_j=r_l,t_j=t_l$
through sources $0,7$, contrary to (F.29.9).
Two boundary types force $t_j=t_l$ through
source $7$ and then $r_j=r_l$ through source $2$.
For the same type $X_5$ or $X_6$, sources $0,7$
determine the pair $(r_j,t_j)$ after cancelling
the common units $u$ or $c$.
The different types $X_5,X_6$ are separated
by source $9$, with scalars $c\ne1$ and $1$.

A core subspace in factor $j$ isomorphic to
type $X_5$ in factor $l$ must be one-dimensional.
Source $9$ and (F.29.13) force
$(c-1)^2(c-(t_ju)^{-1})=0$, hence
$t_j=(cu)^{-1}\in k_1$.
Inverse Frobenius makes $t=(cu)^{-1}$ and all
$t_l=t_j$. Source $0$ then forces $r_j=r_l$,
contradicting (F.29.9).
For a core subspace versus type $X_6$, source
$8$ and (F.29.13) force
$t_l^{-1}\in\{1,c/u\}$, so
$t_l\in\{1,u/c\}\subset k_1$.
Again all $t_j=t_l$, and source $0$ contradicts
(F.29.9). These comparisons cover every pair
of the four types, including the same-type
pairs. The inverse of an intertwining
isomorphism covers the opposite orientation.

\emph{Chart $c\ne1$, $r=1$: two core coordinates and four quotients.}
Here $c=stu$, $k_1=\mathbb F_p(u,c)$ and
$k_0=k_1(t)$. In the same principal quotient
$p=7$, let $b=(1,1,s,st,c,c)$ and use
\[
 \begin{gathered}
 \pi(v)=(bv,v_1,v_5,v_6),\qquad
 \iota(x,y)=(0,-sx-sty,x,y,0,0),\\
 \sigma(z_\partial,z_1,z_5,z_6)
       =(z_1,z_\partial-z_1-c(z_5+z_6),0,0,z_5,z_6),\\
 J=(\iota e_1,\iota e_2,e_2,e_1-e_2,e_5-ce_2,e_6-ce_2),\\
 J^{-1}(v)=(v_3,v_4,bv,v_1,v_5,v_6).
 \end{gathered}                                                \tag{F.29.14}
\]
The identities $\pi\sigma=I_4$,
$\ker\pi=\operatorname{im}\iota=W$ and
$\iota(v_3,v_4)+\sigma\pi(v)=v$ prove the
maps and both basis inverses explicitly.
All eleven source matrices again have the
displayed invariant-core and diagonal-quotient
form. Their required entries are
\[
 \begin{array}{c|c|cccc}
 l&G_l&\lambda_{\partial,l}&\lambda_{1,l}
                         &\lambda_{5,l}&\lambda_{6,l}\\\hline
 1&*&1&u^{-2}&1&1\\
 2&tI_2&t&t^{-1}&t/u&ct\\
 5&I_2&1&c^2&t^{-1}&u^{-1}\\
 7&tI_2&t^{-1}&t&t/u&ct\\
 8&*&1&1&1&t^{-1}\\
 9&*&1&1&c&1
 \end{array}
                                                               \tag{F.29.15}
\]
The smaller core satisfies the exact quadratic
identities
\[
 (G_8-I_2)(G_8-(c/u)I_2)=0,
 \qquad (G_9-I_2)(G_9-(tu)^{-1}I_2)=0.
                                                               \tag{F.29.16}
\]
They, the basis and every quotient coordinate
are independently checked from all eleven
literal source words after $r\mapsto1$.

The five types are now $W,\partial,X_1,X_5,X_6$.
Their source-$2$, source-$7$ scalar pairs are,
respectively,
$(t_j,t_j)$, $(t_j,t_j^{-1})$,
$(t_j^{-1},t_j)$, $(t_j/u,t_j/u)$,
and $(ct_j,ct_j)$.
The same-type comparisons determine $t_j=t_l$,
which is impossible for $j\ne l$.
A core type versus boundary or $X_1$ would
force $t_l^2=1$ through these two sources.
A boundary or $X_1$ type versus $X_5$ or $X_6$
would likewise force $t_j^2=1$ on the first
side. These conclusions contradict
$t_j\notin k_1$.
Types $X_5,X_6$ are separated by source $9$.
A core type versus $X_5$ would force
$c=(t_ju)^{-1}$ by (F.29.16) and source $9$,
putting $t_j\in k_1$.
A core type versus $X_6$ would force
$t_l\in\{1,u/c\}\subset k_1$ by source $8$.
Both are impossible.

The only pair left is boundary in factor $j$
versus $X_1$ in factor $l$. Sources $1,5$ force
$u^2=c^2=1$. Since $c\ne1$ and $u\ne1$,
the factorization $(a-1)(a+1)=a^2-1$ in the
field gives
\[
 c=u=-1,\qquad p\ne2,\qquad
 s=t^{-1},\qquad k_1=\mathbb F_p,\qquad
 t\notin\mathbb F_p.
                                                               \tag{F.29.17}
\]
In characteristic two $a^2=1$ has only the
root $1$, so this remaining case is absent.

In (F.29.17) the coordinate order of (F.29.14)
is $(x,y,z_\partial,z_1,z_5,z_6)$.
Sources $2,7$ have core $tI_2$ and quotient
diagonals $(t,t^{-1},-t,-t)$ and
$(t^{-1},t,-t,-t)$, respectively.
The nonzero factors $t-t^{-1}$ and $2t$ follow
from $t\notin\mathbb F_p$ and $p\ne2$.
A boundary-character eigenvector must therefore
have its $X_1,X_5,X_6$ coordinates zero, by
source $2$. Its boundary coordinate cannot be
zero, because then source $7$ would compare
$t$ on the core with $t^{-1}$ on that vector.
Normalize this nonzero coordinate to one.
Source $7$ then uniquely forces the first
vector in the following exact identities:
\[
 \begin{gathered}
 v_\partial=(1-2t,2,1,0,0,0)^{\mathsf T},
       \qquad T_7v_\partial=t^{-1}v_\partial,\\
 v_1=(-t,0,0,1,0,0)^{\mathsf T},
       \qquad T_2v_1=t^{-1}v_1,\\
 (T_1-I_6)v_\partial=(T_1-I_6)v_1
                         =(3t-1,-2,0,0,0,0)^{\mathsf T}.
 \end{gathered}                                                \tag{F.29.18}
\]
The boundary scalar at source $1$ is $1$.
Its required eigenvector equation contradicts
the nonzero second coordinate $-2$ in
(F.29.18). Thus the boundary-character line
does not occur as a submodule in this case.
For completeness, source $7$ first kills the
other quotient coordinates of an $X_1$ line;
source $2$ uniquely forces $v_1$ after its
$X_1$ coordinate is normalized. The same
residual excludes that line too.
The coefficient specialization for these
identities is $r\mapsto1$, $s\mapsto t^{-1}$,
$u\mapsto-1$. A monomial $s^a t^b u^d$ maps
to $(-1)^d t^{b-a}$, and the pivot becomes
$\delta_7=-2$. The separate literal three-word
audit verifies (F.29.7) with this exact pivot
for sources $1,2,7$, before using its proved
odd-characteristic invertibility.
The same argument transports by Frobenius
to every factor. It excludes the last possible
cross-type comparison and completes this chart.

\emph{Chart $c=1$, $r\ne1$: the actual pivot $X_4$.}
Put $h=rst=u^{-1}\ne1$.
The character is now literally
$(r,s,t,h^{-1},1,1,1)$.
Its coefficient map is
\[
 \begin{aligned}
 \mathbb Z[t_1^{\pm1},\ldots,t_7^{\pm1}]&\longrightarrow
       A_3=\mathbb Z[r^{\pm1},s^{\pm1},t^{\pm1}],\\
 (t_1,t_2,t_3,t_4,t_5,t_6,t_7)&\longmapsto
                         (r,s,t,(rst)^{-1},1,1,1).
 \end{aligned}                                                \tag{F.29.19}
\]
The exponent image is
$r^{a_1-a_4}s^{a_2-a_4}t^{a_3-a_4}$.
This is also the homomorphism
$A_4/(rstu-1)\to A_3$ sending $u$ to
$(rst)^{-1}$ and fixing $r,s,t$.
The reverse map sends $r,s,t$ to their classes;
the relation $rstu=1$ makes the two composites
the identity, including on $u$.
The maps are applied to the unlocalized word
identities, followed here by localization at
$1-h$ for the actual pivot $X_4$ in (F.29.6).
The scalar fields are $k_1=\mathbb F_p(h)$
and $k_0=k_1(r,t)$ because $s=h/(rt)$.

The six quotient coordinates are now those
of $X_1,X_2,X_3,X_5,X_6,X_7$, in that order.
Put $b=(1,r,rs,1,1,1)$. The exact maps are
\[
 \begin{aligned}
 \pi(v)&=(bv,v_4,v_5,v_6),\\
 \iota(x,y)&=(-rx-rsy,x,y,0,0,0),\\
 \sigma(z_\partial,z_5,z_6,z_7)
              &=(z_\partial-z_5-z_6-z_7,0,0,z_5,z_6,z_7).
 \end{aligned}                                                \tag{F.29.20}
\]
They satisfy $\pi\sigma=I_4$,
$\ker\pi=\operatorname{im}\iota=W$ and
$\iota(v_2,v_3)+\sigma\pi(v)=v$.
Thus the basis maps are explicitly
\[
 \begin{gathered}
 J=(\iota e_1,\iota e_2,e_1,e_4-e_1,e_5-e_1,e_6-e_1),\\
                    J^{-1}(v)=(v_2,v_3,bv,v_4,v_5,v_6).
 \end{gathered}                                                \tag{F.29.21}
\]
Every actual source has invariant two-dimensional
core and diagonal quotient coordinates
$\partial,X_5,X_6,X_7$. The literal table is
\[
 \begin{array}{c|c|cccc}
 l&G_l&\lambda_{\partial,l}&\lambda_{5,l}
                         &\lambda_{6,l}&\lambda_{7,l}\\\hline
 0&rI_2&1&r&r&r\\
 2&*&rt&rht&rt&rt\\
 7&tI_2&t^{-1}&ht&t&t\\
 8&I_2&1&1&t^{-1}&h\\
 10&I_2&1&t^{-1}&h&1
 \end{array}
                                                               \tag{F.29.22}
\]
All entries use the original source indices;
in particular $10=\texttt{old\_-1\_4}$.

The five possible types are
$W,\partial,X_5,X_6,X_7$, with the classification
proved from the displayed maps above.
Source $0$ separates boundary from every
other type since $r_j\ne1$.
Two core types determine $(r_j,t_j)$ by
sources $0,7$; two boundary types do so by
sources $7,2$. For the same meridian type,
sources $0,7$ determine the same pair, with
the common $h$ cancelled for type $X_5$.
Every same-type comparison across distinct
factors therefore contradicts (F.29.9).

For the different types, core versus $X_7$
is excluded by source $8$, with scalars $1,h$.
Core versus $X_6$ is excluded by source $10$.
Core versus $X_5$ in factor $l$ would force
$t_l=1$ by source $10$. Inverse Frobenius then
gives all $t_j=1$, and source $7$ compares
$1$ with $h\ne1$.
Types $X_5,X_7$ are excluded by source $8$;
types $X_6,X_7$ by source $10$.
For type $X_5$ in factor $j$ versus $X_6$
in factor $l$, source $8$ forces $t_l=1$,
hence all $t_j=1$, and source $10$ then
compares $1$ with $h$.
This proves every cross-type comparison,
in either orientation. No extra condition
on $t-1$ or on the characteristic has entered
this chart.

\emph{Chart $c=1$, $r=1$: an essential line and actual seed evaluations.}
Here the character is
$(1,s,t,(st)^{-1},1,1,1)$, $h=st\ne1$,
$k_1=\mathbb F_p(h)$ and $k_0=k_1(t)=k_1(s)$.
The coefficient map is $r\mapsto1$ in
(F.29.19), sending $h$ to $st$ and retaining
the principal pivot $X_4$.
Since the field extension is proper, both
$s,t\notin k_1$. In particular
$s-1,t-1,t+1,h-1$ are nonzero, in every
characteristic, as are the units $s,t$.

Keep the six coordinates of the preceding
chart and put $b=(1,1,s,1,1,1)$.
The common kernel of the boundary and the
four character-one meridian covectors is
the line generated by $v_0=(0,1,-s^{-1},0,0,0)$.
The following basis and inverse are exact:
\[
 \begin{gathered}
 J=(v_0,e_2,e_5-e_2,e_4-e_2,e_1-e_2,e_6-e_2),\\
 J^{-1}(v)=(-sv_3,bv,v_5,v_4,v_1,v_6),\\
 J(z_0,z_\partial,z_6,z_5,z_1,z_7)
  =(z_1,z_0+z_\partial-z_6-z_5-z_1-z_7,
                                  -s^{-1}z_0,z_5,z_6,z_7).
 \end{gathered}                                                \tag{F.29.23}
\]
Substitution proves both composites to be the
identity. In these coordinates the actual
source matrices and one required lift identity are
\[
 \begin{gathered}
 T_l=\begin{pmatrix}\nu_l&a_l\\0&\Lambda_l\end{pmatrix},
 \qquad \Lambda_l\text{ diagonal},\\
 (\nu_0,\ldots,\nu_{10})=(1,1,t,1,s,1,1,t,1,s,1),\\
 \widetilde v_0=S_4v_0,\qquad
 A_2\widetilde v_0=t\widetilde v_0+st^2\delta.
 \end{gathered}                                                \tag{F.29.24}
\]
All eleven full matrices and the displayed
principal correction are independently verified.

For $j=1,2,5,7,9$, the additive commutators
$T_4T_j-T_jT_4$ kill the first line and have
image in that line. Their first rows on the
last five coordinates form $gH$, where
\[
 \begin{gathered}
 g=(s-1)(t-1),\\
 H=\begin{pmatrix}
 0&0&0&(st+1)(st+s-1)&-t^{-2}\\
 0&0&1&-(t+1)&0\\
 0&s/t&-t^{-2}&0&0\\
 (t+1)(st-t+1)/(st^3)&0&s&0&0\\
 -(s-1)/(st^2)&0&0&-(s-1)&0
 \end{pmatrix}.
 \end{gathered}                                                \tag{F.29.25}
\]
The independent signed-permutation expansion
of the determinant gives
\[
 \det H=\frac{(s-1)(t-1)(t+1)}{t^6},
 \qquad
 \det(gH)=\frac{(s-1)^6(t-1)^6(t+1)}{t^6}\ne0.
                                                               \tag{F.29.26}
\]
The factor proof preceding (F.29.23) justifies
every nonzero factor, including in characteristic
two where $t+1=t-1$.
Every nonzero subspace invariant under the
actual eleven sources contains
$L=k_0[v_0]$. Indeed, a vector with nonzero
last five coordinates is sent by some displayed
commutator to a nonzero first-line vector,
by (F.29.26). That vector remains in the
invariant subspace and spans $L$. A nonzero
subspace already contained in that line
equals $L$. This proves the essential-line
statement directly from the source operators.

Let $d_0$ have generator values $\widetilde v_0$.
For the actual negative seed
$a=u_{\texttt{cross\_0\_-1\_1}}$ of index $13$,
the word calculation and (F.29.24) give
\[
 \begin{gathered}
 \chi(a)=h^{-1},\qquad d_0(a)=h-1,\\
 \kappa=B_2(a)a^{-1}\in K,\qquad \chi(\kappa)=1,
                    \qquad d_0(\kappa)=1-h.
 \end{gathered}                                                \tag{F.29.27}
\]
In fact $d_0(B_2(a))=t(h-1)+ht(h^{-1}-1)=0$.
For a character-one word the principal
correction in (F.29.24) vanishes. Thus
\[
 \kappa_n=B_2^n(\kappa)\in K,\qquad
 \chi(\kappa_n)=1,\qquad
                       d_0(\kappa_n)=t^n(1-h)\quad(n\geq0).
                                                               \tag{F.29.28}
\]
The source preserves $K$, and the seed belongs
to the original eight negative seeds, so these
are actual elements of $K$. Independently, the
first two derivative rows
$D(a)(A_2-I_7)$ and $D(a)(A_2-I_7)A_2$ are
checked against (F.29.27)--(F.29.28), including
their vanishing on the principal vector.

\emph{The common restriction map and its exact kernel.}
Across all four charts consider the actual map
\[
 \Phi:\bigoplus_{j=0}^{m-1}Q_j
       \longrightarrow H^1_{\rm cont}(K,k_{0,\chi|K}),
       \qquad ([d_j])_j\longmapsto\sum_j[d_j|K].
                                                               \tag{F.29.29}
\]
It is well-defined because the characters
agree on $K$, and it commutes with every
source operator $[d]\mapsto[d\circ B_l]$.
Each individual restriction is injective
by F.24, since $\chi_j(R)=C\ne\{1\}$.

Choose a smallest-dimensional nonzero
source-invariant subspace $T\leq\ker\Phi$
under the supposition that the kernel is nonzero.
Each nonzero coordinate projection $p_j|T$
is injective, since its kernel is invariant.
Its image is a nonzero invariant subspace
with no proper nonzero invariant subspace.
Two nonzero projections therefore give the
exact mutually inverse intertwiners
\[
 p_l\circ(p_j|T)^{-1}:p_j(T)\longrightarrow p_l(T),
 \qquad p_j\circ(p_l|T)^{-1}:p_l(T)\longrightarrow p_j(T).
                                                               \tag{F.29.30}
\]
In the first three charts every pair of
possible types was excluded above for $j\ne l$.
There is therefore only one nonzero projection;
$T$ lies in a single factor, contradicting
that factor's F.24 restriction injectivity.
This proves $\Phi$ injective in those charts.

In the last chart the essential-line argument
puts $L_j$ inside each nonzero projection
image. Minimality makes that image equal to
$L_j$, so $T\leq\bigoplus_jL_j$.
Frobenius transport of (F.29.28) gives, on
these lines, the evaluation matrix at
$\kappa_0,\ldots,\kappa_{m-1}$ and its determinant
\[
 ((1-h)t_j^n)_{0\leq n,j<m},
 \qquad (1-h)^m\prod_{j<l}(t_l-t_j)\ne0.
                                                               \tag{F.29.31}
\]
All these evaluations descend to the common
target cohomology, since the words have
character one. The coefficient inverse is
explicit: write
\[
 \prod_{l\ne j}\frac{X-t_l}{t_j-t_l}
       =\sum_{n=0}^{m-1}h_{jn}X^n,
 \qquad z_j=(1-h)^{-1}\sum_{n=0}^{m-1}h_{jn}y_n.
                                                               \tag{F.29.32}
\]
It inverts $y_n=\sum_j(1-h)t_j^nz_j$;
evaluation of the displayed polynomial at
each $t_l$ proves the inverse composites.
Every denominator is nonzero by $h\ne1$
and the proved distinctness of the $t_j$.
Thus evaluation after $\Phi$ is injective
on $\bigoplus_jL_j$, contradicting the nonzero
$T$ in its kernel. This proves the injectivity
also in the last chart.

Zero summed restriction of cocycles now
forces each cohomology class to be zero,
so each cocycle is $b_j(\chi_j-1)$.
Its summed restriction is
$(\chi|K-1)\sum_jb_j$. Since $C\ne\{1\}$,
this vanishes exactly for $\sum_jb_j=0$.
The converse follows by direct restriction.
The principal vector in every factor is
nonzero, and the displayed parameter-sum
kernel has dimension $m-1$.
This proves (F.29.3), including the rank
$7m-(m-1)=6m+1$.

For the exact trace reduction of F.26, the
tuple $(d^{q_1^j})_j$ gives
\[
 \operatorname{Tr}_{k_0/k_1}(d|K)=0
 \quad\Longleftrightarrow\quad
 d=b(\chi-1)\ \hbox{with}\ \operatorname{Tr}_{k_0/k_1}(b)=0.
                                                               \tag{F.29.33}
\]
The forward implication is the proved cocycle
kernel, and the reverse is direct evaluation,
using $\chi(K)\subset k_1$.
F.26 sends every proper finite affine marginal
image through its explicit conjugation and
coefficient homomorphism to a nonprincipal
element of this trace kernel over the original
scalar field $k_0$. Those maps retain $\chi$,
so they retain (F.29.1), $C$, $k_1$ and $k_0$.
Equation (F.29.33) excludes that element.
This proves (F.29.2) for every remaining
negative representation, including those
with an original coefficient field larger
than $k_0$.

\emph{Integral evidence and the actual sign transport.}
The three symbolic producers retain all full
source matrices and basis maps for the generic
four-parameter chart and the two pivot-$X_4$
charts. Their exact check counts are $754$,
$691$ and $956$.
The separately structured three-chart verifier
imports neither a producer nor an algebra
package. It reconstructs all eleven old-source
and all eight negative-seed braid words in
each of those three integral Laurent rings.
Its $3246$ checks include (F.29.7), both basis
inverses, every invariant inclusion and diagonal
quotient, the comparison tables, the integral
annihilators, the five printed commutator rows,
their $120$-term determinant expansion, the
source-$2$ principal correction and the actual
seed evaluations.
The additional $r=1,c\ne1$ verifier uses the
same reviewed integer Laurent arithmetic,
reconstructs all eleven original words again,
and records $996$ checks of the new basis,
every quotient character and (F.29.16).
Root's separate exceptional audit reconstructs
sources $1,2,7$ directly and checks the exact
coefficient map and both residual identities
in (F.29.18), with $20$ additional recorded
certificate assertions.

These are integral coefficient identities.
The only nonmonomial principal localizations
are $1-c$ and $1-h$ in their respective
charts; their nonvanishing was proved before
each quotient was used. All other determinant
and interpolation factors have their explicit
nonzero proofs above. The exceptional constant
$2$ is inverted only after (F.29.17) proves
odd characteristic. The evidence manifest
binds the source words, programs, outputs,
coefficient inputs and completed dependencies
to this exact locus and statement.

Finally, a positive representation precomposes
with $\mathcal B_+:V_-\to V_+$ to have precisely
the negative character just proved in the
basis $Y_i^+=\mathcal B_+(X_i)$.
This continuous isomorphism and its inverse
carry $\mathcal K_-,R_-$ onto their positive
counterparts by F.23.13. Equality of the two
negative images therefore gives equality of
the two original positive images through
those exact maps. This proves (F.29.2) for
both signs.
\end{proof}

\begin{claim}{F.30}{Proved}{Finite affine product-one images when the first scalar generates the scalar-field extension}
Use the actual bases $Y_i^-=X_i$ and
$Y_i^+=\mathcal B_+(X_i)$ under the based
isomorphism of F.23.13.
For a finite field $k$ and continuous
$\rho:V_\pm\to\operatorname{Aff}_1(k)$, let
$\chi$ be its scalar character and put
$C=\chi(R_\pm)$,
$k_1=\mathbb F_p(C)$,
$k_0=\mathbb F_p(\chi(V_\pm))$ and
$r=\chi(Y_1^\pm)$.
The exact locus
\[
             \prod_{i=1}^7\chi(Y_i^\pm)=1,
                    \qquad k_0=k_1(r)
                                                               \tag{F.30.1}
\]
has the image equality
\[
                         \rho(\mathcal K_\pm)=\rho(R_\pm).
                                                               \tag{F.30.2}
\]
All characteristics, finite extension degrees
and larger original coefficient fields are included.
On the part $C\ne\{1\}$, put
$q_1=|k_1|$, $m=[k_0:k_1]$ and
$\chi_j=\chi^{q_1^j}$, $0\leq j<m$.
The summed continuous $H^1$ restriction for
these $m$ characters is injective. Its cocycle
kernel consists exactly of the principal tuples
whose parameters sum to zero, and its rank is
$6m+1$ on $7m$ coordinates.
\end{claim}
\passage{Part F root-adjoining intersection question; Section 30, folios 194-196, scans 383-387; Section 31, folios 198-199 verso, scans 391-394; original F7/F19 residual subgroup-membership reduction}


\begin{proof}
\emph{Actual fields and the completed scalar cases.}
Use the negative groups $K=\mathcal K_-\leq R=R_-\leq V=V_-$
first. Write their product-one character as
$(r,s,t,u,v,w,c^{-1})$, where $c=rstuvw$.
The normal-generation equality of F.21 maps
through the actual marginal surjection of F.22
to $\overline{\langle\!\langle K\rangle\!\rangle}_R=R$.
Its image under the continuous abelian scalar
map gives $\chi(K)=\chi(R)$.
The actual deleted generators $X_4,\ldots,X_7$
and (F.30.1) therefore give
\[
 \begin{gathered}
 C=\chi(K)=\chi(R)=\langle u,v,w,c^{-1}\rangle,
 \qquad k_1=\mathbb F_p(u,v,w,c),\\
 k_0=\mathbb F_p(r,s,t,u,v,w)=k_1(r).
 \end{gathered}                                                \tag{F.30.3}
\]
All six parameters are nonzero.
For $C=\{1\}$, the image of $R$ is an abelian
translation group, so the same actual
normal-generation map proves (F.30.2).
For $C\ne\{1\}$ and $k_0=k_1$, F.25 gives
the image equality. F.24 supplies the individual
restriction injection and the stated cocycle
kernel with $m=1$.

We now treat $C\ne\{1\}$ and $k_1\subsetneq k_0$.
Equation (F.30.3) implies $r\notin k_1$, hence
$r-1\ne0$. This gives an actual retained-generator
principal pivot, independently of the value of $c$.

\emph{Exact coefficient and principal-quotient maps.}
The word derivations use the coefficient homomorphism
\[
 \begin{aligned}
 \mathbb Z[t_1^{\pm1},\ldots,t_7^{\pm1}]
    &\longrightarrow A=
       \mathbb Z[r^{\pm1},s^{\pm1},t^{\pm1},u^{\pm1},v^{\pm1},w^{\pm1}],\\
 (t_1,t_2,t_3,t_4,t_5,t_6,t_7)
    &\longmapsto(r,s,t,u,v,w,(rstuvw)^{-1}).
 \end{aligned}                                                \tag{F.30.4}
\]
Thus an exponent vector $(a_i)$ maps to
$r^{a_1-a_7}s^{a_2-a_7}t^{a_3-a_7}
 u^{a_4-a_7}v^{a_5-a_7}w^{a_6-a_7}$.
These are derivations on the specified discrete
free words. Their evaluations in a finite field
give the actual cocycle rule.
Every seven-tuple of generator values defines
a homomorphism to the finite affine group
and extends uniquely to $V=\widehat F_7$.
Consequently the continuous cocycle space is
exactly $k_0^7$.

Let $\delta=(r,s,t,u,v,w,c^{-1})^{\mathsf T}-\mathbf1$.
Its first coordinate is $\delta_1=r-1\ne0$.
The actual principal quotient uses the maps
\[
 \begin{aligned}
 P_1(z)&=(z_i-\delta_i(r-1)^{-1}z_1)_{i=2}^7,
       &S_1(x)&=(0,x_1,\ldots,x_6)^{\mathsf T},\\
 P_1S_1&=I_6,
       &S_1P_1(z)&=z-\delta(r-1)^{-1}z_1.
 \end{aligned}                                                \tag{F.30.5}
\]
They identify
$Q_\chi=H^1_{\rm cont}(V,k_{0,\chi})$
with six coordinates corresponding, in order,
to the original $X_2,\ldots,X_7$.
The universal maps are defined over $A[(r-1)^{-1}]$;
the proved nonvanishing justifies every
finite-field use of that localization.

Let $P_0=X_1\cdots X_7$ be the actual boundary word.
Its character is one, and its derivative and
induced quotient covector are
\[
 \begin{gathered}
 D(P_0)=(1,r,rs,rst,rstu,rstuv,c),\\
 b=D(P_0)S_1=(r,rs,rst,rstu,rstuv,c),\\
 D(P_0)\delta=0,\qquad D(P_0)=bP_1.
 \end{gathered}                                                \tag{F.30.6}
\]
The vanishing is the telescoping identity
$D(P_0)\delta=\chi(P_0)-1$; substituting
(F.30.5) proves the last displayed equality.

For $\beta=(s,st,stu,stuv,stuvw)$, define
\[
 \begin{aligned}
 \pi:Q_\chi&\longrightarrow k_0,
       &\pi(z)&=bz,\\
 \iota:k_0^5&\longrightarrow Q_\chi,
       &\iota(x)&=(-\textstyle\sum_{i=1}^5\beta_ix_i,x_1,\ldots,x_5),\\
 \sigma:k_0&\longrightarrow Q_\chi,
       &\sigma(a)&=r^{-1}a e_1.
 \end{aligned}                                                \tag{F.30.7}
\]
One has $\pi\sigma=I_1$,
$\ker\pi=\operatorname{im}\iota=:W_\chi$ and
$\iota(z_2,\ldots,z_6)+\sigma\pi(z)=z$.
In particular the coefficient basis and its
inverse are explicitly
\[
 \begin{gathered}
 J=(e_2-se_1,e_3-ste_1,e_4-stue_1,
                 e_5-stuve_1,e_6-stuvwe_1,r^{-1}e_1),\\
                  J^{-1}(z)=(z_2,z_3,z_4,z_5,z_6,bz).
 \end{gathered}                                                \tag{F.30.8}
\]
The displayed decomposition proves both
basis composites without an assumed invariant complement.

\emph{The actual source operators needed for separation.}
Retain the eleven F.23 automorphisms
$B_l(a)=u_l\alpha_l(a)u_l^{-1}$, $0\leq l\leq10$.
They and their inverses preserve $K$.
Their pure Artin action fixes every abelian
character, and their cocycle matrices $A_l$
therefore satisfy $A_l\delta=\delta$.
Set $\overline A_l=P_1A_lS_1$.

Every pure $\alpha_l$ fixes $P_0$. The cocycle
formula for conjugation gives
$d(B_l(P_0))=\chi(u_l)d(P_0)$ because
$\chi(P_0)=1$. Combining this with (F.30.6)
gives $\pi\overline A_l=\chi(u_l)\pi$.
Thus $W_\chi$ is invariant under all actual
source operators, and $\pi$ is their exact
one-dimensional boundary quotient.

Put $\widetilde\iota=S_1\iota$.
The literal original-word calculation gives
the integral identity
\[
 (r-1)(A_0\widetilde\iota-r\widetilde\iota)
       =\delta\,\operatorname{row}_1(A_0\widetilde\iota).
                                                               \tag{F.30.9}
\]
Apply $P_1$ and cancel the proved nonzero
$r-1$. This gives the actual equality
$\overline A_0\iota=r\iota$ on the five-dimensional
invariant subspace. The same word calculation
gives the boundary entries
\[
 \begin{array}{c|c|c}
 l&\overline A_l|W_\chi&\chi(u_l)\text{ on }Q_\chi/W_\chi\\\hline
 0&rI_5&1\\
 2&*&rt\\
 7&*&t^{-1}
 \end{array}.
                                                               \tag{F.30.10}
\]
The stars stand for the actual core operators;
only the displayed entries are used.
The indices are the original zero-based
$0=\texttt{old0}$,
$2=\texttt{old\_1\_1}$ and
$7=\texttt{old\_-1\_1}$, respectively the
first, third and eighth source positions.
The certificate compares all entries of the
three full seven-coordinate source actions
with direct evaluations of those labelled
decorated Artin words.

\emph{Frobenius transport and all three type comparisons.}
For $0\leq j<m$, put
$r_j=r^{q_1^j}$, $s_j=s^{q_1^j}$ and
$t_j=t^{q_1^j}$. The parameters $u,v,w,c$
are fixed by these powers.
The coefficient map
\[
 d\longmapsto d^{q_1^j},\qquad
 k_{0,\chi}\longrightarrow k_{0,\chi_j}
                                                               \tag{F.30.11}
\]
is additive and semilinear. It preserves the
cocycle rule, principal cocycles and continuity,
and commutes with each actual source precomposition.
The inverse is the power $q_1^{m-j}$, with
the identity at $j=0$.
It transports every displayed inclusion,
quotient, section and inverse basis to
$Q_j=H^1_{\rm cont}(V,k_{0,\chi_j})$ and $W_j$.
All characters agree on $K$, since
$\chi(K)=C\subset k_1^\times$.

The elements $r_0,\ldots,r_{m-1}$ are distinct
and each differs from $1$.
An equality of two would make a nonidentity
$k_1$-Frobenius power fix the generator $r$
and hence all of $k_0=k_1(r)$.
For $0<a<m$, the polynomial $X^{q_1^a}-X$
has degree smaller than $|k_0|$, so cannot
vanish on the whole field.
Also a conjugate equal to $1$ would give
$r=1\in k_1$, contrary to the proper extension.

Take a nonzero source-invariant subspace
$U\leq Q_j$ with no proper nonzero invariant
subspace. A nonzero invariant intersection
with $W_j$ gives $U\leq W_j$.
For zero intersection, $\pi_j|U$ is injective
and identifies $U$ with the one-dimensional
boundary character. These are the two
possibilities through the exact maps (F.30.7).

For distinct $j,l$, the following comparisons
exclude every source-intertwining isomorphism
between such subspaces. Two subspaces in
$W_j,W_l$ have source-$0$ scalars $r_j,r_l$;
an intertwiner $F$ would satisfy
$(r_j-r_l)F=0$, hence $F=0$.
A core subspace and a boundary subspace
have source-$0$ scalars $r_j$ and $1$,
and the same argument excludes an isomorphism.
For two boundary subspaces, source $7$
would force $t_j^{-1}=t_l^{-1}$, so $t_j=t_l$.
Source $2$ would then force
$r_jt_j=r_lt_l$, hence $r_j=r_l$,
contrary to the proved distinctness.
All character values cancelled here are
nonzero. Inverse intertwiners cover both
orientations of the mixed comparison.
These are all three unordered type comparisons.

\emph{Summed restriction and the finite affine image.}
Consider the well-defined actual restriction map
\[
 \Phi:\bigoplus_{j=0}^{m-1}Q_j
      \longrightarrow H^1_{\rm cont}(K,k_{0,\chi|K}),
      \qquad ([d_j])_j\longmapsto\sum_j[d_j|K].
                                                               \tag{F.30.12}
\]
It commutes with every actual source operator.
Every individual restriction is injective
by F.24, because $\chi_j(R)=C\ne\{1\}$.
Under the supposition that $\ker\Phi$ is
nonzero, choose a nonzero invariant subspace
$T$ inside it with smallest dimension.
Each nonzero coordinate projection $p_j|T$
is injective, and its image has no proper
nonzero invariant subspace.
Two nonzero projections give the explicit
mutually inverse intertwiners
\[
 p_l\circ(p_j|T)^{-1}:p_j(T)\longrightarrow p_l(T),
 \qquad p_j\circ(p_l|T)^{-1}:p_l(T)\longrightarrow p_j(T).
                                                               \tag{F.30.13}
\]
The three comparisons exclude this for
$j\ne l$. Hence there is only one nonzero
projection, so $T$ lies in one $Q_j$.
There $\Phi$ is the injective individual
restriction, a contradiction.
This proves the summed $H^1$ injection.

Zero summed restriction of cocycles therefore
forces each cocycle to be principal.
Since all characters agree on $K$ and
$C\ne\{1\}$, the exact kernel and its dimension are
\[
 \left\{\bigl(b_j(\chi_j-1)\bigr)_j:
                      \sum_jb_j=0\right\},\qquad
                     \dim\ker=m-1.
                                                               \tag{F.30.14}
\]
The converse inclusion is direct restriction.
Each principal vector is nonzero, and each
continuous cocycle space has seven generator
coordinates. Thus the stated rank is
$7m-(m-1)=6m+1$.

Apply this result to the tuple
$(d^{q_1^j})_j$ of one cocycle in the trace
kernel of F.26. The cocycle $d$ is principal,
say $d=b(\chi-1)$.
Its Frobenius powers have parameters
$b^{q_1^j}$, uniquely because their principal
vectors are nonzero. Their parameter sum
is exactly the field trace. Consequently
\[
 \operatorname{Tr}_{k_0/k_1}(d|K)=0
 \quad\Longleftrightarrow\quad
 d=b(\chi-1),\qquad \operatorname{Tr}_{k_0/k_1}(b)=0.
                                                               \tag{F.30.15}
\]
The reverse implication follows directly from
$\chi(K)\subset k_1$.
The explicit conjugation and coefficient
homomorphism of F.26 sends every proper
finite affine marginal image to a nonprincipal
element of this trace kernel over $k_0$.
Those maps preserve the scalar character
and both scalar fields. They therefore
preserve the product-one condition and
$k_0=k_1(r)$ of (F.30.1).
Equation (F.30.15) excludes their resulting
nonprincipal element. This proves (F.30.2)
for the negative marginal over every original
finite coefficient field, including fields
larger than $k_0$.

\emph{Integral evidence and actual positive-sign transport.}
The coefficient producer applies precisely
(F.30.4) to the frozen F.24 universal matrices
for the three original sources.
The separate verifier imports neither that
producer nor an algebra package. It uses the
reviewed integer Laurent arithmetic to
reconstruct the three actual decorated Artin
words and compares every entry of their full
seven-coordinate actions with the producer.
Its $45$ exact checks include both inverse
maps in (F.30.8), the boundary kernel and
section, the principal-vector identities,
(F.30.9) and every boundary value in (F.30.10).
The identities have integer coefficients;
the only nonmonomial principal localization
is $r-1$, whose nonvanishing was proved from
the exact field equality before its use.
The evidence manifest binds the statement,
programs, original words, coefficient inputs,
outputs and completed dependencies.

For a positive representation, precomposition
with the actual continuous isomorphism
$\mathcal B_+:V_-\to V_+$ sends the values
at $Y_i^+$ to precisely the negative
coordinates used above.
F.23.13 supplies its inverse and carries
$\mathcal K_-,R_-$ onto their positive
counterparts. Thus it preserves the scalar
images defining $k_0,k_1$, the first scalar
$r$ and all conditions (F.30.1).
Equality of the two negative images gives
equality of the original two positive images
through these specified maps. This proves
(F.30.2) for both signs.
\end{proof}

\begin{claim}{F.31}{Proved}{Finite affine marginal image equality on the full product-one character locus}
Use the actual bases $Y_i^-=X_i$ and
$Y_i^+=\mathcal B_+(X_i)$ under the based
isomorphism of F.23.13.
For every finite field $k$ and continuous
$\rho:V_\pm\to\operatorname{Aff}_1(k)$ whose
scalar character $\chi$ belongs to the exact locus
\[
                         \prod_{i=1}^7\chi(Y_i^\pm)=1,
                                                               \tag{F.31.1}
\]
the actual marginal images satisfy
\[
                         \rho(\mathcal K_\pm)=\rho(R_\pm).
                                                               \tag{F.31.2}
\]
All characteristics, all scalar-field extension
degrees and all larger original coefficient
fields are included.

Put $C=\chi(R_\pm)$,
$k_1=\mathbb F_p(C)$ and
$k_0=\mathbb F_p(\chi(V_\pm))$.
On the part $C\ne\{1\}$, set
$q_1=|k_1|$, $m=[k_0:k_1]$ and
$\chi_j=\chi^{q_1^j}$ for $0\leq j<m$.
The summed continuous $H^1$ restriction
for these $m$ characters is injective.
The corresponding map on cocycles has
exactly the principal tuples whose
parameters sum to zero as its kernel,
and has rank $6m+1$ on $7m$ coordinates.
\end{claim}
\passage{Part F root-adjoining intersection question; Section 30, folios 194-196, scans 383-387; Section 31, folios 198-199 verso, scans 391-394; original F7/F19 residual subgroup-membership reduction}


\begin{proof}
\emph{Actual scalar fields and completed scalar cases.}
Work first with
$K=\mathcal K_-\leq R=R_-\leq V=V_-$.
The negative character has the unique form
\[
 \begin{gathered}
 \chi=(r,s,t,u,v,w,c^{-1}),\qquad c=rstuvw,\\
 C=\chi(K)=\chi(R)=\langle u,v,w,c^{-1}\rangle,\\
 k_1=\mathbb F_p(u,v,w,c),\qquad
 k_0=k_1(r,t),\qquad s=\frac{c}{rtuvw}.
 \end{gathered}                                                \tag{F.31.3}
\]
All six displayed parameters are nonzero.
Here the equality of the scalar images is
the image of
$\overline{\langle\!\langle K\rangle\!\rangle}_R=R$
under the actual abelian scalar homomorphism;
F.21 and the marginal surjection of F.22
give precisely this normal-generation equality.
The actual deleted generators $X_4,\ldots,X_7$
give the displayed generating set for $C$.
The two descriptions of $k_0$ follow from
the last identity in (F.31.3), in both directions.

For $C=\{1\}$ the image of $R$ is an
abelian translation group, and the same
normal-generation map proves (F.31.2).
For $C\ne\{1\}$ and $k_0=k_1$, F.25
gives the image equality, while F.24
gives the asserted individual restriction
injection and cocycle kernel.
For the rest of the proof take
$C\ne\{1\}$ and $m>1$.

\emph{Coefficient maps and actual source actions.}
Let $P_0=X_1\cdots X_7$.
The coefficient homomorphism on the original
seven-variable derivations is exactly
\[
 \begin{aligned}
 \mathbb Z[t_1^{\pm1},\ldots,t_7^{\pm1}]&\longrightarrow
 A=\mathbb Z[r^{\pm1},s^{\pm1},t^{\pm1},u^{\pm1},v^{\pm1},w^{\pm1}],\\
 (t_1,t_2,t_3,t_4,t_5,t_6,t_7)&\longmapsto
          (r,s,t,u,v,w,(rstuvw)^{-1}).
 \end{aligned}                                                \tag{F.31.4}
\]
An exponent vector $(a_i)$ therefore maps to
$r^{a_1-a_7}s^{a_2-a_7}t^{a_3-a_7}
 u^{a_4-a_7}v^{a_5-a_7}w^{a_6-a_7}$.
The second chart applies the additional
coefficient homomorphism $r\mapsto1$.
Every computation below is an integral
Laurent identity on these specified rings.

These derivations are evaluated on the actual
discrete free words, with
$d(xy)=d(x)+\chi(x)d(y)$ and
$d(x^{-1})=-\chi(x)^{-1}d(x)$.
Every seven-tuple of values in $k_0$ gives
a homomorphism from the discrete free group
to its finite affine group, hence extends
uniquely and continuously to $V=\widehat F_7$.
The continuous cocycle space is thus exactly
$k_0^7$.
Its principal line is generated by
$\delta=(r,s,t,u,v,w,c^{-1})^{\mathsf T}-\mathbf1$.

Retain the eleven original F.23 automorphisms
$B_l(x)=u_l\alpha_l(x)u_l^{-1}$,
where throughout this proof $0\leq l\leq10$
is the original zero-based source index.
Thus $l=0$ is \texttt{old0}; indices $1,\ldots,5$
are \texttt{old\_1\_0},\ldots,\texttt{old\_1\_4},
and indices $6,\ldots,10$ are
\texttt{old\_-1\_0},\ldots,\texttt{old\_-1\_4}.
The source position in a one-based list is
always $l+1$.
Both $B_l$ and its inverse preserve $K$.
The pure Artin action fixes every abelian
character and the actual boundary word.
Write $A_l$ for precomposition on the seven
cocycle coordinates. Then
\[
 \begin{gathered}
 A_l\delta=\delta,\qquad
 D=(1,r,rs,rst,rstu,rstuv,c),\qquad D\delta=0,\\
 DA_l=\mu_lD,\qquad \mu_l=\chi(u_l),\\
 (\mu_0,\ldots,\mu_{10})=
 (1,1,rt,1,1,1,u^{-2},t^{-1},v^{-2},w^{-2},c^2).
 \end{gathered}                                                \tag{F.31.5}
\]
For the boundary identity, the exact cocycle
formula gives
$d(B_l(P_0))=\chi(u_l)d(P_0)$,
because $\chi(P_0)=1$ and $\alpha_l(P_0)=P_0$.
The displayed character values are evaluated
from the actual decorations in the eleven
labelled Artin words.

The word recurrence used for the finite
certificate retains both each meridian
image $x_i$ and its decoration $a_i$.
At a positive crossing of positions $j,j+1$ it sends
$(x_j,x_{j+1};a_j,a_{j+1})$ to
$(x_jx_{j+1}x_j^{-1},x_j;x_ja_{j+1},a_j)$;
at a negative crossing it sends them to
$(x_{j+1},x_{j+1}^{-1}x_jx_{j+1};\allowbreak
 a_{j+1},x_{j+1}^{-1}a_j)$.
The eighth meridian is evaluated at the
identity. At the end $u_l=a_8^{-1}$,
and its conjugation gives the seven rows
of $A_l$. The certificate checks the
label permutation, the erased meridian,
all seven preserved character values and
the full boundary product after each word.
These are the original source maps used
in every matrix below.

\emph{The retained principal chart for $r\ne1$.}
The principal quotient and its section are
\[
 \begin{aligned}
 P_1(z)&=(z_i-\delta_i(r-1)^{-1}z_1)_{i=2}^7,
 &S_1(x)&=(0,x_1,\ldots,x_6)^{\mathsf T},\\
 P_1S_1&=I_6,
 &S_1P_1(z)&=z-\delta(r-1)^{-1}z_1.
 \end{aligned}                                                \tag{F.31.6}
\]
They identify
$Q_\chi=H^1_{\rm cont}(V,k_{0,\chi})$
with the coordinates of the original
$X_2,\ldots,X_7$ in that order.
The inverse formulas prove the identification
and its exact kernel, namely $k_0\delta$.
Put $b=DS_1=(r,rs,rst,rstu,rstuv,c)$.
The boundary quotient $\pi(z)=bz$ is
well-defined since $D\delta=0$.
Its kernel is the invariant subspace $W_\chi$
of dimension five, using the explicit maps
\[
 \begin{gathered}
 \iota(x_1,\ldots,x_5)=
 (-sx_1-stx_2-stux_3-stuvx_4-stuvwx_5,x_1,\ldots,x_5),\\
 \sigma(a)=r^{-1}ae_1,\qquad
 \pi\sigma=I_1,\qquad
 z=\iota(z_2,\ldots,z_6)+\sigma\pi(z),\\
 J=(e_2-se_1,e_3-ste_1,e_4-stue_1,
              e_5-stuve_1,e_6-stuvwe_1,r^{-1}e_1),\\
 J^{-1}(z)=(z_2,z_3,z_4,z_5,z_6,bz).
 \end{gathered}                                                \tag{F.31.7}
\]
The displayed decomposition proves both
inverse basis identities; it does not
require $\sigma(k_0)$ to be invariant.
In this basis the actual source block is
\[
 T_l=J^{-1}P_1A_lS_1J=
       \begin{pmatrix}G_l&*\\0&\mu_l\end{pmatrix},
                    \qquad G_0=rI_5.
                                                               \tag{F.31.8}
\]
The scalar identity for $G_0$ is also the
literal-word identity (F.30.9).
The present certificate checks every entry
of all eleven blocks directly from the
original words. Its exact quotient check is
\[
 (r-1)(A_lS_1J-S_1JT_l)
       =\delta\,\operatorname{row}_1(A_lS_1J).
                                                               \tag{F.31.9}
\]
The section has zero first row, so applying
$P_1$ to this identity and cancelling
the proved nonzero $r-1$ gives exactly
(F.31.8), with no coordinate substitution.

\emph{Eight exact moving-root rows on $W_\chi$.}
The following table specifies a list of base
roots $a_{l,h}$ and one moving root $b_l$.
Let $Q_l(Z)=\prod_h(Z-a_{l,h})$.
Repeated base roots are retained as repeated
factors of this polynomial.
\[
 \begin{array}{c|c|c|c}
 l&\text{source position}&(a_{l,h})_h&b_l\\\hline
 6&7&(1,(vw)^{-1})&c/t\\
 8&9&(1,c/u)&(tw)^{-1}\\
 9&10&(1,c/v)&(tu)^{-1}\\
 10&11&(1,(uw)^{-1})&(tv)^{-1}\\
 1&2&(1,(vw)^{-1},(r^2u^2)^{-1})&c/t\\
 3&4&(1,c/u,(r^2v^2)^{-1})&(tw)^{-1}\\
 4&5&(1,c/v,(r^2w^2)^{-1})&(tu)^{-1}\\
 5&6&(1,(uw)^{-1},(c/r)^2)&(tv)^{-1}
 \end{array}.
                                                               \tag{F.31.10}
\]
For the displayed rows below the exact
integral identities are
\[
 \begin{gathered}
 Q_l(G_l)(G_l-b_lI_5)=0,\qquad q_lG_l=b_lq_l,\\
 q_lQ_l(G_l)=\tau_lq_l,
       \qquad \tau_l=Q_l(b_l)=\prod_h(b_l-a_{l,h}).
 \end{gathered}                                                \tag{F.31.11}
\]
The last equality follows by polynomial
evaluation from the second, and is also
checked directly. No conclusion about
diagonalizability is needed.

Here are all the rows used in (F.31.11).
To keep their coordinates explicit, set
$A=rst$, $B=t-1$, $H=rt-t+1$ and
$L=ru+(r-1)v-r$. These are the displayed
Laurent expressions, with $c=Auvw$.
Each displayed column is the transpose
of the named row:
\[
 \begin{gathered}
 q_6^{\mathsf T}=
 \begin{pmatrix}
 -A(c-1)H\\-B(c-1)(AH+1)\\
 -AuB(c-1)(H-u)\\-AuB(c-1)\bigl(v(H-u)+u\bigr)\\
 -cB\bigl((c-1)(H-u)+cu/v\bigr)
 \end{pmatrix},\qquad
 q_8^{\mathsf T}=
 \begin{pmatrix}w-1\\B(w-1)\\uB(w-1)\\-uvB\\0\end{pmatrix}.
 \end{gathered}                                                \tag{F.31.12}
\]
\[
 \begin{gathered}
 q_9^{\mathsf T}=
 \begin{pmatrix}
 A(u-1)\\B\bigl(A(u-1)-1\bigr)\\AuB(u-1)\\AuvB(u-1)\\cB(u-1)
 \end{pmatrix},\qquad
 q_{10}^{\mathsf T}=
 \begin{pmatrix}
 v-1\\B(v-1)\\-uB\bigl((r-1)v+1\bigr)\\
 -uv(r-1)B(v-1)\\-uvw(r-1)B(v-1)
 \end{pmatrix}.
 \end{gathered}                                                \tag{F.31.13}
\]
\[
 \begin{gathered}
 q_1^{\mathsf T}=
 \begin{pmatrix}c-1\\B(c-1)\\uB(c-1)\\uvB(c-1)\\cuvwB\end{pmatrix},\qquad
 q_3^{\mathsf T}=
 \begin{pmatrix}
 stv(w-1)\\B(w-1)(stv+1)\\-stuv(r-1)B(w-1)\\
 -stuvB\bigl((w-1)L+ru\bigr)\\-stuvwB(w-1)L
 \end{pmatrix}.
 \end{gathered}                                                \tag{F.31.14}
\]
\[
 \begin{gathered}
 q_4^{\mathsf T}=
 \begin{pmatrix}
 -A(u-1)H\\-B\bigl(AH(u-1)-1\bigr)\\
 -AuB(u-1)H\\-AuvB(u-1)H\\-cB(u-1)H
 \end{pmatrix},\qquad
 q_5^{\mathsf T}=
 \begin{pmatrix}
 -Aw(v-1)\\-AwB(v-1)\\-B\bigl(Auw(v-1)-1\bigr)\\
 -B(v-1)(c-1)\\-cwB(v-1)
 \end{pmatrix}.
 \end{gathered}                                                \tag{F.31.15}
\]
With rows in the exact order shown, the
four determinant identities are
\[
 \begin{aligned}
 \det(q_6,q_8,q_9,q_{10},q_1)
   &=cu^2vw(r-1)B^4(c-1)(rtv+uv-u),\\
 \det(q_6,q_8,q_9,q_{10},q_3)
   &=c^2u^2(r-1)B^4(w-1),\\
 \det(q_6,q_8,q_9,q_{10},q_4)
   &=rtc^2u^2(r-1)B^4(u-1),\\
 \det(q_6,q_8,q_9,q_{10},q_5)
   &=-cuvw(r-1)B^4(v-1)F,\\
 F&=(c-1)(rt+u)-cu/v.
 \end{aligned}                                                \tag{F.31.16}
\]
The certificate multiplies the original
word matrices, checks every entry of
(F.31.11), compares all eight rows with
(F.31.12)--(F.31.15), and expands each
determinant by its signed permutations.
The resulting identities hold in $A$,
before reduction in any characteristic.

\emph{The exact common fixed field and the generic comparisons.}
Let $\sigma(a)=a^{q_1}$ on $k_0$.
This automorphism has order exactly $m$:
its $m$th power is the identity, and for
$0<a<m$ the polynomial $X^{q_1^a}-X$
has degree smaller than $|k_0|$ and cannot
vanish on all of $k_0$.
Write $r_j=\sigma^j(r)$ and $t_j=\sigma^j(t)$.
The coefficient map $d\mapsto d^{q_1^j}$
is additive and semilinear, preserves the
actual cocycle rule and principal line,
commutes with each source precomposition,
and has inverse $d\mapsto d^{q_1^{m-j}}$.
It transports each exact basis map above
to $Q_j=H^1_{\rm cont}(V,k_{0,\chi_j})$
and its core $W_j$.
It also proves $k_0=k_1(r_j,t_j)$.

Take distinct $j,l$ and source-invariant
subspaces of $W_j,W_l$ with an intertwining
isomorphism. Source $0$ first forces
$r_j=r_l$, since it acts there by
$r_jI$ and $r_lI$.
Choose the unique $a\in\{1,\ldots,m-1\}$
congruent to $l-j$ modulo $m$, and define
the actual subfield
\[
       L_0=\{x\in k_0:\sigma^a(x)=x\}.
                                                               \tag{F.31.17}
\]
It contains $k_1$ and $r_j$.
The element $t_j$ lies outside $L_0$:
otherwise $\sigma^a$ would fix both
$r_j,t_j$ and $k_1$, hence their whole
field $k_0$, contradicting its proved
nonidentity. In particular
$t_j\ne t_l$ and $t_j-1\ne0$.

All the base roots in (F.31.10), evaluated
at the $j$th or $l$th character, coincide
and lie in $L_0$.
Each moving root is $\eta/t_j$ or
$\eta/t_l$ for the same nonzero
$\eta\in k_1$.
Thus the two moving roots are distinct,
both lie outside $L_0$, and every factor
of $\tau_{l',j}$ and $\tau_{l',l}$ is
nonzero, for every source $l'$ in that table.
This proves the exact coefficient-field
inclusions used in the following cancellation.

Identify the two subspaces through the
specified intertwiner. The common
source-$l'$ endomorphism satisfies the
two annihilators
$(Z-b_{l',j})Q_{l'}(Z)$ and
$(Z-b_{l',l})Q_{l'}(Z)$.
Subtracting gives
$(b_{l',l}-b_{l',j})Q_{l'}(Z)=0$ on it.
The nonzero scalar just proved shows
that $Q_{l'}(G_{l',j})$ kills the first
subspace. By (F.31.11) and
$\tau_{l',j}\ne0$, the row $q_{l',j}$
kills that subspace as well.
This argument permits all equalities
among the base roots, including every
characteristic-dependent coincidence.

The first four rows in (F.31.16) therefore
vanish on the subspace, as do all four
possible last rows.
At least one of those four square row
matrices is invertible. Indeed, for
$c\ne1$ use the first determinant:
its last factor cannot vanish, because
$r_jt_jv+uv-u=0$ would put
$t_j=(u-uv)/(r_jv)$ in $L_0$.
Its other factors are nonzero by
$r_j\ne1$, $t_j\notin L_0$, and $c\ne1$.
For $c=1$, choose the second determinant
when $w\ne1$, the third when $w=1$
and $u\ne1$, and the fourth when
$w=u=1$. In this last case $v\ne1$
because $C\ne\{1\}$, and its factor
$F$ is exactly $-u/v\ne0$.
All omitted prefactors in this discussion
are displayed Laurent monomial units.
The invertible row matrix forces the
subspace to be zero. This excludes any
nonzero common core subspace for distinct
Frobenius indices on the entire $r\ne1$ chart.

For a nonzero source-invariant subspace
of $Q_j$ with no proper nonzero invariant
subspace, its intersection with $W_j$
either equals the whole subspace or is
zero. In the latter case $\pi_j$ is an
injective intertwiner onto the boundary
character $\mu_j$.
A core subspace and a boundary subspace
cannot be isomorphic: their source-$0$
scalars are $r_j\ne1$ and $1$.
Two boundary characters at distinct
indices cannot coincide. Source $7$
would give $t_j=t_l$, and source $2$
would then give $r_j=r_l$. Since
$k_0=k_1(r_j,t_j)$, this would make
$\sigma^{l-j}$ the identity.
The inverse of an intertwining isomorphism
covers the reversed mixed comparison.
These are all three unordered types on
this chart, obtained through the exact
boundary quotient rather than an assumed
invariant complementary line.

\emph{The retained chart for $r=1$.}
Now $k_0=k_1(t)$ by (F.31.3).
The proper extension implies $t\notin k_1$,
so $t-1\ne0$.
Use the actual third principal coordinate:
\[
 \begin{aligned}
 P_3(z)&=(z_i-\delta_i(t-1)^{-1}z_3)_{i\in I},
       &I&=(1,2,4,5,6,7),\\
 S_3(x)&=(x_1,x_2,0,x_3,x_4,x_5,x_6)^{\mathsf T},\\
 P_3S_3&=I_6,
       &S_3P_3(z)&=z-\delta(t-1)^{-1}z_3.
 \end{aligned}                                                \tag{F.31.18}
\]
Thus $Q_\chi$ has the coordinates of the
original generators in the order $I$.
Since $\chi(X_1)=1$, evaluation at $X_1$
also descends through the principal line.
Put $b=(1,1,st,stu,stuv,c)$ and define
the exact two-coordinate quotient
\[
 \begin{gathered}
 \pi(z)=(bz,z_1),\qquad
 \iota(x_1,x_2,x_3,x_4)=
 (0,-stx_1-stux_2-stuvx_3-cx_4,x_1,x_2,x_3,x_4),\\
 \sigma(a,h)=(h,a-h,0,0,0,0),\qquad
 \pi\sigma=I_2,\qquad
 z=\iota(z_3,z_4,z_5,z_6)+\sigma\pi(z),\\
 J=(e_3-ste_2,e_4-stue_2,e_5-stuve_2,
                         e_6-ce_2,e_2,e_1-e_2),\\
 J^{-1}(z)=(z_3,z_4,z_5,z_6,bz,z_1).
 \end{gathered}                                                \tag{F.31.19}
\]
These formulas prove both basis inverses
and identify the invariant core
$W_\chi=\ker\pi$ with $k_0^4$.
Its quotient coordinates are the actual
boundary evaluation followed by the
actual $X_1$ evaluation.
The original-word identities are
\[
 \begin{gathered}
 T_l=J^{-1}P_3A_lS_3J=
 \begin{pmatrix}G_l&*\;&*\\0&\mu_l&0\\0&0&\nu_l\end{pmatrix},\\
 (t-1)(A_lS_3J-S_3JT_l)
     =\delta\,\operatorname{row}_3(A_lS_3J),\\
 (\nu_0,\ldots,\nu_{10})=
 (1,u^{-2},t^{-1},v^{-2},w^{-2},c^2,1,t,1,1,1).
 \end{gathered}                                                \tag{F.31.20}
\]
Here $\mu$ is (F.31.5) with $r=1$.
For each word the first row of $A_l$
is exactly $(\nu_l,0,\ldots,0)$;
this proves the $X_1$ quotient before
passing to cohomology. The cleared
identity proves the full block via
$P_3$ and the nonzero scalar $t-1$.

\emph{Four moving rows and all generic type comparisons for $r=1$.}
Use the first four sources of (F.31.10),
namely $l=6,8,9,10$, with their same
base and moving roots and $r=1$.
Their core matrices are now $4$ by $4$.
For the rows below they satisfy precisely
\[
 Q_l(G_l)(G_l-b_lI_4)=0,\qquad
 q_lG_l=b_lq_l,\qquad q_lQ_l(G_l)=Q_l(b_l)q_l.
                                                               \tag{F.31.21}
\]
The actual rows are
\[
 \begin{gathered}
 q_6^{\mathsf T}=
 \begin{pmatrix}
 -(st+1)(c-1)\\stu(u-1)(c-1)\\
 stu(c-1)(uv-u-v)\\c(cu-cu/v-c-u+1)
 \end{pmatrix},\qquad
 q_8^{\mathsf T}=
 \begin{pmatrix}1-w\\-u(w-1)\\uv\\0\end{pmatrix},\\
 q_9^{\mathsf T}=
 \begin{pmatrix}-stu+st+1\\-stu(u-1)\\-stuv(u-1)\\-c(u-1)\end{pmatrix},\qquad
 q_{10}^{\mathsf T}=
 \begin{pmatrix}1-v\\u\\0\\0\end{pmatrix},\\
              \det(q_6,q_8,q_9,q_{10})=u^2c^2.
 \end{gathered}                                                \tag{F.31.22}
\]
This determinant is a Laurent unit in
every characteristic and at every allowed
value of $u,v,w,c$.

All $t_j$ are distinct and lie outside
$k_1$, since $k_0=k_1(t)$ and the
$k_1$-Frobenius has order $m$.
The moving roots $\eta/t_j$ therefore
lie outside $k_1$, and differ for
different indices; all base roots lie
in $k_1$. Subtracting the annihilators
as above, then using (F.31.21) and its
nonzero $Q_l(b_l)$, shows that a common
core subspace is killed by all four
rows of (F.31.22). Its unit determinant
excludes every such nonzero subspace.

The same rows exclude any core subspace
isomorphic to a boundary or $X_1$
character from any Frobenius factor.
Such a subspace is a line, and for
each of the four negative sources its
scalar $\lambda_l$ lies in $k_1$:
the boundary values are
$(u^{-2},v^{-2},w^{-2},c^2)$ and
the $X_1$ values are all $1$.
The annihilator on the line gives
$(\lambda_l-b_{l,j})Q_l(\lambda_l)=0$.
The first factor is nonzero, so
$Q_l(G_{l,j})$ kills the line.
Equation (F.31.21) and $Q_l(b_{l,j})\ne0$
make all four rows vanish on it,
contradicting (F.31.22).
Inverse intertwiners cover both orientations.

A simple invariant subspace of $Q_j$
either lies in $W_j$, or projects
injectively under (F.31.19) to a simple
subspace of its two-dimensional diagonal
quotient. A nonzero projection of the
latter to either one-dimensional
coordinate is injective; its character
is therefore $\mu_j$ or $\nu_j$.
This argument uses only the displayed
coordinate maps and does not assume
semisimplicity of a larger source module.
For distinct indices, $\mu_j=\mu_l$
is excluded by source $7$, and
$\nu_j=\nu_l$ is excluded by source $2$.
These core comparisons and the two
equal-type comparisons leave only
the cross-type equality
$\mu_j=\nu_l$.
Sources $1,3,4,5$ force exactly
\[
                  u^2=v^2=w^2=c^2=1.
                                                               \tag{F.31.23}
\]
In characteristic two this would make
all four parameters equal to $1$ and
give $C=\{1\}$, already excluded.
Thus the remaining field has odd
characteristic, its four parameters
are signs not all $1$, $k_1=\mathbb F_p$,
and $t_j\notin\mathbb F_p$.

\emph{All fifteen sign strata: exact absence of a boundary eigenline.}
For each sign tuple write
$\epsilon=cuvw$.
The coefficient homomorphism is explicitly
\[
 \begin{aligned}
 \mathbb Z[s^{\pm1},t^{\pm1},u^{\pm1},v^{\pm1},w^{\pm1}]
    &\longrightarrow\mathbb Z[t^{\pm1}],\\
 s^at^bu^dv^ew^f
    &\longmapsto\epsilon^a u^dv^ew^f t^{b-a},
       \qquad s\longmapsto\epsilon/t,
 \end{aligned}                                                \tag{F.31.24}
\]
where on the right $u,v,w,c$ are the
specified integer signs.
Addition of exponent vectors multiplies
both their signs and their powers of
$t$, so this is an exact ring
homomorphism, including negative exponents.
It is applied to the complete actual
word matrices and all basis maps.

Let $Q$ have the four rows in (F.31.22)
after this map. Its entries are integers,
and $\det Q=u^2c^2=1$.
For $l=6,8,9,10$ let
$\tau_l=Q_l(b_l)$ and define an extended
six-coordinate row by the exact identity
\[
 \widetilde q_l=
       \tau_l^{-1}(q_l,0,0)Q_l(T_l),\qquad
 \widetilde q_l|_{k_0^4}=q_l.
                                                               \tag{F.31.25}
\]
The second equality follows from
(F.31.21), and its cleared integral
form is checked before any division.
Every base root here is a sign, while
each moving root is a sign times $t^{-1}$.
All factors of $\tau_l$ are nonzero
for $t\notin\mathbb F_p$.
The inverse in (F.31.25) is thus defined
in every remaining field, including
all odd characteristics.

On a boundary eigenline all four
negative-source scalars are $1$.
Consequently $Q_l(T_l)$, and hence
each $\widetilde q_l$, vanishes there.
The last two quotient scalars at
source $2$ are $t$ and $t^{-1}$.
They are distinct because $t\notin\mathbb F_p$.
A boundary eigenline projecting to
the boundary quotient therefore has
last coordinates $(a,0)$ with $a\ne0$;
the nonzero boundary coordinate also
follows from the already proved
exclusion of a boundary line in the core.
Scale by the unique $a^{-1}$ to obtain
last coordinates $(1,0)$.

If $E$ is the four-row matrix with
rows $\widetilde q_l$, its first four
columns are $Q$. The exact constraints
then force the first four coordinates
to be the unique vector
$x=-Q^{-1}E_{\mathrm{boundary}}$.
Both inverses of $Q$ are integral,
as $\det Q=1$.
The table gives this forced vector
and a nonzero coordinate of the
actual residual
$T_l(x,1,0)^{\mathsf T}-\mu_l(x,1,0)^{\mathsf T}$.
The source index $l$ is zero-based
as before; the coordinate $a$ in this
table is one-based in the six-coordinate
basis (F.31.19).
\[
\begin{array}{c|c|c|c|c}
 (u,v,w,c)&x&l&a&\text{residual coordinate}\\\hline
 (1,1,1,-1)&(0,0,0,0)&1&4&-2(t+1)/t\\
 (1,1,-1,1)&(0,0,-2,0)&3&3&2(t+1)/t\\
 (1,1,-1,-1)&(0,0,2,0)&1&3&-4\\
 (1,-1,1,1)&(0,0,0,0)&5&2&-2(t+1)/t\\
 (1,-1,1,-1)&(0,0,0,0)&1&2&4\\
 (1,-1,-1,1)&(0,0,-2,0)&3&2&4/t\\
 (1,-1,-1,-1)&(0,0,2,0)&1&2&4\\
 (-1,1,1,1)&(-2,0,0,0)&4&1&2(t+1)/t\\
 (-1,1,1,-1)&(-2,0,0,-4)&1&1&4\\
 (-1,1,-1,1)&(2,0,2,0)&3&1&4(2t-1)/t\\
 (-1,1,-1,-1)&(2,0,6,4)&1&1&4\\
 (-1,-1,1,1)&(-2,-4,0,0)&4&1&2(3t-1)/t\\
 (-1,-1,1,-1)&(-2,-4,0,-4)&1&1&-4\\
 (-1,-1,-1,1)&(2,4,2,0)&3&1&4(2t-1)/t\\
 (-1,-1,-1,-1)&(2,4,6,4)&1&1&-4
\end{array}.
                                                               \tag{F.31.26}
\]
All fifteen tuples other than
$(1,1,1,1)$ occur exactly once.
The constants $2,4$ and the parameter
$t$ are units in these fields.
The factor $t+1$ cannot vanish for
$t\notin\mathbb F_p$.
The factor $2t-1$ has nonzero leading
coefficient in every odd characteristic,
and its vanishing would put $t$ in
the prime field. For $3t-1$, characteristic
three gives the nonzero constant $-1$;
every other odd characteristic again
gives a prime-field root.
Thus every displayed residual is
nonzero in every remaining field.

The certificate checks both inverses of
each specialized $Q$, the cleared
identities (F.31.25), the four equations
$E(x,1,0)^{\mathsf T}=0$, and each
displayed residual directly in
$\mathbb Z[t^{\pm1}]$.
Uniqueness from $Q^{-1}$ makes the
table exhaustive for all possible
normalized boundary eigenvectors.
There is consequently no boundary
eigenline on any of the fifteen strata.
This excludes a pair of isomorphic
subspaces of types $\mu_j,\nu_l$;
the reverse orientation is excluded
by the inverse intertwiner.
Together with the earlier comparisons,
this accounts for all six unordered
type pairs $W/W$, $W/\mu$, $W/\nu$,
$\mu/\mu$, $\nu/\nu$ and $\mu/\nu$.

\emph{Summed restriction through the actual invariant maps.}
All characters $\chi_j$ restrict to
the same character on $K$, because
$\chi(K)=C\subset k_1^\times$.
The actual restriction therefore defines
\[
 \Phi:\bigoplus_{j=0}^{m-1}Q_j
  \longrightarrow H^1_{\rm cont}(K,k_{0,\chi|K}),\qquad
       ([d_j])_j\longmapsto\sum_j[d_j|K].
                                                               \tag{F.31.27}
\]
It commutes with each of the eleven
source operators, since their actual
automorphisms preserve $K$ and all
the characters. Every individual
restriction is injective by F.24,
as $\chi_j(R)=C\ne\{1\}$.

Suppose that $\ker\Phi$ were nonzero,
and choose inside it a nonzero invariant
subspace $T$ of smallest dimension.
Finite dimension makes this choice
available. Each nonzero coordinate
projection $p_j|T$ is injective:
its kernel is an invariant subspace
of $T$, and cannot be all of $T$.
Its image has no proper nonzero
invariant subspace, by pulling such
a subspace back through this inverse.
Two nonzero coordinate projections
give the explicit mutually inverse
intertwiners
\[
 \begin{gathered}
 p_l\circ(p_j|T)^{-1}:p_j(T)\longrightarrow p_l(T),\\
 p_j\circ(p_l|T)^{-1}:p_l(T)\longrightarrow p_j(T).
 \end{gathered}                                                \tag{F.31.28}
\]
The complete comparisons on the
$r\ne1$ and $r=1$ charts exclude
such isomorphisms for $j\ne l$.
Hence exactly one projection is
nonzero. The subspace $T$ then lies
in one factor, where $\Phi$ is the
injective individual restriction,
a contradiction. This proves the
summed $H^1$ injection.

Zero summed restriction of cocycles
now forces every cocycle to be
principal. Since the common restricted
character has nontrivial image $C$,
the exact kernel is
\[
 \left\{\bigl(b_j(\chi_j-1)\bigr)_{j=0}^{m-1}:
                      \sum_{j=0}^{m-1}b_j=0\right\}.
                                                               \tag{F.31.29}
\]
Its converse inclusion follows directly
by restriction. Each principal vector
is nonzero, so this space has dimension
$m-1$. The cocycle domain has dimension
$7m$ by its actual generator coordinates,
and the claimed rank is $6m+1$.

For one cocycle $d$, apply this result
to the actual coefficient tuple
$(d^{q_1^j})_{j=0}^{m-1}$.
It gives the exact trace-kernel identity
\[
 \operatorname{Tr}_{k_0/k_1}(d|K)=0
 \quad\Longleftrightarrow\quad
 d=b(\chi-1),\qquad
                    \operatorname{Tr}_{k_0/k_1}(b)=0.
                                                               \tag{F.31.30}
\]
Indeed the forward implication makes
$d$ principal. Its Frobenius powers
then have the unique parameters
$b^{q_1^j}$, whose sum is the trace;
(F.31.29) gives the stated vanishing.
The reverse implication uses
$\chi(K)\subset k_1$ and the trace's
$k_1$-linearity.

The explicit conjugation, translation
projection and trace-pairing coefficient
maps in F.26 send every proper finite
affine marginal image to a nonprincipal
element of this trace kernel over the
actual field $k_0$.
Those maps retain the scalar character
and both scalar fields, so they retain
(F.31.1). Equation (F.31.30) excludes
the resulting nonprincipal element.
This proves (F.31.2) for the negative
marginal, for every original finite
coefficient field including fields
larger than $k_0$.

\emph{Integral certificate, continuity and sign transport.}
The independently structured verifier
imports none of the four new producers
and no symbolic algebra package.
It reconstructs all eleven literal
decorated Artin words on each of the
two coefficient charts using the
previously reviewed integer Laurent
word arithmetic. It checks the full
seven-coordinate matrices and the
quotient blocks by the cleared
intertwining identities (F.31.9) and
(F.31.20), independently of the
producer's Laurent-division algorithm.
Its $1603$ exact checks include every
displayed basis inverse, all twelve
moving-root annihilators and direct
left-eigenrow identities, the five
row determinants, and all fifteen
coefficient maps, forced lifts and
residual obstructions.
The nonmonomial principal pivots are
exactly $r-1$ and $t-1$ on their
respective charts. Their nonvanishing,
and every further factor cancelled
in the proof, were established in
the actual finite fields before use.
The evidence manifest binds the
source, programs, original words,
coefficient inputs, outputs and
completed dependencies.

Every cocycle constructed from generator
values is the unique continuous extension
of a homomorphism to a finite affine
group. The source automorphisms are
continuous by F.23, and all maps between
the finite coefficient fields are
continuous. Their cocycle compositions
and the restrictions to the actual
closed subgroup $K$ are therefore
continuous throughout the argument.
For the positive marginal, precompose
with the actual based isomorphism
$\mathcal B_+:V_-\to V_+$ of F.23.13.
Its inverse is supplied there, and
it sends $X_i$ precisely to $Y_i^+$
and carries $\mathcal K_-,R_-$ onto
their positive counterparts.
It therefore preserves the seven
specified character values, their
product, both scalar fields and
all maps defining the summed restriction.
The proved negative image equality
and kernel statements give the
positive statements through these
explicit mutually inverse maps.
This proves the claim for both signs.
\end{proof}

\begin{claim}{F.32}{Proved}{Finite affine marginal image equality when the first character is one and the boundary character is nontrivial}
Use the original bases $Y_i^-=X_i$ and
$Y_i^+=\mathcal B_+(X_i)$ of F.23.13.
For every finite field $k$ and continuous
$\rho:V_\pm\to\operatorname{Aff}_1(k)$ whose
scalar character belongs to the exact locus
\[
 \chi(Y_1^\pm)=1,\qquad
                 h=\prod_{i=1}^7\chi(Y_i^\pm)\ne1,
                                                               \tag{F.32.1}
\]
the actual marginal images satisfy
\[
                         \rho(\mathcal K_\pm)=\rho(R_\pm).
                                                               \tag{F.32.2}
\]
All finite characteristics, scalar-field
extension degrees and larger original
coefficient fields are included.

Put $C=\chi(R_\pm)$,
$k_1=\mathbb F_p(C)$ and
$k_0=\mathbb F_p(\chi(V_\pm))$.
For $C\ne\{1\}$, let $q_1=|k_1|$,
$m=[k_0:k_1]$ and
$\chi_j=\chi^{q_1^j}$, $0\leq j<m$.
The summed continuous $H^1$ restriction
for these characters is injective.
The corresponding cocycle restriction
has exactly the principal tuples with
parameter sum zero as its kernel,
and has rank $6m+1$ on $7m$ coordinates.
\end{claim}
\passage{Part F root-adjoining intersection question; Section30, folios194-196, scans383-387; Section31, folios198-199 verso, scans391-394; original F7/F19 subgroup-membership residual}


\begin{proof}
\emph{The actual scalar fields and coefficient maps.}
Work first with the original negative
marginal $K=\mathcal K_-\leq R=R_-\leq V=V_-$.
Write its character in the unique form
\[
 \begin{gathered}
 \chi=(1,s,t,u,v,w,z),\qquad h=stuvwz\ne1,\\
 C=\chi(K)=\chi(R)=\langle u,v,w,z\rangle,\qquad
 k_1=\mathbb F_p(u,v,w,z),\\
 k_0=k_1(s,t)=k_1(h,t),\qquad
                       s=\frac{h}{tuv wz}.
 \end{gathered}                                                \tag{F.32.3}
\]
Every displayed scalar is nonzero.
The equality of scalar images is the
image of the actual normal-generation
equality
$\overline{\langle\!\langle K\rangle\!\rangle}_R=R$
from F.21 and the marginal map of F.22
under the abelian scalar homomorphism.
The deleted generators $X_4,\ldots,X_7$
give exactly the four generators of $C$.
The identities for $h$ and $s$ prove
both inclusions in the displayed
equality of fields.

For $C=\{1\}$, the image of $R$ is
an abelian translation group, so this
same normal-generation map proves
(F.32.2). For $C\ne\{1\}$ and $m=1$,
F.25 gives (F.32.2), and F.24 gives
the asserted individual restriction
injection and cocycle kernel.
We henceforth take $C\ne\{1\}$ and $m>1$.

The original Laurent coefficient map is
\[
 \begin{aligned}
 \mathbb Z[t_1^{\pm1},\ldots,t_7^{\pm1}]
   &\longrightarrow
       A=\mathbb Z[h^{\pm1},t^{\pm1},u^{\pm1},v^{\pm1},w^{\pm1},z^{\pm1}],\\
 (t_1,t_2,t_3,t_4,t_5,t_6,t_7)
   &\longmapsto (1,h/(tuv wz),t,u,v,w,z).
 \end{aligned}                                                \tag{F.32.4}
\]
After the first character is set to one,
the change from the independent variables
$(s,t,u,v,w,z)$ to $(h,t,u,v,w,z)$ is
the exact Laurent-ring isomorphism
\[
 \begin{gathered}
 s^at^bu^cv^dw^ez^f
    \longmapsto h^at^{b-a}u^{c-a}v^{d-a}w^{e-a}z^{f-a},\\
 h^at^bu^cv^dw^ez^f
    \longmapsto s^at^{b+a}u^{c+a}v^{d+a}w^{e+a}z^{f+a}.
 \end{gathered}                                                \tag{F.32.5}
\]
Both composites fix every monomial.
Adding exponent vectors proves that
both maps preserve multiplication;
extension by integer linearity proves
the asserted inverse ring maps.
In $A$ the letter $s$ will always denote
the specific Laurent monomial in (F.32.3).

\emph{Boundary coordinates for the actual principal quotient.}
The continuous cocycle space for $\chi$
is $k_0^7$, in original generator order.
Indeed every such tuple defines a
homomorphism of the discrete free group
to the finite affine group, and that
homomorphism has a unique continuous
extension to $V=\widehat F_7$.
Its principal line is $k_0\delta$, where
\[
 \begin{gathered}
 \delta=(0,s-1,t-1,u-1,v-1,w-1,z-1)^{\mathsf T},\\
 D=(1,1,s,st,stu,stuv,stuvw),\qquad
                            D\delta=h-1,\\
 N=I_7-\frac{\delta D}{h-1},\qquad H=\ker D.
 \end{gathered}                                                \tag{F.32.6}
\]
The row $D$ evaluates a cocycle on
the original word $P_0=X_1\cdots X_7$.
The nonzero scalar $h-1$ makes
$\delta$ nonzero and defines $N$.
Direct multiplication gives
$N\delta=0$, $DN=0$ and $N^2=N$.
For every $a\in k_0^7$ one has
\[
       a=Na+\delta\frac{Da}{h-1},\qquad
       \ker N=k_0\delta,\qquad \operatorname{im}N=H.
                                                               \tag{F.32.7}
\]
Consequently $[a]\mapsto Na$ and
$b\mapsto[b]$, $b\in H$, are mutually
inverse maps from the actual
$Q_\chi=H^1_{\rm cont}(V,k_{0,\chi})$
to $H$. This proves the precise
principal-quotient identification.

Let $E$ take the last six coordinates,
and let $J:k_0^6\to k_0^7$ have columns
\[
 \begin{gathered}
 J=(e_2-e_1,e_3-se_1,e_4-ste_1,
                         e_5-stue_1,e_6-stuve_1,e_7-stuvwe_1),\\
 DJ=0,\qquad EJ=I_6,\qquad JE=I_7-e_1D,\\
 ENJ=I_6,\qquad JEN=N.
 \end{gathered}                                                \tag{F.32.8}
\]
Thus $J$ and $E|_H$ are explicit
inverse coordinate maps on $H$.
The actual principal quotient in six
coordinates is $EN$, with section $J$.
These coordinates have the original
order $X_2,\ldots,X_7$.

Retain each of the eleven actual
F.23 source automorphisms
$B_l=\operatorname{Ad}(u_l)\alpha_l$.
Here and below $0\leq l\leq10$ is
the original zero-based source index:
$0$ is \texttt{old0}, $1,\ldots,5$
are \texttt{old\_1\_0},\ldots,\texttt{old\_1\_4},
and $6,\ldots,10$ are
\texttt{old\_-1\_0},\ldots,\texttt{old\_-1\_4}.
Its position in a one-based list is
always $l+1$.
Both $B_l$ and its inverse preserve
the actual subgroup $K$.
The pure Artin part preserves every
abelian character and fixes $P_0$.

Let $A_l$ be the seven-coordinate
cocycle action of $B_l$, $M_l$ that
of $\alpha_l$, $a_l$ the row evaluating
$d(u_l)$, and $\lambda_l=\chi(u_l)$.
The actual cocycle formulas give
\[
 \begin{gathered}
 A_l=\lambda_lM_l-\delta a_l,\qquad
 M_l\delta=\delta,\qquad DM_l=D,\\
 a_l\delta=\lambda_l-1,\qquad A_l\delta=\delta,\\
 U_l=A_l+\delta a_l=\lambda_lM_l,
             \qquad NA_l=U_lN,\\
 T_l=EU_lJ,\qquad U_lJ=JT_l.
 \end{gathered}                                                \tag{F.32.9}
\]
For the middle intertwiner, $N\delta=0$
cancels the displayed decoration term,
and $M_l\delta=\delta$, $DM_l=D$
give $NM_l=M_lN$.
Thus $T_l$ is the exact source action
on the actual quotient (F.32.8).
Its entries belong to $A$; the only
division in identifying that quotient
was the already nonzero $h-1$.

For an entirely integral check put
$B=(h-1)I_7-\delta D$. Then
\[
 \begin{gathered}
 B\delta=0,\quad DB=0,\quad B^2=(h-1)B,\quad
 BA_l=U_lB,\\
 DU_l=\lambda_lD,\qquad U_l\delta=\lambda_l\delta,\\
 (\lambda_0,\ldots,\lambda_{10})=
 (1,1,t,1,1,1,u^{-2},(ht)^{-1},v^{-2},w^{-2},z^{-2}).
 \end{gathered}                                                \tag{F.32.10}
\]
Every scalar in the last line is
evaluated from the decoration of the
same labelled original word.

\emph{The invariant five-dimensional kernel and its quotient.}
Because $\chi(X_1)=1$, evaluation
$d\mapsto d(X_1)$ vanishes on the
principal line. On the coordinates
(F.32.8) its negative is the row
$b=(1,s,st,stu,stuv,stuvw)$.
Set $W_\chi=\ker b$ and use the exact
maps and basis
\[
 \begin{gathered}
 \iota(x_1,\ldots,x_5)=
 (-sx_1-stx_2-stux_3-stuvx_4-stuvwx_5,x_1,\ldots,x_5),\\
 \pi(\zeta)=b\zeta=-d(X_1),\qquad \sigma(a)=ae_1,\qquad
 \pi\sigma=I_1,\\
 \zeta=\iota(\zeta_2,\ldots,\zeta_6)+\sigma\pi(\zeta),\\
 F=(e_2-se_1,e_3-ste_1,e_4-stue_1,
                         e_5-stuve_1,e_6-stuvwe_1,e_1),\\
 F^{-1}(\zeta)=(\zeta_2,\zeta_3,\zeta_4,\zeta_5,\zeta_6,b\zeta).
 \end{gathered}                                                \tag{F.32.11}
\]
Both composites of the displayed basis
maps are identity. In particular no
additional scalar is inverted here.
Direct evaluation of the original words
gives $d(B_l(X_1))=\nu_l d(X_1)$ and
\[
 \begin{gathered}
 S_l=F^{-1}T_lF=
       \begin{pmatrix}G_l&*\\0&\nu_l\end{pmatrix},\qquad
                              G_0=h^{-1}I_5,\\
 (\nu_0,\ldots,\nu_{10})=
 (1,u^{-2},t^{-1},v^{-2},w^{-2},z^{-2},1,t/h,1,1,1).
 \end{gathered}                                                \tag{F.32.12}
\]
This proves invariance of $W_\chi$
for every actual source, and identifies
its one-dimensional quotient by the
actual map $\pi$.
The integral certificate checks the
full original first row of $A_l$,
the full block identity, and the cleared
seven-coordinate identity
$BA_lJF=(h-1)JFS_l$.

\emph{Eight moving-root identities and explicit evaluation rows.}
For each row of the following table
let $Q_l(X)$ be the product of $X-a$
over all listed base roots $a$, with
their displayed multiplicities.
Let $b_l$ be its moving root.
\[
\begin{array}{c|c|c}
 l&\text{base roots of }Q_l&b_l\\\hline
 1&1,(vw)^{-1}&(tz)^{-1}\\
 3&1,(uz)^{-1}&(tw)^{-1}\\
 4&1,(vz)^{-1}&(tu)^{-1}\\
 5&1,(uw)^{-1}&(tv)^{-1}\\
 6&1,(vw)^{-1},(h/u)^2&(tz)^{-1}\\
 8&1,(uz)^{-1},(h/v)^2&(tw)^{-1}\\
 9&1,(vz)^{-1},(h/w)^2&(tu)^{-1}\\
 10&1,(uw)^{-1},(h/z)^2&(tv)^{-1}
\end{array}.
                                                               \tag{F.32.13}
\]
The following identities hold over $A$,
before any specialization or cancellation:
\[
 \begin{gathered}
 Q_l(G_l)(G_l-b_lI_5)=0,\qquad
 q_lG_l=b_lq_l,\\
 q_lQ_l(G_l)=Q_l(b_l)q_l,\qquad
 q_l=\widehat q_l\Delta_t,\qquad
               \Delta_t=\operatorname{diag}(1,t-1,t-1,t-1,t-1).
 \end{gathered}                                                \tag{F.32.14}
\]
The diagonal identity defines $q_l$
without dividing by $t-1$.
Here are all eight rows explicitly.
The first four are
\[
 \begin{gathered}
 \widehat q_1^{\mathsf T}=\frac1{uv wz}
 \begin{pmatrix}
 -h(z-1)\\-h(z-1)\\-hu(z-1)\\-huv(z-1)\\-uv w(hz-h-z)
 \end{pmatrix},\qquad
 \widehat q_8^{\mathsf T}=\frac1{uv w}
 \begin{pmatrix}
 -h(w-1)\\-h(w-1)\\-hu(w-1)\\-uv(hw-h-w)\\-uv w(h-1)(w-1)
 \end{pmatrix},\\
 \widehat q_9^{\mathsf T}=\frac1{uv wz}
 \begin{pmatrix}
 h(u-1)\\h(u-1)-uv wz\\hu(u-1)\\huv(u-1)\\huv w(u-1)
 \end{pmatrix},\qquad
 \widehat q_{10}^{\mathsf T}=
                 \begin{pmatrix}1-v\\1-v\\u\\0\\0\end{pmatrix}.
 \end{gathered}                                                \tag{F.32.15}
\]
The other four use the explicit
polynomials
\[
 \begin{aligned}
 P_4={}&h^2\{u(vw+wz-w-z)-vw\}
                  +hu(-vw-2wz+w+z)+uwz,\\
 P_6={}&(z-1)\{h^2[vw(u-1)+uw(z-1)-uz]
                         +huz(1-2w)\}+uwz^2.
 \end{aligned}                                                \tag{F.32.16}
\]
Their rows are
\[
 \begin{gathered}
 \widehat q_3^{\mathsf T}=\frac1{uwz}
 \begin{pmatrix}
 -h(w-1)\\-(h+uwz)(w-1)\\0\\hu(uw-w+1)\\huw(u-1)(w-1)
 \end{pmatrix},\\
 \widehat q_4^{\mathsf T}=\frac1{u^2wz}
 \begin{pmatrix}
 h^2(u-1)\\
 h^2(u-1)+huv wz(u-1)-u^2v wz\\
 -hu(u-1)[h(u-1)-u]\\
 -hu(u-1)\{h[(u-1)v+u(z-1)]-u(v+z-1)\}\\
 -u(u-1)P_4
 \end{pmatrix},\\
 \widehat q_5^{\mathsf T}=\frac1{uwz}
 \begin{pmatrix}
 h(v-1)\\h(v-1)\\-hu[hz+(v-1)(z-1)]\\
 -u(v-1)[h^2z+hv(z-1)+h(w-1)z-wz]\\
 -uw(v-1)[h^2z+hv(z-1)+h(w-2)z-(v+w-1)z]
 \end{pmatrix},\\
 \widehat q_6^{\mathsf T}=\frac1{u^2wz^2}
 \begin{pmatrix}
 h^2(z-1)\\h(h+uv wz)(z-1)\\-h^2u(u-1)(z-1)\\
 -hu(z-1)\{h[u(v+z-1)-v]-uz\}\\-uP_6
 \end{pmatrix}.
 \end{gathered}                                                \tag{F.32.17}
\]
Every denominator in these displays
is a Laurent monomial. The direct
left-eigenrow identities in (F.32.14)
are checked from the actual matrices;
they do not require a rank assertion
to persist under specialization.
The last polynomial evaluation in
(F.32.14) also follows by multiplying
the direct left-eigenrow identity
through each factor of $Q_l$.

For $a=3,4,5,6$, let $\widehat R_a$
be the square matrix with rows
$\widehat q_1,\widehat q_8,\widehat q_9,
\widehat q_{10},\widehat q_a$, in that order.
Let $R_a$ use the corresponding actual
rows $q_l$. The exact determinants are
\[
 \begin{gathered}
 R_a=\widehat R_a\Delta_t,\qquad
                 \det R_a=(t-1)^4\det\widehat R_a,\\
 \begin{array}{c|c}
 a&\det\widehat R_a\\\hline
 3&h u(h-1)(w-1)/(wz)\\
 4&-h(h-1)(u-1)(hw-h-w)/w^2\\
 5&-h(h-1)(v-1)\\
 6&-h(h-1)(z-1)(hw-h-w)/(w^2z)
 \end{array}.
 \end{gathered}                                                \tag{F.32.18}
\]
This is a column-scaling identity
over $A$, including $t=1$.
Its use as an invertibility statement
will occur only after the relevant
field argument proves $t\ne1$.

\emph{Comparison of core subspaces in distinct Frobenius factors.}
Let $\sigma:k_0\to k_0$ be the actual
$k_1$-Frobenius $a\mapsto a^{q_1}$.
Its order is $m$. Put
$h_j=\sigma^j(h)$ and $t_j=\sigma^j(t)$;
the four deleted scalars are fixed.
The coordinates and matrices above
specialize to $Q_j=Q_{\chi_j}$ and
$W_j\subset Q_j$.
For each $j$ the two-way field equality
in (F.32.3), after this automorphism,
gives
\[
                       k_0=k_1(h_j,t_j).
                                                               \tag{F.32.19}
\]

Suppose that nonzero invariant subspaces
$S_j\subset W_j$ and $S_{j'}\subset W_{j'}$,
with $j\ne j'$, have a linear isomorphism
intertwining all eleven sources.
Source zero is $h_j^{-1}I_5$ and
$h_{j'}^{-1}I_5$, respectively, so
$h_j=h_{j'}$. Use the actual automorphism
$\tau=\sigma^{j'-j}$, with negative powers
interpreted by the inverse automorphism,
and put
\[
 L=\{a\in k_0:\tau(a)=a\},\qquad
       k_1\subset L,\quad h_j\in L,\quad
                         t_j\notin L.
                                                               \tag{F.32.20}
\]
Here $L$ is a field because $\tau$
preserves the field operations and
inverses. It is proper, since
$0<|j'-j|<m$ and $\sigma$ has order $m$.
If $t_j$ belonged to $L$, (F.32.19)
would put all of $k_0$ in $L$, a
contradiction. The same reasoning shows
$t_{j'}\ne t_j$: equality would make
$\tau$ fix both generators in (F.32.19).
Also $t_{j'}=\tau(t_j)$ cannot belong
to $L$, since applying $\tau^{-1}$ to
an element of $L$ fixes that element.

For every source in (F.32.13), all base
roots lie in $L$ and agree for the
two factors. Its moving roots are
$\eta/t_j$ and $\eta/t_{j'}$ for one
nonzero $\eta\in k_1$. They are distinct,
and neither belongs to $L$; otherwise
their inverses multiplied by $\eta$
would put $t_j$ or $t_{j'}$ in $L$.
Consequently each factor of $Q_l(b_{l,j})$
is nonzero, including any repeated
base-root factor.

Transport the source action on $S_{j'}$
back to $S_j$ by the given isomorphism.
For a fixed source it is the same
operator on $S_j$, and both annihilators
in (F.32.14) apply to it.
Subtracting these two polynomial
identities gives
$(b_{l,j'}-b_{l,j})Q_l(G_{l,j})=0$
on $S_j$, so $Q_l(G_{l,j})$ vanishes
there. The direct row evaluation in
(F.32.14) and the proved nonzero
$Q_l(b_{l,j})$ now give
$q_{l,j}|_{S_j}=0$ for all eight rows.
This deduction uses only exact
intertwiners and polynomial identities.

At this comparison $t_j\notin L$,
so $t_j-1\ne0$; also $h_j-1\ne0$
by (F.32.1). Every displayed monomial
factor is nonzero. The determinant
cover (F.32.18) is exhaustive as follows.
For $w\ne1$ use $a=3$.
For $w=1$ and $v\ne1$ use $a=5$.
For $w=v=1$ and $u\ne1$ use $a=4$:
the remaining factor $h_jw-h_j-w$
is exactly $-1$.
For $w=v=u=1$, nontriviality of $C$
forces $z\ne1$, and $a=6$ has the
same factor $-1$.
Each case provides a nonzero actual
determinant $\det R_a$.
Thus the eight rows cannot all vanish
on a nonzero subspace, a contradiction.
There are no isomorphic nonzero
invariant core subspaces across
distinct Frobenius factors.
Every argument and every nonzero
factor holds also in characteristic two.

\emph{All other simple-subspace types.}
A simple invariant subspace of $Q_j$
either lies in $W_j$, or intersects
$W_j$ trivially. In the latter case
the actual quotient map $\pi$ restricts
to an injection into its one-dimensional
quotient. Its image is the entire
one-dimensional space, so it has
exactly the character $\nu_j$ of
(F.32.12), with the inverse of this
restricted quotient as the inverse map.
This proves the classification of
possible simple-subspace types used
here without a semisimplicity assumption.

A core subspace from any factor cannot
be isomorphic to such a quotient
character from any factor: source zero
is respectively $h_j^{-1}$ and $1$,
which are unequal because $h_j\ne1$.
For two quotient characters, equality
at source $2$ gives $t_j=t_l$, and
equality at source $7$ then gives
$h_j=h_l$. Equation (F.32.19) makes
$\sigma^{l-j}$ the identity on $k_0$,
which forces $j=l$ in the indicated
index range. Inverse intertwiners cover
both orientations of each comparison.
Together with the core comparison,
these are all three unordered type
pairs: core/core, core/quotient and
quotient/quotient.

\emph{Summed restriction and the exact affine consequence.}
Every $\chi_j$ restricts to the same
character on $K$, because its image
there is $C\subset k_1^\times$.
There is consequently the actual map
\[
 \Phi:\bigoplus_{j=0}^{m-1}Q_j
       \longrightarrow H^1_{\rm cont}(K,k_{0,\chi|K}),\qquad
                   ([d_j])_j\longmapsto\sum_j[d_j|K].
                                                               \tag{F.32.21}
\]
The same eleven actual source
automorphisms act on both sides, and
restriction commutes with them because
they preserve $K$ and the characters.
Each individual restriction is injective
by F.24 and $C\ne\{1\}$.

A nonzero kernel of $\Phi$ would
contain a nonzero invariant subspace
$T$ of least dimension. Each nonzero
coordinate projection $p_j|_T$ is
injective: its kernel is invariant
and is a proper subspace of $T$.
Its image is simple, because the
inverse projection pulls any proper
nonzero invariant subspace back to one
in $T$. Two nonzero projections give
the explicit inverse intertwiners
\[
 \begin{gathered}
 p_l\circ(p_j|_T)^{-1}:p_j(T)\longrightarrow p_l(T),\\
 p_j\circ(p_l|_T)^{-1}:p_l(T)\longrightarrow p_j(T).
 \end{gathered}                                                \tag{F.32.22}
\]
The exhaustive three comparisons exclude
these for $j\ne l$. Exactly one
projection is therefore nonzero, and
$T$ lies in a single factor. There
$\Phi$ is the injective individual
restriction, giving the final
contradiction. This proves the summed
$H^1$ injection.

Zero summed cocycle restriction now
makes every cocycle principal.
Their common restricted character has
nontrivial image $C$, so the exact
kernel on cocycles is
\[
 \left\{\bigl(b_j(\chi_j-1)\bigr)_{j=0}^{m-1}:
                              \sum_{j=0}^{m-1}b_j=0\right\}.
                                                               \tag{F.32.23}
\]
Conversely each tuple in this display
restricts with sum zero. Each principal
line has a nonzero generating vector;
the displayed kernel has dimension
$m-1$. The cocycle domain has its
original $7m$ coordinates, so its
restriction rank is $6m+1$.

Apply this statement to the actual
Frobenius coefficient tuple
$(d^{q_1^j})_{j=0}^{m-1}$. It gives
the exact trace-kernel identity
\[
 \operatorname{Tr}_{k_0/k_1}(d|K)=0
 \quad\Longleftrightarrow\quad
 d=b(\chi-1),\qquad
                      \operatorname{Tr}_{k_0/k_1}(b)=0.
                                                               \tag{F.32.24}
\]
For the forward direction, summed
restriction makes $d$ principal. Its
Frobenius powers have the unique
parameters $b^{q_1^j}$, whose sum is
the trace, and (F.32.23) gives the
stated vanishing. For the reverse
direction use $\chi(K)\subset k_1$
and the trace's $k_1$-linearity.

The explicit affine conjugation,
translation projection and trace-pairing
coefficient maps of F.26 take every
proper marginal affine image to a
nonprincipal element of this trace
kernel over the actual scalar field
$k_0$. Those maps preserve the scalar
character, both scalar fields, its
first coordinate and its boundary
product. They therefore retain the
exact locus (F.32.1).
Equation (F.32.24) excludes the
resulting nonprincipal cocycle, and
proves (F.32.2) for every original
finite coefficient field, including
fields larger than $k_0$.

\emph{Literal-word certificate, continuity and the positive marginal.}
The independent verifier imports none
of the four new producers and no
symbolic algebra package. It reconstructs
all eleven literal F.23 decorated
Artin words directly on the coefficient
chart (F.32.4), using the previously
reviewed integer Laurent word arithmetic.
At a positive crossing the meridian
and decoration pair is transformed by
$(x_j,x_{j+1};a_j,a_{j+1})\mapsto
(x_jx_{j+1}x_j^{-1},x_j;x_ja_{j+1},a_j)$;
at a negative crossing it is transformed
by $(x_{j+1},x_{j+1}^{-1}x_jx_{j+1};
a_{j+1},x_{j+1}^{-1}a_j)$.
The eighth meridian is evaluated at
the identity, and the final decoration
gives $u_l=a_8^{-1}$.
The verifier checks the label permutation,
the erased meridian, every preserved
character and the full boundary product.

Its $1088$ exact checks include the
full seven-coordinate actions and
decoration derivatives, the inverse
coefficient maps (F.32.5), all basis
inverses, the cleared intertwiners
(F.32.10)--(F.32.12), the eight direct
left-eigenrow identities and annihilators,
and all four actual determinant and
column-scaling identities.
They are integer Laurent identities,
so their stated field specializations
apply in every characteristic.
The nonmonomial division $h-1$ in
the principal quotient and all nonzero
factors used in the comparison have
been proved valid in the exact fields
before use. The evidence manifest
binds the original words, all coefficient
data, the programs and outputs, and
the completed dependencies.

All finite coefficient maps are
continuous. The original source
automorphisms are continuous by F.23;
the cocycles have the unique continuous
finite-affine extensions already described.
Their restrictions to the actual
closed subgroup $K$ and all displayed
coefficient compositions are therefore
continuous throughout the proof.
For the positive marginal use the
actual based isomorphism
$\mathcal B_+:V_-\to V_+$ of F.23.13
and its given inverse. It sends $X_i$
exactly to $Y_i^+$ and carries the
actual subgroups $\mathcal K_-,R_-$
onto their positive counterparts.
Precomposition by these inverse maps
preserves all seven character values,
the boundary product, both scalar
fields, and the maps in (F.32.21).
It therefore transports the proved
negative image equality and exact
kernel assertions to the positive
marginal. This proves the claim.
\end{proof}

\begin{claim}{F.33}{Proved}{Exact paired restriction kernel for opposite first characters}
Use the actual bases $Y_i^-=X_i$ and
$Y_i^+=\mathcal B_+(X_i)$ of F.23.13.
Let $k$ be a finite field of odd characteristic,
and let $\chi,\chi':V_\pm\to k^\times$ be the
continuous characters with values
\[
 \begin{aligned}
 (\chi(Y_i^\pm))_{i=1}^7&=(r,s,t,u,v,w,z),\\
 (\chi'(Y_i^\pm))_{i=1}^7&=(-r,s,t,u,v,w,z),\qquad
             \langle u,v,w,z\rangle\ne\{1\}.
 \end{aligned}                                                \tag{F.33.1}
\]
All seven displayed values are nonzero; $r^2=1$
is included. The two characters have the same
restriction $\eta$ to the actual subgroup
$\mathcal K_\pm$. The summed continuous cocycle
restriction has the exact kernel
\[
 \begin{split}
 \ker\bigl[Z^1_\chi(V_\pm,k)\oplus
 Z^1_{\chi'}(V_\pm,k)
   &\longrightarrow Z^1_\eta(\mathcal K_\pm,k)\bigr]\\
   &=\{(b(\chi-1),-b(\chi'-1)):b\in k\}.
 \end{split}                                                  \tag{F.33.2}
\]
It has dimension one, and the restriction
map has rank $13$ on the original $14$
generator coordinates. The induced summed
$H^1$ restriction is injective.

More generally, the explicit $64$ by $13$
Laurent matrix used below has rank $13$
over every field, in every characteristic,
for every nonzero tuple $(R,s,t,u,v,w,z)$
with $(u,v,w,z)\ne(1,1,1,1)$.
\end{claim}
\passage{Part F root-adjoining intersection question; Section30, folios194-196, scans383-387; Section31, folios198-199 verso, scans391-394; actual F7/F19 subgroup-membership residual}


\begin{proof}
\emph{The original subgroup and source words.}
First work with
$K=\mathcal K_-\le R_-\le V_-=V$.
The original retained quotient
$\pi:V\to\widehat F(X_1,X_2,X_3)$ kills
$X_4,\ldots,X_7$, and $R_-=\ker\pi$.
The character $\chi'/\chi$ is exactly
$\epsilon\pi$, where
$\epsilon(X_1)=-1$ and
$\epsilon(X_2)=\epsilon(X_3)=1$.
Thus $\chi|R_-=\chi'|R_-$, proving their
agreement on $K$ by this actual quotient map.
The normal-generation equality of F.21,
transported through the original marginal
map of F.22 and then through the abelian
scalar map, gives
\[
       \chi(K)=\chi(R_-)=\langle u,v,w,z\rangle.
                                                               \tag{F.33.3}
\]
Indeed conjugation does not change an
abelian scalar image, and the finite image
of $K$ is already closed. In particular
$\eta\ne1$.

Retain the eight actual negative collision
seeds in the following order:
\[
 \begin{gathered}
 \omega_0,\ldots,\omega_7:\quad
 (0,-1,1),(0,-1,-1),(2,-1,1),(2,-1,-1),\\
 (3,-1,1),(3,-1,-1),(4,-1,1),(4,-1,-1).
 \end{gathered}                                                \tag{F.33.4}
\]
The triples are the original F.23 labels
$\mathrm{cross}_{e,-1,\varepsilon}$.
Their scalar characters are respectively
$u,u,v,v,w,w,z,z$.
Write $\beta_j$ for the actual negative
old-source automorphism with zero-based
index $j$, including its original
decoration and orientation. The $64$ words
used in the proof are, in order,
\[
 \begin{gathered}
 \omega_0,\ldots,\omega_7,\qquad
 (\beta_j(\omega_0),\ldots,\beta_j(\omega_7))
       \quad\text{successively for}\\
                         j=2,9,8,10,1,0,7.
 \end{gathered}                                                \tag{F.33.5}
\]
These are the original F.24 source rows.
Every displayed word belongs to $K$:
the seeds lie in $K$, and each original
old-source conjugation preserves $K$.
No larger braid group or replacement
source map is used. Since the old sources
are pure, their action preserves every
abelian character, so every word in
(F.33.5) retains its seed character.

\emph{The coefficient ring and the principal-line quotient.}
Set
\[
 \begin{gathered}
 A=\mathbb Z[R^{\pm1},s^{\pm1},t^{\pm1},u^{\pm1},
                         v^{\pm1},w^{\pm1},z^{\pm1}],\\
                         B=A[\xi]/(\xi^2-R).
 \end{gathered}
                                                               \tag{F.33.6}
\]
The element $\xi$ is a unit, with
$\xi^{-1}=\xi/R$. The map from the
original seven-variable Laurent ring
to $B$ sends its variables to
$(\xi,s,t,u,v,w,z)$. Its inverse sends
$\xi$ to the original first variable
and $R$ to its square, fixing the other
six variables. Both composites fix all
generators. The monic quadratic relation
gives the free $A$-basis $1,\xi$;
explicitly, for every integer $e$,
\[
       \xi^e=R^{(e-\bar e)/2}\xi^{\bar e},
       \qquad \bar e\in\{0,1\},\quad
                         \bar e\equiv e\pmod2.
                                                               \tag{F.33.7}
\]
This includes negative exponents and
defines the exact coefficient map.

For each word in (F.33.5), write its
original universal derivation row
uniquely as $a+\xi b$, with $a,b\in A^7$.
Let $E$ be the $64$ by $14$ matrix whose
row is $(b,a)$, in that order.
The exact Fox fundamental identity for
the word, whose character has no $\xi$
component, gives
\[
 -b_1+a_1+\sum_{j=2}^7b_j(\theta_j-1)=0,
       \qquad(\theta_2,\ldots,\theta_7)=(s,t,u,v,w,z).
                                                               \tag{F.33.8}
\]
Consequently $Ep=0$ for the actual vector
\[
 p=(-1,s-1,t-1,u-1,v-1,w-1,z-1;\,1,0,0,0,0,0,0)^T.
                                                               \tag{F.33.9}
\]

Here is an explicit quotient with its
inverse, over $A$ itself. For
$(x_1,\ldots,x_7;y_1,\ldots,y_7)$ put
\[
 \begin{aligned}
 \lambda&=y_1,&q_1&=x_1+\lambda,\\
 q_j&=x_j-(\theta_j-1)\lambda\quad(2\le j\le7),&
 q_{6+j}&=y_j\quad(2\le j\le7).
 \end{aligned}                                                \tag{F.33.10}
\]
The target coordinates are $(q_1,\ldots,q_{13};\lambda)$.
The inverse is
\[
 \begin{aligned}
 x_1&=q_1-\lambda,&y_1&=\lambda,\\
 x_j&=q_j+(\theta_j-1)\lambda\quad(2\le j\le7),&
 y_j&=q_{6+j}\quad(2\le j\le7).
 \end{aligned}                                                \tag{F.33.11}
\]
Substitution proves both inverse
identities coordinate by coordinate.
This subtracts $\lambda p$ and records
the unique removed parameter. By
(F.33.8), $E(x,y)=Nq$, where the row of
the $64$ by $13$ matrix $N$ is
\[
              (b_1,\ldots,b_7,a_2,\ldots,a_7).
                                                               \tag{F.33.12}
\]
There is no division in this quotient.

\emph{Four exhaustive exact row certificates.}
The field specializations are covered by
the four disjoint charts
\[
 \begin{array}{c|l|c}
 \mathsf U&u\ne1&\delta=u-1\\
 \mathsf V&u=1,\ v\ne1&\delta=v-1\\
 \mathsf W&u=v=1,\ w\ne1&\delta=w-1\\
 \mathsf Z&u=v=w=1,\ z\ne1&\delta=z-1.
 \end{array}                                                  \tag{F.33.13}
\]
Setting the preceding variables to one
is the displayed Laurent ring
homomorphism for that chart.
Every row operation is either multiplication
by a Laurent monomial, division by the
active $\delta$ after exact divisibility
has been checked in every entry, or
the cleared elimination
$L\mapsto pL-cP$, where $p$ is the
pivot entry of $P$ and $c$ is the
corresponding entry of $L$.
The inverse to each divided normalization
is multiplication by its recorded
monomial inverse and power of $\delta$.
Thus all resulting rows lie in the
specialized original row space.

The producer permits tests for factors
$R-1$ and $2$ during its search, but
every recorded division count for those
factors is zero, in all operations and
all unit witnesses. They are not inverted
in this proof. The only actual divisions
are by character monomials and the
nonzero active difference in (F.33.13).
The respective operation counts are
$674,483,374,232$; the respective
active-difference division counts are
$17,15,15,9$.

The following table prints every single
pivot as $(\text{row},\text{column};\text{coefficient})$.
Both row and column indices are zero-based,
and the column order is exactly (F.33.12).
\[
\begin{array}{c|c|c|c|c}
 &\mathsf U&\mathsf V&\mathsf W&\mathsf Z\\\hline
1&(11,4;stu)&(0,3;1)&(0,3;1)&(0,3;1)\\
2&(8,9;1)&(1,9;1)&(1,9;1)&(1,9;1)\\
3&(13,5;stu)&(3,10;s)&(2,4;1)&(2,4;1)\\
4&(14,12;1)&(11,4;st)&(3,10;1)&(3,10;1)\\
5&(0,3;1)&(13,5;st)&(13,5;st)&(4,5;1)\\
6&(3,10;su)&(14,12;1)&(14,12;1)&(5,11;1)\\
7&(16,8;s)&(18,7;-1)&(20,7;-1)&(14,12;1)\\
8&(6,6;1)&(6,6;-1)&(4,0;-stz)&(39,7;-1)\\
9&(40,0;1)&(5,11;1)&(28,8;1)&(62,0;st)\\
10&(1,2;s)&(10,0;-st)&(6,6;st)&(7,8;-1)\\
11&(48,7;1)&(34,8;-1)&(5,11;1)&(6,6;st)\\
12&(5,11;1)&\text{block below}&(29,2;s)&(47,2;s)\\
13&(49,1;1)&\text{block below}&(52,1;-1)&(55,1;1).
\end{array}                                                   \tag{F.33.14}
\]
Every entry is a signed Laurent monomial,
and every pivot row is zero in its
preceding pivot columns. The remaining
two columns of chart $\mathsf V$ are
$1,2$. In the actual recorded order
of rows $35,51$, their exact block is
\[
 M=\begin{pmatrix}
 -(t-1)(st+1)&s(st-t+1)\\
 st-s+1&-s(s-1)
 \end{pmatrix},\qquad \det M=-s^2t.
                                                               \tag{F.33.15}
\]
These two full rows vanish in all eleven
previous pivot columns. Direct multiplication
gives $M\operatorname{adj}(M)=
\operatorname{adj}(M)M=-s^2t I_2$.
Thus the adjugate gives the two remaining
coordinate rows with unit coefficient
$-s^2t$. In the reversed row order the
block is $SM$, where
$S=\left(\begin{smallmatrix}0&1\\1&0\end{smallmatrix}\right)$;
$S^2=I_2$ and $\det(SM)=s^2t$.
This is the exact comparison between
the two row-order descriptions.

The fifty single pivots and this one
two-column block prove rank $13$ in
all four charts. The proof is over
the integers, with only the stated
nonzero field divisions. It therefore
establishes the matrix assertion in
every characteristic, including $2$,
and for every $R\ne0$, including $R=1$.

\emph{The two characters and their exact kernel.}
Now specialize $R=r^2$ in the given odd
characteristic field. A derivation row
$a+\xi b$ becomes $a+rb$ for $\chi$
and $a-rb$ for $\chi'$.
For their generator-value columns $c,c'$,
the actual summed evaluation is
\[
 \begin{gathered}
 (a+rb)c+(a-rb)c'=bx+ay,\\
 \begin{aligned}
 x&=r(c-c'),&y&=c+c',\\
 c&=(y+x/r)/2,&c'&=(y-x/r)/2.
 \end{aligned}
 \end{gathered}                                                \tag{F.33.16}
\]
The displayed two maps are inverse;
their only divisors are the nonzero
elements $2,r$. Together with
(F.33.10)--(F.33.12), rank $13$ gives
$\ker E=kp$. The inverse map in
(F.33.16) sends $\lambda p$ to
\[
 \left(\frac{\lambda}{2r}(\chi(X_i)-1)_{i=1}^7,
       -\frac{\lambda}{2r}(\chi'(X_i)-1)_{i=1}^7\right).
                                                               \tag{F.33.17}
\]
Thus vanishing on the actual $64$ words
already forces exactly the generator
values in (F.33.2).

Every assignment of the seven generator
values extends uniquely to a continuous
crossed homomorphism with the fixed
scalar character: it is the unique
homomorphism from the original free
profinite group to the finite group
$\operatorname{Aff}_1(k)$ with those
generator images. Hence the generator
identity (F.33.17) is an identity of
continuous cocycles on $V$.
Conversely the two principal cocycles
in (F.33.2) sum to zero on all of $K$,
because $\chi|K=\chi'|K$.
This proves both kernel inclusions.
The vector in (F.33.2) is nonzero for
$b\ne0$, since $\eta\ne1$; its dimension
is one and the rank is $14-1=13$.

For the cohomology statement, the two
principal cocycles restrict to principal
$\eta$-cocycles, so the summed map is
well-defined on $H^1$. A kernel class
has representatives whose sum on $K$
is $b(\eta-1)$. Subtract $b(\chi-1)$
from the first representative. The new
representatives have zero summed
restriction, so (F.33.2) makes both
principal on $V$. Their original
cohomology classes are therefore zero.

\emph{Literal-word audit and the positive marginal.}
The independent word verifier imports
no producer arithmetic and no symbolic
algebra package. It uses generic
crossed-homomorphism multiplication
$(\alpha,D)(\beta,E)=(\alpha\beta,D+\alpha E)$
and inverse $(\alpha^{-1},-\alpha^{-1}D)$
to reconstruct the original nineteen
negative F.23 braid words.
At a positive crossing it applies
\[
 (x_j,x_{j+1};c_j,c_{j+1})\longmapsto
 (x_jx_{j+1}x_j^{-1},x_j;x_jc_{j+1},c_j),
                                                               \tag{F.33.18}
\]
and at a negative crossing it applies
$(x_{j+1},x_{j+1}^{-1}x_jx_{j+1};
c_{j+1},x_{j+1}^{-1}c_j)$.
The eighth meridian is the identity;
the actual decoration is $c_8^{-1}$.
The verifier checks every pure label
permutation, decoration character and
derivative, all eleven full seven-coordinate
actions, all $64$ original rows, the
principal identities, and every recorded
normalization, elimination, pivot and
adjugate. Its $252309$ exact checks pass.
The separate root replay has $32427$
checks, including the full inverse
coordinate maps, and the separate
actual-divisor audit checks all $51$
unit witnesses. The evidence manifest
binds these inputs, programs and outputs.

For the positive marginal retain the
actual based maps of F.23.13,
\[
 \mathcal B_+(g)=w_+\Pi_+(g)w_+^{-1},\qquad
 \mathcal B_+^{-1}(g)=\Pi_+^{-1}(w_+^{-1}gw_+).
                                                               \tag{F.33.19}
\]
Their composites cancel to the identity.
The original filling squares identify
their restrictions $V_-\leftrightarrow V_+$
and $R_-\leftrightarrow R_+$, and the
actual source comparison gives
$\mathcal B_+(\mathcal K_-)=\mathcal K_+$.
Precomposition sends a positive character
and cocycle to
$(\chi\mathcal B_+,d\mathcal B_+)$;
precomposition with $\mathcal B_+^{-1}$
is its inverse. Since $Y_i^+=\mathcal B_+(X_i)$,
this preserves all seven displayed
values, the paired coordinate maps,
and the actual restriction sums.
It also sends every negative word in
(F.33.5) to its exact positive image.
The equality of cocycles, both kernel
inclusions and the cohomology injection
therefore transport to the positive
marginal. All groups and coefficient
maps used here are continuous, with
finite coefficient targets throughout.
This proves the claim.
\end{proof}

\clearpage\bibliographystyle{plain}\bibliography{LM4_FINITE_GROUP_ACTIONS_references}
\end{document}
